\newcommand{\weightedpin}[3]{%
  \href{https://github.com/lsernau/grh-lean/blob/72fead102f663d5c7a407964a13dc0cf9d977207/#1\#L#2}{#3}}

\subsection*{Admissible data at a square scale}

\teamanchor{proof.weighted-data}%
Fix \(M=18480\), a prime \(d\equiv9239\pmod M\), and \(0<\varepsilon<1/100\).
Thus \(M\mid 2(d+1)\) and \(M\nmid d+1\), as required by the compact
parameters. Use the least maximizing square \(x\) of each sufficiently
late block in the \hyperlink{proof.weighted-square-selection}{test-scale
selection}. The data construction works once \(x\) is large in terms of
\(d\) and \(\varepsilon\).

Take \(m\) and \(r\) to be the least powers of two and three at least
\(\sqrt{x}\) and \(x^5\), respectively. Thus \(\sqrt{x}\le m\le2\sqrt{x}\)
and \(x^5\le r\le3x^5\). The number of lifts is \(m\), the definite
source rank is \(\kappa=r+1-m\), and the enlargement parameter is \(t=5r\).
Set \(s=22\log_2 m\), so \(2^s=m^{22}\) and
\(s\le (11/\log2)\log x+22\le32\log x\) for large \(x\).
Write \(c=\lfloor\log_3 s\rfloor+5\) and \(N_1=M3^c2^s\).
The CM field is \(E=K(\mu_{N_1})\) for \(K=\mathbf Q(\sqrt{-d})\). Its infinite
places \(\Phi\), represented by the embeddings extending a fixed embedding of
\(K\), are indexed by \((\mathbf Z/N_1\mathbf Z)^\times\). Indeed \(d\) is
ramified in \(K\) and does not divide \(N_1\), so \(K\) and
\(\mathbf Q(\mu_{N_1})\) intersect in \(\mathbf Q\). The tops \(\mathcal T\)
are the places of class \(1\) modulo \(M\). Reduction of the unit groups
is surjective, and \(2\) and \(3\) already divide \(M\). Thus
\(T=\#\mathcal T=3^c2^s>0\) and \(f=\#\Phi=\varphi(M)T\).
Write \([u]\) for the place indexed by \(u\) modulo \(N_1\), with the
identity place used when \(u\) is a nonunit. Set \(\log b=x^\varepsilon\)
and \(\log x_1=x^{\varepsilon/6}\). Then
\begin{equation}
\label{eq:weighted-data-scales}
  \log b=x^\varepsilon,\qquad \log x_1=x^{\varepsilon/6},\qquad
  (10^6\log x)^2\le\log x_1,\qquad xb\le a<xb+N_1,
\end{equation}
where \(a\) is the least multiple of \(N_1\) at least \(xb\).

\teamanchor{proof.weighted-scale-constant}%
For any fixed \(C_0\ge0\), the condition
\(10^{12}C_0\le\log\log b=\varepsilon\log x\) holds eventually. The
inequality \(100\log\log b\le\log x\) uses \(\varepsilon<1/100\).
We also have \(10^6\log x\le\log x_1\), \(10^6\log x_1\le\log b\), and
\(6\log\log x_1=\log\log b\). In particular, \(b\) dominates every
fixed power of \(x\) and of \(x_1\). The required fixed constant includes
\(M+2d+6\), \(|\ell_d|\), the finite sum of \(|\log|\Phi'_M(\xi)||\) over
primitive \(M\)-th roots, and the fixed constant
\(C_d=C_{\rm vac}+2d+6\) for the
\hyperlink{proof.weighted-cm-moved-vacuum}{moved vacuum bound}.
The \hyperlink{proof.weighted-cm-jet-transport}{normalized Taylor transport}
uses the same CM period. These constants are fixed before \(x\).
\footnote{\weightedpin{grh-compact/src/New/Data/AdmissibleData.lean}
{12-L45}{Admissibility and the fixed constant}.}

\teamanchor{proof.weighted-thick-places}%
To choose the thick places, put \(A=M3^c\), \(a_0=1+A\), and \(L_0=324s\).
The integer \(a_0\) is a unit modulo \(N_1=A2^s\), since \(A\) is even.
We have \(v_2(A)=4\)
and \(v_2(a_0^{2^j}-1)=4+j\) for \(j\ge0\), by successive squaring since
\(a_0^{2^j}+1\equiv2\pmod4\). Thus the order of \(a_0\) modulo
\(2^{s+1}\mid N_1\) is exactly \(n=2^{s-3}\).
Apply the \hyperlink{proof.thick-sample-construction}{digit-and-sphere construction}
with modulus \(n\) and coefficient bound \(L_0\). Here
\(s=22\log_2 m\asymp\log x\). With \(u=\log(n/(4L_0))\asymp\log x\) and
\(v=\log(8L_0)=O(\log\log x)\), the dimension
\(q_0=\lfloor\sqrt{u/v}\rfloor\) tends to infinity and
\(e^{u/q_0}/L_0\to\infty\), as required for the digit construction.
In its counting bound, \(\log q_0=O(\log\log x)=o(\sqrt{uv})\).
It gives labels \(0<e<n/(4L_0)\) whose number is at least
\(nL_0^{-1}\exp(-16\sqrt{\log n\log L_0})\). For three distinct labels
it excludes
\(j(e_2-e_1)\equiv i(e_3-e_1)\pmod n\) when \(0<i<j\le L_0\).

Take \(\mathcal J\) to be the places \([-a_0^e]\) on these labels. Their
images under reduction \(\phi\) modulo \(2^{s+1}\) are distinct, because
the labels lie in an interval shorter than the order \(n\). They all
have class \(-1\) modulo \(A\), hence modulo every divisor of \(A\).
If \(\sigma,\tau,\rho\in\mathcal J\) satisfy
\(\phi(\tau/\sigma)^j=\phi(\rho/\sigma)^i\) for \(0<i<j\le L_0\), their
labels satisfy the displayed congruence. If \(e_2\ne e_1\), neither
\(e_3=e_1\) nor \(e_3=e_2\) is possible, since a nonzero multiple of
\(e_2-e_1\) by \(j\) or \(j-i\) has absolute value less than \(n/4\).
The exclusion for distinct labels therefore forces \(e_2=e_1\), giving
\(\phi(\tau/\sigma)=1\). This supplies the congruence and three-term
conditions of the datum. The pinned construction uses the same cyclic
coset and reduction.\footnote{\weightedpin
{grh-compact/src/New/Data/DataExist.lean}{83-L96}{Placement in the cyclic coset}
and \weightedpin{grh-compact/src/New/Data/AdmissibleData.lean}{32-L37}{the thick-place conditions}.}

Since \(T=8\cdot3^c n\) and \(3^c,L_0=O(\log x)\), the cardinality
\(J=\#\mathcal J\) satisfies
\[
 \log(T/J)\le\log(8\cdot3^cL_0)+16\sqrt{\log n\log L_0}=o(\log x).
\]
In particular \(J\ge T/\log b>0\) for sufficiently large \(x\) at fixed
\(\varepsilon\).

\paragraph{The split level primes.}
\teamanchor{proof.weighted-level-prime-existence}%
For a prime \(d>3\) with \(d\equiv3\pmod4\), put \(K=\mathbf Q(\sqrt{-d})\).
Let \(H\ge3\) be an integer and \(B'=4dH^2\). There is an absolute
effective \(C_{\rm p}>0\) such that, if \(X\ge(B')^{C_{\rm p}}\) and
\(1\le v_0<v_1\le2\) with \(v_1-v_0\ge1/2\), there is a prime
\(P\in(v_0X,v_1X]\) with \(P=\pi\bar\pi\) for
\(\pi\in\mathcal O_K\) satisfying \(\pi\equiv1\pmod{H\mathcal O_K}\).

Take the ray class field \(L/K\) of \(H\mathcal O_K\). The Artin map
identifies the ideals prime to \(H\), modulo principal ideals with a
generator congruent to one, with \(\operatorname{Gal}(L/K)\)
\cite[Ch.~V, Theorems~3.5--3.6]{f1milne-cft}.
Conjugation preserves the ray subgroup, so uniqueness of the class
field gives that \(L/\mathbf Q\) is Galois. A rational prime \(P\)
prime to \(dH\) that splits completely in \(L\) has a prime
\(\mathfrak p\) of norm \(P\) in the trivial ray class of \(K\).
Thus \(\mathfrak p=(\pi)\) with \(\pi\in\mathcal O_K\) and
\(\pi\equiv1\pmod{H\mathcal O_K}\). Its norm gives
\(\pi\bar\pi=P\equiv1\pmod H\).

Apply Thorner and Zaman with base field \(\mathbf Q\), conjugacy class
\(\{1\}\) in \(\operatorname{Gal}(L/\mathbf Q)\), and abelian subgroup
\(\operatorname{Gal}(L/K)\) with fixed field \(K\). Its Hecke characters
have conductors dividing \(H\mathcal O_K\), so their largest conductor
norm \(Q_{\rm TZ}\) is at most \(H^2\). The absolute discriminant \(d\)
and degree \(2\) of \(K\) give the theorem's parameter
\(D_{\rm TZ}=4dQ_{\rm TZ}\le B'\).
For \(y\ge D_{\rm TZ}^{c_2}\) the count \(\Pi(y)\) of unramified rational
primes at most \(y\) splitting completely in \(L\) is
\[
 \Pi(y)=\frac{A(y)}{[L:\mathbf Q]}(1+e(y)),\qquad
 e(y)=O\!\left(e^{-c_3\log y/\log D_{\rm TZ}}
                    +e^{-\sqrt{c_3\log y/2}}\right)
\]
with absolute effective constants
\cite[Theorem~1.4, p.~1042, and \S2]{f6thorner-zaman-chebotarev}.
Here \(A(y)=\operatorname{Li}(y)\) in the ordinary case. If there is
an exceptional real zero \(\beta_1\) with
\(0<1-\beta_1<1/(8\log D_{\rm TZ})\), then \(\beta_1>1/2\) and
\(A(y)=\operatorname{Li}(y)-\operatorname{Li}(y^{\beta_1})\), since the
exceptional character has value one on the identity class.

Put \(q_X=1-X^{\beta_1-1}>0\) in the exceptional case and \(q_X=1\)
otherwise. On \(X\le y\le2X\) one has
\[
 A'(y)\ge q_X/\log(2X),\qquad A(y)\ll Xq_X/\log X .
\]
In the exceptional case the derivative is
\((1-y^{\beta_1-1})/\log y\). Also, integrating from \(y^{\beta_1}\)
to \(y\) gives \(A(y)\le2y(1-y^{\beta_1-1})/\log y\).
Since \(y\le X^2\), the factor \(1-y^{\beta_1-1}\) is at most \(2q_X\).
In the ordinary case the upper bound follows by splitting the
logarithmic integral at \(\sqrt y\).
The main-term increment \(A(v_1X)-A(v_0X)\) is at least
\(Xq_X/(2\log(2X))\). The two endpoint errors in
\([L:\mathbf Q](\Pi(v_1X)-\Pi(v_0X))\) have total absolute value
\[
 O\!\left(\frac{Xq_X}{\log X}
       \left(e^{-c_3\log X/\log B'}+e^{-\sqrt{c_3\log X/2}}\right)\right).
\]
The common positive factor \(Xq_X/\log X\) cancels in the comparison.
Since \(B'\ge252\), choose \(C_{\rm p}\ge\max(1,c_2)\) large enough
to make the error smaller than the main-term increment for
\(X\ge(B')^{C_{\rm p}}\). The count difference is positive. Each prime
in it exceeds \(X\ge B'>dH\), so the ray class argument supplies its
generator.

\teamanchor{proof.weighted-level-primes}%
For the level primes, put \(H=(Md)^{2(s+6)}\) and \(B'=4dH^2\).
Since \(d\) is fixed and \(s=O(\log x)\), we have
\(\log B'=O_d(\log x)\). The condition \(X\ge(B')^{C_{\rm p}}\)
in the \hyperlink{proof.weighted-level-prime-existence}{split-prime existence result}
holds uniformly for \(X\ge b^{10}\), since \(\log b=x^\varepsilon\).
At \(X=b^{10}\), use the two disjoint intervals \((X,3X/2]\) and
\((3X/2,2X]\) to choose distinct \(g_0,g_1\) with generators
\(\pi_i\in\mathcal O_K\), \(\pi_i\bar\pi_i=g_i\) and
\(\pi_i\equiv1\pmod{H\mathcal O_K}\). Both primes lie in
\((b^{10},2b^{10}]\subseteq[b^{10},b^{11}]\) for large \(x\).
Then \(\pi_G=\pi_0\pi_1\) has norm \(g=g_0g_1\). Both primes are
\(1\) modulo \(H\), and \(N_1d\mid H\) since \(2,3\mid Md\) and
\(c=\lfloor\log_3 s\rfloor+5\le s+5\) in the definition of \(N_1\).
In particular \(g_i\equiv1\pmod{N_1d}\) and each splits completely in \(E\).

For every \(U\in[100\log b,2^s]\), apply the same result to
\((e^U,2e^U]\). Since \(e^U\ge b^{100}\), the same threshold in \(x\)
works for all \(U\). This gives a prime \(p\equiv1\pmod H\)
with \(U<\log p\le U+\log2\), and
\(\log2\le\sqrt{\log b}\) for large \(x\). Together with the scales,
powers and thick set above, these choices give an admissible datum \(D\).
The two \(g_i\) specify the principal-level type conditions. The extra
split prime \(N\) is chosen below from the window for the weights of
this \(\mathcal J\) and specifies ramification on the columns.

\teamanchor{proof.weighted-signed-asks}%
Set \(w(\sigma)=a/x\) for \(\sigma\in\mathcal J\) and \(w(\sigma)=a/x_1\)
otherwise, and write \(\bar w=f^{-1}\sum_{\sigma\in\Phi}\log w(\sigma)\).
Put \(\chi_d=(\cdot/d)\) and \(\ell_d=\gamma-\Re(L'/L)(1,\chi_d)\).
The quantities used in the volume calculation are
\begin{align}
  \operatorname{Ask}(u,\tau)
    &=\sum_{1\le\nu\le a}\Lambda(\nu)(1-\chi_d(\nu))
      \bigl(u/\nu-w([-\nu]\tau)\bigr)_+,\notag\\
  \operatorname{MainAsk}&=a(\log a-1-\ell_d-\bar w),\notag\\
  \mathrm{rem}&=\sum_{\tau\in\mathcal T}
       \bigl(\operatorname{Ask}(a,\tau)-\operatorname{MainAsk}\bigr).
\label{eq:weighted-remainder}
\end{align}
The integral weight at stage zero is \(w_0(\sigma)\), obtained by rounding
\(w(\sigma)\) up to a multiple of \(N_1\) and adding the half-gap \(r^8\).
With \(\operatorname{rd}_E\) the root discriminant of \(E\), put
\begin{align*}
  \Omega&=(2\pi/\operatorname{rd}_E)e^{\ell_d}, &
  S(u)&=\log(u!)-u\log(\Omega/g)\quad(u\in\mathbf Z_{\ge0}),\\
  \mathcal U&=\sum_{\sigma\in\Phi}
     \bigl(S(w_0(\sigma))-\operatorname{Ask}(w(\sigma),\sigma)\bigr), &
  h_{\mathrm{main}}&=m\kappa\mathcal U,\\
  \mathcal B&=fa(\log b)^3, & L&=\mathcal B\log b.
\end{align*}
\teamanchor{proof.weighted-ramification-prime}%
Put \(G=\log(g\operatorname{rd}_E/(2\pi))\) and set
\[
  \lambda=m\bar w+(m-1)G-\tfrac{15}{2}\log b.
\]
Since \(1\le w(\sigma)\le b^2\), \(g\in[b^{20},b^{22}]\), and
\(\log\operatorname{rd}_E=O(\log x)\) by \(E\subseteq\mathbf Q(\mu_{N_1d})\),
we have \(2\log b\le G\le27\log b\) and
\(m\log b\le\lambda\le29m\log b\) for large \(x\).
In particular \(100\log b\le\lambda\le2^s=m^{22}\).
Use the \hyperlink{proof.weighted-level-primes}{prime window} with
\(U=\lambda\) and choose \(N\equiv1\pmod H\) with
\[
  \lambda<\log N\le\lambda+\log2\le\lambda+\sqrt{\log b}.
\]
Then \(N\equiv1\pmod{N_1d}\), \(N\) splits completely in \(E\),
and \(N-1\ge b^9\) for large \(x\). Since \(m\to\infty\), the lower
bound \(m\log b\) for \(\log N\) also separates \(N\) from the two
primes \(g_i\le b^{11}\). As \(4\mid N_1\), the sum of two squares
for \(N\) gives \(\pi\in\mathbf Z[i]\subset\mathcal O_E\) with
\(\pi\bar\pi=N\).\footnote{\weightedpin
{grh-compact/src/New/Seesaw/IndexTypes.lean}{130-L178}{Ramification prime and its level window}.}

\teamanchor{proof.weighted-datum-characters}%
\paragraph{Roots of unity and local frames.}

Put \(F=E^+\) and \(L_{\rm cyc}=\mathbf Q(\zeta_{N_1})\). The intersection \(L_{\rm cyc}\cap K\) is trivial over
\(\mathbf Q\) because \(d\) ramifies in \(K\) and is prime to \(N_1\). Hence the restrictions
of the \(\Phi\)-embeddings to \(L_{\rm cyc}\) run once through \(\operatorname{Gal}(L_{\rm cyc}/\mathbf Q)\). The
roots of unity of \(E\) form exactly \(\mu_{N_1}\). Indeed \(E\subseteq\mathbf
Q(\zeta_{dN_1})\), where \(dN_1\) is even. A root of unity in this cyclotomic field has order
dividing \(dN_1\), by comparing the degree at that conductor with the degree at its least
common multiple with the root's order. If \(E\) contained a root with order divisible by
\(d\), it would contain \(\mathbf Q(\zeta_{dN_1})\), which has degree \(d-1>2\) over \(L_{\rm cyc}\).

The field \(E/F\) is unramified at every finite place. At a place of residue characteristic
\(p\), use a root of unity of order \(3\) or \(5\) prime to \(p\). Its reduction still has
that order, while complex conjugation inverts it. If conjugation stabilizes the place, it acts
nontrivially on the residue field. At a nonsplit place the norm on local integral units is
onto. The norm on the residue fields is onto, and the trace is onto on each successive
principal-unit quotient. Successive lifting and completeness prove the assertion also at the
dyadic places. At a split place the local norm on units is the product of the two coordinates
and is onto.

For every diagonal space in the grids use the indicated signs as entries \(\pm1\) in its
rational Hermitian form over \(E/F\), with the standard integral coordinate lattice. Choose
one local unit of norm \(-1\) at each finite place of \(F\), and multiply all negative
coordinates by it in the finite frame. These are simultaneous isometries from the standard
finite Hermitian lattice to the various signed lattices, in every rank. At split places use
the chosen matrix in one half and the inverse conjugate transpose in the other. This
identifies the compact stabilizers and reductions in the comparisons of equal ranks. At a
nonsplit place the resulting integral unitary group becomes \(\mathrm{GL}_n\) over the
unramified quadratic integer ring, and the level remains hyperspecial wherever the data impose
no split-place congruence.

For the original row, the rank is \(\kappa+2\varrho+1=r\), where \(\varrho=m/2-1\). It has a
definite core of rank \(\kappa\), \(\varrho\) hyperbolic pairs and one negative line. An
integral frame for \(\langle1,-1\rangle\) is given by \(e=(1,1)\) and \(e^*=(c_0,c_0-1)\),
with \(c_0=(1+\sqrt{-d})/2\). Both norms are zero and the paired entry is one, since
\(c_0+\bar c_0=1\). The change of basis has determinant \(-1\). This frame is valid over the
integers at two as well. The row chain adds blocks of size \(t\) with the signatures specified
by the sign word in the manuscript. For the equal-rank inner forms \(W_j\) the numbers of
positive and negative entries are \(\kappa+j-1\) and \(m-j\) for \(1\le j\le m\). They sum to
\(r\), and \(W_m\) is definite. In the two grids the sizes of the added blocks are \(t+z\) and
\(t-z\) with \(z=0,1\), respectively. The entry signs and doubled real parameters can then be
read in the ordered blocks already specified in the manuscript. The lower edge of \(Z\) used
for the core descent is untwisted. The different family of edges in the full \(Z\) grid that
increases the first coordinate can use a twist.

\paragraph{The chosen halves, different and splitting character.}

For \(g_i\), use the places \(w\) dividing \((\bar\pi_i)\) in \(E\). Exactly one of \(w,\bar
w\) above a fixed real-base place divides \((\bar\pi_i)\), with valuation one, since
\(\pi_i\bar\pi_i=g_i\) and \(g_i\) splits completely. There are \(f\) such places for each
\(i\). Write \(S_g\) for their union and \(\mathfrak M=(\bar\pi_G)\), where
\(\pi_G=\pi_0\pi_1\). Its support is \(S_g\), so the chosen principal level is the kernel of
these reductions. The \(\pi\)-half over the separate prime \(N\) is constructed in the same
way from \(\pi\bar\pi=N\). If \(Q\) is the product of the distinct chosen halves at new split
pairs, the principal modulus is \(\mathfrak M^eQ\), with \(e\) as specified in the edge. It is
prime to its conjugate. The base core uses \(\mathfrak M^{b_1}\), where \(b_1\) is the
rounding of \(a/x\) up to a multiple of \(N_1\). The terminal reduction is identity at \(S_g\)
and has its indicated parabolic type at \(Q\). Thus the chosen principal congruence at at
least one odd \(g_i\) is still imposed in every terminal use of freeness.

Put \(P(X)=\Phi_{N_1}(X)\) and \(\delta=P'(\zeta_{N_1})\). The \(K\) and \(L_{\rm cyc}\) discriminants
are relatively prime. Locally at every rational prime one integer ring is étale, so their
tensor product is the full integer ring in the compositum. The tower formula for differents
gives \(\mathfrak D_{E/K}=(\delta)\) and \(\mathfrak D_{E/\mathbf Q}=(\delta\sqrt{-d})\).
Reciprocity of \(P\) gives \(\bar\delta/\delta=-\zeta_{N_1}^{\,2-f}\), with
\(f=\varphi(N_1)\). Thus \(\varpi=i\zeta_{N_1}^{(2-f)/2}\delta\sqrt{-d}\) is a purely
imaginary generator of the absolute different; \(4\mid N_1\) and \(f\) is even.
Differentiating \(X^{N_1}-1=P(X)((X^{N_1}-1)/P(X))\) at \(\zeta_{N_1}\) shows that \(\delta\)
divides \(N_1\) in \(\mathcal O_E\). Choose the finite idele \(\alpha\) with component
\(\delta^{-1}\) at the places dividing \(c_0\) and component \(1\) elsewhere. Every rational
prime dividing \(N_1\) divides \((d+1)/4=c_0\bar c_0\). For odd primes this follows from
\(M\mid2(d+1)\), and for \(2\) it follows from \(d\equiv7\bmod8\). Since \(c_0+\bar c_0=1\),
exactly one of \(c_0,\bar c_0\) is divisible by a place over any of these primes. The
conjugate of \(\delta\) differs from it by a unit, so \(\alpha\bar\alpha\delta\) is a finite
unit. For \(j_{\rm Weil}=-1/\sqrt{-d}\) this gives \((j_{\rm
Weil}\alpha\bar\alpha)\widehat{\mathcal O}_E=\mathfrak D_{E/\mathbf Q}^{-1}\). Also
\(\Im\tau(j_{\rm Weil})>0\) and \((1-j_{\rm Weil}^{-1})/2=c_0\), and \(N_1\alpha\) is
integral. The skew scalar for the Weil pairing is \(-1/\sqrt{-d}\), with \(\sqrt{-d}=i\sqrt
d\) in the chosen embedding of \(K\).

There is a unitary splitting character \(\upsilon\) of \(E^\times\backslash\mathbb
A_E^\times\), unramified at all finite places, whose infinite part is \(\prod_{\tau\in\Phi}
z_\tau/|z_\tau|\) and whose restriction to \(\mathbb A_F^\times\) is the quadratic character
\(\epsilon_{E/F}\). To verify the prescribed unit compatibility let \(\zeta=\zeta_{N_1}\) and
\(v=i(\zeta-1)\). It is a unit with \(v/\bar v=\zeta\), and \(N_{E/K}(v)=1\), since
\(N_{L_{\rm cyc}/\mathbf Q}(\zeta-1)=\Phi_{N_1}(1)=1\) when \(N_1\) has at least two prime divisors. If
\(a/\bar a\) is an integral unit for \(a\in E^\times\), all of its conjugates have modulus
one, hence \(a/\bar a=\zeta^j\) for some \(j\). It follows that \(a/v^j\in F^\times\). In
particular every unit of \(E\) is the product of \(v^j\) and an integral unit of \(F\). For a
unit of \(F\), the sign of its norm to \(\mathbf Q\) is \(+1\), because the global character
\(\epsilon_{E/F}\) is trivial on principal ideles and is unramified at every finite place.
Thus the proposed infinite part is trivial on integral units of \(E\).

On the group of principal fractional ideals of \(E\) prescribe the inverse of the proposed
infinite part evaluated at a generator. On the extensions of fractional ideals of \(F\)
prescribe \(\epsilon_{E/F,f}\). The extension-of-ideals map is injective. If such an extended
ideal is \((a)\), then \(a/\bar a\) is an integral unit and \(a=v^j b\) with \(b\in
F^\times\). As \(v\) has norm one over \(K\), the infinite part on \(a\) equals the sign
product on \(b\). This is the inverse of \(\epsilon_{E/F,f}(b)\). The two prescriptions
therefore agree on their intersection. A character of their product subgroup extends to the
whole fractional ideal group with values in the circle, because circle values admit roots. The
quotient by principal ideals is finite, so one can make the extension on a finite set of
generators, taking compatible roots. This gives the required \(\upsilon_f\). Its values on
finite integral ideles are one. At a nonsplit place an \(F\)-uniformizer is also an
\(E\)-uniformizer, and restriction to \(\epsilon_{E/F}\) gives value \(-1\). The character
\(\upsilon\) has \(\upsilon(\bar t)=\upsilon(t)^{-1}\), since \(t\bar t\) is a norm and
\(\epsilon_{E/F}\) kills norms. Accordingly the base change of its restriction
\(\upsilon_{U(1)}\) to the norm-one torus is \(\upsilon^2\). The character \(\upsilon_{U(1)}\)
has real weight one at every \(\tau\) and trivial finite tame part.

\paragraph{Global stage characters and the signed catalogue.}

Put \(h_2=r^8\), \(h=2h_2\), \(M_\omega=(Md)^{s+6}\) and \(w_j(\tau)=w_0(\tau)+jh\). Then
\(H=M_\omega^2\) from the prime construction and \(dN_1\mid M_\omega\). At each chosen \(w\),
fix a residue character of order \(M_\omega^2\) of the cyclic group \(k_w^\times\), available
because \(g_i\equiv1\bmod H\). Let \(\omega_w\) be its \(M_\omega\)-th power, and use the
trivial character outside the chosen halves. At a split norm-one pair the opposite matrix
coordinate is determined by the conjugate inverse. The chosen local \(\omega_w\) has order
\(M_\omega\) and kills the reduction of every norm-one integral unit of \(E\), since these
units are in \(\mu_{N_1}\) and \(N_1\mid M_\omega\). In particular \(\omega_w^a\) is
nontrivial at all the chosen places for \(0<|a|\le 8m\).

For any integer \(A\) the angular norm-one character \(\prod_{\tau\in\Phi} t_\tau^{-A
w_0(\tau)}\) also kills all norm-one integral units of \(E\). Indeed \(w_0(\tau)\equiv
h_2\bmod N_1\), independent of \(\tau\). For a root of unity \(\zeta'\in\mu_{N_1}\), the
angular value is \(N_{L_{\rm cyc}/\mathbf Q}(\zeta')^{-A h_2}=1\). Thus the prescribed real and finite
compact values of weight \(-A w_0\) and tame part \(\omega^a\), for integers \(A,a\), agree on
the rational points in the compact level. Define their product on
\(U(1)(F_\infty)\prod_{v<\infty}U(1)(\mathcal O_v)\), where the product of groups is an open
subgroup of \(U(1)(\mathbb A_F)\). It descends to the image of that open subgroup in the
rational quotient. Extend it to the rational quotient with values in the circle. This
extension is continuous because it is continuous on the open subgroup, and its local tame
parts and weights are unchanged. The quotient by the open subgroup is finite. The ideals of
the norm-one finite ideles are all of the form \(\mathfrak a/\bar{\mathfrak a}\): at a split
pair use its prescribed exponent at one half and zero at the other, and at a nonsplit place
the norm-one exponent is zero. Principal ideals of the form \(\mathfrak a=(u)\) map to
\((u/\bar u)\), the ideal of a rational norm-one point. The ideal class group therefore
surjects onto this quotient. For weight zero one may choose a finite-order extension. No
coefficient-place choices enter this construction.

Choose \(\mu_\flat\) with weight \(-w_0\) and trivial tame part, and \(\chi_0\) with weight
zero and tame part \(\omega\). Keeping these two choices fixed, put

\[
 \mu_j^0=\mu_\flat(\upsilon_{U(1)}^{-h}\chi_0)^j.
\]

Then \(\mu_j^0\) has weight \(-w_j\) and tame part \(\omega^j\). On the actual edge the
remaining rank twist is the specified floor or ceiling power of \(\upsilon_{U(1)}\). A trial
finite-order character \(\psi\) has weight zero and support at the new split pairs, so
replacing \(\mu_j^0\) by \(\mu_j^0\psi\) leaves all real parameters unchanged. The
auxiliary-ray construction requires the places in \(Q\) and their conjugates to avoid
\(g_0,g_1,N,\ell\); its exact-support condition at each required power makes \(\mathfrak M^e
Q\) the same at that power. These are genuine torus characters, as the ray construction kills
the real ideles and the decomposition groups above \(\ell\).

The unsigned family at stage \(j\) uses

\[
 \nu_j=\mu_j^0\upsilon_{U(1)}^{\,j+(t-11m-1)/2},
 \qquad 1\le z'\le 7m/2+\kappa/2 .
\]

It has the same character at each point \(z'\); its tame part is \(\omega^j\), and the weight
is the negative of \(w_j-j-(t-11m-1)/2\). The \(\kappa/2\) largest points are the sharp-line
tests. They satisfy \(m<z'\le r\) eventually, and the positive weight magnitude \(c\) lies in
\([w_j+r-t,w_j]\), since \(j+(t-11m-1)/2\) is nonnegative and at most \(t-r\). The remaining
unsigned points use the positive sign in the signed interpolation. For \(s_1=+1,-1\) put
\(e_+=m\), \(e_-=m/2\) and use

\[
 \nu_{j,s_1}=\upsilon_{U(1)}^{-e_{s_1}}\chi_0^{\,s_1 j},
 \qquad 1\le z'\le e_{s_1}.
\]

The corresponding sign uses the trial power \(\psi^{s_1}\) at that stage. The ramified flat
pair in the rank-one jet is \((\upsilon_{U(1)}^{-2}\chi_0,2)\). In a grid square at \((a,b)\)
the product character before the rank normalization is \(\mu_{a+1}^0\mu_{b+1}^0\) with tame
exponent \(a+b+2\); the trial power occurs from the \(b+1\) edge in \(V\) or the \(a+1\) edge
in the full \(Z\) grid. For a triangle square whose shorter chain has length \(n\), the
corresponding character is \(\upsilon_{U(1)}^{-(n+1)h}\chi_0^{n+1}\), multiplied
by the trial power \(\theta^{n+1}\) in \(V\). The triangle in \(Z\)
uses its trivial-twist specialization. In \(\Pi_{\nu,S}\) the
positive-infinity character is the inverse of the resulting
negative-weight critical-pair character. The needed lower points have \(z'\le t\)
in those squares. Their negative weight magnitudes dominate \(t\) and are at most \(4b\) at
the large scales under consideration. This follows from \(w_0=O(b)\), \(w_0(\tau)\ge a/x_1\),
\(w_j\le2b\) for \(j\le2m\), and \(t<h\ll b/m\). For the explicit signed and flat families the
lower points satisfy \(1\le z'\le c\) directly. Every ramified family just listed has a
nonzero tame exponent of absolute value at most \(2m\). Thus each signed Euler factor at the
fixed helper inertia has a nontrivial root of unity of order dividing \(q_w-1\).

Finally the change of a Green label by the stage tame character preserves the allowed set.
Squarefreeness and factor degrees are unchanged by scaling roots by an element of the residue
field. The forbidden roots form the subgroup of solutions of \(x^{M_\omega}=1\), which is
invariant under this scaling. A primitive irreducible root also remains primitive. In the
cyclic group of order \(q_w^r-1\), the multiplier has order dividing \(M_\omega\).
Its exponent is divisible by \((q_w^r-1)/M_\omega\), which still contains
every prime factor of \(q_w^r-1\) because \(M_\omega^2\mid q_w-1\).
Adding it to a primitive exponent preserves
coprimality. This verifies the particular local label changes used in the lower row and the
scalar-inertia test.
\footnote{\weightedpin{grh-compact/src/New/Seesaw/IndexTypes.lean}
{35-L111}{Field, finite frames and stage characters},
\weightedpin{grh-compact/src/New/Rows/Separation.lean}
{43-L65}{The original auxiliary ranks and weights}, and
\weightedpin{grh-compact/src/New/Exchange/TriangleArrows.lean}
{22-L63}{The triangle rank and character data}.}

\subsection*{Projected Poincar\'e series and the seesaw}

\teamanchor{proof.weighted-projection}%
Let \(\mathcal C\) be the column space of algebraic forms on the definite
unitary group \(U(X_0)\) of rank \(\kappa\), at the prescribed weight
and stage-zero level. This level imposes the chosen principal level and
requires the last row to be congruent to \((0,\ldots,0,1)\) modulo \(\pi\).
The initial Poincar\'e series takes values in bounded continuous
automorphic forms on the indefinite group \(U(Y_0)\) of rank \(r\),
fixed by the prescribed real type kernel and of principal level
at the chosen places above \(g_0,g_1\). This space has the
unitary action \(\rho\) of \(\mathcal G=\prod_v\mathrm{GL}_r(k_v)\), where
\(v\) runs over the chosen places above \(g_0,g_1\) and \(k_v\) is the
residue field. The central idempotent \(e_0\) for
the chosen types at the places over \(g_0,g_1\) gives the row space
\(W=\operatorname{im}\rho(e_0)\). Write \(C=\dim_{\mathbf C}\mathcal C\)
and \(R=\dim_{\mathbf C}W\). The split cuspidal types give
cuspidality on \(W\) by the constant-term argument in the
\hyperlink{proof.weighted-good-transfer}{trace comparison}.
\footnote{See the
\weightedpin{grh-compact/src/New/Columns/ColumnDefs.lean}{38-L73}{definite column model}
and the \weightedpin{grh-compact/src/New/Trace/GoodIdempotent.lean}{103-L121}{good row space}.}

Evaluation of the columns in double-coset coordinates and in the local
Weyl-to-Bessel coordinates gives the input to a Poincar\'e series map
\(\mathcal P\) on \(U(Y_0)\). Its seed vectors use the product of
the full real Bessel vectors, the class cutoff and the finite
principal-coset indicator.
Let \(\mathcal R=\prod_{v\mid\pi}(1-z_v)\) on \(\mathcal C\), where
\(\pi\in\mathcal O_E\) satisfies \(\pi\bar\pi=N\) and \(z_v\) is the scalar
translation chosen at \(v\).
Write \(\mathcal C_{\rm ram}=\operatorname{im}\mathcal R\) for the ramified
column space.

\paragraph{Trivial stabilizers at the chosen level.}
\teamanchor{proof.weighted-free-level}%
On a definite group over \(E^+\), the stabilizer of any adelic representative
at a compact integral level is finite. The finite coordinates put its
matrix entries in a fixed fractional lattice and definiteness bounds them
at infinity, where the lattice is discrete. At a chosen place over an odd
prime \(p=g_i\) the split coordinates are over \(\mathbf Q_p\). The
principal congruence group \(1+pM_n(\mathbf Z_p)\) is torsion-free. For
\(I+p^hA\ne I\) with \(h\ge1\) and \(A\) having a unit entry, a power prime
to \(p\) leaves the least entry valuation of its difference from \(I\) equal
to \(h\). Its \(p\)-th power raises this valuation to \(h+1\), by the
binomial expansion. Conjugating a finite class stabilizer into this
principal level at the chosen place proves it is trivial. Thus the column
level, the enlarged level below, and the terminal level all have trivial
class stabilizers. At the terminal definite space \(Y_m\), the residue
conditions are identity at the chosen places and block upper triangular at
the twist places, so the principal condition still holds for every twist.
At any other integral level the stabilizer injects under reduction at the
chosen places by applying the same calculation to its kernel.\footnote{\weightedpin
{grh-compact/src/New/Data/ChosenPlaceLevels.lean}{41-L79}{Chosen principal levels}
and \weightedpin{grh-compact/src/New/Seesaw/GradingsOfY.lean}{96-L104}{the terminal residue conditions}.}

\paragraph{The central action on the column classes.}
\teamanchor{proof.weighted-central-orbits}%
At each \(v\mid\pi\), the scalar \(z_v\) lifts a generator of
\(k_v^\times=\mathbf F_N^\times\), acts by the inverse conjugate scalar
at \(\bar v\), and is trivial elsewhere. The places \(v\mid\pi\) choose
one from each conjugate pair and number \(f\). The translations lie in
the chosen principal level since \(N\ne g_i\). The column level \(K_0\)
and these central elements generate \(K^+\). The last-row congruence
shows that \(K^+/K_0\simeq\prod_{v\mid\pi}\mathbf F_N^\times\), since
\(\prod_vz_v^{a_v}\) belongs to \(K_0\) precisely when \(N-1\) divides
each \(a_v\). A translation fixing a class would supply a rational
element of the class stabilizer at \(K^+\). Triviality of that stabilizer
makes the action of this quotient free.

On an orbit, choose representatives \(a_fz\) from one fixed \(a_f\).
The orbit values are then copies of the same coefficient space on which
the \(z_v\) act by cyclic shifts, preserving its local Weyl gauges.
Apply the \hyperlink{proof.diamond-augmentation-kernels}{augmentation
calculation for cyclic shifts} to this orbit over the coefficient ring.
In particular the integral ramified columns at a coefficient embedding
are the image of the integral columns under \(\mathcal R\), and
\[
   \dim\ker\mathcal R
   =C\bigl(1-(1-(N-1)^{-1})^f\bigr)\le fC/(N-1).
\]
The triangle inequality for the cyclic shifts gives a squared supremum
norm loss of at most \(4^f\) under every coefficient automorphism.\footnote{\weightedpin
{grh-compact/src/New/Columns/ChainValuation.lean}{94-L99}{Gauges and conjugate norms}
and \weightedpin{grh-compact/src/New/Columns/ChainValuation.lean}{157-L202}{the central scalar action}.}

\paragraph{Constant terms of the normalized elliptic translates.}
\teamanchor{proof.weighted-elliptic-torsion-values}%
The base theta model has the odd sigma section \(s_\Lambda\) and its
translate \(\theta\) by an order-two point \(t_0\), normalized by
\(\theta(0)=1\). The two components of \(t_0\) under the two idempotents
of \(\mathcal O_K/2\mathcal O_K\) are nonzero. This is the condition that
neither \(\omega t_0\) nor \(\bar\omega t_0\) be a period for
\(\omega=(1+\sqrt{-d})/2\). Write \(a_1\) for the coefficient of the
marked cubic \(y^2+a_1xy=x^3+x\), with \(t_0=(0,0)\). Its discriminant
and invariant are
\[
  \Delta=a_1^4-64,\qquad j=\frac{(a_1^4-48)^3}{a_1^4-64}.
\]
The CM invariant \(j\) is an algebraic integer, so \(a_1\) is algebraic.
One can see integrality of \(j\) here from the cyclic modular polynomial
for \(k+\omega\) with a sufficiently large integer \(k\). The kernel is
cyclic since \(k+\omega\) is primitive in \(\mathcal O_K\) and the period
ideal is invertible. Its norm \(n=k^2+k+(d+1)/4\) is between \(k^2\) and
\((k+1)^2\) for large \(k\). On the diagonal the integral modular
polynomial is monic up to sign. Indeed, in the product over the cyclic
sublattices of index \(n\), each factor
\(j(\tau)-j((u\tau+b)/v)\), \(uv=n\), has distinct leading cusp
exponents \(-1\) and \(-u/v\). Its leading coefficient is a root of
unity. The integral leading coefficient of the product is then \(\pm1\).
Since the cyclic isogeny is an endomorphism, the diagonal polynomial
vanishes at the given \(j\).

The actual \(\Delta\) is a unit at every finite coefficient embedding.
At odd \(p\), \(|a_1^4|_p>1\) would give \(|j|_p>1\). If
\(|\Delta|_p<1\), then \(|a_1^4-48|_p=|16|_p=1\), also impossible.
At \(p=2\), take a good model of the given \(j\) after a finite coefficient
extension. For example take \(y^2+Bxy+y=x^3\) with
\((B^4-24B)^3=j(B^3-27)\). A root \(B\) of this monic integral
equation is integral, and \(\widetilde B^3=1\) in the residue field would
give \(1=0\) in the displayed equation. Thus \(B^3-27\) is a unit.
The CM endomorphisms of the generic curve extend over the good elliptic
model by the N\'eron mapping property. The two idempotents modulo \(2\)
split its finite flat group of 2-torsion into rank-two factors. On the
geometric special fiber their tangent dimensions add to one, since
the differential of doubling is zero. A rank-two factor is either
\'etale with tangent dimension zero or connected with positive tangent
dimension. Hence exactly one factor is \'etale. The generic point
\(t_0\) has a nonidentity component in that factor, whose specialization
is injective on geometric points. In particular \(t_0\) specializes
away from the origin of the good model.

Translate this integral point to \((0,0)\) in the good Weierstrass model.
Its order two makes the coefficient of \(y\) zero, leaving an equation
\(y^2+A_1xy=x^3+A_2x^2+A_4x\) with integral coefficients. Its discriminant
is \(A_4^2((A_1^2+4A_2)^2-64A_4)\), so \(A_4\) is a unit. Remove
\(A_2\) by \(y\mapsto y+s x\) with \(s^2+A_1s=A_2\) and scale by a
fourth root of \(A_4\). Both changes are integral with unit scaling.
Two Weierstrass models of this normalized shape representing the same
marked curve have change of scale \(u\) with \(u^4=1\). Indeed the
translations in \(x,y\) vanish at the marked point and the coefficient
of \(x\) is one in both models. This gives \(|\Delta|_2=1\) for the
actual normalized cubic. In particular \(a_1\) is integral everywhere.
\footnote{\weightedpin{grh-compact/src/Lib/CM/Sigma/CurveTheta.lean}{81-L118}{Marked model and theta normalization}, and
\weightedpin{grh-compact/src/Lib/CM/LocalHeight/IntegralParameter.lean}{102-L117}{The discriminant assertion on that model}.}

For a coefficient embedding into \(\mathbf C_p\) write
\(H(z)=\max(|s_\Lambda(z)|_p,|\theta(z)|_p)\), where absolute values are
taken after the embedding. The two sections have disjoint zero sets, so
\(H(z)>0\). Put \(u=s_\Lambda(z)\), \(v=\theta(z)\). The
sigma duplication and theta addition formulas give
\[
 \theta(2z)=v^4-u^4,
 \qquad s_\Lambda(2z)^2=(uv)^2\bigl(4(v^4+u^4)+(a_1uv)^2\bigr).
\]
For the second formula, use \(v^2/u^2=\wp(z)-\wp(t_0)\) in the
equation for \(\wp'(z)^2\) and the sigma duplication formula. The first
is the even theta addition formula at equal arguments. These analytic
identities extend through the zeros of \(u\) and \(v\).
Normalize \(u,v\) by one of maximal absolute value. If both output
coordinates have smaller absolute value than the indicated fourth
power, their residues satisfy \(U^4=V^4\) with both \(U,V\) units and
\(8+a_1^2(U/V)^2=0\). Squaring the latter equation contradicts that
\(\Delta\) is a unit. The displayed formulas therefore give
\(H(2z)=H(z)^4\) at every coefficient embedding.

For any section \(f\) of degree \(k>0\) for the multiplier of
\(\theta^k\), the quotient \(f/\theta^k\) is a meromorphic function on
the elliptic curve with pole of order at most \(k\) at \(t_0\). Use the
ordinary Weierstrass coordinates evaluated at \(z+t_0\), whose poles
are at \(t_0\). Their monomials with pole orders \(2i\) and
\(2i+3\), together with the constant, give the basis with these pole
bounds by Riemann--Roch. Write \(u^2=x(z+t_0)v^2\) and \(G=y(z+t_0)v^3\).
After multiplying by \(\theta^k\), the basis elements are
monomials in \(u^2,v,G\) of degrees \(2,1,3\), respectively, where the
cubic identity is \(G^2+a_1u^2vG=u^6+u^2v^4\). By the ultrametric
inequality \(|G|_p\le H^3\). It follows that \(|f(z)|_p/H(z)^k\) is
bounded by the maximum absolute value of its coefficients in this basis.

Let \(y\) be torsion modulo \(\Lambda\) and put
\[
 A_y(z)=\exp\bigl(-\pi(z\bar y+|y|^2/2)/\operatorname{area}(\Lambda)\bigr),
 \qquad R_y(z)=\frac{|A_y(z)|_p H(z+y)}{H(z)}.
\]
Raise the two twisted translates to an even multiple of the order
of \(y\). Their multipliers then agree with that of the same power
of \(\theta\), so the preceding basis gives a finite bound for
\(R_y(z)\), uniformly in \(z\). The Appell--Humbert multiplier changes
\(R_y\) by the absolute value of a root of unity when \(y\) is shifted
by a period, and this absolute value is one. In fact the extra phase
is an exponential of \(\pi i\) times a rational Riemann pairing because
\(y\) is torsion. Also \(R_y(z)^4=R_{2y}(2z)\) by the duplication of
\(H\) and \(A_y(z)^4=A_{2y}(2z)\). Hence the positive finite suprema of
\(R_{2^j y}\) are successive fourth powers of one another on a finite
orbit. Each equals one. At \(z=0\), \(H(0)=1\), and the actual
translated value \(A_y(0)\theta(y)\) is integral at every coefficient
embedding.\footnote{\weightedpin{grh-compact/src/Lib/CM/Sigma/CurveTheta.lean}{291-L359}{The analytic identities}, and
\weightedpin{grh-compact/src/Lib/CM/LocalHeight/IntegralParameter.lean}{151-L214}{Local-height inputs and Tate's trick}.}

The complex conjugate bound requires only the action on these torsion
values. Let \(n\) be the order of \(y\) and choose \(k\) divisible by
\(n\). The algebraic degree-\(k\) theta line is \(\mathcal O(k[t_0])\).
Identify its canonical section with \(\theta^k\) and normalize its
evaluation at zero to be one. Translation
by \(y\) preserves this line up to isomorphism. The analytic operator
\(f(z)\mapsto A_y(z)^k f(z+y)\) is a scalar multiple of an algebraic
theta-group lift. Its \(n\)-th power is a root of unity times the identity,
by the period multiplier, so that scalar is algebraic. After a complex
automorphism, the transported operator and the corresponding analytic
operator for the transported torsion point differ by a root of unity since
they lift the same translation and both have finite order. Evaluating on
the normalized section at zero makes \((A_y(0)\theta(y))^k\) algebraic.
Taking \(k\)-th roots of the absolute values proves the equality of the
moduli with those in the transported normalized model. The period lattice
and marked half period up to
lattice translation are specified by the cubic invariants and its
invariant differential. There are only finitely many conjugate models,
since \(a_1\) is algebraic, and the lattice translation phases have
absolute value one. On each model the norm
\(e^{-\pi|z|^2/(2\operatorname{area}(\Lambda))}|\theta(z)|\) descends to
the compact elliptic torus. Its maximum, over these finitely many
models, gives one constant \(B_d\ge1\) for all torsion values and all
their complex conjugates.

\paragraph{The rank-one theta product and its finite projection.}
\teamanchor{proof.weighted-vacuum-projection}%
Write \(\nu=N_1\) in this paragraph and let \(\zeta\) have exact order
\(\nu\). Put \(\eta=\sqrt{2\pi/\sqrt d}\). The rank-one standard vacuum
uses the lattice \(L=r_\infty(\mathfrak a_0)\subset\mathbf C^\Phi\),
where \(r_\infty(\gamma)=(\eta\tau(\gamma))_{\tau\in\Phi}\) and
\(\mathfrak a_0=E\cap\alpha\widehat{\mathcal O}_E\) with the prescribed
finite idele \(\alpha\). The latter has component \(P'(\zeta)^{-1}\)
at places dividing \(\omega\), and component one elsewhere, with
\(P=\Phi_\nu\). The relative different is \((P'(\zeta))\). Indeed the
cyclotomic and quadratic discriminants are relatively prime, so the
compositum has \(\mathcal O_E=\mathcal O_K[\zeta]\). This follows locally
by tensoring with the unramified factor at each rational prime. At any
place where \(P'(\zeta)\) has positive valuation exactly one of
\(\omega,\bar\omega\) lies in the place, since
\(\omega\bar\omega=(d+1)/4\) is divisible by every prime dividing
\(\nu\) and \(\omega+\bar\omega=1\). Consequently
\(\mathfrak a_0\bar{\mathfrak a}_0=\mathfrak D_{E/K}^{-1}\).
Thus \(L\) is self-dual for the alternating form of the standard
Fock Hermitian form. In the \(E\) coordinates this is the trace pairing
\(\operatorname{Tr}_{E/\mathbf Q}(\bar x y/\sqrt{-d})\), using
\(\mathfrak D_{E/\mathbf Q}=\sqrt{-d}\,\mathfrak D_{E/K}\).

Write \(\Lambda=\Omega_0\mathcal O_K\) for the period lattice of the
principal marked model just used. Define the complex
linear forms and product
\[
  \ell_i(z)=\frac{\Omega_0}{\nu\eta}
           \sum_{\tau\in\Phi}z_\tau\overline{\tau(\zeta^i)},
  \qquad S(z)=\prod_{i=0}^{\nu-1}\theta(\ell_i(z))^\nu.
\]
For \(\gamma\in\mathfrak a_0\), we have
\(\ell_i(\nu r_\infty(\gamma))=\Omega_0\operatorname{Tr}_{E/K}
(\gamma\bar\zeta^i)\in\Lambda\), by the inverse different.
The row orthogonality of the finite Fourier matrix at the distinct
\(\tau(\zeta)\), and
\(\operatorname{area}(\Lambda)=|\Omega_0|^2\sqrt d/2\), show that
the sum over \(i\) of the pullbacks of the elliptic Hermitian form,
each counted \(\nu\) times, is the standard Fock Hermitian form.
The multipliers match too. On the elliptic lattice the theta
multiplier is \(\exp(\pi i E_\Lambda(\bar\omega l,l))\), where
\(E_\Lambda\) is the Riemann form. To check this formula, on either
idempotent line modulo \(2\), the sigma sign cancels the sign from
translating by the mixed point \(t_0\). Hence the multiplier is one
on \(\omega\Lambda\) and \(\bar\omega\Lambda\). The Appell--Humbert law
and \(l=\omega l+\bar\omega l\) give the formula. The \(\ell_i\)
commute with the common scalar \(\bar\omega\), so multiplying the
elliptic phases gives the canonical multiplier of the standard
vacuum on \(\nu L\). In particular \(S\) is a section on this lattice
and \(S(0)=1\).

Let \(\Theta_{\rm vac}\) be the entire theta function of the standard
vacuum on \(L\). Use the trace-one element \(\omega\) in the adelic realization. For
\(\gamma\in E\) the rational translation formula takes
the form
\[
 \Theta_{\varrho(\operatorname{finVec}(\alpha^{-1}\gamma))F}
   = U_{-r_\infty(\gamma)}\Theta_F,\qquad
 \chi_f(\alpha^{-1}\gamma)=\chi_\infty(-r_\infty(\gamma))^{-1},
\]
where \(\operatorname{finVec}\) places the vector in the finite adeles,
\(\varrho\) is the Heisenberg translation, and \(\chi_f,\chi_\infty\)
are the finite and canonical archimedean multipliers, and \(U_r\) is
archimedean Heisenberg translation with its Gaussian factor. These identities
use invariance under the rational Heisenberg translations and the
product formula for the additive character, with the archimedean
Gaussian divided out. If \(\Theta_F=0\), then \(F\) vanishes at the
real-only arguments. Rational covariance moves these to arguments
with dense finite components, so continuity gives \(F=0\).

The theta space on \(\nu L\) has dimension \(\nu^{2f}\), since
\((\nu L)^\vee=\nu^{-1}L\). The translates of \(\Theta_{\rm vac}\)
indexed by \(\nu^{-1}L/L\) are eigenvectors with distinct characters
of \(L/\nu L\), for the translations divided by their vacuum
multiplier. The characters come from the perfect Riemann pairing
of these two finite groups. The translates are therefore a basis.
By the rational translation formula the expansion of \(S\) in this
basis realizes it as the entire theta function of one vector \(F_s\)
in the finite part of the actual standard model. To obtain the vacuum,
choose a translate of \(S\) whose coefficient on \(\Theta_{\rm vac}\)
is nonzero, and apply the sum of the normalized translations in
\(L/\nu L\). This kills the nontrivial characters. It gives a nonzero
scalar multiple of the vacuum as a sum of \(\nu^{2f}\) finite
Heisenberg translates of \(F_s\) with root-of-unity coefficients.

At any argument \(x\) of the finite adeles,
\(F_s(\operatorname{finVec}x)=\Theta_{\varrho(\operatorname{finVec}x)F_s}(0)\).
For \(x=\alpha^{-1}\gamma\), the rational formula gives the product of
elliptic translated values at the arguments \(-\ell_i(r_\infty(\gamma))\).
They are torsion modulo \(\Lambda\) because their traces lie in \(K\).
The raw finite translation is locally constant on \(F_s\), so density
allows every finite adelic argument to be replaced by a rational one
with the same value. Compositions of finite translations add only
root-of-unity phases: for each argument of the finite additive
character some integer multiple belongs to its open kernel. This
gives the integrality and conjugate bounds on the finite adeles for
every summand of the fixed vacuum projection.\footnote{\weightedpin{grh-compact/src/New/Data/WeilDataOfE.lean}{198-L219}{The prescribed finite idele},
\weightedpin{grh-compact/src/New/Curve/FrameSections.lean}{83-L126}{The product on the actual lattice}, and
\weightedpin{grh-compact/src/Lib/Weil/Construction/FockThetaFunctions.lean}{25-L67}{Rational translation and finite realization}.}

\paragraph{The uniform bound for scaled vacua.}
\teamanchor{proof.weighted-vacuum-constant}%
The geometric input on this side concerns a single base \(K\)-CM elliptic
theta model. For the standard rank-one space over \(E\), it asks for one
nonzero scalar \(c_s\) at each \(s\), common to every finite adelic
evaluation. All values of \(c_s\) times the vacuum are integral at every
coefficient place, and every conjugate over \(K\) has modulus at most
\(\exp(C_{\rm vac}N_1^2)\), with \(C_{\rm vac}\) independent of \(s\).
The standard rank-\(n\) formula uses the product of the same rank-one
normalization in the \(n\) coordinates. A complex automorphism either
fixes \(K\) or does so after composition with complex conjugation, which
does not change the modulus.

The special theta product has \(N_1^2\) of the preceding elliptic torsion
factors in rank one. A nonzero multiple of the vacuum is a sum of
\(N_1^{2f}\) Heisenberg translates of the product, with coefficients
roots of unity. Thus the archimedean bound for the scaled vacuum is
\[
  N_1^{2f}B_d^{N_1^2}
  \le \exp\bigl((2+\log B_d)N_1^2\bigr),
\]
since \(f\le N_1\) and \(\log N_1\le N_1\). The same finite sum proves
integrality using the normalized elliptic torsion values at every
coefficient embedding.\footnote{\weightedpin{grh/src/New/Derivatives/Vacuum/IntegralVacuum.lean}{56-L100}{Torsion transport and finite projection},
\weightedpin{grh/src/New/Derivatives/Vacuum/IntegralVacuum.lean}{195-L275}{Rank-one product and bounded integral vacuum}, and
\weightedpin{grh-compact/src/New/Curve/IntegralVacuum.lean}{41-L86}{The common scalar and its tensor powers}.}

\paragraph{Conjugating the normalized Taylor series.}
\teamanchor{proof.weighted-cm-jet-transport}%
Use the marked elliptic curve and the standard Fock coordinates of
\hyperlink{proof.weighted-vacuum-projection}{the vacuum construction}.
Write \(\eta=(2\pi/\sqrt d)^{1/2}\), \(\nu=N_1\), and \(A_\sigma\)
for the area of the marked conjugate elliptic lattice \(\Lambda_\sigma\)
in its algebraic differential coordinate. For \(p\in\mathbf C^\times\)
write \(T_pF\) for the Taylor series of \(\Theta_F(pZ)\).
We use the scale
\[
 \mathrm{per}=\frac{s_0}{u_0}\frac{\eta}{\Omega_0},\qquad
 s_0^f=\chi_{\mathrm{fin}}(\alpha)^{-1},\qquad
 A_\sigma|\sigma(u_0)|^2=\pi e^{\ell_d}\frac{2\pi}{\sqrt d}.
\]
Here \(\chi_{\mathrm{fin}}=\upsilon_{\mathrm{fin}}|\cdot|_f^{-1/2}\)
is the finite multiplier character of the Weil model. The root \(s_0\)
and the area equalizer \(u_0\) are chosen once. These are the
normalizations of the weighted jet at every stage. We give the area
and transport calculations explicitly.

The ideal class group of \(K\) has odd order \(h_K\) by genus theory.
The conjugate \(j\)-invariants over \(K\) run through its \(h_K\)
ideal classes by the maximal-order theorem of complex multiplication.
Choose a marked conjugate curve for each class \(c\), with lattice
\(\Lambda_c\) and discriminant \(\Delta_c\) in the differential
coordinate of its normalized cubic. The good-reduction argument
for the base cubic applies to its conjugates. In particular each
\(\Delta_c\) is a unit at all finite places. Write
\(\Lambda_c=u_c\mathfrak d_c\Lambda\), with
\(\mathfrak d_c=\mathfrak a_{c^2}\mathfrak a_c^{-1}\) for a system of
ideal-class representatives \(\mathfrak a_c\), and \(\Lambda\) the
base lattice. These ideals represent the classes \(c\) and
\(\prod_c\mathfrak d_c=\mathcal O_K\).
For \(v\in\mathfrak d_c\) the scalar \(u_cv\) is the differential of
an isogeny to the \(c\)-curve. For \(v'\in\mathfrak d_c^{-1}\) the
scalar \(u_c^{-1}v'\) is the differential of an isogeny in the reverse
direction. These differentials are algebraic and integral on the
specified good models at every finite place, by extension of
isogenies to the elliptic schemes. Taking finite sums of products
of the \(v\)'s, and separately of the \(v'\)'s, proves that
\(\prod_c u_c\) is a unit at every finite place. Applying \(\sigma\)
to the same isogenies gives both lattice inclusions, hence
\(\Lambda_c^\sigma=\sigma(u_c)\mathfrak d_c\Lambda_\sigma\).
The discriminant-normalized area of any lattice \(L'\) is
\(|\Delta(L')|^{1/6}\operatorname{area}(L')\).
It depends only on the shape. The ideal-class mean
\hyperlink{proof.weighted-cm-area-mean}{area calculation} below is the
same Kronecker limit calculation as
\eqref{eq:cm-elliptic-area-mean} in the earlier route.
Consequently any root \(u\) of
\[
 u^{12h_K}=\prod_c u_c^{12}\Delta_c
\]
has the required area property and is a unit at every finite place.
The comparison uses these properties for the already fixed \(u_0\)
and the multiplier identity for the already fixed \(s_0\) in
\(\mathrm{per}\).

\teamanchor{proof.weighted-cm-area-mean}%
In detail, for a basis proportional to \(1,\tau\), with \(y=\Im\tau\),
the discriminant-normalized area divided by \(\pi\) is
\(4\pi y|\eta_{\rm D}(\tau)|^4\), where \(\eta_{\rm D}\) is Dedekind's
eta. For the ideal lattice of this shape the partial ideal zeta
function is
\[
 \frac{(2/\sqrt d)^z}{|\mathcal O_K^\times|}
 \sum_{(j,k)\ne(0,0)}\frac{y^z}{|j\tau+k|^{2z}} .
\]
Kronecker's first limit formula gives its ratio of constant term
to residue at \(z=1\) as
\(\log(2/\sqrt d)+2\gamma-2\log 2
-2\log(\sqrt y|\eta_{\rm D}(\tau)|^2)\).
The residues are the same in every class and the sum is
\(\zeta(z)L(z,\chi_d)\). The mean log normalized area divided by
\(\pi\) is therefore \(\ell_d+\log(2\pi/\sqrt d)\).
Also, for a finite idele \(a_f\), choose \(k>0\) such that the
fractional ideal of \(a_f^k\) is principal, say generated by \(a\).
The splitting character is unramified at the finite places and
its archimedean component is \(\prod_\tau\tau(a)/|\tau(a)|\).
Triviality on \(E^\times\) and the product formula give
\(\chi_{\mathrm{fin}}(a_f)^k=\prod_{\tau\in\Phi}\tau(\bar a)\in K\).
Thus the modulus of \(\chi_{\mathrm{fin}}(a_f)\) is preserved under
all \(\sigma\) fixing \(K\). Since
\(\mathfrak a_0\overline{\mathfrak a_0}=\mathfrak D_{E/K}^{-1}\),
\(|\alpha|_f^{1/f}=\operatorname{rd}_E/\sqrt d\), so
\(|\sigma(s_0)|^2=\operatorname{rd}_E/\sqrt d\).

We next identify the full Taylor series after conjugation. The
following elliptic calculation is the transport of normalized
torsion germs used in the earlier fitted kernels. For the odd
reduced sigma section \(f_\Lambda\) with \(f_\Lambda'(0)=1\), the
exponent subtracted from the logarithm of the usual sigma function
is \(s_2 z^2/2\), where
\(\eta_\Lambda(\lambda)=s_2\lambda+\pi\bar\lambda/\operatorname{area}(\Lambda)\)
is the zeta quasiperiod. The de Rham class
\([\wp(z)\,dz]+s_2[dz]\) has periods proportional to
\(\bar\lambda\). It is the eigenspace of a nonreal CM endomorphism
for the eigenvalue on antiholomorphic differentials. In the
algebraic de Rham basis \([\wp(z)\,dz],[dz]\) this determines
\(s_2\) algebraically and commutes with \(\sigma\) fixing \(K\).
The Taylor coefficients of the usual sigma function follow from
the algebraic Weierstrass coefficients and
\((\log \sigma_\Lambda)''=-\wp\), with the odd,
derivative-one normalization. Thus \(f_\Lambda\) transports
coefficientwise.

For a torsion lift \(x\) choose \(m_0\ge2\) such that \(m_0x/2\)
is a period. Set \(u=x/(2m_0)\), \(M_0=m_0^2\), and let \(U_x\)
be reduced translation. The functions
\(D_j(z)=f_\Lambda(jz)/f_\Lambda(z)^{j^2}\) are algebraic elliptic
functions by their prescribed torsion divisors and leading
coefficient \(j\). If \(u\) is a period, \(x\) is an even period
and acts trivially. Otherwise the multiplier, oddness, and
\(M_0u\in\Lambda\) give, with \(B_u=(U_u f_\Lambda)^{4M_0}\),
\[
 \begin{split}
 U_xf_\Lambda(2m_0 z)&=B_u(z)D_{2m_0}(z+u),\\
 B_u(z)^{M_0}&=f_\Lambda(2M_0z)/D_{2M_0}(z+u),\\
 B_u(0)&=-D_{M_0-1}(u)/D_{M_0+1}(u).
 \end{split}
\]
The second quotient is a unit series. The indicated constant
selects its unique \(M_0\)-th root, so translating and conjugating
the series agrees once the deeper torsion lifts are tracked.
The even section is \(U_{t_0}f_\Lambda/(U_{t_0}f_\Lambda)(0)\)
for the marked half-period \(t_0\). Its conjugate uses the
transported marked point. The composition scalar of two reduced
translations is \(\exp(\pm\pi i E_\Lambda(x,y))\) for their alternating
Riemann pairing \(E_\Lambda\), with the sign determined by the order
of the two translations. On rational lifts it is a Weil pairing
at a sufficiently divisible torsion order, with one argument
possibly an integral multiple of the chosen smaller torsion label.
For example if \(mx,my\in\Lambda\), the pairing of
\(x/(2m)\) and \(y/(2m)\) at order \(2m^2\) gives this scalar or its
inverse. These deeper lifts can be tracked together with the
division-polynomial arguments. Choosing a different period lift of
the marked point multiplies \(U_{t_0}f_\Lambda\) by a constant
which cancels in its normalization.

For each \(\sigma\) fixing \(K\), choose
\(\Lambda_\sigma=\lambda_\sigma\mathfrak b_\sigma\) with
\(\mathfrak b_\sigma\) a fractional ideal of norm one.
Such a representative is \(\mathfrak c/\overline{\mathfrak c}\)
since every ideal class is a square. Thus
\(|\lambda_\sigma|^2=2A_\sigma/\sqrt d\).
The action on torsion commutes with \(\mathcal O_K\), as the
differentials of the CM endomorphisms are fixed by \(\sigma\).
The Tate modules of \(K/\mathcal O_K\) and \(K/\mathfrak b_\sigma\)
are the respective completions of these two rank-one lattices.
Their induced isomorphism is multiplication by an idele
\(\beta_\sigma\in\mathbb A_{K,f}^\times\) with
\(\beta_\sigma\widehat{\mathcal O}_K=\widehat{\mathfrak b}_\sigma\).
Use representatives \(v\in K\) as close to \(\beta_\sigma\)
as needed for any finite list of deeper torsion lifts. This
represents the image of \(\Omega_0 x\) by \(\lambda_\sigma v x\)
simultaneously for that list. Compatibility with the Weil
pairings gives \(\beta_\sigma\bar\beta_\sigma
=\chi_{\mathrm{cyc}}(\sigma)\) as an element of
\(\widehat{\mathbf Z}^{\,\times}\). Indeed, in the norm-one ideal
frame the alternating pairing on rational lifts is the base
pairing multiplied by \(v\bar v\), and the equality of Weil
pairings at all orders determines its residue units at every
rational prime.

For standard rank \(r_0\), use the product \(S(Z)\) from the
vacuum construction in all \(r_0\) coordinates. Let \(F\) be
its vector. The finite translation projection expresses the
vacuum as a finite combination of translates of \(F\).
These translates therefore span the finite part.
Put
\[
 S_\sigma(Z)=
 \prod_{k=1}^{r_0}\prod_{i=0}^{\nu-1}
 \theta_\sigma\left(
     \frac{\lambda_\sigma}{\nu\eta}
         \sum_{\tau\in\Phi}Z_{\tau,k}\overline{\tau(\zeta^i)}
              \right)^\nu .
\]
Choose a positive integer \(e\) with \(e\in\mathfrak b_\sigma\).
On \(\nu e L\), in each copy of the standard lattice \(L\), the
arguments are periods. The Hermitian form is the standard form
by the same Fourier orthogonality since
\(A_\sigma=|\lambda_\sigma|^2\sqrt d/2\).
The multiplier computation for the even marked section uses the
same two idempotent lines modulo \(2\) and gives the canonical
multiplier on \(\nu e L\). As for \(S\), the finite translation
basis on this lattice realizes \(S_\sigma\) as \(\Theta_{F_\sigma}\)
on the same standard finite part.

At a rational finite argument \(\alpha^{-1}\gamma\) the
translation of \(S\) translates its \(i,k\)-th elliptic factor
by \(\Omega_0 x_{ik}\), with
\(x_{ik}=-\operatorname{Tr}_{E/K}(\gamma_k\bar\zeta^i)/\nu\).
Applying \(\mathrm{per}\) to the complex variables leaves the
linear coefficients \((s_0/u_0)\overline{\tau(\zeta^i)}/\nu\).
Their conjugates are the coefficients for \(S_\sigma\) at
\[
 t_\sigma=\sigma(s_0/u_0)\eta/\lambda_\sigma,
 \qquad
 |t_\sigma|^2
    =\frac{\operatorname{rd}_E}{\sqrt d}
       \frac{\pi}{A_\sigma|\sigma(u_0)|^2}=\Omega^{-1},
\]
after the permutation \(\tau\mapsto\sigma\circ\tau\).
Take \(v\) representing all the elliptic translates and their
deeper torsion arguments and sufficiently close to \(\beta_\sigma\)
that the finite translates of \(F_\sigma\) at
\(\alpha^{-1}v\gamma\) and at \(\beta_\sigma\alpha^{-1}\gamma\)
coincide. This is possible by local constancy. The conjugation
of the elliptic germs then gives
\[
 \operatorname{rename}_\sigma\bigl(\sigma T_{\mathrm{per}}
       \varrho(\operatorname{finVec}x)F\bigr)
 =T_{t_\sigma}\varrho(\operatorname{finVec}(\beta_\sigma x))F_\sigma .
\]
The equality holds at all finite arguments \(x\) by density of
the rational arguments and local constancy on both sides.
Taylor series at nonzero scale are injective on finite vectors,
by entire continuation and rational translation covariance.
The formula thus extends by semilinearity to a map \(L_\sigma\)
on the finite part, with
\(\operatorname{rename}_\sigma(\sigma T_{\mathrm{per}}G)
=T_{t_\sigma}L_\sigma G\).
It respects the finite translations through \(x\mapsto\beta_\sigma x\).
Indeed the Heisenberg cocycle and the finite normalization
phases are finite trace-character values, and
\(\sigma(\psi_f(z))=\psi_f(\beta_\sigma\bar\beta_\sigma z)\).
The covariance follows first on the translates of \(F\), then
on their span.

\paragraph{The moved vacuum and its entire majorant.}
\teamanchor{proof.weighted-cm-moved-vacuum}%
Write \(M_f=\widehat{\mathcal O}_E^{\,r_0}\) in the standard finite
coordinates, \(v_0\) for the unit vacuum and \(T(x)\) for the
vacuum-normalized finite translations. For \(g\) finite unitary
the common fixed line of \(T(x)\), \(x\in gM_f\), is
\(\mathbf C\omega(g)v_0\). This follows from the corresponding
statement for \(M_f\), where averaging \(T(x)\) is the vacuum
projector, and from \(T(gx)\omega(g)=\omega(g)T(x)\).
Write \(\beta_\sigma=z a_f\) with \(z\bar z=1\) and
\(a_f\in\widehat{\mathcal O}_E^\times\). At a nonsplit pair the
valuation of \(\beta_\sigma\) is zero, and at a split pair its
two valuations are opposite, so such a factorization is
obtained by retaining just these valuations in \(z\).
With the same symbol for the central unitary scalar matrix,
\(L_\sigma v_0=c_\sigma\omega(z)v_0\) with \(c_\sigma\ne0\).
For every \(g\),
\(L_\sigma\omega(g)v_0=c_{\sigma,g}\omega(gz)v_0\).
To compare the scalars, the pairing against \(v_0\) on
\(T(x)v_0\) is \(1_{M_f}(x)\). Pairing the image against
\(\omega(z)v_0\) gives \(c_\sigma 1_{M_f}(x)\).
By semilinearity this compares the pairings on the entire finite
part. Apply it to \(\omega(g)v_0\). The coefficient
\(\langle v_0,\omega(g)v_0\rangle=1/\overline{\Delta_f(g)}\)
is nonzero, where \(\Delta_f(g)=\chi_{\mathrm{fin}}(\det u_f)\)
and \(u_f M_f=M_f\cap g^{-1}M_f\).
For this normalized Weil formula, one averages over the
translates \(T(gx)v_0\) indexed by
\(M_f/(M_f\cap g^{-1}M_f)\). The averaging scalar is
\(\Delta_f(g)/|\Delta_f(g)|^2\), and only the identity coset
survives in the pairing. Both \(\Delta_f(g)\) and
\(\overline{\Delta_f(g)}\) have unchanged modulus under \(\sigma\)
by the principal-ideal calculation above. Thus
\(|c_{\sigma,g}|=|c_\sigma|\).

Let \(v_{\rm int}=c_s^{r_0}v_0\) be the scaled vacuum already
chosen in the definite chain. The constant term of the
transport, after finite translation, is the conjugate value
of \(v_{\rm int}\) at the corresponding finite point. Because
\(\beta_\sigma\) is invertible, these values bound
\(L_\sigma v_{\rm int}\) at all finite arguments by
\(\exp(C_{\rm vac}\nu^2 r_0)\).
The squared Hilbert norm of a theta vector is its normalized mean
squared value on the adelic vector quotient by the rational
vectors. In a Schr\"odinger polarization this follows by integrating
the squared rational summation first in the dual translation
variable, which kills the cross terms, then in the remaining
translation variable. Absolute values descend to this quotient,
and the finite-only arguments have dense image by archimedean
approximation with rational vectors. Hence the Hilbert norm
is bounded by the finite supremum. Since the unit vacuum
and its unitary moves have Hilbert norm one,
\(|\sigma(c_s^{r_0})c_\sigma|\le \exp(C_{\rm vac}\nu^2 r_0)\).

For a unit moved vacuum one has the uniform pointwise bound
\[
 \sup_y|(\omega(g)v_0)(y)|\le \exp(O_d((f r_0+1)^2)).
\]
Here is a proof retaining the measure and lattice of this model.
Set \(w=\varrho(y)\omega(g)v_0\). The squared value at \(y\)
is \(\sum_{\gamma\in E^{r_0}}\langle w,\varrho(\gamma)w\rangle\)
with the prescribed rational embedding. In a rational
Schr\"odinger polarization this identity is Poisson summation
in the dual rational variable of the matrix coefficient, followed
by rational summation in the other variable. The self-dual
measures give quotient mass one. The translations by \(y\)
affect the coefficients only by phases. Unitarity and the
vacuum matrix coefficient bound their moduli by
\[
  1_{\{\alpha^{-1}g^{-1}\gamma_f\in M_f\}}\,
  \exp\left(-\frac{\pi}{\sqrt d}
                \sum_{\tau,k}|\tau(\gamma_k)|^2\right).
\]
Two distinct vectors in the indicated lattice have distance
at least one in the unscaled complex coordinates. For their
difference \(\gamma_0\), the finite conditions, unitarity,
and \(\alpha\bar\alpha\widehat{\mathcal O}_E
=\mathfrak D_{E/K}^{-1}\widehat{\mathcal O}_E\)
put \(\sum_k\bar\gamma_{0,k}\gamma_{0,k}\) in the inverse
relative different. Its trace to \(K\) is a positive rational
integer equal to the squared distance. In real dimension
\(D_0=2fr_0\), a unit-separated set has at most
\((2R+1)^{D_0}\) points in a ball of radius \(R\). Bounding
\(e^{-c r^2}\le e^{(D_0+1)^2/(4c)} e^{-(D_0+1)r}\) and summing
over integer shells proves the bound.

Consequently the entire theta of the transported vector satisfies
\begin{equation}
 \label{eq:weighted-cm-envelope}
 \left|\Theta_{L_\sigma\varrho(\operatorname{finVec}x)
                      \omega(g)v_{\rm int}}(Z)\right|
 \le \exp\left(C_d(\nu r_0+1)^2+\frac{\|Z\|^2}{2}\right).
\end{equation}
One can take \(C_d=C_{\rm vac}+2d+6\). In the squared
moved-vacuum bound the count over integer shells is at most
\(3^{D_0}(1-e^{-1})^{-1}\) after the factor
\(\exp((D_0+1)^2\sqrt d/(4\pi))\) is removed.
Taking the square root and using \(f\le\nu\) gives the
term \(2d+6\) in this choice. Include \(C_d\) in the fixed
constant of the scale condition \(10^{12}C_0\le\log\log b\).
The factor \(\exp(\|Z\|^2/2)\) removes the unit real Gaussian
from the theta value. The finite translations move the
evaluation point and introduce a phase. Each extra determinant
character in a definite pair is algebraic of its prescribed
weight, with modulus one after conjugation. To see this on a
finite norm-one element, pass to a power whose ideal is
principal, generated by \(a\). The product \(a\bar a\) is a unit.
On doubling that power one can therefore use \(a/\bar a\),
a rational norm-one generator, and an integral compact factor.
The compact character value is a
root of unity by smoothness, and the rational norm-one point
contributes only powers of conjugates of complex modulus one.
Thus \eqref{eq:weighted-cm-envelope} applies to each moved
vacuum on the standard Hom space of the definite pair.
For a complex automorphism not fixing \(K\), compose on its
output with complex conjugation before taking the norm.
The Weyl Gram matrix has real integral entries, so this
preserves its norm as well as the scalar moduli.

\paragraph{The integral terminal branching map.}
\teamanchor{proof.weighted-terminal-branching}%
Use \(n=n_m-1\), and let \(\beta_{\tau,h}\) be the number of
kept columns of upper height \(h\). Let \(\gamma_{\tau,h}\)
be the remaining upper columns, whose height drops by one.
Thus the lower height-\(h\) count is
\(\mu_{\tau,h}=\beta_{\tau,h}+\gamma_{\tau,h+1}\).
For the weighted shape \(\beta_{\tau,h}=0\) for \(h<\kappa\)
and for \(h=n\). The second assertion follows since the
positive kept heights are \(\kappa\bmod t+1\), whereas
\(n=\kappa+m(t+1)-1\). The determinant exponents on
the two blocks therefore agree. Write \(J\) for the total
number of lower boxes in heights \(h\ge\kappa\).
There is a nonzero integer branching matrix in the evaluation
coordinates, with Weyl operator norm at most
\(\max(1,J)^{3J}\). We construct it at the outset so the same
matrix enters the seesaw and the two length and index estimates.

We recall the evaluation-coordinate norm. For a shape \(\nu'\)
let \(b_i(X)\) be the standard products of flag minors, and
write \(e_{\nu'}(X)=(b_i(X))_i\). Integral straightening expands
every ordered product of flag minors of this shape with integer
coefficients in the \(b_i\). In the tensor product of exterior
powers with its standard orthonormal exterior bases, the
coordinate vector \(w_{\nu'}(X)=\bigotimes_{\text{columns }h}
(X_1\wedge\cdots\wedge X_h)\) is therefore \(C_{\nu'}e_{\nu'}(X)\)
for an integer matrix \(C_{\nu'}\). It contains the identity
rows indexed by the standard products. In particular its
integral image in the tensor space is saturated. The kernel
of the Weyl form is
\[
 w_{\nu'}(X)^{\mathsf T}w_{\nu'}(Y)
   =e_{\nu'}(X)^{\mathsf T}G_{\nu'}e_{\nu'}(Y),
 \qquad G_{\nu'}=C_{\nu'}^{\mathsf T}C_{\nu'},
\]
as follows also by Cauchy--Binet. This identifies the tensor
metric on the generated irreducible with the Weyl metric used
on the definite values. The highest evaluation \(X=1\) has
norm one. The associated polynomial of \(\xi\) is
\(e_{\nu'}(X)^{\mathsf T}G_{\nu'}\xi\).
The normalized Fischer form on this polynomial model agrees
with the Weyl form by unitary invariance and the highest vector
normalization. If \(H_{\nu'}\) is the squared Fischer norm of
the highest product of minors, this shows that the integral
matrix of Fischer pairings of the \(b_i\), multiplied by
\(G_{\nu'}\), is \(H_{\nu'}\) times the identity. Consequently
\(H_{\nu'}\) times any integral product of minors is the
associated polynomial of an integral Weyl vector. For a diagram
with \(Q\) boxes the monomial factorial is at most \(Q!\).
The sum of absolute highest-product coefficients is at most the
product of the column-height factorials, also at most \(Q!\).
This gives \(H_{\nu'}\le (Q!)^3\).

At each place insert the last unit vector \(e_*\) at the front
of every shifted column and include the kept columns in the
upper exterior space. Tensoring gives an isometry \(I\) of
the lower tensor space into the upper tensor space. It is
equivariant for \(\operatorname{diag}(g,s)\) with the extra
character \(s^{\sum_h\gamma_{\tau,h}}\).
Its bilinear kernel on evaluation vectors is
\[
 K(X,Y)=\langle Iw_\mu(X),w_{\rm upper}(Y)\rangle_{\rm bil}.
\]
It is the value at \(\operatorname{diag}(X^{\mathsf T},1)Y\)
of the product of upper minors with the added row first in
the shifted columns and last in the kept columns. This is the
same product of minors as in the earlier terminal highest
vector, up to the order signs of the shifted rows.
It is a nonzero polynomial at \(X=1\).

For a matrix \(Y\) with last row \((y,0,\ldots,0)\), expand
each shifted minor in that row. Apart from the monomial
\(y^{\sum_h\gamma_{\tau,h}}\), the result at this place is the
lower flag product at \(Y_{\rm old}^{\mathsf T}X\), where the
height-\(h\) shifted rows are \(2,\ldots,h+1\) and the kept
rows are \(1,\ldots,h\). Here \(Y_{\rm old}\) is the \(n\)-by-\(n\)
block in the old rows and the first \(n\) columns.
Set \(p=\kappa-1\). After permuting the rows to put
\(2,\ldots,p+1\) first, each column of height \(h\) contains its
first \(\min(p,h)\) standard rows. The remaining boxes of this
lower diagram form a diagram \(\mu^\flat\) on the last
\(n-p\) rows. Its total size over all places is at most \(J\).

To verify that the denominator for this product is controlled
by the remaining boxes, use the tensor realization. In each
long column wedge the first \(p\) standard vectors onto the
tail column, and in a short column insert its usual highest
exterior vector. This is an integral isometry of the tail
tensor space into the lower tensor space. The lower Weyl
space intersects this inserted tensor space in exactly the
inserted tail Weyl space. Indeed the prescribed tensor basis
vectors, with these leading rows, are the simultaneous
weight space of diagonal actions on the first \(p\) rows for
the weights of the highest vector on these rows. The first
row must appear in every column. Inductively the next row
must appear in every column tall enough to contain it. This
gives exactly the indicated leading rows. The projection
to these weights sends \(w_\mu(X)\) into the image of the
tail Weyl space. For \(X\) with nonsingular first \(p\)-by-\(p\)
block this follows by subtracting combinations of its first
\(p\) columns from later columns to clear the upper \(p\)
entries. The initial-column wedges are unchanged. Their
projections are the leading minor scalars times the wedges
of initial columns of the lower-right Schur complement.
These \(X\) form a dense set. Conversely the image of the
tail Weyl space is generated by the lower block orbit of the
lower highest tensor. Thus orthogonal projection of an
inserted exterior basis vector to the lower Weyl space
equals the insertion of its projection to the tail Weyl space.
The integral lattices in both Weyl spaces are saturated in
their tensor spaces by the identity rows of the \(C_{\nu'}\).
The preceding Fischer identity for the tail diagram now
shows that \(c=H_{\mu^\flat}\) times the lower flag product
is the associated polynomial of an integral lower Weyl
vector. The row permutation and substitution by
\(Y_{\rm old}^{\mathsf T}\) act integrally on this polynomial
model. In the tensor description the transpose substitution
acts on the vector through the exterior powers of
\(Y_{\rm old}\). It preserves the generated Weyl space and has
integer polynomial coefficients in the ambient tensor basis.
Saturation, with the coordinate left inverse given by the
identity rows, gives the same integrality in the Weyl model.
We have \(c\le\max(1,J)^{3J}\).

For unrestricted generic \(Y\), multiply the first column by
its last entry \(y\), and replace column \(i>1\) by
\(yY_i-Y_{*,i}Y_1\). These are integral upper triangular
column substitutions.
The bottom row is now zero off the first column. Every flag
minor of height \(h\) has been multiplied by \(y^h\).
It follows that a monomial in the \(y\)'s times \(cK(X,Y)\)
has every coefficient in \(Y\) in the integral lower Weyl
polynomial lattice in \(X\). Comparing coefficients with
their exponents shifted by this monomial removes the factor
without introducing any denominator. Finally, the coefficient
vectors of the \(e_{\rm upper}(Y)\) span the full upper
coordinate lattice over \(\mathbf Z\). For a diagonal term
order the leading monomials of the standard products are
distinct with coefficient one, giving a unitriangular
coefficient submatrix. It follows that \(cI^*\) restricted
to the upper Weyl space, with \(I\) restricted to the lower
Weyl space in the adjoint, has integer matrix \(B\) in the
evaluation coordinates. Its norm is at most \(c\).

The last-coordinate character, including the common
determinant shift \(e_m(\tau)\), is
\(s_\tau^{b_\tau}\), \(b_\tau=\sum_h\gamma_{\tau,h}+e_m(\tau)\).
The ordinary multiplicity-one branching rule follows by
removing the entries \(n+1\) from a semistandard tableau of
the upper partition, leaving a horizontal strip. This proves
that the constructed nonzero map spans the branching line
used in the terminal period. Use one such matrix in the
restriction \(\prod_\tau\tau(\bar\pi)^{-b_\tau}B\).
Changing the generator once changes the whole restricted
definite chain by one common scalar, absorbed in \(c_*\) in
\eqref{eq:weighted-seesaw}, retaining the tensor
\(\Psi_\vartheta\) for each twist. This also holds for the
paired-chain expression, since both expressions apply the
restriction to the same terminal form.
At a coefficient embedding above \(p\), the finite
terminal gauge uses the inverse upper representation of
the frame \(\operatorname{diag}(g,s_\tau)\).
In the resulting value \(B\) has matrix norm at most one
over the valuation ring and contributes the last-coordinate
factor \(s_\tau^{-b_\tau}\).
Off \(N\) the \(\tau(\bar\pi)\)'s are units and the ordinary
frame uses \(s_\tau=1\). At \(N\) the adapted upper frame has
\(s_\tau=N\) when the place induced by \(\tau\) contains
\(\bar\pi\), and \(s_\tau=1\) otherwise.
Thus \(|\tau(\bar\pi)|_p=|s_\tau|_p\), and the restriction's
normalizing character agrees in absolute value with the
factor in the upper gauge. The lower gauge is therefore
bounded by the upper gauge at that adapted frame. This is
valid for each embedding, including its composition with
conjugation.

At a top, the \(a\) height-one columns all shift, and
\(e_m(\tau)\ge-w_m(\tau)\ge-2b\) at every place, so
\(b_\tau\ge a1_{\tau\in\mathcal T}-2b\).
Since \(\pi\bar\pi=N\), the restriction costs at most
\(-Ta\log N+2fb\log N+6J\log\max(1,J)\)
in the logarithm of the squared norm at every embedding.
Here \(J=O(fn_m^2r^8)\), since the large height-\(\kappa\)
shifted part moves to height \(\kappa-1\) and only the
kept part and the small higher columns remain at heights
at least \(\kappa\). Hence the final term is \(O(x^{70})\)
on the weighted range.

An \(m\)-fold twist \(\vartheta\) is a weight-zero character of
\(U(1)(\mathbb A_{E^+})\) whose tame support \(Q\) is disjoint from its
conjugate. Its norm \(q_\vartheta=\prod_{v\in Q}\mathrm N(v)\) is prime
to \(gN\), and the \(j\)-th power of its tame character is nontrivial
at every \(v\in Q\) for \(1\le j\le m\). The trivial twist has
\(Q=\varnothing\). Let \(\iota_\vartheta\) be the chain of \(m\) theta
lifts through \(Y_j\) with the corresponding twist powers and
\(\Theta_\vartheta=\iota_\vartheta|_W\) its restriction, realized at the
final level.\footnote{\weightedpin{grh-compact/src/New/Seesaw/IndexTypes.lean}{184-L199}{The twist and its level}
and \weightedpin{grh-compact/src/New/Exchange/LValueRatio.lean}{19-L43}{The powers used in the chain}.}

\paragraph{The seesaw identity.}
There is one scalar \(c_*\in\mathbf C^\times\) on the chosen datum \(D\),
and for each twist \(\vartheta\) there is a good finite datum
\(\Psi_\vartheta\), both independent of the ramified column. They give
\begin{equation}
\label{eq:weighted-seesaw}
  \iota_\vartheta\bigl(\mathcal P(c_*\phi)\bigr)
       = B^\mu_{\vartheta,\Psi_\vartheta}(\phi)
       \qquad(\phi\in\mathcal C_{\rm ram}).
\end{equation}
The support is in the admissible translations for \(Nq_\vartheta\),
the coefficients are integral at all coefficient places prime to
\(q_\vartheta\), and the total absolute value
of the conjugated coefficients is at most \(e^{x^{60}}\) at the
trivial twist, for every \(\sigma\in\operatorname{Aut}(\mathbf C)\).
The tensor is jointly invariant at the column levels with last
component the identity. Its chain also has the paired-chain expression
using the projections onto the local vacua at the places over \(N\).
These last two properties allow the definite estimates in Theorem~A.

\teamanchor{proof.weighted-seesaw-inputs}%
\paragraph{The functions to which the finite cuts apply.}
Let \(F=E^+\), write \(G_j=U(Y_j)\), and write \(B_j\subset G_j\)
for the Bessel subgroup of the fixed grading, with string length
\(\max(m-1-j,0)\). At the bottom let \(\eta\) be the inverse of
the fixed central lattice scalar in the embedded core. The bottom
level \(J_g\) in lattice coordinates is \(\eta K(g)\eta^{-1}\),
with \(K(g)\) the principal level in the filling chart. Put
\(J_j=J_j(\vartheta)\) for the level of
the lower lift at stage \(j\). Thus \(J_j\) is the kernel of the
chosen and twist-place reductions inside the lattice stabilizer
\(K_j^0\subset G_j(\mathbb A_{F,f})\). Write \(T_j\) for
the unnormalized lower theta lift, with the harmonic at infinity and
the finite inclusion datum of the lower edge. The actual arrow is
\(L_j=[K_j^0:J_j]T_j\). Use Haar measure of mass one on \(K_j^0\).

Let \(P_0\) be the periodization of the seed of \(\phi\) over
\(B_0(F)\). The seed in a class and Bessel coordinate is the
prescribed finite coset indicator of \(J_g\) times the full real
Bessel coordinate and the class cutoff \(\beta_c\). The Poincar\'e series
\(\mathcal P\phi\) uses the same seed periodized over \(G_0(F)\),
with the class and Weyl-to-Bessel coordinates of the column
right-translated by the fixed central lattice scalar.
Each seed coset is a translate of \(J_g\) by a point of the
embedded definite core, including the central lattice scalar. Its
translates under \(B_0(F)\) therefore have finite support in
\(B_0(\mathbb A_{F,f})J_g\). On the core the rational unipotent
translates of the cutoff sum to one. For \(y_0\in J_g\) the fixed
bottom normal form reads
\[
 P_0\bigl(\jmath_0(k)(\xi_0,y_0)\bigr)=
  \mathsf A\bigl(\rho_{X_0}(k_\infty^{-1})\phi(k_f)\bigr)
                                      \cdot\mathsf W(\xi_0).
\]
Here \(\jmath_0\) embeds the core, \(\rho_{X_0}\) is its conjugate
Weyl action, \(\mathsf A\) is the fixed change to Bessel
coordinates, and the dot is bilinear with the product vector of
real Bessel functions \(\mathsf W\). The central factor in the
initial column coordinates cancels the inverse central factor in
the seed coset dictionary. This supplies the same initial
coefficient linear form for all twists.

\paragraph{The normalized real Bessel vectors.}
\teamanchor{proof.weighted-real-bessel-vectors}%
Fix a real place of \(F\), with chosen complex embedding \(\Phi_v\) of
\(E\). All the constructions in this paragraph are made at that place.
Write \(G_1\) for the unitary group on the normalized frame, with
\(\varrho\) hyperbolic pairs, a positive core \(X_0\) of rank \(k\) and a
last vector \(z\) of norm \(-1\). Put \(n=\varrho+1\) and
\(r=k+2\varrho+1\). Use the compact-dual Haar \(dy\) from the degree
calculation, probability Haar on \(H=U(X_0)\), and Lebesgue measure
\(du\) on the free entries of the flag unipotent \(U\). Use the
Euclidean-unitary sign frame \(p_v\), with \(p_v^*J_1p_v=J\),
where \(J_1\) is the normalized flag Gram matrix and
\(J=\operatorname{diag}(\varepsilon_i)\). Read the row coefficient
of \(y\) at \(p_v^{-1}yp_v\). Use \(ds=du\,dh\) on \(S=UH\).
Write \(d\dot y\) for the quotient Haar specified by \(dy=du\,d\dot y\).
Put \(\operatorname{size}(y)=\inf_{s\in S}(1+\|sy\|)\) in the operator norm.
We use
\(\langle a,b\rangle=a^*b\). Let \(\sigma\) be the sum of the first
superdiagonal entries on the frame \((e_0,\ldots,e_{\varrho-1},z)\), and put
\(\operatorname{wave}(c)=\exp(c-\bar c)\).
The basis scales at the frame entries are
\[
 s_{va}=\kappa_v^{\varrho-a}/\sqrt N,\qquad
 \kappa_v=g\Phi_v(\varpi)/(2\pi i)\in\mathbf R^\times .
\]
Put \(u=N_1/M\). Every prime dividing \(u\) divides \(M\), so
\(\Phi_{N_1}(X)=\Phi_M(X^u)\). The different generator satisfies
\[
 |\Phi_v(\varpi)|
    =\sqrt d\,u\,|\Phi'_M(\Phi_v(\zeta)^u)|.
\]
The derivatives on primitive \(M\)-th roots form a fixed
finite set of nonzero values. With \(u\) polynomial in \(x\)
and \(b^{20}\le g\le b^{22}\), this gives
\(1\le|\kappa_v|\le b^{23}\) uniformly for large \(x\).%
\footnote{\weightedpin{grh-compact/src/New/Data/Different.lean}
{346-L366}{The different generator and the cyclotomic derivative}.}
Conjugation by the inverse of the basis-scale matrix therefore divides
the rational simple-root sum at \(v\) by \(\kappa_v\).
Write \(\sigma_U\) for the sum before this scaling. In these coordinates
\(\operatorname{wave}(\sigma)\) is the local factor of
\(\chi_\infty(u)=\psi_\infty(\sigma_U(u)/(g\varpi))\).
For the standard adelic additive character write
\(\chi_f^\pm(u)=\psi_f(\pm\sigma_U(u)/(g\varpi))\).
On rational \(u\) the product formula gives
\(\chi_f^+(u)^{-1}=\chi_f^-(u)=\chi_\infty(u)\).
\footnote{\weightedpin{grh-compact/src/New/Filling/Setup.lean}{126-L132}{The basis scale},
\weightedpin{grh-compact/src/New/Filling/Setup.lean}{319-L324}{Conjugation in the normalized frame}, and
\weightedpin{grh-compact/src/Lib/DiscreteSeries/Bessel/RealMeasures.lean}{192-L211}{Haar in free entries}, and
\weightedpin{grh-compact/src/Lib/DiscreteSeries/Bessel/UnipotentCoords.lean}{233-L244}{The quotient measure}.}

Let \(m_\lambda(y)\) be the matrix coefficient of a unit highest vector
in the lowest compact type of the row discrete series at \(v\).
Here \(\lambda\) is its Harish--Chandra parameter in the ordered unitary
sign frame specified below.
On \(H\) the column type is the symmetric power of exponent
\(a_v=\kappa(v,1)\) in the column diagram. Its highest weight
\(\Lambda'\) has entry \(a_v\) at the designated coordinate and
zero elsewhere. The rank-\(k\) weight sums the columns of height strictly
less than \(k\). These have only the height-one contribution, and
the determinant exponent at stage zero is zero.
\footnote{\weightedpin{grh-compact/src/New/Curve/LiveSum.lean}{73-L85}{Column numbers and initial determinant exponent},
\weightedpin{grh-compact/src/Lib/General/Algebra/ZerosCostDegree.lean}{25-L28}{The diagram-weight sum}, and
\weightedpin{grh-compact/src/New/Filling/FillingData.lean}{56-L93}{Actual parameter and compact coefficient}.}

\paragraph{The row coefficient on the actual sign frame.}
\teamanchor{proof.weighted-row-coefficient}%
In that frame \(G_1=U(p,q)\), \(p=k+\varrho\), \(q=\varrho+1\).
Let \(K=U(p)\times U(q)\) with probability Haar. Order the diagonal
coordinates by decreasing \(\lambda_i\), retaining the sign on each
one. The first \(k-1\) entries, numbered from zero, are
\(k/2-1-i\). The remaining entries are \(-B_v-jh\),
\(1\le j\le r-k+1=2(\varrho+1)\), where \(B_v=w_0(v)-h_2\), \(h_2=r^8\) and
\(h=2h_2\). Here \(k\) is even and \(r\) is odd. The rounding
defining \(B_v\) gives \(B_v\ge a/\max(x,x_1)\ge k\) eventually,
since \(a\ge xb\), \(\log b=x^\varepsilon\),
\(\log x_1=x^{\varepsilon/6}\) and \(k\le r\le3x^5\).
The first block is positive for the form and
the signs on the second block alternate. In particular \(\lambda\)
is regular, \(\lambda_i\in(r-1)/2+\mathbf Z\), and opposite-sign
entries differ by at least \(h\).
\footnote{\weightedpin{grh-compact/src/New/Seesaw/IndexTypes.lean}
{124-L126}{The undoubled row parameter},
\weightedpin{grh-compact/src/New/Seesaw/IndexTypes.lean}
{279-L285}{Its indexing on the real row},
\weightedpin{grh-compact/src/New/Data/AdmissibleData.lean}
{129-L132}{The increments in the tail of the parameter}, and
\weightedpin{grh-compact/src/New/Exchange/SteepParameters.lean}
{10-L29}{Signs and lowest weight on the two blocks}.}

Take the upper triangular Borel for this order. Write \(\rho_G\) for
the half-sum of its roots and \(\rho_K\) for the compact half-sum.
Set
\[
 \chi=\lambda-\rho_G,\qquad \mu=\lambda+\rho_G,\qquad
 \Lambda=\chi+2(\rho_G-\rho_K).
\]
The entries of \(\chi\) are integers. They decrease weakly since
successive entries of \(\lambda\) differ by at least one.
Thus \(\chi\) is a unitary character of the diagonal compact torus,
dominant on the positive roots. The Harish--Chandra discrete series
\(\pi_\lambda\) has a single copy of the compact type of highest
weight \(\Lambda\). To apply the semisimple theorem, first restrict
the parameter and roots to \(SU(p,q)\). Its lowest type has the
restriction of weight \(\Lambda\). On \(SU(p,q)\times U(1)\) give
the circle the character \(z^{\sum_i\lambda_i}\).
The sum is the integer \(\sum_i\Lambda_i\), so for \(\zeta^r=1\)
the pair \((\zeta^{-1}I,\zeta)\) acts trivially, as one sees on the
lowest type. The product representation descends to \(U(p,q)\).
On the derived group its infinitesimal character is the restriction
of \(\chi+\rho_G\). In Sun's convention the inducing line is tensored
with the top exterior power of the upper nilradical before applying
the derived functor in degree \(\dim(\mathfrak n\cap\mathfrak k)\).
The parameter \(\chi+\rho_G=\lambda\) pairs strictly positively with
every root in that nilradical. The good-range theorem thus gives an
irreducible unitary module of infinitesimal character \(\lambda\).
Its bottom weight is \(\chi+2(\rho_G-\rho_K)=\Lambda\), as in Sun's
one-dimensional inducing example. Within a sign block, the difference
of the \(\rho_G-\rho_K\) entries at indices \(i<j\) counts the opposite-sign
indices between them, so \(\Lambda\) is compact dominant.
On the derived group both modules
now have the Casimir and compact type of the Vogan--Zuckerman
criterion, whose inducing Levi is the compact Cartan. Its split
direction is zero. Hence their Harish--Chandra modules agree there.
The central circle acts by the inducing character throughout, since
all roots kill the center. This identifies
\(\pi_\lambda\) with the unitary realization of
\(A_{\mathfrak b}(\chi)\) on \(U(p,q)\).
\cite[Theorems~(8.1) and (8.5)]{ar4-atiyah-schmid}%
\footnote{\href{https://www.math.utah.edu/~ptrapa/math-library/vogan/vogan-unitarizability-annals-1984.pdf\#page=2}
{Vogan, Theorems~1.2--1.3}, and
\href{https://www.numdam.org/item/CM_1984__53_1_51_0.pdf\#page=29}
{Vogan--Zuckerman, Proposition~6.1}, with the character and bottom
weight of their (5.1)--(5.2). For the induction convention, see
Sun~\cite[\S\S5.1 and 6.3]{real-branch-sun}.}

Let \(\tau_\Lambda\) denote this compact type, \(w\) its highest
unit vector, \(P_w=ww^*\) and \(d_\Lambda=\dim\tau_\Lambda\).
For an integral weight \(\gamma\) write
\[
 M_\gamma(Y)=\prod_{t\in\mathbf Z}
 \frac{\det Y_{\{i:\gamma_i\ge t\},\,\{i:\gamma_i\ge t\}}}
      {(\det Y)^{{\bf1}_{t\le0}}}.
\]
The empty minor is one. Only finitely many factors differ from one
on invertible matrices. On \(K_{\mathbf C}=\mathrm{GL}_p(\mathbf C)
\times\mathrm{GL}_q(\mathbf C)\),
\(M_\Lambda(a)=w^*\tau_\Lambda(a)w\). Indeed each sign block contributes
the coefficient of a highest wedge at every successive difference
of its dominant weight, times the appropriate determinant power.
Their tensor highest vector lies in the irreducible summand of weight
\(\Lambda\), and the product of these wedge coefficients is the stated
minor product.

For a matrix \(Y\) with positive definite Hermitian part put
\[
 \mathcal E_\lambda(Y)
   =d_\Lambda\int_K
     \frac{\tau_\Lambda(u)P_w\tau_\Lambda(u)^*}
          {M_{-\mu}(u^*Yu)}\,du .
\]
Every principal compression is accretive and invertible, so the
integrand is holomorphic in \(Y\), locally uniformly for \(u\in K\).
For \(Y>0\), identify \(Y\) with \(Y^{1/2}U(r)\).
On this whole cone Sun's Definition~3.2 and Theorem~6.3 give
\(d_\Lambda E_{\mathfrak b}(j_{\mathfrak b}\circ\psi_\chi)\).
Indeed on the complexified torus modulo its compact part the inducing
function is \(\psi_\chi(tT)=\chi(tt^*)\), by the torus action and
its value one at the identity. Its embedding \(j_{\mathfrak b}\)
sends the identity of that line to \(P_w\). The factor
\(\delta_{\mathfrak n}^+\) from (9) multiplies this function by
\(\delta_{\mathfrak n}(tt^*)\). The torus coordinate \(tt^*\) of
\(u^{-1}Y^{1/2}U(r)\) is the positive diagonal \(D\) in
\(u^*Yu=N D N^*\), \(N\) upper unitriangular. Its trailing principal
determinants are the trailing products of \(D_i\). Thus
\((\delta_{\mathfrak n}\chi)(D)=D^\mu=1/M_{-\mu}(u^*Yu)\),
including the full determinant factor when \(\det Y\ne1\).
This is the same torus calculation as in the proof of
\eqref{eq:real-branch-lowest-decay}. The torus bottom space is the line
\(\mathbf Cw\), and (26) has factor \(d_\Lambda\) for probability Haar.
For a positive Hermitian \(z\in G_1\), take \(x=z^{1/2}\in\exp\mathfrak p\)
there. The result is the compression of \(\pi_\lambda(z)\) to
\(\tau_\Lambda\). Compact left and right translation now give, for
\(p_1,q_1\in K\),
\begin{equation}
 \label{eq:weighted-row-coefficient-polar}
 m_\lambda(p_1zq_1^{-1})=
 \frac{\displaystyle\int_K
   M_\Lambda(p_1u)\overline{M_\Lambda(q_1u)}
   \,M_{-\mu}(u^*zu)^{-1}\,du}
 {\displaystyle\int_K|M_\Lambda(u)|^2\,du}.
\end{equation}
The denominator equals \(d_\Lambda^{-1}\) by compact Schur
orthogonality. This gives \(m_\lambda(g)=w^*\pi_\lambda(g)w\) on
all of \(G_1\), including \(m_\lambda(1)=1\). It is of positive
type, since the matrix with entries \(m_\lambda(g_i^{-1}g_j)\)
is the Gram matrix of the vectors \(\pi_\lambda(g_j)w\).
\cite[Definition~3.2, Theorem~6.3 and (25)--(26)]{real-branch-sun}

The same integral gives the local formula needed for expansion.
For \(Y>0\), equation (5) and Theorem~3.3 of Sun give
\[
 \mathcal E_\lambda(aYa^*)=
 \tau_\Lambda(a)\mathcal E_\lambda(Y)\tau_\Lambda(a)^*
 \qquad(a\in K_{\mathbf C}).
\]
Indeed \(a\) acts by congruence on the positive Hermitian model of
\(\mathrm{GL}_r(\mathbf C)/U(r)\). On the compact complexification one has
\(\tau_\Lambda(a)^*=\tau_\Lambda(a^*)\).
Both sides of the resulting identity with independent variables
\((a,Y,b)\) are holomorphic where \(Y,aYb\) are accretive.
Near \((1,1,1)\), the conditions \(Y=Y^*\), \(b=a^*\) describe a
maximal totally real submanifold, and its positive part contains
a relative neighbourhood of that point. Uniqueness of holomorphic
continuation gives
\[
 \mathcal E_\lambda(aYb)=
     \tau_\Lambda(a)\mathcal E_\lambda(Y)\tau_\Lambda(b)
\]
throughout the connected product of operator-norm balls
\(\|Y-1\|<1/2\), \(\|a-1\|,\|b-1\|<1/16\) in the indicated
spaces. There \(\|aYb-1\|\le(17/16)^2(3/2)-1<1\).
In particular \(w^*\mathcal E_\lambda(Y)w\) is unchanged by a
lower unitriangular left factor and an upper unitriangular right
factor within these bounds and within the two sign blocks. The
upper factor fixes \(w\), and the adjoint of the lower factor
does so as well.

This local expression agrees with the splitting used to define the
coefficient. To see the agreement on the whole real group, write
\(g=p_1zq_1^{-1}\) as above and set \(J=\operatorname{diag}(\varepsilon_i)\).
The same-sign diagonal part of \(g\) is
\(a=p_1S_0q_1^{-1}\), where \(S_0=(z+JzJ)/2=(z+z^{-1})/2\).
Hence \(S_0\) is positive block diagonal and commutes with \(z\).
Put \(Z_0=S_0^{-1}z>0\). The unitary matrices \(S_0^{it}\) fix
\(Z_0\) by conjugation. Equivariance on the positive cone shows that
\(\mathcal E_\lambda(Z_0)\) commutes with their action, and hence
with \(\tau_\Lambda(S_0^{1/2})\) by differentiation. Equivariance
under \(S_0^{1/2}\) gives
\(\mathcal E_\lambda(z)=\tau_\Lambda(S_0)\mathcal E_\lambda(Z_0)\).
It follows that
\[
 m_\lambda(g)=w^*\tau_\Lambda(a)\mathcal E_\lambda(a^{-1}g)w.
\]
For \(\|g-1\|<1/16\), diagonal and off-diagonal block extraction
are contractions in operator norm, by averaging \(g-1\) and
\(J(g-1)J\) with either sign. Thus \(\|a-1\|<1/16\) and
\(\|a^{-1}g-1\|<1/15\). The polarized identity then gives
\(m_\lambda(g)=w^*\mathcal E_\lambda(g)w\).
\footnote{\weightedpin{grh-compact/src/Lib/DiscreteSeries/FlenstedJensen/RepresentationDefs.lean}
{95-L122}{The weights, compact integral, specified split and real kernel}, and
\weightedpin{grh-compact/src/Lib/DiscreteSeries/FlenstedJensen/PsiBasic.lean}
{21-L41}{The polar form and inversion}.}

For the filling Haar \(dy\) put
\(\delta=\deg_v=(\int_{G_1}|m_\lambda(y)|^2\,dy)^{-1}\).
With the convention \(\langle u,v\rangle=u^*v\),
Schur orthogonality in \(\pi_\lambda\) gives
\[
 \delta\int_{G_1}m_\lambda(xya)\overline{m_\lambda(y)}\,dy
       =m_\lambda(x)m_\lambda(a).
\]
For the right action \(R(g)F(y)=F(yg)\), use the kernel
\(e_\lambda(g)=\delta\,\overline{m_\lambda(g)}\). The decay
estimate below makes \(e_\lambda\) integrable. Schur orthogonality
gives \(e_\lambda^*=e_\lambda\), \(e_\lambda*e_\lambda=e_\lambda\),
and \(e_\lambda*\delta_g*e_\lambda=m_\lambda(g)e_\lambda\).
Indeed the last convolution at \(x\) has integrand
\(\delta^2\overline{m_\lambda(y)}m_\lambda(x^{-1}yg)\).
Thus \(Q_\lambda=R(e_\lambda)\) is an orthogonal projection and
\(Q_\lambda R(g)Q_\lambda=m_\lambda(g)Q_\lambda\). In \(\pi_\lambda\)
the projection is \(P_w\), again by Schur orthogonality. The column
coefficient is realized independently on the symmetric power.
If \(j\) is its distinguished coordinate and \(v_0=e_j^{\otimes a_v}\)
has norm one there, then
\(m_\eta(h)=(h_{jj})^{a_v}\) and
\(\langle v_0,\rho_H(h^{-1})v_0\rangle=\overline{m_\eta(h)}\).
On the common block \(X\) of zero weights the row highest vector is
fixed by \(U(X)\). Indeed the raising operators there kill \(w\),
and the squared norms of the corresponding lowerings equal the
weight differences, all zero. The diagonal torus acts trivially.
The column vector is fixed there as well. These are the row and
column coefficients, with their indicated conjugations, in the
stationary average.

\paragraph{Global real coefficient estimates.}
\teamanchor{proof.weighted-bessel-global-decay}%
In the sign coordinates write \(D(g)=(g+JgJ)/2\)
for the sign-block diagonal part, and \(O(g)=g-D(g)\). We use the
operator norm \(\|\cdot\|\) and distinguish it from the Frobenius norm
\(\|\cdot\|_{\mathrm F}\). Fix a nonnegative integer \(c\) no greater
than any opposite-sign Harish--Chandra gap \(|\lambda_i-\lambda_j|\).
The bound underlying both norms is
\begin{equation}
 \label{eq:weighted-bessel-tail-determinant}
 |m_\lambda(g)|^2 |\det D(g)|^c\le 1 .
\end{equation}

To prove it, let \(z\) be positive Hermitian in \(G_1\). Then
\(z^{-1}=JzJ\) and \(\det z=1\). Denote its principal minor on \(B\)
by \(\Delta_B(z)\), with empty minor one. Jacobi's identity gives
\(\Delta_{B^c}(z)=\Delta_B(z)\). The opposite-sign differences of \(\mu=\lambda+\rho_G\) have absolute
value at least \(c\), since \(\rho_G\) follows the same order as \(\lambda\).
In the weight minors from \eqref{eq:weighted-row-coefficient-polar},
the determinant denominators are one for this \(z\).
Complementary minors give
\(M_{-\mu}(z)=M_\mu(z)\), and translating all thresholds gives
\(M_{\mu+c}(z)=M_\mu(z)\).

Suppose \(c>0\). Put \(S_t=\{i:\mu_i\ge t\}\),
\(T_t=\{i:\mu_i+c\ge t\}\), and \(Y_t=T_t\setminus S_t\).
The entries in \(Y_t\) have weights in the integer interval
\([t-c,t-1]\), so each nonempty \(Y_t\) lies in one sign block.
The determinant inequality for two principal sets of a positive
Hermitian matrix gives
\[
 \Delta_{Y_t}(z)\le\Delta_{T_t}(z)\Delta_{S_t^c}(z)
                       =\Delta_{T_t}(z)\Delta_{S_t}(z).
\]
Indeed, after taking the Schur complement of the intersection of the
two sets, the determinant of the resulting two-block matrix is at most
the product of its two block determinants. This follows by taking the
Schur complement of one block and using monotonicity of the determinant
on positive definite matrices. Here the union is the full set, whose
minor is one. For each residue class of \(t\) modulo \(c\), the sets
\(Y_t\) in that class partition the indices, since index \(i\) occurs
for \(\mu_i<t\le\mu_i+c\). On each sign block the product of the
partition minors is at least the block determinant, by the same
determinant inequality applied to disjoint sets. Consequently
\[
 \det D(z)^c\le\prod_t\Delta_{Y_t}(z)
            \le M_{\mu+c}(z)M_\mu(z)=M_{-\mu}(z)^2 .
\]

Write \(g=pz\) by polar decomposition, with \(p\) in the maximal
compact group \(K=U(p_+)\times U(p_-)\). For the polar integral of the
unit lowest-type coefficient, the numerator is
\(M_\Lambda(pu)\overline{M_\Lambda(u)}\), the denominator is
\(M_{-\mu}(u^*zu)\), and the outside factor is
\(\bigl(\int_K|M_\Lambda(u)|^2\,du\bigr)^{-1}\) for probability Haar.
The determinant of \(D(u^*zu)\) equals \(\det D(z)=|\det D(g)|\).
Cauchy--Schwarz in \(u\), with left invariance of probability Haar on
\(K\), proves \eqref{eq:weighted-bessel-tail-determinant}. When \(c=0\)
the estimate follows from unitarity of the row representation.

The block diagonal part of \(g^*Jg=J\) gives
\(D(g)^*D(g)=1+O(g)^*O(g)\). Taking determinants, the product of one
plus the nonnegative eigenvalues is at least one plus their sum.
Also \(D(z)=(z+z^{-1})/2\), and the eigenvalues of \(z\) occur with
their reciprocals. If \(\|z\|>1\), its largest eigenvalue and its
reciprocal contribute at least \((\|z\|/2)^2\) to the determinant of
\(D(z)\), and all the other eigenvalues of \(D(z)\) are at least one.
If \(\|z\|=1\), then \(z=1\). Since \(\|g\|=\|z\|\), the same bound
holds for \(|\det D(g)|\). Hence
\begin{equation}
 \label{eq:weighted-bessel-tail-global}
 |m_\lambda(g)|^4\le (1+\|O(g)\|_{\mathrm F}^{2})^{-c},\qquad
 |m_\lambda(g)|\le (2/\|g\|)^c .
\end{equation}
These inequalities use \(c=A\) in the off-sign comparison for a
parameter of gap \(500r^3+A\). They also permit simultaneously all the
operator-norm powers up to that gap used below.
\footnote{The polar integral is Sun's lowest-layer formula
\cite[\S6.3, (25)--(26)]{real-branch-sun} on the matched sign frame.
The determinant deductions correspond to
\weightedpin{grh-compact/src/Lib/DiscreteSeries/FlenstedJensen/PsiDecay.lean}{18-L54}{Weighted minors and power decay},
\weightedpin{grh-compact/src/New/DiscreteSeries/TwoFactorSetup.lean}{102-L114}{The off-sign estimate}, and
\weightedpin{grh-compact/src/Lib/General/LinearAlgebra/BlockDiagonalDet.lean}{19-L30}{The operator norm and block determinant}.}

\paragraph{Derived vectors and the local coefficient mass.}
Write \(\pi\) for the row discrete series, \(u_\lambda\) for its unit
highest vector, and \(L\) for a bound on \(|2\lambda_i|\). For a smooth
vector let \(E_{ab}=d\pi(e_{ab})\), using complexification of the real
unitary Lie algebra. The adjoint on the smooth vectors is
\(E_{ab}^{\dagger}=\varepsilon_a\varepsilon_b E_{ba}\), as follows
from \(X\mapsto-JX^*J\), the real structure on the complexified Lie
algebra. Let \(C_G\) be the sum of \(E_{ab}E_{ba}\) over all pairs and
\(C_K\) the sum over equal-sign pairs. Then
\[
 \sum_{a,b}\|E_{ab}u\|^2=\langle u,(2C_K-C_G)u\rangle .
\]
Let \(V_j\) be the space generated from the lowest \(K\)-type by
products of at most \(j\) derived operators. It is a finite-dimensional
\(K\)-module, as a quotient of the sum of tensor products of the
lowest type with at most \(j\) adjoint factors. Each coordinate of its
highest weights has absolute value at most
\((L+r-1)/2+j\). Indeed, the lowest highest weight is
\(\lambda+\rho_G-2\rho_K\), whose correction has coordinates of
absolute value at most \((r-1)/2\). The weights of that type lie in the
convex hull of its compact Weyl orbit, and each adjoint weight shifts
each coordinate by at most one.

On a highest compact weight \(\nu\), \(C_K\) has eigenvalue
\(\|\nu+\rho_K\|^2-\|\rho_K\|^2\). On the row representation \(C_G\)
has eigenvalue \(\|\lambda\|^2-\|\rho_G\|^2\). Put
\(N_j=L+r+1+j\). We have \(|(\nu+\rho_K)_i|\le N_j\) and
\(\|\rho_G\|^2\le rN_j^2\). Thus the last sum of squares is bounded
by \(3rN_j^2\|u\|^2\) on \(V_j\). Cauchy--Schwarz on the matrix
entries and \(\|X\|_{\mathrm F}\le\sqrt r\,\|X\|\) give
\begin{equation}
 \label{eq:weighted-bessel-tail-casimir}
 \|d\pi(X)u\|\le (r+1)^4(L+r+1+j)\|X\|\|u\|,
 \qquad u\in V_j,
\end{equation}
with \(d\pi(X)u\in V_{j+1}\). This also holds for the compact
representation, with its own matrix size and doubled parameter bound.

If \(\|g-1\|<1/2\), the convergent logarithm satisfies
\(X=\log g\in\mathfrak u(J)\) and \(\|X\|\le2\|g-1\|\).
The real Lie algebra condition follows from
\(g^{-1}=Jg^*J\) by taking logarithms locally. Integrate the derivative
of \(\pi(g_0\exp(tX))u_\lambda\) and apply
\eqref{eq:weighted-bessel-tail-casimir} at \(j=0\). For
\(\|g-1\|\ge1/2\), the coefficient difference is at most two by
unitarity. In both cases, with \(N_0=L+r+1\), we have
\begin{equation}
 \label{eq:weighted-bessel-tail-lipschitz}
 |m_\lambda(g_0g)-m_\lambda(g_0)|\le100 N_0^5\|g-1\| .
\end{equation}
In particular \(|m_\lambda|\ge1/2\) on
\(B_G((200N_0^5)^{-1})\), where
\(B_G(a)=\{g\in G_1:\|g-1\|\le a\}\).
\footnote{The earlier window route gives the same Casimir and tensor
argument in the proof of \eqref{eq:real-branch-derived} and the
exponential path before \eqref{eq:real-branch-local-scale}.
\weightedpin{grh-compact/src/Lib/DiscreteSeries/Rigidity/CasimirGrowth.lean}{64-L83}{Filtration growth} and
\weightedpin{grh-compact/src/Lib/DiscreteSeries/FlenstedJensen/TailBounds.lean}{14-L47}{The local Lipschitz bound}.}

\paragraph{Volume and radial integration.}
\teamanchor{proof.weighted-bessel-tail-haar}%
Let \(j:S_0\to G_1\) be a proper homomorphism and \(ds_0\) right
Haar. Write \(V_D\) for the mass of
\(\{s:\|j(s)-1\|\le D^{-1}\}\). For \(D,R_0\ge1\), any
\(g_0\in G_1\) satisfies
\[
 \operatorname{vol}_{S_0}\{s:\|j(s)g_0\|\le R_0\}
       \le (3D R_0^2)^{2r^2}V_D .
\]
Take a maximal \((DR_0)^{-1}\)-separated subset of the indicated
matrices. Disjoint ambient norm-balls of half this radius show that
the number of centres is at most \((1+2D R_0^2)^{2r^2}\), since
\(M_r(\mathbf C)\) has real dimension \(2r^2\). The nearest centre
\(j(s_i)g_0\) has inverse of norm at most \(R_0\), since
\(g^{-1}=Jg^*J\). Right division places \(ss_i^{-1}\) in the ball
defining \(V_D\). Right invariance then gives the inequality.
For \(a>4r^2\), apply the identity
\(u^{-a}=\int_u^\infty a R_0^{-a-1}\,dR_0\) to
\(u=\|j(s)g_0\|\ge1\), and integrate by Tonelli. The integral of
\((2/u)^a\) is at most
\(2^a(3D)^{2r^2}\frac{a}{a-4r^2}V_D\).
In particular, with \(D=1\) and integral \(a\ge8r^2+1\),
\begin{equation}
 \label{eq:weighted-bessel-tail-subgroup}
 \int_{S_0}(2/\|j(s)g_0\|)^a\,ds_0\le 16^a V_1 .
\end{equation}
Properness makes the closed unit neighbourhood of the identity in
this comparison have finite mass. For the closed subgroup \(S\) and any \(x_0\in G_1\), use the
proper conjugation \(j(s)=x_0^{-1}sx_0\) and \(g_0=x_0^{-1}y\).
The multiplier for conjugation by \(H\) on \(du\) is a
continuous character into \(\mathbf R_{>0}\). Its compact
image is trivial, so \(ds=du\,dh\) is right invariant.
Here \(V_1\) is the mass of \(B_x=\{s:\|x_0^{-1}sx_0-1\|\le1\}\).
The bound is uniform in \(y\).

For \(S_0=G_1\) and \(j=\mathrm{id}\), integrate in the compact-dual
Haar \(dg\). The formal degree here satisfies
\(\delta\int_{G_1}|m_\lambda(g)|^2\,dg=1\). The local mass from
\eqref{eq:weighted-bessel-tail-lipschitz} therefore implies
\(\delta\operatorname{vol} B_G((200N_0^5)^{-1})\le4\).
With \(D=200N_0^5\) and \(a\ge8r^2+1\), the layer integral above
gives the needed formal-degree weighted estimate
\begin{equation}
 \label{eq:weighted-bessel-tail-degree}
 \delta\int_{G_1}(2/\|g\|)^a\,dg
    \le 8(2000 N_0^6)^{2r^2}\,2^a .
\end{equation}
The quotient measure used below is \(d\dot y\) with
\(dg=du\,d\dot y\), so the subgroup and group applications involve
the same real Haar as the Bessel seed and its Schur estimates.
Left translation by \(H\) preserves this quotient measure by the
same conjugation calculation in the \(U\)-fiber.
\footnote{Compare the earlier volume proof
\eqref{eq:real-branch-volume}--\eqref{eq:real-branch-power-integral}.
\weightedpin{grh-compact/src/Lib/DiscreteSeries/Haar/RegularRepL2.lean}{35-L77}{Right-translate covering and decay integral},
\weightedpin{grh-compact/src/Lib/Adelic/QuotientMeasure/Packing.lean}{95-L125}{Layer integral and ambient packing},
\weightedpin{grh-compact/src/Lib/DiscreteSeries/FlenstedJensen/TailBounds.lean}{49-L66}{The degree-weighted integral}, and
\weightedpin{grh-compact/src/New/DiscreteSeries/BesselExistence.lean}{44-L77}{The subgroup and its unfolding}.}

\paragraph{The Bessel average in the column model.}
Use the full Schur matrices \(S_v(h)\) of the prescribed column
coordinate space and its positive Gram \(Q=B^*B\).
In an orthonormal model the action is
\(\rho_H(h)=B S_v(h)B^{-1}\). The column basis and its dual are related by
\(Q S_v(h)=S_v(h^{\mathsf T})^{\mathsf T}Q\), so
\(C=B^*\) satisfies
\[
 C\rho_H(h)C^{-1}=S_v(h^{\mathsf T})^{\mathsf T}.
\]
Let \(V=\mathbf C^{I_v}\) be this irreducible unitary column model
of dimension \(d\). Use its unit highest vector \(v_0\) from the
coefficient calculation. In the bilinear period the column
uses the dual of the displayed transport. Write \(h_s\) for the
compact component of \(s\in S\).
For a point \(x_0\in G_1\) to be specified, set
\[
 Y(y)=\int_S m_\lambda(x_0^{-1}sy)\,
         \operatorname{wave}(-\sigma(s))\rho_H(h_s^{-1})v_0\,ds .
 \tag{*}
\]

Right reindexing gives \(Y(uy)=\operatorname{wave}(\sigma(u))Y(y)\)
and \(Y(hy)=\rho_H(h)Y(y)\). The subgroup estimate
\eqref{eq:weighted-bessel-tail-subgroup} and the operator norm
bound \eqref{eq:weighted-bessel-tail-global} give absolute
convergence and boundedness for every \(y\). The same Haar gives
\[
 \int_{U\backslash G_1}\int_S F(x_0^{-1}sy)\,ds\,d\dot y
   =\int_{G_1}F(g)\,dg \qquad(F\ge0),
\]
since \(H\) normalizes \(U\), preserves \(du\), and has mass one.
This gives \(Y\in L^1(d\dot y)\) in norm.
By the degree-weighted integrability
\eqref{eq:weighted-bessel-tail-degree}, \(e_\lambda\) is in \(L^1\).
Apply the Schur identity with \(a=1\) inside the average
to obtain \(Q_\lambda Y=Y\). The interchange uses the uniform
subgroup bound. Translation of \(e_\lambda\) is continuous
in \(L^1\), which gives continuity of the bounded average.
\footnote{\weightedpin{grh-compact/src/New/DiscreteSeries/BesselExistence.lean}
{34-L79}{Graph and Haar},
\weightedpin{grh-compact/src/New/DiscreteSeries/BesselExistence.lean}
{177-L222}{Average and reproducer}, and
\weightedpin{grh-compact/src/New/Filling/ColumnCoords.lean}
{36-L46}{Bilinear Weyl-to-Bessel coordinates}.}

\paragraph{The local coefficient after removal of its phase.}
\teamanchor{proof.weighted-real-bessel-gaussian}%
Write \(\ell=2\lambda\) in the sign coordinates. With the
positive root order fixed above, \(|(2\rho_K)_i|\le r\) and
\[
 2(\rho_G)_i=\#\{j:\ell_j<\ell_i\}-\#\{j:\ell_j>\ell_i\}.
\]
Thus \(\mu\) has the same order as \(\ell\).
For an integral weight \(w\), put \(P_t=\operatorname{diag}({\bf1}_{w_i\ge t})\).
The minor product above has the equivalent finite expression
\[
 M_w(Y)=\det(Y)^c \prod_{c<t\le d}\det\bigl(1+P_t(Y-1)\bigr),
 \qquad c\le\min_i w_i,\quad d\ge\max_i w_i .
\]
Here \(c,d,t\) are integers. The expression is independent of the chosen
endpoints. The local coefficient formula proved above uses the probability
Haar \(dk\) on \(K\), with normalization
\(D_K=\int_K |M_\Lambda(k)|^2\,dk\).

Let \(g\in U(J)\), \(\|g-1\|_{\rm F}\le\varepsilon_{\rm loc}\), and suppose that
\[
 T\ge\max\{1,r^2,r^2\log(r+L+1)\},
 \qquad |\ell_i|\le L,\qquad \varepsilon_{\rm loc} T\le10^{-12}.
\]
Write \(Z=\log g\) using the series at \(1\). Then
\(\|Z\|_{\rm F}\le2\varepsilon_{\rm loc}\) and \(Z^*=-JZJ\), by
\(g^*=Jg^{-1}J\) and the same local logarithm for \(g^{-1}\).
In particular \(\operatorname{tr}(\Lambda Z)\) is imaginary. Define
\[
 Q_w(A)=\sum_{i,j}|w_i-w_j|\,|A_{ij}|^2,\qquad
 q_g=Q_{-\mu}(Z)/4,\qquad \kappa_g=e^{-q_g}.
\]
For the row coefficient with opposite-sign differences
\(|\lambda_i-\lambda_j|\ge500r^3+A\) as in the local integral
formula, we will show
\begin{equation}\label{eq:weighted-bessel-gaussian}
 \bigl|m_\lambda(g)e^{-\operatorname{tr}(\Lambda Z)}-\kappa_g\bigr|
 \le \kappa_g/10+e^{-100T}.
\end{equation}

In each sign block, order coordinates by decreasing \(\ell_i\).
Let \(Z_-\) and \(Z_+\) be the strictly lower and strictly upper
parts of these blocks, respectively. The triangular covariance in
the local coefficient formula replaces \(g\) by
\[
 \widehat g=e^{-Z_-}e^Ze^{-Z_+},\qquad
 m_\lambda(g)=D_K^{-1}\int_K
       \frac{|M_\Lambda(k)|^2}{M_{-\mu}(k^*\widehat g k)}\,dk .
\]
The two unipotent factors fix the highest line and its dual on the
appropriate sides with scalar \(1\). Each exponent has Frobenius norm
at most \(2\varepsilon_{\rm loc}\), so the matrices in the replacement are within
\(e^{6\varepsilon_{\rm loc}}-1<1/16\) of \(1\) in Frobenius norm. This is inside
the covariance neighborhood, and all denominator compressions are
accretive.

We give the remainder estimate for a compressed determinant. For an
orthogonal projection \(P\) and matrices \(A_1,A_2,A_3\) with
\(\|A_i\|_{\rm F}\le d_0\le 1/100\), set
\[
 \begin{aligned}
 T_P(A,B)&=\operatorname{tr}\bigl(PA(1-P)B\bigr),\\
 S_P&=\tfrac12\sum_i T_P(A_i,A_i)+\sum_{i<j}T_P(A_i,A_j),
 & \zeta&=\max_i\|[P,A_i]\|_{\rm F}.
 \end{aligned}
\]
The series for the logarithm gives
\begin{equation}\label{eq:weighted-bessel-compressed-log}
 \log\det\bigl(1+P(e^{A_1}e^{A_2}e^{A_3}-1)\bigr)
 =\operatorname{tr}\bigl(P\textstyle\sum_i A_i\bigr)+S_P+e_P,
 \qquad |e_P|\le 4000d_0\zeta^2 .
\end{equation}
The logarithm here is the trace of the matrix logarithm near \(1\).
Indeed, expand the function
\[
 t\longmapsto\operatorname{tr}\log\bigl(1+P(e^{t A_1}e^{t A_2}e^{t A_3}-1)\bigr).
\]
The terms linear in \(t\) give the trace displayed above.
In degree two, the terms from the product of exponentials have
coefficients \(1/2\) on \(A_i^2\) and \(1\) on \(A_iA_j\) for \(i<j\).
The quadratic term of the logarithm subtracts the same traces with
\(P\) inserted between the two matrices, giving \(S_P\).
For the higher degrees split every \(A_i\) into its diagonal and
off-diagonal blocks for \(P\). The words using only diagonal blocks sum
to the logarithm of the determinant of the product on \(\operatorname{im}P\).
This is linear in \(t\). A trace with one off-diagonal factor vanishes.
One with at least two such factors in degree \(j\) has modulus at most
\(\zeta^2 d_0^{j-2}\) by applying the Frobenius trace inequality at two
of them and the operator norm at the others. Both block parts have
operator norm at most \(d_0\). Before the block splitting, the sum of
absolute word coefficients in degree \(j\) is at most the coefficient
of \(x^j\) in \(-\log(2-e^{3x})\). This series has nonnegative
coefficients and value less than \(2\) at \(x=1/6\). Including at most
\(2^j\) terms per word, the terms of degrees \(j\ge3\) are bounded by
\(2\zeta^2\sum_{j\ge3}12^j d_0^{j-2}\le4000d_0\zeta^2\).

Take \(P=P_t\) for \(w=-\mu\), \(A_1=-Z_-\), \(A_2=Z\), \(A_3=-Z_+\),
and now put \(P'=kPk^*\). Write \(d_0=2\varepsilon_{\rm loc}\),
\(u=\|[P,Z]\|_{\rm F}\) and \(v=\|[P,k]\|_{\rm F}=\|P'-P\|_{\rm F}\).
The minor of \(k^*\widehat g k\) for \(P\) has the compressed determinant
in \eqref{eq:weighted-bessel-compressed-log} for \(P'\).
Entry masking and \(\|A_i\|_{\rm op}\le d_0\) give
\(\|[P,A_i]\|_{\rm F}\le u\) and
\(\|[P',A_i]\|_{\rm F}\le u+2d_0v\). The diagonal of \(P'-P\) has
absolute sum \(v^2\), by summing the squares of the entries of \(k\)
across the cut for \(P\). Since \(k\) preserves each sign block and
\(\sum_i A_i\) is diagonal within them, the linear trace for \(P'\)
differs from \(\operatorname{tr}(PZ)\) by at most \(d_0 v^2\).
Replacing \(P\) by \(P'\) moves each block
\(PA_i(1-P)\) in Frobenius norm by at most \(2d_0v\), by subtracting
the two products. The same holds for the blocks in the reverse
direction. The trace inequality then gives
\[
 |S_{P'}-S_P|\le18d_0v(u+2d_0v).
\]
To evaluate \(S_P\), let \(i\) be inside the cut and \(j\) outside
it. The order of \(\mu\) is the order of \(\ell\), so \(\ell_i<\ell_j\).
At equal signs the contribution through \(i,j\) is
\(-Z_{ij}Z_{ji}/2\), as the \((1,2),(1,3),(2,3)\) terms
contribute \(-1,+1,-1\) times \(Z_{ij}Z_{ji}\), and the \((2,2)\) term
contributes \(1/2\) times it. At opposite signs only \((2,2)\)
contributes. The relation \(Z_{ji}=-\varepsilon_i\varepsilon_j
\overline{Z_{ij}}\) makes both contributions \(|Z_{ij}|^2/2\).
Thus \(S_P=u^2/4\). The three error bounds, with \(2d_0\le1\),
give an error at most \(10^4d_0(u^2+v^2)\) away from
\(\operatorname{tr}(PZ)+u^2/4\).

For each pair \(i,j\), exactly \(|w_i-w_j|\) cuts separate the two
indices, and \(\sum_{c<t\le d}P_t=\operatorname{diag}(w_i-c)\).
Also \(\det(\widehat g)=e^{\operatorname{tr} Z}\).
Multiplying the minor expansions therefore gives
\[
 M_{-\mu}(k^*\widehat g k)
 =\exp\bigl(q_g-\operatorname{tr}(\Lambda Z)-R_k\bigr),
 \qquad
 |R_k|\le 10^6\varepsilon_{\rm loc}\bigl(q_g+Q_\Lambda(k)+r^2\bigr).
\]
Here \(Q_{-\mu}(k)\le Q_\Lambda(k)+2r^2\) by
\(|\mu_i-\Lambda_i|\le r\) and \(\|k\|_{\rm F}^2=r\).
The linear trace correction is at most \(2\varepsilon_{\rm loc} r^2\).

For a cut \(P_t\) and unitary \(k\), the square modulus of the
compressed determinant on \(\operatorname{im}P_t\) is
\(\det(1-B^*B)\) there, where \(B=(1-P_t)kP_t\).
Its eigenvalue factors are in \([0,1]\) and their product is at most
\(e^{-\|B\|_{\rm F}^2}\). The two off-diagonal blocks of \(k\) have equal
squared Frobenius norm, as follows from the row and column norm
identities on \(\operatorname{im}P_t\).
Consequently \(|M_\Lambda(k)|^2\le e^{-Q_\Lambda(k)/2}\).
For a lower bound put \(U_*=r+L+1\) and \(h_*=(20U_*^3)^{-1}\).
Since \(|\Lambda_i|\le 2U_*\), and every compression of a matrix
\(\|k-1\|_{\rm op}\le h_*\) has least singular value at least \(1-h_*\),
\[
 |M_\Lambda(k)|\ge (1-h_*)^{4rU_*}\ge 1-4rU_*h_*\ge 1/2.
\]
An \(h_*\)-separated set in \(K\), for the operator norm, has at most
\((1+4/h_*)^{2r^2}\) elements by packing disjoint balls of radius
\(h_*/2\) in the ambient real matrix space. A maximal such set
covers \(K\) by right translates of the \(h_*\)-ball about \(1\).
Thus \(D_K\ge\frac14(1+4/h_*)^{-2r^2}\ge e^{-20T}\).

After removal of its phase the coefficient differs from \(e^{-q_g}\)
by the average of \(e^{-q_g}(e^{R_k}-1)\) for the probability
measure \(|M_\Lambda(k)|^2\,dk/D_K\).
If \(q_g+Q_\Lambda(k)+r^2\le50000T\), then
\(|R_k|\le1/20\) and \(|e^{R_k}-1|\le1/10\).
Otherwise, using \(|R_k|\le(q_g+Q_\Lambda(k)+r^2)/4\), we have
\[
 |M_\Lambda(k)|^2 e^{-q_g}|e^{R_k}-1|
 \le e^{-q_g-Q_\Lambda(k)/2+|R_k|}
 \le e^{-5100T}.
\]
Integrating the two bounds and using the lower bound for \(D_K\)
proves \eqref{eq:weighted-bessel-gaussian}.
\footnote{\weightedpin{grh/src/New/Appendix/DiscreteSeries/OneMinor.lean}
 {573-L628}{The projection estimate in logarithm coordinates},
 \weightedpin{grh/src/New/Appendix/DiscreteSeries/GaussianShape.lean}
 {325-L392}{The same Gaussian argument in those coordinates}, and
 \weightedpin{grh-compact/src/New/DiscreteSeries/GaussianShape.lean}
 {41-L98}{The compact proof, with quadratic form evaluated on \(g-1\)}.}

For the actual compact coefficient the column model gives
\(m_\eta(h)=h_{pp}^{a_v}\), with \(p\) the coordinate of the
nonzero weight. If \(Z_H=\log h\), take \(P\) to be the projection onto
that coordinate. Apply \eqref{eq:weighted-bessel-compressed-log} with
only \(A_2=Z_H\) nonzero. Its quadratic term is
\(-\|[P,Z_H]\|_{\rm F}^2/4\) because \(Z_H^*=-Z_H\).
After multiplication by \(a_v\) the error is at most \(32000\varepsilon_{\rm loc} q_H\),
where \(q_H=a_v\|[P,Z_H]\|_{\rm F}^2/4\).
The same split at \(q_H=50000T\) gives
\[
 \left|m_\eta(h)e^{-\operatorname{tr}(\Lambda'Z_H)}-e^{-q_H}\right|
 \le e^{-q_H}/10+e^{-100T}
 \qquad(\|h-1\|_{\rm F}\le\varepsilon_{\rm loc}).
\]
This also covers \(a_v=0\), when the coefficient is \(1\).

\paragraph{A stationary graph on the prescribed frame.}
\teamanchor{proof.weighted-real-bessel-stationary}%
Use \(B_v=w_0(v)-h_2\) from the row parameter. The bound
\(xb\le a<xb+N_1\), with \(x_1\ge x\), gives
\(xb/x_1\le B_v<b+N_1/x+N_1\le2b\).
The lower bound dominates \(2nh+4r\) for large \(x\).
The compact weight is \(\Lambda=\lambda+\rho_G-2\rho_K\).
On the initial positive block \(X\), the root correction at
index \(i\) is \(1-k/2+i\), so its \(k-1\) entries are zero.
The other entries are
\(u_i=-B_v-(i+1)h+(\rho_G-2\rho_K)_{k-1+i}\) for
\(0\le i<2n\). The root correction has modulus at most \(2r\),
so \(h>4r\) makes the entries strictly decreasing.
They lie in \([-3B_v,-B_v]\) by the bound on \(B_v\),
and their form signs are \((-1)^i\).
The column weight is zero on \(X\) and equals \(a_v\) on the remaining
coordinate \(p\) of the definite core. Set \(t_i=u_i-a_v\).
If \(c_*^2,c_i^2\) are the absolute derivatives of the nodal polynomial
at \(0,t_i\), respectively, Lagrange interpolation for
\(0>t_0>\cdots>t_{2n-1}\) gives
\[
 \sum_i\frac{(-1)^i t_i^j}{c_i^2}
 =\begin{cases}
 c_*^{-2}&j=0,\\
 0&1\le j<2n,\\
 -1&j=2n.
 \end{cases}
\]
Indeed the leading coefficient is the sum of nodal values divided by
the signed derivatives, which have signs \(+1,(-1)^{i+1}\).
Use \(u_*=(c_*/c_i)_i\), \(V_j=(t_i^{j+1}/c_i)_i\).
The first vector has form norm \(1\) and is orthogonal to \(V_j\)
for \(j\le\varrho=n-1\). The \(V_j\) satisfy
\(h(V_i,V_j)=0\) for \(i+j<2\varrho\) and \(-1\) for equality.
The Vandermonde matrix shows independence of \(u_*,V_0,\ldots,V_\varrho\).
The orthogonal sum with \(X\) has the Gram of the core and the frame
\(f_0,\ldots,f_\varrho=(e_0,\ldots,e_{\varrho-1},z)\) in the split
coordinates. By Witt extension there is an isometry \(q=x_0p_v\)
from the sign coordinates to these coordinates. Here \(p_v\)
is the fixed unitary sign frame, \(qV_j=f_j\), and \(q\) maps \(X\)
and \(u_*\) to the two designated parts of the definite core.

Let \(D=q\Lambda q^{-1}\). For \(i<\varrho\),
\(Df_i=a_v f_i+f_{i+1}\), by multiplication by \(u_j=a_v+t_j\)
in the sign coordinates. On the definite core the compression of
\(D\) is \(\Lambda'\), since \(\Lambda u_*=a_vu_*+c_*V_0\).
The group \(S\) is upper block triangular along the split flag,
with identity diagonal blocks except for its compact action on the core.
Only the middle diagonal block and the strictly lower entries of \(D\)
contribute to \(\operatorname{tr}(D(s-1))\).
The shifts on \(f_i\), paired with the adjacent entries of \(s\),
contribute \(\sigma(s)=\sum_{i<\varrho}s_{f_i,f_{i+1}}\).
The remaining lower entries are their reflected entries on the dual
flag, by \(D^*J_1=J_1D\), where \(J_1\) is the split Gram matrix.
To see their contribution, write \(J_1 f_i=\varepsilon'_i f_i^\#\),
with \(f_i^\#=e_i^*\), \(\varepsilon'_i=1\) for \(i<\varrho\),
and \(f_\varrho^\#=z\), \(\varepsilon'_\varrho=-1\).
The reflected entry of \(D\) has the factor
\(\varepsilon'_i\varepsilon'_{i+1}\), and the paired entry of \(s\)
has the same factor times
\(\overline{(s^{-1})_{f_i,f_{i+1}}}\), by \(s^{-1}=J_1s^*J_1\).
On adjacent frame levels the latter conjugate equals
\(-\overline{s_{f_i,f_{i+1}}}\). The middle block \(h_s\oplus1\)
acts trivially on \(z\), so this includes the last adjacent pair.
Since the other strict lower entries are zero, the reflected sum is
\(-\overline{\sigma(s)}\).
For \(g_s=p_v^{-1}x_0^{-1}s x_0p_v\) we obtain
\begin{equation}\label{eq:weighted-bessel-stationary-phase}
 \begin{aligned}
 \operatorname{tr}(g_s-1)&=\operatorname{tr}(h_s-1),\\
 \operatorname{tr}\bigl(\Lambda(g_s-1)\bigr)
   -\operatorname{tr}\bigl(\Lambda'(h_s-1)\bigr)
 &=\sigma(s)-\overline{\sigma(s)}=:\phi(s).
 \end{aligned}
\end{equation}
The character \(\phi\) is additive because each frame matrix is upper
unitriangular. The common \(U(X)\) inside \(S\) has the same compact
action in both sign coordinates, by the choice of \(q\).
It fixes both highest vectors. For the row vector this also follows
from the compact coefficient \(M_\Lambda\), as all the weight entries
on \(X\) are zero and each weight cut either contains \(X\) or is
disjoint from it.

The same stationary point has the required quotient size. Since the
distinct nodes are integers, \(c_i\ge1\), and the columns \(p_vV_i\)
through \(i=2\varrho\) have Euclidean lengths at most
\[
 M_0=\sqrt{2n}(3B_v+a_v)^{2n-1}.
\]
Gram--Schmidt on the first \(n\) columns is upper unitriangular
without rescaling. Each resulting column \(\widetilde V_j\) has
length at most \(M_0\), and pairing it with the old column
\(p_vV_{2\varrho-j}\) gives \(-1\), so its length is at least
\(M_0^{-1}\). The frame change lifts to \(U\). In fact, writing
\(E=(e_i)\), \(E^*_{\rm f}=(e_i^*)\) as matrices of frame and dual
columns, a frame action \(E\mapsto EA,\ z\mapsto z+Ev\)
with \(A\) upper unitriangular extends by
\[
 E^*_{\rm f}\longmapsto
 \bigl(E^*_{\rm f}+zv^*+Ev v^*/2\bigr)(A^*)^{-1}
\]
and by the identity on the core. The paired blocks preserve the
split form and have the upper triangular levels of \(U\).
After this change the isotropic duals
\(J_1\widetilde V_j/\|\widetilde V_j\|^2\) for \(j<\varrho\)
have length at most \(M_0\), pair to the identity with the
\(\widetilde V_i\), and are orthogonal to \(\widetilde V_\varrho\).
They can be installed by one further left multiplication in \(U\)
fixing the frame. Indeed, in the new split coordinates the target
dual columns have the form \(E^*_{\rm f}+X_0C+EB\), where \(X_0\)
is the matrix of core columns and \(B+B^*+C^*C=0\).
This sends \(X_0\) to \(X_0-EC^*\) and fixes \(E,z\), giving an
isometry in the subgroup of \(U\) fixing the frame.

For the resulting inverse matrix \(G=(s_0x_0)^{-1}\), \(s_0\in U\),
\(GG^*=J_1+2GP_-G^*\), with \(P_-\) the negative projection of \(J_1\).
Its negative basis is \((e_i-e_i^*)/\sqrt2,z\).
The bounds on the frame and dual columns give
\(\|GP_-\|_{\rm F}^2\le(2\varrho+1)M_0^2\) and hence
\(\|G\|_{\rm op}^2\le4nM_0^2\).
Inversion preserves the operator norm on \(U(J_1)\).
It follows that
\[
 \operatorname{size}(x_0)\le8nM_0\le b^{5n}.
\]
For the last inequality use \(3B_v+a_v\le2b^2\), \(b\ge8n\),
\(8n\sqrt{2n}\le b^2\), and \(2^{2n-1}\le b^n\).
These left translations bound the quotient size of \(x_0\).
The average uses the stationary representative \(x_0\) from
\eqref{eq:weighted-bessel-stationary-phase}.
\footnote{\weightedpin{grh-compact/src/New/Exchange/SteepParameters.lean}
 {9-L29}{The actual steep parameter},
 \weightedpin{grh-compact/src/New/DiscreteSeries/BesselExistence.lean}
 {104-L175}{The graph position}, and
 \weightedpin{grh-compact/src/New/RealPlace/FrameNorm.lean}
 {96-L126}{The quotient size}.}

\paragraph{Nonvanishing of the stationary average.}
\teamanchor{proof.weighted-real-bessel-nonvanishing}%
Put \(\beta=b^3,\ q_0=(r+k)^2,\ T=q_0\log\beta\),
\(\epsilon=10^{-13}/T\), and write
\[
 \theta(s)=m_\lambda(g_s)\overline{m_\eta(h_s)}e^{-\phi(s)},
 \qquad \omega_j(g)=(2/\|g\|_{\rm op})^j.
\]
The actual parameters satisfy
\(\beta\ge10^{12}\), \(r+k+L+1\le\beta\) with \(L=3a\), and the row has
gap at least \(500r^3+A\) for \(A=r^7\ge10^{30}T^3\).
Indeed \(B_v\le2b\), \(a_v\le a\le e x b\), \(k\le r\le3x^5\)
and \(b\ge x^{10^{12}}\) give \(3B_v+a_v\le b(6+e x)\le2b^2\)
and \(r+k+3a+1\le b^3\) for sufficiently large \(x\).
The row gap is at least \(2r^8\).
Since \(\log b\le x^{1/100}\) and \(r\ge x^5\), we have
\(T^3/r^7\le1728x^{3/100-5}\).
\footnote{\weightedpin{grh-compact/src/New/Data/AdmissibleData.lean}
 {23-L32}{Scales},
 \weightedpin{grh-compact/src/New/Data/AdmissibleData.lean}
 {90-L95}{Gap},
 \weightedpin{grh-compact/src/New/Data/Weights.lean}
 {23-L41}{The upper bound for \(B_v\)}, and
 \weightedpin{grh-compact/src/New/Filling/FillingData.lean}
 {56-L96}{Actual Bessel parameters}.}

Suppose first the off-sign blocks of \(g_s\) have combined
Frobenius norm at most \(\epsilon\). With \(A_1=g_s-1\), the row
identity for \(g_sJg_s^*=J\) gives
\(-2\operatorname{Re}(A_1)_{ii}=\sum_j\varepsilon_i\varepsilon_j
|(A_1)_{ij}|^2\).
The real weighted trace for \(\Lambda\) agrees by
\eqref{eq:weighted-bessel-stationary-phase} with that for
\(\Lambda'\) on \(h_s-1\), which is nonpositive. With weights
\(-\Lambda_i\) the sum of the unsigned squares is therefore at most
\(6B_v\epsilon^2\). In particular the rows of \(A_1\) outside \(X\)
have total squared norm at most \(6\epsilon^2\).
Write \(g_s=\left(\begin{smallmatrix}a&b'\\ c&d'\end{smallmatrix}\right)\)
for \(X\oplus X^\perp\), with the form \(1\oplus J'\).
The equations \(a^*a=1-c^*J'c\) and \(a^*b'=-c^*J'd'\) give
\[
 \|b'\|_{\rm F}
 \le\frac{\sqrt6\epsilon(1+3\epsilon)}{\sqrt{1-6\epsilon^2}}
 \le3\epsilon,\qquad
 \|(a^*a)^{1/2}-1\|_{\rm F}\le6\epsilon^2 .
\]
Remove the unitary factor of \(a\) by a common left element of \(U(X)\).
The new first factor is within \(10\epsilon\) of \(1\) in Frobenius
norm by the three block bounds. The new second factor is within
the same distance since its squared norm of the difference from
\(1\) is \(-2\operatorname{Re}\operatorname{tr}(h-1)\), equal to
the sign-weighted sum of squares for the new first factor. The
move preserves \(\theta\) and \(\|g_s\|_{\rm op}\).

At the moved point, let \(X_s=(g_s-1,h_s-1)\), and denote by
\(\mathcal L\) the difference of the two weighted traces in
\eqref{eq:weighted-bessel-stationary-phase}. The homomorphism of the graph
into the matrix pair and the additivity of \(\phi\) imply
\(\mathcal L(X_s^j)=0\) for \(j\ge2\).
In detail, \(X_s^j\) is the \(j\)-th finite difference of
\((g_s,h_s)^i\) at \(i=0\), and
\(\mathcal L((g_s,h_s)^i-1)=i\phi(s)\).
Thus the weighted difference of the two logarithms is \(\phi(s)\).
Their diagonal traces are imaginary, so the two phases of the
coefficients cancel with \(e^{-\phi(s)}\).
Using \eqref{eq:weighted-bessel-gaussian} and its compact version
with \(\varepsilon_{\rm loc}=10\epsilon\), write the product as
\((\kappa_1+e_1)(\kappa_2+e_2)\) with
\(0\le\kappa_i\le1\), \(|e_i|\le\kappa_i/10+E_0\), \(E_0=e^{-100T}\).
Its difference from \(\kappa_1\kappa_2\) is at most
\(0.21\kappa_1\kappa_2+3.2E_0\).
Thus \(|\theta(s)|-2\operatorname{Re}\theta(s)\le 10E_0\).
The moved first factor has operator norm at most \(2\), so
\(\omega_{8q_0}(g_s)\ge1\) also at the original point.

For the other points the off-sign bound gives
\(|\theta(s)|^4\le(1+\epsilon^2)^{-A}\), so \(|\theta(s)|\le e^{-1250T}\)
using \(\log(1+\epsilon^2)\ge\epsilon^2/2\).
Since \(16q_0\le A\), the global coefficient bound also gives
\(|\theta(s)|\le\omega_{16q_0}(g_s)\).
Their geometric mean bounds the modulus by
\(e^{-625T}\omega_{8q_0}(g_s)\).
Combining this bound with the near estimate gives the following
inequality on the whole stationary graph.%
\footnote{\weightedpin{grh-compact/src/New/DiscreteSeries/TwoFactorSetup.lean}
 {102-L148}{Off-sign decay and localization}, and
 \weightedpin{grh-compact/src/New/DiscreteSeries/RestrictionLemma.lean}
 {127-L189}{The two-factor comparison}.}
\begin{equation}\label{eq:weighted-bessel-stationary-cone}
 |\theta(s)|\le 2\operatorname{Re}\theta(s)
            +32\beta^{-50q_0}\omega_{8q_0}(g_s).
\end{equation}

Finally use \(B_x\) as above and
\(B_{\rm sm}=\{s:\max(\|g_s-1\|_{\rm op},\|h_s-1\|_{\rm op})\le h_0\}\),
where \(h_0=(200\beta^5)^{-1}\).
On \(B_{\rm sm}\) the Frobenius differences are at most
\(\sqrt r\,h_0\). The inequality \(\log\beta\le\beta\) gives
\(T\sqrt r\,h_0\le(200\beta^{3/2})^{-1}\le10^{-12}\).
Both quadratic forms in the Gaussian estimates are at most
\(\beta^2h_0^2\), since
\(\max_{i,j}|\mu_i-\mu_j|\le L+r\), \(a_v\le L\),
and the logarithms have Frobenius norms at most \(2\sqrt r\,h_0\).
Hence the modulus of each factor is at least \(1/2\) there.
The paired matrices of \(B_x\) have norms at most \(2\).
An \(h_0/2\)-separated set for the maximum operator norm in the two
matrix spaces has size at most
\((1+1600\beta^5)^{2(r^2+k^2)}\le\beta^{11q_0}\)
by disjoint ambient balls of radius \(h_0/4\).
Division by a covering centre on the right sends the covered points
to \(B_{\rm sm}\), since inversion preserves each operator norm.
The measure \(ds\) is right Haar, so
\(\operatorname{vol}(B_x)\le\beta^{11q_0}\operatorname{vol}(B_{\rm sm})\).
This ball \(B_x\) has positive finite mass. Its interior contains
the identity and it is compact, since the graph is closed and the
second projection takes values in the compact group \(H\).

Integrate \eqref{eq:weighted-bessel-stationary-cone} and use the
subgroup tail estimate \eqref{eq:weighted-bessel-tail-subgroup} with
\(8q_0>8r^2\). The bound on \(B_{\rm sm}\) gives
\[
 \operatorname{vol}(B_x)\le
 8\beta^{11q_0}\operatorname{Re}\int_S\theta(s)\,ds
 +128\beta^{-39q_0}16^{8q_0}\operatorname{vol}(B_x).
\]
The second term is at most half the left side. Thus
\begin{equation}\label{eq:weighted-bessel-nonvanishing}
 0<\operatorname{vol}(B_x)\le
 \beta^{20q_0}\operatorname{Re}\int_S\theta(s)\,ds .
\end{equation}
The scalar integral is \(\langle v_0,Y(x_0)\rangle\) in the unitary
column model, giving the required nonzero average in the same Haar.
In particular the subgroup tail bound with \(M=8r^2+1\) gives
\(\int_S\omega_M(x_0^{-1}sy)\,ds
 \le b^{M+60q_0}\|Y(x_0)\|\) for every \(y\), since \(b\ge16\).

\paragraph{Truncation in the same quotient measure.}
\teamanchor{proof.weighted-bessel-tail-truncation}%
Keep the nonzero average \(Y\) at the stationary point \(x_0\).
Set \(M=8r^2+1\), \(C_0=16^M\) and \(R=b^{6n+2}\). 
The graph bound
\(\operatorname{vol}(B_x)\le b^{60(r+k)^2}\|Y(x_0)\|\), together
with \eqref{eq:weighted-bessel-tail-subgroup} and \(b\ge16\), gives
\[
 \int_S\omega_M(x_0^{-1}sy)\,ds\le P_0\alpha,
 \qquad P_0=b^{M+60(r+k)^2},\qquad \alpha=\|Y(x_0)\|>0 .
\]

Let \(B_1=B_G(1)\) and put
\[
 \begin{aligned}
  \Phi(y)&=\int_S\omega_{2M}(x_0^{-1}sy)\,ds,\qquad
  I_M=\int_{G_1}\omega_{2M}(g)\,dg,\\
  d\mu_0(g)&=\frac{C_0P_0\alpha}{\operatorname{vol}(B_1)}
                           \omega_M(g)\,dg .
 \end{aligned}
\]
The positive function \(\Phi\) is left \(S\)-invariant by Haar
reindexing. Unfolding with the free-entry \(du\) and probability
\(dh\) gives \(\int_{U\backslash G_1}\Phi\,d\dot y=I_M\).
Also \(\|Y(y)\|\le\Phi(y)\), by the coefficient bound at exponent
\(2M\). This gives \(Y\in L^1(d\dot y)\), and its boundedness gives
\(Y\in L^2(d\dot y)\).

For any \(s\in S\) and \(y\in G_1\), let \(t=x_0^{-1}sy\).
If \(tg\in B_1\), then \(\|tg\|\le2\) and
\(\|g\|\le2\|t\|\). The original orbit size satisfies
\(\operatorname{size}(yg)\le2\operatorname{size}(x_0)<R\), using
\(\operatorname{size}(x_0)\le b^{5n}\). To see this, use left
\(S\)-invariance and
\(\operatorname{size}(yh)\le\|h\|\operatorname{size}(y)\) for
\(h\in G_1\), which follows by taking the infimum and using
\(\|h\|\ge1\). Moreover
\(\omega_M(g)\ge C_0^{-1}\omega_M(t)\).
Left invariance of \(dg\) makes the mass of \(t^{-1}B_1\) equal to
that of \(B_1\), so
\[
 \mu_0\{g:\operatorname{size}(yg)\le R\}\ge P_0\alpha\omega_M(t).
\]
Multiply by \(\omega_M(t)\), integrate over \(s\), and use the
uniform estimate for its integral to conclude
\[
 \Phi(y)\le\mu_0\{g:\operatorname{size}(yg)\le R\},\qquad
 \mu_0(G_1)\le C_0^2P_0\alpha .
\]
For the second estimate use \eqref{eq:weighted-bessel-tail-subgroup}
with \(S_0=G_1\), \(j=\mathrm{id}\), \(g_0=1\), and \(a=M\).

Apply the reproducing identity at \(x_0\) and bound
\(|m_\lambda(g^{-1})|\le\omega_{2M}(g)\), since inversion
preserves the operator norm. Cauchy--Schwarz with this weight gives
\[
 \alpha^2\le \delta^2 I_M\int_{G_1}\|Y(x_0g)\|^2
                                      \omega_{2M}(g)\,dg
    =\delta^2 I_M\int_{U\backslash G_1}\|Y(y)\|^2\Phi(y)\,d\dot y .
\]
The equality unfolds the same group Haar, using the \(S\)-invariance
of the squared norm. Since \(\Phi\le C_0^2P_0\alpha\), this gives
\(\alpha\le \delta^2 I_M C_0^2 P_0\|Y\|_2^2\). The group estimate
\eqref{eq:weighted-bessel-tail-degree} gives
\[
 \delta I_M C_0^2\le8(2000N_0^6)^{2r^2}2^{10M}\le b^{99r^2} .
\]
Indeed \(2000\le b\) and \(N_0\le b^4\) give the factor
\(b^{50r^2}\) for the power of \(2000N_0^6\). The remaining factor
is \(2^{80r^2+13}\le b^{10r^2}\), since \(b\ge2000>2^{10}\).
Thus the stated power has room to spare.

Set \(c_0=\sqrt d/\|Y\|_2\), which is finite by continuity and
\(Y(x_0)\ne0\). Define the normalized full vector, its continuous
scalar cut and their positive majorant by
\[
 \begin{gathered}
  W=c_0Y,\qquad W^\flat(y)=\min(1,\max(0,R-\operatorname{size}(y)))W(y),\\
  A(y)=c_0\Phi(y),\qquad \mu_{\rm cut}=c_0\mu_0 .
 \end{gathered}
\]
The size is continuous by the right multiplication inequality and
its application with the inverse. The cutoff preserves the left
Bessel laws and vanishes for size at least \(R\). Its support modulo
\(U\) is compact. In fact, for each point with size at most \(R\)
choose a translating element of \(S=HU\) within one of the infimum,
and discard the Euclidean-unitary left factor in \(H\). This gives
a representative in a fixed compact group norm-ball. We have
\(\|W^\flat\|_2^2\le\|W\|_2^2=d\). The matrix
\(\int W(y)W(y)^*\,d\dot y\) commutes with the irreducible compact
column action, since \(H\) normalizes \(U\) and preserves \(du\) and
\(d\dot y\). Its trace is \(d\), so it is the identity on the
orthonormal column space. The full vector retains the row reproducer.

The majorant estimates are
\begin{equation}
 \label{eq:weighted-bessel-cut-majorant}
 A(y)\le\mu_{\rm cut}\{g:\operatorname{size}(yg)\le R\},\qquad
 \mu_{\rm cut}(G_1)\int A\,d\dot y\le b^{2r^3}d .
\end{equation}
For the product estimate use
\[
 c_0^2 C_0^2 P_0\alpha I_M\le d(\delta I_M C_0^2 P_0)^2
                         \le d b^{2r^3} ,
\]
where \(M+60(r+k)^2+99r^2\le347r^2+1\le r^3\) at the actual
ranks. If \(W(y)\ne W^\flat(y)\), then
\(\operatorname{size}(y)+1>R\). For every \(s\) this implies
\(\|x_0^{-1}sy\|>2b\), since otherwise
\(\operatorname{size}(y)+1\le2b^{5n+1}+1\le R\). Apply the power
bound at exponent \(2M+3r^3f\). It applies because \(2M\le500r^3\)
and \(3r^3f\le r^7\), with the row gap at least \(500r^3+r^7\).
For both \(W^\natural=W,W^\flat\) this yields
\begin{equation}
 \label{eq:weighted-bessel-cut-decay}
 \|W^\natural(y)\|\le
 \begin{cases}
  A(y),&W(y)=W^\flat(y),\\
  b^{-3r^3f}A(y),&W(y)\ne W^\flat(y).
 \end{cases}
\end{equation}
On the set where the vectors differ, the difference also has norm
at most \(\|W(y)\|\), since the cut multiplies by a scalar in
\([0,1]\). Now \(A\le\mu_{\rm cut}(G_1)\), so
\[
 \|W-W^\flat\|_2^2\le b^{-6r^3f}\int A^2\,d\dot y
        \le b^{2r^3-6r^3f}d\le b^{-10r^2}d .
\]

All the magnitude bounds hold on the chosen large squares. There
\(x^5\le r\le3x^5\), \(k\le r\), \(n\le r\), \(L=3a\le3exb\),
and \(f\le e^{C'_0}x^{11}(\log b)^{1/5}\). The scale
\(\log b=x^\varepsilon\) satisfies \(100\log\log b\le\log x\).
In particular \(f\le e^{C'_0}x^{11+1/500}\). Together
these give \(b\ge2000\), \(r+L+1\le b^4\), \(3r^3f\le r^7\) and
the required row gap from the actual gap at least \(2r^8\). They
also give
\[
 (r+1)^4(L+r+1+2n^2)\le b^5 ,
\]
which bounds every factor of \eqref{eq:weighted-bessel-tail-casimir}
on the required normalized derivative words. The
\hyperlink{proof.weighted-bessel-smooth-words}{derivative Schur argument}
uses the full \(W\) in this same quotient. The absolute
periodization uses \eqref{eq:weighted-bessel-cut-majorant} on the
full seed, and the relative trace uses both estimates on the full
and cut link tensors, before their separate arithmetic class cutoff.
\footnote{\weightedpin{grh-compact/src/New/RealPlace/BesselCut.lean}{66-L157}{The positive measure and the cut on the same vector},
\weightedpin{grh-compact/src/New/Filling/BesselSystem.lean}{85-L109}{The pair and its Schur inputs},
\weightedpin{grh-compact/src/New/RealPlace/BesselOrthonormal.lean}{8-L26}{Orthonormality in the quotient},
\weightedpin{grh-compact/src/New/Trace/GoodIdempotent.lean}{143-L149}{The row gap}, and
\weightedpin{grh-compact/src/New/Data/Magnitudes.lean}{129-L179}{Bounds on the actual parameters}.}

\paragraph{Schur bounds for the right derivatives.}
On the right-regular Hilbert space of \(U\)-equivariant vectors
\(Q_\lambda W=W\), and
\(Q_\lambda R(g)Q_\lambda=m_\lambda(g)Q_\lambda\).
Thus the cyclic representation of \(W/\sqrt d\) is isometric to
the row discrete series. Match the Gram matrices of the two sets
of translates of the unit vectors and complete. The orthogonal
projection onto the cyclic space commutes with right translation.
For \(V\) in that space and square-integrable \(F\) in the ambient
right-regular space, Schur orthogonality gives
\[
 \int_{G_1}|\langle F,R(g)V\rangle|^2\,dg
       \le \|F\|_2^2\|V\|_2^2/\delta.
\]
For matrix-unit words \(w\) with \(|w|\le 2n^2\),
\eqref{eq:weighted-bessel-tail-casimir} gives
\(\|W_w\|_2\le P_1^{|w|}\sqrt d\) for \(P_1\ge b^5\).
For each scalar test \(\phi\in L^1(G_1)\cap L^2(G_1)\),
duality and Cauchy--Schwarz in the Schur relation give
\[
 \int_{U\backslash G_1}\left\|\int_{G_1}\phi(g)W_w(yg)\,dg\right\|^2
    d\dot y \le P_1^{2|w|}\frac{d}{\delta}\int_{G_1}|\phi(g)|^2\,dg .
\]

\teamanchor{proof.weighted-bessel-smooth-words}%
Here is an exact smoothing on the same cyclic space. Let \(K\) be the
maximal compact of the row group \(G_1\). Average a nonnegative smooth
compactly supported approximate identity of integral one over
conjugation by \(K\), using probability Haar. The resulting functions
\(\eta_j\) are again an approximate identity, since \(K\) is compact.
On the cyclic representation generated by \(W\), the operators
\(R(\eta_j)\) commute with \(K\). The lowest \(K\)-type occurs there
with multiplicity one and contains \(W\) under the Gram isometry.
Consequently \(R(\eta_j)W=c_jW\), with \(c_j\to1\) by strong
continuity. Choose one \(c_j\ne0\) and put \(\eta=\eta_j/c_j\).
The equality
\[
             W(y)=\int_{G_1}\eta(a)W(ya)\,da
\]
holds pointwise. Indeed both sides are continuous and the quotient
Haar has full support.

For \(X\) in the real Lie algebra, define
\(\ell_X\eta(a)=\left.\frac{d}{dt}\right|_{0}\eta(\exp(-tX)a)\).
Changing variable by left translation in the last integral gives
the derivative of \(W(y\exp(tX))\) with test \(\ell_X\eta\).
Iterating in the same order gives the word test
\(\ell_{X_1}\cdots\ell_{X_j}\eta\). All these tests are smooth and
compactly supported. Since \(W\) is bounded and continuous, the
integrals give bounded continuous representatives of the Hilbert-space
derivatives. This also follows in \(L^2\) by differentiating the
left translates of \(\eta\) in \(L^1(G_1)\). Extend the derivatives
complex-linearly to the matrix units \(E_{ab}\) in the normalized
filling frame. The map \(y\mapsto p_v^{-1}yp_v\) sends those units
to \(p_v^{-1}E_{ab}p_v\) in the row frame, each of operator norm one.
Writing \(W_{(a,b)w}\) for prepending such a derivative to \(w\), one
obtains along a smooth curve with \(\gamma'(t)=\gamma(t)X\)
\[
 \frac{d}{dt}W_w(\gamma(t))
      =\sum_{a,b} X_{ab}W_{(a,b)w}(\gamma(t)).
\]
Right differentiation preserves the two left Bessel transformation laws.
Thus the word and convolution estimates apply to the full pointwise
vector used on the positive Gram cell. The trace estimate for the
radial cut \(W^\flat\) uses its order-zero norm bound and its support.
\footnote{\weightedpin{grh-compact/src/Lib/DiscreteSeries/Rigidity/KernelFixes.lean}
{60-L69}{Minimal idempotent},
\weightedpin{grh-compact/src/New/RealPlace/BesselCut.lean}
{23-L66}{Derivative Schur bound}.}

\paragraph{The real argument in the period and trace.}
If \(D_v\) is the diagonal basis-scale
matrix, \(p_v\) the sign frame and \(x_v\) the point defining the
adelic realization, the real argument in the seed is
\[
 y_v=D_v^{-1}x_v g_v p_v^{-1}
     =(D_v^{-1}x_v p_v^{-1})(p_v g_v p_v^{-1}).
\]
The inverse sign-frame equivalence preserves the compact-dual
Haar, and the remaining factor is a left translation. This map
intertwines rational left multiplication with its action in the
normalized real frame. At each place write \(\mu_v\) for the
finite positive measure \(\mu_{\rm cut}\). The seed on its finite class coset is
\(\beta_c(y)W_i(y)\), with \(\beta_c\) the class cutoff and
\(W_i\) the full Bessel coordinate. The orbital estimates use
the radial cut, with its separate pointwise gain
\eqref{eq:weighted-bessel-cut-decay}.
\footnote{\weightedpin{grh-compact/src/New/RealPlace/BesselOrthonormal.lean}{8-L26}{Orthonormality},
\weightedpin{grh-compact/src/New/Filling/FillingData.lean}{94-L108}{The common choice of the vectors and kernel},
\weightedpin{grh-compact/src/Lib/General/LinearAlgebra/FramedCoefficient.lean}{121-L147}{The measure-preserving sign frame},
\weightedpin{grh-compact/src/New/Filling/PoincareSeries.lean}{61-L107}{The actual real argument and seed},
\weightedpin{grh-compact/src/New/Filling/PoincareSeries.lean}{173-L196}{The change of measure in the trace}, and
\weightedpin{grh-compact/src/New/Seesaw/BesselGeneric.lean}{53-L95}{The realization on the initial Bessel group}.}

\paragraph{Absolute periodization at the weighted level.}
\teamanchor{proof.weighted-poincare-realization}%
Write \(\Gamma=G_0(F)\), \(U\) for the unipotent of \(B_0\), and
\(U_c=U(F)\cap h_cK(g)h_c^{-1}\) for the rational left
stabilizer of the class within \(U(F)\) in lattice coordinates.
Take \(W_v\), the integrable type kernel \(e_v\) and the
finite positive measure \(\mu_v\) from the
\hyperlink{proof.weighted-real-bessel-vectors}{real construction}.
The kernel fixes the Bessel coordinates by right convolution,
and the change of basis intertwines the compact core action
with the transpose Schur action. For the absolute folding
argument use the finite positive translate majorant
\eqref{eq:weighted-bessel-cut-majorant} for the full vector.
Write \(\mathrm{InCut}_v=\{y:\operatorname{size}(y)\le R\}\),
with \(R=b^{6(\varrho+1)+2}\).
The estimate counts the right translates \(yh\) in this set
against the finite \(\mu_v\). The pointwise gain
\eqref{eq:weighted-bessel-cut-decay} enters the full-to-cut trace
comparison.
This is the \hyperlink{proof.period-bessel-seed-fold}{earlier
absolute folding argument} in the coordinates of the weighted seed.

For \(a=(c,i)\) let \(s_a\) be the class-cutoff seed before
periodization. In the finite coordinate of the filling write \(\eta\) for the inverse of
the central lattice scalar, embedded in the core. Its support
coset is \(D_c=\eta^{-1}h_cK(g)\) in the finite chart
\(x_f\mapsto\eta^{-1}x_f\eta\), with \(x_f\) written in the lattice
basis. On this coset the real factor is \(\beta_c W_i\).
Equivalently, in unconjugated lattice coordinates the support is
\(h_cK(g)\eta^{-1}\). The right level is \(K(g)\) in the
filling chart and \(\eta K(g)\eta^{-1}\) in lattice coordinates.
Left multiplication by \(U_c\) preserves the support, with the left
action transported by the same conjugation in the filling chart.
The norm of a Bessel coordinate is invariant under
\(U_\infty\), since the character has modulus one. Moreover
\(\sum_{u\in U_c}\beta_c(uy)=1\), with nonnegative summands.
Grouping \(\sum_{\gamma\in\Gamma}|s_a(\gamma x)|\) by \(U_c\)
thus removes the cutoff and gives a sum of absolute Bessel
coordinates over the \(U_c\)-cosets satisfying the finite condition.

Here is a uniform bound for this sum. The real size cut has a compact
set \(C_\infty\) of representatives modulo \(U_\infty\). Choose a
Bessel translate within one of the infimum in the size cut. Its
inverse norm equals its norm in the sign frame. The resulting
closed matrix ball in the group is compact, and the core factor
removed from the translating element ranges over a compact group.
The group \(U_c\) is a cocompact
lattice in \(U_\infty\). Indeed, in a central filtration of \(U\)
its rational subgroup cut out by the fixed compact open level has
full lattices in the additive successive quotients. Taking
representatives successively gives a compact set \(C_U\) with
\(U_\infty=U_c C_U\). For fixed \(x\) and \(h=(h_v)_v\), each coset
whose real coordinate after right translation by \(h\) is in the cut
therefore has a representative \(\gamma\) with
\(\gamma x(h,1)\) in a fixed adelic compact \(C\).
In the rational filling coordinates its real part lies in
\(C_U C_\infty\) and its finite part in \(h_cK(g)\eta^{-1}\).
The fixed changes of real and finite frame carry compact sets to
compact sets. In these coordinates the motions by \(\Gamma\) are left
translations. For any two representatives in this compact, their
quotient belongs to \(\Gamma\cap CC^{-1}\), a finite set. Hence the
number of contributing cosets is bounded by its cardinality,
independently of \(x,h\). Tonelli and the product measure
\(\otimes_v\mu_v\) now bound the absolute seed sum uniformly in \(x\).

For the sum with signs, the rational \(U_c\)-character at infinity
is one: its character at finite places is one on the conjugated
level and the global additive character is trivial on rational
points. The same grouping therefore leaves the finite-coset
indicator times \(W_i\) in each coset. The product kernel
\(e=\prod_v e_v\) fixes each such summand. The uniform absolute bound
allows integration against \(|e|\) before summing, so \(R(e)\) fixes
the bounded, measurable periodization. It is continuous at infinity
because translates of \(e\) vary continuously in \(L^1\), and
invariance under the indicated right level gives continuity in
the finite variables. It is
left \(\Gamma\)-invariant by reindexing. These properties put the
series in the indicated space of forms, with its actual Petersson
norm. Its finite linear combination for any column defines
\(\mathcal P\). Since \([\Gamma\backslash G_0(\mathbb A_F)]\) has
finite measure, Tonelli on this quotient gives
\(\int_{G_0(\mathbb A_F)}|s_a|<\infty\). For a bounded test \(z\)
and either \(\Delta=B_0(F)\) or \(\Gamma\), unfolding therefore gives
\[
 \int_{\Delta\backslash G_0(\mathbb A_F)}
       z(g)\sum_{\gamma\in\Delta}s_a(\gamma g)\,dg
       =\int_{G_0(\mathbb A_F)}z(g)s_a(g)\,dg
\]
when \(z\) is left \(\Delta\)-invariant. This proves the two initial
pairings are equal for the prescribed transposed test.
In the Petersson convention linear in the second argument, the
functional \(z\mapsto\int_{G_0(\mathbb A_F)}\overline{s_a(g)}z(g)\,dg\)
on the prescribed type has Riesz vector \(\mathcal P_a\), the
periodization just constructed. Its value is
\(\langle\mathcal P_a,z\rangle\) by unfolding. Multiplying the
functional by the positive trace normalization multiplies this
vector by the same factor.

To identify the initial linear form, write the finite core point
\(k_f\eta\) in its unique column double class. On the principal
coset of that class a rational Bessel element contributing to the
seed has core projection in its stabilizer. The stabilizer is
trivial by the free level calculation, so the remaining rational
translations are exactly \(U_c\). Their cutoffs sum to one. To
reach the same coset from a finite unipotent translate, approximate
that translate by a rational unipotent in the prescribed finite
open coset, successively in the additive quotients of its central
filtration. If \(u\) is the rational approximant to the finite
left translate, the compensating rational action is \(u^{-1}\).
The unbarred \(P_0\) then acquires
\(\chi_\infty(u)^{-1}=\chi_f^+(u)\), and its conjugate in the
Petersson reading acquires
\(\chi_f^-(u)=\psi_f(-\sigma_U(u)/(g\varpi))\).
Both characters are trivial on the remaining principal unipotent
level, where the simple-root sum lies in \(g\varpi\) times the
additive annihilator. Thus the bottom formula above holds with
the Weyl-to-Bessel map \(\mathsf A\) equal to the tensor product of
the given compact intertwining coordinate maps. Its contraction
with \(W(\xi_0)\) is bilinear. The value of the translated column at
\(k_f\eta\) is the original value at \(k_f\), since the central
translation defining its coordinates is by \(\eta^{-1}\).
In particular the resulting coefficient is independent of the
twist. The positive rescaling
\((\mu_f(K(g))\nu)^{-1/2}\) in the trace calculation applies to
these same coordinate columns and this same \(\mathcal P\).
\footnote{\weightedpin{grh-compact/src/New/Filling/NearTerms.lean}{125-L151}{Approximation, the finite character and the class stabilizer},
\weightedpin{grh-compact/src/New/Filling/ColumnCoords.lean}{36-L53}{Compact intertwiners and free class coordinates},
\weightedpin{grh-compact/src/New/Seesaw/BesselGeneric.lean}{53-L95}{The value on the Bessel subgroup}, and
\weightedpin{grh/src/New/Spanning/Unfolding/PoincareSeries.lean}{471-L557}{reproduction after the unipotent sum}.}

\paragraph{Successive periods of the graded string.}
Inductively let \(P_{j+1}\) be the generic function obtained from
\(P_j\) at the outer point \(u_j=(\xi_j,y_j)\) by the constant term
at the string map \(Y_j\to Y_{j+1}\). In the mixed Weil model this
means applying the Heisenberg translation by the string map and
taking the zero coefficient in the nonzero graded directions. The
one-step identity, in the measures of the lower chain, is
\begin{equation}
\label{eq:weighted-seesaw-step}
 \int_{[B_j]}P_j(bu_j)(L_j^{\mathrm t}z)(bu_j)\,db
   =[K_j^0:J_j]\int_{B_{j+1}(F)\backslash G_{j+1}(\mathbb A_F)}
                        P_{j+1}(h)z(h)\,dh .
\end{equation}
Here \([B_j]=B_j(F)\backslash B_j(\mathbb A_F)\) and
\(L_j^{\mathrm t}\) is the transposed lift on conjugate forms for the
bilinear product integral. This formula includes absolute convergence
for the tests of the given types and preservation of the generic
condition for the next step. In the mixed model the nonmatching
rational terms are killed by the unipotent character. The surviving
string-map orbit has the next Bessel group as stabilizer. The factor
\([K_j^0:J_j]\) is the normalization of \(L_j^{\mathrm t}\).
Transposition itself follows by Fubini on the two automorphic
quotients, with the same scalar on the two sides, as in the theta
kernel interchange in \cite{seesaw}.
The interior and boundary calculations below establish this
identity for the normalized arrows, including genericity and
absolute convergence at the indicated types.
\footnote{\weightedpin{grh-compact/src/New/Seesaw/GradingsOfY.lean}
{74-L105}{The normalized arrows}. For the orbit method see
Bryan Wang Peng Jun, \emph{Global theta lifting and automorphic
periods associated to nilpotent orbits},
\href{https://arxiv.org/pdf/2311.14234\#page=17}{Proposition 3.2
and Theorem 4.5}.}

\paragraph{The last three string lengths.}
\teamanchor{proof.weighted-last-seesaw-arrows}%
For a string of length \(s\), its ordered basis has grades
\(s-1,s-3,\ldots,1-s\). The core has grade zero.
For \(s=1\) and \(s=0\) all grades are zero. The two final arrows
have lengths \(1\to0\to0\). Their string maps are zero, and the
positive-degree spaces in the constant terms are zero. For the
first arrow the Bessel group is the definite group fixing its one
vector, and for each zero string the Bessel group is the whole
definite unitary group. Accordingly the one-step formula on these
arrows follows by Fubini: the kernel in \(P_{j+1}\) is exactly the
theta kernel at \((bu_j,h)\), with both quotients compact. Its
factor is \([K_j^0:J_j]\) from \(L_j^{\mathrm t}\). The definition of
\(P_{j+1}\) uses precisely the \(db\) on the left of the one-step
formula. In particular the different choice of the grading constant
at \(Y_m\) has no effect on this calculation.\footnote{\weightedpin
{grh-compact/src/Lib/Weil/Periods/DegreeAndUnipotent.lean}{40-L73}{The grading and Bessel subgroup},
\weightedpin{grh-compact/src/Lib/Weil/Periods/FiniteBesselGroup.lean}{18-L32}{the string map},
\weightedpin{grh-compact/src/New/Seesaw/GradingsOfY.lean}{191-L209}{the last two graded spaces}.}

When \(m\ge4\), the preceding arrow has lengths \(2\to1\).
In its mixed expansion the rational negative-degree entries form
one column \(v=w(e_+)\), where \(e_+\) is the source vector of
grade one. All target grades are zero. The central unipotent
\(e_-\mapsto e_-+t e_+\), with \(\bar t=-t\), acts on the term by
the character of the Hermitian norm of \(v\). The prescribed
generic character supplies the inverse character at
\(v_*=c z'\), where \(z'\) is the sole target string vector and
\(c\) is the common lower grading constant. Indeed the string map
sends \(e_+\) to \(c z'\). Write \(T\) for this map and
\(\dagger\) for Hermitian adjoint. The off-diagonal entries of the
source hyperbolic Gram are one. Thus the lower left entries of
\(T^\dagger T\) and \(w^\dagger w\) are \(-Nc\bar c\) and
\((v,v)\). Since \(u-1=tE_{+,-}\), the character quotient
\(\beta(Tu,T)/\beta(wu,w)\) is
\(\psi((j_{\rm Weil}/2)t(-Nc\bar c-(v,v)))\), by
\(\beta(x,y)=\psi((j_{\rm Weil}/2)\operatorname{tr}(y^\dagger x))\).
Integrating
with its probability Haar measure on the compact additive quotient
for \(t\) inside \(db\) kills exactly the
terms with unequal norms. The norm \(-Nc\bar c\) is nonzero.
By rational Hermitian Witt extension the remaining vectors are the
orbit of \(c z'\), with stabilizer \(U(X')=B_{j+1}\).
Unfolding the rational orbit gives the integral over
\(B_{j+1}(F)\backslash G_{j+1}(\mathbb A_F)\), with counting measure
on the rational points. On compact sets of outer variables the
same Schwartz seminorm bounds are uniform, giving continuity.
Rational invariance follows by reindexing by the rational
stabilizer. These are all the genericity conditions on the
one-string target, whose unipotent group is trivial.

For the \(2\to1\) arrow the
\hyperlink{proof.weighted-mixed-fourier-expansion}{mixed expansion}
has precisely the single column just described. The bound in the
\hyperlink{proof.weighted-seesaw-absolute-orbits}{absolute orbit
calculation} applies with weight zero on the target. It permits the
norm-character selection and the unfolding by counting measure on
its rational orbit. The target has no other negative-degree entries.%
\footnote{\weightedpin{grh-compact/src/Lib/Weil/Periods/BesselIntegral.lean}{26-L35}{The character in adjoint coordinates},
\weightedpin{grh-compact/src/Lib/Weil/Periods/UnfoldingStep.lean}{33-L59}{the selected frequency},
\weightedpin{grh-compact/src/Lib/Weil/Periods/UnfoldingStep.lean}{89-L113}{rational Witt comparison}, and
\weightedpin{grh/src/Lib/Weil/Periods/Unfolding.lean}{512-L572}{orbit integral against the rapid test}.}

\paragraph{The rational string orbit at the interior steps.}
\teamanchor{proof.weighted-graded-string-orbit}%
Write \(Y=Y_j\), \(Y'=Y_{j+1}\), \(s=m-1-j\ge3\),
and write \(\lambda,\lambda'\) for the two gradings. The string
bases have grades \(s-1-2i\) and \(s-2-2i\). The core spaces have
grade zero. Opposite string vectors pair to one, with middle norm
\(-N\) when there is a middle vector. The map \(T:Y\to Y'\)
has degree \(-1\) and kills the core. On string vectors that have
a successor in \(Y'\), its coefficients \(t_i\) are \(1,-Nc,c\)
in positive, zero and negative source grades if \(s\) is odd, and
\(c,1\) in positive and negative source grades if \(s\) is even.
The last source vector maps to zero. Thus \(A=T^\dagger T\) is
the string shift \(e_i\mapsto a_i e_{i+1}\) with \(a_i\ne0\)
for \(i<s-1\), and zero on the core. In the next grading
\(TT^\dagger\) is the same shift prescribed there, since the
two constants \(c\) agree.
\footnote{\weightedpin{grh-compact/src/Lib/Weil/Periods/FiniteBesselGroup.lean}
{14-L51}{String maps and genericity}, and
\weightedpin{grh-compact/src/Lib/Weil/Periods/DegreeOfT.lean}
{34-L43}{The two shift identities}.}

\paragraph{The mixed Fourier expansion in the two gradings.}
\teamanchor{proof.weighted-mixed-fourier-expansion}%
Keep the actual weighted datum, the outer point \(u_j\), and one
finite Schwartz vector and real Fock polynomial of the arrow.
Write \(M=\operatorname{Hom}_E(Y,Y')\) in the rational frames and
give a matrix entry from source grade \(a\) to target weight \(b\)
degree \(b-a\). Use either target weight system \(d=0\) or
\(d=\lambda'\). Write \(M_-,M_0,M_+\) for strictly negative, zero
and strictly positive degree. The Hermitian trace pairing on \(M\)
pairs degrees \(a\) and \(-a\), since the Grams of the two strings
pair opposite grades and the cores have grade zero. In particular
\(M=M_-\oplus M_0\oplus M_+\) is a hyperbolic splitting with
retained space \(M_0\).

Let \(K=M_+\) and write \([K]=K(E)\backslash K(\mathbb A_E)\) with
probability Haar. In the fixed oscillator write \(\Theta\) for the
theta functional and \(\rho\) for Heisenberg translation, with
\(\rho(v)\rho(w)=\beta(v,w)\rho(v+w)\) and
\(\beta(v,w)=\psi((j_{\rm Weil}/2)\operatorname{tr}(w^\dagger v))\).
The function \(A_f(z)=\Theta(\rho(z)f)\) on \(K(\mathbb A_E)\)
is periodic modulo \(K(E)\). Indeed \(\beta\) is trivial on
\(K\times K\), and the rational Heisenberg translations fix
\(\Theta\). For each rational \(w\), put
\begin{equation}
\label{eq:weighted-mixed-Fourier-proof}
 C_d(w;f)
  =\int_{[K]}\Theta(\rho(z)\rho(w)f)\,dz
  =\int_{[K]}\beta(z,w)^2 A_f(z)\,dz .
\end{equation}
The second equality commutes \(\rho(w)\) to the left and uses
rational invariance. By adelic trace duality and the opposite-degree
pairing, \(\chi_w(z)=\beta(w,z)^2\) runs once through the characters
of \([K]\) when \(w\) ranges over \(M(E)/M_{\ge0}(E)\). Thus the
displayed integral is the Fourier coefficient in the \(\chi_w\)
convention. This checks both the sign of the phase and the counting
on the rational labels.

For the fixed finite data the function \(A_f\) is invariant under
an open compact subgroup of \(K(\mathbb A_{E,f})\), since Heisenberg
translation is smooth at the finite places. Strong approximation
for the additive group identifies the resulting quotient of
\([K]\) with a real torus: its lattice is the intersection of
\(K(E)\) with that open compact subgroup in the finite coordinates.
In particular the characters that occur form a lattice of real
frequencies. If its real dimension is \(D\), application of
\((1-\Delta)^b\) for \(2b>D\) bounds the absolute sum of coefficients
by a constant times the maximum of finitely many real derivatives
of \(A_f\). Indeed the individual bound is the same maximum times
\((1+4\pi^2\|\xi\|^2)^{-b}\) at frequency \(\xi\). Fourier inversion
at zero proves \(\Theta(f)=\sum_w C_d(w;f)\) with absolute convergence.

For the agreement chart, let \(S\) consist of the entries with
source grade \(a>\max(0,\lambda'(b))\), where \(b\) now denotes the
target basis index. The dual coordinates to \(S\) span a rational
subspace \(K_S\) of \(K\) in both gradings. Requiring \(w\) to
agree with \(T\) on \(S\) is exactly
\(\chi_w|_{[K_S]}=\chi_T|_{[K_S]}\). Summing the absolutely
convergent series on this subtorus against \(\chi_T^{-1}\) gives,
in either model, the same probability mean
\begin{equation}
\label{eq:weighted-mixed-agreement-mean}
 \sum_{w:\ w|_S=T|_S} C_d(w;f)
   =\int_{[K_S]}\beta(z,T)^2\Theta(\rho(z)f)\,dz .
\end{equation}
For \(z=y^\dagger\), the phase on the right is
\(\psi(j_{\rm Weil}\operatorname{tr}(yT))^{-1}\), since
\(j_{\rm Weil}\) is anti-invariant under conjugation and \(\psi\)
is given by the trace. This is the adjoint-coordinate reading of
the common mean.\footnote{\weightedpin
{grh-compact/src/Lib/Weil/Periods/BesselIntegral.lean}{18-L36}
{The original rational-coordinate Heisenberg action},
\weightedpin{grh-compact/src/Lib/Weil/Periods/UnfoldingStep.lean}
{138-L224}{The two graded quotients and the agreement sum}.}

\paragraph{The normalized zero term.}
\teamanchor{proof.weighted-seesaw-zero-term}%
In a hyperbolic splitting \(M=L\oplus U\oplus K\), choose the
Schr\"odinger model retaining \(U\) and taking the mixed variable
\(x\) along \(L\). Use the same pair of rational dual frames at
all places, so that \(K\) acts by \(\rho(z)F(x)=\beta(x,z)^2F(x)\)
and \(L\) translates \(x\). The retained vector action is the
oscillator of \(U\). In these coordinates the rational theta
functional is a fixed nonzero scalar \(\kappa\) times the sum of
the retained theta functional over the rational \(x\). This
follows from the rational Heisenberg invariance in the ordinary
Schr\"odinger realization: invariance under rational modulations
restricts the support of the functional to the rational lattice,
and invariance under rational translations gives its counting
measure. Use the ratio of the original theta functionals for
\(\kappa\). The average over \([K]\) kills every \(x\ne0\), so the
normalized zero term is \(Zf=\kappa F(0)\). For a retained
Heisenberg translate \(y\) the same argument gives
\[
 \Theta_U(\rho_U(y)Zf)
   =\int_{[K]}\Theta(\rho(y)\rho(z)f)\,dz .
\]
This fixes the normalization of \(Zf\) and gives
\[
 C_d(w;f)=\Theta_U(Z\rho(w)f).
\]
At each place evaluate the variables dual to the same rational
isotropic space. In this chart a finite Bruhat--Schwartz vector
remains locally constant and compactly supported after evaluation.
On a pure tensor the real and finite evaluations separate with the
fixed scalar. For the paired arrows use the real adapted vacuum
that defines the finite map and assign its fixed value to that map.
Other real inputs give one retained real vector, independent of the
finite input. The actual \(w\) is in rational coordinates, so the
finite shift keeps the original frame and the factor
\(\alpha_{\rm W}^{-1}\).

For \(g\) stabilizing \(K\), let \(a=g|_K\) in the frame of \(K\)
and \(g_0\) be its induced action on \(K^\perp/K\). We claim
\begin{equation}
\label{eq:weighted-seesaw-normalized-zero-term}
 Z\omega(g)=\upsilon(\det a)^{-1}|\det a|_E^{1/2}\,
                 \omega_U(g_0)Z .
\end{equation}
To calculate the scalar, first use the covariance of the
Heisenberg representation. In the mixed \(x\) coordinate,
invariance under all \(K(\mathbb A_E)\)-modulations characterizes
the distributions that factor through \(F(0)\). The characters
separate every nonzero \(x\); at a real place their derivatives
also rule out derivatives of evaluation. At a finite place
locally constant test functions on compact sets give the same
support assertion. The covariance with the retained Heisenberg
translations then determines the operator on \(F(0)\) up to
scalar, by irreducibility of the retained oscillator for the
Heisenberg group. This argument also works for the continuous
Schwartz operators, since their kernels commuting with modulation
and translation are scalar multiples of the diagonal kernel.

The scalar on the unipotent radical of the stabilizer of \(K\)
is a continuous character. It is trivial. At each place, scalar
dilation \(t\) on \(K\), the identity on \(U\), and \(t^{-1}\)
on \(L\), for \(t\) in the base field, contracts the radical as
\(t\) tends to zero. Conjugating leaves a character unchanged.
Elements supported at finitely many places are dense in the
adelic radical, giving the same conclusion globally.
The Levi action on the retained factor is \(\omega_U(g_0)\):
the orthogonal sum has tensor vacuum and the normalized vacuum
coefficients multiply. It remains to calculate the scalar for
\(g_a=\operatorname{diag}(a^{-*},1,a)\) on \(L,U,K\).
Covariance with translation along \(L\) sends \(F(x)\) to
\(c(a)F(a^*x)\). Unitarity in the mixed \(L^2\) measure gives
\(|c(a)|=|\det a|_E^{1/2}\).

Use the source's adapted vacuum convention for the phase. In the
adapted real frame the hyperbolic mixed vacuum is a positive
Gaussian up to a constant, by the unitary Fourier transform of
the Gaussian. In the adapted finite frame it is constant on the
integral lattice and zero off it. Indeed the integral \(K\)-moves
force self-dual lattice support by modulation, and the integral
\(L\)-moves force constancy there by translation. The integral
hyperbolic frame and the conductor in these coordinates are those
of the standard vacuum, including at the places over two.
Consequently the vacuum matrix coefficient of \(g_a\) is
\(c(a)\) times a strictly positive real number: it pairs this
nonnegative profile with its \(a^*\)-dilate, whose support meets
it in a neighborhood of zero.

Let \(k\) be the fixed adapted move from the standard vacuum
frame. By unitarity, the coefficient of the adapted vacuum at
\(g_a\) is the standard vacuum coefficient at \(k^{-1}g_a k\),
namely \(\overline{\Delta(k^{-1}g_a k)}^{-1}\) in the Hermitian
form linear in its second argument. In the following block
calculation \(\Delta\) is computed in this standard frame.
At a real place the negative block of the transformed \(g_a\) is
\((a^{-*}+a)/2=a^{-*}(I+a^*a)/2\), and \(\Delta_\infty\)
is its conjugate determinant. Its phase is therefore
\(\upsilon_\infty(\det a)^{-1}\). At the finite places, write
\(D_f(A)=\chi_f(\det u_A)\), where the columns of \(u_A\)
span \(\widehat{\mathcal O}_E^p\cap A^{-1}\widehat{\mathcal O}_E^p\)
and \(\chi_f=\upsilon_f|\cdot|_E^{-1/2}\). This is independent of
the integral basis since \(\chi_f\) is unramified. The identity
\(A(\widehat{\mathcal O}_E^p\cap A^{-1}\widehat{\mathcal O}_E^p)
=\widehat{\mathcal O}_E^p\cap A\widehat{\mathcal O}_E^p\)
gives \(D_f(A^{-1})=\chi_f(\det A)D_f(A)\).
Smith factors show \(D_f(A^*)=\overline{D_f(A)}\), using
\(\chi_f(\bar d)=\overline{\chi_f(d)}\).
The integral hyperbolic frame therefore gives
\[
 \Delta_f(g_a)=D_f(a^{-*})D_f(a)
             =\chi_f(\overline{\det a})\,|D_f(a)|^2 .
\]
Thus \(\Delta(g_a)\upsilon(\det a)>0\), and the same phase
calculation and the modulus above prove
\eqref{eq:weighted-seesaw-normalized-zero-term}. For the middle
pair use \((uw)^rv^{-(2r-1)}\) on the positive-degree part and
\(v^{-1}\) on the removed plane. These ideles have modulus one,
so their two contributions change the exponents by
\((-r,2r,-r)\) at \((u,v,w)\). The terminal zero string uses no
positive-degree factor.\footnote{\weightedpin
{grh-compact/src/Lib/Weil/MixedModel/ZeroTerm.lean}{265-L283}
{The actual adapted signed and integral frames},
\weightedpin{grh-compact/src/Lib/Weil/Construction/PairDefs.lean}
{202-L218}{The fixed vacuum coefficient and translated evaluation},
\weightedpin{grh-compact/src/Lib/Weil/Construction/FiniteMultiplier.lean}
{69-L118}{The fixed finite multiplier},
\weightedpin{grh-compact/src/Lib/Weil/MixedModel/ZeroTermLaw.lean}
{67-L106}{The normalized parabolic factor},
\weightedpin{grh-compact/src/New/Seesaw/EdgePair.lean}
{22-L41}{The finite part of the actual paired evaluation}.}

For rational \(\gamma\in B_j(E^+)\) and \(\delta\) preserving the
filtration defined by the target weights,
\(C_d(w;\omega(\gamma,\delta)f)\) is the coefficient at
\(\delta^{-1}w\gamma\). In the mean in
\eqref{eq:weighted-mixed-Fourier-proof} this follows by rational
invariance of \(\Theta\) and change of variables on \([K]\).
For adelic moves raising the respective grades strictly, the action
on \(M\) is the identity modulo a degree-raising operator. Its
action on \(K\) has determinant one and its retained action is the
identity. Hence the same substitution follows by
\eqref{eq:weighted-seesaw-normalized-zero-term}; the determinant
characters of the arrow are one on these unipotents.
When a move changes \(w\) to \(w+q\) with \(q\in K(\mathbb A_E)\),
the Heisenberg law and \(\rho(q)\)-invariance of \(Z\) give
\(C_d(w+q;f)=\beta(w,w+q)C_d(w;f)\). In particular the source
degree-two test changes the integrand by
\(\beta(Tu,T)/\beta(wu,w)\). On the target the shift
\(v'^{-1}T-T\) is strictly positive when \(v'-1\) raises grades
by two. The same rule gives
\(\beta(T,v'^{-1}T)=\psi((j_{\rm Weil}/2)\operatorname{tr}
(TT^\dagger v'))\). Thus the retained \(T\)-coefficient has the
required next genericity, using the original common constant \(c\).


\paragraph{The real types of the fixed lower-chain tests.}
\teamanchor{proof.weighted-seesaw-signed-types}%
At a real place put \(R_j=r+jt\) and \(\epsilon_j=j\bmod 2\).
In doubled Harish--Chandra parameters the entries on \(Y_j\)
are the following entries minus \(\epsilon_j\):
\[
 \begin{array}{ll}
 k-2-2i & (0\le i<k-1),\\
 2h_2-2w_{a+1} & (0\le a<m),\\
 t-1-2w_\ell-2b & (1\le\ell\le j,\ 0\le b<t).
 \end{array}
\]
Only the middle list can have negative signs. Its sign at \(a\)
is \((-1)^{a+\min(a,j)}\). Write \(w=w_{j+1}\) and put
\(\gamma=\lfloor R_j/2\rfloor-w\),
\(\gamma'=\lfloor R_{j+1}/2\rfloor-w\).
These are the two actual real weights of the rank-twisted
character, since \(\vartheta\) has real weight zero.
The Fock rows are on \(Y_{j+1}\) and the columns are the
conjugate side on \(Y_j\). Their determinant exponents are
\(\gamma\) and \(R_{j+1}-\gamma'\).
Consequently, for the actual doubled parameters \(l_j,l_{j+1}\),
the centered row and column entries are
\[
 l_{j+1}-2\gamma+R_j,\qquad -l_j+2\gamma'-R_{j+1}.
\]
Since \(r,t\) are odd, \(R_j\bmod2=\epsilon_{j+1}\) and
\(R_{j+1}\bmod2=\epsilon_j\). The centers of the common
entries are therefore \(x\) and \(-x\), with \(x\) the
unshifted doubled entry plus \(2w\). Every \(x\) in the first
and the old added lists is \(>t\) and keeps its sign. The stem
centers are \(2h_2+2(j-a)h\), with \(h=2h_2\).
They are \(>t\) for \(a\le j\) and \(<-t\) for \(a\ge j+1\).
Exactly the latter signs reverse. The new block has positive
sign and centers \(t-1-2b\), all of absolute value \(<t\).
This proves the signed correspondence in the paired Fock
vacuum, including the final definite row.

The parameter is regular at both nodes, since the same-sign
strings have distinct entries two apart and the centers of
different lists are disjoint. For opposite signs the doubled
stem-to-stem gap is at least \(2h\). A negative stem index
has \(a\ge j+1\), so its distance from any old added-block
center is at least \(3h\), before subtracting \(t-1\).
The distance from the compact first list is larger still.
On the target, the new positive string is at a distance of at least
\(3h-(t-1)>2h\) in doubled units from any negative stem entry, since
the closest possible negative stem has \(a=j+2\).
Thus the opposite-sign gap of the undoubled parameter is at
least \(h\). The datum has \(100m\le k\le r\) and \(t=5r\),
so \(R_j+R_{j+1}<r^2\). With \(h=2r^8\) the required
gap \(500(R_j+R_{j+1})^3\) follows.

Here is the highest polynomial on this matching. For a common
row with center \(x_i\) and sign \(a_i\), its corresponding
column has sign \(a'_i=a_i\operatorname{sgn}(x_i)\).
For any row \(i\), subtract its vacuum weight from its lowest
compact weight, and call the result \(d_i\). The formula
\(\Lambda=\lambda+\rho_{\rm n}-\rho_{\rm c}\) gives
\[
 2d_i=x_i+\sum_u a_i a_u\operatorname{sgn}(x_u-x_i)
                     -\sum_v a_i a'_v .
\]
The row sum includes the \(t\) new entries; the column sum is
over the old entries. For a new row numbered \(b\) in the added
string, the matched terms cancel, and the remaining sum is
\(2b-t+1\). Thus its relative weight is zero.
For a common row add its formula to the column formula for
its mate. The two center terms cancel. A second matched pair
with centers of the same sign gives zero in the two sums.
For opposite signs, \(\operatorname{sgn}(x_u-x_i)
=\operatorname{sgn}(x_u)\), and the sign reversal gives zero
again. The new rows have centers between the old positive and
negative centers, so their contribution also cancels. Hence the
two relative weights are negatives of each other. The sign of
\(d_i\) on a common row is the sign of \(x_i\), since the sum
of the two ranks bounds the correction and \(|x_i|\) exceeds
that sum in the actual datum.

For each pair of row and column signs and each height \(q\ge0\)
form the minor on the matched pairs of that signature with
\(|d_i|>q\). Multiply the nonempty minors. For either sign of
\(d_i\) the rows chosen are an initial segment in the compact
order appropriate to that variable sign, and so are the paired
columns in the swapped order. Indeed the vacuum-subtracted
weights are monotone among rows of the same sign in the
Harish--Chandra order. For neighboring such rows this follows
from the displayed formula: their highest-weight difference
is at least half their doubled gap minus one. Rows with the
other variable sign are outside the initial segment. The
compact raising differential operators therefore give
determinants with a repeated row or column. The diagonal
weights count the occurrences with their variable signs and
give the desired lowest compact weights. Each row and each
column uses variables of only one sign, so all cross-sign
Laplacians kill the product. It is nonzero on the matrix
with the matching entries one and all other entries zero.
This supplies the joint highest line of the two actual lowest
compact types. The chosen polynomial is kept in that line,
whose multiplicity one follows from the joint harmonic
decomposition of Howe \cite[Theorem~2.2(1)]{p2paul-real-theta}.
\footnote{\weightedpin{grh-compact/src/New/Exchange/GridParameters.lean}
{30-L65}{The actual signs, weights and doubled lists},
\weightedpin{grh-compact/src/Lib/DiscreteSeries/Theta/OscillatorDefs.lean}
{49-L58}{The centered and relative weights},
\weightedpin{grh-compact/src/Lib/DiscreteSeries/Theta/JointHarmonicExists.lean}
{42-L93}{The joint highest line in the actual oscillator}.}


\paragraph{Rapidity and the real line in the two theta integrals.}
\teamanchor{proof.weighted-seesaw-rapid-test}%
For the real correspondence divide the two determinant twists
out of the covered oscillator. In the convention with the
source acting by the right inverse, the source column
vacuum after removing its determinant center is
\(-a'_i(p'-q')/2\), where \((p',q')\) is the target
signature. Passing to the source action by the inverse
transpose negates the torus weights and preserves the
Hermitian sign spaces. The source vacuum then has
weights \(a'_i(p'-q')/2\) and the target vacuum has weights
\(a_i(p-q)/2\) for source signature \((p,q)\), precisely
Paul's covers on the two signed groups. The determinant
centers removed on those two sides are
\(\gamma'-R_{j+1}/2\) and \(\gamma-R_j/2\).
Thus the smaller genuine parameter is \((x_i/2)\) with
the source signs, and the larger one has these entries with
the target signs and the positive string of length \(t\).
Both are integral on the genuine covers, whose ranks have
opposite parities. The source action by the inverse transpose
also accounts for the dual in the bilinear integral. Equivalently
one can use the determinant-twisted splitting after
\(\chi_V(z)=(z/|z|)^{2\gamma'-R_{j+1}}\) and
\(\chi_W(z)=(z/|z|)^{2\gamma-R_j}\) in the left-action convention.
Their restrictions to \(\mathbb R^\times\) have the required
parities. The proof of Ichino's Lemma~7.10 uses this inverse
transpose and \((\chi_W,\chi_V^{-1})\) in the right-action
model. Its vacuum shifts are the ones just computed.
The actual global kernel keeps its two original characters and
the same finite vector.

Paul's Theorem~2.7(b) applies because the absolute values of the
common centers after halving are at least \((t+1)/2\),
the target reverses exactly the signs of negative entries,
and all \(t\) new entries are positive-labeled. In reverse,
Proposition~3.4(ii) deletes the maximal positive-labeled
central string, which is maximal because both adjacent
integers are absent. It also gives first occurrence in that
Witt tower. For the forward tower use Atobe's criterion in
Ichino's Theorem~6.1(i).
The smaller centered parameter has no zero and is integral,
so \(k_\pi=-1\). The counts \(r_\pi,s_\pi\) are the numbers
of positive and negative labels after reversing the negative
entries. The signature of \(Y_{j+1}\) is
\((r_\pi+t,s_\pi)\), which uses \(l=0\) and
\(t_{\rm Ichino}=(t-1)/2\).
The two near sets are empty by the strict gaps. Removing a
hyperbolic plane makes \(l<0\) while keeping the offset
\(t_{\rm Ichino}\), so every smaller member in that tower
vanishes. In the zero negative-index case this condition is
vacuous. For transposition the kernel is the same function
with its arguments interchanged. The conjugate target input
pairs with the larger representation on the oscillator side,
and the output belongs to the conjugate of the smaller one.
If a skew-Hermitian model is used for the other member of
the pair, identify the two roles by the same trace-zero
multiple of the forms as in the symmetric tower formula.
\cite[Theorem~2.7(b) and Propositions~2.9 and~3.4(ii)]{p2paul-real-theta}
\cite[Theorem~6.1(i), Remarks~6.2 and~7.3 and Lemma~7.10]{p2ichino-real-theta}

The input at each arrow is a bounded automorphic vector fixed by
the rank-one real discrete-series kernel. At fixed level it is
in \(L^2\). Formal-degree orthogonality puts it on the lowest
compact line of copies of the prescribed discrete series, and
its real translates generate those copies. A compactly
supported smooth kernel acts
as the identity on this line: average a sufficiently small
smooth bump on the two sides against its rank-one compact
idempotent and normalize the nonzero resulting coefficient.
Thus the input is smooth. For the constant-term test use the
following discrete-series assertion on \(SU(p,q)\).
A square-integrable, bi-compact-finite and center-finite
matrix coefficient belongs to the Harish--Chandra Schwartz
space. A center-finite function in that Schwartz space has
zero integrals along every proper parabolic unipotent
radical, including after left and right translation. The
normalized coefficient here is bi-compact-finite, has one
infinitesimal character and is square integrable by
formal-degree orthogonality. The compact central factor
separates as a scalar character and every such radical lies
in the derived group, so the assertion applies on \(U(p,q)\).
\footnote{\href{https://repository.ias.ac.in/30833/1/336.pdf\#page=77}
{Harish-Chandra, \emph{Discrete series for semisimple Lie groups II}
(1966), \S33, Corollary~1 of Lemma~65, pp.~77--78}, and
\href{https://djvu.online/file/KpNCz2R803FHi}
{Wallach, \emph{Real Reductive Groups I}, \S7.2.2, pp.~233--234}.}
Take the
automorphic constant term of the convolution reproducing the
input. The constant term as a function of the real convolution
variable is left invariant by the relevant conjugated real
unipotent radical. Disintegrating along that radical makes the
integral of the discrete-series kernel zero. Boundedness and
integrability of the kernel justify the interchange. The
integrability follows from the opposite-sign gaps by the
discrete-series integrability criterion. In type \(A_{R-1}\)
the root bound in that criterion is
\(\frac14\sum_\beta|(\alpha,\beta)|=R-1\) for squared root
length two. Work on the finite cover by \(SU(p,q)\) times
the compact center. Rapid decay follows on each finite-level
component by the cuspidal decay theorem for these smooth,
center-finite forms.
\cite[Part~III, \S3, Corollary to Theorem~1]{p2milicic-integrability}
\cite[Theorems~2.23--2.24]{p2miller-schmid-decay}

For the fixed finite vector and finite-level representatives, the
Schwartz estimate for the ordinary theta kernel gives
\[
 |\Theta(\omega(g,h)(\phi\otimes f_f))|
     \le C\,p(\phi)(\|g_\infty\|\|h_\infty\|)^A
\]
on reduced real representatives. Here \(p\) is a fixed real Schwartz
seminorm, \(A\) is fixed and the group norms dominate the matrix
and inverse norms. This follows by the lattice sum and the
polynomial seminorm bound for the Weil action given in the
absolute-orbit calculation below. A right differential operator
in \(h_\infty\) replaces \(\phi\) by its corresponding oscillator
derivative, leaving \(p\) and \(A\) fixed. Rapid decay on the input
justifies differentiation under the theta integral and gives
uniform moderate growth of the lifted function. Rational invariance
comes from the theta sum, and an open compact stabilizer of
the given finite vector on the output side gives its level.
Write \(\mathcal T(\phi,z)\) for the bilinear theta integral and
\(R\) for right translation. Right Haar measure and the commuting
actions give
\[
 \mathcal T(\omega(a,1)\phi,R(a)z)=\mathcal T(\phi,z),
 \qquad
 \mathcal T(\omega(1,b)\phi,z)=R(b)\mathcal T(\phi,z).
\]
Fix a target point and the other local vectors in a cuspidal
constituent of the input. Varying the oscillator vector and
the source vector at one real place gives a continuous
bilinear pairing invariant under simultaneous source
translation: moving a right source translation off the
oscillator gives its inverse on the form. The smooth-vector
seminorm bound in the cuspidal decay theorem supplies continuity
in the source vector, and the displayed bound supplies continuity
in the Schwartz vector. In its polynomial
and compact-finite vectors a nonzero pairing would give a
surjection onto the contragredient source module, by
irreducibility. The polynomial vectors are dense in the
Schwartz realization, so the same vanishing criterion
applies to the restricted Schwartz vector in the constant
term. By Howe's correspondence all such pairings vanish
against the smaller Witt-tower members at a real place. The
constant-term formula for a unitary theta lift expresses the
proper-parabolic constant terms as those lower lifts. When
the target is lowered by \(2u\) dimensions, the isotropic map
block consists of the maps into its removed \(u\)-plane. The
input element \(g\) acts there with determinant
\((\det g)^{-u}\) by precomposition. The zero-term law
\eqref{eq:weighted-seesaw-normalized-zero-term} therefore
changes the exponent in the inverse input character from
\(\gamma'\) to \(\gamma'-u\). The genuine center
\(\gamma'-R_{j+1}/2\) is preserved when \(R_{j+1}\)
decreases by \(2u\). In the swapped arrow lower the original
source by \(2u\). On the removed isotropic map block the
input \(h\) has determinant \((\det h)^u\) by postcomposition.
This changes \(\gamma\) to \(\gamma-u\), preserving
\(\gamma-R_j/2\). All these determinants have adelic norm
one, and the base character \(\mu_{j+1}^0\vartheta^{j+1}\)
is unchanged. This checks the input splitting in the two
tower applications, including the conjugate orientation.
The lower vanishing now gives zero proper constant terms in
each direction. Use the growth reduction in the proof of
Miller--Schmid's Theorem~2.24, \S5 after (5.11), directly on
this smooth uniformly moderate function. On a Siegel set choose
a maximal-parabolic radical for a largest split simple-root
coordinate. The zero mean bounds the function on each compact
unipotent quotient by a sum of its unipotent derivatives.
Moving those derivatives to the right gives factors
\(e^{-\alpha(\log a)}\), with each root \(\alpha\) in that
radical bounded below by the chosen simple-root coordinate up
to a constant. Right derivatives have the same zero means and
the same growth exponent. Iteration gives any negative power
of the reduction height. This proves rapid decay and puts the
lift in the cuspidal \(L^2\)-space before identifying its
infinitesimal character there.
\cite[Proposition~3.1, including its symmetric version]{p2wu-unitary-tower}
\footnote{Use \href{https://www.math.utah.edu/~ptrapa/math-library/howe/howe-transcending-1989.pdf}
{Howe, Theorem~1}, on the smooth real oscillator.
See also \href{https://arxiv.org/pdf/math/0610860v3\#page=15}
{Khare--Larsen--Savin, \S4.1}, for the same translated-coefficient
vanishing used in the unfolded constant term.}

On the cuspidal target take an irreducible constituent of the
lift. Pair the theta integral against this constituent, fixing the
vectors at all other places. The pairing is continuous in the local
oscillator and the two local smooth vectors by the same Schwartz
estimate and the input and output cuspidal seminorm bounds.
On polynomial oscillator vectors the source and target compact
equivariance put the pairing in the compact-finite parts of the
two distribution duals. These are the usual Harish--Chandra
contragredients by admissibility. With the conjugation on the
output test this gives a nonzero quotient onto the tensor of
the contragredient source and the output representation. Real
Howe duality identifies the output with the paired discrete series.
The actual joint highest polynomial restricts the output to the
highest line of the multiplicity-one lowest compact type by
compact equivariance. In the transpose it is the conjugate
line of the smaller representation, with the reversed compact
positive order. The formal-degree matrix coefficient of the
specified unit vector in that line, in the original Haar
measure, fixes the line exactly by Schur orthogonality.
This argument applies to each finite-level
cuspidal input by linearity on its fixed type, and to the
actual finite inclusion vector without further normalization.
It proves rapidity of the test and the required real line
for its successive transposes. The sole index in the transpose
remains \([K_j^0:J_j]\).

\paragraph{Absolute sums in the orbit integrals.}
\teamanchor{proof.weighted-seesaw-absolute-orbits}%
Let \(z\) be the rapid automorphic test of the conjugate target
type. First take \(d=0\). The full absolute sum of mixed terms
is invariant by reindexing under the rational target group, and
under simultaneous rational source substitution with the labels.
Choose a compact set of representatives for the source Bessel
quotient, which is compact because its core is definite and its
unipotent fibers have compact rational quotients. Keep \(u_j\)
fixed. At the finite places use the finitely many adelic class
representatives of the target at a common level of the finite
datum and \(z\). The finite translates on the compact source
representatives take finitely many values. Hence one finite open
compact subgroup and one torus lattice suffice in the bound
following \eqref{eq:weighted-mixed-Fourier-proof}.

For the real bound choose reduced target representatives at
those finite points. Each derivative on the fixed compact torus
is a Heisenberg derivative applied to the real Schwartz vector.
A lattice sum in a fixed Schr\"odinger polarization is bounded
by a Schwartz seminorm of order greater than its real lattice
dimension. The real metaplectic action on a fixed seminorm is
bounded by a power of the norm of its matrix and inverse:
the compact factors act boundedly, and the split diagonal action
dilates the variables and their derivatives with its determinant
factor. Heisenberg derivatives are polynomial differential
operators in this chart. The source representatives and the
torus representatives range over compact sets. Thus for some
\(C,A\) the absolute sum on these representatives is at most
\(C\,\operatorname{ht}(h)^A\), where \(\operatorname{ht}\ge1\)
is the reduction height of the target. The determinant phases
have modulus one. The constants may depend on the fixed datum,
test, and outer point, as allowed by the tested period.
Since \(z\) is rapid and the target quotient has finite volume,
the full absolute sum multiplied by \(|P_j(ru_j)z(h)|\) is
integrable in the two original quotient measures. The same
seminorm bound on compact sets of \(h\) works in either grading
and is locally uniform in the outer point.

Use Tonelli on this absolute sum, with counting measure on the
rational orbit labels. If \(\Gamma\) is the source rational group
and \(S_\xi\) the stabilizer of a label, summing
\(k_{\gamma\xi}(r)=k_\xi(\gamma^{-1}r)\) on
\(\Gamma/S_\xi\) then integrating over \(\Gamma\backslash B_j\)
is integration of \(k_\xi\) on \(S_\xi\backslash B_j\).
This is the quotient integration formula with the same Haar and
counting measures on the two discrete subgroups. For a nontrivial
character test \(k_\xi(ur)=c_uk_\xi(r)\), integrate next through
the closed subgroup of left translations acting on the function
by unit scalars. It contains \(S_\xi\) and \(u\), and is closed
by continuity and evaluation at a point where the function is
nonzero. On each fiber right multiplication by \(u\) is defined
and scales Haar by a positive factor. It also gives the unit
scalar \(c_u\), so the fiber integral vanishes if \(c_u\ne1\).
This argument on the stabilizer quotient uses integration in
stages and needs no invariance of that quotient under left
multiplication by \(u\).

For the target chart unfold the original \(d=0\) translates first.
The absolute integral of their sum over the target rational
quotient bounds the integral of the absolute chart sum over
\(B_{j+1}(E^+)\backslash G_{j+1}(\mathbb A_{E^+})\).
At each fixed target point use
\eqref{eq:weighted-mixed-agreement-mean} and the locally absolute
actual grading to reduce the chart to \(T\). This proves absolute
integrability of the surviving period on the larger quotient,
with its quotient measure. The factor \([K_j^0:J_j]\) is still the
normalization of the transposed lift. The case \(2\to1\) is the
same argument with the single column and trivial target grading,
followed by the norm-character selection already computed for
that arrow.\footnote{\hyperlink{proof.period-grid-prefixes}
{Polynomial theta growth and rapid decay in the earlier route},
\weightedpin{grh-compact/src/Lib/Weil/Periods/UnfoldingStep.lean}
{313-L390}{The two absolute orbit integrations},
\weightedpin{grh-compact/src/Lib/Adelic/QuotientMeasure/CocompactIntegrable.lean}
{128-L177}{Unfolding and cancellation on a discrete stabilizer quotient}. The fixed-source passages also include
\weightedpin{grh-compact/src/Lib/Weil/Periods/UnfoldingStep.lean}
{138-L206}{the mixed coefficients},
\weightedpin{grh-compact/src/Lib/Weil/Periods/UnfoldingStep.lean}
{227-L261}{the substitution phase},
\weightedpin{grh-compact/src/Lib/Adelic/QuotientMeasure/CocompactIntegrable.lean}
{167-L177}{the twist on the stabilizer quotient},
\weightedpin{grh-compact/src/Lib/Weil/Periods/UnfoldingStep.lean}
{308-L390}{the absolute orbit integrations},
\weightedpin{grh-compact/src/Lib/Weil/Lift/Arrow.lean}
{39-L48}{rapid tests and transpose}, and
\weightedpin{grh-compact/src/New/Seesaw/GradingsOfY.lean}
{74-L95}{the actual lower-chain arrows}.}

Here is the calculation of the surviving labels. For a vector
\(v\) of strictly positive source grades and a vector \(x\)
such that \((x,v)=0\), use the Eichler matrix
\[
 u=1+v x^\dagger-x v^\dagger-\tfrac12(x,x)v v^\dagger .
\]
This is \(1+D+D^2/2\) for \(D=v x^\dagger-x v^\dagger\),
with \(D^\dagger=-D\) and \(D^3=0\). In the tests below its
nonconstant terms raise grades by at least two, and
\(w(u-1)\) is strictly positive in the indicated model. When
\((x,x)\ne0\) in these tests, \(Tv\) and \(wv\) are isotropic
(the zero vector is allowed). Thus the last term makes zero
contribution to the two traces. The phase comparison, with
\(x\) multiplied by arbitrary adelic scalars, gives
\[
                      (wx,wv)=(Tx,Tv).
\]
Indeed the two remaining rank-one traces are conjugates with
opposite signs, and multiplication by \(j_{\rm Weil}\) and
varying the scalar tests the difference by the nondegenerate
adelic additive pairing.

For weights \(d=0\), first take \(v,x\) in the strictly positive
source space \(Y_{>0}\). Their adjoints are supported on the
negative columns, so all the degree conditions hold. This gives
equality of the two Gram forms on \(Y_{>0}\). To prove
injectivity there, if \(wv=0\) for a nonzero \(v\in Y_{>0}\),
choose its lowest nonzero grade \(a\) and take \(x\) to be dual
to the successor of the string vector of grade \(a\).
Thus \(x\) has grade \(2-a\), and the terms of \(u-1\) raise
grades by at least two. The term involving \(wx\) lies on
negative columns and the terms involving \(wv\) vanish. If
\(x\) has grade zero then \(Tv\) has strictly positive grades.
The comparison gives \((x,Av)=0\), whereas the coefficient of
the successor of the lowest entry of \(v\) in \(Av\) is
nonzero. Now rational Witt extension applies to the isometric
embeddings of \(Y_{>0}\) given by \(w\) and \(T\), including
its possibly degenerate Gram form. A rational target unitary
move therefore makes the entire restriction to \(Y_{>0}\)
agree with \(T\).

For the passage to the actual weights, keep a coarser chart of
labels: require equality with \(T\) only on entries whose
source grade \(a>0\) is strictly larger than the target grade.
The chart is stable under the rational source \(B_j(F)\).
Its translates by \(G_{j+1}(F)\) cover all labels not deleted
in the weights-zero model. Two translates intersect only when
their quotient belongs to \(B_{j+1}(F)\). To see this, for
each nonnegative target string vector \(f_i\) of grade \(b\),
take the source vector \(e_i\) of grade \(b+1\). Its image in
the chart is \(t_i f_i\) modulo grades \(>b\). Descending through
these \(b\), a rational target move carrying one chart label to
another must fix \(f_i\) modulo grades \(>b\). These are the
defining conditions of \(B_{j+1}\), and they also imply that
every element of \(B_{j+1}\) preserves the chart. Thus one
unfolds the chart translates to \(B_{j+1}(F)\backslash
G_{j+1}(\mathbb A_F)\). The sum on this chart agrees in the two models by
\eqref{eq:weighted-mixed-agreement-mean}, applied to these prescribed
entries. Use the original weights-zero absolute integral before
the target chart unfolding and then use the actual-grading
coefficient at each unfolded target point.

In the actual weights, the remaining entries to compare have
target grade \(b<0\) and source grade \(b<a\le0\). Work in
increasing \(b\). For a source basis vector \(x\) of grade \(a\),
\(wx\) has no grades below \(b\) by induction and the degree of
\(T\). Let \(v=e_i\) be the source vector of grade \(1-b\)
mapping to \(t_i f_i\), where \(f_i\) is the target vector
dual to the row of grade \(b\). In the chart, \(wv\) is
\(t_i f_i\) plus terms of grades \(>-b\). In the Eichler
test the term \(wv x^\dagger\) has degree at least \(a-b\),
and \(wx v^\dagger\) has degree at least one. The correction
term also has positive degree. Both \(Tv,wv\) have positive
grades. The identity \((wx,wv)=(Tx,Tv)\) now reads off
the coefficient in the specified row of grade \(b\), because
all higher grades in \(wv\) pair with grades below \(b\).
This proves \(w-T\) has no negative degree entries, and hence
only the class \(T\) contributes in the actual model.
The source Bessel group preserves that rational class, since
it fixes the half-string modulo higher grades and the
nonconstant core block is killed by \(T\).

The remaining coefficient is exactly the integral defining
\(P_{j+1}\). It is continuous by compactness of the rational
source Bessel quotient and the same seminorm bound locally in
the outer variables. The rational target Bessel group fixes
the frequency class. If \(v'-1\) raises target grades by at
least two, the shift from \(T\) to \(v'^{-1}T\) is strictly
positive. The Heisenberg law now gives the character
\(\psi((j_{\rm Weil}/2)\operatorname{tr}(TT^\dagger v'))\);
the equality of \(c\) identifies it with the next prescribed
shift. This proves the genericity and the period formula
\eqref{eq:weighted-seesaw-step}. Rational orbit unfolding uses
the original quotient measures, and the only lift factor in
that formula is \([K_j^0:J_j]\).
\footnote{\weightedpin{grh-compact/src/Lib/Weil/Periods/UnfoldingStep.lean}
{43-L113}{Eichler tests, Witt extension and the chart stabilizer},
\weightedpin{grh-compact/src/Lib/Weil/Periods/UnfoldingStep.lean}
{208-L224}{Partial inversion in the two models}, and
\weightedpin{grh-compact/src/Lib/Weil/Periods/UnfoldingStep.lean}
{265-L290}{The next generic character}.}

\paragraph{The paired tensor at the outer points.}
Pair edges \(2i,2i+1\), \(0\le i<m/2\). For
\(y=(y_0,\ldots,y_{m-1},1)\), translate the finite inclusion datum
of arrow \(j\) by \((y_j,y_{j+1})\). Take the paired zero term
in \(\operatorname{Hom}(Y_{2i+1},Y_{2i}^{-}\oplus Y_{2i+2})\)
at the two translated string maps, with finite tensor
\(E_{i,\vartheta}(y)\). Put
\(E_\vartheta(y)=\bigotimes_{i<m/2}E_{i,\vartheta}(y)\).
The next calculations realize the paired terms as coefficients
of these definite lifts. The ramified support calculation then
bounds the possible adapted profiles. The integral orbit
comparison turns those bounds into finite cuts.

\paragraph{The integral orbit of the initial half-string.}
\teamanchor{proof.weighted-adapted-line-orbits}%
Set \(p_i=e_i\) for the string basis with \(0\le i<s\) and
\(i^\vee=s-1-i\). An adapted basis \(b\) is block upper triangular
for the order which gives the core index \(2s\) and \(e_i\) the
index \(4i+1\). Its columns span \(x_Fy\widehat{\mathcal O}_E^n\),
where \(x_F\) is the finite frame and \(y\) the stage point.
These bases exist at every finite place by triangular elimination
over its valuation ring along this ordered flag, and the choices
are integral frames at almost every place. Denote the string
diagonal by \(\beta_i\). For \(i\le i^\vee\) the column \(b_i\)
equals \(\beta_i e_i\) modulo earlier string vectors, with no core
entry. For \(i,j\) in this first half,
\[
       (b_i,b_j)=\overline{\beta_i}\,\beta_j(e_i,e_j).
\]
Indeed \(i+j\le s-1\), and an earlier vector cannot pair with the
other index or an earlier one, by the anti-diagonal string Gram
matrix. Its only possibly nonzero entry on these columns is the
norm at the middle vector in an odd string.

Suppose \(y,y'\) give equal valuations of \(\beta_i,\beta'_i\) at
every finite place for \(i\le i^\vee\). Multiplying the corresponding
columns of \(b'\) by the units \(\beta_i/\beta'_i\) makes these
diagonals equal. In the two integral orthonormal models the frames
\[
  B_i=(x_Fy)^{-1}b_i,\qquad B'_i=(x_Fy')^{-1}b'_i
                    \quad(i\le i^\vee)
\]
are primitive, being part of invertible integral matrices. They
have equal Gram matrices by the displayed identity. We can
therefore choose an integral unitary \(k\) carrying \(B_i\) to
\(B'_i\) simultaneously, including when their Gram matrix is
degenerate.

For completeness, at each finite place of \(F\) the quadratic
algebra is split or unramified. Here \(E\) contains \(i\) and
\(\omega=(1+\sqrt{-7})/2\), so
\(i\omega-\overline{i\omega}=i\) is a unit. The order generated
locally by \(i\omega\) has unit discriminant. Modulo the uniformizer
the isometry between the primitive frames extends by Witt's
theorem for Hermitian spaces. In the split algebra, the first
frame specifies independent vectors and the conjugate frame
independent covectors. Their pairings fix the induced map on
the quotient by their common covector kernel. Extend the given
map first on the intersection of the vector span with that kernel,
then to the whole kernel and finally on complementary vectors
using the prescribed quotient map. This gives the simultaneous
change of basis for the two halves.

To lift an approximate integral isometry \(g\) matching the frames
modulo \(\varpi^a\), \(a\ge1\), write its correction as
\(g(1+\varpi^a T)\). Over the residue algebra the two equations are
\[
  T+T^*=H,\qquad T B=D,\qquad
  H=(1-g^*g)/\varpi^a,\quad D=g^{-1}(B'-gB)/\varpi^a .
\]
The adjoint here is in the integral orthonormal frame.
Equality of the original Gram matrices gives
\(B^*HB=B^*D+D^*B\) after reduction. Extend the prescribed map
\(B\mapsto D\) to any endomorphism \(T_0\). The Hermitian error
\(H-T_0-T_0^*\) vanishes on the product of the span of \(B\) with
itself. Choose a vector space complement. This error is \(T_1+T_1^*\)
with \(T_1B=0\): in sesquilinear form coordinates assign the mixed
block on the side with second argument in the complement, and
split the complementary Hermitian block using an element of
residue trace one. The same construction works over the split
residue algebra and in characteristic two. Use \(T=T_0+T_1\)
and pass to the inverse limit. Taking the product of these local
solutions gives \(k\in K_j^0\). This is the lifting calculation for
\hyperlink{proof.primitive-frame-densities}{primitive Hermitian
frames}, with the additional columns prescribed at each depth.

The rational-coordinate matrix \(x_Fy'k(x_Fy)^{-1}\) now carries
\(b_i\) to \(b'_i\). Its action on \(e_i\) is the identity modulo
the span of the earlier \(e_t\), by induction up the half-string.
These are the defining conditions of \(B_j\); they include the
middle vector. Thus \(y'ky^{-1}\) belongs to
\(B_j(\mathbb A_{F,f})\), proving the orbit assertion used for the
finite cut. Its separate support input is the bound on the
valuations for a nonzero unfolded ramified column.
\footnote{\weightedpin{grh-compact/src/New/Seesaw/EqualLinesOrbit.lean}{29-L59}{Adapted frames and Gram entries},
\weightedpin{grh-compact/src/New/Seesaw/EqualLinesOrbit.lean}{81-L106}{the half-string orbit}, and
\weightedpin{grh-compact/src/Lib/Adelic/Witt/SpaceOperations.lean}{37-L59}{integral extension for primitive frames}.}

\paragraph{A compact edge in the weighted normalization.}
\teamanchor{proof.weighted-definite-edge-coefficient}%
On the edge \(X_j\to X_{j+1}\) let \(V,V'\) be its coefficient
spaces, with real actions \(\rho,\rho'\) in the conjugate Weyl
convention. For the two ranks \(n<n'\) take the prescribed
diagram, retaining the columns of heights \(<n'\). There are
no columns of heights strictly between \(n\) and \(n'\), since
the heights above one have the form \(k+i(t+1)\) and the
consecutive ranks differ by \(t+1\). At each embedding the
incoming determinant exponent is the incoming real splitting
weight \(\lfloor n'/2\rfloor-w_{j+1}\) plus the multiplicity
\(\kappa(n)\) of source-height columns. This is zero at \(j=0\)
by the formula for \(\kappa(k)\). At \(j>0\) it equals
\(\lfloor n_{j-1}/2\rfloor-w_j\), by the constant weight
increment \(h\), the column count \(h-(t+1)\) and
\(\lfloor n'/2\rfloor-\lfloor n_{j-1}/2\rfloor=t+1\).
The outgoing exponent is \(\lfloor n/2\rfloor-w_{j+1}\).
The twist has real weight zero, so both real actions are fixed
for the entire family. These identities identify the action
on the rectangular Schur coefficient polynomials
\(P_{\ell,v}(Z)=\ell(S_\kappa(Z)v)\):
\[
 \Omega(x,y)P_{\ell,v}=P_{\ell\circ\rho'(y^{-1}),\,\rho(x)v}.
\]
Indeed the Fock substitution acts on the two sides of \(Z\),
and the maximal minors of height \(n\) contribute the
\(\kappa(n)\) extra powers of the source determinant. The
remaining determinant powers are those of the real splitting
characters with the given rank twists.

Let \(K^{X}_j\) be the actual source column level, and let
\(\mathcal J_{\vartheta}\) be the Schur jet matrix at scale
\(c_{\rm jet}=\overline{\pi_G}\,\mathrm{per}\ne0\).

\leavevmode\teamanchor{proof.weighted-fock-jet-coefficient}%
Write \(\Theta\) for the theta functional and \(\mathsf v(P)\otimes f\)
for the insertion of the real Fock polynomial \(P\) in a finite
datum \(f\). In the matrix coordinates at all real places write
\(\alpha!=\prod_i\alpha_i!\). The generating series at scale \(c\)
and its differential pairing are
\[
 T_c(f)(Z)=\sum_\alpha
  \frac{c^{|\alpha|}}{\alpha!}\,
   \Theta(\mathsf v(Z^\alpha)\otimes f)\,Z^\alpha ,
 \qquad
 P(\partial)T_c(f)(0)=\Theta(\mathsf v(P(c\,\cdot))\otimes f).
\]
Indeed, for \(P=\sum_\alpha P_\alpha Z^\alpha\), the derivative
multiplies the \(\alpha\)-coefficient by \(P_\alpha\alpha!\).
The pairing is complex bilinear. For the finite part of the
positive definite model this series gives the ordinary theta jets
as well. In the \(d_0\) complex Fock coordinates put
\(d\mu(w)=\pi^{-d_0}\,dw\), with \(dw\) Lebesgue measure, and
write \(\mathsf H(w)\) for real Heisenberg translation in these
coordinates. Set
\(G_f(w)=e^{-\|w\|^2/2}\Theta(\mathsf H(w)f)\).
The integrated action gives
\(\Theta(\mathsf v(P)\otimes f)=\int P(\bar w)G_f(w)\,d\mu(w)\).
The real vacuum projection fixes \(f\). The translated vacuum
at \(w\) is \(e^{-\|w\|^2/2}e^{-\|y\|^2/2}
\exp(\sum_i y_i\overline{w_i})\). Dividing by the untranslated
vacuum value \(e^{-\|y\|^2/2}\) at \(y\) gives the entire theta function as
\(\int \exp(\sum_i y_i\overline{w_i})G_f(w)\,d\mu(w)\).
On the line \(y=tc z\), its \(m\)-th derivative at zero inserts
\((c\sum_i z_i\overline{w_i})^m\). The Heisenberg translates have
bounded theta value, so Gaussian decay dominates these derivatives
uniformly on compact \(t\)-disks. The multinomial formula gives
\(m![t^m]T_c(f)(t z)\).

The entries of \(\mathcal J_\vartheta(f)\) are the jets of
\(T_{c_{\rm jet}}(f)\) paired with the entries of \(S_\kappa(Z)\).
Each entry is homogeneous of degree
\(D_j=\sum_{\tau}\sum_{h<n'} h\kappa_\tau(h)\), including all
retained columns of source height \(n\). The finite action changes
\(f\) to \(f'=\omega_f(g_f,h_f)f\): it commutes with the real
Heisenberg operators used for the insertion. The real action on the
polynomial is the covariance calculated above. Thus, for every
\(\ell\in(V')^*\) and \(v\in V\),
\[
 c_{\rm jet}^{D_j}\,
 \Theta\!\left(\omega((x,g_f),(y,h_f))
      (\mathsf v(P_{\ell,v})\otimes f)\right)
 =\ell\!\left(\rho'(y^{-1})\,\mathcal J_\vartheta(f')\,\rho(x)v\right).
\]
This identity holds at each finite translate and for every retained
degree, including degree zero.
\footnote{\weightedpin{grh/src/Lib/Weil/Construction/ThetaJets.lean}{84-L128}
{The factorials and the Schur jet matrix},
\weightedpin{grh/src/Lib/Weil/MixedModel/RealTheta.lean}{163-L186}
{the Gaussian integral of the translated vacuum},
\weightedpin{grh/src/New/Derivatives/Vacuum/TightFrames.lean}{473-L506}
{its directional derivatives}, and
\weightedpin{grh/src/Lib/Weil/DualPair/WeilAction.lean}{662-L673}
{the action on the real insertion and finite datum}.}

The \hyperlink{proof.compact-rectangular-average}{compact average}
applies to \(B_{\ell,v}=P_{\ell,v}\), using the covariance just
computed and the nonzero Schur minor on the first \(n\) rows.
It gives one \(\ell(P,t_0)\) for an arbitrary real Fock polynomial
\(P\) and \(t_0\in V^*\). For \(P=P_{\ell_0,v_0}\), the source
average conjugates the rank-one endomorphism \(v_0t_0\) of \(V\).
It commutes with the source action and has trace \(t_0(v_0)\), so
\(\ell(P_{\ell_0,v_0},t_0)=t_0(v_0)\ell_0/\dim V\).

Combine the coefficient identity with the compact average.
For all finite data \(f\), smooth algebraic columns \(\phi\),
twists and target points \(h=(y,h_f)\), integrating the theta
kernel with arbitrary Fock datum \(P\otimes f\) against
the coefficient \(t_0(\rho(x^{-1})\phi(g_f))\) gives
\[
 \frac{c_{\rm jet}^{-D_j}}{[K_{X_j}^0:K^{X}_j]}\,
   \ell(P,t_0)\bigl(\rho'(y^{-1})(\mathscr L_{j,\vartheta}(f)\phi)(h_f)\bigr).
\]
Here \(\mathscr L_{j,\vartheta}\) is the prescribed algebraic
class sum, normalized to give \(K^{X}_j\) volume one. In fact
the finite integrand after the compact average is exactly
\(\mathcal J_{\vartheta}(\omega_f(g_f,h_f)f)\phi(g_f)\).
It is left rational invariant and has an open stabilizer. At
any finer level \(K'\subset K^{X}_j\) fixing it, a class has
mass \(\operatorname{vol}(K')\) divided by the order of its
rational stabilizer. The analytic measure gives \(K_{X_j}^0\)
volume one, so its class masses are
\([K_{X_j}^0:K^{X}_j]^{-1}\) times the algebraic masses.
The same ratio therefore holds at every refinement and is
independent of the twist and finite data. The formula gives
a single output functional uniform in all of them. This
compact argument applies to both edges of the last pair as
well as the edges of all preceding pairs.
\footnote{\weightedpin{grh-compact/src/New/Exchange/CrudeModularData.lean}
{27-L44}{The actual column weights, levels and jet scale},
\weightedpin{grh-compact/src/Lib/Weil/CompactLift/BranchingPeriod.lean}
{90-L139}{The transformation and homogeneity of the Schur jet},
\weightedpin{grh/src/Lib/Weil/CompactLift/ThetaKernelDefs.lean}
{343-L366}{The exact oscillator coefficient of the scaled jet},
and \weightedpin{grh/src/Lib/Weil/CompactLift/DefiniteThetaLift.lean}
{273-L321}{The uniform compact-edge interface}.}

\paragraph{A middle pair of graded arrows.}
\teamanchor{proof.weighted-middle-double-lift}%
Use \(Y,Y',Y''\) for a pair of lengths \(2r+1,2r,2r-1\)
and respective full ranks \(d,d',d''\). The lower ranks
increase by \(t\) at each step, so \(d''-d=2t\) and the outer
ranks have the same parity. On \(\operatorname{Hom}(Y',Y^-\oplus Y'')\)
the first and second data combine by adjoint and flip in the
first summand. The oscillator functional on the combined data
is a fixed scalar times the product of the two theta functionals.
In a product Schr\"odinger model the rational sums multiply,
and the scaled adjoint identifies the first rational lattice
sum with that in the flipped form. The intertwining scalar
depends on the two chosen models. The transported Heisenberg
frequency is the stacked map \(S=(T^\dagger,T')\), with
\[
              S^\dagger S=-TT^\dagger+T'^\dagger T'=0 .
\]
Each term on the right is the shift on \(Y'\) by the common
constant \(c\). By averaging the positive-degree maps, the
product of the two constant terms thus becomes the zero term
of the combined datum at \(S\), with the reciprocal of the
fixed theta scalar. The positive-degree variables agree under
the adjoint because it preserves map degree when reversing
source and target, as their Gram dualities negate the grades.

The zero-degree maps are from \(X'\) to
\(X^-\oplus X''\oplus H\), where \(H\) is the plane of the two
outer middle vectors \(z,z''\). Indeed the outer nonzero grades
are even and the middle string grades are odd. In \(Y^-\oplus
Y''\) the vector \(z+z''\) is isotropic, and
\((z-z'')/(2N)\) is its dual isotropic vector. To integrate over
the source Bessel quotient, the source unipotent has only
degrees at least two. Its character cancels the phase of the
first constant term at \(T\). Its fiber average therefore
leaves only the source core variable. If \(u'\) is in the
middle unipotent, the two frequency maps move to
\(u'^{-1}T,T'u'\), so the new stacked frequency is \(Su'\).
The difference \(S(u'-1)\) has degrees at least zero and its
degree-zero part maps \(X'\) into the isotropic line
\(E(z+z'')\). More precisely its \(X'\) row along this line is
\[
                c\,(u'-1)_{\text{string grade }1,\,X'} .
\]
Here the grade-one vector maps under \(S\) to \(c(z+z'')\),
as follows by taking the adjoint across the middle entries
of the string Grams. All other entries of \(S(u'-1)\) have
positive degree, because \(S\) has degree \(-1\) and
degree-zero maps have just the source and target indicated.
The Heisenberg phase from composing with the shift by \(S\)
is one by \(S^\dagger S=0\). The positive-degree part
disappears in the first zero term, and it pairs trivially with
the retained degree-zero part.

The displayed grade-one row is additive on the unipotent:
quadratic terms raise grades by at least two. It is onto,
both rationally and adelically, by the Eichler matrix with
the grade-one vector and any core vector. Thus the normalized
unipotent quotient measure pushes to probability Haar on the
additive rational quotient of this row. This is exactly the
second mixed zero term, along \(\operatorname{Hom}(X',E(z+z''))\).
Both Bessel quotients are compact. Their fiber masses over the
definite core quotients are fixed constants, so this pair of
averages introduces a single scalar fixed by the graded
spaces and measures.

For the action of the core triple let \(u,v,w\) now denote
the three determinants. With \(\upsilon\) the fixed splitting
character, the underlying Weil action before the zero terms
has extra normalized \(\upsilon\)-exponents
\[
 d'-\lfloor d'/2\rfloor,\qquad
 \lfloor d/2\rfloor-\lfloor d''/2\rfloor-d,\qquad
 \lfloor d'/2\rfloor
\]
at \(u,v,w\), respectively. The \(-d\) and \(d'\) come from
the adjoint flip, and the two base characters are evaluated at
the same determinants before and after restriction. On the
positive-degree maps the determinant is
\((uw)^r v^{-(2r-1)}\): each of the \(r\) negative middle
grades admits both outer core blocks, and each of the
\(2r-1\) positive outer grades admits the contragredient
middle core block. All remaining string blocks are fixed.
On the removed plane the isotropic block has determinant
\(v^{-1}\). Use the normalized parabolic law
\eqref{eq:weighted-seesaw-normalized-zero-term} for both blocks.
Their determinants are norm-one ideles, so the modulus is one.
The exponents change by \((-r,+2r,-r)\). These are exactly
the changes obtained by using the smaller ranks
\((d-2r-1,d'-2r,d''-2r+1)\), by the equal parity of \(d,d''\).
This proves equivariance with the normalized triple Weil
action of the three cores. The shift by \(S\) commutes with
their simultaneous core action, since the string maps kill
the cores and their images lie on the strings.

At any fixed tuple of real outer variables the moved data
split into a real vector and the translated finite tensor.
The two zero terms split in the same way: in adapted mixed
models they evaluate the removed Fourier variables at zero,
separately at the real and finite places. Their normalized
finite map is precisely \(E_{i,\vartheta}(y)\), including the
finite \(\alpha_{\rm W}^{-1}S\) shift. The fixed choice of
polarization identifies the retained oscillator with the two
definite oscillator spaces, and acts by the same finite
intertwiner for all input data. For decomposable polynomial
data the two compact integrals on the cores give the two
successive weighted class sums by the compact-edge calculation
above. A polynomial in the two sets of variables is a finite
sum of such products. For arbitrary real vectors the same
identity follows by Hermite approximation: the value of each
test integral is continuous in the real Schwartz vector, and
the family of evaluations furnished by the finite-dimensional
last coefficient space is a closed subspace of the product
of the scalar test spaces. Consequently one coefficient
functional works for all finite tensors, twists, columns and
evaluation points with those fixed real variables. The same
argument for the final pair uses \(z/\pi\) instead of the plane.
It has no positive-grade or plane term, and changing ranks to
\((d-1,d',d''+1)\) preserves the displayed exponents by
equal parity.

To state the induction, let \(\Phi_{i,\vartheta,y}\) be the
column after the first \(i+1\) pairs with the finite tensors
\(E_{0,\vartheta}(y),\ldots,E_{i,\vartheta}(y)\).
Fix the real outer variables and a real point \(x\) at stage
\(2i+2\). For some functional \(\ell_{\xi,x}\), all twists,
columns, finite tuples with \(y_0\in J_g\), and core points
\(k\) satisfy
\[
 P_{2i+2}\bigl(\jmath_{2i+2}(k)(x,y_{2i+2})\bigr)
 =\ell_{\xi,x}\bigl(\rho_i((\operatorname{emb}_i k)_\infty^{-1})
                 \Phi_{i,\vartheta,y}((\operatorname{emb}_i k)_f)\bigr).
\]
Here \(\rho_i\) is the outgoing conjugate Weyl action.
The bottom Bessel evaluation is the initial functional.
The two compact integrals take its current value on the
embedded core at the given outer point and produce the
functional for the next pair. The finite component of that
same outer point is used on both adjacent edge data. The
level-independent scalar of the pair is absorbed in the
functional. This proves the claimed coefficient identity
on the whole family by induction. The zero-term calculation
also gives equivariance of each finite evaluation under
independent finite core moves at its three vertices. In
the terminal pair the third group is \(U(Y_m)\), acting on
\(X_m\) by the block embedding fixing the extra line.
\footnote{\weightedpin{grh-compact/src/New/Seesaw/EdgePair.lean}
{141-L224}{The positive-grade action and the finite evaluation},
\weightedpin{grh-compact/src/New/Seesaw/EdgePair.lean}
{235-L256}{The plane and last line},
\weightedpin{grh-compact/src/Lib/Weil/MixedModel/ZeroTermLaw.lean}
{95-L106}{The parabolic multiplier}, and
\weightedpin{grh/src/New/UpperBound/Seesaw/UnfoldingStages.lean}
{319-L471}{The middle row and the double integral}.}

\paragraph{One valuation step on the ramified columns.}
\teamanchor{proof.weighted-ramified-line-support}%
Use additive valuations \(\operatorname{ord}_w\). For the
finite Weil idele \(\alpha_{\rm W}\), let
\(a=\alpha_{\rm W}/\bar\gamma\), where \(\gamma\) generates the
chosen finite modulus for \(g\), i.e. it has the valuations
of \(\overline{\pi_G}\). The modulus generator \(b\) for an
edge with twist has
\(\operatorname{ord}_w b=\operatorname{ord}_w\gamma+{\bf1}_{w\in Q}\),
where \(Q=Q_\vartheta\) is the chosen half of the twist places.
The weighted frames have \(\operatorname{ord}_w c
=\operatorname{ord}_w(a\bar a)\). This uses the self-dual
finite normalization of \(\alpha_{\rm W}\) for the different
and \(c=-(j_{\rm Weil}g\varpi)^{-1}\). With
\(n=\lfloor(s+2)/2\rfloor\) at string length \(s\), the standard
profile \(\beta_i^0\) is that of
\[
 (a/\bar a)^n
 \begin{cases}
  1 & s\ \text{odd},\\
  a^{-1} & s\ \text{even},\ 2i<s,\\
  \bar a & s\ \text{even},\ 2i\ge s .
 \end{cases}
\]
Writing \(i^\vee=s-1-i\), these entries satisfy
\(\operatorname{ord}_w\beta_i^0+
 \operatorname{ord}_{\bar w}\beta_{i^\vee}^0=0\).
In particular the odd gluing column
\(e/\bar\pi+u_j z/N\) changes the finite frame but leaves
the adapted middle line \(u_j z\) in this profile.
\footnote{\weightedpin{grh-compact/src/New/Trace/GoodIdempotent.lean}
{56-L71}{The finite idele and lattice scale},
\weightedpin{grh-compact/src/New/Seesaw/PairTensorData.lean}
{191-L205}{The valuations of the grading constant},
and \weightedpin{grh-compact/src/New/Seesaw/LatticeAndTorsionDefs.lean}
{13-L26}{The actual finite lattice}.}

The finite inclusion datum is fixed by the Heisenberg shifts
of \(b\widehat{\mathcal O}_E\)-vectors. Indeed its finite-cell
expansion has shifts in \(\bar b^{-1}\widehat{\mathcal O}_E\);
the translates by \(b\widehat{\mathcal O}_E\) commute with
these shifts by self-duality and fix the vacuum. Suppose an
edge constant term at \(T\) is nonzero at finite points with
rational-coordinate lattice bases \(R,R'\). The mixed
zero-term support argument gives, for an integral matrix \(L\),
\[
          (\bar b/\alpha_{\rm W})T-R'LR^{-1}
                  \quad\text{of degree at least zero}.
\]
To see this, set \(g\) to be the finite unitary action on
the Hom frame, so the moved datum is fixed by the lattice
\(b g\widehat{\mathcal O}_E\). Pair its intersection with the
positive isotropic part against \(\alpha_{\rm W}^{-1}T\) in
finite coordinates. A nonzero zero term requires trivial
commutator on that intersection. The intersection is a
primitive submodule of the full lattice, so each integral
functional there extends to the full lattice. By self-duality
the frequency agrees on these pairings with an element of
\(\bar b^{-1}g\widehat{\mathcal O}_E\). The annihilator is
exactly the maps of nonnegative degree, giving the formula
after conversion to rational coordinates.

Take adapted bases \(B,B'\) with string diagonals \(\beta,\beta'\).
In the source use the dual basis in reverse string order,
\(\widetilde B=J^{-1}B^{-*}P_s\), with \(J\) the source Gram
and \(P_s\) the permutation reversing its string and fixing the
core. This is again a lattice basis, since \(R^* J R=1\).
Both \(\widetilde B\) and \(B'^{-1}\) preserve the grade
filtrations. On a string vector of grade \(\lambda_i\) the
degree-\(\lambda_i\) string entry of \(\widetilde B\) is
\((J_{i,i^\vee}\overline{\beta_{i^\vee}})^{-1}\).
For the middle vector the degree-zero block may also have a core
component, killed by \(T\). In the target odd middle, the
adapted order puts the middle vector before the core, so
\(B'^{-1}\) sends this row component to itself times
\((\beta'_i)^{-1}\) modulo higher grades. Thus the degree
\(-1\) entry at matched indices in \(B'^{-1}(\bar b/\alpha_{\rm W})
T\widetilde B\) is
\[
       \frac{\bar b}{\alpha_{\rm W}}\,
       \frac{t_i}{J_{i,i^\vee}\overline{\beta_{i^\vee}}\beta'_i}.
\]
It is integral by the support formula, since
\(B'^{-1}R' L R^{-1}\widetilde B\) is integral.
The occurrence of \(J_{i,i^\vee}=-N\) at the odd source middle
cancels the factor \(-N\) in \(t_i\). Writing \(A_w=\operatorname{ord}_w a\)
and \(C_w=\operatorname{ord}_w c=A_w+A_{\bar w}\), the coefficient
\(t_i/J_{i,i^\vee}\) has valuation \(0\) on positive odd
source grades and on negative even source grades, and \(C_w\)
on the remaining grades. Inserting the standard profile
shows that its change from source to target is exactly that
valuation minus \(A_w\): \(n\) stays the same for odd to even
and decreases by one for even to odd. Accordingly the
deviations \(e_i(w)=\operatorname{ord}_w(\beta_i/\beta_i^0)\)
satisfy%
\footnote{\weightedpin{grh-compact/src/New/Seesaw/EdgePair.lean}
{43-L54}{The support annihilator},
\weightedpin{grh-compact/src/New/Seesaw/LinesStep.lean}
{53-L94}{The dual source matrix at the matched entries}, and
\weightedpin{grh/src/New/UpperBound/Seesaw/UnfoldingStages.lean}
{943-L1001}{The twist valuation and Heisenberg invariance}.}
\[
               e'_i(w)\le -e_{i^\vee}(\bar w)+{\bf1}_{\bar w\in Q}.
\]

The integral Gram of the adapted target frame gives
\(e'_i(w)+e'_{i^\vee}(\bar w)\ge0\)
whenever the target Gram pairing is a unit. At a target odd
middle and \(w\mid N\) it yields the lower bound \(-1\)
instead. To exclude equality \(-1\), name the place in the
conjugate pair dividing \(\pi\) as \(w\),
and put \(i=(s'-1)/2\) for this target length \(s'\).
At \(w\) replace the target lattice by a diagonal lattice with
core part standard, string diagonals \(\beta'_t(w)\) for
\(t\le i\), and \(1/\overline{\beta'_{t^\vee}}(w)\) for
\(t>i\). Its component at \(\bar w\) is its Gram dual in
the split factors. This defines a finite group point, since
at a split place the factor at \(w\) may be changed freely
and determines the factor at \(\bar w\). The initial half
valuations are unchanged at \(w\) and \(\bar w\):
for \(t<i\) use the opposite diagonal entries, and at
\(i\) use
\(\operatorname{ord}_w\beta'_i+\operatorname{ord}_{\bar w}\beta'_i=-1\).
Outside \(w,\bar w\) the lattice point is unchanged. The
integral half-frame orbit calculation therefore permits a
finite Bessel move of the target point \(x'\) to this split
lattice, up to its integral lattice stabilizer. Some such
move \(x''=d_Bx'\) then satisfies
\[
     (x'')^{-1}\jmath_j(\nu_w I)x''\in J_j ,
\]
where \(j\) is the target stage and \(\nu_w\) is the
norm-one scalar of the ramification operator. It is integral,
supported on \(w,\bar w\), and preserves the split lattice.
The level at these places is maximal because \(N\) is prime
to the chosen and twist moduli.

In this odd-string target \(j\) is even. The nonzero stage is
therefore read by the preceding paired definite chain. In
reading it at \(x''\), write the inverse Bessel move as a
unipotent times a core move. The unipotent gives only the
prescribed generic phase, since all its grades are even. Fix
the real variables and use the even-stage coefficient formula
at \(x''\) and at the associated core point. Simultaneous
translation by \(\nu_w I\) on all definite vertices fixes
each of the edge data, because it acts trivially on the Hom
space and the two determinant characters are unramified
on these units. Equivariance of the class sum therefore
transfers the scalar translate of the input column to the
same scalar translate of the final coefficient. By the
finite triple equivariance, translating at the final
vertex alone is the action induced by replacing \(x''\)
by \(\jmath_j(\nu_w I)x''\). The displayed level inclusion
makes the finite tensor invariant under this change, since
both edge data descend to the prescribed finite level cosets.
The paired chain hence has the same final value on
\(\psi\) and \(R(\nu_w I)\psi\). Every ramified column is the
difference of two such columns at this \(w\), so the original
stage value vanishes. This contradiction rules out the exceptional
Gram bound \(-1\).

Start with the standard profile at stage zero, where finite
support lies in \(B_0(\mathbb A_{F,f})J_g\). At each nonzero
stage value the defining integral has a nonzero integrand at
some preceding Bessel point, giving both a nonzero stage
value and a nonzero edge constant term. Induction with the
upper inequality therefore bounds every \(e'_i(w)\) from above
by \(j\) on \(Q\cup\bar Q\) and zero elsewhere. Applying the
Gram inequality and its ramified middle correction to the
opposite index gives the corresponding lower bound.
This proves the support condition for all finite places and
outer variables of the cuts.
\footnote{\weightedpin{grh-compact/src/New/Seesaw/UnfoldingStages.lean}
{37-L67}{Even-stage coefficients on the entire family},
\weightedpin{grh-compact/src/New/Seesaw/UnfoldingStages.lean}
{70-L111}{Ramified support and the central scalar},
\weightedpin{grh-compact/src/New/Seesaw/LiftCharacters.lean}
{171-L176}{Unramified rank twists above \(N\)},
and \weightedpin{grh-compact/src/New/Columns/ChainValuation.lean}
{183-L209}{Ramified columns as scalar differences}.}

\teamanchor{proof.weighted-seesaw-cuts}%
\paragraph{Cutting and integrating the outer variables.}
Take \(A_0=J_g\) and \(A_j=K_j^0\) for \(j>0\), and set
\(S_j=B_j(\mathbb A_{F,f})\), \(U_j=S_j\cap A_j\). The first
support calculation gives a compact set \(D_{j,\vartheta}\) whose
left \(S_j\)-saturation covers all the finite support in stage
\(j>0\). Indeed the valuation differences take only finitely
many values. Choose a representative for each occurring profile,
keeping its components at the twist places and their conjugates
and setting the other components to one. The profile is unchanged
by this replacement outside the twist places, by equality with the
standard profile there. Take the union of the sets \(yK_j^0\) for
these representatives. For \(j=0\) use \(D_{0,\vartheta}=J_g\).

For each \(S_j\)-orbit on \(G_j(\mathbb A_{F,f})/J_j\) meeting
\(D_{j,\vartheta}/J_j\), choose one of its points there and form its
\(U_j\)-orbit. Denote the union of these finite orbits by
\(\mathcal F_j\). The image of \(D_{j,\vartheta}\) in the discrete
quotient is finite by compactness. Each \(U_j\)-orbit is also finite.
Writing \(s_t=S_j\cap yJ_jy^{-1}\) for the stabilizer of \(t=yJ_j\),
the number of cosets of \(U_j\) whose translates of \(\mathcal F_j\)
contain \(t\) equals
\begin{equation}
\label{eq:weighted-seesaw-multiplicity}
  m_j(t)=[s_t:s_t\cap U_j].
\end{equation}
In fact, in the orbit of \(t\in\mathcal F_j\), the part of the cut
is \(U_jt\). The elements \(s\) with \(s^{-1}t\in U_jt\) are
precisely \(s_t U_j\). This also shows that the multiplicity is
constant under left \(U_j\)-translation.

The representatives \(y\) of the cut are integral off the twist
primes. At a twist prime over \(\ell\), the component
of \(J_j\) lies in the principal congruence subgroup, a pro-\(\ell\)
group. For any fixed \(y\), a common \(q_\vartheta\)-power of all
elements in \(yJ_jy^{-1}\) belongs to \(K_j^0\). To see this at a
twist prime, take the normal core of the open intersection with the
lattice stabilizer inside the conjugated pro-\(\ell\) subgroup.
Its quotient is a finite \(\ell\)-group. Away from these primes the
conjugation is by an integral matrix, so it preserves the lattice
stabilizer. The same exponent works in the restrictions to \(s_t\)
and to all conjugates there of \(s_t\cap U_j\). Cauchy's theorem
for the finite quotient by the normal core then shows that
\(m_j(t)\mid q_\vartheta^a\) for some \(a\ge 0\). At stage zero
\(s_t\subset U_0\), since \(y\in J_g\) and \(J_0\subset J_g\), so
\(m_0(t)=1\). Also \(\mathcal F_0=J_g/J_0\), since two points of
\(J_g/J_0\) in the same \(S_0\)-orbit are in the same \(U_0\)-orbit.
For the trivial twist all the multiplicities are one,
and every \(\mathcal F_j\) has at most
\begin{equation}
\label{eq:weighted-seesaw-cut-count}
  C_g=g^{f n_m^2},\qquad n_m=\dim_E X_m,
\end{equation}
points, by reduction at the chosen places from the integral level
in rank at most \(n_m\).

Normalize Haar measure on \(S_j\) by \(\operatorname{vol}(U_j)=1\),
and take the complementary real Haar factor to give the fixed
measure on \([B_j]\). The Bessel subgroup here is a unitary group
acting on a unipotent group. The determinant of its action on the
unipotent Lie algebra is an \(F\)-rational character of the unitary
group, hence trivial. The Bessel subgroup and its finite and real
factors are unimodular. For a function \(f_j\) left invariant under
\(B_j(F)\), integrate first over \([B_j]\) and then over
\(B_j(\mathbb A_F)\backslash G_j(\mathbb A_F)\). The latter
quotient has separate real and finite factors. On the finite orbit
\(S_jyJ_j\) its mass is \(\operatorname{vol}(J_j)/\operatorname{vol}_{S_j}(s_t)\).
Indeed the left \(S_j\)-average of the indicator of \(yJ_j\) is
\(\operatorname{vol}_{S_j}(s_t)\) on that orbit. Since the cut has
\([U_j:U_j\cap s_t]\) points in this orbit, the mass is exactly
the sum of \(\operatorname{vol}(J_j)/m_j(t)\) over those points.
It follows that when the inner period at \((\xi_j,t)\) is
\([K_j^0:J_j]a(\xi_j)b(t)\), the outer integral is
\begin{equation}
\label{eq:weighted-seesaw-cut-integral}
  \left(\int_{B_j(F_\infty)\backslash G_j(F_\infty)}
                   a(\xi_j)\,d\xi_j\right)
                 \sum_{t\in\mathcal F_j}\frac{b(t)}{m_j(t)} .
\end{equation}
Here \(F_\infty=F\otimes_{\mathbf Q}\mathbf R\) and the real
measure is independent of the twist. The equality uses
\([K_j^0:J_j]\operatorname{vol}(J_j)=1\). When a finite factor
\(b(t)\) is nonzero, the inner period and absolute integrability
ensure that \(a\) is left invariant and integrable on the real
quotient. Assign the real integration functional the value zero
on functions outside this domain. This convention only affects
terms with zero finite factor and gives a functional defined
using the fixed real measure for all other outer variables.

\paragraph{The real coefficient of the last double lift.}
\teamanchor{proof.weighted-terminal-double-lift}%
Assume the incoming coefficient on \(U(X_{m-2})\) has been obtained.
The two compact integrals for the
\hyperlink{proof.weighted-last-seesaw-arrows}{last arrows} now take
place on \([U(X_{m-2})]\) and \([U(X_{m-1})]\).
The only remaining string vector has norm \(+N\) in the flipped
source. Sending the extra unit vector of \(X_m\) to \(z/\pi\) is a
rational isometry, since \(\pi\bar\pi=N\). The retained group at
the last evaluation is \(U(Y_m)\) embedded by this isometry and the
identity on \(Y_m\). There is no positive-grade constant term or
translation by a string map on these two arrows.

Fix the real outer variables \(\xi\) and hence the real Schwartz
vector of this pair before varying its finite tensor, twist, or
column. First take a real Fock polynomial which is a product of a
polynomial on each edge, and a decomposable finite pair. In
Schr\"odinger coordinates on the orthogonal sum, the Schwartz
function is the product, and summing over rational vectors is the
product of its two rational sums. The first variable is changed
by adjoint and flip. Rational vectors correspond bijectively
under that change. The intertwiner between the two choices of
polarization contributes one scalar fixed by the choices. The
determinant characters agree after reassigning the extra line.
Indeed let \(d,d',d''\) be the three original ranks. The additional
splitting-character exponent at the middle determinant, including
the adjoint flip, is \(\lfloor d/2\rfloor-\lfloor d''/2\rfloor-d\).
The new ranks are \(d-1,d',d''+1\), and \(d''-d=2t\) from
\(d=n_{m-2}+1,\ d''=n_m-1\). The equal parity of \(d,d''\)
leaves this exponent unchanged. The other two splitting exponents
are \(d'-\lfloor d'/2\rfloor\) and \(\lfloor d'/2\rfloor\).
Block embedding preserves both determinants, and all factors of
the two base characters use the same three determinants.
This identifies the action used in the thick pair with the product
of the two normalized edge actions. The two resulting compact
integrals are successive definite lifts. Apply
\hyperlink{proof.weighted-definite-edge-coefficient}{the compact
edge calculation} to their actual kernels. Evaluating the incoming
functional on the first integral gives a functional on the
intermediate type. Applying the second gives one on the last
type, uniform over finite tensors, columns, twists and output
arguments.

Finite sums of such products contain every polynomial in the two
sets of variables. Polynomial vectors, with their Gaussian, are
dense in the real oscillator Schwartz space by the Hermite
expansion. For fixed finite tensor, twist, column and output point,
both compact integrals are continuous in the real Schwartz vector.
Indeed absolute theta summation on a compact set of representatives
is bounded by a fixed Schwartz seminorm. Take the product of these
spaces of values over \emph{all} the tests. The subspace obtained
from linear forms on the last finite-dimensional coefficient
space is finite-dimensional and therefore closed in this
Hausdorff product. Passage to the limit gives one linear form for
the whole family. Thus the preceding compact calculation supplies
the \(l_\xi\) on the last pair even for the real Schwartz vectors
arising in the outer integration. Compact branching then gives the
terminal Petersson pairing below with its stated factor.
\footnote{\weightedpin{grh-compact/src/New/Seesaw/DoubleLiftValue.lean}{22-L54}{The common coefficient of the double lift},
\weightedpin{grh-compact/src/Lib/Weil/DualPair/TwoEdgesOnePair.lean}{90-L95}{the three determinant characters},
\weightedpin{grh-compact/src/Lib/Weil/DualPair/TwoEdgesOnePair.lean}{110-L138}{the theta product and its intertwining},
\weightedpin{grh-compact/src/New/Seesaw/EdgePair.lean}{226-L253}{the two rational isometries}, and
\weightedpin{grh-compact/src/New/Columns/ColumnDefs.lean}{61-L101}{the last block embedding and normalized branching}.}

\paragraph{The terminal pairing and the common scalar.}
Let \(\mathscr L_\vartheta(v)\phi\) denote the definite chain on
the tensor \(v\). Denote by \(R_m\) the block restriction followed
by the fixed branching map of the weighted datum and the factor
\(\prod_{\tau\in\Phi}\tau(\bar\pi)^{-b_\tau}\), where \(b_\tau\)
is the prescribed branch exponent, then converted to an automorphic
form by the inverse of the terminal algebraic realization. To pair
the stage-\(m\) matrix coefficient against a conjugate form, integrate over the
compact real group \(U(Y_m)(F_\infty)\). If \(v_0\) is the fixed
highest vector of the test type, the average of
\(v_0\otimes l_\xi\) under the two coefficient representations
is an intertwiner from the restricted coefficient representation on
\(X_m\) to the one on \(Y_m\). The standard interlacing branching
rule for \(U(n_m)\downarrow U(n_m-1)\) has multiplicity one, so the
intertwiner is a scalar multiple of the fixed branching map.
The Petersson pairing of the corresponding forms gives the same
finite class integral, with a fixed constant by Schur orthogonality.
The inverse of the displayed branching factor and this Petersson
constant depend only on \(D\). Thus for every final test form \(z\)
the terminal period is
\[
   a_m(\xi)\,\langle z,R_m\mathscr L_\vartheta(E_\vartheta(y))\phi\rangle,
\]
for a scalar function \(a_m\) independent of \(z,\phi,\vartheta,y\).
Here and below the Petersson pairing is linear in the second
argument. The actual terminal type, branching map and scaling
factor are those in the weighted column data.
\footnote{See
\weightedpin{grh-compact/src/New/Columns/ColumnDefs.lean}{61-L101}{the branching map and normalized restriction} and
\weightedpin{grh-compact/src/New/Filling/DegreeBounds.lean}{15-L39}{the terminal pairing in the lower normalization}.
For the compact averaging identity, see also
\weightedpin{grh-compact/src/Lib/Weil/CompactLift/BranchingPeriod.lean}{145-L172}{the intertwining mean}.}

Apply \eqref{eq:weighted-seesaw-step} and
\eqref{eq:weighted-seesaw-cut-integral} successively, starting at
\(j=m-1\) and ending at \(j=0\). At each step transpose the last
arrow in the pairing, holding the previous variables fixed. Its
local period has the form in \eqref{eq:weighted-seesaw-cut-integral},
with the real factor given by the preceding real integration and
the finite factor given by the preceding finite sum. The finite
support required for each such integration was uniform in the
previous variables. At the bottom use the integral of the
Poincar\'e series against the transposed test. The result, on
testing against \(z\), is the definite chain on the finite tensor
\begin{equation}
\label{eq:weighted-seesaw-cut-tensor}
  \mathscr Z_\vartheta=\sum_{t_j\in\mathcal F_j\ (0\le j<m)}
                    \left(\prod_{j<m}m_j(t_j)^{-1}\right)
                        E_\vartheta(t_0,\ldots,t_{m-1},1)
\end{equation}
followed by \(R_m\), times one scalar \(a_*\) from the real
integrals. It is independent of \(z,\phi,\vartheta\). All summands
are defined on the cosets. The chosen inclusion data are fixed by
the two principal levels of each edge, and so changing the finite
representatives leaves each translated datum unchanged. Taking all
final tests identifies the forms, with the same conversion to
algebraic forms on both sides.

For level elements \(b_j\) of the definite groups, the inclusions
of the grade-zero cores carry \(b_j\) into \(U_j\) for \(j<m\).
At zero the lower level is \(J_g\), and at the intermediate stages
the inclusions preserve the lattice on the core. For \(b_m=1\),
the paired intertwining property makes simultaneous left
translation of the \(t_j\) by these inclusions act on the kernel
tensor by the joint action on its edges. The cuts and their
multiplicities are invariant under these translations. Therefore
\(\mathscr Z_\vartheta\) has the required joint invariance.
\footnote{The core level inclusions use
\weightedpin{grh-compact/src/New/Seesaw/StandardLatticeLevel.lean}{22-L33}{the lattice of the embedded core}
and \weightedpin{grh-compact/src/New/Seesaw/SeesawTheorem.lean}{38-L47}{the stage-zero level}.}

\teamanchor{proof.weighted-seesaw-finite-datum}%
\paragraph{Admissible finite tensors.}
Use the normalized integral vacuum on each definite edge. For
\(y_j\) integral off \(q_\vartheta\), the calculation below writes
\(E_\vartheta(y)\) as a scalar \(C_*\), depending only on \(D\),
times a combination of tensors \(V_t^{\rm pair}\) of paired
vacuum translates. The tuples \(t\) are admissible for
\(q_\vartheta\) before adding \(N\) to the modulus. The
coefficients become integral at every coefficient place away
from \(q_\vartheta\) after multiplication by \(d_*=C_g^{2m}\),
and at the trivial twist the sum of their absolute values under
each conjugate is at most one.

\paragraph{Finite zero terms on the two graded frames.}
\teamanchor{proof.weighted-paired-vacuum-calculation}%
Use the normalized finite oscillator operations for the given
mixed polarization. The mixed zero-term covariance sends
translations on the isotropic and retained frame to their
retained component. The translate of a finite vacuum under a
unitary move \(g\) is a sum over \(gL/(L\cap gL)\) with the
common multiplier \(\Delta_f(g)/|\Delta_f(g)|^2\). Here
\(|\Delta_f(g)|^2=[L:L\cap g^{-1}L]\). On the split side use
the Schr\"odinger integral-box indicator and the adjoint flip.
\footnote{\weightedpin{grh-compact/src/New/Seesaw/EdgePair.lean}
{39-L41}{Translation covariance of the mixed zero term},
\weightedpin{grh-compact/src/Lib/Weil/Construction/FiniteMultiplier.lean}
{154-L183}{The multiplier and the lattice index}, and
\weightedpin{grh-compact/src/New/Seesaw/AdmissibleEvaluation.lean}
{106-L148}{Split indicator and adjoint flip}.}

Write \(L\) for
the corresponding self-dual lattice, \(\rho(t)\) for its
vacuum-normalized Heisenberg translations, and \({\bf v}\) for the
plain vacuum. We first compute the zero term \(Z\) for an isotropic
frame \(k_0\) with retained isometric embedding \(i_0\).
Suppose \(B=i_0g+k_0 C\) is integral, where \(g\) is unitary on the
retained space. Then, for one scalar \(C_0\) independent of \(t\),
\[
 Z(\rho(t){\bf v})=\zeta_t C_0\,\omega(g)\rho(B^*t){\bf v}_0,
 \qquad \zeta_t=0\ \text{or a root of unity}.
\]
Indeed translation by \(Bv\), for \(v\) integral, fixes the
vacuum. Under \(Z\) the isotropic component disappears and the
retained component translates by \(gv\). The jointly fixed line
for \(gL_0\) is \(\mathbf C\omega(g){\bf v}_0\), by conjugating the
standard vacuum projector. If a translated zero term is nonzero,
its commutators with translations by integral vectors of
\(K=\operatorname{im}k_0\) have trivial character. Hence
\(t\in L+K^\perp\). Locally this is
\((L\cap K)^\perp_{\psi}=L+K^\perp\): split off the primitive
submodule \(L\cap K\) in \(L\) and extend its integral functional
using self-duality. Write \(t-l=i_0w+k_0b\) with \(l\in L\).
Then \(w=gB^*(t-l)\) and \(B^*l\) is integral. Commuting the
translations and removing the integral parameters gives the displayed
formula, since all the finite additive character values have finite
order. This reasoning also applies when \(C_0=0\).
\footnote{\weightedpin{grh-compact/src/New/Seesaw/EdgePair.lean}{39-L75}{Zero-term support and its vacuum line}.}

Here are the matrices that make \(B\) integral in each actual pair.
Put \(\beta=\overline{\pi_G}\), \(C=\) the chosen idele generating
\((\beta)\), \(\delta=C/\bar C\), and \(a=\alpha/\bar C\), for the
Weil idele \(\alpha\). Let \(\epsilon\) be the finite integral unit
of norm \(-1\) in the flip, and put
\[
 \lambda=\bar\epsilon\,\bar\alpha/\alpha,\qquad
 u_j=(a/\bar a)^{\lfloor(m-j+1)/2\rfloor}.
\]
Thus \(\lambda u_j=\bar\epsilon\delta u_{j+2}\) for the even source
indices \(j<m\), and \(u_m=1\). Let \(e,e''\) be the last definite
source and target coordinate in the pair. Put \(D_\delta=1\) for a
middle pair. In the last pair it is diagonal with \(\delta\) at
\(e''\) and one elsewhere. It is one above \(N\).

Let \(p,\bar p\) be the two idempotents of the chosen half of the
places over \(N\) and its conjugate. With row vectors understood as
columns when forming \(\rho_0\rho_0^*\), put
\[
 \rho_0=e+\bar\epsilon\bar p\,e'',\quad
 A=p(\pi-1)+\bar p(\bar\pi^{-1}-1),\quad
 \gamma=1+A\rho_0\rho_0^*,\quad
 g=(D_\delta\gamma)\otimes 1.
\]
Here \(\rho_0^*\rho_0=1\) and \(1+A\) has norm one. For column
\(x=\gamma_\ell\) put
\(c=A\pi^{-1}(\bar p\,1_{\ell=e}+\epsilon p\,1_{\ell=e''})\).
It is supported above \(N\), so \(\delta c=c\).
The entries of \(x\) off \(e,e''\) are integral and
\(\bar\pi x\) is integral. Multiplication on the three idempotent
pieces \(p,\bar p,1-p-\bar p\) gives
\[
 c-\pi^{-1}x_e=-\varpi_0\,1_{\ell=e},\qquad
 \bar\epsilon c-\pi^{-1}x_{e''}=-\varpi_0\,1_{\ell=e''},
 \qquad \varpi_0=p+(1-p)\pi^{-1}\in\widehat{\mathcal O}_E .
\]

In an odd string with middle vector \(z\), the two pertinent
lattice columns are \(e/\bar\pi+u_jz/N\) and \(u_jz\).
Thus a core vector \(x\) with additional component \(u_j H z\)
is integral when \(\bar\pi x_e\), its other core coordinates, and
\(H-x_e/\pi\) are integral. In a middle pair the vectors
\(z,z''\) have norms \(+N,-N\) in the flipped sum. The isotropic
vector \(z+z''\) is paired with \((z-z'')/(2N)\). Append
\(\lambda u_j c(z+z'')\) to the embedded column. Removing the
source flip gives \(H=c\) there, and on the target
\(H=\bar\epsilon c\). Both required differences are integral by
the two congruences. In the last pair the extra vector embeds as
\(z/\pi\). Its coefficient is \(\delta x_{e''}/\pi\), which after
removing the flip and dividing by \(u_j\) is
\(-\epsilon x_{e''}/\pi\). The required difference
\(-\epsilon x_{e''}/\pi-x_e/\pi\) is integral, by adding the first
congruence to \(\epsilon\) times the second. The other target
coordinates lie in its standard lattice. Tensor these columns with
the dual core projection of the odd-index intermediate space. Its
string has even length, so its core is an integral direct summand.
This gives the desired integral \(B\) on the Hom space. For the
specified stacked shift \(t_0=\alpha^{-1}(T^*,T')\) in oscillator
coordinates, \(B^*t_0=0\): both string maps in this column vanish
on the intermediate core. Removing the shift therefore changes at
most the phase and the zero indicator.
\footnote{\weightedpin{grh-compact/src/New/Seesaw/LatticeAndTorsionDefs.lean}{13-L26}{The graded lattice columns},
\weightedpin{grh/src/New/UpperBound/Seesaw/NeighbourLattice.lean}{343-L435}{the two congruences}, and
\weightedpin{grh-compact/src/New/Seesaw/AdmissibleEvaluation.lean}{163-L215}{the tensor with the dual core and the stacked shift}.}

We compute the coefficient and support of the resulting finite
expansion. At \(\bar p\) use the Schr\"odinger coordinate \(v\)
column by column. Put
\[
 n(y)=\epsilon\overline{y/N}\,e-\overline{y/N}\,e'',
\]
where \(y/N\) denotes the selected residue lift supported on \(p\).
The \(e\)-entry of \(\gamma v+n(y)\) on \(\bar p\) is
\(\bar\pi^{-1}v_e+\epsilon\overline{y/N}\).
For integral \(v_e\) exactly one residue class \(y\) in each
column makes this entry integral. Existence of such a residue
forces \(v_e\) integral, on multiplying by \(\bar\pi\).
Every other change in this column is \(\rho_{0,i}\) times the
change in the \(e\)-entry. Thus \(\gamma v+n(y)\) belongs to the
integral box for exactly one \(y\) when \(v\) belongs to the
integral box, and for no \(y\) otherwise. This proves the finite
indicator identity. All translations \(n(y)\) are supported on
one split isotropic side, on which the semicharacter and all
mutual commutators are one. The first component after inverse
adjoint and flip is \(-y/N\), by
\(\epsilon\bar\epsilon=-1\); the second is \(-\bar y/N\).
The parameters already present are integral over \(N\), since
\(\beta\) is a unit there and the twist modulus \(q\) is coprime
to \(N\). Remove these integral components on the original vacuum
before applying the box identity. The remaining parameters
commute with every \(n(y)\). Consequently the residue sum has
equal coefficients in each pair; its overall scalar is independent
of the residue.
\footnote{\weightedpin{grh-compact/src/New/Seesaw/AdmissibleEvaluation.lean}{106-L148}{The integral box in the selected polarization},
\weightedpin{grh/src/New/UpperBound/Seesaw/NeighbourVacuum.lean}{118-L200}{commutation and the primes over \(N\)}, and
\weightedpin{grh/src/New/UpperBound/Seesaw/NeighbourVacuum.lean}{208-L271}{the adjoint flip of each translate}.}

Away from \(N\) the column correction \(c\) is zero and
\(\gamma=1\). Give the source block of a stacked translation
denominator \(\beta\) and the target block \(\bar\beta\):
multiplication by that denominator is required to give an
integral parameter away from \(q\). On the retained components
the integral adjoint \(B^*\) preserves this condition with the
new denominators \(d_b(D_\delta)_{bb}\), where \(d_b\) is \(\beta\)
or \(\bar\beta\) according to that block. Indeed every nonzero
block of the embedding retains its side and its valuation except
for the terminal extra coordinate. That coordinate comes from
the source and has
\(v(\bar\beta\delta)=v(\beta)\) at every finite place.
More explicitly, a nonzero entry of \(B^*\) at a place \(w\)
uses a conjugated embedding entry at \(\bar w\). The respective
denominators at \(w\) have equal valuations as just computed,
so integrality follows termwise. Above \(N\) all denominators
are units. Applying \(D_\delta\) itself to the parameter then
gives \(d_b\) as the denominator. Under the inverse adjoint and
flip the source parameter is conjugated and multiplied by an
integral unit, so both ordinary edge parameters have denominator
\(\bar\beta=\pi_G\).

Write \(T(\xi)=\psi_0(c_0\xi^*\xi/2)\rho_{\rm Weyl}(\xi)\) for the
normalized translations used above. Here \(v_0\) is the vacuum on
\(L_0\) and the underlying operators satisfy
\(\rho_{\rm Weyl}(\xi)\rho_{\rm Weyl}(\eta)
=\psi_0(\eta^*\xi/2)\rho_{\rm Weyl}(\xi+\eta)\). The group-law
discrepancy is \(\psi_0(\overline{c_0}\eta^*\xi)\), using
\(c_0+\overline{c_0}=1\) and \(\psi_0(z^*)=\psi_0(-z)\). It is one on
\(L_0\), since \(c_0\) is integral and \(\psi_0\) is trivial on the
integral finite adeles. The unitary change \(D_\delta\) preserves
\(\xi^*\xi\) and \(\eta^*\xi\) exactly, including at the dyadic places.
Thus the same normalized translations are a group action on
\(D_\delta L_0\), conjugate to the action on \(L_0\). For
\(H=L_0\cap D_\delta L_0\), every \(T(\xi)\) with \(\xi\in H\) fixes
\(v_0\). The probability-Haar average on \(D_\delta L_0\) applied to
\(v_0\) is therefore the finite average over \(D_\delta L_0/H\).
Its image lies on the moved vacuum line. The Weyl commutator and
self-duality of \(L_0\) give zero vacuum pairing for
\(\xi\notin L_0\), so the identity coset gives a nonzero coefficient
\(\|v_0\|^2/h\), where \(h=[D_\delta L_0:H]\). Absorb the fixed
scalar multiplying the average into \(C_0\). The coefficients are
\(1/h\) up to roots of unity.
\footnote{\weightedpin{grh-compact/src/Lib/Weil/Construction/FiniteHeisenberg.lean}
{41-L70}{Self-duality and the finite phase} and
\weightedpin{grh-compact/src/Lib/Weil/Construction/Projections.lean}
{210-L220}{The normalized average on the vacuum}.}
The quotient maps injectively into \(\bar\beta^{-1}L_0/L_0\).
Each diagonal factor satisfies
\(\bar\beta(D_\delta)_{bb}\in\widehat{\mathcal O}_E\) and also
\(d_b(D_\delta)_{bb}\in\widehat{\mathcal O}_E\). This shows that
\[
 h\mid |\mathcal O_E/(\beta)|^{(n_j+n_{j+2})n_{j+1}}
\]
and that the averaged translations preserve the required
denominators. The sum of the conjugated absolute coefficients of
the average is at most one, including when the zero term vanishes.

Finally the finite inclusion datum on each lower edge is the
characteristic function of its specified residue matrix. Finite
Fourier inversion on its matrix module expresses it by the
vacuum translations with coefficients roots of unity divided by
the cardinality of the module. Its local pieces away from \(q\)
are the ones over the chosen divisor of \((\beta)\), with one
matrix entry for each source and target index. Thus the part of
the denominator prime to \(q\) divides
\(|\mathcal O_E/(\beta)|^{\dim Y_j\dim Y_{j+1}}\).
The conjugated absolute coefficient sum is one. The moving
level element is integral away from \(q\). Its moved-vacuum
formula has a finite-index denominator supported over \(q\).
Its Weil multiplier is \(\chi_{\rm fin}(\det b)\) for the
saturating lattice \(bL=L\cap g^{-1}L\).
At every \(\ell\nmid q\), \(\det b\) and the determinant for the
additional splitting characters are units at every place over
\(\ell\). The algebraic character values are then \(\ell\)-adic
units under every coefficient embedding: pass to a power with
principal ideal, write its generator times a finite integral
unit, and evaluate its infinity type on the generator, all of
whose conjugates are \(\ell\)-units. The finite norm factor uses
only residue cardinalities over \(q\). At the trivial twist the
moving elements are integral everywhere. Their vacuum multipliers
are one and the extra smooth character on the compact determinant
group has finite image.

On each pair the two lower matrix exponents are
\((\dim_E Y_j+\dim_E Y_{j+2})\dim_E Y_{j+1}\) in sum. Each
\(\dim_E Y_k\) and each \(n_k\), for \(0\le k\le m\), is at most
\(n_m\). Adding the retained exponent \((n_j+n_{j+2})n_{j+1}\)
therefore gives at most \(4n_m^2\). Thus the combined denominator
divides \(C_g^4\), where
\(C_g=|\mathcal O_E/(\beta)|^{n_m^2}=g^{f n_m^2}\).
Tensoring \(m/2\) pairs gives \(C_g^{2m}\), uniformly in the
twist and finite cut points, and conjugated coefficient sum at
most one at the trivial twist. Replacing the plain vacua by the
chosen nonzero integral vacua divides the common scalar by their
fixed product of normalization factors. All the other scalars
absorbed above are attached to the fixed pair, its frame or its
polarization. This gives the single \(C_*\) needed before the
outer cut sum.
\footnote{\weightedpin{grh-compact/src/New/Seesaw/AdmissibleEvaluation.lean}{136-L162}{The shifted weights and finite index},
\weightedpin{grh-compact/src/New/Seesaw/IntegralAway.lean}{21-L36}{the determinant factors at the coefficient places},
\weightedpin{grh-compact/src/New/Seesaw/GoodDatumExists.lean}{23-L81}{the actual two matrices and the exponent \(2m\)}, and
\weightedpin{grh/src/New/UpperBound/Seesaw/GoodDatum.lean}{542-L559}{conversion to the normalized integral vacua}.}

By \eqref{eq:weighted-seesaw-multiplicity}, the cut coefficients
are integral away from \(q_\vartheta\). At the trivial twist the
cut has total mass at most \(C_g^m\) by
\eqref{eq:weighted-seesaw-cut-count}. Multiplying the finite
expansions by the cut coefficients and summing consequently gives
\(\mathscr Z_\vartheta=C_*\sum_t r_t V_t^{\rm pair}\), with
\(d_*r_t\) integral away from \(q_\vartheta\) and
\(\sum_t|\sigma(r_t)|\le C_g^m\) at the trivial twist. Choose
\(c_*=d_* /(a_*C_*)\) if \(a_*C_*\ne0\). Then the tensor for
\(\iota_\vartheta\mathcal P(c_*\phi)\) has coefficients
\(d_*r_t\) in these paired vacua. If \(a_*C_*=0\), the unfolded
pairing vanishes for every column and twist, so one can take
\(c_*=1\) and the zero finite datum.

In each pair the operator used in \(V_t^{\rm pair}\) sums over
residue vectors at the chosen places over \(N\). It adds
\(-y/N\) on the last column of the first edge and
\(-\bar y/N\) on the last row of the second. Composition of
Heisenberg translations contributes a root of unity and adds the
translation parameters. Thus expanding the pairs gives a linear
combination of ordinary edge vacuum tensors whose translations
are admissible for \(Nq_\vartheta\), still with coefficients
integral away from \(q_\vartheta\). Each pair has at most
\(N^{f n_m}\) summands. After collecting identical tuples the
conjugated coefficient sum is at most
\[
   C_g^{3m} N^{f n_m m/2}\le e^{x^{60}} .
\]
For the inequality, \(f=O(x^{11}\log x)\), \(n_m=O(x^{11/2})\),
\(m=O(\sqrt{x})\), \(\log g=O(\log b)\), and
\(\log N=O(m\log b)\) by the level window. The expansion is an
identity of the joint tensor, which retains its invariance.

Finally let \(P_N\) be the projector onto the vacuum at all the
places above \(N\), for each edge. The original paired translations
have integral starting parameters there. On a translate of a
vacuum, the projector multiplies the translate by the indicator
of integrality of its parameter at those places. This follows by
averaging the Heisenberg commutator character over the standard
self-dual lattice. In a translated pair only the zero residue
vector can survive both projections, since \(N\) is split and the
chosen places contain one from each conjugate pair. Thus projecting
each edge of a paired tensor and applying the same residue sum to
each pair returns that tensor. Apply this to the sum and use the
linearity in the joint tensor of the definite chain. This gives
exactly the required paired-chain expression at \(N\) for the
ordinary coefficient function \(\Psi_\vartheta\) obtained above.
This proves \eqref{eq:weighted-seesaw} from the fixed-data inputs.
\footnote{The operators on the two edges are
\weightedpin{grh-compact/src/New/Seesaw/NamesAtSplitLevel.lean}{342-L352}{the residue sum and vacuum projection}.
Their defining application is
\weightedpin{grh-compact/src/New/Seesaw/GoodDatumExists.lean}{84-L116}{projection of a paired translate and expansion},
with the stated goodness conditions in
\weightedpin{grh-compact/src/New/Columns/ChainValuation.lean}{260-L279}{the definite-chain datum}.}

\teamanchor{proof.weighted-projection-transfer}%
\paragraph{Projection at the chosen types.}
For \(s_0\in\mathcal G\) put
\[
   A_{s_0}\phi=|\mathcal G|\,\rho(e_0)\rho(s_0)\mathcal P(c_*\phi)
       \qquad(\phi\in\mathcal C_{\rm ram}).
\]
The following transfer applies to any principal-level form \(u\)
on \(Y_0\), including the auxiliary rows in the height comparison.
Choose a lift of \(s\in\mathcal G\) to integral matrices in the
split charts. At every stage use the block diagonal lift with that
same first block, the identity on later blocks and the identity at
the twist places. The inclusion matrix \(I_j:Y_j\hookrightarrow
Y_{j+1}\) in the finite inclusion datum is fixed by
\(M\mapsto \widetilde s_{j+1}^{-1}M\widetilde s_j\). On the split
lattice model this is the change of arguments in the finite Weil
action. The modulus factors have absolute determinant one. The
tame scalar at a selected place is the ratio of the two values of
the edge residue character on the determinants. Both determinants
are those of the same first block, so the ratio is one for
\(\mu_{j+1}^0\vartheta^{j+1}\). The factors at the other places are
unchanged. Thus the pair of lifts fixes the lower edge datum.
Changing the source variable in the theta integral commutes with
right translation by these simultaneous lifts. This proves
\(\iota_\vartheta(\rho(s)u)(h)=\iota_\vartheta(u)(h\widetilde s_m)\).
The determinant ratios are one at every edge. The conversion to
algebraic forms at the last stage
commutes with finite right translation.
\footnote{The finite model for the present inclusion vector and the
determinant ratio are recorded in
\weightedpin{grh-compact/src/New/Seesaw/GradingsOfY.lean}{34-L63}{the split inclusion datum}.
The target actions and lifts of the first block are those of
\weightedpin{grh-compact/src/New/Seesaw/GradingsOfY.lean}{135-L163}{the terminal block level}.}

For each approved irreducible character \(\chi\) of \(\mathcal G\),
Schur orthogonality gives its central projector. Consequently,
writing \(a_t=|\mathcal G|e_0(t)\),
\[
       a_t=\sum_{\chi\ {\rm approved}}\chi(1)\chi(t^{-1}) .
\]
Each character value is a sum of roots of unity, hence \(a_t\)
is an algebraic integer. For any field automorphism \(\sigma\) of
\(\mathbf C\), applying \(\sigma\) to matrices of the irreducible
representations permutes the irreducible characters and preserves
their degrees. Therefore
\begin{equation}
\label{eq:weighted-projection-coefficients}
 \sum_t|\sigma(a_t)|^2
    =|\mathcal G|\sum_{\chi\ {\rm approved}}\chi(1)^2
    \le |\mathcal G|^2,
 \qquad \sum_t|\sigma(a_t)|\le|\mathcal G|^{3/2} .
\end{equation}
The last bound is Cauchy--Schwarz. The convolution and the covariance
prove the evaluated identity used in Theorem~A
\begin{equation}
\label{eq:weighted-seesaw-evaluated}
 (\Theta_\vartheta A_{s_0}\phi)(h)
     =\sum_{t\in\mathcal G} a_t\,
      B^\mu_{\vartheta,\Psi_\vartheta}(\phi)(h\widetilde{(ts_0)}_m).
\end{equation}
It has the same scalar \(c_*\) before projection for every translate
and twist. More generally, for any principal-level \(u\),
\[
 \iota_\vartheta(|\mathcal G|\rho(e_0)\rho(s_0)u)(h)
    =\sum_{t\in\mathcal G} a_t\,\iota_\vartheta(u)(h\widetilde{(ts_0)}_m).
\]

At any coefficient place the Weyl unit ball at \(h\) is a module
over its valuation ring. The gauge uses the action of \(h^{-1}\)
on the coefficient vector at the induced local places. Each
\(\widetilde{(ts_0)}_m\) and its inverse are integral at every
finite place. The corresponding Weyl operators preserve the
integral Weyl lattice, so the gauge of a given vector is the same
at \(h\) and at \(h\widetilde{(ts_0)}_m\). If all values of a scalar
multiple of the unprojected lift have gauge at most one at their
arguments, their translates and the combination with coefficients
\(a_t\) have gauge at most one at \(h\). This works with the same
scalar and coefficient place. At the trivial
twist, if all the conjugated unprojected values have squared Weyl
norm at most \(e^B\), the triangle inequality in the Weyl norm
and \eqref{eq:weighted-projection-coefficients} give an upper bound
\(|\mathcal G|^3 e^B\) for every squared projected value. In the
class sum the weight of any class is the reciprocal of its finite
stabilizer order. With \(H_m\) the number of classes of the final
trivial-twist level, the logarithmic loss is at most
\begin{equation}
\label{eq:weighted-projection-norm-loss}
                  3\log|\mathcal G|+\log H_m\le x^{99}.
\end{equation}

The \hyperlink{proof.weighted-short-class-sum}{definite class count}
below gives \(\log H_m\le x^{36}\). Since \(\log|\mathcal G|\le fr^2\log g\),
this proves the last inequality at all sufficiently large \(x\).

\subsection*{Short forms (Theorem A)}

\teamanchor{proof.weighted-short}%
Throughout this step, \(D\), the Poincar\'e maps \(A_s\) with their common
seesaw scalar \(c_*\), and the level prime \(N\) are fixed. Put
\(n_j=\kappa+j(t+1)\), the rank of \(X_j\), and \(C=\dim_{\mathbf C}\mathcal C\).
We use \(\|v\|_{\mathrm W}\) for the positive Weyl norm of a value and
\(\|\Theta_\vartheta y\|_\sigma\) for the class-sum norm of a lifted row
after applying \(\sigma\) to its coefficients. At a coefficient embedding
\(j:\mathbf C\hookrightarrow\mathbf C_p\) we write \(\mathcal O_j\) for the
valuation ring pulled back to \(\mathbf C\) and \(v_j(p)=1\).
Integrality of a form at \(j\) means that its Weyl gauge is at most one
at every adelic argument, with the lattice transported at that argument.

\paragraph{A spanning family of short integral columns.}
\teamanchor{proof.weighted-short-spanning}%
At stage zero there are \(a\) columns of height one at each
\(\tau\in\mathcal T\) in the weight diagram. All other columns of the
diagram have height at least \(\kappa\). Thus the coefficient representation
on \(X_0\) has one symmetric-power factor of degree \(a\) for each top
and a trivial factor at every other place, with determinant exponent zero.
In the degree-\(a\) factor write \(v(z)=(z^\alpha)_{|\alpha|=a}\) for the
Weyl evaluation vector. The positive Weyl norm is invariant under the
unitary group and gives \(v(e_1)\) squared norm one. Unitary transitivity
on the unit sphere and homogeneity give
\[
   \|v(z)\|_{\mathrm W}^2=\Big(\sum_i |z_i|^2\Big)^a .
\]
It follows by comparing coefficients in \(z,\bar z\) that the Gram
matrix in these coordinates has diagonal entries \(a!/\alpha!\).
In products of factors the squared norms multiply.

The definite class lattices \(L_g=E^\kappa\cap g\widehat{\mathcal
O}_E^\kappa\) are self-dual for the standard Hermitian form.
The \hyperlink{proof.weighted-free-level}{free-level calculation}
applies to the chosen principal level within each definite stage.
Evaluation at the class representatives therefore identifies the
columns with a direct sum of copies of the coefficient representation.

Write \(Q=dN_1\). Then \(Q\) belongs to the different of \(E\).
Indeed \(E\subset\mathbf Q(\zeta_Q)\), and
\(\Phi_Q'(\zeta_Q)\) divides \(Q\zeta_Q^{-1}\), by differentiating
the factorization of \(Z^Q-1\). Since \(\mathbf Z[\zeta_Q]\) is the
ring of integers, \(Q\) kills its codifferent modulo that ring.
The trace from \(\mathbf Q(\zeta_Q)\) to \(E\) maps integers to integers,
so the codifferent of \(E\) is contained in the cyclotomic codifferent.
Intersecting with \(E\) proves the assertion.

Endow the real embedding of \(L_g\) into \(\mathbf C^{f\kappa}\) with
the Euclidean norm, using one embedding from each conjugate pair.
Its Euclidean dual is \(2\mathfrak D_E^{-1}L_g\) in the same coordinates
because the Euclidean pairing is half the Hermitian trace pairing.
Every nonzero \(v\in L_g\) has \(\|v\|\ge1\). Indeed its squared length
is the sum of the conjugates of the nonzero totally positive integer
given by its Hermitian square. Thus every nonzero vector of the
Euclidean dual has length at least \(2/Q\).

Here is the weak transference estimate needed for this lattice.
For a full Euclidean lattice \(L\) of rank \(D_0\), Minkowski's first
theorem on \(L\) and \(L^*\), and the cube in the unit ball, give
\(\lambda_1(L)\lambda_1(L^*)\le D_0\).
If \(\lambda_1(L^*)\ge\rho>0\), project along a primitive shortest
vector \(b\in L\). The projected lattice has dual \(L^*\cap b^\perp\)
and minimum in its dual at least \(\rho\).
Apply induction in dimension to the projected distance, then lift
the approximate vector and round its remaining coordinate modulo \(b\).
The total distance is at most
\((D_0-1)^2/\rho+D_0/(2\rho)\le D_0^2/\rho\), with the same argument
in dimension one starting from the zero-dimensional lattice.
For \(L_g\) the resulting radius is \(R_{\rm tr}=(2f\kappa)^2Q/2\).
Choose successive \(E\)-independent elements of \(L_g\) by taking a
target vector of Euclidean length \(2R_{\rm tr}\) orthogonal to the
archimedean image of the \(E\)-span already chosen and approximating
it within \(R_{\rm tr}\). Each new length is at most \(3R_{\rm tr}\)
and the new vector is outside that span. The resulting frame
\(Y_g\in\mathrm{GL}_\kappa(E)\) satisfies
\[
     g^{-1}Y_g\in M_\kappa(\widehat{\mathcal O}_E),\qquad
     |\tau((Y_g)_{hi})|\le P_{\rm frame}=6Q(f\kappa)^2
\]
at all complex embeddings. The same proof works for \(n_j\) in place
of \(\kappa\).

\teamanchor{proof.weighted-classwise-evaluation}%
Use a standard positive definite group of rank \(q\) with self-dual
class lattices. At a free class represented by \(h\), the level
preserves the transported integral gauge. Let \(v_\tau(X)\) be the
vector of Weyl coordinate polynomials of the diagram at \(\tau\),
with the determinant character treated separately.
Choose \(Y_h\) by the same frame construction in rank \(q\).
For a nonzero determinant exponent \(a_\tau\), choose also the frame
\(Y_{h^{a_\tau}}\) for \(h^{a_\tau}\). Use the class values
\[
 \bigotimes_\tau v_\tau\bigl(\Phi_\tau(\overline{Y_h})X_\tau\bigr)
                 \prod_{\tau:a_\tau\ne0}
                     \Phi_\tau(\overline{\det Y_{h^{a_\tau}}}).
\]
Here the entries of \(X_\tau\) run over \(D+1\) distinct roots of unity,
where \(D\) is any bound for the number of boxes of the diagram
at each \(\tau\). For the symmetric-power evaluation one can take
\(X_\tau\) to be a single column. Extend each value by automorphy on
its class and put zero elsewhere. Freeness makes this extension
well-defined.

In the transported gauge at \(h\), the conjugate on the frame matches
the coefficient reading \(\Phi_\tau\circ\operatorname{conj}\).
The polynomial arguments use \(h^{-1}Y_h\), and the extra factors
use \(\det(h^{-a_\tau}Y_{h^{a_\tau}})\) at their respective embeddings.
Both are integral in these frames. The class transport preserves
integrality. The values at \(h\) span the coefficient space since the
Weyl coordinate polynomials are linearly independent, substitution
by \(Y_h\) is invertible, and evaluation on the roots-of-unity grid
is injective on polynomials of degree at most \(D\) in each entry.
Indeed univariate interpolation in each variable proves this last
assertion in succession. The determinant factors are nonzero scalars.
Conjugate norm bounds on the representatives hold on their whole
classes since rational unitary transport remains unitary at every
complex embedding.

For the stage-zero space take the single-column evaluation at each
top, of degree \(a\) and determinant exponent zero. The symmetric-power
norm formula above bounds the squared norms uniformly by
\((\kappa^3P_{\rm frame}^2)^{aT}\). Apply the ramifying product
\(\mathcal R\) to this spanning family. Each \(1-z_v\) preserves integrality and multiplies a
squared uniform norm bound by at most \(4\), by the triangle inequality
and compactness of the finite translation in the adelic lattice.
The number of factors is \(f\), as in the
\hyperlink{proof.weighted-central-orbits}{central-orbit calculation}.
Consequently one set \(S\) spans \(\mathcal C_{\rm ram}\), is integral
at every coefficient embedding, and obeys
\begin{equation}
\label{eq:weighted-A-short-column}
    \|\sigma\phi(g)\|_{\mathrm W}^2\le
          4^f(\kappa^3P_{\rm frame}^2)^{aT}\le e^{\mathcal B}
          \quad(\phi\in S,\ \sigma\in\operatorname{Aut}(\mathbf C)).
\end{equation}
For the last inequality use \(T\le f\), \(\log P_{\rm frame}=O(\log x)\) and
\(\log x=o((\log b)^3)\).

\paragraph{Class sums, conjugate bounds and the signed height.}
\teamanchor{proof.weighted-short-class-sum}%
Use the finite datum \(\Psi_\vartheta\) of the fixed seesaw.
For the chain integral, the quotient Haar measure at each stage
assigns one to the class of a free representative when the
level has volume one. The combined integrand over the \(m\) stages,
after summing over the support of \(\Psi_\vartheta\), is invariant
under independent right translations by those levels. This follows
by moving the input and output translations through each kernel;
the intermediate actions cancel and joint invariance fixes the sum.
Thus the chain is the sum over all tuples of classes of these
combined values. In particular no denominators arise from passing
from a joint integral to the joint class sum. If each individual
term propagates an upper bound for the pointwise Weyl gauge, the
ultrametric inequality gives the bound for the whole sum whenever
the coefficients of \(\Psi_\vartheta\) are integral. For complex
norms the triangle inequality instead costs the number of class
tuples and the sum of the absolute values of the coefficients.
The same reasoning applies after changes to the kernel data that
commute with the indicated source and target levels.

\teamanchor{proof.weighted-definite-class-count}%
For these standard positive definite groups with self-dual class
lattices, let \(q\le 2n_m\) and suppose the level
contains the principal congruence subgroup modulo the rational integer
\(M\). The short integral frame in rank \(q\) gives a Gram matrix
whose entries are algebraic integers of modulus at most
\(q(6Q(fq)^2)^2\) at all embeddings. Nonzero integral differences
have maximum modulus at least one by their rational norm. Packing
in the archimedean box therefore gives \(\exp(O_d(fq^2\log x))\)
possible Gram matrices. If \(L'\) is spanned by the frame over
\(\mathcal O_E\), then \(L'\subset L_h\subset (L')^\vee\).
The index \([(L')^\vee:L']\) is the absolute norm of the Gram
determinant and its logarithm is \(O_d(fq\log x)\) by Hadamard.
A finite abelian group of order \(J_0\) has at most
\(J_0^{1+\log_2 J_0}\) subgroups, by choosing a sequence of generators
each increasing the generated subgroup's order by a factor at least
two. Frames with the same Gram matrix are identified by a rational
unitary map, so these subgroups bound the possible class lattices for
each Gram matrix. The classes at principal level modulo \(M\)
surject onto the classes at the given level. Entrywise reduction
modulo \(M\) bounds the principal level index by \(M^{2fq^2}\).
Thus the logarithmic class count is
\[
 O_d\bigl(fq^2\log x+(fq\log x)^2+fq^2\log M\bigr).
\]
For the definite stages of Theorem~A take \(M=gN\), including the
final class-sum level on \(Y_m\) at the trivial twist. The bounds
\(f=O(x^{11}\log x)\), \(q\le n_m=O(x^{11/2})\) and
\(\log(gN)=O(m\log b)\) give a logarithmic count at most \(x^{36}\)
for large \(x\).

\paragraph{A Gaussian estimate for factorial jet pairings.}
\teamanchor{proof.shared-gaussian-jet-pairing}%
Let \(Z_\rho\) be complex rectangular matrix variables, grouped into
blocks \(Z_b\) of \(D_b\) complex entries. Write \(s_{\rho,i}\) for
the singular values of \(Z_\rho\), in decreasing order. Let \(P\) be
a polynomial of degree \(d_b\), homogeneous in each block, and
suppose
\[
 |P(Z)|^2\le A\prod_{\rho,i}s_{\rho,i}^{2\lambda_{\rho,i}},
 \qquad d_b=\sum_{\rho\in b,i}\lambda_{\rho,i},
 \qquad \lambda_{\rho,i}\ge0 .
\]
Let \(H\) be entire, with \(|H(Z)|\le M e^{\|Z\|^2/2}\), where the
norm is the Hilbert--Schmidt norm on all entries. With
\(P=\sum p_I Z^I\), \(H=\sum h_I Z^I\), and
\(\phi(u)=u\log u-u\) for \(u>0\), \(\phi(0)=0\), one has
\begin{equation}
\label{eq:shared-gaussian-jet-pairing}
 \left|\sum_I p_I I!h_I\right|^2
 \le A M^2 e^{\sum_{\rho,i}\phi(\lambda_{\rho,i})}
        \prod_{b:\,d_b>0} e(2d_b+2)^{2D_b}.
\end{equation}
Conjugating the coefficients of \(P\) preserves its bound, because
\(Z\) and \(\bar Z\) have the same singular values. Thus the same
estimate holds with \(\overline{p_I}\) in the pairing.

Evaluate \(H\) at zero in each block of degree zero. In every
other block use \(\epsilon_b=(2d_b+1)^{-1}\) and the Gaussian
measure \(\pi^{-D_b}e^{-(1+\epsilon_b)\|Z_b\|^2}\,dZ_b\).
The squared norm of a monomial of degree \(d_b\) in this measure
is its factorial times \((1+\epsilon_b)^{-D_b-d_b}\).
Monomial orthogonality, Cauchy--Schwarz and the bound for \(H\)
give \(M^2\prod_b(1+\epsilon_b)^{D_b+d_b}\epsilon_b^{-D_b}\)
times the squared norm of \(P\) in the Gaussian measure of
parameter one. In its integral, the identity
\(\|Z_b\|^2=\sum_{\rho\in b,i}s_{\rho,i}^2\) leaves the factor
\(e^{-\epsilon_b\|Z_b\|^2}\) after bounding
\(r^\lambda e^{-(1-\epsilon_b)r}\) by
\((1-\epsilon_b)^{-\lambda}e^{\phi(\lambda)}\).
The residual integral gives \(\epsilon_b^{-D_b}\).
Since \((1+\epsilon_b)/\epsilon_b^2\le(2d_b+2)^2\) and
\(((1+\epsilon_b)/(1-\epsilon_b))^{d_b}\le e\), this proves
\eqref{eq:shared-gaussian-jet-pairing}.

For a rectangular Schur substitution in unitary norms with
normalized highest vectors, the scalar matrix coefficient between
unit vectors has the bound in this estimate with \(A=1\) and
\(\lambda_{\rho,i}\) the row lengths. In fact in the exterior-power
column model, diagonal substitution contributes at most the product
of the first \(h\) singular values in each column of height \(h\).
A coordinate inclusion is isometric since it is unitary-equivariant
and preserves the normalized highest vector. Singular value
decomposition reduces a rectangular matrix to this diagonal and
inclusion. Scaling a block variable by a scalar \(c_\rho\) contributes
\(|c_\rho|^{2\sum_i\lambda_{\rho,i}}\) to \(A\).

\paragraph{The archimedean estimate for a definite edge.}
\teamanchor{proof.weighted-cm-jet-length}%
Put \(r_0=n_j n_{j+1}\), \(D_j=fr_0\), and
\(u=\log(g/\Omega)\). For one moved, normalized vacuum,
\hyperlink{proof.weighted-cm-jet-transport}{Taylor transport}
at the actual coefficient scale
\(\overline{\pi_G}\mathrm{per}\) reduces the pairing with a
Schur matrix coefficient to the Taylor pairing of an entire
function \(H(Z)\) with the conjugate polynomial, with the places
permuted. In the polynomial the argument is scaled by
\(c_\sigma'=\sigma(\overline{\pi_G})t_\sigma\), with
\(|c_\sigma'|^2=g/\Omega\).
The function \(H\) has the majorant \(M_j e^{\|Z\|^2/2}\),
\(M_j=\exp(C_d(\nu r_0+1)^2)\), by
\eqref{eq:weighted-cm-envelope}.

Choose unit input and dual test vectors in the conjugate Weyl
spaces, and use their coefficientwise preimages, with the place
permutation, for the scalar test on the original jet. The Hermitian
dual is taken in the actual conjugate metric. The test polynomial
has Schur row lengths \(\lambda_{j,\tau,i}\) and includes the
scale \(c_\sigma'\) above at each place.
In \eqref{eq:shared-gaussian-jet-pairing} use one block of \(D_j\)
entries and degree \(d_j=\sum_{\tau,i}\lambda_{j,\tau,i}\).
The squared scale factor is \(A=e^{u d_j}\). Thus for \(d_j>0\)
the squared operator norm is at most
\[
 M_j^2(2d_j+2)^{2D_j}
    \exp\left(1+\sum_{\tau,i}
        [\phi(\lambda_{j,\tau,i})+u\lambda_{j,\tau,i}]\right).
\]
For \(d_j=0\) use \(H(0)\). Since
\(\phi(\lambda)\le\log(\lambda!)\), the principal exponent
is bounded by \(\sum_{\tau,i}S(\lambda_{j,\tau,i})\).
The sum over stages of the additional logarithmic costs
\(O_d((N_1 n_jn_{j+1}+1)^2+D_j\log(2d_j+2)+1)\) is
\(O(x^{70})\). Indeed \(N_1,f=O(x^{11}\log x)\),
\(n_m=O(x^{11/2})\), \(m=O(x^{1/2})\), and the stated row
length bound gives \(d_j\le 2afn_j\).
The class-sum and finite coefficient factors are charged
separately in the short-lift argument.

The weight at height \(\kappa\) is
\(w_1(\tau)-\lfloor(\kappa+t+1)/2\rfloor\).
The other heights \(h> \kappa\), separated by \(t+1\), have
multiplicity at most \(2r^8\). There are at most \(m\) of them
in each diagram. It follows that the row lengths are at most
\[
  a\,1_{\{\tau\in\mathcal T,\ i=0\}}
       +w_0(\tau)1_{\{i<\kappa\}}+O(mr^8).
\]
At a top \(w_0(\tau)=a/x_1+O(N_1+r^8)\).
Since \(g/\Omega>1\) eventually, the increment
\(S(v+1)-S(v)=\log(v+1)+\log(g/\Omega)\) is positive
and \(O(\log b)\) for \(v\le2a\).
All the displayed lengths are at most \(2a\). Consequently
\[
 \sum_{j,\tau,i} S(\lambda_{j,\tau,i})
   \le m\bigl(T S(a)+\kappa\sum_\tau S(w_0(\tau))\bigr)
       +O\bigl(mfa\log b/x_1+x^{70}\bigr).
\]
The correction in the first row of a top costs \(O((a/x_1+mr^8)\log b)\),
and all remaining corrections are at most \(O(mr^8\log b)\) per row.
The same error bound dominates the sum of the remaining single-step
Gaussian terms, the logarithmic class counts, and the squared-norm
loss \(2x^{60}\) from \(\Psi_1\).

The \hyperlink{proof.weighted-terminal-branching}{terminal branching
calculation} gives the squared logarithmic norm cost
\[
    -Ta\log N+2fb\log N+O(x^{70}).
\]
Its positive terms and the class-sum and jet errors above are
\(o(\mathcal B)\), using \(\log N=O(m\log b)\), \(a\ge xb\) and
\(x_1\gg x^A\) for every fixed \(A\).

Let \(\gamma_p\) be as in the prime-power decomposition of the
height. More explicitly put
\begin{align*}
   \operatorname{Ask}_p(v,\tau)
   &=\frac{1-\chi_d(p)}{2}
      \sum_{\substack{e\ge1\\e\text{ odd or }p=d}}
          (v/p^e-w([-p^e]\tau))_+,\\
   \gamma_p
   &= m\Bigl(\sum_{\tau\in\mathcal T}\operatorname{Ask}_p(a,\tau)
       +\kappa\sum_\tau\operatorname{Ask}_p(w(\tau),\tau)\Bigr).
\end{align*}
Both asks and gains vanish for split primes and \(p\ge x_1\).
For \(v\le a\), summation of \(2\operatorname{Ask}_p(v,\tau)\log p\)
gives \(\operatorname{Ask}(v,\tau)\) by the weights of prime powers.
Thus on all the values of the unprojected lift \(\iota_1\mathcal P(c_*\phi)\)
the logarithmic squared-norm bound is
\begin{equation}
\label{eq:weighted-A-chain-bound}
                 A+h_0+2\sum_p\gamma_p\log p+\mathcal B .
\end{equation}
Here \(A\) bounds the logarithmic squared norms of the input values,
and the height with the same signed \(\mathrm{rem}\) and levels is
\[
   h_0=h_{\rm main}-m\,\mathrm{rem}
       +T\bigl(m(S(a)-\operatorname{MainAsk})-a\log N\bigr).
\]

Concavity of \(\log\) on the intervals from \(1\) to \(a\) gives
\(\log(a!)\le a\log a-a+1+\tfrac12\log a\).
Hence \(S(a)-\operatorname{MainAsk}\le
a(\bar w+G)+1+\tfrac12\log a\).
The lower endpoint of the level window is
\(m(\bar w+G)-G-\tfrac{15}{2}\log b\).
The terms \(ma(\bar w+G)\) cancel against this choice of level.
Since \(G=O(\log b)\), \(T\le f\), and \(m/a\to0\), this proves
\(h_0\le h_{\rm main}-m\,\mathrm{rem}+\mathcal B\).
For \(G\) one can use \(Q\) in the different to get
\(\operatorname{rd}_E\le Q\).

The \hyperlink{proof.weighted-projection-transfer}{projection formula}
preserves the integral gauge at each coefficient embedding. Its
logarithmic squared-norm cost at the trivial twist is at most \(x^{99}\).

\paragraph{Collision counts and the loss budget.}
\teamanchor{proof.weighted-local-collisions}%
Call \(p\) small non-split when \(p<x_1\) and \(\chi_d(p)\ne1\).
At such a prime write \(c_p=\lceil\log_p m\rceil\) and
\begin{align*}
 \ell_p&=\frac{fa(\log b)^2}{p}
            \begin{cases}
                 m & (p\equiv \pm1\bmod m),\\
                 1 & (\text{otherwise}).
            \end{cases}\\
 e_p&=
 \begin{cases}
   3\ell_p & (p\text{ small non-split}),\\
   4mfb & (p=N),\\
   0 & (\text{otherwise}).
 \end{cases}
\end{align*}
These cases are disjoint. At the short residue classes write
\(p=mc\pm1\) with \(c\ge1\), whence \(m/p\le2/c\).
Bound the sum of \(\log p/p\) at the other small primes by the
harmonic sum over all integers up to \(x_1\), multiplied by
\(\log x_1\). The two short residue classes have the same order of
bound after using \(c\le 2x_1\). Therefore
\[
 \sum_{p<x_1} e_p\log p=O\bigl(fa(\log b)^2(1+\log x_1)^2\bigr).
\]
Also \(4mfb\log N=O(m^2fb\log b)\). Since
\((\log x_1)^2=(\log b)^{1/3}\) and \(m^2b\le4a\), these errors
total \(o(\mathcal B)\), in particular
\begin{equation}
\label{eq:weighted-A-loss}
                  \sum_p e_p\log p\le\mathcal B .
\end{equation}

For \(p\) small non-split, \(p\) is a unit modulo \(N_1\). To see this,
the primes dividing \(M\) are \(2,3,5,7,11\). For each of the odd
ones, \(d\equiv-1\pmod p\) and \(\chi_d(p)=(-d/p)=1\).
For \(2\) use \(d\equiv7\pmod8\). At \(p=d\) we have
\(d\equiv-1+M/2\pmod M\), so \(d^2\equiv1\pmod M\) and
\(-d^e\) is never \(1\bmod M\) for positive \(e\).

For an inert \(p\), write \(u_p=[-p]\in(\mathbf Z/N_1\mathbf Z)^\times\)
and let \(o_p\) be the order of its reduction modulo \(M\).
When \(e\) is odd, a top \(\sigma\) has \([-p^e]\sigma\) a top
exactly when \(o_p\mid e\).
We call an inert \(p\) served when \(p\not\equiv\pm1\bmod m\)
and \(o_p\) is odd and at most \(c_p\).
For any set \(I\) of infinite places, define
\(\operatorname{coll}_p(I)\) as the sum of \(a/p^e\) over the live
pairs \(e\ge1\) and \(\sigma\) for which both \(\sigma\) and
\([-p^e]\sigma\) are tops outside \(I\).
Here live means that \(e\) is odd or \(p=d\).
If all such \(e\) are at least \(e_0\), the sum is at most
\(2fa/p^{e_0}\) by the geometric series.
At an unserved prime this gives
\[
        m\operatorname{coll}_p(\varnothing)\le\ell_p .
\]
Indeed there is no collision for \(p=d\) or for even \(o_p\).
For \(p\equiv\pm1\bmod m\) use the bound with \(e_0=1\) and the
extra \(m\) in \(\ell_p\). In the remaining unserved case
use \(e_0=c_p+1\) and \(p^{c_p}\ge m\).

For a served \(p\), consider the finite places induced by the tops
at the given coefficient embedding \(j\). Write \(\lambda(\sigma)\)
for the place induced by \(j\circ\Phi_\sigma\).
Decomposition in the compositum \(E\) gives
\[
 \lambda(\sigma')\in
 \{\lambda(\sigma),\overline{\lambda(\sigma)}\}
   \quad\Longleftrightarrow\quad
      \sigma'\in\langle u_p\rangle\sigma .
\]
In fact the Frobenius generator, in the product Galois group of \(K\)
and the cyclotomic factor, is \((\operatorname{conj}_K,[p])\),
and complex conjugation is \((\operatorname{conj}_K,[-1])\).
Their product acts as \(u_p\) on the cyclotomic label. Its odd
powers interchange the two conjugate places, and its even powers
fix each. The places in a pair are distinct. Indeed the order of \(u_p\)
is divisible by \(m\), by reduction modulo \(2^s\). To see the
divisibility, put \(v=v_2(p^2-1)\). Since \(p\not\equiv\pm1\bmod m\)
and exactly one of \(p-1,p+1\) has \(2\)-valuation one,
\(v\le\log_2 m\). Successive squaring gives
\(v_2(p^{2^k}-1)=v+k-1\) for \(k\ge1\), so the order modulo \(2^s\)
is divisible by \(2^{s-v+1}\), hence by \(m\). In particular the
cyclotomic order of \(p\) is even and the Frobenius decomposition
group cannot contain complex conjugation: if \([p]^k=[-1]\)
for odd \(k\), then \(u_p^k=1\), contradicting the even order.

Since \(o_p\) is odd, multiplication by \(u_p^{o_p}\) pairs top
places inducing conjugate primes. Choose one prime from each
such pair, giving the marked set \(\Lambda\).
Let \(I=\{\sigma:\lambda(\sigma)\in\Lambda\}\).
For every live collision, one of the two places belongs to \(I\),
so \(\operatorname{coll}_p(I)=0\).
An orbit through a top reduces modulo \(M\) to a cyclic group
of odd order. It cannot contain \(-1\), so it contains no thick
place. Thus multiplication by \(u_p\) bijects the two halves of
each marked pair, preserving every \(w_j(\sigma)\).
If \(p\ge m\), then \(c_p=1\) and \(o_p=1\), so the entire marked
orbit consists of tops. These facts apply for each coefficient
embedding, including its composition with complex conjugation.

\paragraph{The lower exponent in one step.}
\teamanchor{proof.weighted-local-jet-count}%
Write \(B_{j,\sigma}\) for the number of boxes in columns of heights
different from \(1\) and less than \(n_{j+1}\) at step \(j\), and
\(r_{j,\sigma}\) for their number of columns.
The height-\(\kappa\) term and the polynomial bounds on higher
multiplicities give, uniformly in \(j<m\) and \(\sigma\),
\[
  \kappa w(\sigma)\le B_{j,\sigma}\le\kappa w(\sigma)+x^{60},
       \qquad r_{j,\sigma}\le w(\sigma)+x^{60}.
\]
For the lower bound, \(w_1(\sigma)\) includes \(3r^8\) in addition
to the rounded-up nominal weight, exceeding the term
\(\lfloor(\kappa+t+1)/2\rfloor\) subtracted at height \(\kappa\).
For the upper bounds all added boxes come from rounding by at
most \(N_1\), the stage-one gap, and \(O(m)\) polynomial-sized
columns of height at most \(n_m\).

For a prescribed \(I\), put \(t_\sigma=a1_{\sigma\in\mathcal
T\setminus I}\) and write \(\sigma'=[-p^e]\sigma\) for a live exponent.
The coupling count in the Schur jet estimate at \((\sigma,e)\) is
\[
 K_{\sigma,e}
  =\min(t_\sigma/p^e,r_{j,\sigma'}+t_{\sigma'})
     + r_{j,\sigma}/p^e+\min(B_{j,\sigma}/p^e,B_{j,\sigma'}).
\]
Let \(B_j=Ta+\sum_\sigma B_{j,\sigma}\), \(H_p=\lfloor\log_p B_j\rfloor\),
and \(\omega_p=(1-\chi_d(p))/2\). Put
\(\mathscr E_p=\{e\ge0:e\text{ odd or }p=d\}\).
The Schur jet exponent before the charges is
\[
 J_{p,j}(I)=\omega_p B_j
          \sum_{\substack{e\in\mathscr E_p\\e>0}}p^{-e}
      -\omega_p\!\sum_{\substack{e\in\mathscr E_p\\e\le H_p}}
                    \sum_\sigma K_{\sigma,e}
      -fn_{j+1}^3(H_p+1).
\]
Here \(K_{\sigma,0}\) uses \(\sigma'=[-1]\sigma\). The ramified
normal form permits a quadratic term and equal Frobenius powers in
the exponential, so the coupling set includes zero when \(p=d\).
The first term is the factorial gain less the CM contribution per
box, using the geometric series over the positive live exponents.

\paragraph{The normalized elliptic germs.}
\teamanchor{proof.weighted-local-cm-germs}%
Use the marked model of \(\theta\) above, over the valuation ring
\(\mathcal O\) of \(\mathbf C_p\), with \(v_p(p)=1\).
Let \(X_0(z)=\wp(z+t_0)-\wp(t_0)\) and
\(Y_0(z)=(\wp'(z+t_0)-a_1X_0(z))/2\).
Here \(z\) is the invariant differential coordinate and the subscript
only denotes the marked elliptic coordinates. Thus \(Y_0=z+O(z^2)\)
and \(X_0=z^2+O(z^3)\). The cubic and its derivative give
\[
 X_0\in\mathcal O[[Y_0]],\qquad
 Y_0'=1+3X_0^2-a_1Y_0 .
\]
For the first assertion solve the cubic for \(X_0\) at the marked
point, where its derivative in \(X_0\) is a unit. Consequently
\(\mathcal O[[Y_0]]\subset\mathbf C_p[[z]]\) is closed under
\(d/dz\).

For \(0\ne\delta\in\mathcal O_K\) and \(\delta y\in\Lambda\), put
\(n_\delta=\delta\bar\delta\). Then
\begin{equation}
\label{eq:weighted-local-theta-isogeny}
 (U_y\theta)(\bar\delta z)=\theta(z)^{n_\delta}G_{\delta,y}(z),
 \qquad G_{\delta,y}\in\mathcal O[[Y_0]] .
\end{equation}
The two sides have the multiplier of a degree-\(n_\delta\) section.
The canonical multiplier
\(\exp(\pi i E_\Lambda(\bar\omega l,l))\) pulls back under
\(\bar\delta\) to its \(n_\delta\)-th power. The translation
phase is one since \(E_\Lambda(\bar\delta l,y)=E_\Lambda(l,\delta y)\)
is integral. The isogeny height ratio
\(H(\bar\delta z)/H(z)^{n_\delta}\) is bounded. This follows from the
section basis of the marked cubic, applied to the pullbacks of
\(\theta\) and \(s_\Lambda^2\), of degrees \(n_\delta\) and
\(2n_\delta\). Under doubling the ratio becomes its fourth power.
Since doubling on the complex uniformizing line is onto, its finite
positive supremum is one. The torsion-translation height calculation
above then gives
\(\lvert(U_y\theta)(\bar\delta z)\rvert_p\le H(z)^{n_\delta}\).

To read this bound in the formal ring, write the quotient by
\(\theta^{n_\delta}\) as \(P(X_0)+Y_0Q(X_0)\) in the section basis.
For any integral \(x\) whose reduction avoids the roots of
\(4x^3+a_1^2x^2+4x\), the two points above \(x\) have integral
\(Y_0\) values with unit difference. The discriminant polynomial
has nonzero reduction, including at \(p=2\) where \(a_1\) is a unit.
At these two points \(H/\lvert\theta\rvert_p=1\), so the two
quotient values give \(P(x),Q(x)\in\mathcal O\). If a polynomial had
a coefficient of negative valuation, divide by a coefficient of
least valuation and choose a residue outside the finite excluded
set and the roots of the resulting nonzero polynomial. This proves
that \(P,Q\) have integral coefficients and gives
\eqref{eq:weighted-local-theta-isogeny}. The same argument for an
even section of degree \(n_\delta\) gives a polynomial in \(X_0\)
alone of degree at most \(\lfloor n_\delta/2\rfloor\).

Put \(L=-\log\theta\) as a formal series, with \(L(0)=L'(0)=0\).
Logarithmic differentiation of the translated sigma function gives
\(L''=L_2+X_0\) for a constant \(L_2\).
The logarithmic derivative of \(G_{\bar a,0}\) is integral in
\(\mathcal O[[Y_0]]\) for \(0\ne a\in\mathcal O_K\), since
its constant term is one. Applying \(d/dz\) and evaluating at zero
gives \((a\bar a-a^2)L_2\in\mathcal O\). Subtract the readings for
\(a=\omega,\bar\omega\). As \(\omega+\bar\omega=1\), this proves
\(\sqrt{-d}\,L_2\in\mathcal O\).
\footnote{\weightedpin{grh-compact/src/Lib/CM/LocalHeight/IsogenyIntegrality.lean}
{82-L117}{The section-basis integrality argument} and
\weightedpin{grh-compact/src/Lib/CM/KroneckerLimit/LogTheta.lean}
{141-L178}{The logarithmic derivative and the constant term}.}

Now take \(p\) small non-split. Let \(k_p=2\) at inert \(p\) and
\(k_p=1\) at \(p=d\), set \(q_0=p^{k_p}\), \(Q=q_0-1\), and use
\(\pi_0=p\) or \(\sqrt{-d}\), respectively. Then
\(\pi_0\bar\pi_0=q_0\) and \(v_p(\pi_0)=\omega_p\).
All these \(p\) are at least \(13\). The \(Q\)-th roots of unity
are in \(K_p\) with norm \(\zeta^{p+1}\).
For \(a\) in the closure of \(\mathcal O_K\), the formal action
\(z\mapsto az\) sends \(Y_0\) into \(\mathcal O[[Y_0]]\).
For \(a\in\mathcal O_K\) this follows by extending the
endomorphism over the good elliptic model. General \(a\) follows
by coefficientwise limits in the fixed invariant coordinate.
Make the following rotation-equivariant choice of parameter
\[
 T_0(z)=\frac1Q\sum_{\zeta^Q=1}\zeta^{-1}Y_0(\zeta z)
          =z+O(z^{q_0}).
\]
This is an integral change of formal coordinate, and it and its
inverse \(Z(T_0)\) have only degrees \(1\bmod Q\).
The reduction of \([\pi_0]\) is purely inseparable of degree \(q_0\).
At the ramified prime its differential is zero and its degree is
\(p\). At an inert prime, an ordinary special fiber would have a
one-dimensional geometric étale \(p\)-torsion space on which the
unital action of \(\mathcal O_K/p=\mathbf F_{p^2}\) is impossible.
Thus the reduction of \([p]\) has inseparable degree \(p^2\).
Here the CM endomorphism is taken on the good model at the origin
corresponding to \(z=0\) under translation by the marked point.
Its first nonzero coefficient in the reduced \(T_0\) parameter
is in degree \(q_0\). If \(Z_{q_0}=[T_0^{q_0}]Z\), this says that
\((\pi_0-\pi_0^{q_0})Z_{q_0}\) is a unit.

Choose \(\varpi^Q=-1/Z_{q_0}\) and put \(\xi=\varpi^Q\).
We have \(v_p(\xi)=\omega_p\). Also \(Z'\) is integral since
\(dT_0/dz\) is an integral unit in the \(T_0\) parameter.
For \(j\ge2\), the coefficient in degree \(1+jQ\) of
\(\varpi^{-1}Z(\varpi X)\) has valuation at least
\(j\omega_p-v_p(1+jQ)\ge\omega_p\).
For \(2\le j\le4\), \(1+jQ\) is prime to \(p\). For \(j\ge5\),
\(1+jQ<p^{(j-1)\omega_p}\) follows at \(j=5\) from \(p\ge13\)
and then by induction. Thus this series is integral and is
\(X-X^{q_0}\) modulo \(\xi\). Its compositional inverse
\(S(X)=T_0(\varpi X)/\varpi\) is integral and satisfies
\[
          S(0)=0,\qquad S-S^{q_0}=X\pmod{\xi}.
\]

Let \(L^\flat\) be the part of \(L\) in degrees \(p+1\bmod Q\).
The same rotation average proves \(W_0=\exp(L^\flat-L)\in
\mathcal O[[T_0]]\). In fact
\(\exp(a\bar a L(z)-L(az))\) is an integral unit of constant term
one by \eqref{eq:weighted-local-theta-isogeny}, and this persists
under the coefficientwise limit \(a\to\zeta\).
Taking the \(-\zeta^{-(p+1)}\)-th power preserves integrality
by the integral binomial series, since the exponent is in
\(\mathbf Z_p\). Multiplying these powers over all rotations
and taking the \(1/Q\)-th power proves the assertion.

For the exponent, first use
\eqref{eq:weighted-local-theta-isogeny} for \(\pi_0\) and the
even section \(\theta(\pi_0 z)\). The quotient is \(R(X_0)\),
where \(\deg R<q_0/2\), and it equals
\(\exp(q_0 L-L(\pi_0 z))\). In degrees \(n<q_0\) all
coefficients of \(q_0 L-L(\pi_0 z)\) have valuation at least
\(\omega_p\). Indeed \(X_0\) in these degrees has integral
coefficients, since \(T_0=z\bmod z^{q_0}\). For \(3\le n<q_0\),
\(v_p(n(n-1))\le\omega_p\), and the degree two uses the integral
bound on \(\sqrt{-d}L_2\). The factor \(q_0-\pi_0^n\) for \(n\ge2\)
has valuation at least \(2\omega_p\). The exponential in the same
range is one modulo \(\pi_0\), by
\(v_p(r!)\le (r-1)/(p-1)\). Because \(X_0=z^2+\cdots\),
successively reading the coefficients of degrees \(2j<q_0\)
shows that \(R-1\) has all coefficients divisible by \(\pi_0\).
Its logarithm is then in \(\pi_0\mathcal O[[T_0]]\), again by
the factorial bound, since \(R(0)=1\).
Taking degrees \(p+1\bmod Q\) gives
\[
  q_0 L^\flat-L^\flat(\pi_0 z)
   =\pi_0 T_0^g(\lambda+T_0^Q H_0(T_0)),\qquad
  g=\begin{cases}p+1 & (k_p=2),\\2 & (k_p=1),\end{cases}
\]
where \(\lambda\in\mathcal O\) and \(H_0\in\mathcal O[[T_0]]\).
The projector reads the same congruence classes in \(T_0\)
because \(T_0(\zeta z)=\zeta T_0(z)\).

Write \(c_\varpi=\varpi^{p+1}/p\) and
\(\Xi(X)=-c_\varpi^{-1}L^\flat(\varpi X)\).
Dividing the last identity by \(-q_0 c_\varpi\) shows that
\((1-\pi_0^n/q_0)[X^n]\Xi\) is integral and congruent to
\(\lambda_p[X^n]S^g\) modulo \(\xi\), with
\(\lambda_p=-\pi_0\varpi^g\lambda/(q_0c_\varpi)\in\mathcal O\).
The factor multiplying \(\lambda\) has valuation zero.
For \(n\ge3\), \(1-\pi_0^n/q_0\) is one modulo \(\xi\).
Degrees zero and one vanish, and degree two vanishes in the inert
case. At the ramified prime \(1-\pi_0^2/q_0=2\) and
\([X^2]S^2=1\). Therefore \(\Xi\) is integral and
\begin{equation}
\label{eq:weighted-local-cm-germ}
 \Xi\equiv
 \begin{cases}
   \lambda_p S^{p+1} & (k_p=2),\\
   \lambda_p(S^2-X^2/2) & (k_p=1)
 \end{cases}
 \pmod{\xi}.
\end{equation}
In the inert case \(\lambda_p\) is a unit. For this, set
\(h_j=[z^{p-1}]X_0^j\) and reduce modulo the maximal ideal
of \(\mathcal O\). The coefficients in degree \(p-1\) of
\((X_0^jX_0')'\) vanish, since
\([z^p](X_0^jX_0')\) is integral. The cubic gives
\[
 (X_0^j X_0')'=(4j+6)X_0^{j+2}
                +(j+1)a_1^2 X_0^{j+1}+(4j+2)X_0^j .
\]
If \(h_1\) vanished on reduction, then \(h_0=0\) and the
recurrence would force \(h_{(p-1)/2}=0\), since \(4j+6\) is prime
to \(p\) for \(0\le j\le(p-5)/2\). But \(h_{(p-1)/2}=1\).
Now \([X^{p+1}]\Xi=-h_1/(p+1)\).

Finally choose a positive integer \(n\) prime to \(p\) with
\(ny\in\Lambda\). Apply \eqref{eq:weighted-local-theta-isogeny} at
\(z=\varpi X/n\), using \(\delta=n\). This gives
\[
 (U_y\theta)(\varpi X)
    =V_y(X/n)\exp\!\bigl(n^2c_\varpi\Xi(X/n)\bigr),
 \qquad V_y\in\mathcal O[[X]] .
\]
The factor \(V_y\) is the product of the rescaled \(G_{n,y}\)
and \(W_0^{n^2}\); both are integral after scaling because
\(T_0(\varpi X)=\varpi S(X)\). The unit variable rescaling
and the exponent scalar depend only on \(n\).
We have \(v_p(n^2c_\varpi)=1/(p-1)-\omega_p\) in both cases.
At a ramified prime \(X=S-S^p\) modulo \(\xi\), so
\(\lambda_p(S^2-X^2/2)=
\lambda_p(2S^{p+1}-S^{2p}+X^2/2)\) modulo \(\xi\).
This is the ramified shape used in the tensor calculation, with
\(2p=p+p^{k_p}\).
\footnote{\weightedpin{grh-compact/src/New/Curve/FormalParameter.lean}
{125-L150}{Compact parameter and translate statement},
\weightedpin{grh/src/New/Derivatives/FormalGroup/FrobeniusCongruence.lean}
{143-L274}{The degree bound and Frobenius congruence}, and
\weightedpin{grh/src/New/Derivatives/FormalGroup/Exponent.lean}
{416-L487}{The two coefficient readings}.}

\paragraph{The normalized products and the source gauge.}
\teamanchor{proof.weighted-local-source-gauge}%
The same one-variable model applies to the standard-data module
in the placed Weyl coordinates. Write \(\nu=N_1\), \(s_1=s_0/u_0\),
\(\mathfrak a_0=E\cap\alpha\widehat{\mathcal O}_E\), and use
\(L=r_\infty(\mathfrak a_0)\) and \(\eta\) from the vacuum construction.
We first construct a tight frame with frame constant \(M_p'^2\)
in \(\mathfrak a_0\), where \(M_p'\) is an integer prime to \(p\).
List each vector with its multiplicity. For \(p\nmid\nu\) use
\(M_p'=\nu\) and the powers of \(\zeta\), each repeated \(\nu\)
times, as in the original vacuum calculation.

For \(p\mid\nu\) write \(A_0=(d+1)/4\), \(\nu'= \nu/p^{v_p(\nu)}\),
\(r_0=\nu/p\), \(\zeta_p=\zeta^{r_0}\). All primes dividing
\(\nu\) divide \(A_0\). There is an integer \(z_0\) with
\(A_0z_0^2+z_0+1\equiv0\bmod\nu\). Indeed
\(u\mapsto A_0u^2+u\) is injective modulo \(\nu\), since the
difference at \(u,v\) is \((u-v)(1+A_0(u+v))\) with the second
factor a unit. Put \(t_1=\bar\omega+A_0z_0\) and
\(c_1=t_1\nu'/\nu\). Use the list
\[
  w_{jkh}=\zeta^j\bigl(1_{k=0}\zeta_p^h
                   +c_1(\zeta_p^h-\zeta_p^k)\bigr),
      \qquad 0\le j<r_0,\quad 0\le k,h<p.
\]
These vectors lie in \(\mathfrak a_0\). On the denominator
side of \(\alpha\) at a prime over \(p\), this uses
\[
  v_p\bigl(\Phi_\nu'(\zeta)\bigr)+v_p(\zeta_p-1)
       =v_p(\nu).
\]
In fact the prime-to-\(p\) cyclotomic factor is unramified and
the different of the \(p^u\)-part is generated by
\(\Phi_{p^u}'(\zeta_{p^u})\), of valuation \(u-1/(p-1)\),
by differentiating
\((X^{p^u}-1)/(X^{p^{u-1}}-1)\).
On the opposite side \(\bar\omega\) vanishes modulo the prime
and \(\bar t_1\) is a unit. Since
\(t_1\bar t_1=A_0(A_0z_0^2+z_0+1)\), \(t_1\) is divisible
by \(p^{v_p(\nu)}\) there. At every other prime
\(\nu'/\nu\) is a unit.

For two embeddings \(\sigma,\tau\), the sum over \(h,k\) of
\(\Phi_\sigma(w_{jkh})\overline{\Phi_\tau(w_{jkh})}\)
vanishes if their \(\zeta_p\) values differ, by geometric sums.
If the values agree, that sum is
\[
 p(1+c_1+\bar c_1+2p c_1\bar c_1)
      \Phi_\sigma(\zeta)^j\overline{\Phi_\tau(\zeta)^j}.
\]
The sum over \(j\) then vanishes unless \(\sigma=\tau\).
The common diagonal value is
\[
 M_p'=\nu(1+c_1+\bar c_1+2p c_1\bar c_1)
     =\nu+\nu'(1+2A_0z_0)
           +2p(\nu')^2 t_1\bar t_1/\nu .
\]
It is a positive integer congruent to \(\nu'\) modulo \(p\).
Repeating this list \(M_p'\) times gives the desired identity
\(\sum_i \Phi_\sigma(w_i)\overline{\Phi_\tau(w_i)}
=(M_p')^2\delta_{\sigma\tau}\).
\footnote{\weightedpin{grh-compact/src/New/Curve/TightFrameExists.lean}
{87-L148}{The compact frame construction with multiplicities}.}

Use the forms
\[
   L_{i,c}(Z)=\frac{\Omega_0}{M_p'\eta}
         \sum_\tau\overline{\Phi_\tau(w_i)}Z_{\tau,c}
\]
in each matrix coordinate \(c\). With \(\prod_{i,c}\theta(L_{i,c})\),
the trace-dual identity \(\mathfrak a_0\bar{\mathfrak a}_0
=\mathfrak D_{E/K}^{-1}\), the frame identity and the canonical
multiplier give the same construction on \(M_p' L\) as the
\hyperlink{proof.weighted-vacuum-projection}{original vacuum construction}
on \(\nu L\). Let \(F\) be its finite Weil vector. Its evaluation
at zero is one. This product recovers the previously normalized
vacuum with integral coefficients, as follows. Expand \(F\) in the
translates of the normalized vacuum indexed by
\((M_p')^{-1}L/L\) in each coordinate, and for each translate
apply the inverse translation followed by the projection to the
vacuum. This projection averages the multiplier-normalized
translations from \(L/M_p'L\). Multiply its output by the value
at zero of the corresponding translated vacuum and sum. The
result is \(F(0)\) times the normalized vacuum. Each value used
as coefficient is integral by the established vacuum bound, and
the averaging denominator is a power of \(M_p'\).
Every phase is a root of unity, and each shift in the formula
is integral after multiplication by \(M_p'\) in the raw finite
coordinates. This also proves the expansion with the density of
rational finite translations used in the original construction.
\footnote{\weightedpin{grh-compact/src/Lib/Weil/Construction/ConjugateJetDefs.lean}
{27-L36}{The recovery identity for the fixed vacuum}, and
\weightedpin{grh-compact/src/New/Curve/FrameSections.lean}
{83-L126}{The product at the actual finite translates}.}

A datum in \(\mathrm{Std}_\delta\) is an integral linear combination
of finite translates of this normalized vacuum whose vectors become
integral above \(p\) after multiplication by \(\delta\).
For each term of the vacuum recovery choose a rational finite
translation \(\alpha^{-1}\gamma\) with the same value as the
actual translating operator on \(F\). This is possible by smoothness
and weak approximation, while retaining integrality above \(p\)
after multiplication by \(M_p'\delta\).
After multiplication by some integer \(n'\) prime to \(p\),
\(n'M_p'\delta\gamma\) belongs to \(\mathfrak a_0\).
Under the rational Heisenberg formula the elliptic torsion argument
is \(-\Omega_0\operatorname{Tr}_{E/K}(\gamma\bar w_i)/M_p'\).
Thus \(n\delta y_{i,c}\in\Lambda\) with \(n=n'(M_p')^2\)
prime to \(p\). This works for a common \(n\) through each finite sum by
taking a common multiple of the finitely many denominators prime to
\(p\). In the normalized jet of scale \(c_0\,\mathrm{per}\), the
linear coefficients after reading by \(j_p\) are
\begin{equation}
\label{eq:weighted-local-frame-coefficients}
 z_{\tau i}=j_p\!\left(\frac{c_0s_1}{M_p'}
                               \overline{\Phi_\tau(w_i)}\right).
\end{equation}
Here \(c_0\) is the scalar in the kernel scale, so it is
\(\bar\pi_G\) in the chain of short lifts. In particular
for the small non-split \(p\) prime to the twist we use \(\delta=1\).
Such \(p\) is prime to \(\nu\) and off \(g_0,g_1,N\).
Consequently \(s_1,c_0,M_p'\) are all \(p\)-adic units and
\eqref{eq:weighted-local-cm-germ} applies with a common exponent
to the elliptic factors of each standard-data summand. This
supplies the placed products and the coefficient
\(\beta=j_p(c_0s_1/M_p')\) in the tensor calculation below;
the common \(1/n\) in the germ may be included in \(\beta\).

For every \(p\), \(j_p(s_1\overline{\Phi_\tau(w_i)})\) is integral.
To see this, the finite splitting-character valuation formula on
\(\alpha\) reads the product of the local absolute values of
\(\alpha\) through \(j_p\Phi_\tau\operatorname{conj}\). This
formula follows by taking a power of the ideal of \(\alpha\)
which is principal and using the infinity type on its generator.
The absolute values of \(\alpha\) in the product are all equal.
They are one for \(p\nmid\nu\), and for \(p\mid\nu\) every
weight reading lies on the same chosen quadratic side and the
relative different valuation is constant there.
Since \(s_0^f=\chi_{\mathrm{fin}}(\alpha)^{-1}\) and \(u_0\)
is a unit, \(|j_p(s_1)|_p\) is the inverse of that common value.
Now \(w_i\in\mathfrak a_0\) proves the claim.
\footnote{\weightedpin{grh-compact/src/Lib/ClassField/HeckeCharacter/SplittingCharOnU1.lean}
{70-L82}{Valuation of the finite splitting character},
\weightedpin{grh-compact/src/New/Curve/TightFrameExists.lean}
{28-L55}{The resulting period valuation}, and
\weightedpin{grh-compact/src/New/Curve/ConjugatesOfJets.lean}
{50-L84}{The placed product expansion}.}

For the ordinary source-gauge bound, \(d(d/dz)\) preserves the
algebra \(\mathcal O[[Y_0]][L',\theta]\).
It sends \(L'\) to \(d(L_2+X_0)\) and \(\theta\) to
\(-dL'\theta\), with \(dL_2\in\mathcal O\) by the logarithmic
derivative calculation. Every member of the algebra has integral
constant term. It follows from \eqref{eq:weighted-local-theta-isogeny}
that all derivatives of \((U_y\theta)(\bar\delta dz)\) at zero
are integral when \(\delta y\in\Lambda\). Replace \(\delta\) by
\(n\delta\) for each product above, which changes no absolute
value. The multivariable chain rule and
\eqref{eq:weighted-local-frame-coefficients} show that a homogeneous
differential polynomial of degree \(B\) with integral coefficients
has value bounded by
\[
               \left|j_p\!\left(\frac{c_0}{\bar\delta d}\right)\right|_p^B .
\]
The integral Schur jet polynomials have this degree with \(B\)
equal to the number of boxes below the target rank. This gives the
bound on coefficient vectors at the identity.

For arbitrary source and target arguments, use rational unitary
changes to make both arguments integral at all places above \(p\).
Weak approximation follows from the Cayley chart and approximation
in its vector space. Rational equivariance transports the two
Weyl gauges exactly. Integral unitary changes preserve
\(\mathrm{Std}_\delta\). Indeed their action on a vacuum is a
sum of finite Heisenberg translates indexed by the quotient of the
standard lattice by its intersection with the inverse image of that
lattice. The coefficient is the finite Weil multiplier divided by
its complex squared norm. The translating vectors are integral
above \(p\), and the coefficient is a \(p\)-adic unit since the
local determinant factors of the integral automorphisms are units.
The two additional character values in the paired action are units
for the same reason. Heisenberg covariance carries the integral
condition on \(\delta x\) through these changes.
With a marked diagonal inserted on the target, the same calculation
leaves the diagonal action on the target tableau and determinant
character. Each tableau entry then contributes the absolute value
of its diagonal coordinate, which is the charge calculation in
the marked chain.
\footnote{\weightedpin{grh-compact/src/Lib/Weil/Construction/Projections.lean}
{210-L220}{The vacuum under an integral change},
\weightedpin{grh-compact/src/New/Curve/ConjugatesOfJets.lean}
{24-L50}{Stability of the standard-data module}, and
\weightedpin{grh-compact/src/New/Curve/ConjugatesOfJets.lean}
{114-L153}{The placed gauges and the scale per box}.}

Away from \(N\), the good datum \(\Psi_\vartheta\) for a twist
of level norm prime to \(p\) gives
standard data with \(\delta=\pi_G\) and \(c_0=\bar\pi_G\).
The ratio per box is \(\lvert j_p(1/d)\rvert_p\).
At \(N\), the last-column cut uses the same source-gauge
calculation after clearing the last-column translation at the
places over the quadratic place induced by \(j_N\).
Take \(\delta=\omega-c\), where \(c\in\mathbf Z\) has the same
residue as \(\omega\) at that place. This clears the simple poles
there, and \(\bar\delta\) is a unit at \(j_N\) because
\(\omega-\bar\omega=\sqrt{-d}\) and \(N\nmid d\).
Also \(\bar\pi_G\) is a unit at \(j_N\), so the derivative gauge
has unit bound per box. Before that derivative bound, the
residue-modulus averaging is cleared by \(N^{2f n_j n_{j+1}}\).
It contributes at most \(N^{fb}\) per lift. The subsequent
paired Fourier move and its character cost have their separate
bounds in the level-prime calculation, yielding the displayed
\(N^{3mfb}\) bound under the starting gauge condition.
\footnote{\weightedpin{grh-compact/src/New/Columns/ChainValuation.lean}
{281-L319}{The allowable translated vacua},
\weightedpin{grh-compact/src/New/Columns/LocalStatement.lean}
{244-L250}{The ordinary branch on the same chain}, and
\weightedpin{grh-compact/src/New/Columns/TwistedHeckeType.lean}
{43-L74}{The \(N\)-adic cut and its averaging denominator}.}

\paragraph{The tensor and Schur calculations for the charged jets.}
\teamanchor{proof.weighted-local-jet-algebra}%
Use \(v_p(p)=1\) and the valuation ring \(\mathcal O\) of
\(\mathbf C_p\). Every small non-split prime in this application
is at least \(13\), since \(2,3,5,7,11\) split. The normalized
elliptic factors have the common exponent and residue shape of
\eqref{eq:weighted-local-cm-germ}, with
\(v_p(\varpi)=\rho_p\), \(v_p(\xi)=\omega_p\), and
\(v_p(\alpha_p)=1/(p-1)-\omega_p\).
The \hyperlink{proof.weighted-local-source-gauge}{placed product}
gives the coefficients \eqref{eq:weighted-local-frame-coefficients}.
Include the common unit change of variable from the torsion
denominator in \(\beta\). The frame constant and the torsion
denominator are prime to \(p\) in this case.

For one product let \(L_{i,c}=\sum_\sigma z_{\sigma i}X_{\sigma,c}\)
be its linear arguments, with \(c\) a matrix coordinate.
The frame identity gives
\(\sum_i \Phi_\tau(w_i)\overline{\Phi_{\sigma'}(w_i)}
      =l\,1_{\tau=\sigma'}\),
where \(l\) is the tight-frame constant. In the residue calculation the
\(z_{\sigma i}\) are a common integral scalar \(\beta\) times
\(j_p(\overline{\Phi_\sigma(w_i)})\). Any cyclotomic denominator of
\(w_i\) divides a power of \(N_1\) and is a unit here. Raising
\(j_p(\overline{\Phi_\sigma(w_i)})\) to the \(p^e\)-th power modulo
\(\xi\) gives \(j_p(\Phi_{[-p^e]\sigma}(w_i))\) for
\(e\in\mathscr E_p\). For inert \(p\) this uses the odd Frobenius
power to cancel conjugation on \(K\). At \(p=d\), the two embeddings
of \(K\) have equal reductions modulo \(\xi\) because
\(v_p(\sqrt{-d})=v_p(\xi)=1/2\). On the cyclotomic factor the
statement is the Frobenius action on roots of unity of order prime to
\(p\). These readings give, modulo \(\xi\),
\[
 \sum_{i,c} L_{i,c}^{p^e}L_{i,c}
      = l\beta^{p^e+1}\sum_{\sigma,c}
               X_{\sigma,c}^{p^e}X_{[-p^e]\sigma,c}.
\]
In particular the matrix coordinate \(c\) agrees at both ends of a
coupling. Substitutions with zero constant term that kill the
couplings on the right through degree \(B_j\) also kill the exponent
of the product modulo \(\xi\) in these degrees. To see this,
iterate \(S=X+S^{p^k}\) in the truncated ring to write
\(S(L)=\sum_{r\ge0} L^{p^{kr}}\). Each term in the sums of
\(S(L)^{p+1}\) and \(S(L)^{p+p^k}\) is a Frobenius power of
\(\sum_{i,c}L_{i,c}^{p^e}L_{i,c}\) with a live \(e\).
When \(k=2\) the two indices of its Frobenius powers differ by
an odd integer. When \(k=1\) their difference can be zero, and
the term \(u_0L^2\) also uses \(e=0\). Terms with \(e>H_p\)
have degree greater than \(B_j\). Because
\(v_p(\alpha_p\xi)=1/(p-1)\) and
\(v_p(r!)\le r/(p-1)\), the exponential after any such
substitution has integral coefficients through degree \(B_j\).
The integral factors and their common unit rescaling preserve this
property. It passes to the expansion of a standard datum by the
ultrametric inequality.

For the Schur calculation use \(n=n_j\), \(n'=n_{j+1}\), and let
\(\kappa_{\sigma,h}\) be the number of columns of height \(h<n'\) in
the jet. Fix a target tableau \(\eta\) and integral \(n\)-row matrices
\(Y_\sigma\). The flag test is the product \(P_Y(Z)\) of the minors of
\(Z_\sigma Y_\sigma\) on the rows of the target columns and
the first \(h\) columns. The coefficient of a source tableau in the
expansion of this product in the source flag basis is the
corresponding Schur polynomial. This expansion also includes the
common determinant power at the source when full source columns occur.
Its coefficient estimates follow from the estimates of
\(P_Y(\partial)\Theta(0)\) for all integral \(Y\). Indeed, the flag
basis polynomials reduce to linearly independent polynomials over
the residue field. With a diagonal leading-term order their leading
monomials record the multiset in each row of the standard tableau,
which determines its filling. Multiplying by a common determinant
power preserves independence. If a coefficient had smaller valuation
than the claimed bound, division by a coefficient of minimal
valuation would give a nonzero polynomial over the infinite residue
field, with a nonzero value at the reduction of some integral \(Y\).
This would contradict the bound for that flag test.

Write \(y_{\sigma,a}\in\mathcal O^n\) for column \(a\) of \(Y_\sigma\).
We can change the first \(h\) directions at \((\sigma,h)\)
by an integral \(h\)-by-\(h\) matrix \(M_{\sigma,h}\) at controlled
determinant cost. A congruence on each changed row of a matrix,
given by a linear form on that row with modulus \(q_0\in\mathcal O\),
costs at most \(v_p(q_0)\). For one form pick a row with a
nonzero value of minimal valuation, subtract multiples of it from
the other rows, then multiply the chosen row by \(q_0\). If all
values vanish the identity suffices. Previously imposed congruences
survive integral row operations.

For a live \(e\), the desired pairing of directions \(u,v\) at
\(\sigma,\sigma'=[-p^e]\sigma\) is
\(\sum_a u_a^{p^e}v_a\equiv0\pmod{\xi}\).
When the restrictions are put on \(u\), use linear forms with
\(p^e\)-th roots of the coefficients and modulus
\(\xi^{1/p^e}\). Taking the \(p^e\)-th power modulo \(\xi\)
recovers the pairing. When the restrictions are put on \(v\),
use modulus \(\xi\). The following choices cover every
pair of heights except those with first height one and either
\(\sigma\in I\) or both second height one and \(\sigma'\in I\).
\begin{itemize}
\item For first height one outside \(I\), either require its sole
change coefficient to be a multiple of \(\xi^{1/p^e}\) or impose
one pairing congruence on each relevant partner height against
the initial height-one source direction. The two costs are
\(t_\sigma/p^e\) and \(r_{j,\sigma'}+t_{\sigma'}\).
\item For first height different from one and second height one,
put one congruence on the first height against the initial
height-one direction of the partner. This costs \(r_{j,\sigma}/p^e\).
\item For both heights different from one, clear all the
coefficients in each row of the change matrix on one side. Take
the smaller of the two costs \(B_{j,\sigma}/p^e\) and \(B_{j,\sigma'}\).
\end{itemize}
The costs in this list are in units of \(v_p(\xi)=\omega_p\)
and count each change matrix once for every column of its height.
In the last case clearing the coefficients means requiring each
entry of the matrix row to be divisible by the relevant modulus.
In the other cases where a pairing is imposed, the opposite
height-one change is a scalar times its initial direction.
All these congruences can be imposed together using the same
matrices, so that
\[
  \sum_{\sigma,h}\kappa_{\sigma,h}\,v_p(\det M_{\sigma,h})
     \le \omega_p\sum_{\substack{e\in\mathscr E_p\\e\le H_p}}
                     \sum_\sigma K_{\sigma,e}.
\]

For an exempt pair, use the prescribed charges \(g_{\sigma,h}(b)\)
with a positive charge at one end for each common matrix row \(b\).
In that row take independent variables for the changed directions
at all heights, multiplying the direction from height \(h\) by
\(p^{g_{\sigma,h}(b)}\). The substitution made from these rows kills
the live couplings modulo \(\xi\). For a nonexempt pair the
direction pairing vanishes by construction. For an exempt pair it
is multiplied by at least one factor of \(p\). Apply the preceding
product-integrality argument to this substitution and the common
factor \(\varpi\). Applying the product of determinant derivatives
in the new variables multiplies the original flag test by
\(\varpi^{B_j}p^{c(\eta)}
  \prod_{\sigma,h}\det(M_{\sigma,h})^{\kappa_{\sigma,h}}\),
where \(c(\eta)\) is the sum of charges in the target tableau.
This follows by the chain rule and multilinearity of each determinant
in the \(h\) direction slots. The number of variables is at most
\(fn'^3\), since it is the sum over \(\sigma,h<n'\) of \(n'h\).
Every derivative monomial has degree \(B_j\). Its factorial
multiplier has valuation at least
\(B_j/(p-1)-fn'^3(H_p+1)\), by the base-\(p\) digit formula
for each factorial. Subtracting the valuations of the scaling
factors proves that every jet entry has valuation at least
\[
 (1/(p-1)-\rho_p)B_j
    -\omega_p\!\sum_{\substack{e\in\mathscr E_p\\ e\le H_p}}
                 \sum_\sigma K_{\sigma,e}
    -fn'^3(H_p+1)-c(\eta)=J_{p,j}(I)-c(\eta).
\]
For the final equality, if \(p\) is inert then
\(1/(p-1)-1/(p^2-1)=\sum_{e>0,\ e\ {\rm odd}}p^{-e}\).
If \(p=d\), then \(1/(p-1)-\rho_p=\frac12\sum_{e>0}p^{-e}\).

For the nonheight-one part of the exponent, the difference
\(B_{j,\sigma}/p^e-\min(B_{j,\sigma}/p^e,B_{j,\sigma'})\)
is bounded below by
\(\kappa(w(\sigma)/p^e-w(\sigma'))_+-x^{60}\).
If \(\sigma\in\mathcal T\setminus I\) and \(\sigma'\notin\mathcal
T\setminus I\), the first minimum in \(K_{\sigma,e}\) similarly
costs at most \(x^{60}\) beyond the ask for these height-one boxes.
If both are tops outside \(I\), charging \(a/p^e\) covers the
whole shortfall from this ask. If \(\sigma\in\mathcal T\cap I\),
the first minimum is zero and the height-one part contributes
all of \(a/p^e\). For \(e=1\), its excess over
\((a/p-w(\sigma'))_+\) is at least \(a/x_1\).
Outside the tops the height-one part is zero.

The remaining term \(r_{j,\sigma}/p^e\) in the coupling count
is at most \((a/x+x^{60})/p\). The asks vanish beyond \(H_p\)
since \(B_j\ge a>x_1\). The tail of the first sum in \(J_{p,j}(I)\)
is nonnegative. Finally \(H_p+1=O(\log b)\) uniformly in \(p\),
since \(\log B_j=O(\log b)\) and \(p\ge2\).
The per-place error per exponent, including the factorial
rounding penalty, is bounded by
\[
  2x^{60}+(a/x+x^{60})/p+n_{j+1}^3=O(a/(pm))
\]
uniformly for \(p<x_1\). Here \(x_1=b^{o(1)}\) and \(a\ge xb\).
At \(p=d\) the zeroth coupling accounts for the quadratic term,
as in the \hyperlink{proof.generator-valuation-discriminant}
{discriminant-place calculation in the earlier route}.
The positive ask begins at exponent one.
For \(\sigma'=[-1]\sigma\), at most one place is thick and at most
one is a top, since their respective classes are \(-1\) and \(1\)
modulo \(M\). The first minimum is at most \(r_{j,\sigma'}\) if
\(t_\sigma\ne0\), and the last minimum is bounded by the number
of boxes at the thin place. It follows that
\[
   K_{\sigma,0}
       \le 2(a/x+x^{60})+\kappa a/x_1+x^{60}
       = o(a/(dm)).
\]
Here \(d\) is fixed, \(m\le2\sqrt{x}\), and \(x_1\) dominates every
fixed power of \(x\).
Summing over \(e\le H_p\) absorbs this error in
\(fa(\log b)^2/(pm)\) for large \(x\).
We obtain
\begin{equation}
\label{eq:weighted-A-jet-count}
 J_{p,j}(I)\ \ge\
 \frac{\gamma_p}{m}-\operatorname{coll}_p(I)
       +\omega_p\#(\mathcal T\cap I)\frac a{x_1}
       -\frac{fa(\log b)^2}{pm}.
\end{equation}

\paragraph{The marked operators and divisibility off the level.}
\teamanchor{proof.weighted-local-hecke-chain}%
Fix a coefficient embedding \(j_p:\mathbf C\hookrightarrow\mathbf C_p\), a
twist of level norm prime to \(p\), and the good datum of the same
seesaw. Use \(M_p\) for the module of columns integral at \(j_p\).
For \(p\ne N\), use the
\hyperlink{proof.weighted-local-source-gauge}{source gauge
calculation} on the translations allowed in \(\Psi_\vartheta\).
It gives the factor \(|j_p(1/d)|_p\) per box, which is one
when \(p\ne d\).
The final restriction preserves the integral gauge off \(N\).
the branching matrix is integral and its extra scalar is a unit.
This proves the required bound at primes outside
\(\{N\}\cup\{p<x_1:\chi_d(p)\ne1\}\), since the fixed \(d\) lies
below \(x_1\) for large \(x\).
\footnote{\weightedpin{grh-compact/src/New/Columns/ChainValuation.lean}
{307-L331}{Standardness at the allowable places}
and \weightedpin{grh-compact/src/New/Curve/ConjugatesOfJets.lean}
{129-L143}{Standard-data bound}.}

At a small non-split prime, using \(I=\varnothing\) in
\eqref{eq:weighted-A-jet-count} and the plain class sum gives
\[
 v_{j_p}\big(B^\mu_{\vartheta,\Psi}(\phi)\big)
 \ \ge\ \gamma_p-\ell_p-m\operatorname{coll}_p(\varnothing)
             \qquad(\phi\in M_p\cap\mathcal C_{\rm ram}).
\]
Here the valuation of a form denotes the lower bound for all its
integral-gauge coordinates at every argument.
At an unserved prime this is at least \(\gamma_p-2\ell_p\),
with no index cost. At any small non-split prime the trivial
scalar correction \(p^{s_p}\), where
\(s_p=\lceil2mfa/p\rceil\), gives valuation at least
\(\gamma_p-\ell_p\) on the image of \(p^{s_p}M_p\), by the
geometric-series bound for unrestricted collisions.
The estimates before restricting to \(Y_m\) apply to every
integral column, whether or not it is ramified. This lets us
add the two images below before restricting.

For \(p\) served, use the marked set chosen above and put
\[
 A_\lambda=a\#\{\tau\in\mathcal T:\lambda(\tau)=\lambda\},\qquad
 \Delta_p=c_p f(n_m+c_p),\qquad
 Y_\lambda=p^{A_\lambda}T_{\lambda,c_p}.
\]
Here \(T_{\lambda,i}\) is the spherical operator for the double coset
of \(\operatorname{diag}(pI_i,I)\) at \(\lambda\).
It acts at a split place of \(E\) over \(p\), away from the
chosen and level places. The marked place is spherical for each
stage kernel because its vacuum and characters are unramified,
and it lies away from the twist primes.

\paragraph{The split Hecke polynomial on the spherical kernel.}
\teamanchor{proof.weighted-local-hecke-vacuum}%
The \hyperlink{proof.principal-core-span}{spherical string in the earlier
chain} can be read on the whole spherical vector space by a
finite-field calculation. Let \(w=\lambda\) be the marked split
place with residue order \(q=q_\lambda\), and let \(n<n+k\) be the
source and target ranks. Give each maximal compact subgroup volume
one. For rank \(v\), write
\[
  P_{v,w}(X)=\sum_{i=0}^v(-X)^i q^{\binom i2}T_{w,i}.
\]
The double coset for \(T_{w,i}\) consists of the minuscule
modifications of index \(q^i\). A representative \(g_U\) can be
chosen integral with valuation of determinant \(i\) and reduction
kernel \(U\), as \(U\) runs through the \(i\)-planes of the
residue vector space. To see the coset parametrization, lift a basis
adapted to \(U\) and multiply the \(U\)-coordinates by a uniformizer.
The lattice \(g_U^{-1}\mathcal O_w^v\) is the inverse image of \(U\)
in \(\varpi_w^{-1}\mathcal O_w^v/\mathcal O_w^v\).
It determines the left compact coset, and \(g_U^{\mathsf t}\)
gives the corresponding right compact coset.

Use the split matrix realization of the same local Weil kernel,
with matrix variable \(A\) of size \((n+k)\)-by-\(n\).
The vacuum has integral support in this matrix variable.
The action on the target replaces \(A^{\mathsf t}\) by
\(A^{\mathsf t}g_U^{-1}\) and multiplies the integral-box
indicator by \(\alpha^{\dim U}\). On the source at the
conjugate place it replaces \(A\) by \(Ag_U^{-1}\) and
multiplies by \(\zeta^{\dim U}\).
Here \(\alpha=q^{n/2}\chi_S(\varpi_w)\) and
\(\zeta=q^{(n+k)/2}\chi_T(\varpi_w)\).
These scalars are the splitting characters times
\(|\det g_U|^{-n/2}\) and \(|\det g_U|^{-(n+k)/2}\),
respectively. In each case the translated indicator is supported on
integral \(A\) with \(U\) in the kernel of the corresponding residue
matrix. This follows by testing the linear map on
\(g_U^{-1}\mathcal O_w^v\).

For a residue matrix of rank \(r\), summing the two Hecke polynomials
therefore gives \(\prod_{i<n+k-r}(1-q^i\alpha X)\) on the target
and \(\prod_{i<n-r}(1-q^i\zeta X)\) on the conjugate source.
Indeed the coefficient of \((-Y)^i\) in
\(\prod_{a<D}(1-q^aY)\) is \(q^{\binom i2}\) times the number of
\(i\)-planes in a \(D\)-space. The identity follows by fixing a
hyperplane and counting \(i\)-planes lying in it and those meeting
it in an \((i-1)\)-plane, with \(q^{D-i}\) extensions for each
such \((i-1)\)-plane. Factoring the first \(k\) terms separates
the two polynomials at each residue matrix. In the theta pairing,
change variables by each source representative to move its
inverse action on the kernel to the source vector. The source
double coset at the conjugate place is the inverse double coset
at \(w\), since the two actions on the split unitary group are
inverse transpose. In the resulting paired terms \(g_U\) can be
replaced by \(g_U^{\mathsf t}\): they lie in the same double coset,
the source vector is fixed by the compact subgroup on the right,
and the lift is invariant under its action on the left by the
source vacuum. The resulting equality for the actual stage lift
\(L\) and every spherical vector \(u\) is
\[
 P_{n+k,w}(X)Lu
 =\prod_{i<k}(1-q^i\alpha X)\,
       L\big(P_{n,w}(q^k\alpha\zeta^{-1}X)u\big).
\]
The datum and the lifted vector are fixed by the relevant compact
subgroups because the characters are unramified there.
\footnote{\weightedpin{grh-compact/src/Lib/Weil/Unramified/HeckeTransfer.lean}
{58-L115}{Split matrix action and the same Hecke normalization}.}

In the marked chain the rank twists give the same base-change
character \(\chi_{e+1}\) on the two sides of arrow \(e\).
Its value at a uniformizer has \(p\)-adic valuation zero.
For this assertion, raise the prime to a power principal in the ray
class group for the character conductor. The character identity on
its generator computes the valuation, divided by that power, as
the signed sum of its weight over the two halves of the marked
pair. The bijection by \(u_p\) preserves the weights and gives zero.
Finite-order factors also have valuation zero. It follows that
\(q^k\alpha/\zeta=q^{k/2}\).
For the chain \(C_{j+1}\) of the first \(j+1\) lifts use the
variable \(X=-Uq^{(n_{j+1}-n_e)/2}\) at stage \(e\).
The coefficient of \(U^i\) at that stage is
\[
  y_e(i)T_{\lambda,i},\qquad
  y_e(i)=q^{(n_{j+1}-n_e)i/2+\binom i2}.
\]
Iterating the polynomial equality gives
\(P_{n_{j+1},w}(-U)C_{j+1}
 =G(U)C_{j+1}P_{n_0,w}(-Uq^{(n_{j+1}-n_0)/2})\).
The product \(G\) contains the factors
\[
  1+Uq^{i+(n_e+n_{j+1}-n_{e+1})/2}\chi_{e+1}(\varpi_w),
  \qquad 0\le e\le j,\quad 0\le i<n_{e+1}-n_e .
\]
It has integral coefficients and constant term one. The recurrence
for its reciprocal gives integral coefficients. Extracting degree
\(c_p\) from \(G(U)^{-1}P_{n_{j+1},w}(-U)\) transfers the
operator \(y_0(c_p)T_{\lambda,c_p}\) across the chain.
With \(q=p^{2h_\lambda}\), one has \(h_\lambda\le f\) by
\([E:\mathbf Q]=2f\), and
\[
 v_p(y_0(c_p))=
    h_\lambda\bigl((n_{j+1}-n_0)c_p+c_p(c_p-1)\bigr)
 \le\Delta_p .
\]
Multiplying the extracted identity by \(p^{\Delta_p}/y_0(c_p)\)
keeps the coefficients integral. They depend on the ranks and
the chain characters at \(\lambda\), which agree for every term
of \(\Psi_\vartheta\).

Before moving the operators across the chain, consider a Hecke
translate of a stage datum at each \(\lambda\in\Lambda\), of index
at most \(c_p\). Its diagonal representative has \(p\) in a set
\(J_\lambda\) of at most \(c_p\) coordinates and \(1\) in the
others. Right and left compact factors preserve the gauges.
In the transported gauge the two corresponding powers are
reciprocal at paired embeddings. For \(\sigma\in I\), charge each
height-one box whose index is outside \(J_{\lambda(\sigma)}\).
At the conjugate half charge every box with index in the same
\(J_{\lambda(\sigma)}\). These charges cover each exempted
coupling in the jet estimate since live partners induce
conjugate primes. At the marked half, the diagonal change of
gauge pays one power of \(p\) for each box in \(J_\lambda\).
For height one this fills the complement of the charged boxes,
costing a total of \(A_\lambda\) at that prime. For the other
columns it costs at most \(c_p\) per column because a tableau
column has distinct entries. At the conjugate half the reciprocal
gauge cancels the charges in \(J_\lambda\). The determinant
exponents cancel between the two halves by the same balancing
of stage weights. Thus a single step increases the gauge by
at most
\[
     p^{\,\sum_{\lambda\in\Lambda}A_\lambda+
              c_p\sum_{\sigma\in I}r_{j,\sigma}-J_{p,j}(I)}.
\]
This uses the local covariance of the normalized Schur kernel in
the transported Weyl gauge.
\footnote{\weightedpin{grh-compact/src/New/Curve/ConjugatesOfJets.lean}
{114-L136}{The gauge with diagonal powers at the target}
and \weightedpin{grh-compact/src/New/Columns/LocalStatement.lean}
{100-L143}{The choices of the charges and the diagonal action}.}

The coefficient extraction permits the operator \(Y_\lambda\)
to be moved from the initial column across the first \(j+1\)
lifts to a modification of the \(j\)-th datum. In that modification
replace the \(\lambda\)-translate by a finite linear combination
of translates of index at most \(c_p\) and multiply by
\(p^{A_\lambda-\Delta_p}\). Its coefficients are integral.
Thus the gauge cost \(A_\lambda\) in the last display is replaced
by \(\Delta_p\). For distinct marked primes, the operators commute
with each other and the modifications off their own places.
Their target Hecke operators commute with the level actions on
the vacuum data, so the modified joint sum still has the plain
class-sum estimate. Equivariance of the lift moves the modified
target action from the kernel into the target form. With
\(Y=\prod_{\lambda\in\Lambda}Y_\lambda\), induction on the chain
length now shows that applying all the modifications, one per
stage, gives the original chain applied to \(Y^m\phi\).
The resulting total exponent of the gauge is at most
\[
        \sum_{j<m}\left(
          \#\Lambda\,\Delta_p+c_p\sum_{\sigma\in I}r_{j,\sigma}
                                        -J_{p,j}(I)\right).
\]

To check its sign, all \(\sigma\in I\) are thin.
If \(p<m\), then \(c_p(a/x_1+x^{60})=O(a/(pm))\)
uniformly because \(x_1\) dominates every fixed power of \(x\).
If \(p\ge m\), then \(c_p=1\) and these \(\sigma\) are tops.
In both cases \(\#\Lambda\,\Delta_p=O(fa/(pm))\) uniformly, by
the polynomial bounds on \(f,n_m,c_p\) and \(x_1=b^{o(1)}\).
Therefore
\[
 \#\Lambda\,\Delta_p+c_p\sum_{\sigma\in I}r_{j,\sigma}
 \ \le\ \frac{\ell_p}{m}+\#(\mathcal T\cap I)\frac{a}{x_1}
\]
for all large \(x\).
Here \(\omega_p=1\) and \(\operatorname{coll}_p(I)=0\).
Substituting \eqref{eq:weighted-A-jet-count} proves that
the output on \(Y^m M_p\) has valuation at least
\(\gamma_p-2\ell_p\).

\paragraph{Recombination and the served-prime index.}
\teamanchor{proof.weighted-local-index-dispatch}%
We use the following elementary determinant rules, valid over
any valuation ring. A recombination of a tuple is a matrix
with entries in that ring acting on the tuple, with cost
\(v_p(\det)\). If \(M\) is a finite free module and \(A\)
is an integral square matrix invertible over the fraction field,
any tuple in \(M\) can be recombined into \(AM\) at cost at most
\(v_p(\det A)\). To see this, put the tuple in the columns of
\(X\). Perform elementary row and column operations over the
valuation ring to diagonalize \(A^{-1}X\), after initially
multiplying by a common denominator. Such diagonalization
follows by moving an entry of least finite valuation to the
first position, using it to clear its column and row, and
inducting. Scale the columns at each negative-valuation
diagonal entry by its inverse. The sum of the negative
valuations in absolute value equals the negative valuation of
the corresponding maximal minor in the diagonal form. All
minors of the original matrix have valuation at least
\(-v_p(\det A)\), by complementary minors for \(A^{-1}\)
and Cauchy--Binet. The same holds after integral row and column
operations. This proves the cost bound.
The argument uses only finite matrices, including when
\(\mathcal O_{j_p}\) is nondiscrete.

If \(X\) and \(Y\) preserve \(M\), costs for recombining into
\(XM+p^s M\) and \(YM+p^s M\) add for \(XYM+p^s M\).
Indeed first recombine into \(x_i=Xu_i+p^s v_i\).
Next perform the second recombination on the tuple \(u_i\)
and apply that same matrix to the tuple \(x_i\).
The new error is in \(p^s M\), and the determinants multiply.
Similarly one can pull back a recombination by evaluating a
linear map on a tuple and using its recombination matrix on
the original tuple.

\paragraph{The marked operator in the actual column coordinates.}
\teamanchor{proof.weighted-local-marked-edges}%
Fix a served \(p\), a coefficient embedding \(j_p\), and its
marked set \(\Lambda\). Use \(k=\kappa\), \(r_p=c_p\),
and \(C=\dim_{\mathbf C}\mathcal C\). For \(\lambda\in\Lambda\),
let \(L_\lambda\) be the length of the top cycle in
\(\{\lambda,\bar\lambda\}\). This length is the order of \(u_p^{o_p}\),
where \(u_p=[-p]\), so \(\operatorname{ord}(u_p)=o_p L_\lambda\).
It is even, since \(m\) is a power of two, \(m\mid
\operatorname{ord}(u_p)\) and \(o_p\) is odd. Write
\(\tau_i=u_p^{o_pi}\tau_0\), with \(\tau_0\) inducing \(\lambda\).
Read \(W_\lambda=\mathcal O_{E,\lambda}\) through
\[
 g_i=j_p\Phi_{\tau_i}\quad (i\text{ even}),\qquad
 g_i=j_p\Phi_{\tau_i}\operatorname{conj}\quad (i\text{ odd}).
\]
The full arithmetic Frobenius at this inert prime is
\(\Phi_{u_p}\circ\operatorname{conj}\) in these coordinates.
It acts by \([p]\) on the cyclotomic factor and by conjugation
on \(K\). Since \(o_p\) is odd, the step from \(g_i\) to
\(g_{i+1}\) reads its \(o_p\)-th power.
Put \(P=p^{o_p}\) and \(q_\lambda=P^{L_\lambda}\).
The embeddings satisfy \(g_{i+1}(a)\equiv g_i(a)^P\bmod p\)
for \(a\in W_\lambda\). This is congruence modulo \(p\) in
\(\mathcal O_{j_p}\), since \(p\) is unramified in \(E\) and
the congruence before the embedding is in \(p W_\lambda\).
Also \(\#(W_\lambda/p W_\lambda)=q_\lambda\).
Indeed \(4\mid\operatorname{ord}(u_p)\) here, so the order
of \(\Phi_{u_p}\circ\operatorname{conj}\) is
\(\operatorname{ord}(u_p)\). The two marked halves have
\(L_\lambda/2\) letters each, so \(A_\lambda=aL_\lambda/2\).
\footnote{\weightedpin{grh-compact/src/New/Columns/MarkedPrimeType.lean}
{19-L55}{The cycle and the congruence modulo \(p\) on the marked ring}, and
\weightedpin{grh/src/New/UpperBound/Divisibility/MarkOperator.lean}
{277-L332}{The residue field and alternating embeddings}.}

Let \(X\) be the definite class set at the stage-zero level and
\(K_\lambda=\mathrm{GL}_k(W_\lambda)\). The group
\(K_\lambda\) may be used to choose representatives of the finite
parahoric cosets in the completed local group. Put
\[
 d_\lambda=\operatorname{diag}(pI_{r_p},I),\qquad
 \mathscr P=K_\lambda\cap d_\lambda K_\lambda d_\lambda^{-1}.
\]
Choose the class representatives \(\rho_x\) integral above \(p\),
using rational unitary changes as in the source-gauge calculation.
Their components at \(\lambda,\bar\lambda\) can be made the
identity by a right level element supported on this split pair.
They give \(M_p\) by arbitrary integral Weyl values on \(X\),
since the class stabilizers are trivial and the level is integral
above \(p\). The level contains the supported maximal compact
at \(\lambda,\bar\lambda\), as \(p\) is away from the chosen
and level primes. Haar mass one gives the Hecke action as a sum
over the right compact cosets of \(K_\lambda d_\lambda K_\lambda\),
indexed by \(b\in K_\lambda/\mathscr P\).

An edge \(e=(x,b)\) replaces the \(\lambda\)-component of \(\rho_x\)
by \(b d_\lambda\), with conjugate inverse transpose at
\(\bar\lambda\) and no change elsewhere. Let this new adele be
\(\gamma_e\rho_{o(e)}u_e\), with \(\gamma_e\) rational unitary and
\(u_e\) in the stage-zero level. At \(\lambda\) this gives
\(\gamma_e^{-1}=h_e d_\lambda^{-1}b^{-1}\) for
\(h_e\in K_\lambda\). At the other places over \(p\) outside
the paired places, the rational part is integral invertible.
Both edge parametrizations used below are bijective. One is the
original \((x,b)\). For the other, suppose two edges have the
same \(o(e)\) and the same coset
\((h_e^{\mathsf t})^{-1}\mathscr P\).
Then \(d_\lambda h_e^{-1}h_{e'}d_\lambda^{-1}\) is integral
invertible, by the definition of \(\mathscr P\) and transposition.
Hence \(\gamma_e\gamma_{e'}^{-1}\) is integral invertible at
the split pair. Away from the pair its comparison between
\(\rho_x\) and \(\rho_{x'}\) is in the level because the origins
are equal. The supported maximal compact puts the comparison
in the level everywhere. Thus \(x=x'\), and freeness gives
\(\gamma_e\gamma_{e'}^{-1}=1\). The equation at the split pair
then gives \(b^{-1}b'\in\mathscr P\). This proves injectivity
and the two finite sets have equal size.

Let \(J\) index the degree-\(a\) monomials in \(k\) variables in
each letter, and let \(H\) index the Weyl tableaux outside the
cycle. The integral stage-zero Weyl basis consists of these
\(J\times H\) indices. For the letters define the polynomial
substitution action \(\rho\) by the symmetric powers of \(g_i(h)\)
at even \(i\) and of \(g_i(h)^{-\mathsf t}\) at odd \(i\).
In the monomial substitution matrix the entry at
\((\alpha,\beta)\) is the coefficient of \(z^\alpha\) in the
substitution of \(z^\beta\) by the columns of the base matrix.
Thus the coefficient action on the Weyl evaluation vector is
the transpose of the substitution matrix for the transposed
base matrix. This is a bijection of integral tableau and monomial
indices, using the height-one columns, and requires no change
by reciprocal multinomial coefficients.

Let \(\mathrm{dep}(\beta)\) be the degree in the variables
outside the first \(r_p\) at even letters and in the first \(r_p\)
at odd letters. Set \(D_p=\operatorname{diag}(p^{\mathrm{dep}(\beta)})\).
The summand of \(Y_\lambda=p^{A_\lambda}T_{\lambda,c_p}\)
along \(e\), on the cycle and the other tableaux respectively, is
\[
                 \bigl(\rho(h_e)D_p\rho(b^{-1})\bigr)^{\mathsf t}
                              \otimes U_e,\qquad
                 U_e\in \mathrm{GL}_{|H|}(\mathcal O_{j_p}).
\]
Indeed at an even letter, transpose followed by conjugation of
\(\gamma_e\) gives \(\gamma_e^{-1}\). The \(p\) per box changes
\(d_\lambda^{-1}\) to \(p d_\lambda^{-1}\).
At an odd letter the transpose reads
\(\gamma_e^{\mathsf t}\) through the redirected embedding and keeps
\(d_\lambda\). The off-cycle factors are integral invertible
because their top readings see only places outside the pair,
and the other stage-zero factors are trivial.
The formula proves that \(Y_\lambda M_p\subseteq M_p\) and
identifies the marked operator on all the columns.
\footnote{\weightedpin{grh-compact/src/New/Columns/TwistedHeckeType.lean}
{153-L239}{Evaluation at the actual classes and coefficients},
\weightedpin{grh/src/New/UpperBound/Divisibility/HeckeEdges.lean}
{403-L449}{The reverse edge parametrization}, and
\weightedpin{grh/src/New/UpperBound/Divisibility/HeckeEdges.lean}
{769-L864}{Operator, Weyl indexing and the source lattice}.}

\paragraph{The truncated colength of the marked operator.}
\teamanchor{proof.weighted-local-marked-colength}%
The colength calculation uses the
\hyperlink{proof.weighted-local-depth-sum}{residue depth sum} and
the adjoint pairing below. Put
\[
 \mathfrak c=L_\lambda a\left(P^{-2r_p}+\frac{r_p^2}{kP}\right),
 \qquad c_0=k^2 o_p L_\lambda,\qquad A_0'=L_\lambda c_0 .
\]
Write \(\mathcal O=\mathcal O_{j_p}\) in this paragraph, choose
\(w\in\mathcal O\) with \(w^2=p\), and put
\(\Delta=\operatorname{diag}(w^{\mathrm{dep}(\beta)})\).
Let \(M_*\) be the integral space of edge coefficients \(v\) for
which \(\Delta\rho(g)\Delta^{-1}v\) is integral for every
\(g\in K_\lambda\). Over the residue field, let \(\mathcal D_j\)
be the space of polynomials in the specified degrees all of whose
\(\rho(g)\)-translates have zero coefficients in depths below \(j\).

\paragraph{The depth sum for the marked residue cycle.}
\teamanchor{proof.weighted-local-depth-sum}%
For the served cycle \(0<2r_p\le k\). Use the residue field \(\mathbb F\) of \(\mathcal O_{j_p}\) and
the marked residue field \(\mathbf F_{q_\lambda}\), with embeddings
\(\sigma_i=\overline{g_i}\) and \(\sigma_{i+1}=\sigma_i^P\).
Here and below the latter notation means raising each output to
the \(P\)-th power. Write \(R=r_p\) within this paragraph, and
\(\ell=L_\lambda\), so \(q_\lambda=P^\ell\) and \(\ell\) is even.
Use variable indices \(0\le b<k\) in each letter \(i\bmod\ell\).
Call \(v=(b,i)\) deep if \(b\ge R\) at even \(i\) or \(b<R\)
at odd \(i\). Give a deep variable the weight \(1/\kappa_v\), where
\[
                  \kappa_v=
       \begin{cases}P & (b<2R),\\ P^{2R} & (b\ge2R),\end{cases}
\]
and give every other variable weight zero. Denote the set of
degree-\(a\)-in-each-letter monomials by \(J\).
We prove that if \(0\ne F\in\mathcal D_j\), then \(F\) contains
a monomial of weight at least \(j\).

Replace all shallow variables by independent indeterminates
\(u_{b,i}\). The resulting polynomial \(f\) in the deep variables
is nonzero over \(L=\mathbb F(u_{b,i})\).
Let a matrix \(A(y)\) over \(\mathbf F_{q_\lambda}\) have arbitrary
entries \(y_{cb}\) in the positions
\[
       (c,c+R),\quad (c,b)\text{ with }2R\le b<k,
                                \qquad 0\le c<R,
\]
and zero entries elsewhere. Choose lifts with this same support
over \(W_\lambda\), so their squares vanish and \(g=I+A(y)\)
is invertible there. The corresponding action of \(g^{-1}\)
on a deep variable, after specializing the shallow variables,
is \(Z_v\mapsto Z_v-z_v(y)\) with
\[
 \begin{aligned}
 z_{b,i}(y)&=\sum_{c<R}u_{c,i}\sigma_i(y_{cb})
                       &&(i\text{ even},\ b\ge R),\\
 z_{c,i}(y)&=-\sum_{b\ge R}u_{b,i}\sigma_i(y_{cb})
                       &&(i\text{ odd},\ c<R).
 \end{aligned}
\]
Unspecified \(y_{cb}\) in the sums are zero.
Indeed the base matrix of \(g^{-1}\) in the substitution
is \(I-A(y)\) at even letters and \(I+A(y)^{\mathsf t}\)
at odd letters. All the shifted coefficients involve only
shallow variables. Since \(gF\) has depth at least \(j\) and
\(F=g^{-1}(gF)\), the polynomial \(f\) belongs to the \(j\)-th
power of the point ideal of \(z(y)\) in \(L[Z_v]\) for every
such \(y\).
\footnote{\weightedpin{grh/src/New/UpperBound/ElementaryDivisors/Depth.lean}
{398-L455}{The split unipotent translations}, and
\weightedpin{grh/src/New/UpperBound/ElementaryDivisors/Depth.lean}
{574-L605}{Vanishing at the translated points}.}

For each letter \(i\) and distant column \(b\ge2R\) use the
\(R\) integers \(0\le h<2R\) for which \(i-h\) is even.
Frobenius gives the simultaneous equations
\[
  z_{b,i-h}(y)^{P^h}
              =\sum_{c<R}u_{c,i-h}^{P^h}\sigma_i(y_{cb}).
\]
Their coefficient matrix is invertible over \(L\). A monomial
in its determinant has for each \(c\) exactly one factor
\(u_{c,i-h}^{P^h}\) and therefore determines \(h\) from its
exponent, even if some \(i-h\) coincide modulo \(\ell\).
Distinct permutations give distinct monomials. For the paired
position \((c,c+R)\) at even \(i\), divide \(z_{c+R,i}(y)\)
by \(u_{c,i}\). At odd \(i\), subtract the already recovered
distant entries from the equation for \(z_{c,i}(y)\) and divide
by \(-u_{c+R,i}\). In both cases the coefficient is nonzero.
It follows that every \(\sigma_i(y_{cb})\) is the evaluation at
\(z(y)\) of a polynomial with coefficients in \(L\) and all
monomial weights below one. The powers used for distant
columns have weight \(P^h/P^{2R}<1\), and for the paired
positions the degree-one terms have weight \(1/P\).
In particular the points \(z(y)\) are distinct, since the
entries of \(y\) can be read at any fixed \(i\).

For each deep variable \(v=(b,i)\), raising \(z_v(y)\)
to the power \(\kappa_v=P^h\) gives a linear combination of
the entries at letter \(i+h\), with coefficients in \(L\).
Therefore for each \(v\) there is a polynomial \(H_v\) all of
whose monomials have weight below one such that
\[
                 G_v=Z_v^{\kappa_v}-H_v(Z)
\]
vanishes at all these points. There are
\(q_\lambda^{R(k-2R+1)}\) points. The product of the
\(\kappa_v\) is exactly this number. In each even letter there
are \(R\) near and \(k-2R\) distant deep variables, and in
each odd letter there are \(R\) near and no distant deep
variables. Their product thus has exponent
\(\ell R(k-2R+1)\) in \(P\).
\footnote{\weightedpin{grh/src/New/UpperBound/ElementaryDivisors/Depth.lean}
{284-L357}{The determinants and paired-entry interpolation}, and
\weightedpin{grh/src/New/UpperBound/ElementaryDivisors/Depth.lean}
{493-L570}{The equations and distinct graph points}.}

For \(j>0\), impose the jets of order below \(j\) at all these
points. The Chinese remainder theorem gives a surjective
jet map, because the \(j\)-th powers of distinct point ideals
are comaximal. Reduce products \(Z^\alpha G^\gamma\)
by substituting \(Z_v^{\kappa_v}=G_v+H_v\) whenever
\(\alpha_v\ge\kappa_v\). In each expanded term the weight of
the part \(Z^\alpha\) strictly decreases, and
\(\lvert\gamma\rvert+\sum_v\alpha_v/\kappa_v\) never increases.
The reduction terminates because the weights are rational
with fixed denominators. Modulo the common jet kernel the
products with \(0\le\alpha_v<\kappa_v\) and
\(\lvert\gamma\rvert<j\) therefore span. Their number is
\(\prod_v\kappa_v\) times the number of monomials of degree
below \(j\) in the deep variables, exactly the dimension
of the jet target. Thus they are independent there.
If every monomial of \(f\) had weight below \(j\), reduction
would express \(f\) itself as a linear combination of those
products. Since its jets vanish, every coefficient in that
expression would be zero, which contradicts \(f\ne0\).
The case \(j=0\) is automatic. Specialization preserved the
weights of the deep parts, so the costly monomial in \(f\)
comes from one in \(F\).
\footnote{\weightedpin{grh/src/Lib/General/Algebra/MultivariatePolynomials.lean}
{182-L245}{Weighted division and interpolation of jets}.}

For each \(j\), projection to the coefficients of monomials of
weight at least \(j\) is therefore injective on \(\mathcal D_j\).
After summing these dimension bounds for \(1\le j\le M\),
each monomial is counted at most its weight. The mean exponent
in any variable is \(a/k\) by symmetry within its letter.
Thus
\[
 \sum_{j=1}^M\dim_{\mathbb F}\mathcal D_j
    \le |J|\frac{a}{k}
       \left(\frac{\ell R}{P}
                  +\frac{\ell(k-2R)}{2P^{2R}}\right)
    \le |J|\,\ell a\left(P^{-2R}+\frac{R^2}{kP}\right).
\]
This is the required \(\mathfrak c |J|\), on the residue
representation of the same integral columns.
\footnote{\weightedpin{grh/src/New/UpperBound/ElementaryDivisors/Depth.lean}
{608-L659}{The weighted-monomial step}, and
\weightedpin{grh/src/New/UpperBound/ElementaryDivisors/Depth.lean}
{752-L796}{The depth sum}.}

For the adjoint pairing take the diagonal matrix
\(B_{\beta\beta}=p^{A_0'}/\prod_i\binom a{\beta_i}\)
on \(\mathcal O^J\), where
\(\binom a{\beta_i}=a!/\prod_b\beta_{bi}!\).
It has integral complementary diagonal
\(B'\) with \(BB'=p^{A_0'}I\).
Indeed \(a<q_\lambda\) and Legendre's formula give
\(v_p(\binom a{\beta_i})\le k\log_p q_\lambda\le c_0\),
by bounding the carry at each power of \(p\) by \(k\).
The tensor symmetric pairing gives
\(B\rho(g^{\mathsf t})=\rho(g)^{\mathsf t}B\),
by realizing polynomials in each letter as symmetric tensors
with the induced bilinear form. The monomials in that form have
squared weights \(\binom a{\beta_i}^{-1}\).

\paragraph{The edge average at the served prime.}
\teamanchor{proof.weighted-local-edge-average}%
Use the marked Frobenius cycle, source gauge and monomial pairing of
the marked-colength calculation. For this paragraph put
\[
 \begin{gathered}
 L=L_\lambda,\quad n=a,\quad r=r_p,\quad \ell_0=L/2,\quad
 F=o_pL,\quad P=p^{o_p},\quad q=P^L=p^F,\\
 A_0'=Lc_0,\quad c_0=k^2F,\quad c=5A_0',\quad T_0=3kc .
 \end{gathered}
\]
Here \(\ell_0\) is a number of letters. The residue field of the
marked ring \(W_\lambda\) is \(\mathbf F_q\). The maps
\(\iota_s:W_\lambda\to\mathcal O=\mathcal O_{j_p}\) reduce to
embeddings of \(\mathbf F_q\) in the residue field \(\mathbf k\) of
\(\mathcal O\), and
\(\iota_{s+1}\equiv\iota_s^P\bmod p\mathcal O\).
The action \(\rho\) on the degree-\(n\) monomials in each letter
uses \(\iota_s(g)\) at even letters and
\(\iota_s(g)^{-\mathsf t}\) at odd letters.
Let \(G\) be the block diagonal subgroup of \(K_\lambda\) with blocks
of sizes \(r,k-r\).
We use the integral pairing
\(\langle v',v\rangle=v'^{\mathsf t}Bv\) of that calculation,
so \(\rho(g^{\mathsf t})\) is adjoint to \(\rho(g)\).
The depth \(\mathrm{dep}\), \(\Delta=\operatorname{diag}
(w^{\mathrm{dep}(\beta)})\), \(w^2=p\), and
\(M_*=\{v:\Delta\rho(g)\Delta^{-1}v\in\mathcal O^J
\text{ for all }g\in K_\lambda\}\) use the same marked reading.

For an adjacent root \((b,b+1)\) and a letter \(i\) define the linear
defect on real exponent vectors by
\[
 d_{i,b}(\beta)=\beta_{u_i,i}
   +\frac{\sum_s P^{[s-i]_L}\epsilon_s
        (\beta_{b+1,s}-\beta_{b,s})}{q-1},
 \qquad
 \epsilon_s=(-1)^s,\qquad
 u_i=\begin{cases}b+1&i\text{ even},\\ b&i\text{ odd}.\end{cases}
\]
The notation \([s-i]_L\) uses the representative in \([0,L-1]\).
Put \(d_i=d_{i,r-1}\). Define \(\mathscr C\) by requiring all
coefficients of \(\rho(g)v\) at \(d_i(\beta)>3T_0\), for all
\(g\in G\) and \(i\), to lie in \((w^c)\).
Define \(\mathscr D\) over \(\mathbf k\) by the same requirements
with zero as the coefficient ideal. We will construct an integral
matrix \(V\) and prove
\[
 \begin{gathered}
     M_*+w^c\mathcal O^J\subseteq\mathscr C,\qquad
     V\mathcal O^J\subseteq\mathscr C,\\
     \langle v',Vv\rangle\equiv p^{A_0'}\langle v',v\rangle
                   \pmod{w^c}\quad(v'\in M_*,\ v\in\mathcal O^J),\\
     w^{c-1}v\in\mathscr C\ \Longrightarrow\ \bar v\in\mathscr D,\qquad
     [K_\lambda:\mathscr P]\dim_{\mathbf k}\mathscr D\le(1+\theta_0)|J|.
 \end{gathered}
\]
These are the simultaneous inputs to the marked-edge block
argument at the selected \(s_p\).
The size bounds used below are
\[
 Lkn(\ell_0+2)+n+k<P^{\ell_0},\qquad
 2(k+1)\le T_0,\qquad 8LkT_0P^{2r}\le\theta_0 n .
\]
For the served primes, \(P^{\ell_0}=q^{1/2}\ge13^{m/2}\)
dominates \(a\) times the polynomial factors in \(x\).
Also \(T_0\) is polynomial in \(x\), \(P^{2r}=b^{o(1)}\),
\(\theta_0^{-1}=m^2 f\) is polynomial in \(x\), and \(a\ge xb\).
These give the displayed bounds uniformly using the same scale as
the marked-colength calculation.
\footnote{\weightedpin{grh/src/New/UpperBound/ElementaryDivisors/ElementaryDivisorTheorem.lean}
{617-L640}{Edge conditions at the integral and residue precisions}, and
\weightedpin{grh/src/New/UpperBound/ElementaryDivisors/ElementaryDivisorTheorem.lean}
{672-L725}{The full statement and its construction}.}

\paragraph{The root filtration on the marked columns.}
\teamanchor{proof.weighted-local-root-filtration}%
Use the same marked ring \(W_\lambda\), coefficient ring \(\mathcal O\),
\(P\), \(q=P^L\), and embeddings \(\iota_s\). We use that \(P\ge5\)
and that these embeddings satisfy the Frobenius congruences modulo
\(p\mathcal O\), where \(p=w^2\). For coefficients in \(\mathcal O\)
write \(\operatorname{ord}_w(c)\) for the greatest integer \(e\) for
which \(c\in(w^e)\), with order \(+\infty\) at zero. Equivalently
\(\operatorname{ord}_w(c)=\lfloor2v_p(c)\rfloor\) for \(c\ne0\).
For the residue-field argument take the same substitutions reduced
to \(\mathbf k\), set \(w=0\), and give each nonzero coefficient
order zero. For variable weights \(h(z)\), a polynomial has
level at least \(B\) if each coefficient at \(z^\alpha\) has order
at least \(B-h\cdot\alpha\). Sums preserve a common lower level
and products add the levels. In particular a substitution sending
each variable to a polynomial of at least its weight preserves
levels. If a second such substitution differs on each variable
by terms of strictly greater weight, its difference on a
polynomial of level \(B\) has level strictly greater than \(B\).
For this last assertion, expand the difference of each monomial
as a sum with at least one changed factor. All sums are finite.

Fix a root \((b,b+1)\) and letter \(i\). Denote the two variables
at \(s\) by \(X_s,Y_s\), in that order by increasing row at even
letters and by decreasing row at odd letters. Put
\(P_s=P^{[s-i]_L}\). For an exponent array \(\alpha\) write
\[
 A_\infty(\alpha)=q\sum_s P_s\alpha_{X_s},\qquad
 A_0(\alpha)=q\alpha_{Y_i}+\sum_{s\ne i}P_s\alpha_{Y_s}.
\]
The defect \(d_{i,b}(\alpha)\) of the marked edge is
\((A_0(\alpha)-A_\infty(\alpha)/q)/(q-1)\). Suppose the increments
\(\delta_s=h(Y_s)-h(X_s)\) are in \([1/2,1]\). Suppose also that every
translate of \(F\) by the \(\mathrm{SL}_2(W_\lambda)\) on this root
has level at least \(e+h\cdot\alpha\), where \(e\) is a nonnegative integer and
\(F_\alpha\notin(w^{e+1})\). We show \(d_{i,b}(\alpha)\le0\).

Put \(\zeta=\delta_i/q>0\) and \(\eta_s=-1\) when \(s\) has the
parity of \(i\), and \(+1\) otherwise. Make the substitution
\(\mathscr L\) over \(\mathcal O[t]\), as follows:
\[
 \begin{array}{ll}
 X_i\mapsto t^qX_i,\quad Y_i\mapsto X_i,&\\
 X_s\mapsto t^{qP_s}X_s-\eta_s t^{P_s}Y_s,\quad
 Y_s\mapsto Y_s-\eta_s X_s &(s\ne i).
 \end{array}
\]
The other variables are fixed. In \(\mathscr L F\) take the
coefficient of the monomial obtained from \(z^\alpha\) by merging
\(Y_i\) into \(X_i\); reduce this polynomial in \(t\) modulo
\(w^{e+1}\) and denote it by \(\Phi(t)\).

For the expansion at infinity give \(t\) weight \(-\zeta\)
and \(X_s\) weight \(h(X_s)+qP_s\zeta\), keeping the other
weights. The terms \(t^{qP_s}X_s\) and the unaltered variables,
with \(Y_i\mapsto X_i\), preserve the original weights. The
two cross terms at \(s\ne i\) have excesses respectively
\(\delta_s-P_s\zeta\ge1/2-1/P>0\) and
\(qP_s\zeta-\delta_s\ge P/2-1>0\).
At the merged monomial, the weight of the term with \(t\)-degree
\(j\) is \(h\cdot\alpha+\zeta(A_\infty(\alpha)-j)\), because the
new weights of \(X_i,Y_i\) agree. The level lower bound
\(e+h\cdot\alpha\) kills coefficients with \(j>A_\infty(\alpha)\)
modulo \(w^{e+1}\). At \(j=A_\infty(\alpha)\) the cross terms also
vanish modulo \(w^{e+1}\). The merged exponent and this \(t\)-degree
uniquely recover \(\alpha\) in the substitution keeping only
the weight-preserving terms. Indeed the only forgotten exponents
were the division between \(X_i\) and \(Y_i\), and the exponent
of \(X_i\) is recovered from the \(t\)-degree after subtracting
the already fixed contributions from \(s\ne i\).
The leading coefficient is therefore \(F_\alpha\bmod w^{e+1}\)
and \(\deg\Phi=A_\infty(\alpha)\).

For the expansion at a finite point take \(x\in W_\lambda\),
\(u=\iota_i(x)\) and \(u_s=\iota_s(x)\). Give \(t\) weight
\(\zeta\), keep the weights of \(X_s\), and subtract
\(q\zeta\) from \(h(Y_i)\) and \(P_s\zeta\) from \(h(Y_s)\) at
\(s\ne i\). In the root rows take
\(\left(\begin{smallmatrix}0&1\\-1&x'\end{smallmatrix}\right)\),
with \(x'=-x\) for even \(i\) and \(x'=x\) for odd \(i\).
It sends \(X_s\) to \(Y_s\) and \(Y_s\) to
\(X_s+\eta_s u_sY_s\), up to the sign \(-1\) on each input
variable with row index \(b\). This follows by taking the
columns of the matrix in the even letters and of the inverse
transpose in the odd letters. Apply \(\mathscr L\) after this
root action and shift \(t\) to \(t+u\). For \(X_s\) the result is
the signed \(\mathscr L(Y_s)\). For \(Y_s\) at \(s\ne i\) it is
the signed \(\mathscr L(X_s)\) plus the two signed errors
\[
 \big((t+u)^{qP_s}-t^{qP_s}-u_s\big)X_s,\qquad
 -\eta_s\big((t+u)^{P_s}-t^{P_s}-u_s\big)Y_s .
\]
At \(i\) only the first error occurs, with \(P_i=1\) and \(u_i=u\).
The two errors follow by expanding
\(\mathscr L(X_s)(t+u)+\eta_s u_s\mathscr L(Y_s)(t+u)\);
the constant terms cancel with these signs.
Every parenthesis has coefficients in \((p)\) by the Frobenius
congruences, including for the exponent \(qP_s\) since the cycle
returns after \(L\) steps.

In these finite-point weights \(\mathscr L(Y_s)\) has weight at
least \(h(X_s)\), since \(\delta_s-P_s\zeta>0\) for \(s\ne i\),
and \(\mathscr L(X_s)\) has weight at least \(h(Y_s)\) by
\(qP_s\zeta\ge\delta_s\) for \(s\ne i\). At \(i\) the two weights
are exact. Each error has coefficient order at least two,
while the finite-point \(t\)-weight is nonnegative. For the
\(X_s\)-error use \(h(X_s)+2\ge h(Y_s)\), and for the \(Y_s\)-error
use \(2\ge P_s\zeta\). Thus acting, substituting and shifting is
level-preserving. Apply this to the inverse root translate of
\(F\), whose level is at least \(e+h\cdot\alpha\) by hypothesis.
It gives that level for \((\mathscr L F)(t+u)\). At the merged
monomial the term of \(t\)-degree \(j\) weighs
\(h\cdot\alpha+\zeta(j-A_0(\alpha))\), since the adjusted
weights of \(X_i,Y_i\) again agree. The coefficients with
\(j<A_0(\alpha)\) vanish modulo \(w^{e+1}\). Hence
\((t-u)^{A_0(\alpha)}\) divides \(\Phi\).

Choose one \(x\) from each of the \(q\) residue classes of
\(W_\lambda\). Differences of distinct \(u\)'s are units and
remain units in the coefficient quotient. These powers of
linear polynomials are pairwise comaximal, so their product
divides \(\Phi\). The divisor is monic, so \(q A_0(\alpha)\le
\deg\Phi=A_\infty(\alpha)\), including over the possibly
nonreduced quotient. This proves the required extreme-term
inequality for the original exponent and either coefficient ring.

For the between-block root, the depth weight rises by one
from \(X_s\) to \(Y_s\). For each \(i\) replace it by
\(\mathrm{dep}-d_{i,r-1}/3\), with \(d_{i,r-1}\) here viewed
as a weight on individual variables. The rise of the defect
weight on this pair is \(1+2/(q-1)\) at \(i\) and
\(2P_s/(q-1)\) at \(s\ne i\). Both lie in \([0,3/2]\), so the
new increments lie in \([1/2,1]\). If all root translates have
original depth level at least \(u\), they also have new level
at least \(u\). Otherwise choose a term of least new level
below \(u\) over the root orbit. The minimum exists because
support degrees are bounded and finite coefficient orders are
nonnegative integers. Its defect is nonpositive by the
extreme-term inequality, so its original depth level is at
most its new level, a contradiction. In particular,
\(d_{i,r-1}(\beta)\ge3j\) forces the depth-scaled coefficient
into \((w^{u+j})\) for integers \(u,j\ge0\), with the same
integer exponent when \(\mathcal O\) is nondiscrete.
For a within-block root in the residue field, the signed
height rises by one from \(X_s\) to \(Y_s\). The same
extreme-term inequality with \(e=0\) applies to a term of
least signed height over the root orbit. These are the two
root inputs used in halfway divisibility and the Levi flag
selection below.
\footnote{\weightedpin{grh/src/New/UpperBound/ElementaryDivisors/RootLemma.lean}
{509-L589}{Auxiliary substitution and its expansion at infinity},
\weightedpin{grh/src/New/UpperBound/ElementaryDivisors/RootLemma.lean}
{623-L761}{The finite expansions}, and
\weightedpin{grh/src/New/UpperBound/ElementaryDivisors/RootLemma.lean}
{873-L1067}{The degree bound and both filtration conclusions}.}

\paragraph{Halfway divisibility.}
For the root \((r-1,r)\), the deep indicator rises by one in
the root direction at each letter. Use the
\hyperlink{proof.weighted-local-root-filtration}{root filtration}
on this marked action.%
\footnote{\weightedpin{grh/src/New/UpperBound/ElementaryDivisors/RootLemma.lean}
{984-L1049}{The depth shift under the root group}, and
\weightedpin{grh/src/New/UpperBound/ElementaryDivisors/Edges.lean}
{700-L739}{The divisibility statement for integral powers of \(w\)}.}

For \(v'\in M_*\) choose an integer \(u\) large enough that
\(H=w^u\Delta^{-1}v'\) is integral. For every \(g\in K_\lambda\),
the coefficients of \(\Delta\rho(g)H\) are divisible by \(w^u\).
Apply the estimate to \(\rho(g)H\) and all its root-group
translates. Since \(\Delta\) commutes with \(\rho(g)\) for \(g\in G\),
cancelling \(w^u\) gives
\[
 d_i(\beta)\ge 3c\quad\Longrightarrow\quad
                 (\rho(g)v')_\beta\in(w^c)\qquad(g\in G).
\]
In particular \(M_*\subseteq\mathscr C\). For \(g\in G\), the
adjoint relation and \(g^{\mathsf t}\in G\) give
\[
 d_i(\beta)> T_0\quad\Longrightarrow\quad
 \langle v',\rho(g)e_\beta\rangle
          =B_{\beta\beta}(\rho(g^{\mathsf t})v')_\beta\in(w^c).
\]
The first containment for \(w^c\mathcal O^J\) follows from the
definition of \(\mathscr C\). If \(w^{c-1}v\in\mathscr C\), then
cancelling \(w^{c-1}\) puts each tested coefficient of \(v\) in
\((w)\), which reduces to zero in \(\mathbf k\).
This gives the claimed implication to \(\mathscr D\).

\paragraph{An integral functional for one matrix.}
Let \(\xi=e_r-e_{r-1}\) on the row indices and put
\[
 D=Lkn(\ell_0+2)P^{\ell_0},\qquad H_0=2T_0P^{\ell_0}.
\]
Then \(D+k<q\). In the \(j\)-th copy of a generic matrix \(X\),
let \(g\in G\) act by \(X\mapsto X^g=\iota_j(g^{-1})X\iota_j(g)\).
The row and column weights of \(X^\alpha\) are
\(\sum_{b,d}\alpha_{bd}\xi_b\) and \(\sum_{b,d}\alpha_{bd}\xi_d\).
We construct \(I_j\) on polynomials of degree at most \(D\)
with the following properties. It is invariant, vanishes on
monomials of row weight \(>H_0\), and agrees with
\(E_j=p^{c_0}\operatorname{ev}_1\) modulo \((w^c)\) on a
polynomial whose entire \(G\)-orbit has those coefficients in
\((w^c)\).

Fix a homogeneous polynomial \(\psi\) of degree \(d\le D\).
The diagonal monomial pairing
on the entries of \(X\) has coefficient
\(\eta_\alpha=p^{c_0}/\binom{d}{\alpha}\) at \(X^\alpha\).
It is integral. For \(d<p^F\), Legendre's formula bounds each
floor-function contribution to the valuation of the multinomial
by \(k^2-1\), and the contributions at powers at least \(p^F\)
vanish. Units in the denominator are invertible in \(\mathcal O\).
After restricting \(X\) to the two diagonal blocks, form an
operator \(D_d(\psi)\) on \(U_d=(\mathcal O^k)^{\otimes d}\)
by pairing \(\psi\) with the entries of \((X^{\mathsf t})^{\otimes d}\).
Polynomials with an off-block variable contribute zero.
The diagonal pairing in degree \(d\) is proportional over the
fraction field to the factorial pairing
\((X^\alpha,X^\alpha)=\prod_{b,e}\alpha_{be}!\).
Applying a linear substitution to the first argument therefore
applies its transpose to the second, as follows by the
constant-coefficient differential operator pairing at zero.
For \(X\mapsto\iota_j(g^{-1})X\iota_j(g)\) the transpose substitution
is \(X\mapsto\iota_j(g)^{-\mathsf t}X\iota_j(g)^{\mathsf t}\).
Consequently, writing \(g_j=\iota_j(g)\), we have
\[
 D_d(\psi(X^g))=
          g_j^{\otimes d}D_d(\psi)(g_j^{-1})^{\otimes d},
 \qquad
 \operatorname{tr}_{U_d}D_d(\psi)=p^{c_0}\psi(1).
\]
For the trace, a diagonal monomial in \(X\) of degree \(d\)
occurs on \(\binom d\alpha\) diagonal tensor entries, giving the
stated scalar.

The coefficient matrices of the block-restricted tensor polynomial
\(X^{\otimes d}\) lie in the \(\mathcal O\)-span of
\(g_j^{\otimes d}\) for \(g\in G\). Indeed, evaluation vectors
of the degree-\(d\) block monomials at the invertible block matrices
over \(\mathbf F_q\) span over \(\mathbf k\). An annihilator would
give a polynomial vanishing on those invertible matrices.
Multiplying by the product of the two determinants gives a
polynomial vanishing on the entire \(\mathbf F_q\) grid, of degree
at most \(d+k<q\). Interpolation one variable at a time
makes that polynomial zero, and cancellation gives a zero
annihilator. Choosing a residue evaluation basis and lifting its
points to \(G\) yields a matrix with unit determinant over
\(\mathcal O\). Inverting this matrix expresses the tensor
coefficient matrices in the indicated span. The same span
contains the coefficients of \((X^{\mathsf t})^{\otimes d}\), by
transposing the matrices used for evaluation.

Let \(U_{d,+}\) be generated by the \(G\)-translates of those
basis tensors \(e_{\mathbf b}\) with
\(\sum_{\nu=1}^d \xi_{b_\nu}>H_0\).
The span result shows finite generation, since all translates
are linear combinations of the finitely many applications
of coefficient matrices to these tensors. It also shows that
\(D_d(\psi)\) preserves \(U_{d,+}\).
A finitely generated torsion-free module over a valuation
domain is free. In this case one can prove it by eliminating
from the generator matrix a nonzero entry dividing all its
entries, then continuing on the remaining block.
Use the trace difference
\[
 I_j(\psi)=\sum_{0\le d\le D}
   \left(\operatorname{tr}_{U_d}D_d(\psi_d)
                  -\operatorname{tr}_{U_{d,+}}D_d(\psi_d)\right),
\]
where \(\psi_d\) is the homogeneous part of \(\psi\).
Conjugation invariance follows from the operator conjugation
identity and stability of \(U_{d,+}\).
For a monomial of row weight \(>H_0\) the operator is zero
on every low-weight basis tensor. The transpose makes each
matrix input index of \((X^{\mathsf t})^{\otimes d}\) a row
index of the corresponding monomial in \(X\).
The nonzero columns map to \(U_{d,+}\) since their inputs
are among its generators and the operator preserves it.
The two traces are equal by trace cyclicity for the maps
through \(U_{d,+}\).

Suppose next that all conjugates of \(\psi\) have high-row
coefficients in \((w^c)\).
On a high-weight basis tensor \(D_d(\psi_d(X^g))\) uses only
the high-row part of \(\psi_d(X^g)\). Dividing that part by
\(w^c\) gives a polynomial whose operator preserves \(U_{d,+}\).
The conjugation identity, with \(g\) or \(g^{-1}\), now gives
\(D_d(\psi_d)U_{d,+}\subseteq w^c U_{d,+}\) on the generating
translates. This proves divisibility of the restricted trace
even if the lattice \(U_{d,+}\) is unsaturated.
It proves \(I_j(\psi)\equiv E_j(\psi)\pmod{w^c}\).
Finally both \(I_j\) and \(E_j\) kill monomials whose row and
column weights differ. Conjugation by the diagonal
\(\big((1+p)^{\xi_b}\big)_b\in G\) acts on the monomial by
\((1+p)\) to the power of the column minus the row weight.
No nonzero power of \(1+p\) equals \(1\) in characteristic zero.
\footnote{\weightedpin{grh/src/New/UpperBound/ElementaryDivisors/InvariantAverage.lean}
{223-L253}{The weight test under diagonal conjugation}, and
\weightedpin{grh/src/New/UpperBound/ElementaryDivisors/InvariantAverage.lean}
{919-L968}{Trace normalization and an integral invariant functional}.}

\paragraph{Contraction at the marked precision.}
Take \(L\) matrices \(X_s\) with independent entries. The
substitution matrix \(U(X)\) on the monomial indices is the tensor
product of the symmetric powers of the following letter matrices.
In letter \(s\) form the products of
\(X_{s-d}^{[P^d]}\), \(0\le d<m_s\), in descending order of \(d\),
where brackets denote entrywise powers. Take the row at \(b\)
from the product of length \(m_s=\ell_0+\mathrm{deep}(b,s)\).
For odd \(s\), take this row from the transposed adjugate of
the product of that length. Thus \(U(1,\ldots,1)=1\). Each
entry has degree at most \(D\) in each copy. A factor has
degree at most \(P^{\ell_0}\), there are at most \(\ell_0+1\)
factors in a row product and an adjugate entry multiplies at
most \(k-1\) entries. The symmetric powers and tensor product
then give the degree bound.
\footnote{\weightedpin{grh/src/New/UpperBound/ElementaryDivisors/FrobeniusFactors.lean}
{67-L96}{The row-length convention and transposed adjugates}.}

For \(g\in G\) write \(X^g=(X_j^g)_j\).
The matrix
\[
                    U_g(X)=\rho(g)\,U(X^g)\,\rho(g^{-1})
\]
has the same description by generic factors, with each factor
replaced by
\(\iota_s(g)(X_{s-d}^g)^{[P^d]}\iota_s(g^{-1})\).
For the verification, conjugate the full row-length products
and use that the two row-length projectors commute with the
block diagonal matrices. At odd letters the transposed
adjugate changes the conjugating matrix to its inverse
transpose. This follows from the polynomial adjugate identity,
so it holds for generic factors.

\teamanchor{proof.weighted-local-signed-frobenius}%
We prove the coefficient estimate for this \(U_g\). Fix a
copy \(j\) and use the row vector \(\xi=e_r-e_{r-1}\).
For an exponent \(\nu\) in the copies \(X_s\), write \(R(\nu_s)\)
and \(C(\nu_s)\) for its row and column weights against \(\xi\).
Assign the entry \((b,b')\) of \(X_u\) the weight
\[
 \omega_{u,b,b'}=\frac{\epsilon}{P^{\ell_0}}
 \begin{cases}
  \xi_b, & u=j,\\
  P^{[u-j]_L}(\xi_b-\xi_{b'})/q, & u\ne j,
 \end{cases}
 \qquad \epsilon\in\left\{\tfrac12,-\tfrac12\right\}.
\]
Use the coefficient order \(\operatorname{ord}_w\) above.
The weight of \(X^\nu\) is \(\epsilon W_j(\nu)\), where
\[
 W_j(\nu)=P^{-\ell_0}\left(R(\nu_j)
        +\sum_{u\ne j}\frac{P^{[u-j]_L}}{q}
                               (R(\nu_u)-C(\nu_u))\right).
\]

Consider the generic factor
\(\iota_s(g)(X_{s-d}^g)^{[P^d]}\iota_s(g)^{-1}\)
at letter \(s\), with \(0\le d\le\ell_0\). Entrywise
\(P^d\)-th power is a ring map modulo \(p\), and it carries
the readings of \(g\) and \(g^{-1}\) at \(s-d\) to those at \(s\).
The two conjugations therefore cancel modulo \(p\). The entry
\((b,b')\) is congruent to \(X_{s-d,b,b'}^{P^d}\) modulo \(p\)
as a polynomial, and all its monomials have degree \(P^d\).
Since every \(\lvert\omega_{u,b,b'}\rvert\le P^{-\ell_0}\),
their weights lie in \([-1,1]\). The error coefficients have
order at least two, so the entry has level at least
\(P^d\omega_{s-d,b,b'}\), the level of the pure power.

Put \(u_0=1/(2P^{\ell_0})\), and define the real numbers
\[
 \lambda_s(d,b)=\frac{\epsilon \xi_b}{P^{\ell_0}}\,
                             \frac{P^{d+[s-d-j]_L}}q .
\]
For any \(s,d\) these sum to zero over \(b\), and
\(\lvert\lambda_s(0,b)\rvert\le u_0\).
The difference
\(\lambda_s(d+1,b)-\lambda_s(d,b')\) agrees with
\(P^d\omega_{s-d,b,b'}\) unless \(s-d=j\). In fact if
\([s-d-j]_L>0\), decrementing that representative by one
leaves the exponent \(d+[s-d-j]_L\) unchanged. If \(s-d=j\),
advancing \(d\) replaces the representative zero by \(L-1\).
The difference then equals \(P^d\omega_{j,b,b'}\) minus
\(\epsilon\xi_{b'}P^d/(P^{\ell_0}q)\), whose absolute value
as an error is at most \(u_0\). Thus the factor entry has level
at least \(\lambda_s(d+1,b)-\lambda_s(d,b')-u_0\).

Use any fixed allowable length \(m_s\le\ell_0+1\) for a
letter product. Along each matrix product summand these
potential differences telescope. Every entry in row \(b\)
has level at least
\(\lambda_s(m_s,b)-(m_s+1)u_0\).
A transposed adjugate entry in row \(b\) is a signed sum of
products taking one entry from every row other than \(b\).
The sum of their potentials is \(-\lambda_s(m_s,b)\).
The lower level for a row chosen from either the ordinary
product or the transposed adjugate is consequently
\[
 (-1)^s\lambda_s(m_s,b)-k(\ell_0+2)u_0 .
\]
For odd letters this calculation uses the full product at
that fixed length before selecting the adjugate row.


For the prescribed length \(m_s=\ell_0+\mathrm{deep}(b,s)\),
the first term of this level is the variable weight in
\(\epsilon(1-1/q)d_{j+\ell_0}\). To see this, both are zero
off the two root rows. On the shallow row the exponent
contributes \(P^{[s-j-\ell_0]_L}/q\) with a negative sign.
On the deep row the power in the numerator is
\(P^{1+[s-j-\ell_0-1]_L}\): it equals
\(P^{[s-j-\ell_0]_L}\) off letter \(j+\ell_0\), and \(q\)
at that letter. This is exactly the extra term in the
defect formula.

In the degree-\(n\) symmetric power, each coefficient at
a row exponent uses the row levels with the multiplicities
of that exponent. This follows by expanding the substitution
of each input monomial as a product of linear forms in the
output variables. Multiplying over the letters gives
\[
 \operatorname{ord}_w\big([X^\nu](U_g)_{\beta\gamma}\big)
       +\epsilon W_j(\nu)
 \ \ge\
 \epsilon(1-1/q)d_{j+\ell_0}(\beta)
       -\frac{Lnk(\ell_0+2)}{2P^{\ell_0}}.
\]
The scale condition bounds the last quotient by \(1/2\).
If the coefficient survives modulo \(w^c\), its order is less
than \(c\). Use both signs of \(\epsilon\) to get
\(\lvert(1-1/q)d_{j+\ell_0}(\beta)-W_j(\nu)\rvert\le2c+1\).
There are \(n\) variables counted with multiplicity in each
letter, so \(\lvert d_{j+\ell_0}(\beta)\rvert\le2n\) by the
geometric sum in the defect formula. Since \(2n\le q\),
this proves
\[
 [X^\nu](U_g)_{\beta\gamma}\notin(w^c)
      \quad\Longrightarrow\quad
 \big|d_{j+\ell_0}(\beta)-W_j(\nu)\big|\le 2c+2\le T_0=3kc .
\]
All the coefficients here use the original monomials and the
alternating action in the marked edge calculation.
\footnote{\weightedpin{grh/src/New/UpperBound/ElementaryDivisors/FrobeniusFactors.lean}
{470-L569}{Conjugation modulo \(p\) and the coefficient level},
\weightedpin{grh/src/New/UpperBound/ElementaryDivisors/FrobeniusFactors.lean}
{269-L389}{Matrix products, cofactors and symmetric powers}, and
\weightedpin{grh/src/New/UpperBound/ElementaryDivisors/FrobeniusFactors.lean}
{622-L781}{Cyclic potentials and the two signed bounds}.}
The conjugated factors and the same signed bound are recorded in the
fixed proof.%
\footnote{\weightedpin{grh/src/New/UpperBound/ElementaryDivisors/FrobeniusFactors.lean}
{503-L509}{Conjugation on both sides of each entrywise power}, and
\weightedpin{grh/src/New/UpperBound/ElementaryDivisors/FrobeniusFactors.lean}
{660-L781}{Gauss estimate with the signed weights}.}

Contract each copy of \(U(X)\) by \(I_j\) to obtain \(V\).
For the range and pairing, choose a copy \(j\).
Contract the other copies by \(L_u\in\{I_u,E_u\}\) and denote
by \(F_{j,g}(X_j)\) the vector obtained from \(U_g(X)\rho(g)v\).
The invariance of the other \(L_u\) gives
\[
 F_{j,1}(X_j^g)=\rho(g^{-1})F_{j,g}(X_j).
\]
If \([X_j^\alpha]F_{j,g,\beta}\) survives modulo \(w^c\),
at least one contributing monomial of \(U_g\) survives
and has \(R(\nu_u)=C(\nu_u)\) for every contracted copy \(u\ne j\).
The coefficient estimate therefore reads
\[
          \big|d_{j+\ell_0}(\beta)-P^{-\ell_0}R(\alpha)\big|
                       \le T_0.
\]
With all other \(L_u=I_u\), the coefficient
\((\rho(g)Vv)_\beta\) is \(I_j(F_{j,g,\beta})\).
If \(d_{j+\ell_0}(\beta)>3T_0\), \(I_j\) kills all surviving
monomials, since their row weights exceed \(2T_0P^{\ell_0}=H_0\).
This proves \(V\mathcal O^J\subseteq\mathscr C\).

For the pairing put \(\psi=\langle v',F_{j,1}\rangle\).
In \(\psi(X_j^g)\), the coefficient at a monomial of row weight
\(>H_0\) is the sum over \(\beta\) of
\([X_j^\alpha]F_{j,g,\beta}\langle v',\rho(g^{-1})e_\beta\rangle\).
Each summand whose first factor survives has
\(d_{j+\ell_0}(\beta)>T_0\). Its second factor is divisible by
\(w^c\) by halfway divisibility. Thus
\(I_j(\psi)\equiv E_j(\psi)\pmod{w^c}\).
Replace the \(E_j\)'s by \(I_j\)'s one copy at a time. The
initial all-\(E_j\) contraction is \(p^{Lc_0}I\), which proves
the pairing congruence with the stated \(p^{A_0'}\).
\footnote{\weightedpin{grh/src/New/UpperBound/ElementaryDivisors/AveragingMap.lean}
{468-L564}{Range calculation and pairing after each replacement}.}

\paragraph{Packing the small-defect residue space.}
Work over \(\mathbf k\). Let \(M=q^j\) for a positive integer \(j\).
For each array \(u\) of residues in \([0,M-1]^{kL}\), take a form
\(f_u\in\mathscr D\).
Choose a nonnegative \(\eta(u)\equiv u\bmod M\) in each coordinate
by adding
\(M(k-\lfloor\sum_b u_{b,s}/M\rfloor)\) at row \(r-1\) in
letter \(s\). Its degree in that letter lies in
\([Mk,M(k+1))\). Form the polynomial
\[
                        f(z)=\sum_u z^{\eta(u)} f_u(z^M).
\]
Packing is injective by the disjoint exponent classes modulo \(M\).
The reduced substitutions commute with \(z\mapsto z^M\),
since the entries of the substitution matrices lie in the
embedded \(\mathbf F_q\). A nonzero term of any Levi translate of
\(f\) therefore has exponent \(\alpha=M\beta+\gamma\), where
\(\beta\) occurs in a Levi translate of some \(f_u\) and \(\gamma\)
is a nonnegative exponent in the translate of a tag
\(z^{\eta(u)}\). Those tags retain their degrees in each letter.
The defect formula gives
\(\lvert d_i(\gamma)\rvert\le2M(k+1)\), by summing the geometric
weights in the second term. Thus%
\footnote{\weightedpin{grh/src/New/UpperBound/ElementaryDivisors/DefectCount.lean}
{199-L274}{Packing with disjoint exponent residues}, and
\weightedpin{grh/src/New/UpperBound/ElementaryDivisors/DefectCount.lean}
{284-L366}{Degrees and defects after a Levi translate}.}
\[
 d_i(\alpha/M)\le3T_0+2(k+1)\le4T_0,\qquad
             n+k\le\sum_b\alpha_{b,s}/M\le n+k+1.
\]

For a nonzero \(f\), choose a nonzero term of least signed height
among all Levi translates. The variable \(z_{b,s}\) has height
\(b\) at even letters and \(-b\) at odd letters.
Let \(\mathscr I\) be the inverse image of the lower triangular
subgroup modulo \(p\).
We can find this coefficient in a representative of
\((G\cap\mathscr I)\backslash G\).
Indeed \(y\in G\cap\mathscr I\) substitutes a diagonal unit
times each variable plus variables of strictly higher height.
Writing the chosen translate as \(yx\) with \(x\) a coset
representative, the coefficient of least height can only come
from the same term in the translate by \(x\), with a nonzero
diagonal multiplier. In the translate by \(x\), for any adjacent
root inside a block, all translates by its root-\(\mathrm{SL}_2\)
have no coefficient at smaller height.
The signed height rises by one along the root at each letter.
The \hyperlink{proof.weighted-local-root-filtration}{residue root
estimate} for a term of least height gives
\(d_{i,b}(\alpha)\le0\) for every within-block root and all \(i\).
The between-block root already has \(d_i(\alpha/M)\le4T_0\).
\footnote{\weightedpin{grh/src/New/UpperBound/ElementaryDivisors/RootLemma.lean}
{1051-L1067}{The root estimate for a term of least height}, and
\weightedpin{grh/src/New/UpperBound/ElementaryDivisors/DefectCount.lean}
{455-L489}{Choice of a Levi flag representative}.}

Define the cyclic linear map
\[
               (A_1v)_{b,s}=\epsilon_s(Pv_{b,s+1}-v_{b,s})
\]
on real arrays. Substituting in the defect formula gives
\(d_{i,b}(A_1v)=(A_1v)_{u_i,i}+v_{b+1,i}-v_{b,i}\).
Indeed the geometric sum of successive differences telescopes
with the wraparound factor \(q-1\).
Let \(\mu_b=P^{2(r-1-b)}\) for \(b<r\), and \(\mu_b=0\)
for \(b\ge r\).
Apply \(A_1\) to the array which equals \(\mu\) at even letters
and \(P\mu\) at odd letters. The resulting vector \(\sigma\) is
\((P^2-1)\mu\) at even letters and zero at odd letters,
so is nonnegative integral with degrees at most \(P^{2r}-1\)
in each letter. The displayed defect identity gives
\(d_{i,b}(\sigma)\le0\) for every root. For \(b=r-1\) it gives
\(-1\) at even letters and \(-P\) at odd letters.
For \(b<r-1\), the even defect is zero since
\(\mu_b=P^2\mu_{b+1}\), and the odd defect is nonpositive.
Both defects vanish when \(b\ge r\).

Put
\[
 \begin{split}
   D_s&=n+k+1+4T_0\sum_b\sigma_{b,s},\qquad
   E_0=4T_0P^{2r},\\
   \Omega&=\{\beta\ge0:d_{i,b}(\beta)\le0\text{ for all }i,b\},\\
   \mathcal B&=\{\beta\ge0:D_s-1\le\sum_b\beta_{b,s}\le D_s
                                    \text{ for all }s\}.
 \end{split}
\]
All \(D_s\le n+E_0\), and \(\alpha/M+4T_0\sigma\)
belongs to \(\Omega\cap\mathcal B\).
For each of the chosen Levi flag representatives and each
point of this intersection on \(M^{-1}\mathbf Z^{kL}\),
record the coefficient of the translated packed polynomial
at the exponent \(M(\beta-4T_0\sigma)\). Use zero if this
array has a negative coordinate. The least-height argument
proves injectivity of this linear coefficient map.
The domain has dimension \(M^{kL}\dim_{\mathbf k}\mathscr D\).
As \(M=q^j\to\infty\), the grid count divided by \(M^{kL}\)
converges to the volume, since the region is a bounded
polytope and its boundary has measure zero. Consequently%
\footnote{\weightedpin{grh/src/New/UpperBound/ElementaryDivisors/DefectCount.lean}
{497-L581}{The defect telescope and the shift}, and
\weightedpin{grh/src/New/UpperBound/ElementaryDivisors/DefectCount.lean}
{694-L801}{Coefficient sampling and the volume limit}.}
\[
        \dim_{\mathbf k}\mathscr D
        \le [G:G\cap\mathscr I]\operatorname{vol}
                                      (\Omega\cap\mathcal B).
\]

\paragraph{The flag factor in the volume.}
\teamanchor{proof.weighted-local-defect-volume}%
Let
\(\mathcal F_{k,q}=\prod_{j=1}^k\frac{q^j-1}{q-1}\).
We show \(\mathcal F_{k,q}\operatorname{vol}(\Omega\cap\mathcal B)
=\operatorname{vol}(\mathcal B)\).
From the defect identity, \(A_1v\in\Omega\) is equivalent
to interlacing at every link. At even \(s\), the arrays
\(x_b=Pv_{b,s+1}\) and \(y_b=v_{b,s}\) satisfy
\(x_b\ge y_b\) and \(y_b\ge x_{b+1}\) for consecutive indices.
At odd \(s\) use \(x_b=v_{b,s}\) and \(y_b=Pv_{b,s+1}\)
with the same inequalities. Either link forces its two
arrays to be in decreasing order.

Away from the finitely many tie hyperplanes, define
\(M_s(b,d)=\mathbf 1_{\epsilon_s(Pv_{b,s+1}-v_{d,s})>0}\).
For decreasing arrays \(x,y\), the matrix with entries
\(\mathbf 1_{x_b>y_d}\) has determinant one exactly when
\(x_b>y_b>x_{b+1}\), with only \(x_b>y_b\) at the last row,
and determinant zero otherwise. To see this, each row is a
suffix of ones and the suffixes are nested. A nonzero
determinant forces their \(k\) lengths to be distinct and
nonzero, hence \(k,k-1,\ldots,1\), and then the matrix is
upper triangular with diagonal one. This applies by
transposition to an odd link.
When the rows of \(v\) are independently reordered in
each letter, the sign from that reordering occurs once
as a row sign and once as a column sign in \(\prod_s\det M_s\),
so cancels. After sorting, the product is the indicator
that all interlacings hold.
There is precisely one ordering giving interlacing when
that indicator is one. Therefore \(\prod_s\det M_s\)
equals the sum of the indicators that \(A_1v\) satisfies
the interlacing after the independent letter reorderings.

For a tuple of permutations \(\tau=(\tau_s)_s\), set
\[
             (A_\tau v)_{b,s}
                 =\epsilon_s(Pv_{b,s+1}-v_{\tau_s(b),s}).
\]
Expanding the determinant product expresses it as the signed
sum of \(\mathbf 1_{A_\tau v>0}\) over all \(\tau\), with
sign \(\prod_s\operatorname{sgn}(\tau_s)\).
The bounds on the sums of the letters are the same for
every \(\tau\) and for every reordering of the rows within
letters. Multiply by the indicator of these bounds and
integrate. The interlacing interpretation gives
\((k!)^L\operatorname{vol}(\Omega\cap\mathcal B)/|\det A_1|\),
whereas the expanded determinants give
\(\sum_\tau\operatorname{sgn}(\tau)\operatorname{vol}(\mathcal B)/
|\det A_\tau|\).

To compute these Jacobians, let \(D_\epsilon\) multiply by
\(\epsilon_s\) in letter \(s\), let \((Sv)_{b,s}=v_{b,s+1}\),
and let \((R_\tau v)_{b,s}=v_{\tau_s(b),s}\). Then
\[
 A_\tau=P D_\epsilon S(I-P^{-1}H),\qquad H=S^{-1}R_\tau .
\]
The unsigned permutation \(H\) shifts around the cycle and
permutes the rows. Its return
permutation on one letter is a product \(\pi\) of the
\(\tau_s\), up to the convention for representing
permutation matrices. A cycle of length \(j\) of \(\pi\)
corresponds to a cycle of length \(Lj\) of \(H\).
Thus
\[
            |\det A_\tau|=P^{kL}|\det(I-q^{-1}\pi)| .
\]
The signs of the tuples equal the signs of their return
permutations. Each \(\pi\) occurs \((k!)^{L-1}\) times.
The two integrals show that the ratio
\(\operatorname{vol}(\Omega\cap\mathcal B)/
\operatorname{vol}(\mathcal B)\) depends only on \(k,q\)
and is the same as for one even letter with step \(q\).

In that one-letter calculation the cone condition on
\(\beta=A_1v\) is \(\beta_b\ge q\beta_{b+1}\).
The triangular linear map
\[
    y_b=(1+q+\cdots+q^b)(\beta_b-q\beta_{b+1}),
    \qquad \beta_k=0
\]
maps the cone bijectively to the nonnegative orthant and
preserves the sum of the coordinates by telescoping.
Its determinant is \(\mathcal F_{k,q}\). It therefore
gives the asserted volume ratio on a slab of unit width.
\footnote{\weightedpin{grh/src/New/UpperBound/ElementaryDivisors/PieriCone.lean}
{143-L207}{The Jacobian for the cyclic step map}, and
\weightedpin{grh/src/New/UpperBound/ElementaryDivisors/DefectCount.lean}
{898-L920}{The interlacing cone and its exact volume factor}.}

Modulo \(p\), \(\mathscr P\) is the lower two-block parabolic
and \(\mathscr I\) the lower Borel. The quotient of the
two-block parabolic by its full unipotent radical is the
product of the two general linear groups. Hence choices
of a parabolic flag and of the full flags within the two
blocks give a full flag, so
\([K_\lambda:\mathscr P][G:G\cap\mathscr I]=\mathcal F_{k,q}\).
The slab volume is at most
\(\prod_s\bigl(D_s^{k-1}/(k-1)!\bigr)\),
by the volumes of the positive simplices with degree bounds
\(D_s\) and \(D_s-1\).
Since \(2LkE_0\le\theta_0 n\), it follows that
\[
 \begin{split}
 [K_\lambda:\mathscr P]\dim_{\mathbf k}\mathscr D
 &\le\operatorname{vol}(\mathcal B)\\
 &\le \exp\!\left(L(k-1)E_0/n\right)
                               \left(\frac{n^{k-1}}{(k-1)!}\right)^L\\
 &\le (1+\theta_0)\binom{n+k-1}{k-1}^{L}=(1+\theta_0)|J|.
 \end{split}
\]
Here \(e^u\le1+2u\) for \(0\le u\le1/2\).
This completes the construction and the dimension bound used
in the integral edge block at \(s_p\).
The separate sum of the depth-space dimensions enters the
two half-restriction determinants.

We give the passage from these estimates to the truncated
determinant. On \(\mathcal O^{H\times X\times J}\), form the two
restriction matrices \(J_1,J_2\). They send the value at the
origin to \(\Delta\rho(h_e^{\mathsf t})\) times that value
on the forward edge, and the value at the terminus to
\(\Delta\rho(b^{-1})\) times that value on the reverse edge.
They have full column rank. Write \(J_i=Q_i d_i\) by Smith
reduction, with \(Q_i\) having an integral left inverse and
\(d_i\) square. Smith reduction here uses a least-valuation
entry at each finite step. This works over the
nondiscrete coefficient valuation ring.

For either end, every source and parahoric coset occurs. If
\(u\in\mathscr P\), then \(\Delta\rho(u)\Delta^{-1}\) is integral
invertible. Indeed the entries from a deep variable to a shallow
variable are divisible by \(p=w^2\) in each letter, balancing the
one power of \(w\) in the half-depth scale. This is true for \(u^{-1}\)
as well. Consequently a vector in the saturation of \(J_i\)
has values in \(M_*\) in every edge coordinate. For the same
reason, if \(t_\nu\) is a Smith divisor of \(J_i\), the corresponding
primitive source column has all normalized translated coefficients
divisible by \(t_\nu\). If \(v_p(t_\nu)>(j-1)/2\), its reduction
lies in \(\mathcal D_j\) at each source. These columns are independent
on reduction. Thus for each \(j\) their number is at most
\(|H||X|\dim\mathcal D_j\). Their valuations are bounded by half the
maximum depth, by testing the identity translate on the primitive
source column. Summing over \(j\) gives
\[
                 v_p(\det d_i)\le\frac12\mathfrak c C,
                              \qquad C=|H||X||J|.
\]

On the edge coordinates let \(D_U\) act by \(U_e\) on \(H\)
and by the identity on \(J\). Extend \(B,B',V\) diagonally
over the edges and the \(H\)-coordinates, and let \(R_2\)
be a left inverse of \(Q_2\). The actual paired middle
matrix and its near inverse on the edge module are obtained from
\[
     A=Q_2^{\mathsf t}D_U B,\qquad
     K_1=V B'D_U^{-1}R_2^{\mathsf t}.
\]
The image of \(K_1\) has every value in \(\mathscr C\).
Since values of \(Q_2\) lie in \(M_*\), the averaging congruence
gives \(AK_1\equiv p^{2A_0'}I\bmod w^c\). Here the \(U_e\)
cancel because they act in \(H\) and commute with the
operators on \(J\).

Let \(R_1Q_1=I\), put \(\epsilon=w^{c-1}\), and apply Smith
reduction to \((I-Q_1R_1)K_1\) in the free complement of \(Q_1\).
Keep the primitive column for each divisor of valuation at
most \(v_p(\epsilon)\). Each such column multiplied by \(\epsilon\)
belongs to \(\mathscr C\) edge by edge, by dividing the corresponding
column image and using that the original image is in
\(\mathscr C\). The columns of \(Q_1\) have the same property.
All these primitive columns are independent modulo the maximal
ideal. The residue-dimension bound shows that the number \(h\)
of kept columns satisfies
\[
                 C+h\le |H||X|[K_\lambda:\mathscr P]\dim\mathscr D
                       \le (1+\theta_0)C .
\]
The other terms in the complementary image are divisible
by \(\epsilon\). Apply \(A\) to the decomposition of \(K_1\)
into its projection on \(Q_1\), the kept terms and the discarded
terms. Since \(v_p(\epsilon)>2A_0'\), the result is
\[
 p^{2A_0'}\mathcal O^C\subset (AQ_1)\mathcal O^C+
               \operatorname{span}_{\mathcal O}\{\text{\(h\) vectors}\}.
\]
The inclusion follows because the matrix on the left after moving
the error is \(p^{2A_0'}\) times a matrix congruent to the identity
modulo the maximal ideal.

If \(p^u\mathcal O^C\subset B_0\mathcal O^C+\operatorname{span}(h
\text{ vectors})\), at most \(h\) Smith divisors of \(B_0\) have
valuation greater than \(u\). To see this, diagonalize \(B_0\)
and project to those coordinates. In this block the inclusion,
applied to each column of \(p^u I\), gives
\(p^u I=D_0L_0+U_0S_0\), with \(D_0\) the diagonal block
and \(U_0\) the \(h\)-column generator matrix.
The equality \(p^{-u}U_0S_0=I-p^{-u}D_0L_0\) makes the
left side an integral invertible matrix, because
\(p^{-u}D_0\) is zero modulo the maximal ideal. Therefore
the rank of \(U_0\) is at least the size of the block.
It follows that the
truncated determinant cost of \(AQ_1\) at any cutoff \(s\) is at
most \(2A_0'C+\theta_0sC\). The costs of \(d_1\) and \(d_2\)
add to it under multiplication, with the same \(s\). The truncated
valuation of each diagonal divisor \(t_\nu\) is
\(\min(v_p(t_\nu),s)\), because the lattice is enlarged by
\(p^s\mathcal O^C\). The tuple recombination rule gives
subadditivity.

Finally \(d_2^{\mathsf t}AQ_1d_1=J_2^{\mathsf t}D_U BJ_1\)
is the actual edge matrix of \(Y_\lambda\) multiplied at the
source by \(B\). This follows edge by edge from the adjoint
identity and \(\Delta B\Delta=D_pB\). Its image is therefore
contained in \(Y_\lambda M_p\). We conclude that any integral
tuple can be recombined into \(Y_\lambda M_p+p^sM_p\)
at determinant valuation at most
\[
                      C(\mathfrak c+2A_0'+\theta_0s).
\]
This recovers the truncated statement from the two prescribed
estimates on the actual edge module.
\footnote{\weightedpin{grh/src/New/UpperBound/ElementaryDivisors/Edges.lean}
{612-L668}{Smith divisors of the two restriction matrices},
\weightedpin{grh/src/New/UpperBound/ElementaryDivisors/ElementaryDivisorTheorem.lean}
{776-L861}{The weighted pairing and actual operator}, and
\weightedpin{grh/src/New/UpperBound/ElementaryDivisors/ElementaryDivisorTheorem.lean}
{1011-L1141}{The middle estimate and truncation}.}

\paragraph{Uniform sizes and the served-prime budget.}
\teamanchor{proof.weighted-local-marked-budget}%
Take \(\theta_0=1/(m^2 f)\). The two size inequalities in the
averaging estimate hold uniformly for served \(p\).
We have \(L_\lambda\le f=O(x^{11}\log x)\),
\(k\asymp x^5\), \(1\le o_p\le r_p=O(\log x)\), and \(2r_p\le k\).
Also \(p<x_1\), so
\[
 \log(P^{3r_p})=O\bigl((\log x)^2\log x_1\bigr)=o(\log b).
\]
All the other factors of the degree lower bound are at most
polynomial in \(x\), including \(\log q_\lambda\).
Thus the required degree is \(b^{o(1)}\), and \(a\ge xb\).
On the other side, \(m\mid o_p L_\lambda\) gives
\(q_\lambda\ge13^m\), which exceeds \(a^4\) eventually.
In particular \(A_0'\), which is bounded by a polynomial in \(x\),
is at most \(L_\lambda aP^{-2r_p}\le\mathfrak c\) for large \(x\).
Since \(p^{r_p}\ge\max(m,p)\), \(P\ge p\) and
\(m r_p^2/k=o(1)\), we get
\[
          3\mathfrak c\le \frac{6L_\lambda a}{mp}
                              =\frac{12A_\lambda}{mp}.
\]
The one-operator cost is bounded by
\(C(12 A_\lambda/(mp)+s/(m^2 f))\).
For comparison, the compact release gives the conclusion on \(M_p\)
with the bound
\(C(128 A_\lambda\log b/(mp)+s/x^{14}+x^{40})\).
Both are bounds for the absolute determinant cost of each
tuple, with \(C\) the dimension of the complete stage-zero
column space.
\footnote{\weightedpin{grh/src/New/UpperBound/Divisibility/MarkOperator.lean}
{336-L414}{The readable normalized expense}, and
\weightedpin{grh-compact/src/New/Columns/MarkedPrimeType.lean}
{12-L18}{The compact bound on the same lattice}.}

At \(s=s_p=\lceil2mfa/p\rceil\), recombination with products
of the integral operators adds the costs. For \(Y=\prod_{\lambda\in\Lambda}
Y_\lambda\) there are \(m\) occurrences of every operator in
\(Y^m\). Since \(\sum_\lambda A_\lambda\le fa\) and
\(\#\Lambda\le f\), the resulting cost at the original cutoff
is no more than
\[
                 C\left(\frac{12fa}{p}+\frac{s_p}{m}\right)
                 \le C\left(\frac{14fa}{p}+\frac1m\right)
                 \le C e_p .
\]
Here \(e_p=3 fa(\log b)^2/p\) at a served prime.
The marked chain bounds the output on \(Y^mM_p\) and the
scalar correction bounds the output on \(p^{s_p}M_p\).
As the bound holds for every tuple, we may apply it to each
assigned translate block of integral ramified columns and use
linearity to retain ramification of the recombined vectors.
The same \(q\) from \(\beta_p=\lceil\gamma_p-e_p\rceil\)
works on the two conjugate readings, using the respective
marked sets. Thus the paired height cost is
\(F_p=2|\mathcal S|C e_p\log p\) for the fixed selected rows.
Its sum stays within the allowance for
\eqref{eq:weighted-A-paired-index}.

\paragraph{The index at the level prime.}
\teamanchor{proof.weighted-local-fourier-index}%
Use \(p=N\) and a fixed coefficient embedding \(j_N\).
Write \(\lambda(\tau)\) for the place induced by \(j_N\Phi_\tau\)
and let \(P_N=\{\overline{\lambda(\tau)}:\tau\in\Phi\}\).
Since \(N\) splits completely, \(P_N\) contains exactly one place
of each conjugate pair. These are the places used to measure the
coefficient action in the Weyl gauge, which is taken through
\(j_N\Phi_\tau\operatorname{conj}\).
Embed the local general linear groups at \(P_N\) in the unitary
group using inverse transpose at their conjugates.
For \(S\subseteq P_N\), let \(d_S\) be the resulting adelic
diagonal with last entry \(N\) at \(S\), all other entries \(1\).
Put \(d_0=d_{\{w\in P_N:w\mid\pi\}}\) and
\(u_0=d_{P_N}^{-1}\).

For an adelic diagonal \(d'\), let \(M(d')\) consist of the
columns whose Weyl gauge in the lattice at \(g d'\) is at most
one for every \(g\). Thus \(M(1)=M_N\).
After conjugation by \(d_0\), the local column level at \(P_N\)
becomes the group of integral invertible matrices with last
column congruent to the last standard basis vector.
At \(w\mid\pi\) the original condition is on the last row;
conjugating by the last-coordinate diagonal with entry \(N\)
converts it to the last-column condition. At \(\bar w\mid\pi\)
the condition is on the last row of the inverse transpose, which
already says that the original matrix fixes this last column
modulo \(N\). The same block calculation shows that
\((d_0d_S^{-1})^{-1}K_0(d_0d_S^{-1})\) is integral above \(N\)
for \(S=P_N\) and for \(S=\{w\in P_N:w\mid\pi\}\), where \(K_0\)
is the stage-zero column level.

For each fixed \(S\) use weak approximation to move every class
representative \(c\) by a rational unitary matrix so that
\(l=c d_0d_S^{-1}\) and \(d_S^{-1} l d_S\) are both integral
invertible above \(N\). This is an open condition containing
the identity on \(l\). The required weak approximation follows
from the Cayley map \(A\mapsto(1+A)(1-A)^{-1}\)
on the space of skew-Hermitian matrices over \(E^+\).
Its image contains the dense open set of local unitary
matrices with \(1+g\) invertible. Weak approximation in this
vector space meets the finitely many prescribed local opens
at the places over \(N\).
Now \(M(d_0d_S^{-1})\) is determined by ordinary integral values
at these representatives, because the shifted arguments there
are the integral points \(l\). The gauge for \(M(d_0)\) is the
gauge at \(c d_0=l d_S=d_S(d_S^{-1}l d_S)\), hence agrees with
the gauge at \(d_S\) on each representative. Testing only the
representatives suffices since the rational actions cancel in
transported gauges and \(d'^{-1}k d'\) is integral invertible
for \(k\in K_0\). Thus in the same representative coordinates
the two lattices are related by the rational diagonal action
of \(d_S\) at the given coefficient embedding \(j_N\).
In the symmetric-power factors its exponents
are the powers of the last coordinate in the monomials at
tops measured in \(S\). Permuting the \(\kappa\) coordinates
shows that their mean is \(a/\kappa\) in each factor.
The inclusion \(M(d_0)\subset M(d_0d_S^{-1})\) has
determinant valuation at most
\[
                              C T a/\kappa .
\]
The finite-matrix argument above supplies the same bound for
recombining any tuple from the larger lattice into \(M(d_0)\).

In the conjugated local representation let \(R_0(k)\) denote
translation by \(d_0 k d_0^{-1}\), using the indicated embedding
of the local product group. The actual finite Fourier operator is
\[
  F=R_0(u_0)\!
    \sum_{\substack{y\in((\mathbf Z/N)^{P_N})^\kappa\\
                    y_{\mathrm{last}}\text{ a unit}}}
          \psi_N(y_{\mathrm{last}})\,R_0(k_y),
\]
where \(k_y\) is integral invertible and has last column and
last row congruent to \(y\).
The coefficients \(\psi_N\) are roots of unity.
To see that \(F\) preserves the column level, take \(k\) in the
conjugated last-column stabilizer. Then \(u_0^{-1} k u_0\)
is integral invertible with last row congruent to the last
standard basis row. It acts bijectively on the indexing vectors
\(y\), preserving their last coordinate. The matrix formed by
multiplying \(k_y\) by \(u_0^{-1} k u_0\) on the left has
the same last column modulo \(N\) as the representative of the
new index. Their quotient on the right fixes the input column.
Thus left multiplication by \(k\) permutes the sum defining \(F\).
Away from \(N\) the column level commutes with the sum.

If \(\phi\in M(d_0)\), evaluate the \(y\)-summand at \(g\).
Its input argument is \(g d_0u_0 k_y d_0^{-1}\), and the lattice
after the \(d_0\) shift is the one at \(g d_0u_0 k_y\).
This agrees with the transported lattice at \(g d_0u_0\)
because \(k_y\) is integral invertible. The ultrametric
inequality gives \(F M(d_0)\subset M(d_0u_0)\).
Start with a tuple in \(M(1)\) and recombine into \(M(d_0)\).
Each of the subsequent \(m/2\) applications of \(F\)
lands in \(M(d_0u_0)\) and allows another recombination
back into \(M(d_0)\). Pull the matrices back to the input
tuple by linearity. We obtain an integral recombination whose
columns have \(F^{m/2}\)-images in \(M(d_0)\) and whose determinant
valuation is at most
\[
         (m/2+1) C T a/\kappa \ \le\ 4mfb C=e_N C
\]
for large \(x\), since \(T\le f\) and \(a/(\kappa b)=O(x^{-4})\).
This proves the index estimate for the test used at \(N\).
\footnote{\weightedpin{grh-compact/src/New/Seesaw/NamesAtSplitLevel.lean}
{93-L115}{The actual operator on the residue vectors} and
\weightedpin{grh-compact/src/New/Seesaw/LiftCharacters.lean}
{128-L169}{The level-prime preimage test in the fixed code}.}

\paragraph{The paired residue calculation.}
\teamanchor{proof.weighted-local-fourier-weil}%
Use \(R_N=(\mathbf F_N)^{P_N}\) and \(e\) for the last basis
vector in each rank. In rank \(n_j\), take \(d_j\) with last
entry \(N\) at \(P_N\cap\{w:w\mid\pi\}\) and last entry \(1\)
elsewhere, as for \(d_0\). Write \(b_y\) for the indicator that the
last column of an integral matrix is \(y\) modulo \(N\),
\(a_y\) for the corresponding last-row indicator, and
\(c_y(S)=\psi_N(e^{\mathsf t} S y)\) on integral matrices.
All are extended by zero outside the integral matrices.
In the split Schr\"odinger model the \(N\)-part of the glued
pair is a sum of translates by \(-y/N\) in the last column on
the \(\pi\)-half and by its conjugate in the last row.
At \(w\in P_N\) over \(\pi\), the translated first box tests
integrality of the other columns and \(T_{\rm last}-y/N\).
Equivalently it tests \(b_y(Td_{j})\). The other translate
multiplies the second box by \(\psi_N\) of its last row
against \(y\). At \(w\in P_N\) over \(\bar\pi\), the second
box instead tests \(S_{\rm last}-y^{\mathsf t}/N\),
using conjugation to identify the two \(\mathbf F_N\) coordinates,
and the first box gives the character of its last column
against \(y\). On this support the phase is
\(\psi_N((N S_{\rm last}\bmod N)\cdot(T_{\rm last}\bmod N))\).
In each integral first box on this half, \(z\) is its last
column and the phase is the same dot product. The product
over all places thus gives
\[
                     \sum_z b_z(Td_j)c_z(d'_{j+2}S),
\]
where \(d'_{j+2}\) has last entry \(N\) on
\(P_N\setminus\{w:w\mid\pi\}\). The other matrix coordinates
have the same integral tests in both expressions. In the
Schr\"odinger formula the self-pairing in the phase vanishes
for a translate supported on one split half. The remaining
linear term is the character \(\psi_N\) of the original
translation. It is unchanged by integral changes of the
lift of \(y/N\).

Write \(L_j^{\rm raw}[f]\) and \(R_j^{\rm raw}(g)\) for the lift and
right translation in the unshifted coordinates at \(P_N\).
With \(C_s(g)=\chi_{j,n_{j+1}}(g)\) and
\(C_t(k)=\chi_{j,n_j}(k)\), the local Weil action of \((g,k)\)
on the residue function is \(C_s(g)^{-1} C_t(k) f(k^{-1}Sg)\).
Here \(\chi_{j,r}\) is the base-change character for arrow \(j\)
at the determinant times \(|\det|^{-r/2}\) at \(P_N\).
Its character is unramified there because the chosen places and
the twist conductor are prime to \(N\).
Thus equivariance gives
\[
 L_j^{\rm raw}[f(k^{-1}\,\cdot\,g)]
   = C_s(g) C_t(k)^{-1}
      R_{j+1}^{\rm raw}(k)L_j^{\rm raw}[f]R_j^{\rm raw}(g^{-1}).
\]
Applying the lift to the Weil translate by \((d_j,d_{j+1})\)
defines the placed map
\[
 L_j[f]=R_{j+1}^{\rm raw}(d_{j+1})
            L_j^{\rm raw}[f]R_j^{\rm raw}(d_j^{-1}).
\]
Write \(R_j(g)=R_j^{\rm raw}(d_j g d_j^{-1})\).
Put \(B_j'=L_j[b_e]\) and \(u_j=d_{P_N}^{-1}\) in rank \(n_j\).
For integral invertible \(A,B\), split Weil covariance says
\[
   R_{j+1}(B)L_j[f(B\,\cdot\, A)]R_j(A)=L_j[f].
\]
Both determinant characters on the compact groups are one.
For \(u_{j+1}^{-1}S u_j\), write \(S\) in blocks with its
last row and column separated. Conjugation divides the
top part of the last column by \(N\), multiplies the
left part of the last row by \(N\), and fixes the bottom
right entry. Hence \(a_e(u_{j+1}^{-1}S u_j)=b_e(S)\),
including their support conditions. The full covariance gives
\[
   R_{j+1}(u_{j+1}) L_j[a_e]=\gamma_j B_j' R_j(u_j),
   \qquad
   \gamma_j=\prod_{w\in P_N}N^{(n_{j+1}-n_j)/2}.
\]
The same unramified base-change character occurs in numerator
and denominator at each place. Its values at the two
last-coordinate determinants cancel. The remaining
absolute-value powers are exactly the displayed rank difference.
In particular \(|j_N(\gamma_j)|_N\le1\).

When a component of \(y_{\rm last}\in R_N\) vanishes, the
sum testing that last entry is invariant under a source
central unit at that place. Its lift kills a vector in the
image of \(1\) minus this unit's translation. The ramifying
differences give this condition separately at every \(w\in P_N\).
Expanding \(c_e\) along its last column and last row, respectively,
therefore gives for ramified \(\phi\)
\[
  R_{j+1}(u_{j+1})L_j[c_e]\phi
       =F_{j+1}B_j'\phi
       =\gamma_j B_j'F_j\phi ,
\]
where \(F_j=R_j(u_j)\sum_{y_{\rm last}\in R_N^\times}
 \psi_N(y_{\rm last})R_j(k_y)\).
In the last-column expansion use \(k_y e=y\), and in the
last-row expansion use \(e^{\mathsf t}k_y=y^{\mathsf t}\).
If \(y=(v,t)\) with \(t\) a unit, one such residue matrix is
\[
        k_y=\begin{pmatrix}I+t^{-1}vv^{\mathsf t}&v\\
                             v^{\mathsf t}&t\end{pmatrix}.
\]
Its determinant is \(t\), by the Schur complement.
The operators \(F_j\) commute with the central translations.
The maps \(B_j'\) intertwine them, since simultaneous
multiplication by a central unit on both sides fixes the
residue argument. These two maps thus preserve the
ramifying-image conditions.

In the first arrow of a pair use \((g,k)=(d_{2i},1)\) in the
unshifted covariance formula. In the second use
\((g,k)=(1,(d'_{2i+2})^{-1})\). This removes the two
complementary diagonals from the matrix arguments and gives the
scalar
\[
 \chi_{2i,n_{2i+1}}(d_{2i})\
 \chi_{2i+1,n_{2i+1}}(d'_{2i+2}).
\]
Indeed the source determinant factor in the first arrow is
read in rank \(n_{2i+1}\), and the target factor in the second
arrow in rank \(n_{2i+1}\). The intermediate translations
by \(d_{2i+1}\) cancel when the maps are written in placed
coordinates. At the exterior stages there are the \(d_{2i}\)
and \(d_{2i+2}\) shifts of the placed maps, with the last
factor \(R_{2i+2}(u_{2i+2})\) since
\(d_{2i+2}d'_{2i+2}=u_{2i+2}^{-1}\).
Each \(w\in P_N\) contributes to exactly one of the two
determinants. The \(N\)-valuation of its character value there
is the signed infinity weight at the matching place, by the
same ray-class argument as in the marked transfer. Its absolute
value as an exponent is at most \(2b\) by the stage weight bound.
The absolute-value
power in \(\chi_{j,r}\) has \(N\)-adic size at most one for
these diagonals. Thus the scalar for this pair has \(N\)-adic
size at most \(N^{2fb}\).

In the remaining composition
\(R_{2i+2}(u_{2i+2})\sum_y L_{2i+1}[c_y]L_{2i}[b_y]\),
the terms with \(y_w=0\) at any \(w\) kill the ramified input.
Every other \(y\) is \(ke\) for some integral invertible \(k\).
Covariance gives \(L_{2i}[b_y]=R_{2i+1}(k)B'_{2i}\) and
\(L_{2i+1}[c_y]R_{2i+1}(k)=L_{2i+1}[c_e]\).
The sum is therefore the number of these vectors times
\(F_{2i+2} B'_{2i+1} B'_{2i}\) on the input.
Move the \(F\) through all preceding \(B'_j\) using the
same commutation identity. Induction over the pairs gives
a single scalar times
\[
                       B'_{m-1}\cdots B'_0 F_0^{m/2}\phi .
\]
The counts of residue vectors are integers and every \(\gamma_j\)
is \(N\)-adically bounded by one. The scalar is shared by all
terms of the good datum\footnote{\weightedpin{grh-compact/src/New/Seesaw/NamesAtSplitLevel.lean}
{140-L162}{Covariance and the conjugated support} and
\weightedpin{grh-compact/src/New/Seesaw/NamesAtSplitLevel.lean}
{425-L513}{Split determinant characters and the glued data}.}
and has \(N\)-adic size at most \(N^{mfb}\), which suffices for
the allowance \(N^{2mfb}\).

The gauge bound for \(B'_j\) follows from
\hyperlink{proof.weighted-local-source-gauge}{the placed source
gauge at \(N\)} and
\hyperlink{proof.near-reference-kernel}{the finite Fourier
calculation for residue functions}.
The last-column indicator on the integral box has a
Fourier expansion on the residue matrices at \(P_N\), with
coefficients integral after multiplication by
\(N^{2f n_j n_{j+1}}\). Indeed the inversion denominator is
the order \(N^{f n_j n_{j+1}}\) of the matrix group under
addition, and each coefficient numerator is an integral
sum of roots of unity. Naming the residue function first
projects the datum at \(N\) onto its vacuum part. For a Heisenberg
translate, this projection is zero if its translation vector fails
the integral-box test there. Each additive character on the
box is then a Heisenberg translation of the projected vacuum, with
possible simple poles only at \(\overline{P_N}\).
Choose \(\delta\in\mathcal O_K\)
zero modulo the prime of \(K\) induced by \(j_N\) and
one modulo its conjugate. All places of \(\overline{P_N}\)
are over the first prime, and hence multiplication by
\(\delta\) clears the poles. The conjugate \(\bar\delta d\)
is a unit at \(j_N\). In the normalized standard-data
estimate the per-box ratio
\(|j_N(\bar\pi_G)/j_N(\bar\delta d)|_N\) is one.
After removing the Fourier denominator the gauge cost is
at most \(N^{2f n_j n_{j+1}}\le N^{fb}\).
The last-column indicator is stable under the two conjugated
stabilizers. For \(A,B\) fixing the respective last vectors
modulo \(N\), \(B^{-1}S A e=e\) in the residue field exactly
when \(S e=e\), and integral invertibility preserves the
integral box. Thus the shifted lattices and the class-sum
calculation use the same levels and gauges.
\footnote{\weightedpin{grh-compact/src/New/Curve/ConjugatesOfJets.lean}
{129-L147}{The standard-data gauge at the original kernel scale}
and \weightedpin{grh-compact/src/Lib/Weil/MixedModel/NameModulus.lean}
{179-L187}{Residue characters as vacuum translates}.}

For the final restriction the \(d_m\)-shift contributes
\(N^{-b_\tau}\) in the last-coordinate action at embeddings with
\(\overline{\lambda(\tau)}\mid\pi\).
By the branching identity these factors match, in the valuation,
the scalar \(\prod_\tau\Phi_\tau(\bar\pi)^{-b_\tau}\) in the
restriction. The integer branching map is integral in the
remaining rank. Consequently the usual output Weyl gauge is
at most \(N^{3mfb}\) on the columns passing the test, and the
index bound just proved pays for the required local submodule.

\paragraph{The common scalar and the selected rows.}
\teamanchor{proof.weighted-short-paired-index}%
Take the one rational scalar
\[
  \beta_p=\lceil\gamma_p-e_p\rceil,\qquad
  q=\prod_p p^{-\beta_p}>0,\qquad B_s=qA_s .
\]
The product has finite support. At an embedding above \(p\) the
factor \(q/p^{-\beta_p}\) is a unit. At a small non-split prime,
\(\beta_p\le\gamma_p-2\ell_p\) for large \(x\) because \(\ell_p\ge1\)
uniformly. At \(N\) we have
\(\beta_N=\lceil-4mfb\rceil\le-3mfb\).
At the remaining primes \(\beta_p=0\).
Thus the local estimates give, for every finite tuple of integral
ramified columns and every twist of level norm prime to \(p\),
an integral coefficient matrix of determinant valuation at most
\(C e_p\) for which the recombined definite lifts multiplied by
\(q\) are integral at the specified embedding. By the common
seesaw and the integral projection this also holds for the
corresponding \(\Theta_\vartheta B_s\)-images at any fixed \(s\).

At the archimedean embeddings the same \(q\) satisfies
\[
       2\log q+2\sum_p\gamma_p\log p
       \ \le\ 2\sum_p e_p\log p\ \le\ 2\mathcal B .
\]
Use the spanning set \(S\) from \eqref{eq:weighted-A-short-column}
in \eqref{eq:weighted-A-chain-bound} and add the projection cost.
Since \(h_0\le h_{\rm main}-m\,\mathrm{rem}+\mathcal B\) and
\(5\mathcal B+x^{99}\le L\) for large \(x\), the result is
\begin{equation}
\label{eq:weighted-short-forms}
   \|\Theta_1 B_s\phi\|_\sigma^2\le e^h,\qquad
              h=h_{\rm main}-m\,\mathrm{rem}+L
       \quad(\phi\in S,\ s\in\mathcal G,\
                 \sigma\in\operatorname{Aut}(\mathbf C)).
\end{equation}

Take a finite set of translates \(\mathcal S\) and assign each row
of a family \(y\) in their images of integral ramified columns
one representation \(y_i=B_{s_i}\phi_i\), with \(s_i\in\mathcal S\).
For a prescribed \(\iota=(p,j_\iota)\), take a twist \(\theta\) with
level norm prime to \(p\). At each of \(j_\iota\) and
\(j_\iota\circ\operatorname{conj}\), group the rows by \(s_i\),
recombine each block of columns, and put its matrix on the
diagonal in the given row order. The resulting matrices \(U_+\)
and \(U_-\) over the respective valuation rings make the lifts
for the same \(\theta\) integral at the two places. The
determinant valuation of each matrix is at most
\(|\mathcal S|C e_p\). Put
\[
    E_p=|\mathcal S|C e_p\log p,\qquad
    \delta=|\mathcal S|C\mathcal B .
\]
The paired index bound is
\begin{equation}
\label{eq:weighted-A-paired-index}
 -\log|j_\iota(\det U_+)|_p
 -\log|(j_\iota\circ\operatorname{conj})(\det U_-)|_p\le F_p,
       \qquad F_p=2E_p .
\end{equation}
All finite sums of \(F_p\) are bounded by \(2\delta\) by the summed
allowance for \(e_p\). The maps, row preimages and costs precede the
choice of the coefficient place. The matrices are chosen for the
embedding and twist used there.

\subsection*{Enough forms (Theorem B)}

\teamanchor{proof.weighted-rank}%
Work with the datum used in Theorem~A. Write $G=\mathcal G=G_0\times G_1$, where
$G_i$ is the product at the chosen places over $g_i$. Let $K(g)$ be the row
principal level in the common finite adelic unitary group of rank $r$ and
$K_r$ its full lattice stabilizer. Thus $K_r/K(g)=G$. Let $S=U\rtimes H_0$
be the Bessel subgroup on $Y_0$, with $U$ unipotent and $H_0=U(X_0)$ fixing
the frame. Use the class representatives $h_c$ in $H_0(\mathbf A_f)$ and
the Bessel coordinates of the filling datum. The column coordinate
isomorphism first translates by the central lattice-scale element and then
takes the Weyl-to-Bessel coordinates of the class values. We give the
resulting coordinate space its Euclidean inner product. Its basis consists
of $e_{c,i}$ for a class $c$ and a tuple $i$ of Bessel-coordinate indices.
Each of the index sets is nonempty, so $C>0$.\footnote{\weightedpin
{grh-compact/src/New/Seesaw/BesselGeneric.lean}{12-L31}{Filling datum and coordinate isomorphism}
and \weightedpin{grh-compact/src/New/Filling/ColumnCoords.lean}{48-L53}{Class coordinates}.}

\paragraph{Unfolding the normalized pairings.}
\teamanchor{proof.weighted-bessel-unfolding}%
Use the product of the filling Haar measures at infinity and the common
finite Haar measure $d\mu_f$ on the row group. The quotient by rational
points uses counting measure. At infinity write $dy$ as the disintegration
along $U_\infty$, with the filling Haar measure $du$ on $U_\infty$ and
quotient measure $d\dot y$ on $U_\infty\backslash G_\infty$. The product
of full Bessel vectors is $W(y)$, with
$\int W_i(y)\overline{W_j(y)}\,d\dot y=\delta_{ij}$. It transforms on
$U_\infty$ by its Bessel character and on $H_{0,\infty}$ by the unitary
Bessel-coordinate representation. The unipotent cutoff $\beta_c$ satisfies
$\sum_{u\in U_c}\beta_c(uy)=1$ for
$U_c=U(F)\cap h_c K(g)h_c^{-1}$, where $F=E^+$. Its integral along each
$U_\infty$ orbit is $\operatorname{vol}(U_c\backslash U_\infty)=\nu$, the
same number for each $c$. Take the signed frame scales \(s_{va}\)
from the \hyperlink{proof.weighted-real-bessel-vectors}{real Bessel frame}.
Put $\varrho=m/2-1$, $k=\kappa$, $d_U=\varrho(\varrho+k)$ and
\begin{equation*}
 \begin{aligned}
   D_U&=\prod_{v\in\Phi}\prod_{a=0}^{\varrho}
                     (s_{va}^2)^{k+2(\varrho-a)},\\
   \nu&=\frac{(g^{2d_U+\varrho}/N^k)^f\,|\operatorname{disc}E|^{\varrho/4}}
              {\operatorname{vol}(\widehat{\mathcal O}_E)^{d_U}D_U},
   &\operatorname{vol}(\widehat{\mathcal O}_E)&=
                2^f|\operatorname{disc}E|^{-1/2}.
 \end{aligned}
\end{equation*}
Here $g=g_0g_1$ and $\varpi$ is the chosen generator of the different
with $\bar\varpi=-\varpi$. For the finite correlation kernel use
$\chi_f^+(u)=\psi_f(\sigma_U(u)/(g\varpi))$, with
$\sigma_U(u)$ the sum of the entries just above the diagonal of the
frame block. It is trivial on the principal congruence subgroup of $U$
and has inverse $\chi_\infty(u)$ on rational unipotents.
\footnote{\weightedpin
{grh-compact/src/New/Filling/Setup.lean}{128-L160}{Frame and covolume normalization},
\weightedpin{grh-compact/src/New/Data/Different.lean}{346-L350}{The generator of the different},
\weightedpin{grh-compact/src/New/Filling/UnipotentCovolume.lean}{24-L107}{Formal covolume calculation},
and \weightedpin{grh-compact/src/Lib/Adelic/Adeles/RationalPoints.lean}{155-L157}{Additive Haar factor}.}

\paragraph{Covolume of each unipotent fibre.}
\teamanchor{proof.weighted-unipotent-covolume}%
The computation of \(\nu\) uses successive outer rows of the unipotent
group. Use rational coordinates with frame \(e_0,\ldots,e_{\varrho-1},z\)
and dual isotropic vectors \(e_0^*,\ldots,e_{\varrho-1}^*\).
In the finite class lattice, the basis matrix \(B_c\) is the product
of the Bessel lattice change and the action of the class representative
in \(H_0\). The class lattice is the orthogonal sum of its middle
block and the hyperbolic planes.
The matrix is the identity on the hyperbolic pairs. In the middle
block its determinant has finite idele modulus \(N^f\). Indeed the
Bessel change replaces one vector \(x\) of \(X_0\) by
\(x/\bar\pi+z/N\), where \(\pi\bar\pi=N\). The determinant of the
class representative has norm one. The finite modulus is unchanged
by conjugation, so \(|\bar\pi|_f^2=|N|_f=N^{-2f}\).
The finite class lattice is self-dual. In the two coordinates changed
by the Bessel basis the Gram matrix is
\(\left(\begin{smallmatrix}0&-1\\-1&-N\end{smallmatrix}\right)\), which
has determinant \(-1\), and the class representative is an
isometry.\footnote{\weightedpin{grh-compact/src/New/Filling/EichlerFlag.lean}
{15-L47}{Bessel lattice change} and
\weightedpin{grh-compact/src/New/Filling/LevelBridge.lean}
{55-L60}{The basis of each class lattice}.}

At row \(a\) remove the outer hyperbolic plane from the subspace
remaining after rows \(0,\ldots,a-1\). Write \(J_a\) for the form on
the space \(Y_a\) after this removal, of dimension
\(d_a=k+1+2(\varrho-1-a)\) over \(E\). The row subgroup is formed
by the matrices
\[
 n_a(l,t)=
 \begin{pmatrix}1&l&t\\0&I&-J_a^{-1}l^*\\0&0&1\end{pmatrix},
 \qquad t+\bar t=-lJ_a^{-1}l^* .
\]
Here \(l\) is a row on \(Y_a\). Projection to \(Y_a\) gives the next
unipotent group. This projection splits both over the reals and on
the class lattice at principal level since its basis respects the
decomposition into the outer pair and \(Y_a\).
Let \(B_{c,a}\) be the restriction of the class basis matrix to
\(Y_a\). In this basis the conditions on the row and its conjugate
column are equivalent to
\((l/g)B_{c,a}\in\widehat{\mathcal O}_E^{\,d_a}\), by self-duality.
For such an \(l\), the scalar \(lJ_a^{-1}l^*/g^2\) is in \(\mathcal O_F\).
Take \(\omega=-\zeta_3\in\mathcal O_E\), which has trace one since
the field contains the \(5040\)-th roots of unity. Thus
\(t=-\omega lJ_a^{-1}l^*\) satisfies the equation and is integral
modulo the prescribed level in the sense \(t\in g\mathcal O_E\).
All other admissible \(t\) differ from it by \(gi\mathcal O_F\), since
the datum contains \(i\) and \(\mathcal O_E\cap F=\mathcal O_F\).
Consequently every row in the displayed lattice occurs, and its
central fibre is a translate of the same skew lattice.
\footnote{\weightedpin{grh-compact/src/New/Filling/BoxPartition.lean}
{224-L263}{Principal level and the same row translations}.}

In a row subgroup the Haar measure is Lebesgue in the free complex
row and in \(\operatorname{Im}t\). Multiplication shifts the row
additively and shifts the central coordinate by a function of the row.
In the product \(n_a(l,t)\operatorname{diag}(1,h,1)\) the free row
is \((lh,t)\), giving Jacobian one at fixed \(h\) since \(|\det h|=1\).
Successive disintegration therefore gives the product measure in the
free entries. This is also the Haar measure used for
the Bessel quotient. On the group fixing the frame, the free complex
entries in \(X_0\) and the independent real entries of the skew
matrix are the defining coordinates of that Haar measure.
The quadratic correction to the skew matrix gives a translation on
each fibre over the rectangular block. Adding the frame entries by
the upper unitriangular lift preserves the product measure. At fixed
frame entries, right multiplication by that lift is triangular in
the remaining free entries of each row, with diagonal one.
\footnote{\weightedpin{grh-compact/src/Lib/DiscreteSeries/Bessel/UnipotentCoords.lean}
{233-L244}{Definition of the quotient measures} and
\weightedpin{grh-compact/src/Lib/DiscreteSeries/Bessel/RealMeasures.lean}
{192-L211}{The free-entry Haar}.}

The standard \(\mathcal O_E\)-lattice in the complex embeddings has
covolume \(2^{-f}\sqrt{|\operatorname{disc}E|}\).
For the row lattice cut out by \(lB_{c,a}\in\widehat{\mathcal O}_E^{\,d_a}\)
the covolume is this quantity to the power \(d_a\) times
\(|\det B_{c,a}|_f=N^f\). The determinant factor follows by elementary
divisors at the finite places. The imaginary lattice obtained from
\(i\mathcal O_F\) has covolume \(\sqrt{|\operatorname{disc}F|}\).
Both \(i\) and \(\zeta_3\) generate \(E\) over \(F\). Their quadratic
polynomial discriminants \(-4\) and \(-3\) show that the relative
discriminant is the unit ideal. Thus
\(\sqrt{|\operatorname{disc}F|}=|\operatorname{disc}E|^{1/4}\).
In rescaled real coordinates the complex entry in column \(j\) is
multiplied by the scale of \(j\) divided by that of \(e_a\), and the
central entry by \(s_{va}^{-2}\). The two scales at an inner
hyperbolic pair are inverses. The complex entries in \(X_0\) have
column scale one and the entry in column \(z\) has squared column
scale \(s_{v\varrho}^2=1/N\). Hence the real Jacobian at \(v\) on
row \(a\), including the level factor \(g\), is
\[
 g^{2d_a+1}s_{v\varrho}^2\,
       (s_{va}^2)^{-k-2(\varrho-a)} .
\]
In multiplying over the rows, \(N^{f\varrho}\) from the finite
determinants cancels the factor
\(\prod_v s_{v\varrho}^{2\varrho}=N^{-f\varrho}\).
The extra factor for \(a=\varrho\) in \(D_U\) contributes
\(\prod_v s_{v\varrho}^{2k}=N^{-kf}\) when the remaining product is
written over \(D_U\). As \(\sum_a d_a=d_U\), this gives the displayed
\(\nu\) on each class. Integration of the partition of unity on a
fundamental domain gives the same number on each unipotent fibre.

For a set $T$ of rational terms define the blocks
\begin{equation*}
 A_T(c,s,c')_{ij}=\sum_{\substack{\gamma\in T\\
            h_c^{-1}\gamma_f h_{c'}\in K_r\\
            \overline{h_c^{-1}\gamma_f h_{c'}}=s}}
    \int\beta_c(y)\beta_{c'}(\gamma_\infty^{-1}y)
          W_i(y)\overline{W_j(\gamma_\infty^{-1}y)}\,dy.
\end{equation*}
In this formula the bar on a finite group element denotes reduction in
the filling frame. Use the
\hyperlink{proof.gram-unfolding}{Petersson unfolding}
\eqref{eq:relative-gram-unfolding} with both \(\Delta_i=\{1\}\):
the first seed is \(s_{(c,i)}\), and the second is \(s_{(c',j)}\)
right-translated by a lift \(k\in K_r\) of \(s\) in that frame.
This right translation acts by \(\eta k\eta^{-1}\) in the lattice
coordinates. Absolute integrability
follows from the uniform absolute sum and \(L^1\) seed bound in the
\hyperlink{proof.weighted-poincare-realization}{weighted periodization}.
Taking the conjugate gives the order of factors in \(A_T\).
In lattice coordinates the first seed's finite support is
\(x_f=h_c k_1\eta^{-1}\) with \(k_1\in K(g)\). The argument of the
second seed then has finite part \(\gamma_f^{-1}h_c k_1 k\eta^{-1}\).
It lies in \(h_{c'}K(g)\eta^{-1}\) exactly when
\(h_c^{-1}\gamma_f h_{c'}\in k_1kK(g)=kK(g)\), using normality
of \(K(g)\) in \(K_r\). This criterion holds across the first
support, whose Haar measure is \(\mu_f(K(g))\) in the transported
normalization. The real integrand is the product of the two
cutoffs and two full Bessel coordinates displayed above.
Write \(\mathcal Q\) for the unscaled Poincar\'e map
\(\mathcal P\) in the Euclidean column coordinates, and
\(P=(\mu_f(K(g))\nu)^{-1/2}\mathcal Q\). Then
\begin{equation*}
 \langle\rho(s)Pe_{c',j},Pe_{c,i}\rangle=A_{\rm all}(c,s,c')_{ij}/\nu.
\end{equation*}
We use an inner product linear in its second entry. Thus the matrix
entries of $K(s)=P^*\rho(s)P$ are the conjugates of the displayed
entries. If \(i_{c_*}\) takes the coordinates of \(c_*\phi\) on
\(\mathcal C_{\rm ram}\), the same projection as in
\eqref{eq:weighted-seesaw-evaluated} gives
\[
 B_s=q|\mathcal G|\,\rho(e_0)\rho(s)\mathcal Q i_{c_*}.
\]
The coordinate isomorphism and the nonzero scalars preserve preimage
dimensions for the rank calculation.
\footnote{\weightedpin{grh-compact/src/New/Filling/PoincareSeries.lean}
{73-L107}{Finite frame and seed}, \weightedpin{grh-compact/src/New/Filling/PoincareSeries.lean}
{127-L154}{Convergence and unfolding}, and
\weightedpin{grh-compact/src/New/Filling/PoincareSeries.lean}
{176-L196}{Petersson normalization}.}

To read the core component of a finite \(S\)-term, let \(I\) be
the span of the isotropic half-frame. The action on the middle
lattice on \(I^\perp/I\) is the core component, since \(U\)
acts trivially there. At the \(\pi\)-half of \(N\), the gluing
column is \(e/\bar\pi+z/N\). Its projection onto \(X_0\) completes
the unmarked columns to the integral core lattice because
\(\bar\pi\) is a unit. Intersection with \(X_0\) instead requires
the coefficient on that column to be divisible by \(\pi\),
in order to cancel its \(z\)-coordinate. Thus a core isometry
preserving the middle lattice is integral with integral inverse.
Applying it to the unmarked columns and subtracting the gluing
column from its image gives the last-row condition modulo
\(\pi\). Conversely that condition makes these images integral
in the glued lattice, and it also holds for the inverse.
The inverse conjugate transpose gives the condition at the
opposite half. Away from \(N\) the denominators are units and
the principal reductions restrict to the ordinary ones on the
core. Hence the core pullback of \(K(g)\) is the column level
\(K_0\) from the \hyperlink{proof.weighted-projection}{definite
column construction}. If an element of \(S\) belongs to \(K(g)\),
its middle action is in this \(K_0\); removing it leaves an
element of \(U\cap K(g)\). The same calculation for the core
lattice alone is used in the \hyperlink{proof.weighted-marked-disc-sum}
{marked discriminant comparison}.

At $s=1$, a term in $S(F)$ forces $c=c'$ and $\gamma\in U_c$. Indeed its
$H_0$ component identifies the two column classes at column level, and
the class stabilizers there are trivial. The rational unipotent phase at
infinity is trivial on $U_c$ by the finite congruence and the product
formula for the additive character. The sum of the second cutoff over
$U_c$ is one. Integration of the first cutoff along $U_\infty$ then gives
$\nu I_C$ for this block matrix. More generally, write $J$ for the image
of the integral points of $S$ in $G$, and $T_S(s)$ for the normalized
matrix obtained from the $S(F)$ terms. For $t\in J$, each class $c$ has
exactly one class $c'$ with an $S(F)$ term over $(c,t,c')$ modulo the
unipotent actions. The existence follows by decomposing the $H_0$ part
using the column classes and approximating the unipotent part by a
rational unipotent in its prescribed finite open coset. Multiplying one
term by the inverse of another reduces uniqueness to the case $s=1$.
It follows that $T_S(t)$ permutes the class fibres by unitary operators.
It vanishes outside $J$ and $T_S(1)=I_C$. Translating a rational term by
one such $S(F)$ element permutes the rational terms over the corresponding
class and residue labels. To see that this transport preserves the
cutoff normalization, group the absolutely summable terms into
$U_c$--$U_{c'}$ double orbits. In an orbit represented by $\gamma$,
summing the second cutoff over $U_{c'}$ gives one and summing the
first cutoff over $U_c$ modulo
$U_c\cap\gamma U_{c'}\gamma^{-1}$ unfolds the first integral to the
quotient by that intersection. The two Bessel factors have trivial
phase under the respective class unipotents. Thus the orbit integral
is independent of the choices of the two partitions. Applying the
Bessel transformation law to the transported orbit gives
\begin{equation*}
  K(ts)=T_S(t)K(s),\qquad T_S(ts)=T_S(t)T_S(s)\quad(t\in J).
\end{equation*}
For $s$ supported in $G_i$, only the fixed class fibres contribute to
$\operatorname{tr}T_S(s)$. A term in such a fibre has its $H_0$ part
integral and congruent to one above $g_{1-i}$, so freeness at that single
level forces this part to be one. Conversely an integral unipotent
reduction acts on all fibres by its phase, by rational unipotent
approximation. For a rational representative \(u\), the factor
in the amplitude is \(\overline{\chi_\infty(u^{-1})}=\chi_\infty(u)\).
The operator entries are conjugate to the amplitude entries.
Their factor is therefore \(\chi_\infty(u)^{-1}=\chi_f^+(u)\).
Thus \(\operatorname{tr}T_S(s)=C\chi_U(s)\) for these \(s\),
where \(\chi_U\) is the reduction of \(\chi_f^+\) on integral \(U\),
extended by zero.\footnote{\weightedpin{grh-compact/src/New/Filling/NearTerms.lean}
{145-L174}{Bessel terms and the one-prime freeness}, \weightedpin
{grh-compact/src/New/Filling/AmplitudeSums.lean}{163-L211}{Covariance and identity term},
and \weightedpin{grh/src/New/Spanning/Unfolding/Moments.lean}
{696-L773}{Trace on a one-factor translate}.}

\paragraph{The share of allowed finite types.}
\teamanchor{proof.weighted-good-share}%
The good idempotent $e_0$ is the product of two central orthogonal
idempotents $e^{(i)}\in\mathbf C[G_i]$. At a chosen place $w$ above
$g_i$, an allowed label is a monic squarefree polynomial of degree $r$
with nonzero constant term and no root among $x\in k_w$ with
$x^{M_\omega}=1$, where $M_\omega=(Md)^{s+6}$. A good tuple uses allowed
labels everywhere and at least one irreducible label whose root generates
$\mathbf F_{q^r}^{\times}$, where $q=g_i$. Use the Green irreducible
representations for these labels, with the choice of characters in the
datum. Use the \hyperlink{proof.green-unipotent-criterion}{degree formula}
\eqref{eq:squarefree-green-degree} for this label with irreducible-factor
degrees \(d_1,\ldots,d_l\) and residue size \(q=g_i\). It gives
\[
 \frac{\dim(\tau_c)^2}{|\mathrm{GL}_r(k_w)|}
   = q^{-r}\frac{\prod_{j=1}^r(1-q^{-j})}
                  {\prod_{h=1}^{l}(1-q^{-d_h})^2}
   \ge q^{-r}(1-r/q),
\]
since \(\sum_h d_h=r\) and each numerator factor is at least \(1-1/q\).
The corresponding idempotent
has coefficient $\dim(\tau_c)\chi_{\tau_c}(u^{-1})/|\mathrm{GL}_r(k_w)|$
at $u$. Its identity coefficient is its share in the regular
representation.\footnote{\href{https://www.ams.org/tran/1955-080-02/S0002-9947-1955-0072878-2/S0002-9947-1955-0072878-2.pdf}
{Green, Theorems 2, 9 and 12--14}, and \weightedpin
{grh-compact/src/New/Trace/BadShare.lean}{29-L59}{Primitive and allowed labels}.}

There are at most $(1+M_\omega)q^{r-1}$ monic degree-$r$ polynomials
with a forbidden linear factor, including $X$. Those divisible by a
nonconstant monic square number at most $\sum_{1\le j\le r/2}q^{r-j}$,
which is at most $2q^{r-1}$. The regular share of all allowed labels
is therefore at least $1-(M_\omega+3+r)/q$. The number of primitive
labels is $\varphi(q^r-1)/r$, since a primitive root lies in no proper
subfield. If $n$ has $a$ distinct prime factors, its $j$-th prime factor
is at least $j+1$, so $\varphi(n)/n\ge1/(a+1)\ge1/(1+\log_2 n)$.
Consequently the primitive share is at least $c_0/(r^2\log q)$ for an
absolute $c_0>0$ at these scales. There are $f$ factors in $G_i$, with
$f\asymp x^{11}\log x$, $r\asymp x^5$, $M_\omega=x^{O_d(1)}$ and
$b^{10}\le q\le b^{11}$. For $x$ sufficiently large the union bound
for forbidden factors and the product bound for missing all primitive
factors give
\begin{equation}
 \label{eq:weighted-bad-share}
  \beta_i:=(1-e^{(i)})(1)\le
   \frac{f(M_\omega+3+r)}{q}+\exp\!\left(-\frac{c_0 f}{r^2\log q}\right)
   \le b^{-8}/2.
\end{equation}
The first term is $b^{-10+o(1)}$, and $f/(r^2\log q)$ is at least a fixed
positive multiple of $x^{1-\varepsilon}\log x$.

\paragraph{Transfer of traces on the good types.}
\teamanchor{proof.weighted-good-transfer}%
Work over $F=E^+$ with base field $\mathbf Q$ for restriction of scalars.
The Hermitian space $W_j$ has diagonal entries consisting of
$\kappa+j-1$ copies of $1$ and $m-j$ copies of $-1$.
The finite adelic frames on $Y_0$ and on each $W_j$ identify their finite
groups, compact subgroups, and reduction maps with the same $G_f,K_r,K(g)$
and $G$. The spaces of interest each use a single archimedean discrete
series at each real place, with the common regular Harish--Chandra
parameter up to permutation and the signature of the group in question.
The archimedean convolution kernel is $e_\pi=d_\pi\overline{m_\pi}$,
where $m_\pi(u)=\langle v_\pi,\pi(u)v_\pi\rangle$ for its specified unit
vector and $d_\pi$ is the formal degree. The parameter gaps ensure \(e_\pi\in L^1\). By the
\hyperlink{proof.bessel-real-projector}{formal-degree projection}
\eqref{eq:formal-degree-projector}, evaluation at \(v_\pi\)
identifies the multiplicity space with the \(e_\pi\)-fixed space
and preserves the commuting finite Hecke action. In particular the
two spaces give the same finite Hecke trace.
Let $p$ be a common level-preserving Hecke operator away from the chosen
places and their conjugates. Its realizing element is a finite linear
combination of translations which are trivial there. Sandwich each
translate by the averaging idempotent of the fixed compact subgroup away
from the chosen places. This gives a finite-away convolution operator with
the same action on principal-level forms. Use Haar probability on this
compact subgroup and the same finite test function on every inner form.
\footnote{\weightedpin{grh-compact/src/Lib/Adelic/Unitary/SpaceDefs.lean}
{267-L282}{Common finite frame}, \weightedpin{grh-compact/src/New/Trace/GoodIdempotent.lean}
{99-L121}{Spaces and reduction action}, \weightedpin
{grh-compact/src/Lib/DiscreteSeries/FlenstedJensen/RepresentationDefs.lean}
{119-L123}{Matrix coefficient kernel}, \weightedpin
{grh-compact/src/Lib/DiscreteSeries/CuspKernel/SchurOrthogonality.lean}
{108-L135}{Integrability and coefficient}, and \weightedpin
{grh-compact/src/New/Selmer/FitTwists.lean}{115-L139}{Common Hecke algebra}.}

For each good tuple, insert the tensor product of its finite type
idempotents. At a chosen split place $v$ over each $g_i$, its primitive
label gives a cuspidal representation $\tau$ of $\mathrm{GL}_r(k_v)$ by
Green's construction. Inflate $\tau$ to $K_v=\mathrm{GL}_r(\mathcal O_v)$,
extend its central character unitarily to $F_v^\times$, and compactly
induce from $F_v^\times K_v$. The resulting depth-zero representation
is supercuspidal. In the induced model a vector supported on the inducing
subgroup has its diagonal coefficient supported there, with restriction
to $K_v$ equal to the corresponding coefficient of $\tau$. Choose a unit
basis and use the contragredient coefficient where required to get
$\chi_\tau(u^{-1})$. Multiply by $1_{\mathcal O_v^\times}(\det u)$.
On $F_v^\times K_v$ this cuts out exactly $K_v$. Each type idempotent
at these two places, inflated as a convolution function with probability
Haar on $K_v$, is therefore a sum of truncated diagonal supercuspidal
coefficients times the common scalar $\dim\tau$. The split places are
distinct, being over distinct rational primes.\footnote{\href
{https://www.e-periodica.ch/cntmng?pid=com-001\%3A1996\%3A71\%3A\%3A10\#page=18}
{Moy--Prasad, Proposition 6.6, for the depth-zero compact induction},
and \href
{https://archive.intlpress.com/site/pub/files/_fulltext/journals/ajm/2004/0008/0004/AJM-2004-0008-0004-a019.pdf\#page=16}
{Harris--Labesse, \S1.9, for the truncated coefficients}.}

Every good form is cuspidal by the
\hyperlink{proof.finite-type-constant-term}{local constant-term test}.
The discrete-series vectors used in the real projection are smooth.
For each proper rational parabolic use a chosen split place with a
primitive label. A lattice basis adapted to the local flag makes
the compact part of the radical reduce to that of a proper finite
parabolic. The cuspidal Green type at this place has zero invariants
under this radical. The test gives a zero constant term on each good
tuple and hence on the good space.
Conversely the fixed $K_\infty$-finite lines of
the cuspidal spectrum are bounded by rapid decay, so their
principal-level vectors belong to the bounded automorphic model of
the given kernel. Thus the cuspidal trace computes the trace on the
full good space.\footnote{\href{https://arxiv.org/pdf/1106.2149\#page=12}
{Miller--Schmid, Theorem 2.24, for rapid decay}.}
The real unitary group has compact centre. By Clozel--Delorme,
a discrete-series representation \(\pi_v\) admits a compactly
supported smooth \(K\)-finite test with trace \(1\) on \(\pi_v\)
and \(0\) on every other basic representation. The basic tests
use normalized induction from limits of discrete series with
arbitrary complex unramified parameter. The characterization of
discrete trace functionals in \S5.2 then gives zero on every
proper parabolic induction, including the standard modules at
their Langlands parameters. Use the
Harish--Chandra parameter of the selected real row type for
\(\pi_v\), normalizing the test in the transported Haar.
The standard-module expansion of any irreducible character preserves
its infinitesimal character.
Thus a nonzero trace requires the designated infinitesimal character.
For the decreasing regular Harish--Chandra parameter $\lambda_v$, the
differences and parity of its coordinates make
$\lambda_v-\rho_{\rm roots}$ dominant integral. This is the highest
weight of the finite-dimensional algebraic coefficient system with
infinitesimal character $\lambda_v$. For $p,q>0$, restrict to
$SU(p,q)$, using the finite central covering by $SU(p,q)\times U(1)$.
The restriction of an irreducible unitary representation is irreducible
on $SU(p,q)$ since the center acts by a scalar. Its infinitesimal
character is regular integral on this derived group. By Salamanca-Riba's
Theorem~1.2 the factor has a dominant inducing weight in a compatible
positive root order. It vanishes on the Levi roots and is nonnegative
on the nilradical roots. Its shift by half the positive roots and
the shifted algebraic coefficient weight are dominant representatives
of the same regular infinitesimal character, so they coincide. The
Vogan--Zuckerman formula then gives nonzero cohomology with the
indicated coefficient dual. The
central circle acts trivially after tensoring with the dual, since
the coefficient weight has the same central character. Relative
cohomology therefore gives the same nonvanishing on $U(p,q)$.
The definite cases follow by compact representation theory.
Harris--Labesse, by Proposition~1.8.1, can lift
each contributing cuspidal representation to a cuspidal similitude
representation. In Theorem~3.1.5(B) the source is a cohomological
cuspidal representation of the unitary group with two
supercuspidal split components. It gives temperedness at split,
unramified spherical and archimedean places and discrete series
at the archimedean places.
Here the real cohomology was checked above and the two truncated
local coefficient tests force the two split supercuspidal
components by Lemma~1.9.1. The pseudo-coefficient trace on that spectrum is
one precisely for the specified member. The selected archimedean
kernel takes one line in each such discrete-series copy. The tests
give the finite Hecke traces on the prescribed spaces using the full
cuspidal spectrum over $F$.

Use group Haar measures from transported invariant top-degree forms and
transport Haar on the corresponding compact Cartan tori. Normalize the
pseudo-coefficients with respect to these measures. They have matching
stable integrals with the relative Kottwitz sign, by the real comparison
used in the proof of Theorem~3.1.6 of Harris and Labesse. That sign is one at each infinite
place here, since $r$ is odd and $pq$ is even for every signature $(p,q)$
of rank $r$. The finite parts match under the common frames, so the
cuspidal trace identity of their Proposition~1.10.6 applies to each
product of truncated coefficients. In the case of that proposition
used here, both tests transfer to the same test on the quasi-split
group, with the matching regular stable orbital integrals just checked.
In addition,
one split test is a truncated supercuspidal coefficient with a
supercuspidal Jacquet--Langlands transfer, a second split test is
cuspidal, and the tests are cuspidal at the nonisomorphic real
places. The first split transfer is the identity on the common
general linear group, the second primitive coefficient is
cuspidal, and the real pseudo-coefficients give the required
cuspidal tests. This yields equality of the cuspidal traces in
the transported Haar. Taking the sum over the coefficients
and over the good tuples gives equality of the traces of $p$ on the good
spaces. The argument also applies to $p^a$ for every $a\ge1$, so equality
of dimensions and Newton identities give equality of the determinants
of the common operators on the good spaces. This applies to the common
Hecke operators used in the single-lift comparison in Theorem~C.
\footnote{\href{https://www.numdam.org/item/ASENS_1990_4_23_2_193_0.pdf\#page=21}
{Clozel--Delorme, \S5.2, pp.~212--213 and Corollary},
\href{https://www.numdam.org/item/CM_1988__68_3_251_0.pdf\#page=4}
{Salamanca-Riba, Theorem 1.2, p.~253}, \href{https://www.numdam.org/item/CM_1984__53_1_51_0/}
{Vogan--Zuckerman, Theorem 5.5}, \href{https://www.numdam.org/article/CM_1979__39_1_11_0.pdf\#page=12}
{Shelstad, \S4} and \href{https://www.numdam.org/article/CM_1979__39_1_11_0.pdf\#page=33}
{Lemma 5.3 and \S6}, \href
{https://archive.intlpress.com/site/pub/files/_fulltext/journals/ajm/2004/0008/0004/AJM-2004-0008-0004-a019.pdf\#page=15}
{Harris--Labesse, Proposition 1.8.1}, \href
{https://archive.intlpress.com/site/pub/files/_fulltext/journals/ajm/2004/0008/0004/AJM-2004-0008-0004-a019.pdf\#page=19}
{Proposition 1.10.6}, and \href
{https://archive.intlpress.com/site/pub/files/_fulltext/journals/ajm/2004/0008/0004/AJM-2004-0008-0004-a019.pdf\#page=27}
{Theorems 3.1.5--3.1.6 and Proposition 3.1.7}.}

\paragraph{The dimension bound from the definite inner form.}
\teamanchor{proof.weighted-good-mass}%
Take $p=1$ to carry $R$ to the good space of the definite $W_m$. It is
bounded by the full principal-level dimension there. Let
$\operatorname{mass}_r=\sum_a|\Gamma_a|^{-1}$ be the mass of the genus of
the standard rank-$r$ lattice at $K_r$. Refining each class at $K_r$ to
$K(g)$ acts on its $|G|$ level choices by $\Gamma_a$. The action is free
by the principal-level stabilizer bound. Hence the number of new classes
is $|G|\operatorname{mass}_r$. At each class the value space is the full
algebraic representation of the compact real group with the common
parameter. For the local parameters $\lambda_{v,1}>\cdots>\lambda_{v,r}$
written without the factor two, its dimension is
$\prod_{a<b}(\lambda_{v,a}-\lambda_{v,b})/\prod_{a=0}^{r-1}a!$ by Weyl's
dimension formula.

\paragraph{The compact-dual degree in the filling Haar.}
\teamanchor{proof.weighted-compact-dual-degree}%
Here the real filling Haar is the compact-dual normalization in
the formal-degree formula. At a real place write the group as
\(U(J)\) with \(J=\operatorname{diag}(I_p,-I_q)\), \(p+q=r\).
The differential of the Cayley chart
\(Y\mapsto(1+J^{-1}Y/2)(1-J^{-1}Y/2)^{-1}\)
identifies the filling density with independent-entry Lebesgue measure
on \(JX\), divided by \(\operatorname{vol}U(r)\) in those coordinates.
The \hyperlink{proof.compact-unitary-coordinate-volume}{compact unitary
volume calculation} gives \(\operatorname{vol}U(r)=
\prod_{j=1}^r 2\pi^j/(j-1)!\). It also shows that the tangent identification
with the compact dual preserves the absolute independent-entry
density. Thus the filling Haar corresponds to compact probability
Haar. If \(p=0\) or \(q=0\) the group is already compact and Schur
orthogonality gives the degree equal to the representation dimension.
\footnote{\weightedpin{grh-compact/src/Lib/DiscreteSeries/Haar/CompactDualMeasure.lean}
{22-L48}{The filling Haar in the Cayley chart}.}

For \(p,q>0\), use the semisimple group \(S=SU(p,q)\) with
\(K_S=S(U(p)\times U(q))\). On \(S\) start with Haar induced by
the metric volume of \(S/K_S\) and probability Haar on \(K_S\).
Use the corresponding dual metric on
\(X^c=SU(r)/K_S\) with the same probability Haar on \(K_S\).
The discrete atom in the Plancherel formula is the Schur formal
degree for the same Haar measure. The formula in these measures gives
\(\dim\tau/\operatorname{vol}(X^c)\) for a discrete series
with the infinitesimal character of the finite-dimensional
representation \(\tau\). The same quotient disintegration
gives total mass \(\operatorname{vol}(X^c)\) on the compact
group. Divide both Haar measures by this mass. The compact Haar
is now a probability measure and the formal degree on \(S\) is
\(\dim\tau\).
\footnote{\href{https://arxiv.org/pdf/math/0009039\#page=8}{Olbrich,
\emph{\(L^2\)-invariants of locally symmetric spaces}, Section~2}
and \href{https://arxiv.org/pdf/math/0009039\#page=15}{Corollary~5.2}.}

The maps \(SU(p,q)\times U(1)\to U(p,q)\) and
\(SU(r)\times U(1)\to U(r)\) by scalar multiplication have the
same kernel \(\{(\zeta^{-1}I_r,\zeta):\zeta^r=1\}\).
Push forward the measures just compared, using probability Haar
on each central circle. Both maps have the same local determinant
under the tangent identification and have \(r\) sheets, so the
comparison is unchanged. The measure on \(U(p,q)\) is the filling
Haar by the Cayley calculation. An irreducible unitary
representation stays irreducible on \(SU(p,q)\) since the central
circle acts by scalars. Its matrix coefficient norm squared is
unchanged in the product integral, so the normalized formal
degree remains \(\dim\tau\). In this application \(\tau\) is the
restriction of the compact \(U(r)\) representation of highest
weight \(\lambda_v-\rho\), extended algebraically to
\(\mathrm{GL}_r(\mathbf C)\) before restriction. It has the required
infinitesimal character on \(SU(p,q)\). The weighted Bessel system
takes \(\deg_v\) to be the formal degree of this real row type.
We consequently obtain exactly the same Weyl dimension in the
orbital energy and the definite row count. The finite factor
\(|G|\operatorname{mass}_r\) still comes from the free refinement
of classes, with the residues in counting measure.
\footnote{\weightedpin{grh-compact/src/New/Filling/FillingData.lean}
{56-L76}{The real row parameter and formal degree} and
\weightedpin{grh-compact/src/New/Filling/FillingData.lean}
{94-L108}{The Bessel system using this degree}.}

Thus, with $\deg_{\rm Bes}=\prod_v\deg_v$,
\begin{equation}
 \label{eq:weighted-row-mass}
    R\le |G|\deg_{\rm Bes}\operatorname{mass}_r.
\end{equation}

\paragraph{Orbital energy and the fixing sets.}
\teamanchor{proof.weighted-orbital-energies}%
For a set $\Omega\subseteq G$ and a set $T$ of rational terms stable
under rational unipotents on both sides,
write $\mathscr E_T(\Omega)=\nu^{-2}\sum_{s\in\Omega,c,c'}
\|A_T(c,s,c')\|_{\rm HS}^2$. This is a sum in counting measure. In
particular $\mathscr E_{\rm all}(\Omega)=\sum_{s\in\Omega}\|K(s)\|_{\rm HS}^2$.
Put $\mathscr M=C^2/(\deg_{\rm Bes}\operatorname{mass}_r)$,
$\mathscr D=|J|C$ and $\tau=b^{kf/4}$. Write
\(A^{\rm cut}_T=A_T^\flat\) for the same blocks with the cut
Bessel vectors in place of the full vectors. The link sum uses
these cut vectors, and \(\mathscr E\) uses the full vectors. The
field contains the $5040$-th roots of unity since $5040\mid N_1$. The
admissibility size ratios below are those of the filling datum attached
to $D$.\footnote{\weightedpin{grh-compact/src/New/Filling/FillingAdmissible.lean}
{39-L81}{Admissibility on the constructed datum}.}

For a subset $A$ of the chosen places above $g_0$, prescribe at each
$w\in A$ a nonzero vector of the Bessel frame span, in either the
standard module at $w$ or the contragredient module at $\bar w$.
Write $F_A\subseteq G$ for the subgroup fixing the prescribed vectors
and $\Omega_A$ for the portion of $F_A$ having all leading frame minors
of sizes $1,\ldots,\varrho+1$ nonzero at one and the same chosen
place above $g_0$. The last minor is the determinant of the augmented
frame pairing. The open-orbit parametrization uses it to give a
nonzero terminal pivot. When
$A$ is the full set take $\Omega_A=F_A$. Define
$\mathscr M_A=b^{-6k|A|}\mathscr M$. For \(X\) the span of the
isotropic frame, the link of \(\gamma\) is the image of
\(X^\perp\cap(\gamma X)^\perp\) in \((X^\perp/X)^2\) under
\(y\mapsto([y],[\gamma^{-1}y])\). The frame discriminant is
\(|\det Z|^2-N^2|\det Z_\varrho|^2\) for the frame-pairing
matrix \(Z\) and its leading block \(Z_\varrho\).
The link energy squares the sum within one link at fixed
\((c,s,c')\) before summing over these labels and links.
\footnote{\weightedpin{grh-compact/src/New/Filling/EichlerFlag.lean}
{39-L63}{Frame discriminant and link}, and
\weightedpin{grh-compact/src/New/Filling/SchurTest.lean}
{17-L25}{Counted energy and link energy}.}

\paragraph{Summing the finite denominator shells.}
\teamanchor{proof.weighted-finite-denominator-shells}%
Enlarge the finite set of places to a conjugation-stable set containing
the support of \(g\) and the denominators of the finite packet under
consideration. Let \(A\) be the product of these completions and \(R\)
its ring of integers. For the restricted character
\(\psi=\psi_f(\,\cdot\,/\varpi)\), the annihilator of \(R\) is \(R\).
Taking the inverse character gives the same annihilator.
Thus the normalization \(\operatorname{vol}(R)=1\) is self-dual,
and the finite Plancherel constant is one. Write
\(\operatorname{md}(c)\) for the modulus of multiplication by \(c\).
Put \(n=\varrho+1\). Write \(M\) for the group Haar measure of the
unitary level, \(K\) for its integral group, \(K(g)\) for its
principal congruence group and \(Z\) for the intermediate subgroup
imposing the prescribed residue conditions. Put
\[
 v_U=\operatorname{vol}_U(U\cap K(g)),\qquad
 M_Z=M(Z)M(K(g))/v_U^2 .
\]
Here the measures on the successive unipotent coordinates are those
of the finite unfolding. The finite vector has one coordinate for
each right coset \(K(g)a\) in \(Z\). At \(y\) that coordinate is
\(\chi_g(ya^{-1})\), where \(\chi_g(uk)=\psi(\sigma(u)/g)\)
for \(u\in U,\ k\in K(g)\), extended by zero elsewhere.
In particular there is no factor \(1/[Z:K(g)]\) in its inner
product. Indeed on \(uy\in Z\) there is a unique coset containing
\(uy\). Integrating first over the \(u\)'s giving that coset,
each nonempty fibre has measure \(v_U\), and the two characters
give
\[
 v_U\langle W_Z(y),W_Z(yx)\rangle
   =\int_U 1_Z(uy)\chi_g\bigl((uy)x(uy)^{-1}\bigr)\,du .
\]
The sign here uses the inner product linear in its second argument.

Let \(dM_0^+\) be the upper measure at floor zero. The ordered
factorization is \(dM=du*dM_0^+\), with \(du\) the unipotent
Haar. At a point \(y\) put \(B(y)=\int_U1_Z(uy)\,du\). Whenever
this fibre is nonempty, its translate in \(U\) is contained in
\(U\cap K\), whose volume is one in the integral unipotent
coordinates. Thus \(0\le B(y)\le1\) and
\(\int B(y)\,dM_0^+(y)=M(Z)\).
The coset formula bounds the pairing by \(B(y)/v_U\).
If it pairs nontrivially with \(W_Z(y')\), then \(y'\in U zK(g)\)
for any chosen \(z\in Uy\cap Z\). Indeed the common nonzero
coset coordinate gives two representatives differing on the right
by \(K(g)\), which is normal in \(K\). Unfolding \(M(zK(g))\)
makes the measure of \(U zK(g)\) in \(M_0^+\) equal to
\(M(K(g))/v_U\). Since \(B(y)^2\le B(y)\), the pairing norm is
\[
 \mathcal N(W_Z,W_Z)=
 \iint|\langle W_Z(y),W_Z(y')\rangle|^2\,dM_0^+(y')\,dM_0^+(y)
 \le \frac{M(Z)M(K(g))}{v_U^3}.
\]
Consequently the common restricted integral uses
\(v_U\mathcal N(W_Z,W_Z)\le M_Z\).
\footnote{\weightedpin{grh-compact/src/New/Filling/StandardLevel.lean}
{218-L226}{Unipotent coordinate volumes and unfolding},
\weightedpin{grh-compact/src/New/Filling/TopFloor.lean}
{102-L109}{The common coset} and
\weightedpin{grh-compact/src/New/Filling/TopFloor.lean}
{146-L205}{Coset pairing and upper-measure normalization}. The fixed source also records \weightedpin{grh-compact/src/New/Filling/TopFloor.lean}{134-L198}{the pairing of the unnormalized coset vector}.}

\leavevmode\teamanchor{proof.weighted-position-shell-inputs}%
Use three estimates on the same finite or complex local floor and
the same pair of head gauges. The
\hyperlink{proof.weighted-restricted-position-integral}{restricted
positive-weight integral} is proved after the scalar-volume bound.
The capped estimate and the reconstruction from the dual coordinates
are proved here. The numerical applications to the two local measures
follow in their respective passages.

\paragraph{The absolute orbital majorant with a capped coordinate volume.}
\teamanchor{proof.weighted-capped-orbital-majorant}%
Use the local level at the finite product or at \(\mathbf C\), with
level scalar \(g\), box \(K\), vector ball \(\mathcal B\), gauge
constant \(\Lambda\ge2\), and volume exponent \(s\). Write \(n=\varrho+1\).
Number the hyperbolic floors from \(0\) to \(\varrho-1\).
For \(l\le\varrho\), \(H_l\) fixes the first \(l\) hyperbolic pairs.
After the marked floor the remaining group is the core group
\(H_{\varrho+1}\), and \(H_{\varrho+2}=\{1\}\).
The group \(Q_l\) acts upper unitriangularly on the first \(l\)
frame vectors, with \(Q_{\varrho+2}=U\). The
earlier unipotent radical is \(\mathcal N_{<l}\), with
\(Q_l=H_l\mathcal N_{<l}\). Nearness at \(l\) means \(Q_l y\cap K\ne
\varnothing\), and is left \(Q_l\)-invariant. The groups \(Q_l\) and
\(H_l\) decrease with \(l\), and \(U,H_l\subseteq Q_l\).
The Weyl elements \(w_p,w_q\) of two heads of length \(j\le\varrho\)
commute with \(H_j\). The ratio \(d=w_q^{-1}w_p\) is diagonal on
the earlier hyperbolic pairs and normalizes \(U\).

Fix an allowed position \(x\). For \(a<\varrho\) let \(L_a(y,x)\)
be the scalar values of the dual coordinate \(h(e_a,w)\) on the
vectors \(w\) with \(h(e_b,w)=0\) for \(b<a\) and
\[
 y^{-1}w\in\mathcal B,\qquad
 (y^{-1}w-(yx)^{-1}w)/g\in\mathcal B .
\]
Use the local additive Haar on these scalar sets. Put
\(r_a(y,x)=\operatorname{vol}L_a(y,x)\) and
\(P_l(y)=\prod_{a<l}r_a(y,x)^s\). Each scalar set contains a
neighbourhood of zero and lies in the image of a compact set.
Indeed the two ball conditions are open, contain zero, and the
first bounds \(w\) in \(y\overline{\mathcal B}\); projection onto the
dual coordinate is open on the subspace with the earlier coordinates
zero. The same projection shows that the condition
\(c\in L_a(y,x)\) is open in \((y,x,c)\), so \(r_a\) is
measurable. The first \(a\) dual-zero conditions and the \(a\)-th
dual coordinate are preserved by \(Q_{a+1}\) on that subspace.
Changing \(w\) by this group element preserves the two pullbacks
after changing \(y\) by the same left multiplication. Thus \(r_a\)
is left \(Q_{a+1}\)-invariant, \(0<P_l<\infty\), and
\[
 P_l(\tau y)=P_l(y)\quad(\tau\in Q_l),\qquad
 P_{l+1}(y)=P_l(y)r_l(y,x)^s\quad(l<\varrho).
\]

Write \(\theta_l\) for the actual floor measure, supported on \(H_l\),
and \(M_j=(\eta_{j-1}*\theta_{j-1})*\cdots*(\eta_0*\theta_0)\)
for the preceding stack, where \(\eta_l\) is the Heisenberg measure.
Convolutions are read with the highest floor on the left.
Let \(W\) be a bounded measurable cut Bessel vector of frequency
\(1/g\), so \(W(y)\ne0\) implies \(Uy\cap K\ne\varnothing\).
Let \(W'\) be bounded and measurable with the same character on
\(U\), supported where \(y\) is near at \(j\). Use the two kernels
\[
 k(y,z)=\langle W(y),W(z)\rangle,\qquad
 k'(y,z)=\langle W(y),W'(dz)\rangle .
\]
For either kernel \(b\) put
\[
 \begin{aligned}
 b_l(y,z)&=\int b(\tau y,\tau z)\,
             d(\theta_{\varrho+1}*\cdots *\theta_l)(\tau),\\
 m_l^b(y)&=\int |b_{l+1}(\tau y,\tau yx)|\,d\theta_l(\tau).
 \end{aligned}
\]
The empty convolution has unit mass at \(1\). In the weighted
square sums defined in the
\hyperlink{proof.weighted-one-floor-plancherel}{floor calculation},
write \(A_l^b(y)=P_l(y)|b_{l+1}(y,yx)|^2\),
\(B_l^b(y)=\int(A_l^b+B_{l+1}^b)(\tau y)\,d\theta_l(\tau)\)
for \(l\le\varrho\), and \(B_n^b=0\). Let \(Z_l^b\) integrate
\(B_l^b\) against \(\theta_{l-1}*\cdots*\theta_0\).
With \(C_\Gamma=\Lambda^{10^6n^3}\) and
\(J_j=\prod_{a<j}\operatorname{fl}_a(p_a^{-1})\), the square
product is \(\mathcal D_j(x)=C_\Gamma^3 J_j^2 Z_0^k Z_j^{k'}\).
The positive majorant and the active set at \(x\) are
\[
 \begin{split}
 \mathcal F_j^+(x)&=\int m_j^k(y)m_j^{k'}(w_p^{-1}y)\,dM_j(y),\\
 \mathcal U_j(x)&=\{y:k(y,yx)\ne0,\ w_p^{-1}y\text{ near at }j,\
                                w_q^{-1}(yx)\text{ near at }j\}.
 \end{split}
\]
For the weighted floor form with multiplier
\(\omega(\sigma\tau^{-1})\) of modulus at most one, use \(W'=W\).
At fixed \(y,x\) the two integrations against
\(\theta_j\) have upper arguments \(\sigma y,\sigma yx\) and
\(\tau w_p^{-1}y,\tau w_p^{-1}yx\), because the heads commute
with \(H_j\). The triangle inequality and Tonelli bound the modulus
by \(m_j^k(y)m_j^{k'}(w_p^{-1}y)\). The absolute integral over
\(dM_j\,dM\), times the head density \(\kappa_j\), therefore
majorizes the form whenever the exterior integral is finite.
The real mixed application uses the displayed positive majorant
directly with its choice of \(W'\).

Suppose that for each \(y\) the \(\theta_j\)-mass of
\(\{\tau:\tau y\text{ is near at }j+1\}\) is at most \(V\ge1\).
Take \(V\) finite in the estimates below. Either local level
has a finite bound from the size-seven argument, so an infinite
mass bound can be replaced by that finite bound when needed.
For \(0\le\sigma<1\), assume every \(y\in\mathcal U_j(x)\)
has an \(a\) with \(j\le a<\varrho\) and \(r_a(y,x)\le\sigma\).
For \(\sigma\ge1\) there is no extra volume condition. Then
\begin{equation}
 \label{eq:weighted-capped-position}
 \mathcal F_j^+(x)^2\le \sigma^s V^2\mathcal D_j(x).
\end{equation}
The majorant vanishes if \(\mathcal U_j(x)\) is empty.
This covers the case \(\sigma=0\) under the cap hypothesis,
and the case \(j=\varrho\) when that hypothesis is imposed.

\smallskip
\noindent\textit{The measured one-floor bounds.}
\leavevmode\teamanchor{proof.weighted-shell-floor-mass}%
The map \(\tau\mapsto\tau^{-1}e_l\) takes \(\theta_l\) to the
quadric measure at floor \(l\le\varrho\) used in the measured
floor coordinates. Here \(e_\varrho\) denotes the last marked
vector. On a lower floor the Weyl parameter \(c\) gives the
last hyperbolic coordinate and scales the \(d_l\) horizontal
coordinates by \(c\) up to conjugation and signs. Its density
\(\operatorname{md}(c)^{d_l}dc\) is the Jacobian of this change.
The skew coordinate has the translated skew Haar of the quadric
chart. The terminal floor measure has this pushforward by its
definition. The \(H_l\)-action preserves both the ambient additive
Haar and the Hermitian value. It therefore takes the continuous
family of quadric measures to another disintegration for the
same norm. Uniqueness almost everywhere, continuity against
compactly supported tests and full support of the skew measure
give equality at the fixed Gram value, including zero. The stabilizer of \(e_l\) in
\(H_l\) is \(H_l\cap Q_{l+1}\).
The left quotient by this closed stabilizer is a standard
Borel space. The orbit map induces an injective Borel map on
the quotient with Borel inverse on its image. Thus invariant
measurable functions of \(\tau\) descend measurably to the
quadric orbit point.
In particular, if a measurable function is invariant on the left
by \(Q_{l+1}\), its integral at \(\tau y\) against \(\theta_l\)
is invariant when \(y\) is replaced by \(\zeta y\) with
\(\zeta\in Q_l\). To see this, factor \(\zeta\) in
\(H_l\mathcal N_{<l}\), conjugate the unipotent factor to the
left, and use that \(\mathcal N_{<l}\subset Q_{l+1}\) is
normalized by \(H_l\). The \(H_l\) factor on the right changes
the orbit point by its inverse action and preserves its measure.
The same argument works with invariance under \(Q_{l+1}\cap H_b\)
and \(\zeta\in Q_l\cap H_b\) for \(b\le l\), using the factorization
within \(H_b\). At \(l=\varrho+1\) it is the invariance of the
core-group measure. These identities apply to integrable complex
functions by taking their real and imaginary parts.

For \(l<\varrho\) the inverse moment needed at one floor is
\begin{equation}
 \label{eq:weighted-shell-one-floor-inverse}
 \int_{\{\tau:\tau y\ {\rm near\ at}\ l+1\}}
        r_l(\tau y,x)^{-s}\,d\theta_l(\tau)
 \le \Lambda^{10^4 n}.
\end{equation}
If \(y\) is not near at \(l\), the integration set has zero
measure. Otherwise take \(\zeta y=y_0\in K\) with \(\zeta\in Q_l\)
and use the preceding invariance, since the near indicator and
the reciprocal volume are left \(Q_{l+1}\)-invariant. If
\(\xi\tau y_0\in K\) with \(\xi\in Q_{l+1}\), the projection of
\(y_0(\xi\tau y_0)^{-1}e_l\) to the remaining floor space is
\(\tau^{-1}e_l\). Indeed \(\xi^{-1}e_l-e_l\) lies in the earlier
frame span and \(\tau\in H_l\) fixes that span. Both box
matrices have size one. Inversion adds four and multiplication
adds the two sizes and one, putting the orbit point in size seven.

Use the
\hyperlink{proof.weighted-inverse-weight}{inverse-volume moment}
at \(y_0\) with no extra row equations. Its set \(D\) consists of
orthonormal coordinates \(v\) for which the pullbacks of
\(E_0v\) at \(y_0,y_0x\) are a close pair, where \(E_0\) embeds
the orthonormal part. If \(J\) is the split Gram matrix, set
\(u=E_0^*J\tau^{-1}e_l\), the orthonormal projection of the
orbit point. The vector \(\tau E_0v\) has the required earlier
dual coordinates zero and its two pullbacks are unchanged.
Its \(l\)-th dual coordinate is \(h(\tau^{-1}e_l,E_0v)=u^*v\).
Therefore
\(\{u^*v:v\in D\}\subset L_l(\tau y_0,x)\).
The inverse volume to the power \(s\) is bounded by
\(\phi(u)^s\), with a measurable inverse-volume majorant
\(\phi\) supplied by that moment. On the size-seven quadric the
paired-coordinate fibre integral costs \(\Lambda^{1287n}\) and
the orthonormal coordinate has size at most eleven. Integrating
\(\phi(u)^s\) costs at most \(\Lambda^{1188n}\). Their product is
at most \(\Lambda^{10^4n}\), proving
\eqref{eq:weighted-shell-one-floor-inverse}.
The same size-seven argument with integrand one, including
the terminal floor, bounds the \(\theta_l\)-mass of a translated
near set by \(\Lambda^{1287n}\operatorname{vol}(\text{size }11)^{d_{\rm orth}}\).
Here \(d_{\rm orth}\) is the number of orthonormal coordinates.

\smallskip
\noindent\textit{Supports, squares and the downward induction.}
The last mass bound and the mass at most one on the core
factor over the cut support justify the upper averages. More
explicitly, for a product of upper-floor factors to put the
first argument of \(b\) in \(UK\), the partial product through
each floor \(l\) must put it in the near set at \(l+1\), since
the remaining factors are in \(H_{l+1}\subset Q_{l+1}\).
Successive integration bounds the measure of this set by the
product of the floor mass bounds. Boundedness of \(W,W'\) now
gives absolute convergence. A nonzero upper average or pivot
therefore yields some point of the corresponding \(H_l\)-orbit
where the unaveraged kernel is nonzero. In particular
\(b_{l+1}(\tau y,\tau yx)\) vanishes off the near-\((l+1)\)
set. Since \(P_l(\tau y)=P_l(y)\),
\[
 \int |b_{l+1}(\tau y,\tau yx)|^2 d\theta_l(\tau)
       \le B_l^b(y)/P_l(y).
\]
Cauchy--Schwarz on the near set at \(j+1\) thus gives
\[
 (m_j^b(y))^2\le V B_j^b(y)/P_j(y).
\]

For \(l\ge j\) the simultaneous left action of \(Q_l\cap H_j\)
preserves \(b_l\). At the unaveraged top this follows from the
two Bessel characters, which cancel because the ratio \(d\)
commutes with \(H_j\). Descending induction uses the
right-invariance identity for the floor measures.
For \(\nu\in\mathcal N_{<l}\) the simultaneous left action on
\(b_{l+1}\) only introduces a unit character. To check this,
move \(\nu\) through each \(\tau\in H_{l+1}\) in the upper
average. The conjugate stays in \(\mathcal N_{<l}\) and has
the same simple-root entries. Indeed \(H_{l+1}\) fixes the
frame through the entry indexed by \(l\), and the later
simple-root entries of this radical are zero. For the translated
kernel the second conjugate is also conjugated by the diagonal
ratio \(d\), which multiplies each simple-root entry by its fixed
ratio. Both characters are consequently independent of \(\tau\).
Factoring \(Q_{l+1}\subset Q_l=H_l\mathcal N_{<l}\) now gives
left \(Q_{l+1}\)-invariance of \(A_l^b\). Its support lies
in the near set at \(l+1\). The definition of \(B_l^b\) and the
right-invariance identity for \(\theta_l\) show by downward
induction that \(B_l^b\) is invariant on the left by \(Q_l\)
and zero away from the near set at \(l\).

Put \(C_1=\Lambda^{10^4n}\). Suppose some \(l\ge j\) and some
base \(y\) have the following property. For each
\(\tau\in H_l\) with \(k(\tau y,\tau yx)\ne0\), a coordinate
at an index \(a\) with \(l\le a<\varrho\) has volume at
most \(\sigma\). For \(0<\sigma<1\) we claim
\begin{equation}
 \label{eq:weighted-capped-pivot}
 m_l^k(y)^2
   \le \sigma^s C_1^{\varrho-l} B_l^k(y)/P_l(y).
\end{equation}
At \(l=\varrho\) the hypothesis makes the raw kernel zero on
the whole \(H_\varrho\)-orbit. Thus its pivot is zero.
Suppose \(l<\varrho\) and write \(z=\tau y\) for
\(\tau\in H_l\). When \(r_l(z,x)\le\sigma\),
\[
 |k_{l+1}(z,zx)|^2
      \le \sigma^s A_l^k(z)/P_{l+1}(z).
\]
When \(r_l(z,x)>\sigma\), that volume is unchanged along
the \(H_{l+1}\)-orbit of \(z\). The cap hypothesis on the
\(H_l\)-orbit therefore forces a later small coordinate whenever
the raw kernel is nonzero on this smaller orbit. The triangle
inequality \(|k_{l+1}(z,zx)|\le m_{l+1}^k(z)\) and the induction
hypothesis bound the square by
\(\sigma^s C_1^{\varrho-l-1} B_{l+1}^k(z)/P_{l+1}(z)\).
In both cases it is bounded by
\[
 r_l(z,x)^{-s}\,
 \frac{\sigma^s C_1^{\varrho-l-1}}{P_l(y)}
                  \bigl(A_l^k(z)+B_{l+1}^k(z)\bigr).
\]
Cauchy--Schwarz on the near-\((l+1)\) set uses
\eqref{eq:weighted-shell-one-floor-inverse} in its first factor
and the defining integral of \(B_l^k(y)\) in its second.
This proves \eqref{eq:weighted-capped-pivot}.

For the product of the two pivots, take a \(y\) where it is
nonzero. The second pivot gives \(\tau\in H_j\) with
\(W(\tau w_p^{-1}y)\) and \(W'(\tau w_q^{-1}yx)\) both nonzero.
The cut condition for \(W\) and the near-\(j\) support of \(W'\)
imply both translated near conditions. Since the Weyl elements
commute with \(H_j\), they hold also with any \(\rho_0\in H_j\)
in place of \(\tau\), by multiplying on the left by
\(\tau\rho_0^{-1}\in Q_j\). A nonzero first pivot then gives
an active point by its raw kernel witness. This proves the
vanishing assertion and treats \(\sigma=0\) directly, even for
a vacuous mass bound \(V=\infty\). If \(0<\sigma<1\), the
first pivot satisfies \eqref{eq:weighted-capped-pivot}, with
\(C_1^{\varrho-j}\le C_\Gamma\). The second pivot costs \(V\).
For \(\sigma\ge1\) use the bound \(V\) on each pivot.
In both branches, using \(V\ge1\), the square of the product is
bounded by
\[
 \sigma^s V^2 C_\Gamma\,
       \frac{B_j^k(y)}{P_j(y)}
       \frac{B_j^{k'}(w_p^{-1}y)}{P_j(w_p^{-1}y)} .
\]

\smallskip
\noindent\textit{Cauchy--Schwarz in the two Weyl coordinates.}
\leavevmode\teamanchor{proof.weighted-shell-two-weyl}%
Write \(I_j(y)=1/P_j(y)\) on the near set and zero elsewhere.
On the supports of the two weighted square sums the last
product is unchanged on replacing its reciprocals by \(I_j\).
Cross the reciprocal weights in Cauchy--Schwarz with \(dM_j\).
The resulting two integrals are
\[
  \int B_j^k(y)I_j(w_p^{-1}y)\,dM_j(y),\qquad
  \int B_j^{k'}(w_p^{-1}y)I_j(y)\,dM_j(y).
\]
Their estimates use \(M_j=\eta_{<j}*(\theta_{j-1}*\cdots*\theta_0)\),
by moving each new floor measure past the earlier Heisenberg
convolution, which is invariant under conjugation by that floor
group. The first integral uses the written
\hyperlink{proof.weighted-inverse-weight}{preceding inverse-weight bound}
with the head \(\bar p\), since \(w_p^{-1}=w_{\bar p}\) and
conjugation preserves all local moduli. The cost is
\(C_\Gamma J_j Z_j^k\).

In the second use \(y=w_p z\) in the integral of the left
\(H_j\)-invariant function \(B_j^{k'}(z) I_j(w_p z)\).
For this change of variables, write the full Haar as
\(M=S_j*M_j\) with the upper stack \(S_j\) on \(H_j\).
That stack is right \(H_j\)-invariant. Induct from the core
group. At a lower floor the higher stack convolved with the
current Heisenberg measure is right invariant under the stabilizer
\(Q_{l+1}\cap H_l=H_{l+1}\mathcal N_l\), by conjugation invariance
of the Heisenberg Haar under \(H_{l+1}\). The right-invariance
identity for \(\theta_l\) now gives the assertion on \(H_l\).
At the marked floor the stabilizer is just \(H_{\varrho+1}\).
For any positive integrable function \(a\) on the group, put
\(A(z)=\int a(hz)\,dS_j(h)\). For a nonnegative left
\(H_j\)-invariant function \(\Phi\), unfolding gives
\(\int\Phi\,dM_j=\int\Phi\,a/A\,dM\), as \(A\) is positive and
finite \(M_j\)-almost everywhere. In this identity a left
translation by \(w_p\) on Haar replaces \(a\) by \(a(w_p\,\cdot)\)
and its denominator by \(A(w_p\,\cdot)\), since \(w_p\) commutes
with \(H_j\). This proves invariance of the integral of such
functions under the Weyl change. The inverse-weight estimate
on the second integral now costs \(C_\Gamma J_j Z_j^{k'}\)
with head \(p\). Finally \(Z_j^k\le Z_0^k\): integrating the
recursion for \(B_l^k\) against the preceding \(\theta\)-stack
gives \(Z_l^k\) as \(Z_{l+1}^k\) plus a nonnegative piece.
The two costs and the pointwise factor give
\eqref{eq:weighted-capped-position}.
\footnote{\weightedpin{grh-compact/src/New/RealPlace/LevelGauge.lean}
{63-L95}{Near sets, coordinate weights and the floor orbit},
\weightedpin{grh-compact/src/Lib/DiscreteSeries/Bruhat/BigCell.lean}
{434-L440}{Measurable descent through the orbit point},
\weightedpin{grh-compact/src/Lib/General/Measure/Place.lean}
{123-L131}{Continuous quadric disintegration},
\weightedpin{grh-compact/src/New/RealPlace/LevelGauge.lean}
{157-L194}{Invariance and factorization of the stack},
\weightedpin{grh-compact/src/New/RealPlace/WeightPlancherel.lean}
{20-L52}{The size-seven fibre estimates},
\weightedpin{grh-compact/src/New/RealPlace/LocalTheorem.lean}
{98-L115}{The descending cap},
\weightedpin{grh-compact/src/New/RealPlace/LocalTheorem.lean}
{347-L380}{Crossed Cauchy--Schwarz and the two translated supports},
and \weightedpin{grh/src/New/Spanning/Frames/Gains.lean}
{335-L421}{The floor-mass and small-volume arguments in the readable edition}.}

\smallskip
\noindent\textit{Coincident pivots at the bottom.}
\leavevmode\teamanchor{proof.weighted-bottom-cap}%
At \(j=0\) with \(W'=W\), both heads are empty and \(M_0\)
is the point mass at \(1\). Thus \(\mathcal F_0^+(x)=m_0^k(1)^2\),
\(P_0(1)=1\), \(Z_0^k=B_0^k(1)\) and
\(\mathcal D_0=C_\Gamma^3 B_0^k(1)^2\).
Here all points are near at \(0\), since \(H_0=Q_0\) is the
whole group and the box is nonempty. If every raw nonzero
pair at position \(x\) has a coordinate of volume at most
\(e<1\), apply \eqref{eq:weighted-capped-pivot} once to their
common pivot and square the bound. This gives
\begin{equation}
 \label{eq:weighted-coincident-pivot-cap}
 \mathcal F_0^+(x)^2
 \le e^{2s}\Lambda^{2\cdot 10^4n\varrho} B_0^k(1)^2
 \le e^{2s}\mathcal D_0(x).
\end{equation}
At \(e=0\) the raw support is empty and the same conclusion
follows by vanishing. This is the input on the exterior
of the principal coset in the
\hyperlink{proof.weighted-finite-bottom}{finite bottom integral}.
\footnote{\weightedpin{grh-compact/src/New/Filling/TopFloor.lean}
{398-L415}{The coincident pivot on the complementary set}. The fixed source also records \weightedpin{grh-compact/src/New/RealPlace/LocalTheorem.lean}{99-L124}{the one-pivot cap}, and
\weightedpin{grh-compact/src/New/RealPlace/LocalTheorem.lean}{365-L383}{the crossed capped estimate}.}

\paragraph{Reconstruction from the dual coordinates.}
\leavevmode\teamanchor{proof.weighted-triangular-reconstruction}%
Work over the local coefficient ring \(A\) with its involution. Order the
split basis as \(e_0,\ldots,e_{\varrho-1}\), the middle block, and
\(e^*_{\varrho-1},\ldots,e^*_0\). The Gram matrix \(J\) pairs each
\(e_a\) with \(e_a^*\) with value one and has no other entries on these
rows. Write \(n=\varrho+1\) and \(r=\dim_A V\). The Bessel group \(U\)
is block upper unitriangular in this order. The group \(Q_j\) sends
each \(e_a\), \(a<j\), to itself plus a combination of earlier \(e_b\)'s.
It contains \(U\). For \(h\in Q_j\) the row at \(e_a^*\), \(a<j\), of
\(h\) differs from the identity row only in the columns \(e_b^*\),
\(b<a\). Indeed \(Jh=(h^{-1})^*J\), and \(h^{-1}e_a\) has the
stated upper-triangular form. Consequently \(h\) preserves the
\(a\)-th dual coordinate of a vector whose earlier dual coordinates
vanish.

In a gauge, products add sizes, a sum of at most \(r+1\) entries
of size \(i\) has size \(i+1\), and negation and conjugation preserve
size. The elements \(0,1\) belong to every scalar size set.
Multiplication by \(1\) gives monotonicity in the size.
All matrix and vector sizes below are entrywise. Thus a matrix
product or matrix--vector product costs one additional size beyond
the sum of the two sizes. Since \(J,J^{-1}\) have size one,
\(K^{-1}=J^{-1}K^*J\) has size at most \(i+4\) when \(K\) is a
unitary matrix of size \(i\).

First consider a head \(p\) of length \(j\le\varrho\) in a gauge \(S\).
This means \(p_a\in A^\times\) for \(a<j\) and the inverse of
\(P_l=\prod_{a<l}p_a\) has size one for each \(l\le j\).
The Weyl element \(w_p\) sends \(e_a\) to \(p_a e_a^*\) and
\(e_a^*\) to \(\bar p_a^{-1}e_a\) on these pairs.
Suppose that \(s z\) and \(h w_p^{-1}z\) lie in the box, with
\(s,h\in Q_j\). Their quotient \(s w_p h^{-1}\) has size seven in \(S\).
In its first \(a+1\) primal columns, for \(a<j\), right multiplication
by \(h^{-1}\) changes the columns \(p_b s e_b^*\), \(0\le b\le a\),
by an upper unitriangular matrix. Take the minor on the rows
\(e_0^*,\ldots,e_{a-1}^*,i\) and these columns. Its determinant is
\[
             P_{a+1}(s e_a^*)_i .
\]
In the unchanged columns the first \(a\) rows form a lower-triangular
block with diagonal \(p_0,\ldots,p_{a-1}\) and have zero last entries.
Expansion in the last column gives the identity, also when a row
is repeated. Expansion by one row at a time puts a \(t\times t\)
determinant with size-seven entries in size \(8t\), starting
with the empty determinant \(1\). At \(i=e_a^*\) the minor is
\(P_{a+1}\), so all \(P_l\), \(l\le j\), have size at most \(9n\).
Using the inverse partial products gives size \(10n\) for
\(s e_a^*\), \(p_a=P_{a+1}/P_a\) and \(p_a^{-1}=P_a/P_{a+1}\).

For the matrix estimate let \(0\le\ell\le\varrho\), and let \(\nu\) be
unitary and block upper unitriangular for the first \(\ell\) split
pairs, with the remaining coordinates treated as one middle block.
Let \(g=\bar g\) be a unit of size one in the estimating gauge
\(\Gamma\), and set
\[
                  D_\nu=(I-\nu)/g,\qquad
                  \widetilde D_\nu=(I-\nu^{-1})/g.
\]
Suppose for each \(a<\ell\) that \(f_a(e_b^*)=\delta_{ab}\) for
\(b\le a\), and both \(f_a\) and \(D_\nu f_a\) have size \(c\).
Consider any matrix row for which the entries on the nondual columns
of the \(\ell\) pairs (the primal columns and the entire middle block)
are known to have size \(d\), and whose evaluations at all \(f_a\)
also have size \(d\). It suffices to recover the dual columns in
decreasing \(a\). To find the column at \(e_a^*\), subtract the
evaluation at \(f_a-e_a^*\) from the evaluation at \(f_a\). The vector
\(f_a-e_a^*\) has size \(c+1\) and vanishes in the as yet unknown
columns \(e_b^*\) for \(b\le a\). If the current size bound is \(d'\),
the product costs at most \(d'+c+2\) and the subtraction gives
\(d'+c+3\). All entries of the row have size \(d+(c+3)\ell\).

Apply this rule to the rows of \(D_\nu\) in the middle block and at
the dual vectors of the first \(\ell\) pairs. The upper-triangular
pattern gives zero in all their nondual columns. These rows have
size \(c+(c+3)\ell\). Since
\(J\widetilde D_\nu=D_\nu^*J\), the nondual columns of
\(\widetilde D_\nu\) have size \(c+4+(c+3)\ell\). For such a column
of \(J\), only the rows of \(D_\nu\) just bounded can contribute to
\(D_\nu^*J\). The two multiplications by Gram matrices cost four.
Now use \(\nu f_a=f_a-gD_\nu f_a\) as test vectors for
\(\widetilde D_\nu\). Their initial dual coordinates are still
\(\delta_{ab}\), their size is at most \(c+2\), and
\(\widetilde D_\nu(\nu f_a)=-D_\nu f_a\) has size \(c\).
A second application of the rule gives
\begin{equation}\label{eq:weighted-triangular-inverse}
        \widetilde D_\nu\quad\text{of size}\quad c+4+(2c+8)\ell .
\end{equation}

For the relative position \(x\), use the pair \((z,zx)\).
Let \(p,q\) be heads of length \(j\le\varrho\) in \(S\).
Assume \(w_p^{-1}z\) and \(w_q^{-1}zx\) are near at \(j\), so their
left \(Q_j\)-cosets meet the box. Choose \(u,v\in U\) so that
\(K_1=u z\), \(K_2=vzx\) lie in the box and
\(\Delta=(K_2-K_1)/g\) has size two in \(\Gamma\).
When \(j>0\), assume every scalar of size \(10n\) in \(S\), divided
by \(g\), has size one in \(\Gamma\).
For every \(j\le a<\varrho\), assume that \(L_a(z,x)\) contains a unit
\(c_a\) of \(A\) with \(c_a^{-1}\) of size one in \(\Gamma\).

Take \(\nu=v u^{-1}\) and \(\ell=\varrho\) in
\eqref{eq:weighted-triangular-inverse}. For \(a<j\), set \(f_a=u e_a^*\).
The head calculation at the two box points gives size \(10n\) in \(S\)
for \(u e_a^*\) and \(v e_a^*\). Hence \(f_a\) has size at most two
in \(\Gamma\) and \(D_\nu f_a=(u-v)e_a^*/g\) has size at most two there.
For \(a\ge j\), take a vector \(w\) realizing \(c_a\) in \(L_a(z,x)\).
Write \(b=z^{-1}w\), \(b'=(zx)^{-1}w\) and \(e=(b-b')/g\).
Both \(b\) and \(e\) are in the level ball and have size one.
Set \(f_a=c_a^{-1}u w=c_a^{-1}K_1 b\). It has the required initial
dual coordinates and has size at most four. Moreover
\[
     D_\nu f_a=c_a^{-1}(K_2e-\Delta b).
\]
The two terms inside the parentheses have sizes at most three
and four. Subtraction and multiplication by \(c_a^{-1}\) give
size at most six. All the test vectors therefore work with \(c=6\).
Equation~\eqref{eq:weighted-triangular-inverse} gives size
\(10+20\varrho\) for \(\widetilde D_\nu\). Since \(K_1^{-1}\nu^{-1}K_2=x\),
\begin{equation}\label{eq:weighted-triangular-relative}
              (x-I)/g=K_1^{-1}(\Delta-\widetilde D_\nu K_2).
\end{equation}
The right side has size at most \(19+20\varrho\le99n\).
For \(j=0\) the head scaling condition and the first set of tests
are empty. For \(j=\varrho\) the condition on remaining \(L_a\)'s
is empty. The proof includes \(\varrho=0\), when \(U\) is the identity.
\footnote{The fixed source uses the same minors and two triangular
substitutions in \weightedpin{grh-compact/src/New/RealPlace/Size.lean}
{90-L164}{the size calculations}. Its main application is
\weightedpin{grh-compact/src/New/RealPlace/FlatCharts.lean}
{211-L233}{reconstruction at the relative position}.}

For the finite test use \(W=W'=W_Z\), the preceding unipotent
factor \(v_U\), and Plancherel constant \(c=1\).
Fix \(0\le j\le\varrho\) and let \(p,q\) be heads for the gauge
\(\alpha\), with \(\alpha/g\in R\) when \(j>0\). Put
\(\Lambda=2/\operatorname{md}(\alpha)\).
The orthonormal part has \(d_{\rm orth}=k-6\varrho-13\) coordinates
and the level has \(3\varrho+6\) additional hyperbolic pairs, so
the volume exponent is \(s_{\rm f}=k-12\varrho-23\).
\teamanchor{proof.weighted-finite-hyperbolic-frame}%
Put \(\tau=-\zeta_3\in R\), so \(\tau+\bar\tau=1\).
In the marked integral block with Gram matrix
\(\left(\begin{smallmatrix}0&-1\\-1&-N\end{smallmatrix}\right)\),
take \(u=-e_o\), \(v=z-N\tau e_o\). They have norm zero and pairing
one, and \(z=v-N\tau u\). For two remaining norm-one slots \(a,b\),
take a unit \(c\) with \(c\bar c=-1\), using the split or
unramified norm on units. Then \(a+cb\) and
\(\tau a-\bar\tau cb\) have norm zero and pairing one.
Their change-of-basis determinant is \(-c\). The marked
pair has determinant \(-1\). Applying the construction on
\(3\varrho+6\) disjoint pairs leaves the stated number of
orthonormal coordinates and preserves the integer-box measures.
Put \(m_0=3\varrho+5\). There are at most \(2f\) local factors
above any rational prime, so the elementary ideal product gives
\(\zeta_R(m_0)^n\le\zeta(m_0)^{2fn}\le1+2^{-n}\) at the selected
large scales. This also bounds the ideal sums with larger
exponents needed in the finite shell calculation.
For floor mass use the integral gauge \(\alpha=1\), whose constant
is two. The size-seven quadric fibre is bounded by
\(2^{99n(7+6)}=2^{1287n}\). Size eleven is still \(R\), of
additive volume one. Hence the floor mass bound is
\(V_{\rm f}=2^{1287n}\le\Lambda^{9999n}\) in any head gauge.
On the shell \(I_x=cR\), with
\(I_x=\{c\in R:c(x-I)/g\text{ is integral}\}\) and
\(h=[R:I_x]\), the gauge of depth \(\beta=\bar c c\) makes
\((x-I)/g\) of size one. Its constant is \(2h^2\).
\footnote{\weightedpin{grh-compact/src/New/RealPlace/FlatCharts.lean}
{123-L136}{Measured quadric fibres},
\weightedpin{grh-compact/src/New/RealPlace/WeightPlancherel.lean}
{47-L52}{The floor mass at size seven} and
\weightedpin{grh-compact/src/New/Filling/GlobalToLocal.lean}
{46-L55}{The dimensions of the actual finite level}.}
In the self-dual integral coordinates each unipotent fibre at the
principal level has volume \(v_U=\prod_{a<\varrho}{\rm fl}_a(g)\).
Indeed the dilation by \(g\) has modulus \({\rm fl}_a(g)\) on
the horizontal and skew coordinates at floor \(a\).
Write \(\mathcal D=\mathcal D_j(x)\) for the square product in the
\hyperlink{proof.weighted-capped-orbital-majorant}{capped majorant},
on this same pair of finite coset vectors.

\leavevmode\teamanchor{proof.weighted-triangular-finite}%
On the finite product \(A=\prod_w E_w\) with \(R=\prod_w\mathcal O_w\),
the box is \(K\) and the level ball is \(R^r\). A nonzero inner product
of \(W_Z(z)\) and \(W_Z(zx)\) gives a common nonzero coset coordinate
\(K(g)b\). By the definition of the phase there are \(u,v\in U\) with
\(K_1b^{-1},K_2b^{-1}\in K(g)\), where \(K_1=uz\) and \(K_2=vzx\).
Thus \(K_1,K_2\in K\) and \(K_1^{-1}K_2\in K(g)\), since \(b\in Z\)
normalizes \(K(g)\). Consequently
\(\Delta=K_1(K_1^{-1}K_2-I)/g\) has integral entries.

For \(L_a(z,x)\), put \(w'=u w\). The earlier zero dual coordinates
and the \(a\)-th value are preserved in both directions by the
triangular action of \(U\). The condition \(z^{-1}w\in R^r\) is
equivalent to \(w'\in R^r\), since \(K_1\in\mathrm{GL}_r(R)\).
The second ball condition is the integrality of
\(g^{-1}(z^{-1}-(zx)^{-1})u^{-1}w'\). The permissible \(w'\)'s
form an \(R\)-submodule of \(R^r\). Taking \(w'\) to be a sufficiently
divisible multiple of \(e_a^*\) clears all the coefficients in the
second condition at each place and gives an \(a\)-th coordinate in
\(A^\times\). The image under the \(a\)-th coordinate is therefore an
ideal of \(R\) with a nonzero component at every place. Write it
\(d_aR\) with \(d_a\in A^\times\cap R\) by choosing a uniformizer
power in each component. In the normalization \(\operatorname{vol}(R)=1\),
\(\operatorname{vol}L_a(z,x)=\operatorname{md}(d_a)\).

Let \(d=\prod_{j\le a<\varrho}d_a\). For the head gauge of depth
\(\alpha\), size \(i\) means \(\alpha^i c\in R\) and
\(\Lambda=2/\operatorname{md}(\alpha)\). Use the estimating gauge of
depth \(\alpha'=\alpha^{10n+1}\bar d d\).
This is integral, invertible in \(A\) and self-conjugate. For each
remaining coordinate \(\alpha'/d_a\in R\). When \(j>0\),
\(\alpha g^{-1}\in R\) in the head gauge. If \(c\) has size \(10n\)
there, \(\alpha'c/g=(\bar d d)(\alpha g^{-1})(\alpha^{10n}c)\)
is integral. These facts and the integral \(\Delta\) give all the
hypotheses of the \hyperlink{proof.weighted-triangular-reconstruction}
{reconstruction at the same active pair}. Hence
\((\alpha')^{99n}\in I_x\), the denominator ideal of \((x-I)/g\).
If \(h=[R:I_x]\), taking moduli gives
\[
 \operatorname{md}(d)^{198n}
 \le h^{-1}\operatorname{md}(\alpha)^{-99n(10n+1)}
 \le h^{-1}\Lambda^{99n(10n+1)} .
\]
For \(\sigma^{400n^2}=h^{-1}\Lambda^{99n(10n+1)}<1\), all the
remaining volumes being greater than \(\sigma^2\) would contradict
this inequality. Indeed there are at most \(\varrho<n\) factors,
and their product to power \(198n\) would exceed
\(\sigma^{400n^2}\). For the empty product the left side is one.\footnote{The fixed source also records \weightedpin{grh-compact/src/New/Filling/TopFloor.lean}{216-L242}{the finite pivot test}, and
\weightedpin{grh-compact/src/New/RealPlace/FlatCharts.lean}{210-L233}{the estimating gauge}.}

Take \(e\ge0\) and a nonnegative integer \(t\).
Let \(0\le f\) be bounded by the absolute majorant, and suppose
that when \(e<1\) it satisfies \(f^2\le e^{2s_{\rm f}}\mathcal D\)
on the positions. The majorant is supported there by the
common-coset test. Put \(e_0=\min(1,e)\),
\(B=99n(10n+1)\), and use the threshold
\(\sigma^{400n^2}=h^{-1}\Lambda^B\) of the
\hyperlink{proof.weighted-triangular-finite}{finite reconstruction}.
Its small-coordinate conclusion gives cap \(\sigma^2\) if
\(\sigma<1\), including an empty active locus at \(j=\varrho\).
With cap one for \(\sigma\ge1\),
\eqref{eq:weighted-capped-position} now supplies the
pointwise factor \(E=\Lambda^{9999n}\min(e_0,\sigma)^{s_{\rm f}}\).
For \(t+10^8n^4\le k\), we have
\(s_{\rm f}\ge t+400n^2(200000n^2)\). Putting \(t\) powers
on \(e_0\) and \(400n^2(200000n^2)\) powers on \(\sigma\) gives
\[
 \min(e_0,\sigma)^{s_{\rm f}}
 \le e_0^t h^{-200000n^2}\Lambda^{200000n^2 B}.
\]
This includes \(\sigma\ge1\). If \(e_0=0\), the hypothesis
gives a zero integral on each shell since the weighted
self-pairings are finite there almost everywhere.
The \hyperlink{proof.weighted-restricted-position-integral}{restricted
integral} uses \(M_Z\) and \((2h^2)^{50000n^2}\). This spends
\(100000n^2\) powers of \(h\), leaving \(h^{-100000n^2}\).
The actual level has \(\sum_I[R:I]^{-\mu}\le2\)
for some \(\mu<9n\), so the residual sum is at most two.
Using \(\Lambda\ge2\), all factors beyond \(e_0^tM_Z\) are bounded
by \(\Lambda^{B_{\rm f}}\), where
\[
 B_{\rm f}=200000n^2 B+1500000n^3+9999n+50000n^2+2
    \le10^9n^4.
\]
Consequently
\[
 \kappa_j\int f\,dM\le\min(1,e)^t M_Z\Lambda^{10^9n^4}.
\]
Taking \(e=1\) proves absolute integrability of the finite majorant.
\footnote{\weightedpin{grh-compact/src/New/Filling/StandardLevel.lean}
{84-L89}{The ideal zeta bound} and
\weightedpin{grh-compact/src/New/Filling/TopFloor.lean}
{271-L340}{Denominator summation}. The fixed source also records \weightedpin{grh-compact/src/New/RealPlace/LocalTheorem.lean}{322-L382}{the common estimates at the finite gauges}, and
\weightedpin{grh-compact/src/New/Filling/TopFloor.lean}{273-L382}{the finite shells with an inserted frequency}.}

\paragraph{Plancherel on one floor.}
\teamanchor{proof.weighted-one-floor-plancherel}%
Let \(A\) be \(\mathbf C\) at a real place of the base field, or
\(\prod_{w\in T}E_w\) at the finite places, with the involution of \(E\).
Write \(d\) for the number of orthonormal coordinates of the local
floor basis, and \(h(v,v')=v^*Jv'\). In that basis write a vector
as \((p,w,z_0)\), where \(p\) consists of one coordinate from each
hyperbolic plane, \(z_0\) of its partners and \(w\in A^d\). Thus
\[
   h(v',v)=(z'_0)^*p+(w')^*w+(p')^*z_0 .
\]
The order in \(p\) pairs the second coordinate of the distinguished
hyperbolic plane with its first. Write \(B_i\) for the vectors all
of whose ambient entries have gauge size \(i\), and \(n=\varrho+1\).

Use the following measure condition at size \(i\). In floor coordinates
the measure \(q\) is the image of \(dp\,dw\,d\zeta(z)\) by a map
\((p,w,z)\mapsto(p,w,z_0(p,w,z))\). The \(\zeta\)-mass of the
points in the coordinate box of size \(i+4\) is at most \(\mu_i(p)\)
for every \(p,w\), and \(\int\mu_i(p)\,dp\le M_i(q)\). The increase
by four covers passage from ambient to floor coordinates.
For the floor quadric measure and for the additive volume of the
floor one may use
\[
                         M_i(q)=\Lambda^{99n(i+6)} .
\]
Indeed in the distinguished plane the quadric chart has coordinates
\((\bar t^{-1}\tau(\gamma-h(v,v))+t\beta,t)\), where
\(\tau+\bar\tau=1,\ \bar\beta=-\beta\). Its measure is
\(dt\,dv\,d\beta\), with the normalized skew Haar \(d\beta\).
Fix \(p,w\) and subtract one admissible set of residual coordinates
from all the others. The terms \(\gamma\) and \(w^*w\) cancel.
The resulting residual parameters belong to the homogeneous
hyperbolic cone chart at the same \(p\), with coordinates of size
\(i+5\). Translation preserves their measure. The integral in \(p\)
is therefore bounded by the mass of this cone box. For additive volume the ordinary
product box in the paired coordinates gives the same upper bound.
\footnote{\weightedpin{grh-compact/src/Lib/General/Measure/Place.lean}
{116-L125}{The chart and its skew Haar measure},
\weightedpin{grh-compact/src/Lib/DiscreteSeries/Bruhat/BigCell.lean}
{263-L295}{The floor Gram and additive volume}, and
\weightedpin{grh-compact/src/New/RealPlace/FlatCharts.lean}
{94-L148}{The two measured fibre bounds}.}

Here are the cone masses for the gauges in use. Let \(e\) be the
number of hyperbolic planes including the distinguished plane.
Over \(\mathbf C\) the box radius is \(M=L^{i+5}\), \(\tau=1/2\)
and \(\beta\) is purely imaginary. Since \(h(v,v)\) is real,
the bound on the first coordinate forces
\(|\beta||t|\le M\). Its skew length is at most \(2M/|t|\).
The integral of \(|t|^{-1}\) over \(|t|\le M\) is \(2\pi M\).
The other \(2(e-1)\) coordinates contribute at most
\((\pi M^2)^{2(e-1)}\). This fits \(M_i\) with
\(\Lambda=L^2,\ e\le n,\ L\ge4\).

At a finite product put \(R=\prod_{w\in T}\mathcal O_w\), with
\(\operatorname{vol}(R)=\nu(R^-)=1\). In the unit box write the
hyperbolic form as \(\operatorname{Tr}(x^*y)\), \(x,y\in R^e\),
where \(\operatorname{Tr}(z)=z+\bar z\). For almost every \(x\)
the ideal generated by the entries of \(x^*\) is \(tR\) for a unit
of \(A\). At fixed \(x\), the image of \(dy\) under
\(\delta\operatorname{Tr}(x^*y)\) is normalized Haar on
\(\delta\operatorname{Tr}(tR)\), with \(\delta\) the skew unit.
This image contains \(t\bar t R^-\), since
\(\tau+\bar\tau=1\), \(\tau\in R\) and \(\delta\in R^\times\).
Its density relative to the skew Haar is at most
\(\operatorname{md}(t)^{-1}=N(tR)\). Indeed multiplication by
\(t\bar t\) has modulus \(\operatorname{md}(t)\) on skew elements.
Averaging the quadratic level over a neighborhood of zero gives
this same upper bound at fixed \(x\). On integrating in \(x\), bound
the content index by
\(\sum_{I\subset R}N(I){\bf1}_{x^*\in I^e}\), over open integral
ideals. This gives \(\zeta_R(e-1)\le2\) by the rank of the finite
chart. Shrinking the neighborhood yields the unit-box mass on the
zero quadric by continuity on compact open tests. For the box
\(\alpha^{-(i+5)}R\) in each coordinate, \(\alpha\) self-conjugate,
dilation multiplies this bound by
\(\operatorname{md}(\alpha)^{-(i+5)(2e-1)}\). The rescaling back
to the unit box acts on the \(2e\) ambient variables by
\(\alpha^{i+5}\) and on the skew level variable by
\(\alpha^{2(i+5)}\). Here \(e\le 1+9n+\varrho\) and
\(\Lambda=2/\operatorname{md}(\alpha)\), which gives \(M_i\).
\footnote{\weightedpin{grh-compact/src/New/RealPlace/ComplexGauge.lean}
{157-L184}{The complex cone box},
\weightedpin{grh-compact/src/Lib/General/Measure/Place.lean}
{123-L132}{Continuous quadric disintegration}, \weightedpin
{grh-compact/src/Lib/Adelic/Adeles/FinitePlaceDensity.lean}
{213-L258}{The finite content moments and cone mass}, and \weightedpin
{grh-compact/src/New/Filling/StandardLevel.lean}
{146-L156}{The self-conjugate finite gauge}.}

Suppose \(q_1,q_2\) satisfy this condition at sizes \(i_1,i_2\).
For \(f\) integrable and square integrable on \(q_1\), supported in
\(B_{i_1}\), a unit \(\lambda\in A^\times\), and a fixed shift \(y_0\),
one has
\[
 \int_{B_{i_2}}\left|\int\psi(\lambda h(v'+y_0,v))f(v)\,dq_1(v)\right|^2
   dq_2(v')
 \le M_{i_1}(q_1)M_{i_2}(q_2)
        \left(\frac c{\operatorname{md}(\lambda)}\right)^d
        \|f\|_{L^2(q_1)}^2 .                                      \tag{F}
\]
To see this, fix both bases \(p,p'\). In the input residual variables
integrate \(f\) against the character of
\(\lambda((p')^*z_0+h(y_0,v))\), calling the result \(a_{p,p'}(w)\).
By Cauchy--Schwarz,
\[
 \int |a_{p,p'}(w)|^2dw
 \le \mu_{i_1}(p)\iint |f(p,w,z_0(p,w,z))|^2dw\,d\zeta(z).
\]
The output residual variables contribute the character of
\(\lambda(z'_0)^*p\), of modulus one. At fixed \(p'\), the absolute
output is bounded by the integral in \(p\) of the absolute Fourier
transforms of \(a_{p,p'}\) in the \(w\) variables. Minkowski's
inequality, ordinary Plancherel in \(w\), and Cauchy--Schwarz in
\(p\) bound its squared \(L^2(dw')\) norm by
\(M_{i_1}(q_1)(c/\operatorname{md}(\lambda))^d\|f\|_2^2\),
uniformly in \(p'\). Integrating the output residual bound and then
\(\mu_{i_2}(p')\) proves (F).

Over \(\mathbf C\), \(d=k-1\) and the character is
\(\exp(2\mathrm i\operatorname{Im}z)\). Relative to the ordinary
\(2\pi\) Fourier transform on the real plane, the dual-variable
Jacobian in one coordinate is \(|\lambda|^2/\pi^2\).
Thus \(c=\pi^2,\ \operatorname{md}(\lambda)=|\lambda|^2\).
Over the finite product \(d=k-6\varrho-13\). The character has the
integral product ring as its own annihilator, that ring has mass
one, and conjugation preserves this Haar measure. Hence \(c=1\)
and \(\operatorname{md}(\lambda)=\prod_{w\in T}|\lambda_w|_w\).%
\footnote{\weightedpin{grh-compact/src/Lib/DiscreteSeries/Bruhat/BigCell.lean}
{464-L482}{The place Fourier inequality},
\weightedpin{grh-compact/src/Lib/DiscreteSeries/Bruhat/ComplexPlace.lean}
{28-L48}{Lebesgue measure and the real floor basis},
\weightedpin{grh-compact/src/New/Filling/GlobalToLocal.lean}
{44-L49}{The finite floor basis}, and \weightedpin
{grh-compact/src/New/Filling/GlobalToLocal.lean}{102-L143}
{The self-dual character on the product}.}

Number the frame \((f_1,\ldots,f_n)=(e_0,\ldots,e_{\varrho-1},z)\)
in order. For the inserted character take both measures in (F) to be the
quadric measure on floor \(j\le\varrho\). At a given position
and preceding translate, the two upper kernels \(K_{j+1},K'_{j+1}\)
occur with one factor on each of the two orbits. The orbit map
\(\sigma\mapsto t\sigma^{-1}f_{j+1}\) takes \(d\theta_j\) to the
quadric measure for every \(t\in H_j\). The upper kernel descends
to this orbit, by invariance under its stabilizer.
\teamanchor{proof.weighted-local-plancherel-support}%
When both kernels contribute, \(y\) and \(w_p^{-1}y\) are near at \(j\).
Choose \(tu\in H_j\mathcal N_{<j}\) with \(tuy\in K\), and put
\(y_0=ty\), \(\nu=tut^{-1}\). Then \(\nu y_0\in K\), and
\(w_p^{-1}y_0\) is still near. The head minors give size \(10n\)
for the \(\nu e_a^*\) and \(p_a^{\pm1}\). The triangular inverse
estimate with level scalar one and \(c=10n+1\) bounds
\(I-\nu^{-1}\) by size \(c+4+(2c+8)j\). Adding the identity and
multiplying by \(\nu y_0\) gives size \(60n^2\) for \(y_0\).
Multiplication by \(w_p^{-1}\) gives size \(99n^2\).
The empty head uses the box representative itself. For active
\(\sigma\in H_j\), choose \(\xi'\in Q_{j+1}\) with \(\xi'\sigma y\in K\).
Project \((ty)(\xi'\sigma y)^{-1}f_{j+1}=t\sigma^{-1}
(\xi')^{-1}f_{j+1}\) onto the remaining floor. This gives
\(t\sigma^{-1}f_{j+1}\) because \((\xi')^{-1}f_{j+1}-f_{j+1}\)
lies in the earlier frame fixed by \(t,\sigma\). The two orbit
supports therefore have sizes \(60n^2+6,99n^2+6\).
The inserted entry is the pairing of these two
points and has size at most \(200n^2\). A frequency agreeing with
\(\lambda\) as a character on that box gives exactly the same
insertion. Apply Cauchy--Schwarz to the first orbit and (F) to the
second, using \(-\lambda\) when moving the complex conjugation.
The squared double integral is bounded by
\((c/\operatorname{md}(\lambda))^d\) times the product of the two
squared \(L^2(d\theta_j)\) norms and \(\Lambda^{10^6n^3}\).
\footnote{\weightedpin{grh-compact/src/New/RealPlace/LevelGauge.lean}
{93-L120}{The exact push-forward of the orbit measure},
\weightedpin{grh-compact/src/New/RealPlace/FlatCharts.lean}
{194-L209}{Support sizes for a common center}, and
\weightedpin{grh-compact/src/New/RealPlace/WeightPlancherel.lean}
{137-L188}{Factorization of the upper kernel through its orbit}.}

For unequal last pivots put \(j=a+1\le\varrho\) and assume that
\(\xi=(p_a^{-1}-q_a^{-1})/g\) is a unit of \(A\). On the
Heisenberg group of floor \(a\) write an element as \(n(v,\beta)\),
with \(v\in V_{a+1}\) and \(\beta\) skew. Its measure is the image
of \(dv\,d\beta\). For \(\sigma\in H_{a+1}\), the only possible
nonzero term of the simple-root sum of
\(\sigma n(v,\beta)\sigma^{-1}\) is
\(-h(v,\sigma^{-1}f_{a+2})\). The Weyl ratio \(w_q^{-1}w_p\)
multiplies that simple-root entry by \(p_a/q_a\).
The translated kernel law at level \(g\) therefore inserts the
frequency \(\xi'=(p_a/q_a-1)/g=-p_a\xi\) on this root sum.

Choose an active translate as origin. After translation and
conjugation by the floor isometry from the support estimate, the
active difference, modulo the preceding unipotent group, has its
dual-frame column divided by \(u=p_a^{-1}\) of size \(h_1\).
If the centered input translate has size \(h_0\), the orbit point
support has size \(h_0+6\). In these coordinates \(v/u\in B_{h_1}\).
At fixed \(v\), two admissible skew coordinates have difference
divided by \(u\) of size \(h_1+1\), by subtracting the two central
entries of that column. If \(S\) is the central slice and
\(\delta\) is the skew unit, the measure-preserving split
\((\beta,\beta')\mapsto\tau\delta^{-1}\beta+\beta'\) shows that
\[
  \nu(S)^2\le\operatorname{md}(u)\,\Lambda^{2(h_1+2)} .
\]
Indeed \(\tau\delta^{-1}\) has gauge size one. The differences of
two points in the image of \(S\times S\)
belong to \(u\) times a scalar box of size \(h_1+3\), whose volume
before scaling is at most \(\Lambda^{h_1+4}\).
Changing \(v=u v'\) contributes \(\operatorname{md}(u)^{d_a}\),
where \(d_a=\dim_A V_{a+1}=k-1+2(\varrho-a)\).
Thus the total scaling factor from this box is
\(\operatorname{fl}_a(u)=\operatorname{md}(u)^{d_a+1/2}\).%
\footnote{\weightedpin{grh-compact/src/Lib/DiscreteSeries/Bruhat/BigCell.lean}
{211-L222}{The simple-root sum in horizontal coordinates},
\weightedpin{grh-compact/src/Lib/General/Measure/Place.lean}
{174-L180}{Measure of a skew slice from its differences},
\weightedpin{grh-compact/src/New/RealPlace/LevelGauge.lean}
{283-L293}{Gauge sizes for the fixed scalars}, and
\weightedpin{grh-compact/src/New/RealPlace/FlatCharts.lean}
{150-L173}{The centered Heisenberg box}.}

Expand the translated \(K'_j\) as the \(\theta_j\)-integral of
\(K'_{j+1}\). Moving \(n\) through \(\sigma\) supplies the character
just computed. Translation of the origin gives a fixed shift \(y_0\)
in (F) and the horizontal change supplies frequency
\(-\xi'\bar u\). Apply (F) with the input floor quadric and output
ordinary volume of \(V_{a+1}\). Since
\(\operatorname{md}(\xi'\bar u)=\operatorname{md}(\xi)\), its
contribution to the squared integral, relative to the next
\(L^2(d\theta_j)\) norm, is at most
\[
    \left(\frac c{\operatorname{md}(\xi)}\right)^d
      \operatorname{fl}_a(p_a^{-1})\,
      \Lambda^{100n(h_0+h_1+31)} .
\]
Both \(u'y\) and \(w_{p,j}u'y\) are near at \(j=a+1\) on the
active support. Apply the common-center calculation to the
second point with transform \(w_{p,j}^{-1}\). It gives
\(t_{u'}\in H_j\) with \(t_{u'}u'y\) of size \(99n^2\).
For a fixed active \(u_0\), the \(e_a^*\)-column of
\((t_{u_0}u_0y)(t_{u'}u'y)^{-1}\) is that of
\(t_{u_0}u_0(u')^{-1}t_{u_0}^{-1}\) because the centering
isometries fix \(e_a^*\). Inversion and multiplication give
size \(198n^2+5\). Subtracting the identity costs one, and
division by \(p_a^{-1}\) costs \(10n\). Thus \(h_0=99n^2\)
and \(h_1=198n^2+10n+6\) suffice. Translations of the origin
and conjugation preserve the Heisenberg Haar. The sets of
active points are the measurable near sets of the integral.
These choices put the last gauge power below
\(\Lambda^{10^6n^3}\).
\footnote{\weightedpin{grh-compact/src/New/RealPlace/Size.lean}
{79-L101}{Matrix size under products and inversion},
\weightedpin{grh-compact/src/New/RealPlace/Size.lean}
{156-L164}{The active pivot size},
\weightedpin{grh-compact/src/New/RealPlace/WeightPlancherel.lean}
{190-L220}{Translation, the horizontal frequency and volume}, and
\weightedpin{grh-compact/src/New/RealPlace/LocalTheorem.lean}
{436-L452}{The column comparison for the last pivot}.}

We record how the two estimates use the same weighted norms.
For a fixed position write \(P_l\) for the product of the first
\(l\) floor-value volumes to the local volume exponent. For either
kernel put \(A_l=P_l|K_{l+1}|^2\),
\(B_l(g)=\int(A_l+B_{l+1})(\sigma g)\,d\theta_l(\sigma)\)
with \(B_n=0\), and \(Z_l\) for the integral of \(B_l\) against
\(\theta_{l-1}*\cdots *\theta_0\), with the identity mass at \(l=0\).
Since \(P_l(\sigma g)=P_l(g)\)
on \(H_l\), the next squared orbit norm is at most \(B_l(g)/P_l(g)\).
Put \(C_\Gamma=\Lambda^{10^6n^3}\) and
\(J_l=\prod_{b<l}\operatorname{fl}_b(p_b^{-1})\).
The \hyperlink{proof.weighted-inverse-weight}{inverse-weight
calculation} bounds the integral of an invariant amplitude
divided by \(P_l(w_{p,l}g)\) against these preceding Heisenberg
measures. The bound is \(C_\Gamma J_l\) times the integral
against the remaining measure, on the stated near support.

For the inserted entry, apply Cauchy--Schwarz to the two factors
\((B_j/P_j)^{1/2}\) in the preceding translate. Applying the
inverse-weight estimate in the two Weyl coordinates and using
\(Z_j\le Z_0\) gives \(C_\Gamma^3J_j^2 Z_0Z'_j\) times the
Fourier factor. For unequal pivots, write \(h=w_{p,j}^{-1}g\).
Before Cauchy--Schwarz the squared integrand on its support factors as
\[
  \frac{A_a(g)}{P_j(h)}
       \ \frac{P_j(h)}{P_a(w_{p,a}h)}\,|K'_j(h)|^2 .
\]
Here \(P_a(w_{p,a}h)=P_a(g)\) by floor invariance.
Retain the indicators that \(g,h\) are near at \(j\).
The first integral costs \(C_\Gamma J_j\) and uses the \(A_a\)
summand of \(Z_0\). In the second, the function of \(h\), with its
near-point indicator, is left \(H_j\)-invariant. Left transport
by \(w_{p,j}^{-1}\) therefore preserves its integral against
the ordered stack measure \((\eta_{j-1}*\theta_{j-1})*
\cdots*(\eta_0*\theta_0)\), where \(\eta_l\) is the Heisenberg
measure at \(l\). In this integral isolate the last
Heisenberg measure. The quotient of weights is constant along
this integration. The horizontal estimate above costs
\((c/\operatorname{md}(\xi))^d C_\Gamma\operatorname{fl}_a(p_a^{-1})\).
The ensuing \(B'_j/P_j\) cancels the numerator \(P_j\).
The inverse-weight estimate at length \(a\) costs \(C_\Gamma J_a\)
and leaves \(Z'_j\). Thus both bounds for the squared integral of
the floor modulus have the required normalization
\[
   (c/\operatorname{md}(\lambda))^d\,\mathcal D_j(x),
   \qquad \mathcal D_j(x)=C_\Gamma^3J_j^2 Z_0Z'_j ,
\]
with \(\lambda=\xi\) for unequal pivots. At the real places the
level in this difference is \(g=1\). For finite places all measures
and weights refer to the same subgroup and local cut.
\footnote{\weightedpin{grh-compact/src/New/RealPlace/LevelGauge.lean}
{257-L276}{The positive weighted sums},
\weightedpin{grh-compact/src/New/RealPlace/LevelGauge.lean}
{226-L232}{Weyl transport on the ordered stack},
\weightedpin{grh-compact/src/New/RealPlace/LocalTheorem.lean}
{347-L363}{Cauchy--Schwarz in the two Weyl coordinates}, and
\weightedpin{grh-compact/src/New/RealPlace/LocalTheorem.lean}
{418-L458}{Weighted integration at an unequal pivot}.}

\paragraph{Integration against the preceding inverse weights.}
\teamanchor{proof.weighted-inverse-weight}%
Fix a position \(x\). Let \(I_l(h)\) be \(1/P_l(h)\) when
\(Q_lh\) meets the box \(K\), and zero otherwise. Here \(P_l\)
is the product of the first \(l\) floor-value volumes to the power
\(s\). Write \(\omega_{<l}\) for the convolution of the first \(l\)
Heisenberg measures. For a head \(p\) with all inverse initial
products of size one, we prove, for \(l\le\varrho\),
\[
 \int \Phi(h) I_l(w_{p,l}h)\,d(\omega_{<l}*\mu)(h)
 \le \Lambda^{10^6 n^3}\prod_{a<l}\operatorname{fl}_a(p_a^{-1})
                    \int\Phi\,d\mu .
\]
Here \(\Phi\ge0\) is measurable, invariant on the left by \(Q_b\)
and supported where \(Q_bh\) meets \(K\), for some \(b\ge l\).
The measure \(\mu\) may be any countable sum of finite measures.
The preceding unipotents lie in \(Q_b\), so \(\Phi\) passes
through their integration. For each remaining \(h\), some element
of \(Q_b\subset Q_l\) takes \(h\) to \(g_0\in K\). Factor that
element in \(H_l\mathcal N_{<l}\). The measure \(\omega_{<l}\)
is invariant under translations by \(\mathcal N_{<l}\) and under
conjugation by \(H_l\). Also \(H_l\) commutes with \(w_{p,l}\) and
preserves \(I_l\) on the left. Thus it suffices to bound
\(\int I_l(w_{p,l}\nu g_0)\,d\omega_{<l}(\nu)\) uniformly on the
box. At length zero the bound is one.
\footnote{\weightedpin{grh-compact/src/Lib/DiscreteSeries/Bruhat/BigCell.lean}
{169-L187}{Factorization of the partial unitriangular group},
\weightedpin{grh-compact/src/Lib/DiscreteSeries/Bruhat/BigCell.lean}
{443-L462}{The Heisenberg convolution}, and
\weightedpin{grh-compact/src/New/RealPlace/WeightPlancherel.lean}
{95-L115}{Reduction to the compact box}.}

The Heisenberg convolution has the factorization
\(\omega_{<a+1}=\eta_a*\omega_{<a}\) by invariance of the earlier
measure under the new floor group. Fix \(\nu\in\mathcal N_{<a}\)
and put \(n_v=n(v,\beta)\) for the
new Heisenberg element at \(a<l\). The points \(w_{p,a+1}n_v\nu g_0\)
and \(w_{p,a}\nu g_0\) differ by left multiplication by an element
of \(H_a\). Their first \(a\) weights agree, and nearness at
\(a+1\) for the first implies nearness at \(a\) for the second.
On the active set \((n_v\nu)^{-1}\) returns \(n_v\nu g_0\) to
\(K\). Apply the head-column bound with the head \(\bar p\), whose
Weyl element is \(w_{p,a+1}^{-1}\). The scalar \(\bar p_a\) and
the column \((\nu^{-1}n_v^{-1})e_a^*\) both have size \(10n\).
For \(u=(-\bar p_a)^{-1}\), the column of
\((\nu^{-1}n_v^{-1}-1)/u\) therefore has size \(j_0=20n+1\).
Thus \(v/u\) lies in the additive floor box of size \(j_0\). At
fixed \(v\), differences of admissible skew coordinates divided by \(u\)
have size at most \(j_0+1\). The measure-preserving combination
\((b,b')\mapsto\tau\delta^{-1}b+b'\) sends two copies of this skew
slice into a translate of \(u\) times a box of size \(j_0+3\).
Its squared mass is at most \(\operatorname{md}(u)\Lambda^{j_0+4}\).
The change \(v=u v'\) in \(d_a=k-1+2(\varrho-a)\) horizontal
coordinates therefore costs, with the skew slice, at most
\(\operatorname{fl}_a(u)\Lambda^{j_0+2}\), where
\(\operatorname{fl}_a(u)=\operatorname{md}(u)^{d_a+1/2}\).

Let \(E_0\) embed the \(d\) orthonormal coordinates. Define \(D\)
by the \(\varrho\) row constraints \((\nu E_0w)_{e_i}=0\) and
the condition that \((g_0^{-1}E_0w,(g_0x)^{-1}E_0w)\) is a
close pair. The latter requires the first vector and the
difference of the two divided by the level \(g\) to be in the
level ball. In fact \(\nu E_0w=E_0w\) on \(D\). A preceding
unipotent fixes all the other coordinates of a vector with zero
earlier dual coordinates. The Heisenberg and Weyl formulas now give
\[
   w_{p,a+1}n_v\nu E_0w
       =E_0w-p_a h(v,E_0w)e_a^* .
\]
This vector has zero earlier dual coordinates and its pullbacks
at the two translated points give the prescribed close pair.
Its \(e_a^*\)-coordinate is \(h(v/u,E_0w)\). The new floor-value
set thus contains \(z^*D\), where \(z=E_0^*J(v/u)\).
\footnote{\weightedpin{grh-compact/src/Lib/DiscreteSeries/Bruhat/FloorGroups.lean}
{43-L58}{The action on a hyperbolic pair},
\weightedpin{grh-compact/src/Lib/DiscreteSeries/Bruhat/BigCell.lean}
{119-L134}{The horizontal term of the Heisenberg action},
\weightedpin{grh-compact/src/New/RealPlace/Size.lean}
{156-L164}{The active head column}, and
\weightedpin{grh-compact/src/New/RealPlace/WeightPlancherel.lean}
{116-L135}{The transported test vector}.}

We need a measurable \(\phi\ge\operatorname{vol}(z^*D)^{-1}\) with
\(\int_{B_t^d}\phi(z)^s\,dz\le\Lambda^{99n(t+1)}\) at every size
\(t\). Over \(\mathbf C\), \(d=s+1+3n,\ s=k-3n-2\). Add the row
constraints forcing the negative projections of \(g_0^{-1}E_0w\)
and \((g_0x)^{-1}E_0w\) to vanish. The negative space has dimension
\(n\), so there are at most \(3n\) rows altogether. Choose \(s+1\)
orthonormal columns \(O\) in their common kernel. On their image both
pullbacks preserve Hilbert norm, since they preserve the Hermitian
form and take values in the positive eigenspace. Thus
\(O\{v:\|v\|\le1/9\}\subset D\), since the difference has norm
at most \(2/9\), the level is \(g=1\), and the ball has radius
\(1/4\). The image \(z^*D\) contains the disk of
radius \(\|O^*z\|/9\). Take \(\phi(z)=36/\|O^*z\|^2\).
After unitary completion of \(O\), enclose the size-\(t\) box in
the ball of radius \(S=\sqrt d\,L^t\). In \(s+1\) complex
dimensions the integral of \(\|v\|^{-2s}\) over that ball is
\(\pi^{s+1}S^2/s!\). The other \(3n\) coordinates contribute at
most \((\pi S^2)^{3n}\). The resulting moment is bounded by
\(\pi e^{36\pi}(\pi d)^{3n}d\,L^{2t(3n+1)}\), which is at most
\(\Lambda^{99n(t+1)}\) for \(\Lambda=L^2,\ 4\le L,\ d\le L\).
\footnote{\weightedpin{grh-compact/src/New/RealPlace/ComplexGauge.lean}
{123-L151}{Projection to the common kernel}, and
\weightedpin{grh-compact/src/New/RealPlace/ComplexGauge.lean}
{185-L206}{The disk in the floor values}.}

At the finite places \(d=k-6\varrho-13,\ s=k-12\varrho-23\).
The finite inverse moment is on the open position set
\(\{y^{-1}v_0yk_0:v_0\in U,\ k_0\in K(g)\}\) at fixed \(y\).
Openness follows from \(K(g)\) being open. The common-coset test
puts its position in this set.
Write \(x=y^{-1}v_0y k_0\) with \(v_0\in U,\ k_0\in K(g)\),
and \(C=yg_0^{-1}E_0\). In addition to the first \(\varrho\)
rows require the dual-frame coordinates of \(Cw\), and the
first \(\varrho\) coordinates of \((v_0^{-1}-1)Cw\), to vanish.
Then \(v_0^{-1}Cw=Cw\), since a Bessel unipotent can change only
the first frame coordinates when the dual coordinates vanish.
Clear the finitely many row denominators in each factor and use
column operations over its discrete valuation ring. The resulting
matrix in \(\mathrm{GL}_d(R)\) has its last
\(h_f=d-3\varrho=s+3\varrho+10\) columns \(Y\) in this common
kernel. Hence \(YR^{h_f}\subset D\). The first pulled-back vector
is integral, and the second is its image under \(k_0^{-1}\).

Take \(\phi(z)=\operatorname{vol}\{c:cY^*z\in R^{h_f}\}\).
For almost all \(z\) the entries of \(Y^*z\) generate a
fractional ideal \(tR\). Thus \(\phi(z)=\operatorname{md}(t)^{-1}\)
because \(z^*YR^{h_f}=\bar t R\). Its inclusion in \(z^*D\)
gives the required majorization.
On \(R^d\), integral changes of coordinates preserve volume.
Summing over the possible content ideals bounds the moment by
\(\sum_{I\subset R}N(I)^{s-h_f}=\zeta_R(h_f-s)\le2\).
On \(\alpha^{-t_0}R^d\) the cost of scaling is
\(\operatorname{md}(\alpha)^{-t_0(d-s)}\). Since
\(d-s\le10n\), this fits the desired power of the finite \(\Lambda\).
\footnote{\weightedpin{grh-compact/src/New/Filling/StandardLevel.lean}
{156-L182}{The integral row constraints and stable columns}, and
\weightedpin{grh-compact/src/Lib/Adelic/Adeles/FinitePlaceDensity.lean}
{193-L230}{Content and its integral moment}.}

In either case use the volume-preserving floor basis for \(v/u\).
The box of ambient size \(j_0\) has all floor coordinates of size
\(j_0+4\). The paired coordinates have total volume at most
\(\Lambda^{99n(j_0+6)}\), since each scalar box costs at most
\(\Lambda^{j_0+5}\) and there are at most \(2(1+9n+\varrho)\)
such coordinates. On the orthonormal coordinates use the moment
at size \(j_0+4\). Together with the skew and horizontal factors
this costs, in one step, at most
\[
 \operatorname{fl}_a(p_a^{-1})
       \Lambda^{\,j_0+2+99n(2j_0+11)}
 \le \operatorname{fl}_a(p_a^{-1})\Lambda^{10^4n^2}.
\]
Integrate the earlier weight and iterate for the \(l\le n\)
floors. This proves the asserted inverse-weight inequality on both
local levels.
\footnote{\weightedpin{grh-compact/src/New/RealPlace/FlatCharts.lean}
{150-L181}{The skew slice and projection to orthonormal coordinates},
and \weightedpin{grh-compact/src/New/RealPlace/LevelGauge.lean}
{294-L301}{The two geometric volume estimates}.}

\paragraph{The scalar volume in the position shells.}
\teamanchor{proof.weighted-scalar-volume}%
The same disks and integral columns also bound the positive floor
weight. In the construction of \(D\), allow any system of at most
\(\varrho\) homogeneous row equations in \(w\). Their number is the
only property used to construct \(O\) and \(Y\). For fixed \(z\), let
\(C_z\) consist of the scalars \(c\) for which each \(w\in D\) admits
two close pairs \((a,a-ge),(b',b'-gf')\) with
\[
  (cz)^*w
     =g^{-1}\bigl(h(a,b')-h(a-ge,b'-gf')\bigr)
     =h(e,b')+h(a,f')-g h(e,f').
\]
Here \(g=\bar g\), and \(a,b',e,f'\) lie in the level ball.
At the finite places the last expression lies in \(R\), since the
Gram matrix and \(g\) are integral. Testing each column of \(Y\) and
conjugating gives \(cY^*z\in R^{h_f}\). The latter scalar set has
volume \(\phi(z)\), so \(\operatorname{vol}(C_z)\le\phi(z)\).

Over \(\mathbf C\), \(g=1\) and the Gram matrix has operator norm one.
The three terms have total absolute value at most \(3/16\).
If \(r_z=\|O^*z\|>0\), the disk in \(z^*D\) contains \(r_z/9\).
Testing this value gives \(|c|\le 27/(16r_z)\le3/r_z\), whence
\(\operatorname{vol}(C_z)\le9\pi/r_z^2\le36/r_z^2=\phi(z)\).
When \(r_z=0\) the majorant is infinite. Thus the same
\(s\)-th moment controls the inverse floor volume and this
scalar volume. The positive-weight integral uses the case of zero
row constraints before integration over the positions.
\footnote{\weightedpin{grh-compact/src/New/RealPlace/LevelGauge.lean}
{294-L301}{Both scalar volume conditions},
\weightedpin{grh-compact/src/New/Filling/StandardLevel.lean}
{169-L191}{Testing the integral columns},
\weightedpin{grh-compact/src/New/RealPlace/ComplexGauge.lean}
{185-L214}{Testing the projected disk}, and
\weightedpin{grh-compact/src/New/RealPlace/WeightPlancherel.lean}
{61-L89}{The positive floor weight in the position integral}.}

\paragraph{Integration on a restricted set of positions.}
\leavevmode\teamanchor{proof.weighted-restricted-position-integral}%
We prove the first of the
\hyperlink{proof.weighted-position-shell-inputs}{three position estimates}.
The two gauges use the same level, scalar \(g=\bar g\), box \(K\),
close pairs, floor measures and Haar measure \(M\). Write
\(\mathcal P=\{x\in{\rm positions}:(x-I)/g\text{ has size }1
\text{ in the second gauge}\}\). All sizes in the proof up to the
final head calculation refer to this gauge, of constant \(\Lambda'\).
Put \(d=d_{\rm orth}\), \(d_a=k-1+2(\varrho-a)\), \(n=\varrho+1\),
and
\[
 C_a={\rm fl}_a(g)c^d(\Lambda')^{4600n}\quad (a<\varrho),
 \qquad C_\varrho=1 .
\]
Here \({\rm fl}_a(g)=\operatorname{md}(g)^{d_a+1/2}\). For two
points \(h,y\) let \(P_l(h,y)\) be the product of the first \(l\)
floor-value volumes to the power \(s\). The \(a\)-th value is
\(\zeta_{e_a^*}\) for a vector \(\zeta\) with earlier dual coordinates
zero for which \((h^{-1}\zeta,y^{-1}\zeta)\) is a close pair.
These sets have jointly open membership in the points and the scalar,
as one can vary the selected scalar coordinate in a witness vector
and keep the others fixed. Their volumes are measurable, positive and finite,
since each contains a neighborhood of zero and lies in the image of
the bounded set \(h\overline{\mathcal B}\) under the coordinate map.

Write \(\eta_a\) for the Haar measure of the Heisenberg group
\(\mathcal N_a\), with \(\eta_\varrho=\delta_1\), and
\(\omega_l=\eta_0*\cdots *\eta_{l-1}\), the measure of
\(\mathcal N_{<l}\). The coordinates \(n(v,\beta)\) give
\(d\eta_a=dv\,d\beta\). Translation shears the scalar coordinate
and translates the vector coordinate. Conjugation by \(H_{a+1}\)
acts on \(v\) by an isometry of additive modulus one and fixes
\(\beta\). Using the factorization
\(\mathcal N_{<a+1}=\mathcal N_{<a}\mathcal N_a\) and
\(\mathcal N_a\subset H_a\), induction shows that \(\omega_l\)
is invariant on both sides by
\(\mathcal N_{<l}\) and under conjugation by \(H_l\).

\emph{The positive weight.}
For \(Q_lh\cap K\ne\varnothing\) and \(l\le\varrho\) we claim
\[
 \int {\bf1}_{\mathcal P}(h^{-1}\nu y)P_l(h,\nu y)\,d\omega_l(\nu)
 \le \prod_{a<l}C_a .                                             \tag{1}
\]
Factor a transporter taking \(h\) into \(K\) as \(\gamma u\)
with \(\gamma\in H_l,\ u\in\mathcal N_{<l}\).
Replace \(h,y\) by \(\gamma u h,\gamma y\), and change
\(\nu\) to \(\gamma u\nu\gamma^{-1}\). The position is unchanged.
Both points of the transported pair are multiplied by \(\gamma u\).
This preserves the first \(l\) value sets, since \(Q_l\) preserves
each successive dual coordinate on the subspace where the earlier
ones vanish. It suffices to prove (1) for \(h\in K\), by induction
on \(l\). For \(l=0\) its left side is an indicator.

In the step \(a\longmapsto a+1\), a nonzero integral allows us to
choose an element \(u_0\in\mathcal N_{<a+1}\) with
\(h^{-1}u_0y\in\mathcal P\). Right invariance permits replacement
of \(y\) by \(u_0y\), so write \(x_0=h^{-1}y\in\mathcal P\).
Factor \(\omega_{a+1}=\omega_a*\eta_a\). Whenever
\(x_1=h^{-1}\nu n(v,\beta)y\in\mathcal P\), the matrix
\[
 \frac{(hx_0)(hx_1)^{-1}-I}{g}
 =h\left(\frac{x_0-I}{g}-\frac{x_1-I}{g}\right)(hx_1)^{-1}
\]
has size \(14\). Indeed \(x_1=I+g((x_1-I)/g)\) has size \(3\),
products add the sizes plus one, and \(X^{-1}=J^{-1}X^*J\) costs
four. Thus the two factors after the first multiplication have sizes
\(4\) and \(9\). The element \((hx_0)(hx_1)^{-1}=n^{-1}\nu^{-1}\)
is an earlier unipotent times \(n^{-1}\). Its \(e_a^*\)-column
gives the horizontal coordinates of \(n^{-1}\) and its \(e_a\)-entry
unchanged by that earlier factor. Consequently \(v/g\) has size
\(14\). At fixed \(v\), differences of admissible \(\beta\)'s divided
by \(g\) have size \(15\), by subtraction of the two central entries.
The skew-slice calculation in the
\hyperlink{proof.weighted-one-floor-plancherel}{one-floor proof}
then gives, for a nonnegative measurable function \(b_1\) supported
on such slices and bounded there by \(B(v/g)\) with \(B\) measurable,
\[
 \int b_1(n')\,d\eta_a(n')
 \le {\rm fl}_a(g)(\Lambda')^{16}
       \int_{B_{14}} B(v')\,dv'.                                   \tag{2}
\]
In fact the translate of the image of two copies of the skew slice
under \((b,b')\mapsto\tau\delta^{-1}b+b'\) lies in \(g\) times a
scalar box of size \(17\). Thus its squared mass is at most
\(\operatorname{md}(g)(\Lambda')^{18}\). The horizontal change of
variables costs \(\operatorname{md}(g)^{d_a}\). These give (2),
also for a smaller measurable support.

Apply the \hyperlink{proof.weighted-scalar-volume}{scalar-volume bound}
at \(h,x_0\) with no row constraints. Thus \(w\in D\) when
\((h^{-1}E_0w,y^{-1}E_0w)\) is a close pair. Set
\(z=E_0^*J(v/g)\). For every value \(b=\zeta_{e_a^*}\) at
\((h,\nu n y)\) and every \(w\in D\), unitarity gives
\[
 \begin{split}
  g^{-1}\!\left(h_J(h^{-1}\zeta,h^{-1}E_0w)
          -h_J((\nu n y)^{-1}\zeta,y^{-1}E_0w)\right)
   &=g^{-1}h_J(\zeta,E_0w-\nu n E_0w)\\
   &=(b z)^*w ,
 \end{split}
\]
where \(h_J(v,v')=v^*Jv'\). Indeed \(n E_0w=E_0w-h_J(v,E_0w)e_a\),
and the earlier unipotent adds only earlier frame vectors, orthogonal
to \(\zeta\). The two pairs on the left are close. Hence the volume
of the new value set is at most the same \(\phi(z)\) for every
such \(\nu\). Integrating the earlier \(P_a\) with its position
indicator costs at most \(\prod_{b<a}C_b\) by induction, with \(ny\)
as second point. Apply (2) to the resulting majorant \(\phi^s\)
on the remaining active slices. The additive floor box has paired
coordinate fibre mass at most \((\Lambda')^{99n(14+6)}\).
Its orthonormal coordinates have size \(14+4\), where the moment
of \(\phi\) costs at most \((\Lambda')^{99n(14+5)}\).
Since \(16+99n(20+19)\le4600n\) and \(c\ge1\), this proves (1).
\footnote{\weightedpin{grh-compact/src/New/RealPlace/FlatCharts.lean}
{150-L193}{The Heisenberg slice and the position difference} and
\weightedpin{grh-compact/src/New/RealPlace/WeightPlancherel.lean}
{54-L89}{The positive floor weight}.}

\emph{One integral of a squared pairing.}
Let \(\Theta=\theta_n*\cdots *\theta_0\),
\(\lambda_l=\theta_l*\cdots *\theta_0\), and \(k\) be either raw
pairing \(k(h,y)=\langle W(h),W'(D_j y)\rangle\),
\(D_j=w_{q,j}^{-1}w_{p,j}\), or the self pairing with \(D_0=1\)
and \(W'=W\). Use the corresponding starting floor \(j\) or \(0\).
Write \(K_l(h,y)\) for its diagonal average in
\(\theta_n*\cdots *\theta_l\), so \(K_{n+1}=k\).
The cut on the first vector implies \(K_l(h,y)=0\) unless
\(Q_l h\) meets \(K\), since each higher transporter belongs to
\(H_l\subset Q_l\) and \(U\subset Q_l\).
These are measurable, absolutely convergent averages. To see this,
at each floor \(b\le\varrho\) first translate a point near \(b\)
into the box.
The orbit point of an active translate then has size \(7\), and
its quadric mass there is at most
\(M_7\,\operatorname{vol}(B_{11}^{\rm scalar})^d\) by the fibre
estimate. The change of representative preserves this integral by
invariance of the quadric and its orbit parametrization. The terminal
measure gives mass at most one to the translates meeting the cut.
Iteration and boundedness of both vectors give uniform finite
absolute integrals on every upper stack, so Fubini gives
\(K_l=\int K_{l+1}(\sigma\,\cdot,\sigma\,\cdot)\,d\theta_l(\sigma)\).

On \(\mathcal N_{<l}\) the two unipotent transformations of
\(K_{l+1}(h,y)\) have the unitary phases coming from the two
Bessel characters. Indeed one conjugates them by the averaging
element of \(H_{l+1}\), which fixes the first \(l+1\) frame vectors,
then uses the Bessel law in the raw pairing. In particular
\(\lvert K_{l+1}(h,\nu y)\rvert=\lvert K_{l+1}(h,y)\rvert\)
for \(\nu\in\mathcal N_{<l}\). For \(b\ge j\) the twist commutes
with \(H_b\). Thus the amplitude
\(K_{b+1}(\sigma h,\sigma y)\), \(\sigma\in H_b\), descends to
the orbit \(\sigma^{-1}f_{b+1}\). Here simultaneous translations
in the stabilizer \(H_b\cap Q_{b+1}\) preserve the amplitude.
The unipotent phases cancel, and right change of representative
in successive floor integrals preserves the invariant quadric
measure, downwards from right invariance of the terminal measure.
The orbit factor is measurable. In fact \(H_b\), its ambient vector
space and \(\mathbf C\) are standard Borel spaces and the orbit map
is Borel. For any rational sublevel set of the real or imaginary
part of the amplitude, its image and the image of its complement
are disjoint analytic sets because the amplitude is constant on
fibres. Borel separation gives a Borel set agreeing with the
sublevel set on the image. The countable family of these sets
recovers each component there by rational cuts and extends it
measurably to the ambient vector space. This supplies the measurable
amplitude on the quadric in (F).

For \(0\le l\le n\) use the measurable functions
\[
 m_l(h,y)=\int {\bf1}_{\mathcal P}(h^{-1}\nu y)\,d\omega_l(\nu),
 \qquad A_l(h,y)={\bf1}_{\{m_l(h,y)>0\}}\,|K_{l+1}(h,y)|^2 .
\]
The cut by positive mass avoids choosing a representative in a
projected support. We prove for \(j\le a\le\varrho\) that
\[
 \int A_a(h,n'y)\,d\eta_a(n')
 \le C_a\int A_{a+1}(\sigma h,\sigma y)\,d\theta_{a+1}(\sigma).
                                                                    \tag{3}
\]
If \(m_{a+1}(h,y)=0\), the left side vanishes by Tonelli and
\(\omega_{a+1}=\omega_a*\eta_a\).
Otherwise \(m_{a+1}(\sigma h,\sigma y)=m_{a+1}(h,y)>0\)
on \(H_{a+1}\) by conjugation invariance. Thus the right side
uses the full square. At \(a=\varrho\) the left integral is at
the identity and (3) follows by Cauchy--Schwarz against
\(\theta_n\) on the part of its support meeting the cut, of mass
at most one.

Suppose \(a<\varrho\) and the left integral is positive. Choose
one active \(n_0\in\mathcal N_a\) with
\(h^{-1}\nu_0 n_0y\in\mathcal P\), \(\nu_0\in\mathcal N_{<a}\).
Nonvanishing implies that \(h\) is near \(a+1\). Factor a
transporter taking \(h\) into \(K\) as
\(\xi=R u_*\), where \(u_*\in\mathcal N_{<a}\) and
\(R\in H_a\cap Q_{a+1}\), and factor \(R=t r\) with
\(t\in H_{a+1}\), \(r\in\mathcal N_a\).
For any active \(n'\) choose \(\nu\) with
\(h^{-1}\nu n'y\in\mathcal P\). Compare these two positions
at the center \(\xi h\) by the size-\(14\) calculation above.
In the variable \(m=R n'n_0^{-1}R^{-1}\) their quotient is
\[
 (\xi\nu_0 n_0y)(\xi\nu n'y)^{-1}=\nu' m^{-1},
 \qquad \nu'\in\mathcal N_{<a},
\]
since \(H_a\) normalizes the earlier unipotents. Therefore (2)
holds for the active \(m=n(v,\beta)\). This change of variables
preserves \(d\eta_a\). Its horizontal coordinate in \(n'\) is
\(t^{-1}v+v_0\), where \(v_0\) is that of \(n_0\).

Expand \(K_{a+1}(h,n'y)\) at the next floor. Moving \(n'\)
past \(\sigma\in H_{a+1}\) inserts
\(\psi(-g^{-1}h_J(t^{-1}v+v_0,\sigma^{-1}f_{a+2}))\).
Indeed the conjugated simple-root entry is the negative of this
pairing before division by \(g\). The twist acts trivially on
\(\mathcal N_a\) since \(a\ge j\), and only the second argument
was translated. Thus with \(v=g v'\) the frequency in (F) is
\(-1\), its shift is \(t v_0/g\), and its input orbit variable is
\(t\sigma^{-1}f_{a+2}\). This orbit carries the floor quadric
measure. Where \(K_{a+2}(\sigma h,\sigma y)\ne0\), that variable
has size \(7\). In detail \(t^{-1}\xi\in\mathcal N_{<a+1}\), and
conjugating this element by \(\sigma\) preserves that subgroup.
Consequently \((\sigma t^{-1})(\xi h)\) is near \(a+2\).
Its translate into the box and the size-one matrix \(\xi h\)
give the column of \((\sigma t^{-1})^{-1}\) on \(f_{a+2}\)
of size \(1+5+1=7\) after projecting off the earlier frame
vectors fixed by \(H_{a+1}\).
Applying (2) and then (F) with quadric input at size \(7\) and
additive output at size \(14\) bounds the left side of (3) by
\[
 {\rm fl}_a(g)c^d
   (\Lambda')^{16+99n(7+6)+99n(14+6)}
     \int |K_{a+2}(\sigma h,\sigma y)|^2\,d\theta_{a+1}(\sigma).
\]
The exponent is at most \(4600n\). This proves (3), including
the case of a null active slice.
\footnote{\weightedpin{grh-compact/src/New/RealPlace/LevelGauge.lean}
{93-L120}{Orbit push-forward in the stated measure},
\weightedpin{grh-compact/src/Lib/DiscreteSeries/Bruhat/BigCell.lean}
{434-L440}{Measurable descent to the orbit},
\weightedpin{grh-compact/src/New/RealPlace/FloorKernels.lean}
{35-L63}{Successive averages and stabilizer invariance}, and
\weightedpin{grh-compact/src/New/RealPlace/WeightPlancherel.lean}
{190-L220}{The centered Heisenberg Fourier integral}.}

Set \(\mu_n=\Theta\) and \(\mu_l=\eta_l*\mu_{l+1}\) downwards.
In the ordered stack move the earlier Heisenberg measures past
the higher floor factors using their conjugation invariance.
Since \(\eta_\varrho=\eta_n=\delta_1\), this gives
\(M=\omega_l*\mu_l\) for each \(l\).
When \(s_1\in H_l\), the measures \(\mu_l\) and its
left translate by \(s_1\) therefore give the same integral to
every nonnegative left \(\mathcal N_{<l}\)-invariant function.
Indeed their convolutions on the left with \(\omega_l\) are both
\(M\), by left invariance of \(M\) and conjugation invariance of
\(\omega_l\). On invariant functions such a factor is unique by the
quotient integral formula. Explicitly, choose a strictly positive
integrable \(b\) on \((G,M)\) and put
\(b_l(y)=\int b(\nu y)\,d\omega_l(\nu)\).
This is positive, finite almost everywhere for either factor, and
left \(\mathcal N_{<l}\)-invariant. Set \(b/b_l=0\) on the
exceptional null set. For an invariant nonnegative
\(F_1\), integrating \(F_1 b/b_l\) against \(M\) gives
\(\int F_1\,d\mu_l\) by the convolution identity, and likewise
for the translated factor.

Write \(T_l=\iint A_l(h,y)\,d\mu_l(y)\,d\lambda_l(h)\).
Integrate (3) against \(d\lambda_a(h)\,d\mu_{a+1}(y)\).
For fixed \(\sigma h\), \(A_{a+1}(\sigma h,\cdot)\) is invariant
under \(\mathcal N_{<a+1}\), by right invariance of \(\omega_{a+1}\)
and the character law for \(K_{a+2}\). The preceding identity
allows the simultaneous translate \(\sigma y\) to be replaced by
\(y\) in its integral. It follows that \(T_a\le C_a T_{a+1}\),
and \(T_n\le\iint |k(h,y)|^2\,d\Theta(y)\,d\Theta(h)\).

Expanding the finite sum defining \(Z_j(k;x)\) gives one term for
each \(j\le l\le\varrho\), namely the \(\lambda_l\)-integral of
\(P_l(h,hx)|K_{l+1}(h,hx)|^2\). In its integral over
\(x\in\mathcal P\), put \(hx=\nu y\) using \(M=\omega_l*\mu_l\)
and Haar invariance. The square is constant as \(\nu\) varies.
If \(m_l=0\), the indicator vanishes almost everywhere in \(\nu\).
A zero square also gives zero. Otherwise \(h\) is near
\(l+1\), hence near \(l\), and (1) bounds the term by
\((\prod_{a<l} C_a)T_l\). Iterating (3) proves
\[
 \int_{\mathcal P} Z_j(k;x)\,dM(x)
 \le n\left(\prod_{a<\varrho} C_a\right)
          \iint |k(h,y)|^2\,d\Theta(y)\,d\Theta(h).                 \tag{4}
\]
There are at most \(n\) summands. The choices of centers in
(1) and (3) were made only for fixed outer variables in pointwise
inequalities. The integrals use the measurable \(m_l,A_l\).

\emph{The head factors.}
Put \(\mathcal N(V,V')=\iint|\langle V(h),V'(y)\rangle|^2
\,d\Theta(y)\,d\Theta(h)\). The raw self pairing in (4) has squared
integral \(\mathcal N(W,W)\). For the second pairing it is
\[
 J\,\mathcal N(W,W'),\qquad
 J=\prod_{a<j}\operatorname{md}(p_a/q_a)^{d_a+1}.
\]
Indeed the \(a\)-th Weyl parameter in \(\Theta\) has density
\(\operatorname{md}(b)^{d_a}\,db\), and its Weyl factor commutes
with all higher floors. Write \(D_a=w_{q,a}^{-1}w_{p,a}\) and
\(r_a=p_a/q_a\). In that factor
\(D_{a+1}w_a(r_ab)^{-1}=D_aw_a(b)^{-1}\) by the two-by-two
hyperbolic block. Substituting \(r_ab\) contributes
\(\operatorname{md}(r_a)^{d_a+1}\). Iterating proves the formula.

Let \(N_0^2\ge\mathcal N(W,W)\mathcal N(W,W')\).
With \(S\subset\mathcal P\), \(f\) nonnegative measurable and
\(f^2\le E^2\mathcal D_j\) there, apply Cauchy--Schwarz in \(x\)
to the two square roots in \(\mathcal D_j\), extending their squared
integrals to \(\mathcal P\). Formula (4) bounds \(\kappa_j\int_S f\)
by
\[
 E N_0\, n\left(\prod_{a<\varrho} C_a\right)
    \Lambda^{1500000 n^3}
    \kappa_j\left(\prod_{a<j}{\rm fl}_a(p_a^{-1})\right)J^{1/2}.
\]
Since \(\kappa_j=\prod_{a<j}\operatorname{md}(p_aq_a)^{d_a/2}\),
the last three factors starting with \(\kappa_j\) multiply to
\(\operatorname{md}((\prod_{a<j}q_a)^{-1})^{1/2}\).
The inverse prefix has head-gauge size one, hence size two and
modulus at most \(\Lambda^2\). The argument requires only \(q\)
to be a head. As \(n\le2^n\le(\Lambda')^n\) and
\(4600n\varrho+n\le 50000n^2\), the resulting estimate is
\[
 \kappa_j\int_S f\,dM
 \le E N_0\,c^{d\varrho}
        \left(\prod_{a<\varrho}{\rm fl}_a(g)\right)
        \Lambda^{1500000n^3+1}(\Lambda')^{50000n^2}.
\]
This includes \(j=0\) with empty head product. For the finite
pairing norms in the applications, a zero norm makes the
corresponding square integral vanish by (4), and \(E=0\)
already gives \(f=0\) almost everywhere on \(S\).
The proof also bounds the upper integral of \({\bf1}_S f\)
for an arbitrary subset \(S\subset\mathcal P\) and arbitrary
nonnegative \(f\) satisfying the pointwise comparison, since
\({\bf1}_{\mathcal P} E\sqrt{\mathcal D_j}\) is its measurable
majorant in the same Cauchy--Schwarz calculation.
\footnote{\weightedpin{grh-compact/src/New/RealPlace/LevelGauge.lean}
{141-L155}{The Weyl modulus in the upper measure},
\weightedpin{grh-compact/src/New/RealPlace/LevelGauge.lean}
{191-L225}{The ordered factorization and quotient integration},
\weightedpin{grh-compact/src/New/RealPlace/LocalTheorem.lean}
{117-L175}{The integral of one squared pairing}, and
\weightedpin{grh-compact/src/New/RealPlace/LocalTheorem.lean}
{322-L344}{The restricted-position integral}.}

\paragraph{Which finite frequency survives the summation.}
\teamanchor{proof.weighted-finite-frequency}%
For the same \(A,R,Z\) and \(j\), write \(D_l(p)=p_0\cdots p_{l-1}\).
On a nonzero finite orbital majorant all \(D_l(p),D_l(q)\),
\(1\le l\le j\), are integral: they are the leading minors of
the two frames in integral representatives of the active cell.
Let \(u=\prod_{1\le l\le j}D_l(p)\) and
\(v=\prod_{1\le l\le j}D_l(q)\). Their quotients by any of the
respective partial products or by any single pivot are integral.
For example \(u/p_i=(u/D_{i+1}(p))D_i(p)\).
Thus \(\alpha=g\overline{uv}\,uv\) is an integral self-conjugate
scale clearing both heads, \(1/g\), and all \(1/p_i-1/q_i\).
For an exponent \(K_1\) define
\(\Phi_{K_1}(p)=\operatorname{md}(u)^{-K_1/2}\) on this support;
use zero when any required partial product is nonintegral.
For \(4E_1\le K_1\), multiplicativity and
\(\operatorname{md}(g),\operatorname{md}(uv)\le1\) give
\[
 \Lambda^{E_1}\le
  (2/\operatorname{md}(g))^{K_1}\Phi_{K_1}(p)\Phi_{K_1}(q).
\]

Let \(\xi\in A^\times\). The two local tests needed for the sieve
are the additive character of \(\xi\) times the next diagonal
Gram entry, and, at \(j=i+1\), the constant weight with
\(\xi=(1/p_i-1/q_i)/g\).
Write \(\beta=\alpha^{200n^2}\). For the character test choose
\(\eta\) componentwise to be either \(\xi\) or \(\beta\), taking
one of smaller valuation at each component. Then
\(\xi/\eta,\beta/\eta\in R\) and \(\eta-\xi\in\beta R\).
On the arguments \(z\) with \(\beta z\in R\) the characters of
\(\xi z\) and \(\eta z\) agree. Put \(w=\beta/\eta\).
For the other test take \(\eta=\xi,\ w=\beta\); here both
\(\alpha/g\) and \(\alpha(1/p_i-1/q_i)\) are integral.
In both cases \(w,\xi w,\xi/\eta\in R\), and, with
\(a_0=\min(1,\operatorname{md}(\eta)^{-1/2})\),
\[
 a_0^2\le \operatorname{md}(w)/\operatorname{md}(\beta).
\]
The \hyperlink{proof.weighted-one-floor-plancherel}{one-floor
Plancherel bounds} for these tests show that their
absolute pivot integrand squared is at most
\(\operatorname{md}(\eta)^{-|\iota|}\mathcal D(x)\), where
\(|\iota|\) is the number of orthonormal coordinates. For the first
test the Fourier transform is taken in the next frame coordinate, on
which frequencies agreeing on \(\beta^{-1}R\) give the same character.
For the second it is taken in the Heisenberg coordinate, whose two
translated characters differ by \((1/p_i-1/q_i)/g\).
Since \(s_{\rm f}\le|\iota|\), these are exactly the gains required by
the shell bound whenever \(\operatorname{md}(\eta)^{-1/2}<1\).
\footnote{\weightedpin{grh-compact/src/New/RealPlace/LocalTheorem.lean}
{384-L447}{The two local Fourier bounds} and
\weightedpin{grh-compact/src/New/Filling/TopFloor.lean}
{245-L273}{Head denominators in the integral gauge}.}

Put \(d(\xi)=[R:R\cap\xi^{-1}R]\). Since \(wR\subset
R\cap\xi^{-1}R\) and \(\xi/\eta\in R\), we have
\(\operatorname{md}(w)\le d(\xi)^{-1}\) and
\(a_0\le\operatorname{md}(\xi)^{-1/2}\).
Use the shell estimate at \(t+12\), keeping \(t\) powers for
the frequency and six pairs for its denominator. It gives
for the floor integral \(I_j\) with either of these weights
\[
 \frac{|I_j|}{M_Z}\le
 \operatorname{md}(\xi)^{-t/2}d(\xi)^{-6}
 (2/\operatorname{md}(g))^{K_1}\Phi_{K_1}(p)\Phi_{K_1}(q).
\]
Indeed the residual power of \(\Lambda\) is at most
\(E_1=10^9 n^4+1200n^2\), including the six factors of
\(\operatorname{md}(\beta)^{-1}\). The assumptions
\(t+10^{10}n^4\le k,\ K_1\ge10^{10}n^4\) give \(4E_1\le K_1\)
and the required shell range. In the sieve \(t=k-n^7,\ K_1=n^6\)
satisfy these inequalities for sufficiently large \(n\). The finite
set contains all places over the rational level \(g=g_0g_1\), so
\(\operatorname{md}(g)=g^{-2f}\). From \(g\le b^{22}\) we obtain
\((2/\operatorname{md}(g))^{n^6}\le b^{45fn^6}\le b^{fn^9}\).
This gives the specified finite frequency bound in both uses of
the floor recursion.\footnote{\weightedpin
{grh-compact/src/New/Filling/TopFloor.lean}{343-L382}
{Finite Fourier estimate} and \weightedpin
{grh-compact/src/New/Filling/FiniteSystem.lean}{235-L266}
{The two uses at the sieve exponents}.}

\paragraph{The contribution of the principal level.}
\teamanchor{proof.weighted-finite-bottom}%
Put \(\delta_R=\operatorname{vol}(R\setminus R^\times)\) and write
the floor-zero form as
\[
 \mathcal B(Z)=\int\left|\int_Z\chi_g(yxy^{-1})\,dM(y)\right|^2\,dM(x).
\]
On \(K(g)\) each conjugated character equals one, since \(K(g)\)
is normal in \(K\). The contribution of this coset is exactly
\(M(Z)^2M(K(g))\).

\leavevmode\teamanchor{proof.weighted-triangular-principal}%
For the finite bottom use \(j=0\) and the depth-one gauge in which
every size set is \(R\). The Weyl elements are the identity and
\(Q_0\) is the whole isometry group, so both points are near at zero.
The head scaling is empty. A nonzero coset-vector pairing supplies
the integral box representatives and integral divided difference
as above. If every \(L_a(z,x)\) contained a unit of \(R\), those
inverses would also be integral, and the
\hyperlink{proof.weighted-triangular-reconstruction}{reconstruction}
would give integral \((x-I)/g\). Since \(g\in R\), this makes \(x\)
integral and its unitary inverse is integral as well. This would put
\(x\) in \(K(g)\).

For a position outside \(K(g)\), therefore, every active pair has
some \(L_a(z,x)\) with no unit of \(R\). By the
\hyperlink{proof.weighted-triangular-finite}{finite ideal calculation},
it is an integral ideal \(d_aR\) disjoint from \(R^\times\), with
\([R:d_aR]\ge2\). Thus the cap on its additive volume is
\(\min(\operatorname{vol}(R\setminus R^\times),1/2)\).
This is the cap for the exterior term. The contribution of \(K(g)\)
itself comes from the principal phase identity.


In the integral gauge, take \(j=0,\ p=q=1\) and
\(e=\min(\delta_R,1/2)\). The
\hyperlink{proof.weighted-triangular-principal}{principal reconstruction}
gives the cap \(e\) at every nonzero pair outside \(K(g)\).
Here the two absolute pivots coincide. The
\hyperlink{proof.weighted-bottom-cap}{coincident-pivot bound}
\eqref{eq:weighted-coincident-pivot-cap} gives
\(f_0^2\le e^{2s_{\rm f}}\mathcal D\) for the floor-zero majorant.
The modulus of the phase integral is \(v_U\) times the modulus of the
floor-zero upper pairing, so its square is at most \(v_U^2 f_0\).
Apply the shell bound to \(1_{K(g)^c}f_0\) with
\(t'=t+1+n+10^9n^4\). If \(t+10^{10}n^4\le k\), then
\(t'+10^8n^4\le k\) and
\(2^{10^9n^4}e^{t'}\le2^{-n}\delta_R^{t+1}\).
Restoring \(v_U^2M_Z=M(Z)M(K(g))\) yields
\[
 \mathcal B(Z)\le M(Z)M(K(g))
                 \bigl(M(Z)+2^{-n}\delta_R^{t+1}\bigr).
\]
The zero cap also gives this estimate when \(\delta_R=0\).
\footnote{\weightedpin{grh-compact/src/New/Filling/TopFloor.lean}
{133-L138}{The integral on the principal coset},
\weightedpin{grh-compact/src/New/Filling/TopFloor.lean}
{199-L205}{Phase integral and the upper pairing} and
\weightedpin{grh-compact/src/New/Filling/TopFloor.lean}
{384-L415}{The exterior integral}. The fixed source also records \weightedpin{grh-compact/src/New/Filling/TopFloor.lean}{198-L210}{the normalized phase integral}.}

\paragraph{The real correlation forms for the open term.}
\teamanchor{proof.weighted-real-correlation-forms}%
Work at one real place and write \(F_\flat=W^\flat\). For
\(p\in(\mathbf C^\times)^{\varrho+1}\), let \(w_p\) be the Weyl element
which sends \(e_a\) to \(p_a e_a^*\) and \(e_a^*\) to
\(\bar p_a^{-1}e_a\) for \(a<\varrho\).
Let \(\sigma_c\) fix the hyperbolic pairs and act on the plane
of a specified positive unit vector and \(z\) by
\[
 \begin{pmatrix}\bar u&\sqrt{|u|^2-1}\\-\sqrt{|u|^2-1}&-u\end{pmatrix},
 \qquad u=c/\min(1,|c|) .
\]
It is the identity on the other positive coordinates.
For \(\gamma=w_p\sigma\) with \(\sigma\in U(X_0\oplus\mathbf C z)\),
write
\[
 J_F(\gamma)=\int_{G_1} F_\flat(y)F_\flat(\gamma^{-1}y)^*dy .
\]
This Haar is the filling Haar on the group. To check absolute
convergence, put \(K_R=\{y:\|y\|\le R\}\).
Both factors can be nonzero only on \(U K_R\cap\gamma U K_R\).
Conjugation by \(\gamma\) sends \(U\) to the opposite unipotent.
In a factorization of a point of this intersection the product
of the inverse lower and the upper unipotent factors lies in
\(\gamma K_R K_R^{-1}\). Elimination on lower-times-upper
unitriangular matrices recovers both factors with pivots one.
Thus the supports have compact intersection, locally uniformly
in these parameters. This also gives continuity of the correlation
by dominated convergence.

Define
\[
 \begin{aligned}
 c_G&=\frac{\operatorname{vol}U(k)}{\operatorname{vol}U(r)}
                   \frac{2\pi^k}{(k-1)!},&
 \Theta&=\pi^{(2\varrho+1)k}
            \frac{\operatorname{vol}U(k)}{\operatorname{vol}U(r)},\\
 d_a&=k-1+2(\varrho-a)\quad (a<\varrho),&
 \Xi(p)&=\frac{\sqrt{c_G}}{d}
 (|p_\varrho|^2-1)_+^{(k-1)/2}
 \prod_{a<\varrho}|p_a|^{d_a}\,J_F(w_p\sigma_{p_\varrho}).
 \end{aligned}
\]
The top form is \(J_n(p,q)=\langle\Xi(p),\Xi(q)\rangle_{\rm HS}\),
with the Hilbert--Schmidt inner product linear in the second variable.
It is continuous and positive semidefinite when the lists have
a common terminal entry. For any one \(\sigma\) with
\(h(z,\sigma z)=p_\varrho=q_\varrho\), one may replace
\(\sigma_{p_\varrho}\) in both arguments by \(\sigma\).
When the terminal modulus is one the scalar is zero.
When it is larger than one, aligning the positive part of
\(\sigma z\) by \(H\) gives a factorization
\(\sigma=h_1\sigma_{p_\varrho}h_2\) with \(h_1,h_2\in H\).
Indeed the remaining isometry fixes \(z\).
Then \(J_F\) in both arguments is multiplied on the left by
\(\rho_H(h_1)\) and on the right by \(\rho_H(h_2)\).
These changes preserve the inner product.

For a rational packet with lower pivots \(p_a\) and terminal pairing
\(t=h(z,\sigma z)\) in the rational frame of norm \(-N\), the input
to this real form is the list
\[
 p'_{va}=s_{va}^2\Phi_v(p_a)\quad(a<\varrho),\qquad
 p'_{v\varrho}=\Phi_v(t)/N .
\]
Indeed the basis change conjugates each Weyl block by
\(\operatorname{diag}(s_{va},s_{va}^{-1})\), with the inverse on the
left, and scales the terminal pairing by \(s_{v\varrho}^2\).
Equivalently the leading minor of order \(j\) is multiplied by
\(\prod_{a<j}s_{va}^2\) under \(\Phi_v\).
The condition \(|\Phi_v(t)|>N\) makes the real terminal modulus larger
than one. These substitutions apply to both members of each pair
at the same \(t\), with the same \(F_\flat\).
\footnote{\weightedpin{grh-compact/src/New/Filling/BesselSystem.lean}{45-L74}{The correlation and its normalization},
\weightedpin{grh-compact/src/New/RealPlace/BandLimits.lean}{15-L64}{The top form at a common terminal pairing},
\weightedpin{grh-compact/src/New/Filling/Setup.lean}{375-L381}{Scaling the actual leading minors}, and
\weightedpin{grh-compact/src/New/Filling/GlobalToLocal.lean}{70-L74}{Scaling the Weyl blocks}.}

Here are the lower forms and their estimates in these measures.
For \(i\le\varrho\), let \(G_i\) denote the group fixing the earlier
hyperbolic pairs. At a lower floor \(i<\varrho\), its measure
\(\theta_i\) is the law of \(w_i(c)^{-1}u\), with
\(|c|^{2d_i}dc\) and Lebesgue Haar \(d\eta_i(u)\) on the corresponding
Heisenberg unipotent. At the terminal floor use the law of
\(\sigma_c^{-1}h\) with
\(2\pi^k(|c|^2-1)_+^{k-1}dc/(k-1)!\) and probability Haar on \(H\).
This is polar integration on the positive vector in
\(\sigma_c^{-1}z\), whose squared length is \(|c|^2-1\).
At floor \(n\) use probability Haar on \(H\).
Write \(M_i\) for the ordered convolution of \(\eta_a*\theta_a\) for
\(a<i\), highest floor on the left, and \(\mathcal M=M_{n+1}\).
There are no Heisenberg factors at or above \(\varrho\).
Also write \(B_i^D(z,z')\) for the average of
\(\langle F_\flat(vz),F_\flat(Dvz')\rangle\) against the convolution
of \(\theta_n,\ldots,\theta_i\), in that order.
For truncated Weyl products \(w_p^{(i)},w_q^{(i)}\) of length \(i\),
put \(D=(w_q^{(i)})^{-1}w_p^{(i)}\). Define
\[
 \begin{split}
 \mathcal I_i(p,q;\omega)=&
 \prod_{a<i}|p_a q_a|^{d_a}
 \int_{x,g}\int_{\sigma,\tau}
 \omega(\sigma\tau^{-1})
 B_{i+1}^{1}(\sigma g,\sigma gx)
 \\[-2pt]
 &\hspace{9mm}\cdot
 \overline{B_{i+1}^{D}\bigl(\tau(w_p^{(i)})^{-1}g,
                           \tau(w_p^{(i)})^{-1}gx\bigr)}
 \,d\theta_i(\sigma)d\theta_i(\tau)\,dM_i(g)d\mathcal M(x).
 \end{split}
\]
The superscript \(1\) means the identity operator.

\paragraph{The real Gram cell and its Haar measure.}
\leavevmode\teamanchor{proof.weighted-real-gram-measure}%
Write \(m=\varrho\), \(n=m+1\), \(W=X_0\simeq\mathbf C^k\)
and \(r=2m+k+1\). Use the Euclidean-orthonormal split basis
in the order \((E,W,\mathbf C z,D)\), with \(E=(e_a)_{a=1}^m\),
\(D=(e_a^*)_{a=1}^m\), and take \(x_0\) to be the designated basis
vector of \(W\). The Hermitian involution \(J\) interchanges \(E,D\),
acts by \(1\) on \(W\) and by \(-1\) on \(z\).
Let \(F=(E,z)\), \(F'=(D,x_0)\), \(p=\operatorname{diag}(0_m,1)\),
and \(p_0=1-p\). The mixed Gram map is \(Z(y)=F^*JyF\).
Write \(D_m(Z)\) for the leading principal minor of order \(m\).
The cofactor formula gives, for invertible \(Z\),
\[
 \mathfrak g(Z)=|\det Z|^2-|D_m(Z)|^2=|\det Z|^2\delta(Z),
 \qquad \delta(Z)=1-| (Z^{-1})_{nn}|^2 .
\]
Thus the open Gram cell \(\Omega_{\rm Gr}=\{\mathfrak g>0\}\) is the
domain on which \(Z\) is invertible and \(\delta>0\). Write
\(\Omega_G=\{y\in G_1:Z(y)\in\Omega_{\rm Gr}\}\).

\paragraph{The section and its stabilizer.}
\teamanchor{proof.weighted-real-gram-section}%
For \(Z\in\Omega_{\rm Gr}\), put
\[
 K_Z=(Z^{-1})^*pZ^{-1},\qquad
 T_Z=-p-\tfrac12 p_0K_Z(1+p),\qquad
 \Lambda=p_0+\sqrt{\delta(Z)}\,p .
\]
The section \(\gamma_Z\) used below sends \(F,F'\) to
\[
 X_Z=(F T_Z+F'\Lambda)Z,\qquad
 X'_Z=\bigl(F(Z^{-1})^*-X_Z T_{Z^*}\bigr)\Lambda^{-1},
\]
and fixes \(W\ominus\mathbf C x_0\) pointwise. These are the
section matrices in the filling coordinates.
Indeed \(F^*JF=-p,\ (F')^*JF'=p,\ F^*JF'=p_0\).
Consequently \(F^*JX_Z=Z\). The Gram of \(X_ZZ^{-1}\)
has blocks \(- (K_Z)_{11},-(K_Z)_{1n},-(K_Z)_{n1}\)
and \(\delta(Z)-1=-(K_Z)_{nn}\), so \(X_Z^*JX_Z=-p\).
Using \(pT_{Z^*}=-p\) gives \(X_Z^*JX'_Z=p_0\). For the
last block, \(T_{Z^*}+T_{Z^*}^*=-2p-K_{Z^*}+pK_{Z^*}p\)
gives
\[
 (X'_Z)^*JX'_Z=\Lambda^{-1}(p-pK_{Z^*}p)\Lambda^{-1}=p .
\]
The form on the span of \(F,F'\) is nondegenerate, so the
Gram identities give \(\gamma_Z\in G_1\). The displayed formulas
are smooth on \(\Omega_{\rm Gr}\).\footnote{\weightedpin{grh-compact/src/New/RealPlace/GramSectionBounds.lean}
{25-L59}{The positive-discriminant domain and the section matrices}, and
\weightedpin{grh-compact/src/New/RealPlace/GramSection.lean}
{24-L61}{The section in the filling group}.}

The pointwise stabilizer of \(F\) is \(S=U^0H\), with unique
factorization \(u(B,A)h\), where \(H=U(W)\) acts trivially on the
other blocks and
\[
 u(B,A)=
 \begin{pmatrix}
   1&-B^*&0&A-\tfrac12 B^*B\\
   0&1&0&B\\
   0&0&1&0\\
   0&0&0&1
 \end{pmatrix},\qquad
 B\in M_{k,m}(\mathbf C),\quad A^*=-A .
\]
In fact the fixed frame forces the \(D\) and \(z\) rows
to agree with the identity. Pairings of the remaining columns
force the \(W\)-block \(h\) to be unitary, the \(E,W\)-block to
be \(-B^*h\), and the Hermitian part of the \(E,D\)-block to be
\(-B^*B/2\). The free entries of \(u(B,A)\) in the \(E,W\)-block
are \(-B^*\), and those in the \(E,D\)-block are the strictly
upper entries and the imaginary diagonal. At fixed \(B\) the
upper entries differ from those of \(A\) by a translation,
and the imaginary diagonal is the same as in \(A\).
Thus \(du=dB\,dA\) is the prescribed free-entry measure.
The \(U^0\) product adds the two \(B\)'s and adds the two \(A\)'s
with a skew Hermitian shear. Left and right multiplication
therefore preserve \(du\), as does conjugation by \(h\), which
sends \(B\) to \(hB\) and fixes \(A\). This also gives the
bi-invariant product Haar on \(S\).\footnote{\weightedpin{grh-compact/src/Lib/DiscreteSeries/Bessel/UnipotentCoords.lean}
{25-L80}{The prescribed frame and free entries}, and
\weightedpin{grh-compact/src/Lib/DiscreteSeries/Bessel/UnipotentCoords.lean}
{217-L244}{The frame stabilizer and unipotent Haar}.}

\paragraph{The normalized frame density.}
Use independent-entry Lebesgue measure on skew Hermitian
matrices, with complex off-diagonal entries and real imaginary
diagonal entries. Write \(d\omega_G\) and \(d\omega_H\) for the raw
Haar densities whose values at the identity use these coordinates
on \(JY\) and \(Y\), respectively. The filling measure is
\(dy=d\omega_G/\operatorname{vol}U(r)\). In fact the Cayley
map of a skew Hermitian matrix \(V\) is \(g=2C^{-1}-1\)
with \(C=1-J^{-1}V/2\). Its differential in left-translated
coordinates is \(Jg^{-1}dg=C^{-*}(dV)C^{-1}\), where
\(C^{-*}=(C^{-1})^*\). On Hermitian or skew Hermitian
matrices of size \(j\), congruence by a matrix \(C\) has
absolute real determinant \(|\det C|^{2j}\).
Unitary factors preserve Euclidean volume, whose density is a
constant times independent-entry Lebesgue measure. Positive
diagonal singular values
\((s_i)\) contribute \(s_i^2\) in the real diagonal and
\((s_is_j)^2\) in each complex off-diagonal position. This proves
the determinant formula and gives the filling density on the
Cayley chart with the prescribed divisor \(\operatorname{vol}U(r)\).
Its complement, where \(g+1\) is singular, has Haar measure
zero. For each \(g\) there are finitely many phases \(t\) in one
period with \(1+e^{it}g\) singular, so central invariance and
Tonelli give this null-set assertion.

For \(H=U(k)\), the last-column map has stabilizer \(U(k-1)\).
At the base column \(x_0\), the off-diagonal entries and imaginary
diagonal in that column give the Euclidean sphere tangent density
with absolute determinant one. The other entries give the same
raw Haar on \(U(k-1)\). The unitary action preserves sphere area,
so these densities disintegrate \(d\omega_H\). The sphere area
follows by polar integration of the Gaussian on \(\mathbf C^k\).
The compact volumes satisfy the resulting recurrence
\[
 s_k=\frac{2\pi^k}{(k-1)!},\qquad
 \operatorname{vol}U(k)=s_k\,\operatorname{vol}U(k-1),
 \qquad dh=\frac{d\omega_H}{\operatorname{vol}U(k)} .
\]
Start with \(\operatorname{vol}U(0)=1\). Here \(s_k\) is the
area of the unit sphere of \(W\) and \(dh\) is probability Haar.
Its image under \(h\mapsto hx_0\) is unit
sphere probability. These are the compact volumes in the filling
normalization.\footnote{\weightedpin{grh-compact/src/Lib/DiscreteSeries/Haar/UnitaryNorm.lean}
{25-L34}{Independent skew Hermitian entries}, and
\weightedpin{grh-compact/src/Lib/DiscreteSeries/Haar/CompactDualMeasure.lean}
{19-L48}{Cayley chart, density and compact-dual divisor}.}

On full-rank \(r\times n\) complex matrices \(X\), the map
\(q(X)=X^*JX\) is a submersion onto its image in the space of
Hermitian matrices. In its differential \(X^*J\,dX\) can be
arbitrary. On that Hermitian space use the density \(d\lambda(q)\)
in the independent complex off-diagonal entries and the real
diagonal entries \(q_{ii}/2\). Disintegrating ambient Lebesgue
\(dX\) against this density gives a smooth density \(d\mu\) on the
fibre \(q=-p\). The action of \(G_1\) preserves \(q\) and \(dX\),
since \(|\det y|=1\), and therefore preserves \(d\mu\).

At the base frame \(F\), put \(P=F^*J\,dX\). Then
\(dq=P+P^*\). For \(i<j\), passage from
\((P_{ij},\overline{P_{ji}})\) to the \(ij\) entries of
\((P-P^*)/2\) and \(dq\) has complex determinant \(1\).
On the diagonal the real part of \(P_{ii}\) gives \(dq_{ii}/2\)
and the imaginary part gives the independent skew entry.
Hence \(d\mu\) at \(F\) is Lebesgue on these skew entries together
with all entries of the \(F^\perp\)-block of \(J\,dX\), where
\(F^\perp\) denotes the Euclidean orthogonal complement.
The differential of \(y\mapsto yF\) reads just these entries
in the \(F\)-columns of \(JY\). Its kernel is the remaining
skew Hermitian block on \(F^\perp=(D,W)\).
On the tangent of \(u(B,A)h\) this remaining block is
\(\begin{psmallmatrix}A&-B^*\\ B&Y_H\end{psmallmatrix}\).
Its independent entries have density \(du\,d\omega_H\).

It follows that \(d\omega_G\) disintegrates over the frame orbit
with quotient density \(d\mu\) and right stabilizer density
\(du\,d\omega_H\). In detail, the differential just computed makes
the orbit open in the full-rank fibre and supplies smooth local
lifts \(y_X\) with \(y_XF=X\). Integrating a test on \(y_X S\)
against the raw stabilizer Haar and then \(d\mu(X)\) defines
a left-invariant smooth density. The two tangent densities at the
identity give \(d\omega_G\), fixing its scale. Inversion on \(S\)
preserves its Haar, so the right factor can be written
\((vh)^{-1}\) with density \(\operatorname{vol}U(k)\,dv\,dh\).

\paragraph{The density on the positive-discriminant cell.}
For \(X=yF\), use \(Z=F^*JX\) as part of the ambient coordinates.
It records the \(D\)-rows and the negative of the \(z\)-row.
At fixed \(Z\), right-multiply \(X\) by \(Z^{-1}\).
Write the \(E\)- and \(W\)-blocks of \(XZ^{-1}\) as
\((L,a)\) and \((B,x)\), with \(a\in\mathbf C^m\) and
\(x\in W\) the last columns. The other two blocks are fixed,
and the constraint is
\[
 q'=Z^{-*}q(X)Z^{-1}=
 \begin{pmatrix}
 L+L^*+B^*B & a+B^*x\\
 a^*+x^*B & \|x\|^2-1
 \end{pmatrix}=-K_Z .
\]
Returning to \(X\) has ambient Jacobian \(|\det Z|^{2(m+k)}\)
on the \(m+k\) free complex rows. Returning from \(q'\) to \(q\)
has Jacobian \(|\det Z|^{2n}\) on Hermitian matrices. Both are
changes with \(Z\) fixed, so their ratio is also the Jacobian in
the joint coordinates with \(dZ\).
In the upper-left block the change from \(L\) to its skew part
\(A=(L-L^*)/2\) and \(L+L^*\) has the same determinant as
at \(F\). The off-block \(a\) is translated by \(B^*x\).
For the last real coordinate use \(q'_{nn}/2=(\|x\|^2-1)/2\).
Polar integration at \(\|x\|^2=\delta(Z)>0\) leaves
\(\delta(Z)^{k-1}d\sigma(\xi)\), where \(x=\sqrt{\delta(Z)}\,\xi\)
and \(d\sigma\) is unit sphere area. Therefore the frame density
in these coordinates is
\[
 d\mu(X)=|\det Z|^{2(k-1)}\delta(Z)^{k-1}\,
         dZ\,dB\,dA\,d\sigma(\xi)
       =\mathfrak g(Z)^{k-1}\,dZ\,du\,d\sigma(\xi).
\]
All coordinates of \(X\) are recovered since
\[
 L=-\tfrac12\bigl((K_Z)_{11}+B^*B\bigr)+A,\qquad
 a=-(K_Z)_{1n}-B^*x.
\]
For a lift \(k_\xi\in H\) of the column \(\xi=k_\xi x_0\),
the frame \(u(B,A)k_\xi\gamma_ZF\) has exactly these blocks.
Thus the coordinates cover the whole positive-discriminant
part of the fibre. Two lifts differ by right multiplication by
\(H_0=U(W\ominus\mathbf C x_0)\), which commutes with \(\gamma_Z\).
For \(y=u k_\xi\gamma_Z s^{-1}\), the simultaneous change
\((k_\xi,s)\mapsto(k_\xi\ell,s\ell)\), \(\ell\in H_0\),
is precisely the redundancy: the frame determines \(Z,u,\xi\),
and for a fixed lift the right factor is unique.

The two disintegrations, with sphere probability supplied by
probability Haar \(dk_0\) on the left \(H\), now give
\[
\begin{aligned}
 \int_{\Omega_G} f(y)\,dy
   &=\int_H \mathcal I(k_0)\,dk_0,\\
 \mathcal I(k_0)
   &=\int_{\Omega_{\rm Gr}} c_G\,\mathfrak g(Z)^{k-1}
       \int_H\int_{U^0}\int_{U^0}
       f\bigl(u k_0\gamma_Z(vh)^{-1}\bigr)\,du\,dv\,dh\,dZ,\\
 c_G&=\frac{\operatorname{vol}U(k)}{\operatorname{vol}U(r)}
                  \frac{2\pi^k}{(k-1)!}.
\end{aligned}
\]
Here the right-factor integral is independent of the chosen lift
by right invariance on \(S\). For every nonnegative Borel
\(f\), Tonelli permits these extended integrals.
In particular \(\inf_{k_0\in H}\mathcal I(k_0)\le\int_{G_1} f\,dy\).
The infimum is over a single compact variable outside the
Gram and stabilizer integrals. This is the lower bound used in
the \hyperlink{proof.weighted-real-cell-contraction}{real cell
contraction}.\footnote{\weightedpin{grh-compact/src/New/Filling/BesselSystem.lean}
{18-L20}{The constant \(c_G\)},
\weightedpin{grh-compact/src/New/RealPlace/SobolevCell.lean}
{17-L34}{The lower bound used on the cube}, and
\weightedpin{grh/src/New/Spanning/Bessel/MeanValue.lean}
{581-L617}{The mean-value bound with both compact averages}.}

\paragraph{The Cayley-coordinate Jacobian.}
The column lift in the compact chart takes \((t,b)\in
\mathbf R\times\mathbf C^{k-1}\) to the Cayley transform
of \(\begin{psmallmatrix}0&b\\-b^*&it\end{psmallmatrix}\).
Its last column is
\[
 \xi=(b/\Delta,\ 2/\Delta-1),\qquad
 \Delta=1+\|b\|^2/4-it/2 .
\]
The inverse on the sphere off \(-x_0\) reads
\(\Delta=2/(1+\xi_k)\), \(b=\Delta\,\xi_{<k}\),
\(t=-2\operatorname{Im}\Delta\). The sphere relation gives
\(\operatorname{Re}\Delta=1+\|b\|^2/4\), so this is a smooth
section on that chart. To compute its density, first use
\((b,\Delta)\) as complex ambient coordinates in the displayed
map. Their Jacobian is \(4|\Delta|^{-2(k-1)-4}\).
On the sphere the absolute derivative of
\((\|\xi\|^2-1)/2\) in \(\operatorname{Re}\Delta\) is
\(2/|\Delta|^2\). After dividing by this derivative and using
\(d(\operatorname{Im}\Delta)=dt/2\) in absolute density, the
sphere measure is \(|\Delta|^{-2k}\,db\,dt\).
This is the left Cayley weight \(w_L\), since \(\Delta\) is the
determinant of its Cayley denominator.

On the right use the full compact Cayley chart, with weight
\(w_R=|\det(1-Y_R/2)|^{-2k}\) for skew Hermitian \(Y_R\).
Its inverse exists on the set where \(h+1\) is nonsingular.
With these two choices the group chart is injective. The
frame recovers \(Z,u,\xi\), the column chart recovers its left
lift, and \(y^{-1}u k_\xi\gamma_Z=vh\) recovers the right
parameters by the stabilizer factorization. These inverse formulas
are smooth on the chart image. The omitted sphere point and the
omitted right compact set have fibre measure zero.
The group chart thus pulls back \(d\omega_G\) with density
\(\mathfrak g(Z)^{k-1}w_Lw_R\) in the indicated Lebesgue
coordinates, and \(dy\) with this density divided by
\(\operatorname{vol}U(r)\). Splitting a skew Hermitian matrix
into independent entries, or isolating its last column,
preserves these coordinate densities. Integrating the left
weight gives \(s_k\) and integrating the right gives
\(\operatorname{vol}U(k)\), as in the normalized formula above.\footnote{\weightedpin{grh-compact/src/New/RealPlace/CellNeighbourhood.lean}
{49-L131}{Normalized Jacobian and injectivity in the compact chart}
(in particular \weightedpin{grh-compact/src/New/RealPlace/CellNeighbourhood.lean}
{49-L87}{the Jacobian} and
\weightedpin{grh-compact/src/New/RealPlace/CellNeighbourhood.lean}
{88-L123}{injectivity}), and
\weightedpin{grh-compact/src/Lib/DiscreteSeries/Haar/SphereFibreIntegral.lean}
{10-L35}{The integrated column Cayley weight}.}

\paragraph{The ordered Haar measure and its quotient.}
\leavevmode\teamanchor{proof.weighted-ordered-haar}%
Write \(V_i\) for the space after removing the preceding Witt
pairs and \(G_i^{\rm fl}=U(V_i)\), so \(G_0^{\rm fl}=G_1\)
in the filling notation. Write \(d\omega_i\) for the Haar
density whose value at the identity is independent-entry
Lebesgue measure on the skew Hermitian matrix \(J_iX\),
\(X\in\mathfrak u(V_i)\). The split frame has
\(J_i^2=1\). The Cayley chart used in the
\hyperlink{proof.weighted-compact-dual-degree}{filling measure}
identifies \(dy=d\omega_0/\operatorname{vol}U(r)\).
At the final compact group \(H=U(X_0)\), the prescribed
probability Haar is \(d\omega_n/\operatorname{vol}U(k)\).
These densities are both left and right invariant. For example,
the conjugation action on the Lie algebra has absolute real
determinant one since its complexification is conjugation
on the full matrix algebra.

At a Witt floor write the fixed frame vector as \(e\), its
opposite as \(f\), with \(h(e,f)=1\), and identify
\(V_i=\mathbf C e\oplus V_{i+1}\oplus\mathbf C f\).
The stabilizer \(S_i\) of \(e\) in \(G_i^{\rm fl}\) is
\(G_{i+1}^{\rm fl}N_i\), where \(N_i\) is the Heisenberg group
of that floor. Indeed \(g\in S_i\) has \(f\)-coefficient one
in \(gf\), and \(h(gf,gf)=0\) gives
\[
 gf=f+a+(it-h(a,a)/2)e
 \qquad(a\in V_{i+1},\ t\in\mathbf R).
\]
There is exactly one \(n(a,t)\in N_i\) with this value at \(f\).
Then \(n(a,t)^{-1}g\) fixes the pair. The compact or indefinite
unitary action of \(G_{i+1}^{\rm fl}\) on \(a\) preserves
Lebesgue measure and fixes \(t\). The group law in \(N_i\)
is translation in \(a\) and translation plus a linear shear
in \(t\). Thus \(d\omega_{i+1}*d\eta_i\) is bi-invariant
Haar on \(S_i\).

The orbit map \(g\mapsto g^{-1}e\) identifies \(S_i\backslash
G_i^{\rm fl}\) with the nonzero isotropic quadric. Use the
invariant measure on this quadric obtained by disintegrating
complex Lebesgue measure against \(d(h(v,v)/2)\). This is the
real place's skew coordinate \(\beta=(i/2)h(v,v)\), with
Lebesgue measure on \(\operatorname{Im}\beta\).
In the chart \(v=c(f+a+(u+it)e)\), \(c\ne0\), one has
\(h(v,v)/2=|c|^2(u+h(a,a)/2)\). The ambient Jacobian is
\(|c|^{2d_i+2}\), since \(d_i=\dim_{\mathbf C}V_{i+1}\).
Dividing by \(|c|^2\) gives
\(|c|^{2d_i}dc\,da\,dt\) on the quadric. The omitted
complex hyperplane has real dimension smaller than the smooth
quadric and hence has zero quadric measure. For
\(g=w_i(c)^{-1}n(a,t)\) its orbit value is
\(c n(a,t)^{-1}f\). Inversion in these Heisenberg coordinates
sends \((a,t)\) to \((-a,-t)\), which preserves the density.
Thus this invariant quadric measure is exactly the image
of the specified \(\theta_i\). At the point \(e\), use the
chart with \(e,f\) exchanged. Its quadric density at
\(c=1,a=t=0\) is one.

Integrating over the stabilizer on the left and this quotient
on the right gives a Haar measure. For right invariance, a
right translation takes the quotient to itself by its
unitary action and changes a chosen lift by left
multiplication in \(S_i\); the inner integral is unchanged
by right translation within \(S_i\). The measure has a
smooth positive local density, by a smooth local section.
To determine that density at the identity, put \(Y=J_iX\).
The differential of the orbit map there is \(-J_iYe\).
The complex entries in the \(e\)-column of \(Y\) and the
imaginary diagonal entry give the quotient coordinates
\((c,a,t)\) at \(e\) with absolute real determinant one.
Indeed \(J_i\) interchanges \(e,f\) and is unitary on
\(V_{i+1}\), and the \(e\)-diagonal is sent to the imaginary
\(f\)-coordinate. The kernel has this column zero.
Its remaining independent entries are exactly those of
the middle skew Hermitian block, the column coupling the
middle block to \(f\), and the imaginary \(f\)-diagonal.
These give the stabilizer density
\(d\omega_{i+1}*d\eta_i\), with absolute determinant one.
The infinitesimal action of \(n(a,t)\) on \(f\) is
\(a+it e\), so its last column in \(Y\) reads
\(J_{i+1}a+it f\). Hence
\[
                  d\omega_{i+1}*d\eta_i*d\theta_i=d\omega_i
                  \qquad(i<\varrho).
\]

At the terminal floor use \(V_\varrho=X_0\oplus\mathbf C z\),
\(h(z,z)=-1\) and the stabilizer of \(z\) is \(H\).
On \(h(a+cz,a+cz)=-1\), disintegration against \(d(h/2)\)
and ordinary polar integration in \(a\in\mathbf C^k\)
give the density
\(2\pi^k(|c|^2-1)_+^{k-1}dc/(k-1)!\) and the uniform
sphere probability in \(a\) of squared length \(|c|^2-1\).
That sphere probability is the image of the right compact
factor of \(\theta_\varrho\). The terminal boost sends \(z\)
to a vector with \(z\)-coordinate \(-c\) and first component
of squared length \(|c|^2-1\). The change \(c\mapsto-c\)
preserves the density, so this is precisely the orbit law. Near \(z\), write
\(c=e^{i\phi}\sqrt{1+\|a\|^2}\). The ambient polar
factor \(|c|\,d|c|\,d\phi\) in the last coordinate
becomes \(d(h/2)\,d\phi\) in absolute density, giving
\(da\,d\phi\) on the hyperboloid. The differential
\(-J_\varrho Yz\) identifies these coordinates at \(z\)
with the off-diagonal \(z\)-column and the imaginary
\(z\)-diagonal with absolute determinant one. The remaining
middle entries are the stabilizer density \(d\omega_n\).
This gives
\(d\omega_n*d\theta_\varrho=d\omega_\varrho\).
In particular no extra scalar comes from either probability
compact factor. Iterating the two floor formulas yields
\[
 \mathcal M=\frac{d\omega_0}{\operatorname{vol}U(k)}
     =\lambda_0\,dy,\qquad
 \lambda_0=\frac{\operatorname{vol}U(r)}{\operatorname{vol}U(k)}>0,
 \qquad
 \lambda_0\frac{\operatorname{vol}U(k)}{\operatorname{vol}U(r)}\le1.
\]

For the quotient formula, each later floor group fixes the
preceding pair and conjugates \(n(a,t)\) by a unitary change
of \(a\) with \(t\) fixed. Therefore \(\eta_i\) commutes
under convolution with the measures supported in the later
floor groups. Move these factors to the left in increasing
order:
\[
 \mathcal M=(\eta_0*\cdots *\eta_{\varrho-1})*
                            (\theta_n*\cdots *\theta_0).
\]
The first factor is Haar on \(U\). Successively fixing
the opposite Witt vectors writes each \(u\in U\) uniquely
in these ordered factors, with each later group acting on the
preceding Heisenberg group by the same measure preserving
conjugation. To fix the scale, the prescribed free entries
are in the row of \(e_i\), in the columns of the middle block
and the later Witt pairs, together with the imaginary part in
the column of \(f_i\). At the identity the \(i\)-th
Heisenberg factor gives the row \(-a^*J_{i+1}\) in those
complex entries and \(t\) in the real entry, with all other
free rows zero. The real determinant has absolute value one,
so the first factor is exactly \(du\). Put
\(\Theta_0=\theta_n*\cdots *\theta_0\).
Comparing the identity \(\mathcal M=du*\Theta_0\) with
\(dy=du\,d\dot y\) gives the image measure
\(\lambda_0\,d\dot y\) of \(\Theta_0\) on \(U\backslash G_1\).
Explicitly, use a measurable section of this quotient and
multiply a nonnegative quotient test by a function on \(U\)
of Haar integral one in the section coordinates. Integration
over \(U\) gives the test, so both iterated integrals give
the asserted quotient identity.

The lower bound for the full Gram chart is used separately
in the real cell contraction.
\footnote{\weightedpin{grh-compact/src/Lib/DiscreteSeries/Haar/UnitaryNorm.lean}
{25-L34}{Independent skew Hermitian Lie coordinates},
\weightedpin{grh-compact/src/Lib/DiscreteSeries/Haar/CompactDualMeasure.lean}
{19-L48}{Filling and compact densities},
\weightedpin{grh-compact/src/Lib/DiscreteSeries/Bruhat/ComplexPlace.lean}
{79-L155}{The terminal boost, polar density and compact factor},
\weightedpin{grh-compact/src/Lib/DiscreteSeries/Bessel/RealMeasures.lean}
{192-L210}{The free-entry unipotent measure}, and
\weightedpin{grh-compact/src/New/RealPlace/ComplexGauge.lean}
{43-L110}{The ordered and quotient Haar comparison on this datum}.}

Set%
\footnote{\weightedpin{grh-compact/src/New/RealPlace/LevelGauge.lean}{43-L62}{The floor measures},
\weightedpin{grh-compact/src/Lib/DiscreteSeries/Bruhat/ComplexPlace.lean}{52-L97}{The real floor densities},
\weightedpin{grh-compact/src/New/RealPlace/LocalTheorem.lean}{30-L67}{Weighted pairings on each floor},
and \weightedpin{grh-compact/src/New/RealPlace/ComplexGauge.lean}{31-L110}{The quotient and Haar comparison}.}
\[
 a_H=\operatorname{vol}U(k)/\operatorname{vol}U(r),\qquad
 J_i(p,q)=\frac{a_H}{(\lambda_0d)^2}\mathcal I_i(p,q;1).
\]

\paragraph{Common-entry descent on the real level.}
\leavevmode\teamanchor{proof.weighted-real-common-entry}%
Use the real cut vector \(F_\flat\) and the ordered measures above.
For \(0\le i\le\varrho\), use \(G_i^{\rm fl}=U(V_i)\) for the
group fixing the first \(i\) hyperbolic pairs. Write \(f_i\) for
the next frame vector, which is \(e_i\) for \(i<\varrho\) and
\(z\) for \(i=\varrho\). Its stabilizer in \(G_i^{\rm fl}\) is
\(S_i=G_i^{\rm fl}\cap Q_{i+1}\).
Let \(D=w_q^{-1}w_p\), with both Weyl products truncated at \(i\).
These products and \(D\) commute with \(G_i^{\rm fl}\).
The two upper pairings \(B_{i+1}^{1}\) and \(B_{i+1}^{D}\)
are invariant under simultaneous left multiplication by \(S_i\)
in their two arguments, by the
\hyperlink{proof.weighted-restricted-position-integral}{stabilizer
invariance of the upper averages}.

The floor measure \(\theta_i\) pushes forward to the invariant
quadric measure under \(\tau\mapsto\tau^{-1} f_i\).
It therefore gives the same integral to a left \(S_i\)-invariant
observable of \(\tau\) when \(\tau\) is replaced by \(\tau s\),
for any \(s\in G_i^{\rm fl}\). Indeed the observable descends
to the orbit and the substitution acts there by \(s^{-1}\).
The entry \(Z_i(v)=h(f_i,v f_i)\) is invariant under both left
and right multiplication by \(S_i\), since \(S_i\) fixes \(f_i\).
Write the ordered disintegration
\[
 d\theta_i
   =\bigl(c\mapsto v_i(c)^{-1}\bigr)_*
         \bigl(\rho_i(c)\,dc\bigr)*d\zeta_i .
\]
For \(i<\varrho\), \(v_i(c)\) is the Weyl element on the
next pair with entry \(c\), \(\zeta_i=\eta_i\) is its
Heisenberg measure, and \(\rho_i(c)=|c|^{2d_i}\) with
\(d_i=k-1+2(\varrho-i)\). At \(i=\varrho\), \(v_i(c)\) is
the terminal boost and \(\zeta_i=\theta_n\) is probability Haar
on \(H=U(X_0)\). The density from the skew-coordinate polar
calculation in the ordered Haar comparison is
\[
 \rho_\varrho(c)
   =\frac{2\pi^k}{(k-1)!}(|c|^2-1)_+^{k-1}.
\]
The skew base measure is \(d(\beta/2)\) for the Hermitian level
\(\beta=h(v,v)\). On the support
of \(\rho_i\) we have \(Z_i(v_i(c))=c\) almost everywhere.
In both cases \(\zeta_i\) is carried by \(S_i\).

Put \(\omega(v)=\operatorname{wave}(\xi Z_i(v))\).
In the integral \(\mathcal I_i(p,q;\omega)\), hold \(x,y,\sigma\)
fixed and first integrate over \(\tau\). Its integrand has the
three factors
\[
 \omega(\sigma\tau^{-1})\,
 B_{i+1}^{1}(\sigma y,\sigma yx)\,
 \overline{B_{i+1}^{D}(\tau w_p^{-1}y,\tau w_p^{-1}yx)} .
\]
The observable of \(\tau\) is left \(S_i\)-invariant, by the
right invariance of \(Z_i\) under \(S_i\) and the simultaneous
left invariance of the mixed pairing. Use the orbit change above
to replace \(\tau\) by \(\tau_0\sigma\) inside the integral
and then put \(\tau_0=v_i(c)^{-1}u\), \(u\in S_i\), using the
displayed disintegration. The phase argument becomes
\(\sigma\tau^{-1}=u^{-1}v_i(c)\), hence its entry is \(c\).
Put \(y'=u\sigma y\). The first pairing is unchanged when
\(\sigma y,\sigma yx\) are replaced by \(y',y'x\).
Since both \(u\) and \(\sigma\) lie in \(G_i^{\rm fl}\),
\[
 \tau w_p^{-1}y
    =v_i(c)^{-1}u\sigma w_p^{-1}y
    =(w_pv_i(c))^{-1}y' .
\]
The same change gives \((w_pv_i(c))^{-1}y'x\) in the
other argument. The phase is \(\operatorname{wave}(\xi c)\).
The new measure of \(y'\) is precisely
\(d(\zeta_i*\theta_i*M_i)\).

Apply the change first to the product of the absolute values
of the two pairings, using nonnegative integrals. Its initial
integral against \(\mathcal M\), \(M_i\) and the two copies of \(\theta_i\)
is the integral of the absolute majorant
\(\mathcal F_i^+\), finite by the
\hyperlink{proof.weighted-real-shell-adapter}{real shell estimate}
at \(t_1=1\). Thus the changed integral is absolutely convergent,
and Fubini also applies with its complex phase. Below the terminal
floor the new measure is \(M_{i+1}=\eta_i*\theta_i*M_i\).
Appending the same \(c\) to the two heads changes the Weyl
products to \(w_pv_i(c),w_qv_i(c)\), keeps the ratio \(D\),
and multiplies the head density by \(|c|^{2d_i}=\rho_i(c)\).
Write \(p\oplus c\) for replacing the entry at \(i\) by \(c\),
with the suffix ignored at level \(i+1\). The successive-average
identity for the two upper pairings
identifies their product with the form at the next floor and
weight one. Consequently the normalized forms satisfy
\begin{equation}\label{eq:weighted-real-common-entry}
 \int_{\mathbb C}J_{i+1}(p\oplus c,q\oplus c)
           \operatorname{wave}(\xi c)\,dc
 =\frac{a_H}{(\lambda_0 d)^2}
        \mathcal I_i(p,q;\operatorname{wave}(\xi Z_i)).
\end{equation}
For \(i<\varrho\) the appended entry at \(c=0\) can be chosen
as the unit \(1\), with the displayed phase still evaluated at
\(c=0\). This point has measure zero. The terminal case of
\eqref{eq:weighted-real-common-entry} uses the Gram computation
below.\footnote{\weightedpin{grh-compact/src/New/RealPlace/LevelGauge.lean}
{93-L120}{The orbit pushforward and its invariance},
\weightedpin{grh-compact/src/New/RealPlace/LocalTheorem.lean}
{89-L96}{Stabilizer invariance of the frame entry},
\weightedpin{grh-compact/src/New/RealPlace/FloorKernels.lean}
{38-L61}{The absolutely convergent upper averages} and
\weightedpin{grh-compact/src/New/RealPlace/LocalTheorem.lean}
{253-L320}{The two-translate change of variables}.}

\paragraph{The terminal Gram in the same measures.}
\leavevmode\teamanchor{proof.weighted-real-terminal-gram}%
At \(i=\varrho\) both \(\eta_\varrho\) and \(\eta_n\) are
point masses at the identity, and \(\theta_n=\zeta_\varrho\).
Thus \(\zeta_\varrho*\theta_\varrho*M_\varrho=\mathcal M=\lambda_0\,dy\).
The upper averages at floor \(n=\varrho+1\) are over probability
Haar on \(H\). The cut vector satisfies
\(F_\flat(hy)=\pi_H(h)F_\flat(y)\) with \(\pi_H(h)\) unitary.
As \(D\) commutes with \(H\), both displayed upper averages
therefore equal their raw pairings.
For \(|c|>1\) put \(\gamma_p=w_pv_\varrho(c)\) and
\(\gamma_q=w_qv_\varrho(c)\).
The commutation with the terminal group gives
\(D\gamma_p^{-1}=\gamma_q^{-1}\).
In the remaining integral in \(x,y'\) put \(y''=y'x\), using
left invariance of Haar for the \(x\)-integral.
The product of measures is \(\lambda_0^2\,dy'\,dy''\).
The two raw factors become
\[
 \langle F_\flat(y'),F_\flat(y'')\rangle
 \overline{\langle F_\flat(\gamma_p^{-1}y'),
                         F_\flat(\gamma_q^{-1}y'')\rangle}.
\]

For each terminal \(\gamma=w_p v\), \(v\in G_\varrho^{\rm fl}\),
the products of coordinates used in \(J_{F_\flat}(\gamma)\)
have compact support in the integration variable. In fact, if
both \(F_\flat(y)\) and \(F_\flat(\gamma^{-1}y)\) are nonzero,
take \(u\in U\) with \(u y\) in the cut box. Since
\(G_\varrho^{\rm fl}\subset Q_\varrho\), the second nonzero
value makes \(w_p^{-1}y\) near the box at \(\varrho\).
The \hyperlink{proof.weighted-triangular-reconstruction}{head
minor calculation} bounds \(u e_a^*\) for every \(a<\varrho\)
in any gauge in which \(p\) is a head. Use the test vectors
\(e_a^*\) in \eqref{eq:weighted-triangular-inverse}, with
\(\nu=u\) and \(g=1\). The vectors and
\((I-u)e_a^*\) are bounded, so the entries of \(u^{-1}\) are
bounded in that gauge. Hence so are the entries of
\(y=u^{-1}(uy)\). The isometry group is closed in matrix
space, so this bounds the support in a compact set.
The same gauge works locally in the nonzero head variables.
This gives absolute integrability of the tensors and
continuity by dominated convergence in the boost and heads.

Recall
\[
 J_{F_\flat}(\gamma)_{a,b}
   =\int F_{\flat,a}(y)
             \overline{F_{\flat,b}(\gamma^{-1}y)}\,dy.
\]
Expanding the two raw factors, summing the finite coordinates,
and applying Fubini on the compact support gives
\(\lambda_0^2\langle J_{F_\flat}(\gamma_p),
J_{F_\flat}(\gamma_q)\rangle\), with the inner product linear
in its second argument. The normalization of the form
cancels \(\lambda_0^2\). Write
\[
 c_G=a_H\,\frac{2\pi^k}{(k-1)!},\qquad
 \Xi(p;c)=
 \frac{\sqrt{c_G}}d
 (|c|^2-1)_+^{(k-1)/2}
 \prod_{a<\varrho}|p_a|^{d_a}
 J_{F_\flat}(w_pv_\varrho(c)).
\]
Define \(J_n(p,q)\) on all unit heads by the inner product of
\(\Xi(p;p_\varrho)\) and \(\Xi(q;q_\varrho)\). These top values
form a positive semidefinite Gram kernel. The density
\(\rho_\varrho(c)\) multiplied by the earlier
head density supplies exactly the two scalar factors in
\(\langle\Xi(p;c),\Xi(q;c)\rangle\). Thus the new inner integral
at \(i=\varrho\), with its density and normalization,
is \(J_n(p\oplus c,q\oplus c)=\langle\Xi(p;c),\Xi(q;c)\rangle\).
Both sides vanish for \(|c|\le1\), since \(k>1\). The padding
unit \(1\) at \(c=0\) also gives a zero terminal vector. The density and correlation
formula prove \eqref{eq:weighted-real-common-entry} also at
the terminal floor.

On the actual terminal group let \(\gamma\in G_\varrho^{\rm fl}\)
with \(h(z,\gamma z)=c\). Its positive component of \(\gamma z\)
has squared length \(|c|^2-1\), since \(h(\gamma z,\gamma z)=-1\).
For \(|c|>1\), \(H\) acts transitively on the sphere of that
length and fixes \(z\). Take \(h_1\in H\) sending
\(v_\varrho(c)z\) to \(\gamma z\); then
\(h_2=(h_1v_\varrho(c))^{-1}\gamma\) fixes \(z\) and lies in
\(H\). Hence \(\gamma=h_1v_\varrho(c)h_2\).
Changing variables in the tensor integral and using the unitary
law for \(F_\flat\) gives for \(a,b\in H\)
\[
 J_{F_\flat}(a\gamma b)
   =\bigl(\pi_H(a)\otimes\overline{\pi_H(b^{-1})}\bigr)
                                      J_{F_\flat}(\gamma).
\]
The \(w_p,w_q\) commute with \(h_1\), so replacing the boost by
\(\gamma\) applies the same unitary operator to both tensors.
At \(|c|=1\) the polar prefactors give zero for either choice.
This is exactly the top Gram on the prescribed terminal
element and head data. In particular the real inserted-frequency
bound on the right of \eqref{eq:weighted-real-common-entry}
applies to the same sequence of forms whose top belongs to
the open packet. The unequal-head bound applies directly at its
own floor, and the bottom estimate concerns \(J_0\) in this
same normalization.%
\footnote{\weightedpin{grh-compact/src/Lib/DiscreteSeries/Bruhat/ComplexPlace.lean}
{112-L158}{Terminal orbit, Cartan decomposition and polar law},
\weightedpin{grh-compact/src/New/RealPlace/LocalTheorem.lean}
{183-L187}{Terminal stack},
\weightedpin{grh-compact/src/New/RealPlace/LocalTheorem.lean}
{207-L227}{Compact support and the tensor identity},
\weightedpin{grh-compact/src/New/RealPlace/BandLimits.lean}
{60-L84}{The terminal Gram in the sieve} and
\weightedpin{grh-compact/src/New/Filling/BesselSystem.lean}
{64-L77}{Normalized correlations on the terminal element}. The fixed source also gives
\weightedpin{grh-compact/src/New/RealPlace/LocalTheorem.lean}{267-L320}
{the common-entry descent},
\weightedpin{grh-compact/src/New/RealPlace/BandLimits.lean}{65-L84}
{the terminal descent in these measures}, and
\weightedpin{grh-compact/src/New/RealPlace/ComplexShells.lean}{106-L112}
{absolute convergence of the majorant}.}

\leavevmode\teamanchor{proof.weighted-real-shell-adapter}%
Put \(\rho=R=b^{6n+2}\) and \(L=L_1\ge\rho\). The real level
has \(g=1\), with \(k-1\) orthonormal coordinates and exponent
\(s_0=k-3n-2\). Use the ordered measure comparison above,
\(dM=\lambda_0\,dy\) and \(dM_0^+=\lambda_0\,d\dot y\) on
left \(U\)-invariant integrands, with \(a_H\lambda_0\le1\).
For the first vector use \(F_\flat\). For the second take a
Bessel vector \(W'\) with the same character, supported near
floor \(j\), with squared quotient norm at most \(\epsilon_1^2d\).
The pointwise inner-product inequality gives
\[
 \mathcal N(F_\flat,F_\flat)\le(\lambda_0d)^2,\qquad
 \mathcal N(F_\flat,W')\le \epsilon_1^2(\lambda_0d)^2,
\]
since the squared quotient norm of \(F_\flat\) is at most \(d\).
Thus the factor \(N_0=\epsilon_1(\lambda_0d)^2\) becomes
\(a_H\epsilon_1\) after multiplication by the form constant.

In Lebesgue coordinates \(z=x+iy\), the oscillator character is
\(\operatorname{wave}(\xi z)=
 e^{2\pi i((\operatorname{Im}\xi)x+(\operatorname{Re}\xi)y)/\pi}\).
The change to Euclidean frequency gives Plancherel constant
\(\pi^2\) in one complex coordinate. Multiplication by \(\xi\)
has modulus \(|\xi|^2\). Hence the shared one-floor estimates
have squared factor \((\pi^2/|\xi|^2)^{k-1}\), both for the
inserted diagonal frequency and, when \(j=i+1\le\varrho\),
for the Heisenberg frequency \(\xi=p_i^{-1}-q_i^{-1}\).
The floor-modulus product is one at \(g=1\).

Fix \(t_1\ge1\) and \(0\le f\le\mathcal F_j^+\). When \(t_1>1\)
assume in addition \(f^2\le((\pi\rho^{24})^k/t_1^{s_0})^2\mathcal D\).
Put \(T_i=t_1L^{11ni}\) and let \(B_i\) be the positions where
every entry of \(x-I\) has absolute value at most \(T_{i+1}^{99n}\).
The sets \(B_i\) increase to the whole position set. In the annulus
\(A_i=B_i\setminus B_{i-1}\), with \(B_{-1}=\varnothing\),
the restricted integral uses a second gauge of base \(T_{i+1}^{99n}\)
and constant \(T_{i+1}^{198n}\).
\leavevmode\teamanchor{proof.weighted-triangular-real}%
At the real place \(A=\mathbf C\), \(g=1\), the cut box is the
operator-norm ball of radius \(R_{\rm cut}\) and a gauge of base
\(L\ge R_{\rm cut}\) has scalar size \(i\) given by \(|c|\le L^i\).
With \(R_{\rm cut}=R\) and \(t=t_1\), let the nested position
sets have entrywise radii \(T_{i+1}^{99n}\) for \(x-I\). On annulus
\(i\ge1\), a position is outside the preceding set of radius
\(T_i^{99n}\). The head gauge has base \(L\) and the estimating
gauge has base \(T_i\). When \(j>0\), the head scaling holds because
\(L^{10n}\le T_i\). At an active pair both values of the cut vector
are nonzero, so the two \(U\)-translates into the cut box exist.
Their difference has size two. The required Weyl translates are
near by the active-pair condition.

If each \(L_a(z,x)\), \(j\le a<\varrho\), had area greater than
\(\pi/T_i^2\), then in each there would be \(c_a\) outside the open
disk of radius \(1/T_i\). This is a nonzero complex number, with
\(|c_a^{-1}|\le T_i\). The
\hyperlink{proof.weighted-triangular-reconstruction}{reconstruction}
would give \(|(x-I)_{bc}|\le T_i^{99n}\) at every entry, contradicting
the annulus condition. The same implication rules out an active pair
on that annulus when \(j=\varrho\). On the initial set \(i=0\) the
application instead uses the supplied pointwise frequency bound when
\(t>1\), and the unrestricted cap when \(t=1\).

Thus the small cap on \(A_i\), \(i\ge1\), is \(\pi/T_i^2\).
On \(A_0\), take the supplied pointwise bound when \(t_1>1\)
and the bound with cap one when \(t_1=1\).

The cut gauge \(\Gamma_\rho\) supplies the current floor mass.
Its size-seven fibre and the \(k-1\) size-eleven coordinates give
\[
 V\le \rho^{2574n}(\pi\rho^{22})^{k-1}
   \le\rho^{2574n+24(k-1)}.
\]
The constant of the cut gauge is \(\rho^2\), and
\(\rho\ge4\) gives the last inequality. With cap
\(\pi/T_i^2\) the pointwise factor from
\eqref{eq:weighted-capped-position} is bounded by
\(\pi^{s_0/2}T_i^{-s_0}\rho^{2574n+24(k-1)}\).
All three cases therefore allow
\(E_i=(\pi\rho^{24})^k\rho^{9999n}T_i^{-s_0}\).
In the restricted integral the head constant is \(L^2\), and
\(T_{i+1}^{198n\cdot50000n^2}
=T_i^{9900000n^3}L^{108900000n^4}\).
The normalized bound on \(A_i\) is thus
\[
 \Theta\epsilon_1\rho^{24k+9999n}
 L^{108900000n^4+3000000n^3+2}
 T_i^{9900000n^3-s_0},
 \qquad\Theta=a_H\pi^{(2\varrho+1)k}.
\]
The power of \(\pi\) uses \(k+2(k-1)\varrho\le(2\varrho+1)k\).
As \(s_0>10^9n^3\), summing the last power of \(T_i\)
costs at most \(2t_1^{9900000n^3-s_0}\).
Use \(\rho\le L\), \(L\ge2\), and
\(108900000n^4+3000000n^3+9999n+3\le10^9n^4\).
This gives the bound, including convergence at \(t_1=1\),%
\footnote{\weightedpin{grh-compact/src/New/RealPlace/ComplexGauge.lean}
{29-L53}{The real level and its ordered Haar comparisons},
\weightedpin{grh-compact/src/New/RealPlace/ComplexGauge.lean}
{154-L172}{Complex sizes and their constant},
\weightedpin{grh-compact/src/Lib/DiscreteSeries/Bruhat/ComplexPlace.lean}
{27-L49}{Fourier normalization and orthonormal coordinates} and
\weightedpin{grh-compact/src/New/RealPlace/ComplexShells.lean}
{52-L110}{Mixed absolute integrals on the annuli}. The fixed source also records \weightedpin{grh-compact/src/New/RealPlace/WeightPlancherel.lean}{18-L55}{the floor volume weights},
\weightedpin{grh-compact/src/New/RealPlace/LocalTheorem.lean}{389-L452}{the diagonal and unequal-head Fourier formulas},
\weightedpin{grh-compact/src/New/RealPlace/ComplexShells.lean}{51-L109}{the complex shell sum}, and
\weightedpin{grh-compact/src/New/RealPlace/ComplexShells.lean}{69-L105}{the real gauge and cap}.}
\begin{equation}\label{eq:weighted-real-shell-integral}
 \frac{a_H}{(\lambda_0d)^2}\kappa_j\int f\,dM
 \le \Theta L_1^{10^9n^4}R^{24k}t_1^{9900000n^3-s_0}\epsilon_1 .
\end{equation}

Let \(\ell_K(u)=u^K\min(1,u)^6\), and put
\[
 H_j(p)=\prod_{a<j}\ell_{n^6}(|p_0\cdots p_a|^{-1}),\qquad
 \mathfrak b(\xi)=\Theta b^{2000n^9}\ell_{k-n^7}(|\xi|^{-1}).
\]
Here \(\xi\ne0\). In the chosen real data \(n\ge10^8\),
\(k\ge n^7\), and \(s_0>10^9n^3\).
Take \(t_1=\max(1,\rho^{24}|\xi|)\) and \(\epsilon_1=1\).
For \(t_1>1\), \((\pi/|\xi|)^{k-1}\le
(\pi\rho^{24})^k/t_1^{s_0}\), since \(k-1\ge s_0\).
If there is no active pair anywhere, the absolute majorant
vanishes, so the desired bound holds. Otherwise an active pair
passes both near-frame tests, and each direct partial product is a
leading minor in the
\hyperlink{proof.weighted-triangular-reconstruction}{head calculation}
and has modulus at most
\(R_*=\rho^{9n}\). The common head radius is the maximum
of \(\rho\) and the inverse moduli of the partial products.
For \(0<a\le R_*\), \(R_*^{n^6+6}\ell_{n^6}(1/a)\)
dominates both \(1\) and \(a^{-n^6}\). Applying this to each
partial product gives
\[
 L^{10^9n^4}\le L^{n^6}
   \le R_*^{(n^6+6)(2j+1)}H_j(p)H_j(q).
\]
Now \(s_0-9900000n^3\ge k-n^7+6\) and
\(\rho^{24}/t_1\le|\xi|^{-1}\), \(t_1^{-1}\le\min(1,|\xi|^{-1})\).
Therefore \(\rho^{24k}t_1^{9900000n^3-s_0}\) is at most
\(\rho^{24n^7}\ell_{k-n^7}(|\xi|^{-1})\).
With \(\rho=b^{6n+2}\), the resulting power is bounded by
\(b^{2000n^9}\). This proves \(\mathfrak b(\xi)H_j(p)H_j(q)\)
for the inserted integral at floor \(j\) and for the
unequal-head test at \(j=i+1\le\varrho\).
\footnote{\weightedpin{grh-compact/src/New/RealPlace/Size.lean}
{148-L164}{The leading partial products},
\weightedpin{grh-compact/src/New/RealPlace/ComplexShells.lean}
{38-L49}{The common head radius} and
\weightedpin{grh-compact/src/New/RealPlace/ComplexShells.lean}
{112-L144}{The two uses of the band bound}. The fixed source also records \weightedpin{grh-compact/src/New/Filling/EichlerFlag.lean}{65-L91}{the selected real frame}, and
\weightedpin{grh-compact/src/New/RealPlace/ComplexShells.lean}{113-L144}{the two frequency uses}.}

At the empty head write
\(Q_{A,B}(x)=\int_{U\backslash G_1}
\langle A(y),B(yx)\rangle\,d\dot y\), and
\(E_\flat=F_\flat-W\) for the full vector \(W\).
Both upper averages give \(\lambda_0\), and the position
measure supplies a third factor. Hence
\[
 J_0=\frac{a_H\lambda_0}{d^2}
         \int |Q_{F_\flat,F_\flat}(x)|^2\,dx .
\]
For a compactly supported scalar test \(h\), the Schur estimate
of \(W\) yields
\[
 \left|\int h(x)Q_{F_\flat,W}(x)\,dx\right|
 \le \|F_\flat\|_{L^2(d\dot y)}(d/\delta)^{1/2}\|h\|_{L^2(dx)}
 \le d\delta^{-1/2}\|h\|_{L^2(dx)},
\]
where \(\delta=\deg_v\). Fubini on this support follows by
Cauchy--Schwarz in the quotient and right invariance. Thus
\(L^2\) duality, together with \(a_H\lambda_0\le1\), bounds
the normalized integral of \(|Q_{F_\flat,W}|^2\)
by \(1/\delta\). The inequality
\[
 |Q_{F_\flat,F_\flat}|^2
 \le |Q_{F_\flat,W}|^2+
        2|Q_{F_\flat,F_\flat}Q_{F_\flat,E_\flat}|
\]
leaves the mixed absolute majorant. The difference \(E_\flat\)
is bounded and has the same Bessel character. At floor zero
every point is near, by translating to the identity in the box.
In \eqref{eq:weighted-real-shell-integral} take \(L_1=R\),
\(t_1=1\) and \(\epsilon_1=b^{-5r^2}\) from the squared cut norm.
The mixed contribution is at most
\(2\Theta R^{10^9n^4+24k}b^{-5r^2}\).
The Weyl degree and the compact sphere factors give
\(\delta\Theta\le b^{4r^2}\): the factorials of indices \(k\)
through \(r-1\) cancel and \(\pi\) has nonpositive power.
In the actual parameters \(10^9n^4<nk\),
\((6n+2)(n+24)<7n^2\), and \(7n^2k+11\le r^2\).
As \(b\ge2\), the mixed term is at most \(b^{-10}/\delta\).
This proves \(J_0\le(1+b^{-10})/\deg_v\) for the same forms
whose highest member is the normalized top Gram.
\footnote{\weightedpin{grh-compact/src/New/RealPlace/BandLimits.lean}
{85-L124}{The Schur term and the mixed majorant},
\weightedpin{grh-compact/src/New/Filling/BesselSystem.lean}
{82-L96}{Squared norms of the cut pair}, and
\weightedpin{grh-compact/src/New/Filling/FillingData.lean}
{59-L66}{Actual weight size and real rank},
\weightedpin{grh-compact/src/New/Filling/FillingData.lean}
{110-L118}{The Weyl-factor bound for the band constant}.}

Thus the actual top packet Gram has the positive top form,
the diagonal descent and the numerical bottom and frequency
bounds for the real part of the sieve.
\footnote{\weightedpin{grh-compact/src/New/RealPlace/BandLimits.lean}{85-L125}{Mixed Schur and cut-error calculation at the bottom},
\weightedpin{grh-compact/src/New/Filling/FillingData.lean}{94-L116}{Weyl degree and band constant},
and \weightedpin{grh-compact/src/New/Filling/GlobalSieve.lean}{167-L184}{The real and finite forms on the same packet}.}

\paragraph{Rational open orbits and their finite packets.}
\teamanchor{proof.weighted-open-packets}%
The next three local arguments apply to the finite sets \(T\)
and fixing subgroups \(Z\) used in the rational packet
comparison below. We use a single coset vector for the
terminal fibre and the two floor identities.

\paragraph{The finite orbital tensors and their terminal fibres.}
\leavevmode\teamanchor{proof.weighted-finite-packet-fibre}%
Fix the conjugation-stable set \(T\) of the finite packet calculation.
Write \(A_T=\prod_{w\in T}E_w\), \(R_T=\prod_{w\in T}\mathcal O_w\),
and \(z\) for the marked vector in the filling lattice basis.
Use \(h(v,w)=v^*Jw\), so \(h(z,z)=-N\).
The additive measure gives \(R_T\) mass one, and \(\psi_T\) is
the restriction of \(\psi_f(\,\cdot\,/\varpi)\) to \(A_T\).
The inverse different calculation makes \(R_T\) its own annihilator,
so this is the self-dual measure for either sign of the character.
Take \(K_0=K(g)_T\),
\(K=K(1)_T\), and \(K_0\subseteq Z\subseteq K\).
For \(Q=K_0\backslash Z\) use counting measure. Its vector is
\[
 W_a(y)=\chi_g(ya^{-1})\quad(a\in Q),\qquad
 \chi_g(uk)=\psi_T\bigl(g^{-1}\operatorname{sr}(u)\bigr)
                 \quad(u\in U,\ k\in K_0),
\]
with zero value off \(UK_0\). Here \(\operatorname{sr}\) is the
sum of simple-root entries. It is additive on \(U\) and belongs
to \(gR_T\) on \(U\cap K_0\), so the value is well defined.
Right \(K_0\)-invariance of \(\chi_g\) makes the formula independent
of the chosen representative of \(a\). Normality of \(K_0\) in
\(K\) also gives \(W(yk)=W(y)\) for \(k\in K_0\).
In particular \(W\) is locally constant, bounded by \(\sqrt{|Q|}\),
and \(W(y)\ne0\) implies \(uy\in K\) for some \(u\in U\).
These are the same coordinates used in the
\hyperlink{proof.weighted-finite-denominator-shells}{finite shell bound}.

Let \(G\) be the isometry group over \(A_T\), let \(G_i\subset G\)
fix the first \(i\) hyperbolic pairs for \(0\le i\le\varrho\), and
let \(H\subset G_\varrho\) fix \(z\).
Put \(t(l)=h(z,l^{-1}z)\), \(e(l)=l^{-1}z\) on \(G_\varrho\).
The terminal floor measure \(\theta_\varrho\) pushes forward by
\(e\) to the norm-\(-N\) quadric measure \(\delta\).
Let \(dh\) be right Haar on \(H\) with \(H\cap K\) of mass one.
The terminal group measure and the full group measure are
\[
 \nu_T=dh*\theta_\varrho,\qquad
 \mu=\nu_T*M_\varrho,\qquad
 M_i=(\eta_{i-1}*\theta_{i-1})*\cdots*(\eta_0*\theta_0),
 \quad M_0=\delta_1 .
\]
Here \(\eta_i\) is the Heisenberg measure and \(\theta_i\)
the next-vector floor measure. Thus the terminal stabilizer in
\(\nu_T\) is integrated with its actual Haar measure.
Indeed \(l,l'\) over \(e\) differ by a left factor of \(H\).
For a nonnegative test \(F\), the fibre integral
\(\int_H F(hl)\,dh\) is independent of \(l\) by right invariance.
Right multiplication of \(l\) by \(s\in G_\varrho\) acts by
\(s^{-1}\) on \(e\), preserving \(\delta\). This proves the right
invariance of \(\nu_T\). For compact \(C\subset G_\varrho\), any
nonempty fibre in \(C\) has Haar mass at most
\(\operatorname{vol}_H(H\cap CC^{-1})\), and its image lies
in the compact set \(e(C)\). The quadric measure is locally finite
there by the regular norm chart. Hence \(\nu_T(C)<\infty\).

The same orbit argument at each lower floor yields the ordered
Haar \(\mu\). In detail, the stabilizer of the next isotropic
vector in \(G_i\) is generated by the Heisenberg group at \(i\)
and \(G_{i+1}\). The latter preserves horizontal volume by its
isometry action and fixes the skew centre. Its right Haar followed
by \(\eta_i\) is therefore right invariant on this stabilizer.
Integrating that Haar along \(\theta_i\) gives right Haar on \(G_i\)
by the invariant quadric pushforward. Compact sets project to
compact sets of regular orbit points. The same fibre bound as
above gives local finiteness, and an open set meeting the group
has positive measure by the local quadric chart.
The full group is unimodular. Indeed its adjoint action preserves
the trace pairing \(\operatorname{tr}(XY)\) on its skew-adjoint
Lie algebra over the local fixed field. This pairing is
nondegenerate, as extension to a splitting field gives the trace
pairing on all matrices. Thus the absolute adjoint determinant
is one, also in a split pair. The left translate of a right
invariant volume has this same density, so \(\mu\) is left Haar
as well.\footnote{\weightedpin{grh-compact/src/New/Filling/EichlerFlag.lean}
{15-L47}{The form in the filling lattice basis},
\weightedpin{grh-compact/src/New/RealPlace/LevelGauge.lean}
{23-L61}{The ordered measures},
\weightedpin{grh-compact/src/New/RealPlace/LevelGauge.lean}
{93-L125}{Quadric pushforward and isometry invariance} and
\weightedpin{grh-compact/src/New/Filling/StandardLevel.lean}
{218-L258}{The finite level volumes and Haar scale}.}

For unit heads \(p\), let \(w_p\) use the first \(\varrho\)
entries, \(d_i=k-1+2(\varrho-i)\) for \(i<\varrho\), and put
\[
 \kappa_p=\prod_{i<\varrho}\operatorname{md}_T(p_i)^{d_i/2},
 \qquad
 C_p(l)_{ab}=\kappa_p\int_G W_a(y)
             \overline{W_b(lw_p^{-1}y)}\,d\mu(y).
\]
This is the tensor used in the rational packet calculation below. We record
the support needed in these integrals. For fixed heads, the
\hyperlink{proof.weighted-triangular-reconstruction}{triangular
reconstruction} bounds \(y\) uniformly in \(l\in G_\varrho\)
whenever \(W(y)\) and \(W(l w_p^{-1}y)\) are both nonzero.
Indeed take \(u\in U\) with \(uy\in K\). The second condition
puts \(w_p^{-1}y\) in \(Q_\varrho K\), where \(Q_\varrho\)
is the group unitriangular on the earlier frame. The head-minor
calculation bounds all columns \(u e_i^*\). Applying the triangular
inverse formula to \(e_i^*\) bounds \(u^{-1}\), hence \(y\).
One gauge works in a neighbourhood of a finite list of heads and
their conjugates, since there are only finitely many inverse partial
products to put in its first size ball. Its bounded matrix set in
the closed isometry group is compact. Now write \(y'=lw_p^{-1}y\).
The Weyl blocks give \(w_{\bar p}=w_p^{-1}\), and \(G_\varrho\)
commutes with them. Thus \(y=l^{-1}w_{\bar p}^{-1}y'\).
The same bound with the conjugate head bounds \(y'\).
Finally \(l=y'y^{-1}w_p\) lies in a compact set as well.
At a split pair the involution interchanges the two components
when forming \(\bar p\). The conjugation-stable gauge controls both.

In particular the four-vector product used to compute
\(\langle C_p(l),C_q(l)\rangle\) is bounded and has common compact
support in \(l,y_1,y_2\) for nearby \(p,q\). There is also uniform
local constancy. For \(p'\) near \(p\) the head modulus is unchanged.
For \(s\in G_\varrho\) sufficiently near \(1\),
\(y^{-1}w_p s w_{p'}^{-1}y\in K_0\) throughout the compact set
just obtained. This follows from continuity and compactness, since
\(K_0\) is open. The right \(K_0\)-invariance of \(W\) now gives
\(C_{p'}(ls)=C_p(l)\) for every \(l\in G_\varrho\). Outside that
compact \(y\)-set both integrands vanish whenever their first
vector is nonzero. For \(a\in H\cap K_0\) we also have
\(W(ay)=W(y)\), since \(a\) normalizes \(U\) and fixes the
simple-root character. Consequently
\[
                     C_p(al)=C_p(l)=C_p(la).
\]
For the second equality change the integration variable by left
multiplication by \(a^{-1}\) and commute \(a\) with \(w_p\).\footnote{\weightedpin{grh-compact/src/Lib/DiscreteSeries/Bruhat/FloorGroups.lean}
{200-L214}{The Weyl blocks and the conjugate head},
\weightedpin{grh-compact/src/New/RealPlace/Size.lean}
{148-L174}{Head columns and the uniform support bound},
\weightedpin{grh/src/New/Spanning/Frames/TopCorrelation.lean}
{438-L462}{Applying the support bound in both directions} and
\weightedpin{grh-compact/src/New/Filling/TopFloor.lean}
{423-L431}{The finite tensor in the level measure}.}

For a finite family \(p\) with a common last entry \(t\), let
\(R\) be any positive semidefinite matrix on its index set.
For \(d=(a,b)\in Q^2\), write \(v_d(l)_p=C_p(l)_{ab}\).
With the inner product linear in its second argument,
\[
       \sum_{p,q}R_{pq}\langle C_p(l),C_q(l)\rangle
                  =\sum_{d\in Q^2}v_d(l)^*R\,v_d(l)\ \ge 0 .
\]
Integrate this expression on \(t(l)-t\in L\) for a compact open
additive subgroup \(L\subset A_T\). For arbitrary heads set
\[
 \mathcal K_L(p,q)
   =\int_{t(l)-p_\varrho\in L}
         \langle C_p(l),C_q(l)\rangle\,d\nu_T(l).
\]
Thus the integral of the positive form is
\(\Re\sum_{p,q}R_{pq}\mathcal K_L(p,q)\). At a common last entry this is the
Gram matrix of the functions \(1_{t+L}(t(l))C_p(l)\) in
\(L^2(d\nu_T;\mathbf C^{Q^2})\).
The integral is jointly locally constant in the two unit heads,
since the tensors agree uniformly for nearby heads and changing
\(p_\varrho\) inside the same coset of \(L\) preserves the aperture.

On a terminal fibre put
\[
 E_R(e)=\inf_{l^{-1}z=e}\int_H
           \sum_d v_d(hl)^*R\,v_d(hl)\,dh .
\]
The infimum ranges over \(G_\varrho\), with value \(+\infty\) for
an empty fibre. On the orbit its value equals every displayed
integral by right invariance of \(dh\). The function on that orbit
is measurable. In fact there are continuous local sections there:
at a fixed point \(e_0\), for \(v\) on its quadric near \(e_0\),
put \(a=e_0+v\), \(D=h(a,e_0)\), and
\[
       r_v x=x-a\,h(a,x)/D,\qquad s(v)=r_v r_{e_0}^{-1}
                       \quad (x\in V_\varrho).
\]
The element \(D\) is a unit of \(A_T\) near \(2h(e_0,e_0)=-2N\).
Since \(D+\bar D=h(a,a)\), direct expansion shows that \(r_v\)
preserves \(h\). It sends \(e_0\) to \(-v\), so \(s(v)e_0=v\)
and \(s(e_0)=1\). Extend by the identity on the hyperbolic
pairs. The same construction works about any point of the
norm quadric, so each orbit is open in the quadric.
Integration along a section and then against \(\delta\) now gives
\[
  \int_{h(z,e)-t\in L}E_R(e)\,d\delta(e)
                =\Re\sum_{p,q}R_{pq}\mathcal K_L(p,q).
\]
The pushforward description of \(\delta\) gives full measure to
the orbit in this identity.

Here is the consequence on individual labels. Suppose a finite
list \((l_j,d_j)\) lies over \(e_0\), and
\(l_jl_{j'}^{-1}\in K_0,\ d_j=d_{j'}\) imply \(j=j'\).
For \(v\) near \(e_0\) on the quadric, \(v=a_0y\) with
\(a_0\in H\cap K_0\), take \(s=s(v)^{-1}\). The elements
\(l_jsa_0\) lie over \(y\) and have exactly the original
\(C_p(l_j)\) for every head on the list. This uses the uniform
constancy for right multiplication by \(s\), followed by the
right \(H\cap K_0\)-invariance. For any \(l\) over \(y\) put
\(h_j=l_jsa_0l^{-1}\in H\). In the integral over \(H\), retain
coordinate \(d_j\) on \((H\cap K_0)h_j\). It is constant there by
left invariance of the tensor under \(H\cap K_0\). For a given
coordinate these cosets are disjoint, since
\(h_jh_{j'}^{-1}=l_jl_{j'}^{-1}\). Positivity on every coordinate
gives throughout this neighbourhood
\[
 E_R(y)\ \ge\ \operatorname{vol}_H(H\cap K_0)
             \sum_j v_{d_j}(l_j)^*R\,v_{d_j}(l_j).
\]
For \(y\) off the orbit the same inequality holds with infinite
left side. This provides the neighbourhood and the individual
coset mass for the one-vector count. No averaging factor occurs
in the coordinates of \(W\).\footnote{\weightedpin
{grh-compact/src/New/Filling/TopFloor.lean}{453-L508}
{Local constancy and the moving labelled cosets},
\weightedpin{grh-compact/src/New/Filling/TopFloor.lean}
{558-L592}{The common-entry Gram and the aperture}, and
\weightedpin{grh-compact/src/New/Filling/FiniteSystem.lean}
{185-L223}{Application on the rational labels}.}

\paragraph{The Fourier multiplier of the finite terminal aperture.}
\leavevmode\teamanchor{proof.weighted-finite-terminal-aperture}%
Set \(\mu_0=\mu(K_0)\), \(v_U=\operatorname{vol}_U(U\cap K_0)\), and
\(m_Z=\mu(Z)\mu_0/v_U^2\).
With the filling modulus,
\(v_U=\prod_{i<\varrho}\operatorname{md}_T(g)^{d_i+1/2}\) and
\(\operatorname{md}_T(g)=g^{-2f}\). Thus
\(v_U^{-1}=g^{f\varrho(2k+2\varrho+1)}=U_0\).
Write \(\mathcal I_i(p,q;\omega)\) for the finite floor form,
including its head density \(\kappa_i(p,q)=
\prod_{a<i}\operatorname{md}_T(p_aq_a)^{d_a/2}\).
Write \(Z_i(x)=h(f_i,x f_i)\), with \(f_i=e_i\) for \(i<\varrho\)
and \(f_\varrho=z\). We claim
\begin{equation}\label{eq:weighted-finite-terminal-form}
 \int_{G_\varrho}\psi_T(\xi t(l))\,
          \langle C_p(l),C_q(l)\rangle\,d\nu_T(l)
       =\mathcal I_\varrho(p,q;\psi_T(\xi Z_\varrho)).
\end{equation}

To prove the identity expand the two tensors. With
\(w_p=w_{p,\varrho}\), \(w_q=w_{q,\varrho}\), the four factors give
\[
 \kappa_p\kappa_q\,
 \langle W(y_1),W(y_2)\rangle
 \overline{\langle W(lw_p^{-1}y_1), W(lw_q^{-1}y_2)\rangle}.
\]
The integrations in \(y_1,y_2\) use \(d\mu\,d\mu\).
The compact support proved above allows all three integrations
in \(l,y_1,y_2\) in either order, even with absolute values.
Put \(y_2=y_1x\), using left Haar in \(x\). In the \(y_1\)-integral
put \(y_1=h_1\sigma y\) by
\(d\mu=dh_1*d\theta_\varrho*dM_\varrho\). For fixed \(h_1,\sigma\)
replace \(l\) by \(h_2\tau(h_1\sigma)^{-1}\), using right
invariance of \(\nu_T\) and then \(d\nu_T=dh_2*d\theta_\varrho\).
Both \(h_1,h_2\) fix \(z\), giving
\[
 t(l)=h(z,\sigma\tau^{-1}z),\qquad
 l w_p^{-1}y_1=h_2\tau w_p^{-1}y,\qquad
 l w_q^{-1}y_2=h_2\tau w_q^{-1}yx .
\]
With \(D=w_q^{-1}w_p\), the last expression is also
\(D h_2\tau w_p^{-1}yx\). Integrating the first raw pairing
over \(h_1\) gives the upper pairing of
\((u,v)\mapsto\langle W(u),W(v)\rangle\) at
\((\sigma y,\sigma yx)\). Integrating the second over \(h_2\)
gives the conjugate of the upper pairing of
\((u,v)\mapsto\langle W(u),W(Dv)\rangle\) at
\((\tau w_p^{-1}y,\tau w_p^{-1}yx)\).
These upper integrals are over \(dh\) at this terminal stage.
The two remaining \(\theta_\varrho\)-integrals, with the multiplier
at \(\sigma\tau^{-1}\), are precisely the defining two translates
of \(\mathcal I_\varrho\). Its density here is
\(\kappa_\varrho(p,q)=\kappa_p\kappa_q\). This proves
\eqref{eq:weighted-finite-terminal-form} with the specified
conjugation and finite stabilizer measures.\footnote{\weightedpin{grh-compact/src/New/RealPlace/LocalTheorem.lean}
{30-L76}{The head density, the two kernels and the floor form},
\weightedpin{grh-compact/src/New/Filling/TopFloor.lean}
{510-L556}{Terminal unfolding for the unnormalized vector} and
\weightedpin{grh/src/New/Spanning/Frames/TopCorrelation.lean}
{596-L677}{The changes of variables with both core Haar measures}.}

Use \(p\oplus_i c\) for replacing the entry at \(i\) by \(c\).
In a unit head use \(1\) when \(c\) has a zero component. Such
\(c\)'s form a null set for \(dc\), being a finite union of
coordinate hyperplanes in \(A_T\). For the terminal form put
\(J_{\varrho+1}=\mathcal K_L/(|L|m_Z)\), where \(|L|=\operatorname{vol}(L)\).
Fubini is applicable when integrating it over the common last
entry \(c\). Indeed integration over \(c\) of the absolute
majorant \(1_L(t(l)-c)|\langle C_p(l),C_q(l)\rangle|\)
costs \(|L|\) times an integrable function of \(l\).
The substitution \(u=t(l)-c\) yields
\begin{align}
 &\int_{A_T}\psi_T(\xi c)\,
       J_{\varrho+1}(p\oplus_\varrho c,q\oplus_\varrho c)\,dc
       \notag\\[-2pt]
 &\hspace{15mm}
   =\frac{\mathcal I_\varrho(p,q;\psi_T(\xi Z_\varrho))}{m_Z}\,
          \frac{1}{|L|}\int_L\psi_T(-\xi u)\,du .
 \label{eq:weighted-finite-aperture}
\end{align}
The last factor is exactly \(1_{L^\perp}(\xi)\), where
\(L^\perp=\{\xi:\psi_T(\xi L)=1\}\). If the character is
nontrivial on \(L\), translation by a point where it differs
from one forces the integral to vanish. In particular the factor
equals one at zero frequency and has modulus at most one at
every frequency.\footnote{\weightedpin
{grh-compact/src/New/Filling/FiniteSystem.lean}{262-L280}
{The aperture and its signed character average}.}

\paragraph{Descent of the finite common entry.}
\leavevmode\teamanchor{proof.weighted-finite-common-entry}%
For \(i\le\varrho\) put \(J_i(p,q)=\mathcal I_i(p,q;1)/m_Z\).
For \(i<\varrho\), write \(P=w_{p,i}\), \(P'=w_{q,i}\) for the
truncated Weyl elements and \(D=(P')^{-1}P\). Each commutes with
\(G_i\). Denote the two upper pairings just used, now at floor
\(i+1\), by \(B^1,B^D\). The
\hyperlink{proof.weighted-shell-floor-mass}{quadric pushforward}
and the \hyperlink{proof.weighted-restricted-position-integral}
{stabilizer invariance of upper averages} apply to this finite
level. Both pairings are invariant under simultaneous left
translation by \(S_i=\operatorname{Stab}_{G_i}(f_i)\).
For the finite floor, \(\theta_i\) is the law of
\(w_i(c)^{-1}u\) for
\(\operatorname{md}_T(c)^{d_i}\,dc\,d\eta_i(u)\).
Here \(u\) fixes \(f_i\), and \(Z_i(w_i(c))=c\).

In \(\mathcal I_i(p,q;\psi_T(\xi Z_i))\) the upper pairings occur as
\[
  \psi_T(\xi Z_i(\sigma\tau^{-1}))\,
  B^1(\sigma y,\sigma yx)\,
  \overline{B^D(\tau P^{-1}y,\tau P^{-1}yx)} .
\]
Keep \(\sigma\) fixed. The expression in \(\tau\) descends to
\(\tau^{-1}f_i\), since \(B^D\) is invariant under simultaneous
left translation of both arguments by \(S_i\), and \(Z_i\) is
invariant under right multiplication by \(S_i\).
Thus \(\tau\) may be replaced in the
integral by \(\tau_0\sigma\) using invariance of the quadric
measure. Put \(\tau_0=w_i(c)^{-1}u\) and \(y'=u\sigma y\).
Then \(Z_i(\sigma\tau^{-1})=Z_i(u^{-1}w_i(c))=c\).
The first pairing becomes \(B^1(y',y'x)\) and the second uses
\((P w_i(c))^{-1}y',(P w_i(c))^{-1}y'x\).
Appending \(c\) to both heads leaves \(D\) unchanged:
\(w_i(c)\) commutes with the preceding blocks. The resulting
measure in \(y'\) is \(\eta_i*\theta_i*M_i=M_{i+1}\).
The new head density is
\(\kappa_{i+1}(p\oplus_i c,q\oplus_i c)
   =\kappa_i(p,q)\operatorname{md}_T(c)^{d_i}\).
With weight one, each upper pairing is the successive average of
the next upper pairing, by the ordered convolution of its floor
measures. Their product is therefore the two-translate integral
defining \(\mathcal I_{i+1}\) at these heads.

For the absolute values the initial integral is bounded by the
\hyperlink{proof.weighted-finite-denominator-shells}{finite
denominator-shell estimate} with gain \(1\) and exponent \(0\).
At fixed inputs each successive upper average converges absolutely
by the measured floor-mass bound on its support. The same change
of variables for the nonnegative integrands gives integrability
in \(c\), then Fubini applies to the complex integral. We have proved
\begin{equation}\label{eq:weighted-finite-common-entry}
  \int_{A_T}\psi_T(\xi c)\,J_{i+1}(p\oplus_i c,q\oplus_i c)\,dc
       =\frac{\mathcal I_i(p,q;\psi_T(\xi Z_i))}{m_Z}
                  \qquad (i<\varrho).
\end{equation}
At the bottom the Weyl elements are \(1\), the preceding stack
is \(\delta_1\), and the two upper pairings coincide.
The \hyperlink{proof.weighted-finite-bottom}{unfolded phase
pairing} gives
\[
 J_0=\frac{1}{v_U^2m_Z}
       \int_G\left|\int_Z\chi_g(yxy^{-1})\,d\mu(y)\right|^2d\mu(x)
     =\frac{\mathcal B(Z)}{\mu(Z)\mu_0}.
\]
Thus the bound from each prescribed standard or dual vector
stabilizer acts on this very bottom form. In
\eqref{eq:weighted-finite-common-entry} and
\eqref{eq:weighted-finite-aperture} the right side at \(\xi=0\)
is \(J_i\). At a unit frequency \(\xi\) the
\hyperlink{proof.weighted-finite-frequency}{inserted-character
bound} applies, with the two head weights at \(i\).
For \(i<\varrho\) the unequal-pivot bound applies directly to
\(J_{i+1}(p,q)\) when \(g^{-1}(p_i^{-1}-q_i^{-1})\in A_T^\times\),
with both head weights at \(i+1\).
Together with common-last-entry positivity and continuity, these
are all the finite inputs of the adelic sieve for this \(Z,L\).\footnote{\weightedpin{grh-compact/src/New/RealPlace/LocalTheorem.lean}
{258-L318}{Successive averaging and lower-floor descent},
\weightedpin{grh-compact/src/New/Filling/TopFloor.lean}
{336-L353}{The majorant and the two unit-frequency bounds},
\weightedpin{grh-compact/src/New/Filling/GlobalToLocal.lean}
{245-L253}{The unipotent principal volume} and
\weightedpin{grh-compact/src/New/Filling/EichlerFlag.lean}
{78-L91}{The finite hypotheses used in the sieve}.}

\paragraph{The residue conditions and their finite levels.}
\leavevmode\teamanchor{proof.weighted-finite-residue-stabilizers}%
Use \(n=\varrho+1\) and \(r=k+2\varrho+1\) on the given filling
datum. Let \(h\) be the number of column classes and \(d_B\) the
number of product Bessel coordinates, so \(C=h d_B\). For proper \(A\), a residue in \(\Omega_A\) fixes the prescribed
vectors and has all leading frame minors of orders \(1,\ldots,n\)
nonzero at one chosen place above \(g_0\). The class representatives
fix the frame, and taking a minor of the integral term position
\(h_c^{-1}\gamma_f h_{c'}\) commutes with reduction in the filling
lattice basis. Hence every term over this residue has all these minors
nonzero over \(E\). The last minor includes the marked anisotropic
column. The full fixing set and the two complementary residue sets
use the empty choice of retained links in the crude comparison.

For a finite calculation take a conjugation-stable set \(T\) containing
the places of \(g\) and \(N\) and the supports of the denominators of
the filling basis, the finitely many term positions, and their inverses.
Enlarge \(T\) so that the selected pivot values, terminal entries
and their inverses are integral off \(T\).
Write \(K\) for the integral group at \(T\) and \(K_0=K(g)_T\) for its
principal level. Membership in the global principal level means
membership in \(K_0\) at \(T\) and integrality with integral inverse
off \(T\). At a split pair \(\{w,\bar w\}\) the unitary isometry is
determined by its first matrix. In dual coordinates its other
component is the inverse adjoint. The reduction of the first matrix
is onto \(\mathrm{GL}_r(\mathbf F_q)\), since entrywise lifts of an
invertible residue matrix have unit determinant. At a prescribed
pair \(q=g_0\), let \(Z_w\) be the inverse image of the specified
vector stabilizer. The frame is injective on residue vectors and
both actions are transitive on nonzero vectors. Thus
\([K_w:Z_w]=q^r-1\). Take \(Z=\prod_{w\in A} Z_w\times K_{\rm rem}\)
over the distinct pairs and the remaining conjugation-stable
factor. We have \(K_0\subseteq Z\subseteq K\), and the components
at \(T\) of all lifts of residues in \(\Omega_A\) belong to \(Z\).

Write \(H_T\) for the terminal stabilizer \(H\) in the local
packet measure. Use the
\hyperlink{proof.weighted-finite-common-entry}{normalized
finite bottom}
\[
 J_0^T=\frac{\mathcal B(Z)}{\mu(Z)\mu_0},\qquad \mu_0=\mu(K_0),
 \qquad
 \mathcal B(Z)=\int\left|\int_Z\chi_g(yxy^{-1})\,d\mu(y)\right|^2d\mu(x).
\]
For \(x\in K_0\) the phase in the inner integral is one by
normality. Its contribution to \(\mathcal B(Z)\) is
\(\mu(Z)^2\mu_0\). On a prescribed split pair, the unit-pivot charts
have integral horizontal and skew coordinates of volume one.
At the last floor use the marked pair above and the unit
coordinate \(a=h(u,e)\) on the norm-\(-N\) fibre. Scale \(v\) by
\(a\) and \(u\) by \(\bar a^{-1}\), then apply the Heisenberg
translation fixing \(u\). Division by the unit \(a\) keeps the
horizontal coordinates integral, and the norm equation leaves
an integral skew coordinate. The same integral chart therefore
has mass at least \((1-q^{-1})^2\) on each of the \(n\) floors.
The integral terminal stabilizer has mass one. It follows that
\(\mu(K_w)\ge(1-q^{-1})^{2n}\ge2^{-n}\) for the chosen \(q\).
The nonunits of the split ring
have additive measure \(\delta_w\le 2/q\le b^{-9}\).
Apply the \hyperlink{proof.weighted-finite-bottom}{finite bottom
estimate} with \(t_0=k-10^{10}n^4\) to \(Z_w\) and the principal
lower bound to \(K_w\). This gives
\[
 \frac{\mathcal B(Z_w)/\mu(Z_w)}{\mathcal B(K_w)/\mu(K_w)}
 \le\frac1{q^r-1}+\frac{2^{-n}}{\mu(K_w)}\delta_w^{\,t_0+1}
 \le 2b^{-10r}+b^{-8k}\le b^{-6k}.
\]
Here \(t_0+1\ge8k/9\) on the sufficiently large selected squares.
The split pairs and complementary factor give a direct product
of groups, and the phase is a product. Rescaling the product
Haar by a positive scalar rescales \(\mathcal B\) by its cube
and each subgroup mass by that scalar. Thus the quotient of
the two ratios is unchanged, and Fubini multiplies the individual
bounds. On the whole semilocal group the finite bottom estimate
at exterior exponent zero uses \(\delta\le1\) and
\(\mu(K)\le1+2^{-n}\). For this upper bound unload the preceding
integral unipotents, of volume one. At each orbit choose box
representatives at that floor and the next. The indicated
column of their quotient modulo the earlier frames is integral,
so the translated orbit vector remains in the unit box. The
finite cone calculation with \(e\) split planes bounds this
fibre by \(\zeta_R(e-1)\) uniformly in its norm level. Indeed
the content divisor bound uses trace one in the skew coordinate
after the orthonormal coordinates are fixed. There are at least
\(m_0+1\) planes at each floor, including the auxiliary terminal
planes, and the orthonormal boxes have mass one. The flag
disintegration therefore gives \(\mu(K)\le\zeta_R(m_0)^n\)
with \(\mu_H(H\cap K)=1\). This is at most \(1+2^{-n}\).
Since \(2^{1-n}\le b^{-6}\), this yields
\[
 0\le J_0^T\le e_T=(1+b^{-6})b^{-6k|A|}.
\]
The estimate holds throughout the finite exhaustion.\footnote{\weightedpin
{grh-compact/src/New/Filling/CrudeBoundDefs.lean}{27-L40}
{Residue action and test}, \weightedpin
{grh-compact/src/New/Filling/TopFloor.lean}{133-L149}
{Principal term and unnormalized vector}, \weightedpin
{grh-compact/src/New/Filling/TopFloor.lean}{199-L215}
{Floor-zero pairing and box lower bound}, \weightedpin
{grh-compact/src/New/Filling/StandardLevel.lean}{224-L237}
{The measure of the integral box}, and \weightedpin
{grh-compact/src/New/Filling/FixedLines.lean}{19-L143}
{Fixed vectors and factorization of Haar}.}

\paragraph{Unfolding on the labelled cosets.}
\leavevmode\teamanchor{proof.weighted-finite-labelled-packets}%
For the rational open-cell parametrization write \(F=E^+\),
\(\gamma_p=w_p\sigma\), \(t=h(z,\sigma z)\) and
\(y=(\sigma z)_{X_0}\). The nonzero leading minors provide the
successive pivots. At each floor the middle and skew coordinates
of the image of the pivot vector, divided by its nonzero opposite
coefficient, specify the first Eichler translation. The opposite
column and the unit pairing specify the second. The link in the Weyl
position is the graph of \(\sigma^{-1}\) on \(X^\perp/X\).
For a fixed rational open term and its terminal representative
\(\sigma\), this elimination determines the two unipotents and
all the pivots uniquely. The block fixes \(\sigma\) and can have
many heads. The block labels \((c,s,c',\Lambda)\) inject into
\((c,c',s,\sigma)\). The identities
\(\sigma z=y-(t/N)z\), \(Ny^*y=t\bar t-N^2\) and
\(\operatorname{disc}(\gamma_p)=|\prod_{a<\varrho}p_a|^2
(t\bar t-N^2)\) give \(|t_v|>N\) at every real place on the
open links. We keep the sum over heads inside each block and group
the block labels by \(t\).

For a block choose a lift \(\tilde s\in K_r\), put
\(D_\gamma=h_c^{-1}\gamma_f h_{c'}\), and let
\(l=(D_\sigma|_T)^{-1}\). The class representatives commute
with the Weyl operators, giving \(D_{\gamma_p}|_T=w_pl^{-1}\).
In the filling basis \(l^{-1}z\) is the lattice-coordinate vector
of \(h_c^{-1}y-(t/N)z\). Use the finite correlation \(C_p(l)\) of
the \hyperlink{proof.weighted-finite-packet-fibre}{actual vector
\(W_Z\)}, with factor
\(\varkappa(p)=\prod_{a<\varrho}\operatorname{md}_T(p_a)^{e_a/2}\),
\(e_a=k-1+2(\varrho-a)\). For each \(b\in Q=K_0\backslash Z\)
choose \(a=K_0\tilde s_T b\), the unique coset with
\(\tilde s_T^{-1}ab^{-1}\in K_0\). Right multiplication of the
integration variable by \(a\) in the \((a,b)\) entry gives
\(\chi_g(x_T)\overline{\chi_g((D_{\gamma_p}^{-1}x\tilde s)_T)}\).
Here \(x\) can be taken to equal one off \(T\). Write
\(\zeta=D_{\gamma_p}^{-1}x\tilde s\), which is integral off \(T\)
with integral inverse.

For the two nonzero phases the rational unipotents solving
\[
 (h_c^{-1}u_fh_c)x\in K(g),\qquad
 (h_{c'}^{-1}u'_fh_{c'})\zeta\in K(g)
\]
form one left coset of \(U_c\) and \(U_{c'}\), respectively.
Indeed the local support gives finite adelic unipotent solutions,
taking the identity off \(T\). The rational points of \(U\) are
dense in its finite adeles by the affine unipotent coordinates
and approximation with denominators supported at finitely many
places. The conditions are open and the ratio of two solutions
is in the respective class unipotent. A vanishing phase has no
solutions. On the two nonempty cosets the unbarred characters
in the correlation are \(\chi_f^+(u)^{-1}\) and
\(\chi_f^+(u')^{-1}\). For example \(x\) is the inverse of
\(h_c^{-1}u_fh_c\) times a principal-level element.
The simple-root sum is invariant under the middle stabilizer
and the integral factors off \(T\) have trivial phase.
The product formula therefore gives the two unbarred values
\(\chi_\infty(u)\) and \(\chi_\infty(u')\).

In the cut real correlation \(B_\infty(\gamma_p)\) insert the
partitions \(\sum_u\beta_c(uY)=1\) and
\(\sum_{u'}\beta_{c'}(u'\gamma_p^{-1}Y)=1\) on these two cosets.
The real phase law converts the integrand to the amplitude of
\(\delta=u\gamma_pu'^{-1}\) after \(Y\mapsto uY\). For fixed \(u,u'\)
the two finite conditions say that \(D_\delta\in K_r\) has
residue \(s\), \(h_c^{-1}u_fh_c\) is integral off \(T\), and
\(x_T\in(h_c^{-1}u_fh_c)_T^{-1}K_0\). In fact if \(k_1,k_2\)
are the products in the two conditions then
\(D_\delta=k_1\tilde s k_2^{-1}\). Conversely the condition on
\(x_T\) defines \(k_1\in K(g)\), and
\(k_2=D_\delta^{-1}k_1\tilde s\in K(g)\) by its reduction.
The finite coset contributes \(\mu_0\). Open-cell uniqueness and
the \hyperlink{proof.weighted-poincare-realization}{absolute
integral majorant for amplitudes} justify the interchange of the
sum and the two integrals, since the absolute integrals sum to
at most \(\mu_0\) times that majorant. Thus
\[
 C_p(l)_{ab}B_\infty(\gamma_p)
       =\varkappa(p)\mu_0\mathcal A_{p,T}(c,s,c',\sigma),
\]
where \(\mathcal A_{p,T}\) is the indicated restricted sum of
rational cut amplitudes. This holds for every \(b\in Q\) with
its specified \(a\).\footnote{\weightedpin
{grh-compact/src/Lib/DiscreteSeries/Bruhat/BigCell.lean}{510-L554}
{Rational cell}, \weightedpin
{grh-compact/src/New/Filling/GlobalToLocal.lean}{159-L213}
{Terminal coordinates and absolute sums}, and \weightedpin
{grh-compact/src/New/Filling/GlobalToLocal.lean}{298-L376}
{Unfolding with both class cutoffs}.}

\paragraph{The labelled multiplicity at the terminal quadric.}
\leavevmode\teamanchor{proof.weighted-finite-labelled-mass}%
For fixed \((c,y,t)\) the terminal vector \(e\) is common to the
corresponding labels. The pairs \((K_0l,(a,b))\) are distinct
as the label and \(b\in Q\) vary. Indeed equality of the cosets for
labels \(i,j\) gives \(h_{c'_i}^{-1}\sigma_i^{-1}\sigma_j h_{c'_j}\)
in \(K_0\) at \(T\). The two term positions are integral off \(T\),
so this element is in the global principal level. Equality of the
terminal vector and of the fixed frames places \(\sigma_i^{-1}\sigma_j\)
in \(H_0(F)\). The \hyperlink{proof.weighted-bessel-unfolding}{column
level} is the pullback of the principal row level to \(H_0\).
The class double cosets agree, and freeness at that level gives
\(\sigma_i=\sigma_j\). Finally equality of
\((a,b)\) forces equality of the two residues, both reductions of
\(ab^{-1}\). This gives the injection into the labelled cosets with
all \(|Q|\) choices per block.

Use \(\ell_{va}=|\kappa_v|^{2(\varrho-a)}/N\), \(d_B\) the number
of product Bessel coordinates, and \(G_t=\prod_v(|t_v|^2-N^2)^{k-1}\).
The \hyperlink{proof.weighted-real-correlation-forms}{real terminal
Gram formula} on the scaled head and entry \(t_v/N\)
gives the same Gram matrix \(R\) for all labels with that \(t\).
Its vectors are \(\eta_p B_\infty(\gamma_p)\), where
\[
 \lambda_0=\frac{c_G^{f/2}}{d_B}\prod_v\prod_{a=0}^{\varrho}\ell_{va}^{e_a},
 \qquad
 \eta_p=\lambda_0\sqrt{G_t}\prod_v\prod_{a<\varrho}|\Phi_v(p_a)|^{e_a}.
\]
The terminal scale is \(\ell_{v\varrho}=1/N\) and contributes
\(N^{-(k-1)}\) at each \(v\). The product formula for each
\(T\)-unit gives \(\varkappa(p)\eta_p=\lambda_0\sqrt{G_t}\).
Thus at the labelled coordinate \(v_p=C_p(l)_{ab}\) the form
\(v^*Rv\) equals \(\mu_0^2G_t\lambda_0^2\|\sum_p\mathcal A_{p,T}\|^2\).

\paragraph{The one-vector count on the marked terminal quadric.}
\teamanchor{proof.weighted-open-quadric-count}%
Write \(V=\operatorname{vol}(\widehat{\mathcal O}_E)\) for the
global finite additive measure and
\[
 c_* = |D_E|^{1/4} V^k
            \left({\pi^k\over (k-1)!\,N^k}\right)^f .
\]
For a finite list of distinct \((c,y)\) with
\(Ny^*y=t\bar t-N^2\) and \(h_c^{-1}y\) integral off \(T\), let
\(a_{c,y}\ge0\) be lower bounds for a nonnegative function \(F_0\)
on \(H_T\cap K_0\)-orbits of a neighbourhood of the corresponding
terminal points. For all sufficiently small \(L\),
\[
  |L|\sum_{(c,y)}a_{c,y}
   \ \le\ (1+b^{-6})h c_*G_t
     \int_{h(z,e)-t_T\in L} F_0(e)\,d\delta_{-N}(e).
 \tag{*}
\]
To see the constant, choose a compact locally constant test
attaining the specified lower bounds and minorizing \(F_0\)
on the norm quadric. Take the maximum of the finitely many
\(a_{c,y}\) on sufficiently small compact open neighbourhoods
saturated under \(H_T\cap K_0\). Such saturations are again compact
and open. Apply Weil's one-vector formula to \(q_X(y)=y^*y\) on
the definite \(E\)-space \(X_0\), of rank \(k>2\). This is case
(I\({}_2\)), \(\epsilon=1/2\), with one input vector. At a split
place over \(g_0\), the two halves are a vector and a covector.
Nonzero vectors and covectors of fixed contraction form one
orbit, including contraction zero for \(k\ge2\). This gives the
full-rank null orbit and the transitivity in Weil's Theorem~5.
The quotient is compact, the Schwartz theta sum bounded there
by the finite class decomposition and the lattice theta sums.
The compact-level test is the saturated finite test at \(T\)
times the integral boxes outside it. At infinity take a Schwartz
test equal to one on the spheres of norm
\(\alpha_v=(|t_v|^2-N^2)/N>0\).
With probability Haar on the quotient each of the \(h\) classes
has mass \(1/h\) because the rational column stabilizers are trivial.

In Weil's (41), use the fibre gauge given by disintegration
of the ambient additive gauge against the fixed-field additive
gauge. With real Lebesgue in the chosen coordinates, the finite
integer boxes have covolume-one additive measures
\(V=2^f|D_E|^{-1/2}\) over \(E\) and \(|D_F|^{-1/2}\) over \(F\).
These give the normalized rational gauges, since multiplying a
gauge by a rational scalar has adelic product one. The quotient
of the normalization factors is
\(V^k\sqrt{|D_F|}=V^k |D_E|^{1/4}\), using finite unramifiedness
of \(E/F\). At each real place polar integration against
\(d\alpha\) gives \(\pi^k\alpha_v^{k-1}/(k-1)!\). The class sum
thus has scalar \(hV^k|D_E|^{1/4}\) times the real fibre masses
and the finite fibre integral whose integer boxes have mass one.
Off \(T\), Fourier inversion can first be performed on a congruence
slab around \(\alpha\). At a place of \(F\) of residue size \(Q\),
square the absolute integral for a character of conductor \(Q^a\)
and subtract the two variables. The perfect unramified or split
trace pairing, including above \(2\), forces the difference into
the \(a\)-th prime-power lattice. The squared integral is at most
\(Q^{-2ak}\) with the integer box of mass one. Since there are at
most \(Q^a\) frequency classes of this level, absolute summation
bounds the density by \((1-Q^{1-k})^{-1}\), uniformly in \(\alpha\).
Their product is at most \(\zeta_F(k-1)\le\zeta_E(k/2-1)\), and
\(\zeta_E(s)\le\zeta(s)^{2f}\) makes its relative error
\(o(b^{-6})\) on these data. Let \(B_{\rm lat}\) have the class
lattice columns in the rational middle space. The chart
\((y,\tau)\mapsto B_{\rm lat}^{-1}(y-\tau z/N)\) has determinant
\(\bar\pi/N=1/\pi\) up to sign, of absolute modulus \(N^f\).
In volume-one finite coordinates the ambient chart measure is
therefore \(N^f\,dy\,d\tau\). Dividing by the unchanged fixed-field
norm measure for \(q_X(y)-\tau\bar\tau/N\) gives
\(d\tau\,dy|_{q_X=\tau\bar\tau/N-N}=N^{-f}\,d\delta_{-N}\).
The skew coordinate is changed by an integral skew unit and keeps
the integer-lattice measure. Together with
the \(N^{-(k-1)f}\) from \(\alpha\) this gives precisely \(c_*\).
For a compact locally constant minorant the slab integrals at
the regular norm fibre tend to its quadric integral. Indeed a
finite partition of the compact support into local norm-coordinate
charts shows that the minorant's fibre integral is locally constant
as \(\tau\) varies near \(t_T\). The outside-density bound is
uniform in this perturbation. Its slack within \(1+b^{-6}\)
gives \((*)\) on all sufficiently small slabs for the finite list.
\footnote{\href{https://archive.ymsc.tsinghua.edu.cn/pacm_download/117/5988-11511_2006_Article_BF02391774.pdf\#page=77}
{Weil, Theorem~5 and (41), pp.~76--77}, with case (I\({}_2\)) in
\S27 and the measures in \S\S44, 53, and
\weightedpin{grh-compact/src/Lib/SiegelWeil/OneVector/OneVector.lean}
{38-L48}{One-vector mass formula}, \weightedpin
{grh-compact/src/Lib/SiegelWeil/Mass/TorusMajorant.lean}{68-L112}
{Quadratic character and denominator ideals}, \weightedpin
{grh-compact/src/Lib/Adelic/Adeles/FinitePlaceDensity.lean}{325-L364}
{Density off the selected places}, \weightedpin
{grh-compact/src/New/Filling/StandardLevel.lean}{38-L58}
{Covolume in the terminal chart}, and \weightedpin
{grh-compact/src/New/Filling/FiniteSystem.lean}{43-L147}
{Comparison on a neighbourhood of the terminal points}.}

\paragraph{The normalized packet comparison on the common fibre.}
The \hyperlink{proof.weighted-finite-packet-fibre}{neighbouring-fibre
lower bound} for these pairs supplies the factor \(\mu_H(H_T\cap K_0)\)
on each occupied coset. Apply the
\hyperlink{proof.weighted-open-quadric-count}{one-vector quadric count}
with these lower bounds, grouping by distinct \((c,y)\) for fixed \(t\).
In rational coordinates \(D_\sigma z=h_c^{-1}y-(t/N)z\).
Since \(D_\sigma\), the basis change and its inverse, and \(t/N\)
are integral off \(T\), the required \(h_c^{-1}y\) is integral
there. The norm and strict positivity conditions were verified in
the rational open cell.
For the test \(F_0\) in that count use the nonnegative fibre
form \(E_R\) from the local argument. Its integral over the
aperture is the integral in the terminal comparison.
For sufficiently small compact open additive \(L\), the integral of
the positive fibre form over the aperture is at most
\(\Re\sum_{p,q}R_{pq}\mathcal K_L(p,q)\) by the terminal
disintegration. A compact open additive subgroup inside the
intersection of the finitely many neighbourhoods works for all
selected \(t\).

Normalize \(\mu_H(H_T\cap K)=1\) as in the local terminal measure.
Identity extension outside \(T\) identifies \(\bar H\) with
\((H_T\cap K)/(H_T\cap K_0)\). Thus \(\mu_H(H_T\cap K_0)=|\bar H|^{-1}\).
At the level primes the integral Bessel subgroup has integral
unipotent and middle factors in the flag basis. Their reductions
intersect trivially by their action on the diagonal frame blocks.
The local factors extend integrally by the identity off those
primes, so \(|J|=|\bar H|U_0\), with the unipotent index
\(U_0=v_U^{-1}=g^{f\varrho(2k+2\varrho+1)}\) in the local
normalization. Also \(\mu(Z)=\mu_0|Q|\). These identities give
\[
 m_Z=\mu(Z)\mu_0 U_0^2
     =|J|U_0\,\mu_H(H_T\cap K_0)\mu_0^2|Q|.
\]
In the one-vector inequality divide by \(|L|\), cancel \(G_t>0\)
and apply this identity. The result for the chosen labels is
\[
 \lambda_0^2\sum_{\rm labels}\Big\|\sum_p\mathcal A_{p,T}\Big\|^2
 \le (1+b^{-6})\Pi\,
      \Re\sum_{p,q}R_{pq}\frac{\mathcal K_L(p,q)}{|L|m_Z},
 \qquad \Pi=|J|U_0h c_*.
\]
The terminal kernel here uses the same unnormalized coset vector
and Haar as the finite floor identities.\footnote{\weightedpin
{grh-compact/src/New/Filling/StandardLevel.lean}{70-L83}
{Haar of the integral stabilizer}, \weightedpin
{grh-compact/src/New/Filling/GlobalToLocal.lean}{245-L297}
{Unipotent order and reduction indices}, and \weightedpin
{grh-compact/src/New/Filling/FiniteSystem.lean}{150-L223}
{Comparison on the common terminal fibre}.}

\paragraph{The rational heads and the frequency budget.}
\teamanchor{proof.weighted-open-sieve}%
\teamanchor{proof.weighted-open-poisson-budget}%
We give the summation over rational pivots, since their weights
vary with the actual scale and with the places of \(T\).
Put \(c_i=V/\prod_v\ell_{vi}^2\). Let \(O_T=\prod_{w\in T}\mathcal O_w\),
\([y]_T=\prod_{w\in T}|y|_w^{-1}\), and
\(D_T(y)=[O_T:O_T\cap y^{-1}O_T]\). Put
\(L_h(s)=s^h\min(1,s)^6\).
For frequencies \(y\in\mathcal O_{E,T}\), integral away from \(T\),
we shall use
\[
 \sum_{0\ne y\in\mathcal O_{E,T}}
  {[y]_T^{h/2}\over D_T(y)^6}
     \prod_v L_h(|c_vy_v|^{-1})
       \ \le\ C_1^f B^{6f}\prod_v X_v^h
       \qquad(|c_v|^{-1}\le X_v\le B,\ B\ge1).
\]
Here \(C_1\) is absolute. The product formula cancels the first
\(h\) powers, since \([y]_T^{1/2}\le\prod_v|y_v|\).
The remaining real factors are at most
\(B^{6f}\prod_v\max(1,|y_v|)^{-6}\).
Write \(\mathfrak d(y)\) for the denominator ideal. Its norm is
at most \(D_T(y)\). For a fixed denominator \(\mathfrak d\),
its elements lie in at most \(N\mathfrak d\) cosets of
\(\mathcal O_E\). In each coset a product of complex squares of
side \(1/2\) contains at most one element, since every conjugate
of a nonzero integral difference would otherwise have modulus less
than one. Summing the real weights over the grid costs at most an
absolute constant per place. Summing the remaining denominator
weights costs at most \(\zeta_E(2)\le\zeta(2)^{2f}\). This also
proves the bound for all finite subsums uniformly in \(T\).
\footnote{\weightedpin{grh/src/New/Spanning/MainTerm/Frequencies.lean}
{194-L270}{Packing on individual denominator cosets}, and \weightedpin
{grh-compact/src/New/Filling/GlobalSieve.lean}{43-L80}
{The compact frequency weight}.}

Remove the two head weights from the real frequency bound. At a
nonzero frequency the remaining factor is
\[
      \Theta_\infty b^{2000n^9} L_{k-n^7}(|\xi_v|^{-1}),
        \qquad
      \Theta_\infty=\pi^{(2\varrho+1)k}
          {\operatorname{vol}U(k)\over\operatorname{vol}U(r)} .
\]
At the finite places the
\hyperlink{proof.weighted-finite-frequency}{two local Fourier bounds}
use the next diagonal Gram entry and the difference of two last
pivots, respectively. With their
\hyperlink{proof.weighted-finite-denominator-shells}{denominator summation},
the finite frequency factor is bounded by
\[
   b^{n^9 f}[y]_T^{(k-n^7)/2}/D_T(y)^6.
\]
Here the frequency exponent is \(k-n^7\), the head exponent is \(n^6\),
and \((2/\operatorname{md}_T(g))^{n^6}=(2g^{2f})^{n^6}
\le b^{n^9 f}\) since \(g\le b^{23}\).
It gives the same bound in the off-diagonal step at
\(y=(p^{-1}-q^{-1})/g\) for distinct last rational pivots \(p,q\).
In this case the real frequency is \(g y_v/\ell_{vi}\).
In both cases the inverse real frequency scale is at most
\[
       X_v={|\kappa_v|^{2\varrho+1}\over gN},
       \qquad \max(1,X_v)\le b^{O(n)}.
\]
Indeed \(|\kappa_v|=g|\Phi_v(\varpi)|/(2\pi)\ge1\).
The packing bound therefore bounds the sum of the frequency factors,
or one such off-diagonal factor, by
\(\epsilon_*=\Theta_\infty^f\prod_v X_v^{\,k-n^7}b^{O(fn^9)}\).
All implied constants here are absolute.

For a head of length \(i\), its successive cumulative products
\(d_j=p_0\cdots p_j\), \(j<i\), determine the head. The finite head
weight is zero unless every \(d_j\) is integral at \(T\), and equals
\(\prod_j[d_j]_T^{n^6/2}\) otherwise. Away from \(T\) these \(d_j\)
are units. For each \(j\) apply the packing bound with \(h=n^6\),
\(c_v=\prod_{a\le j}\ell_{va}\) and \(X_v=B=N^i\). These choices
are valid since \(\ell_{va}\ge N^{-1}\). The real head weights are
exactly the corresponding products of \(L_{n^6}\). As \(N\le b^{60n}\),
summing the heads independently via the \(d_j\)'s costs
\(b^{O(fn^9)}\), also after squaring. Moreover
\(\ell_{vi}\le b^{46n}\) and \(V=(2/\operatorname{rd}_E)^f\).
Consequently the squared sum of head weights divided by
\(\prod_{a<i}c_a\) is bounded by a common
\(\Gamma=b^{O(fn^9)}\).

Let \(\bar u\) be the mean of the logarithms of the nominal weights
and put \(G_*=\log(g\operatorname{rd}_E/(2\pi))\).
The different generator has \(\prod_v|\Phi_v(\varpi)|
=\operatorname{rd}_E^f\). Since \(m=2\varrho+2\), the specified
window for \(N\) yields
\[
 A_*={1\over f}\sum_v\log X_v
    =(m-1)G_*-\log(gN)
    \le -m\bar u-\tfrac{25}{2}\log b,\qquad
 |A_*|=O(m\log b).
\]
Here \(g\ge b^{20}\). In Weyl's dimension product the factors for
the unchanged initial \(k-1\) parameters cancel. For each last
index \(l=k-1,\ldots,r-1\) there are \(l\) remaining differences,
each at most \(w(v)e^{1/x}\), with denominator \(l!\).
Thus
\[
   f^{-1}\log\deg_{\rm Bes}
       \le mr(\bar u+1/x)-\sum_{j=k-1}^{r-1}\log(j!).
\]
The sphere product for \(\operatorname{vol}U(r)/\operatorname{vol}U(k)\)
gives \(\log\Theta_\infty\le\sum_{j=k}^{r-1}\log(j!)\). It follows that
\[
 {1\over f}\log\!\left(\deg_{\rm Bes}\Theta_\infty^f
                         \prod_v X_v^{k-n^7}\right)
 \le -\tfrac{25}{2}k\log b
         +O\!\left(m^2\log b+mr/x+n^7m\log b\right)
 \le -10 k\log b
\]
for all sufficiently large selected squares.
Indeed \(m\asymp\sqrt{x}\), \(k\asymp x^5\) and
\(0\le\bar u\le2\log b\). Also \(n^9=o(k)\). Hence
\((2\varrho+1)\Gamma\epsilon_*\le b^{-6kf-10}/\deg_{\rm Bes}\).\footnote{\weightedpin
{grh-compact/src/New/Filling/DegreeBounds.lean}{52-L78}
{Weyl differences on the prescribed real type},
\weightedpin{grh-compact/src/New/Data/Different.lean}{354-L363}
{Product of the different generator}, and \weightedpin
{grh-compact/src/New/Seesaw/IndexTypes.lean}{130-L160}
{The selected ramification window}.}

\paragraph{Ordered Poisson summation on the actual heads.}
\leavevmode\teamanchor{proof.weighted-finite-ordered-poisson}%
Multiply the normalized finite forms by the real forms on the scaled
heads. Denote the product forms by \(J_i\), \(0\le i\le n\).
At the top with common terminal value the finite form is
\(\mathcal K_L/(|L|m_Z)\) and the real form is \(R\). Their entrywise
product is positive semidefinite. Descent integrates one common
coordinate at frequency zero by
\eqref{eq:weighted-finite-aperture} and
\eqref{eq:weighted-finite-common-entry} and the real descent.
The multiplier \(1_{L^\perp}\) is one at zero and its modulus
is at most one at all frequencies. The product system thus
uses the same two local frequency bounds.

Let \(\mathcal P\subset E^\times\) be a finite set of \(T\)-units
including the terminal values. Use the factors \(c_i\)
defined in the frequency budget. Let \(\Lambda_i(p)\) be the product of the finite and real
head weights for \(p\in\mathcal P^i\) with the real scales applied.
Write \(H_i=\sum_{p\in\mathcal P^i}\Lambda_i(p)\). Use the bounds
from the \hyperlink{proof.weighted-open-poisson-budget}{denominator
packing and frequency budget},
\(H_i^2\le\Gamma\prod_{a<i}c_a\) and a common \(\epsilon_*\) for
the sum of nonzero frequency gains and for each unequal-last-entry
gain. These are uniform in \(T\) and \(L\).

For each \(i\) let \(F_i(z)\) be the sum of
\(J_{i+1}(p\oplus z,q\oplus z)\) over the two length-\(i\) heads,
with the real scalings on those heads. This function is real
nonnegative by the top Gram and successive zero-frequency
integrations. It is integrable by the common-entry identities.
At a rational sample the heads and varied entry are units in the
product of fields. Continuity at the top and Fatou's lemma at each
descent give lower semicontinuity at each such sample. The
zero-coordinate sets in the integrating variables have measure zero.

Pull \(F_i\) back to the adeles using the real scalings \(\ell_{vi}\)
and projection at \(T\), extending by zero outside those adeles
integral off \(T\). The complex scaling has real determinant
\(\prod_v\ell_{vi}^2\). The finite box off \(T\) in the product
adelic measure has volume \(V\), since the local measures used at
\(T\) give the local integers mass one. The pullback has integral
\(c_i\int F_i\). The adelic quotient by \(E\) has volume one:
the real covolume of \(\mathcal O_E\) is \(2^{-f}\sqrt{|D_E|}\),
and the finite integral box has measure \(V=2^f/\sqrt{|D_E|}\).
Index the Fourier coefficients by the inverse characters
of \(z\mapsto\psi_{\mathbb A_E}(yz/\varpi)\), \(y\in E\).
The Fourier integrand thus uses \(\psi_{\mathbb A_E}(yz/\varpi)\).
The surviving frequencies have \(y\in\mathcal O_{E,T}\) and occur at
\[
 \xi_T=y_T,\qquad \xi_v=\frac{2\pi\mathrm i\,\Phi_v(y)}
                                  {\Phi_v(\varpi)\ell_{vi}} .
\]
The finite factor in the pairing is \(\psi_f(y\,\cdot\,/\varpi)\),
and the factor at \(v\) on the scaled argument \(z_v\) is
\(\operatorname{wave}(\xi_v z_v)\), where
\(\operatorname{wave}(z)=\exp(z-\bar z)\).
If \(y\) fails integrality off \(T\), an integral translation at
those places leaves the pullback invariant and acts by a nontrivial
character, so its coefficient vanishes. At nonzero surviving \(y\)
the product frequency is a unit in the product of local fields.
The absolute coefficients sum to at most \(c_i\epsilon_*H_i^2\)
by the product bounds. Negating the frequency index preserves the
finite integral annihilators and \(L^\perp\), the real determinants,
and all the absolute frequency estimates. The same self-dual Haar
and local masses therefore apply with this indexing.

Periodize the nonnegative pullback on \(\mathbb A_E/E\).
Tonelli's theorem gives an \(L^1\) function with these Fourier
coefficients. Their absolute convergence gives a continuous
function agreeing with the periodization almost everywhere by
uniqueness of Fourier coefficients. Near zero each finite sum over
the translates at the sampled rationals is bounded by the
periodization. Lower semicontinuity of this sum at zero and
density of the full-measure equality set show that its value at
zero is at most the continuous value at zero. For \(V_i\subseteq
\mathcal P\) this gives
\[
 \sum_{t\in V_i}F_i(t)
 \le c_i\left(\Re\sum_{p,q\in\mathcal P^i}J_i(p,q)+\epsilon_*H_i^2\right),
\]
where \(F_i(t)\) is evaluated at the scaled rational point.
At \(i=0\) the zeroth coefficient uses the bottom
\(E_0=e_T(1+b^{-10})^f/\deg_{\rm Bes}\). At \(i>0\) the summands
of \(J_i\) with equal last pivots give the preceding sum with
common entry set \(\mathcal P\). Unequal last pivots \(u,u'\)
use \(y=(u^{-1}-u'^{-1})/g\), nonzero and integral off \(T\).
The real frequencies are \(g\Phi_v(y)/\ell_{v,i-1}\) and the
finite frequency is \(y_T\). Their bound costs a further
\(\epsilon_*H_i^2\). Induction now gives
\[
 \sum_{t\in V_i}F_i(t)
 \le \left(\prod_{a=0}^{i}c_a\right)
                  \bigl(E_0+(2i+1)\Gamma\epsilon_*\bigr).
\]
The \hyperlink{proof.weighted-open-poisson-budget}{numerical budget}
on this datum gives
\((2\varrho+1)\Gamma\epsilon_*\le b^{-6kf-10}/\deg_{\rm Bes}\).
At \(i=\varrho\) the left side is
\(\Re\sum_{t,p,q}R_{pq}\mathcal K_L(p,q)/(|L|m_Z)\).\footnote{\weightedpin
{grh-compact/src/New/Filling/EichlerFlag.lean}{77-L148}
{Product forms and lower semicontinuity}, \weightedpin
{grh-compact/src/New/Filling/GlobalSieve.lean}{87-L213}
{Adele scaling and ordered sums}, and \weightedpin
{grh/src/New/Spanning/MainTerm/GlobalSieve.lean}{367-L446}
{Explicit Fourier transform and one-sided summation}.}

\paragraph{The packet constant and the row mass.}
\teamanchor{proof.weighted-open-mass}%
Set \(s_U=\varrho(k+\varrho)\) and
\(F_{\rm sc}=\prod_{v,a}\ell_{va}^{k+2(\varrho-a)}\).
For the actual column classes and the class-unipotent covolume,
\[
 h=|\bar H|(N^k-1)^f\operatorname{mass}_k,\qquad
 \nu F_{\rm sc}=U_0 N^{-kf} V^{-s_U}|D_E|^{\varrho/4}.
\]
The first identity follows by refining the integral genus of rank
\(k\). At the places of \(N\) the lattice basis prescribes a
nonzero dual vector in the residue space of dimension \(k\).
Each of the \(f\) split pairs has index \(N^k-1\).
At \(g\) the principal refinement has index \(|\bar H|\) and its
stabilizers are trivial, so the refined mass counts classes.
The second identity is the
\hyperlink{proof.weighted-unipotent-covolume}{class-unipotent
covolume} in the real frame coordinates used for \(P\).

Every factor in the packet allowance can now be evaluated. With
\(\mathscr M=C^2/(\deg_{\rm Bes}\operatorname{mass}_r)\) one has
\[
 \frac{\Pi V^n}
      {\mathscr M\deg_{\rm Bes}\,c_G^f(\nu F_{\rm sc}/d_B)^2}
 =(1-N^{-k})^{-f}\frac{\operatorname{mass}_r}{\operatorname{mass}_k}
   \prod_{j=k}^{r-1}
       \left({(2\pi)^{j+1}\over
                   \operatorname{rd}_E^{\,j+1/2}j!}\right)^f .
\]
Indeed \(C=h d_B\), \(\Pi=|\bar H|U_0^2h c_*\), and the powers of
\(N\) reduce to \((1-N^{-k})^{-f}\).
The remaining factor before the unitary volumes is
\(|D_E|^{1/4-\varrho/2}V^{2s_U+n+k}\).
Substitute \(V=(2/\operatorname{rd}_E)^f\) and
\(|D_E|=\operatorname{rd}_E^{2f}\). The power of the root
discriminant is \(-f(r^2-k^2)/2\); the powers of \(2\) before
the sphere product use \(2s_U+n+k-1=\sum_{j=k}^{r-1}j\).
Finally \(\operatorname{vol}U(r)/\operatorname{vol}U(k)
=\prod_{j=k}^{r-1}2\pi^{j+1}/j!\), giving the displayed identity.

The masses in this ratio are those of the standard positive definite
self-dual genera in ranks \(j\). The relative discriminant of \(E/F\)
is the unit ideal by the
\hyperlink{proof.weighted-unipotent-covolume}{covolume calculation}.
Apply the \hyperlink{proof.unitary-compact-dual-mass}{self-dual mass
formula} \eqref{eq:unitary-self-dual-mass} and its consecutive
critical-value factor \eqref{eq:unitary-critical-mass-value}. Since
\(|D_F|=\operatorname{rd}_E^f\), their quotient gives
\[
 \frac{\operatorname{mass}_{j+1}}{\operatorname{mass}_j}
  =\left(\frac{\operatorname{rd}_E^{j+1/2}j!}{(2\pi)^{j+1}}\right)^f
     L_F(j+1,\eta_{E/F}^{\,j+1}).
\]
The absolute logarithm of the last factor is at most
\(f\log\zeta(j+1)\), by the absolutely convergent Euler product.
This is at most \(f\,2^{-j-1}(1+2/j)\) by bounding the zeta tail
beyond its first term by the integral from \(2\). For \(j\ge4u+12\)
and \(8f\le2^u\) we obtain
\begin{equation}
 \label{eq:weighted-consecutive-mass}
 \left|\log\frac{\operatorname{mass}_{j+1}}{\operatorname{mass}_j}
    - f\log\frac{\operatorname{rd}_E^{j+1/2}j!}{(2\pi)^{j+1}}\right|
       \le 2^{1-u}.
\end{equation}
Take \(u=\varrho\) throughout \(k\le j<r\). The logarithm of
the packet ratio above is bounded above by
\[
       2f/N+2(r-k)\,2^{-\varrho}
             =o(b^{-9}).
\]
Here \(\log N\ge m\log b\), \(f\) grows polynomially in \(x\),
\(\varrho=m/2-1\asymp\sqrt{x}\), and \(\log b=x^\varepsilon\).
Thus the ratio is at most \(1+b^{-9}\) on all sufficiently large
selected squares.\footnote{\weightedpin
{grh-compact/src/New/Filling/FillingAdmissible.lean}{12-L37}
{The column-level index and its class count},
\weightedpin{grh-compact/src/New/Filling/GlobalToLocal.lean}{245-L297}
{Finite unipotent mass and the image of the middle subgroup},
\weightedpin{grh-compact/src/Lib/Adelic/Transference/MassBasic.lean}
{188-L213}{Norm-one vector count in the formal proof}, and
\weightedpin{grh-compact/src/Lib/SiegelWeil/Mass/MassRatio.lean}
{16-L48}{Consecutive mass ratios from the norm-one count}.}

Combining the packet and sieve inequalities requires
\[
 {1\over\lambda_0^2}\prod_{a=0}^{\varrho}c_a
      ={V^n d_B^2\over c_G^f F_{\rm sc}^2}.
\]
This uses the exponent \(k-1+2(\varrho-a)\) in \(\lambda_0\)
at every \(a\), including the terminal index. The mass comparison
now bounds the finite packet sum, divided by
\(\nu^2\mathscr M b^{-6k|A|}\), by
\[
 (1+b^{-6})(1+b^{-9})
    \bigl((1+b^{-10})^f(1+b^{-6})
                    +b^{-6k(f-|A|)-10}\bigr)\le 1+b^{-4}.
\]
Indeed \(f\le b/1000\), \(b\ge1000\), and \(|A|\le f\).

\paragraph{The comparison on the original column space.}
\leavevmode\teamanchor{proof.weighted-finite-bessel-comparison}%
The packet bound is independent of the finite choices. Exhaust \(T\) with the finitely many labels and heads fixed. The
off-\(T\) indicator in \(\mathcal A_{p,T}\) tends to one, with
the absolutely summable rational amplitudes as dominator.
Next exhaust the pivot values inside each block using that
absolute sum. Finally take the supremum over finite
lists of blocks, after squaring within each link. We obtain
\[
 \sum_{s\in\Omega_A,\,c,c',\,\Lambda}
  \bigl\|A^{\rm cut}_{\{\operatorname{disc}\ne0,\,
                         \operatorname{link}=\Lambda\}}
                         (c,s,c')\bigr\|_{\rm HS}^{\,2}
        \le(1+b^{-4})\nu^2 b^{-6k|A|}\mathscr M.
\]
The \hyperlink{proof.weighted-crude-energy}{crude energy comparison}
below uses this to bound its retained-link term before accounting
for the cutoff error.

For the complementary energies the finite comparison matrix on
the \(C\) original column coordinates is
\[
 \begin{aligned}
 (T_S(s))_{(c,i),(c',j)}
   &=\nu^{-1}\overline{A_{S(F)}(c,s,c')_{ij}},\\
 (K(s)-T_S(s))_{(c,i),(c',j)}
   &=\nu^{-1}\overline{A_{S(F)^c}(c,s,c')_{ij}}.
 \end{aligned}
\]
These are the full-vector blocks and the conjugation is that of
the original Petersson unfolding. The finite tensor \(C_p(l)\)
has coordinates on \(Q\times Q\) and occurs with the cut real
correlation in the open-link inequality. In the original
unfolding the finite row condition contributes the principal
coset volume \(\mu_f(K(g))\), which cancels the same scalar in
\(P=(\mu_f(K(g))\nu)^{-1/2}\mathcal Q\) for the Poincar\'e map
\(\mathcal Q\). The \hyperlink{proof.weighted-bessel-unfolding}{same
Petersson unfolding} gives the identity block on \(S(F)\) and
the unitary transport matrices \(T_S(t)\) for \(t\in J\).
The \hyperlink{proof.weighted-complement-crude}{complementary
energy comparison} uses these same coordinates to bound
\(\Delta_i=\mathscr E_{S(F)^c}(JG_i)/(|J|C)\).\footnote{\weightedpin
{grh-compact/src/New/Filling/PoincareSeries.lean}{176-L196}
{Original Poincar\'e normalization}, \weightedpin
{grh-compact/src/New/Filling/AmplitudeSums.lean}{163-L216}
{Class transport and identity block}, and \weightedpin
{grh-compact/src/New/Filling/Moments.lean}{105-L213}
{Fixing-set and complementary consumers}.}

\paragraph{The two energy ratios on the selected datum.}
\teamanchor{proof.weighted-energy-ratios}%
Here \(r=k+m-1\) and \(\bar w=f^{-1}\sum_v\log w(v)\).
Use \(U_0=g^{f\varrho(2k+2\varrho+1)}\)
for the order of the reduced unipotent group and \(d_B\) for the
number of Bessel coordinates per column class. The class count and the decomposition
of \(J\) give the exact ratio
\[
 \frac{\mathscr M}{\mathscr D}
   = \frac{(N^k-1)^f\operatorname{mass}_k\,d_B}
      {U_0\,\deg_{\rm Bes}\operatorname{mass}_r}.
\]
Indeed the number of column classes is
\(|\bar H|(N^k-1)^f\operatorname{mass}_k\), and
\(|J|=|\bar H|U_0\). These are the class and unipotent counts
used in the
\hyperlink{proof.weighted-open-mass}{open packet argument}.

For the real Bessel coordinates at \(v\), the \(U(X_0)\)-type
in the construction is \(\operatorname{Sym}^{\ell_v}\mathbf C^k\)
with \(0\le\ell_v\le a\). Thus its dimension multiplied by
\((k-1)!\) is \(\prod_{j=1}^{k-1}(\ell_v+j)\le(a+k)^{k-1}\).
Here \(a+k\le 2xb\le b^{5/4}\) for all sufficiently large selected
squares, using \(a<xb+N_1\) and the polynomial sizes of \(N_1,k\).
Consequently%
\footnote{\weightedpin{grh-compact/src/New/Filling/FillingData.lean}
{64-L84}{The symmetric power in the Bessel space}, and
\weightedpin{grh-compact/src/New/Filling/DegreeBounds.lean}
{42-L50}{The dimension bound on the filling datum}.}
\[
 0\le f^{-1}\log d_B+\log((k-1)!)\le \tfrac54 k\log b .
\]

In the Weyl product for \(\deg_v\), the factors involving only
the first \(k-1\) parameters cancel against their denominators.
Here \(k=r+1-m\) is even. For \(i\le k-2\) in zero-based
indices, \(\lambda_i=k/2-1-i\ge1-k/2\).
For \(l\ge k-1\) put \(j=l+2-k\), so
\(\lambda_l=h_2-w_j(v)\). The parameter bounds
\(w_j(v)-h_2\ge w(v)+2jh_2\) and \(h_2=r^8\ge k\) yield
\[
    \lambda_i-\lambda_l\ge w(v)+2h_2+1-k/2\ge w(v).
\]
Every remaining difference is at least one and at
most \(w(v)e^{1/x}\), including the differences from the initial
parameters. The denominator contributed by the parameter of index
\(j\), starting from zero, is \(j!\). It follows that
\[
 f^{-1}\log\deg_{\rm Bes}
 = mk\bar w-\sum_{j=k-1}^{r-1}\log(j!)
          +O(m^2\log b+mr/x).
\]
For the lower bound before taking this error, use
\(m(k-1)\bar w\). For the upper bound use \(mr(\bar w+x^{-1})\).
\footnote{\weightedpin{grh-compact/src/New/Seesaw/IndexTypes.lean}
{124-L126}{The row parameters},
\weightedpin{grh-compact/src/New/Data/Weights.lean}
{29-L54}{The nominal and actual weights}, and
\weightedpin{grh-compact/src/New/Filling/DegreeBounds.lean}
{52-L78}{The Weyl product on these parameters}.}

Sum the consecutive mass estimate \eqref{eq:weighted-consecutive-mass}
over \(k\le j<r\). The exponents of
\(\operatorname{rd}_E\) and \(2\pi\) in its logarithm are both
\((m-1)k+O(m^2)\), with opposite signs. Hence
\[
 f^{-1}\log\frac{\operatorname{mass}_r}{\operatorname{mass}_k}
 =(m-1)k\log\frac{\operatorname{rd}_E}{2\pi}
   +\sum_{j=k}^{r-1}\log(j!)+O(m^2\log b+m2^{-\varrho}/f).
\]
Also \(f^{-1}\log U_0=(m-2)k\log g+O(m^2\log b)\).
Insert these expressions in the exact ratio and use
\[
 \log N=m\bar w+(m-1)\log\frac{g\operatorname{rd}_E}{2\pi}
            -\tfrac{15}{2}\log b+O(\sqrt{\log b})
\]
for the selected prime \(N\). The weight terms cancel and the two
factorial sums leave \(\log((k-1)!)\). Thus
\[
 f^{-1}\log(\mathscr M/\mathscr D)
 =k\bigl(\log g-\tfrac{15}{2}\log b\bigr)
   + f^{-1}\log d_B +\log((k-1)!)+\mathcal E ,
\]
where
\[
 |\mathcal E|
  \ll k\sqrt{\log b}+m^2\log b+mr/x+m2^{-\varrho}/f+N^{-k}.
\]
Since \(m\asymp\sqrt{x}\), \(k\asymp x^5\), and
\(\log b=x^\varepsilon\), this error is at most
\(\tfrac18 k\log b\) on all sufficiently large selected squares.
The lower bound \(g\ge b^{20}\) and the Bessel dimension estimate
therefore give
\[
     b^{11kf}\le\frac{\mathscr M}{\mathscr D}
       \le (g/b^{49/8})^{kf}\le \tfrac12(g/b^6)^{kf}.
\]
The last comparison uses \(b^{kf/8}\ge2\). This proves
\(\tau^{44}\mathscr D\le\mathscr M\) and
\(2(b^6/g)^{kf}\mathscr M\le\mathscr D\), with
\(\tau=b^{kf/4}\) on the same datum. The auxiliary factors used
in the crude bound also satisfy \(2\tau^2(b/g)^{kf}\le1\),
\(N^f\le\tau\), and \(\tau\ge b^{n+31}\), where \(n=\varrho+1\).
For these use \(g\ge b^{20}\), \(N\le b^{60n}\) and
\(k\ge1000(2n)^9\).
\footnote{\weightedpin
{grh-compact/src/Lib/SiegelWeil/Mass/MassRatio.lean}
{16-L48}{Consecutive mass ratios from the norm-one count},
\weightedpin{grh-compact/src/New/Seesaw/IndexTypes.lean}
{131-L162}{The window of the selected ramification prime}, and
\weightedpin{grh-compact/src/New/Filling/Setup.lean}
{291-L318}{The energy inequalities used in the two applications}.}

\paragraph{Compatible frame orbits.}
\teamanchor{proof.weighted-compatible-orbits}%
The character test and the flag argument used for
\hyperlink{proof.gram-orbit-expansion}{the earlier Gram expansion}
apply to the weighted frame as follows. For a rational term
\(\gamma\) over \((c,s,c')\), put
\(B=U\cap\gamma U\gamma^{-1}\) and
\(\Lambda=U_c\cap\gamma U_{c'}\gamma^{-1}\).
After joining the cutoff integrals in its class-unipotent double
orbit, the quotient of integration is
\(\Lambda\backslash G_\infty\).
The intersection \(B\) is a unipotent group over \(F\), and
\(\Lambda\) is its rational intersection with an open compact subgroup
at the finite places. Unipotent strong approximation gives a compact
quotient \(\Lambda\backslash B_\infty\). The integrand transforms on
\(B_\infty\) by the ratio of the two Bessel characters.
This character is trivial on \(\Lambda\). If the two finite
characters disagree on any rational element of \(B\), the product
formula gives a nontrivial ratio on \(B_\infty\).
Disintegrating the absolutely convergent integral over
\(\Lambda\backslash B_\infty\) then gives zero by translation on
this compact quotient. We can therefore restrict the rational
orbits to the pairs with agreeing characters.

For such a pair let \(f_0,\ldots,f_{\varrho}\) be \(e_0,\ldots,
e_{\varrho-1},z\), and write \(f'_i=\gamma f_i\).
Let \(X_i,X_i'\) be their spans through index \(i-1\).
Suppose that a common vector has leading coefficients
\(c,c'\ne0\) at indices \(i,i'<\varrho\). Denote the vector by \(u\).
For \(y\) perpendicular to \(X_{i+1}+X_{i'+1}'\), the isometry
\[
 v\longmapsto v+h(v,y)u-h(v,u)y+\ell\, h(v,u)u,
 \qquad \ell+\bar\ell=-h(y,y),
\]
belongs to both flag unipotents. Here \(h\) is linear in its first
argument. The sum of its simple entries for the first flag is
\(c\,h(f_{i+1},y)\). Earlier simple entries vanish by the
orthogonality imposed on \(y\), and later entries vanish by the level
of \(u\). The second flag gives \(c'\,h(f'_{i'+1},y)\).
The agreement is tested with \(\chi_f^+=\psi_f(\sigma_U/(g\varpi))\)
on both flags. Inverting both characters gives the same equality.
The character equality holds
on replacing \(y\) by every \(E\)-multiple of \(y\). The field
\(E\) is dense in its finite adeles, so a nonzero difference of the
two simple-entry sums would contradict the nontriviality of this
additive character. Thus
\[
 c f_{i+1}-c'f'_{i'+1}\in X_{i+1}+X'_{i'+1}.
\]
Thus a common vector at the next pair of levels has the same leading
coefficients. A common vector with the terminal index on only one
side would have negative norm on that side and zero norm on the
other. Iteration shows that the levels are equal. At the terminal
index both distinguished vectors have norm \(-N\), so
\(c\bar c=c'\bar c'\). A common vector already at the terminal level
gives that equality directly.

Write \(X=X_{\varrho}\), \(X'=X_{\varrho}'\) and \(I=X\cap X'\).
At a real place \((X'+Ez')^\perp\) is positive semidefinite with
radical \(X'\), where \(z'=\gamma z\).
If a common vector occurs, let \(j\) be the first level with
one. The lifted common vectors from \(j\) through \(\varrho-1\)
are independent, and the intersection \(I\) meets \(X_j\) and
\(X'_j\) trivially. Hence \(X=X_j\oplus I\) and
\(X'=X'_j\oplus I\). The terminal common vector gives
\(z'-\zeta z\in X+X'\) with \(\zeta\bar\zeta=1\).
A vector of \(X_j\) perpendicular to \(X'_j\) is also perpendicular
to \(I\) and hence to \(X'\). The terminal relation makes it
perpendicular to \(z'\). Its isotropy then puts it in \(X'\), by
the positive semidefinite description. Thus it is zero and the
leading frame minor of size \(j\) is nonzero. The convention for
size zero is \(1\).

If \(X+Ez\) and \(X'+Ez'\) have zero intersection, their cross
pairing \(Z\) is nonsingular. A vector in its left kernel has
nonpositive norm as an element of \(X+Ez\). Since it is in
\((X'+Ez')^\perp\), it must have zero norm and belong to \(X'\).
It also belongs to \(X\), and hence is zero. For a pair with
zero discriminant the equation
\(|\det Z|^2=N^2|\det Z_{\varrho}|^2\) now gives a nonsingular
pairing between \(X\) and \(X'\). Their joint span is hyperbolic.
Project \(z,z'\) orthogonally onto its complement. We subtract
only an element of \(X\) from \(z\) and an element of \(X'\) from
\(z'\), so the projections have norm \(-N\).
The vanishing discriminant says that their joint Gram matrix is
singular, using the block Gram determinant of the two augmented
frames. The complement has negative index one. Subtracting a
multiple of the first projected vector from the second then gives
an isotropic vector in a positive definite space, so the two
projections are proportional. This gives a nonzero intersection of
the augmented spans. Consequently every compatible pair with
zero discriminant has a first common level \(j\le\varrho\),
with common leading coefficients \((\zeta,1)\) at every
level from \(j\) on and with the leading minor of size \(j\) nonzero.
\footnote{\weightedpin{grh-compact/src/New/Filling/CrudeBoundDefs.lean}
{53-L95}{Levels and common vectors} and
\weightedpin{grh-compact/src/New/Filling/CrudeEnergy.lean}
{22-L82}{Character vanishing and the degenerate levels}.}

At depth zero the relations make \(\gamma\) upper triangular on the
entire augmented frame, with constant frame diagonal \(\zeta\).
The class representatives fix the frame, and the lattice change
preserves its lines with the same scale on both sides of a term
position. Thus the multiplier on its first vector is still
\(\zeta\). Integrality of the position gives
\(\zeta\in\mathcal O_E\). For a position in \(JG_i\), its component
away from \(G_i\) fixes that first vector, so \(\zeta\) reduces to
one at a chosen place above \(g_{1-i}\). At this split place the
completion is \(\mathbf Q_{g_{1-i}}\) and its principal units are
torsion-free since \(g_{1-i}\) is odd. The norm-one algebraic
integer \(\zeta\) is a root of unity, so it equals one.
The upper unitriangular frame condition then puts \(\gamma\) in
the Bessel subgroup \(S(F)\). This is the required depth-zero
exclusion for the complementary energy.

\paragraph{The class cutoffs in the frame box.}
\teamanchor{proof.weighted-box-cutoff}%
Write \(a=\varrho\), \(n=a+1\), \(B=b^{2n^2}\), and
\(\theta=b^{6n+2}\). Let \(\Gamma_c=U_c\) be the rational class-unipotent
group at level \(g\), and \(\Gamma_c^0=U_c^0\) its subgroup fixing the
frame \(\mathbf f=(e_0,\ldots,e_{a-1},z)\) pointwise. At the real places
use the normalized frame with scales \(s_{vh}=\kappa_v^{a-h}/\sqrt N\)
for \(0\le h\le a\). The set \(\mathrm{Box}\) consists of the points \(y\)
whose inverse frame vectors have length at most \(B\) at every place
and for which the local size is at most \(\theta\). Both conditions
are left invariant under \(U_\infty^0 H_{0,\infty}\).

The rectangular coordinates in the
\hyperlink{proof.weighted-unipotent-covolume}{class-unipotent fibre}
give an element of \(\Gamma_c\) with rational frame action
\(1+gN x E_{ih}\) for every \(i<h\le a\) and \(x\in\mathcal O_E\).
Use the single free entry \(gN x\) at column \(f_h\) in row \(i\).
The hyperbolic coordinates are integral in the class basis and the
marked coordinate has denominator dividing \(N\), so the row
divided by \(g\) is integral in that basis. The skew entry can be
chosen as in the fibre computation.
In the real frame the same entry is multiplied by \(s_{vh}/s_{vi}\),
of absolute value at most one. In the ordinary Euclidean norm on
\(\mathbf C^f\), the trace dual of \(\mathcal O_E\) is
\(2\mathfrak D_E^{-1}\). Its shortest length is at least \(2/b\),
since \(\mathfrak D_E=(\varpi)\), \(|\Phi_v(\varpi)|\le b\), and a
nonzero algebraic integer has some conjugate of absolute value at
least one. The \hyperlink{proof.weighted-short-spanning}{elementary
covering-radius estimate}
\(\mu(L)\le (2f)^2/\lambda_1(L^*)\) therefore gives a covering
radius at most \(b^3\) for \(\mathcal O_E\), using \(f\le b/1000\).

At an arbitrary \(y\), choose \(u_v\in U_v\) with
\(\|u_vy_v\|_{\rm op}\le\operatorname{size}(y_v)\), by the definition
of the size with the additive one. Starting with the frame matrices
of \(u_v\), multiply on the right by the preceding rational shears,
processing row indices from largest to smallest. A shear at \((i,h)\)
changes that entry by \(gN(s_{vh}/s_{vi})\Phi_v(x)\), and changes no
entries in lower rows. Approximation in \(\mathcal O_E\) at once over
all \(v\) leaves absolute value at most \(gN b^3\le b^{60n+26}\) in
each processed entry. For the resulting \(q\in\Gamma_c\), the inverse
frames of \(q^{-1}y\) have lengths at most
\(b^{60n+27}\operatorname{size}(y_v)\) at \(v\). This follows by
applying \((u_vy_v)^{-1}\), whose operator norm equals that of
\(u_vy_v\), to the columns of the frame matrix of \(u_vq_v\).

Let \(\mathcal F_c(y)\) be the left \(\Gamma_c^0\)-cosets of
\(t\in\Gamma_c\) for which \((ty)_v^{-1}\mathbf f\) obeys these last
length bounds. This is a finite nonempty set. Finiteness follows by
applying \(y_v\) to the bounded columns. The resulting matrices have
bounded frame coefficients at every embedding, and the rational
coefficients along the isotropic rows are integral in the class
chart, which fixes those rows and the frame. Equal frame matrices
give the same coset. Put
\[
 \omega_c(y)=\frac{{\bf1}_{[1]\in\mathcal F_c(y)}}{|\mathcal F_c(y)|}.
\]
This function is measurable and belongs to \([0,1]\). Every element
of \(U_\infty^0H_{0,\infty}\) fixes the inverse frame and the size in
this test, since \(t_v^{-1}\mathbf f\) is in the frame span.
For \(t\in\Gamma_c\) the cosets counted at \(ty\) are the cosets
counted at \(y\) right translated by \(t^{-1}\). Thus
\(\sum_{t\in\Gamma_c^0\backslash\Gamma_c}\omega_c(ty)=1\).
When \(\omega_c(y)>0\) and the cut holds at \(v\), the inverse frame
lengths at that place are at most \(b^{66n+29}\le B\).
Choose a measurable partition cutoff \(p_c\) for the discrete left
action of \(\Gamma_c^0\) on \(G_\infty\). For example, the indicator
of a measurable fundamental section is such a cutoff. Take
\(\beta_c=p_c\omega_c\). Summing first within the left cosets shows
\(\sum_{t\in\Gamma_c}\beta_c(ty)=1\). This is the partition of unity in
the Poincar\'e unfolding, so it gives the same periodization and the
same fibre integral \(\nu\).\footnote{\weightedpin
{grh-compact/src/New/Filling/Setup.lean}{239-L255}
{Class groups and cutoff data}, \weightedpin
{grh-compact/src/New/Filling/Setup.lean}{330-L357}
{Discreteness and the two cutoff factors}, and \weightedpin
{grh-compact/src/New/Filling/BoxPartition.lean}{451-L508}
{Frame reduction and coset partition}.}

\paragraph{Separation of the rational frames in the cutoff.}
\teamanchor{proof.weighted-cut-frame-packing}%
Let \(n=\varrho+1\), \(B=b^{2n^2}\), and write \(f_0,\ldots,f_{\varrho}\)
for the augmented frame, with \(f_{\varrho}=z\).
Consider two rational terms in the same residue count after their
class-unipotent reductions into the cutoff box. Write \(\gamma\)
for their relative rational isometry, and \(y_v\) for the real
coordinates of the first reduced term. Use the supremum, over \(v,a\),
of Euclidean distance between corresponding vectors
\(y_v^{-1}f_a\) and \(y_v^{-1}\gamma_v f_a\). If this distance is
less than \((NB)^{-1}\), the two rational frames coincide.

Indeed the class representatives preserve the frame. The relative
integral term position consequently makes every frame pairing
\(h(f_b,\gamma f_a)\) an algebraic integer: in the class lattice
chart these are entries of the product of its integral Gram matrix
and an integral isometry. The entries of \(\gamma f_a\) along the
first \(\varrho\) coordinate rows are integral as well, since the
class change acts on the middle block and fixes those coordinate
rows and the marked frame. In the normalized real frame,
the scale at \(f_a\) is \(s_{va}=\kappa_v^{\varrho-a}/\sqrt N\),
with \(1\le|\kappa_v|\le b^{23}\).
The real Gram matrix has operator norm one. Pairing the difference
of the two vectors with \(y_v^{-1}f_b\), of length at most \(B\),
bounds each conjugate of
\(h(f_b,\gamma f_a)-h(f_b,f_a)\) by \(NB\) times their distance.
The factor \(N\) covers \(1/|s_{va}s_{vb}|\).
An algebraic integer whose absolute values at all embeddings are
less than one vanishes, by its norm. All these pairings agree.

The pairings with the isotropic frame and with \(z\) fix the dual
frame coordinates and the \(z\)-coordinate. Unitarity now makes
the component of \(\gamma f_a-f_a\) in the definite middle space
have norm zero. Thus this difference lies in the isotropic span.
For fixed \(a\), recover its coordinates there starting with the
last one. In the normalized frame the coefficient at \(e_i\) is
\((s_{va}/s_{vi})\Phi_v((\gamma-I)_{i,a})\), where the index \(a\)
on the right denotes the frame column. As \(y_v\) is in the cut,
choose \(u_v\in U_v\) with
\(\|u_v y_v\|_{\mathrm{op}}\le b^{6n+2}\). The \(i\)-th row of
the upper unitriangular action of \(u_v\) reads precisely that
coefficient once the later coefficients vanish. Hence
\[
 |\Phi_v((\gamma-I)_{i,a})|
 \le b^{23\varrho+6n+2}
       \,\|y_v^{-1}\gamma_v f_a-y_v^{-1}f_a\|
 <1 .
\]
Here \(b^{23\varrho+6n+2}\le NB\) for the selected large squares.
Integrality and the same norm argument finish the descending
induction. This proves the asserted separation.

Translate one frame to the origin in the real subspace containing
the differences, of dimension \(D\le2rnf\) in the unrestricted
count. The resulting centres lie in a norm ball of radius \(2B\).
Their open balls of radius \(1/(2NB)\) are disjoint, so comparison
of Lebesgue volumes bounds their number by
\((1+4NB^2)^D\le b^{6n^2D}\); here \(N\le b^{60n}\) and \(n\)
tends to infinity. For two reduced terms with the same source
class and residue, the relative finite position is the quotient
of their two integral positions and belongs to \(K(g)\). When
their frames coincide, class freeness gives the same target
class and the relative term is a class unipotent fixing the frame.
The boxing corrections are themselves class unipotents.

Fix the source class, residue and \(t\), and take any right translate
\(x\in G_\infty\). For each target class \(c'\), group the rational
terms by right multiplication by \(U_{c'}\). This preserves the
finite position modulo \(K(g)\) and the condition
\(\gamma^{-1}tx\in\mathrm{InCut}\), because the cut is left
\(U_\infty\)-invariant. In each orbit the weights
\(\beta_{c'}(\gamma^{-1}t)\) at the unshifted arguments sum to one
for the full orbit and to at most one for a restricted term set.
Use the partition at the shifted argument to choose a representative
with \(\omega_{c'}(\gamma^{-1}tx)>0\). Its inverse frame
\((tx)^{-1}\gamma\mathbf f\) lies in the box. Distinct orbits
have distinct boxed frames by class freeness and the relative
position argument above. The separation and volume comparison
therefore bound the total weight by \(b^{12rn^3 f}\) for each
right translate in the positive majorant.\footnote{\weightedpin
{grh-compact/src/New/Filling/NearTerms.lean}{252-L317}
{Integrality, class comparison and the fixed frame},
\weightedpin{grh-compact/src/New/RealPlace/FrameNorm.lean}
{25-L45}{Reading a frame vector after a unipotent translation}, and
\weightedpin{grh-compact/src/New/Filling/BoxPartition.lean}
{510-L555}{Separation and the weighted frame count}.}

\paragraph{The full and cut relative traces.}
\teamanchor{proof.weighted-cut-energy}%
Use the Bessel vectors and the class cutoffs of the normalized
unfolding. Write $a=\varrho$, $n=a+1$, $r=k+2a+1$, and
$d_W=\prod_v |I_v|$, so $C=|\mathrm{Cl}|d_W$. For a family $B$ of
matrix blocks put
\[
 \|B\|_\Omega=\nu^{-1}
     \biggl(\sum_{s\in\Omega,c,c'}\|B(c,s,c')\|_{\rm HS}^2\biggr)^{1/2}.
\]
All the residue sums use counting measure. The norm of the tensor
formed from two product Bessel vectors is the product of their norms.
For the cut pair \(W_v,W_v^\flat\) use the
\hyperlink{proof.weighted-bessel-tail-truncation}{real majorant and decay}
\eqref{eq:weighted-bessel-cut-majorant}--\eqref{eq:weighted-bessel-cut-decay}.
Write \(M_v=A_v\) and \(m_v=\mu_v\) for the \(U_v\)-invariant
majorant and its finite positive measure at \(v\).
The norm in these bounds is the one on the full vectors
and their cut, and the integral uses the same quotient Haar.
Here \(\mathrm{InCut}_v=\{\operatorname{size}\le R\}\) as in the
\hyperlink{proof.weighted-poincare-realization}{absolute periodization}.

Write $M(t)=\prod_v M_v(t_v)$ and let $K_\gamma(t)$ and
$K_\gamma^\flat(t)$ be the full and cut integrands before their two
class cutoffs. If the two tensors differ, some local vector differs
at $t_v$ or at $(\gamma^{-1}t)_v$. Hence
\[
 \|K_\gamma(t)-K_\gamma^\flat(t)\|_{\rm HS}
 \le 2b^{-3r^3 f}M(t)M(\gamma^{-1}t).
\]
At fixed $c,s,t$ and right translate $x\in G_\infty$, the sum of
$\beta_{c'}(\gamma^{-1}t)$ over $(c',\gamma)$ with
$\gamma^{-1}tx\in\mathrm{InCut}$ is at most $b^{12rn^3 f}$ by the
\hyperlink{proof.weighted-cut-frame-packing}{weighted frame count}.
This holds for each right translate \(x\) in the positive measure.
Integrate this count against $\prod_v m_v$ using the majorant of
$M(\gamma^{-1}t)$. The integral of $\beta_c(t)M(t)$ is
$\nu\prod_v\int_{U_v\backslash G_v}M_v$ by the class cutoff partition
and the unipotent covolume. Thus for every $c,s$ and any term set $T$,
\[
 \sum_{c'}\|A_T(c,s,c')-A_T^\flat(c,s,c')\|_{\rm HS}
 \le 2\nu d_W b^{(12rn^3-r^3)f}
 \le \tfrac12\nu b^{-k^3 f}.
\]
For the last inequality use $d_W\le b^{2kf}$,
$r^3-k^3\ge r^2$, and $12rn^3+2k+2\le r^2$.
The latter follows from $r\ge24n^3$, $k+1\le r$ and $r\ge4$.
The square of the sum in $c'$ bounds the sum of the squares. Since
$|G|\le g^{r^2 f}\le b^{23r^2 f}$ and
$|\mathrm{Cl}|\le \mathscr D=|J|C$, we get
\begin{equation}
 \label{eq:weighted-cut-energy-error}
 \|A_T-A_T^\flat\|_\Omega^2\le b^{-40}\mathscr D.
\end{equation}
Here $2k^3-23r^2\ge 2k^2(k-46)\ge40$ by $r\le 2k$ and $k\ge47$.
The bound also proves convergence of the difference with these
weights. Absolute convergence of the two amplitude sums permits
their use on every term set in the unfolding.

\paragraph{The Haar measures on the two real unipotent groups.}
Work at a fixed infinite place in the normalized frame. Write \(m=\varrho\),
\(n=m+1\), \(k=\dim X_0\), and \(r=k+2m+1\), and order the blocks as
\((e,X_0,z,e^*)\). Use the Euclidean operator norm, so \(g^{-1}=Jg^*J\) gives
\(\|g^{-1}\|=\|g\|\). Write \(F=(e_0,\ldots,e_{m-1},z)\) for the
frame matrix. In the \hyperlink{proof.weighted-real-gram-section}{Gram
section coordinates} the frame stabilizer is \(U^0H\), with
\(H=U(X_0)\) and \(U^0\) parametrized there by \((B,A_{\rm Gr})\).
For the tube write \(A=-B^*\) and \(S=A_{\rm Gr}\). This linear
change has absolute real determinant one, and the \(e,e^*\)
block is \(S-AA^*/2\).

Multiplication of \(u(A,S)\) and \(u(B,V)\) adds \(A,B\) and replaces the
second coordinate by \(S+V+(BA^*-AB^*)/2\). Inversion negates \(A,S\).
Thus \(du_0=dA\,dS\) is left and right Haar, with \(dS\) Lebesgue in the
entries \(S_{ab}\in\mathbf C\) for \(a<b\) and \(\operatorname{Im}S_{aa}\in
\mathbf R\). To compare the normalization, the real coordinates of the
pinned skew matrix are sent to these coordinates by
\[
 (c_{ab},c_{ba})\longmapsto
 \bigl((c_{ab}-c_{ba})/\sqrt2,\ (c_{ab}+c_{ba})/\sqrt2\bigr)
 \quad(a<b),\qquad c_{aa}\longmapsto\operatorname{Im}S_{aa}.
\]
All these maps have absolute determinant one. At fixed \(A\) the free
entries of \(S-AA^*/2\) differ from those of \(S\) by a translation.
Consequently the prescribed Haar measure on \(U^0\) is \(du_0=dA\,dS\).\footnote{\weightedpin{grh-compact/src/Lib/DiscreteSeries/Haar/UnitaryNorm.lean}{25-L34}{Real skew coordinates},
 \weightedpin{grh-compact/src/Lib/DiscreteSeries/Bessel/UnipotentCoords.lean}{174-L185}{The frame-fixing group},
 and \weightedpin{grh-compact/src/Lib/DiscreteSeries/Bessel/UnipotentCoords.lean}{228-L244}{The Haar prescriptions}.}

The frame action \(uF=FQ(u)\) of \(U\) has kernel \(U^0\) and image the
upper unitriangular \(n\)-by-\(n\) group. For
\(Q=\left(\begin{smallmatrix}a&v\\0&1\end{smallmatrix}\right)\)
use the lift fixing \(X_0\), taking \(e\) to \(ea\), \(z\) to \(z+ev\),
and \(e^*\) to \((e^*+zv^*+evv^*/2)a^{-*}\). Here \(C^{-*}=(C^*)^{-1}\)
for any invertible matrix \(C\). These columns have the original
\(J\)-pairings, as follows from \(a^*a^{-*}=1\) and the cancellation
\(vv^*/2+vv^*/2-vv^*=0\) in the dual Gram matrix. The lift belongs to
\(U\) by its triangular form. Its conjugation action on \(U^0\) sends
\((A,S)\) to \((aA,aSa^*)\). Both maps preserve the displayed Lebesgue
measures since \(a\) is unipotent and the corresponding linear maps
are unipotent. Multiplication on either side in the upper unitriangular
group is triangular in the entries ordered by superdiagonal distance
and has diagonal one. The product \(dQ\,du_0\) under
\((Q,u_0)\mapsto u_0\widetilde Q\) is therefore Haar on \(U\)
and agrees with the prescribed product measure. It also gives
Lebesgue in the full free entries. In fact the \(e,e^*\) block of
\(u(A,S)\widetilde Q\) is
\((S-AA^*/2+vv^*/2)a^{-*}\). For fixed \(A,Q\), right multiplication
by \(a^{-*}\) is a shear, triangular in descending column order in
each free upper row, with coefficient one on each complex entry
and on the imaginary diagonal.
The remaining free entries are \(A\) and the strict upper entries of \(Q\).
Here \(dQ\) is \(d\Re\,d\Im\) in each free complex entry.\footnote{\weightedpin{grh-compact/src/Lib/DiscreteSeries/Bessel/UnipotentCoords.lean}{65-L83}{Free entries},
 \weightedpin{grh-compact/src/Lib/DiscreteSeries/Bessel/UnipotentCoords.lean}{187-L213}{The frame-action lift},
 and \weightedpin{grh-compact/src/Lib/DiscreteSeries/Bessel/RealMeasures.lean}{197-L211}{The full free-entry measure}.}

Write \(dy,dh,c_G\) for the filling Haar, the probability Haar
on \(H\), and the constant in the
\hyperlink{proof.weighted-real-gram-measure}{Gram density formula}.
Write \(d\dot y\) for the quotient disintegration of \(dy\) by the
just computed \(du\) on \(U\).
The local data are read after the diagonal basis change. At place \(v\)
the frame scales are \(s_{va}=\kappa_v^{m-a}/\sqrt N\), with the
reciprocal scales on \(e_a^*\) and the unscaled \(X_0\). The integral
coordinates and rational term are transported before taking operator
norms, free entries, and Gram pairings. In particular the term
\(\gamma\) of the real calculation is the rescaled image of the
rational term. Each leading principal Gram determinant of order \(j\) gains
\(\prod_{a<j}|s_{va}|^2\) from row and column scaling. The last scale is
\(1/\sqrt N\), so the rational discriminant
\(|D_n|^2-N^2|D_m|^2\) has common factor
\((\prod_{a<n}|s_{va}|^2)^2\) in the rescaled real discriminant.
This is also the frame and Haar used by the full vector in the normalized
Poincaré map.\footnote{\weightedpin{grh-compact/src/Lib/DiscreteSeries/Haar/CompactDualMeasure.lean}{34-L48}{Filling Haar},
 \weightedpin{grh-compact/src/New/Filling/BesselSystem.lean}{16-L20}{The Gram density constant},
 \weightedpin{grh-compact/src/New/Filling/Setup.lean}{130-L148}{Basis scales},
 \weightedpin{grh-compact/src/New/Filling/Setup.lean}{164-L167}{Rational discriminant},
 \weightedpin{grh-compact/src/New/Filling/Setup.lean}{203-L215}{The rescaled real argument},
 \weightedpin{grh-compact/src/New/Filling/Setup.lean}{375-L382}{Principal determinant change},
 and \weightedpin{grh-compact/src/New/Filling/NearTerms.lean}{225-L236}{The scaled discriminant}.}

\paragraph{A compact tube with prescribed fiber bound.}
\teamanchor{proof.weighted-real-cell-tube}%
For \(u=u(A,S)\) set \(L(u)=\max(\|A\|,\|S\|^{1/2})\).
The group law gives \(L(uu')\le L(u)+L(u')\), since the central
correction has norm at most \(\|A\|\|B\|\). Also \(L(u^{-1})=L(u)\).
The dilations \((A,S)\mapsto (\lambda A,\lambda^2 S)\) multiply \(du_0\)
by \(\lambda^{D_0}\), \(D_0=2m(k+m)\le r^2\), and multiply \(L\) by
\(\lambda\). This includes \(m=0\), when \(U^0=\{1\}\) with mass one.

For \(T\ge8r\), put \(R=r^2T^2\), \(\Omega=\{L\le R\}\) and
\(A_T=\Omega B(T)\), where \(B(T)=\{y\in G_1:\|y\|\le T\}\).
The two factors are compact. The first is a closed bounded box
in the displayed coordinates and the second is a closed matrix norm ball
in \(U(J)\). The box has finite positive \(du_0\)-measure. On this box,
\(\|u\|\le1+2R+3R^2/2\le5R^2\), so
\[
 A_T\subset B(5r^4T^5)\subset B(T^9)\subset B(T^{20n}).
\]
If the fiber over \(y\) is nonempty, fix \(u_0y=\omega_0z_0\) with
\(\omega_0\in\Omega\) and \(z_0\in B(T)\), and set \(q=\omega_0^{-1}u_0\).
For another \(uy=\omega z\) in the fiber, \(uq^{-1}=\omega zz_0^{-1}\).
Here \(zz_0^{-1}\in U^0\) and its norm is at most \(T^2\).
For any \(u(A,S)\) of norm at most \(T^2\), its \(A\)-block and the
skew part \(S\) of its \(e,e^*\)-block have norms at most \(T^2\),
so \(L(u)\le T^2\). Right translation by \(q^{-1}\), followed by
homogeneous scaling of the box, gives
\[
 \int_{U^0}{\bf1}_{A_T}(uy)\,du_0
 \le (1+1/r^2)^{D_0}\operatorname{vol}_{U^0}(\Omega)
 \le e\,\operatorname{vol}_{U^0}(\Omega)
 \le T\,\operatorname{vol}_{U^0}(\Omega).
\]
The inequality holds also for an empty fiber. The readable tube proof
uses the homogeneous gauge in the unrecentered central block, and the
compact proof uses dilated matrix norm balls. The recentered box gives
the radius and fiber bounds used here with the same Haar.\footnote{\weightedpin{grh/src/Lib/DiscreteSeries/Bessel/UnipotentHaarMeasure.lean}{718-L775}{Homogeneous balls in the readable edition},
 \weightedpin{grh/src/Lib/DiscreteSeries/Bessel/UnipotentHaarMeasure.lean}{890-L933}{Its tube and right-translated slices},
 and \weightedpin{grh-compact/src/Lib/DiscreteSeries/Bessel/RealMeasures.lean}{159-L190}{The compact tube interface},
 and \weightedpin{grh-compact/src/Lib/DiscreteSeries/Bessel/RealMeasures.lean}{159-L195}{The tube and fibre bound in the original reading}.}

\paragraph{Geometry of the actual real frame box.}
\teamanchor{proof.weighted-real-cell-contraction}%
Fix \(\gamma\) with \(Z=F^*J\gamma F\) and
\(\mathfrak g=|\det Z|^2-|D_m(Z)|^2>0\). Use the notation
\(\operatorname{Box}(\rho)=\{y:\|y^{-1}f_a\|\le\rho\text{ for all }a\}\)
and \(\operatorname{size}(y)=\inf_{s\in UH}(\|sy\|+1)\). Suppose
\(\operatorname{Box}(\rho)\cap\gamma\operatorname{Box}(\rho)\ne\emptyset\),
with \(\rho\ge b\ge8r\). This gives \(|Z_{ab}|\le\rho^2\) by taking the
pairings after a common isometry. The determinant bound
\(\mathfrak g\le(n\rho^2)^{2n}\) implies \(M_2\ge\rho\) for
\[
 M_2=\rho^{9n}/\mathfrak g,\quad P_2=M_2^{9n},\quad T_2=P_2^5,
 \quad P=P_3=T_2^n,\quad c_2=c_G(\mathfrak g/2)^{k-1}.
\]
Also \(M_2^{-1}\le\mathfrak g\) and \(P\ge b^5\). The sign condition
holds at each real place of a rational term of nonzero discriminant.
Indeed the Gram matrix on \((F,\gamma F)\) has determinant
\((-1)^n\mathfrak g\) by swapping the two block rows and using the rank-one
Schur complement. The identity extends to singular \(Z\) by continuity.
When the pair Gram is nonsingular its negative index is \(n\): the
ambient negative index is \(n\), and inside the pair span the negative
vector \(z\) and the \(m\)-dimensional totally isotropic subspace in
\(z^\perp\) force at least \(n\).

More generally, let \(M\ge8n\), \(|Z_{ab}|\le M^2\),
\(\mathfrak g(Z)\ge M^{-1}\) and \(L_*\ge M^{9n}\). For
\(\|Z'-Z\|<L_*^{-6}\), the two matrix norms are at most \(M^3\).
By column multilinearity the determinant of a principal block of order
\(j\ge1\) changes by at most \(j M^{3(j-1)}\|Z'-Z\|\).
The total change in \(\mathfrak g\) is at most
\(4nM^{6n-3}L_*^{-6}\le M^{-1}/2\) (the order-zero minor is constant).
Hence \(\mathfrak g(Z')\ge\mathfrak g(Z)/2\).
The singular values and the last cofactor now show that
\[
 \|Z'\|,\quad\|(Z')^{-1}\|,\quad
 \bigl(1-|((Z')^{-1})_{nn}|^2\bigr)^{-1/2}
 \ \le\ L_* .
\]
Indeed \(|\det Z'|\ge(2M)^{-1/2}\), which bounds the inverse norm
by \(\sqrt{2M}\,M^{3n-3}\). The last quantity is
\(|\det Z'|/\sqrt{\mathfrak g(Z')}\), at most \(\sqrt{2M}\,M^{3n}\).
The same inverse estimate applies without a neighborhood whenever the
entries and the discriminant of the tested matrix satisfy the original
bounds with \(M\).\footnote{\weightedpin{grh-compact/src/New/Filling/BesselSystem.lean}{28-L45}{Size, frame box and mass},
 and \weightedpin{grh-compact/src/New/RealPlace/CellNeighbourhood.lean}{135-L163}{The pinned perturbation bounds}.}

Use the \hyperlink{proof.weighted-real-gram-section}{explicit
Gram section}, with \(E=p,\ D=p_0,\ Q=Z^{-1},\ \Delta=\Lambda\),
\(A=T_Z,\ A'=T_{Z^*}\), and \(S_Z=X_Z=\gamma_ZF\). Its column
identities give
\[
 \gamma_ZF'=(FQ^*-S_ZA')\Delta^{-1},\qquad
 \gamma_Z^{-1}F(Z^*)^{-1}=FA'+F'\Delta .
\]
We use these identities to bound the section and its derivative.

On a region where the three displayed bounds hold with \(L_*=P\ge8\),
we have \(\|A\|,\|A'\|\le2P^2\), \(\|S_Z\|\le3P^3\), and
\(\|\gamma_Z\|\le 11P^6\le P^{12}\). For a real matrix coordinate
direction, \(dQ=-Q(dZ)Q\) has norm at most \(P^2\). Differentiation
gives \(\|dA\|,\|dA'\|\le2P^3\), \(|ds|\le P^3\),
\(|d(s^{-1})|\le P^5\), and \(\|dS_Z\|\le6P^4\).
In the formula for \(\gamma_ZF'\), the derivative before differentiating
\(\Delta^{-1}\) has norm at most \((P^2+18P^6)P\), and differentiation
of \(\Delta^{-1}\) costs at most \((P+6P^5)P^5\).
Thus the derivative of \(\gamma_Z\) itself is at most
\(32P^{10}\le4P^{18}\) in each real coordinate.\footnote{\weightedpin{grh-compact/src/New/RealPlace/GramSectionBounds.lean}{16-L61}{Section formulas and the last cofactor},
 \weightedpin{grh-compact/src/New/RealPlace/GramSectionBounds.lean}{30-L49}{The section matrices},
 \weightedpin{grh-compact/src/New/RealPlace/GramSectionBounds.lean}{89-L114}{The derivative bounds},
 and \weightedpin{grh-compact/src/New/RealPlace/GramSection.lean}{17-L61}{The section as an isometry}.}

The pair \((F,\gamma F)\) has the same Gram as \((F,\gamma_ZF)\).
Their Gram is nonsingular and has signature \((n,n)\). The isometry
matching the two pairs extends by a unitary map between the positive
orthogonal complements. It thus gives \(s_1\in U^0H\) with
\(s_1\gamma_ZF=\gamma F\), and \(s_2=\gamma^{-1}s_1\gamma_Z\)
also fixes the frame. Hence \(\gamma=s_1\gamma_Zs_2^{-1}\).
Choose \(h,v\) independently with laws \(dh\) and
\(\left.du_0\right|_\Omega/\operatorname{vol}_{U^0}(\Omega)\),
where \(\Omega\) is the tube box for \(T=T_2\).
Set \(t=s_2(vh)^{-1}\) and \(A_2=A_{T_2}\).

For these choices, suppose \(y\in\operatorname{Box}(\rho)\) and
\(\operatorname{size}(\gamma t y)\le\rho\). Put \(y_0=(vh)^{-1}y\).
Then \(y_0^{-1}F=y^{-1}F\), and
\(\operatorname{size}(\gamma_Zy_0)=\operatorname{size}(\gamma t y)\)
by the left \(UH\)-invariance of size. By its definition, some
\(u\in U\) satisfies \(\|u\gamma_Zy_0\|\le\rho\): approximate the
infimum of \(\|s\gamma_Zy_0\|+1\) to get \(\|s\gamma_Zy_0\|<\rho\),
then remove the Euclidean-unitary left \(H\)-factor. If
\(u^{-1}F=F C'\), then \(Z'=F^*J u\gamma_Z F=(C')^*Z\).
Since \(C'\) is upper unitriangular, left multiplication by \((C')^*\)
preserves every leading principal determinant. By the two frame bounds,
\(|Z'_{ab}|\le M_2^2\), and therefore \(\|(Z')^{-1}\|\le P_2\).
We have
\[
 y_0^{-1}\gamma_Z^{-1}F(Z^*)^{-1}
   =(u\gamma_Zy_0)^{-1}F((Z')^*)^{-1}.
\]
The right side has norm at most \(M_2P_2\) because \(F\) is an orthonormal
column inclusion. By the inverse section formula and \(\|A'\|\le2P_2^2\),
each \(\|y_0^{-1}e_a^*\|\) is at most
\(M_2P_2+2\sqrt n M_2P_2^2\le P_2^3\).
In the negative eigenspace of \(J\) the orthonormal basis
\((e_a-e_a^*)/\sqrt2,z\) gives
\(\|y_0^{-1}|_{J=-1}\|\le \sqrt{2n}P_2^3\).
For any \(J\)-unitary matrix \(X=\left(\begin{smallmatrix}a_1&b_1\\c_1&d_1\end{smallmatrix}\right)\)
in the positive-negative splitting, \(XJX^*=J\) gives
\(a_1a_1^*=1+b_1b_1^*\) and \(c_1c_1^*=d_1d_1^*-1\).
Thus if the negative restriction has norm at most \(L\ge1\), its full
norm is at most \(\sqrt{1+L^2}+3L\le5L\).
It follows that
\(\|v^{-1}y\|=\|hy_0\|=\|y_0^{-1}\|\le5\sqrt{2n}P_2^3\le T_2\).
In particular \(y\in A_2\).\footnote{\weightedpin{grh-compact/src/New/RealPlace/GramSection.lean}{64-L80}{Double-coset factorization},
 \weightedpin{grh-compact/src/New/RealPlace/FrameNorm.lean}{33-L46}{Removing the compact factor from the size constraint},
 and \weightedpin{grh-compact/src/New/RealPlace/OpenCellBound.lean}{76-L100}{The support inclusion}.}

\paragraph{Evaluation at the section and the full-vector contraction.}
Use the full normalized \(W\), \(d=|I_v|\), and \(\delta=\deg_v\).
The input from the Bessel estimates is
\(\int_{U\backslash G_1}\|W\|^2d\dot y=d\), with the unitary left
\(UH\)-law for \(W\) and its right derivative words \(W_w\).
For \(P\ge b^5\) their Schur estimate, in the same Haar, is
\[
 \int_{U\backslash G_1}\Big\|\int_{G_1}\psi(y)W_w(gy)dy\Big\|^2 d\dot g
 \le P^{2|w|}(d/\delta)\int_{G_1}|\psi|^2\,dy,\qquad |w|\le2n^2.
\]
For the geometric calculation use the
\hyperlink{proof.weighted-real-gram-measure}{Gram inequality}
with a single left compact variable outside all integrations.
The density on the open cell is
\(c_G\mathfrak g(Z')^{k-1}dZ'\,dh\,dv_0\,du_0\), where \(dv_0=du_0(v)\)
and \(dZ'\) is coordinate Lebesgue measure.

First note an unfolding bound. Let \(\Phi\ge0\) be Borel, with
\(\int_{U^0}\Phi(ug)du_0\le K\) for each \(g\). Suppose that whenever
\(\Phi(g)\) and \(\Phi(sg)\) are nonzero, with \(s\in U\), every
free entry of \(Q(s)\) has modulus at most \(R_*\ge1\). For any
nonnegative left-\(U\)-invariant \(w\) one has
\[
 \int_{G_1}\Phi(g)w(g)\,dg
    \le K(2R_*)^{2n^2}\int_{U\backslash G_1}w(g)\,d\dot g.
\]
On each orbit meeting the support, use a point in the support as origin
by right translation in its \(U\)-integral. In the decomposition
\(s=u_0\widetilde Q\) the \(U^0\)-fiber has integral at most \(K\) and
the upper entries range over discs of radius \(R_*\). Their product
volume is \((\pi R_*^2)^{n(n-1)/2}\le(2R_*)^{2n^2}\).
Apply this on the quotient using the stated disintegration.

Let \(\psi\) be integrable and supported in \(B(P^{20})\).
For \(v\in V_P=U^0\cap B(P^{20})\), \(h\in H\), evaluate the convolution
at \(g=\gamma_{Z'}(vh)^{-1}\). The cube
\(Z+[0,P^{-8}]^{2n^2}\) in the real and imaginary entries lies in
\(\mathcal Q=\{Z':\|Z'-Z\|<P^{-6}\}\), since its operator radius is at
most \(2nP^{-8}\). The discriminant throughout this region is at least
\(\mathfrak g/2\). In a cube direction \(X=g^{-1}\partial g\) has norm
at most \(4P^{70}\) from the section bounds and \(\|vh\|\le P^{20}\).
The derivative of \(W_w(gy)\) uses \(y^{-1}Xy\). For each term in
\((y^{-1}Xy)_{ab}\), absorb
\((y^{-1})_{ac}y_{db}/P^{40}\) in \(\psi\) and retain the coefficient
\(P^{40}X_{cd}\), of modulus at most \(P^{112}\). The absorbed factor
has modulus at most one on the support.

Include all such modified tests together with their corresponding
derivative words in the family of convolutions. A first derivative
of any family member is a sum of at most \(r^4\) members of one higher
word length, with the stated coefficient bound. Differentiation under
the integral follows from the continuous derivatives on the compact
sets in this calculation and the integrability of the test.
For a one-dimensional interval of length \(\varepsilon=P^{-8}\),
the fundamental theorem of calculus and Cauchy--Schwarz give
\[
 \|f(0)\|^2\le \tfrac{2}{\varepsilon}\int_0^\varepsilon\|f(x)\|^2 dx
                +2\varepsilon\int_0^\varepsilon\|f'(x)\|^2 dx .
\]
Integrate in \(h,v\) and the other current coordinates, and at each
step bound the resulting integrals by their maximum over the permitted
members. This costs at most
\(2P^8(1+r^8P^{224})\le P^{241}\) per coordinate, and uses only first
derivatives of family members and words of length at most \(2n^2\).

To bound each resulting cube integral, insert a nonnegative Borel
factor \(\beta(g)\) with \(\int_{U^0}\beta(ug)du_0=1\).
For instance take \(e^{-\|g\|^2}\) divided by its \(U^0\)-orbital
integral. This denominator is positive and finite, since
\(\|ug\|\ge\|u\|/\|g^{-1}\|\) and \(L(u)\le\|u\|\), with the polynomial
volume growth proved above. It is constant on the left \(U^0\)-orbit.
Multiply by \({\bf1}_{\{g:F^*JgF\in\mathcal Q\}}\) to define \(\Phi\).
The Gram and the squared norm of each convolution are left
\(U^0H\)-invariant. Thus the \(u\)-integral in the Gram lower bound
removes \(\beta\). Every chart integral for fixed \(h_0\) is at
least \(c_2\) times the corresponding cube integral over
\(H\times V_P\), so the same lower bound holds for its infimum and
therefore for the Haar integral.

For the reverse estimate on this Haar integral, the unfolding bound
applies with \(K=1\) and \(R_*=P^2\). Indeed, for \(s\in U\) and two
Gram matrices from \(\mathcal Q\),
\[
 F^*JgF=Q(s)^*F^*JsgF .
\]
The Gram norm and inverse norm bounds give \(\|Q(s)\|\le P^2\).
Hence unfolding and Schur bound the Haar integral by
\((2P^2)^{2n^2}P^{4n^2}(d/\delta)\int|\psi|^2\).
With the coordinate evaluation cost this is bounded by the power
\(P^{2000n^2}\).\footnote{\weightedpin{grh-compact/src/New/Filling/BesselSystem.lean}{52-L62}{The derivative Schur input},
 \weightedpin{grh-compact/src/New/Filling/BesselSystem.lean}{85-L109}{The full vector and its parameters},
 \weightedpin{grh-compact/src/New/RealPlace/SobolevCell.lean}{16-L54}{The two Haar inequalities},
 \weightedpin{grh-compact/src/New/RealPlace/SobolevCell.lean}{116-L132}{The cube calculation},
 \weightedpin{grh-compact/src/New/RealPlace/SobolevCell.lean}{36-L56}{Unfolding in the free upper entries},
 \weightedpin{grh-compact/src/Lib/General/Analysis/ZpowSum.lean}{66-L91}{Successive coordinate estimates}, and
 \weightedpin{grh-compact/src/New/RealPlace/SobolevCell.lean}{57-L134}{Derivative words and evaluation at the section}.} We obtain
\[
 c_2\int_H\int_{V_P}
 \Big\|\int\psi(y)W(\gamma_Z(vh)^{-1}y)dy\Big\|^2\,dv_0\,dh
 \le P^{2000n^2}(d/\delta)\int|\psi|^2\,dy .
\]

Now \(A_2\subset B(P^{20})\) and \(\Omega\subset V_P\).
For a bounded measurable \(\phi\), apply the section estimate to
\(\psi_j(y)={\bf1}_{A_2}(y)\phi(y)\overline{W_j(y)}\) and sum
over the \(d\) columns. The left Bessel actions of \(s_1\) and \(t\)
are unitary and independent of \(y\), so
\[
 \begin{aligned}
 \mathbf E_{h,v}
 \left\|\int_{A_2}\phi(y)W(\gamma t y)W(t y)^*dy\right\|_{\rm HS}^{2}
 &\le \int |\phi|^2\,d\nu_2,\\
 d\nu_2 &=
 \frac{P^{2000n^2}d}{\delta c_2\operatorname{vol}_{U^0}(\Omega)}
             {\bf1}_{A_2}(y)\|W(y)\|^2dy .
 \end{aligned}
\]
Use the same unfolding bound with \(\Phi={\bf1}_{A_2}\).
Now \(K=T_2\operatorname{vol}_{U^0}(\Omega)\) by the tube estimate,
and \(R_*=P^{40}\) since \(\|s\|\le\|sy\|\|y^{-1}\|\) for two points in
\(A_2\). Thus
\(\int_{A_2}\|W(y)\|^2dy\le
 T_2\operatorname{vol}_{U^0}(\Omega)(2P^{40})^{2n^2}d\).
The tube volume cancels. Using \(P=M_2^{45n^2}\) and \(M_2\ge2\) gives
\[
 \nu_2(G_1)
 \le M_2^{93600n^4+2n^2+45n}\frac{d^2}{\delta c_2}
 \le \left(\frac{\rho^{9n}}{\mathfrak g}\right)^{10^5n^4}
        \frac{|I_v|^2}{\deg_v c_G(\mathfrak g/2)^{k-1}} .
\]

In the classwise application below \(\gamma\) is read in its
normalized frame at every infinite place. The radius is
\(\rho=b^{2n^2}\), which is at least the analytic cutoff radius
\(b^{6n+2}\) since \(\varrho\ge 10^5\). Use the local tubes and
sampling laws above in the
\hyperlink{proof.weighted-large-discriminant-unfolding}{class
unfolding with a large discriminant}.\footnote{\weightedpin{grh-compact/src/New/RealPlace/OpenCellBound.lean}{13-L75}{Contraction and mass},
 \weightedpin{grh-compact/src/Lib/DiscreteSeries/Bruhat/BigCell.lean}{41-L67}{Product of contractions},
 \weightedpin{grh-compact/src/New/Filling/Setup.lean}{243-L255}{The actual partition and box radius},
 \weightedpin{grh-compact/src/New/Filling/SchurTest.lean}{69-L99}{Packet-constant comparison},
 \weightedpin{grh-compact/src/New/Filling/BlockBound.lean}{330-L361}{The class application},
 \weightedpin{grh-compact/src/New/RealPlace/OpenCellBound.lean}{13-L100}{The original contraction and support on the same data},
 \weightedpin{grh-compact/src/New/Filling/Moments.lean}{102-L120}{The fixing-set comparison},
 and \weightedpin{grh-compact/src/New/Filling/Moments.lean}{186-L218}{The complementary comparison}.}

\paragraph{The sizes of cut blocks.}
\teamanchor{proof.weighted-cut-block-sizes}%
At fixed $(c,s,c')$ choose a representative of each
$U_c^0\gamma U_{c'}^0$ class, where $U_c^0$ fixes the whole frame.
Only classes with both real arguments in the frame box contribute.
Write $j$ for the level, so $0\le j\le n$, and set
$d_j=\min(j,a)$. For the position $\kappa=h_c^{-1}\gamma h_{c'}$ in
the adelic lattice chart let $I_j(\kappa)$ be the index of the
$\widehat{\mathcal O}_E$-span of the adjoints of the first $d_j$ rows
of the frame pairing in the span of its first $a$ rows. Write
\[
 \begin{split}
 D_j(\gamma)&=|N_{E/\mathbf Q}\det Z_j(\gamma)|/N^{2f(j-a)_+},\\
 x(\gamma)&=\prod_{v\mid\infty}|\delta(\gamma)|_v,\\
 B_j(\gamma)&=\frac{\max(x(\gamma),N^f D_j(\gamma))}
                         {I_j(\kappa) I_j(\kappa^{-1})},\\
 M_0&=|J|\,|\mathrm{Cl}|\,N^{-kf},\\
 \mathcal D_j&=\begin{cases}
            \tau^{13-32j}\mathscr D&j\le a,\\
            \mathscr M&j=n.
          \end{cases}
 \end{split}
\]
Here $\delta$ is the frame discriminant and a leading minor of
size zero is one. The indices are finite and positive. For a stable
set of level $j$, the cut amplitude norm of its block is bounded by
the sum, over its retained near representatives in the block, of
$\nu u_j(\gamma)$, where
\begin{equation}
 \label{eq:weighted-class-size}
 u_j(\gamma)=\tau^2\left(
 \frac{\mathcal D_j}{M_0 g^{2fd_j k}B_j(\gamma)^{k+1}}
 \right)^{1/2}.
\end{equation}
Only compatible residues are included in this sum. This means that
the reduced unipotent phase takes equal values at $\bar u,\bar u'$
whenever $\bar u s=s\bar u'$. The two-sided covariance multiplies
the block by the ratio of these values and gives zero when they differ.
This is a condition on the residue in addition to the character test
on a rational stabilizer.\footnote{\weightedpin
{grh-compact/src/New/Filling/AmplitudeSums.lean}{180-L194}
{Two-sided covariance and vanishing for incompatible residues}.}

\paragraph{Integral frame translations within a class.}
\teamanchor{proof.weighted-thin-class-lattice}%
Use \(Z=Z(\gamma)=(h(f_i,\gamma f_h))_{i,h}\) in the rational frame, and
let \(Y\) have the first \(j\) rows of \(Z\) followed by the corresponding
rows of the identity. It is invertible since its determinant is the
leading minor of order \(j\). For \(h\ge j\), the shared flag vectors in
the level condition can be combined so that the coefficients on the
moved frame from index \(j\) onward are exactly those of \(e_h\) in the
coordinate space of the frame. Indeed these trailing coefficient
vectors are unitriangular. The resulting vector belongs to the
original frame span, so its first \(j\) pairings vanish, since these
first \(j\) frame vectors are isotropic. Its moved-frame coefficients
are therefore \(Y^{-1}e_h\). In particular \(\gamma\mathbf f Y^{-1}e_h\)
has zero positive middle and dual-isotropic components. If \(j=n\)
and \(h=a\), the dual component still vanishes and the \(z\) component
is \(-1/N\), by the equation \(Z Y^{-1}e_a=e_a\). Its positive middle
component has squared length at most \(1/N\) at every embedding.
In fact, preservation of the Hermitian square gives that length
as \(1/N-N|\Phi_v((Y^{-1}e_a)_a)|^2\).

Let \(C_{c,0}\) be the finite coordinate map from \(X_0\) into the
middle block of the class lattice of \(c\). The intersection lattice
\(L_c=\{x\in X_0(E):C_{c,0}x\in\widehat{\mathcal O}_E^{k+1}\}\)
has integral Hermitian squares because the full class lattice is
self-dual. A vector of \(L_c\) whose squared length is strictly less
than one at every infinite place must be zero by the norm of its
Hermitian square.

Here is the translation calculation in a form also used for the
frame-value count. Let \(p,p_0\) be terms of level \(j\) with the same
\(Z\) and source class \(c\), and let \(D_{\rm mid}\) be the positive middle
block of \((p-p_0)\mathbf f\). Choose isotropic row indices \(J_0\)
containing all indices less than \(d_j=\min(a,j)\), and write
\(\xi=Z_{J_0,:}\). Suppose that in the finite chart
\(C_{c,0}D_{\rm mid}=g(A_0\xi+B_0Y)\) with \(A_0,B_0\) integral. In each finite
valuation ring the intersection of the image of \(C_{c,0}\) with the
integral middle module is saturated, hence admits an integral
projection onto it. Apply this projection to the two coefficient
matrices. By density of \(E\) in its finite adeles, choose a rational
matrix \(L_0\) close enough to the preimage of the projected \(A_0\)
that \(L_0\) and all columns of
\(Q=g^{-1}(D_{\rm mid}-gL_0\xi)Y^{-1}\) belong to \(L_c\). The two integrality
conditions define an open neighbourhood of that preimage.
For \(h\ge j\) the column of \(Q\) vanishes, since both shared vectors
have zero positive middle block and \(\xi Y^{-1}e_h=0\) by their
dual-isotropic blocks. The omitted index below \(j\) can only be
\(h=a\) when \(j=n\), and then \(\xi Y^{-1}e_a=0\) as well. Its column
is integral with squared length at most
\(4/(g^2N)<1\) at every place, by the two bounds just proved, and is
zero as well. Since \(Y_{h,:}=\xi_{h,:}\) for the selected rows \(h<j\),
we obtain \(D_{\rm mid}=gL\xi\) with columns of \(L\) in \(L_c\). All the other
rows below the isotropic block in \((p-p_0)\mathbf f\) vanish since
the frame pairings give its dual-isotropic and marked rows.

For the thin class \(\mathcal C=\Gamma_c^0\gamma\Gamma_{c'}^0\), both
unipotent factors fix the frame, so every term has the same level,
\(Z\) and residue. The right factor fixes the moved columns before
applying the left factor. In the finite class chart, the change in
each middle row of a moved frame is a linear combination of the
dual-isotropic rows of the original moved frame, with coefficients
in \(g\widehat{\mathcal O}_E\). These are the rows of the isotropic
frame pairing. Let \(R_a\) and \(R_j\) be the finite spans of the
adjoints of the first \(a\) and \(d_j\) rows of \(Z\), respectively.
They have rank \(d_j\). For \(j\le a\) the shared vectors give zero
columns from \(j\) onward after multiplying the isotropic rows by
\(Y^{-1}\). At \(j=n\) this rank follows from the invertible full
pairing. In either case the first \(d_j\) rows are independent.
The quotient \(gR_a/gR_j\) has order \(I_j(\kappa)\). Colour the
terms by the differences of their \(k+1\) adjoint middle rows in
this quotient, giving at most \(I_j(\kappa)^{k+1}\) colours. In one
colour choose a base term \(p_0\). The translation calculation with
\(J_0=\{0,\ldots,d_j-1\}\) and \(B_0=0\) gives
\(D_{\rm mid}=gL(p)Z_{J_0,:}\). For equal \(L(p)\) the moved frames therefore
differ only in the isotropic rows.\footnote{\weightedpin
{grh-compact/src/New/Filling/BoxPartition.lean}{413-L423}
{Shared flag vectors}, \weightedpin
{grh-compact/src/New/RealPlace/VolumeKernel.lean}{55-L122}
{The definite-coordinate lift}, and \weightedpin
{grh-compact/src/New/Filling/ThinClasses.lean}{21-L39}
{The actual double-coset congruence}.}

\paragraph{The real volume integral.}
\teamanchor{proof.weighted-thin-volume-kernel}%
At a fixed real place write \(J_1\) for the Hermitian involution
defining the normalized form, and use \(\mathbf f\) for the normalized
frame matrix. The columns of \(F_-=(1-J_1)\mathbf f\) span the maximal
negative space. Put \(\Delta(y)=|\det(F_-^*y^{-1}\mathbf f)|\).
It is positive. A vector in \(y^{-1}\operatorname{span}(\mathbf f)\)
has nonpositive Hermitian square, whereas a nonzero vector whose
projection to the negative space is zero has positive square.
Moreover \(\Delta\) is left invariant under \(U^0_v\), which fixes
\(\mathbf f\).

Fix \(y\) of size at most \(\theta\), and let \(\zeta\) be an
isometry with the given normalized pairing matrix \(Z_v\) and level
\(j\) normal form. In the following calculation \(\iota\) includes
the first \(j\) frame coordinates, \(F_j=\mathbf f\iota\), and
\(z_j=\iota^*Z_v\iota\). Define
\[
 B_0=(F_-^*y^{-1}\mathbf f)^{-1}F_-^*y^{-1},\qquad
 t=\zeta F_j z_j^{-1},\qquad p=Z_v\iota z_j^{-1}.
\]
Then \(B_0\mathbf f=1\), \(p=\mathbf f^*J_1t\) and \(\iota^*p=1\).
Write \(M=B_0t z_j-\iota\). For the first \(j\) standard columns,
\(B_0\zeta\mathbf f\) and \(1+Mp^*\) have equal values.
For every remaining index, the shared flag vector takes the form
\(\mathbf f x=\zeta\mathbf f x'\) with both coefficient vectors
supported through that index and last coefficients of modulus one.
Here \(p^*x=t^*J_1\zeta\mathbf f x'=z_j^{-*}F_j^*J_1\mathbf f x'=0\),
because the first \(j\) vectors are isotropic when there are
remaining indices. Thus \(B_0\zeta\mathbf f x'=(1+Mp^*)x=x\).
The two triangular change-of-columns matrices have determinant of
modulus one. The identity \(\det(1+Mp^*)=\det(1+p^*M)\) now gives
\[
 \frac{\Delta(\zeta^{-1}y)}{\Delta(y)}
    =|\det B_0\zeta\mathbf f|=|\det z_j|\,|\det(p^*B_0t)|.
\]
At \(j=0\) this is one, with empty determinants.

Choose \(u\in U_v\) with \(\|uy\|_{\rm op}\le\theta\). Let \(T\) be
\(\theta\) times the first \(j\) square block of its frame matrix,
and set
\[
 C=T p^*B_0,\qquad
 D=Cy,\qquad
 Y_1=J_1tT^*-C^*,\qquad
 W_1=\operatorname{pr}_{X_0}Y_1,\qquad
 \Theta=-CtT^*.
\]
Since \(\mathbf f^*Y_1=0\), we have \(Y_1^*J_1Y_1=W_1^*W_1\).
The rows of \(D\) lie in the negative subspace, so
\(C J_1 C^*=D J_1D^*=-DD^*\). Also
\(Q=\theta^{-1}(uy)^{-1}F_j\) is a contraction and
\(QT=y^{-1}F_j\), whence \(DQ=1\). Thus \(DD^*\succeq1\).
Expansion gives
\[
 \Theta+\Theta^*=W_1^*W_1+DD^*-Tt^*J_1tT^*\succeq 1+W_1^*W_1,
\]
as \(t^*J_1t=z_j^{-*}F_j^*J_1F_j z_j^{-1}\preceq0\).
If \(H=(\Theta+\Theta^*)/2\), conjugating \(\Theta\) by \(H^{-1/2}\)
gives the identity plus \(i\) times a Hermitian matrix. Its determinant
has modulus at least one, so \(|\det\Theta|\ge\det H\).
For an inclusion \(e:\mathbf C^d\to\mathbf C^j\) of coordinate columns,
\(W_1W_1^*\succeq (W_1e)(W_1e)^*\). The preceding determinant identity
and \(|\det T|=\theta^j\) give
\[
 \frac{\Delta(y)}{\Delta(\zeta^{-1}y)}
 \le \frac{(2\theta^2)^j}
 {|\det z_j|\det(1+(W_1e)^*(W_1e))}.
\]

Suppose the specified positive middle block of \(\zeta F_j\) has the
form \(V_0+A R e^*z_j\), where \(V_0\) is fixed and \(R\) is an
invertible \(d\)-by-\(d\) matrix. Then \(W_1e=A R(e^*T^*e)+C_0\) for a
fixed \(C_0\), since \(J_1\) is the identity on \(X_0\). The affine map
on each row has real Jacobian \(|\det R|^2\theta^{2d}\).
To integrate the determinant to the power \(-(k+n)\), integrate one
row \(w\) at a time. For the positive matrix \(H_0\) from the remaining
rows, \(\det(H_0+w^*w)=\det H_0(1+wH_0^{-1}w^*)\). Changing variables
by \(H_0^{1/2}\) lowers the exponent of \(\det H_0\) by one, with
factor at most \(\pi^d/d!\), since for an integer \(q\ge d+1\)
\[
 \int_{\mathbf C^d}(1+\|w\|^2)^{-q}\,dw
     =\frac{\pi^d}{(q-1)\cdots(q-d)}\le\frac{\pi^d}{d!}.
\]
Here the equality is the radial beta integral, with the empty product
at \(d=0\). The successive exponents are \(k+n,\ldots,n+1\).
We obtain a measurable majorant \(K_v(A)\) for
\((\Delta(y)/\Delta(\zeta^{-1}y))^{k+n}\) with
\begin{equation}
 \label{eq:weighted-thin-local-volume}
 \int K_v(A)\,dA\le
 \left(\frac{(2\theta^2)^j}{|\det z_j|}\right)^{k+n}
 \left(\frac{\pi^d}{d!\theta^{2d}|\det R|^2}\right)^k .
\end{equation}
The measure in each complex coordinate is ordinary Lebesgue
measure, with unit disk area \(\pi\). The case \(j=0\) uses the
unit mass of the zero-dimensional parameter space.\footnote{\weightedpin
{grh-compact/src/New/RealPlace/VolumeKernel.lean}{129-L235}
{The volume determinant and the compression to the selected columns}, and
\weightedpin{grh-compact/src/Lib/General/LinearAlgebra/HuaIntegral.lean}
{187-L209}{Row integration}.}

\paragraph{The one-sided Schur estimate.}
\teamanchor{proof.weighted-thin-schur}%
For cut vectors put \(\beta_c^0={\bf1}_{\mathrm{Box}}\beta_c\) and
\(V(y,z)=\prod_v(\Delta(y_v)/\Delta(z_v))^{k+n}\).
For the class \(\mathcal C=\Gamma_c^0\gamma\Gamma_{c'}^0\) define
\[
 F(y)=\sum_{p\in\mathcal C}\beta_{c'}^0(p^{-1}y)V(y,p^{-1}y),\qquad
 \sum_{p\in\mathcal C} S(c,c',p)
    =\int\beta_c^0(y)\|W^\flat(y)\|^2 F(y)\,dy.
\]
Here \(S(c,c',p)\) denotes the integral with the single summand of
\(F\). All rational arguments in real integrals mean their normalized
real images. The function \(F\) is invariant under left
\(\Gamma_c^0\), by reindexing \(p\) and the invariance of \(\Delta\).
In the last integral use \(\beta_c^0={\bf1}_{\mathrm{Box}}p_c\omega_c\)
and disintegrate along \(U_\infty^0\). At a section point \(y\), all
factors except \(p_c\) and \(F\) are constant on the fibre.
The cutoff \(p_c\) unfolds that fibre to
\(\Gamma_c^0\backslash U_\infty^0\). Right translation by any
\(s\in U_\infty^0\) preserves its quotient Haar, since a unipotent
group is unimodular. Consequently, if for every boxed \(y\) and every
\(u\in U_\infty^0\) one has \(\int_{\mathcal X} F(u s(x)y)\,dx\le M_*\),
where \(s(x)\in U_\infty^0\) and \(\operatorname{vol}\mathcal X\ge1\),
Tonelli bounds the last displayed integral by
\(M_*\int\beta_c(y)\|W^\flat(y)\|^2\,dy\le M_*\nu d_W\).
This uses the full class partition on \(U_\infty\) and the cut squared
norms \(\int_{U_v\backslash G_v}\|W_v^\flat\|^2\le |I_v|\).

Take \(d=d_j\), \(\mathcal X=\prod_v\{x_v\in\mathbf C^{k\times d}:
\max|x_{v,ih}|<1\}\), of volume \(\pi^{kdf}\). In the \(v\)-factor put
\[
 R_v=\operatorname{diag}_{h<d}\frac{\pi g}{b s_{vh}}.
\]
Let \(s_v(x_v)\) fix \(z\) and have blocks on \(X\oplus X_0\oplus X^*\)
\[
 \begin{pmatrix}1&Q_v&-Q_vQ_v^*/2\\0&1&-Q_v^*\\0&0&1\end{pmatrix},
 \qquad Q_v=\iota_{a,d}(x_vR_v)^*,
\]
where \(\iota_{a,d}\) includes the first \(d\) indices of the isotropic
block. The matrix preserves the split form and fixes the frame.
Its inverse adds \(Q_v^*\) times the dual-isotropic block to the
positive middle block. For a colour with base term \(p_0\), the
isometry \(s_v(x_v)^{-1}u_v^{-1}p_v\) has pairing \(Z_v\) and the
same shared flag relations. On its first \(j\) frame columns the
positive middle block is \(V_0+(x_v+\lambda_v(p))R_v e^*z_j\), with
\[
 V_0=\operatorname{pr}_{X_0}(u_v^{-1}p_{0,v}F_j),\qquad
 \lambda_v(p)=(b/\pi)\Phi_v(L(p)),\qquad e=\iota_{j,d}.
\]
Indeed the rational translation \(gL(p)Z_{J_0,:}\) is rescaled on its
frame columns by \(s_{vh}\), and the corresponding row of \(Z_v\)
has the additional factor \(s_{vi}\) at row \(i\). Hence its coefficient
after rescaling is \(g\Phi_v(L(p))\operatorname{diag}(s_{vi}^{-1})\).
This is \(\lambda_v(p)R_v\). Also \(u_v^{-1}\) does not change the
positive middle block of the difference, whose dual-isotropic
coordinates are zero. The volume ratio at \(u s(x)y\) is therefore
bounded by \(\prod_v K_v(x_v+\lambda_v(p))\) with the same kernels
throughout the colour.

Translate the integration variable for the summand indexed by \(p\)
by \(\lambda(p)\). At a fixed translated argument, at most one
matrix \(L(p)\) can occur with a nonzero boxed weight. In fact an
intersection of two translated unit boxes would imply that every
coordinate of \(\Phi_v(L(p)-L(q))\) has modulus less than \(2\pi/b\).
In each column the squared length would be less than
\(64k/b^2\le 64/b<1\) at every place. The lattice norm argument above
then gives equality of the matrices. For that fixed matrix all
nonzero summands have the same original \(x\), and their boxed
inverse frame tuples differ in \(\prod_v(y_v^{-1}X_v)^n\), of real
dimension at most \(2anf\). This follows from the isotropic row
difference and the fact that \(u s(x)\) fixes \(X\) pointwise.
The \hyperlink{proof.weighted-cut-frame-packing}{weighted frame count}
bounds the sum of their \(\beta_{c'}^0\) weights by
\(b^{12an^3 f}\le\tau\). Equal frames give relative elements in
\(\Gamma_{c'}^0\) by class freeness at the same residue, so this count
uses the partition weights on one class lattice. Now integrate the
common product kernel and sum over the colours. The required estimate
for \(M_*\) is at most \(I_j(\kappa)^{k+1}\tau\prod_v\int K_v\).

Write \(\epsilon=j-d\), \(D_v=|\Phi_v(\det Z_j(\gamma))|\),
\(P_0=\prod_v D_v\), and \(w_v=\prod_{h<d}|\kappa_v|^{a-h}\).
We have \(\epsilon\in\{0,1\}\), \(w_v\ge1\) and \(P_0\ge1\)
because the nonzero rational leading minor is an algebraic integer.
The normalized local block satisfies
\[
 |\det z_j|=w_v^2N^{-j}D_v,
 \qquad |\det R_v|=(\pi g/b)^d N^{d/2}/w_v,
 \qquad \sqrt{D_j(\gamma)}=P_0/N^{\epsilon f}.
\]
Thus \eqref{eq:weighted-thin-local-volume} is a product of the factor
\(((2\theta^2)^j/|\det z_j|)^n\) and the \(k\)-th power of
\[
 \frac{2^j\theta^{2\epsilon}}{\pi^d d!}
       \frac{b^{2d}N^\epsilon}{g^{2d}D_v}.
\]
The first fraction is at most \((g/b^6)^d\). If \(\epsilon=0\), use
\(2\le\pi\) and \(g\ge b^{19}\). If \(\epsilon=1\), then \(d=a\),
and \(2\theta^2\le b^{12a+17}\le b^{13a}\le(g/b^6)^a\).
The \(n\)-th power factors have product at most
\[
 \frac{N^{\epsilon f}}{P_0}
 \,(2\theta^2)^{jnf}N^{(jn-\epsilon)f}
 \le \frac{\tau}{\sqrt{D_j(\gamma)}}.
\]
For the last inequality the leftover exponent of \(b\) is at most
\((72n+5)n^2f\) by \(N\le b^{60n}\) and \(j\le n\), below
\(kf/4\) since \(k\ge1000(2n)^9\). Together with packing and
disintegration this proves
\begin{equation}
 \label{eq:weighted-thin-schur}
 \sum_{p\in\mathcal C}S(c,c',p)
 \le \nu d_W\tau^2(b^4g)^{-d_j k f}
      \left(\frac{I_j(\kappa)}{\sqrt{D_j(\gamma)}}\right)^{k+1} .
\end{equation}
All the estimates also hold for the zero level with the empty
minor and the zero-dimensional shift.\footnote{\weightedpin
{grh-compact/src/New/Filling/SchurTest.lean}{101-L169}
{The reciprocal test and quotient averaging}, \weightedpin
{grh-compact/src/New/Filling/SchurTest.lean}{101-L150}
{The boxed volume weight}, \weightedpin
{grh-compact/src/New/Filling/ThinClasses.lean}{43-L123}
{The class integral and scale factors}, and \weightedpin
{grh-compact/src/New/Filling/BoxPartition.lean}{510-L556}
{The weighted count on the same frame fibres}.}

\paragraph{The cut amplitude in the small discriminant range.}
Both arguments of a nonzero cut integrand with nonzero cutoff
belong to \(\mathrm{Box}\). For a single term \(p\), insert the
reciprocal weights \(V(y,p^{-1}y)\) and \(V(p^{-1}y,y)\) in
Cauchy--Schwarz and change variables by \(p^{-1}\) in the second
integral. This bounds the square of its absolute cut amplitude by
\(S(c,c',p)S(c',c,p^{-1})\). Apply Cauchy--Schwarz to the sum over
\(p\), and then use \eqref{eq:weighted-thin-schur} for the two
orientations. Inversion sends \(Z\) to \(Z^*\), so it preserves
\(j,D_j\), and the finite position becomes \(\kappa^{-1}\).
If \(x(\gamma)\le N^fD_j(\gamma)\), divide the resulting squared
bound by \(\nu^2u_j(\gamma)^2\) from \eqref{eq:weighted-class-size}.
With \(M_0=|J|\,|\mathrm{Cl}|\,N^{-kf}\) and
\(\mathscr D=|J|\,|\mathrm{Cl}|d_W\), the ratio is at most
\[
 \frac{d_W\mathscr D\,N^f b^{-8d_j k f}}{\mathcal D_j}\le1.
\]
Use \(d_W\le\tau^8\), \(N^f\le\tau\), and
\(b^{-8d_j k f}=\tau^{-32d_j}\). For \(j\le a\) take
\(\mathcal D_j=\tau^{13-32j}\mathscr D\), and at \(j=n\) take
\(\mathcal D_n=\mathscr M\ge\tau^{44}\mathscr D\).
If the cut absolute amplitude in this range is nonzero at a term,
there is a point witnessing that the term is near. Nearness is
preserved by both frame-fixing class-unipotent factors, so only
near representatives enter the bound.\footnote{\weightedpin
{grh-compact/src/New/Filling/ThinClasses.lean}{127-L149}
{The small-discriminant application}.}

\paragraph{The real cell mass on a top-level class.}
\leavevmode\teamanchor{proof.weighted-cell-mass-budget}%
Let \(\gamma\) belong to the large range \(x(\gamma)>N^fD_j(\gamma)\)
of the fixed block. Its rational frame discriminant \(\delta\)
is nonzero. Put
\(a=\varrho,\ n=a+1,\ B=b^{2n^2},\ s_0=10^5n^4\).
Write \(x=\prod_v|\Phi_v(\delta)|\) and
\(L_v=\prod_{i=0}^a\ell_{vi}\), \(L=\prod_vL_v\), where
\(\ell_{vi}=|\kappa_v|^{2(a-i)}/N\).
The two frame-row indices at level \(n\) are one. Put
\(u=u_n(\gamma)\) from \eqref{eq:weighted-class-size} in the large
range \(B_n(\gamma)=x\).

To see the sign of the discriminant, multiply the rational Gram
matrix \(Z\) by \(y=\operatorname{adj}(Z)e_n\), where \(e_n\) is
the last coordinate vector. This gives \((\det Z)e_n\), and the
last entry of \(y\) is the leading minor \(D_a\).
In the isometric image of the frame combination \(Fy\), all dual
isotropic coordinates therefore vanish, and its terminal coordinate
is \(-\det Z/N\). Equality of its norms in the two frames gives
\(\Phi_v(\delta)=N\|(\gamma Fy)_{X_0,v}\|^2\).
The frame products on the retained finite coset are integral.
Thus \(x\) is the absolute norm of a nonzero algebraic integer
from the totally real field, and \(x\ge1\). Scaling each row and
column of \(Z\) to the real frame gives
\(\mathfrak g_v=L_v^2|\Phi_v(\delta)|\). In this calculation the
last scale satisfies \(\ell_{va}=1/N\), which accounts for the
\(N^2\) on the rational minor \(D_a\).

For the compact-cell estimates take \(\gamma\) to be a retained near
representative. Let \(M_v\) denote the upper bound for the mass in the
\hyperlink{proof.weighted-real-cell-contraction}{local contraction}
at \(\gamma_v\), formal degree \(\deg_v\), and
\(d_v=|I_v|\). It is
\[
 M_v=\left(\frac{B^{9n}}{\mathfrak g_v}\right)^{s_0}
       \frac{d_v^2}{\deg_v c_G(\mathfrak g_v/2)^{k-1}} .
\]
The product of the \(d_v\) is \(d_B\). Use the
\hyperlink{proof.weighted-open-mass}{packet normalization inequality}
for \(\Pi,V,F_{\rm sc},\nu,\mathscr M\).
Substitution of \(\prod_v\mathfrak g_v=L^2x\) gives
\[
 \frac{4\prod_vM_v}{\nu^2u^2}
 \le \frac{4(1+b^{-9})}{\tau^4}\,
 x^{2-s_0} B^{9nf s_0}\,
 \frac{F_{\rm sc}^{\,2}}{L^{2(k-1+s_0)}}\,
 \frac{2^{(k-1)f}M_0g^{2fak}}{\Pi V^n}.
\]
Here \(M_0=|J|hN^{-kf}\) and
\(\Pi=|J|g^{fa(2k+2a+1)}h\,c_*\).
The last fraction is
\[
  g^{-fa(2a+1)}
  \left(\frac{(k-1)!\operatorname{rd}_E^{\,k+n-1/2}}
                   {2^{n+1}\pi^k}\right)^f .
\]
It is at most \(b^{(k+n)f/1000}\), by
\((200fk\operatorname{rd}_E)^{1000}\le b\).
The quotient containing \(F_{\rm sc}\) is the product of
\(\ell_{vi}^{\,2(2(a-i)+1-s_0)}\). All these exponents are
negative and \(\ell_{vi}\ge N^{-1}\), so the quotient is at most
\(N^{2nf s_0}\). By \(N\le b^{60n}\) and \(x\ge1\), the exponent
of \(b^f\) contributed by these factors and \(B^{9nf s_0}\)
is at most \((120n^2+18n^3)s_0\). This is at most \(k/2\)
for \(n\ge10\) and \(k\ge1000(2n)^9\), as on the chosen
scales. Since \(\tau^4=b^{kf}\), the displayed ratio is at most one.

For the cutoff on the same representative, the two boxed frames
give real pairings of modulus at most \(B^2\) by Cauchy--Schwarz.
Passing to the rational frame multiplies an entry by at most
\(N\), since \(|\kappa_v|\ge1\). Thus each entry is at most
\(NB^2\le b^{5n^2}\) on the chosen scales.
The determinant expansion, \(n\le b\), and
\(0\le\Phi_v(\delta)\le|\Phi_v(\det Z)|^2\) give
\(x\le b^{11n^3f}\). Using
\(\mathscr M\ge\tau^{44}|J|h d_B\) in the same class size now yields
\[
  u^2\ge
   \frac{\tau^{48}N^{kf}d_B}{g^{2fak}x^{k+1}}
   \ge b^{-2k^3f}.
\]
Indeed \(g\le b^{22}\) and \(44ak+11n^3(k+1)\le2k^3\).
Consequently \(\sqrt{\prod_v M_v}\le \nu u/2\) and
\(\nu b^{-k^3f}/2\le\nu u/2\) for the two estimates on the
same class.\footnote{%
\weightedpin{grh-compact/src/New/Filling/NearTerms.lean}{215-L236}
{The rational and real discriminants},
\weightedpin{grh-compact/src/New/Filling/NearTerms.lean}{294-L305}
{The near frame and its leading minors},
\weightedpin{grh-compact/src/New/Filling/SchurTest.lean}{63-L99}
{The two class allowances on the same near representative}.}

\paragraph{The cut amplitude in the large discriminant range.}
\teamanchor{proof.weighted-large-discriminant-unfolding}%
For \(x(\gamma)>N^f D_j(\gamma)\) one has \(j=n\), since the
discriminant vanishes at other levels. Both row indices at \(n\)
are one, and \(B_n(\gamma)=x(\gamma)\). The discriminant computation
in the \hyperlink{proof.weighted-cell-mass-budget}{cell mass comparison}
shows that every normalized local discriminant \(\mathfrak g_v\)
is positive and \(x(\gamma)\ge1\). The sign and integrality arguments
there apply to all terms in the range, before using nearness.
In particular the full normalized pairing and hence \(Z\) are
invertible.

Fix a stable set of terms in this range and a block \((c,s,c')\).
Stability holds for this restriction since the two leading
determinants and the discriminant are unchanged by the upper
unitriangular frame actions on the left and on the right. Choose
one representative \(\rho\) of each double class for
\(\Gamma_c\) on the left and \(\Gamma_{c'}^0\) on the right. Choose
it among the fixed representatives for the smaller left group
\(\Gamma_c^0\), preferring a near one whenever it exists.
The larger left group acts freely on these right cosets. To see
this, if \(u\in U\) fixes \(\gamma\mathbf f\), it fixes the first
\(a\) columns of \(\gamma\mathbf f Z^{-1}\), whose dual-isotropic
block is the identity. In \(X\oplus (X_0\oplus Ez)\oplus X^*\)
an element of \(U\) is upper block triangular, acts trivially on
\(X^\perp/X\), and has contragredient diagonal blocks on \(X,X^*\).
Fixing these columns forces both those diagonal blocks to be one
and the block from \(X^*\) to the middle to vanish. Preservation of
the form then makes the middle-to-\(X\) block vanish. Fixing the
columns now also makes the \(X^*\)-to-\(X\) block vanish. This proves
\(\Gamma_c\cap\gamma\Gamma_{c'}^0\gamma^{-1}=\{1\}\).

On the pair \((\gamma,y)\) let \((u,v)\in\Gamma_c\times\Gamma_{c'}\)
act by \((u\gamma v^{-1},uy)\). Counting measure in the terms and
real Haar in \(y\) are preserved. The weight
\[
 \phi(\gamma,y)={\bf1}_{\gamma\text{ chosen}}\,
                       \omega_{c'}(\gamma^{-1}y)
\]
sums to one along the action. For each left \(\Gamma_{c'}^0\)-coset
of \(v\) there is exactly one pair mapping \(\gamma\) to the chosen
representative, by the preceding freeness. Its weight is
\(\omega_{c'}(v\gamma^{-1}y)\), constant on that coset, and their
sum is one. The original cutoff product also sums to one under
the action. Both kernels are invariant by their class-level
Bessel characters. Inserting the sum of either weight into the
other weighted integral, using Tonelli for absolute values, and
changing variables by the action proves equality of the two
integrals for each invariant absolutely integrable integrand.
This applies to the full and cut kernels and to the norm of their
difference. In particular the sum of absolute errors with the new
weights is at most \(\nu b^{-k^3 f}/2\). Here the absolute error
bound is the integral of the majorant in the
\hyperlink{proof.weighted-cut-energy}{full-to-cut estimate}, before
the triangle inequality on the individual blocks.

If a chosen representative has a nonzero cut integrand, its
right argument is boxed by the property of \(\omega_{c'}\), and the
left argument is in the cut. The left partition gives a
\(u\in\Gamma_c\) moving that left argument into \(\mathrm{Box}\).
Thus the larger double class contains a near term \(u\rho\), and
its smaller double class has a near representative. By the choice
rule the selected representative is near.

For a near selected \(\gamma\) apply the
\hyperlink{proof.weighted-real-cell-contraction}{real cell contraction}
at every place with frame bound \(B\) and \(\mathfrak g_v>0\). It gives
compact sets \(A_v\), probability laws on the permitted translates
\(t_v\in U_v^0H_{0,v}\), and positive measures of masses at most
\(M_v\) from the \hyperlink{proof.weighted-cell-mass-budget}{cell mass comparison}.
For a product translate change
variables to the old left argument \(\gamma t y\). The weight is
now \(\omega_{c'}(t y)=\omega_{c'}(y)\). If the cut integrand is
nonzero, \(y_v\) has inverse frame of size at most \(B\) and
\(\operatorname{size}(\gamma_vt_vy_v)\le\theta\le B\), hence
\(y_v\in A_v\). The cut integral is therefore the integral on
\(A=\prod_v A_v\). Tensoring the local contractions and applying
them to \(\omega_{c'}|_A\) shows that the average squared norm of
the full-vector integral on \(A\) is at most
\(\prod_v M_v\), since \(0\le\omega_{c'}\le1\).
To tensor, first use finite sums of simple tensors. Approximate
the bounded weight in \(L^2\) for the product of the local sums of
cell measure and restricted Haar. This gives approximation in both
product measures separately. Convergence in restricted Haar identifies
the displayed integral for each fixed translate by continuity of the
full vectors on the compact domain. The product contraction gives
convergence in the averaged \(L^2\) norm and hence the claimed bound.
Choose a translate satisfying the support assertion and this
squared norm bound. The change from the full vector to the cut
on \(A\) costs at most the absolute error for this representative
under \(\phi\), since the latter integrates the norm on the whole
group. By the \hyperlink{proof.weighted-cell-mass-budget}{cell mass
comparison}, the full integral on \(A\) has norm at most
\(\nu u_n(\gamma)/2\) and the error at most
\(\nu b^{-k^3 f}/2\le \nu u_n(\gamma)/2\). This proves the large-range
bound by summing over the selected near representatives, which
are a subset of the representatives for the smaller double
classes in \eqref{eq:weighted-class-size}. Together with the
small-range bound and the residue test, it gives the cut block
estimate on every stable term set of a fixed level.\footnote{\weightedpin
{grh-compact/src/New/Filling/AmplitudeSums.lean}{116-L158}
{Cutoff replacement and the absolute difference}, \weightedpin
{grh-compact/src/New/RealPlace/OpenCellBound.lean}{13-L100}
{The product-cell inputs}, and \weightedpin
{grh-compact/src/New/Filling/BlockBound.lean}{301-L394}
{Choice of representatives and the two discriminant ranges}, and
\weightedpin{grh-compact/src/New/Filling/SchurTest.lean}{69-L99}
{The packet-constant comparison on this class}.}

\paragraph{A finite-module count for inverse indices.}
Write \(\widehat O=\widehat{\mathcal O}_E\), and for an integral
finite idele \(d\) write \([d]=|\widehat O/d\widehat O|\).
For a finite adelic lattice \(M\) and a submodule \(K\) put
\(I_K(M)=[M:M\cap K]\) when the index is finite.
Let \(Y_\alpha\) be integral matrices of size \(\sigma\times(r_0+s)\),
and let \(m\) be an integral matrix whose \(r_0\) columns span a
primitive direct summand of \(\widehat O^{r_0+s}\).
Suppose \(K\) contains \(d\widehat O^\sigma\).
For each \(\alpha\) let \(u_\alpha\) be a nonnegative integer
divisible by \(I_K(Y_\alpha m\widehat O^{r_0})\).
If at most \(x\) indices \(\alpha\) have any one matrix of column
classes modulo \(K\), then for \(p_0\ge r_0\)
\begin{equation}
 \label{eq:weighted-finite-inverse-index}
 \sum_\alpha u_\alpha^{-p_0}\le [d]^{\sigma(\sigma+s)}x .
\end{equation}
The weight at \(u_\alpha=0\) is zero.

Complete \(m\) to a basis of \(\widehat O^{r_0+s}\) and express
the columns in this basis. Put
\(Q=\widehat O^\sigma/(\widehat O^\sigma\cap K)\), of order at
most \([d]^\sigma\). Group the matrices by the submodule
\(A\subset Q\) generated by their first \(r_0\) column classes
and by the ordered last \(s\) column classes. Each submodule of
\(Q\) can be generated by \(\sigma\) elements. Indeed, over
each discrete valuation ring in the finite support of \(Q\),
lift the submodule to the free module of rank \(\sigma\), use
elementary divisors, and project a basis back to the submodule.
Combining the local generators gives the desired \(\sigma\)
generators. There are therefore at most
\(|Q|^{\sigma+s}\) groups. In a group the index
\(I_K(Y_\alpha m\widehat O^{r_0})\) is \(|A|\).
There are at most \(|A|^{r_0}\) ordered first-column classes.
Each occurs at most \(x\) times, and its summands have weight at
most \(|A|^{-p_0}\). This proves \eqref{eq:weighted-finite-inverse-index}.\footnote{\weightedpin{grh-compact/src/New/RealPlace/VolumeKernel.lean}
{16-L48}{The finite congruence count}.}

\paragraph{The sum of inverse discriminants on the marked class lattices.}
\teamanchor{proof.weighted-marked-disc-sum}%
Put \(a=\varrho\), \(n=a+1\), \(k=\dim_E X_0\), and
\(f=[E:\mathbf Q]/2\). Use one complex embedding at each of the \(f\)
infinite places and the usual complex Euclidean norm in the
orthonormal coordinates of \(X_0=E^k\). For a finite adelic lattice
\(\Lambda\) in \(X_0\otimes_E\mathbb A_{E,f}\) and an \(E\)-subspace \(V\)
of \(X_0\), let \(d_\Lambda(V)\) be
\(\lvert\det(m_V^*m_V)\rvert_f^{-1/2}\), where \(m_V\) is a basis of
\(\Lambda\cap V_{\mathbb A_f}\), \(*\) is conjugate transpose, and the
finite absolute value is the product of the normalized local
absolute values. A change of basis over \(\widehat O\) leaves this
quantity unchanged. With \(r_0=\dim_E V\), the covolume of
\(V\cap\Lambda\) in its real span in the \(f\) complex embeddings is
\[
       (2^{-f}\sqrt{|D_E|})^{r_0}d_\Lambda(V).
\]
To see this, write \(m_V=xA\) with \(x\) an \(E\)-basis of \(V\)
and \(A\) invertible over the finite adeles. Finite elementary
divisors multiply the covolume of \(\mathcal O_E^{r_0}\) in
the coefficient space by \(|\det A|_f^{-1}\).
At a complex embedding the real volume multiplier of \(x\) is
\(\det(x_v^*x_v)\). The product formula gives
\(\prod_v\det(x_v^*x_v)=|\det(x^*x)|_f^{-1/2}\).
The normalized discriminant is their product with \(|\det A|_f^{-1}\).
In particular passing to a sublattice of finite index \(I\)
in the same subspace multiplies the discriminant by \(I\).

First let \(\Lambda=u\widehat O^k\) with \(u\) a finite adelic
isometry of \(X_0\). This lattice is self-dual. For the primitive
lattice \(M=\Lambda\cap V_{\mathbb A_f}\), projection from
\(\Lambda\) to \(V_{\mathbb A_f}\) has image \(M^\vee\), where the
dual is taken for the Hermitian form. Indeed a functional on \(M\)
extends to \(\Lambda\) by a finite adelic integral complement,
and self-duality realizes that functional by a member of \(\Lambda\).
The projection has kernel \(\Lambda\cap V^\perp_{\mathbb A_f}\).
It follows that the index of
\((\Lambda\cap V_{\mathbb A_f})\oplus(\Lambda\cap V^\perp_{\mathbb A_f})\)
in \(\Lambda\) is \([M^\vee:M]=d_\Lambda(V)^2\).
Interchanging \(V,V^\perp\) gives \(d_\Lambda(V)=d_\Lambda(V^\perp)\).

Let \(s=\operatorname{codim}_E V\). In the real dimension \(2fs\)
of \(V^\perp\), denote the Euclidean successive minima of
\(\Lambda\cap V^\perp\) by \(\lambda_h\). If \(s>0\), choose
independent \(E\)-vectors \(x_1,\ldots,x_s\) there with
\(\|x_i\|\le\lambda_{2f(i-1)+1}\), where \(\|\cdot\|\) is the
Euclidean norm summed over the complex embeddings. At step \(i\)
the preceding \(E\)-span has real dimension \(2f(i-1)\), so one
of the \(2f(i-1)+1\) independent lattice vectors at that minimum
lies outside it. The Euclidean bound for the product of the
minima follows by stretching the unit ball along an orthonormal
flag with scales \(\lambda_h\). Choose the flag so that its
\((h-1)\)-space contains all lattice vectors of norm strictly below
\(\lambda_h\). A lattice point in the open stretched ball with
last nonzero coordinate \(h\) would have norm \(<\lambda_h\).
Half of the stretched ball thus injects, up to sets of measure
zero, into the quotient by the lattice. Its volume is at most
the lattice covolume. Writing \(\operatorname{rd}_E=|D_E|^{1/(2f)}\)
and using the cube of side \(2/\sqrt{2fs}\) in the unit ball gives
\[
    \left(\prod_{i=1}^s\|x_i\|\right)^{2f}
      \le \prod_{h=1}^{2fs}\lambda_h
      \le (8fs\operatorname{rd}_E)^{fs}d_\Lambda(V).
\]
For \(s=0\) the empty product has value one.

Every nonzero \(x\in X_0\cap\Lambda\) has
\(\|x\|\ge1\), since its squared Hermitian length is a nonzero
totally positive algebraic integer. The same bound for differences
gives unit spacing. Packing half-unit balls yields at most
\((1+2R)^{2fk}\) points of norm at most \(R\).
Integration of this count for \(R\ge1\) gives
\[
  \sum_{0\ne x\in X_0\cap\Lambda}\|x\|^{-2f(k+1)}
       \le 3^{2fk}(k+1)\le 5^{2f(k+1)} .
\]
The ordered tuple chosen for \(V\) determines \(V^\perp\) and
hence \(V\). Raising the minima bound to \(k+1\) and summing the
products of inverse vector lengths over all tuples now gives
\begin{equation}
 \label{eq:weighted-unmarked-disc-sum}
 \sum_{\operatorname{codim}_E V=s}d_{u\widehat O^k}(V)^{-k-1}
       \le (200fk\operatorname{rd}_E)^{f(k+1)s}.
\end{equation}
Here \(s\le k\). For \(s>k\) the sum is empty.\footnote{\weightedpin{grh-compact/src/New/Filling/GramLiouville.lean}
{46-L79}{The normalized discriminant and index},
\weightedpin{grh-compact/src/New/Filling/NearTerms.lean}
{22-L89}{The Euclidean minima and inverse discriminants}.}

Now use the actual class lattices. Let \(p=\pi\) be the integral
finite idele of the marked class basis. Then \(p\bar p=N\) and \([p]=N^f\).
Let \(o\) be the marked coordinate of \(X_0\), \(K^0\) its definite
unitary standard lattice stabilizer, and \(K_H\) the column level.
For the class represented by \(h_c\) the definite part is
\[
   \Lambda_c^+
     =h_c\{v\in\widehat O^k:v_o\in p\widehat O\}.
\]
Indeed in the marked basis the new vector is
\(e_o/\bar p+z/N\). On \(X_0\) its coefficient is \(\bar p\,v_o\),
and the coefficient of \(z\) is \(-v_o/p\). Thus
\(d_c(V)=d_{\Lambda_c^+}(V)\).
The preimage of the full ambient lattice stabilizer in \(K^0\)
is \(K_p=\{h\in K^0:e_o^*h\equiv e_o^*\bmod p\}\).
In the marked coordinates the differences \(h-I\) must be
integral after dividing row \(o\) by \(p\) and column \(o\) by
\(\bar p\). The row congruence for \(h^{-1}\) follows from
the row congruence for \(h\). Conjugate transpose therefore
gives the required column congruence for \(h\).
At each finite place at least one of \(p,\bar p\) is a unit,
so both divisions preserve the required integrality.
Conversely the marked chart recovers the row congruence
by multiplying back the column scale \(\bar p\).
The column level satisfies \(K_H\subset K_p\). The admissible datum\footnote{\weightedpin{grh-compact/src/New/Filling/EichlerFlag.lean}
{24-L37}{The marked basis},
\weightedpin{grh-compact/src/New/Filling/LatticeChart.lean}
{34-L63}{The definite stabilizer and its row condition},
\weightedpin{grh-compact/src/New/Filling/Setup.lean}
{303-L318}{The level and the index hypotheses}.}
gives
\[
             [K^0:K_p]\ge (N^k/b)^f .
\]

Fix a representative \(u\) at level \(K^0\), and write
\(y_\phi\) for representatives of \(K^0/K_p\).
On \(u\widehat O^k\) the row defining the marked lattice of
\(uy_\phi\) is \(v\mapsto e_o^*y_\phi^{-1}u^{-1}v\bmod p\).
These rows are distinct modulo \(p\) in the \(u\)-basis.
Equality of the rows for \(\phi,\phi'\), followed by right
multiplication by \(y_{\phi'}\), would give
\(y_\phi^{-1}y_{\phi'}\in K_p\).
For a fixed \(V\) of codimension \(s\), its primitive lattice
in \(u\widehat O^k\) has \(k-s\) basis vectors. The index
\(I_\phi(V)\) of the marked lattice in this primitive lattice
is the order of the image of the row on it modulo \(p\).
Apply \eqref{eq:weighted-finite-inverse-index} with one output
coordinate, modulus \(p\), \(x=1\), and \(p_0=k+1\). It gives
\(\sum_\phi I_\phi(V)^{-k-1}\le N^{f(s+1)}\).
Since \(d_{uy_\phi}^+(V)=d_{u\widehat O^k}(V)I_\phi(V)\),
\eqref{eq:weighted-unmarked-disc-sum} and the index bound give
\[
 \frac{1}{[K^0:K_p]}
      \sum_\phi\sum_{\operatorname{codim}_E V=s}d_{uy_\phi}^+(V)^{-k-1}
 \le N^{-kf} b^f N^{f(s+1)}
                          (200fk\operatorname{rd}_E)^{f(k+1)s}.
\]
This average is also the uniform average over the column-level
classes above the class of \(u\). In fact the finite rational
stabilizer \(u^{-1}U(X_0)(E^+)u\cap K^0\) acts on \(K^0/K_H\)
with trivial stabilizers by freeness of the column level.
The sum over subspaces is invariant under rational transport.
Thus all orbits have the same size, and passing from \(K^0/K_H\)
to \(K^0/K_p\) has constant multiplicity \([K_p:K_H]\).
This proves the assertion about the average, including its class
multiplicity.

Use \(N\le b^{60n}\),
\((200fk\operatorname{rd}_E)^{1000}\le b\), and
\(k\ge1000(2n)^9\). The extra exponent of \(b\) divided by \(f\)
is at most \(1+60n(s+1)+(k+1)s/1000\le k(s+1)/4\).
With \(\tau=b^{kf/4}\), summing over the column-level classes\footnote{\weightedpin{grh-compact/src/New/Filling/BlockBound.lean}
{18-L75}{The marked discriminant sum over all classes}.} gives
\begin{equation}
 \label{eq:weighted-marked-disc-sum}
  \sum_c\sum_{\operatorname{codim}_E V=s}d_c(V)^{-k-1}
       \le |\mathrm{Cl}|\,N^{-kf}\tau^{s+1}.
\end{equation}

\paragraph{The inverse frame-value sum on a fixed label.}
\teamanchor{proof.weighted-frame-index-sum}%
Fix a source class \(c\), a level \(j\) with \(0\le j\le a+1\),
and the \(J\)-coordinate \(t\) in the fixed set bijection
\(G\simeq G/J\times J\). Take any set \(\mathcal A\) of pairs
\((c',\gamma)\) of this label represented by near terms of level \(j\)
chosen in their \(U_c^0\)-\(U_{c'}^0\) double cosets.
Write \(\mathsf F=(e_0,\ldots,e_{a-1},z)\) as a column matrix
and \(Z=\mathsf F^*h\gamma\mathsf F\). For a fixed \(V\subset X_0\)
of codimension at most \(s\), let
\[
  M_\gamma(V)=
       \{((\gamma\mathsf F)^*h)v:v\in\Lambda_c^+\cap V_{\mathbb A_f}\},
 \qquad R_\gamma=\langle Z_{0,:}^*,\ldots,Z_{a-1,:}^*\rangle_{\widehat O}.
\]
For any nonnegative integer \(u_\gamma\) on each pair in
\(\mathcal A\) divisible by \(I_{gR_\gamma}(M_\gamma(V))\),
with the infinite index recorded as \(0\), we prove
\begin{equation}
 \label{eq:weighted-frame-index-sum}
  \sum_{(c',\gamma)\in\mathcal A} u_\gamma^{-k}
       \le b^{2000n^6(s+1)f}.
\end{equation}
The zero convention for the weight is as in
\eqref{eq:weighted-finite-inverse-index}.

In the class chart the paired frames are integral and
the class basis preserves the frame vectors. Consequently \(Z\)
has algebraic integer entries. At a near witness \(y_v\), both
\(y_v\) and \(\gamma_v^{-1}y_v\) are in the real frame box,
with \(\|y_v^{-1}f_h\|\) and
\(\|(\gamma_v^{-1}y_v)^{-1}f_h\|\le B=b^{2n^2}\).
The signed frame scales are \(s_{vh}=\kappa_v^{a-h}/\sqrt N\)
with \(1\le|\kappa_v|\le b^{23}\).
After undoing the two scales the Gram pairing at \(v\) is bounded
by \(NB^2\le b^{5n^2}\). Two distinct integral \(Z\)'s differ
by an algebraic integer with absolute value at least one at
some complex embedding and coordinate. Packing in the product
of the entry disks thus bounds the number of \(Z\)'s by
\((1+2b^{5n^2})^{2fn^2}\le b^{12fn^4}\).

Fix one such \(Z\). Let \(P\) be its first \(a\) rows and \(Y\)
its first \(j\) rows padded with the remaining rows of the
identity. Put \(\delta=\det Y\ne0\), \(M=g_0^{n^4}\), and
\[
         d=M g\delta\bar\delta,\qquad
         K=gR_\gamma+gY^*\widehat O^n .
\]
These objects are constant in the fixed \(Z\) family.
Multiplication by the adjugate of \(Y^*\) shows
\(d\widehat O^n\subset gY^*\widehat O^n\).
Also \(I_K(M_\gamma(V))\) divides \(u_\gamma\), because
\(gR_\gamma\subset K\) and the index modulo \(K\) is finite.
The weight is zero in the infinite-index case of the
original \(gR_\gamma\) index.
The Leibniz bound at the real places gives
\(\lvert\delta\rvert_v^2\le b^{11n^3}\).
Using \(g\le b^{23}\), \(g_0\le b^{11}\), \(n\ge10^5\), and
the absolute norm of a principal integral idele gives
\begin{equation}
 \label{eq:weighted-frame-modulus}
            [d]=\prod_v|\Phi_v(d)|^2\le b^{24fn^4}.
\end{equation}

For the application of \eqref{eq:weighted-finite-inverse-index},
use the adjoint of the middle part of the moved frame in the
class chart as the matrix \(Y_\gamma\), of size \(n\times(k+1)\).
Use the chart's integral unimodular middle Gram matrix times an
integral basis of \(\Lambda_c^+\cap V_{\mathbb A_f}\), in that
chart, as the matrix \(m\). Its columns generate a primitive
direct summand. Their ambient codimension is at most \(s+1\),
and \(Y_\gamma m\) generates precisely \(M_\gamma(V)\).

Consider a nonempty set of the matrices \(Y_\gamma\) whose
columns agree modulo \(K\). Take a fixed comparison term in the
set. The \hyperlink{proof.weighted-thin-class-lattice}{integral
translation calculation} with \(\xi=P\) writes the change
needed in the positive middle rows as \(gLP\), with the columns of
\(L\) in the definite class lattice. Its congruence hypothesis is
agreement modulo \(gR_\gamma+gY^*\widehat O^n\).
Extend \(L\) by zero on the \(z\)-row. If \(h_M\) is the middle
Gram matrix and \(\omega=-\zeta_3\) has trace one, the rational
unipotent matrix
\[
 \begin{pmatrix}
 I&-gL^*h_M&-\omega g^2 L^*h_ML\\
 0&I&gL\\ 0&0&I
 \end{pmatrix}
\]
is unitary and fixes the augmented frame. The first property
follows from the trace of \(\omega\), and the second from the
orthogonality of \(L\) to \(z\). Self-duality of the class lattice
makes the three off-diagonal blocks integral multiples of \(g\)
in its chart. Thus the matrix gives a left translate in \(U_c^0\).
It adds \(gLP\) in the middle and leaves the dual rows fixed.
The marked row is fixed by the pairings, so after this translate
the moved frame agrees below the \(e\)-rows with the comparison
frame.
Both \(Z\) and the vectors
\(\gamma\mathsf F\operatorname{adj}(Y)e_h\) for \(h\ge j\)
are unchanged by any such translate, since the latter vectors
belong to the standard augmented frame span.

We bound the remaining \(a\)-by-\(n\) integral \(e\)-rows
by grouping their residues modulo \(d\), after the middle
shears. To prove injectivity of these residues, we first bound
the shared vectors. For \(h\ge j\) write
\(\gamma\mathsf F\operatorname{adj}(Y)e_h=\mathsf F w_h\),
where \(w_h\in E^n\). At the witness \(y_v\) we have
\[
 y_v^{-1}\mathsf F(\Phi_v(w_h)/s_v)
  =(\gamma_v^{-1}y_v)^{-1}\mathsf F
    (\Phi_v(\operatorname{adj}(Y)e_h)/s_v),
\]
where division by \(s_v\) is entry by entry.
The right side has length at most
\(nB\sqrt N\,b^{5n^3}\) by the adjugate bound and the box.
Choose a left unipotent \(v_0\) with
\(\|v_0y_v\|_{\rm op}\le b^{7n}\), using the cutoff radius
\(b^{6n+2}\). Its action on the augmented frame is upper
unitriangular with entries of size at most \(b^{7n}B\), since
\(\|v_0f_h\|\le\|v_0y_v\|_{\rm op}\|y_v^{-1}f_h\|\).
The inverse on this span has entries bounded by
\((n-1)!(b^{7n}B)^{n-1}\le b^{3n^3}\).
Multiplying by \(v_0y_v\) and this inverse, then undoing
the scales \(|s_{vh}|\le b^{23n}\), gives
\(\lvert\Phi_v((w_h)_e)\rvert\le b^{10n^3}\) in every
\(e\)-coordinate. In particular twice this bound is strictly
less than \(M\), since \(g_0\ge b^{10}\).\footnote{\weightedpin{grh-compact/src/New/Filling/BlockBound.lean}
{90-L120}{Bounds on the shared vectors},
\weightedpin{grh-compact/src/New/Filling/BoxPartition.lean}
{411-L420}{The common frame vectors}, and
\weightedpin{grh-compact/src/New/Filling/NearTerms.lean}
{295-L317}{The sizes and integrality of frame products}.}

Suppose two aligned frames of the same \(e\)-row residue
differ in their \(e\)-rows by \(dX_1\) with \(X_1\) integral.
For \(h\ge j\) the difference of the two shared vectors
in the \(e\)-rows is \(dX_1\operatorname{adj}(Y)e_h\).
Its quotient by \(M\) is integral and has absolute value
strictly less than one at every infinite place. Hence
\(X_1\operatorname{adj}(Y)e_h=0\).
The frames have equal internal Gram matrix and agree below
the \(e\)-rows, including the \(e^*\)-block \(P\).
As \(d=\bar d\), the equality of the internal Gram matrices gives
\(X_1^*P=-P^*X_1\).
Let \(\alpha\) consist of the first \(a\) columns of
\(\operatorname{adj}(Y)\) and put \(A_1=X_1\alpha\).
For \(j\le a\), the vanishing of the other columns
and \(\operatorname{adj}(Y)Y=\delta I\) give \(A_1P=\delta X_1\).
For \(j=a+1\), let \(\beta\) be the last column of the
adjugate. Then \(P\beta=0\), so \(P^*X_1\beta=0\) by the Gram
relation. Since \(P\) has full row rank, \(X_1\beta=0\).
This gives \(A_1P=\delta X_1\) also at the top level.
The integral matrix
\[
 S_1=\bar\delta A_1-\delta A_1^*-(P\alpha)^*A_1
\]
is skew Hermitian by the Gram relation. Using \(A_1P=\delta X_1\)
and \(A_1^*P=-(P\alpha)^*X_1\) gives \(S_1P=\delta\bar\delta X_1\).
The unipotent shear adding \(MgS_1\) times the \(e^*\)-rows to
the \(e\)-rows therefore matches the two frames. It fixes the
standard frame, and is integral and congruent to one modulo
\(g\) on the class lattice, since the class chart is the
identity on these hyperbolic pairs. This puts it in \(U_c^0\).
For the resulting equal frames, the relative position
\(\kappa^{-1}\kappa'\) of the two integral positions fixes the
standard frame. Its residue is in \(J\), so the \(G/J\)-coordinates
of the two residues agree. Their fixed \(t\) gives equal exact
residues, hence \(\kappa^{-1}\kappa'\in K(g)\).
The relative rational isometry \(\eta\) fixing the frame has
upper triangular blocks in \(X\oplus(X_0\oplus Ez)\oplus X^\vee\).
On its middle quotient it acts on \(X_0\) as an isometry
\(h_\eta\) and fixes \(z\). In the finite class chart of
\(\kappa^{-1}\kappa'\), the diagonal middle block is the
marked lift of \(h_{c'}^{-1}h_\eta h_{c''}\), where \(c',c''\)
are the two target classes. Since the position is integral
and congruent to one modulo \(g\), that lift belongs to
the principal level, so \(h_{c'}^{-1}h_\eta h_{c''}\in K_H\).
The class condition gives \(c'=c''\), and freeness of the
column level gives \(h_\eta=1\).
The remaining rational relative term belongs to \(U_{c'}^0\).
The original pairs therefore have the same chosen double-coset
representative.

There are at most \([d]^{an}\) \(e\)-row residue matrices.
Thus in the application of \eqref{eq:weighted-finite-inverse-index}
we can take \(x=[d]^{an}\), \(\sigma=n\), \(p_0=k\), and
complementary rank at most \(s+1\). Summing over \(Z\), the
exponent of \(b\) is at most
\[
        f\{12n^4+24n^4[n(n+s+1)+an]\}
                \le 2000n^6(s+1)f .
\]
This proves \eqref{eq:weighted-frame-index-sum}.\footnote{\weightedpin{grh-compact/src/New/RealPlace/VolumeKernel.lean}
{79-L122}{Fixed-label uniqueness and the middle-row shear},
\weightedpin{grh-compact/src/New/Filling/BlockBound.lean}
{127-L219}{The inverse frame-value sum and the skew shear}.}

\paragraph{Pair weights for the congruent frame terms.}
\teamanchor{proof.weighted-pair-count}%
For two near terms $\gamma,\gamma'$ of level $j$ in the same
block let $L$ be the intersection of their links. Each link is a
maximal isotropic for $h_M\oplus(-h_M)$ on $M^2$ with
$M=X_0\oplus Ez$. Put $i=k+1-\dim_E L$. Let $V$ be the first
projection of $L\cap X_0^2$ and put $l=k-\dim_E V$. The first
projection is injective on $L\cap X_0^2$ by positivity of $X_0$.
The terminal pairings on either link have rank at most one in the
degenerate levels and at most two in the top level. Thus
$l\le i+(j-a)_+$, and $i=0$ means that the two links coincide.

Let $d_c(V)$ be the Hermitian discriminant of the intersection of
$V$ with the definite lattice of $c$. If $m$ is an adelic basis of
this primitive lattice in the coordinates of $X_0$, set
$d_c(V)=|\det(m^*m)|_{\mathbb A_f}^{-1/2}$. Define
$F_\gamma(V)$ as the index
$[M_\gamma:M_\gamma\cap R_\gamma]$, where $M_\gamma$ consists of
the frame values of the integral vectors of $V$ and $R_\gamma$
is the span of the adjoints of the first $a$ rows of the position's
frame pairing. Define $d_\gamma$ as the index of the lattice
generated by the two integral isotropic frames in its saturation.
Use the corresponding definitions for $\gamma'$. These indices
are taken in the lattices of $c,c'$ in the given term position.
The two ordered weights are $\omega(\gamma,\gamma')$ and its
transpose, with
\[
 \omega(\gamma,\gamma')=
 \frac{d_c(V)^{-k-1}
       (g^{2fd_j} F_\gamma(V))^{-k}d_{\gamma'}^{-k}}
      {M_0(2\tau^2)^l}.
\]
Use zero weight off the congruent near pairs of level $j$.

\paragraph{Summing the ordered pair weights.}
\teamanchor{proof.weighted-pair-index-assembly}%
For a congruent pair in the same block \((c,s,\ c')\)
let \(V\) be the definite projection of the common link,
as in the \hyperlink{proof.weighted-pair-count}{definition of the pair weight}.
Write \(l=\operatorname{codim}_E V\).
For either link and \(y\in V\), its lift in \(X^\perp\cap(\gamma X)^\perp\)
can be written \(y+x\), with \(x\in X=\langle e_0,\ldots,e_{a-1}\rangle\),
and its pullback lies in \(y'+X\) for \(y'\in X_0\).
Both projections have zero \(z\)-coordinate. Thus the lift is
orthogonal to \(\gamma\mathsf F\), including its terminal vector.
It follows that \(M_\gamma(V)\) takes values in the adelic
span of \(R_\gamma\), so
\(F_\gamma(V)=I_{R_\gamma}(M_\gamma(V))\) is finite and positive.
The rank of \(R_\gamma\) is \(d_j=\min(j,a)\). The nonzero leading
minor gives the lower bound, and for \(j\le a\) the shared
columns give \(PY^{-1}e_h=0\) for \(h\ge j\).
Therefore \([R_\gamma:gR_\gamma]=g^{2fd_j}\), and
\[
       I_{gR_\gamma}(M_\gamma(V))\mid g^{2fd_j}F_\gamma(V).
\]

For the second sum, keep \(\gamma\) fixed and restrict to
\(V_0=V\cap(\gamma\operatorname{span}\mathsf F)^\perp\),
of codimension at most \(l+n\). If \(y\in V_0\), then for the
lift \(y+x\) just used, \(\gamma^{-1}x\) is isotropic and
orthogonal to \(\mathsf F\). Its image in the positive definite
quotient \(X^\perp\cap z^\perp/X\) must therefore be zero.
It follows that \(\gamma^{-1}y\in y'+X\). The second link has
a lift of the same pair \((y,y')\), so
\((\gamma'\gamma^{-1}-1)y\in X+\gamma'X\).
For integral \(y\) in the source chart, equality of the exact
residues makes \(g^{-1}(\gamma'\gamma^{-1}-1)y\) integral.
Pairing it with the moved augmented frame gives
\(-g^{-1}(\gamma'\mathsf F)^*hy\) by the defining orthogonality
of \(V_0\).
The unsaturated two-frame lattice has image precisely
\(R_{\gamma'}\) under this pairing. Its saturated quotient has
order \(d_{\gamma'}\). Thus
\[
      I_{gR_{\gamma'}}(M_{\gamma'}(V_0))\mid d_{\gamma'} .
\]
The assertions hold on the finite adelic lattices by extending
the rational subspaces and using the same congruence there.

For fixed \(c,V,t\) sum first over \(\gamma'\) with \(\gamma\)
held fixed, then over \(\gamma\), using the two displayed
divisibilities in \eqref{eq:weighted-frame-index-sum}.
The second sum uses \(V_0\) on that same \(\gamma\);
the first-term sum uses \(V\). The product of bounds is
\[
 b^{2000n^6(2l+n+2)f}\le \tau^{l+1}.
\]
The inequality follows from \(k\ge1000(2n)^9\) and
\(2l+n+2\le3n(l+1)\). All pairs in the same block carry one
label \(t\), and there are \(|J|\) labels. Each ordered pair
is counted with its actual \(V\) and no further multiplicity.
With \(M_0=|J|\,|\mathrm{Cl}|\,N^{-kf}\), insertion of
\eqref{eq:weighted-marked-disc-sum} bounds the sum of the ordered
weights at codimension \(l\) by
\[
   \frac{|J|}{M_0(2\tau^2)^l}
     \tau^{l+1} |\mathrm{Cl}|N^{-kf}\tau^{l+1}
       =2^{-l}\tau^2 .
\]
Summing \(l\ge0\) gives, at each level \(j\),
\begin{equation}
 \label{eq:weighted-pair-weight-sum}
 \sum_{c,c',\gamma,\gamma'}\omega(\gamma,\gamma')\le2\tau^2 .
\end{equation}
This applies to the entire family of congruent near pairs,
and hence to each restricted residue set in the fixing-set
and complementary estimates. Apply the \hyperlink{proof.weighted-bordered-link-lattice}
{bordered lattice calculation} below to these same pairs for
the gain from their common-link codimension and the terminal
divisor on equal top links.\footnote{\weightedpin{grh-compact/src/New/Filling/BoxPartition.lean}
{329-L395}{Definitions of the discriminants, indices and weights},
\weightedpin{grh-compact/src/New/Filling/BlockBound.lean}
{229-L300}{The summation for the ordered pairs}.}

\paragraph{The bordered Gram determinant on the class lattice.}
\teamanchor{proof.weighted-bordered-link-lattice}%
Let all residues in \(\Omega\) fix nonzero frame vectors at places
with total norm at least \(L_*\ge1\), counting chosen and conjugate
places separately. Keep the notation of the pair count. Let \(X\) be the isotropic
frame span, \(M=X_0\oplus Ez\) the middle space, \(h\) the ambient
Hermitian Gram matrix, and \(\Lambda_c\) the self-dual adelic
class lattice. Gram entries are written \(v^*h w\). Intersections
of a rational subspace with an adelic lattice use its adelic span.
For a term \(\gamma\) of level \(j\), put \(d=\min(j,a)\),
\(t=a-d\), \(A=X+\gamma X\), and \(I=X\cap\gamma X\). The
columns of
\[
 F=(e_0,\ldots,e_{a-1},\gamma e_0,\ldots,\gamma e_{d-1})
\]
are a basis of \(A\). For \(j\le a\), the leading minor of size
\(j\) implies that the first \(j\) isotropic vectors are
independent modulo \(I\). The shared vectors at positions
\(j,\ldots,a-1\) give an independent family in \(I\) with
unit triangular last \(t\) coordinates. The same holds in
the moved frame after rescaling the common leading coefficient.
For \(j=a+1\), invertibility of the full frame pairing gives
\(I=0\), since a vector in \(I\) would pair to zero with the entire
moved frame. Thus \(\dim I=t\) in both cases. The map
\[
 \ell_\gamma:A^\perp\longrightarrow M^2,\qquad
 v\longmapsto ([v],[\gamma^{-1}v])
\]
has kernel \(I\) and image \(L_\gamma=\operatorname{link}(\gamma)\).
Its image is isotropic for \(h_M\oplus(-h_M)\) and has
dimension \(k+1\), so it equals its orthogonal complement for that
form.

Write \(S_0=(e_d^*,\ldots,e_{a-1}^*)\), \(p=(z,\gamma z)\)
and \(P_0=(S_0,p)\). The bordered matrix is
\[
 B_\gamma=(F,P_0)^*h(F,P_0)+\operatorname{diag}(0,N1_{t+2}).
\]
The notation \([u]\) for a nonzero finite idele means its inverse
finite modulus, so \([u]=|\mathrm{Nm}_{E/\mathbf Q}(u)|\) on \(E^\times\).
Group the columns as \((\gamma e_{<d},S_0,\gamma z)\) and
\((e_{<a},z)\). The second diagonal block is zero, since the
added \(N\) cancels \(h(z,z)=-N\). The cross block \(R\) has the
first \(d\) rows of \(Z(\gamma^{-1})\), the last \(t\) isotropic
identity rows and the last row of \(Z(\gamma^{-1})\).
Hence \([\det B_\gamma]=[\det R]^2\).
For \(j=a+1\), \(R=Z(\gamma^{-1})\). Otherwise replace the
last row of \(R\) by the last identity row to obtain \(Y\).
Then \(\det Y=\det Z_j(\gamma^{-1})\). A shared terminal vector
can be given terminal coordinate one in the first frame and zero
in its last \(t\) isotropic coordinates by subtracting shared
isotropic vectors. Its leading coefficient in the other frame
has absolute value one at every complex place. Its first-frame
coordinate vector is \(Y^{-1}e_{\mathrm{last}}\), because it is
orthogonal to \(\gamma e_{<d}\). The row replacement therefore
multiplies \(\det Y\) by \(-N\) times a scalar of norm one.
In the notation of the cut block sizes, this gives%
\footnote{\weightedpin{grh-compact/src/New/Filling/LinkPairing.lean}
{22-L50}{The bordered determinant in the weighted frame}.}
\begin{equation}
 \label{eq:weighted-bordered-det}
 [\det B_\gamma]=(N^{2f}D_j(\gamma))^2 .
\end{equation}

Let \(Q=\tfrac12\operatorname{diag}(1_k,N,1_k,N)\) on \(M^2\),
and let \(\mathcal L_\gamma=\ell_\gamma(\Lambda_c\cap A^\perp)\).
For an adelic basis \(v\) of a lattice in this positive form write
\(\operatorname{disc}_Q(v)=[\det(v^*Qv)]^{1/2}\).
We compute the discriminant of \(\mathcal L_\gamma\).
Choose a rational basis \((u_I,u_L)\) of \(A^\perp\) with
\(u_I\) a basis of \(I\), \(S_0^*h u_I=1_t\) and
\(S_0^*h u_L=0\). The projection on the last \(t\) isotropic
coordinates of \(I\) is invertible by the shared vectors, so
this choice is possible. The columns \(\ell_\gamma(u_L)=v\)
are a basis of the link. In the flag \(I\subset A^\perp\),
choose an integral adelic basis of the saturated intersection
with \(\Lambda_c\), of the form
\(\mathcal U=(u_I,u_L)\left(\begin{smallmatrix}A_0&B_0\\0&E_0\end{smallmatrix}\right)\).
Choose the basis over each valuation ring by extending a basis of
the primitive first lattice. The associated changes are unimodular
outside a finite set. The columns \(vE_0\) are a basis of
\(\mathcal L_\gamma\).

Let \(\iota=[\Lambda_c\cap A:\langle F\rangle]\), with the spans
taken over \(\widehat{\mathcal O}_E\). The complementary-minor
identity on this lattice gives
\[
 [\det B_\gamma]=[N]^{t+2}\iota^2
  [\det(\mathcal U^*(h+N^{-1}hP_0P_0^*h)\mathcal U)].
\]
To see the identity, first replace \(F\) by a primitive basis
\(F_{\rm sat}\), extracting \(\iota^2\). Primitivity and
self-duality make \(F_{\rm sat}^*h\) a surjection on the integral
ambient lattice. In the unimodular basis \((T,\mathcal U)\)
provided by a section \(F_{\rm sat}^*hT=1\), let \(J_0\) be
the ambient Gram matrix and \(B\) the coordinates of \(P_0\).
The bordered Gram matrix for \(F_{\rm sat}\) is the principal
minor deleting the \(\mathcal U\)-indices of
\[
 \begin{pmatrix}
  J_0^{-1}&B\\ B^*&B^*J_0B+N1
 \end{pmatrix}.
\]
The Schur complement gives the full determinant and upper
inverse block, respectively,
\[
 \frac{N^{t+2}}{\det J_0},\qquad J_0+N^{-1}J_0BB^*J_0.
\]
Jacobi's identity gives the formula since \([\det J_0]=1\).

In the rational basis \((u_I,u_L)\) the last Gram matrix before
the adelic basis change is \(\operatorname{diag}(N^{-1}1_t,v^*Qv)\).
Indeed the \(S_0\)-pairings give the first block. The two middle
components of \(\ell_\gamma\) have the same \(h_M\)-pairing on
the link, and the \(z,\gamma z\) terms change their negative
coordinates to positive ones. The factor
\([N]^{t+2}[N^{-1}]^t=[N]^2\) cancels \(N^{4f}\) in
\eqref{eq:weighted-bordered-det}, giving
\[
 \operatorname{disc}_Q(\mathcal L_\gamma)[\det A_0]\iota
     =D_j(\gamma).
\]

The class chart fixes \(e_i,z,e_i^*\). Its other columns use the
middle lattice basis with its marked coordinate. The map
\(x\mapsto(\gamma f)^*h x\) on the integral unmoved isotropic
lattice has image \(R_\gamma\) and kernel \(I\cap\Lambda_c\).
For the kernel, \(\gamma^{-1}x\) is isotropic and orthogonal to the
unmoved augmented frame. Its class in the
positive definite quotient of \(X^\perp\cap z^\perp\) by \(X\)
is zero. This proves the kernel identity over \(E\) and hence
in the adelic chart. Modulo the image of the first \(d\)
unmoved isotropic vectors, the quotient is the last \(t\)
coordinates modulo \(A_0\widehat{\mathcal O}_E^t\). Its index
is \(I_j(\kappa)=[\det A_0]\). In the other direction, insert
the whole two-frame lattice between \(\langle F\rangle\) and
\(\Lambda_c\cap A\), extracting its saturation index \(d_\gamma\).
Pairing with the unmoved augmented frame kills \(X\), and maps
the remaining index onto the quotient of the columns for all
\(\gamma e_i\) by those for \(i<d\). Those columns are the
adjoints of the isotropic rows at \(\kappa^{-1}\). Thus
\(I_j(\kappa^{-1})d_\gamma\) divides \(\iota\), and%
\footnote{\weightedpin{grh-compact/src/New/Filling/LinkPairing.lean}
{52-L119}{Integral lifts and the three index factors}, and
\weightedpin{grh-compact/src/Lib/General/Algebra/LatticeComplement.lean}
{66-L102}{The complementary-minor identity over the finite adeles}.}
\begin{equation}
 \label{eq:weighted-bordered-lifts}
 \operatorname{disc}_Q(\mathcal L_\gamma)
       I_j(\kappa)I_j(\kappa^{-1})d_\gamma\le D_j(\gamma).
\end{equation}

\paragraph{Separating the common-link and terminal factors.}
Let \(L=L_\gamma\cap L_{\gamma'}\), \(L_0=L\cap X_0^2\), and
let \(V\) be its first projection. Use \(Q\)-orthogonal
complements for the rational flag \(L_0\subset L\subset L_\gamma\).
An adapted adelic basis of \(\mathcal L_\gamma\) is upper triangular
in this flag. Its discriminant is the product of the discriminants
of the three diagonal lattices. The first projection identifies
the form on \(L_0\) with the definite form on \(V\). An integral
lift of a vector of \(L_0\) has the form \(x+x_0\), with
\(x\in\Lambda_c\cap V\), \(x_0\) in the integral isotropic
lattice, and is orthogonal to the entire moved frame. Hence
\((\gamma f)^*h x=-(\gamma f)^*h x_0\in R_\gamma\). The first
diagonal lattice therefore has discriminant at least
\(d_c(V)F_\gamma(V)\), by the finite index of this frame-value
condition on \(\Lambda_c\cap V\).

Let \(z_0\) take pairings with \(z\) on the two components of
the link. On \(L_\gamma\) it is represented for \(Q\) by two
vectors \(\epsilon_1=(u-z,yz)\), \(\epsilon_2=(\bar y z,u'-z)\)
in \(L_\gamma\). At top level obtain \(\epsilon_1\) by applying
\(\gamma f\) to \(NZ(\gamma)^{-1}e_{\mathrm{last}}\) and taking
its link image. Then \(y=N\det Z_a/\det Z\) and
\(\delta(\gamma)=\det Z\,\overline{\det Z}(1-y\bar y)\). At lower levels use
a shared terminal vector, giving \(y\bar y=1\). The inverse
construction gives the other vector with terminal multiplier
\(\bar y\), by isotropy for the difference form. This isotropy
also gives \(\epsilon^*Q=z_0\) on the link. Their Gram is
\(N\left(\begin{smallmatrix}1&-\bar y\\-y&1\end{smallmatrix}\right)\).
By positivity its rank is the rank of \(z_0\) on the link, one
below the top level and two at the top.

For the middle diagonal lattice, of rank \(r_0=\dim(L/L_0)\),
this shows that \(l=k-\dim V=i-1+r_0\le i+(j-a)_+\). Its
terminal pairings agree with those of the full integral lifts.
Thus they are integral and have a nonzero minor of
size \(r_0\). Let \(L'=L_*\) when \(j=a+1\) and \(L=L_\gamma\),
and \(L'=1\) otherwise. In the first case \(r_0=2\), and the
terminal determinant has index at least \(L_*\) by the
fixed-vector determinant test below. In the other cases the
nonzero integral minor has index at least one. Gram
Cauchy--Schwarz on the corresponding \(r_0\) columns of
\(\epsilon\) bounds that index by their discriminant times
the discriminant of the middle diagonal lattice. The former
factor is \(1\) for \(r_0=0\), \(N^f\) for \(r_0=1\), and
\(x(\gamma)/D_j(\gamma)\) for \(r_0=2\). The last equality
uses \(r_0=2\) only at top level and the factor \(N^{2f}\)
in \(D_j(\gamma)\).

For the determinant test, work at one of the specified split
places and use the reduction of the lattice chart. If the
terminal determinant were a unit, the two forms defined by the
integral lifts, which vanish on \(X\) and \(\kappa X\), would
force both terminal coefficients of every common vector of the
two augmented frames to be zero. Suppose a nonzero common vector
has last nonzero isotropic coefficients \(x_b,y_{b'}\) on the
two sides. For any linear form \(\lambda\) vanishing on
the first \(b+1\) and moved first \(b'+1\) isotropic vectors,
the rank-one translation \(1+u\,v\lambda\), \(u\) in the
residue field, belongs to both flag unipotents. At the conjugate
place the corresponding inverse adjoint fixes the augmented
frame, so its simple-root sum is zero. Lift the two translations
separately, using the corresponding flag vectors and linear
forms in their respective coordinates. The lifts can be supported
on this one split pair. The reductions intertwine at both places
by inverse adjoint, including when the fixed vector is on the
conjugate side. The locally self-dual phase has conductor
\(\mathcal O_{E,w}\) at the working place \(w\), where \(g\) has
valuation one. Division by \(g\) therefore gives a nontrivial
additive character of its residue field. Comparing the phases
for all \(u\) gives
\(x_b\lambda(f_{b+1})=y_{b'}\lambda(\kappa f_{b'+1})\).
The annihilator relation gives a common vector with these same
leading coefficients at the next indices. Iterating reaches
a nonzero terminal coefficient. This contradicts the unit
determinant since a fixed nonzero frame vector is already a
common vector. This applies at each chosen or conjugate place
separately, so the product of their norms divides the ideal
index of the determinant.

Write \(\mathcal Y_\gamma\) for the third diagonal lattice
in the complement of \(L\) in \(L_\gamma\). Combining the first
and the middle factors with \eqref{eq:weighted-bordered-lifts}
gives exactly
\[
 L' d_c(V)F_\gamma(V)d_\gamma
               \operatorname{disc}_Q(\mathcal Y_\gamma)
       \le B_j(\gamma).
\]
The pairing for \(h_M\oplus(-h_M)\) between the complements of
\(L\) on the two sides is nondegenerate. It equals the pairing
of the integral lifts by \(h(1-\gamma\gamma'^{-1})\), since
components in the common link pair to zero. Congruence of the
two term positions makes this pairing divisible by \(g\).
The difference form is twice a \(Q\)-unitary change, so the
nonzero pairing determinant in codimension \(i\) gives
\(\operatorname{disc}_Q(\mathcal Y_\gamma)
\operatorname{disc}_Q(\mathcal Y_{\gamma'})\ge(g/2)^{2if}\)
by Gram Cauchy--Schwarz. Multiplication and taking square roots
give the pair charge on the actual term positions%
\footnote{\weightedpin{grh-compact/src/New/Filling/BoxPartition.lean}
{150-L210}{Terminal determinant at a fixed-vector place},
\weightedpin{grh-compact/src/New/Filling/LevelBridge.lean}
{142-L148}{Conductor of the self-dual phase},
\weightedpin{grh-compact/src/New/Filling/GramLiouville.lean}
{47-L79}{Discriminant and index in the definite lattice}, and
\weightedpin{grh-compact/src/New/Filling/LinkPairing.lean}
{165-L315}{The three factors of the link and the congruent-pair gain}.}
\begin{equation}
 \label{eq:weighted-pair-gauge}
 \sqrt{B_j(\gamma)B_j(\gamma')}
 \ge (g/2)^{if} L' d_c(V)
        \sqrt{F_\gamma(V)F_{\gamma'}(V)d_\gamma d_{\gamma'}} .
\end{equation}

\paragraph{The charge for a pair of cut terms.}
\teamanchor{proof.weighted-crude-energy}%
To use this gain, retain a set $\mathcal R$ of links in the top
level, and retain none at a lower level. An omitted pair either
lies at a lower level, has distinct top links, or has both terms
in a top link outside $\mathcal R$. For such a pair inserting
\eqref{eq:weighted-pair-gauge} into \eqref{eq:weighted-class-size}
and keeping the extra factor $d_c(V)^{-1}$ gives
\[
 u_j(\gamma)u_j(\gamma')\le
 \frac{\tau^4\mathcal D_j(2\tau^2)^l(2/g)^{ikf}}{(L')^k}
 \sqrt{\omega(\gamma,\gamma')\omega(\gamma',\gamma)} .
\]
The extra powers of the positive integer indices and of
$(g/2)^{if} L'$ can be dropped because these quantities are at
least one. Put \(\eta=2\tau^2(b/g)^{kf}\). The
\hyperlink{proof.weighted-energy-ratios}{energy ratios on this datum}
give \(\eta\le1\) and
\(\tau^{18}\eta\mathscr M=2(b^6/g)^{kf}\mathscr M\le\mathscr D\).
Fix
$C_*\ge\mathscr D$, with $C_*\ge\tau^{20}\mathscr D$ if depth zero
occurs, and with $C_*\ge\tau^{12}\mathscr M/L_*^k$ if top links
are omitted. Then the coefficient of the geometric mean is at
most $C_*/\tau^3$ in every omitted pair.

For $j=0$ the bound follows from $\tau^{17}\mathscr D\le C_*/\tau^3$.
For $1\le j\le a$, the factor $\eta^i\le1$ leaves at most
$\tau^{17-32j}\mathscr D$. For $j=n,i\ge1$, use $l\le i+1$
and $2\tau^6\eta\mathscr M\le 2\tau^{-12}\mathscr D$.
For $j=n,i=0$ use $2\tau^6\mathscr M/L_*^k\le 2\tau^{-6} C_*$.
The inequalities use $b\ge1000$ and $\tau\ge2$.

At fixed $(c,s,c')$ sum the amplitudes within each link before
expanding the squared norm. Retained diagonal link squares stay
in that sum. Bound the remaining inner products by the products
of the norms, then use \eqref{eq:weighted-class-size} and the
bound for omitted pairs. A pair with fixed $c,c'$ determines its
residue $s$, so it enters the sum over $\Omega$ once. By AM--GM and
\eqref{eq:weighted-pair-weight-sum}, the total omitted energy at
level $j$ is at most
\[
 \frac{C_*}{2\tau^3}\sum_{c,c',\gamma,\gamma'}
       \bigl(\omega(\gamma,\gamma')+\omega(\gamma',\gamma)\bigr)
 \le\frac{2C_*}{\tau}
 \le \frac{b^{-30} C_*}{4^{n-j}}.
\]
Here $\tau\ge b^{n+31}$ on the datum. This calculation uses the
same near representatives and class charts throughout.\footnote{
\weightedpin{grh-compact/src/New/Filling/CrudeEnergy.lean}{98-L160}{The per-level pair charge}.}

For a stable set $T$, partition the degenerate terms with agreeing
rational characters by their least level $j<n$. The least-level
sets are stable because each level and the character-agreement
condition are invariant under the two rational unipotent actions.
Terms with disagreeing rational characters have zero amplitude.
The norm of the entire cut degenerate part is at most
$b^{-15}\sqrt{C_*}\sum_{j=0}^{n-1}2^{-(n-j)}\le b^{-15}\sqrt{C_*}$.
If the retained cut link squares at the top level sum to at most
$\nu^2 w_*$, the top norm is at most $\sqrt{w_*+b^{-30}C_*}$.
The triangle inequality and \eqref{eq:weighted-cut-energy-error}
therefore give
\begin{equation}
 \label{eq:weighted-crude-norm}
 \|A_T\|_\Omega\le \sqrt{w_*}+3b^{-15}\sqrt{C_*}.
\end{equation}
Young's inequality with parameter $b^{-4}$ gives the stated full
energy bound $(1+b^{-4})^2w_*+b^{-12}C_*$. Indeed the two
coefficients supplied by Young are $1+b^{-4}$ and
$9(1+b^4)b^{-30}$, which are no larger than the stated bounds.

\paragraph{The full energy on the required residue sets.}
On \(\Omega_A\) take \(C_*=\mathscr M_A\), using
\(\tau^{44}\mathscr D\le\mathscr M\) from the
\hyperlink{proof.weighted-energy-ratios}{datum ratios} and \(|A|\le f\).
For proper \(A\) retain all top links and set
\(w_*=(1+b^{-4})\mathscr M_A\) by the
\hyperlink{proof.weighted-finite-bessel-comparison}{cut open-link estimate}.
When \(A\) is the full set, retain no links, take \(w_*=0\)
and use \(L_*\ge g_0^f\), which gives
\(\tau^{12}\mathscr M/L_*^k\le\mathscr M_A\).
The crude estimate with these constants gives
\begin{equation}
 \label{eq:weighted-open-energy}
    \mathscr E_{\rm all}(\Omega_A)\le(1+4b^{-4})\mathscr M_A.
\end{equation}
The constant follows from \((1+z)^3+z^3\le1+4z\)
at \(z=b^{-4}\).\footnote{\weightedpin
{grh-compact/src/New/Filling/Setup.lean}{296-L318}{Admissibility ratios}.}

\paragraph{The complement on the two required residue sets.}
\teamanchor{proof.weighted-complement-crude}%
Set \(\Omega_i^\perp=JG_i\) and \(T=S(F)^c\), and retain no links.
The coordinate away from $G_i$ reduces to a Bessel-subgroup
element. It fixes the first isotropic frame vector at every chosen
place and its conjugate over $g_{1-i}$, so $L_*=g_{1-i}^{2f}$.
The depth-zero calculation for
\hyperlink{proof.weighted-compatible-orbits}{compatible frame orbits}
puts each such term in \(S(F)\).

The ratio of the two level primes is at most $b$. Thus
\[
 \frac{\tau^{12}\mathscr M}{L_*^k}
 \le (b^4/g)^{kf}\mathscr M
 \le \tfrac12 b^{-2kf}\mathscr D\le\mathscr D,
\]
The middle inequality uses
\(2(b^6/g)^{kf}\mathscr M\le\mathscr D\) from the
\hyperlink{proof.weighted-energy-ratios}{datum ratios}. Take $C_*=\mathscr D$ and
$w_*=0$ in \eqref{eq:weighted-crude-norm}. In the normalized unfolding, the amplitude is that of
\(P^*\rho(s)P-T_S(s)\) on the given map \(P\). Thus
\begin{equation}
 \label{eq:weighted-complement-energy}
    \mathscr E_{S(F)^c}(\Omega_i^\perp)
       \le9b^{-30}\mathscr D\le b^{-12}|J|C .
\end{equation}
The fixed-component index in the pinned theorem is $1-i$.\footnote{
\weightedpin{grh-compact/src/New/Filling/Moments.lean}{186-L218}{Application on the Bessel-subgroup complement}.}

\paragraph{The finite Bessel phase.}
\teamanchor{proof.weighted-finite-bessel-phase}%
Write \(\chi_U^\natural(s)=|G|^{-1}\sum_{u\in G}\chi_U(usu^{-1})\).
We bound the finite phase in counting measure. At a chosen place
above \(g_i\), put \(q=g_i\), \(a=\varrho\ge1\), \(n=k+2a+1\), and
\(V=\mathbf F_q^n\). At the split place the residual group is
\(\mathrm{GL}(V)\). Order the frame and its dual as
\[
 e_1,\ldots,e_a,\quad [z,X_0'],\quad e_a^*,\ldots,e_1^* .
\]
Here \(X_0'\) is the orthogonal complement of \(z\) in the middle
block over the residue field. The lattice change preserves \(z\),
of norm \(-N\), and \(N\) is a unit at this place.
The image of the integral unipotents is \(1+\mathfrak u\), with
\(\mathfrak u\) the strictly block upper triangular matrices for
block sizes \((1^a,k+1,1^a)\). To lift a residue, lift the strictly
upper blocks in the chosen split component. At its conjugate component
use the inverse transpose transported by the lattice Gram matrix.
These integral lifts are block unitriangular and can be set to the
identity away from the chosen pairs of places.

The unipotent character has the simple entries joining
\(e_1,\ldots,e_a,z\) at the chosen place. At the conjugate place the
inverse transpose moves those entries to the successive links joining
\(z,e_a^*,\ldots,e_1^*\). In fact the adjacent block of the inverse is
the negative of the original adjacent block. The inverse Gram matrix
interchanges the paired isotropic vectors, and the middle boundary
uses the covector \(h(z,\cdot)\), which takes the nonzero value
\(-N\) on \(z\). Both residue additive characters are nontrivial because
\(\psi_f(-x/\varpi)\) is self-dual on the local integers, and the
level prime has valuation one. Relative to one fixed nontrivial
additive character \(\psi\) of \(\mathbf F_q\), each link has a
nonzero coefficient. Conjugating by a diagonal change of basis along
the chain makes them all one. Such a change preserves
\(\mathfrak u\) and the squared norm of the conjugacy average.
With \(J_0\) the resulting lower shift along this chain, the phase is
\(\psi(\operatorname{tr}(J_0X))\) at \(1+X\), \(X\in\mathfrak u\),
and zero elsewhere.\footnote{\weightedpin
{grh-compact/src/New/Filling/LevelBridge.lean}
{117-L151}{Conjugate reduction and self-dual character}, \weightedpin
{grh-compact/src/New/Filling/Setup.lean}
{386-L396}{The finite phase}, and \weightedpin
{grh/src/New/Spanning/Unfolding/Moments.lean}
{725-L738}{The extended character and its average}.}

Let \(\mathfrak p\) be the block upper triangular algebra including
its diagonal blocks, and let \(\chi^\natural\) denote the local
conjugacy average. For
\(F(X)=1_{\mathfrak u}(X)\psi(\operatorname{tr}(J_0X))\) on
\(\operatorname{End}(V)\), use the counting transform
\(\widehat F(A)=\sum_X F(X)\psi(-\operatorname{tr}(AX))\).
An entry in a strict upper block of \(X\) pairs with the transposed
entry of \(A\), so \(\mathfrak u^\perp=\mathfrak p\) for the trace
pairing. Thus \(\widehat F=|\mathfrak u|1_{J_0+\mathfrak p}\).
Fourier transform commutes with conjugacy averaging. The average
of \(F\) is supported on nilpotents, and \(X\mapsto 1+X\) identifies
it there with \(\chi^\natural\), with \(\chi^\natural(1)=F(0)=1\).
Plancherel, with \(A\) uniform in \(\operatorname{End}(V)\), gives
\[
 \|\chi^\natural\|_{\ell^2(\mathrm{GL}(V))}^2
   =|\mathfrak u|^2\,\mathbb E_A p(A)^2,\qquad
 p(A)=\mathbb P_{g\in\mathrm{GL}(V)}(g^{-1}Ag\in J_0+\mathfrak p).
\]

In an ordered basis \((b_1,\ldots,b_a,z,x_1,\ldots,x_k,t_1,\ldots,t_a)\)
write \(b_{a+1}=z\), \(c_i=t_{a+1-i}^*\) for \(1\le i\le a\),
and \(c_{a+1}=z^*\), using the dual basis. The lower blocks of
\(J_0+\mathfrak p\) prescribe the first \(a\) columns and last \(a\)
rows below the diagonal blocks. Equivalently,
\[
 A b_i-b_{i+1}\in\langle b_1,\ldots,b_i\rangle,\qquad
 c_i A-c_{i+1}\in\langle c_1,\ldots,c_i\rangle
 \quad(1\le i\le a).
\]
Starting at \(v=b_1\) and \(\lambda=c_1\), the two lists are unit
triangular changes of \(v,Av,\ldots,A^av\) and
\(\lambda,\lambda A,\ldots,\lambda A^a\). Such changes preserve the
evaluation matrix with only its last diagonal entry equal to one.
Hence the dual-basis conditions \(c_i(b_j)=0\) except for
\(c_{a+1}(b_{a+1})=1\) are equivalent to
\[
 \lambda A^h v=0\quad(0\le h<2a),\qquad \lambda A^{2a}v=1 .
\]
These moments force independence of both sequences through power
\(2a\). A dependence with highest index \(h\) would give a
contradiction after multiplying by \(A^{2a-h}\) and evaluating
against the other initial vector or covector. For each \(v\) with
independent sequence there are \(q^k\) choices of \(\lambda\), since
the moments prescribe \(2a+1\) independent evaluations.
Each pair gives \(q^{a(a+1)}\) choices of the two unit triangular
changes. The \(x_i\) lift an ordered basis of
\(\bigl(\bigcap_{j=1}^{a+1}\ker c_j\bigr)/\langle b_1,\ldots,b_a\rangle\),
which has dimension \(k\), with \(q^a\) choices of lift per vector.
Each \(t_i\) can be any vector in its prescribed affine
\(c_j\)-evaluation fibre of size \(q^{n-a-1}\). These choices complete
the basis and their counts are constant. It follows that \(p(A)\) is
proportional to the fraction \(d(A)\) of vectors with independent
sequence through power \(2a\). The affine codimension gives
\(\mathbb E_A p(A)=|\mathfrak u|^{-1}\).
For a uniform \(A\) and
uniform \(v\), expose \(v,Av,\ldots,A^{2a}v\) in order. After
independence up to any given step, the next vector is the image of a
new independent input to the uniform linear map, and hence is
uniform in \(V\). Thus
\[
 \mathbb E_A d(A)=\prod_{j=0}^{2a}(1-q^{j-n}),
 \qquad
 \|\chi^\natural\|_{\ell^2(\mathrm{GL}(V))}^2
 =\frac{\mathbb E_A d(A)^2}{(\mathbb E_A d(A))^2}
 \le \prod_{j=0}^{2a}(1-q^{j-n})^{-1}.
\]
All \(2f\) split factors are independent in the residual group \(G\).
Both the phase and its conjugacy average factor over them. The
product of the denominators in this bound is at least
\(1-2f/(q_{\min}^k(q_{\min}-1))>1/2\), using
\(q_{\min}=\min(g_0,g_1)\ge b^{10}\) and \(f\le b\) for large \(x\).
We have proved
\[
 \sum_{s\in G}|\chi_U^\natural(s)|^2\le2 .
\]

\paragraph{The projected covariance.}
\teamanchor{proof.weighted-projected-moment}%
To use \eqref{eq:weighted-open-energy}, let $\Theta(s)=\|K(s)\|_{\rm HS}^2$.
It is nonnegative and of positive type since it is the Hilbert--Schmidt
pairing of $PP^*$ with $\rho(s)PP^*\rho(s)^*$. Choose $A$ and fixed
vectors to maximize $b^{6k|A|}\sum_{s\in F_A}\Theta(s)$.
If $A$ is full, \eqref{eq:weighted-open-energy} suffices. Otherwise take
one missing place, with residue field $\mathbf F_q$. Let
$M_A=\sum_{s\in F_A}\Theta(s)$. Maximality bounds the $\Theta$-mass of
the stabilizer in $F_A$ of any prescribed nonzero vector of the frame
or dual frame there by $Y=b^{-6k}M_A$.

For a leading $l$-by-$l$ frame block of $s$ at this place, write it
as $A_1^\mathsf{t}\pi(s)A_0$, where $\pi$ acts on $\mathbf F_q^r$ and
$A_0,A_1$ are the frame and paired-frame inclusions of $\mathbf F_q^l$.
For a nontrivial additive character $\psi$ of $\mathbf F_q$, the left
kernel size equals $\sum_y\mathbb E_z\psi(y^\mathsf{t}A_1^\mathsf{t}\pi(s)A_0z)$
and is at least $q$ for a singular block. The two axes of this character
sum cost at most $2M_A$ after multiplying by $\Theta$ and summing.
For $y,z\ne0$, use the delta basis for $\pi(s)A_0z$ and the orthonormal
Fourier basis for the contragredient orbit of $A_1y$. Tensor each with
vectors $u_s$ having $\langle u_a,u_b\rangle=\Theta(a^{-1}b)$.
The squared norm of either sum of tensor vectors over $s\in F_A$ equals
$|F_A|$ times the $\Theta$-mass of the corresponding stabilizer, by expanding
the squared norm and grouping by $a^{-1}b$. The mixed pairing of the two
sums equals $|F_A|q^{-r/2}$ times the corresponding phase sum.
Cauchy--Schwarz therefore bounds that sum by $q^{r/2}Y$. Summing in
$y$ and averaging in $z$ gives
\begin{equation*}
 q\sum_{\substack{s\in F_A\\\det(A_1^\mathsf{t}\pi(s)A_0)=0}}\Theta(s)
       \le 2 M_A + q^{l+r/2}Y\le 3M_A .
\end{equation*}
The last bound uses $l\le\varrho+1$, $r=k+2\varrho+1$, $q\le b^{11}$
and $k\ge1000(2\varrho+2)$, which gives
$11(l+r/2)\le 6k$ through $l=\varrho+1$. The union bound over all
these minors costs at most $3(\varrho+1)/q\le b^{-4}$ of $M_A$,
since $\varrho+1\le k\le b$ and $q\ge b^{10}$. Every residue in
$F_A\setminus\Omega_A$ is singular at one of these minors at the
missing place. Thus $(1-b^{-4})M_A\le(1+4b^{-4})\mathscr M_A$.
The maximal choice bounds the unrestricted sum by
$b^{6k|A|}M_A$. In particular $\sum_s\Theta(s)\le(1+9b^{-4})\mathscr M$.
\footnote{\weightedpin{grh/src/New/Spanning/SecondMoment/PositiveType.lean}
{148-L183}{Finite Fourier argument}, and \weightedpin
{grh-compact/src/New/Filling/Moments.lean}{220-L257}{Positive-type moment assembly}.}

For the rank argument use the projected covariance
\begin{equation*}
   H=|G|^{-1}\sum_s P_W\rho(s)PP^*\rho(s)^*P_W .
\end{equation*}
\teamanchor{proof.weighted-covariance-identity}%
The \hyperlink{proof.full-trace-identity}{relative trace identity}
gives the squared Hilbert--Schmidt norm
\(|G|^{-1}\sum_s\Theta(s)\) of the ambient covariance. Take
\(T_s=\rho(s)P\) and group the sum by the relative translate.
Compression decreases the norm. Applying
\eqref{eq:weighted-row-mass} to the same \(W\) gives
\begin{equation}
 \label{eq:weighted-projected-covariance}
  \frac{R\operatorname{tr}(H^2)}{C^2}\le 1+9b^{-4}.
\end{equation}

\paragraph{The two complementary estimates needed for rank.}
\teamanchor{proof.weighted-induced-bad-trace}%
The small share of a bad finite type reduces the precision needed in
one of the complementary orbital estimates. Put
\[
 \Delta_i=\frac{1}{|J|C}\sum_{s\in JG_i}
       \|K(s)-T_S(s)\|_{\rm HS}^2 .
\]
Fix \(i\) and put \(D_i=J\cap G_i\),
\(\tau_i=T_S|_{D_i}\) and \(N_i=[G_i:D_i]\).
On \(D_i\) the comparison term is a unitary representation on the
column space \(\mathcal C\). Its extension by zero to \(G_i\) is
\(T_S|_{G_i}\). The covariance \(K(ts)=\tau_i(t)K(s)\) implies
\(P^*\rho(t)=\tau_i(t)P^*\) on the span of the
\(\rho(s)P\mathcal C\), and both sides vanish on its orthogonal
complement. Taking adjoints and replacing \(t\) by \(t^{-1}\) gives
\(\rho(t)P=P\tau_i(t)\).

Let \(\pi\) be the induced representation on the orthogonal sum of
the column fibres over \(G_i/D_i\). We use
\([ht,v]=[h,\tau_i(t)v]\) and left translation. The map
\(S_i[h,v]=\rho(h)Pv\) intertwines this representation with \(\rho|_{G_i}\).
Put \(L_i=S_i^*S_i\). The \((hD_i,h'D_i)\) block of \(L_i-I\)
in coset representatives is \(K(h^{-1}h')-T_S(h^{-1}h')\).
The Hilbert--Schmidt norm of this block is unchanged by left and
right multiplication of its argument by \(D_i\). Each summation
over \(G_i/D_i\) is consequently a sum over \(G_i\) divided by \(|D_i|\).
Left \(J\)-covariance gives
\[
 \frac{\|L_i-I\|_{\rm HS}^2}{N_iC}
  =\frac{1}{|D_i|C}\sum_{h\in G_i}\|K(h)-T_S(h)\|_{\rm HS}^2
  =\Delta_i .
\]

For any self-adjoint central idempotent \(a\in\mathbf C[G_i]\),
the identity-coset block of \(\pi(a)(L_i-I)\) is
\(\sum_{h\in G_i}a(h)(K(h)-T_S(h))\).
The diagonal coset blocks are unitarily conjugate by equivariance.
Their squared norms sum to at most \(\|L_i-I\|_{\rm HS}^2\), since
\(\pi(a)\) is an orthogonal projector. Thus
\[
 \Big\|\sum_{h\in G_i}a(h)(K(h)-T_S(h))\Big\|_{\rm HS}^2
 \le C\Delta_i .
\]

Take \(a_i=1-e^{(i)}\), \(A_i=\pi(a_i)\) and
\(p_i=\operatorname{rank}(A_i)/(N_iC)\).
The compression of \(\pi(h)\) to the identity coset is \(T_S(h)\).
As all diagonal coset blocks of \(A_i\) have the same trace,
\[
 p_i=C^{-1}\sum_{h\in G_i}a_i(h)\operatorname{tr}T_S(h)
     =\sum_{h\in G_i}a_i(h)\chi_i^\natural(h),
\]
where the one-factor trace formula gives the phase and \(\chi_i^\natural\)
is its conjugacy average on \(G_i\).
Likewise the identity-coset trace of
\(A_iL_i=S_i^*\rho(a_i)S_i\) is \(\operatorname{tr}(P^*\rho(a_i)P)\).
Hilbert--Schmidt Cauchy--Schwarz gives
\[
 \left|\frac{\operatorname{tr}(P^*\rho(a_i)P)}{C}-p_i\right|
 =\frac{|\operatorname{tr}(A_i(L_i-I))|}{N_iC}
 \le\sqrt{p_i\Delta_i}.
\]
This uses the counting normalization of the actual complementary
orbital sums and the phase of the Bessel term at the same
place.\footnote{\weightedpin{grh/src/New/Spanning/Unfolding/Moments.lean}
{434-L458}{Normalized complementary blocks} and
\weightedpin{grh/src/New/Spanning/Unfolding/Moments.lean}
{725-L773}{Trace of the comparison term on a single factor}.}

\paragraph{The identity contribution to the finite share.}
\teamanchor{proof.weighted-phase-share}%
Put \(q_{\min}=\min(g_0,g_1)\) and
\(u=2f/[q_{\min}^k(q_{\min}-1)]\). On \(G_i\), \(\chi_i^\natural(1)=1\).
The product of the chain probabilities in the
\hyperlink{proof.weighted-finite-bessel-phase}{finite phase calculation}
is at least \(1-u\). The product over just the factors of \(G_i\)
has the same lower bound. In counting measure, with \(\delta_1\)
the indicator of the identity, this gives
\[
 \|\chi_i^\natural-\delta_1\|_2^2
    =\|\chi_i^\natural\|_2^2-1\le\frac{u}{1-u}.
\]
Since \(a_i\) is a self-adjoint idempotent,
\(\|a_i\|_2^2=a_i(1)=\beta_i\).
Separate the identity in the formula for \(p_i\) and apply
Cauchy--Schwarz to the remaining terms. By \eqref{eq:weighted-bad-share},
\[
 p_i\le\beta_i+\sqrt{\beta_i u/(1-u)}\le b^{-8}
\]
for sufficiently large selected squares, since \(\beta_i\le b^{-8}/2\),
\(q_{\min}\ge b^{10}\) and \(f\le b\).

\paragraph{The projected diagonal bound.}
\teamanchor{proof.weighted-diagonal-bad}%
Taking \(a\) to be the unit in the induced block estimate gives
\(\|P^*P-I_C\|_{\rm HS}^2\le C\min_i\Delta_i\).
The two commuting bad projections dominate \(I-P_W\) in sum.
The positive-part trace is the supremum of the trace against
\(0\le Q\le I\), and the trace norm of \(I_C-P^*P\) is at most
\(\sqrt C\|I_C-P^*P\|_{\rm HS}\). Hence
\begin{equation}
 \label{eq:weighted-projected-deficit}
 d_+:=\frac{\operatorname{tr}((I_C-P^*P_WP)_+)}{C}
 \le \min_i\sqrt{\Delta_i}
      +\sum_{i=0}^1\bigl(p_i+\sqrt{p_i\Delta_i}\bigr)
 \le 5b^{-4},
\end{equation}
where the last bound uses \(p_i\le b^{-8}\) and the supplied
\(\Delta_i\le b^{-12}\) for both factors. For sufficiently large
selected squares it also suffices that \(\min_i\Delta_i=o(b^{-3})\)
and \(\max_i\Delta_i=o(b^5)\), and that
\(\mathscr E_{\rm all}(G)\le(1+o(b^{-3/2}))\mathscr M\).
The bound \(p_i\le b^{-4}\) would give the sufficient condition
\(\max_i\Delta_i=o(b)\). All asymptotics refer to the fixed precision
and quadratic prime. In either case the diagonal deficit, half the
excess covariance and the ramification loss then sum to at most
\(b^{-3/2}\). The same preimage estimate selects the required number
of translates and the same short-row construction applies.\footnote{
\weightedpin{grh-compact/src/New/Filling/Moments.lean}{186-L218}
{Complementary orbital estimate for each factor} and
\weightedpin{grh-compact/src/New/Filling/FillingTheorem.lean}
{17-L47}{Application to the ramifying operator and projected rows}.}

\paragraph{Selection using the projected maps.}
\teamanchor{proof.weighted-frame-compression}%
Use \(d_+\) from \eqref{eq:weighted-projected-deficit} and
put \(\epsilon=(R\operatorname{tr}(H^2)/C^2-1)_+\).
If \(d_+<1\), then \(R>0\), since \(P_WP=0\) would give \(d_+=1\).
Let \(T_s=P_W\rho(s)P\), and for \(F\subseteq W\) let \(Q_s\) project
onto \(S_s=T_s^{-1}F\). In the
\hyperlink{proof.cut-frame-energy}{positive trace compression estimate}
\eqref{eq:positive-trace-compression}, take \(L_j=T_s\), \(Z=F\)
and \(A=|G|H\). Centrality of \(P_W\) gives
\(T_s^*T_s=P^*P_WP\) for each \(s\), so \(n_j=Cd_+\). Dividing by
\(|G|\) and using the arithmetic--geometric mean gives
\[
 |G|^{-1}\sum_s\dim S_s-Cd_+
 \le \sqrt{\dim F\,\operatorname{tr}(H^2)}
 \le C\sqrt{(1+\epsilon)\frac{\dim F}{R}}
 \le\frac C2\left(1+\epsilon+\frac{\dim F}{R}\right).
\]
If $k_{\rm ram}=C^{-1}\dim\ker\mathcal R$, then
$\dim \mathcal R^{-1}S_s\le \dim S_s+Ck_{\rm ram}$ by rank--nullity.
Thus the needed hypothesis is the combined bound
$d_++\epsilon/2+k_{\rm ram}\le\eta_0$, where we take $\eta_0=b^{-3/2}$.
Equations \eqref{eq:weighted-projected-deficit} and
\eqref{eq:weighted-projected-covariance} and the \hyperlink{proof.weighted-central-orbits}{central orbit bound}
$k_{\rm ram}\le f/(N-1)$ imply it for sufficiently large $x$.
Here $N-1\ge b^9$. On the ramified columns the nonzero scalars from
normalization, the seesaw and Theorem~A preserve these preimages.
Consequently for every $F\subset W$ there is a translate with
\begin{equation}
 \label{eq:weighted-preimage}
  \dim(B_s\mathcal R)^{-1}F
              \le C\left(\frac{1+\dim F/R}{2}+\eta_0\right).
\end{equation}

Start with $F=0$. While the codimension $D_F=R-\dim F$ exceeds
$4\eta_0R$, add the image of the selected translate. The rank gain
is $\Delta\ge C(D_F/(2R)-\eta_0)>0$. The potential
$5R\log(1+D_F/(\eta_0R))$ decreases by at least
$5R\Delta/(D_F+\eta_0R)\ge C$. Thus a set of translates $\mathcal S$
has $|\mathcal S|C\le5R\log(1+1/\eta_0)$ and images spanning a space
of dimension at least $(1-4\eta_0)R$. Since $\eta_0=b^{-3/2}$, the budget is at most $5R\log(1+b^2)$.

\paragraph{The common family of rows.}
\teamanchor{proof.weighted-rows}%
Choose $n$ independent rows $y=(y_1,\ldots,y_n)$ from the short images
of Theorem~A which span this space. Record the index allowances for
these exact row preimages, and use Hadamard on their untwisted lifts.
This gives one family with
\begin{equation}
 \label{eq:weighted-selected-rows}
  n\ge (1-4\eta_0)R,\qquad
  \delta\le5R\log(1+b^2)\mathcal B,\qquad
  \operatorname{covol}_\sigma^2(\Theta_1y)\le e^V
    \quad\text{for every }\sigma,\qquad V=nh .
\end{equation}

\subsection*{Twists and determinants (Theorem C)}

\teamanchor{proof.weighted-fit}%
\teamanchor{proof.weighted-fit-local-output}%
We first derive the comparison at a prescribed coefficient place.
The determinant argument will fix a single untwisted frame and
untwisted multipliers, then choose the place from their product.
Write \(\iota=(\ell,j_\iota)\) for a rational prime and a field
isomorphism \(j_\iota\colon\mathbf C\simeq\overline{\mathbf Q}_\ell\),
with \(|\ell|_\ell=\ell^{-1}\). The local construction for \(D\) at
this pair gives one twist \(\theta\) of added level norm prime to
\(\ell\). Each \(\theta^j\), \(1\le j\le m\), has the same added
conductor support as \(\theta\). Use this same twist also at the
paired embedding \(j_\iota\circ\operatorname{conj}\). The
\hyperlink{proof.weighted-fit-ray-primes}{ray construction} makes
the inverse fitted avatar trivial on the decomposition groups
above \(\ell\), so its trial powers preserve the local row there.
Put \(B_0=\mathcal B/(m\kappa)\). The comparisons use nonnegative
allowances \(\beta_{j,\ell}\) for the two determinant tests and
\(v_\ell=9c_\ell\log\ell\) for the ramified ratios, where
\[
 c_\ell=v_\ell(G)+2^{-\ell},\qquad
 G=\prod_{v\ {\rm chosen}}|\operatorname{GL}_{rt}(\mathbf F_{q_v})|.
\]
In \(v_\ell(G)\) we use the normalized additive valuation. The
allowances are fixed before choosing any twist. Their sums over
finite sets of primes are bounded by \(B_0\) at each stage for
\(\beta_{j,\ell}\), and by \(x^{999}\) for \(v_\ell\).

\begin{enumerate}
\item Write \(\vartheta=\theta^j\) and \(Z_{j,\vartheta}\) for
the image of the single lift on good forms on \(U(W_j)\). Let
\(g_{j,\vartheta}\) be the comparison map from the tables, and
\(k_j\) its positive integral hook factor. On \(Z_{j,\vartheta}\) the
specific compressed Gram operator
\[
 K_{j,\vartheta}
   = k_j\,\operatorname{pr}_{Z_{j,\vartheta}}\,
             g_{j,\vartheta}g_{j,\vartheta}^*\big|_{Z_{j,\vartheta}}
\]
satisfies
\[
 |j_\iota(\det K_{j,\vartheta})|_\ell\le
        \exp(R r\beta_{j,\ell}).
\]
Here \(\dim Z_{j,\vartheta}=R\), and the unitary source rank is \(r\).

\item The \(\kappa/2\) auxiliary lines at stage \(j\) use the lower
members \(p=p_j(z)\) of the unsigned character pairs of lift rank \(2z\),
where \(7m/2<z\le 7m/2+\kappa/2\). They use one fixed character as
the lift rank varies.
Let \(Z_{p,\vartheta}\)
be the line of the lift of the constant function from \(U(1)\), and
let \(h'\) be the number of \(U(1)\)-classes at the chosen level.
Use the two integral test forms \(\Xi,\Xi'\) from the jet test,
at \(j_\iota\) and \(j_\iota\circ\operatorname{conj}\) with the same
\(p,\vartheta\) and the same fixed height witness for the test.
Their projected pairing satisfies
\[
  \left|j_\iota\!\left(
     h'\langle\Xi,\operatorname{pr}_{Z_{p,\vartheta}}\Xi'\rangle
                  \right)\right|_\ell\le\exp(\beta_{j,\ell}).
\]
Both forms are integral at every embedding above \(\ell\). Their
Hecke tests use the tuple idempotent labelled \(X-1\) at the chosen
places and the kernel of the descent map from \(U(2z)\) to \(U(1)\).

\item For the ramified pairs and signs used in this comparison,
write \(\rho_{p,s}(\vartheta)\) for the twisted over untwisted
Hecke-value ratio. The local fit gives
\(\log|j_\iota(\rho_{p,s}(\vartheta)^{-1})|_\ell\le v_\ell\).
The list consists of the auxiliary flat pair \(q_2\) with sign \(+1\)
in the jet test, the lower members of the unsigned stage-\(j\) pairs
up to \(7m/2\) with sign \(+1\), the lower members of the two signed
families up to \(m\) and \(m/2\) with signs \(+1\) and \(-1\) respectively,
and the lower members with sign \(+1\) in the scalar square
calculations for the chain and the completion. Each of the latter
calculations uses one product over the lower points of its specified
ramified pair, at its own specified stage.
\end{enumerate}
The table map and the jet test are the ones in the single-lift proof.
The square calculations use the grid-pair and triangle-pair
characters on \(D\).\footnote{\weightedpin
{grh-compact/src/New/Rows/LiftVolumeHeight.lean}{47-L73}{Compressed Gram determinant},
\weightedpin
{grh-compact/src/New/Rows/HeckeValueHeight.lean}{154-L201}{Two coefficient readings of the jet test},
\weightedpin
{grh-compact/src/New/Rows/HeckeValueHeight.lean}{223-L237}{Rank-one and ramified factors of the product},
\weightedpin
{grh-compact/src/New/Exchange/RowLadder.lean}{16-L25}{Grid pair},
\weightedpin
{grh-compact/src/New/Exchange/TriangleArrows.lean}{125-L135}{Triangle pair}, and
\weightedpin
{grh-compact/src/New/Exchange/CostOfSquares.lean}{54-L63}{Square factors}.}

For the determinant comparison we use the
\hyperlink{proof.weighted-fit-joint-det}{joint fitted product}
and the signed ramified bounds in item~3. The checks at one
prescribed \(\iota\) supply both with its one twist. The
\hyperlink{proof.weighted-selected-coefficient}{norm selector}
uses them only if the absolute rational norm is less than one.

The local construction uses the \hyperlink{proof.weighted-fit-rows}
{normalized Galois rows} and the
\hyperlink{proof.weighted-fit-integral-rows}{balanced lengths}.
With the prescribed finite length \(M_\ell\), the reciprocity and
specialization below give lengths at most \((10fr_i+1)c_\ell\)
and the advertised signed inverse-ratio bounds on the same
first \(m\) powers.

\paragraph{Paired ray coordinates.}
\teamanchor{proof.weighted-fit-ray-primes}%
Fix a coefficient prime \(\ell\), a ray order \(q_0=\ell^N\) with \(N\ge1\), and
a conjugation-stable finite set \(T_0\) of places whose decomposition
groups the next coordinate must kill. Include the places over \(\ell\).
Suppose the proposed Frobenius is in the closed normal subgroup
generated by inertia at finitely many places away from \(T_0\).
Include these places and the other exclusions in a finite set \(B\).
The \hyperlink{proof.twist-ray-quotient}{earlier ray quotient} can be
formed after enlarging \(T_0\) to a conjugation-stable set \(T\) which
kills the ideal class group. Choose the extra prime-ideal generators
outside \(B\) and its conjugate, using Chebotarev in the Hilbert class
field. Adjoin \(\mu_{q_0}\) and the \(q_0\)-th roots of a finite set of
generators of the \(T\)-units. At each place supporting the proposed
Frobenius these are roots of units of order prime to the residue
characteristic. The resulting Galois extension is unramified there.
Thus its trivial Frobenius prescription and the finite-quotient
prescription for the proposed Frobenius agree on the intersection.
The Chebotarev set has positive Dirichlet density over \(E\).
The places of absolute degree at least two have density zero,
since their sum of \(q_P^{-s}\) is bounded at \(s=1\).
Thus one can choose a place \(P\) of absolute degree one outside
the exclusions and the ramified rational primes. It is split
over the real field and satisfies \(q_P\equiv1\bmod q_0\).

The ray group of the localized ring of integers at \(P\) is a
quotient of \(k_P^\times\) by the images of the \(T\)-units. Their
images are \(q_0\)-th powers, so the cyclic quotient of order \(q_0\)
descends. Subtract its conjugate. Since \(P\ne\overline P\), this
gives a character \(h_P\) onto \(\mathbb Z/q_0\mathbb Z\) on inertia at
\(P\), the inverse character at \(\overline P\), and a trivial
character on the decomposition groups at \(T\). Normalize the
coordinate to give \(1\) on the chosen tame generator \(\delta_P\).
Change the new Frobenius lift by a power of \(\delta_P\) to give
\(h_P(\phi_P)=0\). It has the same image in the finite quotient used
for its prescription, since \(P\) can be chosen outside that quotient's
ramification. The \hyperlink{proof.twist-picard-ray-map}{earlier
transfer calculation} shows that the character kills the real
ideles and gives a character of the norm-one torus with the
specified inverse avatar. The construction works also at \(\ell=2\).
\footnote{\weightedpin
{grh-compact/src/New/Selmer/Corepresentable.lean}{74-L132}{The new ray coordinate and its Frobenius},
and \weightedpin
{grh-compact/src/New/Selmer/TwistTheorem.lean}{31-L40}{The helper prime}.}

\paragraph{Reciprocity for a tied Selmer group.}
\teamanchor{proof.weighted-fit-paired-reciprocity}%
Use the finite local conditions of one balanced row \(\rho\) on
\(O^n\), where \(O\) is a coefficient DVR finite over
\(\mathbb Z_\ell\). Write \(c\) for complex conjugation on the
absolute Galois group \(G_E\), and \(\chi\) for the cyclotomic
character of \(G_E\). Let \(a\mid\ell^{M_\ell}\) be a nonunit
and \(V=(O/a)^n\).
At the conjugation-stable set of fixed places \(S\) use the prescribed
modules of local cocycles containing coboundaries. At an existing
auxiliary pair use the same slope \(t\) for the two relations
\[
 z(\delta_P)=t z(\phi_P),\qquad
 z((c\delta_P)^{-1})=t z((c\phi_P)^{-1}).
\]
Here \(cg\) means conjugation of \(g\) by complex conjugation.
The ties contain coboundaries. Tame inertia acts trivially at
an auxiliary place, and a nonzero slope is allowed only where
both Frobenius readings act trivially on \(V\).
Away from these places require unramifiedness. Normalize the
global cocycles by \(z(\gamma)=0\), where the distinguished inertia
element \(\gamma\) has \(\chi(\gamma)=1\),
\(\rho(\gamma)=sI\), \(\rho((c\gamma)^{-1})=sI\), and \(s^2-1\)
is a unit. This normalization is available and unique on each
cohomology class, since the coboundary at \(\gamma\) is
\((s-1)\) times its defining vector. Denote the resulting module by
\(X\).

Assume \(X\) contains a cocycle \(z\) having a unit coordinate in
some value. Suppose that at every admissible new pair the two
Frobenius values of this \(z\) are the two tame values of a single
cocycle with the old conditions and ramification allowed at the new
pair. We show that \(\operatorname{length}_O X\le\beta\), the
balanced-row allowance. The contrapositive gives the new pair used
in the slope argument.

Take the images of the prescribed finite local cocycle modules in
local \(H^1\). If there is a choice of decomposition group in the
local length bound, use one attaining the least length at each
place. The global Selmer group with these local conditions contains
the image of \(X\). At an auxiliary place the local module of tied
cocycles injects into \(V\) by evaluation at Frobenius. Indeed the
tie determines the value at a generator of tame pro-\(\ell\)
inertia, wild inertia gives zero, and the representation is
unramified there. Write \(Y\) for the dual Selmer group with the
annihilator local conditions on
\(V^{\mathrm D}=\operatorname{Hom}(V,\mu_{\ell^\infty})\).
The Greenberg--Wiles count gives
\[
 \operatorname{length}_O X
 \ \le\ \operatorname{length}_O Y+\sum_{v\in S}
 \operatorname{length}_O Z_v-(|S|+f)\operatorname{length}_O V ,
\]
where \(Z_v\) is the chosen local cocycle module and \(f\) is the
number of complex places of \(E\). The finite-place contribution
with local cocycle module \(Z_v\) is
\(\operatorname{length}_O Z_v-\operatorname{length}_O V\), after
subtracting the local coboundaries and invariants. This is
nonpositive at the tied places. Global invariants vanish on both
modules by the scalar at \(\gamma\). The complex places give
\(-f\operatorname{length}_O V\). These are the terms obtained by
taking orders in Poitou--Tate and inserting the local and global
Euler characteristics. In the length identity the local Euler
characteristic contributes \(-[E_v:\mathbb Q_\ell]\operatorname{length}_O V\)
at each \(v\mid\ell\), and the global Euler characteristic
contributes \(f\operatorname{length}_O V\) with the opposite
sign. Their sum is the stated complex-place term because
\([E:\mathbb Q]=2f\). This applies also at \(\ell=2\) over this
totally imaginary field.\footnote{\href
{https://math.stanford.edu/~conrad/BSDseminar/refs/MilneADT.pdf}
{Tate's duality and Euler characteristics in Milne, \emph{Arithmetic
Duality Theorems}, 2nd ed., I.2.8, I.4.10, and I.5.1}.
Compare the Selmer order calculation in
\href{https://annals.math.princeton.edu/1995/141-3/p01}
{Wiles, Proposition 1.6, p.~472}.}

Let \(U\) be the common kernel of the actions on \(O^n/\ell^{M_\ell}\),
its conjugate, and \(\mu_{\ell^{M_\ell}}\). Let \(H\) be the closed
normal subgroup generated by the commutators \([g,\gamma]\) and their
conjugates under \(c\). The two scalar conditions give \(H\subset U\).
Let \(\delta_0\) be the tame generator at the helper prime. It lies
in \(U\) as well. For each \(h\in H\) take a new pair with Frobenius
arbitrarily close to \(\delta_0h\), using the preceding ray
construction. The target is in the normal closure of inertia at the
helper and the fixed places away from \(\ell\). The coordinate just
constructed kills the groups at the nonhelper auxiliary places.
The old coordinate values at the target are \(1\) at the helper
and zero at the other places, since old coordinates kill commutators.
Require the new Frobenius to lie in the open subgroup \(U\).
The pair may therefore be assigned any slope in this row and zero
in the other rows.

Write \(x\) for a dual cocycle representing an element of \(Y\).
Identify the roots of unity of order \(\ell^{M_\ell}\) with their
additive cyclic group via a fixed generator. On \(U\) the
restrictions of \(x,z\) are homomorphisms, and the restriction of \(x\)
is independent of its cocycle representative. At the new pair,
\(x\) is unramified. By the hypothesis on \(z\), let \(z'\) be a
relaxed cocycle with the two tame values equal to the two
Frobenius values of \(z\). The global sum of local invariants of
\(tz'\cup x\), for \(t\in O\), vanishes. All terms outside the
new pair vanish by the local conditions and unramifiedness.
With a normalized tame generator \(\delta\) the remaining term
at \(P\) is \(x(\phi_P)(t z'(\delta))\), up to the common sign
of the local cup convention. Indeed \(x\) is unramified and both
actions at the new place are trivial. The local invariant of
the cup of an unramified character with the Kummer class of a
uniformizer evaluates that character at Frobenius, with the
sign determined by the order of the factors.
At \(\overline P\) the corresponding normalized generator
is \((c\delta)^{-1}\) since \(c\) inverts the roots of unity.
The actual chosen generator and its conjugate inverse are the
same multiple of these normalized generators in tame
pro-\(\ell\) inertia, so the relation holds on the two actual
values together. Since \(z((c\phi_P)^{-1})=-z(c\phi_P)\),
prescribing the values of \(x,z\) on the Frobenius and its
conjugate yields
\[
 x(\delta_0 h)(t z(\delta_0 h))
   -x(c(\delta_0h))(t z(c(\delta_0h)))=0\qquad(h\in H).
\]
Both prescriptions use finite quotients because the cocycles
have finite target. At those quotients the change of Frobenius
lift by the new inertia is trivial.

For \(h,h'\in H\) put
\[
 B_t(h,h')=x(h')(t z(h))-x(ch)(t z(ch')) .
\]
Taking the mixed difference of the last equality at \(hh',h,h',1\)
gives \(B_t(h,h')+B_t(h',h)=0\) by additivity on \(U\).
Let \(h_\gamma=\gamma h\gamma^{-1}\). Equivariance on \(U\) gives
\[
 B_t(h_\gamma,h')=B_{st}(h,h'),\qquad
 B_t(h',h_\gamma)=B_{s^{-1}t}(h',h).
\]
Indeed \(x(ghg^{-1})(v)=x(h)(\chi(g)\rho(g)^{-1}v)\) and
\(z(ghg^{-1})=\rho(g)z(h)\). Apply antisymmetry with
\(h_\gamma\) and \(st\), then subtract antisymmetry with \(h\)
and \(t\). Since \(B_t\) is additive in \(t\), it follows that
\(B_{(s^2-1)t}=0\). The unit condition gives \(B_t=0\) for all
\(t\), including at coefficient prime \(2\).

The values \(z(H)\) span \(V\) over \(O\). In fact
\(z([g,\gamma])=(1-s)z(g)\), so one such value has a unit
coordinate, and its orbit spans \(V\) because \(\rho(G_E)\)
spans the full matrix algebra over \(O\). The equality \(B_t=0\)
now gives an \(O\)-linear map \(\eta:V\longrightarrow(V^{\mathrm
D})^c\) with \(\eta(z(h))=x(ch)\). A linear relation among
the \(z(h)\) gives a linear combination of the \(x(ch)\)
vanishing on every \(t z(ch')\) by \(B_t=0\). This proves
that the combination vanishes everywhere. The same conjugation
identities show that \(\eta\) is \(G_E\)-equivariant. If \(\eta=0\),
then \(x(H)=0\). Comparing \(x(u)\) with
\(x(\gamma u\gamma^{-1})\) for \(u\in U\) gives \(x(U)=0\),
since \([u,\gamma]\in H\) and \(s-1\) is a unit. The class
inflates from \(H^1(G_E/U,V^{\mathrm D})\). In this quotient
\(\gamma\) is central and acts by \(s^{-1}\) on \(V^{\mathrm D}\).
For any cocycle on the quotient, comparing its values on
\(\gamma g\) and \(g\gamma\) writes it as a coboundary,
using \((s^{-1}-1)^{-1}\). Thus \(Y\) injects into the module
of equivariant maps.

Let \(g_*\) be the element from the balanced-row bound, with
\(\rho(g_*)=\alpha I\) and \(\rho((cg_*)^{-1})=\alpha'I\).
Equivariance gives \(\eta((\alpha-\chi(cg_*)\alpha')V)=0\).
An equivariant map is determined by its value on one vector with
unit coordinate, because the orbit of that vector spans \(V\).
The length of their module is therefore at most
\(n\,v_O(\alpha-\chi(cg_*)\alpha')\). The balanced inequality is
\[
 \sum_{v\in S}\operatorname{length}_O Z_v
   +n\,v_O(\alpha-\chi(cg_*)\alpha')
 \ \le\ (|S|+f)\operatorname{length}_O V+\beta .
\]
Combining the two bounds proves
\(\operatorname{length}_O X\le\beta\). This supplies the
unseen-class implication in the finite-length induction for the
actual crystalline relation modules. The ray order can still
be chosen after the bound on the number of auxiliary pairs.
\footnote{\weightedpin
{grh-compact/src/New/Selmer/Corepresentable.lean}{222-L256}{The balanced bound and the ties},
\weightedpin
{grh-compact/src/New/Selmer/TiltingLemma.lean}{89-L198}{The compact unseen-class implication},
and \weightedpin
{grh/src/New/Congruences/Selmer/Mirror.lean}{541-L701}{The corresponding equivariant map}.}

\paragraph{The slope at two successive coefficient lengths.}
For one row and a new admissible pair let \(Y\) be the relaxed
module modulo \(a=\varpi^s\), and \(A,B:Y\longrightarrow(O/a)^{2n}\)
its paired tame and Frobenius evaluations. The old tied module
is \(\ker A\), because vanishing on both tame generators restores
unramifiedness at the two new places. Let \(D\) be the sum of the kernels of \(uA-B\)
over \(u\in O\), so that \(B(D)\subseteq A(D)\). The inverse
image \(P\) of \(A(D)\) in \(O^{2n}\) is free of rank \(2n\).
Choose a lift \(\sigma:P\to D\) of reduction through \(A\), and an
endomorphism \(T\) of \(P\) lifting \(B\sigma\). Both lifts exist
because \(P\) is free. If \(u\) avoids the roots of the reduction
of \(\det(U-T)\), then \(u-T\) is invertible on \(P\).
Applying \(uA-B\) to \(\sigma (u-T)^{-1}P\) proves that
\((uA-B)D=A(D)\). Since its kernel on \(Y\) is contained
in \(D\), it follows that
\[
 \operatorname{length}_O\ker(uA-B)
 \le \operatorname{length}_O(\ker A\cap B^{-1}A(Y)).
\]
There are at most \(2n\) exceptional residue classes at each length.
The balanced residue-field bound allows one \(u\) that is a unit
and works at both \(s\) and \(s-1\). The new slope in this row is
\(u^{-1}\).

For a fixed finite list of rows, apply induction on \(s\), starting
at \(s=0\) with the helper prime and zero slopes. A pencil with \(d\)
pairs has length at most the length at the original places plus
\(2nsd\) in each row, by evaluating at the paired tame generators.
Thus a bound for the number of pairs at \(s-1\), independent of
ray order, gives such a bound \(C_s\) for the number of pairs plus
the sum of the lengths at \(s\) on a pencil already working at
\(s-1\). Among pencils satisfying that sum bound which still
work at \(s-1\), minimize the sum of lengths at \(s\).
If some row exceeds its \(\beta\), its module contains a cocycle
with a unit coordinate, since otherwise division by \(\varpi\)
would put the whole module in the one at \(s-1\) by homogeneity
of the cocycle equations and local relations. The paired
reciprocity implication supplies a new admissible pair with
\(B(\ker A)\not\subseteq A(Y)\). The new slope makes the
length drop strictly at \(s\) and prevents an increase at \(s-1\).
Give the other rows slope zero at this pair, so they retain
the previous unramified condition there. The strict integral
length drop pays for adding one pair and preserves \(C_s\),
contrary to minimality. This proves the simultaneous tied bounds
at \(s\) with a bound for the number of pairs independent of the
ray order. The argument reaches \(s=v_O(\ell^{M_\ell})\).
\footnote{\weightedpin
{grh-compact/src/New/Selmer/TiltingLemma.lean}{204-L224}{The slope estimate},
and \weightedpin
{grh-compact/src/New/Selmer/TwistTheorem.lean}{31-L94}{Induction with one slope at two lengths}.}

\paragraph{Finite interpolation of the signed values.}
\teamanchor{proof.weighted-fit-signed-values}%
Fix a critical pair \(p=(\nu,z)\) from the ramified signed list.
Its tame part is nontrivial at each fixed chosen place and trivial
away from those places. For \(z\ge2\), put \(n=2z\), \(k_\tau=c_\tau-z\) and
\(k=\sum_\tau k_\tau\), where \(-c_\tau\) is the weight of
\(\nu\). Fix the modulus \(I\) of paired auxiliary places,
write \(q=N(I)\), and let \(G\) be the common multiple of the
orders of the class stabilizers. The section at each added place
requires the first coordinate to equal \(1\) in the residue
field. Keep this section even if the character on that modulus
is imprimitive. For every \(\ell\nmid q\) the fixed integer
\[
             D_z=g^{2fn}(gd)^k
\]
clears its point kernel in the integral Weyl gauges
(\(f=\#\Phi\)). In fact the Fourier expansion of the residue
indicator on a box of \(n\) entries puts \(g^{2fn}\) times the
normalized vacuum section in the standard-data module of level
\(g\). The added part of the modulus has norm prime to \(\ell\).
After moving to integral frames, the
\hyperlink{proof.weighted-local-source-gauge}{integral CM jet}
at scale \(gd\,\mathrm{per}\) applies in degree \(k\).
The point kernel uses \(\bar\pi_G\,\mathrm{per}\), and
\(\bar\pi_G\) is integral, so multiplying by \((gd)^k\)
suffices. Returning to the original frames gives the asserted
Weyl gauges.

At the chosen modulus write the lift as
\(\Xi_\vartheta=\sum_a\vartheta(u_a)^{-1}K_a\), summing over the
source classes and the residue lifts. Its squared norm is a sum
of \(\vartheta(u_a/u_b)\langle K_a,K_b\rangle\), with the inner
product linear in the second entry. The common class-level sum
over the target of the pairings of the cleared kernels is integral
after multiplication by \(G\). Indeed the Weyl pairing is integral
in the point gauges and the only class weights are the reciprocal
stabilizer orders. Integrality at the conjugate reading covers the
conjugated entry of the pairing. Division by \(q^{n-1}\) is a unit
at the same places. Thus
\[
       F(I,\vartheta)=G D_z^2\|\Xi_\vartheta\|^2/q^{\,n-1}
\]
is the evaluation of a single finite sum in the ray group ring
with coefficients integral away from \(q\). In the
\hyperlink{proof.weighted-rank-one-norm}{rank-one norm calculation}
the modulus fixes the source support and the target index. At an
added place the support in source units is the kernel of reduction,
on which every allowed tame character is trivial. The same local
calculation therefore contributes \(q_v^{n-1}(1-q_v^{-n})\)
also when the character at \(v\) is imprimitive. The spherical
Euler factors and the factors at the fixed places are those in
the cited calculation. After dividing by \(q^{n-1}\) one obtains
a fixed nonzero multiple of the incomplete Hecke value at \(z\).
For exact additional conductor the quotient by the trivial value
is \(\rho_{p,+}(\vartheta)\).

\smallskip
\noindent\textit{The rank-two determinant period.}
\teamanchor{proof.weighted-fit-signed-determinant}%
For \(z=1\), the test is the determinant period \(\Lambda_\vartheta\)
on \(U(2)\). Put \(\alpha=\nu\vartheta\upsilon_{U(1)}\) and
\(Q(X)=\prod_\tau\det(X_\tau)^{c_\tau-1}\), of total degree
\(h=2\sum_\tau(c_\tau-1)\). Each \(c_\tau\ge1\).
Here is the clearing factor for this test. Set \(\delta=g^2\)
and let \(v_0\) be the scalar that normalizes the rank-one vacuum.
The scalar
\[
              N_p=\delta^{8f}\,v_0^4
                        (\delta d\,\mathrm{per})^{h}
\]
is nonzero and independent of the modulus and the reading. The
subscript \(p\) indexes the critical pair \((\nu,1)\).
The four entries of the residue matrix give a Fourier
denominator dividing \(\delta^{8f}\) at \(\ell\nmid q\), since
the modulus ideal contains \(q\delta\). Multiplying the vacuum
by \(v_0^4\) and by this integer puts it in the standard-data
module of level \(\delta\). Integral finite translates preserve
that module. Pairing its Taylor series at scale
\(\delta d\,\mathrm{per}\) with the integral homogeneous polynomial
\(Q\) multiplies the kernel by \((\delta d\,\mathrm{per})^h\).
Thus \(N_p\) times the kernel is integral in every such reading
at integral arguments.

For \(x=(u,A)\), where \(u\) represents a rational double coset
at maximal compact level and \(A\) is an invertible residue matrix
for the fixed modulus, write \(g_x=u\widetilde A\). Let \(K_1\)
be the untwisted kernel for the indicator of the identity matrix
at this modulus, and set
\[
    b_x=\frac{G}{|\operatorname{Aut}(u)|}\,N_p K_1(g_x).
\]
These numbers are integral away from \(q\). For each coefficient
prime, rational translation can move \(g_x\) into the integral
subgroup there. For definite \(U(2)\) this approximation follows
already from the Cayley parametrization: its inverse is defined
on a dense open subset of each required local coset, and weak
approximation on the space of skew-Hermitian matrices gives a
simultaneous rational parameter. The determinant character from
a rational translate cancels the determinant weight of \(Q\);
the kernel is unchanged. The character of \(K_1\) is
\(\nu\upsilon_{U(1)}\), with its original weight and tame part.
Expanding the residue character as translates of the identity
indicator now gives
\[
       G N_p\Lambda_\vartheta
               =\sum_x b_x\,\vartheta(\det g_x).
\]
In this identity the modulus is fixed and the character varies.
The square \(F(I,\vartheta)=|G N_p\Lambda_\vartheta|^2\) has
coefficients \(b_x\overline{b_y}\) on
\(\det(g_x)/\det(g_y)\). Each conjugate is integral because the
preceding test holds for all coefficient embeddings.

For an exact twist the analytic input here is the
\hyperlink{proof.weighted-complementary-period-norm}{normalized
complementary exchange} on \(W=I_2\), with larger space \(I_2\)
and smaller space \(0\). The rank difference is two and the
half-difference is one. The finite partners are
\(\chi(\det X)\), extended by zero on singular matrices, and
the function \(1\) on a single point. On the fixed places
\(\chi\) has the original nontrivial tame component. On the
disjoint added places the exact twist has nontrivial tame
component. At each real place the exchange weight is
\((c_\tau-1)+1=c_\tau\ge1\), as required in rank two.
The partner identity follows from the
block partner formula with the rectangular block empty.
The real polynomial on the large side is exactly \(Q\);
on the small side the empty determinant is one.
Consequently that exchange, evaluated at the identity, gives
\[
      \gamma_\vartheta\Lambda_\vartheta=1,\qquad
      |\gamma_\vartheta|^2=
          C\big/\Pi_{S_\vartheta}((\nu\vartheta)^{-1},1).
\]
The mass of the zero-dimensional quotient is absorbed in the
same fixed constant \(C\) by normalizing the first identity.
Here \(S_\vartheta\) consists of the fixed chosen places and
the added ones. By definition
\(\Pi_{S}(\xi,1)=(\prod_{v\in S}q_v)\,|L^{S}(1,\xi)|^2\),
where the omitted Euler places include the conjugates.
Dividing by the identity at the trivial twist cancels the mass,
\(G\), \(N_p\) and the contribution of the fixed conductor.
This proves
\[
       F(I,\vartheta)/F(\varnothing,1)
                     =q\,\rho_{p,+}(\vartheta)^2 .
\]
Evaluating the finite sums at the inverse character gives the
minus reading. The norm relation on enlarging an imprimitive
modulus follows from the residue computation below.%
\footnote{\weightedpin{grh-compact/src/New/Exchange/TriangleLevelExponent.lean}{64-L126}{The normalized one-column test},
\weightedpin{grh-compact/src/New/Rows/Carrier.lean}{154-L173}{The complementary space of rank zero},
\weightedpin{grh-compact/src/New/Selmer/ValueFamilies.lean}{30-L80}{Both interpolation families},
and \weightedpin{grh/src/New/UpperBound/Kernels/DefiniteKernelSize.lean}{114-L190}{The common degree-dependent CM scale}.}

Add an unramified pair \(P,\overline P\). Write
\(Y=(\nu\vartheta)_{\rm bc}(\varpi_P)\) for the unitary base-change
value. At \(z\ge2\) omission of the
two Euler factors gives
\((1-Yq_P^{-z})(1-Y^{-1}q_P^{-z})\).
At \(z=1\) the added residue condition is invertibility. On
integral matrices \(X\) one has
\[
 {\boldsymbol 1}_{\ker\overline X=0}
    =\sum_{U\subseteq\ker\overline X}
             (-1)^{\dim U}q_P^{\binom{\dim U}{2}} .
\]
This is the finite \(q\)-binomial theorem evaluated at one.
The subspace condition at \(U\) is given by translating
the split vacuum by the lattice which is the inverse image
of the prescribed residue column condition, using transpose
in the matrix coordinates for the right action. If
\(d=\dim U\), choose bases reducing it to the first \(d\)
coordinates and use a diagonal matrix with \(d\) entries
\(\varpi_P\) and the rest one. The split Weil action gives the
indicator of this column condition multiplied by \(A^d\),
where \(A=q_P\,\alpha_{\rm bc}(\varpi_P)\upsilon(\varpi_P)^{-2}
=q_PY\). Here \(q_P\) is the determinant factor of the
vacuum in rank two, and
\((\upsilon_{U(1)})_{\rm bc}=\upsilon^2\) by restriction to
norm-one ideles. Thus the local sum is the split Hecke
polynomial at \(A^{-1}\) applied to the vacuum. Integrating
over the class sum leaves every right translate unchanged.
The sum of the Hecke coefficients is
\(\prod_{d=0}^1(1-q_P^d A^{-1})\), again by the finite
\(q\)-binomial theorem. Squaring gives
\(\prod_{d=0}^1(1-Yq_P^{d-1})(1-Y^{-1}q_P^{d-1})\).
This residue-indicator calculation gives the norm relation
for the rank-two period.

Pull back each of these finite character sums by the
\hyperlink{proof.twist-picard-ray-map}{earlier transfer}, using
the inverse avatar. Every character argument becomes a
monomial in the installed coordinates, with powers read
modulo \(\ell^N\). In the signed value polynomial the coefficient
field can be enlarged to contain all the cleared coefficients.
At the chosen place \(w\nmid\ell\) the fixed tame character is
nontrivial on a generator. Its value \(\xi\) is a
nontrivial root of unity of order dividing \(q_w-1\).
Evaluating \((T^{q_w-1}-1)/(T-1)\) at \(1\) shows that
\(v_\ell(1-\xi)\le v_\ell(q_w-1)\), since all other
factors are integral. The same applies to the inverse
character. The helper Frobenius can be placed sufficiently
close to that inertia element for all the finite tests,
including the cyclotomic powers in the Euler factors.
Consequently the line polynomial for modulus consisting
only of the helper is bounded at \(u=1\) with the loss of
the helper Euler factors. At a later auxiliary place the
value line has trivial tame coordinate. All the nonhelper
coordinates are \(1\), so the inverse avatar
\(\chi_\vartheta\) reads \(u\) at its Frobenius.
For sign \(s=\pm1\) the base-change value to exponent
\(\epsilon=\pm1\) reads
\[
 j_\iota\bigl(\nu_{\rm bc}(\varpi_P)^\epsilon\bigr)
                    \chi_\vartheta(\phi_P)^{-s\epsilon}.
\]
The omitted factors are therefore at most four linear guards
\(1-\xi' u^{\pm1}\) with \(\xi'\) a unit, up to multiplication
by a unit. Include them in the one-variable specialization
on the helper line. The helper factors cost at most
\(4v_\ell(q_w-1)\), and one can choose the per-guard tolerance
so that eight times the coordinate bound times this tolerance
fits the additional slack. The coordinate bound precedes the
choice of ray order.
Then apply the common specialization to the resulting
value and Selmer polynomials.\footnote{\weightedpin
{grh-compact/src/New/Exchange/TriangleLevelExponent.lean}{64-L131}{The coefficient expansion and rank-one normalization},
\weightedpin
{grh-compact/src/New/Columns/ChainValuation.lean}{24-L83}{The determinant period and its split norm relation},
\weightedpin
{grh-compact/src/Lib/Weil/Unramified/HeckeTransfer.lean}{32-L77}{The split Hecke and vacuum action},
and \weightedpin
{grh-compact/src/New/Selmer/ValueFamilies.lean}{23-L80}{The two value identities}.}
The conductor powers are units at \(\ell\), so the inverse ratio
bounds apply to the same fitted character as the local determinant tests.

\paragraph{Rank-one avatars and the local operator.}
\teamanchor{proof.weighted-fit-rank-one-avatars}%
Let \(\eta_f\) be the finite part of an algebraic Hecke character of
\(E\) with infinity exponents \(m_\sigma\), so on \(E^\times\) it is
\(\prod_\sigma\sigma(x)^{-m_\sigma}\). For the selected coefficient
reading \(j_\iota\), put
\[
 \eta_\ell(\operatorname{Art}^{\rm geom}_E x)
   =j_\iota(\eta_f(x))
      \prod_{w\mid\ell}\prod_{\sigma\mapsto w}
            (j_\iota\sigma)(x_w)^{m_\sigma}.
\]
Here \(\sigma\mapsto w\) means that \(j_\iota\sigma\) induces \(w\), and
its value on \(x_w\) uses its continuous local extension. The right
side kills principal ideles. Choose an open subgroup of the finite
integral units on which \(\eta_f\) is trivial. The quotient by this
subgroup and by \(E^\times\) is a finite ray class group. A power of each
representative of the quotient is therefore principal up to that
subgroup, which shows that the values of \(\eta_f\) are algebraic and
lie in a single number field. The corrected character has values in a
finite \(\ell\)-adic field. Its values on integral ideles are units,
since the smooth character has finite image on the integral units
and each local algebraic factor there is a unit. The finite ray
quotient shows that all values of the corrected character are units.
It is continuous and takes values in a profinite group, so global
reciprocity gives the continuous character displayed above.

Write \(\eta_w^\iota(x)=j_\iota(\eta_f(\operatorname{idele}_w x))\).
For \(w\nmid\ell\) the avatar on the local Weil group is
\(\eta_w^\iota\) via geometric reciprocity. Its inertia has finite
image, and is trivial if \(\eta_f\) is unramified there. For \(w\mid\ell\)
write \(L=E_w\), \(f_w=[k_w:\mathbf F_\ell]\), and
\(\lambda(x)=\prod_{\sigma\mapsto w}(j_\iota\sigma)(x)^{m_\sigma}\).
Then \(\eta_\ell\circ\operatorname{Art}^{\rm geom}_w=\eta_w^\iota\lambda\).
Remove the finite character \(\delta\) of \(G_L\) corresponding to the
restriction of \(\eta_w^\iota\) on units, extended to have value \(1\)
at a chosen uniformizer \(\varpi\). Put \(\xi=\eta_\ell/\delta\).
With Conrad's arithmetic Artin map the algebraic character of
\(\xi\) on units is \(\lambda^{-1}\). Enlarge the coefficient field
to contain all the embeddings of \(L\). Proposition~B.4(i) of
Conrad gives crystallinity of \(\xi\). His convention counts the
cyclotomic character as having weight \(1\). Proposition~B.3
applied to \(\lambda^{-1}\) gives \(m_\sigma\) in that convention,
hence the jump at \(j_\iota\sigma\) is \(-m_\sigma\) here.
Proposition~B.4(ii), for the covariant crystalline module, gives
the scalar
\[
 \xi(\operatorname{Art}^{\rm ar}_w\varpi)^{-1}
              \lambda(\varpi)^{-1}=\eta_w^\iota(\varpi)
\]
for \(\varphi^{f_w}\). On the potentially crystalline periods the
operator for a Weil element \(s\) of geometric degree \(b\) is the
linear operator induced by \(\varphi^{b f_w}s\). For \(\xi\) it has
scalar \(\eta_w^\iota(\varpi)^b\) because \(s\) fixes the crystalline
periods. For \(\delta\), periods may be formed over the finite extension
killing \(\delta\). They have jump zero and on their span the operator
is \(\delta(s)\), since \(\varphi^{b f_w}s\) acts trivially on the
unramified period scalars. The extension cut out by \(\delta\) is
totally ramified, since \(\delta\) has value \(1\) on the symbol of
\(\varpi\), so these period scalars have unramified subfield \(L_0\).
Tensoring gives \(\eta_w^\iota(x)\) when
\(s\) has symbol \(\operatorname{Art}^{\rm geom}_w(x)\), with
\(\operatorname{ord}_w x=b\). This gives every degree-zero inertia
operator and the degree-one geometric Frobenius operator at every
inducing embedding. The crystalline module is the full rank-one
module when the smooth local character is unramified. With a ramified
smooth part, the monodromy is zero after the finite twist, and the
jumps are the same. This calculation applies over each completion
of \(E\), including those of residue characteristic two.%
\footnote{\href{https://math.stanford.edu/~conrad/papers/locchar.pdf\#page=34}
{Conrad, \emph{Lifting global representations with local properties}, Appendix B,
Propositions B.3 and B.4};
\weightedpin{grh-compact/src/Lib/ClassField/Global/Avatars.lean}
{239-L267}{Idelic formula for the avatar};
\weightedpin{grh-compact/src/Lib/PadicHodge/Crystalline/Defs.lean}
{60-L113}{Weil-pair degrees, jumps and linear operator}.}

For a norm-one character \(b\) with weights \(c_\tau\), write
\(b^{\rm bc}(x)=b(x/\bar x)\). Its algebraic character
\[
  \eta_{b;k,h}=b^{\rm bc}\upsilon^{-k}|\cdot|^{h-k/2}
\]
has exponents \(h+c_\tau-k\) and \(h-c_\tau\) at \(\tau\) and
\(\bar\tau\). Indeed \(b^{\rm bc}\upsilon^{-k}\) has angular exponent
\(2c_\tau-k\) and the norm factor gives radial exponent \(2h-k\).
Its smooth finite part is unramified at a split pair when \(b\)
is trivial on the integral norm-one units of that pair, since
\(\upsilon\) and the norm are unramified. For an upper auxiliary pair of ranks \(r',n\), the character
\(a=\upsilon_{U(1)}^{r_-}\) has weight \(r_-\), the number of negative
source signs. Its normalizing avatar is that of
\(\eta_{a^{-1};-r',0}\). Its jumps at \(\tau,\bar\tau\) are
\(r_- - r'\) and \(-r_-\), and its smooth parameter is unramified
at every finite place. The lower block uses
\(\eta_{(\mu_j^0)^{-1};0,(r'-t')/2}\), with \(t'=n-r'\).
It subtracts \((r'-t')/2+\epsilon_\sigma w_j\) from the source jump
when \(\mu_j^0\) has weight \(-w_j\). The same formula with the actual
point weight applies for the rank-one source. Its finite character
on a chosen-place inertia operator is the inverse of the tame part
of \(\mu_j^0\) through reciprocity. A trial stage character
\(\vartheta\), specialized to \(\theta^j\) in the fit, uses the
weight-zero avatar of \((\vartheta^{-1})^{\rm bc}\), with conductor
prime to \(\ell\).
Its inertial action at an added place is the character labelling
the inverse-twist determinant slot on the marked block. Conjugation inverts its
norm-one parameter, so it has the same value at the
inverse-conjugate inertia element. The cyclotomic character is the
avatar of \(|\cdot|\), with geometric Frobenius value \(q_w^{-1}\),
jump \(-1\), and trivial inertia on the degree-zero operators at
\(w\mid\ell\). Its inverse powers give the normalized string and its
first power gives the jump \(-u-1\) for the Tate dual \(V^\vee(1)\).%
\footnote{\weightedpin{grh-compact/src/Lib/ClassField/HeckeL/AutCharLFunction.lean}
{82-L94}{Two infinity exponents};
\weightedpin{grh-compact/src/New/Selmer/AptTupleDefs.lean}
{66-L87}{Normalization and first-entry scalar};
\weightedpin{grh-compact/src/New/Selmer/ValueFamilies.lean}
{86-L96}{The inverse tame avatar}.}

For the conjugate reading \(j_\iota^\#=j_\iota\circ\operatorname{conj}\),
write the \(\ell\)-adic embedding \(j_\iota\sigma\) as
\(j_\iota^\#\bar\sigma\). The same formula uses \(m_{\bar\sigma}\)
and the smooth value read by \(j_\iota^\#\). Both readings use the
same stage character and the same conjugate pair of local tests.

\paragraph{The Hecke transfer on the genuine auxiliary lifts.}
\teamanchor{proof.weighted-fit-stage-hecke-transfer}%
Use the split matrix calculation of
\hyperlink{proof.weighted-local-hecke-vacuum}{the spherical kernel}
on each auxiliary pair. For a genuine stage lift take source
\(W_j\) of rank \(r'=r\), target \(I_n\) with \(n=r+t\), and
\(\mu=\mu_j^0\vartheta\), where the trial stage character
\(\vartheta\) is specialized to \(\theta^j\) in the fit.
For a genuine sharp lift take the
constant source on \(I_1\), so \(r'=1\), with target \(I_{2z}\) and
the point character times the same \(\vartheta\). In both cases put
\(t'=n-r'\) and let \(r_-\) count the negative source signs.
The Weil characters in the source and target slots are
\[
 a=\upsilon_{U(1)}^{r_-},\qquad b=\mu\,\upsilon_{U(1)}^{n/2}.
\]
Here \(n\) is even. These are the characters of the actual
Petersson adjoint and of the rank-one kernel, respectively.
\teamanchor{proof.weighted-genuine-theta-kernels}%
Write the descent kernel as an operator \(K(g,h)\) on coefficient
vectors, with \(g\) on the source and \(h\) a definite target class.
Then the descent of a class function \(\Xi\) is
\(C\Xi(g)=\sum_h K(g,h)\Xi(h)/|\operatorname{Aut}(h)|\).
Inserting this into \(\langle C^*u,\Xi\rangle=\langle u,C\Xi\rangle\)
and moving the finite sum through the source pairing gives
\((C^*u)(h)=\int K(g,h)^*u(g)\,dg\) in the same class-sum metric.
In the stage kernel the adjoint uses the transposed conjugate
coefficient and the prescribed highest-vector datum, giving the
bilinear theta pairing. In the sharp kernel the source is constant
and the target coefficient is evaluated in the one-row jet. The matrix
datum at the finite places is named on the vacuum by a function
of the reductions at \(\mathfrak M Q\). Thus, at a split pair off
this support, the datum is the vacuum. Both characters are
unramified there. In particular the spherical calculation
applies at every off-level split pair needed for the auxiliary
prime tests and at the cofinal split places used to compare
Frobenius coefficients.%
\footnote{\hyperlink{proof.weighted-first-arrow-rallis}{Stage adjunction and datum};
\hyperlink{proof.weighted-rank-one-norm}{Rank-one kernel};
\weightedpin{grh-compact/src/New/Rows/StackedDeterminants.lean}
{133-L151}{The two Weil characters and the finite datum};
\weightedpin{grh-compact/src/New/Exchange/CharacterLiftNorm.lean}
{104-L143}{Stage adjunction on the same source and target};
\weightedpin{grh-compact/src/New/Exchange/CharacterLiftNorm.lean}
{169-L183}{Rank-one application}.}

\teamanchor{proof.weighted-genuine-theta-annihilator}%
At such a place put \(q=q_w\) and let
\[
 A_w=q^{r'/2}\upsilon(\varpi_w)^{2r_--r'},\qquad
 B_w=q^{n/2}\mu^{\rm bc}(\varpi_w),\qquad
 C_w=q^{t'}A_w/B_w.
\]
The identity \((\upsilon_{U(1)})^{\rm bc}=\upsilon^2\) follows
from \(\upsilon(\bar x)=\upsilon(x)^{-1}\); it gives the displayed
\(A_w,B_w\) from the two \(\operatorname{spaceChar}\) factors in the
split calculation. If \(P_w\) is the monic source polynomial and
\(\check P_w(T)=T^{r'}P_w(1/T)\), the eigenvalue of the
reciprocal Hecke operator on the actual lifted vector is
\[
 \check Q_w(T)=\prod_{i=0}^{t'-1}(1-A_wq^iT)\,
                            \check P_w(C_wT).
\]
For the sharp lift, \(\check P_w(T)=1-T\). Equality for all
\(T\) gives separately the eigenvalue of each spherical generator
when the lift is nonzero. A scalar multiple from the class-sum
realization cancels in this identity. The genuine lifts used in
the two determinant tests are precisely the lifts in the
\hyperlink{proof.weighted-first-arrow-rallis}{stage Gram pairing}
and the \hyperlink{proof.weighted-rank-one-norm}{rank-one norm}.
The extra factor in their reciprocal polynomial vanishes at
\(\xi_w=(A_wq^{t'-1})^{-1}\), so the test operator kills the
actual lifted vector, including a zero lift. The general attachment
supplies the Hecke tests on other systems in the deeper levels.

In monic form, if \(\operatorname{sr}_v\) multiplies all roots
by \(v\), the identity becomes
\[
 \operatorname{sr}_{A_w^{-1}} Q_w(X)
   =\prod_{i=0}^{t'-1}(X-q^i)\,
                           \operatorname{sr}_{q^{t'}/B_w}P_w(X),
 \qquad
 \frac{q^{t'}}{B_w}=q^{(t'-r')/2}\mu^{\rm bc}(\varpi_w)^{-1}.
\]
Here \(A_w^{-1}\) is the value of the upper normalizing
character \(\eta_{a^{-1};-r',0}\) from the rank-one calculation.
The lower scalar is the value of
\(\eta_{\mu^{-1};0,(r'-t')/2}\). Both exponents are integral
in the formula defining an algebraic character. Consequently
the untwisted lower normalization is fixed before fitting, and
its change at the trial stage character is the inverse weight-zero
avatar of \(\vartheta\). The string is the polynomial of the
inverse cyclotomic powers at a geometric Frobenius. Comparison
on the Frobenius filter gives the trace identity for the same
lower Hecke system, transferred from \(W_m\) to \(W_j\) as above.
The root identity holds at each prescribed off-level split
place on the given vector, independently of the approximation
used in the trace comparison.%
\footnote{\weightedpin{grh-compact/src/New/Rows/StackedDeterminants.lean}
{155-L190}{The test point and exact reciprocal polynomial};
\weightedpin{grh-compact/src/New/Selmer/RowsExist.lean}
{24-L55}{Trace comparison on that polynomial}.}

\paragraph{The catalogue of normalized rows.}
\teamanchor{proof.weighted-fit-rows}%
For stage \(j\) the source of the auxiliary lift is \(W_j\), of
rank \(r\), and its target is the standard definite space of rank
\(r+t\), \(t=5r\). The common Hecke system is taken on the
\emph{good} principal-level forms on \(W_m\). The rigid-label
transfer compares the traces of the common Hecke algebra on
the good spaces of \(W_j\) and \(W_m\). For any of the finitely
many systems in their sum, interpolate a polynomial in
common Hecke operators that is \(1\) on this system and \(0\)
on the others. Its trace gives equality of the dimensions of
these parts. The nonzero source of each genuine packet in the
stage-\(j\) lift therefore has its Hecke system on \(W_m\).
For each nonzero
isotypic part fix a nonzero component of one good tuple idempotent.
The group-algebra action at the chosen places commutes with every
Hecke operator used here. Thus taking that component preserves the
Hecke system. Its source attachment is taken on the standard model
of \(W_m\). The label compatibility needed is annihilation by the
polynomials for the actual squarefree product type on the indicated
tame inertia, and the Hodge jumps of the actual Harish--Chandra
parameter at \(\ell\). For the rank-one sharp pair one starts with
the constant on \(U(1)\), the trivial source representation and
source label \(X-1\). Its target has rank \(2z\), with \(t'=2z-1\).
Its normalized row is a character obtained from the source by the
same rank normalization. The trivial source character itself would
have jump zero and could not supply the row conditions.
\footnote{\weightedpin{grh-compact/src/New/Trace/GoodIdempotent.lean}{113-L149}{Auxiliary pair and good forms},
\weightedpin{grh-compact/src/New/Trace/GoodFormTransfer.lean}{87-L108}{Transfer on the good Hecke spaces},
\weightedpin{grh-compact/src/New/Selmer/RowsExist.lean}{57-L89}{Actual lower eigenpackets},
and \weightedpin{grh-compact/src/New/Selmer/FitTwistsExist.lean}{24-L34}{The constant source in rank one}.}

\paragraph{The Galois representation of a good lower packet.}
\teamanchor{proof.weighted-fit-good-lower-attachment}%
Take the nonzero good component on \(W_m\) selected for the lower
Hecke system. In its cuspidal decomposition take an irreducible
constituent \(\pi\) on which this component projects nontrivially.
The projection preserves its archimedean type, its Hecke values and
its finite label idempotents. At a split place over each of the two
distinct primes \(g_i\), the primitive Green type forces the local
component to be depth-zero supercuspidal, as in the
\hyperlink{proof.weighted-good-transfer}{good trace transfer}.
The definite archimedean factors are cohomological with the
regular parameters of the prescribed Weyl modules. The odd rank
\(r=3^{a+1}\) and the two split places give the hypotheses of Harris and
Labesse, Theorem~2.2.2. They supply a conjugate-self-dual
cohomological cuspidal \(\Pi\) of \(\mathrm{GL}_r(\mathbf A_E)\).
At a split place \(v\) of \(F\), with \(w,\overline w\) above it, use the projection onto
the \(w\)-matrix as the isomorphism of \(U(W_m)(F_v)\) with
\(\mathrm{GL}_r(E_w)\). The theorem gives \(\Pi_w=\pi_v\) in that
frame and \(\Pi_{\overline w}=\pi_v^\vee\) in the conjugate frame.
In particular this local identification holds at a specified split
place over the coefficient prime.\footnote{\href{https://archive.intlpress.com/site/pub/files/_fulltext/journals/ajm/2004/0008/0004/AJM-2004-0008-0004-a019.pdf\#page=23}{Harris--Labesse, Theorem~2.2.2, p.~675}, and \weightedpin{grh-compact/src/Lib/Adelic/Unitary/SpaceDefs.lean}{148-L159}{The split matrix projection}.}

For a cohomological conjugate-self-dual cuspidal
\(\Pi\) on \(\mathrm{GL}_r(\mathbb A_E)\) with \(E\) CM,
Chenevier--Harris give a continuous semisimple rank-\(r\)
attachment \(R\) for each coefficient prime and coefficient
embedding. At a place \(w\nmid\ell\) its Frobenius-semisimplified
Weil--Deligne representation is dominated by
\(\operatorname{rec}(\Pi_w|\det|_w^{(1-r)/2})\) in the
ordering that fixes the Weil semisimplification and bounds
the monodromy. At \(w\mid\ell\) it is de Rham, with the
labelled jumps of the algebraic coefficient, and if
\(\Pi_w\) is spherical it is crystalline with the
characteristic polynomial of geometric Frobenius in that
same normalized parameter, read on the first linear power
of crystalline Frobenius. The lower \(\Pi\) just constructed
satisfies the three global hypotheses by the definite regular
type and Harris--Labesse base change. Use the chosen \(j_\iota\).
At a split spherical
\(w\), let \(\alpha_1,\ldots,\alpha_r\) be the normalized Satake
parameters of \(\pi_v\), with the uniformizer going to geometric
Frobenius. With the compact subgroup of volume one, the Hecke
eigenvalue of the \(j\)-th minuscule double coset is
\(q_w^{j(r-j)/2} e_j(\alpha_1,\ldots,\alpha_r)\). In the monic
Hecke polynomial the extra factor is \(q_w^{j(j-1)/2}\), so its
roots are \(q_w^{(r-1)/2}\alpha_i\). This is exactly the shift
\(\Pi_w\otimes |\det|_w^{(1-r)/2}\) used in Theorem~3.2.3.
At \(w\nmid\ell\), part~(a\(^{\prime}\)) gives trivial inertia,
zero monodromy and this characteristic polynomial for geometric
Frobenius. At \(w\mid\ell\), part~(c) gives crystallinity and
the same characteristic polynomial for \(\varphi^{[k_w:\mathbf F_\ell]}\).
For a geometric Frobenius lift \(s\), this is the first linear
power in the crystalline Weil operator. Indeed, on crystalline
periods the diagonal action of \(s\) is trivial. The power of
\(\varphi\) is then fixed by the action of \(s^{-1}\) as
\(x\mapsto x^{q_w}\) on the residue field. Over the coefficient
field the Frobenius permutes the components of the unramified
base field, so the polynomial holds in every component.
Cayley--Hamilton gives annihilation on the comparison module for
every inducing embedding and every such lift. Degree-one places of
\(E\) are split over \(F\), so the exact equalities at almost all
of them give the required Frobenius limits. Chebotarev and
Brauer--Nesbitt identify \(R\) with any semisimple
attachment to this lower Hecke system.\footnote{\href{https://webusers.imj-prg.fr/~michael.harris/ConstructionII.pdf\#page=2}{Chenevier--Harris, \S1 geometric reciprocity convention} and \href{https://webusers.imj-prg.fr/~michael.harris/ConstructionII.pdf\#page=11}{Theorem~3.2.3 in the authors' text}, published pp.~66--67.}

Fix a real place of \(F\) and write \(\lambda_1\ge\cdots\ge\lambda_r\)
for the highest weight of the Weyl module used in the definite
algebraic forms. The scalar realization sends an algebraic form
\(f\) to \(L(\tau_\lambda(\overline{g_\infty}^{-1})f(g_f))\).
Right translation acts on the functionals \(L\) through
\((\tau_\lambda\circ c)^\vee\). Since the group is compact unitary,
this has highest weight \(\lambda\) at the chosen embedding
\(\Phi\). The numbers \(h_i=\lambda_i+(r+1)/2-i\) are strictly
decreasing integers. The archimedean
base-change parameter on \(\mathbf C^\times\) has characters
\((z/\overline z)^{h_i}\). Thus its two infinitesimal characters
are the weights \(h_i\) and \(-h_i\), with the latter in reverse
order in the dominant chamber. The coefficient system for cohomology with
\(\Pi_\infty\) has the opposite duality, with highest weights
\(\mu_{\Phi,i}=-\lambda_{r+1-i}\) and
\(\mu_{\overline\Phi,i}=\lambda_i\). By the formula
\(\mu_{\sigma,i}+r-i\) in Chenevier and Harris, the labelled
jump multisets are \(\{i-1-\lambda_i\}_i\) and
\(\{\lambda_i+r-i\}_i\), respectively. These are the jumps of
the lower form before the row twist. The convention gives jump
\(-1\) to the cyclotomic character.\footnote{\href{https://intlpress.com/site/pub/files/_fulltext/journals/cjm/2013/0001/0001/CJM-2013-0001-0001-a002.pdf\#page=9}{Chenevier--Harris, equation (1.6), p.~61}.
For the realization and the jumps see \weightedpin{grh-compact/src/Lib/GLn/SchurModule/SchurOrthogonality.lean}{18-L31}{Scalar realization} and \weightedpin{grh-compact/src/Lib/Weil/Unramified/SphericalHecke.lean}{120-L157}{Hecke polynomial and labelled jumps}.}

The lower principal level imposes its reductions only at split
places. At a nonsplit finite place it is the full stabilizer of
the standard self-dual lattice. The unit conjugate difference in
\(\mathcal O_E\) shows that \(E/F\) is unramified at every finite
place. After the unramified quadratic extension the lattice's
integral unitary model becomes \(\mathrm{GL}_r\) over the integers.
Reductive descent gives a hyperspecial subgroup of the
quasi-split \(p\)-adic group. Its spherical unramified base change is included
in Harris and Labesse, Theorem~2.2.2. Thus parts~(a\(^{\prime}\))
and~(c) also give unramifiedness or crystallinity there, as needed
for the lower row off the split level.\footnote{\weightedpin{grh-compact/src/Lib/Adelic/Adeles/CoprimeValuation.lean}{245-L247}{Unit conjugate difference} and \weightedpin{grh-compact/src/Lib/Adelic/Unitary/CongruenceLevel.lean}{105-L120}{The lower level}.}
The constant lower packet
in rank one has the trivial representation with polynomial \(X-1\)
and jump zero before applying its own row twist.

\paragraph{The tame labels of these lower packets.}
\teamanchor{proof.weighted-fit-good-lower-labels}%
At a chosen split \(w\), write the squarefree label as
\(c_w=\prod_a Q_a\) and put \(d_a=\deg Q_a\). Its finite Green
representation is the parabolic induction of the cuspidal Green
representations \(\tau_a\) attached to the factors. The label
idempotent has nonzero image on the selected \(\pi_v\). By finite
Frobenius reciprocity, \(\pi_v\) contains the inflation of
\(\boxtimes_a\tau_a\) to the parahoric for the blocks of size
\(d_a\). Moy and Prasad, Theorem~6.11(2), place \(\pi_v\) in the
normalized induction of depth-zero supercuspidals \(\pi_a\) with
these finite cuspidal types. Distinct factors give distinct types
in each fixed dimension, so the inducing supercuspidals are pairwise
inequivalent up to unramified twist. Their singleton segments are
unlinked, and Zelevinsky, Theorem~4.2, gives irreducibility of the
induction. The local parameter of \(\Pi_w=\pi_v\) is consequently
the direct sum of the parameters of the \(\pi_a\), with zero
monodromy.\footnote{\href{https://www.ams.org/tran/1955-080-02/S0002-9947-1955-0072878-2/S0002-9947-1955-0072878-2.pdf}{Green, \S\S7--8}, \href{https://www.e-periodica.ch/cntmng?pid=com-001\%3A1996\%3A71\%3A\%3A10\#page=21}{Moy--Prasad, Theorem~6.11, p.~118}, and \href{https://www.numdam.org/item/10.24033/asens.1379.pdf\#page=20}{Zelevinsky, Theorem~4.2, pp.~184--185}.}

Write \(\omega_D=\exp(2\pi i/D)\), and use the faithful
character \(T\) of the algebraic residue field from the label.
For a root \(x\) of \(Q_a\), set \(D_a=q_w^{d_a}-1\) and take
\(b_x\) with \(T(x)=\omega_{D_a}^{b_x}\).
The prescribed Green representation \(\tau_a\) uses the
regular character \(\theta_x(y)=T(y)^{b_x}\) of
\(\mathbf F_{q_w^{d_a}}^\times\). In particular its character on
the regular elliptic element \(y\) in this torus is
\((-1)^{d_a-1}\sum_{j=0}^{d_a-1}\theta_x(y^{q_w^j})\).
The depth-zero local Langlands formula for \(\pi_a\) induces the
Weil character from the unramified extension of degree \(d_a\),
with this restriction on units through geometric reciprocity.
Both the rectifier and the rank normalization are unramified.

With \(D=q_w^{r!}-1\), an inertia element \(s\) of tame exponent
\(b\) has \(T(\mathrm{res}(s(z)/z))=\omega_D^b\) for a
\(D\)-th root \(z\) of a uniformizer. Put \(z_a=z^{D/D_a}\).
Under the geometric Artin map, a unit \(u_a\) representing the
image of \(s\) in abelianized inertia of the unramified
extension has \(\mathrm{res}(u_a)=\mathrm{res}(s(z_a)/z_a)\).
Thus \(\theta_x(\mathrm{res}(u_a))=\omega_{D_a}^{b_x b}=T(x)^b\).
The Frobenius conjugates give the roots \(j_\iota(T(x^{q_w^j})^b)\)
for \(0\le j<d_a\). Taking all factors gives the polynomial for
\(c_w\) at \(s\).\footnote{\weightedpin{grh-compact/src/Lib/LocalGL/Green/GreenTau.lean}{54-L61}{The Green torus character}, \href{https://doi.org/10.1090/S0894-0347-05-00487-X}{Bushnell--Henniart, Part I, \S2.2 and Corollary~3.3}, \href{https://doi.org/10.1112/plms/pdp053}{Part III, Proposition~21}, \href{https://intlpress.com/site/pub/files/_fulltext/journals/mrl/2016/0023/0003/MRL-2016-0023-0003-a001.pdf\#page=5}{Adrian--Roe, Proposition~3.4}, and \weightedpin{grh-compact/src/Lib/ClassField/Local/ArtinMapAtPlace.lean}{189-L198}{The geometric Artin symbol on units}.}

Away from \(\ell\), part~(a\(^{\prime}\)) of Chenevier--Harris
transfers this finite inertia action to \(R\), since monodromy
is zero and Frobenius semisimplification preserves inertia.
The polynomial is the characteristic polynomial of \(R(s)\),
so it also supplies the determinant of \(R(s)\). At a chosen
place over \(\ell\), use the Shin-regular case of
Barnet-Lamb--Gee--Geraghty--Taylor. For a regular algebraic
essentially conjugate-self-dual cuspidal \(\Pi\) over a CM field,
it gives
\[
 \iota\operatorname{WD}(R|_{G_{E_w}})^{F\text{-}\mathrm{ss}}
   \simeq \operatorname{rec}\bigl(\Pi_w|\det|_w^{(1-r)/2}\bigr)
   \qquad (w\mid\ell)
\]
if the weight is Shin-regular. This condition holds in odd rank,
or in even rank if at some embedding the dominant algebraic
weight strictly decreases at an odd index. The present
cuspidal \(\Pi\) is regular algebraic, conjugate self-dual and
has odd rank \(r\). Both representations have the same
semisimple attachment by their almost-everywhere spherical
polynomials. The theorem identifies the Frobenius-semisimplified
Weil--Deligne representation at the place over \(\ell\), in
particular its zero monodromy and its action of inertia.
Potential semistability with zero monodromy gives potential
crystallinity. The theorem holds for any coefficient prime.
In this datum the level primes are odd, so at \(\ell=2\) the
chosen-place types are already covered by the prime-to-\(\ell\)
clause. Every inertia Weil pair here has Frobenius
exponent zero, so its linear operator is the inertia action on
the potentially crystalline periods. Cayley--Hamilton gives
the required polynomial on this comparison module as well.
The conjugate places follow from the dual local matching and
the same geometric reciprocity. The polarization
\(R^c\simeq R^\vee\otimes\varepsilon_\ell^{1-r}\) follows by
comparing the normalized spherical polynomials at conjugate
places and using Chebotarev. On a geometric Frobenius,
\(\varepsilon_\ell^{1-r}\) gives \(q_w^{r-1}\), which takes the
inverse of each normalized root to the normalization of the
inverse Satake parameter.\footnote{\href{https://www.numdam.org/item/10.24033/asens.2212.pdf\#page=7}{Barnet-Lamb--Gee--Geraghty--Taylor, \emph{Local-global compatibility for \(\ell=p\), II}, Theorem~1.1, p.~170}, with Shin-regularity on the preceding page.}

These results supply the unshifted labelled lower column. The
stage twist and the later row bounds apply to this column as
before. The same lower representation gives the genuine
block-plus-string representation wherever the theta polynomial
has that factorization. The general attachment and its local clauses also apply to the
other packets in the upper catalogue.

\paragraph{Galois attachment on the definite catalogue.}
\teamanchor{proof.weighted-fit-spherical-cancellation}%
Keep the field \(E=K(\mu_{N_1})\) at the selected square and write
\(F=E^+\). For each rank \(a\geq1\) in the catalogue, use the group of the
standard positive definite Hermitian form over \(E/F\). In particular,
this applies to the upper groups of ranks \(6r\) and \(2z\), at the base,
twisted and deeper levels of the coefficient test. The weight includes the
determinant weight of the form at that level. The field contains
\(\mu_{5040}\), and \(E=F\mathbf Q(i)\) and \([F:\mathbf Q]\geq2\).
The latter inequality follows already from
\(\mathbf Q(\zeta_5+\zeta_5^{-1})\subset F\).
At a finite prime of residue characteristic \(p\), choose a root of unity
in \(E\) of order \(3\) or \(5\) prime to \(p\). Its reduction still has
that order, so complex conjugation cannot act trivially on the residue
field if it stabilizes the prime. This proves that \(E/F\) is unramified
at all finite places. The standard self-dual lattice gives a reductive
model by descent from \(\mathrm{GL}_a\) over the unramified quadratic
ring, so its stabilizer is hyperspecial at a nonsplit finite place.
The base, twisted and deeper levels of the upper catalogue keep this
stabilizer at every nonsplit place. The chosen \(g_i\), the new level
places \(Q\) and the Taylor--Wiles primes are all split. Thus when a
nonsplit completion is read at \(\ell\), its actual upper packet has
the spherical vector used in the crystalline comparison.
\footnote{\weightedpin{grh-compact/src/New/Selmer/BaseLevelOfLabel.lean}
{23-L31}{The base level and split places}, and
\weightedpin{grh-compact/src/New/Selmer/BaseLevelOfLabel.lean}
{57-L68}{Local conditions on the prescribed base and deeper forms}.}
At a split \(v=w\overline w\), use the projection \(i_w\) onto the
\(w\)-matrix. In that frame the standard compact is
\(\mathrm{GL}_a(\mathcal O_{E_w})\).

Let \(f\) be a nonzero algebraic eigenform in the catalogue, with
monic polynomials \(P_w\) at almost all places of prime norm. Fix the
coefficient isomorphism \(j_\iota:\mathbf C\longrightarrow
\overline{\mathbf Q}_\ell\) extending the reading under consideration.
Realize \(f\) in scalar automorphic forms by taking matrix coefficients
against the dual of its weight module. The definite rational group
is anisotropic and the quotient is compact. Project onto an irreducible
automorphic constituent \(\pi\) on which the realization is nonzero.
Projection preserves the weight and all the Hecke equations on \(f\).
For a specified local test one may first take the local translate,
mean or finite idempotent occurring in that test, and project the
resulting nonzero vector. The uniqueness below gives the same attachment
for all such choices with this Hecke system.

Labesse's Corollary~5.3 applies since the archimedean factors of \(\pi\)
are discrete series on compact groups and satisfy the sign condition
\((*)\). The local matches are equivalently those of Thorne's
Theorem~4.2, for the very group of the standard form in his equation
(4.1). Write the base change as \(\Pi=\boxplus_i\Pi_i\).
Each \(\Pi_i\) is discrete and individually conjugate self-dual.
The base change has \(\Pi_w=\pi_v\circ i_w^{-1}\) at every split place,
and at a hyperspecial spherical place its local representation is the
spherical base change. At infinity its infinitesimal character is the
base change of that of \(\pi\).
\footnote{\href{https://webusers.imj-prg.fr/~michael.harris/CBCM.pdf\#page=44}
{Labesse, \emph{Changement de base CM et s\'eries discr\`etes}, Corollary~5.3},
and \href{https://doi.org/10.1112/plms.12584}
{Thorne, \emph{A \(p\)-adic approach to the existence of level-raising
congruences}, \S4, Theorem~4.2}. Labesse's local spherical transfer
uses the hyperspecial test and the split transfer uses the indicated
matrix projection.}

\teamanchor{proof.weighted-fit-spherical-shifts}%
By the residual spectrum theorem of M{\oe}glin and Waldspurger,
\(\Pi_i\) is the Speh representation attached to a unitary cuspidal
\(\sigma_i\) of rank \(a_i\) and a positive integer \(d_i\).
Its exponents are \(\nu_{ik}=(d_i-1)/2-k\) for \(0\leq k<d_i\).
Here \(\sum_i a_i d_i=a\). The Speh datum is unique, so
\(\sigma_i^c\simeq \sigma_i^\vee\) for each \(i\).
In fact conjugation and duality replace \(\sigma_i\) by
\((\sigma_i^c)^\vee\) and reverse the same string of exponents.
Write
\[
 t_{ik}=\nu_{ik}+\frac{a_i-a}{2},
 \qquad \pi_{ik}=\sigma_i\otimes|\det|^{t_{ik}}.
\]
At a complex embedding \(\tau\), let \(h\) range over the
infinitesimal entries of \(\sigma_i|\det|^{\nu_{ik}}\).
Together over \(i,k\) these are the entries of \(\Pi\).
They are distinct and lie in \((a-1)/2+\mathbf Z\) at each
embedding. Thus the entries \(h+(a_i-a)/2\) of \(\pi_{ik}\) are
distinct and lie in \((a_i-1)/2+\mathbf Z\).
They give a dominant integral weight after subtracting the
Harish--Chandra shift in rank \(a_i\). This is regular algebraicity.
The representation \(\pi_{ik}\) is cuspidal and polarized by
\[
 |\cdot|_F^{2t_{ik}}\delta_{E/F}^{\,a_i\bmod 2}.
\]
Indeed the absolute value on \(E\) is the pullback by norm and the
quadratic character is trivial on norms and has sign \(-1\) at
each real place. In particular the parity correction here is a
twist of the automorphic cuspidal representation by a power
of the norm. The Galois attachment is taken only after the twist
has made the constituent regular algebraic.

Apply Theorem~2.1.1 of Barnet-Lamb, Gee, Geraghty and Taylor to
each \(\pi_{ik}\) and form the continuous semisimple sum
\[
 R=\bigoplus_{i,k}r_{\ell,j_\iota^{-1}}(\pi_{ik}).
\]
The total local shift for this summand is
\(t_{ik}+(1-a_i)/2=\nu_{ik}+(1-a)/2\). At a hyperspecial
spherical place the base change is unramified. Each cuspidal seed
there is unramified as well. Indeed the spherical representation
is a subquotient of the normalized induction from the translated
seeds. Taking maximal compact invariants is exact, and by the
Iwasawa decomposition the induced fixed space is the tensor of
the fixed spaces of the seeds. At \(w\nmid\ell\) the local
comparison gives unramifiedness. At \(w\mid\ell\) the spherical
clause of the same theorem gives crystallinity, also when \(E_w\)
is ramified over \(\mathbf Q_\ell\). Both statements hold for every
\(\ell\), including \(2\).
\footnote{\href{https://www.numdam.org/articles/10.24033/asens.1595/}
{M{\oe}glin--Waldspurger, \emph{Le spectre r\'esiduel de \(\mathrm{GL}(n)\)},
main theorem and \S1.11},
\href{https://doi.org/10.1017/S1474748020000602}
{Anastassiades--Thorne, proof of Corollary~3.4},
and \href{https://annals.math.princeton.edu/wp-content/uploads/annals-v179-n2-p03-p.pdf\#page=38}
{Barnet-Lamb--Gee--Geraghty--Taylor, \emph{Potential automorphy and
change of weight}, Theorem~2.1.1(2)--(4), p.~538}.}

To calculate the labelled jumps at \(\tau\), a rank-\(a_i\) block with
infinitesimal entry \(H\) has jump \((a_i-1)/2-H\).
This is the formula in Theorem~2.1.1(3) with the coefficient weight
dual to the infinitesimal weight. Substituting
\(H=h+(a_i-a)/2\) gives \((a-1)/2-h\).
At the chosen CM embedding \(\Phi\) over a real place, the scalar
realization acts on the functionals by \((\tau_\lambda\circ c)^\vee\),
of highest weight \(\lambda\) in that frame. Hence the entries there
are \(h_j=\lambda_j+(a+1)/2-j\), and at the conjugate embedding they
are their negatives in reverse order. The resulting lists are
\[
 \{j-1-\lambda_j:1\leq j\leq a\},
 \qquad
 \{\lambda_j+a-j:1\leq j\leq a\}.
\]
This is the labelled jump formula on the actual weight of \(f\).
The cyclotomic character in this convention has jump \(-1\).

\teamanchor{proof.weighted-fit-split-wd}%
\teamanchor{proof.weighted-fit-tame-restriction}%
At a split spherical \(w\), with normalized Satake roots \(\alpha_j\)
of \(\pi_v\circ i_w^{-1}\), the rank-\(a\) monic Hecke polynomial
has roots \(q_w^{(a-1)/2}\alpha_j\), counted with multiplicity.
Indeed the \(j\)-th minuscule Hecke eigenvalue is
\(q_w^{j(a-j)/2}e_j(\alpha)\), and its coefficient in degree \(a-j\)
carries the factor \((-1)^j q_w^{j(j-1)/2}\). The total shift of the
direct sum gives exactly this polynomial at all spherical split
places away from \(\ell\). A place of \(E\) of prime norm is split
over \(F\), since \(E/F\) is unramified. Thus outside a finite set
of places \(w\) of prime norm it equals \(j_\iota P_w\) on every
geometric Frobenius lift. This implies coefficientwise approximation
on the Frobenius filter. More explicitly, on the triples \((w,A,s)\)
with \(A\) above \(w\) and \(s\) geometric, in every nonempty open
subset of \(G_E\) the equality holds on every such triple outside the
same finite set of \(w\). As \(s\) approaches any \(g\in G_E\) on this
filter, the coefficients tend to those of \(\det(X-R(g))\).
If a second continuous semisimple representation has the corresponding
approximation property for each error \(\eta>0\), use a nonempty
open subregion furnished by its Frobenius filter, and then a place of
prime norm outside both finite exceptional sets by Chebotarev.
Its trace differs from that of \(R\) by at most \(\eta\) on a dense
set, and thus everywhere by continuity. Letting \(\eta\) tend to zero
and applying Brauer--Nesbitt proves uniqueness.

Thorne's Corollary~4.3 attaches a semisimple representation to this
same \(\pi\), at any coefficient prime. Its almost-everywhere spherical
polynomials are the same, so it is isomorphic to \(R\).
For every split \(w\), including a specified \(w\mid\ell\), it gives
\[
 \mathrm{WD}(R|_{G_{E_w}})^{F\text{-}\mathrm{ss}}
   \ \simeq\
 j_\iota\mathrm{rec}_{E_w}\!\left(
   (\pi_v\circ i_w^{-1})\otimes|\det|_w^{(1-a)/2}\right).
\]
Here local reciprocity sends a uniformizer to geometric Frobenius.
In particular, if \(f\) has Hecke polynomial \(Q\) at a specified
spherical split \(w\), projection preserves \(Q\) and the characteristic
polynomial on the local parameter is exactly \(j_\iota Q\).
Frobenius semisimplification keeps the inertia action and the
monodromy operator. For the characteristic polynomial at any Weil
element, write it as inertia times a power of Frobenius. Some power
of Frobenius centralizes the finite inertia image, so the unipotent
part of Frobenius centralizes that image as well. Removing the
commuting unipotent factor leaves the characteristic polynomial
unchanged, with every multiplicity.
\footnote{\href{https://doi.org/10.1112/plms.12584}
{Thorne, Notation for the Tate normalization and \S4, Corollary~4.3,
for the split comparison}. The formula
\(\mathrm{rec}^T(\pi)=\mathrm{rec}(\pi\otimes|\det|^{(1-a)/2})\)
also appears in the conventions of
\href{https://www.numdam.org/item/10.1007/s10240-021-00126-4.pdf}
{Newton--Thorne, \emph{Symmetric power functoriality for holomorphic
modular forms, II}, pp.~120--121}.}

For a spherical \(w\nmid\ell\), this is the polynomial of \(R(s)\)
for any geometric Frobenius lift \(s\). For a spherical
\(w\mid\ell\), put \(f_w=[k_w:\mathbf F_\ell]\).
The comparison uses the first linear power \(\varphi^{f_w}\) on
the crystalline periods. For a lift \(s\) of geometric Frobenius
the paired operator there is \(\varphi^{f_w}(s\otimes R(s))\),
where \(s\otimes R(s)\) fixes the crystalline periods.
The components of the unramified base field are permuted by
\(\varphi\), so the characteristic polynomial on each component
is \(j_\iota Q\). Cayley--Hamilton gives the required polynomial
on the comparison space for every prolongation \(A\) and every
compatible inducing embedding. More generally, at a split place
where the local parameter has \(N=0\), the same argument reads its
characteristic polynomials on all Weil elements. In residue
characteristic \(\ell\) the representation is then potentially
crystalline, and the operator of a geometric-degree-\(k\) element
on the potentially crystalline periods is
\(\varphi^{f_w k}(s\otimes R(s))\). At degree zero it is the finite
inertia action. Away from \(\ell\), \(N=0\) gives finite inertia on
the original representation and the Weil action agrees with that
restriction. These conclusions use the local parameter of the
same split component that retains the test.

At almost all split places the conjugate Satake roots are their
inverses, since the conjugate matrix projection is inverse transpose.
The cyclotomic value at a geometric Frobenius is \(q_w^{-1}\).
The two normalized sets of roots therefore give
\(R^c\simeq R^\vee\otimes\varepsilon_\ell^{1-a}\) by uniqueness.
For the determinant let \(\omega_\pi\) be the central character of
\(\pi\) on the norm-one ideles. The algebraic character
\[
 x\longmapsto \omega_\pi((x/x^c)I_a)\,|x|_E^{-\binom a2}
\]
has determinant avatar \(\det R\) by comparing it with the product
of the normalized roots at almost all split places. At a split \(w\)
its value on a unit equals the scalar by which that central matrix
acts. On a spherical uniformizer it equals \((-1)^a Q(0)\).
The infinity exponents are \(\sum_j\lambda_j-\binom a2\) at the
chosen embedding and \(-\sum_j\lambda_j-\binom a2\) at the conjugate
embedding. The separate rank-one avatar calculation applies to this
character and to the row twists, with the same \(j_\iota\). Thus the
spherical attachment applies to the whole catalogue before any
geometric theta or local label identities are used.

\paragraph{Sorting the local exponents on the same packet.}
\teamanchor{proof.weighted-fit-sort-exponents}%
Let \(F=E_w\) with residue size \(q\) and uniformizer \(\varpi\),
normalized by \(|\varpi|=q^{-1}\). Let \(\pi\) be the smooth
unitary local action at this split
place. Write \(I\) for the upper pro-\(p\) Iwahori and
\(P_I\) for its averaging projector. Fix a nonzero \(I\)-fixed
vector \(x\), and let \(V_x\) be the span of its local
translates. We use the finite-dimensional space \(V_x^I\) on the
actual forms, with all the other levels and the weight unchanged.
Suppose that diagonal units act on \(x\) by
\(\prod_i\tau_i(t_i)\). Suppose that each permutation \(\sigma\)
with \(\tau_{\sigma(i)}=\tau_i\) for every \(i\) acts on \(x\)
by a scalar. We will find characters
\(\theta_i:F^\times\to\mathbb C^\times\) with
\(\theta_i|_{\mathcal O_F^\times}=\tau_i\) and a functional \(L\)
nonzero on \(x\) with
\[
 L\bigl(\pi(t)P_I v\bigr)=\xi(t)L(v),\qquad
 \xi(t)=\prod_{i=1}^n\theta_i(t_i)|t_i|^{n-i}
\]
for every diagonal \(t\) with
\(\operatorname{ord}(t_1)\ge\cdots\ge\operatorname{ord}(t_n)\).
We will also find a nonzero pro-Iwahori eigenvector in \(V_x\)
whose torus character comes from a permutation of the same
\(\theta_i\), with \(|\theta_i(\varpi)|\) in nonincreasing
order. The unit values and the uniformizer values are
permuted together.

\emph{The dominant operators and a rank-one identity.}
For a dominant diagonal \(t\) put \(A_t=P_I\pi(t)P_I\).
These operators multiply as \(A_t A_u=A_{tu}\).
To see this, write the pro-Iwahori by elimination as upper
unipotents over \(\mathcal O_F\), diagonal units congruent to
\(1\), and lower unipotents with entries in \(\varpi\mathcal O_F\),
in that order. In a term \(P_I\pi(t i u)P_I\) of the middle
average, move the upper and diagonal parts of \(i\) left across
\(t\), and the lower part right across \(u\). They are still in
\(I\), by dominance, so each term equals \(A_{tu}\).
In particular, these operators commute. They preserve \(V_x^I\),
as do their complex adjoints \(P_I\pi(t^{-1})P_I\).

For adjacent positions \(i,j=i+1\), write \(s\) for their
transposition, also viewed as its permutation matrix, and
\(x_+(c)=1+cE_{ij}\), \(x_-(c)=1+cE_{ji}\).
If \(L\) is any \(I\)-invariant functional, the cosets of
\(I\cap s^{-1}I s\) in \(I\) have representatives \(x_+(c)\)
indexed by \(\bar c\in k_F\). Indeed this intersection adds
the single condition that the \(ij\)-entry reduce to zero.
If \(z\) is fixed by \(x_+(\mathcal O_F)\), it follows that
\(L(\pi(s)P_I z)=L(\pi(s)z)\).
If \(z\) is fixed by \(x_-(\mathcal O_F)\), take \(0\) for
the zero coset and units \(u\) for the other residues. Direct
multiplication of the two-by-two blocks gives
\[
 s x_+(u)=x_+(u^{-1})\,d(u)\,x_-(u^{-1}),\qquad
 d(u)=\operatorname{diag}_{{i,j}}(-u^{-1},u).
\]
The remaining diagonal entries of \(d(u)\) equal \(1\).
Averaging gives the identity
\begin{equation}
\label{eq:weighted-sort-rank-one}
 q L(\pi(s)P_I z)
   =L(\pi(s) z)+\sum_{\bar u\ne0}L(\pi(d(u))z).
\end{equation}
The same calculation applies to the vector-valued map \(P_I\).

Every permutation operator \(P_I\pi(\sigma^{-1})\) is injective
on \(I\)-fixed vectors. Indeed if \(P_I\pi(\sigma^{-1})v=0\)
and \(\sigma\) has adjacent inversion \(i<j\), the vector
\(z=\pi(\sigma^{-1})v\) is fixed by \(x_-(\mathcal O_F)\).
Diagonal units normalize \(I\), so \(P_I\pi(d(u))z=0\).
Equation \eqref{eq:weighted-sort-rank-one} gives \(P_I\pi(s)z=0\),
removing that inversion.
Iteration gives \(P_Iv=v=0\).
Now the rotation \(\omega\) with
\(\omega e_1=\varpi e_n\), \(\omega e_a=e_{a-1}\) for \(a>1\),
normalizes \(I\). For every entry that crosses the end of the
ordered basis, conjugation supplies or removes one power of
\(\varpi\), exactly the change between the lower and upper
conditions. For \(1\le a<n\), the diagonal \(w_a\) with
\(\varpi\) in the first \(a\) positions and \(1\) elsewhere
is a permutation matrix times \(\omega^a\), since those are
the positions that wrap in \(\omega^a\). Thus \(A_{w_a}\) is
injective on \(I\)-fixed vectors.\footnote{\weightedpin
{grh-compact/src/Lib/LocalGL/Iwahori/IwahoriInjective.lean}{48-L101}{Rank-one averaging and dominant injection},
and \weightedpin
{grh/src/Lib/LocalGL/Iwahori/QuadraticRelation.lean}{207-L242}{Quadratic identity}.} Their product \(A_b\), with
\(b=\prod_{a<n}w_a\), is injective and \(b\) is strictly
dominant.

\emph{An exponent nonzero on the original vector.}
Choose a common eigenvector of the commuting \(A_t\) on the
\(\tau\)-space of \(V_x^I\), which contains \(x\).
Its common eigenvalues give a multiplicative function \(\xi\).
Injectivity of \(A_b\) gives \(\xi(b)\ne0\). Every torus element
becomes dominant after multiplication by a large power of
\(b\); division by \(\xi(b)^k\) therefore extends \(\xi\) to
the full torus as a character. Removing the powers of
\(|t_i|\) in the displayed convention gives the \(\theta_i\).
Their unit characters are the prescribed ones because \(A_t\)
on \(I\)-fixed vectors is just \(\pi(t)\) for diagonal units.

A common eigencharacter of commuting operators in finite
dimension over \(\mathbb C\) occurs also on the dual with the
transpose action, and conversely. Indeed joint generalized
eigenspace decomposition for finitely many generators of the
operator algebra gives the same characters after transposition.
In each nonzero joint block, commuting nilpotents have a common
kernel in the space and in its dual. Thus on \(V_x^I\) there is
a nonzero dual eigenfunctional with character \(\xi\). Write it
as \(\langle h,\cdot\rangle\) with \(0\ne h\in V_x^I\), using
linearity in the second argument. Since the adjoints preserve
\(V_x^I\), we have \(A_t^*h=\overline{\xi(t)}h\) there and
on the ambient representation. Therefore
\(L(v)=\langle h,P_Iv\rangle\) satisfies the exponent law
and \(L(h)>0\).

If \(L(x)=0\), it vanishes on all \(\pi(\sigma)x\) for finite
permutations \(\sigma\). In fact moving the diagonal units
past \(\sigma\) shows that a nonzero value would require
\(\sigma\) to preserve the unit characters of \(x\). Its scalar
action would then give zero as well.
For any vector \(y\) fixed by the lower unipotents over
\(\varpi\mathcal O_F\) and the diagonal units congruent to \(1\),
and any dominant \(t\), the same Gauss factorization gives
\(L(\pi(t)y)=L(\pi(t)P_I y)=\xi(t)L(y)\).
In particular this applies to a finite permutation of \(x\),
which is fixed by the principal congruence subgroup.

Write any \(g\in\mathrm{GL}_n(F)\) as \(u t \sigma k\), with
\(u\) upper unitriangular, \(t\) diagonal and \(k\in I\).
An upper triangular basis of the lattice \(g\mathcal O_F^n\)
gives the Iwasawa factorization. Reduction of its compact
factor gives the triangular--permutation--unitriangular
factorization by elimination over \(k_F\); the last factor
lifts to \(I\). For large \(m\), \(b^m u b^{-m}\in I\) and
\(b^m t\) is dominant. The preceding evaluation gives
\(L(\pi(b^m g)x)=\xi(b^m t)L(\pi(\sigma)x)=0\).
This vanishing holds eventually on every finite sum of
translates of \(x\), in particular on \(h\). It contradicts
\(L(\pi(b^m)h)=\xi(b^m)L(h)\ne0\).
\footnote{\weightedpin
{grh/src/Lib/LocalGL/Iwahori/SortedExponents.lean}{421-L454}{Contraction and monomial translates},
and \weightedpin
{grh-compact/src/Lib/LocalGL/Iwahori/SortedExponents.lean}{103-L144}{The same cyclic space and the original vector}.}

\emph{A strict adjacent inversion.}
Let \(L\) now denote an exponent nonzero at
\(\pi(\sigma^{-1})x\) for a current permutation \(\sigma\).
Use \(\theta_i\) for the characters in the current order.
Suppose \(|\theta_i(\varpi)|<|\theta_j(\varpi)|\) for
\(j=i+1\), and take \(s=(i\,j)\). The exchanged character is
\[
 \xi_s(t)=\xi(t\circ s)|t_i/t_j|,
\]
because the modulus exponents at adjacent positions differ
by one. Put \(S=L\pi(s)P_I\) and \(w=w_i\).

First consider a dominant \(t\) with gap
\(d=\operatorname{ord}(t_i)-\operatorname{ord}(t_j)\le1\).
The vector \(\pi(t\circ s)\pi(s)P_Iv\) is fixed by
the lower unipotents with entries in \(\varpi\mathcal O_F\)
and the diagonal units congruent to \(1\).
For a lower elementary matrix this follows by conjugating
back through \(t\circ s\) and \(s\). Its entry still has
order at least one except possibly for the lower \(ji\)
entry, which becomes an upper entry of order at least zero.
The vector \(\pi(s)P_Iv\) has the same fixed groups.
Apply the dominant evaluation above to both vectors, after
multiplying \(t\circ s\) by a sufficiently large power of
\(b\), and divide by \(\xi(b)^m\). This gives
\begin{equation}
\label{eq:weighted-sort-small-gap}
 L\pi(s)\pi(t)P_I v=\xi(t\circ s)S(v).
\end{equation}
For gap zero, \(\pi(t)P_Iv\) is also fixed by
\(x_+(\mathcal O_F)\), so the upper-root case of the averaging
formula gives \(S A_t=\xi_s(t)S\) at that gap.
For \(t=w\), \(\pi(w)P_Iv\) is fixed by
\(x_-(\mathcal O_F)\). Equations \eqref{eq:weighted-sort-rank-one}
and \eqref{eq:weighted-sort-small-gap} give
\begin{equation}
\label{eq:weighted-sort-step}
 q S(A_w v)=\xi(w\circ s)S(v)
       +\xi(w)L(v)\sum_{\bar u\ne0}\xi(d(u)).
\end{equation}
Here diagonal units normalize \(I\), so \(L\pi(d(u))=\xi(d(u))L\).
On \(\ker L\), this is \(S A_w=\xi_s(w)S\).
For a positive gap \(d\) of a general dominant \(t\), the
diagonal \(t/w^d\) is dominant with gap zero.
Multiplicativity and invariance of \(\ker L\) therefore give
\(S A_t=\xi_s(t)S\) on \(\ker L\) for every dominant \(t\).
Since \(A_w-\xi(w)\) maps into \(\ker L\) and commutes with
the dominant operators,
\[
 L'=S(A_w-\xi(w))
\]
is an exponent for \(\xi_s\).

We show \(L'\bigl(\pi((\sigma s)^{-1})x\bigr)\ne0\).
Put \(y=\pi((\sigma s)^{-1})x\), so \(L(\pi(s)y)\ne0\).
The functional \(L\) reads the unit characters of
\(\pi(\sigma^{-1})x\) in the current order, since it is
nonzero there. If the two current unit characters differ,
those of \(y\) disagree with the character of \(L\) and
\(L(y)=0\). The \(I\)-invariance of \(x\) gives either the
whole upper or the whole lower adjacent integral root group
fixing \(y\), depending on the order of \((\sigma s)i\) and
\((\sigma s)j\). By the two cases of the averaging formula,
\(S(y)\) is \(L(\pi(s)y)\) or \(q^{-1}L(\pi(s)y)\), and is
nonzero. If \(L'(y)=0\), the eigenrelation on \(\ker L\)
implies \(\xi_s(w)=\xi(w)\), hence
\(\theta_j(\varpi)=\theta_i(\varpi)\).

If the two unit characters agree, the transposition
\(\sigma s\sigma^{-1}\) preserves the unit type of \(x\).
Thus \(y\) is a nonzero scalar multiple of
\(\pi(\sigma^{-1})x\). The two orders of the adjacent roots
show that both \(x_+(\mathcal O_F)\) and \(x_-(\mathcal O_F)\)
fix \(y\). Now \(S(y)=L(\pi(s)y)\ne0\). The two-by-two
factorization gives \(L(\pi(d(u))y)=L(\pi(s)y)\) for each
nonzero residue, using invariance of \(L\) on the left and
of \(y\) on the right. If \(L'(y)=0\),
\eqref{eq:weighted-sort-step} says
\[
 q\,\xi(w)S(y)=\xi(w\circ s)S(y)+(q-1)\xi(w)S(y).
\]
It follows that \(\xi(w)=q\,\xi_s(w)\), hence
\(\theta_i(\varpi)=q\,\theta_j(\varpi)\).
Both possible conclusions contradict the strict inversion.
Thus the adjacent exchange preserves the required nonzero
value at the permuted vector.

Swapping the strict adjacent inversions in turn gives the
sorted character. Its exponent is nonzero on the
\(I\)-average of a permutation of \(x\), which belongs to
\(V_x^I\). The finite-dimensional dual argument gives a
nonzero common eigenvector of the same character in
\(V_x^I\). This is the sorted vector on the original packet.
The original functional \(L\) still reads the ordered
coordinates on \(x\). For example the first-line operator in
the local ties is \(q^{n-1}A_{w_1}\), so that
\(L(Uv)=\theta_1(\varpi)L(v)\) on \(I\)-fixed vectors.
\footnote{\weightedpin
{grh-compact/src/Lib/LocalGL/Iwahori/SortedExponents.lean}{24-L84}{Adjacent exchange},
\weightedpin
{grh/src/Lib/LocalGL/Iwahori/SortedExponents.lean}{217-L273}{The exchanged character on the kernel},
and \weightedpin
{grh/src/Lib/LocalGL/Iwahori/SortedExponents.lean}{329-L419}{Nonvanishing in both unit-type cases}.}

\emph{The two applications.}
For the determinant character in the local ties, all the
unit characters agree and every permutation acts by its
determinant character. For a nontrivial first-entry
character, the character-preserving permutations fix the
first index and lie in the line parahoric, where they act by
a scalar. In particular the result applies to a nonzero
vector on either space, including a putative difference of
the two sides of the Frobenius trace identity.

For a finite product of determinant types at a split place, let
\(P_{\mathrm{fin}}\subseteq\mathrm{GL}_n(k_F)\) be the
parabolic with the designated consecutive blocks of sizes
\(m_a\), and \(\lambda=\prod_a\chi_a\circ\det\) its Levi
character, inflated across the unipotent radical. At this
stage the \(\chi_a\) are distinct. On a principal-level
vector, the nonzero central character projection for the
finite representation induced from \(\lambda\) means that
its finite cyclic span contains an irreducible constituent
of that induction. Indeed the complex character sum acts by
the corresponding character multiplicity times the common
group order divided by the constituent dimension. Finite
Frobenius reciprocity gives a nonzero vector in the same
cyclic span transforming by \(\lambda\) under \(P_{\mathrm{fin}}\).
Its lift is \(I\)-fixed, has exactly \(m_a\) occurrences of
\(\chi_a\) on the diagonal units, and every permutation
preserving these unit characters preserves the blocks and
acts by a scalar. Apply the sorting argument to that vector.
This gives the sorted eigencharacter, including the required
multiplicities of each old determinant label. All the
operations are local, so their vectors retain the Hecke
system away from \(w\).
\footnote{\weightedpin
{grh-compact/src/Lib/GaloisRep/LocalGlobal/TypeFrobenius.lean}{22-L45}{Character tests at the separating place},
and \weightedpin
{grh-compact/src/Lib/LocalGL/Iwahori/ExponentOfLabel.lean}{15-L56}{The determinant blocks of the old label}.}

\paragraph{The spherical vector for a sorted unramified exponent.}
\teamanchor{proof.weighted-fit-spherical-mean}%
Keep the local cyclic representation of the actual sorted vector
at \(w\). In this paragraph index the coordinates by \(0,\ldots,n-1\).
Put \(K=\mathrm{GL}_n(\mathcal O_w)\), write \(I\) for the upper
pro-\(p\) Iwahori, and write \(\pi\) for the local action. Let \(L\)
be a nonzero exponent on the cyclic representation with character
\[
 \xi(t)=\prod_{i=0}^{n-1}\theta_i(t_i)|t_i|^{n-1-i},
 \qquad \lambda_i=\theta_i(\varpi),
 \qquad |\lambda_0|\ge\cdots\ge|\lambda_{n-1}|.
\]
All the \(\theta_i\) are unramified. If one starts from the sorted
pro-Iwahori eigenvector, take its joint generalized-character
block for the commuting dominant operators on the finite
space of \(I\)-fixed vectors in its own cyclic span. The ideal
generated by the operators minus their eigenvalues is nilpotent
on this block. The quotient by its image is nonzero, so a
nonzero covector on the quotient, extended by zero on the other
blocks, supplies the dual eigenfunctional with character \(\xi\).
Its extension by \(P_I\) is the exponent \(L\).
Invariance under \(I\) and the trivial diagonal unit character
make \(L\) invariant under the upper Iwahori in \(K\).

We will use the dominant evaluation also for a few diagonals
outside the dominant chamber. For a vector \(y\) fixed by the
lower unipotents congruent to \(1\) and the diagonal units
congruent to \(1\), Gauss factorization of \(I\) gives
\(L(\pi(t)y)=\xi(t)L(y)\) for dominant \(t\).
Indeed the lower and diagonal factors in the average fix \(y\),
and dominance carries the upper factor into \(I\) after
conjugation by \(t\). If both \(y\) and \(\pi(a)y\) have
these fixed groups, multiply the diagonal \(a\) by a sufficiently
large power \(b^k\) of a strictly dominant diagonal \(b\).
Applying the evaluation with \(b^ka\) on \(y\) and with \(b^k\)
on \(\pi(a)y\) proves
\begin{equation}
\label{eq:weighted-mean-compensation}
 L(\pi(a)y)=\xi(a)L(y).
\end{equation}
This uses only the nonzero character value \(\xi(b)^k\).

\emph{Descent of vanishing through the parahoric flag.}
For \(0\le m\le n\), let \(P_m\) be the inverse image in \(K\)
of the subgroup preserving the first \(m\) initial coordinate
subspaces, of dimensions \(1,\ldots,m\). Its reduction is block
upper triangular for singleton blocks at the first \(m\)
coordinates and the remaining full block. Thus \(P_0=K\) and
\(P_n\) is the upper Iwahori. Suppose \(L\) is zero on the
\(K\)-fixed vectors. We prove by induction on \(m\) that it is
zero on the \(P_m\)-fixed vectors.

For \(m<n\), let \(C_m\subset P_m\) stabilize the line \(k_w e_m\)
in the reduction, and let \(R_m\subset P_m\) be the subgroup
whose \(m\)-th row is supported on the \(m\)-th column in the
reduction. Then \(C_m\subset P_{m+1}\). Fix a \(C_m\)-fixed vector
\(x\). For \(i\ge m\), put \(s_i=(m\,i)\), with \(s_m=1\), and
\(L_i=L\pi(s_i)\).

The orbit of \(e_m\) under the reduction of \(P_m\) consists of
the columns \(c\in k_w^n\) having a nonzero coordinate at or
after \(m\). For such \(c\) define \(v_c=\pi(p)x\) using a
lift \(p\in P_m\) with \(p e_m=c\) in the reduction.
The choice of lift is immaterial by the line stabilizer.
If the last nonzero coordinate of \(c\) is \(i\), an upper
triangular matrix sending \(e_i\) to \(c\) gives
\(L(v_c)=L_i(x)\). For example, take the matrix obtained
by replacing column \(i\) of the identity by \(c\); its
diagonal entry there is nonzero. Set
\[
 a_i=(q-1)q^i L_i(x).
\]
There are exactly \((q-1)q^i\) columns with last nonzero
coordinate \(i\), and \(\sum_c v_c\) is \(P_m\)-fixed.
The outer induction hypothesis therefore yields
\(\sum_{i\ge m}a_i=0\).

Now set \(z=\sum_{c_m\ne0}v_c\). It is fixed by \(R_m\) because
every such matrix scales the \(m\)-th coordinate of the column
by a nonzero scalar. For \(i\ge m\), applying \(s_i\) exchanges
the condition \(c_m\ne0\) for \(c_i\ne0\).
If the last nonzero coordinate is \(i\), there is no additional
restriction. At every last coordinate \(k>i\), the fraction
of columns with \(c_i\ne0\) is \(1-q^{-1}\). Thus
\begin{equation}
\label{eq:weighted-mean-column-count}
 L_i(z)=a_i+(1-q^{-1})\sum_{k>i}a_k,
 \qquad L_m(z)=a_m/q.
\end{equation}

Let \(d_m\) be diagonal with \(\varpi^{-1}\) at \(m\) and \(1\)
elsewhere. Conjugation \(d_m^{-1}C_m d_m\) is contained in
\(R_m\). It divides the off-diagonal \(m\)-th column by
\(\varpi\), multiplies the off-diagonal \(m\)-th row by
\(\varpi\), and preserves the earlier flag conditions.
Thus \(\pi(d_m)z\) is \(C_m\)-fixed. We have
\begin{equation}
\label{eq:weighted-mean-shift}
 L_i(\pi(d_m)z)=\lambda_i^{-1}q^{n-1-i}L_i(z).
\end{equation}
For this identity move \(d_m\) past \(s_i\), replacing it by
\(d_i\). The vector \(\pi(s_i)z\) is fixed by the principal
congruence subgroup. The lower root groups congruent to \(1\)
also fix \(\pi(d_i s_i)z\). Conjugating back through \(d_i\)
can turn a lower root of order one into one of order zero
only if its column is \(i\) and its row is greater than \(i\).
Conjugating further by \(s_i\) then puts it in column \(m\)
below row \(i\), where \(R_m\) imposes no zero condition.
All other conjugated roots still belong to the principal
congruence subgroup. The diagonal units congruent to \(1\)
also fix both vectors. Equation \eqref{eq:weighted-mean-compensation}
gives the stated shift.

The inner induction increases \(j\) from \(m\) to \(n\).
Its assertion is that, for every \(C_m\)-fixed \(x\), vanishing
of \(L_i(x)\) for all \(i\ge j\) implies \(L_m(x)=0\).
The case \(j=m\) is immediate. To pass from \(j\) to \(j+1\),
assume \(L_i(x)=0\) for \(i>j\) and put \(C=(q-1)q^{n-1}\).
Form the \(C_m\)-fixed vector
\[
 x'=\lambda_j\pi(d_m)z-Cx.
\]
For \(i\ge j\) the tail in \eqref{eq:weighted-mean-column-count} vanishes and
\(q^{n-1-i}a_i=C L_i(x)\).
The shift gives \(L_i(x')=C(\lambda_j/\lambda_i-1)L_i(x)=0\).
The inner induction hypothesis gives \(L_m(x')=0\).
But \(L_m(z)=a_m/q\), so this last equality is
\[
 C\left(\frac{\lambda_j}{q\lambda_m}-1\right)L_m(x)=0 .
\]
The coefficient is nonzero by \(|\lambda_j|\le|\lambda_m|\)
and \(q>1\). This includes equality of moduli and repeated
roots. The inner induction at \(j=n\) gives the vanishing
on every \(C_m\)-fixed vector, hence on \(P_{m+1}\)-fixed
vectors.\footnote{\weightedpin
{grh-compact/src/Lib/LocalGL/Tadic/ChainLemma.lean}{26-L62}{Vanishing along the parahoric flag},
\weightedpin
{grh/src/Lib/LocalGL/Iwahori/SphericalExponents.lean}{266-L318}{The finite column sums},
and \weightedpin
{grh/src/Lib/LocalGL/Iwahori/SphericalExponents.lean}{463-L635}{Diagonal shift and the two inductions}.} The outer induction now gives vanishing at the
\(P_n\)-average of every vector. Invariance under \(P_n\)
forces \(L=0\), a contradiction.

\emph{The mean and its Hecke polynomial.}
On the finite-dimensional \(I\)-fixed subspace of this cyclic
representation let \(h\) represent \(L\) in the invariant Hermitian
form, linear in its second argument. Since \(L=L P_I\), it is
represented by the same \(I\)-fixed \(h\) on the cyclic space.
Its spherical mean \(g=P_Kh\) represents the restriction of \(L\)
to the \(K\)-fixed subspace and is nonzero by the preceding descent.
The compact average stays in the cyclic space. In particular this
is a vector in the packet supplying the sorted exponent.

For \(0\le d\le n\), let \(T_d\) be the unnormalized spherical
Hecke operator of the diagonal with \(d\) copies of \(\varpi\)
and \(n-d\) copies of \(1\).
It sums the lattice translates indexed by
\(\varpi\mathcal O_w^n\subseteq\Lambda\subseteq\mathcal O_w^n\)
of residue codimension \(d\). Partition the corresponding
subspaces of \(k_w^n\) by the pivot set \(S\) of the echelon
basis with last nonzero coordinates as pivots. Write \(A\)
for the complement of \(S\), of cardinality \(d\).
At a pivot \(s\in S\) the basis vector is \(e_s\) plus a
linear combination of the \(e_i\) with \(i\in A\) and \(i<s\).
These coefficients vary freely. The number of subspaces
with this pivot set is
\[
 q^{\#\{(i,s):i\in A,\ s\in S,\ i<s\}}
   =q^{\sum_{i\in A}(n-1-i)-\binom d2}.
\]
Lift that basis by an upper unitriangular matrix \(u\).
A representative for the lattice is \(u a_A\), where
\(a_A\) has \(\varpi\) precisely on the coordinates \(A\).
Indeed the pivot columns lift the echelon vectors and the
remaining columns are \(\varpi e_i\).
Since \(L\) is upper-Iwahori invariant, it reads this translate
as \(L(\pi(a_A)v)\) for \(v\) spherical. The vectors \(v\) and
\(\pi(a_A)v\) satisfy the fixed group conditions in
\eqref{eq:weighted-mean-compensation}.
In the second case, conjugation back by \(a_A\) changes
an entry in a lower elementary matrix by at most one power
of \(\varpi\), so the conjugated matrix is still integral.
Consequently
\[
 L(\pi(a_A)v)=\xi(a_A)L(v)
       =q^{-\sum_{i\in A}(n-1-i)}\prod_{i\in A}\lambda_i\,L(v).
\]
Multiplying by the cell counts and summing gives
\[
 L(T_d v)=q^{-\binom d2}e_d(\lambda_0,\ldots,\lambda_{n-1})L(v),
\]
where \(e_d\) is the elementary symmetric polynomial.

The spherical operators are normal on the \(K\)-fixed subspace.
In the invariant Hermitian form the adjoint comes from the
inverse double coset. Transposition preserves the Haar measure
giving \(K\) volume one and every \(K\)-double coset by the
elementary divisor decomposition. As an anti-involution it
reverses convolution, so the spherical operators commute.
Put \(a_d=q^{-\binom d2}e_d(\lambda_0,\ldots,\lambda_{n-1})\).
The functional equation gives \(T_d^*g=\overline{a_d}g\), and
normality makes \(\|(T_d-a_d)g\|=\|(T_d^*-\overline{a_d})g\|=0\).
The usual Satake parameters are \(q^{-(n-1)/2}\lambda_i\),
where \(q^{1/2}\) is positive.
The factor \(q^{d(n-d)/2}\) for the unnormalized sum on their
degree-\(d\) elementary symmetric function gives precisely \(a_d\).
The monic Hecke polynomial in the reader's convention is
therefore
\[
 \sum_{d=0}^n(-1)^d q^{\binom d2}a_d X^{n-d}
        =\prod_{i=0}^{n-1}(X-\lambda_i).
\]
The formula counts repeated roots with their multiplicities.\footnote{\weightedpin
{grh-compact/src/Lib/LocalGL/Iwahori/HeckeCounts.lean}{156-L168}{Satake normalization of the lattice sums},
and \weightedpin
{grh/src/Lib/GaloisRep/LocalGlobal/IwahoriCompatibility.lean}{142-L201}{The spherical vector on the same Hecke packet}.}
The local translates are supported on the split pair
\(\{w,\bar w\}\), so \(g\) has the same Hecke system outside
that pair. Thus on this vector an unramified exponent can be read from its
spherical polynomial.

\paragraph{The split Iwahori polynomial on the tested channel.}
\teamanchor{proof.weighted-fit-tame-iwahori}%
Take a sorted eigenvector from the
\hyperlink{proof.weighted-fit-sort-exponents}{local exponent calculation},
and project it nontrivially to an irreducible constituent \(\pi_w\).
It has the same Hecke system away from the split pair. With \(I\) the
upper pro-\(p\) Iwahori and \(U\) the upper unipotent group, let \(j_U\)
denote the projection to the unnormalized \(U\)-coinvariants.
For a dominant diagonal \(d\) and an \(I\)-fixed vector \(y\),
\[
  j_U(P_I\pi_w(d)y)=d\,j_U(y).
\]
In fact factor each element of \(I\) into an upper unipotent, a
diagonal principal unit and a lower unipotent in that order. The
upper factor disappears in the coinvariants, and conjugation by
\(d^{-1}\) takes the other two factors into \(I\).
The image of the sorted eigenvector is nonzero. If it vanished, write
the vector as a finite sum of differences \((\pi_w(u)-1)z\) for
\(u\in U\). Averaging over a compact subgroup of \(U\) containing
these \(u\)'s kills it. A sufficiently large power of a strictly
dominant diagonal conjugates that compact subgroup into \(I\).
Averaging after translation by this power then kills the vector,
contrary to its nonzero eigenvalue for the dominant operator.

In the normalized Jacquet module, the diagonal character attached
to this vector has entries
\[
  \alpha_i=\xi_i\,|\cdot|^{(N-1)/2}.
\]
Indeed the averaging eigencharacter is
\(\prod_i\xi_i(d_i)|d_i|^{N-i}\) and normalization divides it by
\(\delta_B(d)^{1/2}=\prod_i|d_i|^{(N+1-2i)/2}\).
The identity on dominant diagonals and their inverses determines the
action on the entire diagonal torus. The principal-unit invariants
in the Jacquet module are finite dimensional by admissibility. A
character that occurs on a vector there also occurs as a quotient,
by taking its simultaneous generalized eigenspace and dividing
by the maximal ideal for that character. Averaging over diagonal
principal units extends this quotient to the Jacquet module.
Normalized Frobenius reciprocity gives
\(\pi_w\hookrightarrow i_B^{\mathrm{GL}_N}(\alpha)\).
Sorting by decreasing \(|\xi_i(\varpi)|\) means that the real
exponents of the \(\alpha_i\) are nondecreasing, since
\(|\varpi|=q^{-1}\). Thus \(\pi_w^\vee\) is the Langlands quotient
of the decreasing normalized induction of the \(\alpha_i^{-1}\).
Within a group of equal exponents the inducing unitary characters
give the tempered induction. The local parameter of \(\pi_w\) is
therefore the sum of the characters \(\alpha_i\) through geometric
reciprocity and has zero monodromy. The Tate normalization multiplies
each summand by \(|\cdot|^{(1-N)/2}\), giving exactly \(\xi_i\).
\footnote{\href{https://www.numdam.org/article/SB_1981-1982__24__201_0.pdf\#page=11}
{Rodier, \emph{Repr\'esentations de \(\mathrm{GL}(n,k)\) o\`u \(k\) est un corps \(p\)-adique}, Bourbaki expos\'e 587, \S\S4.3--4.4, Theorem~3},
and \href{https://doi.org/10.1112/plms.12584}
{Thorne, Notation for normalized induction and geometric reciprocity}.}

For a Weil element \(s\) of nonnegative geometric degree, let \(x\)
be the inverse image of its abelianization under geometric Artin
reciprocity. The
\hyperlink{proof.weighted-fit-split-wd}{split comparison}
gives the characteristic and annihilating polynomial
\(\prod_i(X-j_\iota(\xi_i(x)))\) on the actual operator for \(s\).
In residue characteristic \(\ell\) this is the linear operator
on the potentially crystalline periods. At the separator, which
lies away from \(\ell\), this gives finite inertia on the original
representation. A Frobenius lift and the selected inertia element
commute there. In the Frobenius-semisimplified parameter this follows
from the character decomposition, and the omitted unipotent factor
commutes with inertia. For the mixed operator at \(s=\phi\tau^b\)
use both values from each same \(\xi_i\) in this polynomial.
Sorting permutes the unit and uniformizer values together, so the
polynomial is the one attached to the original exponent functional
in the local tie. The projection to \(\pi_w\) leaves the global
semisimple attachment unchanged by uniqueness.
\footnote{\weightedpin{grh-compact/src/Lib/GaloisRep/LocalGlobal/TypeFrobenius.lean}{22-L45}{The polynomial and the two channel tests}.}

\paragraph{The product label and its inertial polynomial.}
\teamanchor{proof.weighted-fit-tame-product}%
Let \(c=\prod_a Q_a\) be the squarefree Green factor of the
type, with distinct monic irreducible \(Q_a\) and nonzero constant
term. It may be \(1\). With \(x_{\eta_b}\) the names of the
determinant characters, write
\[
  C(X)=c(X)\prod_b(X-x_{\eta_b})^{t_b}.
\]
The distinct factors have disjoint roots before the indicated
repetitions. On a nonzero vector of the type projector, at least
one irreducible constituent has a nonzero image of the same
projector. The restriction there to the finite quotient
\(\mathrm{GL}_N(k_w)\) contains the type given by the parabolic
product of the cuspidal Green representations of the \(Q_a\)
and the determinant representations \(\eta_b\circ\det\) of
the indicated ranks.

Write this local representation \(\pi\) as a Langlands quotient
of the standard induction of essentially square-integrable segment
representations \(\Delta_j\). Taking fixed vectors under the
kernel of reduction is exact over \(\mathbf C\). By Iwasawa
decomposition the fixed space of the standard induction is the
finite parabolic induction of the block fixed spaces of the
\(\Delta_j\). In fact a function in the induction is determined
on the maximal compact, and after imposing reduction-kernel
invariance it is a function on the finite group with the stated
parabolic equivariance. The modulus of the parabolic is one on
its integral subgroup. Since \(\pi\) has the required finite
label, it occurs already in this finite induction.

Each \(\Delta_j\) is an unramified twist of a segment of length
\(m_j\) on a supercuspidal \(\rho_j\). It is a subquotient of
induction of the corresponding unramified twists of \(\rho_j\).
Applying the same calculation shows that any nonzero fixed
space of \(\Delta_j\) requires a nonzero reduction-kernel fixed
space of \(\rho_j\). At depth zero this is the finite cuspidal
\(\bar\rho_j\) of its maximal compact type. Indeed the depth-zero
compact induction is from the maximal compact times the center.
In the Cartan decomposition, a noncentral diagonal has a gap
between consecutive valuations. The conjugated reduction kernel
has determinant valuation zero, so its intersection with the
compact times the center lies in the compact. The reduction
of this intersection contains the unipotent radical across the gap.
Cuspidality kills the corresponding fixed values in the compact induction.
The central double coset gives precisely \(\bar\rho_j\).
Consequently every constituent of the finite fixed space of
\(\Delta_j\) has the single finite cuspidal support
\(\bar\rho_j\), with multiplicity \(m_j\).

The distinct Green factors contribute distinct cuspidal
supports, each once. The determinant block of rank \(t_b\)
has support \(\eta_b\) in rank one, repeated \(t_b\) times.
The separability of \(c\prod_b(X-x_{\eta_b})\) in the actual
label makes all these supports disjoint. The finite parabolic
induction ring factors by cuspidal support. In particular the
multiplicity of the given product is the product of the
multiplicities on these separate supports. For every Green
factor, its occurrence forces its sole segment to have length
one. For a determinant block, its occurrence forces
\(\eta_b\circ\det\) to occur in the induction of the finite
fixed spaces on that support. Here the supercuspidal is of rank
one, so \(\Delta_j\) is a determinant twist of the local Steinberg
of rank \(m_j\). Its fixed space is the same twist of the finite
Steinberg. This follows by applying exact fixed vectors to the
Steinberg quotient of the induction of the trivial Borel
character by the larger parabolic inductions and using the
Iwasawa calculation for each term. By Frobenius reciprocity
after twisting by \(\eta_b^{-1}\circ\det\), the trivial character
can occur there only when each Steinberg factor is trivial.
The finite Steinberg is nontrivial in rank greater than one.
Consequently each segment on this support has length one as
well. In the local correspondence the standard Langlands
quotient therefore has zero monodromy. The
\hyperlink{proof.weighted-fit-good-lower-labels}{Green reciprocity
calculation for the lower packet} applies separately to each
supercuspidal singleton of the standard data. If its residue
label has root \(x\) of degree \(e\), the corresponding Weil
character is induced from the unramified extension of degree
\(e\), with the regular character \(T^a\) on residue units,
where \(T(x)=\exp(2\pi i a/(q^e-1))\). The unramified rectifier
and the rank twist do not affect inertia. For a tame symbol
of exponent \(a_s\) at an unramified degree divisible by every
\(e\), the Frobenius-conjugate inertial values are
\(T(x^{q^j})^{a_s}\) for \(0\le j<e\). For a determinant slot the
same calculation in degree one gives the value of \(\eta_b\)
on the Artin unit, repeated \(t_b\) times. Thus the inertial
multiset has exactly the prescribed multiplicities. The
calculation applies to every irreducible admissible local
factor containing the label.%
\footnote{See \href{https://www.numdam.org/item/10.24033/asens.1379.pdf}{Zelevinsky, \emph{Induced representations of reductive \(p\)-adic groups II}, \S\S6,9} for the segment classification,
\href{https://www.e-periodica.ch/cntmng?pid=com-001\%3A1996\%3A71\%3A\%3A10\#page=21}{Moy--Prasad, Theorem~6.11, p.~118} for the depth-zero support,
\href{https://www.uni-muenster.de/FB10/u/pschnei/publ/pap/temtype.pdf\#page=21}{Schneider--Zink, \S5, Proposition~6} for the finite generalized Steinberg fixed spaces, and
\href{https://maths.dur.ac.uk/users/jack.g.shotton/papers/gln.pdf\#page=37}{Shotton, \S6.5, Proposition~6.9} for the factorization of finite induction by cuspidal support.
For depth-zero local Langlands and the geometric reciprocity convention, see
\href{https://doi.org/10.1090/S0894-0347-05-00487-X}{Bushnell--Henniart, Part~I, \S2.2 and Corollary~3.3},
\href{https://doi.org/10.1112/plms/pdp053}{Part~III, Proposition~21} and
\weightedpin{grh-compact/src/Lib/ClassField/Local/ArtinMapAtPlace.lean}{189-L205}{The geometric Artin symbol}.}

For the values on the chosen inertia element, keep the faithful
residue character \(T\) and the chosen prolongation. Take
\(f_0\ge1\) divisible by all the degrees of the \(Q_a\), and
a \((q^{f_0}-1)\)-st root \(z\) of a uniformizer. Write \(b\)
for the tame exponent of \(s\), so
\[
  T\bigl(\overline{s(z)/z}\bigr)=
                  \exp(2\pi i b/(q^{f_0}-1)).
\]
The \hyperlink{proof.weighted-fit-split-wd}{split comparison}
thus gives the characteristic polynomial
\[
     \prod_{x\,:\,C(x)=0}\bigl(X-j_\iota(T(x)^b)\bigr)
\]
on the actual inertia operator, counting roots of \(C\) with
multiplicity. It also gives the determinant, the finite inertial
type and the same labelled jumps, using the degree-zero operator
on potentially crystalline periods at \(w\mid\ell\).
The finite inertia action is diagonalizable in characteristic
zero, so its distinct-root polynomial also annihilates it.

\paragraph{The labels of the fitted catalogue.}
\teamanchor{proof.weighted-fit-quadratic-inertia}%
At each chosen split place the lower packet has the prescribed
squarefree label of degree \(r'\). The character-corrected
action on the upper block scales its roots by the name of the
inverse prescribed tame character. The remaining \(t'\) slots
have the trivial determinant character. Inducing from the
parabolic and applying the finite central projector gives
the character tests of these product labels on the cyclic
spans. At every added place
\(w\in Q\) there are only determinant slots, of ranks \(r'\)
and \(t'\). The marked determinant reads the inverse fitted
character and the complementary determinant reads the trivial
one. The added conductor is disjoint from the chosen places.
Every prescribed power stays nontrivial there, so the fitted
character gives the full tame part on the marked block and
the two determinant names are distinct. The shifted squarefree
label at the old places omits \(1\) by the good tuple condition,
so its roots are disjoint from the trivial slots. The projections
commute with the other finite
levels, so these statements apply to the whole catalogue at
each allowed auxiliary level.

At a chosen place with \(q=g_i\), each label value at inertia
has order dividing \(q^a-1\) for some \(1\le a\le r\), since a
root of the squarefree factor has irreducible degree at most
\(r'\le r\) and the determinant slots are defined over the
base residue field. The finite-image inertia matrices away
from \(\ell\) are semisimple, and \(q^a-1\) divides the factor
\(|\mathrm{GL}_{rt}(\mathbf F_q)|\) of \(G\). Hence
\(N_x(s)^G=I\). For a chosen tame generator, its exponent is
prime to \(q^{r!}-1\), taking \(f_0=r!\) in the tame symbol.
Thus every nontrivial label root still
gives a nontrivial eigenvalue. This is used for the pure row,
where there are no complementary trivial slots. At a chosen
place above \(\ell\) the degree-zero Weil
operators have the same eigenvalue orders. Their eigenvalues
lie in unramified coefficient extensions of degree at most
\(r\), since here \(q=\ell\). Taking the product of the
minimal polynomials of the labelled values, with multiplicity,
gives the common annihilator of degree at most \((r+t)r\)
used for the crystalline test. The normalizing avatar of the
upper model has unramified finite parameter at these split
places. It shifts the labelled jumps by its own algebraic
weights and preserves the inertial Weil eigenvalues. The
additional character of the lower row scales its inertia
labels by the value of its tame avatar, with the corresponding
shift of its labelled jumps. These assertions transport to
the conjugate chosen place by sending each inertial element
\(s\) to \(c s^{-1}c\), where \(c\) is complex conjugation.
For the unshifted upper attachment the polarizing exponent
is \(1-n\). The upper normalizing avatar \(\eta\) adds jumps
\(-u,-r_-\) at a conjugate pair with \(u+r_-=r'\), and
\(\eta^c\eta=\chi^{r'}\) by its algebraic character formula.
Thus the normalized packet satisfies
\(N_x^c\simeq N_x^\vee\otimes\chi^{1-t'}\).
Away from \(\ell\) the cyclotomic factor is trivial on
inertia. Above \(\ell\) it is crystalline and its degree-zero
Weil action is trivial. Thus on inverse conjugation the
inertial operator is transposed, at the conjugate inducing
embedding in the period interpretation. The same labelled
polynomial and root orders hold there, and the labelled
jumps transport by the polarizing weights.

For \(w\in Q\), the squarefree Green factor is \(1\) and
the product-label calculation uses \(f_0=1\). Put \(a(s)=\psi(s)\), the inverse fitted
avatar, on inertia at \(w\). Local reciprocity identifies it
with the character labelling the marked determinant block.
Since \(w\nmid\ell\), the normalizing avatar is trivial on
inertia. The normalized packet matrix \(N_x(s)\) has finite
order and is therefore diagonalizable. The two determinant
slots give the exact relation
\[
 (N_x(s)-a(s)I)(N_x(s)-I)=0.
\]
If \(a(s)=1\), the matrix is the identity by semisimplicity.
Its determinant is \(a(s)^{r'}\). The inverse fitted avatar has
the same value under inverse conjugation, so the same
quadratic relation holds on \(Q\cup\bar Q\) for the same
normalized packet, including nonlift packets.%
\footnote{\weightedpin{grh-compact/src/New/Selmer/BaseLevelOfLabel.lean}
{23-L74}{Application to the whole catalogue and the two types};
\weightedpin{grh-compact/src/New/Selmer/LabelledColumns.lean}
{36-L62}{Orders at the chosen places};
\weightedpin{grh-compact/src/New/Selmer/LabelledColumns.lean}
{127-L155}{Quadratic relation on both halves};
\weightedpin{grh-compact/src/New/Selmer/HeckeStableSubspaces.lean}
{80-L102}{Inverse conjugation on inertia and the period operators};
\weightedpin{grh-compact/src/New/Selmer/AptTupleDefs.lean}
{66-L93}{Normalizing avatar and labelled shift}.}

\paragraph{The labelled weights on the stage and sharp targets.}
\teamanchor{proof.weighted-fit-stage-hecke-weights}%
At an infinite place of the full stage pair, write \(u=r-r_-\).
Use the positions \(p=0,\ldots,n-1\) in descending order on the
definite target and \(q=0,\ldots,r-1\) on the source. The target
highest weight in these positions is
\[
 \lambda'_p=
 \begin{cases}
   \varsigma_q+w_j-u+p-\lfloor t/2\rfloor,& p=q<u,\\
   0,&u\le p<u+t,\\
   \varsigma_q+w_j-u+p-\lfloor t/2\rfloor,&p=q+t,\ q\ge u.
 \end{cases}
\]
To see this from the diagram, the positive columns give its
adjacent differences up to position \(u\), and the negative
columns give the differences at the opposite end in reverse
order. In the interiors these differences are
\(\varsigma_{q-1}-\varsigma_q-1\).
The positive boundary length is
\(\varsigma_{u-1}+w_j-\lfloor t/2\rfloor-1\), and the first
negative weight after the zeros is
\(\varsigma_u+w_j+\lceil t/2\rceil\) when the negative branch
is present. On the selected square \(h_2>\lceil t/2\rceil\).
The active entries satisfy \(\varsigma_q+w_j\ge h_2\) on the positive
branch and \(\varsigma_q+w_j\le-h_2\) on the negative branch.
At the final stage \(u=r\), the positive column formula uses
the next entry of the same row-parameter function,
\(\varsigma_r=h_2-w_{m+1}\). Hence
\(\varsigma_r+w_m=-h_2\) there and the positive boundary
formula also holds at that stage.
The cutoffs in the two column formulas are therefore active
exactly at their boundary differences. Starting at the zero
block and summing the differences gives the displayed highest
weight, with the determinant exponent cancelling the negative
columns on the zero block. Equivalently its half
Harish--Chandra parameter at an active position is
\(\varsigma_q+w_j-u+r/2\), since \(t\) is odd. The zero
positions have the ordinary consecutive half-sum entries
\((n-1)/2-p\).

The unshifted target jumps at \(\tau,\bar\tau\) are
\(p-\lambda'_p,\lambda'_p+n-1-p\). The upper normalizing
avatar adds \(-u,-r_-\), respectively. At the active
positions this gives \(\lfloor t/2\rfloor-(\varsigma_q+w_j)\)
and \(\lfloor t/2\rfloor+(\varsigma_q+w_j)\). At the zero
positions \(p=u+i\) the two entries are \(i,t-1-i\), giving
the same multiset \(0,\ldots,t-1\) at each embedding. The
active entries agree with the normalized row on the definite
lower packet. Its unshifted jump is
\((r-1)/2-\epsilon_\sigma\varsigma_q\), and the lower
normalizer subtracts \((r-t)/2+\epsilon_\sigma w_j\).

For a sharp point of rank \(n=2z\) with character of weight
\(-c_p(\tau)\), the target has one row of length \(c_p(\tau)-z\) and
determinant exponent zero. The source has weight zero and
\(r_-=0\). The upper shift \((-1,0)\) then gives active
jumps \(z-1-c_p(\tau)\) and \(z-1+c_p(\tau)\) and the complementary
multiset \(0,\ldots,2z-2\) at either embedding. The weight-zero
twist \(\vartheta\) gives the same weights. The upper
weights just computed are the weights of the ambient
spaces as well, so the general attachment reads them on
the base, twisted and deeper packets. The decomposition as
representations applies to genuine lifts by their Frobenius
identity.%
\footnote{\weightedpin{grh-compact/src/New/Trace/GoodIdempotent.lean}
{122-L149}{Actual diagrams and the sign cut};
\weightedpin{grh-compact/src/New/Exchange/AlgebraicPartDefs.lean}
{49-L76}{The same parameters on the two spaces};
\weightedpin{grh-compact/src/New/Exchange/AlgebraicPartDefs.lean}
{91-L120}{The sharp one-row weight};
\weightedpin{grh-compact/src/Lib/Weil/Unramified/SphericalHecke.lean}
{150-L157}{Labelled jumps}.}

For a pair with source rank \(r'\), target rank \(n=r'+t'\)
and trial character \(\vartheta\), the upper representation is
normalized by the avatar with geometric Frobenius value
\(j_\iota(A_w)^{-1}\). The lower factor of the
\hyperlink{proof.weighted-fit-stage-hecke-transfer}{spherical transfer}
gives one fixed twist \(T\) of the source attachment and the inverse
avatar \(\psi\) of \(\vartheta\). Taking the trace coefficient and
using Frobenius density gives, for every genuine lift with that source,
\begin{equation}
\label{eq:weighted-fit-galois-string}
 \operatorname{tr}N_\vartheta(g)=
       \psi(g)\operatorname{tr}T(g)+\sum_{0\le a<t'}\chi(g)^{-a}.
\end{equation}
This needs only equality by limits of the Hecke coefficients at
geometric Frobenii. A separate local assertion at spherical places
states that the Hecke polynomial kills Frobenius. The latter will
serve the Hecke tests below. Equality \eqref{eq:weighted-fit-galois-string} for all powers of \(g\)
determines the characteristic polynomial by Newton's identities
over the coefficient field. The trace-nondegenerate attachments in
use have semisimple group-algebra image. Indeed its radical pairs
to zero with the image under the trace, since every product in the
radical is nilpotent. Thus Brauer--Nesbitt gives the actual
representation over \(\overline{\mathbf Q}_\ell\) as the direct sum
of \(T\psi\) and the displayed characters. One may choose the
integral lattice in this representation to be that direct sum.
The compact congruence argument makes this choice on its genuine
lift packets. Its other packets carry their stable integral
lattices in the same local class.
\footnote{\weightedpin{grh-compact/src/New/Selmer/AptTupleDefs.lean}{66-L82}{Geometric Frobenius normalization},
\weightedpin{grh-compact/src/New/Selmer/RowsExist.lean}{24-L55}{Twist-free row and string},
and \weightedpin{grh-compact/src/New/Selmer/RibetElement.lean}{143-L167}{Chosen direct-sum lattice}.}

Let \(\epsilon_\sigma=1\) on the chosen CM type and \(-1\) on its
conjugate. For the full range the normalized jumps are
\[
 u_q(\sigma)=\lfloor t/2\rfloor-
  \epsilon_\sigma\bigl(\varsigma_q(\sigma)+w_j(\sigma)\bigr),
 \qquad 0\le q<r .
\]
The shift of the lower attachment is
\((r-t)/2+\epsilon_\sigma w_j(\sigma)\). For \(q+1<\kappa\) we
have \(\varsigma_q=\kappa/2-1-q\), so
\(|\epsilon_\sigma u_q+w_j|\le t\). Its jump is above the
\(2h_2^2\) small window in absolute value. On the remaining \(m\)
coordinates write \(i=q+2-\kappa\), \(1\le i\le m\). Then
\(\varsigma_q+w_j=(j-i+1/2)h\), with \(h=2h_2\). The modulus of
the shifted jump is between \(h_2-t/2\) and
\((2m-1)h_2+t/2\). The scale inequalities put it in
\([h_2/2,h_2^2]\), and \(m\kappa\le rm\) gives the required
bound on the number of small jumps. At a conjugate pair of
embeddings the two jumps in each coordinate sum to
\(2\lfloor t/2\rfloor\), hence at most \(r\) are negative.
For a sharp point \(p=(\nu,z)\) the jump is
\(z-1-\epsilon_\sigma c_p(\sigma)\), with \(-c_p\) the weight of
\(\nu\). Since \(m<z\le r\) and \(w_j+r-t\le c_p\le w_j\), the
high-window inequality follows from
\(|w_j-c_p+\epsilon_\sigma(z-1)|\le t-r+z-1\).
There is no small jump and at most one negative jump at the two
embeddings. These computations also identify the Tate-dual jumps
as \(-u-1\) in this convention.
\footnote{\weightedpin{grh-compact/src/New/Seesaw/IndexTypes.lean}{121-L134}{Row parameter},
\weightedpin{grh-compact/src/New/Rows/Separation.lean}{29-L41}{Jump windows},
\weightedpin{grh-compact/src/New/Selmer/FitTwists.lean}{243-L247}{Sharp-pair range},
and \weightedpin{grh/src/New/Congruences/GaloisRepresentations/RowsFromEigensystems.lean}{66-L222}{Expanded weight comparison}.}

\paragraph{Integral rows and scalar estimates.}
\teamanchor{proof.weighted-fit-integral-rows}%
A continuous representation of the compact group \(G_E\) into
\(\operatorname{GL}_d(\overline{\mathbf Q}_\ell)\) descends to a
finite coefficient extension. To see this, exhaust
\(\overline{\mathbf Q}_\ell\) by an increasing sequence of finite
Galois extensions. Such a sequence exists by Krasner's lemma:
approximate the coefficients of a separable minimal polynomial
by rational coefficients and lift a nearby root in the same
finite extension. Krasner's inequality places the original root
in the field of the nearby root. The fields generated by roots
of rational polynomials thus suffice. The preimages of the
general linear groups over these fields are closed subgroups that
cover \(G_E\). One is open by the Baire theorem and has finite
index. Adjoin the entries of the matrices of finitely many
coset representatives. The \(O\)-span of the translates of
\(O^d\) is then an invariant lattice: it contains \(O^d\) and
lies in a bounded multiple of \(O^d\) by compactness. This works
over one sufficiently large finite coefficient DVR for any finite
list of packets and additional scalar or root data. An
unramified coefficient enlargement also gives a residue field
of size greater than \(4d+2\) for each row in the list.
Finite free enlargement preserves normalized length. Above \(\ell\) the
inverse avatar \(\psi\) used here is trivial on the decomposition
groups, so the local relations of \(T\psi\) use the same
coefficient model as those of \(T\).

At a chosen primitive type prime \(q=g_i\ne\ell\) of a full
row, the roots are a primitive \(q\)-Frobenius orbit of length
\(r=3^{a+1}\). Scaling by the label of \(\omega^{-j}\) preserves
primitivity. In a cyclic group of order \(q^r-1\), a factor of
order dividing \(M_\omega\), with \(M_\omega^2\mid q-1\), is a power
of a generator whose exponent is divisible by each prime
dividing \(q^r-1\). Multiplying by that factor preserves generator
order. On reduction modulo the coefficient maximal ideal an
\(\ell\)-power factor can disappear. Write \(s=r/3\).
If the residual orbit had size less than \(r\), it would be
fixed by the \(q^s\)-power map. This would force
\(q^{2s}+q^s+1\) to divide a power of \(\ell\). Since
\(q\equiv1\bmod 9\), this integer is \(3\bmod9\) and exceeds
three. It has the prime \(3\) and a different prime divisor.
Hence the residual orbit still has length \(r\).

Let \(\tau,\phi\) be tame inertia and geometric Frobenius at this
prime. In the residual lattice, the eigenvalue of \(\tau\) and
its \(q\)-power orbit give all \(r\) distinct lines, and
\(\phi\) cycles these lines. Interpolation in the single matrix
\(T(\tau)\) gives the projections on the residual eigenlines.
Multiplying them by powers of \(T(\phi)\) gives all matrix units
up to invertible scalars. Nakayama therefore gives
\(O[T(G_E)]=\operatorname{Mat}_r(O)\) as a linear span, for any
stable lattice over a coefficient DVR containing the eigenvalues.
The rank-one span holds directly. In the full row the orbit
and the tame relation imply that \(T(\phi^{-r})\) commutes with
\(T(\tau)\) and \(T(\phi)\), thus is a scalar over the field,
and a unit scalar in the lattice. In rank one use \(\phi^{-1}\).
For \(d=r'\) the cyclotomic value of \(g_*=\phi^{-d}\) is \(q^d\).

In the full row
\(T(\tau^{(q^r-1)/(q-1)})\) is a scalar of order \(q-1\).
In the sharp line the tame scalar comes from the label
\(\omega^{-j}\), of order \(M_\omega/\gcd(M_\omega,j)\).
For each \(\ell\), one of the prime factors \(2,5\) of the fixed
modulus is different from \(\ell\); its \((s_0+6)\)-th power
divides \(M_\omega\), where \(s_0\) is the parameter of the weighted
cyclotomic datum. Thus the residual scalar has order greater than
\(4d+2\), by \(m(4r+2)<2^{s_0+6}\). Denote this scalar inertia
element by \(\delta\) and its scalar by \(y\). Both \(y-1\) and
\(y+1\) are units. The polarization on traces of \(T\) is
\(\operatorname{tr}T(cgc)=\chi(g)^{1-t'}\operatorname{tr}T(g^{-1})\).
The perfect matrix trace pairing with \(T(G_E)\) gives
\(T(cg^{-1}c)=\chi(g)^{t'-1}\alpha I\) whenever \(T(g)=\alpha I\).
In particular it gives the same scalar \(y\) at
\((c\delta c)^{-1}\), and the polar term for \(g_*\) is
\[
 d\,v_\ell(\alpha-\chi(cg_*c)\alpha')
     =d\,v_\ell(q^{dt'}-1)\le d\,v_\ell(G).
\]
Here \(q^{dt'}-1\) divides the indicated general linear group
order because \(dt'\le rt\).

At each old level place above \(g_i\ne\ell\), the twisted label
lacks root \(1\) and its tame eigenvalues \(\lambda\) have
\(\lambda^{q^a-1}=1\), \(a\le r\). Since
\(\lambda-1\) divides \(q^a-1\) in \(O\) for \(\lambda\ne1\), the
two readings of a chosen pair satisfy separately
\begin{equation}
\label{eq:weighted-fit-helper-determinants}
 v_\ell\det(T(\tau)-1)\le d\,c_\ell,\qquad
 v_\ell\det(T((c\tau c)^{-1})-1)\le d\,c_\ell.
\end{equation}
For the divisibility expand \(X^{q^a-1}-1\) at \(1\) and at
\(\lambda\), and use that the value of the quotient by \(X-1\)
at \(1\) is \(q^a-1\). This applies also to the rank-one label
scaled by \(\omega^{-j}\). The tame relation bounds the local
cocycle length at such a prime by \(d\operatorname{length}_O(O/a)\)
plus the length of \(\ker(T(\tau)-1\bmod a)\). Evaluation at \(\tau\)
accounts for the first term. For a cocycle zero there, the relation
shows that its value at a Frobenius lies in that kernel.
These two evaluations determine the cocycle because the wild
pro-\(g_i\) inertia acts trivially and contributes no continuous
maps to the finite \(\ell\)-module. Smith form bounds the kernel
length by the respective determinant valuation.

At a place over \(\ell\), impose the finite linear relations of
the integral crystalline cocycles of the lattice, as in
\(\mathrm{crysRels}\). For \(a\mid\ell^M\), the space obeying
these relations has length at most
\(\operatorname{rank}_O Z_f^1\,\operatorname{length}_O(O/a)\).
In fact its coordinates give a homomorphism to \(O/a\) on the
\(O\)-module spanned by the coordinate evaluations in the dual
of \(Z_f^1\). This module has rank at most
\(\operatorname{rank}_O Z_f^1\). The local Bloch--Kato rank
formula gives this rank as \(d\) plus the number of negative
jumps over that place. Summing over places above \(\ell\)
counts each relevant embedding once and adds at most \(fd\).
There are at most \(4f\) old level places off \(\ell\). The
excess over the balanced baseline, the polar term and the two
additional determinants \eqref{eq:weighted-fit-helper-determinants} at the helper pair total at most
\((4f+3)d\,c_\ell\le10fd\,c_\ell\). At the chosen reading take
the single finite union of the rows indexed by \((j,\text{good
packet on }W_m)\) and by \((j,p)\) for the sharp lines. Apply
the auxiliary-prime and torsion specialization argument to this
union, with \(M_\ell=\lfloor r b_\ell\rfloor+1\) and the
prescribed signed interpolation families. It gives one
\(\theta\), of conductor prime to \(\ell\), and for each of its
first \(m\) powers gives the row cocycle bound \((10fd+1)c_\ell\)
on the same conductor support and the required signed ramified
ratio bounds. The second coefficient reading acts on the
coefficients of that same twist. Its point test can vary.
\footnote{\weightedpin{grh-compact/src/Lib/General/NumberTheory/StableLattice.lean}{75-L108}{Finite coefficient lattices},
\weightedpin{grh-compact/src/New/Selmer/LabelledColumns.lean}{64-L125}{Residual span, tame eigenvalues and scalar},
\weightedpin{grh-compact/src/New/Selmer/RowBalanced.lean}{119-L175}{Balanced local lengths and both helper terms},
and \weightedpin{grh-compact/src/New/Selmer/TwistForRows.lean}{14-L43}{One local specialization for all rows and powers}.}

\paragraph{The local crystalline defect that the row argument uses.}
\teamanchor{proof.weighted-fit-crystalline-duality}%
Fix a finite normalized target catalogue \(N_x\) of one auxiliary
pair and one tuple label. These are the packets at the base and
deeper Taylor--Wiles levels in the separator. At places over
\(\ell\) they have a common normalized Hodge multiset. Away from
\(g_0g_1\) the local lattices are crystalline. Above \(g_i=\ell\)
they are potentially crystalline with the common tame polynomial
condition from their shifted squarefree labels and \(t'\) trivial
slots. At an inertial element, each relevant root is a
Teichm\"uller root from an extension of residue degree
\(a\le r\). The product of its minimal polynomials over
\(\mathbf Q_\ell\), over the at most \(r'+t'\) label roots with
multiplicity, is a common monic annihilator for the Weil
inertial operator, of degree at most \((r+t)r\). The roots also satisfy \(X^L-1\) for
\(L=\operatorname{lcm}_{1\le a\le r}(\ell^a-1)\).
The auxiliary twist primes and the Taylor--Wiles primes lie
away from \(\ell\). On the \(T\psi\) block the auxiliary twist
is trivial locally. On the other blocks the jumps are \(a\)
for \(\chi^{-a}\). A genuine lift has a lattice equal to this
direct sum, so all these blocks lie in the same local class as
the catalogue. In particular every allowed class jump is
attained by a block of that lift.

Here is the passage from the congruence cocycles to the dual
defect. Write \(L\) for a block lattice and \(A=O/\ell^M\), with
\(M\) allowed to vary by block. Let \(C\) be the \(O\)-span of
\((N_x(g)-1)_x\) in the finite sum of matrix lattices \(B\) with
left translation action. A finite cocycle obeying the relations
among these tuples defines an \(O\)-linear map \(C\to L/\ell^M L\).
The cocycle law makes it equivariant. By Artin--Rees
\(C\cap \ell^{kM}B\subseteq\ell^M C\) for large \(k\), so the
map extends by zero on \(\ell^{kM}B\) to an equivariant map
\(Y=C+\ell^{kM}B\to L/\ell^M L\). The rationalization of \(Y\)
is that of \(B\); locally \(Y\) is a lattice in the same class.
The cocycle \(g\mapsto(N_x(g)-1)_x\) in \(Y\) is a rational
coboundary. Its restrictions above \(\ell\) are therefore
crystalline cocycles in \(Y\).

A linear lift \(p:Y\to L\) of the map exists since \(Y\) is free.
The coboundary \((g p-p)/\ell^M\) is integral and is a rational
coboundary. Thus the induced pairing \(Y\times L^\vee(1)\) is
one of the pairings used in the compact cup-orthogonal
condition. Use a generator of the inverse different of
\(O/\mathbf Z_\ell\) for the trace functional to get perfect
duality after restricting scalars. Put \(A_f,A_f^\vee\) for
the reductions of the two integral crystalline cocycle
modules, and \(H\) for the reduction of the cup-orthogonal
part of the dual one. They contain the usual finite
coboundaries in the first two cases. Local Tate duality gives
\((A_f^\vee)^\perp=A_f\) in finite cocycles. Indeed
Bloch--Kato orthogonality gives the inclusion, and the rank
formula gives the exact orders. For local degree \(n_v\) and
finite module \(V=L/\ell^M L\), the two reduced crystalline
cocycle modules have product of orders \(|V|^{2+n_v}\):
the jump pairs are \(u,-u-1\) and the ranks of the two
coboundary parts account for the \(2\). The finite local
Euler characteristic and perfect duality on \(H^1\) give
the same product for complementary cocycle groups.

The restricted cocycles under consideration pair trivially
with \(H\) by the pairing built from \(p\). The quotient by
their intersection with \(A_f\) thus has length at most
that of \(A_f^\vee/H\). This is bounded by the integral
dual crystalline defect for every reduction modulus \(\ell^M\),
which may vary by block. In the
\hyperlink{proof.weighted-fit-ribet-fitting}{first-order calculation}
the row block uses \(P_0=\ell^{M_\ell}\), and the scalar Tate blocks
use \(P=\ell^\nu\) at the triangular depth. Both applications use
the same stable lattice catalogue for the chosen trial twist.
Since reductions of crystalline
cocycles obey the required relations, the same bound applies
to the quotient by those relations. The global quotient
injects into the direct sum of such local quotients. The
bound needed at each place is
\[
 [E_v:\mathbf Q_\ell]\,\mathrm{loss}_{\mathrm{crys}}(\ell)
       \begin{cases}r',&L=T\psi,\\ 1/t,& L=O(-a),\end{cases}
\]
on the normalized \(O\)-length of
\(Z_f^1(L^\vee(1))/(Z_f^1(L^\vee(1))\cap L^{\perp\mathcal Q})\),
where \(L^{\perp\mathcal Q}\) here denotes the compact
cup-orthogonal cocycles in the dual module.
The coefficient is the established local allowance
\[
 \mathrm{loss}_{\mathrm{crys}}(\ell)=10^6\frac m\kappa
       \left(b_1v_\ell(g_0g_1)+v_\ell(A_{\rm crys})\right),
\]
with \(A_{\rm crys}\) the factorial and short interval product
in the fixed definition of \(b_\ell\).
Summing gives \(2f(r'+t'/t)\mathrm{loss}_{\mathrm{crys}}(\ell)
\le4fr'\mathrm{loss}_{\mathrm{crys}}(\ell)\), with no extra
factor for the finite coefficient extension. This proves the
reduction and summation. The
\hyperlink{proof.weighted-fit-crystalline-chart}{local chart and signed letters}
give the inputs to the
\hyperlink{proof.weighted-fit-crystalline-budget}{one-place allowance}.
The \hyperlink{proof.weighted-fit-tame-cup}{tame-place argument}
proves this bound at \(\ell=g_i\) by integral lifting.
The \hyperlink{proof.weighted-fit-small-unramified-cup}{signed cluster argument}
covers \(\ell\le\kappa t\) with \(\ell\nmid N_1d\).
At an unramified crystalline place in the
\hyperlink{proof.weighted-fit-crystalline-fl-width}{Fontaine--Laffaille width range}
the defect vanishes. The \hyperlink{proof.weighted-fit-small-ramified-cup}{sparse cut argument}
covers the small conductor primes. The large unramified completions
use the \hyperlink{proof.weighted-fit-large-two-sizes}{two-size estimate}
or the \hyperlink{proof.weighted-fit-large-few-low}{few-low-jumps estimate}. The numerical inequality
\(\sum_\ell b_\ell\log\ell\le
B_0\) is the established allowance on the datum.%
\footnote{\weightedpin{grh-compact/src/New/Selmer/BlocksAndCocycles.lean}{13-L40}{Graph lattice},
\weightedpin{grh-compact/src/New/Selmer/RowBalanced.lean}{70-L115}{Passage to local duality},
\weightedpin{grh-compact/src/New/Selmer/OnePlaceDefect.lean}{15-L84}{One-place bound},
and \weightedpin{grh-compact/src/New/Selmer/RibetElement.lean}{74-L114}{Summation on these blocks}.}

\paragraph{Local fields and the cyclotomic letters.}
\teamanchor{proof.weighted-fit-crystalline-chart}%
Work at a fixed place \(v\mid\ell\) of \(E\) in the
\hyperlink{proof.weighted-fit-crystalline-duality}{dual defect
calculation}. Fix its embedding \(e:E\hookrightarrow\overline{\mathbf Q}_\ell\),
and identify the image of the completion group with
\(\Gamma_v=\operatorname{Gal}(\overline{\mathbf Q}_\ell/E_v)\).
Write \(e_v,f_v\) for the ramification and residue degrees, so
\([G_\ell:\Gamma_v]=e_vf_v=[E_v:\mathbf Q_\ell]\).
With the chosen coefficient-field identification \(\iota_\ell\), the
complex embedding attached to \(g\in G_\ell\) is
\(\sigma_g=\iota_\ell^{-1}g^{-1}e\). Two such embeddings are equal
exactly when \(g'g^{-1}\) fixes \(E_v\). Let \(\tau_g\) denote the
infinite place of \(\sigma_g\), indexed by the chosen CM embedding
as in the \hyperlink{proof.weighted-thick-places}{thick set construction}.
Let \(\epsilon_g=1\) for that chosen embedding and \(-1\) for its
conjugate. On the cyclotomic generator \(\zeta_E\), the power read by
\(\sigma_g\) is \(\epsilon_g\tau_g\) modulo \(N_1\), using the
identification of the place group with \((\mathbf Z/N_1)^\times\).

The Kummer coordinate for \(E_v\) comes from a compatible tower
of \(\ell\)-power roots of a uniformizer \(\pi_v\). In the period
ring set \(z=[\underline\pi_v]\), with \(\theta(z)=\pi_v\). The
coefficient ring of the chart is \(\Lambda=W(\mathbf F_{\ell^{f_v}})\).
Its image under \(\theta\) is the integer ring of the maximal
unramified subfield of \(E_v\). The Eisenstein equation of \(\pi_v\)
over this ring is \(z^{e_v}+\ell v(z)\) with \(v(z)\) a unit.
Dividing by \(v(z)\) makes it
\(\mathcal E_v(z)=z^{e_v}v(z)^{-1}+\ell\).
Fixing \(\theta(\Lambda)\) and
\(\theta(z)\) is equivalent to fixing \(E_v\). If \(E_v\) is
unramified one can take \(\pi_v=-\ell\), giving
\(\mathcal E_v(z)=z+\ell\).
For the Kummer shifts, the polynomials \(X^{\ell^n}-\pi_v\)
are Eisenstein over \(E_v\), so \(\Gamma_v\) acts transitively
on their roots. Fix a compatible tower of \(\ell\)-power roots
of unity \(\underline\varepsilon\). For each integer \(a\ge0\)
the sets of elements of \(\Gamma_v\) taking each chosen
\(\pi_v^{1/\ell^n}\) to
\(\varepsilon_n^a\pi_v^{1/\ell^n}\) are nested nonempty closed
sets. Their intersection supplies the shift
\(z\mapsto[\underline\varepsilon]^a z\).

The letters are the Teichm\"uller lifts of the nonzero elements
of \(\mathbf F_{\ell^{f_v}}\). Frobenius on a letter is
\(q\mapsto q^\ell\). Enlarge the coefficient ring \(O\) to
contain the integer ring of \(E_v\) and a primitive root
\(\zeta_O\) of order
\(q_v=\ell^{f_v}-1\). Each embedding \(O\to\mathbf C_\ell\) is
the restriction of some \(g\in G_\ell\), and its letter is
\(F(g)=g(\zeta_O)\). Since \(E_v/\mathbf Q_\ell\) is Galois,
the coefficient embedding determines its reading \(\sigma_g\).
Set \(T(q)=1\) if some \(g\) with
\(F(g)=q\) reads a thick place. Letters without such a reading
have \(T(q)=0\). In the unramified case two readings of the same
letter restrict to the same embedding of \(E_v\), because
\(g'g^{-1}\) fixing the primitive root fixes every
\(q_v\)-th root of unity. Thus at unramified places both
the thick indicator and the CM sign are constant on each
letter fiber. Choose the sign of that reading as the sign of
the letter, with an arbitrary choice on empty fibers.

For later use let \(N_1'=N_1/\ell^{v_\ell(N_1)}\). The prime-to-\(\ell\)
roots of unity in \(E_v\) inject into its residue field, so
\(N_1'\mid q_v\). If \(F(g')=F(g)^{\ell^n}\), the quotient
\(g'g^{-1}\) raises every \(q_v\)-th root to the power \(\ell^n\).
Applying this to the \(N_1'\)-part of \(e(\zeta_E)\) gives
\begin{equation}
 \label{eq:weighted-fit-letter-ratio}
 \tau_g/\tau_{g'}\equiv \epsilon_g\epsilon_{g'}\,\ell^n
       \pmod {N_1'} .
\end{equation}
In fact \(g^{-1}e(\zeta_E)\) on that part equals
\((g'^{-1}e(\zeta_E))^{\ell^n}\), so the ratio is in the
indicated order. This is the relation on the actual
coefficient embeddings that is used in the local defect.

At a ramified place the only possible rational primes divide
\(N_1d\). Each prime dividing \(N_1\) splits in \(K\), so
its absolute local Galois group fixes the image of \(K\)
and all readings at this place have the same CM sign.
At \(\ell=d\), the quotient of any two readings acts on
the \(N_1\)-th roots of unity by a power of \(d\), since
these roots are unramified at \(d\). If two readings
are thick, their ratio of infinite places has class one
modulo \(M\). By \eqref{eq:weighted-fit-letter-ratio},
opposite signs would give \(-d^n=1\bmod M\), excluded by
\(d\equiv M/2-1\bmod M\). Thus the thick readings over \(d\)
also have a common sign. For \(\ell\mid N_1d\), take this common
sign on every letter, choosing one sign if there is no thick
reading. All these primes satisfy \(\ell\le\kappa t\) for large
\(x\), since \(d\) and every prime divisor of \(N_1\) are fixed.
\footnote{\weightedpin{grh-compact/src/New/Selmer/DeterminantBound.lean}
{54-L100}{Normalized embeddings and completion degree},
\weightedpin{grh-compact/src/Lib/PadicHodge/WachKisin/KummerChartExists.lean}
{147-L211}{The chart and its ramification and residue indices},
\weightedpin{grh-compact/src/Lib/PadicHodge/WachGamma/KisinModuleDefs.lean}
{154-L171}{Letters and Teichm\"uller representatives}, and
\weightedpin{grh-compact/src/New/Selmer/ThickSpacing.lean}
{20-L72}{Signed comparison of cyclotomic powers}.}

\paragraph{Spacing in the actual Frobenius orbits.}
\teamanchor{proof.weighted-fit-crystalline-spacing}%
Put \(H=r^8\), the half-gap of the weights, and \(Q=\kappa t\)
with \(t=5r\). Let \(\ell\le Q\) and
\(J_\ell=\lfloor\log_\ell(H^3)\rfloor\). For large \(x\),
\(H^3\le 2^{12s}\), since \(r\le3x^5\) and \(2^s=m^{22}\ge x^{11}\).
Suppose that \(q,q^{\ell^{a_1}},q^{\ell^{a_1+a_2}}\) are all thick
letters, where \(1\le a_1,a_2\le J_\ell\). Pick thick lifts with
places \(\tau_0,\tau_{a_1},\tau_{a_1+a_2}\). Their signed ratios from
the last place, by \eqref{eq:weighted-fit-letter-ratio},
are \(\pm\ell^{a_1+a_2}\) and \(\pm\ell^{a_2}\) modulo \(N_1'\).
Here \(a_1+a_2\le24s<27s\).

For odd \(\ell\), reduce the two ratios using
\(\phi:(\mathbf Z/N_1)^\times\to(\mathbf Z/2^{s+1})^\times\).
Their squares eliminate the CM signs. The thick triple then
satisfies the three-term condition with exponents \(2a_2\) and
\(2(a_1+a_2)\), both at most \(54s\). It forces
\(\phi(\tau_{a_1}/\tau_{a_1+a_2})=1\), hence
\(\ell^{2a_2}\equiv1\bmod 2^{s+1}\). Factoring the difference
of squares and successively squaring gives
\(v_2(\ell^{2a_2}-1)=v_2(\ell-1)+v_2(\ell+1)+v_2(a_2)\).
One of the first two valuations is one, so \(2^s\le a_2(\ell+1)\).
This contradicts \(a_2(\ell+1)=O(sx^{10})=o(x^{11})\).

For \(\ell=2\), use the modulus \(3^c\mid N_1'\) with
\(c=\lfloor\log_3 s\rfloor+5\). All thick places have the
same class modulo \(3^c\). Already the last two places would
give \(4^{a_2}=1\bmod3^c\) after squaring their ratio.
The order of \(4\) modulo \(3^c\) is \(3^{c-1}\), since
\(v_3(4^{a_2}-1)=1+v_3(a_2)\). This order exceeds \(27s\).
In either case,
\begin{equation}
 \label{eq:weighted-fit-letter-spacing}
 T(q)=T(q^{\ell^{a_1}})=1 \quad\Longrightarrow\quad
 T(q^{\ell^{a_1+a_2}})=0
 \qquad(1\le a_1,a_2\le J_\ell).
\end{equation}
The argument works for all letters of the chart, since any
violation would have thick lifts at all three positions.
It also allows repeated positions in a Frobenius cycle.
\footnote{\weightedpin{grh-compact/src/New/Selmer/ThickSpacing.lean}
{36-L59}{Short signed arcs, including the prime two},
and \weightedpin{grh-compact/src/New/Selmer/OnePlaceDefect.lean}
{64-L78}{Applying the spacing to the actual letters}.}

\paragraph{The two jump windows and the local degree.}
\teamanchor{proof.weighted-fit-crystalline-windows}%
Let \(B=b_1\) and \(G=b_2\) be the rounded thick and thin
weights. In the row type take the common stage \(j\le m\)
for all its high jumps. Use the window parameters
\[
 \delta=(2j+1)H-2t,\qquad W=4t,\qquad \eta=H^2+1,
 \qquad
 \beta_g=\begin{cases}B+\delta&\tau_g\text{ thick},\\
                       G+\delta&\tau_g\text{ thin}.
         \end{cases}
\]
For the original class, let \(U(g)\) be its allowed jumps
at the coefficient embedding read by \(g\). Each is attained
by a given row or Tate block. The row jumps have either
\(H/2\le |u|\le H^2\) or
\(|\epsilon_g u+w_j(\tau_g)|\le t\). The high case puts
\(-\epsilon_g u\) in \([\beta_g+t,\beta_g+3t]\).
The Tate jumps \(0,\ldots,t'-1\), \(t'\le t\), fit in the
low window. The actual Tate-dual weights are \(u'=-u-1\).
In the high case \(\epsilon_g u'\) belongs to
\([\beta_g,\beta_g+4t]\), since the translation changes the
signed value by at most one. Its low window is
\([-\eta,\eta]\), which also contains all the dual Tate jumps.
Thus the original sign is \(-\epsilon_g\), and the dual
sign is \(\epsilon_g\). Use for \(U'(g)\) precisely the dual
jumps of the block being tested, a subset of the reflected
set \(\{-u-1:u\in U(g)\}\).

The limit on low jumps for a row of rank \(r'\) gives at most
\(\lfloor r'm/\kappa\rfloor\) jumps of its dual in the low
window, counted with multiplicity. Indeed a dual jump of
absolute value at most \(H^2+1\) has original absolute value
at most \(H^2+2\le2H^2\). The negative members of \(U(g)\)
and the negative jumps of this dual row are all at most
\(-H/2\). The original Tate jumps are nonnegative.
For a tested Tate block instead, \(U'(g)\subseteq[-\eta,0]\).
These are the two alternatives in the crystalline bound.

The letter jump sets are the unions of the embedding jump
sets over each fiber of \(F\), with zero added. At a thick
reading they carry the sign and far window of that letter.
For small \(\ell\le Q\), enlarge the low window to
\([-2G,2G]\). All the thin far jumps lie in this window
because \(G+\delta+W\le2G\). This remains true when a thin
reading shares its letter with a thick one. Thus the
window on letter \(q\) has far center \(B+\delta\) if \(T(q)=1\)
and \(G+\delta\) otherwise, with the sign already chosen.
For \(\ell>Q\) the chart is unramified and \(\ell\ge7\), and
the original low window \([-\eta,\eta]\) suffices. In this
case every embedding on a letter fiber has the same thick
indicator and sign.

The numerical hypotheses for the local estimate follow from
the weighted scales. In fact \(Q=O(x^{10})\), \(W=O(x^5)\),
\(H\asymp x^{40}\), \(m\asymp\sqrt{x}\),
\(G\asymp bx/x_1\), and \(B\asymp b\). Consequently
\[
\begin{gathered}
 99\bigl(Q^2+(W+1)Q+Q+W+m\bigr)\le H,\qquad
 99H^3\le G,\qquad 99(GH^3+G)\le B,\\
 2\log(4B)<m,\qquad
 0\le\delta<\eta,\qquad
 99\bigl((\eta+W)Q+\eta+H\bigr)\le H^3 .
\end{gathered}
\]
Also \(e_v\le[E_v:\mathbf Q_\ell]\le2f\le H\) for large \(x\),
since \(f=\varphi(M)3^c2^s=O_d(x^{11}\log x)\).
The degree used at this place is therefore the completion
degree \(e_vf_v\) and is independent of the enlarged coefficient
ring. At \(\ell=g_i\), the congruence \(g_i\equiv1\bmod N_1d\)
makes \(E_v=\mathbf Q_\ell\). The Hodge windows then use one
embedding of \(E\) at \(v\), and the tame case has degree one.
\footnote{\weightedpin{grh-compact/src/New/Selmer/DeterminantBound.lean}
{62-L79}{Row jumps and their multiplicities},
\weightedpin{grh-compact/src/New/Lift/ExponentTheorems.lean}
{37-L45}{Passage to letter jumps}, and
\weightedpin{grh-compact/src/New/Selmer/OnePlaceDefect.lean}
{24-L84}{The window, chart and degree on the tested dual block}.}

\paragraph{Zero crystalline defect in the Fontaine--Laffaille range.}
\teamanchor{proof.weighted-fit-crystalline-fl-width}%
Let \(v\mid\ell\) be a crystalline place in the weighted catalogue with
\(K=E_v\) unramified over \(\mathbf Q_\ell\). At each embedding
\(\tau:K\hookrightarrow\overline{\mathbf Q}_\ell\), let \(I_\tau\)
be the smallest integral interval containing zero and all the allowed
jumps of the original class at \(\tau\). Suppose
\[
             \max_\tau(\max I_\tau-\min I_\tau)\le\ell-2.
\]
Then the cup-orthogonal defect at this place is zero for each of the
tested row and Tate blocks, at every coefficient precision.

Indeed, use one crystalline character \(\xi\) of \(G_K\) of labelled
jumps \(-\min I_\tau\). Such a character with integral unit values is
obtained from the product of the labelled Lubin--Tate characters with
these signed exponents. Its values can be taken in the coefficient
ring \(O\) containing all images of \(K\), by taking the \(K\)-valued
Lubin--Tate character and its embeddings and inverses. Keep the
convention that the cyclotomic character has jump \(-1\). Twisting by
\(\xi\) puts all the allowed jumps and those of the formerly trivial
quotient in \([0,\ell-2]\). The covariant Fontaine--Laffaille
comparison applies to a finite unramified \(K/\mathbf Q_\ell\)
and a coefficient DVR containing all embeddings of \(K\).
For \(0\le h\le\ell-2\) it sends the filtered
\(\mathcal O_K\otimes_{\mathbf Z_\ell}O\)-modules \(M\) of
finite type to Galois modules. The decreasing filtration
consists of \(\mathcal O_K\)-direct summands, with
\(\mathrm{Fil}^0 M=M\) and \(\mathrm{Fil}^{h+1}M=0\).
The maps \(\varphi_i:\mathrm{Fil}^i M\to M\) are semilinear
for arithmetic Frobenius on \(\mathcal O_K\) and linear on
\(O\), with
\(\varphi_i|_{\mathrm{Fil}^{i+1}M}=\ell\varphi_{i+1}\)
and \(\sum_i\varphi_i(\mathrm{Fil}^i M)=M\).
The functor is exact and fully faithful, identifies the
labelled filtration jumps with the Galois jumps, and its image
is closed under subquotients. Every stable lattice in a crystalline
representation with these jumps belongs to that image. Take
\(h=\max_\tau(\max I_\tau-\min I_\tau)\) after the twist, for
both the graph lattice and the extension of the twisted unit
quotient below.%
\footnote{\href{https://www.ma.imperial.ac.uk/~tsg/Index_files/gswc.pdf\#page=26}{Gee--Herzig--Savitt, Lemma~5.1.6(i)},
\href{https://www.numdam.org/item/10.24033/asens.1437.pdf\#page=36}
{Fontaine--Laffaille, Theorem~6.1 and Proposition~7.15},
and \href{https://math.jhu.edu/~savitt/papers/pdfs/FLlifts.pdf\#page=15}
{Gee--Herzig--Liu--Savitt, \S2.3, pp.~15--17, for the coefficient functor}.}

Take the \hyperlink{proof.weighted-fit-tame-cup}{integral graph}
\(\mathcal G_f\twoheadrightarrow Y\) of the equivariant reduction
\(f:Y\to L/\ell^M L\) of any of the tested pairings. Its
rationalization is \(L[1/\ell]\oplus Y[1/\ell]\), so the same
intervals apply. An integral finite class of \(Y\) gives an extension
of \(O\) by \(Y\) with crystalline generic fibre. Twist the graph and
this extension by \(\xi\). In Fontaine--Laffaille modules write
\(D\twoheadrightarrow D'\) for the graph projection and \(C\) for
the module of the line \(O(\xi)\). In each labelled factor choose
a generator \(e_\tau\) of \(C\) and write \(d_\tau\) for its
filtration degree. With \(\tau'\) the factor reached by divided
Frobenius, \(\Phi_{d_\tau}e_\tau=u_\tau e_{\tau'}\) for a unit
\(u_\tau\). Strictness of the filtered extension by \(D'\) gives
a lift \(x_\tau\in\operatorname{Fil}^{d_\tau}\) of each
\(e_\tau\). On this filtered direct sum the extension is specified by
\(a_{\tau'}=\Phi_{d_\tau}x_\tau-u_\tau x_{\tau'}\in D'_{\tau'}\).
Lift \(a_{\tau'}\) to \(\widetilde a_{\tau'}\in D_{\tau'}\).
On the filtered direct sum of \(D\) and \(C\) set
\(\Phi_{d_\tau}\widetilde x_\tau=u_\tau\widetilde x_{\tau'}+
\widetilde a_{\tau'}\). In degrees \(i\le d_\tau\) use
\(\ell^{d_\tau-i}\) times this value. These maps obey the divided
Frobenius relations. Their images together with the images from \(D\)
generate the entire module because the \(u_\tau\) are units. This
extension by \(D\) maps to the given extension by \(D'\), with the
identity on \(C\). Untwisting proves the surjection
\[
 H_f^1(K,\mathcal G_f)\longrightarrow H_f^1(K,Y).
\]
The projection of a lifted class to \(L\) has reduction equal to
the pushforward under \(f\). Its cup with any integral finite class
in \(L^\vee(1)\) vanishes by rational Bloch--Kato orthogonality and
\(H^2(K,O(1))\simeq O\). Reduction modulo \(\ell^M\) gives
orthogonality for the original pairing. This applies to every
pairing and precision in the definition of the local defect.

In the \hyperlink{proof.weighted-fit-crystalline-windows}{weighted
windows} the upper bound for the labelled width is
\(\beta_\tau+W+\eta\), since only one of the two signs of the far
interval occurs at a given embedding. Hence the test holds whenever
\(\ell\ge B+\delta+W+\eta+2\), and the smaller bound with \(G\)
holds when every label of the completion is thin. Here
\(\delta+W+\eta=o(B)\), so \(\ell>2B\) suffices for all large
squares at crystalline unramified places. In particular it covers
the crystalline case \(\lambda_\ell=0\), since otherwise
\(\ell\le 2B\) would divide \((2B)!\). On the full range
\(3B/2<\ell<2B\) it gives the same zero defect while
\(\lambda_\ell=1\) for large squares. Indeed \(\ell^2>2B\),
the factorial contains one multiple of \(\ell\), and the interval
about \(B\) lies below \(\ell\).

There are also genuine applications to the two-size range on thin
completions. Put \(X=B/(4H^3)\). Here \(B=b_1=b+O(N_1)\) and \(H=r^8\),
so \(\log X=x^\varepsilon+O(\log x)\). The ray modulus \(N_1d\) has
\(\log(4d(N_1d)^2)=O_d(\log x)=o(\log X)\). Thus the
\hyperlink{proof.weighted-level-prime-existence}{split-prime existence result}
gives a prime \(X<\ell\le2X\) with \(\ell\equiv1\pmod{N_1d}\), splitting
completely in \(E\). Here \(\ell>Q\),
\(\lambda_\ell\ge\lfloor2B/\ell\rfloor\ge4H^3\), and
\(G+\delta+W+\eta+2=o(X)\) since \(G/B=O(x/x_1)\).
At a completion corresponding in the chosen coefficient reading to a
top embedding, the unique label is thin. The width test gives zero
defect at that completion. The actual norm-selected prime, when a
test there is needed, retains its own completions and labels.

\paragraph{Charging the local crystalline estimates.}
\teamanchor{proof.weighted-fit-crystalline-budget}%
Let
\[
 \lambda_\ell =
 v_\ell\!\left((2B)!\!
   \left(\prod_{\max(B-30H^3,0)<u\le B+30H^3}u\right)^{H^3}\right),
 \qquad K_\ell=\left\lfloor\frac{2\lambda_\ell}{Q}\right\rfloor .
\]
First use the zero defect in the
\hyperlink{proof.weighted-fit-crystalline-fl-width}{Fontaine--Laffaille width range}.
For the other crystalline places, the following estimates charge
the fixed allowance on the same blocks. If the
tested block has rank \(R\), put \(x_L=1\) for the row and
\(x_L=t\) for a Tate block. Put
\(\mu=\lfloor Rm/(\kappa x_L)\rfloor\). The signed windows
above give the \(\mu\) bound for a row, and there are no
low-jump restrictions to check for a Tate block in the
one-sided case.

For \(\ell\le Q\) and \(\ell\nmid N_1d\), the
\hyperlink{proof.weighted-fit-small-unramified-cup}{small unramified argument}
gives the required bound. At a crystalline prime with
\(\ell\le Q\) and \(\ell\mid N_1d\), use the
\hyperlink{proof.weighted-fit-small-ramified-cup}{sparse-letter estimate}
on the actual letter jump sets.
Here its low half-width is \(2G\), its far width is \(W\),
and its combined width can be taken to be \(9G\) since
\(W+8G\le9G\). The far centers are
\(\widetilde B=B+\delta,\widetilde G=G+\delta\), with
\(\widetilde G+W\le2G\).
For \(J_\ell=\lfloor\log_\ell H^3\rfloor\), the spacing
condition is \eqref{eq:weighted-fit-letter-spacing}. The
sign condition holds as above. In the estimate take its
precision parameter to be \(K_\ell\). Indeed
\(\lambda_\ell\ge\lfloor 2B/\ell\rfloor\) gives
\(K_\ell\ge B/Q^2\) for all sufficiently large \(x\). Since
\(\ell^{J_\ell}\le H^3<\ell^{J_\ell+1}\) and \(e_v\le H\), we have
\[
\begin{gathered}
 11(\ell^{J_\ell}+1)\,9G\le\widetilde B,\qquad
 2e_v\widetilde B\le\ell^{J_\ell+1}K_\ell,\qquad
 8e_v\le\ell^{J_\ell+1},\\
 4\bigl(J_\ell+1+d_v'+\lfloor\log_\ell(2\widetilde B)\rfloor\bigr)\,9G
       \le K_\ell .
\end{gathered}
\]
Here \(d_v'=0\) for \(e_v=1\), and
\(d_v'=v_\ell(e_v)+1\) otherwise. The first three inequalities
follow from \(B/G\gg H^3\) and \(H^2\gg Q^2\); the last follows
from \(B/G\gg mQ^2\) and \(\log(2\widetilde B)=o(m)\).
The sparse estimate yields the cup annihilator exponent
\(10^4 K_\ell(J_\ell+1)\) with \(J_\ell+1\le m\).

At a crystalline \(\ell>Q\) with \(\lambda_\ell\ge H^3\), use
the \hyperlink{proof.weighted-fit-large-two-sizes}{unramified estimate with two sizes}. Keep the same
far centers and take the combined width \(W'=W+4\eta\).
We have \(2B<\ell(\lambda_\ell+1)\), obtained from
\(\lambda_\ell\ge\lfloor 2B/\ell\rfloor\). Since
\(K_\ell\ge\lambda_\ell/Q\) in this range, the weighted scales
give
\[
 5W'\le\widetilde G\le\widetilde B,\qquad
 \widetilde B<K_\ell\ell^2,\qquad
 2\widetilde G<K_\ell\ell,\qquad W'\le K_\ell .
\]
For the two inequalities with \(\ell\), use
\(\lambda_\ell\ell>3B/2\) and \(B/G\gg Q\).
For the width use \(H^3/Q\gg H^2\). The exponent supplied is
\(10^4K_\ell(\lfloor\log_\ell(2\widetilde B)\rfloor+1)\),
also at most \(10^4K_\ell m\).

At a crystalline \(\ell>Q\) with \(\lambda_\ell<H^3\), the factorial factor
implies \(\ell>2B/(\lambda_\ell+1)\ge 2B/H^3>9G\).
The interval factor implies that no integer in
\((B-30H^3,B+30H^3]\) is divisible by \(\ell\).
Put \(A=2(\lambda_\ell+1)\). All nonzero integers within
distance \(6(A+\eta+W)\) of \(0,\pm\widetilde G,\pm\widetilde B\)
are prime to \(\ell\). Near \(\pm\widetilde B\) they lie up to
sign in the interval charged by \(\lambda_\ell\). Near
\(0,\pm\widetilde G\) their absolute values are less than \(\ell\).
The \hyperlink{proof.weighted-fit-large-few-low}{few-low-jumps estimate} applies, since also
\(6(A+\eta+W)<\widetilde G\) and
\(\widetilde B+W+\eta\le(\ell-3)A\).
It gives zero defect for the Tate block. For a row, it gives
zero when \(\lambda_\ell<H/5\), since then \(A<H/2\) and
neither jump set has negative members between \(-A\) and zero.
When \(\lambda_\ell\ge H/5\) its defect bound is
\(11 f_v(\mu(A+1)+R(W+1))\) in normalized length. Here
\(A+1\le3\lambda_\ell\) and \(W+1\le K_\ell\), by \(H/Q\gg W\).

For the exponent estimates, an annihilator \(\ell^k\) of
the defect gives normalized length at most \(2e_vf_vRk\),
by the local crystalline cocycle rank bound. Indeed the
rank is at most \(R+e_vf_vR\), and the length of \(O/\ell O\)
is one after normalization. The crystalline cases above
therefore give length at most
\[
 e_v f_v\,10^5(RK_\ell m+\mu\lambda_\ell)
 \le e_v f_v\,\frac{R}{x_L}\,
     \mathrm{loss}_{\mathrm{crys}}(\ell).
\]
The inequality uses \(x_L\le t\) and \(K_\ell\le2\lambda_\ell/(\kappa t)\).
The crystalline deductions use the cup and length estimates
in their stated ranges on the letter sets and dual block above.
The \hyperlink{proof.weighted-fit-tame-cup}{tame-place argument}
gives the same bound at the two chosen primes.
\footnote{\weightedpin{grh-compact/src/New/Lift/SparseThickLetters.lean}
{175-L243}{Cup estimate at sparse cut windows},
\weightedpin{grh-compact/src/New/Lift/SparseThickLetters.lean}
{272-L286}{Two-size cup estimate},
\weightedpin{grh-compact/src/New/Lift/ExponentTheorems.lean}
{61-L99}{Few low jumps},
\weightedpin{grh-compact/src/New/Selmer/CrystallineRank.lean}
{25-L56}{Cocycle rank and normalized length}, and
\weightedpin{grh-compact/src/New/Selmer/CrystallineBound.lean}
{17-L74}{All crystalline branches and their charges}.}

\paragraph{Placement of cuts for the sparse-letter estimate.}
\teamanchor{proof.weighted-fit-crystalline-cuts}%
Write \(\varphi(q)=q^\ell\) for letter Frobenius and \(J=J_\ell\).
A cut is a letter whose next \(J\) successors are thin.
Fix a letter \(q\). If \(\varphi^{-(J+1)}q\) is a cut, it
lies in the required previous \(2J+1\) positions. Otherwise
some \(\varphi^{-a}q\), \(1\le a\le J\), is thick. If it
is a cut, use it. Otherwise a thick successor occurs within
\(J\) steps of \(\varphi^{-a}q\). By the spacing condition,
\(\varphi^{-(a+J+1)}q\) is then a cut. Indeed any thick
successor of this last letter within \(J\) steps would
lie before \(\varphi^{-a}q\) at distance at most \(J\),
giving three thick letters with successive gaps at most \(J\).
This proves existence of a cut in the specified distance.

Let \(\varphi^{-l}q\) be the nearest strictly preceding cut.
If \(\varphi^{-a}q\) is thick for some \(0<a<l\),
it has a thick successor within \(J\) steps, since it is
not a cut. The preceding argument gives a cut at
\(\varphi^{-(a+J+1)}q\) and thus \(l\le a+J+1\).
On the other hand \(l-a>J\), since \(\varphi^{-l}q\)
has only thin successors for \(J\) steps. Thus
\(a=l-J-1\). The only possible thick positions after the cut,
counted backwards from \(q\), are \(0\) and \(l-J-1\).
If at most one thick position occurs or their signs agree,
all far jumps between the cut and \(q\) outside the enlarged
low window have a common sign. Thin jumps are absorbed by
the enlarged low window.
\footnote{\weightedpin{grh-compact/src/New/Lift/SparseThickLetters.lean}
{248-L270}{The sparse placement in the actual letter orbit}.}

At the conductor-dividing places the thick sign is common.
Mixed-sign clusters at small unramified primes use the
signed Wach bound. In the large-prime calculation with two
sizes the mixed-sign case uses the
\hyperlink{proof.weighted-fit-crystalline-long-gap}{two-letter lift}
under its own size bounds.

\paragraph{The chart and the integral pairing.}
\teamanchor{proof.weighted-fit-cup-chart}%
Fix a completion \(K=E_v\), put \(p=\ell\), and write
\(e_K,f_K\) for its ramification and residue degrees. Denote
the completion of \(\overline{\mathbf Q}_p\) by \(C_p\), its tilt
by \(F\), and put \(A=W(\mathcal O_F)\), \(W_F=W(F)\). Choose
a uniformizer \(\pi\) of \(K\) and compatible \(p\)-power roots
\(\pi_a\), and write \(z=[\pi^\flat]\), where \(\pi^\flat\) is the
element of the tilt formed by this tower. If \(e_K=1\) take
\(\pi=-p\). Put \(S=W(\mathbf F_{p^{f_K}})[[z]]\subset A\).
The homomorphism \(\theta:A\to\mathcal O_{C_p}\) maps \(z\) to
\(\pi\). Scale the monic Eisenstein polynomial by a unit of
the coefficient ring so that its constant coefficient is \(p\).
Write the result as \(E_0(z)=pU_0(z)+a_e z^{e_K}\), with
\(U_0(0)=1\) and \(a_e\) a unit, and set \(E=E_0/U_0\).
Then \(E=p+z^{e_K}u(z)\) for a unit \(u(z)\) in \(S\).
The kernel of \(\theta\) is generated by \(E\). Indeed its
reduction mod \(p\) has tilt valuation \(1\), the valuation
of the reduction of a generator of \(\ker\theta\). Hence the
quotient of the two generators is a unit of \(A\).
In particular \(E=z+p\) in
the unramified chart. Write \(E_a=\varphi^a E\) and
\(\theta_a=\theta\circ\varphi^{-a}\). Then the \(E_a\) generate
distinct prime ideals, since \(\pi_a\) has valuation
\(1/(e_Kp^a)\). The fixer of \(\theta(z)\) and of the Witt
coefficient ring is \(\operatorname{Gal}(\overline K/K)\).
For each integer \(j\) there is a \(g_j\) in this group with
\(g_j z=[\varepsilon]^jz\), where
\(\varepsilon=(1,\zeta_p,\zeta_{p^2},\ldots)\) is a compatible
cyclotomic system. At each finite stage the Eisenstein
polynomial \(X^{p^n}-\pi\) is irreducible over \(K\), so
some automorphism sends \(\pi_n\) to \(\zeta_{p^n}^{j}\pi_n\).
The resulting closed sets of automorphisms are nested and
their intersection is nonempty by compactness.

We use \(\tau=\varphi^{-1}([\varepsilon]-1)\) and
\(\Delta=z\tau\). Thus \(E\tau\) is associated to
\([\varepsilon]-1\). For all \(g\in G_K\) the difference
\(gz-z\) is a multiple of \(\Delta E\). In fact
\(gz=[\varepsilon]^u z\) for some \(u\in\mathbf Z_p\), and
\([\varepsilon]^u-1\) is divisible by \([\varepsilon]-1\)
by the binomial expansion in the complete period ring. It follows
that \(gs-s\) is in \(\Delta E A\) for \(s\in S\).
The product \(\Delta\) is a unit of \(W_F\) and its image is a
unit at each de Rham divisor \(\theta_a\) with \(a\ge0\).
For \(0\le a<q\),
\(E_a\mid\varphi^q\tau\), by factoring successive differences
\([\varepsilon]^{p^j}-1\). These are the scalar chart properties
used below. The Frobenius permutation \(\sigma:i\mapsto i^p\)
on the nonzero residue letters and the change of complex
embedding are as in the
\hyperlink{proof.weighted-fit-crystalline-chart}{local chart calculation}.

Let \(Y\) range over the crystalline catalogue \(\mathcal Q\)
and write \(L_2=L^\vee(1)\) for the tested dual block. At an
arbitrary modulus \(p^M\) take an \(O\)-linear form
\(\psi:Y\otimes_O L_2\to O(1)\) whose coboundary divided by
\(p^M\) is an integral finite cocycle \(c_3\). Every equivariant
map \(Y\to L/p^M L\) gives such a form by lifting linearly,
since its divided coboundary is rationally a coboundary.
Let \(h:O\to\mathbf Z_p\) be any \(\mathbf Z_p\)-linear
functional. The pairing used in the cup is
\(b=h\circ\psi\bmod p^M\), with values in
\(T_M=\mathbf Z/p^M(1)\). It is equivariant. This includes
the scaled block maps and the inverse-different trace functional
of the congruence argument. The zero modulus exponent gives
a zero pairing.

Use the convention \(d_g V=(g-1)V\) for an ordinary vector.
The finite cocycles on \(Y,L_2\) give extensions of their trivial
parameter lattices by \(Y,L_2\). We write \(\widetilde Y\) and
\(\widetilde L_2\) for these extensions, using a module splitting
to write their action as
\((v,c)\mapsto(gv+a_c(g),c)\) and
\((v,c')\mapsto(gv+a'_{c'}(g),c')\). One can take the
parameter lattices to be the entire finite cocycle modules.
The extensions are crystalline, since evaluating the finite
cocycles in a crystalline period ring splits them there. The
third extension \(\widetilde T_3\) is the extension of \(O\)
by \((Y\otimes_O L_2)^\vee(1)\) for \(c_3\). Its equivariant
pairing with \(Y\otimes_O L_2\), with values in \(O(1)\), is
\[
   ((d,c),y\otimes v)\longmapsto p^M d(y\otimes v)+c\psi(y\otimes v).
\]
Substitution of \(p^M c_3(g)=g\psi-\psi\) verifies equivariance.
The allowed third jumps at a letter comprise \(0\) and
\(-u-v-1\) for the corresponding allowed jumps \(u,v\) on
the first two sides. In each case the zero jump includes the
trivial parameter. This construction uses a lattice extension
even when the tested cup pairing has finite coefficients.
\footnote{\weightedpin{grh-compact/src/Lib/PadicHodge/WachKisin/KummerChartExists.lean}
{144-L211}{Kummer chart and local degrees},
\weightedpin{grh-compact/src/Lib/PadicHodge/Herr/GaloisLattices.lean}
{35-L94}{Pairing, coefficient functional and third extension}, and
\weightedpin{grh-compact/src/New/Lift/CupProductBound.lean}
{18-L67}{The three finite extensions}.}

\paragraph{A Witt-vector cochain for cup annihilation.}
\teamanchor{proof.weighted-fit-cup-witt}%
Let \(x,w\in W_F\otimes\widetilde Y\) project to
\(1\otimes p^{k_1}c\), and let \(y\in W_F\otimes\widetilde L_2\)
project to \(1\otimes p^{k_2}c'\).
Call a pair of vectors \emph{safe} if the pairing of their
first components, via the \(W_F\)-bilinear extension \(B\) of
\(p^{k_3}b\), is in \(N=W_M(\mathfrak m_F)(1)\) inside
\(D=W_M(F)(1)\). Here \(M\) denotes the length of truncated
Witt vectors, \(\mathfrak m_F\) is the tilt maximal ideal, and
\(W_M(\mathfrak m_F)\) consists of vectors with every coordinate
in \(\mathfrak m_F\). Suppose that for all \(g,h\in G_K\)
the four pairs
\begin{equation*}
\begin{split}
 &(gx-x,\ g(hy-y)),\qquad (w-x,\ gy-y),\\
 &(gw-w,\ g(\varphi y-y)),\qquad
 (\varphi x-w,\ g\varphi y-\varphi y)
\end{split}
\end{equation*}
are safe. This implies that the cup of \(a_c\) and
\(p^{k_1+k_2+k_3}a'_{c'}\) is zero.

Here is the cochain calculation. Write \(X,W,Y_0\) for the
first components of \(x,w,y\), and \(A_g=p^{k_1}a_c(g)\),
\(C_g=p^{k_2}a'_{c'}(g)\). The first components of their
Galois differences are \(\delta_g X=(g-1)X+A_g\),
\(\delta_g W=(g-1)W+A_g\), and
\(\delta_g Y_0=(g-1)Y_0+C_g\). Frobenius acts on the
scalar factor. Set \(T=\varphi Y_0-Y_0\) and choose
\(q\in D\) with \((\varphi-1)q=B(W,T)\).
Surjectivity of \(\varphi-1\) on \(D\) follows from the
Artin--Schreier equation in the algebraically closed tilt,
then induction on the Witt length. Define
\[
   d_g=B(A_g,gY_0)-B(X,\delta_gY_0)+(g-1)q.
\]
Since \(\varphi\) commutes with Galois and fixes the constant
coefficient vectors, expansion gives
\[
 (\varphi-1)d_g=B(\delta_gW,gT)-B(W-X,\delta_gY_0)
                 -B(\varphi X-W,\delta_g(\varphi Y_0)).
\]
All three terms are in \(N\). Also the inhomogeneous cochain
differential is
\[
 g d_h-d_{gh}+d_g=B(A_g,gC_h)-B(\delta_g X,g\delta_hY_0),
\]
by the two cocycle laws. The last term is in \(N\).
The fixed elements of \(D/N\) under \(\varphi\) are precisely
\(T_M\), embedded by constant Witt vectors. Indeed
\(\varphi-1\) is surjective on \(N\) by the convergent sum
\(-\sum_{j\ge0}\varphi^j\) there, the fixed elements in \(D\)
are the Witt vectors over \(\mathbf F_p\) with the Tate twist,
and these intersect \(N\) trivially. Thus \(d_g\bmod N\) is
a \(T_M\)-valued cochain with differential equal to the claimed
cup. The Witt vectors of length \(M\) carry the product of the
valuation topologies of their coordinates. Their Galois action
is continuous and \(N\) is open, so the reduced cochain is
continuous. If the chart is taken over a subgroup of index
prime to \(p\), restriction followed by corestriction multiplies
the finite cup class by that index, which is invertible on
\(H^2(G_K,T_M)\). For the chart above the subgroup is \(G_K\).
Finally \(T_M\simeq\mu_{p^M}\) maps to
\(\overline K^{\times}\), giving the required continuous
multiplicative coboundary. In the other direction the map on
\(H^2\) is injective by the Kummer sequence and Hilbert 90.
\footnote{\weightedpin{grh-compact/src/New/Lift/CupProductBound.lean}
{49-L122}{Safety and cochain passage}, and
\weightedpin{grh-compact/src/Lib/PadicHodge/WachGamma/KisinModuleDefs.lean}
{80-L126}{Artin--Schreier--Witt calculation}.}

\paragraph{The integral period modules used in the cut calculation.}
\teamanchor{proof.weighted-fit-cup-frame}%
Keep the Kummer chart and work on underlying
\(\mathbf Z_p\)-lattices. For a complete discretely valued
extension \(K/\mathbf Q_p\) with perfect residue field,
Bhatt--Scholze's étale realization gives an equivalence between
finite locally free prismatic \(F\)-crystals on \(\mathcal O_K\)
and finite free \(\mathbf Z_p\)-lattices in crystalline
representations of \(G_K\). The functor is compatible with
restriction to \(\mathcal O_C\), with coefficient endomorphisms
and with the crystalline and de Rham period identifications.
Here \(K\) is a completion of the weighted coefficient datum,
and the crystalline lattices occur in the fixed untwisted
row blocks and in the graphs of the tested maps.
For such a lattice \(T\) let \(\mathcal E_T\) be its prismatic
\(F\)-crystal and let
\(M_T=\mathcal E_T(S,(E))\), a finite free \(S\)-module.
Use \(P_T=A\otimes_S M_T\) inside \(W_F\otimes T\), via
the map of prisms \((S,(E))\to(A,(E))\). Write \(S_E\) and
\(A_E\) for the rings with all \(E_a\), \(a\ge0\), inverted.
The localized modules for the points and their exponent tests
are \(\mathcal N_T=S_E\otimes_S M_T\) and \(A_E\otimes_A P_T\).
The Hodge--Tate map in Bhatt--Scholze's Theorem 5.2 is the
Frobenius-shifted map for which \([\varepsilon]-1\) has
nonzero residue. Passage by \(\varphi^{-1}\) replaces their
inverted scalar by \(\tau\) and their divisor by
\(\ker\theta=(E)\). In this convention
\[
 P_T[1/\tau]=A[1/\tau]\otimes T,
 \qquad
 B_{\rm dR}^+\otimes_A\varphi P_T
     =B_{\rm dR}^+\otimes_K D_{\rm dR}(T[1/p])
\]
in the common period representation at \(\theta\). The
second equality uses Remark 5.4 after the same Frobenius
shift. In particular \(P_T\) at \(\theta\) is the integral
de Rham lattice \(B_{\rm dR}^+\otimes T\), since \(\tau\)
is a unit there. The comparisons are functorial, including
the action of the coefficient ring and of \(G_K\).%
\footnote{Bhatt--Scholze,
\href{https://people.mpim-bonn.mpg.de/scholze/PrismaticCrystals.pdf\#page=15}
{\S5, especially Theorem 5.2, Proposition 5.3, Remark 5.4 and Theorem 5.6},
and \href{https://people.mpim-bonn.mpg.de/scholze/PrismaticCrystals.pdf\#page=31}
{Theorem 7.9 and Remark 7.11}. The integral identification
after inverting the indicated scalar is Bhatt--Morrow--Scholze,
\href{https://arxiv.org/pdf/1602.03148\#page=41}{Lemma 4.26}.}

The Frobenius map on \(M_T\) is an isomorphism after inverting
\(E\). Thus the frames \(\varphi^j M_T\) and \(M_T\) differ
by an invertible matrix over \(S[1/E_0,\ldots,1/E_{j-1}]\).
At the divisor \(E_a\) the
\(\varphi^{-a}\)-reading of \(\varphi^j P_T\) is the constant
de Rham lattice when \(a<j\), since all further intervening
divisors are units at \(\theta\). It is the representation
lattice if \(a\ge j\), since the image of \(\varphi^{j-a}\tau\)
is a unit at \(\theta\) in this range. An invariant basis of
forms adapted to the decreasing filtration reads these two
lattices by orders \(0\) and \(-u\), respectively, where \(u\)
is the jump of the corresponding type. A form of type \(u\)
is a combination over \(K\) of basis forms with jumps
\(\le u\). This also holds on a letter after applying its
Teichm\"uller idempotent. The tensor product filtration
gives the tensor order rule in the cut calculation.

These order tests characterize \(\varphi^j P_T\) inside
\(A_E\otimes_A P_T\), with the profile \(\operatorname{st}_j\).
Indeed in a basis of \(\varphi^j P_T\) an element of the
localized module has coefficients with denominators bounded
by a finite product of the \(E_a\). The lattice test at
\(E_a\) says the order of a numerator is at least the
corresponding denominator exponent. The rule
\(\ker\theta_a=(E_a)\), applied repeatedly in the de Rham
DVR, gives divisibility by that power in \(A\). Distinct
prime factors cancel simultaneously. This proves the standard
profile test, including the profile \(1\) in the Tate lattice
of jump \(-1\). For \(\mathbf Z_p(1)\) the vector
\(\tau^{-1}(1\otimes1)\) generates the prismatic lattice
over \(A\). Its Frobenius factor is \(\tau/\varphi\tau\).
It has the étale comparison and the constant de Rham lattice
on its Frobenius pullback, which determine it by Theorem 5.2
in the shifted convention. Consequently multiplication by
\(\tau\) carries the standard lattice into the integral
Tate lattice, also after extending scalars to \(O\).

There is a second description of \(P_T\). A vector of
\(A[1/\tau]\otimes T\) belongs to \(P_T\) if and only if
all its positive Frobenius translates belong to the constant
de Rham lattice at \(\theta\). For the reverse implication,
take coordinates in a basis of \(P_T\), with denominator a
power of \(\tau\). The translates force divisibility of the
numerators by that power at each prime \(\varphi^{-j}E\),
\(j\ge1\). Divisibility by all these primes implies
divisibility by \(\tau\). In fact
\(\tau=\varphi^{-N}(\tau)\prod_{j=1}^N\varphi^{-j}(\xi)\),
where \(\xi=([\varepsilon]-1)/\tau\) is associated to \(E\).
Modulo \(p\) the valuation of the finite product increases
to that of \(\tau\). Subtract a multiple of \(\tau\) with
the same reduction as the numerator, then divide by \(p\).
All the prime divisibilities survive since none of these
primes divides \(p\). Iteration in the complete ring gives
divisibility by \(\tau\); repeating removes its powers.

For a kernel lattice \(T\) in a cocycle extension
\(\widetilde T\), the underlying sequence of \(\mathbf Z_p\)-modules
is split, and the rational de Rham sequence is strict for
the filtration. Apply this description after clearing the
finitely many \(E_a\)-denominators to get
\[
   (W_F\otimes T)\cap (A_E\otimes_A P_{\widetilde T})
      =A_E\otimes_A P_T.
\]
Filtered strictness also passes the order tests to the
kernel, since the typed forms on the kernel extend to
forms on the extension with the same order bound. The
functorial tensor maps take bilinear images into the
corresponding localized period modules.

For \(x\in M_T\), division of \(gx-x\) by \(\Delta=z\tau\)
leaves a vector of \(P_T\). Use the two maps
\((S,(E))\rightrightarrows(A/(z\tau),(E))\) from the
inclusion and from \(g\in G_K\). The ideal \((z\tau)\) is
stable under \(\delta\), since \(\delta(z)=0\) and
\(\delta(\tau)=((1+\tau)^p-1-\tau^p)/p\in(\tau)\).
The divisor \(E\) stays a nonzero divisor in the quotient,
because it is prime in \(A\) and does not divide \(z\tau\).
Derived \((p,E)\)-completeness follows from the cofiber
sequence for multiplication by \(z\tau\) on \(A\), and
\(p\in(E,\varphi(E))\) descends. Its quotient by \(E\) is
\(\mathcal O_{C_p}/(\pi(\zeta_p-1))\), killed by \(p^2\)
since \(v_p(\pi)+v_p(\zeta_p-1)\le2\), including at \(2\).
This is a bounded prism over \(\mathcal O_K\). The two
maps from \(S\) coincide by the chart congruence. Thus the
crystal base-change identifications agree there, giving
the divisibility for the action through étale comparison.
The scalar chart congruence gives \(gE/E-1\in\Delta A\).
Applying \(\varphi^a\) and using \(\Delta\mid\varphi^a\Delta\)
gives \(gE_a/E_a-1\in\Delta A\). Hence each denominator
transforms by a unit congruent to one modulo \(\Delta\),
and the divisibility remains true on \(\mathcal N_T\).
Division by \(\Delta\) preserves all orders at nonnegative
divisors. The normalized Galois difference is therefore a
kernel vector with the depth profile. Frobenius gives
\(\Phi_q^0\) on the ordinary profiles and \(\Phi_q^1\)
on depths by the factor \(E_a\mid\varphi^q\Delta\) for
\(0\le a<q\). This gives the integral inputs for the cut walk.

\paragraph{Crystalline coordinates and their finite jets.}
\teamanchor{proof.weighted-fit-cup-radius}%
The crystalline comparison gives a scalar retraction on jets
and an explicit index bound for the long gap. Put
\(\Lambda=W(\mathbf F_{p^{f_K}})\). Let \(\mathscr R\) be
the \(p\)-adic completion of the divided-power polynomial
ring over \(\Lambda\) with basis \(y^{[n]}=y^n/n!\).
An element is \(\sum b_n y^{[n]}\) with \(b_n\in\Lambda\)
tending \(p\)-adically to zero. Give it Frobenius extending
Witt Frobenius with \(\varphi(y)=p^{p-1}y^p\). This preserves
the ring and reduces to Frobenius. For \(n>0\), both
\(\varphi(y^{[n]})=p^{(p-1)n}(pn)!y^{[pn]}/n!\) and
\((y^{[n]})^p=(pn)!y^{[pn]}/(n!)^p\) vanish modulo \(p\).
For the latter coefficient the valuation is the positive
base-\(p\) digit sum of \(n\). The specified Frobenius thus
defines a \(\delta\)-structure. The ring is \(p\)-torsion-free
and \(p\)-complete, so \((\mathscr R,(p))\) is a bounded
crystalline prism. The completion of the divided-power
envelope defining \(A_{\rm cris}\) has the same
\(\delta\)-ring property. Its Frobenius is inherited from
\(A\) and sends \(E\) to a multiple of \(p\). It therefore
sends each positive divided power of \(E\) to zero modulo
\(p\), as does taking its \(p\)-th power.

For \(j\ge1\) the two maps from \(S\) have the same
coefficient map \(\varphi^j\) on \(\Lambda\) and send \(z\)
to \(py\) and \(0\), respectively. They commute with
Frobenius and send \(E\) to \(p\) times a unit. Here we use
the normalized generator \(E=E_0/U_0=p+z^{e_K}u(z)\)
from the \hyperlink{proof.weighted-fit-cup-chart}{chart},
with \(U_0\in S^\times\) and \(y\) topologically nilpotent.
The base map
\(\mathcal O_K\to\mathscr R/p\) is reduction to the residue
field followed by its \(p^j\)-power automorphism and
inclusion. Both prism maps induce this same base map.
The second map factors through the crystalline prism
\((\Lambda,(p))\) over \(\mathcal O_K\) with that residue
map. For a basis \(m\) of \(M_T\), choose a basis of the
evaluation on this second prism. The crystal identifies
their pullbacks to \(\mathscr R\) by a matrix
\(V(y)\in\operatorname{GL}(\mathscr R)\).

For sufficiently large \(j\) put \(y=z^{p^j}/p\in pA_{\rm cris}\),
where \(A_{\rm cris}\) is the \(p\)-complete divided-power
envelope of \((E)\). The expansion of \((E-p)^q/q!\) in
the products \((E^i/i!)((-p)^{q-i}/(q-i)!)\) belongs to
\(A_{\rm cris}\). Since \(z^{e_K}=u(z)^{-1}(E-p)\), this
gives divisibility of \(z^{p^j}\) by \(p^3\) whenever
\(v_p(\lfloor p^j/e_K\rfloor!)\ge3\). Substitution on the
divided powers is integral, is continuous for the \(p\)-adic
topology and commutes with Frobenius. This is a map of the
crystalline prisms over \(\mathcal O_K\), where the base map
to \(A_{\rm cris}/p\) has the same \(p^j\)-power residue twist.
The first structural map \(S\to A_{\rm cris}\) is now
precisely \(\varphi^j\). The inclusion with Frobenius
\(\varphi^j:(A,(E))\to(A_{\rm cris},(p))\) is a prism
map over the indicated base. In fact
\(\varphi(E)/p=\delta(E)+(p-1)!E^{[p]}\) is a unit, since
\(\delta(E)\) is a unit and \(E^{[p]}\) is nilpotent modulo
\(p\). The same holds for \(\varphi^j(E)/p\), \(j\ge1\).
The evaluation of \(\mathcal E_T\) there is
\(A_{\rm cris}\otimes_{\varphi^j,A}P_T\) by transitive
base change. After extending to \(B_{\rm cris}\), the
element \(\varphi^j\tau=\varphi^{j-1}([\varepsilon]-1)\)
is invertible. The étale comparison therefore identifies
this pullback with \(B_{\rm cris}\otimes T\). Under that
identification a vector in the frame \(m\) is sent to its
scalar Frobenius translate. The resulting relation is
\begin{equation}\label{eq:weighted-cup-prismatic-jet}
            \varphi^j m=C\,V(z^{p^j}/p).
\end{equation}
Here \(C\) comes from \((\Lambda,(p))\). Its columns are
fixed by \(G_K\) because \(G_K\) fixes the embedding of
\(\Lambda\) and the comparisons commute with the Galois
action. An invariant de Rham form therefore takes values in
\(K\) on \(C\). These are the maps of the crystalline
comparison in Bhatt--Scholze's Proposition 5.3. Their
canonical map from the shifted Fontaine prism corresponds
to \(\varphi\) from \((A,(E))\). The additional Frobenius
iterations account for the residue twist \(p^j\) fixed
above. The left side is the actual frame of the level-\(j\)
module in the common étale realization.

For qualitative estimates at the \(a\)-th divisor put
\(K_a=K(\pi_a)\). The completion of
\(\varphi^{-a}(S)[1/p]\) along \(E\) has coefficient field
\(K_a\) and parameter \(E\). Weierstrass division gives
the residue field since \(\theta(\varphi^{-a}z)=\pi_a\)
and \(\pi_a\) has Eisenstein degree \(e_Kp^a\). Hensel's
lemma lifts this separable residue field to the coefficient
field, and the completion maps into \(B_{\rm dR}^+\).
Both \(V(z^{p^j}/p)\) and its inverse have finite jets in
this scalar completion. For a fixed order \(H\), write
\(z^{p^j n}=z^r u(z)^{-q}(E-p)^q\) with
\(q=\lfloor p^j n/e_K\rfloor\) and \(0\le r<e_K\).
Modulo \(E^H\), the last factor is divisible by
\(p^{q-H+1}\) for \(q\ge H\). The denominator of the
\(n\)-th divided power has valuation \(n+v_p(n!)\le2n\).
For large enough \(j\), \(q-H+1-2n\) tends to infinity.
This gives convergence in each finite jet of the scalar
ring with \(p\) inverted, including at \(p=2\). If
\(\varphi^{j+a}m=mH_{j+a}\) is the iterated height matrix,
then
\[
   \varphi^{-a}m=(\varphi^jm)\varphi^{-a}(H_{j+a}^{-1}).
\]
The last matrix has meromorphic jets over \(K_a\) by the
finite height inclusions. Equation
\eqref{eq:weighted-cup-prismatic-jet} now shows that the
evaluation of \(\varphi^{-a}x\) for \(x\in\mathcal N_T\)
by an invariant form has meromorphic jets over \(K_a\).

In the unramified chart, if \(p\ge5\) and \(j\ge2\), the
evaluation \(y=z^{p^j}/p\) already works at this \(j\),
since \(z=E-p\) and \(v_p((p^j)!)\ge3\). Suppose
\(H\le p^j\delta\) for integers \(H,\delta\ge0\).
Both \(p^{2\delta}V\) and \(p^{2\delta}V^{-1}\) then have
order-\(H\) jets that lift to matrices over \(S\). For
\(n<\delta\), the denominator of the corresponding divided
power has valuation at most \(2n\), paid by \(2\delta\).
For \(n\ge\delta\), the coefficient of \(E^h\) for
\(0\le h<H\) in \((E-p)^{p^jn}\), including this scalar
multiplier and denominator, has valuation at least
\(p^jn-h-2n+2\delta\ge0\). It tends to infinity with \(n\).
When \(H=0\) there is no jet restriction. Take matrices
\(P_1,P_2\) over \(S\) lifting the jets of
\(p^{2\delta}V^{-1}\) and \(p^{2\delta}V\). The columns
\(a\) of \(P_1\) have \(Va\) constant modulo \(E^H\).
The columns of \(p^{4\delta}I-P_1P_2\) have
\(Va\equiv0\bmod E^H\). The \(S\)-span of the two sets
contains \(p^{4\delta}S^{\operatorname{rank}T}\). By
\eqref{eq:weighted-cup-prismatic-jet} this proves that
\(p^{4\delta}\) times the level-\(j\) module over \(S\)
is in the \(S\)-span of vectors all of whose form values
are \(G_K\)-invariant modulo \(t^H\). The same holds at
each letter after applying the Teichm\"uller idempotent.
This is the constant needed in the unramified long gap.
The ordered exponent test there is precisely this level-
\(j\) module over \(S\), by the coefficient retraction below.

\paragraph{The integral coefficient retraction.}
\teamanchor{proof.weighted-fit-cup-retraction}%
Take an \(S\)-linear retraction \(\lambda:A\to S\) as the
coordinate of \(1\) in a topological basis. To construct
the basis, choose a basis of \(A/(p,z)\) over
\(\mathbf F_{p^{f_K}}\) containing \(1\) and lift it to
\(A\), using \(1\) as its own lift. The regular sequence
\(p,z\) gives the polynomial associated graded module over
this residue space. Induction through \((p,z)^n\) and
completeness give unique coefficients in \(S\), with finite
support modulo each \((p,z)^n\). Completeness follows
from Witt \(p\)-adic completeness and \(z\)-adic
completeness of \(A/p=\mathcal O_F\). In fact
\((p,z)^n:p=(p,z)^{n-1}\) by reduction modulo \(p\), so
one can lift an element of the \((p,z)\)-adic completion
by successive \(p\)-adic approximations after the first
lift modulo \(p\). Taking the chosen coordinate defines
\(\lambda\) on all of \(A\). Localize
and extend the scalar map in an \(S\)-basis to obtain a
functorial retraction \(\mathcal R:A_E\otimes_A P_T\to
\mathcal N_T\).

In particular \(S_E\cap A=S\) inside \(A_E\), since the
localized retraction fixes \(S_E\) and maps \(A\) to \(S\).
For a point of \(\mathcal N_T\) with the level-\(j\)
standard profile, its coordinates in the frame
\(\varphi^jm\) are in \(A\) by the order tests and in
\(S_E\) by the finite-height change of basis. Thus this
point already belongs to the level-\(j\) module over \(S\).
In particular the radius bound applies to the points
coming from the cutoff operators.

The retraction preserves form orders after subtracting
any constant in \(K_a\). Evaluate at the \(a\)-th divisor
after expressing the vector in an \(S\)-basis and clearing
the finitely many scalar denominators. Approximate the
form values on the basis and the constant, to the needed
order, by elements of \(\varphi^{-a}S[1/p,1/E]\), using
the meromorphic jet calculation. After clearing powers of
\(p\) and of \(E_a\), the claimed vanishing to a prescribed
order is a divisibility by a power of \(E_a\) of a scalar
\(\sum s_j r_j-s_0\), with \(s_j,s_0\in S\) and
\(r_j\in A\). The equivalence holds in \(A\) because
\(\ker\theta_a=(E_a)\). The divisibility survives
\(\lambda\) by \(S\)-linearity. Undoing the denominators
proves the assertion. In particular the retraction preserves
the exponent profiles and commutes with projection to the
trivial parameter. If a vector of the \(S\)-module maps to
zero in that parameter, the strictness at the localized
\(A\)-level and the natural retraction place it in the
\(S\)-module of the kernel.

\paragraph{Integral operators for raising depth.}
\teamanchor{proof.weighted-fit-cup-depth}%
Use the \hyperlink{proof.weighted-fit-cup-chart}{Kummer chart} and the
\hyperlink{proof.weighted-fit-cup-frame}{integral period modules}
for one of the three extensions, denoted by \(T\) in this
paragraph. Write \(\iota_a:A\to B_{\rm dR}^+\) for the map
through \(\varphi^{-a}\), \(F^n=t^nB_{\rm dR}^+\) for its
de Rham filtration for any integer \(n\), and \(Y_a(x)\) for
the value of a form of type \(u\) at \(i[-a]\) on
\(\varphi^{-a}x\). The form is Galois-equivariant, and its
value is \(\iota_a\)-linear in scalars. Recall two properties
of the crystalline coordinates. The finite-jet calculation
puts the residue of an integral form value of a localized
point in \(K_a=K(\pi_a)\). This also holds after subtracting a
constant in \(K\) and shifting the filtration by a power of
\(\iota_a(E_a)\). The
\hyperlink{proof.weighted-fit-cup-retraction}{coordinate projection}
\(\lambda:A\to S\) induces a retraction \(\mathcal R\) from
the ambient localized module to \(\mathcal N_T\). It is
functorial for maps of the lattices and preserves the congruence
\(Y_a(y)-c\in F^n\) for \(c\in K\). This includes preservation
of every order bound. In particular the correction below is
taken in \(\mathcal N_T\) at every step.

Let \(d=0\) when \(e_K=1\), and \(d=v_p(e_K)+1\) otherwise.
For each pair \((a,n)\) with \(a\ge0\) and \(n\in\mathbf Z\),
there is an additive operator on \(\mathcal N_T\) of cost
\[
              c(a,n)=v_p(n)+d+3a,\qquad v_p(0)=0.
\]
On points over constants it multiplies the projected parameter
by \(p^{c(a,n)}\). It preserves all orders and all previously
valid depths and raises depth \(n\) to \(n+1\) at level \(a\).
Every form at this level whose value is invariant
modulo \(F^n\) has invariant value modulo \(F^{n+1}\) on
the output. The same operator works for every such form.

For \(n\ne0\), write \(\gamma=[\varepsilon]\) and choose
\(g\in G_K\) with \(gz=\gamma^{p^a}z\). Set \(\omega_0=gE/E\),
which is a unit of \(A\) with \(\omega_0-1\in\Delta A\), and
let \(\omega=\varphi^a\omega_0=gE_a/E_a\).
Under \(\theta_a\), \(\gamma^{p^a}\) maps to one, so \(g\)
fixes the residue \(\pi_a\) of \(z\), and hence fixes \(K_a\).
The identity \(\gamma-1=E\tau U\) for a unit \(U\in A\) gives
\[
       \theta\bigl((gz-z)/E\bigr)
                   =p^a\theta(U)\theta(\Delta).
\]
Taylor expansion of the Eisenstein polynomial at \(z\),
followed by division by \(E\) and evaluation at \(\theta\), gives
\[
  \theta(\omega_0-1)=E'(\pi)\,p^a\theta(U)\theta(\Delta).
\]
For a unit normalization \(E=E_0/U_0\) of the polynomial
\(E_0\), division by \(E_0\) in Taylor's formula replaces
\(E_0'(\pi)\) by \(E_0'(\pi)/U_0(\pi)=E'(\pi)\).
The multiplying factor \(U_0/gU_0\) in the Galois quotient
has image one under \(\theta\). The derivative of the
Eisenstein polynomial is nonzero at \(\pi\). The valuations of
its nonzero terms of degrees \(j-1\), \(1\le j\le e_K\),
have distinct fractional parts \((j-1)/e_K\). Its valuation
is at most that of the leading term,
\(v_p(e_K)+(e_K-1)/e_K\), and is zero if \(e_K=1\).
In particular \(v_p(E'(\pi))\le d\).

Here \(v_p(\theta\Delta)=1/e_K+1/(p-1)>1/(p-1)\).
If \(v_p(q-1)>1/(p-1)\), the binomial expansion for \(q^p-1\)
has its unique lowest-valuation term \(p(q-1)\). Raising to a
positive exponent prime to \(p\) preserves the valuation of
\(q-1\), and inversion preserves it as well. Applying these
facts to \(q=\theta(\omega_0)\) proves
\[
 v_p\bigl(\theta(\omega_0^n-1)\bigr)
    =v_p(n)+a+v_p(E'(\pi))+v_p(\theta\Delta).
\]
At \(p=2\) the strict bound \(v_2(q-1)>1\) gives
\(v_2(2(q-1))=1+v_2(q-1)<2v_2(q-1)\), so the same identity holds.
Since \(\Delta\mid\varphi^a\Delta\)
and \(\Delta\mid\omega_0^n-1\), set
\(X=(\omega^n-1)/\Delta\in A\). The valuation of
\(\theta_a(\Delta)\) is \(p^{-a}v_p(\theta\Delta)\), computed
from \(\pi_a\) and \(\zeta_{p^{a+1}}-1\). Thus
\[
  v_p(\theta_a X)
  \le v_p(n)+d+a+(1-p^{-a})\bigl(1/e_K+1/(p-1)\bigr)
  \le c(a,n).
\]
Surjectivity of \(\theta_a\) lifts \(p^{c(a,n)}/\theta_a X\)
to \(\nu\in A\). We have obtained \(\nu X\equiv p^{c(a,n)}\)
modulo \(E_a\), with integral \(\nu\).

The congruence improvement uses
\((B_{\rm dR}/F^n)^{G_K}=\operatorname{im}(K)\) for finite
\(n\). Indeed the successive graded pieces are \(C_p(j)\),
whose invariants vanish for \(j\ne0\) by Tate's theorem and
are \(K\) at \(j=0\) by Ax--Sen--Tate. Starting at the first
nonzero term, eliminate the order-zero term by subtracting its
constant in \(K\), then apply vanishing at the other orders
below \(n\). If \(Y_a(x)=C+y\) with \(C\in K\) and
\(y\in F^n\), then \(\iota_a(E_a)^{-n}y\) has residue in
\(K_a\). It is fixed modulo \(F^1\) by the chosen \(g\).
Since \(gE_a=\omega E_a\), this yields
\[
        gy-\iota_a(\omega^n)y\in F^{n+1}.
\]
The scalar \(\iota_a(\Delta)\) is a unit for the de Rham
filtration. It follows that the value on \((gx-x)/\Delta\)
is congruent to \(\iota_a(X)y\) modulo \(F^{n+1}\).
Consequently the value of
\[
            p^{c(a,n)}x-\nu(gx-x)/\Delta
\]
is congruent to \(p^{c(a,n)}C\) modulo \(F^{n+1}\). The
retraction \(\mathcal R\) preserves this congruence and gives the
desired output in \(\mathcal N_T\).

For \(n=0\) take \(N=p^{d+3a}\) and automorphisms \(g_j\)
with \(g_jz=\gamma^jz\), \(0\le j<N\). Before retraction
use \(\sum_j g_jx=Nx+\Delta\sum_j(g_jx-x)/\Delta\).
If \(Y_a(x)\) is invariant modulo \(F^0\), it is integral
and has residue in \(K_a\). Write this residue in the basis
\(1,\pi_a,\ldots,\pi_a^{p^a-1}\) over \(K\), which exists
because \(X^{p^a}-\pi\) is Eisenstein over \(K\).
The residue of the sum on each term with exponent
\(1\le m<p^a\) is zero by the geometric sum in
\(\zeta_{p^a}^m\), since \(p^a\mid N\). Its constant
term belongs to \(K\). The retraction preserves congruence
to that constant modulo \(F^1\). This also applies when
\(a=0\), where every residue already lies in \(K\).

In both cases, a normalized difference \((g_jx-x)/\Delta\)
lies in the ambient localized module and has zero projection
when \(x\) lies over a constant. It retains the original
orders, since the order tests are Galois-stable and division
by \(\Delta\) leaves their bounds unchanged. It has orders
at least the previous depth bounds as well. Its multiplying
coefficient belongs to \(A\), so after retraction the
correction lies in both order submodules. The latter is
Galois-stable, which proves retention of depth. All operations
are integral and commute with the coefficient-lattice action
by functoriality. No infinite sum of depth operators is used.

Suppose a point over \(p^m c\) has orders \(o\) and depth
\(D\), and the target is depth \(D'\ge D\). Take a finite
set \(I\) of pairs containing every \((a,n)\) with
\(D(a,u)\le n<D'(a,u)\) for an allowed jump \(u\) at
\(i[-a]\). Applying the operators at each level in increasing
order of \(n\) gives depth \(D'\) and costs exactly
\(C_I=\sum_{(a,n)\in I}c(a,n)\) in the multiplier of the
parameter. A single operation serves all the forms and jumps
using its pair \((a,n)\). Multiply by the remaining power
of \(p\) to obtain any exponent at least \(m+C_I\).
The resulting point has orders at least
\[
  \begin{cases}
    \max(o(a,u),D'(a,u)),&o(a,u)>0,\\
    \max\bigl(o(a,u),\min(0,D'(a,u))\bigr),&o(a,u)\le0.
  \end{cases}
\]
In fact each value is congruent to a constant modulo the new
depth. When both bounds are positive the constant must be zero.
When the old order is nonpositive, the constant has order at
least zero, giving the other bound. This applies also to a
localized point whose initial orders below the Frobenius level
are zero. The only additional cost comes from the specified
steps. At level \(a\), the nonzero terms have the usual
\(v_p(n)\) charge and every term, including zero, has the
\(d+3a\) charge. In particular the unramified steps at level
zero have zero cost if their range contains no nonzero multiple
of \(p\). For an interval of real length \(L\) in \([-H,H]\)
with \(H\ge1\), the sum of the \(v_p(n)\) over its nonzero
integers is at most \(L/(p-1)+\lfloor\log_p H\rfloor\).
For each \(h\ge1\) up to \(\lfloor\log_p H\rfloor\), count
at most \(L/p^h+1\) multiples of \(p^h\). The zero term in
the step count incurs only its stated \(d+3a\) charge.%
\footnote{\weightedpin{grh-compact/src/New/Lift/ExponentLattices.lean}
{23-L85}{The scalar Galois steps},
\weightedpin{grh-compact/src/New/Lift/HodgeExpansion.lean}
{32-L75}{The normalized step and the improved orders}, and
\weightedpin{grh-compact/src/New/Lift/LongGap.lean}
{153-L207}{The integral depth operators and their finite cost}.}

\paragraph{Lifting ideals and propagation along Frobenius.}
\teamanchor{proof.weighted-fit-cup-lift-ideal}%
Let \(C\) be the trivial parameter lattice of any side, write
\(c_i\) for the letter component of \(1\otimes c\), and denote
the projection to the parameter by \(\operatorname{pr}\).
For a nonnegative integer modulus \(r\), let \(J_l(i;r)\)
be the sum of \((r)\subset A\) and the ideal of scalars taking
every \(c_i\) into the image of the ambient module at letter
\(i\) with standard orders \(\operatorname{st}_l\). The second
term is an ideal because the ambient order conditions are
\(A\)-linear. Equivalently \(f\in J_l(i;r)\) if one \(b\in A\)
works for all \(c\), with a lift of \((f-rb)c_i\) of standard
orders for each \(c\). If \(f\in S\), applying \(\mathcal R\) gives
lifts in \(\mathcal N_T\) with common witness \(\lambda(b)\in S\),
since the retraction commutes with projection and is the identity
on the \(S\)-localized module. The same argument works for
lifting \((f-r\lambda(b))y\) for any standard localized point
\(y\) in the trivial quotient at \(i\). Such \(y\) is integral
in \(A\otimes C\), by the level-zero standard test for the
trivial lattice. Its letter part is generated over \(A\) by
the \(c_i\), so an ambient lift of \((f-rb)y\) exists before
retraction.

The exact ideal \(J_l(i;0)\) records the cokernel multipliers.
If \(p^c\in J_l(i;p^{c+1})\), then \(p^c(1-pb)\) is in the
exact ideal, and \(1-pb\) is invertible in the \(p\)-complete
ring \(A\) by the geometric series. Thus \(p^c\) is in the
exact ideal and \(J_l(i;p^{c+1})=J_l(i;0)\). If
\(p^c E^s\in J_0(i;0)\), lift \(p^c E^s c_i\) and divide by
\(E^s\) in \(\mathcal N_T\). The result lies over \(p^c c_i\)
with orders \(-u-s\) at level zero and \(-u\) above.
The depths satisfy the same bounds by Galois stability of
the order tests. This keeps the \(p\)-charge at \(c\).

The standard profile test identifies the ambient lattices of
orders \(\operatorname{st}_l\) with the Frobenius translates of
the level-zero lattice, so
\[
  \varphi^l\bigl(J_0(i;r)\bigr)=J_l(\sigma^li;r).
\]
In particular \(p^c z^{pa}\in J_1(\sigma i;r)\) if and only
if \(p^c z^a\in J_0(i;r)\). If the jumps at \(i\) are all at
most \(h\ge0\), multiplication by \(E^h\) sends \(J_0(i;r)\)
into \(J_1(i;r)\), since it raises the level-zero order to
\(h-u\ge0\) and preserves the others.

The operation \(\alpha(g-1)/(\Delta E)\) on a multiplier
\(f\in J_1(i;r)\cap S\) stays in \(J_1(i;r)\) for
\(\alpha\in A\) and \(g\in G_K\). Indeed start with a
localized lift of \((f-rb)c_i\) at standard level one and
\(b\in S\). Its level-zero values are integral for the de
Rham filtration and have residue in \(K_0=K\). The Galois
difference increases that order by one. Dividing by \(\Delta\)
preserves the orders and is integral in the ambient localized
module by the frame calculation. Division by \(E\) uses up the
gain at level zero, leaving the other orders unchanged. The
result after multiplication by \(\alpha\) lies over
\[
 \left(\alpha\frac{gf-f}{\Delta E}-r\alpha\frac{gb-b}{\Delta E}\right)c_i.
\]
Both scalar quotients lie in \(A\) by the chart congruence
on \(S\), and the second is independent of \(c\). This
proves the assertion. Any integral multiple of \(f\) can be
subtracted from the resulting multiplier within the ideal.

To lower a positive power \(z^n\) at level one, take
\(g_1z=\gamma z\), let \(j=v_p(n)\), and put
\(\xi=(\gamma-1)/\tau\). The prime-to-\(p\) part of \(n\)
gives a unit geometric sum in \(A\), since \(\gamma\equiv1\)
in the residue field of the local ring \(A\). The factorization
\[
     \gamma^{p^j}-1=\tau\prod_{h=0}^{j}\varphi^h(\xi)
\]
and \(\xi\sim E\) show that \((\gamma^n-1)/(\tau E)\) is a
unit times \(\prod_{h=1}^{j}E_h\). Here \(E_h\equiv p\)
modulo \(z\). Invert that unit as the coefficient \(\alpha\).
On \(f=p^c z^n\) the operation gives
\(p^{c+j}z^{n-1}\) up to an \(A\)-multiple of \(f\), which
can be subtracted. Iterating and then using Frobenius gives
\[
 \begin{split}
 p^c z^N\in J_1(\sigma i;r),\ pM\le N
 \quad\Longrightarrow\quad
 p^{c+v_p(N!)-v_p((pM)!)}z^M\in J_0(i;r).
 \end{split}
\]
Taking \(M=\lfloor N/p\rfloor\) adds zero cost. Starting and
ending at level zero, one can first apply Frobenius and lower
from \(pN\) to \(pM\); the added cost for \(M\le N\) is
\(N-M+v_p(N!)-v_p(M!)\), by Legendre's formula.

To exchange one factor of \(E\) for \(z\), let
\(f=p^c z^b E^n\) with \(n>0\) lie in \(J_1(i;r)\).
Use the scalar depth calculation at \(a=0\) and this \(n\),
so \(gE=\omega E\), \(\omega^n-1=\Delta X\), and
\(\nu X=p^{d+v_p(n)}+E B\) with \(B\in A\). The chart
identity gives \(gz^b=z^b(1+\tau E Q)\) for \(Q\in A\).
Expanding \(gf-f\) yields
\[
  z\nu\frac{gf-f}{\Delta E}
    =p^{c+d+v_p(n)}z^{b+1}E^{n-1}+(zB+\nu Q\omega^n)f.
\]
Subtract the last multiple of \(f\). Iteration gives
\(p^{c+dn+v_p(n!)}z^{b+n}\in J_1(i;r)\). These two
divided differences apply at every prime, and the \(E\)-exchange
uses the chart bound on the derivative through \(d\).

For \(r=p^{\kappa+1}\), a common power of \(z\) belongs to
\(J_0(i;r)\) at all letters. The period-lattice comparison
after inverting \(\tau\) supplies one power \(\tau^b\) taking
the integral lattice of the extension into the standard
level-zero module. Lifting a constant in the underlying
integral splitting and projecting to its letter gives \(\tau^b\)
in the exact ideal at every letter. Modulo \(p\), the ring
\(A/p\) is a rank-one valuation ring and \(z,\tau\) have
positive finite valuation. Thus some power of \(z\) belongs
to \((p,\tau^b)\); raising it to the power \(\kappa+1\) gives
the claimed common multiplier. Suppose all jumps are bounded
above by \(h\) at every letter, and take the least \(B\) with
this property for \(z^B\). At level one the ideals contain
\(z^B E^h\). The identity \(z^{e_K}\in(E,p)\) implies
\(z^{B+e_K(h+\kappa)}\in J_1(i;r)\) for all \(i\), since
each binomial term is a multiple of \(z^BE^h\) or of
\(p^{\kappa+1}\). The descent with no added cost gives the
common power with exponent \(\lfloor(B+e_K(h+\kappa))/p\rfloor\).
Minimality therefore yields
\[
                       (p-1)B\le e_K(h+\kappa).
\]

Finally suppose \(K/\mathbf Q_p\) is unramified with \(p\ge5\)
and that every jump at every letter is at most \((p-3)(s+t)\),
where \(s,t\) are nonnegative integers. Take \(\kappa=s+t\)
and \(h=(p-3)\kappa\). The last inequality gives
\((p-1)B\le(p-2)\kappa\). Since
\((p-1)v_p(B!)\le B\), it follows that
\[
       (p-2)\bigl(B+v_p(B!)\bigr)\le(p-1)B,
       \qquad B+v_p(B!)\le\kappa.
\]
For \(0\le m\le t\), if \(m\le B\) the level-zero lowering
cost is \(B-m+v_p(B!)-v_p(m!)\le\kappa-m\). If \(m>B\),
multiply by \(z^{m-B}\) within the ideal. In either case
\(p^{\kappa-m}z^m\in J_0(i;p^{\kappa+1})\). The case \(m=0\)
removes the modulus by the unit argument. With \(E=z+p\)
the binomial expansion of \(p^s E^t\) now uses only these
exact multipliers. This proves \(p^s E^t\in J_0(i;0)\) for
every letter and each of the three extensions to which the
jump bound applies.%
\footnote{\weightedpin{grh-compact/src/New/Lift/CokernelKillers.lean}
{57-L202}{Lift ideals, monomial operations and unramified bound},
\weightedpin{grh/src/New/Congruences/Lifting/CupProductBound.lean}
{994-L1178}{Level-zero propagation and the exact multiplier},
and \weightedpin{grh-compact/src/New/Lift/LongGap.lean}
{150-L162}{The localized point over the scaled constant}.}

\paragraph{Exponents and safety at a cut.}
\teamanchor{proof.weighted-fit-cup-walk}%
Use any of the three extensions above, with its allowed letter
sets. A point at \(i\) over a parameter \(p^kc\) has
exponents \(o(a,u)\) if each de Rham form of type \(u\) at
\(\sigma^{-a}i\), evaluated on \(\varphi^{-a}x\), has order
at least \(o(a,u)\). Type \(u\) means that on the original
lattice the form has order at least \(-u\). Equivalently, it belongs
to \(\operatorname{Fil}^{-u}\) of the filtered de Rham dual in that
coefficient factor. Let its depth
\(D(a,u)\) mean that \(gx-x\) has those exponents for every
\(g\in G_K\). We use points in the localized period module of
the comparison, where the exponents and depth test the indicated
vectors before reduction. Write \(i[-a]=\sigma^{-a}i\) and
\[
 (\Phi_q^j o)(a,u)=
   \begin{cases}j,&a<q,\\o(a-q,u),&a\ge q,\end{cases}
 \qquad
 \operatorname{st}_q(a,u)=
   \begin{cases}0,&a<q,\\-u,&a\ge q.\end{cases}
\]
The \hyperlink{proof.weighted-fit-cup-frame}{integral period
comparison} gives the following properties.
Differences of points over the same parameter belong to the
kernel, with the minimum of their two exponent profiles.
On a point of depth \(D\), \((g-1)x/\Delta\) belongs to the
kernel with exponents \(D\). The Frobenius iterate \(\varphi^q x\)
is a point at \(\sigma^q i\), with exponents \(\Phi_q^0 o\) and
depth \(\Phi_q^1 D\). For the latter assertion at \(a<q\),
use the vanishing of \(\varphi^q\Delta\) at \(E_a\) in
\((g-1)\varphi^qx=\varphi^q(\Delta)\varphi^q((g-1)x/\Delta)\).
In the localized period module of the Tate lattice, the
standard level-zero exponents \(1\) at every \(E_a\) imply
that multiplication by \(\tau\) gives an integral \(A\)-vector.
These are the precise integral inputs to the following walk.

Suppose that a point of the third extension over \(p^{k_3}\)
has exponents \(Z(a,w)\) antitone in the jump \(w\), and that
kernel vectors on the two sides have exponents \(D,D'\) at
the same letter. If for every pair of allowed jumps
\begin{equation}\label{eq:weighted-cup-safe-exponents}
       Z(a,-u-v-1)+D(a,u)+D'(a,v)>0,
\end{equation}
then multiplying either kernel vector by \(\Delta\) gives a
safe pair. To see this, use filtered de Rham bases on both
sides. In a bilinear map, the pullback of a form of type \(s\)
uses only pairs of basis forms with \(u+v\le s\), by the tensor
product filtration. Thus the tensor vector has exponents at
least \(1-Z(a,-s-1)\) for type \(s\), by antitonicity and
\eqref{eq:weighted-cup-safe-exponents}. In the final pairing
to the Tate lattice a third type \(w\) can contribute together
with tensor type \(s\) only if \(w+s\le-1\). The inequality
\(Z(a,w)\ge Z(a,-s-1)\) gives order at least \(1\) on each
term. Hence \(\tau\) sends the paired vector into the integral
Tate lattice, and multiplying by \(z=[\pi^\flat]\) sends it
into Witt vectors with coordinates in \(\mathfrak m_F\).
Reduction of the third pairing and application of \(h\) then
give the claimed safety for \(p^{k_3}b\). Distinct letters
pair to zero, since their Teichm\"uller eigenvalues have unit
difference and the pairing is balanced for the action of
\(\zeta_O\). Thus safety on matching letters gives safety
for sums over all letters.

Let \(\mathcal C\) be a set of cut letters meeting each cycle
and let \(q_i\) be the least positive distance from \(i\) to
a preceding cut. At each cut \(i\) take two points \(x_i,y_i\)
over \(p^{k_1}c,p^{k_2}c'\), with exponents \(o_i,o_i'\) and
depths \(D_i,D_i'\). At each cut also take a third point over
\(p^{k_3}\), with antitone exponents \(Z_i\) in the jump.
For each of the two points at \(i\), allow its own distance
\(r\ge0\) to a cut for use of that third point. For the first
point require, at every \(a,u,v\),
\begin{equation}\label{eq:weighted-cup-walk-base}
 Z_i(a,-u-v-1)+D_i(a,u)+\min(-v,o_i'(a,v))>0.
\end{equation}
For the second point require the interchanged condition.
Designate one side at each cut. If the first is designated,
require also \(o_i\ge\operatorname{st}_{q_i}\) and
\begin{equation}\label{eq:weighted-cup-walk-cross}
 (\Phi_r^0 Z_{i[-r]})(a,-u-v-1)+D_i(a,u)
       +\min\bigl(o_i'(a,v),(\Phi_{q_i}^0 o'_{i[-q_i]})(a,v)\bigr)>0.
\end{equation}
If the second is designated, require the interchanged
condition with its corresponding \(r\). This includes the
case \(r=0\) in which the direct third point serves both
inequalities.

At all letters put \(X_i,Y_i\) equal to the direct points at
the latest cut, transported by Frobenius. The distance is
zero at each cut and is \(q_i\) elsewhere. At a cut \(i\),
\((\varphi X)_i\) and \((\varphi Y)_i\) are the transports
from the preceding cut at positive distance \(q_i\). Away
from the cuts they equal \(X_i,Y_i\). Choose \(W_i=X_i\)
where the first side is designated and
\(W_i=(\varphi X)_i\) where the second is designated, with
\(W_i=X_i=(\varphi X)_i\) off the cuts. Summing gives \(x,w,y\)
with the required constant projections.

The first safety pair in the cochain calculation follows
from \eqref{eq:weighted-cup-walk-base} at the latest cut,
using the first Galois difference divided by \(\Delta\).
At levels before that cut's transport distance the three
orders are \(1,0,0\). The other three safety pairs need
only be tested at cuts. If the first side is designated,
\(W_i-X_i=0\). For \((gW_i-W_i,\ g((\varphi Y)_i-Y_i))\)
use \eqref{eq:weighted-cup-walk-cross}, since the second
difference has the minimum of its two exponent profiles.
For \(((\varphi X)_i-W_i,\ (g-1)(\varphi Y)_i)\) use the
second point's condition \eqref{eq:weighted-cup-walk-base}
at the preceding cut and transport it. At \(a<q_i\) the
orders are \(0,1,0\), since \(o_i\ge\operatorname{st}_{q_i}\).
At \(a\ge q_i\) the same inequality gives \(o_i(a,u)\ge-u\),
so the minimum required for the difference follows from the
\(\min(-u,o)\) in \eqref{eq:weighted-cup-walk-base}.
If the second side is designated,
\((\varphi X)_i-W_i=0\). The pair with \(W_i-X_i\) then
uses the interchanged condition \eqref{eq:weighted-cup-walk-cross},
and the pair with \(gW_i-W_i\) uses the transported first
point's condition \eqref{eq:weighted-cup-walk-base}. At the
levels \(a<q_i\), safety again follows from the order \(1\)
of the transported depth and order \(0\) of the difference
on the designated side. This proves the four safety tests
and hence the scaled cup assertion. A larger exponent
works by bilinearity. The same proof with every letter a cut
applies to estimates using only the first preceding step.
\footnote{\weightedpin{grh-compact/src/New/Lift/HodgeExpansion.lean}
{134-L207}{Frobenius and tensor exponents}, and
\weightedpin{grh-compact/src/New/Lift/LongGap.lean}
{198-L313}{Strictness, safe filtered products and the cut calculation}.}

\paragraph{The length of the resulting local quotient.}
\teamanchor{proof.weighted-fit-cup-length}%
Suppose the cup assertion with exponent \(k\) holds for
every \(Y\in\mathcal Q\), every \(M\), every integral pairing
above and every finite cocycle of \(Y\). A tested cocycle
\(a'\in Z_f^1(L_2)\) then satisfies
\(p^ka'\in L_2^{\perp\mathcal Q}\) by the definition of the
orthogonal module. The same conclusion with exponent zero
holds on any fixed \(\mathbf Z_p\)-submodule of finite
cocycles for which the uniform cup assertion applies. Such
a submodule may be used to bound the length even if it was
constructed as a \(\mathbf Z_p\)-module.

For \(R=\operatorname{rank}_O L_2\) and coefficient field
\(F_O=O[1/p]\), put \(V_O=L_2[1/p]\). The finite cocycles
form a free \(O\)-module. They are finitely generated by
the finiteness of \(H^1(K,L_2)\) over \(O\) and that of the
integral coboundaries, and are torsion-free as a submodule
of the continuous cochains. For a de Rham \(F_O\)-representation
\(V\) of \(G_K\), Bloch--Kato's Corollary~3.8.4 gives
\[
 \dim_{F_O}H_f^1(K,V)=\dim_{F_O}H^0(K,V)
       +\dim_{F_O}\bigl(D_{\rm dR}(V)/\mathrm{Fil}^0\bigr).
\]
Their duality in \S3.8 identifies \(H_f^1(K,V)\) and
\(H_f^1(K,V^\vee(1))\) as exact annihilators under the local
Tate pairing. The row block \(V_O\) and its Tate dual here
are crystalline, so these apply. Adding the integral
coboundaries of rank \(R-\dim_{F_O}H^0(K,V_O)\) gives
\[
 \operatorname{rank}_O Z_f^1(L_2)
 =R+\dim_{F_O}\bigl(D_{\rm dR}(V_O)/\mathrm{Fil}^0\bigr)
 \le (1+[K:\mathbf Q_p])R .
\]
Thus an
exponent \(k\) bounds the \(p\)-normalized length of the
defect by \((1+[K:\mathbf Q_p])Rk\le2e_K f_K Rk\).
Indeed the defect is a quotient of
\(Z_f^1(L_2)/p^kZ_f^1(L_2)\), whose \(O\)-length is the
rank times \(k\,\operatorname{len}_O(O/p)\).
For a finite \(O\)-module \(D_0\) the same normalized
length is \(\operatorname{len}_{\mathbf Z_p}(D_0)/[F_O:\mathbf Q_p]\).
This follows by taking the cardinality of both modules from
the residue fields. If a \(\mathbf Z_p\)-submodule of finite
cocycles of index at most \(p^A\) lies in the cup-orthogonal
part, the normalized length is therefore at most
\(A/[F_O:\mathbf Q_p]\). The ramification degree of the
coefficient ring does not give a further factor.

For the finite cocycle application, \(Z_f^1(L_2)\) is
\(p^M\)-saturated in the cochains with coefficients in
\(L_2\): division by \(p^M\) preserves the cocycle identity
and the rational finite condition. Hence its reduction
has cardinality \(|O/p^M|^{\operatorname{rank}_O Z_f^1(L_2)}\).
The corresponding rank for \(L\) and that for \(L_2\)
sum to \((2+[K:\mathbf Q_p])R\), since the complementary
jumps are \(u\) and \(-u-1\). Both reductions contain all
finite coboundaries. The cup of finite rational classes
vanishes after inverting \(p\), and integral cup vanishes
as well because \(H^2(K,O(1))\simeq O\) is torsion-free.
Reduction and the perfect pairing defined by the
inverse-different trace give orthogonality of the two
reduced subgroups. For the finite \(G_K\)-module \(A=L/p^M L\)
and its Cartier dual \(A^{D}=\operatorname{Hom}(A,\mu_{p^\infty})\),
local Tate duality pairs
\(H^j(K,A)\) with \(H^{2-j}(K,A^D)\) perfectly into
\(\mathbf Q/\mathbf Z\). The finite local Euler characteristic
is \(|H^0(K,A)|\,|H^2(K,A)|/|H^1(K,A)|
=|A|^{-[K:\mathbf Q_p]}\). This applies to the underlying
\(\mathbf Z_p\)-modules via the stipulated trace pairing.
The exact product of
cardinalities shows they are full annihilators in finite
cocycles. More explicitly the product of the two
coboundary orders times \(|H^1(K,L/p^M)|\) equals
\(|O/p^M|^{(2+[K:\mathbf Q_p])R}\). It follows that the
quotient of any tested cocycle group by its intersection
with the finite subgroup has order bounded by the index
of the reduction of the integral cup-orthogonal part in
the reduction of \(Z_f^1(L_2)\). This index is bounded by
the integral defect just estimated.%
\footnote{Bloch--Kato,
\href{https://math.stanford.edu/~conrad/BSDseminar/refs/BKTamagawa.pdf\#page=23}
{\S3.8 and Corollary~3.8.4, pp.~355--356}, and
Milne, \cite[Corollary~I.2.3 and Theorem~I.2.8]{f1milne-adt},
for the finite local duality and Euler characteristic.
\weightedpin{grh-compact/src/New/Lift/ExponentTheorems.lean}
{44-L99}{Uniform cup assertion on the universal modules},
\weightedpin{grh-compact/src/New/Selmer/CrystallineRank.lean}
{24-L56}{Local degree and normalized length}, and
\weightedpin{grh-compact/src/New/Selmer/RowBalanced.lean}
{68-L115}{The finite cocycle application}.}
Summation over the actual places uses
\(\sum_{v\mid p}e_vf_v=[E:\mathbf Q]\).

\paragraph{Lifting across thin Frobenius successors.}
\teamanchor{proof.weighted-fit-crystalline-thin-successors}%
Use the \hyperlink{proof.weighted-fit-cup-chart}{Kummer period
coordinates} over a local field \(K/\mathbf Q_p\). Write
\(e_K\) for its ramification index and \(d_K\) for the chart
bound \(d_K=0\) if \(e_K=1\), and \(d_K=v_p(e_K)+1\)
otherwise. Denote Frobenius on letters by \(\sigma\).
Consider one crystalline extension of a trivial lattice by a
crystalline lattice, and include zero in its allowed jumps at each
letter. Suppose all its jumps are at most \(\omega\), and all jumps
at \(\sigma i,\ldots,\sigma^J i\) are at most \(\omega'\). If
\begin{equation}
\label{eq:weighted-fit-thin-successor-cost}
 p^{J+1}\ge 8e_K,\qquad
 k_0\ge
 2Jd_K\omega'+\frac{(2J+4)\omega'}{p-1}
                  +\frac{8e_K\omega}{p^{J+1}},
\end{equation}
then \(p^{k_0}\) lifts every constant at \(i\) to the period module
with the standard exponents of level zero.

Here and below use the
\hyperlink{proof.weighted-fit-cup-lift-ideal}{lifting ideals}
at the fixed modulus \(p^{k_0+1}\). The least exponent \(A\)
such that \(z^A\) belongs to the level-zero ideal at every
letter satisfies \((p-1)A\le e_K(\omega+k_0)\), by the
initial multiplier calculation with global jump bound \(\omega\).

Start with \(z^A\) at \(\sigma^J i\). If a multiplier
\(p^c z^n\) lies in the level-zero ideal at a successor whose jumps
are at most \(\omega'\), multiply by \(E^{\omega'}\) to reach
level one. Exchanging the Eisenstein factors there gives
\(p^{c+d_K\omega'+v_p(\omega'!)}z^{n+\omega'}\).
Descending again costs no additional \(p\)-power and replaces
the exponent by \(\lfloor(n+\omega')/p\rfloor\).
After the \(J\) steps, there is a multiplier \(p^c z^n\)
at level zero at \(i\) with
\[
     c\le J\bigl(d_K\omega'+v_p(\omega'!)\bigr),\qquad
     n\le A/p^J+\omega'/(p-1).
\]
Frobenius transports this multiplier to \(p^c z^{pn}\) at level
one at \(\sigma i\). Descending to exponent zero at \(i\) costs
\(v_p((pn)!)=n+v_p(n!)\le 2n\). Legendre's formula and
\eqref{eq:weighted-fit-thin-successor-cost} bound the total by
\[
\begin{aligned}
 c+2n
 &\le Jd_K\omega'+(J+2)\frac{\omega'}{p-1}
            +\frac{2e_K(\omega+k_0)}{(p-1)p^J}\\
 &\le Jd_K\omega'+(J+2)\frac{\omega'}{p-1}
            +\frac{4e_K\omega}{p^{J+1}}+\frac{k_0}{2}
 \le k_0 .
\end{aligned}
\]
Thus \(p^{k_0}\) belongs to the ideal taken modulo \(p^{k_0+1}\).
The \(p\)-adic unit argument for the lifting ideals gives the
same multiplier in the exact ideal.
\footnote{\weightedpin{grh-compact/src/New/Lift/CokernelKillers.lean}
{87-L188}{Level changes, scalar exchanges and the initial
multiplier}, and
\weightedpin{grh/src/New/Congruences/Lifting/CupProductBound.lean}
{1186-L1242}{Iteration through successors and the resulting bound}.}

\paragraph{The small conductor-dividing places.}
\teamanchor{proof.weighted-fit-small-ramified-cup}%
Fix an actual place \(v\mid\ell\) with \(\ell\mid N_1d\) and
\(p=\ell\le Q\). The fixed block \(L\) is the row of rank \(r'\)
or a scalar Tate block. For any crystalline lattice \(Y\)
in the catalogue class, pair its integral finite cocycle
with one on \(L^\dagger=L^\vee(1)\).
Use an arbitrary integral lift \(\psi:Y\otimes_O L^\dagger\to O(1)\)
of the pairing modulo \(p^M\), followed after reduction by the
given \(\mathbf Z_p\)-linear coefficient functional.
The first two extensions of trivial lattices are those of
the two cocycles. The third has kernel
\((Y\otimes_O L^\dagger)^\vee(1)\) and cocycle
\(g\mapsto(g\psi-\psi)/p^M\). This is integral by equivariance
modulo \(p^M\), and rationally a coboundary, so the third extension
is crystalline. This construction applies, in particular, to the
congruence maps from the
\hyperlink{proof.weighted-fit-crystalline-duality}{row and Tate
graph lattices} at their given moduli.

Keep the letters and Frobenius order of the
\hyperlink{proof.weighted-fit-crystalline-chart}{actual completion}.
Put \(b_*=B+\delta\), \(g_*=G+\delta\), \(V=W=4t\),
\(\eta_0=2G\), \(D=9G\), and \(J=J_\ell\).
At letter \(i\) write \(\beta_i=b_*\) if it is thick,
and \(\beta_i=g_*\) otherwise. The common thick sign over \(v\)
lets us designate one side as the negative side at every letter,
and one as the positive side. Their allowed jump sets are in
\[
 [-\eta_0,\eta_0]\cup[-\beta_i-V,-\beta_i],
 \qquad
 [-\eta_0,\eta_0]\cup[\beta_i,\beta_i+V],
\]
respectively. Thin jumps are included in the low interval even
in a letter fibre that also contains a thick reading.
With \(u\) on the negative side and \(v'\) on the positive side,
the third extension has jumps in the set comprising zero and
\(-u-v'-1\). Since \(v'\ge-\eta_0\) and
\(u\ge-\max(\eta_0,\beta_i+V)\), all three jump sets are bounded
above by \(\beta_i+D\), where \(D\ge V+4\eta_0\).

Call a letter a cut if its next \(J\) successors are thin.
Every letter \(i\) has a cut at positive backward distance at
most \(2J+1\). In fact if \(\sigma^{-(J+1)}i\) is a cut we
can use it. Otherwise some thick letter \(j=\sigma^{-a}i\)
occurs at a distance \(1\le a\le J\). If \(j\) is a cut we
can use it. If it has a thick successor within \(J\), each
of its preceding \(J\) letters must be thin by
\eqref{eq:weighted-fit-letter-spacing}. Then
\(\sigma^{-(a+J+1)}i\) is a cut. This proof also applies if
positions in an orbit repeat. Write \(l_i\) for the least
positive backward distance to a cut.

At every cut use \eqref{eq:weighted-fit-thin-successor-cost}
on the three extensions with \(\omega=2b_*\),
\(\omega'=2D\) and \(k_0=16(J+1)K_\ell\).
Indeed \(\beta_i+D\le2b_*\) everywhere and
\(g_*+D\le2D\) on thin successors. The
\hyperlink{proof.weighted-fit-crystalline-budget}{small-prime
inequalities} give \(p^{J+1}\ge8e_K\),
\(2e_Kb_*\le p^{J+1}K_\ell\), and
\[
 4\bigl(J+1+d_K+\lfloor\log_p(2b_*)\rfloor\bigr)D\le K_\ell.
\]
Using \(p-1\ge1\), the cost on the right side of
\eqref{eq:weighted-fit-thin-successor-cost} is at most
\((2J+10)K_\ell\le k_0\). Thus all three extensions give
points over \(p^{k_0}\) times the respective constants,
with standard exponents and the same Galois-difference depth
\(-u\) on jump \(u\) at every level.

For the first two points raise the depth at level \(a\) to
\(-u+C_a\), where
\[
  A_a=\eta_0\,\mathbf1_{a\le 2J},\qquad
  C_a=\eta_0\,\mathbf1_{a\le 4J+1}.
\]
Here use the \hyperlink{proof.weighted-fit-cup-depth}{integral
depth operators}. At each affected level the required integers
lie in two intervals of radius \(D\), centred at zero and at
the signed opposite of the far centre of the jump window at
\(\sigma^{-a}i\). They have absolute value at most \(2b_*\).
Their total cost at that level, with the cost of zero
defined using \(v_p(0)=0\), is at most
\[
 (4D+2)\bigl(\lfloor\log_p(2b_*)\rfloor+d_K+3a\bigr)
   \le 18K_\ell .
\]
The last inequality uses \(D\ge1\) and \(a\le4J+1\).
The sum over at most \(4J+2\) levels is at most
\(72(J+1)K_\ell\). Hence both points can be taken over
\(p^{25k_0}\) times the corresponding constants. Keep the third
point over \(p^{k_0}\) with exponents and depth \(-w\)
on its jump \(w\).

We check the strict inequalities in the
\hyperlink{proof.weighted-fit-cup-walk}{Frobenius passage
between cuts}. The negative and positive points at cut \(i\)
may be read with exponents
\[
 o_-(a,u)=\max\{\operatorname{std}_{l_i}(a,u),-u-A_a\},
 \qquad
 o_+(a,v')=-v'-A_a,
\]
where \(\operatorname{std}_l(a,u)=0\) for \(a<l\) and \(-u\)
for \(a\ge l\). Indeed the improved exponent once the depth
has been raised is \(-u+C_a\) if \(u<0\), and
\(\min(0,-u+C_a)\) if \(u\ge0\). On the negative-far side
any nonnegative jump is at most \(\eta_0\). When \(a<l_i\)
we have \(A_a=C_a=\eta_0\), so the improved exponent is at
least \(\operatorname{std}_{l_i}(a,u)=0\) in that case as
well. For \(a\ge l_i\) it is at least \(-u\). These facts
give both displayed exponents.

In the third point take \(w=-u-v'-1\), so its exponent is
\(u+v'+1\), an antitone function of \(w\). For each of
the two sides, add its depth to this third exponent and
the minimum of the exponent of the other side and that
side's standard exponent \(-\text{jump}\).
The result is at least \(1+C_a-A_a\ge1\).
Let the negative side lead the passage, since it retains
\(\operatorname{std}_{l_i}\). For its additional
inequality replace the other side's standard exponent
by its exponent transported from the preceding cut.
When \(a<l_i\), the transported exponent is zero and
\(o_+(a,v')=-v'-\eta_0\le0\). When \(a\ge l_i\), the
transported exponent is \(-v'-A_{a-l_i}\). Both required
lower bounds follow from
\[
 C_a\ge A_a,\qquad C_a\ge A_{a-l_i}\quad(a\ge l_i),
\]
the second since \(l_i\le2J+1\). In all comparisons the
third point is taken at the cut itself. The walk permits
either original side to lead. Graded symmetry in cohomology
and bilinearity also allow the scaling power to be placed
on the tested cocycle of \(L^\dagger\).

The cup is therefore killed by \(p^{k_\ell}\) for
\(k_\ell=10^4K_\ell(J+1)\), since the two point exponents
and the third one add to \(51k_0\le k_\ell\).
This applies to every pairing, and hence to the actual
cup-orthogonal quotient. The argument uses only \(p\ge2\),
including the ramified prime two. For
\(R=\operatorname{rank}_O L\) and
\(n_v=[E_v:\mathbf Q_\ell]\) the
\hyperlink{proof.weighted-fit-cup-length}{crystalline cocycle
rank} bounds the normalized length of the defect by
\[
 2n_vRk_\ell
   \le 4\cdot10^4 n_vR\,m\lambda_\ell/Q
   \le n_vR\,\mathrm{loss}_{\mathrm{crys}}(\ell)/t .
\]
Here \(J+1\le m\) and \(K_\ell\le2\lambda_\ell/Q\),
with \(Q=\kappa t\). This is the required allowance for a
Tate block and also fits the row allowance. The bound is
uniform in the modulus exponent \(M\), including the row
modulus \(p^{M_\ell}\) for
\(M_\ell=\lfloor r b_\ell(X)\rfloor+1\) and the scalar
modulus \(p^\nu\) in the first-order polynomial. The local
degree is the same as in the crystalline duality calculation.
\footnote{\weightedpin{grh-compact/src/New/Lift/SparseThickLetters.lean}
{175-L270}{Cut multiplier, points and sparse placement},
\weightedpin{grh-compact/src/New/Selmer/CrystallineBound.lean}
{39-L59}{Small-prime application}, and
\weightedpin{grh-compact/src/New/Selmer/OnePlaceDefect.lean}
{15-L23}{The tested row and Tate blocks}.}

\paragraph{The two sizes at large unramified primes.}
\teamanchor{proof.weighted-fit-large-two-sizes}%
Fix the actual unramified \(E_v/\mathbf Q_p\), where \(p=\ell>Q\)
and \(\lambda_\ell\ge H^3\). Write
\[
 b=B+\delta,\qquad g=G+\delta,\qquad W_1=W+4\eta,
 \qquad k=K_\ell,\qquad \beta_i\in\{g,b\}.
\]
Use the opposite windows at the letters of the chart constructed
above, with far width \(W\) and low half-width \(\eta\).
We use the point notation of
\hyperlink{proof.weighted-fit-cup-depth}{the depth calculation}.
In this paragraph \(E\) denotes the chart's polynomial \(z+p\),
and \(i_{-a}\) is the letter \(a\) Frobenius steps before \(i\).
An admissible pairing has a third side with jumps \(0\) and
\(w=-u-v-1\), where \(u,v\) are jumps of the two sides at the
same letter. All three sets of jumps are bounded above by
\(\beta_i+W_1\). Indeed when both jumps are far, their sum is in
\([-W,W]\), and when only one is far the low jump changes it
by at most \(\eta\).

Put
\[
 \begin{gathered}
 q=\lfloor b/p\rfloor,\qquad
 s_i=\begin{cases}\max(0,q-7k)&\beta_{\sigma i}=b,\\
                    0&\beta_{\sigma i}=g,\end{cases}\\
 s_* =q+12k,\qquad C=g+W_1,\qquad D=2(b+W_1).
 \end{gathered}
\]
Take \(i\) as a cut when \(s_i<\beta_i\).
The inequalities in the allowance calculation give
\[
 \begin{gathered}
 b<kp^2,\qquad 2g<kp,\qquad W_1\le k,\\
 C\le(p-1)k,\qquad
 \beta_i+W_1\le(p-3)s_*,\qquad D\le4kp^2.
 \end{gathered}
\]
For the penultimate inequality use \(b<pq+p\), \(q<kp\),
\(W_1\le k\), and \(p\ge7\) in
\((p-3)(q+12k)-pq=12k(p-3)-3q\).
We also have \(s_i+W_1\le pk\) and
\(\beta_i+s_i+W_1\le D\).
Every thick letter is a cut, and the predecessor of every thin
letter is a cut. Thus the distance \(l=l(i)\) to the nearest
strict predecessor cut is either \(1\) or \(2\), including on
Frobenius orbits of length one or two.

Let \(L_p=\lfloor\log_p D\rfloor\), \(\Delta=4k\) and
\(\kappa_*=120k+24L_p+1\).
At any cut on each of the three sides, \(p^{\kappa_*}E^{s_i}\)
kills the cokernel at level zero. If the successor is thick,
\(s_i+\kappa_*\ge s_*\) because \(s_*-s_i\le19k\). Apply
the unramified calculation of
\hyperlink{proof.weighted-fit-cup-lift-ideal}{the lifting ideals}.
If the successor is thin, \(s_i=0\) and the
\hyperlink{proof.weighted-fit-crystalline-thin-successors}{thin-successor calculation}
with one successor uses \(D\) and \(C\), and its bound is
\[
 6C/(p-1)+8D/p^2\le 38k<\kappa_* .
\]
Choose the third point at each cut over \(p^{\kappa_*}\cdot1\)
with orders \(z_i(0,w)=-w-s_i\) and \(z_i(a,w)=-w\) for
\(a\ge1\). These orders are antitone in \(w\).

\emph{Lifts for the two-letter gap.}
\teamanchor{proof.weighted-fit-crystalline-long-gap}%
We record the part of the quantitative period-frame argument used
here. A localized point of standard orders of level \(l>0\) has
orders \(0\) below \(l\) and \(-u\) on the jump \(u\) from level \(l\)
on. We use \(l=2\), \(p\ge7\), and \(D\le p^l\Delta\).
The quantitative frame at level \(l\) spans \(p^{4\Delta}\) times
these points by points of the same orders with depth at least
\(D\) at level zero. This is the quantitative clause of
\hyperlink{proof.weighted-fit-cup-radius}{the crystalline frame}.
Choose a finite family of such points whose span contains the
\(p^{4\Delta}\)-multiple of the standard module. For a standard
point \(y\), express \(p^{4\Delta}y\) in this family and expand its
scalar coefficients in \(E=z+p\) using \(\Lambda[[z]]=\Lambda[[E]]\),
where \(\Lambda=W(\mathbf F_{p^{f_v}})\). This gives
\[
 p^{4\Delta}y=\sum_{a=0}^{N}E^a\xi_a+E^{N+1}y_{N+1}
\]
with the remainder of standard orders. The coefficients \(\xi_a\)
are \(\Lambda\)-linear combinations of the same points, so they
have standard orders and depth at least \(D\) at level zero.
The chart group fixes \(\Lambda\).

If \(y\) has extra order \(d(u)\le D\) at level zero, then
\(\xi_N\) has order \(D\) there on \(d(u)>N\).
For the constant term reduce the expansion modulo \(E\)
on the indicated jump. In the truncated de Rham value of an
almost invariant coefficient, the positive-order part must vanish
modulo \(t^D\) by the fixed period calculation of
\hyperlink{proof.weighted-fit-cup-depth}{the improved-order rule}.
Thus an order of at least one improves to \(D\).
Subtract the coefficient
and divide by \(E\). This lowers \(d(u)\) by one at level zero
and leaves the constraints at every higher level intact. Repeat
through \(N\). The same argument applies to the projection of
the expansion into the trivial quotient.

Suppose now that \(E^{s_*}\) kills the cokernel at \(i\), that
\(p^{\kappa_*}\) kills it at \(i_{-l}\), and that all jumps on
the side being lifted are at most \(C\) between these two letters,
including \(i\) and excluding \(i_{-l}\). Start with a point over
\(c\) at \(i\) with orders \(-u-s_*\) at level zero and \(-u\)
above. Multiply it by
\[
 F(E)=E^{s_*}\prod_{a<l}\bigl((E-p)^{p^a}+p\bigr)^C .
\]
This is monic of degree \(N=s_*+C\sum_{a<l}p^a\). The new
orders are at least \(C-u\) below \(l\) and \(-u\) above, so
the point is standard of level \(l\). If \(N<D\), the
\(N\)-th coefficient of its expansion, after multiplication
by \(p^{4\Delta}\), has order \(D\) at level zero for
\(u<-N\). In the standard module of the trivial quotient, the
projection of each \(\xi_a\) is constant modulo \(E^D\) by its
depth at level zero and monic division. Compare the coefficient
of \(E^N\) with that of \(p^{4\Delta}F(E)c_i\). Since \(N<D\) and
\(F\) is monic, the difference between the projection of \(\xi_N\)
and \(p^{4\Delta}c_i\) is a multiple of \(E^D\) in that module.
Divide by \(E^D\) there and lift \(p^{\kappa_*}\) times the quotient
using the lifting ideal transported by \(\varphi^l\) from
\(i_{-l}\). Multiplying the lift by \(E^D\) gives order and depth
at least \(D\) at level zero and preserves the standard orders at
all other levels. Subtract it from \(p^{\kappa_*}\xi_N\).
The resulting point over
\(p^{\kappa_*+4\Delta}c_i\) has depth \(D\) at level zero,
\(0\) at \(0<a<l\), and \(-u\) above. Its orders are at
least \(D\) at level zero for \(u<-N\) and standard of level
\(l\) elsewhere. All terms are in the localized Kisin
module, so the removal of the error is an exact lift.%
\footnote{\weightedpin{grh-compact/src/New/Lift/LongGap.lean}
{415-L496}{Taylor coefficient and the exact correction},
\weightedpin{grh/src/New/Congruences/Lifting/LongGap.lean}
{118-L175}{The coefficient calculation}, and
\weightedpin{grh/src/New/Congruences/Lifting/LongGap.lean}
{365-L469}{The lifted coefficient}.}

\emph{Depth choices and the cup.}
For each cut \(i\), call the side with negative far sign at \(i\)
the leading side. Set \(A_a=\eta\) for \(a\le2\) and \(A_a=0\)
for \(a>2\). Use orders \(-u-A_a\) on the other side and
\(\max(\operatorname{std}_{l}(a,u),-u-A_a)\) on the leading
side, where \(\operatorname{std}_{l}\) is the standard order
of level \(l(i)\).

First consider the ordinary point on either side at \(i\).
The cut supplies an initial lift with orders and depth
\(-u-s_i\) at level zero and \(-u\) above. Raise its depth to
\[
 d_i(a,u)=-u+\eta\,\mathbf 1_{a\le5}+s_i\mathbf 1_{a=0}.
\]
At level zero the steps are within \(s_i+W_1\) of \(0\)
or the negative of the far center of that side. At levels
\(1,\ldots,5\) the radius is \(W_1\). There are no other
steps. The sum of \(v_p(n)\) on the nonzero integers of an
interval of radius \(pk\) in \([-D,D]\) is at most
\(2pk/(p-1)+L_p\), by counting the multiples of each \(p^h\).
For the two intervals at any level use \(8k+2L_p\) as a bound.
The extra cost \(3a\) per step at positive levels is at most
\(90k\) at each level, since there are at most \(4W_1+2\)
possible steps. Each of the six levels therefore costs at most
\(98k+2L_p\), and
\(\kappa_*+4\Delta+6(98k+2L_p)\le25\kappa_*\).
This bound pays for the ordinary point and for the additional
steps of the gap point below.

The improved orders of the ordinary point supply the prescribed
nonleading orders. At level zero, if the initial order is
nonpositive, \(s_i<\beta_i\) excludes a far negative jump, and
then \(-u-\eta\le0\). At higher levels the initial order is
\(-u\). If all letters from \(i_{-(l-1)}\) through \(i\) have
the same sign, the leading side has \(u\le\eta\) at all levels
below \(l\). Its improved orders then supply zero below \(l\)
and \(-u\) at levels \(a\ge l\).

The sole remaining case for the leading side has \(l=2\)
and opposite signs at \(i\) and \(i_{-1}\).
The noncut \(i_{-1}\) is thin and is followed by the thick
letter \(i\). It follows that \(s_{i_{-2}}=0\) and
\(g+7k\le q\). On the leading side the jumps at \(i\) are
at most \(\eta\) and those at \(i_{-1}\) are at most \(C\).
The \hyperlink{proof.weighted-fit-cup-lift-ideal}{unramified lifting-ideal
estimate} gives the exact multiplier \(E^{s_*}\) at \(i\), taking
exponents \(0,s_*\). At each letter \(j\) all its upper jumps are
at most \(\beta_j+W_1\le(p-3)s_*\).
The preceding construction applies with \(l=2\), since
\[
 s_*+C(1+p)\le(p+2)q-(6p-6)k<pq\le b<D .
\]
The strict inequality uses \(q<kp\) and \(p\ge7\).
Thus the far negative jumps at \(i\) are \(u<-N\).
Raise the depth to \(D\) at \(a=0\), \(3\eta\) at \(a=1\),
and \(-u+3\eta\,\mathbf1_{a\le5}\) for \(a\ge2\).
The new steps at level one are near zero, and those at
levels two through five are within \(W_1\) of the same
centers as for the ordinary points. The improved orders give
the required leading orders. At \(a=0\) the coefficient
argument gives the order on negative far jumps, and the order
is at least zero on low jumps. At \(a=1\) all leading
jumps are \(\ge-\eta\), so order zero suffices there.
From level two on the required order is \(-u\).
Multiply each point by the remaining power to put both sides
over \(p^{25\kappa_*}\) times their constants.

For the strict inequalities of
\hyperlink{proof.weighted-fit-cup-walk}{the Frobenius walk},
write \(w=-u-v-1\).
For a comparison at the same cut, the other order in the
minimum with \(-v\) is at least \(-v-A_a\).
The ordinary depth therefore gives a sum at least \(1\).
For the gap depth, at \(a=0\) the sum is at least
\(D+u+1-s_i-\eta>0\) by \(-u\le b+W\). At \(a=1\)
it is at least \(u+2\eta+1>0\) by \(u\ge-\eta\).
At higher levels it is again at least \(1\).
For the leading comparison with the preceding cut in the
ordinary case use the third point at \(i\). Below \(l\),
the moving side has \(v\ge-\eta\), which gives
\(\min(-v-A_a,0)=-v-A_a\). At \(a\ge l\) the two
moving orders are bounded below by \(-v-\eta\) if
\(a\le l+2\), and by \(-v\) otherwise. The extra \(\eta\)
in the ordinary depth covers these levels. In the gap case
use the third point from \(i_{-2}\), whose shift is zero.
At level zero its transported order is zero and the
moving orders have minimum at least \(-v-\eta\), as
\(v\ge-\eta\). The sum is at least \(D-v-\eta>0\).
At level one the moving jumps have \(v\le\eta\), so their
minimum with zero is at least \(-2\eta\), less in absolute
value than the depth \(3\eta\). From level two on the
transported third order is \(-w\), and the additional
\(3\eta\) in the gap depth covers both moving orders
until both become standard. This proves the strict
inequalities on every letter, for arbitrary choices of
the two crystalline cocycles and the compact pairing.

The walk now gives the cup bound with exponent \(51\kappa_*\).
Since \(W_1\le b\), \(L_p\le\lfloor\log_p(2b)\rfloor+1\).
With \(k\ge1\), the exponent is bounded by
\[
 51\kappa_*\le10^4 k\bigl(\lfloor\log_p(2b)\rfloor+1\bigr)
 \le 10^4K_\ell m.
\]
For the last inequality use \(2b\le4B\) and
\(\log_p(4B)\le\log_2(4B)<m\).
It acts on the integral dual row or Tate block, so every
reduction of this bound has the original precision of the
compact pairing. Here \(e_v=1\), so the conversion to normalized
length is at most \(2 f_v R\) times this exponent.
Since \(K_\ell\le2\lambda_\ell/(\kappa t)\), the normalized
length is at most \(4\cdot 10^4 f_vRm\lambda_\ell/(\kappa t)\).
This fits \(f_v R\,\mathrm{loss}_{\rm crys}(\ell)/x_L\)
for \(x_L=1\) on the row and \(x_L=t\) on a Tate block.
The case \(\lambda_\ell=H^3\) uses these same inequalities.%
\footnote{\weightedpin{grh-compact/src/New/Lift/SparseThickLetters.lean}
{199-L243}{Depth and exponent at cuts},
\weightedpin{grh-compact/src/New/Lift/SparseThickLetters.lean}
{273-L286}{The two sizes}, and
\weightedpin{grh/src/New/Congruences/Lifting/CutExponentBounds.lean}
{895-L1046}{Nearest cuts and their charges}.}

\paragraph{The loss from low jumps at large unramified primes.}
\teamanchor{proof.weighted-fit-large-few-low}%
Keep \(p=\ell>Q\), now with \(\lambda=\lambda_\ell<H^3\).
Let \(A=2(\lambda+1)\), and keep \(\eta,W,\beta_i\) from
the actual jump windows. The factorial gives \(p A>4B\)
and \(p>2B/H^3>9G\). Since \(A\le2H^3\), the scale
inequalities give
\[
 \beta_i+W+\eta\le(p-3)A,\qquad 6(A+\eta+W)<\beta_i .
\]
For example \((p-3)A>4B-6H^3\), and
\(\beta_i+W+\eta\le B+\delta+W+\eta<2B\).
No nonzero multiple of \(p\) is within \(6(A+\eta+W)\)
of \(0,\beta_i,-\beta_i\).
Near 0 or the thin center the absolute value is less than \(p\).
Near the thick center the absolute value belongs to
\((B-30H^3,B+30H^3]\), since
\(\delta<\eta\) and \(6(A+\eta+W)+\delta<30H^3\).
One multiple there would contribute at least \(H^3\) to \(\lambda\).

Write \(T\) for an arbitrary crystalline side in the original
class and \(S\) for the side with the actual tested dual lattice.
The far sign on \(S\) is \(\epsilon_i\). The third side of a
pairing \(P\) has jump \(w=0\) or \(w=-u-v-1\) for a jump
\(u\) of \(T\) and \(v\) of \(S\). These jumps lie within
\(2\eta+W+1\) of \(0,\beta_i,-\beta_i\), and are bounded
above by \(\beta_i+W+\eta\). The terms with both jumps far
and those with both jumps low give the neighborhood of zero.
With one jump far, the low jump and the translation by one
give the neighborhood of its signed far center. In particular
\(E^A\) kills the level-zero cokernels of all three sides by
\hyperlink{proof.weighted-fit-cup-lift-ideal}{the unramified lifting calculation}.

Write \(r(t)\) for the orders \(t-u\) at level zero and
\(-u\) at higher levels, with \(u\) the jump under test.
An initial point over any constant at a letter has orders and
depth \(r(-A)\). The levels above zero will be kept unless
explicitly changed. The depth can be raised with cost zero
to \(r(2A+\eta)\) on \(T\) or \(S\), and to \(r(A)\) on the
third side. Every step is at level zero, and is in the
three neighborhoods just considered. The unramified step
cost is \(v_p(n)\) at level zero with the cost at \(n=0\)
defined as zero. By the improved-order rule, a point of
depth \(r(2A+\eta)\) obtained in this way on either side has
level-zero order
\[
 o_0(u)=
 \begin{cases}
   2A+\eta-u & u<-A,\\
   \min(2A+\eta-u,0)& u\ge-A .
 \end{cases}
\]
The third side, when deepened to \(r(A)\), has antitone
order \(z_0(w)=A-w\) for \(w<-A\) and
\(z_0(w)=\min(A-w,0)\) otherwise.
We will use \(o_0(u)\ge-u-A\), with
\(o_0(u)\ge-u\) on the far jumps, and \(z_0(w)\ge-w-A\).

\emph{Two cases of zero defect.}
Suppose first that the allowed jumps on both sides miss
\([-A,0)\). Then \(o_0(u)\ge-u\) on both sides. All letters
are cuts separated by one step. At a cut the side with
negative far sign leads and has \(u\le\eta\), so its level-zero
orders are also nonnegative. Leave the third point at
orders \(r(-A)\). The same-cut depth inequality at level zero
has sum at least \(A+\eta+1\), using \(w=-u-v-1\).
For the comparison with the previous cut, use the third
point at the present cut. The moving side's level-zero order
has minimum at least \(\min(-v,0)\) with the shifted
order from the preceding cut. Since its present jumps satisfy
\(v\ge-\eta\), the sum is at least \(A+1\). At higher
levels all bounds are the standard ones and the sum is at least
\(1\). Thus the walk gives orthogonality at exponent zero.

Suppose instead \(S\) has all jumps in \([-\eta,0]\), as for
the tested Tate block. For \(S\), the unramified cokernel
bound with upper jump bound zero supplies a point with
orders and depth \(r(0)\) at every letter. Raise its depth to
\(r(2A+\eta)\) with cost zero. Transport the point from the
previous letter by Frobenius. This gives depth \(1\) at
level zero and \(2A+\eta-v\) at level one. Raise the level-zero
depth to \(2A+\eta-v\) by steps near zero, again with cost zero.
Use these points on \(S\), with orders \(0\) at level zero
and \(-v\) above, and with depth \(2A+\eta-v\) at levels zero
and one and \(-v\) above. Use the preceding points on
\(T\) and a third point of orders \(r(-A)\) at each cut.
Now \(S\) leads at every letter. For the same-cut comparison
at level zero the \(T\)-depth sum is at least
\(A+\eta+v+1\ge A+1\), and the \(S\)-depth sum is at
least \(\eta+1\) since the \(T\)-order is at least \(-u-A\).
Above level zero the sums are at least \(1\).

When comparing the \(T\) points across one step, use
the third point at the current letter if the \(T\)-far
sign is positive there. Then \(u\ge-\eta\). At level
zero the moving order is at least
\(\min(o_0(u),0)\ge-u-A-\eta\), so its sum with the
third order and \(S\)-depth is at least \(1\).
At level one the third order is standard and the moving
order is at least \(-u-A\).
If the \(T\)-far sign is negative, use the third point from
the preceding letter. At level zero this gives order zero
and \(u\le\eta\), so the sum is at least \(A-v>0\).
At level one the transported third order is \(-w-A\), and
the moving order is at least \(-u-A\), giving a sum
at least \(\eta+1\). Above level one the orders and depths
are standard. This proves orthogonality at exponent zero
on the one-sided dual block.%
\footnote{\weightedpin{grh-compact/src/New/Lift/LongGap.lean}
{314-L350}{Free level-zero steps and the two zero cases},
\weightedpin{grh/src/New/Congruences/Lifting/DeepPoints.lean}
{124-L210}{The gap below zero}, and
\weightedpin{grh/src/New/Congruences/Lifting/DeepPoints.lean}
{264-L368}{The one-sided walk}.}

\emph{The price of increasing a single order.}
\teamanchor{proof.weighted-fit-large-jump-colength}%
For a fixed letter \(i\), write \(m_i(u)\) for the jump
multiplicity on that letter. Use the multiplicities summed
over the coefficient embeddings on the letter, as in the
chart. We first record a relation that applies to a kernel
lattice and to an extension lattice.
Suppose \(c_a(u)<c_a(v)\) for every jump \(v<u\) at \(i_{-a}\).
Among \(f_v m_{i_{-a}}(u)+1\) localized points of orders \(c\)
there is a \(\mathbf Z_p\)-linear combination with one coefficient
one which splits as the sum of a point whose orders are raised
by one at level zero and a point whose order is raised by one
only at \((a,u)\).

To see this, use a basis of leading forms for jump \(u\)
modulo forms of smaller jump at \(i_{-a}\).
The number of forms is at most \(m_{i_{-a}}(u)\).
Evaluating \(\varphi^{-a}\) of the points and removing the
power \(t^{c_a(u)}\) gives a matrix over \(B_{\rm dR}^+\).
Upon reduction to \(C_p\), a nonzero minor of largest
absolute value among the full-rank minors expresses every
row in terms of at most \(m_{i_{-a}}(u)\) rows with
integral coefficients. Lift the coefficients to \(A_{\inf}\)
and apply \(\varphi^a\) so that the evaluation at level \(a\)
reads the chosen lifts. The residual point has gained the
single order at \((a,u)\). Forms of smaller jump have
already gained that order by the strict inequality for
the exponents. Apply the
\hyperlink{proof.weighted-fit-cup-retraction}{order-preserving coefficient retraction}
to make the coefficients lie in \(\Lambda[[z]]\).
It fixes the localized points and preserves the orders
of their residual points. Modulo \(E\), each coefficient
is in \(\Lambda\). The corresponding vectors of
remainders lie in at most \(m_{i_{-a}}(u)\) copies of
\(\Lambda\), each of \(\mathbf Z_p\)-rank \(f_v\).
Take a relation over \(\mathbf Q_p\) between these vectors
and divide by a coefficient of least valuation to get the
stated primitive relation. The remaining multiples of \(E\)
give the summand with all level-zero orders raised by one.%
\footnote{\weightedpin{grh-compact/src/New/Lift/LongGap.lean}
{358-L404}{One order and the multiplicity},
\weightedpin{grh/src/New/Congruences/Lifting/HodgeExpansion.lean}
{596-L768}{Leading forms and the primitive relation}.}

Let \(D_i(c)\) be the constants at \(i\) that admit
points of depth \(c\) and of orders at least \(r(0)\).
Assume \(c_b(v)\ge-v\) at every level \(b\) and jump \(v\).
Fix \((a,u)\) with \(c_a(u)<c_a(v)\) for every jump \(v<u\)
at \(i_{-a}\), and let \(c'\) be obtained by raising only
\(c_a(u)\) by one. If no multiple of \(p\) is strictly between
\(c_0(v)\) and \(c_0(v)+p^a\) for each jump \(v\) at \(i\),
then
\[
 \operatorname{length}_{\mathbf Z_p}
       \bigl(D_i(c)/D_i(c')\bigr)
 \le f_v m_{i_{-a}}(u)\bigl(v_p(c_a(u))+3a\bigr).
\]
Here the constants are read in the trivial quotient of
the side. To prove the inequality, let \(X_c\) be the
points defining \(D_i(c)\) and \(M_c\) the kernel points
of orders \(c\). The constant quotient is a quotient of
\(X_c/(X_{c'}+M_c)\). The single step operator gives
\(p^{v_p(c_a(u))+3a}\) times any element of this group equal to
zero. It raises the depth and differs from the indicated
scalar multiplication by a point of \(M_c\), which has no pole.

Take \(f_v m_{i_{-a}}(u)+1\) lifts of elements killed by
\(p\) in \(X_c/(X_{c'}+M_c)\). For each, write its
\(p\)-multiple as a sum of a point of depth \(c'\) and a
kernel point of orders \(c\). Apply the primitive relation
to the kernel points. For the resulting combination \(x\)
of the lifts, \(p x-w\) is of depth \(c'\), where
\(w\) is over zero and has orders \(c\) raised by one
at level zero. An operator \(J\) of cost zero raises
this level-zero depth from \(c_0(v)+1\) to \(c_0(v)+p^a\).
The steps avoid multiples of \(p\) by hypothesis.
The improvement rule gives order \(c_0(v)+p^a\) for
\(Jw\) at level zero. If \(c_0(v)+1>0\) use the positive
initial order, and otherwise the interval avoids zero
and \(c_0(v)+p^a\le0\). Thus \(Jw=E^{p^a}w'\)
with \(w'\in M_c\). On this unramified chart
\(E^{p^a}=E_a+p r\) for \(r\in\Lambda[[z]]\). Put
\(y=rw'-(Jx-x)\in M_c\). Then
\[
 p(x-y)=J(px-w)+E_a w'.
\]
The first summand keeps depth \(c'\), since a cost-zero
operator differs from the identity by a kernel point of
the input depth. The second has orders at least \(c'\),
so its Galois differences also have depth \(c'\).
No \(E_b\) divides \(p\). It follows that \(x-y\) has depth
\(c'\) and still has no pole. This is a primitive relation
on the \(p\)-torsion quotient. Its dimension is therefore
at most \(f_v m_{i_{-a}}(u)\). Filtering a module killed
by \(p^q\) by the kernels of successive powers gives length
at most \(q\) times that dimension. This proves the estimate
also when the individual point modules have infinite length.

When raising several orders at one level to an antitone
target, raise the smallest jump whose current order is still
below its target. Each lower jump is already at least its own
target, which is at least the target at the chosen jump. This
gives the strict inequality needed for each one-step estimate.
Orders at the other levels may be lowered after this
calculation without adding to the length.%
\footnote{\weightedpin{grh-compact/src/New/Lift/FewLowJumps.lean}
{32-L93}{The depth increment and division by \(p\)}, and
\weightedpin{grh/src/New/Congruences/Lifting/FewLowJumps.lean}
{107-L308}{Primitive relations on the quotient}.}

\emph{The colength of the useful constants.}
Apply the calculation to the universal side on
\(Z_f^1(S)\) that sends each parameter to its own cocycle.
Set
\[
 n_i=f_v\sum_{|v|\le\eta}m_i(v),\qquad
 r_i=f_v\sum_v m_i(v)
\]
for the kernel lattice \(S\).
All constants have points of depth \(r(2A+\eta)\) and
orders at least \(r(-\eta)\) by the initial calculation.
Indeed negative far jumps are below \(-A\).
Also \(p^A\) times every constant has a point of depth
\(r(2A+\eta)\) without a pole, by starting with the
cokernel killer \(p^A\) and then making the free steps
from \(r(0)\).

Remove the pole in stages \(r(t)\) to \(r(t+1)\) for
\(-\eta\le t<0\). The quotient of projected constants at
each stage is killed by \(p^A\).
Lift any \(f_v m_i(t)+1\) constants, in particular
representatives of elements of the \(p\)-torsion in the
stage quotient, to points with orders \(r(t)\) and
the assigned depth. The primitive relation for the extension
lattice gives a combination that gains the order at the jump
\(t\) at level zero. Its other summand
gains all level-zero orders, and hence also gains at \(t\).
On jumps \(v<t\) the original order \(t-v\) is already
positive, so the depth improves it to \(2A+\eta-v\).
For \(v>t\), the required \(t+1-v\le0\) follows by the
same improvement rule and the assigned depth.
Thus the combination has orders \(r(t+1)\) and retains
depth \(r(2A+\eta)\). The extension by the trivial
parameter has the same multiplicity at \(t<0\) as its
kernel. The resulting primitive relation on the projected
quotient bounds the dimension of its \(p\)-torsion by
\(f_v m_i(t)\). If \(t\) is absent there is no order to
raise and the same argument gives zero loss. The length of
this stage is at most \(A f_v m_i(t)\).
Consequently \(D_i(r(2A+\eta))\) costs at most \(A n_i\)
in colength.

If \(\epsilon_i=1\), retain this depth. If \(\epsilon_i=-1\),
take the least multiple \(N_i\) of \(p\) at least
\(\beta_i+W+2A\). First impose depth \(N_i\) at level
zero, using the single-depth estimate at \(a=0\). Lower any initially
stronger bound to \(N_i\). All these level-zero depths still exceed
or equal \(-v\) since \(N_i\ge\beta_i+W+2A\).
The hypothesis excluding multiples in an open interval
of length \(p^0=1\) is automatic. On a far negative jump
the steps \([2A+\eta-v,N_i)\) cost zero. A divisible
integer there would be at most \(N_i-p\), hence would
be nonzero and within \(6(A+\eta+W)\) of \(\beta_i\).
On a low jump the cost is at most \(v_p(N_i!)\le A+1\),
since \(N_i<\beta_i+W+2A+p\le(p-1)A+p\).
The cost is at most \((A+1)n_i\).

Next raise the depth at level one to
\[
 W+1-v+2A\mathbf{1}_{|v|\le\eta}
\]
and lower the level-zero depth to \(\beta_i+W+2A\).
The hypothesis on multiples for the single-depth estimate
at \(a=1\) holds because \((N_i,N_i+p)\) contains no
multiple of \(p\). The new depth at level one is
antitone on the jump set. A crossing from a far negative
jump to the low window changes \(-v\) by more than \(2A\).
All steps are within the prescribed three neighborhoods at
\(i_{-1}\), so each has cost \(3\), including the step at
zero if present. They cost at most
\(3(W+1)r_{i_{-1}}+6A n_{i_{-1}}\).
Thus in either sign the final constants with no pole
have colength at most
\[
 (2A+1)n_i+3(W+1)r_{i_{-1}}+6A n_{i_{-1}} .
\]
Intersect the inverse images in \(Z_f^1(S)\) of these
spaces of constants at all letters. Denote the result by
\(V\). The predecessor permutes the letters, so the length
of the quotient by \(V\) is at most
\[
 11\sum_i\bigl((A+1)n_i+(W+1)r_i\bigr).
\]
The construction of \(V\) uses only \(S\) and its jump
windows. In particular it precedes the choice of \(T\)
and the finite pairing.%
\footnote{\weightedpin{grh-compact/src/New/Lift/FewLowJumps.lean}
{101-L153}{The two depths and their length}, and
\weightedpin{grh/src/New/Congruences/Lifting/FewLowJumps.lean}
{485-L727}{Removing the pole and the two increases}.}

\emph{The cup bound on \(V\).}
Use all letters as cuts, and take the \(T\) points of
depth \(r(2A+\eta)\) and the third points with orders
\(z_0(w)\) at level zero and \(-w\) above.
For \(d\in V\), the chosen \(S\) points have the two
depths just computed. Their orders may be taken as \(-v\)
at all levels, except that when \(\epsilon_i=-1\) the
level-zero order can be taken as \(\max(0,-v)\) by the
improved-order rule. At the same cut the inequalities
at level zero for the \(T\) and \(S\) depths give
respectively lower bounds \(A+\eta+1\) and \(\eta+1\)
when the corresponding depth is \(r(2A+\eta)\).
When the \(S\) depth at level zero is \(\beta_i+W+2A\),
the same-cut lower bound there is
\(\beta_i+W+v+1\ge1\). At higher levels the assigned
depths are at least the standard orders, giving sums
at least \(1\).

When \(\epsilon_i=1\), the \(T\) side leads. Its
jumps at level zero are \(\le\eta\) and its orders are
nonnegative there. Use the third point at the current
cut. The moving order at level zero is at least
\(\min(-v,0)\) with \(v\ge-\eta\), giving a sum at least
\(A+1\). Above level zero the sums are at least \(1\).
When \(\epsilon_i=-1\), the \(S\) side leads at \(i\).
Use the third point transported from \(i_{-1}\).
At level zero its order is zero, the \(T\) move has order
at least \(\min(-u-A,0)\), and \(u\le\beta_i+W\).
The sum is at least \(A\).
At level one the three orders in the inequality are
\[
 z_0(-u-v-1),\qquad
 \min(-u,o_0(u)),\qquad
 W+1-v+2A\mathbf{1}_{|v|\le\eta}.
\]
If \(v\) is low, use the lower bounds \(-w-A\) and
\(-u-A\) for the first two terms. If \(v\) is far and
\(u\) low, \(|w|>A\) so \(z_0(w)=A-w\), on either branch
of the definition. The sum is again at least \(W+2\).
If both are far, \(|w|\le W+1\)
gives \(z_0(w)\ge-w-(W+1)\), and the moving order is at
least \(-u\). The resulting sum is at least \(1\).
Above level one the orders are standard with sum \(1\).
The Frobenius walk now gives cup products that are
coboundaries with exponent zero for every integral \(T\)
and every compact pairing on the indicated lattices.

The cup-orthogonal quotient on the actual dual row is
therefore a quotient of \(Z_f^1(S)/V\) as a
\(\mathbf Z_p\)-module. Each coefficient embedding
contributes its actual multiplicity to a single letter.
Thus \(\sum_i n_i\le f_v[O:\mathbf Z_p]\mu\) and
\(\sum_i r_i\le f_v[O:\mathbf Z_p]R\).
For an \(O\)-module of finite length its
\(\mathbf Z_p\)-length divided by \([O:\mathbf Z_p]\)
is its normalized \(O\)-length. Applying this to the
cup-orthogonal quotient itself gives the charge
\[
 11 f_v\bigl(\mu(A+1)+R(W+1)\bigr).
\]
The integral argument uses the given compact pairing
throughout, hence also its reductions at the finite
precision in the dual defect calculation.
For the row, if \(\lambda<H/5\) then
\(A< H/2\) for sufficiently large \(x\)\@. Hence both
allowed sets miss \([-A,0)\), so the defect is zero.
The Tate branch always has zero defect here.
For the remaining row case \(\lambda\ge H/5\),
\(A+1\le3\lambda\) and
\(W+1\le K_\ell\), using \(H/Q\gg W+1\).
These inequalities include the boundary \(\lambda=H/5\)
when it is integral. For a row, \(\mu\le Rm/\kappa\) and
\(K_\ell\le 2\lambda_\ell/(\kappa t)\).
Thus the charge is at most
\(11 f_v (3Rm\lambda_\ell/\kappa+
  2R\lambda_\ell/(\kappa t))\),
within \(f_vR\,\mathrm{loss}_{\rm crys}(\ell)\).%
\footnote{\weightedpin{grh-compact/src/New/Lift/FewLowJumps.lean}
{148-L163}{The walk for the submodule of constants},
\weightedpin{grh-compact/src/New/Lift/ExponentTheorems.lean}
{42-L99}{From letters to normalized length}, and
\weightedpin{grh-compact/src/New/Selmer/CrystallineBound.lean}
{60-L74}{The valuation split}.}

\paragraph{The tame-place defect from integral lifting.}
\teamanchor{proof.weighted-fit-tame-cup}%
Let \(\ell=g_i\), and fix a place \(v\mid\ell\). Then
\(E_v=\mathbf Q_\ell\) and the complex embedding and CM sign are
fixed at this place. Write \(L\) for the actual row block or scalar
Tate block, \(L^\dagger=L^\vee(1)\), and \(\mathcal Q\) for the
class of lattices with the allowed jumps and Weil inertia
polynomials of the catalogue. We use the
\hyperlink{proof.tame-type-lifting}{integral lifting argument at a
tame type prime} on a graph lattice constructed from an arbitrary
pairing in the definition of the cup-orthogonal submodule.

The characteristic-zero Weil inertia of every member of
\(\mathcal Q\) has exponent dividing
\[
 L_{\rm tame}=\operatorname{lcm}_{1\le a\le r}(\ell^a-1).
\]
This order is prime to \(\ell\), so the action of wild inertia is
trivial and the tame action is through the cyclic quotient of
order \(L_{\rm tame}\). Take an inertial element mapping to a
generator of this quotient. At this element the common monic
annihilator from the catalogue has coefficients in \(\mathbf Z_\ell\)
and degree at most \(N=(r+t)r\). Its roots among the
\(L_{\rm tame}\)-th roots of unity give one set of eigenvalues for
all the lattices in \(\mathcal Q\). They determine their tame
inertia characters at the generator. Adjoin the trivial character
to obtain a Frobenius-stable set of cardinality at most
\(N_*=N+1\). Each character has order dividing \(L_{\rm tame}\),
so its period divides \(\operatorname{lcm}(1,\ldots,r)\).
For two roots, their ratio has period dividing the least common
multiple of their periods. If they are in one orbit the orbit
length bounds the ratio period. If they are in distinct orbits,
the product of the orbit lengths bounds it and their sum is at
most \(N_*\). Thus \(q\le\lfloor N_*^2/4\rfloor\), using also
the trivial orbit in the single-orbit case.

The inverse set serves the Tate duals, since the Tate twist is
crystalline and hence has trivial Weil inertia. The common tame
crystalline extension in the earlier argument is obtained by
taking an unramified degree divisible by the character periods
and adjoining a root of the corresponding prime-to-\(\ell\)
Eisenstein equation. The unramified degree can be chosen prime
to \(\ell\), since all periods divide
\(\operatorname{lcm}(1,\ldots,r)\) and \(r<\ell\).

All members of \(\mathcal Q\) have jumps among the row jumps and
the scalar Tate jumps \(0,\ldots,t'-1\), \(t'\le t\). At the
fixed place there are at most \(r+t\) allowed integers.
The \hyperlink{proof.weighted-fit-crystalline-windows}{jump
comparison} puts the large jumps on a common side. Choose the
original orientation if they are positive and the Tate-dual
orientation if they are negative, using either orientation when
there are no large jumps. In the chosen orientation all jumps lie
in \([-\eta_*,H_*]\), where
\(\eta_*=\eta+1\) and \(H_*=\widetilde B+W+1\). Reflection
\(u\mapsto-u-1\) is included in those endpoints. Adjoining the
trivial jump still gives at most \(r+t+1\le N_*+1\) integers.
The prime threshold of the earlier lifting argument holds:
\(1\le\eta_*<H_*\) and
\[
   40(H_*N_*^2+N_*^3)<\ell .
\]
Indeed \(H_*=O(b)\), \(N_*=O(r^2)\) and
\(\ell\ge b^{10}\). We may therefore use the uniform exponent
\[
             k=7\bigl(N_*q+\eta_*(N_*+q)\bigr),
             \qquad q\le N_*^2.
\]
In particular \(k=O(r^{20})\).

For the application, let \(Y\in\mathcal Q\) and let
\(\psi:Y\otimes_O L^\dagger\to O(1)\) be the integral form of
a pairing at modulus \(\ell^M\). Its failure of equivariance
is divisible by \(\ell^M\). Thus reduction gives an equivariant map
\(f:Y\to L/\ell^M L\), using the perfect duality between \(L\)
and \(L^\dagger\). Consider the integral graph
\[
 \mathcal G_f
    =\{(l,y)\in L\oplus Y: l\bmod\ell^ML=f(y)\}.
\]
Projection to \(Y\) is an actual surjection with kernel
\(\ell^ML\), identified with \(L\) by multiplication by
\(\ell^M\). After inverting \(\ell\) the graph is the direct
sum of \(L\) and \(Y\). In particular its characteristic-zero
jumps and inertia types obey the same sets. Its Tate dual has
their reflected jumps and inverse types. Both integral
surjections used below thus have the hypotheses of the old
lifting argument in the stated orientation.

Let \(a\) and \(c\) be integral finite cocycles of \(Y\) and
\(L^\dagger\), respectively. Their reductions modulo every power
of \(\ell\) are continuous, so the cocycles themselves are continuous
in the integral coefficient topology and give classes in \(H_f^1\).
In the original orientation,
\(\ell^k[a]\) lifts to \(H_f^1(\mathbf Q_\ell,\mathcal G_f)\).
Projecting that lift to \(L\) gives an integral finite class
whose reduction in \(L/\ell^ML\) is \(\ell^k f_*[a]\).
Its integral cup with \([c]\) vanishes because the
Bloch--Kato finite conditions are orthogonal after inverting
\(\ell\), and the local invariant identifies
\(H^2(\mathbf Q_\ell,O(1))\) with \(O\), using finite free
scalar extension from \(\mathbf Z_\ell(1)\).
Reducing now proves vanishing for the pairing induced by
\(f\). By bilinearity of this reduced pairing we can place the
factor \(\ell^k\) on either cocycle.

In the dual orientation, dualize the graph exact sequence.
The map \(\mathcal G_f^\dagger\to L^\dagger\) sends a
functional \(d\) to \(l\mapsto d(\ell^M l,0)\). Lift
\(\ell^k[c]\) through this map to an integral finite class.
One may take a cocycle representative \(d\) whose image is
\(\ell^k c\), by correcting by a coboundary through the
surjection of the coefficient lattices. The map
\[
                 d\longmapsto\bigl(y\mapsto-d(0,\ell^M y)\bigr)
\]
is integral and equivariant into \(Y^\dagger\), because
\((0,\ell^M y)\in\mathcal G_f\). The resulting cocycle of
\(Y^\dagger\) is finite. For any lift \(\tilde f(y)\) of \(f(y)\)
the vector \((\tilde f(y),y)\) lies in the graph, and
\[
   d(\ell^M\tilde f(y),0)+d(0,\ell^M y)
          =\ell^M d(\tilde f(y),y).
\]
Consequently the resulting class reduces to the class
induced by \(\ell^k c\) and \(f\) on the dual of \(Y\).
Its integral cup with \([a]\) vanishes by the same finite
orthogonality. Reduction gives the required cup vanishing
also in this orientation.

Finally apply the integral \(\mathbf Z_\ell\)-linear
functional \(O\to\mathbf Z_\ell\) of the given pairing.
It preserves the vanishing modulo \(\ell^M\); the class
also vanishes after mapping the resulting Tate twist into
\(\overline{\mathbf Q}_\ell^\times\). Since the pairing,
its modulus and \(Y\) were arbitrary, \(\ell^k Z_f^1(L^\dagger)\)
lies in the cup-orthogonal submodule. If \(R=\operatorname{rank}_O L\),
the normalized defect is at most \(2Rk\) at this degree-one
place. The allowance pays this since
\[
 2k\ \le\ \frac{\mathrm{loss}_{\mathrm{crys}}(\ell)}{t},
 \qquad
 \mathrm{loss}_{\mathrm{crys}}(\ell)
       \ \ge\ 10^6\,\frac m\kappa\,b_1 ,
\]
using \(v_\ell(g_0g_1)=1\), \(2kt=O(r^{21})\) and
\(b_1\ge b=\exp(x^\varepsilon)\gg r^{22}\).
This gives the required tame bound both
for the row and for each scalar Tate block.
\footnote{\weightedpin{grh-compact/src/Lib/PadicHodge/Herr/GaloisLattices.lean}
{35-L75}{The actual Tate dual, pairing and cup condition},
\weightedpin{grh-compact/src/New/Lift/ExponentTheorems.lean}
{35-L40}{Common tame annihilators},
\weightedpin{grh-compact/src/Lib/PadicHodge/Limit/HodgeTateProjectors.lean}
{51-L65}{Continuity of the finite cocycle reductions}, and
\weightedpin{grh-compact/src/New/Selmer/OnePlaceDefect.lean}
{15-L45}{The same tame-place allowance and tested blocks}.}

\paragraph{The small unramified defect from signed cluster lifting.}
\teamanchor{proof.weighted-fit-small-unramified-cup}%
Let \(\ell\le Q=\kappa t\) and \(\ell\nmid N_1d\). Here
\(E_v/\mathbf Q_\ell\) is unramified, \(\ell\ge13\), and the
catalogue at the fixed place \(v\mid\ell\) is crystalline.
Write \(n_v=f_v=[E_v:\mathbf Q_\ell]\). The coefficient ring
contains the integer ring of \(E_v\). The component at \(q\)
has \(1\otimes\zeta_O=q\otimes1\), so its left-factor label
\(\tau_i:E_v\hookrightarrow O[1/\ell]\) satisfies
\(\tau_i(q)=\zeta_O\). For \(q=F(g)\), this identifies
\(\tau_i=g^{-1}|_{E_v}\), the local embedding in
\(\sigma_g=\iota_\ell^{-1}g^{-1}e\).
The \(\ell\)-power orbit of a starting letter runs through
all these labels. As \(q\) advances to \(q^\ell\), left
Frobenius \(\sigma_{E_v}\) gives
\(\tau_{i+1}\circ\sigma_{E_v}=\tau_i\).

Use the
\hyperlink{proof.crystalline-cluster-lifting}{signed cluster
lifting argument} with
\[
      j_*=J_\ell=\lfloor\log_\ell(H^3)\rfloor,\qquad
      b_*=\widetilde B,\quad w_*=\widetilde G,\quad
      S_*=W,\quad \eta_*=\eta .
\]
The chosen centers \(\widetilde B=b_1+\delta\) and
\(\widetilde G=b_2+\delta\) are integral.
Each coefficient embedding has large jumps in one allowed
signed interval at its own center. At the thick labels this
center is \(b_*\), and elsewhere it is \(w_*<b_*\). Put
\(g_*=w_*+S_*+\eta_*\). The scale inequalities give
\(g_*=O(G)\) and \(b_*\ge C(j_*+1)(g_*+1)\) for any fixed
constant \(C\) from the cluster lifting argument and all
sufficiently large \(x\).

We have \(j_*\ge1\), since \(\ell\le Q\le H^3\), and
\(j_*<n_v\). For the second assertion, all \(N_1\)-th roots
are prime to \(\ell\) here and lie in \(E_v\), so
\(\ell^{n_v}\equiv1\bmod2^{s+1}\). The valuation calculation
in the \hyperlink{proof.weighted-fit-crystalline-spacing}{spacing
argument} gives \(2^s\le n_v(\ell+1)\).
For large \(x\) this implies \(n_v>12s\ge j_*\), since
\(2^s\ge x^{11}\) and \(Q=O(x^{10})\).
If a thick label has a thick predecessor at distance
\(a\in[1,j_*]\), the spacing relation excludes a second thick
predecessor at a smaller distance. It also excludes every
thick label in the \(j_*\) positions immediately before that
predecessor. These are exactly the cyclic predecessor conditions
of the earlier argument. Either sign and either order of two
different signs in a cluster are allowed by its proof.

Apply the lifting to the coefficient projection
\(\mathcal G_f\twoheadrightarrow Y\) of the
\hyperlink{proof.weighted-fit-tame-cup}{integral graph construction}.
Both lattices are crystalline, because \(\mathcal G_f[1/\ell]\)
is the direct sum of \(L[1/\ell]\) and \(Y[1/\ell]\).
Its allowable windows are the same windows for every member
of the catalogue, including when \(L\) is one of the Tate
blocks. The theorem uses the decreasing-filtration jump
convention in which the cyclotomic twist has jump \(-1\).
It can also be applied to the dual graph projection. For
that orientation take \(b_*^\dagger=\widetilde B-1\),
\(w_*^\dagger=\widetilde G+1\), \(S_*^\dagger=W\),
and \(\eta_*^\dagger=\eta+1\). The actual far magnitudes before
reflection lie in \([\beta+t,\beta+3t]\). After \(u\mapsto-u-1\)
they lie in \([\beta+t-1,\beta+3t+1]\), which is contained in
both \([\beta-1,\beta-1+W]\) and \([\beta+1,\beta+1+W]\)
because \(t\ge2\) and \(W=4t\).
These choices preserve the thick set, the predecessor conditions,
and the separation for sufficiently large \(x\).

The earlier proof gives an annihilating exponent for the
cokernel on integral \(H_f^1\), uniformly in the ranks of the
two lattices. Its bound is
\[
 k_{\rm cl}\ll
   \frac{H_{\rm loc}}{(\ell-1)^{j_*+1}}
       +(g_*+1)(j_*+1)+\log_\ell\max(1,H_{\rm loc}),
\]
where \(H_{\rm loc}=b_*+\eta_*+S_*\) if a thick label occurs,
and \(H_{\rm loc}=g_*\) otherwise. The same formula holds
with the parameters for the dual orientation. It is the
annihilator in the proof that is used here, before conversion
to a length bound.
Since \(\ell-1\ge\sqrt{\ell}\) and
\(\ell^{j_*+1}>H^3\), the denominator in the first term
exceeds \(H^{3/2}\). Thus the first term is \(o(B/Q^2)\)
uniformly, since \(H_{\rm loc}=O(B)\) and \(H^{3/2}\gg Q^2\).
The second term is \(O(G\log x)=o(B/Q^2)\) by the thin-to-thick
scale ratio. The same is true of the logarithmic term,
which is \(O(\log b)\). Hence for all sufficiently large
\(x\), the old exponent \(k_{\rm cl}\) is at most
\(K_\ell=\lfloor2\lambda_\ell/Q\rfloor\), using
\(K_\ell\ge B/Q^2\).

Use the integral graph argument over \(E_v\).
The local invariant identifies \(H^2(E_v,O(1))\) with \(O\),
so rational Bloch--Kato orthogonality again gives integral
orthogonality before reduction. The graph therefore gives cup
annihilation by \(\ell^{K_\ell}\) for every pairing in the one-place
defect, including its arbitrary integral coefficient
functional. In particular the rank of \(Y\) plays no role
in the final length. The normalized defect for the tested
block of rank \(R\) is at most \(2n_v R K_\ell\).
This is at most \(n_vR\,\mathrm{loss}_{\mathrm{crys}}(\ell)/t\),
since \(2K_\ell\le4\lambda_\ell/(\kappa t)\).
This proves the required crystalline bound at the small
unramified places on the same row and Tate lattices.
\footnote{\weightedpin{grh-compact/src/New/Selmer/OnePlaceDefect.lean}
{15-L23}{The tested defect}, and
\weightedpin{grh-compact/src/New/Selmer/CrystallineRank.lean}
{43-L56}{Length from a uniform annihilator}.}

\teamanchor{proof.weighted-torsion-specialization}%
\paragraph{Specialization at the first \(m\) powers.}
The old ray construction chooses a common character through its
coordinates on the auxiliary prime pairs. Its primitive-power
estimate applies when the exponent is prime to \(\ell\).
Here the order of the ray coordinates is chosen before the primes,
and the fitted exponents include \(j=\ell\) if \(\ell\le m\).
The following valuation argument handles those exponents as well.
\footnote{\hyperlink{proof.twist-ray-quotient}{Earlier ray-coordinate construction},
\hyperlink{proof.common-twist-specialization}{earlier primitive-power estimate},
and \weightedpin
{grh-compact/src/New/Selmer/RootsOfUnityPoints.lean}{15-L105}{Uniform degree-free specialization}.}

Normalize \(v_\ell(\ell)=1\). For a polynomial \(F\) over the
valuation ring in \(\overline{\mathbf Q}_\ell\) let
\(\|F\|_s=\max_i |[T^i]F|_\ell s^i\). This Gauss norm is
multiplicative. For instance, in a product the coefficient at the
sum of the lowest maximizing indices has a unique contribution
of maximal weighted norm. Fix \(0<s<1\) with
\(|\ell|_\ell<s^{\ell-1}\) and finitely many positive targets \(y_i\).
For all sufficiently large \(N\), if \(\|F_i\|_s\ge y_i\) for every
\(i\), there is one \(\zeta\in\mu_{\ell^N}\) with
\(|F_i(\zeta-1)|_\ell\ge s y_i\) for every \(i\).

Put \(v=v_\ell\), \(\Delta_s=-\log_\ell s>0\), and
\(Q=\prod_i F_i\). For each \(n\ge1\), let \(X_n\) consist of
\(\zeta-1\) for the primitive \(\ell^n\)-th roots \(\zeta\). Its
\(d_n=\varphi(\ell^n)\) distinct elements have valuation \(t_n=1/d_n\).
Consider \(n\ge n_s\), where \(n_s\) is the least \(n\) with
\(t_n\le\Delta_s\). Let \(k_n\) count the roots of \(Q\) of valuation
\(t_n\), with multiplicity. For a polynomial \(F\) with leading
coefficient \(c_F\), multiplicativity and linear factorization give
\[
  -\log_\ell\|F\|_{\ell^{-t}}
      =v(c_F)+\sum_{F(a)=0}\min(t,v(a)).
\]
Each root contributes a nonnegative amount to the difference between
this expression at \(t=\Delta_s\) and at \(t=0\). A root on the
\(n\)-th circle contributes \(t_n\). Since all \(F_i\) are integral,
\[
 \sum_{n\ge n_s} k_n t_n
  \le -\log_\ell\|Q\|_s+\log_\ell\|Q\|_1
  \le -\log_\ell\prod_i y_i .
\]

Suppose every \(x\in X_n\) fails a test. At such an \(x\), a failed
factor has \(|F_i(x)|_\ell<s y_i\le s\|F_i\|_{\ell^{-t_n}}\),
and every other factor is bounded by its Gauss norm at this radius.
Thus \(v(Q(x))+\log_\ell\|Q\|_{\ell^{-t_n}}>\Delta_s\).
This excess equals \(\sum_{v(a)=t_n}(v(x-a)-t_n)\), by
ultrametricity for the roots on the other circles. Truncating each
nonnegative summand at \(\Delta_s\) leaves a sum at least
\(\Delta_s\), including when \(Q(x)=0\).
For a root \(a\) on this circle choose \(x_a\in X_n\) maximizing
\(v(x_a-a)\). At the other points \(v(x-a)\le v(x-x_a)\).
The differences \(x-x_a\) are unit multiples of distinct factors
\(1-\xi\), \(\xi\in\mu_{\ell^n}\setminus\{1\}\). The product of all
these factors is \(\ell^n\), so their valuations are nonnegative
and sum to \(n\). The total truncated contribution of \(a\) on
\(X_n\) is therefore at most \(\Delta_s+n\). Summing in the two
orders gives \(k_n/d_n\ge\Delta_s/(n+\Delta_s)\).
If every point through level \(N\) failed, the sum of these lower
bounds would grow without bound with \(N\), contrary to the fixed
bound \(-\log_\ell\prod_i y_i\). This gives a threshold independent
of the degrees and coefficients of the polynomials.

For several variables fix a coordinate bound \(q\), an integer
\(a_0>0\), a tolerance \(\epsilon_0>0\), and bounds \(B_i\ge0\).
For each \(i\) and \(1\le j\le m\), suppose an integral polynomial
\(F_{i,j}\) has a point \(u_{i,j}\) satisfying
\[
 v_\ell(F_{i,j}(u_{i,j}))\le B_i,\qquad
 |(u_{i,j})_v-1|_\ell^{a_0}\le|\ell|_\ell
 \quad\text{at each coordinate }v.
\]
Choose \(s\) close enough to one that
\[
    |\ell|_\ell<\min(s^{\ell-1},s^{ma_0}),\qquad
    s^q\ge|\ell|_\ell^{\epsilon_0}.
\]
In the next unchosen coordinate replace the variable by
\((1+T)^j\). Some \(j\)-th root \(\alpha\) of \((u_{i,j})_v\) has
\(|\alpha-1|_\ell\le s\). Otherwise the product of all
\(|\alpha-1|_\ell\) would exceed \(s^j\), contrary to the assumed
bound on \(|(u_{i,j})_v-1|_\ell\). Evaluating at \(T=\alpha-1\)
shows that the one-variable Gauss norm is at least the lower bound
already attained with the unchosen coordinate \(u_{i,j}\).
The preceding one-variable argument now replaces that coordinate
by \(\zeta_v^j\) for one \(\zeta_v\in\mu_{\ell^N}\), losing at most
a factor \(s\) for every \(i,j\). Repeat for all coordinates.
The finite list of targets for each coordinate is the original
list scaled by the accumulated powers of \(s\), so one threshold
for \(N\) suffices. This gives
\[
  v_\ell(F_{i,j}((\zeta_v^j)_v))\le B_i+\epsilon_0
                 \qquad\text{for all } i,j.
\]
Monic integral auxiliary test polynomials \(D_b\) can be included
by applying the one-variable argument also to \(\prod_b D_b(1+T)\).
Put \(L_*\ge\sum_b\deg D_b\). The product has Gauss norm at least \(s^{L_*}\).
Taking \(s^{L_*+1}\ge|\ell|_\ell^{\epsilon_0}\) forces
\(v_\ell(D_b(\zeta_v))\le\epsilon_0\) for every \(b,v\), since
each \(D_b(\zeta_v)\) is integral. With \(D_b(X)=X^b-1\) for
\(1\le b\le m\) we can take \(L_*=m(m+1)/2\). Each required power
then has nontrivial inertia on each chosen pair.

\teamanchor{proof.weighted-fit-line-points}%
\paragraph{Finite polynomial relations at a fixed coefficient depth.}

Write \(p=\ell\), let \(O\) be a finite coefficient DVR with \(v(p)=1\), and let \(K_0=M_\ell\)
be the positive integer defining the modulus \(p^{K_0}\). Fix an auxiliary pencil with
coordinates \(h_P:G_E\to\mathbf Z/p^N\) and a finite allowed set of ramification containing
the original places and the auxiliary pairs. For a row \(\rho\colon
G_E\to\operatorname{GL}_a(O)\), define the matrix function

\[
 B_{\mathbf Y}(g)=\rho(g)\prod_P Y_P^{[h_P(g)]},
\]

where the brackets choose the representative in \([0,p^N)\). Impose the crossed-homomorphism
equations with this matrix function, the original unramified and integral linear relations on
the cocycle, and the fixed normalization at the distinguished inertia element. We may also
impose a specified subset of the linear ties at the auxiliary pairs. For a finite extension
DVR \(O'\) and a point \(\mathbf y\in O'^J\) satisfying \(y_P^{p^N}=1\bmod p^{K_0}\) for all
\(P\), these equations define the intended module of continuous cocycles over
\(A'=O'/p^{K_0}O'\). The defining equations have polynomial coefficients even at other
\(\mathbf y\). Such other points need not define a Galois representation.

One finite set of evaluation coordinates suffices for all points satisfying those congruences.
Indeed, first factor out the closed normal subgroup of \(G_E\) generated by the inertia groups
outside the allowed finite set. The action and every unramified crossed homomorphism factor
through the resulting group \(\Pi\). Choose an open normal subgroup \(U\) of \(\Pi\) in the
kernels of all \(h_P\) and of the row representation modulo \(p^{K_0}\). Over \(U\) a cocycle
restricts to a homomorphism into a finite abelian \(p\)-group killed by \(p^{K_0}\). The
largest abelian quotient of \(U\) killed by \(p^{K_0}\) is finite. In fact \(U\) comes from a
finite extension of \(E\), with ramification still confined to a finite set, and its abelian
\(p\)-quotient of bounded exponent is finite by the ray class description. The finite class
group and the local unit groups at the allowed finite places give the bound. At residue
characteristic \(p\) the deep units are finitely generated over \(\mathbf Z_p\), and at the
other allowed places the \(p\)-part of the unit quotient is finite. The kernel of this abelian
quotient is characteristic in \(U\) and normal in \(\Pi\). All the cocycles under
consideration vanish on the kernel and therefore are determined by their values on one fixed
finite set of coset representatives. The original relations are groups of linear equations
with coefficients in \(O\) evaluated on these same coordinates. The crossed-homomorphism
equations are polynomial linear forms in the coordinates by the displayed formula for
\(B_{\mathbf Y}\).

Include the relation rows \(p^{K_0} e_i\) as well. If \(d\) is the number of coordinates, the
\(O\)-module generated by the specialized rows at a point \(\mathbf y_0\in O^J\) has full rank
\(d\). Choose its basis and express each member as a finite \(O\)-linear combination of
specialized relation rows. Use the same combinations of their polynomial rows to form a square
matrix and let \(F(\mathbf Y)\) be its determinant. The matrix at \(\mathbf y_0\) has Smith
exponents at most \(v_O(p^{K_0})\), since the row module contains \(p^{K_0}O^d\). Hence
\(v(F(\mathbf y_0))\) is the normalized length of the intended coefficient module at that
point. At every other admissible point every cocycle lies in the kernel of the specialized
square matrix modulo \(p^{K_0}\). Smith form there bounds its normalized length by
\(v(F(\mathbf y))\), with the infinite bound allowed when \(F(\mathbf y)=0\). If the center
lies in \(O'\), first form the same polynomial relations over \(O'\) and choose the Smith rows
there. All inequalities continue to hold over larger DVRs. This uses \(v(p)=1\) in every
extension, and length is divided by the ramification index in whichever coefficient DVR is
used.

\paragraph{Cost of removing ties on the helper line.}

Let \(P_0\) be the helper and \(J\) the full auxiliary set. The anticyclotomic ray coordinates
give, at \(P\ne P_0\) and at the inverse-conjugate reading, character value \(y_P\) on tame
inertia and \(y_{P_0}\) on the selected Frobenius. Write \(t_P\) for the slope of an
individual row there. A nonzero \(t_P\) implies that \(\rho(\phi_P)\) and
\(\rho((c\phi_Pc)^{-1})\) are the identity modulo \(p^{K_0}\). For a root \(u\) of order
dividing \(p^k\) set

\[
 y_{P_0}=u,\qquad y_P=1+t_P(u-1)\quad(P\ne P_0).
\]

If \(N\ge k+2K_0\), the reductions of these entries define a character modulo \(p^{K_0}\). To
see this, raise the formula for \(y_P\) to \(p^k\) modulo \(p\), obtaining \(1\), then use the
binomial theorem successively modulo higher powers of \(p\). The order of its reduction
divides \(p^{k+K_0}\), which suffices for the ray order. In addition

\[
 |y_P-1|_p^{p^k}\le |p|_p.
\]

For \(u=1\) this is immediate. Otherwise its order is \(p^b\) with \(b\le k\), and
\(v(u-1)=1/\varphi(p^b)\).

For a cocycle \(z\) at a new pair, use its affine local representation \(g\mapsto (B_{\mathbf
y}(g),z(g))\) over \(O'/p^{K_0}O'\). Wild inertia and the prime-to-\(p\) part of tame inertia
have trivial image. Their linear action is trivial, and their translation action would be a
homomorphism into a \(p\)-group. The linear action of the tame pro-\(p\) part is scalar of
exponent dividing \(p^{k+K_0}\). Its affine exponent then divides \(p^{k+2K_0}\). Since
\(q_P\equiv1\bmod p^N\), the tame Frobenius relation makes the images of \(\delta_P\) and
\(\phi_P\) commute. The inverse-conjugate readings have the same property. At either reading,
set \(B=\rho(\phi_P)\) with the matching choice of Frobenius. For \(P\ne P_0\) the affine
equality is

\[
 (uB-I)z(\delta_P)=(y_P-1)z(\phi_P).
\]

It follows that \((uB-I)(z(\delta_P)-t_P z(\phi_P))=0\), since \(t_P(I-B)=0\) in the coefficient
module. The quotient obtained by dropping all these ties injects via the displayed
discrepancies into the indicated kernels, one for each of the two readings of each pair. Smith
form bounds the increase in normalized length by the sum of \(v\det(uB-I)\). Here \(B\) in the
determinant can be the fixed integral characteristic-zero Frobenius matrix. At \(u=1\), the
helper tie is \(z(\delta_{P_0})=0\). The same argument uses the kernel of \(B-I\) for the
unrestricted Frobenius matrix \(B\). This bounds the increase when only that tie is dropped at
\(u=1\).

\paragraph{Initial points and fixed costs.}

Fix all row modules, signed families and the bound \(J_0\) on the number of pencil coordinates
before choosing a ray order. The bound \(J_0\) comes from the integral length descent in the
tied-pencil induction. Work with two separate choices of the finite-relation polynomials. For
each row first choose \(G_i\) for the relations with the nonhelper ties, taking the sharp
center to be the identity tuple. The fully tied untwisted module has normalized length at most
its balanced bound \(\beta_i^{\rm bal}\). Allowing ramification at the helper without its tie
costs at most the sum of the two helper determinant valuations. The approximating inertia
element gives the upper bound \(d_i c_\ell\) for each determinant in rank \(d_i\) by the two
estimates in \eqref{eq:weighted-fit-helper-determinants}. The helper's Frobenius and
inverse-conjugate Frobenius can be chosen in the same determinant-valuation neighborhoods.
Thus

\[
 v(G_i(\mathbf 1))\le \beta_i^{\rm bal}+2d_i c_\ell\le 10 f d_i c_\ell.
\]

This calculation also requires the original untwisted cocycle relations and the normalization
at the same distinguished inertia element used in the tied induction. Test \(G_i\) on the
row's helper line, with the slopes of that row, using \(U-1\) as the polynomial variable. Its
value at \(U=1\) has the bound just obtained.

For a signed family, at the modulus consisting only of the helper set its single coordinate to
\(U\). Its value at \(U=1\) is its original nonzero untwisted value times the omitted Euler
factors at the helper. The helper can satisfy the finite approximation for all the signed
families at once. At the prescribed inertia element each Euler parameter has a nontrivial root
of unity with order dividing \(q_w-1\); the Tate factor there is one. The valuation of its
difference from \(1\) is at most \(v(q_w-1)\), and continuity permits the same bound at the
helper. There are at most four factors in the displayed norm identities. In particular the
weaker allowance \(8c_\ell\) suffices for the total helper cost. At the extra pairs the signed
family uses the helper line with \(y_P=1\) for \(P\ne P_0\). Adding their modulus when the
character on it is unramified multiplies the value by the Euler factors in the signed
interpolation calculation. At these pairs the avatar takes value \(U\) on the prescribed
Frobenius. With the sign \(s_0\) of the signed test and the choice
\(\epsilon\in\{1,-1\}\) of local reading, each factor is
\(1-aU^{-s_0\epsilon}\) for an integral unit \(a\). For unit \(U\) this
is either \(U^{-1}(U-a)\) or \(-a(U-a^{-1})\).

For the row polynomials the extra factors to control are \(\det(U\rho(\phi_P)-I)\) at both
readings of each nonhelper pair. Each is a unit multiple of a monic polynomial since
\(\rho(\phi_P)\) is integral and invertible. The signed factors are the preceding linear
factors. If there are \(a\) signed families, the total degree of the guard factors is at most
\(D_*=2J_0\sum_i d_i+4J_0a\). The one-variable torsion lemma proved above applies to each of
the row line polynomials and the signed helper polynomials with the target given by its
displayed valuation bound at \(U=1\). The Gauss norm on the disc in \(U-1\) is at least the
absolute value at its center. Include also the product of the monic factors after translating
\(U=1+T\). Its Gauss norm at radius \(s\) is at least \(s^{D_*}\). Take \(s\) sufficiently
close to \(1\) that \((D_*+2)(-\log_p s)<c_\ell/2\) and
\(|p|_p<s^{p-1}\). The common lemma then chooses one
\(u\in\mu_{p^k}\). The sum of the valuations of the monic factors is at most \((D_*+1)(-\log_p
s)\), and the additional loss for each line value is at most \(-\log_p s\). Its threshold
\(k\) depends only on the fixed lower targets and the degree bound for the guards. Thus \(k\)
can be fixed before \(N\) and before the primes. The row line at this \(u\) gives the untied
module of normalized length at most \(10 f d_i c_\ell+c_\ell/2\) by the removal-of-ties
calculation. The signed line at this \(u\) gives the full-modulus value with valuation at most
the untwisted valuation plus \(8c_\ell+c_\ell/2\) by the imprimitive norm identities. These
are separate line points for the separate inputs to the common specialization, even though
their helper coordinate \(u\) is common.

Now construct the second row polynomials \(F_i\) for the untied relations, sharp at the
specific points just obtained over the valuation ring of a finite coefficient extension
containing \(O\), \(u\) and the signed coefficients. Use the full-modulus signed polynomials.
The inequalities \(|y_P-1|_p^{p^k}\le |p|_p\) hold for every center. Take \(N\ge k+2K_0\) also
beyond the degree-free threshold in the multivariate lemma above, with \(a_0=p^k\), the
preceding valuation bounds as targets, \(q=J_0\), and losses at most \(c_\ell/2\) in that
lemma. This \(N\) can be chosen before building the pencil, because the targets and the bound
on the number of coordinates are then already fixed. The second lemma gives one torsion tuple
whose first \(m\) powers have every required bound simultaneously. Taking its monic factors
\(X^b-1\), \(1\le b\le m\), with loss per coordinate less than \(c_\ell/J_0\) gives exact
support at every power, and bounds the sum of the inertia-generator valuations at each power
by \(c_\ell\). The ray construction identifies the tuple with one inverse avatar, trivial at
the prescribed local groups above \(\ell\). Its powers give one character and the required
common conductor on both coefficient readings.

The total row bound is \((10fd_i+1)c_\ell\). For each signed family the total increase in
valuation is at most \(9c_\ell\). The conductor norm is a \(p\)-unit. By the exact-conductor
value relation this gives
\(\log|j_\iota(\rho_{p,s}(\theta^j)^{-1})|_\ell\le9c_\ell\log\ell=v_\ell\) for each critical
pair and sign, including the square of the ratio at rank two.
\footnote{\weightedpin
{grh-compact/src/New/Selmer/LinePointPowers.lean}{58-L95}{Finite Selmer polynomial bounds},
\weightedpin
{grh-compact/src/New/Selmer/TwistTheorem.lean}{27-L91}{Pencils with boundedly many primes}, and
\weightedpin
{grh-compact/src/New/Selmer/TwistTheorem.lean}{115-L195}{Line points and common specialization}.}

\teamanchor{proof.weighted-fitting-budget}%
\paragraph{The sum over coefficient primes.}
Put \(h_2=r^8\), take \(b_1\) to be the datum's integral rounding of
\(a/x\), and set \(g=g_0g_1\) and
\[
  A=(2b_1)!\left(\prod_{\max(0,b_1-30h_2^3)<u\le b_1+30h_2^3}
                       u\right)^{h_2^3}.
\]

\paragraph{A common integer for the Tate-twisted cocycles.}
\teamanchor{proof.weighted-fit-odd-class-budget}%
Put \(N_{\rm cyc}=dN_1\), \(F=E^+\), and \(n_E=[E:\mathbb Q]\).
The weighted field \(E=K(\mu_{N_1})\) lies in
\(\mathbb Q(\mu_{N_{\rm cyc}})\). Since \(12\mid N_1\), we have
\(E=F(i)=F(\zeta_3)\). The discriminant of \(E/F\) divides
both \((4)\) and \((3)\) by the orders generated by \(i\)
and by \(\zeta_3\). Hence \(E/F\) is unramified at every
finite place. The restriction and corestriction maps between its
rings of \(\ell\)-integers satisfy
\(\operatorname{cor}\operatorname{res}=2\) and
\(\operatorname{res}\operatorname{cor}=1+c\), where \(c\)
is complex conjugation on \(E\).

The cocycles to be bounded have coefficients in
\((O/\ell^\nu O)(-e)\), with \(0\le e<t\). They are unramified
outside \(\ell\). At \(e=0\) they kill inertia above
\(\ell\), so their length for \(O=\mathbb Z_\ell\) is at most
\(v_\ell(h_E)\) by ordinary class field theory. At \(e>0\)
they vanish on the kernels of the cyclotomic character in
the decomposition groups above \(\ell\), and they obey
\(z(cgc^{-1})=-z(g)\).\footnote{\weightedpin
{grh-compact/src/Lib/ClassField/Places/SelmerDefs.lean}{89-L97}{The local and oddness conditions on the finite cocycles}.}

Take \(e>0\) and \(O=\mathbb Z_\ell\). Embed the coefficients
into \(B=(\mathbb Q_\ell/\mathbb Z_\ell)(-e)\).
At \(v\mid\ell\) the local cocycle factors through the
cyclotomic image of the decomposition group, which is open
because \(E_v/\mathbb Q_\ell\) is finite. At odd \(\ell\)
this image has a procyclic subgroup inside \(1+\ell\mathbb Z_\ell\)
of index prime to \(\ell\). At \(2\) it is itself procyclic
because \(i\in E\). In either case a topological generator
of the procyclic subgroup acts on \(B\) with its action
minus one a nonzero scalar. That scalar acts surjectively on
the divisible module \(B\), so \(H^1\) of the subgroup
vanishes. The prime-to-\(\ell\) quotient at odd \(\ell\)
has vanishing \(H^1\) on the invariant module as well.
Thus the global cohomology class belongs to the strict
localization kernel used in the
\hyperlink{proof.selmer-minus-parts}{earlier critical-parity
comparison}.

Conjugation acts on coefficients by \((-1)^e\), so the
cohomology class of an odd cocycle has eigenvalue
\(\epsilon=(-1)^{e+1}\). Every coboundary as a function on
\(G_E\) is even, since \(\chi(cgc^{-1})=\chi(g)\). An odd
cocycle that is a coboundary in \(B\) is therefore killed
by two. The map to the strict localization kernel is
injective at odd \(\ell\), and has kernel of order at
most two at \(\ell=2\). For the latter assertion the
image of the coboundary map is \(B/B^{G_E}\), again a
divisible group whose two-torsion has order two, since
\(B^{G_E}\) is finite cyclic.

Set \(k=e+1\) and
\[
 H_{L,\ell,k}=H^2_{\mathrm{\acute et}}
    (\mathcal O_L[1/\ell],\mathbb Z_\ell(k))
 \qquad(L=E,F).
\]
At finite coefficients the cohomology over the ring of
\(\ell\)-integers agrees with Galois cohomology unramified
outside \(\ell\). Poitou--Tate duality pairs the finite
localization kernels in degrees one and two. Filtered
direct limits commute with these kernels since the set
of places above \(\ell\) is finite. The finite cohomology
groups in degree one give the Mittag--Leffler condition
for identifying inverse-limit cohomology in degree two.
Consequently the strict localization kernel is the
equivariant Pontryagin dual of the wild kernel
\(A\subset H_{E,\ell,k}\), as in the earlier Selmer
calculation. There are only complex infinite places.
The \(\ell\)-adic groups \(H_{L,\ell,k}\) are finite for
\(k\ge2\).\cite[VI, \S8]{f1weibel-kbook}
\cite[Propositions~6.12--6.13]{f2-rognes-weibel-dyadic}
The signed parts
have the same orders on the two sides, including at \(2\).
Indeed
\[
 |(A^\vee)^\epsilon|
 = |A/(1-\epsilon c)A|
 = |A^\epsilon|
 \le |H_{E,\ell,k}^{\,\epsilon}|.
\]
Here superscripts are the kernels of \(c-\epsilon\).
This yields the required bound on the raw finite cocycles
from the signed part of \(H_{E,\ell,k}\), with the one
possible factor of two just proved.
\cite[Theorem~I.4.10(a) and Proposition~II.2.9]{f1milne-adt}

Use the
\hyperlink{proof.shared-critical-selmer-values}{critical special values}
with \(L=E\), \(F_L=F\) and \(N_L=N_{\rm cyc}\), so that
\(\Lambda_k=\Lambda_k(E)\). This compares the ordinary groups
\(H_{L,\ell,k}\) calculated with the same coefficients as the
wild kernel.

At odd \(\ell\), \(H_{E,\ell,k}^+\simeq H_{F,\ell,k}\)
by transfer, and \(H_{E,\ell,k}\) is the sum of its two
sign parts. Thus the even-\(k\) formula controls the
invariant part and the odd-\(k\) ratio controls the
anti-invariant part, both by
\(v_\ell(\Lambda_k)+v_\ell(w_k(E))\), using
\(w_k(F)\mid w_k(E)\).

For the coefficient prime \(2\), abbreviate \(H_L=H_{L,2,k}\).
If \(k\) is even then \(2H_E^+\subseteq\operatorname{res}H_F\),
so \(|H_E^+|\le |H_F||H_E[2]|\).
If \(k\) is odd, then \(k\ge3\) and
\((1+c)H_E\supseteq 2\operatorname{res}H_F\).
The kernel of restriction is killed by \(2\), and
the kernel of doubling on its image is contained in \(H_E[2]\).
Consequently
\[
 |H_E^-| = \frac{|H_E|}{|(1+c)H_E|}
      \le\frac{|H_E|}{|H_F|}|H_F[2]||H_E[2]|.
\]

The coefficient sequence embeds \(H_L/2H_L\) into
\(H^2_{\mathrm{\acute et}}(\mathcal O_L[1/2],\mu_2)\).
By the Kummer sequence the latter has order at most
\[
 h_L\,2^{\#\{v\mid2\ \mathrm{in}\ L\}+r_1(L)}.
\]
Indeed, the class group of \(\mathcal O_L[1/2]\) is a
quotient of the ordinary class group, and the \(2\)-torsion
of the Brauer group has local invariants only at the real
and dyadic places. Since \(H_L\) is finite,
\(|H_L[2]|=|H_L/2H_L|\).
A common upper bound for the dyadic error, including the
possible factor of two on passing from raw cocycles to
classes, is therefore
\[
 C_2=1+2n_E+\lceil\log_2 h_E\rceil+\lceil\log_2 h_F\rceil .
\]
The bound at \(2\) now uses the critical part of the wild
kernel and applies to the weighted field.

For \(E\) the integer \(M_k=M_k(E)\) from the
\hyperlink{proof.shared-critical-selmer-values}{special value calculation}
satisfies \(v_\ell(\Lambda_k)\le v_\ell(M_k)\) and
\[
 \log|M_k|\ll n_Ek\log(2kN_{\rm cyc}),\qquad
 \log w_k(E)\ll n_Ek\log(2n_Ek).
\]

It follows that the single positive integer
\[
 X=h_E\,2^{C_2}\prod_{k=2}^t |M_k|w_k(E)
\]
bounds every finite cocycle length by \(v_\ell(X)\) for
\(O=\mathbb Z_\ell\). The factor \(2^{C_2}\) is shared
because the same bound is used separately for each \(e\).
For a finite coefficient extension \(O/\mathbb Z_\ell\),
the defining cocycle conditions commute with finite free
base change. In a \(\mathbb Z_\ell\)-basis of \(O\), each
of the coefficient components obeys the identical
conditions at the same modulus \(\ell^\nu\).
Ordinary \(O\)-length is therefore the ramification
index times \(\mathbb Z_\ell\)-length. With
\(v_\ell(\ell)=1\), the normalized bound remains \(v_\ell(X)\).

Finally \(E,F\) have root discriminants at most
\(N_{\rm cyc}\) by the cyclotomic different and the
discriminant tower formula. Minkowski's bound and the
count of ideals of each possible norm give
\(\log h_L\ll n_E^2\log(2N_{\rm cyc})\) for \(L=E,F\).
Summing the displayed estimates gives
\[
 \log X \ll n_E^2\log(2N_{\rm cyc})
          + n_E t^2\log(2n_EtN_{\rm cyc}).
\]

\paragraph{The local length and its sum over primes.}
The normalized local budget and the finite length used for
the Selmer tests are
\[
 b_\ell=h_2^{27}\bigl(c_\ell+v_\ell(X)\bigr)
       +4f\,10^6\frac{m}{\kappa}
            \bigl(b_1 v_\ell(g)+v_\ell(A)\bigr),\qquad
 M_\ell=\lfloor r b_\ell\rfloor+1 .
\]
The construction can use \(\beta_{j,\ell}=b_\ell\log\ell\) at
every stage. Use the compact local argument with the \(X\)
constructed above. The proof of the Tate block estimate uses the
upper bound \(v_\ell(X)\) on each relevant normalized cocycle
length. Thus the Tate block and crystalline defect estimates give
the allowance \(b_\ell=b_\ell(X)\) in this specialization. Both
the finite Selmer tests and the subsequent Hecke denominator proof
use this same \(M_\ell\).
\footnote{\weightedpin
{grh-compact/src/New/Selmer/RibetElement.lean}{204-L223}{Tate block proof using the cocycle allowance},
\weightedpin
{grh-compact/src/New/Selmer/DenominatorBound.lean}{43-L55}{Specialization to the same finite length}, and
\weightedpin
{grh-compact/src/New/Selmer/RibetElement.lean}{134-L147}{Choice of the coefficient ring}.}
The fixed declaration uses the integer from its own uniform bound.
The application here uses the displayed cocycle bound.

For a positive integer \(C\), the sum of
\(v_\ell(C)\log\ell\) over a finite set of primes is at most
\(\log C\). Also
\(\sum_\ell 2^{-\ell}\log\ell
 \le\sum_{n\ge1} n2^{-n}=2\).
Each chosen \(g_i\) has \(f\) places of norm \(g_i\), and
\(|\mathrm{GL}_{rt}(\mathbf F_{g_i})|\le g_i^{(rt)^2}\).
Hence \(\log G\le f(rt)^2\log g\).
On the scales of \(D\), this is \(O(x^{31+\varepsilon}\log x)\).
The sum of the allowances \(v_\ell=9c_\ell\log\ell\) is at most
\(9(\log G+2)\le x^{999}\) on every finite set of primes for large \(x\).

The bounds on \(X\), \(f\), \(N_{\rm cyc}\), \(r\) and \(t\) imply that
\(h_2^{27}(\log G+2+\log X)\) is polynomial in \(x\).
It is \(o(B_0)\) because \(a\ge xb\) and \(b=\exp(x^\varepsilon)\).
For the other term, the interval product in \(A\) has at most
\(60h_2^3\) terms. Since \(b_1=O(b)\), \(h_2\) is polynomial in \(x\)
and \(\log g=O(\log b)\), we have \(\log A=O(b\log b)\).
Consequently its contribution together with that of \(g\) is
\[
 O\!\left(f\frac{m}{\kappa}b\log b\right)
   =O\!\left(\frac{B_0}{(\log b)^2}\right),
\]
using \(a\ge xb\) and \(m^2\le4x\). This proves the claimed
finite-prime budget for large \(x\), uniformly over the coefficient
isomorphisms and over all stages.

\teamanchor{proof.weighted-fit-denominators}%
\paragraph{The split Hecke tests and the packets that pass them.}
\teamanchor{proof.weighted-fit-hecke-packets}%
Fix one of the auxiliary pairs and its twist. Let
\(\xi_w=(A_wq_w^{t'-1})^{-1}\) be the test parameter at an
eligible split place. The same spherical transfer that gives
\eqref{eq:weighted-fit-galois-string} kills all genuine lifts under \(H_w(\xi_w)\).

\paragraph{The local Hecke polynomial.}
\teamanchor{proof.weighted-fit-string-orthogonality}%
At a split place \(w\) write \(q=\operatorname{N}w\) and let \(\pi\)
be the unitary local action of \(\mathrm{GL}_n(E_w)\). Let \(I_w\)
be the upper pro-\(p_w\) Iwahori, where \(p_w=\operatorname{char}k_w\),
and \(P_I\) its orthogonal averaging operator. For a uniformizer \(\varpi\), put
\(t_w=\operatorname{diag}(\varpi,1,\ldots,1)\). On \(I_w\)-fixed
vectors the unnormalized Hecke operator \(U\) of \(I_wt_wI_w\)
is \(q^{n-1}P_I\pi(t_w)\). The \(q^{n-1}\) translations are
indexed by the freely varying first-row entries of
\(I_w/(I_w\cap t_wI_wt_w^{-1})\).

Let \(x\) be a nonzero spherical eigenvector at \(w\), and \(s_j\) the
eigenvalue of the sum of the translations indexed by sublattices
of \(\mathcal O_w^n\) containing \(\varpi\mathcal O_w^n\) and having
residue codimension \(j\). For a lattice \(L\), its translate
\(x_L\) is well defined because \(x\) is spherical. Write \(T_j\)
for the corresponding partial sum over the lattices whose reduction
contains the first line. Thus \(T_0=x\), \(T_n=0\), and
\(s_n\pi(t_w^{-1})x=T_{n-1}\), since the full top spherical
translation is by \(\varpi I_n\). It follows that \(s_n\ne0\).

These sums are \(I_w\)-fixed. For \(0\le j<n\),
\[
 U T_j=q^j(s_{j+1}x-T_{j+1}).
\]
Indeed \(t_w\) takes a lattice containing the first coordinate
line modulo \(\varpi\) to a lattice whose reduction lies in the
tail \(k_w^{n-1}\) and has dimension \(n-1-j\). The upper row
unipotents in \(I_w\) add all linear maps from that tail subspace
to the first line. Each one occurs with relative measure
\(q^{-(n-1-j)}\) under the average. The tail action permutes
the summed subspaces. All subspaces of codimension \(j+1\)
avoiding the first line occur as these graphs.

The adjoint on the \(I_w\)-fixed subspace is
\(U^*=q^{n-1}P_I\pi(t_w^{-1})\). For \(1\le j<n\), applying
\(t_w^{-1}\) to all the lattices of residue codimension \(j\)
gives \(T_{j-1}\) together with the translates at
\[
 \mathcal O_w(\varpi^{-1}e_1+\widetilde c)+L_W,\qquad
 W\subseteq k_w^{n-1},\quad
 \dim W=n-1-j,\quad
 c\in k_w^{n-1}/W ,
\]
where \(L_W\) is the inverse image of \(W\) in the tail lattice.
For subspaces containing the line, this is \(c=0\). For a graph
of a nonzero functional, \(W\) is its kernel and \(c\) is the
coset of the vectors of functional value \(1\). The graphs of
the zero functional give \(T_{j-1}\). The lower column
unipotents with entries in \(\varpi\mathcal O_w\) translate
\(c\) uniformly through its \(q^j\) possibilities, and the
tail action permutes \(W\). Averaging the transforms of
\(T_j\) gives \(q^{-j}\) times the sum over the displayed
family. Thus
\[
 s_j U^*x=q^{n-1}T_{j-1}+q^j U^*T_j,\qquad
 U^*x=q^{n-1}T_{n-1}/s_n,\qquad
 UU^*x=q^{2n-2}x .
\]

Let \(P_j(Z)=\sum_{i=0}^j(-1)^i q^{i(i-1)/2}s_iZ^{j-i}\),
where \(s_0=1\), and let \(P_x=P_n\) be the monic spherical
Hecke polynomial. The forward recurrence gives
\(P_j(U)x=(-1)^j q^{j(j-1)/2} T_j\).
If \(P_x(c)=0\), put \(z=(P_x(Z)/(Z-c))(U)x\). Then \(Uz=cz\).
Expand \(P_x(Z)/(Z-c)=\sum_{j<n}c^{n-1-j}P_j(Z)\).
The adjoint recurrence gives
\[
 cU^*z=q^{-n}P_x(qc)U^*x+q^{n-2}z .
\]
For this calculation the term with \(s_n U^*x\) cancels the
missing last \(T_{n-1}\) term by the displayed endpoint identity.
Taking the inner product of \(U^*x\) with this formula, with the
second argument linear, yields
\[
 c(q^n-1)\langle x,z\rangle=P_x(qc)\langle x,x\rangle .
\]
Every root \(c\) is nonzero since \(s_n\ne0\).
Taking the inner product with \(z\) and using \(Uz=cz\) gives
\[
 (|c|^2-q^{n-2})\|z\|^2
   =\frac{q^{-n}|P_x(qc)|^2}{q^n-1}\,\|x\|^2 .
\]
If \(|c|^2<q^{n-2}\), then \(P_x(qc)=0\).
Whenever \(P_x(qc)=0\), the adjoint identity gives
\(U^*z=q^{n-2}z/c\) and \(\langle x,z\rangle=0\).
It follows that \(z\) is orthogonal to every \(U^i x\), and
hence zero. The cyclic space generated by \(x\) under \(U\)
is the span of the \(T_j\), so both \(U\) and \(U^*\) preserve it.%
\footnote{\weightedpin
{grh-compact/src/Lib/LocalGL/Iwahori/HeckeCounts.lean}{113-L147}{Forward lattice count}
and \weightedpin
{grh/src/Lib/LocalGL/SphericalChainLemma.lean}{94-L203}{Dual lattice count and adjoint identity}.}

The roots of \(P_x\) are also closed under
\(c\mapsto q^{n-1}/\overline c\).
Indeed, the inverse of the \(i\)-th minuscule double coset
is the \((n-i)\)-th minuscule double coset multiplied by
\((\varpi I_n)^{-1}\). The adjunction on spherical vectors
gives \(s_{n-i}=s_n\overline{s_i}\). Substituting the
coefficients of \(P_x\) then gives
\[
 Z^n\overline{P_x(q^{n-1}/\overline Z)}
   =(-1)^n q^{n(n-1)/2}s_n^{-1}P_x(Z).
\]

\paragraph{Cofinality on the fixed denominator packet.}
\teamanchor{proof.weighted-fit-packet-cofinality}%
Fix the auxiliary pair of ranks \(r'\) and \(n\), its
base level and label, the coefficient isomorphism \(\iota\),
and the chosen auxiliary twist \(\vartheta\) with inverse
avatar \(\psi\). Put \(t'=n-r'\). Denote the attached
normalized representations of the finite set of base
packets by \(N_x\), and let \(x_0\) be the genuine packet
under test. Its main block is \(U_0=\psi T\) for the same
row \(T\) as in the Selmer estimate. Its characteristic
polynomial at \(g\) is the product of that of \(U_0(g)\) and
\(\prod_{0\le i<t'}(X-\chi(g)^{-i})\).
This follows from the trace formula by applying it at
the powers of \(g\) and using Newton's identities. The
denominator test uses the fixed base packet space and these
normalized attachments, with the local polynomial
annihilation and the generic Frobenius coefficient
approximation of those attachments. A packet is weak when
each \(N_x(g)\) has eigenvalue \(1\).
The normalized representations have integral
characteristic polynomials. For this assertion, boundedness
of the matrix entries on the compact Galois group bounds
all matrix powers for each \(g\). Applying the powers to
an eigenvector shows its eigenvalue has \(\ell\)-adic norm at
most \(1\). This also applies to the main block \(U_0\).

Write \(q=\operatorname{N}w\) and \(\phi\) for the
geometric Frobenius of an off-level split datum. Let \(A_w\)
be the scalar of the spherical test. Thus the multiplier
in \(N_x(\phi)\) relative to the attached representation is
\(\iota(A_w)^{-1}\), and
\(|A_w|_{\mathbf C}^2=q^{r'}\). Its \(\ell\)-adic image has
norm \(1\), by continuity of the normalizing character
on the compact Galois group. If \(x\) is weak, exact local
annihilation implies \(P_{x,w}(A_w)=0\) for its spherical
Hecke polynomial \(P_{x,w}\). Here the weak eigenvalue \(1\)
is an eigenvalue of the normalized representation, so
the unnormalized eigenvalue is precisely \(\iota(A_w)\).

For the spherical eigenvector on this packet, the
\hyperlink{proof.weighted-fit-string-orthogonality}{local Hecke polynomial}
\(P_x\) is \(P_{x,w}\), with coefficient
\((-1)^i q^{i(i-1)/2}s_i\) in degree \(n-i\).
Starting with the root \(A_w\), the local norm identity gives
the roots \(A_wq^i\) through \(i=\lfloor(t'-1)/2\rfloor\).
At each preceding index \(i\) one has \(2i<t'-2\), so
\[
 |A_wq^i|^2=q^{n-t'+2i}<q^{n-2}.
\]
The reciprocal-conjugate root symmetry sends \(A_wq^i\) to
\(A_wq^{t'-1-i}\) because \(|A_w|^2=q^{n-t'}\).
This gives the distinct exact roots \(A_wq^i\) for \(0\le i<t'\).
In particular the reciprocal test at
\(\xi_w=(A_wq^{t'-1})^{-1}\) vanishes on weak packets.

Use the distinguished row's scalar inertia element
\(\gamma_{\mathrm{sc}}\) at a chosen level prime off \(\ell\)
and write \(T(\gamma_{\mathrm{sc}})=yI\). The label
calculation above gives residual order \(>4r'+2\) for
\(y\). The auxiliary twist and its conjugate conductor
avoid all the chosen level primes, and \(\chi\) is
unramified at this prime. Thus
\(\psi(\gamma_{\mathrm{sc}})=\chi(\gamma_{\mathrm{sc}})=1\).
The powers \(y^{i\ell^\nu}\), \(0\le i\le r'\), have
distinct residues for every \(\nu\), since taking an
\(\ell^\nu\)-power is injective in the residue field.
For any \(g\in G_E\), the residual polynomial
\(\det(1-ZU_0(g))\) has degree at most \(r'\) and constant
term \(1\). Choose \(i\) for which its value at the
residue of \(y^{i\ell^\nu}\) is nonzero. Then
\(\delta=\gamma_{\mathrm{sc}}^i\) lies in \(\ker\chi\), and
\(\det(1-U_0(\delta^{\ell^\nu}g))\) is an \(\ell\)-adic unit.
This unit condition persists throughout an open neighborhood
of \(\delta^{\ell^\nu}g\).

On the same Frobenius datum put
\[
 B_w(X)=\prod_{0\le i<t'}(X-q^i),\qquad
 C_{x,w}(X)=\iota(A_w)^{-n}\iota\bigl(P_{x,w}(A_w X)\bigr).
\]
Here the isomorphism acts on coefficients. The chain just
proved gives the exact factorization \(C_{x,w}=B_w T_{x,w}\)
for each weak packet. The generic coefficient approximation
of the attached representation, with accuracy \(<1\),
gives \(C_{x,w}\equiv\operatorname{charpoly} N_x(\phi)\)
modulo the maximal ideal of the \(\ell\)-adic integers.
The normalization preserves this accuracy since it scales
each coefficient by a unit. In particular \(C_{x,w}\) is
integral. If its prescribed residual comparator is
\(B_w S_w\), with \(|S_w(1)|_\ell\ge1\),
division by the integral monic \(B_w\) makes
\(T_{x,w}\) integral and \(T_{x,w}(1)\) a unit: first
\(B_w S_w\) is integral by the coefficient congruence, so
\(S_w\) is integral, and then the reduced monic factor
cancels to give \(\overline{T}_{x,w}=\overline S_w\).
For every weak packet residually congruent to \(x_0\),
use \(S_w=\operatorname{charpoly} U_0(\phi)\), since
\(\chi(\phi)^{-1}=q\) and the unit test has been imposed.
Each polynomial \(T_{x,w}/T_{x,w}(1)\) is integral with
value \(1\) at \(1\). Their finite product \(e_w\) has the
same value and is divisible by each exact Hecke quotient
in the rescaled variable, which differs from \(T_{x,w}\)
by the nonzero scalar \(\iota(A_w)^n\).

For any neighborhood of \(\delta^{\ell^\nu}g\), intersect
with the open unit condition. For each of the finitely
many weak packets, the generic Frobenius coefficient
approximation gives a nonempty open subset of the
current open set on which all degree-one Frobenius data
off a finite set satisfy that approximation. Refine
successively through the finite packet list. Exclude the
base-level bad places and their conjugates, the coefficient
prime, and the conductor of the fixed twist and its
conjugate. Exclude also any additional finite set required
by the triangle. There is a degree-one Frobenius datum
in the resulting open set by Chebotarev. This is the
cofinality after scalar perturbation required at the
prime-selection step. In the triangular application
\(g\in G_\nu\) and the open cyclotomic and evaluation
conditions are included in the prescribed neighborhood.
The center \(\delta^{\ell^\nu}g\) satisfies them because
\(\delta^{\ell^\nu}\) contributes zero to the additive
evaluations on \(G_\nu\).%
\footnote{\weightedpin{grh-compact/src/New/Selmer/AptTupleDefs.lean}{47-L82}
{The actual scalar inertia and normalization},
\weightedpin{grh-compact/src/Lib/GaloisRep/Main/AttachedDefs.lean}{13-L29}
{Generic coefficient approximation and exact local annihilation},
\weightedpin{grh-compact/src/Lib/ClassField/Chebotarev/FrobeniusFilter.lean}{48-L64}
{Meaning of generic Frobenius data}, and
\weightedpin{grh-compact/src/New/Selmer/DenominatorBound.lean}{98-L128}
{The packet condition and product used in the denominator test}.}

For each packet whose Galois
characteristic polynomial differs from that of \(x_0\) on
residue coefficients, choose a coefficient function \(h\)
and a group element witnessing the difference. Approximate
the group element at generic spherical Frobenii, where the
normalized monic Hecke coefficients give the limits. The
integral Hecke algebra is closed since preservation of the
Weyl integral class values is an intersection of closed
norm inequalities. Thus \(h\) is in this algebra, and
\(h-h(x)\) vanishes at that packet and has unit value at
\(x_0\).
\footnote{\weightedpin{grh-compact/src/Lib/LocalGL/TaylorWiles/HeckeAlgebra.lean}{70-L83}{Closed integral algebra}.}

\paragraph{Separating the test from the unwanted packets.}
\teamanchor{proof.weighted-fit-hecke-separator}%
For the fixed tuple label \(c\), use the base space containing
the image of its idempotent \(e_c\). The finite group character
formula shows that \(G e_c\) preserves the integral smooth
forms at \(j_\iota\): the group order divides \(G\), the
character values are algebraic integers, and the tame factors
are units. Work packetwise on this base level and then at the
indicated Taylor--Wiles level. Take the Ribet normalizer
\[
 B_x=N_x(\gamma a)-I-y_0(N_x(a)-I)
          +\ell^{M_\ell+h_*}\sum_g \lambda_g(N_x(g)-I),
 \qquad b_x=\det B_x ,
\]
where \(\gamma,y_0\) are the scalar data on \(x_0\) in its
short Selmer condition and \(a\) is the cyclotomic separator.
For every genuine packet \(b_x=0\) by the trivial summand in
\eqref{eq:weighted-fit-galois-string}. For every nonweak base packet there is \(g\) with
\(\det(N_x(g)-I)\ne0\). One can choose the finitely supported
nonnegative integers \(\lambda_g\) so \(b_x\ne0\) on all such
packets. Add the terms successively: at the next packet the
determinant as a polynomial in the new integer coefficient
has nonzero leading term, and the determinants on the
previous packets are nonzero at coefficient zero. Avoid
their finitely many roots. Each \(b_x\) is integral by the
stable lattice. The polynomial
\[
 p(T)=\prod_{x\ {\rm nonweak}}(T-b_x)
\]
is integral and \(p(0)\ne0\), with the empty product equal to
\(1\). The base normalizer and \(p(0)\) here are fixed before
choosing the depth of the triangular primes. The locus
\(p(b_x)\ne0\) on the base is contained in
the weak packets. A weak packet can still have \(b_x\ne0\).
For example, each element of the group of diagonal matrices
\(I,\operatorname{diag}(1,-1,-1),\operatorname{diag}(-1,1,-1),
\operatorname{diag}(-1,-1,1)\) has eigenvalue \(1\); the sum
of its three nonidentity elements minus \(3I\) is \(-4I\).
It is the right-hand orthogonality of the geometric test that
controls weak packets outside the genuine lift range.

Write \(l=e_c\Xi'\) and \(w=e_cP_Z\Xi'\). The commuting normal
Hecke actions and \(e_c\) preserve the orthogonal decomposition
by packets, and \(Z\) is an invariant subspace. Thus the
weak-packet components of \(l-w\) lie in
\(\operatorname{im}e_c\cap Z^\perp\cap\ker H_v(\xi_v)\) for
an eligible split \(v\). The right-hand test gives zero pairing
with each component. Its projection to the \(Z\)-part gives
zero as well, so these components vanish. Consequently \(l,w\)
agree after every diagonal operation supported where
\(p(b_x)\ne0\). The form \(w\) is supported on the genuine
packets. This verification allows a complementary subspace
within the same weak Hecke packet as a genuine lift.%
\footnote{\weightedpin{grh-compact/src/New/Selmer/BaseLevelOfLabel.lean}{18-L74}{Integral tuple image and normalized competing columns},
\weightedpin{grh-compact/src/New/Selmer/DenominatorBound.lean}{62-L97}{Normalizer},
and \weightedpin{grh-compact/src/New/Selmer/DenominatorBound.lean}{133-L166}{Tested projection agreement}.}

\paragraph{The triangular and first-order inputs.}
\teamanchor{proof.weighted-fit-first-order}%
Let \(t'\) be the rank difference of the auxiliary pair in this
denominator test. Choose a cyclotomic separator \(\gamma\) with
\(\chi(\gamma)^j-1\mid\ell^{h_0}\) for \(0<j<t'\).
For \(\ell\le h_2^2\) use Frobenius above one of the designated
primes \(g_i\ne\ell\): every \(g_i^j-1\) for \(0<j<t\) divides \(G\),
so one may take \(h_0\le v_\ell(G)\). For \(\ell>h_2^2\) the base
is unramified over \(\mathbf Q_\ell\); use cyclotomic inertia with
residue order \(\ell-1>t-1\) and \(h_0=0\). The assembled auxiliary
primes have \(q_w\equiv1\bmod\ell^\nu\) and avoid the conductor
of the fitted twist.

\paragraph{Uniform detection of the unramified Tate cocycles.}
\teamanchor{proof.weighted-fit-triangle-detection}%
Let \(E\subseteq\mathbf Q(\zeta_N)\) and let \(\chi\) be the
\(\ell\)-adic cyclotomic character on \(G_E\). There is a finite
set \(S_{\rm det}\subset G_E\) that detects every continuous
\(O/\ell^\nu O\)-valued cocycle for \(\chi^{-e}\) unramified off
\(\ell\), for every \(\nu\ge1\), \(e\ge0\), and ring of integers
\(O\) of a finite extension of \(\mathbf Q_\ell\). Unramified
means zero on the inertia groups at the places outside the
allowed set. Detection means that zero values on \(S_{\rm det}\)
force the cocycle to vanish. The same set satisfies
\(\#S_{\rm det}\le y^{20}\) for \(N\le y\) and \(\ell\le y^2\).

Put \(M'=N\ell^2\), \(F=\mathbf Q(\zeta_{M'})\), and \(d=[F:\mathbf
Q]\le M'\). Since \(M'\ge4\), the field \(F\) is totally imaginary.
Take coset representatives for \(G_E/G_F\) and lifts
in \(G_F\) of a basis of the maximal elementary abelian
\(\ell\)-quotient of \(G_F\) unramified away from \(\ell\).
Global Artin reciprocity bounds the dimension of that quotient by
\[
 \dim_{\mathbf F_\ell}\operatorname{Cl}(F)/\ell+
 \sum_{v\mid\ell}\dim_{\mathbf F_\ell}
 \mathcal O_{F_v}^{\times}/(\mathcal O_{F_v}^{\times})^\ell
 \ \le\ \log_\ell h_F+2d.
\]
Indeed, in the inverse limit of the ray class groups for
moduli supported above \(\ell\), the kernel of passage to the ordinary
class group is the image of the product of local unit groups.
At \(v\mid\ell\) the deep principal units are isomorphic by the
local logarithm to a full additive lattice of rank
\([F_v:\mathbf Q_\ell]\). The pro-\(\ell\) unit group has this
rank and cyclic torsion (the local \(\ell\)-power roots of
unity). Summing its generator bounds uses
\(\sum_{v\mid\ell}[F_v:\mathbf Q_\ell]=d\) and
\(\#\{v\mid\ell\}\le d\).

The field discriminant is at most \((M')^d\) in absolute value.
For this upper bound one can use the power basis for
\(\zeta_{M'}\): the identity obtained by differentiating
\(X^{M'}-1=\Phi_{M'}(X)((X^{M'}-1)/\Phi_{M'}(X))\) at
\(\zeta_{M'}\) shows that the norm of \(\Phi'_{M'}(\zeta_{M'})\)
divides \((M')^d\) up to sign. Minkowski gives an integral
ideal of norm at most \((M')^d\) in every ideal class, since
\((4/\pi)^{r_2}d!/d^d\le2^d\) and \(M'\ge4\).
The number of sublattices of index \(n\) in \(\mathbf Z^d\)
is at most \(n^{2d-2}\). In a Hermite form there are at most
\(n^{d-1}\) choices of the diagonal, and for each diagonal
the reduction bounds on the other entries give at most
\(n^{d-1}\) choices. Counting ideals among the sublattices up
to index \((M')^d\) gives \(h_F\le (M')^{2d^2}\).
Consequently \(\#S_{\rm det}\le3M'+2(M')^3\le3(M')^3\).
If \(N\le y\) and \(\ell\le y^2\), then \(M'\le y^5\) and
\(3y^{15}\le y^{20}\).

For detection, fix the module \(M=O/\ell^\nu O\) and
induct on \(k\) for cocycles with values in \(M[\ell^k]\),
starting with the zero module at \(k=0\). If \(z\) vanishes
on \(S_{\rm det}\), so does \(\ell z\), which takes values
in \(M[\ell^{k-1}]\) and is zero by induction. The resulting
\(\ell\)-torsion-valued restriction to \(G_F\) is a
homomorphism, since \(F\) contains \(\mu_\ell\). It is
unramified off \(\ell\) and vanishes on the chosen
generators, hence is zero. For \(b\in G_F\) and a chosen
coset representative \(s\), the identity
\(z(sb)=z(s)+\chi(s)^{-e}z(b)\) gives \(z=0\).
This proves detection for a fixed set, including cocycles
with values in the ramified coefficient rings.
\footnote{For Artin reciprocity and ray class fields use Artin and Tate,
\href{https://doi.org/10.1090/chel/366}{\emph{Class Field Theory}},
Chapter VII \S3 and Chapter VIII \S\S1--2. For the convex body theorem use Minkowski,
\href{https://archive.org/details/geometriederzahl00minkrich}{\emph{Geometrie der Zahlen}}.
Compare the pinned
\weightedpin{grh-compact/src/Lib/Cyclotomic/OddSha/OddShaBound.lean}{16-L57}{detection result}
and \weightedpin{grh/src/Lib/Cyclotomic/OddSha.lean}{182}{readable detection result}.}

\paragraph{Triangular primes from the primitive reciprocity laws.}
\teamanchor{proof.weighted-fit-triangle-assembly}%
Fix the field \(E\), selected coefficient prime \(\ell\),
denominator test and fitted twist of this application.
Here \(E\) is the CM field containing \(\mu_{5040}\).
It is stably unramified in the sense of the primitive law,
and its number of complex places is at least
\(\varphi(5040)/2\ge4\). To see
the unit conjugate difference used by that law, take a
primitive \(\zeta_{28}\in E\), put \(i=\zeta_{28}^7\) and
\(\omega=\zeta_{28}^2-\zeta_{28}^4-\zeta_{28}^8\).
The fourteenth cyclotomic equation gives
\(\omega+\overline\omega=1\), so \(i\omega-\overline{i\omega}=i\)
is a unit in \(\mathcal O_E\).
Use \(N\) equal to the cyclotomic conductor used for \(E\),
\(y=h_2\), and \(t_{\mathrm a}>0\) equal to the rank difference of
the auxiliary pair in the current denominator test. Use the
cyclotomic separator \(\gamma\) and loss \(h_0\) described above.
Thus \(N+t_{\mathrm a}<y\), \(\ell\le y^2\) or \(h_0=0\), and
\[
 \chi(\gamma)^j-1\mid\ell^{h_0}\quad(0<j<t_{\mathrm a}).
\]
Set \(a=\gamma^{-1}\) and \(H=yh_0\).
Here \(y\ge t_{\mathrm a}+2\). For \(0\le e,k<t_{\mathrm a}\)
with \(k\ne e\), every difference
\(\chi(\gamma)^e-\chi(\gamma)^k\) is a unit times a signed
difference \(\chi(\gamma)^j-1\) with \(0<j<t_{\mathrm a}\).
Thus the character-separation product
\[
 \prod_{\substack{0\le k<t_{\mathrm a}\\ k\ne e}}
      \bigl(\chi(\gamma)^e-\chi(\gamma)^k\bigr)\mid\ell^H
       \qquad (0\le e<t_{\mathrm a})
\]
has the required divisibility in \(\mathbf Z_\ell\), including
the empty product if \(t_{\mathrm a}=1\). For \(t_{\mathrm a}>1\)
one also has \(v_\ell(4)\le h_0\). At \(\ell=2\) this uses
\(\mu_4\subset E\) and the nonzero difference at exponent \(1\);
at odd \(\ell\) it holds because \(v_\ell(4)=0\). It follows for
\(0<e<t_{\mathrm a}\) that \(4(\chi(a)^{-e}-1)^2\mid\ell^H\)
in \(\mathbf Z_\ell\), hence also in \(O\).
\footnote{\weightedpin{grh-compact/src/New/Data/Different.lean}{294-L296}{Roots of unity in the datum}
and \weightedpin{grh-compact/src/Lib/Adelic/Adeles/CoprimeValuation.lean}{228-L246}{Stable unramifiedness and the unit conjugate difference}.}
For each \(\nu\ge1\),
put \(A=\mathbf Z/\ell^\nu\) and \(G_\nu=\ker(\chi\bmod\ell^\nu)\).
An even cocycle of type \(e\) has values in \(O/\ell^\nu O\)
with action \(\chi^{-e}\), is invariant under \(g\mapsto g^c\)
for the fixed complex conjugation \(c\), and vanishes on the
intersection of each decomposition group above \(\ell\) with
\(\ker\chi\). Its other possible ramification off \(\ell\) is
at the indicated set of split primes and their conjugates.

The primitive law uses the earlier places of the same type. For a
list of \(s\) earlier type-\(e\) Frobenius data, off \(\ell\),
at degree-one primes over distinct primes of \(E^+\), with
Frobenii in \(G_\nu\), prescribe at most \(s\) elements of
\(G_\nu\). There is a locally constant \(z:G_E\to A\) that
vanishes on these elements and takes a unit value on \(G_\nu\).
Its restriction to \(G_\nu\) is additive and satisfies
\[
 z(gbg^{-1})=\chi(g)^{e+1}z(b),\qquad
 z(b^c)=-z(b)\qquad (g\in G_E,\ b\in G_\nu).
\]
It comes with a finite exceptional set of primes of \(E^+\).
For a further family of \(\nu\)-split Frobenius data
\((w_i,\phi_i)\) over distinct primes outside this set, if
\(z(\phi_i)=0\) except possibly at \(j\), every even
type-\(e\) cocycle \(b\) ramified only at \(\ell,w_i,\bar w_i\)
and vanishing at the earlier listed type-\(e\) Frobenii
satisfies
\[
 n\,b(\tau_j)=0\qquad\text{whenever }2z(\phi_j)\mid n
 \text{ in }A .
\]
Here \(\tau_j\) is the tame generator at the chosen place
above \(w_j\) in its Frobenius datum. We use this assertion
with \(n=\ell^H\), including its factor \(2\) in residue
characteristic \(2\).

\paragraph{Primitive reciprocity on the real subfield.}
\teamanchor{proof.weighted-fit-primitive-reciprocity}%
Put \(K=E^+\) and retain the fixed complex conjugation \(c\). At
every finite place of \(E\) the unit conjugate difference above
excludes the nontrivial element of \(\operatorname{Gal}(E/K)\) from
inertia. Thus \(E/K\) is unramified at all finite places, and every
prime of \(E\) of prime norm lies over a split prime of \(K\).
Write \(\chi_K\) for the cyclotomic character on \(G_K\).
Extend its restriction \(\chi\) from \(G_E\) to the character
\(\kappa(gc^a)=\chi(g)\) of \(G_K\), for \(a=0,1\). The
cyclotomic character is invariant under conjugation by \(c\), so
this is a character, unramified away from \(\ell\) at finite
places. Every real complex conjugation has \(\kappa\)-value \(1\).
At odd \(\ell\), \(\mathbf Z_\ell^\times\) is procyclic. At
\(\ell=2\), the inclusion \(\mu_4\subseteq E\) puts the image
of \(\kappa\) in \(1+4\mathbf Z_2\), also procyclic.

Fix \(\nu\ge1\), let \(A=\mathbf Z/\ell^\nu\), and put
\[
 M=A(\kappa^{-e}),\qquad B=\operatorname{Hom}(M,\mu_{\ell^\nu}),
 \qquad U=G_{E(\mu_{\ell^\nu})}=G_\nu .
\]
Both coefficient modules are trivial on \(U\), and \(G_K/U\)
is the product of the group of order two generated by \(c\)
and the cyclic image of \(\chi\bmod\ell^\nu\) on \(G_E\).
The character of \(B\) is \(\chi_K\kappa^e\bmod\ell^\nu\).
Write \(P_0\) for the \(s\) distinct primes of \(K\) below the
earlier type-\(e\) data. Let \(k\le s\) be the number of
prescribed elements \(u_1,\ldots,u_k\in U\). Let \(S\) consist
of \(P_0\), the primes over \(\ell\) and the real places.

We use local Tate pairings for \(M,B\) and Poitou--Tate
orthogonality of the two global localization images. For any
finite set \(V\) of test primes split in \(E(\mu_{\ell^\nu})/K\)
and disjoint from \(S\), take the following local conditions for
primal cocycles. At \(v\mid\ell\), require zero on
\(D_v\cap\ker\kappa\), with \(D_v\) the chosen decomposition
group in \(G_K\). At \(P_0\) and at the real places require zero
on \(D_v\). Allow all cocycles at \(V\) and unramified cocycles
at other finite places. The conditions include all local
coboundaries. At \(P_0\) the chosen Frobenius is in \(U\), so
the entire local coefficient action is trivial. At \(V\) a
local primal cocycle is determined by its unramified and tame
coordinates. Evaluation on an \(\ell^\nu\)-tame generator
\(\tau_v\) gives the tame coordinate. The local unramified
dual group pairs perfectly with this coordinate. Indeed local
class field theory and the Kummer class of a uniformizer give
the invariant as the product of a tame value and a dual
Frobenius value, divided by \(\ell^\nu\), up to a unit
depending on the chosen generators. At all finite places off
\(\ell\) where the modules are unramified, the unramified
subgroups are mutual annihilators. At a real place the degree
one local group in Tate duality is the usual \(H^1\), including
at \(\ell=2\).%
\footnote{\href{https://math.stanford.edu/~conrad/BSDseminar/refs/TateICM.pdf\#page=2}{Tate, Stockholm ICM (1962), Theorems 2.1, 2.2, 2.4 and 3.1}
for the local and global formulas.
See \href{https://intlpress.com/site/pub/files/_fulltext/journals/mrl/2015/0022/0006/MRL-2015-0022-0006-a005.pdf}{\v{C}esnavi\v{c}ius, Theorem 5.1 and \S6}
for the exact sequence with the real-place conventions.}

Write \(\bar M=M[\ell]\) and
\(\bar B=\operatorname{Hom}(\bar M,\mu_\ell)\).
Choose \(V\) using the two finite groups
\(H^1(G_{K,S},\bar B)\) and \(H^1(G_{K,S},B)\), where
\(G_{K,S}\) is the Galois group unramified outside \(S\).
Finiteness follows from the finite-coefficient global duality
theorem. For each element of the first group with nonzero
restriction to \(U\), choose a witness in \(U\) with nonzero
value. Add a distinct prime to \(V\) with an arithmetic
Frobenius in \(U\) having this value. For each \(u_i\), add a
distinct prime \(v_i\) whose arithmetic Frobenius in \(U\) has
the same values as \(u_i\) under all classes of the second
group. Restrictions to \(U\) are homomorphisms independent of
the choice of cocycle representative, so these conditions
specify open neighborhoods in \(G_K\). Chebotarev supplies
prime-norm primes with chosen places above them and allows any
finite exclusions. Each of these primes splits in the fixed
field of \(U\). Extra primes split in that field may also
be included in \(V\).

Let \(Y_\nu\) be the group of primal global cocycles with the
listed conditions. Its \(\ell\)-torsion is exactly the group
of cocycles with those conditions valued in \(\bar M\).
The dual Selmer conditions for \(\bar B\) are zero at \(V\),
since the primal conditions are full there. In particular
these dual classes are in \(H^1(G_{K,S},\bar B)\).
By the choice of \(V\), they restrict to zero on \(U\).
Inflation--restriction puts them in
\(H^1(G_K/U,\bar B)\). This last group has dimension at most
\(1+h^0(K,\bar B)\). A cocycle on the two-generated group
\(G_K/U\) is determined by its two values on generators, and
its coboundaries have dimension \(1-h^0(K,\bar B)\).

For the \(\bar M\) local subspaces \(L_v\) just specified,
the Selmer dimension formula is
\[
 h^1_L(\bar M)-h^1_{L^\perp}(\bar B)
 =h^0(K,\bar M)-h^0(K,\bar B)
   +\sum_{v\in S\cup V}\bigl(\dim L_v-h^0(K_v,\bar M)\bigr).
\]
In this formula \(h^0(K_v,\bar M)\) counts ordinary
invariants, also at the real places. To see the real-place
term directly, let \(I_M,I_B\) be the two localization images
on \(S\cup V\). Poitou--Tate gives \(I_M^\perp=I_B\).
Taking the dimensions of their intersections with the local
conditions gives the Selmer difference as
\[
 h^1(G_{K,S\cup V},\bar M)
 -\dim\ker(\mathrm{loc}_{\bar B})
 +\sum_{v\in S\cup V}(\dim L_v-h^1(K_v,\bar M)).
\]
By the degree-two terms of duality,
\(\dim\ker(\mathrm{loc}_{\bar B})=h^2(G_{K,S\cup V},\bar M)
-\sum_v h^2(K_v,\bar M)+h^0(K,\bar B)\).
For the rank-one module \(\bar M\), each finite-place
\(h^2-h^1\) is \(-h^0\), with the further subtraction
\([K_v:\mathbf Q_\ell]\) at \(v\mid\ell\). Each real-place
\(h^2-h^1\) is zero by periodicity. The global Euler
characteristic in dimensions is zero since \(\bar M\) is
trivial at all real places of the totally real field \(K\).
The sum of the extra finite-place terms is \(-[K:\mathbf Q]\).
It equals the sum of \(-h^0(K_v,\bar M)\) over the real
places, giving the asserted formula.%
\footnote{\href{https://www.numdam.org/item/SB_1964-1966__9__415_0.pdf\#page=10}{Tate, Bourbaki expos\'e 306 (1965/66), Theorem 2.2}
for the global Euler characteristic, and the local Euler
characteristic in Tate's Stockholm ICM lecture.}

At \(v\mid\ell\) the permitted cocycles factor through the
procyclic group \(\kappa(D_v)\). Evaluation on a generator is
injective, giving \(\dim L_v-h^0(K_v,\bar M)\le0\) after
subtracting the local coboundaries. At \(P_0\) and at each
real place the summand is \(-1\). At \(V\) it is \(1\), since
the two coordinates give all of \(H^1\). Adding back the
global coboundary dimension \(1-h^0(K,\bar M)\) yields
\[
 \dim_{\mathbf F_\ell}Y_\nu[\ell]
 \le |V|+2-s-[K:\mathbf Q]\le |V|-k-1 .
\]
This calculation uses the group killed by \(\ell\) at the
original depth. It therefore controls the number of
generators of \(Y_\nu\) over \(A\).

Let \(D\subseteq A^V\) be generated by the tame evaluation
vectors of \(Y_\nu\) and the \(k\) coordinate vectors of the
primes \(v_i\). The bound shows that \(D\) has at most
\(|V|-1\) generators. Smith form over \(A\) yields a vector
\(r\in A^V\) in its annihilator for the coordinate dot
product, with a unit coordinate. Indeed, in elementary-divisor
coordinates one index has zero cyclic summand, and the dual
coordinate vector is primitive. In particular \(r_{v_i}=0\).
Use the local unramified dual coordinates at \(V\) to
represent \(r\) by a tuple \(\beta_{\mathrm{loc}}\) whose
pairing with a primal local tuple is
\(\ell^{-\nu}\sum_{v\in V} r_v x_v(\tau_v)\bmod\mathbf Z\).
Set its coordinates on \(S\) equal to zero. It annihilates
\(I_M\cap \prod_v L_v\). For these finite cohomology groups
the annihilator of this intersection is
\[
 I_B+\bigl(\textstyle\prod_v L_v\bigr)^\perp.
\]
Consequently a global class \(\beta\) has the same
localizations as \(\beta_{\mathrm{loc}}\) at \(V\), and meets
the dual conditions at \(S\). Since it is unramified at
\(V\), it comes from \(H^1(G_{K,S},B)\). Restrict a
representative to \(U\), identify the value module with \(A\)
by a generator of \(\mu_{\ell^\nu}\), and call the resulting
homomorphism \(z\). Then \(z(u_i)=0\) by the choice of \(v_i\).
A unit coordinate at \(V\) makes \(z\) a unit at its
arithmetic Frobenius. For \(g\in G_K\) and \(b\in U\), the
cocycle identity gives
\(z(gbg^{-1})=\chi_K(g)\kappa(g)^e z(b)\). Its two
specializations give the required twist and oddness.
Extend \(z\) by zero on the open-and-closed complement of
\(U\) in \(G_E\). The resulting function is locally constant.

Take the finite exceptional set to contain the finite places
of \(S\) and \(V\). For any further family in the reciprocity
statement, first take an \(A\)-valued even cocycle \(b\) on
\(G_E\). It extends to a \(\kappa^{-e}\)-cocycle \(b^+\) on
\(G_K\) by \(b^+(gc^a)=b(g)\), using its evenness and
\(\kappa(c)=1\). At \(v\mid\ell\), for
\(g\in D_v\cap\ker\kappa\) one has \(g^2\in
G_E\cap D_v\cap\ker\chi\). Hence
\(2b^+(g)=b(g^2)=0\). At a real place the same calculation
uses the square of a complex conjugation. At \(P_0\), the
unramified restriction is zero on the prescribed geometric
Frobenius, and hence on the decomposition group. At every
remaining finite place where \(b\) is unramified, so is
\(b^+\), since the inertia group lies in \(G_E\).

Thus the local invariant of \((2b^+)\cup\beta\) is zero
at \(S\), by the dual conditions, and at every other
place outside the new ramification set, by unramifiedness.
At a new place the given geometric Frobenius and tame
generator are those of the same decomposition group over
\(K\), because the place splits in \(E\). The dual
Frobenius value is \(z(\phi_i)\), up to switching arithmetic
and geometric Frobenius. The sum of local invariants is
zero. Since all the other \(z(\phi_i)\) vanish, the pairing
at \(j\) gives \(2z(\phi_j)b(\tau_j)=0\) in \(A\), with a
unit change allowed for the given tame generator. This
implies the assertion for any multiple \(n\) of \(2z(\phi_j)\).
For \(O/\ell^\nu O\)-valued cocycles, take coordinates in a
\(\mathbf Z_\ell\)-basis of \(O\). All the local conditions
and equality at the prescribed points hold coordinatewise,
and the same \(z\) works on every coordinate.
Thus no division by \(2\) enters this reciprocity law.%
\footnote{\weightedpin{grh-compact/src/New/Selmer/ReciprocityLaw.lean}{24-L57}
{The required local law and its quantifiers}, and
\weightedpin{grh-compact/src/New/Selmer/ReciprocityLaw.lean}{162-L192}
{Descent and tame transport in the compact proof}.}

Order \(J=\{1,\ldots,t_{\mathrm a}-1\}\times S_{\rm det}\)
lexicographically, using any fixed order on \(S_{\rm det}\).
The first coordinate is the type. For the test point
\(s\in S_{\rm det}\) write \(g_s=(as)^{-1}sa\in\ker\chi\).
At depth \(\nu\) construct the primes, laws \(z_j\) and
exceptional sets in this order. Before the step
\(j=(e,s)\), the simultaneous image of \(G_\nu\) under the
earlier type-\(e\) laws is an \(A\)-submodule of \(A^q\),
where \(q\) is the number of those laws. Any submodule of
\(A^q\) is generated by at most \(q\) elements, by taking its
preimage in \(\mathbf Z_\ell^q\). Prescribe lifts of these
generators as the zeros of the new primitive law. Take an
element where the new law is a unit and multiply by powers
of the lifts to cancel its image under the earlier type-\(e\)
laws. The resulting \(b_0\in G_\nu\) still gives a unit value
for \(z_j\). For an integer \(m_0\) representing
\(\chi(\gamma)\) in \(A\), apply, successively for
\(0\le k<t_{\mathrm a}\), \(k\ne e\), the operations
\[
 b\ \longmapsto\ (\gamma b\gamma^{-1}) b^{-m_0^{\,k+1}} .
\]
Each operation stays in \(G_\nu\) and multiplies the value
of a law of type \(d_0\) by \(m_0^{d_0+1}-m_0^{k+1}\).
The resulting element \(u\) kills all earlier laws. Also
\[
 z_j(u)=z_j(b_0)m_0^{t_{\mathrm a}-1}
     \prod_{k\ne e}(m_0^e-m_0^k).
\]
Lift \(z_j(b_0)\) to a unit in \(\mathbf Z_\ell\) and
lift \(m_0\bmod\ell^\nu\) by \(\chi(\gamma)\). The resulting
lift of \(4z_j(u)\) has valuation at most
\((t_{\mathrm a}-1)h_0+v_\ell(4)\le H\).
Reducing the divisibility in \(\mathbf Z_\ell\) gives
\(4z_j(u)\mid\ell^H\) in the ring \(A\), including
when \(H\ge\nu\) and the right-hand side is zero in \(A\).

Set \(g=g_sg_s^c u(u^c)^{-1}\). Apply the
\hyperlink{proof.weighted-fit-packet-cofinality}{cofinality on the fixed denominator packet}
at the already chosen twist. On the off-level set where the
distinguished unit test and weak-packet comparisons hold,
choose \(\delta\in\ker\chi\) such that the data are cofinal at
\(\delta^{\ell^\nu}g\). At this
element each of the laws has value \(2z_i(u)\), by oddness
and additivity on \(G_\nu\). Each even type-\(e\) cocycle
unramified away from \(\ell\) has value \(2b(g_s)\), by
evenness and additivity on \(G_\nu\). The factor
\(\delta^{\ell^\nu}\) contributes zero to all these values.
There are only finitely many such scalar \(A\)-valued
cocycles by detection. The values of all cocycles with
coefficients \(O/\ell^\nu O\) are then controlled in the
same neighborhood, by taking coordinates in a
\(\mathbf Z_\ell\)-basis of \(O\). Impose these finitely
many scalar conditions, the values of all current and
earlier laws, and membership in \(G_\nu\), which are open
conditions. Choose a degree-one prime \(w_j\) off \(\ell\)
and its Frobenius datum there. Exclude the conductor
primes of the already fixed twist, the primes above the real
primes already chosen or in the current and earlier
exceptional sets, and those ramified in \(E/E^+\). The
same datum gives \(\tau_j\) and \(\phi_j\); in particular
\(\operatorname{N} w_j\equiv1\bmod\ell^\nu\), by the
cyclotomic value \((\operatorname{N} w_j)^{-1}\)
of geometric Frobenius.
\footnote{The existence of degree-one Frobenii in nonempty open sets
off any finite set follows from
\href{https://doi.org/10.1007/BF01206606}{Chebotarev's theorem}.
For the use of the packet conditions see
\weightedpin{grh-compact/src/New/Selmer/DenominatorBound.lean}{98-L108}{the coefficient-place application}.
The compact prime choice is
\weightedpin{grh-compact/src/New/Selmer/TrianglePrimes.lean}{24-L123}{\texttt{exists\_isPrimeTriangle}}.}

For an even cocycle of type \(e\) vanishing at the earlier
\(\tau_i\) and at the earlier \(\phi_i\) of type \(e\),
the earlier conjugate pairs can be removed from its
ramification set. Its restriction to old tame inertia is
additive, its value at the conjugate generator is zero by
evenness, and wild inertia at these primes has residue
characteristic different from \(\ell\). Apply the law \(z_j\)
to the current and later primes. All the later
\(z_j(\phi_i)\) vanish by construction, and
\(z_j(\phi_j)=2z_j(u)\), so \(\ell^H b(\tau_j)=0\).
Here the divisor in reciprocity is \(2z_j(\phi_j)=4z_j(u)\),
whose valuation was charged above.

For the remaining coboundary assertion, suppose \(b\) is
even of type \(e>0\), unramified off \(\ell\), and zero at
every \(\phi_j\) of type \(e\). Then \(2b(g_s)=0\) for
each test point \(s\). Put \(d_e(g)=\chi(g)^{-e}-1\).
The identity \((as)g_s=sa\) and the cocycle law give
\[
 d_e(a)(2b)(s)-d_e(s)(2b)(a)
       =-\chi(as)^{-e}(2b)(g_s)=0 .
\]
The expression \(d_e(a)(2b)(g)-d_e(g)(2b)(a)\) is an
unramified-off-\(\ell\) cocycle, so detection makes it zero
on \(G_E\). The argument works integrally at \(2\),
and is paired with
\(4d_e(a)^2\mid\ell^H\), already proved. Finally
\[
 (\#J+\#S_{\rm det})H=t_{\mathrm a}\#S_{\rm det}\,y h_0
     \le y^{22}h_0 .
\]
For \(h_0=0\) the cost is zero at every prime with a finite
detecting set. The ordered types and their count are
independent of \(\nu\). These are the triangular inputs to
the Tate-block argument on the current packet at the actual
depth. Each prime and its inverse-conjugate reading use the
same variable and the already chosen fitted character.
\footnote{\weightedpin{grh-compact/src/New/Selmer/Corepresentable.lean}{134-L162}{Exact compact triangle interface}
and \weightedpin{grh-compact/src/New/Selmer/TateTwistSteps.lean}{15-L60}{Tate-block use}.
The readable edition uses detecting points in \(\ker\chi\)
together with a generator modulo that kernel. Its
\weightedpin{grh/src/New/Congruences/FirstOrder/TrianglePrimes.lean}{624-L706}{prime construction}
has the same role.}

In the Tate block put \(h_*=H=h_2h_0\) and use \(a=\gamma^{-1}\)
as above.
For a Tate character \(\xi=\chi^{-e}\), the paired differences
at \(\tau,\phi\) read \(c+c\circ\operatorname{conj}\), since
\(\xi(\tau)=\xi(\phi)=1\) modulo \(\ell^\nu\) and the conjugate
arguments are inverses. This even cocycle is subject to the
triangular relation. On the module left after all the prescribed
tame zeros and Frobenius equalities, the odd map
\(c\mapsto c\circ\operatorname{conj}-c\) has kernel killed by
\(C=\ell^{h_*}\) when \(e>0\). In its kernel \(c\) is even.
The triangle applied to \(c+c\circ\operatorname{conj}\) says
\(4(\xi(a)-1)c\) is a coboundary; the Ribet normalizer's
coboundary condition kills it on multiplying by \(\xi(a)-1\).
The factor \(4(\xi(a)-1)^2\) divides \(C\). This works at \(2\)
without a parity projector. The strict local condition away
from the ramified places follows from the local crystalline
relations: at \(\ell\), crystalline cocycles in \(O(-e)\)
vanish on the cyclotomic kernel for \(e>0\), since their rational
classes are coboundaries. For \(e=0\) the finite condition is
unramified and the final everywhere-unramified group is
bounded by the class number. Therefore the quotient used in
the Tate-block cost is bounded by the local defect and the
already established odd or class-number bound. The triangular
prime construction uses the primitive reciprocity law on the
same field and Frobenius conditions as above.
\footnote{\weightedpin{grh-compact/src/New/Selmer/Corepresentable.lean}{134-L162}{Even cocycles, triangle and coboundaries},
\weightedpin{grh-compact/src/New/Selmer/TrianglePrimes.lean}{18-L30}{Prime choice at arbitrary depth},
and \weightedpin{grh-compact/src/New/Selmer/TateTwistSteps.lean}{15-L60}{Entry into the Tate-block bound}.}

The trace ideal over the integral Hecke algebra adjoins one
variable \(X_i\) per pair. Its generators include
\(\ell^\nu\operatorname{tr}N_x(g)\) and
\[
 \operatorname{tr}((N_x(\tau_i)-1)N_x(g))-
 X_i\operatorname{tr}(((X-1)e_i)(N_x(\phi_i))N_x(g)).
\]
The inverse-conjugate reading uses the same variable. Indeed
write \(\bar h=chc\). Polarization gives
\[
 \operatorname{tr}N_x((\bar h^{-1})^a g)=
 \chi(g)^{1-t'}\chi(h)^{a(t'-1)}
 \operatorname{tr}N_x\bigl(h^a(cgc)^{-1}\bigr).
\]
The factor \(\chi(h)^{a(t'-1)}\) is \(1\) modulo \(\ell^\nu\)
for \(h=\tau_i,\phi_i\).
Its error is thus in the trace ideal.

\paragraph{The first-order Ribet polynomial.}
\teamanchor{proof.weighted-fit-ribet-fitting}%
Fix the coefficient reading, auxiliary pair, tuple, fitted twist and
triangular primes of the preceding test. Write \(x_0\) for its
distinguished genuine packet, \(\mathcal X\subseteq Y\) for the genuine
packets, and \(Y\) for the finite catalogue at the deeper level.
Put \(P=\ell^\nu\), \(P_0=\ell^{M_\ell}\), \(C=\ell^{h_*}\), and
\(\Lambda=O[X_j:j\in J]\), where \(O\) is a common finite coefficient
DVR. We take the direct-sum lattice
\(U_x=\rho_x\oplus\bigoplus_{0\le e<t'}\xi_e\) on \(\mathcal X\), with
\(\rho_x=T_x\psi\) and \(\xi_e=\chi^{-e}\). On the full \(Y\) we use
the stable lattices of the normalized attachments. Enlarge \(O\) to
contain the coefficients of every \(e_j\) and the scalar data. The
normalized length of an \(O\)-module of finite length is its length
divided by \(v_O(\ell)\). Thus \(v_\ell(d)=v_O(d)/v_O(\ell)\) for
\(d\in O\).

Let \(S\) be the image of the integral Hecke algebra with the variables
adjoined, read in \(\prod_{x\in Y}\overline{\mathbf Q}_\ell[X_j:j\in J]\).
Its scalars contain the diagonal \(O\), and the normalized
characteristic-polynomial coefficients on every Galois element lie
in \(S\). Indeed the normalized spherical coefficients come from the
integral Hecke operators, and their Frobenius limits lie in the closed
integral Hecke algebra. Denote the trace ideal of the preceding
paragraph by \(\mathfrak I\subset S\). It contains the two copies of
each trace tie, with one \(X_j\) for both copies. Write
\(\tau_{j,+}=\tau_j\), \(\tau_{j,-}=(c\tau_jc^{-1})^{-1}\) and use the
same convention for \(\phi\). The \(e_j\) satisfy
\(e_j(\xi_e(\phi_{j,\pm}))\equiv1\bmod P\) for every \(e\), since
\(\xi_e(\phi_{j,\pm})\equiv1\bmod P\) and \(e_j(1)=1\).
\footnote{\weightedpin
{grh-compact/src/New/Patching/OnePrimeTie.lean}{63-L89}{Hecke coefficients and traces},
\weightedpin
{grh-compact/src/New/Selmer/RibetElement.lean}{52-L70}{The paired traces with a single variable},
and \weightedpin
{grh-compact/src/New/Selmer/RibetFirstOrder.lean}{65-L106}{Trace ideal and required coefficient}.}

\emph{A module of tied rows.}
Write \(A=\Lambda[G_E]\), with \([g]\) the basis element for \(g\),
and let \(U_x\), \(\rho_x\) and \(\eta_e\) denote the induced
representations of \(A\), the last given by
\(\eta_e([g])=\xi_e(g)\). Choose a finite list of group elements
\(g_i\) as follows. The first part consists of the two tame readings
\(\delta_v\) at each twist place and old level place away from \(\ell\).
The second part spans the tuples \((U_x(g)-I)_{x\in Y}\) over
\(O\), which is possible since this span is a submodule of a finite
free \(O\)-module. The last part is a set \(s\) formed from the
detecting set \(S_{\mathrm{det}}\) and all \(\tau_{j,\pm}\). An element
occurring in more than one part has a separate coordinate in each.
Choose constant coordinate vectors \(z(g)\) with
\(\sum_i z(g)_i(U_x(g_i)-I)=U_x(g)-I\) for every \(x\).
For a row \(w\), put \(V(w)=\sum_i w_i([g_i]-1)\).
Write \(\gamma_0,y_0\) for the scalar data of the short Selmer
condition on \(x_0\), and set
\[
 \omega=z(\gamma_0a)-y_0 z(a)+P_0C\sum_g\lambda_gz(g),
 \qquad B=V(\omega).
\]
Here \(a\) and \(\lambda_g\) are the separator and the correction
already chosen in the normalizer. Thus \(b_x=\det U_x(B)\) uses
the base-level normalizer evaluated on the deeper catalogue.

Let \(\mathfrak T\) consist of the elements \(b\in A\) for which
\((\operatorname{tr}U_x(b[g]))_x\in\mathfrak I\) for every \(g\).
Multiplication in \(A\), linearity and cyclicity of trace show that
\(\mathfrak T\) is a two-sided ideal. Define the left row module
\[
 \mathcal N=V^{-1}\bigl(\mathfrak T+(\ker\eta_0)B\bigr).
\]
It contains \(P A^{\{i\}}\), the rows
\[
 \begin{gathered}
 z(gh)-z(g)-[g]z(h),\qquad z(g_i)-\mathbf e_i,\qquad
 ([g]-1)\omega,\\
 z(\tau_{j,\pm})-X_j e_j([\phi_{j,\pm}])z(\phi_{j,\pm}),
 \end{gathered}
\]
and every constant row representing an \(O\)-linear relation among
the tuples \(U_x(g)-I\). These assertions follow by applying \(V\)
and then the \(U_x\). For the last displayed row the trace is exactly
the corresponding generator of \(\mathfrak I\).

For each old tame coordinate \(\delta_v=g_i\) the module also
contains \(o_v([\delta_v])\mathbf e_i\), where
\[
 o_v(Z)=
 \begin{cases}
   \psi(\delta_v)-Z,&\text{at a fitted-twist place},\\
   1+Z+\cdots+Z^{G-1},&\text{at an old level place}.
 \end{cases}
\]
The quadratic relation at the twist places gives
\(o_v(U_x(\delta_v))(U_x(\delta_v)-I)=0\).
At an old level place off \(\ell\), the local attachment has finite
inertial image. In characteristic zero its matrices on inertia are
semisimple. The labelled inertia annihilator restricts their roots
to orders dividing some \(g_i^d-1\) with \(1\le d\le r\).
Since these integers divide \(G\), we have \(U_x(\delta_v)^G=I\).
These assertions apply on all of \(Y\), including the nonlift
packets. All \(\xi_e(\delta_v)\) equal \(1\).
\footnote{\weightedpin
{grh-compact/src/New/Selmer/RibetFirstOrder.lean}{115-L148}{Tied rows and normalizer},
\weightedpin
{grh-compact/src/New/Selmer/RibetElement.lean}{192-L204}{Old tame relations on the whole catalogue},
\weightedpin
{grh-compact/src/New/Selmer/LabelledColumns.lean}{36-L62}{Old level tame annihilators},
\weightedpin
{grh-compact/src/Lib/PadicHodge/Crystalline/GoodCriterion.lean}{38-L43}{Finite inertia and annihilating roots},
and \weightedpin
{grh-compact/src/New/Selmer/LabelledColumns.lean}{128-L151}{Quadratic relation}.}

\emph{The main-block determinant.}
Evaluate \(\mathcal N\) at all \(X_j=0\) through \(\rho_{x_0}\).
Let \(L\subseteq O^{r'|\{i\}|}\) be the span of the ordinary rows
of the resulting block rows. A vector \((w_i)\) modulo \(P_0\)
annihilating \(L\) gives
\(u(g)=\sum_i z(g)_i w_i\), a cocycle for \(\rho_{x_0}\) which
obeys the tuple relations and vanishes on the two readings of each
\(\tau_j\). The row \(([\gamma_0]-1)\omega\) forces
\((y_0-1)\sum_i\omega_iw_i=0\). Since \(y_0-1\) is a unit and
\(\rho_{x_0}(\gamma_0)=y_0 I\), the cocycle rule and \(P_0=0\) on
the coefficient module give
\(\sum_i\omega_iw_i=u(\gamma_0)\). Therefore \(u(\gamma_0)=0\).
The generator relations give \(u(g_i)=w_i\), so this map of
annihilators into cocycles is injective. These cocycles are
continuous because they factor by an \(O\)-linear map on the
finitely generated span of the tuples, which has its lattice
topology by Smith normal form. The tuple relations force
\(u(g)=0\) whenever every \(U_x(g)=I\), giving unramifiedness
outside \(\ell\), the old level, the fitted-twist places and
the triangular pairs. The two tame zeros remove the
triangular pairs from this list. The Artin--Rees graph
construction in the
\hyperlink{proof.weighted-fit-crystalline-duality}{local reduction}
applies to these cocycles and to the scalar cocycles below.

Let \(L_0\) denote the normalized crystalline defect on this
block, at \(P_0\), and \(L_e\) the corresponding defects on the
Tate blocks at \(P\). The preceding local reduction gives
\(L_0+\sum_e L_e\le4fr'\mathrm{loss}_{\mathrm{crys}}(\ell)\).
The short Selmer condition and this reduction bound the length
over \(O\) of the above annihilator by
\[
 b_{\mathrm{main}}
    =v_O(\ell)\bigl((10fr'+1)c_\ell+L_0\bigr)
    < v_O(P_0).
\]
The strict inequality follows from the definition of \(b_\ell\)
and \(M_\ell>r b_\ell\). Smith normal form says that the length of
this annihilator equals the sum of the exponents of \(L\) truncated
at \(v_O(P_0)\). Thus none is truncated and \(L\) has full rank
and length of its cokernel at most \(b_{\mathrm{main}}\).

Full matrix span means \(\rho_{x_0}(\Lambda[G_E])|_{X=0}
=M_{r'}(O)\). Left multiplication by preimages of matrix units
permits placing any ordinary row of \(L\) in any one position of
an evaluated block row, with the other positions zero. Group a
basis of \(L\) into \(r'\)-row blocks and lift them in this way
to rows of a square matrix \(M_0\) over \(A\), all in \(\mathcal N\).
For \(x\in\mathcal X\) write \(D_x=\det\rho_x(M_0)\), the determinant
of the flattened main blocks over \(\Lambda\). Then
\(v_O(D_{x_0}(0))\le b_{\mathrm{main}}\).
\footnote{\weightedpin
{grh-compact/src/New/Selmer/AptTupleDefs.lean}{30-L38}{The short Selmer condition},
\weightedpin
{grh-compact/src/Lib/General/Algebra/BlockFittingBound.lean}{11-L30}{The block determinant from the short annihilator},
and \weightedpin
{grh-compact/src/New/Selmer/RibetFirstOrder.lean}{149-L169}{The main block in the tied module}.}

\emph{The scalar rows and their relations.}
Apply \(\eta_e\) to the rows, discarding the first, old-tame
coordinates. Any \(O\)-linear form on the constant coordinate
space with values in \(O/P\) annihilating an indicated projected
relation module gives a scalar cocycle \(u(g)\) by evaluation
at \(z(g)\). Here and below annihilation means that the form
kills constant vectors whose images belong to the module.
The displayed rows supply the cocycle and tuple relations.
The old-tame coordinates have been discarded, so the generator
relations give \(u(\delta_v)=0\). The cocycles are therefore
continuous and unramified away from \(\ell\) and the new pairs.

If one of these cocycles is a coboundary
\(u(g)=(\xi_e(g)-1)m\), the row \(([g]-1)\omega\) gives
\((\xi_e(g)-1)u(\omega)=0\), where \(u(\omega)\) is the form
applied to \(\omega\). For \(e>0\), the triangular separation
gives \(\xi_e(a)-1\mid C\). Since \(\xi_e(\gamma_0)=1\),
\[
 u(\omega)=
  \left(1-y_0+
   P_0\,\frac{C}{\xi_e(a)-1}\sum_g\lambda_g(\xi_e(g)-1)\right)u(a).
\]
The scalar in parentheses is a unit. It follows that
\((\xi_e(a)-1)u(g)=0\) for every \(g\). For \(e=0\) every
coboundary is zero. These are the hypotheses of the Tate-block
length argument. On every projected module used below, that
argument gives the triangular cost at type \(e\) with constant
\(C\) and length bound \(b_e=v_O(\ell)(L_e+v_\ell(X))\).
In this expression \(X\) is the integer bounding the odd
cocycle modules in the budget calculation, including the
class-number bound for \(e=0\).

If all the last coordinates are also discarded, the cocycle
vanishes on \(s\), hence on both tame readings, so it is unramified
outside \(\ell\). The detecting condition now makes the cocycle
zero. Duality for \(O/P\) shows that the projection of
\(\eta_e(\mathcal N)\) to the middle coordinates contains each
standard basis row. Indeed a constant row outside this projection
has nonzero image detected by an \(O/P\)-valued form, by Smith
normal form applied to the preimage of the projection in the
constant coordinate lattice, which contains \(P\) times that lattice.
For a middle generator such a form would have a nonzero value
on the corresponding \(g_i\). Choose rows \(m_{e,i}\in\mathcal N\)
projecting to the middle unit rows under \(\eta_e\).
They can be made simultaneous. Starting with a zero row for each
\(i\), correct a row \(w\) at the next \(e\) by adding
\(\sum_j(\mathbf e_i(j)-w_j)m_{e,j}\), with \(j\) running over the
middle coordinates. At the already treated characters these
coefficients evaluate to zero. Call the resulting rows \(\nu_i\).

For each \(e\), use the rows \(\eta_e(\nu_i)\) to clear the
middle coordinates and keep the last coordinates \(s\).
Let \(N_e\subseteq\Lambda^s\) be the image of \(\mathcal N\)
under this linear operation after \(\eta_e\), and \(\Phi_e\)
its restriction to constant rows. The image of \(z(g)\) for
\(g\in s\) is the corresponding standard basis vector modulo
\(N_e\). Also
\[
  \Phi_e(z(\tau_{j,\pm}))-
           X_j\Phi_e(z(\phi_{j,\pm}))\in N_e.
\]
Here one can remove \(e_j\) because its value on
\(\xi_e(\phi_{j,\pm})\) is \(1\bmod P\) and \(P\Lambda^s\subseteq N_e\).
Every form on the constant lattice vanishing on
\(\Phi_e^{-1}(N_e)\) still gives one of the cocycles just examined.
\footnote{\weightedpin
{grh-compact/src/New/Selmer/RibetFirstOrder.lean}{170-L201}{Coboundaries and removal of generators},
and \weightedpin
{grh-compact/src/New/Selmer/TateTwistSteps.lean}{15-L60}{Tate-block costs on the tied cocycles}.}

\emph{Fitting determinants for the types.}
Let \(\mathfrak q=(X_j^2:j\in J)\),
\(I_e=\{j:\mathrm{ty}(j)=e\}\),
\(P_e=\prod_{j\in I_e}X_j\), and
\(P_{<e}=\prod_{u<e}P_u\).
Set \(N_e^+=N_e+(\mathfrak q:P_{<e})\Lambda^s\).
For \(j\in I_e\), any earlier \(i\notin I_e\) has
\(\mathrm{ty}(i)<e\) and \(X_i\Lambda^s\subseteq N_e^+\).
We claim there is \(\Gamma_e\in\Lambda\) such that
\[
 \Gamma_e P_e\ \in\operatorname{Fitt}_0(\Lambda^s/N_e^+),
 \qquad
 v_O(\Gamma_e(0))\le b_e+(|I_e|+|s|)v_O(C).
 \tag{*}
\]
Here \(\operatorname{Fitt}_0\) is generated by the full square
determinants of rows of the relation module.

For the proof, fix \(e\), abbreviate \(\Phi_e\) as \(\Phi\),
and write \(Z(H)=\{u(g)=\alpha(z(g)):\alpha\text{ vanishes on }
\Phi^{-1}(H)\}\) for a relation module \(H\supseteq N_e^+\).
We use two elementary consequences of Smith normal form and
alternating multilinearity. First, for an integral vector \(v\),
if all these forms kill \(v\), then \(\Phi(v)\in H\), since the
constant quotient is killed by \(P\) and its \(O/P\)-dual separates
points. Second, if \(a w\in H\) then \(a\) times the Fitting ideal
after adding \(w\) is contained in the Fitting ideal before adding
it. Expand determinants in the enlarged rows. Terms using the new
row twice vanish, and one scalar \(a\) deals with each remaining
term.

The strict-local-zero subset of \(G_E\) in the triangular cost
condition is denoted by \(\mathcal L_e\). Enlarge \(H=N_e^+\)
by the vectors \(\Phi(z(g))\), \(g\in\mathcal L_e\), by adding
one uniformizer-division step at a time. More precisely, when a
required vector \(\Phi(v)\) is outside \(H\), adjoin
\(\varpi^{a-1}\Phi(v)\), where \(a>0\) is least with
\(\varpi^a\Phi(v)\in H\) and \(\varpi\) is a uniformizer.
If \(v=z(g)\) for \(g\in\mathcal L_e\), this costs one in the
Fitting valuation and strictly reduces the image of \(Z(H)\)
modulo the zero conditions on \(\mathcal L_e\). The image
drops because some old form has
\(\varpi^{a-1}u(g)\ne0\) and every new form has that value zero.
The finite relative length \(\lambda_1\) forces this procedure
to terminate with all the prescribed vectors adjoined, and
bounds the total cost.

Next, in order for \(j\in I_e\) adjoin
\(v_j=\Phi(z(\phi_{j,+})-z(\phi_{j,-}))\).
Before adjoining \(v_j\), a form in \(Z(H)\) has equal values
on the two \(\tau_i\) for every \(i<j\). For earlier \(i\in I_e\)
the difference of their ties gives this because \(v_i\in H\).
For earlier \(i\notin I_e\) it follows from \(X_i\Lambda^s\subset H\).
The two Frobenius values also coincide at earlier \(i\in I_e\).
The triangular condition now makes every form kill
\(C\,(z(\tau_{j,+})-z(\tau_{j,-}))\).
Smith duality and the difference of the ties show
\(C X_jv_j\in H\), so adjoining \(v_j\) costs the Fitting
multiplier \(C X_j\).

Call the resulting module \(H'\), and temporarily add the ideal
of all the variables times \(\Lambda^s\), obtaining \(H'_0\).
The forms in \(Z(H'_0)\) vanish on both tame readings of every
pair and agree on the Frobenius readings of type \(e\).
Write \(\lambda_2\) for the length of \(Z(H'_0)\) modulo its
\(C\)-torsion. The triangular cost is \(\lambda_1+\lambda_2\le b_e\).
Adjoin \(C\Phi(z(g))\) for \(g\in s\) to \(H'_0\), using the
same uniformizer-division argument. This time the nonzero test
value is \(\varpi^{a-1} C u(g)\), which vanishes on the
\(C\)-torsion, so the cost is at most \(\lambda_2\).
The enlarged module contains \(C\) times every standard basis
row, yielding a determinant \(C^{|s|}\). After multiplying
by the uniformizer to cover this last cost it lies in
\(\operatorname{Fitt}_0(\Lambda^s/H'_0)\). Modulo the ideal of all
the variables, this Fitting ideal is contained in the image of
\(\operatorname{Fitt}_0(\Lambda^s/H')\), by reducing the entries
of its generating determinants. Lift an element with that
constant coefficient and multiply by \(\prod_{j\in I_e} C X_j\)
and the first uniformizer cost. This proves \((*)\), also for
\(I_e=\varnothing\).
\footnote{\weightedpin
{grh-compact/src/New/Selmer/RibetFirstOrder.lean}{17-L55}{Fitting bound for one type},
\weightedpin
{grh-compact/src/Lib/General/Algebra/FittingDuality.lean}{213-L280}{Fitting multiplication and DVR duality},
and \weightedpin
{grh/src/New/Congruences/FirstOrder/FittingBlock.lean}{137-L210}{The three enlargements}.}

Expansion modulo \((\mathfrak q:P_{<e})\) gives
\[
  \Gamma_e P_e=d_e+\zeta_e,\qquad
  d_e\in\operatorname{Fitt}_0(\Lambda^s/N_e),\quad
  \zeta_e\in(\mathfrak q:P_{<e}).
\]
In the difference between \(\prod_e(\Gamma_eP_e)\) and
\(\prod_e d_e\), telescope with the original factors
\(\Gamma_uP_u\) to the left of the \(e\)-th difference and the
factors \(d_u\) to its right. Each summand has the factor
\(P_{<e}\zeta_e\) and thus belongs to \(\mathfrak q\). Consequently
\[
 \prod_e d_e\equiv \left(\prod_e\Gamma_e\right)\prod_j X_j
                    \pmod{\mathfrak q}.
 \tag{**}
\]

\emph{Combining the blocks in the trace ideal.}
Write each \(d_e=\sum_i c_{e,i}\det B_{e,i}\) with
\(c_{e,i}\in\Lambda\), the rows of \(B_{e,i}\) in \(N_e\).
Lift the rows to \(\mathcal N\). Choose an element \(q_e\in A\)
by multiplying
\(\prod_{u\ne e}([a]-\xi_u(a))\) by the scalar
\(C/\prod_{u\ne e}(\xi_e(a)-\xi_u(a))\). Then
\(\eta_u(q_e)=C\) if \(u=e\) and zero otherwise.
There is also \(k\in A\) acting by \(1\) on every Tate block and
by zero on every \(\rho_x\) for \(x\in\mathcal X\). For each
such \(x\), use its own scalar pair \(\gamma_x,y_x\) from the
short Selmer condition and take the product of
\((y_x-[\gamma_x])/(y_x-1)\). The order of these factors can be
fixed arbitrarily.

For a tuple \(\mathbf i=(i_e)_e\), form a matrix \(M_{\mathbf i}\)
with rows in \(\mathcal N\). Its first rows are the old-tame
polynomial rows, followed by the middle pivot rows \(\nu_i\).
Its last rows are the sums over \(e\) of \(q_e\) times the chosen
lifts of the rows of \(B_{e,i_e}\). Under \(\eta_e\), its
determinant is
\[
 A_*\,\det B_{e,i_e},\qquad
 A_*=\left(\prod_v o_v(1)\right)C^{|s|}.
\]
Indeed the first rows give the displayed old-tame diagonal, the
middle pivot is identity, and clearing that pivot gives the last
block \(C B_{e,i_e}\). Replace \(M_{\mathbf i}\) by
\(M_{\mathbf i}'=M_0+k(M_{\mathbf i}-M_0)\), with left
multiplication entry by entry. On a genuine packet its main
block then has determinant \(D_x\) and its Tate blocks have
the determinants just computed. The element
\(h_{\mathbf i}=k q_0\prod_e c_{e,i_e}\) acts on a genuine packet
only on the trivial Tate slot, there by \(C\prod_e c_{e,i_e}\).
In particular every element of \(\ker\eta_0\) annihilates
\(U_x(h_{\mathbf i})\) by left multiplication on that packet.

We give the trace calculation that takes these matrix
determinants back to the Hecke algebra. For any matrix over \(A\),
including a \(1\)-by-\(1\) matrix, its evaluated flattened
characteristic coefficients belong to \(S\). This follows from
the determinant identity for a sum of matrices. In formal power
series, write \(\det(1-u\sum A_i)\) as the product, over primitive
cyclic words \(w\) in the indices, of
\(\det(1-u^{|w|}A_w)\), where \(A_w\) is the ordered word in
the matrices. Taking formal logarithms expands the traces of all
words, grouped by powers of primitive cyclic words, proving the
identity in characteristic zero. There are finitely many factors
per coefficient. For a matrix over \(\Lambda[G_E]\), take the
summands to be scalar multiples of \(E_{ij}[g]\). A word is zero,
off diagonal with square zero, or diagonal in one block with
a scalar times a single group element there. Its characteristic
coefficients belong to \(S\) by the group-element assertion.
The product identity therefore proves the stated integrality.

Now let \(M\) have rows in \(\mathcal N\). Write
\(V(M_i)=r_i+\lambda_i B\) with \(r_i\in\mathfrak T\) and
\(\lambda_i\in\ker\eta_0\). Work also over \(S[G_E]\), with
the central \(S\)-action on the evaluated matrices. Let \(D\)
be the evaluated flattened determinant of \(M\), and \(W\) a
matrix over \(S[G_E]\) whose evaluated flattening is the adjugate
of that matrix. Such a lift is a polynomial in \(M\) by
Cayley--Hamilton, with the characteristic coefficients in \(S\).
Multiplying \(V(M_i)\) by \(W\) and then by the row \(\omega\)
gives, after evaluation in the packet family,
\[
 (D-\omega W\lambda) B=\omega W r .
\]
For \(h\in A\) put
\(f_{M,h}=(\operatorname{tr}U_x((D-\omega W\lambda)h))_x\).
It belongs to \(S\) by the characteristic coefficient assertion
over \(S[G_E]\), by the same word argument. Multiply the last
identity by the adjugate of \(U_x(B)\) and then by \(U_x(h)\);
taking traces gives \(b f_{M,h}\in\mathfrak I\).
Indeed the adjugate of \(U_x(B)\) also lifts by Cayley--Hamilton,
and the trace of each term containing an \(r_i\) belongs to
\(\mathfrak I\), by cyclicity and the trace ideal property.
On a genuine packet on which \(U_x(\lambda_i h)=0\) for each \(i\),
this formula gives
\(f_{M,h}(x)=\det U_x(M)\operatorname{tr}U_x(h)\).
\footnote{\weightedpin
{grh-compact/src/New/Selmer/RibetCramer.lean}{24-L74}{Characteristic coefficients and two adjugates},
\weightedpin
{grh-compact/src/New/Selmer/RibetFirstOrder.lean}{205-L245}{Combining the blocks and taking coefficients},
and \weightedpin
{grh/src/New/Congruences/FirstOrder/RibetMethod.lean}{208-L274}{The trace multiplier}.}

Take \(f=\sum_{\mathbf i} f_{M_{\mathbf i}',h_{\mathbf i}}\).
Then \(bf\in\mathfrak I\), and \(b\in S\) with \(b_x=0\) on
\(\mathcal X\) because \(B\) vanishes on the trivial Tate slot.
On every genuine packet, the sum of products is
\[
 f(x)=C A_*^{t'}D_x\prod_e d_e .
\]
By \((**)\), every proper squarefree coefficient of \(f(x)\)
below \(\prod_jX_j\) vanishes. Its full squarefree coefficient at
\(x_0\) is
\[
 d=C A_*^{t'}D_{x_0}(0)\prod_e \Gamma_e(0),\qquad
 v_O(d)\le b_{\mathrm{main}}+\sum_e b_e
   +t'\sum_v v_O(o_v(1))
   +(2t'|s|+|J|+1)v_O(C).
\]
The constant coefficient of \(D_{x_0}\) is the only contribution
from that determinant to the full squarefree monomial.

To compare with the prescribed allowance, at the two readings of
the twist places the normalized sum of \(v_O(o_v(1))\) is at
most \(2c_\ell\), since \(v_\ell(u^{-1}-1)=v_\ell(u-1)\) for a
unit and the twist comes with this ramification bound. At the two
readings of old level places it is at most \(8f\,v_\ell(G)\).
There are at most \(4f\) such places before taking the two readings.
The triangular estimate gives
\((|J|+|S_{\mathrm{det}}|)h_*\le h_2^{22}h_0\), and
\(|s|\le |S_{\mathrm{det}}|+2|J|\). Using
\(h_*=h_2h_0\) and \(h_0\le v_\ell(G)\le c_\ell\) therefore gives
\[
 \begin{split}
 v_\ell(G^3d)\ \le\ 
 &4fr'\mathrm{loss}_{\mathrm{crys}}(\ell)+t'v_\ell(X)\\
 &+\bigl(10fr'+4+t'(2+8f)+(4t'+1)h_2^{22}+h_2\bigr)c_\ell .
 \end{split}
\]
Here \(r'\le r\), and \(f,t',r\le h_2\) with \(h_2\ge2\).
The coefficient of \(c_\ell\) is at most
\(18h_2^2+3h_2+4+5h_2^{23}\le h_2^{27}\), and
\(t'\le h_2^{27}\). This proves \(v_\ell(G^3d)\le r'b_\ell\).
The count uses the cost inequality multiplied by \(h_*\) and
applies also when \(h_0=0\). Both trace relations use the selected
coefficient reading and the same fitted twist. The two Galois places
in each conjugate pair share one variable. The determinant estimate
uses the multiplier at the selected reading. The conjugate coefficient
reading enters through the local test forms and integral row recombinations.
\footnote{\weightedpin
{grh-compact/src/New/Selmer/RibetElement.lean}{122-L147}{The coefficient ring and its determinant},
\weightedpin
{grh-compact/src/New/Selmer/RibetElement.lean}{204-L223}{Normalized length, tame factors and the budget},
and \weightedpin
{grh-compact/src/New/Selmer/DenominatorBound.lean}{120-L138}{Use on the distinguished genuine packet}.}

\paragraph{Local operators for the two first-order ties.}
\teamanchor{proof.weighted-fit-integral-ties}%
Work at one selected auxiliary place \(w=w_j\), with residue size
\(q\), at the same coefficient reading and depth as the Fitting
polynomial. Put \(m_1=\ell^\nu\), \(n=r'+t'\), and let
\(\Gamma=\mathrm{GL}_n(k_w)\) act by the residue action \(\rho\)
on the deeper forms. Write \(N_w=\sum_{g\in\Gamma}\rho(g)\)
and \(N\) for the product of the norms at all the selected
auxiliary places. The split local action, written \(\pi\), agrees
with \(\rho\) on the integral lifts of residue matrices. Let
\(H\subseteq\Gamma\) stabilize the first coordinate line.
Its inverse image in \(\mathrm{GL}_n(\mathcal O_w)\) is the
parahoric \(P_H\). The pro-\(p\) Iwahori group \(I_w\)
reduces to the upper unitriangular group. The chosen Artin data give
a uniformizer \(\varpi\) and a unit \(u\) with residue generating
\(k_w^\times\), representing the specified geometric Frobenius
\(\phi\) and tame inertia element \(\tau\), respectively.
Write \(R_x\) for the attachment before normalization at this step,
so \(N_x(h)=a_{\mathrm{nor}}(h)R_x(h)\). The normalizing avatar
is unramified at this auxiliary place and has unit values at
the coefficient place. Thus \(a_{\mathrm{nor}}(\tau)=1\), and
the complex Frobenius factor
\(\alpha_w=j_\iota^{-1}(a_{\mathrm{nor}}(\phi))\) is \(A_w^{-1}\)
in the preceding spherical test.

Use \(t_w\), \(P_I\) and \(U\) from the
\hyperlink{proof.weighted-fit-string-orthogonality}{local Hecke calculation}.
For every character of \(P_H\) through its first diagonal residue,
\(U\) preserves its eigenspace. Indeed for
\(g=\bigl(\begin{smallmatrix}a&b\\ \varpi c&D\end{smallmatrix}\bigr)
\in P_H\), the upper row operation with row \(-bD^{-1}\) is in
\(I_w\) and makes the upper-right block of \(g\) zero. Call the
result \(g'\). Then \(t_w^{-1}g't_w\in P_H\) has the same first
diagonal residue as \(g\). On a vector \(y\) of character \(\eta\)
this gives \(P_I\pi(gt_w)y=\eta(g)P_I\pi(t_w)y\).
Average this equality against \(\eta(g)^{-1}\) over \(P_H\).
The projector so obtained has image fixed by \(I_w\), since
\(\eta\) is trivial there, and therefore
\(P_I\pi(t_w)y\) belongs to the \(\eta\)-space.
\footnote{\weightedpin
{grh-compact/src/Lib/LocalGL/Iwahori/HeckeCounts.lean}{64-L83}{Local coset count},
\weightedpin
{grh-compact/src/Lib/LocalGL/Iwahori/HeckeOnLineStabiliser.lean}{15-L51}{Preservation of the line-parahoric character},
and \weightedpin
{grh-compact/src/New/Patching/OnePrimeTie.lean}{96-L127}{The actual local data and polynomial}.}

Apply the cofactor conclusion of the
\hyperlink{proof.weighted-fit-string-orthogonality}{local Hecke calculation}
to the string \(A_wq^i\), \(0\le i<t'\), on each spherical
base packet on which \(T=Ap(b)\) is nonzero. By the
\hyperlink{proof.weighted-fit-hecke-separator}{base separator}
these packets are weak and residually congruent to \(x_0\).
For \(0\le i<t'-1\) the root \(c=A_wq^i\) has \(P_x(qc)=0\),
so the cofactor obtained by deleting one occurrence of \(c\)
kills \(x\) under evaluation at \(U\). Their greatest common
divisor is
\[
 S_x(Z)=P_x(Z)\Big/\prod_{0\le i<t'-1}(Z-A_wq^i).
\]
It kills the entire \(U\)-cyclic space of \(x\). For \(t'=1\)
the empty product gives \(S_x=P_x\), which kills this space
by the forward recurrence. Adjoint invariance makes \(x\)
orthogonal to the image of \(S_x(U)\) on all \(I_w\)-fixed vectors.

Recall that \(\alpha_w=A_w^{-1}\).
The polynomial \(S_x(A_w Z)\) is a nonzero scalar times
\((Z-q^{t'-1})j_\iota^{-1}(T_{x,w}(Z))\), with \(T_{x,w}\)
the quotient by the full string in the residual Hecke test.
Since \(T_{x,w}\) divides \(e_j\) at the selected reading, put
\[
 Q_j(Z)=(Z-q^{t'-1})e_j(Z).
\]
Here and below the coefficients act on complex forms through
\(j_\iota^{-1}\). The divisibility just proved gives
\(Q_j(\alpha_w U)\) as a polynomial multiple of \(S_x(U)\).
Thus \(x\) is orthogonal to \(Q_j(\alpha_w U)y\)
for every \(I_w\)-fixed \(y\).
\footnote{\weightedpin
{grh/src/Lib/LocalGL/SphericalChainLemma.lean}{352-L518}{The root and cofactor consequences}.}

\emph{Integral weighted sums and their norm.}
Let \(\lambda:k_w^\times\to\mathbf Z/m_1\) be the logarithm
with \(\lambda(\bar u)=1\), which exists since \(m_1\mid q-1\).
Define \(b_H(h)=\lambda(h_{11})\) on \(H\).
Its transfer to \(\Gamma\) is \(b^\Gamma=\lambda\circ\det\).
Indeed for \(g=\operatorname{diag}(z,1,\ldots,1)\) the
transfer is computed on the cycles of the action on lines.
The first line contributes \(z\) before applying \(\lambda\).
Lines in the complementary hyperplane contribute \(1\), and
all other lines have orbit length \(\operatorname{ord}(z)\)
and also contribute \(1\), since \(g^{\operatorname{ord}(z)}=I\).
The transfer is multiplicative before taking the logarithm,
by permuting the lines in the transfer product. Its logarithm
kills each elementary transvection because the residue
characteristic is prime to \(m_1\). Diagonals of the displayed
form, their conjugates and elementary transvections generate
\(\Gamma\), proving the assertion. Choose
\(\gamma_\Gamma=\operatorname{diag}(\bar u,1,\ldots,1)\), so
\(b^\Gamma(\gamma_\Gamma)=1\).

For an integer-valued weight \(f\) on \(\mathbf Z/m_1\) define
\[
 J_f=\sum_{h\in H} f(b_H(h))\rho(h),\qquad
 D_f=\sum_{g\in\Gamma} f(b^\Gamma(g))\rho(g),\qquad
 B_f=Q_j(\alpha_w U)J_f .
\]
The range of \(J_f\) is \(I_w\)-fixed, since the pro-\(p\)
Iwahori residues belong to \(\ker b_H\). The coset sum \(U\)
preserves integrality on this range, since each translation is
at \(w,\bar w\) away from \(\ell\) and preserves the coefficient
gauge at \(\ell\). Its output is \(I_w\)-fixed and hence
fixed by the principal subgroup at \(w\). The old-level
conditions and the conditions at the other auxiliary places
are preserved by commutation. Thus \(J_f,D_f,B_f\) are
integral operators on the actual deeper space, supported
locally at \(w\). They commute with the diagonal Hecke
operators by preservation of the packets. In particular reducing
their integral weight functions modulo \(m_1\) gives valid
operator congruences. The operator \(B_1\) for weight \(1\)
has \(H\)-fixed image by the parahoric invariance of \(U\).

We use the following exact norm identity for any subgroup
\(\Delta\) of the full product of auxiliary residue groups.
On the full principal-level class set, with maximal compact
at the old places, choose a point in each \(\Delta\)-orbit.
Use the multiplier \(G/|\operatorname{Stab}_\Delta|\) on
that point and zero on the other points of its orbit.
Each stabilizer is a subgroup of the reduction of a class
automorphism group at maximal compact level, so its order
divides \(G\). Multiplication by this integral scalar function
preserves the coefficient lattice. It commutes with the
old compact translations and therefore preserves the specified
old-type subspace of the deeper forms. Denote the resulting
operator by \(\Psi_\Delta\). Each chosen orbit point is counted
exactly \(|\operatorname{Stab}_\Delta|\) times in the conjugate
sum, giving
\[
 \sum_{g\in\Delta}\rho(g)\Psi_\Delta\rho(g)^{-1}=G I .
\]
Use \(\Psi_H\) here and take the full-product operator as
\(\Psi\) in the extraction.
\footnote{\weightedpin
{grh-compact/src/Lib/Adelic/CuspForms/BaseLevel.lean}{140-L179}{Old-type and integral form conditions},
\weightedpin
{grh-compact/src/Lib/LocalGL/TaylorWiles/AuxPlaceOperators.lean}{31-L94}{Residue translations and integral polynomials in \(U\)},
\weightedpin
{grh-compact/src/Lib/LocalGL/TaylorWiles/AuxPlaceOperators.lean}{96-L134}{The class multiplier},
and \weightedpin
{grh-compact/src/New/Patching/OnePrime.lean}{11-L50}{Transfer for the line stabilizer}.}

Fix representatives \(c_\eta\) for the lines \(\Gamma/H\),
and write \(g c_\eta=c_{g\eta}h(g,\eta)\).
Put \(A_+=\sum_\eta\rho(c_\eta)\) and
\(A_-=\sum_\eta\rho(c_\eta^{-1})\).
For an integer lift \(v\) of the residue coordinate
\(x\mapsto x\) on \(\mathbf Z/m_1\) set
\[
 F=-A_+ B_v A_- .
\]
This is integral and place-supported. The exact right shift
\(B_f\rho(h)=B_{f(\cdot-b_H(h))}\) yields
\[
 F(\rho(g)-1)\equiv
 A_+ B_1\sum_\eta v(b_H(h(g,\eta)))\rho(c_\eta^{-1})\pmod{m_1}.
\]
Let \(V_g\) be the same last sum preceded by \(\Psi_H B_1\).
By \(H\)-invariance of \(B_1\) and the norm identity,
\(N_w\Psi_H B_1=G A_+ B_1\).
Thus \(G F(\rho(g)-1)\equiv N_w V_g\bmod{m_1}\).

To check the required vanishing, write \(L=\widehat{T}N\).
The norm has image in the base level. On a spherical base
eigenpacket on which \(T\ne0\), the string orthogonality
shows that the image of \(B_1\) is orthogonal to the base vectors.
The other base packets are killed by \(\widehat{T}\).
Using the invariant positive inner product and \(N^*=N\)
gives \(LB_1=0\).
Moreover \(L\rho(g)=L\) for every \(g\) at this auxiliary place,
so the exact equation for \(N_w\Psi_H B_1\) gives
\(|\Gamma|L\Psi_H B_1=0\) before reduction.
Cancel \(|\Gamma|\) on the complex form space to get \(LV_g=0\).
This is the first tie with its exact vanishing condition.
\footnote{\weightedpin
{grh-compact/src/New/Patching/OnePrime.lean}{77-L107}{The first-order norm shift},
and \weightedpin
{grh-compact/src/New/Patching/OnePrimeTie.lean}{128-L152}{Orthogonality on the actual weak spherical packets}.}

\emph{Frobenius on the residue character spaces.}
Here we use the
\hyperlink{proof.weighted-fit-tame-iwahori}{tame local--global compatibility}
for a sorted
pro-Iwahori exponent on the actual attachment. Its exact
input off \(\ell\) is finite inertia and the annihilating
polynomial \(\prod_i(Z-j_\iota(\theta_i(v)))\) at an
element of the decomposition group with Artin preimage \(v\)
and nonnegative geometric Frobenius degree. Apply the
\hyperlink{proof.weighted-fit-sort-exponents}{local exponent calculation}
to a nonzero vector with a diagonal unit character and scalar
action of the permutations preserving it. This gives an exponent
\((\theta_i)_i\) whose
functional is nonzero on the vector. Its ordered coordinates
have the original unit characters, and a simultaneous
permutation of its exponents gives a sorted eigencharacter
on a vector in the same cyclic representation. The
pro-Iwahori fixed part of each such cyclic representation is
finite dimensional here, with the other levels and the
weight fixed. The permutation condition holds on both
character spaces used below. For a determinant character it
holds for all permutations. For a nontrivial first-entry
character it holds for the permutations fixing the first
index. The exponent functional satisfies
\(\mathcal L(Uy)=\theta_1(\varpi)\mathcal L(y)\) on
\(I_w\)-fixed vectors. The factor \(q^{n-1}\) in \(U\) cancels
the first-coordinate power \(|\varpi|^{n-1}\) of the normalized
torus character.
\footnote{\weightedpin
{grh-compact/src/Lib/LocalGL/Iwahori/SortedExponents.lean}{19-L34}{Hypotheses for the local exponent},
\weightedpin
{grh-compact/src/Lib/LocalGL/Iwahori/SortedExponents.lean}{103-L130}{A sorted eigencharacter on the cyclic representation},
\weightedpin
{grh-compact/src/Lib/GaloisRep/LocalGlobal/EigencharCompatibility.lean}{31-L42}{Tame attachment for the sorted eigencharacter},
and \weightedpin
{grh-compact/src/Lib/GaloisRep/LocalGlobal/TypeFrobenius.lean}{22-L45}{Unsorted labels on the actual channel}.}

On a determinant character space \(\xi\circ b^\Gamma\) put
\(\zeta=j_\iota(\xi(1))\). All tame exponents at \(u\) have
that value, so \(N_x(\tau)=\zeta I\) by finite inertia and
the annihilating polynomial at \(\tau\).
On a first-entry space \(\xi\circ b_H\), suppose
\(\xi(n!)\ne1\).
Now \(R_x(\tau)\) is semisimple with possible eigenvalues
\(\zeta,1\), and \(R_x(\phi)\) commutes with it. The latter
assertion follows by applying the local annihilator with
unit preimage \(1\) to the commutator, whose tame exponents
all have value \(1\).
For an exponent \(\theta\) nonzero on a vector in this space,
put \(\alpha_i=j_\iota(\theta_i(\varpi))\) and
\(H_\theta(Z)=\prod_{i>1}(Z-\alpha_i)\).
On the \(\zeta\)-eigenspace of \(R_x(\tau)\), the same local
attachment at \(\phi\tau^b\), with preimage \(\varpi u^b\),
gives
\[
 \bigl(R_x(\phi)-\alpha_1 I\bigr)
               H_\theta(\zeta^b R_x(\phi))=0,\qquad 0\le b<n .
\]
These \(n\) powers of \(\zeta\) are distinct. Any positive
difference of their exponents is at most \(n-1\) and divides
\(n!\). Interpolation in \(\zeta^b\) and \(H_\theta(0)\ne0\)
give \(R_x(\phi)=\alpha_1 I\) on that eigenspace.

Let \(M\) be any matrix over the coefficient field.
For a vector \(y\) in the packet and its first-entry
space, this implies the operator identity
\[
 j_\iota^{-1}\!\bigl(\operatorname{tr}((R_x(\tau)-I)M)\bigr)\, U y
 =
 j_\iota^{-1}\!\bigl(\operatorname{tr}((R_x(\tau)-I)R_x(\phi)M)\bigr)y .
\]
Indeed a nonzero difference would lie in the same character
space by the preservation proved above. An exponent functional
nonzero on that difference would give zero by the scalar
calculation on the \(\zeta\)-eigenspace and by
\(\mathcal L(Uy)=\theta_1(\varpi)\mathcal L(y)\).
Testing the identity with \(R_x(\phi)^bM\) and iterating
substitutes \(\alpha_w U\) for \(N_x(\phi)\) in any polynomial
in this trace identity. These channel calculations apply to
the full catalogue at the deeper level.
\footnote{\weightedpin
{grh-compact/src/Lib/GaloisRep/LocalGlobal/TypeFrobenius.lean}{47-L82}{Frobenius identity from the exponent and Vandermonde}.}

\emph{The second tie.}
For the group argument \(g_0\) of a trace generator, write
\(\widehat r\) for the operator of
\(\operatorname{tr}((N_x(\tau)-I)N_x(g_0))\).
Write \(\widehat r'\) and \(\widehat s'\) for the operators of
\(\operatorname{tr}((N_x(\tau)-I)Q_j(N_x(\phi))N_x(g_0))\) and
\(\operatorname{tr}(Q_j(N_x(\phi))N_x(g_0))\), respectively,
and \(\widehat s\) for the trace operator with
\((Z-1)e_j(Z)\) in place of \(Q_j(Z)\).
These traces are integral Hecke operators by the coefficient
argument above. The difference
\(\widehat s'-\widehat s\) is divisible by \(m_1\) as an
integral operator, since \(q\equiv1\bmod{m_1}\) and the
difference of the polynomials is \((1-q^{t'-1})e_j(Z)\).

On the range of \(D_f\), decompose over \(\mathbb C\) into the
characters \(\xi\circ b^\Gamma\). On each such space
\(N_x(\tau)-I=(\zeta-1)I\), and \(\rho(\gamma_\Gamma)-1\)
acts by \(\xi(1)-1\). Hence exactly
\[
 D_f\widehat r'=(\rho(\gamma_\Gamma)-1)D_f\widehat s',
 \qquad
 N_w\widehat r=(\rho(\gamma_\Gamma)-1)N_w=0.
\]
For \(k_0=n!\), put \(\Delta_{k_0}f(x)=f(x-k_0)-f(x)\).
On a first-entry channel of character \(\xi\),
\(J_{\Delta_{k_0}f}\) is \((\xi(k_0)-1)J_f\).
The channels with \(\xi(k_0)=1\) thus have zero weight.
On the others the preceding Frobenius identity gives, exactly,
\[
 B_{\Delta_{k_0} f}\widehat r=J_{\Delta_{k_0} f}\widehat r'.
\]
The character decompositions verify identities between the
integral operators. Their construction used finite integer
weighted sums.

Choose \(f(x)=v(x)(v(x)+k_0)\), so
\(\Delta_{k_0}f(x)\equiv -2k_0v(x)\bmod{m_1}\).
The transfer formula gives
\(A_+J_v A_-\equiv D_v\bmod{m_1}\), as the coefficient of
\(\rho(g)\) in the left sum reduces to
\(\sum_\eta b_H(h(g,\eta))=b^\Gamma(g)\).
Also \((\rho(\gamma_\Gamma)-1)D_v\equiv-N_w\bmod{m_1}\)
by the weight shift. Combining the identities gives
\[
 2k_0F\widehat r
   \equiv -2k_0D_v\widehat r'
   =-2k_0(\rho(\gamma_\Gamma)-1)D_v\widehat s'
   \equiv 2k_0N_w\widehat s\pmod{m_1}.
\]
Thus the same \(F\) as in the first tie gives the second
with \(\kappa_0=2n!\). The operators \(F_i,N_i\) and the
residue translations commute across distinct auxiliary
places by their disjoint support. In the first tie \(V_g\)
contains the class multiplier and sits on the right of the
norm. These identities supply the hypotheses for the
squarefree coefficient extraction with the previously fixed
depth and the same scalar \(G\).
\footnote{\weightedpin
{grh-compact/src/New/Patching/OnePrime.lean}{108-L118}{Quadratic weight difference and transfer},
and \weightedpin
{grh-compact/src/New/Patching/OnePrimeTie.lean}{153-L181}{The two trace identities on the residue characters}.}

\paragraph{Extracting the diagonal operator from the first-order input.}
\teamanchor{proof.weighted-fit-first-order-extraction}%
Let \(\mathcal E\) be the ring of integral endomorphisms at
the deeper level. For each triangular prime \(i\) use the
operators \(N_i,F_i,V_{i,g}\) from the
\hyperlink{proof.weighted-fit-integral-ties}{local ties},
with \(N=\prod_iN_i\) and the full-product norm operator
\(\Psi\). For an integral diagonal Hecke function \(a\),
write \(\widehat a\) for its operator. The separator is
\(T=Ap(b)\).

Use the integral vector \(l_0=G l\) as input and work in
\(\mathcal E l_0/\ell^\nu\mathcal E l_0\).
For a squarefree multi-index \(d\) use \(\Theta_d\) with
factor \(F_i\) when \(d_i=0\), \(N_i\) when \(d_i=1\),
and zero for any exponent above \(1\). Thus
\(\Theta_{\mathbf 1}=N\). Put \(P=G^2T\) and \(k=|I|+1\),
where \(I\) is the set of triangular primes. The identities
imply
\begin{equation}
\label{eq:weighted-fit-tie-iteration}
 \widehat P^{\,m}\Theta_d z=0
 \quad\text{if}\quad |d|+m\ge k
                \ \text{and}\ N\widehat T z=0.
\end{equation}
For \(m=0\) this follows from the definition of \(\Theta_d\).
For the induction step consider the operator
\(W=\widehat P^{\,m}\Theta_d G\).
If \(d_i=1\), then \(W(\rho_i(g)-1)=0\) by the averaging
factor. If \(d_i=0\), the first tie makes this expression
\(\widehat P^{\,m}\Theta_{d+e_i}V_{i,g}\), which vanishes
on every vector by the induction hypothesis, since
\(\widehat T N V_{i,g}=0\). Hence \(W\rho(g)=W\).
The identity of \(\Psi\) yields \(WG=W\Psi N\).
Apply it to \(\widehat T z\) to obtain \eqref{eq:weighted-fit-tie-iteration} at \(m+1\).

The linear evaluation of a Hecke polynomial
\(\sum a_d X^d\) by \(\sum\Theta_d\widehat a_d\), followed
by multiplication by \(\kappa_0\), kills the trace ideal in
this quotient. For a monomial multiplied by a generator
\(r-X_i s\), its \(i\)-th factor is exactly the last
displayed tie or the zero \(\kappa_0N_i\widehat r\).
The generators with \(\ell^\nu\) already vanish there.
Take depth so that
\(\ell^\nu=w_0\kappa_0(G^2p(0))^k\) with \(w_0\) integral.
This is possible because \(p(0)\ne0\) and the length and
the number of triangular places have been fixed.
There is an integral Hecke function \(v\) with
\[
 P^k-(G^2p(0))^k A^k=bv .
\]
Evaluate \(vbf\) at \(z_0=\Psi l_0\) and multiply by
\(w_0\kappa_0\). The second term in the displayed formula
vanishes in the quotient. Each proper squarefree
coefficient \(f_d\) satisfies \(N\widehat T\widehat f_d z_0=0\):
in the original space this is \(G\widehat T\widehat f_d l_0\),
which vanishes by projection agreement and the lower coefficient
vanishing on \(w\). Thus \eqref{eq:weighted-fit-tie-iteration} kills its term. At the top one
obtains \(w_0\kappa_0 N\widehat{P^k f_{\mathbf 1}}z_0=0\)
in the quotient. Since \(N\Psi l_0=G l_0\) and \(P^k\) is
divisible by \(T\), the element in that last congruence
before reduction is
\(\ell^\nu G\widehat{A^k f_{\mathbf 1}}(G w)\).
Cancellation in the original vector space gives
\(G^2\widehat{A^k f_{\mathbf 1}}w\in\mathcal E l_0\),
which is integral. Its eigenvalue on \(x_0\) is
\(A(x_0)^k f_{\mathbf 1}(x_0)\) before the two factors \(G\),
and \(A(x_0)\) is a unit. The final pairing with the original
left test costs one further \(G\), from finite class stabilizers.
Thus the pairing is integral after \(G^3\), with
\(v_\ell(G^3 A(x_0)^k f_{\mathbf 1}(x_0))\le r'b_\ell\).
\footnote{\weightedpin{grh-compact/src/New/Patching/OnePrimeTie.lean}{185-L304}{Squarefree coefficient extraction at the prescribed depth},
and \weightedpin{grh-compact/src/New/Selmer/DenominatorBound.lean}{133-L166}{The factors from the tuple and class pairing}.}

\paragraph{From diagonal tests to the two determinants.}
For each required \(Z\), use the good idempotent on the full
range and the \(X-1\) tuple idempotent on a line. The diagonal
operators just constructed give the eigenvalue lower bound
\(e^{-r_Z\beta_{j,\ell}}\), including \(G^3\), and integral
pairings on the indicated tests. Here \(r_Z=r\) for the full
range and \(r_Z=1\) for a line. The left test is integral at
all embeddings above \(\ell\). The right test is integral at
\(j_\iota\) and orthogonal to
\(Z^\perp\cap\operatorname{im}(e)\cap\ker H_w(\xi_w)\)
for the indicated split places.

Here is the determinant implication for just the test vectors
in use. Use the resulting \(G^3\)-multiplied diagonal operators.
For a given eigenbasis coordinate, divide the corresponding
linear form in the combination coefficients by one of its
coefficients of largest \(\ell\)-adic norm. This gives a nonzero
linear form over the residue field. Since that field is infinite,
choose integral combination coefficients avoiding the zero sets of
these finitely many forms. Each diagonal entry of the combination
then has norm at least \(e^{-r_Z\beta_{j,\ell}}\). The matrix of
its pairings with the tests has integral entries. Its determinant
is the tested Gram determinant times the product of the diagonal
entries, by writing the projected right-hand tests in the common
eigenbasis. We obtain the bound
\(\exp(\dim Z\,r_Z\beta_{j,\ell})\) for the tested Gram determinant.

\teamanchor{proof.weighted-fit-table-frame}%
Fix a fitted \(\psi\) with \(\ell\nmid\operatorname N(Q_\psi)\).
On the small block \(V^\circ\) take \(D_-=D_{k^-}\),
\(B_-=B_{k^-}\) and \(k_j=h_{k^-}\) in its exact rank. The tables
are the direct sum of these Hermitian coordinate spaces, one for
each member index \(i\) and principal \(\mathfrak M^2\)-class \(c\),
with the unweighted sum metric. The target form metric is the
class sum of the target Gram form with weight
\(1/|\operatorname{Aut}|\). In contrast, the adjunction with the
good source uses the original Petersson measure on \([U(W_j)]\).
These are the metrics defining the two adjoints in the comparison.

For each small-block representative \(c\), choose a rational
unitary \(\delta_c\) with \(\delta_c c\) integral and invertible at
all places above \(\ell\). Weak approximation gives this by density
of the rational points in the finite product of unitary groups,
as one sees on the Cayley open set. The small-block gauge at \(c\)
of a vector \(v\) at a reading \(a\) is then at most one exactly
when the coordinates of \(\rho_-(\delta_c)v\) are \(a\)-integral.
Here \(\rho_-\) is the small-block rational weight action. Indeed
rational transport moves both the representative and the vector,
and integral invertible matrices act in both directions on the
small-block coordinate lattice. The same choice of \(\delta_c\)
works for every coefficient reading above \(\ell\).

For \(a_0=(i,c)\) and each coordinate vector \(e_b\), let
\(f_{a_0,b}\) be the table supported at \(a_0\) with value
\(\overline{\rho_-(\delta_c)^{-1} B_-e_b}\). This conjugation
fixes the integral matrix \(B_-\) and changes the transport to
the one required in the second reading. With the pairing linear
in the second entry, unitarity of \(\rho_-\) and the identity
\(D_-B_-=k_j I\) give
\[
 \langle f_{a_0,b},\zeta\rangle
     = k_j\,\overline{[\rho_-(\delta_c)\bar\zeta_{a_0}]_b}.
\]
The transport of \(\bar f_{a_0,b}\) is the integral vector
\(B_-e_b\). Consequently \(G_\psi f_{a_0,b}\) is integral at
every reading above \(\ell\), by the gauge bound for the divided
point forms. No inversion of \(k_j\) occurs in these tests.

On \(L_\psi=Z_{j,\vartheta}\) with \(\psi=\vartheta\) use the
linear functions \(\ell_{a_0,b}(z)=\langle f_{a_0,b},G_\psi^*z\rangle\).
They separate \(L_\psi\) because \(G_\psi^*|_{L_\psi}\) is
injective and \(D_-B_-=k_jI\). Among the determinants of \(R\)
of these functions on any basis, choose a nonzero determinant of
maximal absolute value at \(j_\iota\), and let \(z_i\) be the
dual basis. Each \(j_\iota(\ell_{a_0,b}(z_i))\) is integral by
Cramer's rule. In the formula for the linear functions this is
the condition that the conjugate table of \(k_jG_\psi^*z_i\)
has integral gauges at \(j_\iota\circ\operatorname{conj}\).
Thus \(k_jG_\psi G_\psi^*z_i\) is integral at \(j_\iota\).
It is tested on the false spaces by the marked-period calculation
and the common-adjoint identity. The first frame is the image by
\(G_\psi\) of the selected test tables. For a selected \(f_i\)
and each \(z_j\), its paired entry is
\[
 \big\langle G_\psi f_i, P_{L_\psi} k_jG_\psi G_\psi^* z_j\big\rangle
       =\ell_i(K_{j,\vartheta} z_j).
\]
The matrix is exactly the matrix of the compressed operator in
the dual basis. The fitted denominator estimate applies to these
two frames, with the first integral at every reading and the
second integral at the selected reading and tested. It yields the
bound \(\exp(Rr\beta_{j,\ell})\) for the determinant.%
\footnote{\weightedpin{grh-compact/src/Lib/Adelic/CuspForms/ClassReps.lean}
{21-L47}{Transport of the p-adic gauge};
\weightedpin{grh-compact/src/New/Rows/LiftVolumeHeight.lean}
{44-L71}{Integral Fischer frame and the fitted determinant test};
\weightedpin{grh-compact/src/New/Rows/ComparisonForms.lean}
{113-L126}{Application to the divided forms}.}

\paragraph{The paired rank-one bound on the same twist.}
\teamanchor{proof.weighted-rank-one-paired-test}%
For \(\psi=\vartheta\), put \(Z=\mathbf C F_\psi\) and let \(e_{\mathrm{triv}}\) be
the rank-one trivial-character tuple idempotent. The normalized
local point tests are initially orthogonal to
\(\operatorname{im}e_{\mathrm{triv}}\cap\ker H
\cap\ker\Theta\), where \(\Theta\) is the descent to \(U(1)\).
Write \(h_1=h'\) for the chosen-level \(U(1)\) class number.
Choose class representatives \(g_c\) integral above \(\ell\).
Let \(C^*\) be the sum of the character-corrected central
operators at \(g_c^{-1}\). Its summands are unitary, commute
with \(e_{\mathrm{triv}}\) and \(H\), and fix \(Z\). Applying
their adjoints to a vector
\(y\in\operatorname{im}e_{\mathrm{triv}}\cap\ker H\cap Z^\perp\)
and descending gives the \(U(1)\) class sum of \(\Theta y\).
Its value at any translate equals its value at \(1\). The
class-sum adjunction identifies this value with the pairing
with \(F_\psi\), up to its nonzero fixed scale, so it vanishes.
It follows that \(C^*\Xi'\) is orthogonal to the required
false-space intersection with \(Z^\perp\). Also it remains
integral because \(g_c\) is integral at \(\ell\) and the
character factors there are units. Since \(P_ZC^*=h_1P_Z\),
the fitted determinant on this line gives
\[
  \bigl|j_\iota(h_1\langle\Xi,P_Z\Xi'\rangle)\bigr|_\ell
      \le \exp(\beta_{j,\ell}).
\]
The two Fourier selections may give different test forms.
The projection and the central operators are those of the
same \(\psi\) in both readings.\footnote{\weightedpin
{grh-compact/src/New/Selmer/RankOneDenominator.lean}{43-L66}{The rank-one pairing with one class number}.}

\teamanchor{proof.weighted-volume}%
On the full row space \(W\), put
\(Z_\vartheta=\det(\Theta_\vartheta^*\Theta_\vartheta)\), using an
orthonormal basis for the fixed source metric. We compare this
determinant for the trivial twist and the twist at a specified
coefficient place. The
single-lift determinant at stage \(j\) is
\(\lambda_j(\psi)=\det(u_\psi^*u_\psi)\) on the good source space.

\paragraph{The two local tests in the determinant of a single lift.}
\teamanchor{proof.weighted-single-compression}%
At stage \(j\), write \(U_j\) for the good forms on \(U(W_j)\), of
dimension \(R\), and \(u_\psi:U_j\to V_\psi\) for its single lift at
the ordinary twist \(\psi\). Denote the target space by
\(V_\psi\) and the image by \(L_\psi\). Use the same finite
Weyl-table space for all \(\psi\) at this stage.

\paragraph{The comparison lattice at a selected place.}
\teamanchor{proof.weighted-comparison-glue}%
Fix \(j\) and a table index \(i\). In the prescribed splitting
\(V_j=V''\oplus V^\circ\), write \(N''=\dim V''\) and let \(E''\)
and \(J^\circ\) be the coordinate inclusions of the \(r\) marked
columns into the respective blocks. At a chosen split place \(v\),
write \(\mathcal O_v\) for its valuation ring and \(\varpi\) for
the local component of the generator of the chosen modulus
\(\mathfrak M\). Lift the \(3m\)-row matrix \(s_i\) from
\(\mathcal O_v/\varpi^2\) and form the stacked matrix
\(A_i=[s_i;\varpi I_{N''}]\). Its rows generate a full lattice
between \(\varpi\mathcal O_v^{N''}\) and \(\mathcal O_v^{N''}\).
Choose a basis matrix \(D_i\) for that lattice. Both \(D_i\) and
\(\varpi D_i^{-1}\) are integral. Let \(D_\circ\) be diagonal with
\(\varpi\) on the \(J^\circ\)-coordinates and \(1\) elsewhere.
The two integral matrices
\[
 \gamma_{i,v}=
 \begin{pmatrix}\varpi D_i^{-1}&E''(J^\circ)^{\mathsf t}\\
                  0&\varpi D_\circ^{-1}\end{pmatrix},
 \qquad
 G_{i,v}=
 \begin{pmatrix}D_i&-D_iE''(J^\circ)^{\mathsf t}\\0&D_\circ\end{pmatrix}
\]
have product \(\varpi I\). Use \(\gamma_{i,v}\) at \(v\), its
conjugate inverse transpose at \(\bar v\), and the identity off
the chosen pairs. This gives a finite adelic unitary element
\(\gamma_i\). The inverse at \(\bar v\) is integral. At \(v\) we
have \(\varpi\gamma_{i,v}^{-1}=G_{i,v}\). These choices use only the
untwisted modulus and the fixed splitting.

For \(X=[X'';X^\circ]\) modulo \(\varpi^2\), the condition
\(G_{i,v}X=\varpi[0;J^\circ]\) implies \(X^\circ=J^\circ U\) with
\(U\equiv I_r\bmod\varpi\). The unique \(Y\equiv I_r\bmod\varpi\)
with \(X^\circ Y=J^\circ\) is \(U^{-1}\). The first block of the
condition is then equivalent to \(D_iX''Y=D_iE''\), or
\(A_iX''Y=A_iE''\), since the two matrices generate the same row
lattice. Conversely these two block conditions with any such \(Y\)
give \(G_{i,v}X=\varpi[0;J^\circ]\). The indicator for the glued
lattice therefore equals the sum of the product indicators for
these \(Y\), with multiplicity one. Chinese remaindering gives
the identity modulo \(\mathfrak M^2\). The upper \(N''\) columns
of \(G_{i,v}\) also have the following property. A change of
basis in \({\rm GL}_{N''}(\mathcal O_v)\) puts the rows of
\(s_i\bmod\varpi\) into the designated \(3m\) coordinates, since
their rank is at most \(3m\). In that basis \(\varpi G_{i,v}\)
vanishes modulo \(\varpi^2\) on all other upper columns.
This gives the column bound on the very lattice used for the
indicator identity.\footnote{\weightedpin
{grh-compact/src/New/Rows/GluedFormTheta.lean}{193-L208}{Split placement of the inverse lattice},
and \weightedpin
{grh-compact/src/New/Rows/GluedFormTheta.lean}{224-L294}{The glued indicator and the marked columns}.}

\paragraph{One normalization of the large-block vectors.}
\teamanchor{proof.weighted-comparison-vectors}%
Fix \(j\). In the large block of rank \(r+9m+1\) use the
prescribed \(3m\)-row residue matrix \(s_i\) and the short
basis \(Y_i\) of the prescribed rank-\((6m+1)\) definite
class.
\teamanchor{proof.weighted-fischer-weyl}%
For the Weyl coordinates, use the integer standard tableau basis
\(b_i(X)\) of products of minors in the initial columns of each
matrix \(X_\tau\). For a column diagram \(\lambda\), write
\(f_\lambda\) for the product of the leading principal minors with
its prescribed multiplicities. Define the Gram matrix \(D_\lambda\)
by the coefficient expansion
\[
 f_\lambda(Y^{\mathsf t}X)
      =\sum_{i,k} (D_\lambda)_{ik} b_i(X)b_k(Y).
\]
Products and transposes in this formula are taken place by place.
Use the bilinear Fischer pairing
\([P,Q]=P(\partial)Q(0)\), with
\([X^\alpha,X^\beta]=1_{\alpha=\beta}\prod_a\alpha_a!\).
Its matrix \(B_\lambda=([b_i,b_k])_{ik}\) is integral. Put
\(h_\lambda=[f_\lambda,f_\lambda]>0\).

The Fischer pairing with \(f_\lambda\) is \(h_\lambda\) times
evaluation at the identity on this flag-polynomial space. Indeed,
only the row multidegree of \(f_\lambda\) contributes to its pairing
or to evaluation at the identity. If a product of basis minors has
this multidegree, the number of entries in the first \(s\) rows of
the ambient matrix is the sum of \(\min(s,h)\) over the columns of
the diagram, where \(h\) is the height of the column. Every column
therefore attains its maximum for each \(s\). Taking \(s=h\)
forces the first \(h\) ambient rows in that minor, giving
\(f_\lambda\) itself. Linear substitution in the Fischer pairing
is adjoint to its transpose, by the chain rule. Hence
\[
 [f_\lambda(Y^{\mathsf t}X),b_k(X)]_X
   =[f_\lambda(X),b_k(YX)]_X=h_\lambda b_k(Y).
\]
Row substitution by \(Y\) preserves the flag-polynomial space by
Cauchy--Binet. Comparing the standard basis coefficients and using
symmetry gives \(D_\lambda B_\lambda=h_\lambda I\). In particular the
polynomial of a Weyl coordinate vector \(v\) is
\(\sum_i(D_\lambda v)_i b_i\), and pairing it with \(b_k\) gives
\(h_\lambda v_k\). This also shows that \(D_\lambda\) is positive
definite for the Hermitian pairing \(\overline v^{\mathsf t}D_\lambda w\).

For the large positive block use the complementary diagram. At a
place of rank \(N''\), the substitution
\[
 P(X)\longmapsto \det(X)^c P(X^{-\mathsf t}w_0),
 \qquad c=\sum_{0<h<N''} k^+_\tau(h),
\]
with \(w_0\) reversing the columns, sends the basis products of the
diagram \(\lambda^\vee\), whose multiplicities are
\(k^+_\tau(N''-h)\), to signed products of the original minors.
Jacobi's identity gives this on each minor, since the complementary
columns to the last \(N''-h\) columns are the first \(h\) columns.
Applying the same substitution in the reverse direction composes
to multiplication by \(\det(w_0)^c\), so the matrix on the integer
bases is unimodular. The integrality here uses determinant
straightening over \(\mathbb Z\). Equivalently, within the flag
space one can read the coefficients of any polynomial with
coefficients in a subring from the tableau basis by elimination.
For the diagonal term order, the leading monomials are distinct
and monic, since their exponents record the ordered entries in
each tableau row.

On applying this substitution to both variables of the Gram
kernel, its argument becomes \(w_0^{\mathsf t}(Y^{\mathsf t}X)^{-\mathsf t}w_0\).
Jacobi's identity takes each leading principal minor of that
argument to the complementary leading minor of \(Y^{\mathsf t}X\)
divided by \(\det(Y^{\mathsf t}X)\). The determinant factors cancel.
Thus there is an integer matrix \(J\) with integer inverse for which
\[
 J^{\mathsf t}D_\lambda J=D_{\lambda^\vee},\qquad
 h_{\lambda^\vee}D_\lambda^{-1}=J B_{\lambda^\vee}J^{\mathsf t}.
\]
Here and below the diagram and its norm over all infinite places
use the products of the local bases and pairings.

Each stacked determinant of height \(h\) contains exactly \(h\)
entries from the active columns and \(N''-h\) added-block entries
in every nonzero term. Let
\[
 d_j^{\mathrm{pol}}=\sum_\tau\sum_{0<h<N''}(N''-h)k_\tau^+(h),
 \qquad D_j=\sum_\tau|\operatorname{DetIndex}_\tau|.
\]
The first number is the common added-block degree and satisfies
\(d_j^{\mathrm{pol}}\le31mD_j\), with \(D_j\le2fb\).
It is also the box count of \(\lambda^\vee\). Over the
ring generated by the entries of the \(\tau(Y_i)\), and with the
added-block variables as coefficients, each determinant product
lies in the standard flag-polynomial space for \(\lambda\).
In the factors attached to standard rows the upper unit rows
use the other active columns. Expansion along those unit rows
and then the first \(h\) variable columns identifies the needed
minors. For a factor attached to the short basis, all the active
columns are variable, and the same expansion applies.
The basis coefficient vector is integral over that ring in the
sense above. Therefore multiplying its Weyl coordinate vector
by \(H=h_{\lambda^\vee}\) clears it over the same ring.
The determinant expansion of \(f_{\lambda^\vee}\) has at most
\(d_j^{\mathrm{pol}}!\) terms counted with multiplicity. Each
monomial of degree \(d_j^{\mathrm{pol}}\) has squared Fischer norm
at most \(d_j^{\mathrm{pol}}!\). The triangle inequality gives
\(1\le H\le(d_j^{\mathrm{pol}}!)^3\), so
\(H\le(d_j^{\mathrm{pol}})^{3d_j^{\mathrm{pol}}}\), with the
empty degree giving \(H=1\). Both the degree and this multiplier
are fixed over table choices and twists.%
\footnote{\weightedpin{grh-compact/src/Lib/GLn/SchurModule/Gram.lean}
{123-L165}{Fischer form and the integral Gram matrix};
\weightedpin{grh-compact/src/Lib/GLn/Branching/Branching.lean}
{66-L120}{Integral complementary-minor isometry};
\weightedpin{grh/src/New/OneLift/ComparisonOperator/LineCoordinates.lean}
{696-L757}{Complementary degree and the clearing integer}.}

For a factor of height \(h\), the stacked determinant has
unit rows in the positive columns from \(h+1\) through
\(r_+=\kappa+j-1\). Subtract multiples of those upper rows
from the lower variable rows to clear the entries in these
columns. In the remaining determinant the first \(h\)
columns are the variables of the Weyl polynomial and have
zero upper block. The complementary columns have entries
from \(0,1,Y_i,z\), with \(z\) the specialized added-block
variables. The column count in this complement is at most
\(12m+1\), since \(r_+-h\le j\) on this diagram.
Expanding along the variable columns gives complementary
minors as the exterior coordinates. Their squared norm
is bounded by the full sum of squared minors of the
complementary columns and hence by the product of their
squared lengths, by Cauchy--Binet and Hadamard. On
\(\max|\sigma(z_a)|\le g^4\) the short basis has the same
entry bound. Each column has squared length at most \(g^9\),
since \(r+11m+1\le g\). The Cartan tensor norm is bounded by
the product of the exterior norms. Thus the squared
Weyl norm of the cleared coordinate vector at any such
specialization is at most \(H^2(g^9)^{(12m+1)D_j}\).%
\footnote{\weightedpin{grh-compact/src/New/Rows/LineCoordinates.lean}
{189-L272}{The actual shapes, determinant coordinates and hook};
\weightedpin{grh/src/New/OneLift/ComparisonOperator/LineCoordinates.lean}
{191-L233}{Complementary-minor bound}.}

By definition the coordinates of the large-block vector
are the class contraction of three definite kernels. The
first is the degree-zero carrier kernel of ranks \(m,3m\)
at unit modulus. The second is the degree-zero kernel
of ranks \(j,r+9m+1\) from the negative block, read via
the opposite-space identity. The third is the added-block
kernel from rank \(12m-j+1\) to \(r+9m+1\) with the
cleared polynomial. The latter two use
\(\mathfrak M^2 Q_\psi\). The scalar for a standard
definite kernel of ranks \(u,v'\) and degree \(d'\) can be
taken to be
\[
 C_{u,v',d'}
    =(g^2)^{2fuv'}\,v^{uv'}(g^2d\,\eta_0)^{d'}.
\]
Here \(v\) and \(\eta_0\) are the vacuum and jet parameters
and \(d\) is the discriminant prime. The normalized
vacuum datum at \(\mathfrak M^2 Q_\psi\) is integral at
every place over \(\ell\nmid\operatorname N(Q_\psi)\)
after including \((g^2)^{2fuv'}\): the integral
standard-data clearer uses \(g^2\operatorname N(Q_\psi)\)
and the omitted norm power is a unit. The jet then uses
the scale \(g^2d\,\eta_0\). This is the weighted standard
datum calculation with the two stars from the negative
block included in the conjugate reading. Coefficients of
the residue data are sums of additive roots of unity,
or indicators, and are integral.

To apply this bound in the transported gauges, move each
finite argument to integral frames at the places above
\(\ell\). A rational change of the representative exists by
weak approximation, using the Cayley chart in the unitary
group. The short-basis columns in these frames are integral
by their construction in the lattice of the corresponding
class. Under the change of representative, the positive
stacked minors transform by row substitutions on the
variable matrix and the associated substitutions on its
constant rows. Indeed the right determinants from the
change of columns multiply each determinant once. The
total exponent of this factor is the number of minors at
that place. The two changes of splitting character have
product \(\upsilon^{r+9m+1}\) in the opposite pair, so the
determinant scalar is the one in the Weyl action. Thus
transporting the output before reading its coefficients
gives the same polynomials on the integral input data.
This evaluates the unit gauge on the actual large-block
vector. Put \(D_{\mathrm{aut}}=\prod_{v\mid\mathfrak M}
|{\rm GL}_{rt}(k_v)|\).
In the double class sum use
\(D_{\mathrm{aut}}^2\) to clear both stabilizer orders.
The fixed normalization \(\beta_j\) is the product of this
integer, \(H\), the three kernel scalars with the negative
one conjugated, and the fixed opposite-space phase.
It is nonzero and depends on \(j\).
\footnote{\weightedpin{grh-compact/src/New/Rows/LineArithmetic.lean}
{30-L105}{The three kernels and their transport into coefficient lattices};
\weightedpin{grh/src/New/UpperBound/Kernels/DefiniteKernelSize.lean}
{119-L186}{Scale of a standard definite kernel}.}

For the complex norm at \(\psi=1\), a normalized kernel
of either degree-zero pair has conjugate modulus
\(\le e^{x_1}\), where \(x_1\) is the secondary cutoff.
For the third apply the weighted Gaussian integral to
the homogeneous polynomial after contraction with an
arbitrary unit vector in the dual Weyl norm. On the
polydisc for its coefficients, the preceding determinant
bound and Cauchy--Schwarz bound this contraction by
\(H(g^9)^{(12m+1)D_j/2}\). Homogeneity extends the estimate
to larger balls, using \((1+R)^{d_j^{\mathrm{pol}}}\le
e^{d_j^{\mathrm{pol}}R}\). With the actual ranks \(\le40\kappa\)
and degree \(\le g\), the Gaussian factor and the residue
and kernel clearers contribute at most \(e^{x_1}\) times
the maximum norm of the finite datum. The latter is at
most \(|{\rm GL}_{3m}(\mathcal O_E/\mathfrak M^2)|\),
the number of terms in the translated trace character.
All class counts and residue group orders here have
logarithms bounded by a fixed power of \(x\).
Applying the triangle inequality to the two class sums
and taking the supremum in the dual norm yields
\[
 \log\big\|y_{1,i}(h)^\sigma\big\|^2
 \le 6x_1+O_d(x^{99})+O(fmb\log b).
\]
Since \(r B_0=rf a(\log b)^3/(m\kappa)\) and
\(m^2\le4x\), the right side is \(o(rB_0)\).
This gives the required \(e^{rB_0}\) bound at each
conjugate, with the same \(\beta_j\) as the finite gauge.
\footnote{\weightedpin{grh-compact/src/New/Rows/MemberKernels.lean}
{143-L167}{Conjugate kernel bound on the polydisc};
\weightedpin{grh-compact/src/New/Rows/LineArithmetic.lean}
{52-L71}{The class counts and dual norm on the actual vector}.}

\paragraph{The comparison forms on these lattices.}
\teamanchor{proof.weighted-comparison-forms}%
Use the \hyperlink{proof.weighted-comparison-vectors}{normalized
vectors} \(y_{\psi,i}(h)\) with their one stage multiplier
\(\beta_j\). Their conductor-away gauges and conjugate squared
norms satisfy the bounds just proved. For the table entry
\(\xi_{i,c}\) at a small-block class \(c\), form the Cartan product
of \(y_{\psi,i}(h)\) and \(\xi_{i,c}\). Let \(D_{i,h,c,\psi}(\xi)\)
be the corresponding point form at
\(\operatorname{diag}(h,c)\gamma_i\), with the original marked
columns and level character. If \(J_i\) is the large-block level
fixing the stacked rows, set
\begin{align*}
 G_\psi\xi&=g^{L_j}\sum_{i,c}\;\sum_{[h]\,\mathrm{at}\ J_i}
                     D_{i,h,c,\psi}(\xi),\\
 L_j&=\sum_\tau
 \left(3m\sum_{a<N''}k_\tau^+(a)
             +\dim V^\circ\sum_{a<\dim V^\circ}k_\tau^-(a)\right).
\end{align*}
Here \(k_\tau^\pm(a)\) count the columns of height \(a\) in the
two parts of the stage weight. The row-fixing condition for
\(\varpi I_{N''}\) puts \(J_i\) inside the principal level
modulo \(\mathfrak M\), so its classes have trivial stabilizers. At primes of the
twist the point form keeps the prescribed determinant character.

The mixed Cartan map used in this product comes from multiplication
of the two polynomial modules, after complementing the negative
diagram. Multiplication is an integral surjection onto the module
for the union of their columns. A standard tableau of the union
can be factored by assigning its columns of each height to the two
diagrams in the prescribed numbers. Each factor is an integer
combination of standard products. The coefficient module of the
union is free, so the dual of the multiplication map has an
integral left inverse. Evaluations span the dual coordinate spaces,
and applying multiplication shows that the dual map sends an
evaluation for the union to the tensor of the two evaluations.
The Gram kernel of the union is the product of the two Gram
kernels, proving that this map is isometric with the normalized
Gram forms. In the complemented factor apply the unimodular
isometry just proved. Its formula on evaluation vectors gives the
contragredient action and the accompanying power
\(-\sum_{h>0} k_\tau^-(h)\) of the determinant. At the identity the
image is the tensor of the highest evaluation for the positive
diagram and the reversed highest evaluation for the negative
diagram. Place their supporting rows in \(V''\) and \(V^\circ\),
respectively, by a permutation of the target coordinates. The
orbits under the two block general linear groups span the two
coefficient modules on those blocks separately. This proves that
their pure tensors lie in the Cartan image. The block extensions
preserve their metrics, since the Gram kernels of the extended
evaluations restrict to those on the blocks. Applying the integral
left inverse to the pure tensors gives the actual integral
product and its asserted norm.%
\footnote{\weightedpin{grh-compact/src/Lib/GLn/Dual/MixedCartan.lean}
{15-L47}{The isometric embedding and determinant exponent};
\weightedpin{grh-compact/src/Lib/GLn/Dual/ClassSum.lean}
{35-L85}{Product on the two blocks}.}

For the gauge bound, transport the coefficient vectors into the
local integral frames at \(h,c\). Lift the local change of columns
to the compact integral group, using the inverse transpose at the
conjugate place. Inserting this matrix between the inverse lattice
and \(h^{-1}\) replaces \(h\) by a compact translate and preserves
the input gauge. In the positive part only the
designated \(3m\) columns of the inverse lattice can have a local
denominator \(\varpi\), by the column bound. In the negative part
use the conjugate inverse on the dual block, for which at most
\(\dim V^\circ\) columns can have this denominator. Every
determinant in a Weyl polynomial is multilinear in the columns.
Thus clearing a column in the positive or negative block
costs at most its number of minor factors in the corresponding
weight. The integral left inverse of the mixed Cartan embedding
carries this bound to the coefficient vector of the product.
This is the integral Weyl model of
\hyperlink{proof.near-level-determinants}{the earlier class lattice},
using its value lattices before projection. Since \(g\in\mathfrak M\),
the factor \(g^{L_j}\) clears all the resulting losses at the two
readings. More precisely, at a place \(a\) the input gauge on the
small block is the gauge of \(\bar\xi_{i,c}\) at
\(a\circ\operatorname{conj}\). The point form is an orbit sum of the
transported coefficient vector. Its terms have no extra denominator
and level characters are local units. Summing in \(h,i,c\) therefore
preserves integrality.\footnote{\weightedpin
{grh-compact/src/New/Rows/LineArithmetic.lean}{25-L41}{The normalized large-block input},
\weightedpin{grh-compact/src/Lib/GLn/Dual/ClassSum.lean}{35-L85}{Cartan product and integral left inverse},
and \weightedpin{grh-compact/src/New/Rows/CartanIntegrality.lean}{17-L68}{Loss at the two block embeddings}.}

The rational Cartan coordinates give the same product norm
after applying a coefficient automorphism to the vectors. At a target
class a point form contributes at most \(J_{\rm aut}\) vectors
of that norm, where \(J_{\rm aut}=\prod_{v\mid\mathfrak M}
|{\rm GL}_{rt}(k_v)|\). Indeed reduction
of a class stabilizer is injective at the chosen levels by freeness,
and \({\rm GL}_{r+t}(k_v)\) embeds in \({\rm GL}_{rt}(k_v)\) for
the ranks under consideration. Each summand comes from a unitary rational
transport and a root of unity, and remains an isometry after
applying the coefficient embedding to the transported matrices.
Apply the triangle inequality to the \(h\)-sum and
Cauchy--Schwarz to the table index sum. If \(C_V\) is the number
of target classes, \(C_J\) bounds the number of \(J_i\)-classes,
and \(I_j\) counts the pairs \((i,c)\), the result is
\[
 \|G_1\xi\|_\sigma^2\le
 C_V J_{\rm aut}^2 C_J^2 I_j\,g^{2L_j}e^{rB_0}
                         \sum_{i,c}\|\xi_{i,c}^\sigma\|^2 .
\]
At the trivial twist all these definite levels contain the
principal congruence subgroup modulo \(g^2\), with ranks
at most \(r+t\). The
\hyperlink{proof.weighted-short-class-sum}{definite class count}
therefore bounds their logarithmic class counts by \(x^{99}\)
for large \(x\). The logarithm of \(I_j\) is bounded by the two
class counts, \(3mN''\log|\mathcal O_E/\mathfrak M^2|\) for the
matrix choices, and \(2mf r^8\log(6m+1)\) for the row choices.
These terms and \(\log J_{\rm aut}\) are \(O(x^{99})\).
The number of boxes for the small negative weight is
\(O(f(r+t)^2r^8)\), so the logarithm of its hook number \(k_j\)
is also \(O(x^{99})\). Finally
\(\sum k_\tau^+=w_j(\tau)-j-\lfloor(t+\kappa-1)/2\rfloor\le2b\)
and \(k_\tau^-(a)\le 2r^8\). It follows that
\[
 2 L_j\log g=O(fmb\log b+x^{99})=o(rB_0).
\]
The last equality uses \(m^2\le4x\), \(a\ge xb\), and
\(B_0=fa(\log b)^3/(m\kappa)\). Therefore \(k_j\|G_1\xi\|_\sigma^2
\le e^{2rB_0}\|\xi^\sigma\|^2\) for sufficiently large scales.
The hook is used once, in the normalization of the compressed
operator.\footnote{\weightedpin
{grh-compact/src/New/Rows/ComparisonForms.lean}{128-L163}{The same norm calculation in class coordinates},
\weightedpin{grh-compact/src/New/Data/ChosenPlaceLevels.lean}{34-L79}{The common stabilizer bound},
\weightedpin{grh-compact/src/New/Rows/LineCoordinates.lean}{123-L138}{The actual table choices},
and \weightedpin{grh-compact/src/New/Rows/LineCoordinates.lean}{189-L192}{The positive weight sum}.}

\paragraph{The real matching of the determinant products.}
\teamanchor{proof.weighted-real-matching}%
Use the ordered positive and negative blocks on the two signed
spaces. Let \(W\) have signature \((p+a,q)\) and \(S\) signature
\((p,q)\), with complementary target \(I_a\) definite. Write
\(N=p+a+q\), and put \(d_\tau=|\operatorname{DetIndex}_\tau|\)
and \(\gamma_\tau=q-d_\tau\). Order the positive indices of \(W\)
by the fixed bijection from the stack of \(p\) and \(a\) rows.
In these coordinates let \(I\) denote the first \(p\) columns
and \(J\) the remaining \(a\). Write \(Z_+\) for the positive
columns in the \(I_a\)-by-\(W\) Fock variables, and \(Y\) for
the positive-positive block of the \(W\)-by-\(S\) variables.
Let \(Y_{--}\) be the negative-negative block of the latter.
For the constant matrices \(\eta_{\tau b}\) of size \(p\) by
\(p+a\), the two polynomials are
\[
 P_\eta=\prod_{\tau,b}\det\binom{\eta_{\tau b}}{Z_{+,\tau}},\qquad
 Q_\eta=\prod_\tau\det(Y_{--,\tau})^{d_\tau+N-2q}
                    \prod_{\tau,b}\det(\eta_{\tau b}Y_\tau).
\]
Let \(\eta^0=(I_p\ 0)\) for all factors, and put \(P_0=P_{\eta^0}\),
\(Q_0=Q_{\eta^0}\). Empty determinants and products have value one.

For these two oscillator actions use splitting weights
\((\gamma,0)\). The compact action of a diagonal sign-block pair
\((x,y)\) substitutes \(y^{\mathsf t}Zx^{-\mathsf t}\) on
equal signs and \(y^{-1}Zx\) on unequal signs, with the blocks
read separately. Its scalar in a \(+\) row block of \(y\) is
\(\det(y_+)^{\gamma_\tau-\#\{x\text{ negative indices}\}}\).
For the rank-\(a\) compact group of the first oscillator and
for the positive-block group acting on the second oscillator
this is \(\det(y_+)^{-d_\tau}\). It is the splitting character
times the vacuum determinant of the opposite-sign variables.
The positive-block group acting through \(x\) on the first
oscillator has scalar one because \(I_a\) is definite and the
second splitting weight is zero.

Denote the first and second real actions by \(\Omega_A\) and
\(\Omega_B\). Integrate in the first Fock space over the circle
with normalized measure. At every place let \(t(z)\) act on
the positive block of \(W\) by \(z^{-2}\) on \(I\) and \(z^{-1}\)
on \(J\), trivially on the negative block. Put
\(u_\tau(z)=z^{-1}I_a\). The operator
\[
 \left(\prod_\tau\det(u_\tau(z))^{d_\tau}\right)
                       \Omega_A(t(z),u(z))
\]
substitutes \(z\) times each variable outside \(Z_J\) and
fixes \(Z_J\), with no scalar. Its circle average sends
polynomials to their restrictions to \(Z_J\). After that,
average the action \(\Omega_A(1,U(a))\) over the product of
compact groups with mass one. For a polynomial in \(Z_J\), the
result transforms under left substitution by the determinant
character \(\prod_\tau\det(u_\tau)^{d_\tau}\). The identity
extends to the full general linear groups by polynomial
continuation from the unitary groups. An invertible square
\(Z_{J,\tau}\) is in the orbit of the identity, so such a
polynomial is a multiple of
\(P_0=\prod_\tau \det(Z_{J,\tau})^{d_\tau}\).
Both compact averaging operators are bounded in the Fock norm
by unitarity, and the polynomial vectors are dense there.
The composite thus sends each real datum \(v\) to
\(\kappa(v)v_{P_0}\), with a scalar independent of the finite
data and all evaluating functionals.

Fix the adelic evaluation point and denote by \(\ell_x\) the
continuous first theta functional at its right translate by \(x\).
Let \(\ell'\) be a linear functional on the second polynomials.
The first functional is invariant under the \(I_a\)-compact
factor, and \(\ell_x(\Omega_A(y,u)w)=\ell_{xy}(w)\) for the
real translations in use. Suppose that
\(\ell_k(v)=\ell'(\Omega_B(1,k)v_{Q_0})\) on the positive compact
block, identifying polynomials and their real Fock vectors on
the right. Apply the two averages and evaluate with \(\ell_k\).
The \(U(a)\) average leaves \(\ell_k\) unchanged.
On the other side, the polynomial of \(\Omega_B(1,kt(z))v_{Q_0}\)
is
\[
 z^{a\sum_\tau d_\tau}
       \left(\prod_\tau\det(k_\tau)^{-d_\tau}\right)
       Q_{\eta^0 k^{\mathsf t}}.
\]
Indeed \(\det(t_\tau)^{-d_\tau}=z^{(2p+a)d_\tau}\) and
\(\eta^0 t_\tau^{\mathsf t}=z^{-2}\eta^0\), so the circle
weight cancels with \(\prod_\tau\det(u_\tau(z))^{d_\tau}\).
On the first side, substitution in the stacked determinant gives
\[
 \Omega_A(k,1)v_{P_0}
   =\left(\prod_\tau\det(k_\tau)^{-d_\tau}\right)
                    v_{P_{\eta^0 k^{\mathsf t}}}.
\]
For example, multiply the whole stacked matrix on the right by
\(k_\tau^{-\mathsf t}\) after replacing the upper rows by
\(\eta^0k_\tau^{\mathsf t}\). Cancelling the common determinant
scalar proves the matching for all equal-row powers on each
place.

To allow the different rows \(\eta_{\tau b}\) in the same
product, let \(F_\eta(Y)=\prod_{\tau,b}\det(\eta_{\tau b}Y_\tau)\).
All these polynomials transform on the right by
\(\prod_\tau\det(u_\tau)^{d_\tau}\) for \(u_\tau\in\mathrm{GL}_p(\mathbb C)\).
The Hermitian Fischer pairing with \(\prod_\tau\det(Y_{I,\tau})^{d_\tau}\)
is its squared norm times evaluation at \(Y_\tau=\binom{I_p}{0}\).
In fact one can set the other rows to zero in this pairing,
and the stated determinant transformation on the right gives
the assertion on the square surviving block. After unitary row
translation these pairings recover evaluation at every tuple
of orthonormal \(p\)-frames. Gram--Schmidt and the right
determinant rule recover evaluation at all full-rank frames.
Those frames are dense, so the equal-row powers span the space
of the \(F_\eta\).

Every linear relation among the \(F_\eta\) is also a relation
among the \(P_\eta\). For a specialization of \(Z_{+,\tau}\)
of full row rank, take a matrix \(A_\tau\) whose \(p\) columns
are a basis of its kernel and complete it to an invertible
matrix \((A_\tau\ B_\tau)\). Multiplication by this matrix
puts the lower-left block of each stacked determinant equal
to zero. Thus the stacked determinant is
\(\det(Z_{+,\tau}B_\tau)\det(\eta_{\tau b} A_\tau)/
\det(A_\tau\ B_\tau)\). Its scalar factor is independent of
\(\eta_{\tau b}\). These specializations are dense, proving
the relation as a polynomial identity. The passage from \(F_\eta\)
to \(Q_\eta\) is multiplication by the fixed negative minor
powers and also preserves all relations. Linearity gives
\(\kappa(v)\ell_1(v_{P_\eta})=\ell'(Q_\eta)\) for the full
family of determinant products.%
\footnote{\weightedpin{grh-compact/src/Lib/Weil/DualPair/PairReal.lean}
{35-L66}{Real substitution and vacuum determinant};
\weightedpin{grh-compact/src/New/Rows/RealMatching.lean}
{38-L113}{The same real matching for independent determinant rows}.}

\paragraph{Determinant partners and the exchange scalar.}
\teamanchor{proof.weighted-det-exchange-scalar}%
For the same real spaces let \(b'=p+q\), and suppose
\(a=b'+2e\) and \(N\ge 2q+2\). Let \(\alpha\) have weight
\(\gamma_\tau=q-d_\tau\) and tame character \(\chi_\alpha\)
on \(\mathcal O_E/\mathfrak m\), where \(\mathfrak m\) is supported
on conjugate-disjoint split places. For an arbitrary function
\(F\) on the rectangular residue matrices use
\begin{align*}
 \Phi(X'\ X'')&=\chi_\alpha(\det X'')F((X'')^{-1}X'),\\
 \Psi\binom{Y'}{Y''}&=\chi_\alpha(\det Y')F(-Y''(Y')^{-1}).
\end{align*}
The expressions are zero where the indicated square block is
noninvertible. These are matching data for \(\chi_\alpha^{-1}\).
To verify this, write the first \(b'\) columns of an integral
invertible matrix \(k\) as \(\binom{A}{C}\) and the last \(a\)
rows of \(k^{-1}\) as \((E\ D)\). Jacobi gives
\(\det D=\det A/\det k\), so the two supports agree on these
columns and rows. Where they are nonzero, \(EA+DC=0\) gives
\(D^{-1}E=-CA^{-1}\). Substitution yields
\(\Phi(E\ D)=\chi_\alpha(\det k)^{-1}\Psi\binom{A}{C}\),
with a primitive frame on each open support. This proof works
over the product residue ring including its squared local
factors. At the remaining places the datum is the vacuum.

In the \hyperlink{proof.weighted-complementary-period-norm}{normalized
complementary-period formula} take
\(\nu=\alpha^{-1}\upsilon^{N/2}\). Its weight at \(\tau\) is
\(N/2-q+d_\tau\), within that formula's stable range. Apply the
formula to the fixed real vector \(v_{Q_0}\). It gives a first
real datum \(v\) independent of the tame character and a nonzero
scalar \(c_\alpha\) for the partner equality, with
\[
 |c_\alpha|^2=C\,\Pi_{\nu,\mathcal S}(e),\qquad
 \Pi_{\nu,\mathcal S}(e)=
     \left(\prod_{w\in\mathcal S}q_w\right)^{e^2}
     \prod_{z=1}^{e}|L^{\mathcal S}(z,\nu)|^2,
 \qquad \mathcal S=\operatorname{supp}\mathfrak m.
\]
The scalar on the integral over \([U(S)]\) is \(c_\alpha\).
Use that integral multiplied by \(c_\alpha\) for \(\ell'\) in the
real matching. Translation of the original formula gives the
required equality on the positive compact block. The transfer
of polynomials gives \(\kappa(v)/c_\alpha\) as the multiplier of
the first theta period to obtain the second one.

For the resulting \(P_\eta\), the compact real action of
\(I_a\) is trivial, because substitution in the \(a\) lower rows
of the stack contributes exactly the determinant character
inverse to its vacuum and splitting phase. At finite integral
points of \(I_a\), the multiplier \(\alpha(\det k)\) cancels
\(\chi_\alpha(\det k)^{-1}\) from replacing the matrix argument
by \(k^{-1}X\). Thus integration over \([U(I_a)]\) is the class
sum at full level with weights \(1/|\operatorname{Aut}|\), since
the real and the standard finite compact factors have volume
one. Changing the second splitting character on the first
kernel from \(1\) to \(\beta\) multiplies it by
\(\beta(\det g)^{-1}\). The same multiplier on the second
kernel replaces its first splitting character \(\alpha\) by
\(\alpha/\beta\). Hence the exact matched scalar has squared
norm \(C'/\Pi_{\nu,\mathcal S}(e)\), with \(C'=|\kappa(v)|^2/C\)
fixed across the twists. Nonvanishing in the applications is
established by their period identities.%
\footnote{\weightedpin{grh-compact/src/New/Rows/Carrier.lean}
{112-L149}{Class measure, character changes and real matching};
\weightedpin{grh-compact/src/New/Exchange/ExchangeTheorem.lean}
{14-L46}{Adelic partner identity and its squared scalar};
\weightedpin{grh-compact/src/Lib/SiegelWeil/Tower/LargeSideConstantTerm.lean}
{22-L31}{The conductor power and squared partial values}.}

\paragraph{The mass and the three characters in the matching.}
\teamanchor{proof.weighted-comparison-matching-reduction}%
Write \(\chi_\psi=\chi^j\psi\), and
\(b_\psi=\mu_j^0\psi\,\upsilon^{j+\lfloor(t+\kappa-1)/2\rfloor}\).
In the matching construction the first two exchanges on
the carrier have complementary rank differences \(m\)
and \(2m\). Their Hecke characters in the norm of the
exchange scalar are \(\chi_\psi\upsilon^{m/2}\) and
\(\chi_\psi^{-1}\upsilon^m\), at \(m/2\) and \(m\).
In the third exchange the ranks are \(r+9m+1\) and \(3m\),
so half the difference is \(z_*=(r+6m+1)/2\).
The common negative block has rank \(m\), the padded
source has rank \(r+12m+1\) after including the carrier
indices, and this rank exceeds \(2m+1\). For the three
exchanges the square blocks have the tameness required
at \(\mathfrak M^2Q_\psi\), since \(\chi^j\) has nontrivial
tame part at both chosen sets and \(\psi\) is nontrivial
at every prime of its conductor. The real weight check
on the third block is
\(|\operatorname{DetIndex}_\tau|+
j+\lfloor(t+\kappa-1)/2\rfloor=w_j(\tau)\).
The product of the carrier character with \(\chi_\psi\)
is \(b_\psi\upsilon^m\). The character entering the third
exchange at \(z_*\) is
\[
 (b_\psi\upsilon^m)^{-1}\upsilon^{3m+z_*}
  =\bigl(\mu_j^0\psi\,
       \upsilon^{j+\lfloor(t-11m-1)/2\rfloor}\bigr)^{-1}.
\]
Here \(2m+z_*=(11m+\kappa)/2\) and both
integers inside the half-rank floors are even. This
identifies the third factor with the first pair of
\(\Gamma_j\) in the reader's convention.
\footnote{\weightedpin{grh-compact/src/New/Rows/Separation.lean}
{43-L65}{The three Hecke pairs};
\weightedpin{grh/src/New/OneLift/ComparisonOperator/ThreeExchanges.lean}
{528-L554}{The character equalities in the two carrier and third exchanges}.}

The finite change of variables can be made directly.
For a square block \(x''\) on the support of the determinant
character, write \(x'=x''z\). Then the invariant finite
datum is \(\chi_\psi(\det x'')\,F(z)\), where
\(F(z)=\sum_{k\in{\rm GL}_{3m}(\mathcal O_E/\mathfrak M^2)}
\operatorname{trchar}_i(zk)\).
If \(x''\) is singular modulo the local modulus the
determinant datum vanishes. Unfold the full-level class
sum by representatives for the row-fixing subgroup.
The allowed integrally invertible \(x''\) lies in exactly
one such coset. On that coset replace \(x''\) by its
marked columns and force the other columns to be
those of the identity. The invariance of \(F\) under
the row-fixing subgroup uses \(s_iA=s_i\) and cyclicity
of the trace. The normalized class measure converts the
sum over the cosets to the class sum at the row-fixing
level. This leaves the finite indicator of the comparison
lattice and the two class variables of its prescribed
large-block vector.

The first exchange uses \(W=I_m\), \(S=0\), \(a=m\),
\(d_\tau=0\), and \(\alpha=\chi_\psi^{-1}\). The second uses
\(W=I_m\oplus S_0\) with \(S_0\) of signature \((2m,m)\),
\(S=I_m^-\), \(a=3m\), \(d_\tau=0\), and
\(\alpha=\chi_\psi\upsilon^m\). The third uses
\(W=I_{3m}\oplus W_j\oplus I_j^-\oplus I_{3m-j}\oplus I_{6m+1}\),
\(S=S_0\), \(a=r+9m+1\), and \(\alpha=b_\psi\upsilon^m\).
Here the original source has \(m-j\) negative signs, so the
first two nondefinite blocks account for all \(m\) negative
indices of \(W\) in the last exchange. The ambient dimensions
are even and at least \(2q+2\) in all three cases. The half-rank
differences are \(m/2\), \(m\) and \(z_*\). The real block
weights just checked are exactly \(q-d_\tau\) in these spaces.
For the first two exchanges the upper-row determinant
polynomial on \(W\) is the empty product. The corresponding
polynomials on \(S\) are \(1\) and
\(\prod_\tau\det(Y_{--,\tau})^{2m}\). For the third, the negative
exponent is \(d_\tau+r+10m+1\). Its positive factors from the
standard rows select the padded minors in \(W_j\oplus I_j^-
\oplus I_{3m-j}\), and the remaining factors have the indicated
rows of \(Y_i\) in \(I_{6m+1}\). This is the prescribed
polynomial before the row-fixing restriction.

The first class identity is normalized to the constant one
on \(U(I_m)\), by dividing by the fixed mass of the rank-zero
quotient. Conjugating and negating its source yields the
constant one on \(U(I_m^-)\). The sign change of the oscillator
conjugates the two characters and multiplies both by
\(\upsilon^m\). It sends the finite determinant datum to the
new determinant character times a factor of modulus one,
since scaling the matrix by the finite norm-minus-one unit
scales its determinant by the \(m\)-th power of that unit.
The fixed phase of the oscillator also has modulus one.
Thus the conjugated first normalizer has the same squared norm.

Insert its class expression for the constant function in the
integral over \(U(I_m^-)\) in the second exchange. For the
finite datum take \(F=1\) in the determinant partner formula.
The product of the two kernels in this integral is the kernel
for the direct sum with the determinant datum on the square
block. This follows by separating the rational lattice sum and
using the product of the two splitting characters. Apply the
second exchange to each class term and split the larger
kernel on \(I_m\oplus S_0\) again. The untwisted finite function
on the \(I_m\) block is the standard vacuum, so this term is
the full-level class coefficient of the carrier. The resulting
class sum on \(I_{3m}\) comes with the product of the first two
normalizers and the modulus-one phase. On the smaller integral
the first character is \(\chi_\psi\upsilon^m\), so its value is
\(\chi_\psi(\det x)\) times the fixed carrier. Changing the second
character on the class expression to \(\chi_\psi\) multiplies
it by \(\chi_\psi(\det x)^{-1}\) and gives the asserted equality.

For the third exchange insert this carrier class sum. Here take
\(F(z)=\sum_k\operatorname{trchar}_i(zk)\), using the reduction
of the rectangular matrix to \(\mathcal O_E/\mathfrak M^2\),
with \(k\) running through \(\mathrm{GL}_{3m}(\mathcal O_E/
\mathfrak M^2)\). On the support of \(\Psi\), right invariance
of \(F\) makes \(F(-Y''(Y')^{-1})=F(Y'')\). The product kernel
therefore has precisely the determinant partner datum on the
third space. Taking the class sum and applying the third
exchange produces the fixed carrier integral with this sum of
translated trace characters. Full level of the carrier and
triviality of the tame part of its other splitting character
give identical integrals for the translates. Thus the scalar
on the row-fixing class sum is the product of the three
normalizers, including the modulus-one phase, divided by
\(\lvert\mathrm{GL}_{3m}(\mathcal O_E/\mathfrak M^2)\rvert\).
The measure at the smaller row-fixing level is the class sum
with the same stabilizer weights, as in the preceding
unfolding. Each squared norm contributes the reciprocal of
the corresponding factor computed above. It follows that
\(|\kappa_\psi|^2=C/\Gamma_j(\psi)\) on the actual comparison
vectors, with one twist-independent \(C\).%
\footnote{\weightedpin{grh-compact/src/New/Rows/Carrier.lean}
{154-L199}{The two carrier exchanges};
\weightedpin{grh/src/New/OneLift/ComparisonOperator/ThreeExchanges.lean}
{686-L790}{Third polynomial and the three squared factors}.}

The Weil phase from
the glued lattice uses the source splitting character
\(\upsilon^{\#\{\text{negative signs of }W_j\}}\), which
is independent of \(\psi\) even on the nonintegral
chosen-place point. Its other factor is a power of the
finite idele absolute value. Thus the common-adjoint
normalization does use an untwisted \(\omega_i\).
\footnote{\weightedpin{grh-compact/src/New/Rows/Carrier.lean}
{52-L103}{Change of variables and row-fixing subgroup};
\weightedpin{grh-compact/src/New/Rows/RealMatching.lean}
{52-L98}{Real Fock matching};
\weightedpin{grh-compact/src/New/Rows/GluedFormTheta.lean}
{159-L165}{Glued-lattice factor and source character}.}

For the large-block detection calculation put \(n_*=2ft^3\) and \(B=5b\).

\paragraph{A first minimum at a half-supremum point.}
\teamanchor{proof.weighted-comparison-peak}%
For a moved Hermitian lattice \(L\) use the Euclidean norm of its
complex coordinates, taking one embedding per conjugate pair. Put
\[
 \lambda(L)=\min\bigl(1,\inf_{0\ne y\in L}\|y\|\bigr).
\]
Suppose a bounded automorphic function \(U\) is holomorphic of
weight \(u\) in an order with the positive indices first. Assume that
\(\max_{\tau,i}|u_{\tau i}|\le B\), \(B\ge2\), and that every negative
entry of \(u_\tau\) is at least one larger than every positive entry.
Let \(p\in\mathcal O_E\) satisfy \(\bar p=-p\) and
\(0<\operatorname{Im}\tau p\le C\) at the chosen embeddings, with
\(C\ge1\). Suppose \(U\) is invariant under the principal level of
\(p\). Then every \(z\) such that
\(|U(z)|>\tfrac12\sup|U|\) satisfies
\(\lambda(L_z)\ge (BC)^{-1/2}\).

Indeed, a shorter nonzero vector has integral Hermitian square. At
each embedding the absolute value of that square is bounded by the
squared Euclidean norm of the corresponding component, hence is
strictly less than one. Its field norm gives isotropy. In the moved
frame write \(v_\tau\) for its real-place coordinates, \(J\) for the
diagonal signature matrix and
\(\eta_\tau=\operatorname{Im}(\tau p)\|v_\tau\|^2\). Each \(\eta_\tau\)
is positive and \(B\sum_\tau\eta_\tau<1\). On the upper half-plane
consider
\[
 M_\tau(s)=I-(\tau p)(s-i)v_\tau v_\tau^*J .
\]
Isotropy gives \(M_\tau(s+1)=M_\tau(1+i)M_\tau(s)\).
The transvection at \(1+i\) comes from the rational vector.
Its finite component is congruent to \(I\) modulo \(p\) in the moved
frame, since the finite coordinates of the vector are integral.
Left rational invariance and the right level therefore make the
holomorphic extension of \(U(z\,\cdot)\) along \(M(s)\) periodic
of period one.

To bound this extension, order the coordinates as above and put
\[
 P_l=1-2\operatorname{Re}((\tau p)(s-i))
                  \sum_{i'\le l} J_{i'i'} |v_{\tau i'}|^2 .
\]
The determinant lemma identifies \(P_l\) with the leading Gram
minor of \(M_\tau(s)^*JM_\tau(s)\) divided by the corresponding
minor of \(J\). As \(l\) increases, the \(P_l\) go monotonically
from \(1\) to \(P=1+(\operatorname{Im}s-1)\eta_\tau\) through the
positive indices and back to \(1\) through the negative indices.
They stay positive. Gram--Schmidt for \(J\) thus gives a unitary
factor followed by an upper triangular factor with squared
diagonal moduli \(P_l/P_{l-1}\). If \(\operatorname{Im}s\ge1\),
the logarithms of these ratios have the signs of \(J\) and sum to
zero. The weight gap bounds their \(u\)-weighted sum by
\(-\log P\). If \(0<\operatorname{Im}s<1\), the size of the weights
bounds the sum by \(-2B\log P\le 4B\eta_\tau\), since
\(\eta_\tau<1/2\). The modulus of the weight character uses one
half of these logarithmic bounds. The extension is consequently
bounded by \(e^2\sup|U|\), and tends uniformly to zero as
\(\operatorname{Im}s\longrightarrow\infty\) by the bound
\(\prod_\tau P^{-1/2}\) in the first case. In the variable
\(q=\exp(2\pi i s)\) it extends holomorphically across \(0\) with
value zero. Schwarz's lemma at \(s=i\), where \(M_\tau(i)=I\),
gives \(|U(z)|\le e^{2-2\pi}\sup|U|\). This excludes the putative
short vector. A nonzero bounded \(U\) has such a half-supremum
point by the definition of the supremum.\footnote{\weightedpin
{grh-compact/src/Lib/DiscreteSeries/Holomorphic/NearMaximum.lean}
{22-L118}{The upper-half-plane transvection, period and signed minors}.
For the earlier scalar Siegel extremum compare
\href{https://faculty.fordham.edu/poor/files/Semihull00.pdf}
{Poor and Yuen, \emph{Linear dependence among Siegel modular forms},
Theorems 1.2--1.3}.}

We use two instances. Write \(G_0=g_0g_1=\pi_G\bar\pi_G\) for
the product of the chosen rational primes, with the prescribed
half \(\mathfrak M=(\bar\pi_G)\), and \(S\) for the Hermitian space of
signature \((2m,m)\). The fixed carrier \(\phi\) on \(U(S)\) has
weight \((0;3m)\). The first two carrier exchanges, at the
untwisted stage character, express it as a finite combination of
theta kernels with determinant character on square matrices.
The finite support forces an invertible residue matrix in the
orientation with \(S\) as source as well, since the matrices are
square. Apply the \hyperlink{proof.weighted-grid-primitive}
{primitive-support decay estimate} with its short isotropic
vector \(d'\) in this source and a full frame \(C'\) in the
fixed target. The corresponding modulation has coefficient
\(C'{}^*Yd'\) up to the invertible Hermitian and frame factors.
It is nonzero on the named support since \(Y\) is injective.
Thus the additive average vanishes there, and repeated differences
of the bounded real translates give every power of the inverse
source height. The definite target runs over finitely many classes.
The level character cancels
the residue determinant character on the full integral group.
The real polynomial \(1\) gives the displayed holomorphic weight.
The carrier is nonzero. In its original definition it is the lift
of \(1\) from \(U(0,m)\) to \(U(S)\), with the unit finite vacuum
and the polynomial \(\prod_\tau\det(X_{\tau,-})^{2m}\).
On the compact real source its source matrix coefficient is
constant, equal to the positive squared norm of the polynomial
vector. In the block-minor formula the source has only negative
indices. Its weight there is the positive target rank \(2m\)
minus the determinant exponent \(2m\), and the accompanying
conjugate factor has weight zero there. The
convergent doubling pairing for \(3m>2m\) sends a zero lift to
zero convolution with this coefficient; its convolution with the
constant \(1\) is the positive compact Haar mass times that squared
norm. The finite vacuum has positive mass at full level and the
source is definite. Thus one can fix a carrier peak \(z_S\),
independently of the good form to be tested.

For a good form \(F\) on \(W_j\) apply the same estimate to
\(U=\overline F\). To see the holomorphic weight, use the injective
principal-space lift \(L\) at the trivial twist, supplied by the
\hyperlink{proof.weighted-auxiliary-injection}{sharp-point pairing on
the full auxiliary base}. In finite dimension
\(L^*L\) is invertible, so \(F\) is the descent of \(Ly\) for some
\(y\). The highest polynomial of the descent is a product of the
joint harmonics of the edge. Its translates on the definite side
give the conjugate holomorphic type \(k\) on \(W_j\). Its weight
has gap at least \(300n_*\) and \(\max |k|\le5b\). The good level
contains the principal level of \(G_0\sqrt{-d}\), since
\(\bar\pi_G\mid G_0\). Its cuspidal realization makes \(F\) bounded.
Thus one may use \(B=5b\) for both forms. For \(\phi\) take
\(p=\sqrt{-d}\), \(C=\sqrt b\). For \(\overline F\) take
\(p=G_0\sqrt{-d}\), \(C=b^{23}\). These imaginary parts are
positive on the fixed CM type, and \(G_0\le b^{22}\).\footnote{\weightedpin{grh-compact/src/New/Rows/NaturalFormNear.lean}
{18-L52}{Carrier weight, full level and the compact-source test};
\weightedpin{grh/src/New/OneLift/ComparisonOperator/ProductsSpan.lean}
{729-L764}{The holomorphic type of the conjugate good form};
\weightedpin{grh/src/New/OneLift/ComparisonOperator/ProductsSpan.lean}
{862-L883}{Level for the second peak}.}
The respective first-minimum bounds are
\[
 \lambda_S(z_S)\ge(5b\sqrt b)^{-1/2},\qquad
 \lambda_{W_j}(z)\ge(5b^{24})^{-1/2}
       \quad\text{at every peak of }\overline F .
\]

\paragraph{A sparse theta kernel on the two moved lattices.}
\teamanchor{proof.weighted-comparison-sparse-kernel}%
Let \(\widetilde W_j=W_j\perp I_j^-\perp I_{3m-j}^+\), and put
\(\nu=n_*=2ft^3\). Fix \(z\) with
\(\lambda_{W_j}(z)\ge (5b^{24})^{-1/2}\), and extend it by the
identity on the padding. Both standard padding lattices have first
minimum one with the cutoff at one. Thus
\[
 \lambda_S(z_S)\lambda_{\widetilde W_j}(z\oplus1)
 \ge (5b^{51/4})^{-1}\ge b^{-13}.
\]
This estimate uses the same padded point in the theta kernel.
Write \(s=3m,\ s'=r+3m\) for the two ranks and \(Q_E=dN_1\).
For their moved Hom lattice the lower bound is
\[
 \lambda_{\operatorname{Hom}}\ \ge\
 \frac{\lambda_S\lambda_{\widetilde W_j}}
 {6Q_E(fs)^2 s s'} .
\]
Here is the transference argument with indefinite signatures.
The Hermitian lattice is self-dual at every finite place. The
Euclidean dual of its archimedean realization is obtained from the
trace dual by the sign matrix and the factor two. The sign matrix
is a Euclidean isometry. Since \(Q_E\) belongs to the different,
the Euclidean dual has minimum at least \(2\lambda_S/Q_E\).
The Euclidean covering-radius argument preceding
\hyperlink{proof.weighted-classwise-evaluation}{classwise evaluation}
gives an \(E\)-basis in the source lattice with lengths at most
\(6Q_E(fs)^2/\lambda_S\). A nonzero integral Hom matrix takes
one basis vector to a nonzero integral target vector. The
Hilbert--Schmidt bound on this evaluation gives the displayed
inequality, even with the harmless factors \(ss'\). Cutting the
first minima at one preserves this lower bound.

The polynomial \(P\) on \((S,\widetilde W_j)\) in the definition of
\(F_i^\natural\) uses the negative \(m\)-row minor with exponent
\(D_\tau+N''+m\) and the positive \(2m\)-row minors with the column
multiplicities \(k_\tau^+(h)\). Here
\(N''=r+9m+1\), \(D_\tau=\Delta'+\sum_h k_\tau^+(h)\), and
\(\Delta'=h_2-\lfloor(t+1)/2\rfloor\).
For each minor the rows and columns are distinct and have the
same sign, including the rows carried into the padding. Thus \(P\)
is nonzero and is homogeneous of degree
\[
 d_P=\sum_\tau\left[m(D_\tau+N''+m)
                       +2m\sum_h k_\tau^+(h)\right].
\]
We have \(D_\tau\le5b\), \(N''\le t\) and \(m,f,t,N_1\) are
bounded by fixed powers of \(x\). The scale inequalities give
\(100(d_P+4\nu)\le b^2\) at the sufficiently large test squares.
Let \(v_P\) be its Fock vector, with the splitting characters
\(\upsilon^m\) and the fixed carrier character, and put
\(C_P(h,u)=\langle v_P,\Omega(h,u)v_P\rangle\).
In particular \(C_P(1,1)=\|v_P\|^2>0\).

Use the element
\(c=N_1/\bar\pi_G^2\in E\). It satisfies
\(c\alpha\,\mathfrak M^2\widehat{\mathcal O}_E
\subset\widehat{\mathcal O}_E\) since \(N_1\alpha\) is integral,
and \(|\tau c|=N_1/G_0\). The real Heisenberg coordinate at \(\tau\)
is scaled by \(\sqrt{2\pi\operatorname{Im}\tau\mathfrak j}\),
where \(\mathfrak j=-(\sqrt{-d})^{-1}\). The Hom minimum gives
\[
 b^6|\tau c|\le
   \sqrt{2\pi\operatorname{Im}\tau\mathfrak j}\,
       \lambda_{\operatorname{Hom}} .
\]
Indeed \(G_0\ge b^{20}\), and
\(6Q_E(ft)^2t^2 N_1/\sqrt{2\pi/\sqrt d}\le b\) at these scales.
The indices of the finite datum below are residue matrices for
\(\mathfrak M^2\). Let \(\psi_T(X)\) be their additive trace
characters, write \(\Theta_T\) for their kernels with polynomial
\(P\), and set
\[
 \Psi=\frac{1}{\#\{T\}}\sum_T
       \overline{\Theta_T(z_S,z\oplus1)}\,\psi_T .
\]
Only the two base points enter these coefficients.

To estimate its kernel, expand each product of two theta values
as a sum of matrix coefficients of rational Heisenberg translates.
In a Schr\"odinger polarization this identity follows by expanding
both theta sums, writing the difference of their rational position
variables as a translation, and applying Poisson summation to the
remaining variable with the self-dual measures. The dual rational
variable gives the modulation coordinate. The Weyl half-phase on
rational pairs is one by the global additive character. Transport
by the common adelic Weil element moves the rational vectors
into the Hom frame, and its scalar phases cancel between the
conjugate pair. The mean over \(T\) kills every term except those
whose \(\alpha^{-1}\)-scaled finite coordinates are integral and
zero modulo \(\mathfrak M^2\). Thus the zero vector contributes
\(C_P(h,u)\). For all other surviving rational \(y\), \(cy\)
is in the moved Hom lattice, and the finite vacuum factor has
absolute value one. The real Heisenberg coordinates \(Y\) of
these \(y\) have norm at least \(L=b^6\). Their pairwise
distances are also at least \(L\), by the same argument on a
difference.

For clarity, in \(e=fss'\) complex Fock variables, let \(A=d_P+e\)
and \(q=2\nu>e\). For a real Weil operator attached to the matrix
\(\sigma\), the translated coefficient at \(Y\ne0\) satisfies
\[
 |\langle v_P,\rho(Y)\omega(\sigma)v_P\rangle|
 \le \|v_P\|^2
       \|\sigma\|^{e+2q}[4(A+q)]^q\|Y\|^{-2q}.
\]
Indeed the vacuum integral bounds the coefficient by
\(\|\sigma\|^e\int |v_P(X)v_P(\sigma X+Y)|\,dX\).
For \(a=(1+q/A)^{1/2}\), homogeneity and the Gaussian give
\[
 |v_P(X)|=a^{d_P}|v_P(X/a)|
             \exp\!\left(-\frac{q\|X\|^2}{2(A+q)}\right).
\]
Use \(\|Y\|^2\le2\|\sigma\|^2
(\|X\|^2+\|\sigma X+Y\|^2)\), which follows from the
triangle inequality and \(\|\sigma\|\ge1\).
Bounding the product of the two dilated functions by half
the sum of their squares costs \(a^{2(d_P+e)}\le \exp(q)\)
in the integral, since \(\sigma\) preserves real volume.
The Gaussian factor is at most
\([4(A+q)\|\sigma\|^2/(\exp(1)\|Y\|^2)]^q\)
by \(t\exp(-t)\le\exp(-1)\).
This proves the bound. In the representation on Hom coordinates
\(\|\sigma\|\le\|h\|\|u\|\).

Disjoint balls of radius \(L/2\) show that at most
\((1+2R/L)^{2e}\) of the points have norm at most \(R\).
Integrating this count against \(2qR^{-2q-1}\) for \(R\ge L\)
gives \(\sum\|Y\|^{-2q}\le(5/L)^{2q}\), since
\(2q\ge2e+2\) and \(3^{2e}\,q/(q-e)\le 5^{2q}\).
Consequently at the same two points the entire error is bounded by
\[
 \left|\Theta_{\Psi,P}(z_Sh,(z\oplus1)u)-C_P(h,u)\right|
 \le b^{-20\nu}\|h\|^{6\nu}\|u\|^{6\nu} C_P(1,1).
\]
Here \(e+2q\le 3q\) and
\(100(d_P+2q)/L^2\le b^{-10}\). The finite trigonometric sum
and the summable majorant justify the averaging.
\footnote{\weightedpin{grh-compact/src/New/Rows/SparseThetaSum.lean}
{96-L108}{Minimum of the moved Hom lattice} and
\weightedpin{grh-compact/src/New/Rows/SparseThetaSum.lean}
{133-L171}{Averaging the actual residue characters};
\weightedpin{grh-compact/src/New/Weil/FockCoeffDecay.lean}
{11-L35}{The homogeneous oscillator vector and the dilation bound};
\weightedpin{grh/src/New/OneLift/ComparisonOperator/ProductsSpan.lean}
{116-L170}{The element clearing the finite coordinates and the sparse scale}.}

\paragraph{The carrier integral in the comparison span.}
\teamanchor{proof.weighted-comparison-carrier-approximation}%
Keep the carrier point \(z_S\) fixed. Put
\(\iota h=h\oplus1\) on the padded space. For a weight \(a\) write
\(m_a\) for its product of normalized holomorphic principal-minor
functions. For the following three weights the two entries apply
to every positive and every negative coordinate of \(S\), respectively:
\[
 k_1=(-\Delta',6m+1),\qquad
 k_2=(\Delta',D_\tau+r+2m),\qquad
 k_3=(0,D_\tau+N''+2m).
\]
At each place the positive minors of \(P\) use the whole block of
\(2m\) positive columns and the negative minor uses the whole block
of \(m\) negative columns. Linear substitution on the source
therefore multiplies each such minor by the determinant of that
block of the complex-linear part. In the Fock coefficient this
gives
\[
 C_P(u,1)=C_P(1,1)\,m_{k_2}(u)\overline{m_{k_3}(u)} .
\]
Indeed the two powers of the positive and negative block determinants
from the polynomial are \(\sum_h k_\tau^+(h)\) and \(D_\tau+N''+m\).
The vacuum for \(\widetilde W_j\) has \(m\) negative coordinates
and \(r+2m\) positive ones. Including its splitting weight \(m-D_\tau\)
gives the displayed weights, up to adding the same integer in all
entries of both weights. Such a shift cancels in the product
because the determinant of a unitary matrix has absolute value one.

Use the same real Haar measure throughout, and put \(B=5b\),
\(H=B^{7\nu}\), \(J=\int_{U(S)(\mathbf R)}|m_{k_3}(u)|^2\,du\).
Then \(0<J<\infty\) and
\[
 \int_{U(S)(\mathbf R)}\|u\|^{6\nu}|m_{k_1}(u)|\,du\le HJ .
\]
Here is the estimate at this rank. The weights \(k_1,k_3\) have
constant compact blocks, their gaps are at least \(300\nu\)
by \(\Delta'+6m+1\ge300\nu\) and \(D_\tau\ge\Delta'\), and
\(\max |k_3|\le B\) by the bound on the column counts.
The squared \(k_3\)-coefficient has integral the reciprocal formal
degree at each place. In a Cartan decomposition the modulus of
the \(k_1\)-coefficient is
\(|\det(u_{++})|^{-(\Delta'+6m+1)}\). The singular values of this
block include the hyperbolic cosines in the Cartan decomposition,
so \(|\det(u_{++})|\ge\max(1,\|u\|/2)\) at that place. Take
\(a_0=10s^3>8s^2\). The local decay powers can be taken as \(a_0\)
at all places and \(a_0+6\nu\) at one of largest norm. Both fit
the gap. Multiplying the local bounds absorbs \(\|u\|^{6\nu}\)
at cost \(2^{6\nu}\). Use the
\hyperlink{proof.weighted-bessel-tail-haar}{formal-degree tail integral}
with \(L=2\lfloor B\rfloor+s\), which bounds the absolute doubled
\(k_3\)-parameter. It gives a ratio to \(J\) of at most
\[
 2^{6\nu+fa_0}
 \left[8\{2000(s+L+1)^6\}^{2s^2}\right]^f\le H .
\]
Indeed the bracket is at most \(B^{17s^2}\) for the scales here,
\(s=3m\ge6\), and \(2fs^3\le\nu\). This proves integrability
on the same group on which the carrier will be reproduced.

Define a function on the real group of \(W_j\) by
\[
 K_0(h)=\int_{U(S)(\mathbf R)}
          \phi(z_Su) C_P(u,\iota h)m_{k_1}(u)\,du .
\]
The coefficient has modulus at most \(C_P(1,1)\) by unitarity,
so \(K_0\) is bounded. Its value at one is
\[
 K_0(1)=\phi(z_S)C_P(1,1)J\ne0 .
\]
In fact \((0,3m)+k_1+k_2=k_3\) since \(N''=r+9m+1\). Hence
\(\phi(z_S\,\cdot)m_{k_1}m_{k_2}\) is holomorphic of weight
\(k_3\). For such a function \(V\), averaging \(V(c_\theta u)\)
against the inverse compact weight of the ordered circle leaves
the lowest term \(V(1)m_{k_3}(u)\). Haar invariance and the same
circle character of \(m_{k_3}\) then give
\(\int V\overline{m_{k_3}}=V(1)J\). Absolute convergence here
follows from the coefficient bound and the preceding integral.
\footnote{\weightedpin{grh-compact/src/New/Rows/NaturalFormNear.lean}
{77-L103}{Real coefficient and carrier convolution};
\weightedpin{grh/src/New/OneLift/ComparisonOperator/ProductsSpan.lean}
{283-L350}{The block substitutions and weights};
\weightedpin{grh-compact/src/Lib/DiscreteSeries/Holomorphic/HolomorphicMinor.lean}
{42-L63}{Circle projection and its real reproducer}.}

For the chosen \(z\) of the sparse-kernel estimate put
\[
 K_z(h)=\int_{U(S)(\mathbf R)}
   \phi(z_Su)\,
   \Theta_{\Psi,P}\bigl(z_Su,(z\oplus1)\iota h\bigr)
     m_{k_1}(u)\,du .
\]
We next verify that its coefficients use precisely the prescribed
finite span. On \(S\) use the additional definite space
\(W^\flat=I_{6m+1}\), whose rank exceeds \(2\dim S\), with characters
\(\upsilon^m,1\). Fix an inclusion \(e\) of the \(3m\) source
coordinates and take the finite indicator \(\delta_e\) of its
matrix modulo \(\mathfrak M^2\). At infinity use
\(Q_0=\prod_\tau\det X_{\tau,e(+),+}^{\,\Delta'}\).
This vector is nonzero and holomorphic of weight \(k_1\):
it is independent of the negative letters and transforms in the
positive columns by the indicated determinant power, with the
negative vacuum weight \(6m+1\). Fix the actual class representative
\(p\) of \(W^\flat\) and its invertible short basis \(Y_p\).
Every real target translate of \(Q_0\) is in the span of the
prescribed products \(Q_{p,R}\) of determinants with rows of
\(\tau Y_p\). Indeed at each embedding its transformed rows
are linear combinations of those rows, and multilinearity in
each row of each determinant gives exactly the row choices \(R\).

For a rapid function \(v\) on \([U(S)]\) of principal level
\(\mathfrak M^2\) write
\[
 \ell_{p,R}(v)=\int_{[U(S)]}v(x)
          \Theta_{\delta_e,Q_{p,R}}(x,p)\,dx .
\]
If these finitely many functionals vanish on \(v\), the
\((\delta_e,Q_0)\)-lift of \(v\) vanishes everywhere.
Indeed the flat target is definite, so its finite class
representatives and their real translates exhaust its adelic
quotient at this level. The preceding polynomial expansion
applies at each representative, and the finite level fixes
\(\delta_e\).

Recall the sharp-data consequence of the convergent doubling
formula. For a zero lift from rank \(s\) to rank \(M>2s\) with
this finite inclusion indicator, convolution of the input by the
real coefficient of its polynomial vanishes identically whenever
that coefficient is integrable. The convergent case of Yamana's
Theorem~2 used here has a split skew-Hermitian space of rank \(2s\),
a Hermitian space \(V\) of rank \(M\) and Witt index \(w\), and a
splitting character \(\chi_V|_{\mathbb A_F^\times}
=\epsilon_{E/F}^{M}\). For \(M-w>s\) and a Schwartz function
\(\Phi\) with polynomial Fock vectors at infinity, the Siegel
section \(f_\Phi(g,z)=|a(g)|^{z-z_0}\omega(g)\Phi(0)\),
\(z_0=(M-s)/2\), gives an Eisenstein series holomorphic at
\(z_0\) with value
\[
 E(g,f_\Phi,z_0)=
 \int_{U(V)(F)\backslash U(V)(\mathbb A_F)}
        \sum_{x\in V^s(F)}(\omega(g)\Phi)(h^{-1}x)\,dh .
\]
The integral is absolutely convergent, \(dh\) is probability
Haar, and the multiplier in this range is \(1\).
Here \(w\le M/2\), so the displayed strict inequality holds.
The doubling construction evaluates this section on the two
copies of the actual source and uses the two lift vectors
to form \(\Phi\). Changing the quotient Haar from probability
to the Petersson convention contributes the common quotient
volume. The graph-orbit argument for this section now unfolds
the pairing of the two theta lifts. Each boundary Lagrangian has a nonzero
intersection with the first copy of the source. Its term
requires a matrix killing a nonzero vector there, which the
primitive-frame support excludes at a chosen split place.
For a definite source these intersections are isotropic and
therefore already zero. The graph Lagrangians give the integral
of the Weil coefficient times the right-translation Petersson
coefficient on the adelic source group. At the chosen split places
two supported integral frames related by a source matrix force
it to be integral and to induce the identity modulo
\(\mathfrak M^2\), by the left inverse of the inclusion.
The finite self-pairing there is constant on the principal
subgroup and has positive mass. At any additional finite set of
places one can cut the vacuum coefficient down to the integral
group in the second argument of the doubling pairing. For
example use \(s\) integral hyperbolic pairs in the target there.
For the positive standard lattice the pair
\((1,\eta),(1,0)-\xi(1,\eta)\) has Gram matrix
\(\left(\begin{smallmatrix}0&1\\1&0\end{smallmatrix}\right)\)
with \(\eta\bar\eta=-1,\ \xi+\bar\xi=1\) in the finite integers
of the Weil datum. In its position coordinates the first vacuum
is the indicator of integral \(s\)-by-\(s\) matrices, up to
the fixed lattice scale. Against it take the normalized indicator
of \({\rm GL}_s(\mathcal O)\) in the same coordinates, retaining
the vacuum on the complement. The source acts on the position
matrix by right multiplication. Its local Haar Jacobian
is one because the determinant of a unitary matrix has norm one
to the base field. On this test the first vacuum
evaluates to one exactly when the acted matrix is integral.
Taking conjugate places together, preservation of the Hermitian
form also makes the inverse integral. On the integral group the
coefficient agrees with the vacuum self-pairing, including its
unitary splitting character on the integral subgroup. This
argument uses the hyperbolic polarization only at the added
places. Increasing their finite
set removes all nonintegral finite translates by dominated
convergence in the absolutely convergent doubling formula.
Every pairing with the zero lift was zero before this limit.
The remaining positive finite mass multiplies the Petersson
pairing of any rapid test with the real convolution. Hence the
convolution vanishes, by density and continuity. For the original
definite source of the carrier the same argument uses the
unit-modulus finite vacuum and \(M=3m>2m\); the target has \(m\)
integral hyperbolic pairs there as well.
\footnote{The ordinary Siegel--Weil value is
\href{https://msp.org/pjm/2013/264-1/pjm-v264-n1-p11-s.pdf\#page=4}
{Yamana, \emph{The Siegel--Weil formula for unitary groups},
Theorem~2, p.~237}, with \(M-\operatorname{Witt}(V)>s\).
Our measured pairings have \(M>2s\), so
Liu \cite[Theorem~2.1(1), p.~859]{tr3liu} supplies this identity
with the probability normalization explicitly.
The graph unfolding and finite restriction are given at
\weightedpin{grh-compact/src/Lib/SiegelWeil/RallisInner/SharpData.lean}
{145-L185}{Convergent doubling for the clean datum},
\weightedpin
{grh-compact/src/Lib/Weil/CompactLift/RallisFormula.lean}
{193-L244}{Finite mass, localization and real convolution},
\weightedpin{grh-compact/src/Lib/SiegelWeil/RallisInner/SharpData.lean}
{235-L276}{Integral frame localization} and
\weightedpin{grh-compact/src/Lib/SiegelWeil/RallisInner/SharpData.lean}
{448-L472}{Sharpness of the inclusion indicator}.}

For \(Q_0\) the real coefficient is \(\|v_{Q_0}\|^2m_{k_1}\).
Indeed it has that holomorphic weight, and its value at one
normalizes the lowest coefficient by the same circle projection
and Hermitian symmetry. Thus vanishing of all \(\ell_{p,R}\)
implies \(\int v(z_Su)m_{k_1}(u)\,du=0\). Linear algebra on
the image of \(v\mapsto(\ell_{p,R}(v))\) now writes this latter
functional as \(\sum_{p,R}a_{p,R}\ell_{p,R}\), with coefficients
independent of \(v\).

Apply this identity to
\(v(x)=\phi(x)\Theta_{\Psi,P}(x,(z\oplus1)\iota h)\).
Its level is \(\mathfrak M^2\): the second splitting character
has trivial tame part and the finite datum factors through that
quotient. It is rapid because \(\phi\) is rapid and the theta
kernel has at most polynomial growth. The product under
\(\ell_{p,R}\) is the kernel on
\(\widetilde W_j\perp W^\flat\), with polynomial
\(P(X_1)Q_{p,R}(X_2)\) and finite function
\(\Psi(X_1)\delta_e(X_2)\). The primitive additive trace
pairing modulo \(\mathfrak M^2\) gives a finite Fourier expansion
of this function, say with coefficients \(b_T\), over exactly
the residue matrices in the definition of \(F_i^\natural\).
Absolute convergence permits termwise integration against the
same \(\phi\). Thus
\[
 K_z(h)=\sum_{p,R,T} a_{p,R}b_T
                  F_{(T,p,R)}^\natural(z(h,1))
          \qquad\text{for every real }h .
\]
Only \(\Psi\) and hence the \(b_T\) may vary with \(z\).
The sparse-kernel estimate, \(\|\iota h\|=\|h\|\), and
\(|\phi|\le2|\phi(z_S)|\) give
\[
 |K_z(h)-K_0(h)|\le\epsilon_0\|h\|^{6\nu},\qquad
 \epsilon_0=2Hb^{-20\nu}|K_0(1)|.
\]
Indeed \( (4H+1)\epsilon_0/|K_0(1)|
  \le 10(5^{14}/b^6)^\nu<1\).
This proves the approximation at the untwisted square level
for every point on the required minimum locus.
\footnote{\weightedpin{grh-compact/src/Lib/Weil/Fock/ReproducingFormula.lean}
{96-L129}{Stable-range reproduction on the given classes and real translates};
\weightedpin{grh/src/New/OneLift/ComparisonOperator/ProductsSpan.lean}
{411-L507}{Finite expansion on the actual flat classes and row bases};
\weightedpin{grh-compact/src/New/Rows/NaturalFormNear.lean}
{103-L146}{Use of the same coefficients for every real translate}.}

On \(W_j\) apply the same reproducing statement with the definite
small block of rank \(t-9m-1>2r\), the given inclusion datum
modulo \(\mathfrak M^2\), and the negative-letter highest vector
of the small weight \(w\). Its translates at a small class are
in the span of the negative Weyl basis there, by equivariance of
the letter substitution. A table supported at this class
with the conjugate of any such coordinate as its entry
produces the corresponding functional. To identify the type
of the approximants on \(W_j\), apply the third large-block
matching identity at twist \(1\), with the factorization
\(I_{3m}s_i=s_i\) of the prescribed member matrix. It writes
\(F_i^\natural\) as \(\kappa_1\) times a finite class sum
of kernels from \(W_j\) to the positive definite large block.
In these kernels the polynomial uses the positive letters of the
untwisted diagram. Its highest vector is holomorphic, and
substitution by all definite target translates spans its Weyl
module. The sum of its staircase weight and the negative-letter
weight \(w\) is \(k\), coordinate by coordinate. Thus each
\(F_i^\natural\) is holomorphic of weight \(k-w\). The finite
data have level \(\mathfrak M^2\) with trivial tame part off
that modulus. The kernels have moderate growth, so their products
with the rapid \(F\) give \(FF_i^\natural\) rapid at the table
level. These conclusions follow by multiplying by \(\kappa_1\),
before its nonvanishing is established.
\footnote{\weightedpin{grh-compact/src/New/Rows/Separation.lean}
{164-L182}{The third identity at twist one on the exact member};
\weightedpin{grh-compact/src/New/Exchange/GridModuli.lean}
{67-L80}{Addition of the two staircase weights};
\weightedpin{grh/src/New/OneLift/ComparisonOperator/ThreeExchanges.lean}
{986-L1021}{The untwisted member equality, product level and weight}.}
Vanishing
of the given table periods therefore gives the real vanishing
integral in the \hyperlink{proof.weighted-comparison-detection-reduction}
{circle isolation argument} for each \(K_z\) above. Each rotation
of a good-form peak retains the same first minimum. Isolation
now forces \(F=0\) for the full good space at each \(j\).

\paragraph{Detection at the large-block products.}
\teamanchor{proof.weighted-comparison-detection-reduction}%
The
\hyperlink{proof.weighted-comparison-carrier-approximation}{carrier approximation}
gives functions \(K_g\) in the span of the prescribed
\(h\mapsto F_i^\natural(g(h,1))\), on the first-minimum
locus, with
\[
 |K_g(h)-K_0(h)|\le \epsilon_0\|h\|^{6n_*},
 \qquad |K_0(1)|>(4B^{7n_*}+1)\epsilon_0 .
\]
The function \(K_0\) is bounded and independent of \(g\).
Let \(k\) be the holomorphic lowest weight of \(U=\bar F\)
and \(w\) the weight of the small kernel. By the untwisted
third exchange in the preceding argument, each \(K_g\)
is holomorphic of weight \(k-w\). If all the table periods
vanish, the preceding small-block reproducing calculation
gives
\[
 \int_{U(W_j)(\mathbf R)}
     F(g(h,1))K_g(h)m_w(h)\,dh=0 .
\]
Here \(m_w\) is the holomorphic principal-minor function
normalized to equal \(1\) at the identity.

Suppose \(F\ne0\) and choose \(g_0\) with
\(|U(g_0)|>\tfrac12\sup|U|\).
By the \hyperlink{proof.weighted-comparison-peak}{first-minimum bound},
\(g_0\) lies on the required locus. Every ordered compact
diagonal rotation \(c_\theta\) preserves the real lattice
norm, so put \(K_\theta=K_{g_0c_\theta}\) there.
Divide \(U(g_0c_\theta(h,1))\) by its unit weight character
on the rotation. Its circle mean is \(U(g_0)m_k(h)\), by
the lowest-term projection of the holomorphic
principal-minor model. Write \(I_k=\int |m_k(h)|^2\,dh>0\).
The weights \(k,w\) have dominant compact blocks, signed
gaps at least \(300n_*\), and \(\max |k|\le B\). The real
minor integral bound on these weights is
\[
 \int_{U(W_j)(\mathbf R)} \|h\|^{6n_*}|m_w(h)|\,dh
          \le B^{7n_*} I_k ,
\]
where the integral is over the real unitary points with
the same Haar measure as the reproducing identity.
Equivalently one may apply the norm integral bound for
a function bounded by \(C\|h\|^{6n_*}\).
In the vanishing identity at \(g_0c_\theta\) replace
\(K_\theta\) by \(K_{g_0}\). Its integral changes by at
most \(4|U(g_0)|\epsilon_0 B^{7n_*}I_k\), since each
approximation differs from \(K_0\) by at most
\(\epsilon_0\|h\|^{6n_*}\). Now the integrand with
\(K_{g_0}\) is continuous in \(\theta\) and has an integrable
majorant. Average only this fixed function times the
rotated translates. The circle mean and real reproducing
identity for the weight \(k\) applied to \(K_{g_0}m_w\)
give \(\overline{U(g_0)}K_{g_0}(1)I_k\) for that average.
Thus \(|K_{g_0}(1)|\le4\epsilon_0 B^{7n_*}\) and
\(|K_0(1)|\le(4B^{7n_*}+1)\epsilon_0\), a contradiction.
The table periods therefore detect every good form on
this fixed untwisted span. The integration uses the
single fixed \(K_{g_0}\).
\footnote{\weightedpin{grh-compact/src/New/Exchange/GridModuli.lean}
{78-L87}{The actual weight inequalities};
\weightedpin{grh-compact/src/New/Rows/SparseThetaSum.lean}
{27-L84}{Minor integral and circle averaging}.}

\paragraph{The common adjoint and the finite tests.}
\teamanchor{proof.weighted-comparison-adjoint}%
The glued indicator identity gives a product of the two block
kernels at \(g\widetilde Y\) for each of its terms. The Fock test
splits into the two prescribed polynomials because the mixed
Cartan tensor is supported on the disjoint blocks. The adelic
point \(\gamma_i\) is the identity at the twist places, so its
Weil scalar \(\omega_i\ne0\) uses only the fixed source character
and the chosen split places. The representatives \(\widetilde Y\)
are principal modulo \(\mathfrak M\) there and trivial at the
twist places. They fix the good source forms. In the Petersson
pairing change \(g\) to \(g\widetilde Y\) in each term. Then they
all have the same value, and the multiplicity is
\(\#\{Y:Y\equiv I_r\bmod\mathfrak M\}\). This calculation uses
the same class normalization as
\hyperlink{proof.period-row-models}{the earlier class formula}.
The cuspidal convergence of the block periods permits the finite
sums to be taken outside the integral.

Write \(F_i^\natural\) for the fixed large-block functions and
\(\mathcal T_i(\xi)\) for the small-block theta linear in the
\(i\)-th table row. The matched large-block identity gives
\(F_i^\natural=\kappa_\psi\) times the large-block sum before its
normalization, with
\(\lvert\kappa_\psi\rvert^2=C/\Gamma_j(\psi)\) for a common \(C\).
The unfolding just described gives, with the second entry linear,
\[
 \overline{\kappa_\psi}\langle F,u_\psi^*G_\psi\xi\rangle
   = C_j\sum_i\overline{\omega_i}
         \int_{[U(W_j)]}\overline{F F_i^\natural}\,
                                      \mathcal T_i(\xi) ,
\]
where \(C_j\ne0\) contains \(g^{L_j}\), the conjugated stage
multiplier and the congruence multiplicity. This right-hand side
is independent of \(\psi\). Use the detection statement for the
actual products \(F_i^\natural\mathcal T_i\). If some
\(\kappa_\psi=0\), the matched identity would give
\(F_i^\natural=0\) for every \(i\), and detection would force the
good source space \(U_j\) to be zero. Its dimension is \(R>0\).
Similarly, if \(F\) is orthogonal to \(u_1^*G_1\) of all tables,
take a table supported on one \(i\). Its right-hand period
vanishes for every small-block table. Detection gives \(F=0\).
Thus the common map \(A=u_1^*G_1\) is onto, and the scalar is
\(\gamma_\psi=\overline{\kappa_1/\kappa_\psi}\). It satisfies
\(\lvert\gamma_\psi\rvert^2=\Gamma_j(\psi)/\Gamma_j(1)\).

\paragraph{Two applications of the marked period.}
\teamanchor{proof.weighted-marked-period-reduction}%
The marked period for the comparison form integrates over
the small block of rank \(t-9m-1\), with the block of rank
\(r+9m+1\) held fixed. A false form \(Y\) is good and
orthogonal to the good-lift range. Its theta descent is good
by equivariance at the chosen places. Adjunction then makes
the descent zero. The reflection sums at these places kill
\(Y\): their normalized sums are projectors on the determinant
character of the rank-one quotient of the corresponding
hyperplane stabilizer. Under the finite parabolic induction
rule this component can occur only in a Green label with
linear factor named by that character. The defining
good-label condition excludes that factor for the
\(\omega^j\) on the actual marked block. At a new twist place
the character is nontrivial by definition of the conductor.
The required rank gap is
\(2r+2(r+9m+1)+3<6r\), which follows from \(r\gg m\).
At a chosen place the small-block principal level modulo
\(\mathfrak M^2\) fixes the small kernel. Conjugation of such
an element by \(\gamma_i\) changes the identity by a matrix
in \(\varpi M_{r+t}(\mathcal O_v)\), because both
\(\gamma_{i,v}\) and \(\varpi\gamma_{i,v}^{-1}\) are integral.
At the conjugate place use the conjugate inverse transpose.
At twist places \(\gamma_i\) is the identity and the
small-block action keeps the marked rows in the fixed block.
These conditions give the same compact stabilizer required
by the period, including its determinant character.
\footnote{\weightedpin{grh-compact/src/New/Trace/GoodIdempotent.lean}
{166-L174}{Reflection sums on the good projector};
\weightedpin{grh-compact/src/New/Rows/MemberOrthogonality.lean}
{122-L138}{Application at the small block}; and
\weightedpin{grh-compact/src/New/Rows/GluedFormTheta.lean}
{172-L190}{The conjugate congruence}.}

For the rank-one point test, keep the first three blocks
fixed and expand the flat lift on the last block into its
rank-one source classes. For each such class the last-block
sum is a marked period of the same form. The compact rank
inequality is \(2+2(12+4+4)+3=45<2z_p\).
Here the false form is in the trivial-label summand at all
chosen places and the descent is zero by its defining test.
The reflection sum on \({\rm GL}_1(k_v)\) is the sum of
\(\omega_v^{-j}\) over \(k_v^\times\) and vanishes on this
summand. At twist places the rank-one character has the
nontrivial tame part of \(\psi\). At chosen places the last
block is at the level fixing \(J_1\), which fixes both the
flat column datum and the full point form's column
\((x_0,J_2,J_3,J_1)^{\mathsf t}\). The determinant character
is read in the second block and stays fixed. The Cartan
contraction against the flat form is equivariant because the
flat coefficient type is trivial and the first twelve
columns are fixed, as in the point-form construction.%
\footnote{\weightedpin{grh-compact/src/New/Selmer/HeckeStableSubspaces.lean}
{31-L47}{The actual rank-one labels};
\weightedpin{grh-compact/src/New/Rows/HeckeValueHeight.lean}
{122-L146}{The marked period on the last block}.}

Both applications use the same marked-period implication.
Suppose that \(Y\) has zero normalized marked descent, the
chosen-place reflection sums vanish on \(Y\), and the tame
character is nontrivial at every new twist place. Assume the
off-level split Hecke equation at
\((a_{\mathrm{Frob},r}(w)q_w^{n-r-1})^{-1}\), with
\(a_{\mathrm{Frob},r}(w)\) the rank-\(r\) Frobenius value of
the source splitting character. The datum is vacuum off the
chosen and new twist pairs. When the fixed block of rank \(b\)
satisfies \(2r+2b+3<n\), the equivariant period over its definite
complement vanishes on the observed residue frames. The proof
uses the normalized Weil actions on these data and the mixed
zero-term formula in the
\hyperlink{proof.weighted-seesaw-zero-term}{seesaw calculation}.

\paragraph{From the marked section to the allowed descent data.}
\teamanchor{proof.weighted-marked-descent-bottom}%
Write \(A=\mathscr M\cup Q\) for the chosen and new twist places of
the CM field \(E\). These are split and pairwise distinct even up to
conjugation. Let \(F=E^+\), \(V=I_n\), and let the lower space \(W\)
have rank \(r\). Use the marked column frame \(I_{\rm mk}=\binom{I_r}{0}\)
in each indicated ordering of the rows. Write
\(Z_v=\binom{Z_{0,v}}{Z_{1,v}}\) in these coordinates. Let
\(\chi_v\) be the residue character of the actual auxiliary pair.
The marked finite function, with vacuum at the other finite pairs, is
\[
 \Psi_0(Z)=
 \prod_{v\in A}1_{M_{n,r}(\mathcal O_v)}(Z_v)
 \prod_{v\in\mathscr M}1_{\overline Z_v=I_{\rm mk}}\,
 \prod_{v\in Q}\chi_v(\det\overline Z_{0,v}),
\]
where the residue character is extended by zero at zero. The level
\(K_A\) fixes the marked columns modulo \(\mathscr M\) and is block
lower triangular modulo \(Q\). Its character on the form is
\(\lambda(k)=\prod_{v\in Q}\chi_v(\det\overline k_{00,v})^{-1}\).
The vacuum condition on a datum means the vacuum at every finite
pair outside \(A\cup\overline A\). At a designated split place use
the Schr\"odinger coordinates \(M_{n,r}(E_v)\), with the other
place of the pair fixed by the unitary relation.

Let \(y\) be the algebraic form whose marked descent vanishes,
with real Weyl type \(\rho\) on \(L\).
For a linear functional \(u\) on \(L\), use the conjugate
scalar coefficient \(Y_u(h)=
\overline{u(\rho(h_\infty^{-1})y(h_f))}\).
For a real Schwartz vector \(P\) and \(g\in
U(W)(\mathbb A_F)\), write \(D(\Phi,y)\) for the descent of \(Y_u\)
at \(g\) with section \(P\otimes\Phi\otimes{\rm vac}\) and splitting
characters \((a,1)\). This pairing is linear in \(\Phi\) and
conjugate-linear in \(y\). The defining marked class sum agrees with
the conjugate of the adelic descent on the highest harmonic and
highest coefficient, multiplied by
\(\dim L\,[K_0:K_A]\). Here \(L\) is the actual archimedean Weyl type
on \(V\), with the highest vector of unit Gram norm. Indeed, on each
class the compact average contracts the two \(L\)-coordinates by
Schur orthogonality with scalar \(1/\dim L\). This is the usual
compact integral of \(\langle u,\rho(k)x\rangle
\overline{\langle u,\rho(k)y\rangle}\). The simultaneous finite
actions have reciprocal phases, and the class measure is
\([K_0:K_A]^{-1}/|\operatorname{Aut}|\) with the original
volume-one compact factors. Erasing the second splitting character
in the descent only removes its character of \(\det g\), which is
nonzero. Thus the highest harmonic gives zero at every \(g\).

The same holds for every real Schwartz vector with the same marked
finite section and for every coefficient \(u\). Identify each
coefficient by the Gram form with \(u=u_\zeta\),
\(u_\zeta(v)=\langle\zeta,v\rangle_L\), linear in \(v\).
Then \(\zeta\mapsto Y_{u_\zeta}\) is linear. For the compact
action \(R(k)\) it takes \(\rho(k)\zeta\) to
\(R(k)Y_{u_\zeta}\), since the coefficient at \(hk\) reads
\(\rho(k)^{-1}\rho(h_\infty^{-1})y(h_f)\).
The bilinear integral is invariant under simultaneous right
actions of \(U(V)(F_\infty)\) on the oscillator and this
coefficient. In the polynomial oscillator it therefore factors
through contraction of \(L\) with \(L^\vee\). At each real place
Howe's compact-member theorem gives irreducibility of the
multiplicity space of \(L^\vee\) for the real \(W\)-side
Harish--Chandra actions. Write \(J\) for the tensor product
of those multiplicity spaces.
In the marked polynomial the positive minors substitute the
transpose of the \(U(n)\)-matrix and the negative minors its
inverse. The first splitting character is \(a=\upsilon^{r_-}\)
with \(r_-\) the number of negative columns of \(W\).
The negative-column diagram contributes determinant exponent
\(-\sum_{i=1}^{n-1}c^-_{\tau,i}\) at \(\tau\), where
\(c^-_{\tau,i}\) counts the negative columns of height \(i\).
These are the determinant powers of \(L\), and the substitutions make
\(\bar\zeta\mapsto P_{\bar\zeta}\) an intertwiner from \(L^\vee\).
It is nonzero on the paired highest line, which is the product of
the staircase minors. The compact contraction of this polynomial
with \(\zeta\) at the paired highest vector therefore gives a
nonzero vector of \(J\). The same vector generates \(J\).
Moving each real \(W\)-action to the evaluation point \(g\), and
then differentiating, gives zero on the generated polynomial
module. The polynomial Fock vectors
correspond to finite Hermite sums, dense in real Schwartz space.
Continuity of the integral on this space follows from the theta
lattice sum over the compact quotient of \(U(V)\). This proves
\(D(\Psi_0,y)=0\) with all \(P,u,g\).%
\footnote{Use Howe \cite[p.~542, (3.9)(a),(d)]{p2flip-howe-real}
for the compact-member theorem.
\weightedpin{grh-compact/src/New/Forms/AuxiliaryPair.lean}
{118-L179}{Harmonic polynomial, compact actions and class normalization}.}

For any \(B=(B_v)_{v\in A}\) with \(B_v\in{\rm GL}_r(E_v)\), embed \(B\)
as an element of \(U(W)(\mathbb A_{F,f})\) on those split pairs.
The source Schr\"odinger action on the marked datum replaces
\(\Psi_0(Z)\) by \(\Psi_0(ZB)\) times a nonzero scalar.
This follows from change of variables \(Z\mapsto ZB\) in the
Schr\"odinger model. The scalar consists of the splitting
character and \(\prod_v|\det B_v|_v^{n/2}\).
Right translation of \(g\) by this element then gives
\(D(\Psi_0(\,\cdot\,B),y)=0\).
For the integral target elements \(k\) at \(A\) the same model
acts as \(\Phi(Z)\mapsto\Phi(k^{-1}Z)\) since \(a\) is an
unramified splitting power at these places. Haar invariance on
\([U(V)]\) gives
\begin{equation}
\label{eq:weighted-marked-finite-covariance}
 D(\Phi(k^{-1}\,\cdot),R(k)y)=D(\Phi,y).
\end{equation}
The character projection by \(K_A\), acting on functions by
\(\overline{\lambda(k)}\Phi(k^{-1}Z)\), thus leaves
\(D(\Phi,y)\) unchanged. The projection uses compact Haar
probability, or the identical average over a finite quotient
fixing the function on its compact support.

Fix a centre \(Z\) and a product of small matrix boxes \(C\) about
it. Suppose first that all \(Z_{0,v}\) are invertible and the
ratios \(Z_{1,v}Z_{0,v}^{-1}\) are integral, in the maximal
ideal at the chosen places. With \(B_v=Z_{0,v}^{-1}\),
the function \(\Psi_0(\,\cdot\,B)\) is nonzero at \(Z\).
Its support is a single \(K_A\)-orbit. After multiplying the
argument by \(B\) on the right, a point of the support is
\(\binom{X_0}{X_1}\) with \(X_0\in{\rm GL}_r(\mathcal O_v)\).
The integral completion
\(\begin{pmatrix}X_0&0\\X_1&I\end{pmatrix}\) takes \(I_{\rm mk}\)
to this point, fixes the marked columns modulo the chosen
place, and is block lower triangular modulo a twist place.
The orbit is compact and open in the product of matrix spaces.
The marked function transforms on it by
\(\Psi_0(kZB)=\overline{\lambda(k)}\Psi_0(ZB)\).
Take \(C\) contained in this orbit. The character projection of
\(1_C\) is supported there and has the same transformation rule.
Transitivity makes it a scalar multiple of \(\Psi_0(\,\cdot\,B)\).
Thus \(D(1_C,y)=0\).

If the ratio condition fails at \(v\), there is an integral row
\(\ell=(\ell_0,\ell_1)\) with
\[
 \ell Z_v=0,\qquad \overline\ell_0\ne0,\qquad
 \overline\ell_1=0\quad(v\in Q).
\]
For singular \(Z_{0,v}\), scale a nonzero left null row of that
block to have integral entries and a unit entry, and append
zeros. Otherwise take a row \(a\) of \(Z_{1,v}Z_{0,v}^{-1}\)
with an entry of maximal modulus at least one at a chosen place
or greater than one at a twist place. Append the negative of the
corresponding coordinate row in the lower block and divide
by the entry of maximal modulus. This gives the three asserted
properties.

For every \(s=1+t\overline\ell_0\in{\rm GL}_r(k_v)\), lift the column
\(t\) integrally and put \(k_t=1+I_{\rm mk}\widetilde t\,\ell\).
The determinant of \(k_t\) is a unit because its residue is
\(\det s=1+\overline\ell_0t\). It fixes \(Z_v\) and hence every
matrix box about \(Z_v\). At a chosen place \(k_t\) differs from
the lift of \(\operatorname{diag}(s,1)\) by a level element of
character one. Writing \(\pi_v(s)\) for the chosen action
twisted by \(\chi_v(\det s)\), we get
\(R(k_t)y=\chi_v(\det s)^{-1}\pi_v(s)y\).
Equation~\eqref{eq:weighted-marked-finite-covariance} makes each
term of the weighted reflection sum
pair with \(1_C\) to \(D(1_C,y)\). The sum kills \(y\), so its
nonzero number of terms times \(D(1_C,y)\) is zero.
At a twist place \(\overline\ell_1=0\), so \(k_t\) belongs to
the level and acts by \(\chi_v(\det s)^{-1}\).
The map \(t\mapsto 1+\overline\ell_0t\) takes all nonzero
residue values among these \(s\), and the character is nontrivial.
Equation~\eqref{eq:weighted-marked-finite-covariance}
again gives \(D(1_C,y)=0\).
Every locally constant compactly supported function at \(A\)
is a finite linear combination of such indicators. In the good
case one can shrink boxes inside the orbit, using compactness of
the support and disjoint refinement of the finitely many boxes.
This proves vanishing on all allowed descents to \(W\).%
\footnote{\weightedpin{grh-compact/src/New/SiegelWeil/AuxPairVanishing.lean}
{12-L56}{Split matrix congruences and reflection sums};
\weightedpin{grh/src/New/OneLift/FalseLifts/BottomOfTower.lean}
{1226-L1491}{Finite character projection and off-orbit annihilation}.}

\paragraph{Integral frames and adapted compact subgroups.}
The field contains a primitive \(5040\)-th root of unity. At
places of residue characteristic different from three a primitive
cube root has distinct reduction from its conjugate. At the places
over three the fourth roots \(i\) and \(-i\) have distinct reductions.
Thus conjugation cannot lie in the inertia subgroup of any
nonsplit quadratic local extension of \(E/F\), and such a local
extension is unramified.
In an unramified quadratic ring every unimodular Hermitian lattice
has an orthonormal basis. In its residue form some vector has
nonzero norm, by nondegeneracy, polarization and the surjective
residue trace. The unit norm map is surjective, first on the
residue field and then on the principal units by lifting with the
residue trace. Rescale the lifted vector to norm one and repeat on
its orthogonal complement. This also extends every integral
orthonormal frame. At a split place such a frame gives a matrix
inclusion and its left inverse over the valuation ring. Extend
its basis by a basis of the kernel of the left inverse and use
the inverse transpose at the paired place.

For an isotropic line at a nonsplit finite place choose a primitive
integral vector \(v\) in it. The unimodular pairing gives an
integral partner \(w\) with \(\langle v,w\rangle=1\).
Replace \(w\) by \(w-cv\), where
\(c+\bar c=\langle w,w\rangle\), to make \(w\) isotropic.
The trace onto the base valuation ring is surjective for the
unramified extension, also at the dyadic places. The orthogonal
complement of the resulting integral hyperbolic plane is unimodular,
so the previous basis argument identifies it with the complement
for the adapted plane. At a split place the isotropic line gives
a line and its annihilating hyperplane in one factor. Integral
bases for the saturated flag give transitivity under the integral
compact. At infinity the positive and negative parts of an
isotropic vector have the same norm, so the product of the two
definite unitary groups is transitive on the isotropic lines.
Factoring the stabilizer of the line into its Levi and unipotent
radical gives the adapted Iwasawa decomposition used below.

\paragraph{Bounded descents up the hyperbolic tower.}
\teamanchor{proof.weighted-marked-descent-bounded}%
Let \(W_j=W\oplus\mathbb H^j\), \(a_j=a\upsilon^j\).
Suppose the allowed descents to \(W_{j-1}\) vanish. For a
fixed allowed section \(\phi\) on \(\operatorname{Hom}(W_j,V)\),
write
\[
  F_j(g)=\int_{[U(V)]}
   \Theta_{W_j,V}(\phi;g,h)\,Y_u(h)\,dh .
\]
We show that \(F_j\) is bounded. All measures on definite groups
are the original quotient measures. The additive means used here
are Haar probability on the compact adelic quotients.

Use the adapted Witt splitting of the last plane. Let \(k\) be
its designated isotropic vector and \(J_j\) the form on \(W_j\).
The isotropic space of maps
\(Z=\{d\,k^*J_j:d\in V\}\subset\operatorname{Hom}(W_j,V)\)
acts by modulation of the mixed variable. Its mean on theta,
\(\Theta_0=\int_{[Z]}\Theta\circ\rho(z)\,dz\), selects the mixed
slice at zero by additive Fourier orthogonality. This is the
normalized zero term of the lower-chain calculation, with the
same Haar and vacuum scales. For each \(g\), decompose \(g=u m k_c\)
relative to the last-plane parabolic, with \(k_c\) in the
adapted maximal compact. The zero slice of \(\omega(g,h)\phi\)
is, up to the nonzero scalar of the plane Levi \(m_1\),
the lower Weil action on the zero slice of
\(\omega(k_c,1)\phi\). The unipotent part acts trivially on
the zero slice. For the Levi, the action \(X\mapsto hXg^{-1}\)
on the isotropic map block has determinant
\(\det(h)\,\overline{\det(m_1)}^{\,n}\).
The normalized zero-term law multiplies by
\(\upsilon\) of this determinant to power \(-1\) and its
absolute norm to power \(1/2\). The \(h\)-factor changes \(a_j\)
to \(a_{j-1}\) since \(\det h\) has norm one. The scalar of
\(m_1\) is \(\upsilon(\det m_1)^n|\det m_1|_E^{n/2}\)
for second character \(1\). In the adapted finite bases the
vacuum is a product over the retained space and the split
plane. Compact action off \(A\cup\overline A\) preserves the
vacuum line, and the zero slice of the split vacuum is a
nonzero multiple of the retained vacuum. Hence the lower
section is allowed. The inductive vanishing says that
\(\Theta_0\) has zero integral against \(Y_u\) at every \(g\).

Here is the uniform estimate for the difference \(\Theta-\Theta_0\).
For a point \(g\) let \(H(g)=\max(1,\lambda_1(g)^{-1})\), with
\(\lambda_1(g)\) the first Euclidean minimum of the moved integral
lattice in normalized archimedean coordinates. We first recall
the polynomial bound for \(\Theta\) in this height and give
the reduction needed for its uniform use. All moved lattices
have one covolume by the determinant norm identity for an
adelic unitary matrix. Those with a fixed positive lower
bound on \(\lambda_1\) have uniformly bounded Euclidean
\(\mathbb Z\)-bases and basis inverses. One proof chooses a
shortest primitive vector \(v\), bounded above by the
volume pigeonhole bound. The orthogonal quotient lattice has
covolume divided by \(\|v\|\) and its shortest vector has length
at least \(\sqrt{3}\|v\|/2\), by rounding a lift by multiples
of \(v\). Induction gives a bounded basis of the quotient and
bounded lifts. The bounded covolume and basis then bound the
inverse. In these bases the matrices of an integral basis of
\(\mathcal O_E\) and the Hermitian Gram entries take finitely
many values. They are integral and bounded at infinity. For a
fixed table of these values, the map between the two bases is a
rational \(E\)-linear isometry. It sends one lattice to the
other, so the moved finite matrices differ modulo the fixed
integral compact. The real coordinates and inverse coordinates
are bounded in this identification. This proves cocompactness
of the bounded-height part modulo rational points. In particular
the definite group has a compact quotient and admits, at every
\(h\), an \(E\)-basis \(d_1,\ldots,d_n\) of integral vectors of
uniformly bounded moved lengths.

Two moved integral vectors of length less than one have
Hermitian product zero. Indeed its value is an algebraic
integer of absolute value less than one at every embedding.
At each real place the span of these short vectors is thus
totally isotropic. For the diagonal signature matrix \(J\),
an orthonormal basis \((v_i)\) of this span and the vectors
\((Jv_i)\) are mutually Euclidean orthogonal isotropic frames
dual for the Hermitian form. The real unitary dilation with
factors \(2,1/2\) on these frames and \(1\) on the complement
doubles the short lengths and shrinks no length by more than
two. Archimedean right translation by its inverse changes the
minimum to at least \(\min(2\lambda_1,1/2)\).
At most \(\log_2 H(g)+O(1)\) such steps bring the point to the
bounded-height region. Thus a rational representative lies in
a fixed compact set times archimedean matrices whose norms and
inverse norms are \(O(H(g))\).

For any fixed finite Schwartz factor \(\phi_f\), the theta
distribution in a rational Schr\"odinger polarization is a
lattice sum and is bounded by a fixed real Schwartz seminorm.
The finite compact representatives give finitely many translates
of \(\phi_f\) since it has an open stabilizer. One can take the
real seminorm \(p_s\) to be the sum of \(L^2\)-norms of products
of at most \(s\) infinitesimal Heisenberg coordinate operators.
For \(s\) large this controls the polynomially weighted
supremum norm used in the lattice sum, by Euclidean Sobolev
embedding. Weyl covariance moves each coordinate operator
through the real Weil action by the corresponding inverse
symplectic matrix. Unitarity then bounds \(p_s\) of its
translate by a fixed power of the matrix and inverse norms
times \(p_s\) of the vector. The matrix on the space of maps
is \(X\mapsto hXg^{-1}\). Rational invariance, the previous
reduction, and compactness for the definite variable therefore
give
\[
 \big|\Theta(\omega(g,h)(f_\infty\otimes\phi_f))\big|
   \le C H(g)^{B}p_s(f_\infty)
\]
uniformly in \(h,g\), with \(B,s,C\) depending on the fixed
finite factor and the spaces. The constant works on every
real vector with this finite factor.

When \(H(g)>2\), a shortest integral vector is isotropic
by the same algebraic-integer argument. Rational Witt extension
takes it to the fixed \(k\) in the last plane. Replace \(g\) by
this rational translate, so that \(k\) has moved length
\(H(g)^{-1}\) and integral finite coordinates. The rank-one
maps \(d_i k^*J_j\) form an \(E\)-basis of \(Z\).
Their moved coordinates are integral at finite places and
\(O(H(g)^{-1})\) at infinity, by multiplying the coordinates
of \(d_i\) by the Hermitian dual coordinates of \(k\).

Take \(\phi=f_\infty\otimes\phi_f\); finite sums of such tensors
suffice. Choose a positive integer \(M\) such that finite
Heisenberg translations by \(M\widehat{\mathcal O}_E\) in
each normalized matrix coordinate fix \(\phi_f\).
This is possible by smoothness of the finite oscillator.
Represent each coefficient in \(\mathbb A_E/E\) by an adele
with finite part in \(M\widehat{\mathcal O}_E\) and bounded
archimedean part. Additive approximation first clears the
finite congruences, after which reduction by \(M\mathcal O_E\)
bounds the archimedean part. For \(z\in[Z]\) use these representatives
in the displayed basis of \(Z\). After applying Weil covariance
to move \(\rho(z)\) across \(\omega(g,h)\), its finite part
fixes \(\phi_f\), and it is a real Heisenberg translation of
size \(O(H(g)^{-1})\). This uses the same normalized product
coordinates at every place.

For an integer \(N\ge1\), rational invariance and the
isotropy of \(Z\) give
\[
 (\Theta-\Theta_0)(\omega(g,h)\phi)
  =\int_{[Z]^N}\Theta\left(
     \prod_{i=1}^N (1-\rho(z_i))\,\omega(g,h)\phi\right)\,d\mathbf z .
\]
Indeed each nonempty subset of the factors averages the
translate by a Haar-uniform point of \([Z]\), and the
alternating coefficients of the nonempty subsets sum to \(-1\).
In this identity rational changes of any representative have no
effect, since the shifts on \(Z\) commute and rational shifts
fix theta. For a small real Heisenberg shift \(v\),
\(\rho(v)-1=\int_0^1\rho(tv)\,d\rho(v)\,dt\).
Each factor costs one order of \(p_s\) and gains \(O(\|v\|)\).
Applying the theta bound with the same finite factor gives
\(O(H(g)^{B-N}p_{s+N}(f_\infty))\). Choose \(N\ge B\).
For \(H(g)\le2\), compact reduction bounds the kernel and the
lower zero-slice integral is still zero. Rational invariance
of the ordinary descent permits the chosen translate when
\(H(g)>2\). Now integrating the difference proves boundedness
of \(F_j\), since \(Y_u\) is bounded on the definite quotient
of finite measure.%
\footnote{\weightedpin{grh-compact/src/New/SiegelWeil/DescentsVanishPlanes.lean}
{24-L104}{Parallel zero-slice and boundedness argument}.
The real seminorm argument also uses the normalized zero term
in \hyperlink{proof.weighted-seesaw-zero-term}{the lower-chain calculation}.}

\paragraph{The Hecke equation in the tower.}
\teamanchor{proof.weighted-marked-descent-hecke}%
Use the original split place \(w\) off \(A\cup\overline A\)
and let \(q=q_w\). Set \(m=r+2j\), \(e=n-m\) and use
the conjugate root
\[
 X=\overline{(a_{\mathrm{Frob},r}(w)q^{n-r-1})^{-1}},
 \qquad |X|=q^{1-n+r/2}
\]
on the conjugate scalar coefficient. The coefficient map is
conjugate-linear and commutes with the finite right actions,
so \(H_w(X)Y_u=0\). The datum is the vacuum at this pair of
places, and both relevant splitting characters are unramified.
Use the Hecke polynomial with the \(K_w\)-double-coset
coefficient \((-X)^i q^{i(i-1)/2}\), Haar measure giving
\(K_w\) mass one. Put
\[
 \alpha=a_{j,{\rm Frob},m}(\overline w),\qquad
 \beta=1_{{\rm Frob},n}(\overline w),\qquad
 X'=q^e\alpha X/\beta .
\]
The character defining \(a_{j,{\rm Frob},m}(v)\) is
\(z\mapsto a_j(z/\bar z)\upsilon(z)^{-m}\), evaluated at
the uniformizer and multiplied by \(q_v^{m/2}\).
The \(\beta\) convention is the same in rank \(n\) with
trivial norm-one character.

For the bilinear descent integral, move the whole
bi-compact double coset through the quotient integral.
The translating elements are inverted with no measure factor
because the split group is unimodular. Inverse conjugate
transpose then identifies the reversed double coset with the
positive double coset at \(\bar w\), where the compact again
has mass one. Evaluate its action on the vacuum at \(\bar w\)
in matrix coordinates \(Z\in M_{n,m}(E_{\bar w})\).
For \(U\subset k_{\bar w}^{\,n}\), choose an integral
\(\gamma_U\) of determinant valuation \(\dim U\) with
kernel \(U\) modulo the uniformizer and elementary divisors
only \(1,\varpi\). The target double-coset representatives
are \(\gamma_U^{\mathsf t}\). The Schr\"odinger multiplier
on the vacuum is
\(\alpha^{\dim U}1_{Z^{\mathsf t}\gamma_U^{-1}\ {\rm integral}}\),
from the determinant character and \(|\det\gamma_U|^{-m/2}\).
The indicator is one exactly when \(Z\) is integral and
\(U\subset\ker\overline Z^{\mathsf t}\).
The source action of \(H_w(X')\) in these opposite coordinates
uses \(\gamma_U^{-1}\) on the right, with
\(U\subset k_{\bar w}^{\,m}\).
Its vacuum multiplier is
\(\beta^{\dim U}1_{Z\gamma_U^{-1}\ {\rm integral}}\),
with the condition that \(Z\) is integral and
\(U\subset\ker\overline Z\).
These representatives come from the lattices between
\(\varpi\mathcal O^a\) and \(\mathcal O^a\), with transpose
interchanging the left and right cosets. For a kernel of
dimension \(d\), summing the coefficients by the subspace
\(q\)-binomial identity gives
\(\prod_{i<d}(1-q^i\alpha X)\), or the same expression with
\(\beta X'\) in place of \(\alpha X\) on the source.
The two kernel dimensions differ by \(e\) for every integral \(Z\).
Since \(\beta X'=q^e\alpha X\), the product on the target
splits at \(e\). The identity on the vacuum and source
equivariance now give the equation on \(F_j\)
\[
 0=\left(\prod_{i<e}(1-q^i\alpha X)\right)H_w(X')F_j .
\]
The spherical property follows from the compact action on the
vacuum at \(w\). The transfer is unchanged for arbitrary data
at other places.

At this stage \(1\le j\le b_{\rm fix}\) and \(n>2r+2b_{\rm fix}+3\).
We have \(e+m/2=n-r/2-j\le n-r/2-1\) and
\[
  m+1+e/2=(n+r+2j)/2+1\le n-r/2-1 .
\]
Thus \(q^e|\alpha X|=q^{1-j}\le1\), so every factor in the
product is nonzero. Also
\(q^{m+1}|X'|=q^{(2r+2j+4-n)/2}\le1\).
For a bounded spherical function \(f\) on rank \(m\), the
nonconstant terms of \(H_w(X')\) have total absolute
coefficient at most
\[
  \prod_{i<m}(1+q^i|X'|)-1
       \le \exp\!\left(|X'|\sum_{i<m}q^i\right)-1<1 .
\]
Here \(\sum_{i<m}q^i\le q^m\), so the exponent is at most
\(1/q\le1/2\). The coefficient of degree zero is the
identity. Taking sup norms in \(H_w(X')f=0\) forces \(f=0\).
It applies to \(F_j\) by the preceding estimate.
Induction therefore gives vanishing for all allowed data on
\(W\oplus\mathbb H^{b_{\rm fix}}\) with character
\(a\upsilon^{b_{\rm fix}}\).%
\footnote{\weightedpin{grh-compact/src/Lib/Weil/Unramified/HeckeTransfer.lean}
{30-L94}{Double cosets and residue-kernel transfer};
\weightedpin{grh-compact/src/New/SiegelWeil/DescentsVanishPlanes.lean}
{15-L53}{Supremum estimate and rank convention}.}

\paragraph{Extraction of the complementary period.}
\teamanchor{proof.weighted-marked-descent-extraction}%
Write \(V=V_1\oplus V_2\) in the actual block ordering, with
\(V_1=I_b\), \(b=b_{\rm fix}\), and \(H=U(V_2)\).
We first use the split smaller group \(\mathbb H^b\).
Let \(f\) be continuous on \([U(V)]\), right invariant under
the integral compact \(K_{\rm off}\) supported off
\(A\cup\overline A\), with every allowed descent to
\(\mathbb H^b\) zero for the character \(\upsilon^b\).
For the standard frame \(x_0=\binom{I_b}{0}\) and an arbitrary
\(h_1\in U(V_1)(\mathbb A_F)\), put
\(X(h)=h^{-1}x_0h_1\).
In the split mixed model the central element with upper-right
block \(B-B^*\) multiplies the summand for a rational frame
\(x\) by \(\psi_E(-\jmath\,\operatorname{tr}(x^*xB))\).
This follows by substituting \(-2\jmath B\) in the radical
action, whose scalar at the mixed variable \(x\) is
\(\psi_E(\tfrac12\operatorname{tr}(z x^*x))\).
The Fourier coefficient of the descent at
\(\jmath I_b\) is the mean with the multiplier
\(\psi_E(\jmath\,\operatorname{tr}B)\) on
\(M_b(\mathbb A_E)/M_b(E)\), with probability Haar.
Additive orthogonality selects exactly the frames of Gram
matrix \(I_b\). Rational Witt extension makes these
frames a single orbit with stabilizer \(H(F)\). Unfolding
gives integration over \(H(F)\backslash U(V)(\mathbb A_F)\)
with its quotient measure. At a fixed value of the source
Levi the Schwartz sums are absolutely integrable. In fact a
compact set of representatives of \([U(V)]\) gives finitely
many finite supports and uniformly bounded real Schwartz
seminorms, and the rational sum on each support is absolutely
and uniformly convergent. This justifies Fourier integration
and unfolding term by term.

Let \(\kappa\) be the scaling idele of the Weil model.
Its finite component is \(\alpha_f\). At \(\tau\) its real
component is \((2\pi\operatorname{Im}\tau(\jmath))^{-1/2}\).
Put \(\Delta=\jmath\kappa\bar\kappa\). In the split mixed model
use the finite vacuum and choose the real Schwartz vector
to be \(\exp(-\sum_{\tau,i,j}|X_{ij,\tau}|^2)\) in the variable
\(x=\kappa X\). At the finite places the two isotropic half
lattices in these coordinates are dual under
\((B,X)\mapsto\psi_E(\Delta_f\operatorname{tr}(B^{\mathsf t}X))\).
Indeed \(\Delta_f=\jmath\alpha_f\bar\alpha_f\) has the
inverse-different valuations. The normalized vacuum is fixed
by Heisenberg shifts in each of these half lattices. To check
the phase here, its finite matrix coefficient on the integral
phase-space lattice is the unit phase
\(\psi_f((c_0\jmath/2)\langle v,v\rangle)\) in the standard
finite frame, with \(c_0=(1-\jmath^{-1})/2\).
Each half lattice is isotropic, so its phase is one, including
at dyadic places. Invariance under the modulating half forces
support on integral \(X\), by self-duality, and invariance under
the translating half makes the function constant there.
The finite vector is nonzero, so the constant is nonzero.
This determines every allowed off-level finite vacuum
component and allows the same choice at the marked pairs.
Thus the combined datum at \(\kappa X\) is, up to a common
nonzero constant,
\[
w(X)=\exp\!\left(-\sum_{\tau,i,j}|X_{ij,\tau}|^2\right)
                       1_{\widehat{\mathcal O}_E^{\,n\times b}}(X_f) .
\]
A Heisenberg shift with mixed coordinates \(yk=0\) and
\(yl=\kappa\bar u\) in the adapted target frames leaves the
mixed argument unchanged. The mixed action multiplies the
datum by the character with argument
\((\jmath/2)\operatorname{tr}((\overline{\kappa\bar u})^{\mathsf t}
 (x+x))\). At \(x=\kappa X\) this is
\(\jmath\kappa\bar\kappa\sum u_{ij}X_{ij}\).
Taking \(u\) supported at infinity and the marked pairs
leaves every off-level vacuum intact. The determinant of
\(h\) contributes no
character in the split model since
\(\upsilon^b(\det h)\upsilon(\det h)^{-b}=1\).
Using the Levi matrix \(\kappa h_1\) on the \(b\)-plane
gives \(X(h)\) above and contributes the fixed nonzero
Levi scalar. Consequently vanishing of these Fourier
coefficients gives
\[
 \int_{H(F)\backslash U(V)(\mathbb A_F)}
    \psi_E\!\left(\Delta\sum u_{ij}X(h)_{ij}\right)
       w(X(h)) f(h)\,dh=0 \qquad\text{for all such } u.
\]

Disintegrate the last displayed integral through
\(H(\mathbb A_F)\backslash U(V)(\mathbb A_F)\), with the measure
induced by the original Haar measures on \(H\) and \(U(V)\).
The function to be tested there is
\(P(h)=\int_{[H]} f(th)\,dt\).
It is continuous by compactness of \([H]\).
The frame map identifies the cosets with their frames \(X(h)\).
By integral transitivity its integral support is compact modulo
\(H(\mathbb A_F)\). The function \(P\) there depends only on
the coordinates of \(X\) at infinity and the marked pairs.
Indeed if these coordinates agree, an element of \(K_{\rm off}\)
carries one full integral frame into the other. Left \(H\) and
right compact invariance then give equality of \(P\).
The induced function on the compact image of the observed
coordinates is continuous. Push the positive measure
 \(w(X(h))\,d\dot h\) to this image. The additive characters
tested above separate its points. Their linear span is a
self-conjugate algebra containing the constant function.
By Stone--Weierstrass
the induced function times this measure is zero.
The image of \(x_0\) is in its support, since \(x_0\) occurs at
\(\operatorname{diag}(h_1,1)\), integrality is open there, and
the Gaussian is positive. Thus
\(\int_{[H]}f(\operatorname{diag}(h_1,t))\,dt=0\).%
\footnote{\weightedpin{grh-compact/src/Lib/Shimura/ComplexTheta/PeriodVanishing.lean}
{21-L104}{Allowed modulation, frame measure and stabilizer}.
The Fourier orbit unfolding is the unitary version of
\href{https://arxiv.org/pdf/1501.01565\#page=6}
{Walls, Proposition 2.1.1}; the vacuum restriction is treated
here by the observed-coordinate measure.}

To apply this to the original \(Y_u\), fix an allowed
small-block datum \(\Phi_2\), its character \(b_2\) on the
source determinant, and the actual \(g,h_1\) of the marked
table. Its finite translate is \(\gamma=xA_{\rm pt}\), where
\(x\) is the divided translation and \(A_{\rm pt}\) the
integral point lift of the prescribed residue frames.
Both are supported on \(A\cup\overline A\), so \(Y_u(h\gamma)\)
has the same
off-level invariance and its allowed top descents still vanish.
This follows by right translation of the datum in the
integral, leaving the off-level components fixed. Choose a
datum \(\Phi_1\) on the first block with fixed finite vacuum
and \(\Theta_{W,V_1}^{a,b_2^{-1}}(\Phi_1;g,h_1)\ne0\).
One can choose it directly in a rational Schr\"odinger
polarization. The finite Weil translate of the vacuum is a
nonzero locally constant compactly supported function. A
rational vector meets its nonzero set by finite additive
approximation. The rational vectors on its finite support
are discrete at infinity because their denominators are
bounded. Choose a real smooth compactly supported section
after the real Weil translation that isolates the chosen
rational vector, and apply the inverse real Weil translation.
The theta sum at the given point is nonzero.

The function
\[
    f(h)=\Theta_{W,V}^{a,1}(\Phi_1\otimes\Phi_2;g,h)
               Y_u(h\gamma)
\]
has zero allowed descents to \(\mathbb H^b\) by the top
vanishing for \(W\oplus\mathbb H^b\).
In the orthogonal sum on the source the two kernels multiply,
and the product character is \(a\upsilon^b\).
Its right \(K_{\rm off}\)-invariance follows from
vacuum invariance and the support of \(\gamma\).
On the block \(V_1\oplus V_2\) the same tensor product on the
target gives the two theta factors with \(b_2^{-1},b_2\),
whose product character is \(1\).
The fixed first factor is nonzero, so the extracted period
gives, upon taking complex conjugates,
\[
 \int_{[H]}
    \overline{\Theta_{W,V_2}^{a,b_2}(\Phi_2;g,t)}\,
    \overline{Y_u(\operatorname{diag}(h_1,t)\gamma)}\,dt=0 .
\]
All these integrals use the given class measures.
Expanding the Cartan contraction in the actual Weyl coordinates
gives a finite sum of these periods. The invariant Gram form and
Cartan equivariance cancel the real compact action on the
sum, so its integral is
\([K_{0,H}:K]^{-1}\sum_c f_c/|\operatorname{Aut}(c)|\)
when the resulting finite contraction is fixed by the actual
chosen level \(K\) of \(H\).
The original class sum therefore vanishes with its specified
weights.
In the comparison application the principal level at
\(\mathscr M^2\) fixes the kernel and
\(\gamma_{i,v}^{-1}(k_v-1)\gamma_{i,v}\) is in
\(\varpi_v M_n(\mathcal O_v)\) at every chosen \(v\).
Here \(\gamma_{i,v}\) is integral by the conjugate side of
the divided-lattice condition and
\(\varpi_v\gamma_{i,v}^{-1}\) is integral by its chosen side.
At the conjugate place use the inverse transpose.
At a chosen place \(A_{\rm pt} I_{\rm mk}\) is the specified
column in the table. The conjugated complementary element
fixes it and hence conjugation by \(A_{\rm pt}\) fixes
\(I_{\rm mk}\). At twist places \(x=1\) and
\(\overline A_{\rm pt}^{\,\mathsf t}I_{e'}=I_{\rm mk}\)
for the prescribed rows \(e'\) of the fixed block.
The complementary motion fixes those rows, so its conjugate
by \(\overline A_{\rm pt}\) fixes \(I_{\rm mk}\) under
transpose. The corresponding determinant in the level
character is one. Thus the conjugation uses the original
divided lattice and table coordinates.
The rank gap is \(2r+2(r+9m+1)+3<6r\), since
\(r+1=\kappa+m\) and \(\kappa>100m+9\).
In the point application the rank is \(r=1\),
the first \(12+4+4=20\) positions are fixed and
\(45<2z_p\). At the chosen places the level fixes \(J_1\)
and therefore the full column \((x_0,J_2,J_3,J_1)^{\mathsf t}\).
The determinant character is read from the second block.
When the first three factors are fixed, the last-block
contraction is on the same flat character and the same
chosen-level classes as the point test. Its real compact
type is trivial, so the Cartan equivariance used in the
class sum is just invariance on this last block.
In both cases the finite translate is \(\gamma=xA_{\rm pt}\), the
specified divided translation followed by the prescribed integral
point lift. The rank-one test uses the trivial divided translation.
This proves the vanishing of the comparison adjoint tables and the testedness
of the rank-one form on their respective false spaces.%
\footnote{\weightedpin{grh-compact/src/New/Rows/MemberOrthogonality.lean}
{57-L138}{Marked point, class sum and divided level};
\weightedpin{grh-compact/src/New/Rows/ComparisonForms.lean}
{165-L192}{Application on the actual comparison adjoint};
\weightedpin{grh-compact/src/New/Rows/HeckeValueHeight.lean}
{122-L146}{Application on the actual four-block test}.}

Apply this implication to the same tables. For a form \(Y\) in
\(\operatorname{im}e_0\cap\ker H\cap L_\psi^\perp\) the marked
period vanishes in each \(\mathcal T_i(G_\psi^*Y)\).
Applying the common adjoint formula at \(G_\psi^*Y\) then gives
\(A G_\psi^*Y=0\). Hence \(G_\psi A^*F\) is orthogonal to
that false-lift space for each \(F\). The tested forms and their
integrality thus refer to the same \(G_\psi\) that compresses
the single-lift determinant.\footnote{\weightedpin
{grh-compact/src/New/Rows/Separation.lean}{101-L155}{Matched large-block identity},
\weightedpin{grh-compact/src/New/Rows/Separation.lean}{195-L236}{Detection on the actual small-block tables},
and \weightedpin{grh-compact/src/New/Rows/ComparisonForms.lean}{165-L192}{Use of marked vanishing on the comparison form}.}

\paragraph{The compressed determinant.}
Use the
\hyperlink{proof.weighted-fit-table-frame}{integral table frames}
and the
\hyperlink{proof.weighted-comparison-matching-reduction}{matched real exchange scalar}.
The common adjoint is onto by the
\hyperlink{proof.weighted-comparison-detection-reduction}{detection on those tables},
and the rank-one calculation uses the same datum's
\hyperlink{proof.weighted-rank-one-depth-reduction}{local depth coefficient}. Set
\(K_\psi=k_j P_{L_\psi}G_\psi G_\psi^*|_{L_\psi}\) and
\(d_j(\psi)=\det K_\psi\). The identity \(u_\psi^*G_\psi=\gamma_\psi A\)
makes \(u_\psi^*\) surjective, so \(u_\psi\) is an isomorphism onto
\(L_\psi\). Multiplying out \(u_\psi^*K_\psi u_\psi\) and taking
determinants gives
\begin{align*}
  \lambda_j(\psi)d_j(\psi)
       &= k_j^R |\gamma_\psi|^{2R}\det(AA^*),\\
  \frac{\lambda_j(1)}{\lambda_j(\psi)}
       &= \frac{d_j(\psi)}{d_j(1)}
           \left(\frac{\Gamma_j(1)}{\Gamma_j(\psi)}\right)^R.
\end{align*}
In particular \(d_j(\psi)\ne0\), and \(G_\psi^*\) has trivial kernel
on \(L_\psi\) by \(G_\psi^*u_\psi=\overline{\gamma_\psi} A^*\).
The finite-place estimate supplied by the fitted determinant test is
\[
  |j_\iota(d_j(\theta^j))|_\ell
       \le \exp(Rr\beta_{j,\ell}).
\]
It is applied to the actual \(G_\psi\) and \(K_\psi\) above, with
\(\psi=\theta^j\). The
\hyperlink{proof.weighted-fit-table-frame}{dual-coordinate selection}
uses the Fock matrix, whose product with the Weyl Gram matrix is
\(k_jI\). The first frame consists of forms
integral at every embedding above \(\ell\). The second consists of
forms integral at \(j_\iota\) and tested against the false lifts.
\footnote{\weightedpin{grh-compact/src/New/Rows/LiftVolumeHeight.lean}{41-L63}{Finite determinant test on the Weyl tables}.}

For the archimedean estimate, the projection to \(L_1\) is
contractive on tested forms for each coefficientwise norm. Indeed
the good idempotent \(e_0\) is self-adjoint for this norm because the
finite group acts unitarily. At a split place off the level write
\(H\) for the test Hecke polynomial operator. Its adjoint for the
coefficientwise metric is the Hecke polynomial at the conjugate place
with its parameter conjugated in that metric. It commutes with \(H\)
and with \(e_0\). Consequently \(H\) is normal for every such metric
and the orthogonal complement of \(P=\operatorname{im}e_0\cap\ker H\)
is the same vector space \(\ker e_0+\operatorname{im}H\) for all of
them. The lift image \(L_1\) has the good type selected by \(e_0\).
It lies in \(\ker H\) by the
\hyperlink{proof.weighted-genuine-theta-annihilator}{split vacuum
calculation on that genuine lift}. A tested form, orthogonal in the
original metric to \(P\cap L_1^\perp\), therefore splits as its
\(L_1\)-projection plus a vector in this common orthogonal complement.
The two terms are orthogonal in the coefficientwise metric as well.
\footnote{\weightedpin{grh-compact/src/New/Selmer/FitTwists.lean}{16-L34}{Hecke adjoint at the conjugate place},
\weightedpin{grh-compact/src/New/Selmer/FitTwists.lean}{63-L113}{Commutation, finite group invariance and tested projection}, and
\weightedpin{grh-compact/src/New/Rows/RowVolume.lean}{72-L82}{Tests vanish on the lift range}.}

Let \(z\in L_1\) be a nonzero eigenvector of \(K_1\), with eigenvalue
\(\mu\), and put \(\xi=G_1^*z\ne0\). Write
\(\sigma^\#=\operatorname{conj}\circ\sigma\circ\operatorname{conj}\),
and choose \(\xi'\) with \((\xi')^{\sigma^\#}=\xi^\sigma\). Both
the form metric and the table metric have rational Gram entries.
With the inner product linear in its second entry, applying \(\sigma\)
to the pairing of \(G_1\xi'\) with
\(\mu z=k_j P_{L_1}G_1\xi\) gives
\[
   \sigma(\mu)\|\xi^\sigma\|^2
    = k_j\big\langle (G_1\xi')^{\sigma^\#},
                         (P_{L_1}G_1\xi)^\sigma\big\rangle.
\]
The vector \(G_1\xi\) is tested by the adjoint identity above. Using
the contractivity and the two bounds for \(G_1\) shows that the
right side has modulus at most \(e^{2rB_0}\|\xi^\sigma\|^2\). Thus
\(|\sigma(\mu)|\le e^{2rB_0}\) and
\(|\sigma(d_j(1))|\le e^{2RrB_0}\). This uses the mixed coefficient
pairing from the \hyperlink{proof.fitted-row-mixed-gram}{base quotient
comparison}, now on the eigenvector test.\footnote{\weightedpin{grh-compact/src/New/Rows/LiftVolumeHeight.lean}{65-L94}{Conjugate eigenvalue estimate}.}

\paragraph{The signed contribution of the rank-one tests.}
\teamanchor{proof.weighted-rank-one-height}%
Fix one rank-one auxiliary pair \(p\) at stage \(j\). Write
\(F_\psi\) for its nonzero trivial-character lift, with the compact
jet and vacuum scales, and \(h_1\) for the class number of its
chosen-level \(U(1)\). Let \(P_\ell\) be the sum over \(\sigma\) of
the local asks at \(u=w(\sigma)\), including the factor
\((1-\chi_d(\ell))/2\). Then
\begin{equation}
\label{eq:weighted-forced-rank-one}
  \sum_\ell 2P_\ell\log\ell
       =\sum_\sigma \operatorname{Ask}(w(\sigma),\sigma).
\end{equation}
Indeed only prime powers contribute. For inert \(\ell\) only the
odd powers survive and for the ramified prime all powers survive.
The supports are finite at \(\ell\le x_1\) by the minimum weight.

For each such \(j,p\), fix the ramified pair \(q_2\), the scalar
\(V_{jp}\ne0\) of the rank-one test, and the flat pair of the same
rank used in its untwisted norm comparison. Set
\[
  S_0=\sum_{\sigma\in\Phi}S(w_0(\sigma)),\qquad
  a_{jp}=I_{jp}/\|F_1\|^2,\qquad
  A_{jp}=a_{jp}V_{jp}\overline{V_{jp}},
\]
where \(I_{jp}\ne0\) is the integral algebraic part at the trivial
twist of the flat pair. Choose \(V_{jp}\) as in the
\hyperlink{proof.weighted-rank-one-character-test}{character
selection}. Its conjugates have modulus at most \(\exp(B_0)\).
The \hyperlink{proof.weighted-rank-one-flat-bounds}{flat-part
comparison} gives the bound \(\exp(B_0-S_0)\) for the
conjugates of \(a_{jp}\).

Given the fit at \(\iota\), the construction gives a tested form
\(\Xi_a\) for each \(a=j_\iota\) and
\(a=j_\iota\circ\operatorname{conj}\), with
\begin{equation}
\label{eq:weighted-rank-one-test}
 |a(V_{jp})|_\ell\ell^{-P_\ell}
  \le\left|a\left(h_1^2\langle\Xi_a,F_\psi\rangle
                         /\rho_{q_2}(\psi)\right)\right|_\ell,
             \qquad \psi=\theta^j.
\end{equation}
The test at the conjugate embedding uses the same \(\psi\), and each
\(\Xi_a\) is integral at all embeddings above \(\ell\).

\paragraph{The rank-one norm at a twist place.}
\teamanchor{proof.weighted-rank-one-norm}%
For a rank-one pair \(p\) of point \(z_p\), the definite target
has rank \(n_p=2z_p\ge4\). The source \(U(1)\) is compact
modulo rational points, the constant source form is cuspidal,
and \(n_p>2\), so the convergent doubling norm formula applies
to the prescribed finite section and compact Fock vector.
Use the split-support calculation of
\S\ref{sec:period-definite-triangle}. At a split place newly in
the conductor, the section has first coordinate \(1\) in the residue field
and \(q^{n_p-1}\) choices for the other residues. Its squared
norm relative to the vacuum is \(q^{-1}\). The matrix coefficient
is supported on source units reducing to \(1\), as both
first coordinates must equal \(1\). The measure of this support
is \((q-1)^{-1}\) with the maximal compact of volume one. The
finite projection on the constant source multiplies the lift
by \(q-1\), because its \(q-1\) summands are equal after translating
the constant function. The target level preserves a hyperplane
of codimension one and has index \((q^{n_p}-1)/(q-1)\), by
transitivity on hyperplanes. Consequently the factor in the
normalized class norm is
\[
 \frac{q^{n_p}-1}{q-1}\,(q-1)^2q^{-1}\frac1{q-1}
                       =q^{n_p-1}(1-q^{-n_p}).
\]
This is the same support and level-index calculation as for the
earlier bottom norm, with a single marked column.\footnote{\weightedpin
{grh-compact/src/New/Exchange/CharacterLiftNorm.lean}{46-L73}{The projected section and its support cardinality},
and \weightedpin{grh-compact/src/New/Rows/StackedDeterminants.lean}{43-L47}{The auxiliary target level}.}

Off the chosen places and the twist, the split vacuum has
overlap \(q^{-n_p|a|/2}\) under a source translation of valuation
\(a\in\mathbb Z\). This follows by intersecting the standard
lattice with its scalar translate in the split Schr\"odinger
model and inserting the square-root determinant in the Weil
action. Its character factor is \(u^a\) for the unitary
parameter \(u\) in the local Rallis splitting. The two geometric
series give
\[
  \sum_{a\in\mathbb Z}u^a q^{-n_p|a|/2}
    =\frac{1-q^{-n_p}}
      {(1-uq^{-z_p})(1-u^{-1}q^{-z_p})}.
\]
The two denominator factors come from the inverse of the base
change of the weight character \(\mu_p(\psi)\) of this rank-one
pair. The Weil splitting in the oscillator overlap has already
been included in \(q^{-n_p|a|/2}\) and in \(u\). At an unmarked
nonsplit place the norm-one group is compact and the unramified
vacuum coefficient is one on it. Thus its normalized local factor
is one. At a fixed chosen place the section has all \(n_p\)
residues prescribed, the target index is the number \(q^{n_p}-1\)
of nonzero vectors, and the coefficient has support in source
units reducing to one. Changing the source integral to its
principal-level class sum multiplies the lift by \(q-1\).
The fixed place thus contributes \((q-1)(1-q^{-n_p})\).

The product of the spherical numerators, the compact real
volume and the real harmonic norm gives a common nonzero
factor. In the unscaled kernel normalization the real harmonic
factor is the hook number times the coefficient dimension.
The jet and vacuum scales multiply it by the same positive
quantity for every twist. The relative norm consequently has
the ratio of the incomplete Hecke values of
\(\mu_p(\psi)^{-1}\), with the factor \(q^{n_p-1}\) at each new
conductor place. Absolute convergence at \(z_p>1\) gives
nonzero Hecke products. In particular \(F_\psi\ne0\)
and
\[
 \rho_p(\psi)^{-1}=
    \frac{\|F_1\|^2}{\|F_\psi\|^2}\operatorname N(Q_\psi)^{\,n_p-1}.
\]
The same calculation applies to the flat pairs of ranks \(4\)
and \(2z_p-20\) used for the rank-one test.
\footnote{\weightedpin
{grh-compact/src/Lib/SiegelWeil/RallisInner/SharpData.lean}{472-L510}{Spherical rank-one factor},
and \weightedpin
{grh-compact/src/New/Exchange/CharacterLiftNorm.lean}{186-L252}{Relative normalization and compact factors}.}

\paragraph{The flat algebraic parts and the untwisted multiplier.}
\teamanchor{proof.weighted-rank-one-flat-bounds}%
For a flat pair of rank \(n=2z\le t\), the diagram is empty
and the coefficient space is one dimensional with rational
Gram matrix \(1\). At the trivial twist the denominator clearer
is \(g^{2fn}\). Set \(D_{\mathrm{flat}}=
D_{\mathrm{aut}}g^{4fn}\), so its integral algebraic part is
\(I=D_{\mathrm{flat}}\|F_1\|^2\). This is integral at every
coefficient place. Indeed the vacuum clearers make
\(g^{2fn}F_1\) integral by the normalized indicator calculation,
and \(D_{\mathrm{aut}}\) clears the stabilizer weights of the
final class sum. The same argument off \(Q_\psi\) applies
to the ramified flat block with its conductor power a unit.
\footnote{\weightedpin{grh-compact/src/New/Exchange/AlgebraicPartDefs.lean}
{223-L255}{Empty shape and scaled lift};
\weightedpin{grh/src/New/OneLift/HeckeValues/AlgebraicPartIntegrality.lean}
{176-L253}{Point kernels, integral class sums and the algebraic part}.}

For the same flat pair, expand the unramified indicator
on matrices modulo \(\mathfrak M\) into finitely many moves
of the normalized vacuum. The coefficients have conjugate
moduli at most the number \(g^{fn}\) of residue matrices.
At order zero, the conjugate Gaussian integral in the
weighted jet bound, with \(L=1\), gives a bound of
\[
 C_n=g^{fn}\exp\bigl(C_d(N_1n+1)^2+fn\log4\bigr)
\]
for every conjugate point kernel. Indeed the Weyl polynomial
is the constant \(1\); the Gaussian majorant has real
dimension \(2fn\) and the factor from its integration is
\(4^{fn}\). The level determinant scalars have modulus one
under the coefficient embeddings. The class counts of
\(U(1)\) and \(U(n)\) are at most \(e^{x^{30}}\) at their
chosen levels. The level comparison for this count uses the
principal level modulo \(\bar\pi_G\pi\), which is contained
in both actual levels. Thus after summing first in the source
group and then in the target group, the mixed pairing between
the \(\sigma\) and \(\operatorname{conj}\sigma\operatorname{conj}\)
readings gives
\[
 |\sigma(\|F_1\|^2)|\le e^{3x^{30}} C_n^2.
\]
This follows by Cauchy--Schwarz with inverse stabilizer orders
as positive rational weights. Here
\(\log D_{\mathrm{aut}}\le f (rt)^2\log g\).
The bounds \(f,N_1=O(x^{11}\log x)\), \(n=O(x^5)\) and
\(\log g=O(x^\varepsilon)\) give
\(\log C_n=o(x^{33})\) and
\(\log D_{\mathrm{flat}}=o(x^{33})\).
Hence \(|\sigma(I)|\le e^{x^{34}}\) for all the flat blocks
on this datum.\footnote{\weightedpin
{grh-compact/src/New/Rows/MemberKernels.lean}{42-L83}
{Conjugate Gaussian bound at order zero};
\weightedpin{grh-compact/src/New/Columns/SizeEstimates.lean}
{67-L85}{The applicable class counts};
\weightedpin{grh-compact/src/New/Rows/MonomialCoefficientVector.lean}
{18-L36}{Flat bound on the normalized lift}.}

To compare the untwisted norm of the main rank-one pair
to the flat pair of the same rank, write
\(m_\tau=c_p(\tau)-z_p\).
The conjugate CM jet transport uses the same scalar and the
same relabelling of places for the two pairs. The transported
conjugate Hecke character uses the same chosen level and
is unramified elsewhere. Apply the rank-one norm computation
at the transported character. The scale of each degree is
\((g/\Omega)^{m_\tau}\) in the squared norm, from the actual
\(\bar\pi_G\) factor and the conjugate CM period. Its
row hook is \(m_\tau!\). The character and vacuum phases
cancel in the mixed conjugate pairing. For \(\sigma\)
fixing \(K\) this gives
\[
 \log|\sigma(\|F_1\|^2)|
    =c_{\sigma,n}
      +\sum_\tau S(m_\tau)+\log d_p
      +\log|L^S(z_p,\nu_{\sigma,p}^{-1})|,
\]
where \(d_p\) is the dimension of its coefficient module and
\(c_{\sigma,n}\) is common to the two untwisted pairs.
The transported character is unitary on the rank-one torus.
Absolute convergence at \(z_p\ge2\) bounds the last logarithm
in modulus by \([E:\mathbf Q]\log\zeta(2)\le4f\).
Since \(d_p\ge1\) and the flat dimension is \(1\), subtracting
the formula for the main norm from that for the flat
algebraic part \(I_{jp}\) gives
\[
 \log\big|\sigma(I_{jp}/\|F_1\|^2)\big|
 \le \log D_{\mathrm{flat}}+8f-\sum_\tau S(m_\tau).
\]
An automorphism of \(\mathbf C\) restricting to conjugation
on \(K\) can be composed with complex conjugation on the
output without changing any modulus. The same bound
therefore holds for all conjugates. As \(g/\Omega\ge1\) and
\(m_\tau\ge w_0(\tau)\), \(S(m_\tau)\ge S(w_0(\tau))\).
Finally \(\log D_{\mathrm{flat}}+8f\le B_0\).
This gives the required bound \(e^{B_0-S_0}\) for the
fixed multiplier with no cost at finite places.
\footnote{\weightedpin{grh-compact/src/New/Exchange/AlgebraicPartConjugates.lean}
{17-L77}{Conjugate transport and the common norm constant};
\weightedpin{grh-compact/src/New/Rows/Separation.lean}
{80-L95}{Comparison with the flat witness at the same rank}.}

\paragraph{Curve coefficients and the twelve-dimensional indicator.}
\teamanchor{proof.weighted-rank-one-depth-reduction}%
Fix the rank-one pair \(p\) at stage \(j\), with degrees
\(n_\sigma=c_p(\sigma)-z_p\). The rank-one inequalities give
\(w_j(\sigma)-t\le n_\sigma\le w_j(\sigma)\).
Since \(j\ge1\), we have \(w_j-w_0=2jr^8>t\). Rounding and \(j\le m\)
give \(w(\sigma)\le w_0(\sigma)\le n_\sigma\le w(\sigma)+x^{42}\)
and \(n_\sigma\le2b\) for large \(x\). We will find a nonzero \(t_p\)
of conjugate modulus at most \(\exp(fmb/\kappa)\) such that for
every prime \(\ell\) and reading \(a_\ell:\mathbf C\to\mathbf C_\ell\)
there is a rank-\(12\) residue indicator at \((\bar\pi_G)\) and
a Taylor monomial of total degree \(n_\sigma\) in each row \(\sigma\).
The coefficient at scale \(\bar\pi_G\eta_0\), multiplied by the
vacuum scale to the twelfth power and by \(\prod_\sigma n_\sigma!\),
has absolute value at least \(\ell^{-P_\ell}|a_\ell(t_p)|_\ell\).

\paragraph{Packing the Frobenius directions.}
\teamanchor{proof.weighted-rank-one-depth-packing}%
Put
\(H=a/x\), \(L=a/x_1\), and \(\Delta=x^{42}\). Thus
\(w_\sigma=H\) on the thick set \(\mathcal J\), and \(w_\sigma=L\)
elsewhere, with \(w_\sigma\le n_\sigma\le w_\sigma+\Delta\) and
\(n_\sigma\le 2b\). Write \(c_* = \lceil 60b/\kappa\rceil\),
\(n_\sigma^- = \max(n_\sigma-c_*,0)\), and
\(D_\sigma=(\prod_{h=1}^{c_*}(n_\sigma^-+h))^2\).
Set \(\omega_\ell=(1-\chi_d(\ell))/2\) and
\[
 P_{\ell,\sigma}=\omega_\ell\!!\!
 \sum_{\substack{e\ge1\\ e\ {\rm odd}\ \text{or}\ \ell=d}}
 \bigl(w_\sigma/\ell^e-w_{[-\ell^e]\sigma}\bigr)_+,
 \qquad P_\ell=\sum_\sigma P_{\ell,\sigma}.
\]
For a nonsplit \(\ell\), the same live powers have
\[
 \omega_\ell\!\!\!\sum_{\substack{e\ge1\\ e\ {\rm odd}\ \text{or}\ \ell=d}}\ell^{-e}
   =\frac{1}{\ell-1}-\rho_\ell,
 \qquad
 \rho_\ell=\begin{cases}
  1/(\ell^2-1)&\ell\text{ inert},\\
  1/(2(\ell-1))&\ell=d.
 \end{cases}
\]

First let \(\ell\le\kappa\) be inert. It is prime to \(N_1\) and
\(\ell\ge13\), since \(2,3,5,7,11\) all split. Put \(u=[-\ell]\) and
let \(\phi\) be reduction modulo \(2^{s+1}\) on the place group. If
\(\phi(u)^v=1\) for \(0<v\le 2s\), then \(2^{s+1}\mid\ell^{2v}-1\).
For any odd \(z\) and positive \(v\), factoring the difference of
squares and then successively squaring gives
\[
 v_2(z^{2v}-1)=v_2(z-1)+v_2(z+1)+v_2(v).
\]
In fact the odd part of \(v\) leaves the valuation unchanged, and
each doubling after \(z^2\) adds one. The smaller of the first two
valuations is one, so \(2^s\mid v(\ell-1)\) or
\(2^s\mid v(\ell+1)\). This is impossible, since
\(2^s=m^{22}\ge x^{11}\) and \(2s(\kappa+1)=O(x^5\log x)\).

Join \(\sigma\) to \(u^i\sigma\) whenever both are thick and
\(0<i<s\) is odd. These directed edges have distinct endpoints and
are disjoint. Indeed a thick triple of the form
\(\sigma,u^i\sigma,u^j\sigma\), with \(0<i<j\le2s\), satisfies the
three-term identity in the \hyperlink{proof.weighted-thick-places}{construction
of the thick set}. It would force \(\phi(u)^i=1\). Two outgoing edges
give such a triple. Two distinct incoming edges of exponents \(i<j\)
give one by starting at the source of the longer edge, with exponents
\(j-i,j\); equal exponents give the same source. Two consecutive
edges give one with exponents \(i,i+j\).

On an edge out of \(\sigma\), choose \(i_\sigma=i\) and
\(y_\sigma=\lfloor n_\sigma/\ell^i\rfloor\). Elsewhere set \(y_\sigma=0\)
and \(i_\sigma=0\). The outgoing cost \(\ell^iy_\sigma\) is at most
\(n_\sigma\). Each target is thick and has only its single incoming
cost, which is at most \((H+\Delta)/\ell\le H\le n_{u^i\sigma}\).
For two active sources \(\sigma,\tau\), the equality
\(u^e\tau=u^{i_\sigma}\sigma\) at a positive odd \(e<i_\sigma\)
would give the same excluded triple, now with exponents
\(i_\sigma-e,i_\sigma\) from \(\sigma\). At \(e=i_\sigma\), the
equality forces \(\tau=\sigma\). Thus the capacity and lowest-term
conditions both hold. At the ramified prime \(\ell=d\), take
\(i_\sigma=y_\sigma=0\) for every \(\sigma\). No live edge there takes
thick to thick, since all thick places have class \(-1\) modulo \(M\),
whereas \(d\equiv M/2-1\pmod M\) and \(4\mid M\) give
\(d^2\equiv1\pmod M\) and \(-d^e\ne1\pmod M\) for every \(e\ge1\).

For a live \(e<s\) that is outside a thick-to-thick edge, the loss
in replacing \(n_\sigma/\ell^e\) by
\((w_\sigma/\ell^e-w_{[-\ell^e]\sigma})_+\) is at most \(2L\).
For a thin source this follows from \(n_\sigma\le L+\Delta\).
For a thick source the target is thin, and the loss is at most
\(L+\Delta/\ell^e\). Here \(\Delta\le L\). At the single edge
of a source the bound \(n_\sigma/\ell^e\le y_\sigma+1\) suffices.
The tail \(e\ge s\) is at most \(2b\sum_{e\ge s}2^{-e}=4b/2^s\).
The factor \(\omega_\ell\) is one in the inert case and one half at
\(d\), so in both cases
\[
 \left(\frac{1}{\ell-1}-\rho_\ell\right)n_\sigma
 \le P_{\ell,\sigma}+y_\sigma+2sL+1+4b/2^s.
\]
In the weighted scales \(b_2\le L+N_1\le2L\), \(s=O(\log x)\),
and \(L=O(bx/x_1)\), with \(x_1\ge x^{10^6}\). For
\(\ell\le\kappa\le 3x^5\) we have
\(c_*/\ell\ge (20/3)b/x^{10}\). This also covers \(\ell=d\) for
large \(x\). Since \(4b/2^s\le 4b/x^{11}\), it follows uniformly
that \(2sL+1+4b/2^s+4b_2\le \lfloor c_*/\ell\rfloor\).
The product of \(c_*\) consecutive positive integers in the
definition of \(D_\sigma\) is divisible by \(c_*!\). Summing over
the places gives the needed small-prime bound\footnote{\weightedpin{grh-compact/src/New/Curve/LiveSum.lean}{16-L24}{The local asks}
and \weightedpin{grh-compact/src/New/Curve/DepthPairs.lean}{12-L92}{The capacity, collision and numerical bounds}.}
\begin{equation}
 \label{eq:weighted-depth-small-packing}
 \sum_\sigma \left(\left(\frac{1}{\ell-1}-\rho_\ell\right)n_\sigma
                       +4b_2-y_\sigma\right)
 \le P_\ell+\sum_\sigma v_\ell(D_\sigma).
\end{equation}

For an inert \(\ell>\kappa\), put \(i_\sigma=1\) and
\(y_\sigma=\min(\lfloor n_\sigma^-/\ell\rfloor,n_{[-\ell]\sigma})\).
The partner map is a permutation, and
\(\lfloor n_\sigma^-/\ell\rfloor\le 2b/\ell<c_*\).
At a target with \(n_\sigma<c_*\) the outgoing cost is zero and the
incoming cost is capped at \(n_\sigma\). At any target with
\(n_\sigma\ge c_*\) its outgoing cost is at most \(n_\sigma-c_*\)
and its incoming cost at most \(c_*\). This gives the capacity in
both cases. The least allowed odd exponent is one, and partners
are distinct at that exponent. Put \(J_\sigma=\prod_{h=1}^{c_*}(n_\sigma^-+h)\).
If \(n_\sigma<c_*\), \(n_\sigma!\) divides \(J_\sigma=c_*!\).
Otherwise \(n_\sigma!=n_\sigma^-!J_\sigma\). In either case,
writing \(q_\sigma=\lfloor n_\sigma^-/\ell\rfloor\), Legendre's
formula gives
\[
 v_\ell(n_\sigma!)\le q_\sigma+v_\ell(q_\sigma!)+v_\ell(J_\sigma).
\]
Here \(q_\sigma\le c_*\), and \(c_*!\mid J_\sigma\). Also
\(n_\sigma^-\le w_\sigma\) since \(c_*\ge\Delta\). Therefore
\[
 q_\sigma-y_\sigma
 \le (w_\sigma/\ell-w_{[-\ell]\sigma})_+
 \le P_{\ell,\sigma}.
\]
Using \(D_\sigma=J_\sigma^2\) yields\footnote{\weightedpin{grh-compact/src/New/Curve/DepthPairs.lean}{45-L59}
{The large-inert choice with natural-number subtraction}.}
\begin{equation}
 \label{eq:weighted-depth-large-packing}
 v_\ell(n!)-\sum_\sigma y_\sigma
     \le P_\ell+\sum_\sigma v_\ell(D_\sigma),
 \qquad n!=\prod_\sigma n_\sigma!.
\end{equation}

\paragraph{The frame of length \(g\) in the rank-one gauge.}
\teamanchor{proof.weighted-rank-one-cm-frame}%
Set \(\nu=N_1\). We use the lattice \(\mathfrak a_0=E\cap
\alpha\widehat{\mathcal O}_E\) and the period factor \(s_1\) from
the \hyperlink{proof.weighted-local-source-gauge}{normalized source gauge}.
The list of powers of \(\zeta\) has frame constant \(\nu\).
For every rational prime \(q\mid\nu\), use the list \(w_{jkh}\)
from that gauge, before repetition. Put \(A_0=(d+1)/4\), choose
\(0\le z_0<\nu\) as in its construction, and write
\(t_1\bar t_1=A_0\nu e_0\), with \(e_0\in\mathbf Z\). Denote
\(\nu'_q=\nu/q^{v_q(\nu)}\). The frame constant of the unrepeated
list is
\[
 C_q=\nu+\nu'_q(1+2A_0z_0)+2q(\nu'_q)^2 A_0 e_0.
\]
This is a positive integer by the same diagonal frame sum. It is
\(0\) modulo each prime divisor of \(\nu\) other than \(q\) and
\(\nu'_q\) modulo \(q\). Thus \(C=\sum_{q\mid\nu}C_q\) is prime to
\(\nu\), where the sum runs over primes. It is of polynomial size
in \(x\), since \(z_0<\nu\), \(e_0=O_d(\nu)\) and \(\nu\) is of
polynomial size. Put \(R=3|n|\), where \(|n|=\sum_\sigma n_\sigma\).
We have \(g\ge b^{20}\) and \(R\le6fb\), so
\(g-R(C+\nu)\ge(\nu-1)C\). The integer on the left is \(a'\nu+b'C\)
with \(a',b'\ge0\), by first choosing \(0\le b'<\nu\) in its
residue class after multiplication by \(C^{-1}\). Repeat each
vector of the powers-of-\(\zeta\) list \(R+a'\) times, and each
vector of the other lists \(R+b'\) times. This gives
\[
 \sum_i c_i\Phi_\sigma(w_i)\overline{\Phi_\tau(w_i)}
       =g\,1_{\sigma=\tau},\qquad c_i\ge3|n|.
\]
The list is in \(\mathfrak a_0\), so for any reading
\(a_\ell:\mathbf C\to\mathbf C_\ell\) all coefficients
\(\mu_{i,\sigma}=a_\ell(s_1\overline{\Phi_\sigma(w_i)})\)
are integral. There is an index \(i\) at which all are units.
If the period factor is a unit, use \(w_i=1\). Otherwise \(\ell\mid\nu\)
and the readings of \(\overline{\Phi_\sigma(w_i)}\) see the side
where \(\alpha=\Phi_\nu'(\zeta)^{-1}\). In the block \(q=\ell\), take
the vector with \((j,k,h)=(0,1,0)\). This is
\(t_1\nu'_\ell(1-\zeta_\ell)/\nu\). The factor \(t_1\) is a unit
on that side. Its valuation is therefore
\(1/(\ell-1)-v_\ell(\nu)=v_\ell(\alpha)\) for every \(\sigma\),
and the period factor cancels it. In particular \(s_1\) and its
conjugate are units at every nonsplit \(\ell\). At these primes
\(g\) is also a unit, since both of its prime factors split.
\footnote{\weightedpin{grh-compact/src/New/Curve/TightFrameExists.lean}{33-L55}
{Period valuations}, \weightedpin{grh-compact/src/New/Curve/TightFrameExists.lean}{87-L148}
{Length and unit directions}, and \weightedpin{grh-compact/src/Lib/Weil/Construction/ConjugateJetDefs.lean}{19-L25}
{Frame identity with multiplicity}.}

\paragraph{Invariant derivatives in the marked elliptic chart.}
\teamanchor{proof.weighted-rank-one-local-derivatives}%
Work in the marked normalized parameter \(Y_0(z)=z+O(z^2)\) of
the \hyperlink{proof.weighted-local-cm-germs}{elliptic germ calculation}.
Let \(\mathcal O\) be the valuation ring of \(\mathbf C_\ell\).
Translation by the marked point identifies this chart with the formal
group at the origin of the good elliptic model. Its integral group
law expresses \(Y_0(z+w)\) as a power series with integral
coefficients in \(Y_0(z),Y_0(w)\). For nonsplit \(\ell\), the same
elliptic germ calculation gives \(T\in \mathcal O[[Y_0]]\) with
\[
 |\varpi|_\ell=\ell^{-\rho_\ell},\qquad
 S(X)=T(\varpi X)/\varpi\in\mathcal O[[X]],\qquad
 \overline S-\overline S^{\,q_0}=X,
\]
where \(q_0=\ell^2\) for inert \(\ell\) and \(q_0=\ell\) at \(d\).
The construction of this local parameter uses only \(\ell\ge13\)
and the indicated splitting type, so applies to every nonsplit \(\ell\).

For split \(\ell\) the same parameter conclusion holds with
\(q_0=\ell\) and \(\rho_\ell=1/(\ell-1)\). To verify this case,
choose \(a_0\in\mathcal O_K\) of valuation one at the prime of \(K\)
defined by the reading and valuation zero at the conjugate prime.
The CM isogeny \([a_0]\) extends to the good elliptic model. Its
degree has \(\ell\)-valuation one by the norm of \(a_0\). Its
reduced differential is zero, so its inseparable degree is exactly
\(\ell\). In an integral formal parameter its reduction starts in
degree \(\ell\). Average \(Y_0(\zeta z)\) with coefficients
\(\zeta^{-1}/(\ell-1)\) over the \((\ell-1)\)-st roots of unity
in \(\mathbf Z_\ell\), giving an integral parameter \(T\). The
integral action for these scalars follows by limits of integer
multiplications. Write \(Z(T)\) for the inverse in the differential
coordinate. It has only degrees \(1\bmod\ell-1\), and \(Z'\) is
integral since the invariant differential is integral in \(T\).
The coefficient of \([a_0]\) at degree \(\ell\), in the \(T\)
coordinate, is \((a_0-a_0^\ell)Z_\ell\), where \(Z_\ell=[T^\ell]Z\).
Thus \(v_\ell(Z_\ell)=-1\). Choose
\(\varpi^{\ell-1}=-1/Z_\ell\). In degree \(1+h(\ell-1)\) for \(h\ge2\),
the coefficient of \(\varpi^{-1}Z(\varpi X)\) has valuation at
least \(h-v_\ell(1+h(\ell-1))\ge1\). The last inequality follows
from \(1+h(\ell-1)<\ell^h\). Thus its reduction is \(X-X^\ell\),
and its integral compositional inverse is the desired \(S\).

Let \(A_m\) be the maximum of
\(\lvert m![z^m]J(Y_0(z))\rvert_\ell\) over \(J\in\mathcal O[[Y_0]]\).
The maximum is attained on one of the finitely many monomials
\(Y_0^h\) with \(0\le h\le m\). The integral addition law gives
\(A_{m+n}\le A_m A_n\): express \(J(Y_0(z+w))\) in the two
parameters and take the coefficient of \(z^m w^n\) times \(m!n!\).
In the differential coordinate this is \((m+n)![z^{m+n}]J\).
Put \(B_m=A_m/\ell^{(\rho_\ell-1/(\ell-1))m}\). The reduced
equation for \(S\) gives
\(\overline S=\sum_{j\ge0} X^{q_0^j}\), by iteration starting at
zero constant term. Therefore at \(m=q_0^j\) the coefficient of
\(T\) in degree \(m\) has absolute value \(|\varpi|_\ell^{1-m}\).
Legendre's formula gives \(B_{q_0^j}\ge|\varpi|_\ell\).
If some \(B_m<1\) with \(m>0\), submultiplicativity and division
of every index by \(m\) would force \(B_n\to0\). Thus for every
\(m\) an integral series \(J_m\) in the marked parameter satisfies\footnote{\weightedpin{grh-compact/src/Lib/CM/LocalHeight/EulerAdditionFormula.lean}{27-L50}
{Integral addition in the marked parameter},
\weightedpin{grh-compact/src/New/Curve/FormalParameter.lean}{125-L150}
{The rescaled parameter}, and
\weightedpin{grh-compact/src/New/Curve/KatzDivisibility.lean}{13-L42}
{Submultiplicativity and sharpness in every degree}.}
\[
 \lvert m![z^m]J_m(Y_0(z))\rvert_\ell
 \ge \ell^{(\rho_\ell-1/(\ell-1))m}.
\]

\paragraph{Twelve factors on the marked frame.}
\teamanchor{proof.weighted-rank-one-depth-products}%
Write \(Z=(Z_\sigma)_\sigma\), \(\mu_i(Z)=\sum_\sigma\mu_{i,\sigma}Z_\sigma\),
and \(\Theta=\prod_i\theta(\mu_i(Z))^{c_i}\), with coefficients read by
\(a_\ell\) and \(\theta(0)=1\). Let \(\mathcal R\) be the image of the
integral formal series ring in the variables \(Y_0(\mu_i(Z))\).
For each multi-index \(a\), choose \(H_a=\prod_iY_0(\mu_i)^{h_i(a)}\)
maximizing \(N(a)=|[Z^a](H_a\Theta)|_\ell\) among monomials of
total degree at most \(|a|\). The same bound applies to the coefficient
of \(J\Theta\) for every \(J\in\mathcal R\) by the ultrametric
inequality. If a product of \(k\) such factors has a coefficient at
\(a\), the coefficient formula bounds it by some
\(\prod_{j=1}^kN(a_j)\) with \(\sum_j a_j=a\). The same bound holds
for any factor obtained by rescaling \(Z\) by a scalar in
\(\mathbf Z_\ell\), because this multiplies each coefficient by
a scalar of norm at most one.

Put \(L_\theta=-\log\theta\). The
\hyperlink{proof.weighted-local-cm-germs}{marked isogeny formula}
implies \(\exp(t^2L_\theta(z)-L_\theta(tz))\) is an integral unit
in \(Y_0(z)\) for \(t\in\mathbf Z_\ell\). Use nonzero integers and
coefficientwise continuity in the marked parameter. Choose
\(t_1,\ldots,t_4\in\mathbf Z_\ell\) with \(\sum_h t_h^2=-2\), by
the four-squares theorem modulo each power of \(\ell\) and compactness.
It follows that \(U_\theta=\Theta(Z)^2\prod_h\Theta(t_hZ)\) is an
integral unit in \(\mathcal R\) with constant coefficient one.
Its inverse belongs to \(\mathcal R\) by expanding the inverse of
an integral formal series with constant coefficient one. Now
\(\Theta^{-1}=(U_\theta^{-1}\Theta)\prod_h\Theta(t_hZ)\), so each
coefficient of \(\Theta^{-1}\) is bounded by a five-fold product of
maxima \(N\) over a splitting of its degree.

For a residual degree \(r'\), apply the single-variable derivative
bound in order \(|r'|\) to the form with unit coefficients. It gives
an element of \(\mathcal R\) with coefficient at \(r'\) at least
\(\ell^{(\rho_\ell-1/(\ell-1))|r'|+v_\ell(r'!)}\) in absolute value,
by the multinomial formula. Here \(v_\ell(r'!)=\sum_\sigma v_\ell(r'_\sigma!)\).
Writing the element as a product with \(\Theta\) times \(\Theta^{-1}\)
gives this bound on the product of six maxima at indices summing to
\(r'\).

At an inert prime, let \(y_\sigma,i_\sigma\) be the packing above and
\(\sigma^*=[-\ell^{i_\sigma}]\sigma\). Put
\[
 A=\sum_\sigma \ell^{i_\sigma}y_\sigma e_\sigma,\qquad
 B=\sum_\sigma y_\sigma e_{\sigma^*},\qquad b'=A+B\le n.
\]
Use \(\theta(\varpi X)=\exp(\lambda\Xi(X))V_0(X)\) from the marked
germ with no torsion translation. Here \(V_0(0)=1\),
\(|\lambda|_\ell=\ell^{1-1/(\ell-1)}\), and
\(\overline\Xi=\lambda_0\overline S^{\,\ell+1}\) with \(\lambda_0\ne0\).
Thus \(F(Z)=\Theta(\varpi Z)=\exp(\lambda\Psi(Z))V(Z)\), where
\(\Psi=\sum_i c_i\Xi(\mu_i)\) and \(V=\prod_i V_0(\mu_i)^{c_i}\)
are integral and \(V(0)=1\). Write \(r_{i,\tau}=\overline{\mu_{i,\tau}}\)
in the residue field. As \(\overline S-\overline S^{\,\ell^2}=X\),
the coefficient at \(Z_\tau^{\ell^e}\) of
\(\partial_\sigma\overline\Psi\), for odd \(e\ge1\), is
\(\lambda_0\sum_i c_i r_{i,\tau}^{\ell^e}r_{i,\sigma}\).
These are the only possible monomials. Odd Frobenius on the
reduction of \(\overline{\Phi_\tau(w_i)}\) acts as the embedding
\(\Phi_{[-\ell^e]\tau}\) on \(w_i\). Frobenius undoes the conjugation
on \(K\) and sends the cyclotomic generator to its indicated power.
This reading applies to the frame vectors, since \(\ell\nmid dN_1\)
and they are integral at \(\ell\). The frame identity gives
\[
 \sum_i c_i r_{i,\tau}^{\ell^e}r_{i,\sigma}
 =\overline{a_\ell(s_1)}^{\,\ell^e+1}\overline g\,
       1_{[-\ell^e]\tau=\sigma}.
\]
Thus every indicated coefficient is nonzero, even when individual
multiplicities vanish in the residue field.

For each partial derivative, the new integral factor in
\(\lambda^{-|B|}F^{-1}\partial^BF\) comes from multiplying by
the corresponding partial of \(\Psi\), modulo the maximal ideal.
Indeed differentiation of any integral factor, or of \(V\), has
integral coefficients and comes with the factor \(\lambda^{-1}\).
Induction gives the reduction
\(\prod_\sigma(\partial_{\sigma^*}\overline\Psi)^{y_\sigma}\).
Give the active directions weight one and the inactive ones weight
\(\ell^{\max i_\sigma}\) when there is an active direction.
In each of the indicated partials the chosen monomial is the unique
one of lowest weight. For active directions this uses the packing's
collision property and the bijection at a fixed exponent. An inactive
direction contributes only exponents at least \(\ell\). The product
has a unit coefficient at \(A\). If all \(y_\sigma=0\), the product
is the constant one and the same conclusion holds.

Take a largest Cauchy summand at \(A\) and write its indices as
\(\gamma+\beta=A\) in the product \(F^{-1}\partial^BF\). The
coefficient of \(\partial^BF\) at \(\beta\) is the coefficient of
\(F\) at \(\beta+B\) times \(B!\binom{\beta+B}{B}\). The binomial
is integral. Undoing the scaling by \(\varpi\) gives
\[
 |[Z^\gamma]\Theta^{-1}|_\ell\,|[Z^{\beta+B}]\Theta|_\ell
 \ge \ell^{(1-1/(\ell-1))|B|+\rho_\ell|b'|+v_\ell(B!)}.
\]
This yields six maxima at indices summing to \(b'\). At \(d\)
use \(A=B=b'=0\) and six degree-zero factors. The corresponding
inequality then has right side one.

For \(\ell\le\kappa\) inert or \(\ell=d\), take the other six
maxima from the residual degree \(r'=n-b'\). Every entry of \(B\)
and \(r'\) is at most \(2b<\ell^{2b_2}\), as \(b_2\ge a/x_1\)
on the weighted scales. Legendre's formula bounds the combined
two factorial losses by \(4fb_2\). Adding the two six-factor
exponents gives at least
\(\rho_\ell|n|+\sum_\sigma y_\sigma-4fb_2\), which by
\eqref{eq:weighted-depth-small-packing} and
\(v_\ell(n!)\le |n|/(\ell-1)\) is at least
\(v_\ell(n!)-P_\ell-\sum_\sigma v_\ell(D_\sigma)\).

For inert \(\ell>\kappa\), all \(i_\sigma=1\), so
\(\rho_\ell|b'|=|B|/(\ell-1)\). The first six maxima have exponent
at least \(\sum_\sigma y_\sigma\). For the residual degree, use
the elements of \(\mathcal R\) with linear terms
\(a_\ell(\bar s_1s_1g) Z_\tau\) obtained by summing
\(c_i a_\ell(\bar s_1\Phi_\tau(w_i))Y_0(\mu_i)\) over \(i\).
Their coefficients as combinations of the generators are integral
by the same period bound at the conjugate reading. The common
linear scalar is a unit. The product of their powers with exponents
\(r'_\tau\), multiplied by \(\Theta\), has a unit coefficient in
degree \(r'\). Use this maximum and pad with degree-zero maxima.
Equation \eqref{eq:weighted-depth-large-packing} gives the same
total exponent. Finally for split \(\ell\), \(P_\ell=0\) and
\(\rho_\ell=1/(\ell-1)\), so the residual construction alone in
degree \(n\), with six more degree-zero maxima, supplies exponent
\(v_\ell(n!)\). In every case there are twelve arrays \(a_1,\ldots,a_{12}\)
with sum \(n\)\footnote{\weightedpin{grh-compact/src/Lib/CM/KroneckerLimit/LogTheta.lean}{140-L151}
{Logarithmic multiplication}, \weightedpin{grh-compact/src/New/Curve/ThetaProducts.lean}{44-L90}
{The actual frame and Frobenius log derivatives}, and
\weightedpin{grh-compact/src/New/Curve/DepthPairs.lean}{94-L146}
{The six-factor coefficient bound}.} and
\begin{equation}
 \label{eq:weighted-depth-twelve-maxima}
 \prod_{j=1}^{12}N(a_j)\ge
 \ell^{v_\ell(n!)-P_\ell-\sum_\sigma v_\ell(D_\sigma)}.
\end{equation}

\paragraph{Recovery of the rank-\(12\) indicator.}
\teamanchor{proof.weighted-rank-one-depth-indicator}%
Use the arrays \(a_b\) of \eqref{eq:weighted-depth-twelve-maxima},
where here \(b\) indexes the twelve columns.
Each \(Y_0^{h_i(a_b)}\theta^{c_i}\) is a degree-\(c_i\) elliptic
section, since \(Y_0\theta^3\) is a section of degree three and
\(3h_i(a_b)\le3|n|\le c_i\). Set \(\delta=\pi_G\), so
\(\delta\bar\delta=g\). With \(\eta=(2\pi/\sqrt d)^{1/2}\), the
product on the standard rank-\(12\) space is
\[
 \prod_{b=1}^{12}\prod_i
  (Y_0^{h_i(a_b)}\theta^{c_i})
  \left(\frac{\Omega_0}{\bar\delta\eta}
        \sum_\tau\overline{\Phi_\tau(w_i)}W_{\tau b}\right).
\]
For a period \(\bar\delta r_\infty(\gamma)\), \(\gamma\in\mathfrak a_0\),
the argument in each column is
\(\Omega_0\operatorname{Tr}_{E/K}(\gamma\bar w_i)\in\Lambda\),
by the trace-dual identity for \(\mathfrak a_0\). The frame identity
and \(|\delta|^2=g\) make the sum of the pulled back Hermitian forms
the standard Fock form. The elliptic canonical multipliers then
multiply to the standard multiplier on
\(L'=\bar\delta L^{12}\), where \(L=r_\infty(\mathfrak a_0)\).
This is the lattice and multiplier calculation of the
\hyperlink{proof.weighted-vacuum-projection}{finite vacuum projection},
now with \(\bar\delta\) as the lattice scale.

The rational finite translation formula identifies the translates of
the vacuum on this space with translates of the standard finite
vector. In fact \(L'^\vee=\delta^{-1}L^{12}\), by self-duality of
\(L\), and \(L'^\vee/L^{12}\) indexes a basis of translates of the
vacuum. The characters of \(L^{12}/L'\) on these vectors are
distinct by the Riemann pairing, and the dimension for the
polarization on \(L'\) is \(g^{12f}\). Hence the product is the
entire function of a vector \(U\) in the finite part of
\(\operatorname{Hom}(E,E^{12})\) with its standard Weil data. Every
finite value \(U(\operatorname{finVec} x)\) is integral at each
reading. For a rational \(\gamma\) this follows by translating the
elliptic sections to the corresponding torsion points and using
the normalized torsion height bound on each
\(Y_0^{h}\theta^{c_i}=(Y_0\theta^3)^h\theta^{c_i-3h}\). The same
bound at arbitrary finite arguments follows by the local constancy
of the finite translation and approximation by rational \(\gamma\).
The Heisenberg phases are roots of unity.

The Taylor scale \(\bar\delta\eta_0=\bar\delta s_1\eta/\Omega_0\)
turns the complex linear forms in the product into precisely the
forms read as \(\mu_i\) above. The twelve groups of variables are
disjoint. Thus if \(L_\alpha\) denotes coefficient extraction at
\(\alpha_{\sigma b}=(a_b)_\sigma\) in this Taylor scale, then
\(\sum_b\alpha_{\sigma b}=n_\sigma\) and
\[
 \bigl|a_\ell(L_\alpha(U))\bigr|_\ell
   =\prod_{b=1}^{12}N(a_b).
\]

To read a coefficient of a residue indicator, use Schur
orthogonality for the finite group of shifts
\(\mathcal A=L'^\vee/L'\), of order \(g^{24f}\). For each class
choose a representative \(r_\infty(\gamma)\) with
\(\gamma\in\delta^{-1}\mathfrak a_0^{12}\), and put
\(x=\alpha_{\mathrm W}^{-1}\gamma\) at the finite adeles, where
\(\alpha_{\mathrm W}\) is the idele defining \(\mathfrak a_0\).
Then \(\delta x\) is integral. Write \(R_x=\varrho(\operatorname{finVec}x)\)
for the induced shift.
The shifts act projectively, since \(L'\) acts by its scalar
multiplier. The vacuum translates give one-dimensional eigenspaces
for \(L^{12}/L'\), and the shifts permute these eigenspaces
transitively. Thus their action is irreducible. Averaging the
operator \(u\mapsto L_\alpha(u)U\) under conjugation yields a
scalar operator; its scalar is its trace \(L_\alpha(U)\) divided
by the dimension \(g^{12f}\). The Heisenberg shifts at \(x\) and
\(-x\) are inverse operators. Applied to the scaled vacuum
\(F_{\rm vac}=v^{12}\operatorname{vac}\), the average gives
\[
 |\mathcal A| L_\alpha(U)F_{\rm vac}
    =\sum_{x\in\mathcal A}g^{12f}L_\alpha(R_xF_{\rm vac})R_{-x}U.
\]
Fix a finite argument \(y_0\) with
\(u_0=F_{\rm vac}(\operatorname{finVec}y_0)\ne0\). Such a point
exists because the standard vacuum has a nonzero finite factor
and its Gaussian has nonzero value at zero. Evaluation at \(y_0\) in the last
identity gives the coefficient comparison with
\(g^{24f}u_0L_\alpha(U)\). All its weights
\(g^{12f}(R_{-x}U)(\operatorname{finVec}y_0)\) are integral.

For \(\mathfrak M=(\bar\delta)\), every \(R_x\operatorname{vac}\)
with \(\delta x\) integral is a sum of standard residue indicators
at \(\mathfrak M\), with root-of-unity coefficients. In the
Schr\"odinger coordinates at these split places, the component
of \(x\) conjugate to the residue coordinates multiplies the
vacuum by the additive character pairing \(\bar x\) with the
residue coordinates. It factors through \(\mathfrak M\), since
\(\bar x\) times \(\bar\delta\) is integral. The components of the
shift outside this part are integral and act on the vacuum by
their semicharacter. Expanding the residue character as its
values times the point indicators gives the assertion, in the
normalized vacuum with no averaging denominator. This is the
finite translation formula in the same standard model as the
entire product. Let \(F_y\) be the vacuum multiplied by the
indicator of the residue fibre \(y\) modulo \(\mathfrak M\) in
these coordinates. The ultrametric inequality supplies a residue
\(y\) for which the rank-\(12\) coefficient satisfies
\[
 \bigl|a_\ell(v^{12}L_\alpha(F_y))\bigr|_\ell
 \ge \bigl|a_\ell(g^{24f}u_0 L_\alpha(U))\bigr|_\ell .
\]

Choose, before any prime or reading,
\[
 t_p=g^{24f}u_0\prod_\sigma D_\sigma.
\]
This is nonzero, and the fixed integral vacuum bound gives
\(\log|\gamma(u_0)|\le O_d(N_1^2)\) for every complex
automorphism \(\gamma\). The other factors contribute at most
\(24f\log g+2fc_*\log(3b)\), since \(c_*=o(b)\) and
\(n_\sigma\le2b\). Their sum is at most \(fmb/\kappa\) for large
\(x\), using \(\log b=o(m)\), \(\log g=O(\log b)\), and the
polynomial bound on \(N_1\). Multiplying the indicator
coefficient by \(n!\) and dividing by \(t_p\), equation
\eqref{eq:weighted-depth-twelve-maxima} gives the required
lower bound \(\ell^{-P_\ell}\). The monomial and the residue
can be chosen individually at each reading. The stage, pair,
degree, modulus, character at the depth point, and \(t_p\)
are fixed.
\footnote{\weightedpin{grh-compact/src/New/Curve/FrameSections.lean}{82-L126}
{The section product in the finite Weil model},
\weightedpin{grh-compact/src/Lib/Weil/Construction/FockThetaFunctions.lean}{61-L112}
{Orthogonality on the lattice with scale \(\bar\delta\)},
\weightedpin{grh-compact/src/Lib/Weil/Lift/Staircase.lean}{148-L166}
{Residue indicators at \((\bar\delta)\)}, and
\weightedpin{grh-compact/src/New/Rows/GluedFormTheta.lean}{117-L143}
{The normalized coefficient on the chosen indicator}.}

\paragraph{Integral point forms for the rank-one tests.}
\teamanchor{proof.weighted-rank-one-point-test}%
Use the partition \(2z_p=12+4+4+(2z_p-20)\) of the rank-one
target. The three flat blocks are taken at the twist \(\psi\),
the trivial twist, and the trivial twist, in that order.
They have ranks \(4,4,2z_p-20\). Since \(m<z_p\le r\), all
three even ranks lie in \([4,t]\) and the last rank is positive.
Let \(\alpha\) be any exponent array on the first block with
\(\sum_b\alpha_{\tau b}=m_\tau\) at each place \(\tau\).
Write \(X^{\tilde\alpha}\) for its extension by degree zero
on the other blocks. The rank-one harmonic model has one
row of length \(m_\tau\) in the jet polynomial at \(\tau\).
\teamanchor{proof.weighted-rank-one-fischer-test}%
Let \(N_\lambda\) be the integral tableau basis in rank \(2z_p\)
for this one-row diagram. Put
\(c_\lambda=[X^{\tilde\alpha}]N_\lambda\). Its entries are
integers. The highest polynomial at each \(\tau\) is a single
variable to the power \(m_\tau\). Thus the Fischer identity for
this Weyl module has hook \(\prod_\tau m_\tau!\). The Fischer
pairing of \(X^{\tilde\alpha}\) with \(N_\lambda\) is
\(\bigl(\prod_{\tau,b}\alpha_{\tau b}!\bigr)c_\lambda\).
The polynomial of the Weyl vector \(c\) is uniquely determined by
its Fischer pairings with the basis. It follows that its harmonic
is
\[
 P_c=\frac{\prod_\tau m_\tau!}{\prod_{\tau,b}\alpha_{\tau b}!}
                    X^{\tilde\alpha}.
\]
For any matrix fixing the first twelve columns of the identity,
the transposed row substitution fixes \(X^{\tilde\alpha}\).
By transpose covariance of the Fischer pairing the Weyl action
fixes \(c\). In this one-row pair all columns come from the
positive sign, so the negative determinant exponent is zero.
The coefficient identity and the fixed-vector identity are
polynomial identities over \(\mathbb Q\). They hold for the same
frames and levels at both coefficient readings, for whichever
monomial and indicator are chosen in each reading.%
\footnote{\weightedpin{grh-compact/src/New/Rows/MonomialCoefficientVector.lean}
{77-L94}{Integral coefficient vector and fixed columns};
\weightedpin{grh-compact/src/New/Exchange/AlgebraicPartDefs.lean}
{215-L223}{Hook of the actual one-row diagram}.}

Fix the rank-\(12\) residue column \(x_0\) modulo \(\mathfrak M\).
In the full point form at \(\operatorname{diag}(1,h_2,h_3,h_1)\),
use the column \((x_0,J_2,J_3,J_1)^{\mathsf t}\) at the chosen
places, where \(J_a\) is the standard column of the indicated
flat block. It is primitive there because \(J_1\) is a unit
column. At a twist place read the marked row of the second
block for the determinant character. The prescribed column
at the chosen places and the prescribed row at the twist
places can each be completed to an invertible residue matrix.
Lift these matrices by Chinese remaindering to one compact
finite adelic element. The class delta at the transported
point pairs with a form by evaluation with no mass factor,
since its stabilizer sum cancels the stabilizer divisor in the
class pairing. Transport through this compact element preserves
the integral gauge. Before this compact transport the point
acts trivially on the first twelve columns, so its gauge on
\(c\) is one or less at every embedding over \(\ell\).

Reduce all four finite data modulo \(\mathfrak M^2 Q_\psi\).
At a chosen place the four indicator constraints are exactly
the four consecutive blocks of the full column constraint.
At a twist place the second block alone contributes
\(\psi_w\) of the marked row. By Chinese remaindering their
product is the datum of the full point form. The real test
is \(P_c\) on the first block and the constant polynomial on
the other three blocks. Choose the splitting character of the
first block such that its product with the three flat splitting
characters is the full target splitting character. The
oscillator representation on the direct sum is the tensor
product in these characters, and the theta sum separates into
the four independent lattice sums. The result is the product
of the four kernels with coefficient
\(\prod_\tau m_\tau!/\prod_{\tau,b}\alpha_{\tau b}!\).
This proves the finite and harmonic identity before taking any
class means.
\footnote{\weightedpin{grh-compact/src/New/Rows/MemberKernels.lean}
{173-L206}{Point form, delta pairing and residue data};
\weightedpin{grh/src/New/OneLift/HeckeValues/TwistedValueHeights.lean}
{313-L431}{The four supported data and their product}.}

Average the point forms over the flat class variables using
the three lifts of the constant function, each multiplied by
its denominator clearer and the coefficient
\(\operatorname N(Q)^{1-d'}\) of its flat rank \(d'\).
As in the finite class formula, multiplication by
\(D_{\mathrm{aut}}\) clears the stabilizer orders in each
class sum. Each flat coefficient is integral away from its
conductor after including the two powers of its denominator
clearer, since the class average uses an integral point kernel
and the flat coefficient pairing has Gram matrix \(1\).
The factor \(\operatorname N(Q)^{1-d'}\) is a unit at \(\ell\).
Thus the resulting \(\Xi\) is integral simultaneously at all
readings over \(\ell\). The scalar averages are exactly
\(A/(h_1v^{d'})\), where \(A\) is the flat integral algebraic
part, \(v\) is the rank-one vacuum scale and \(h_1\) is the
chosen-level class count of the rank-one source. The three
ranks sum to \(2z_p-12\). Substitution gives the normalized
point test with the multiplier \(h_1^3v^{2z_p-12}\).
The marked-period verification supplies its testedness.
\footnote{\weightedpin{grh-compact/src/New/Rows/HeckeValueHeight.lean}
{58-L146}{Flat averages and application on the rank-one auxiliary pair}.}

\paragraph{Character selection for the normalized rank-one test.}
\teamanchor{proof.weighted-rank-one-character-test}%
Keep the same pair and ordinary twist \(\psi=\theta^j\).
Write \(h_1\) for the number of classes of \(U(1)\) at its chosen
level. The \hyperlink{proof.weighted-rank-one-depth-reduction}
{depth calculation} gives a fixed nonzero \(t_p\) and, for
each embedding \(a\) above \(\ell\),
a rank-\(12\) indicator and monomial \(Z^\alpha\) of degrees
\(m_\tau=c_p(\tau)-z_p\). Let \(s_\alpha(g)\) be its kernel at the finite
argument \(g\), with the vacuum factor \(v^{12}\) and with
the coefficient normalization
\(\eta^{\sum_\tau m_\tau}\prod_\tau m_\tau!/\prod_i\alpha_i!\).
Here \(\eta\) is the scaled jet parameter. At \(g=1\) this is
the depth coefficient in the asserted normalization, since an
unscaled derivative multiplies a Taylor coefficient by
\(\prod_i\alpha_i!\). Thus
\(\lvert a(s_\alpha(1)/t_p)\rvert_\ell\ge\ell^{-P_\ell}\).
\footnote{\weightedpin
{grh-compact/src/New/Rows/GluedFormTheta.lean}{27-L38}{Depth on the prescribed twelve-dimensional indicator}.}

Choose \(z\) inverse to the representative of the identity
class, and apply finite Fourier inversion as in
\hyperlink{proof.selected-toric-ratio}{the earlier toric class
calculation}. The class quotient of \(U(1)\) is a finite
abelian group \(C\). Each character of \(C\) lifts to a
weight-zero character, tame on the chosen places and
unramified elsewhere. For each character \(\chi\), form the
average
\[
 B_\chi=\frac1{h_1}\sum_{c\in C}
             \chi(\det(g_cz))\,s_\alpha(g_cz)
\]
on the chosen representatives \(g_c\). Orthogonality gives
\(\sum_{\chi}B_\chi=s_\alpha(1)\). The ultrametric inequality
selects \(\chi\) with
\(\lvert a(B_\chi/t_p)\rvert_\ell\ge\ell^{-P_\ell}\),
including when \(\ell\mid h_1\). The division by \(h_1\)
is in Fourier inversion and costs no power in this lower bound.
In the oscillator kernel the source splitting factor is
\(b(\det g)^{-1}\). Replacing \(b\) by \(b/\chi\) on the
rank-\(12\) block supplies precisely the character in the sum.

Let \(\mathcal X\) be the finite group of the permitted
weight-zero characters on the chosen places. Write \(A_\phi\)
and \(B'_\nu\) for the two fixed untwisted algebraic parts of
the flat pairs of ranks \(4\) and \(2z_p-20\), with characters
\(\phi,\nu\in\mathcal X\). Let \(A_2(\psi)\) be the analogous
algebraic part of the fixed ramified rank-\(4\) pair, so
\(\rho_{q_2}(\psi)=A_2(\psi)/A_2(1)\). Use the same normalized
point form above for \(\chi=\phi\nu\).
The \hyperlink{proof.weighted-rank-one-point-test}{four-block pairing}
with the unscaled full kernel \(K_\psi(g)\) gives
\[
 h_1^3 v^{2z_p-12}\langle\Xi,K_\psi(g)\rangle
 =\frac{\prod_\tau m_\tau!}{\prod_i\alpha_i!}
       A_\phi B'_\nu A_2(\psi)\,
                      K_{12,b/\chi,\alpha}(g),
\]
where \(K_{12,b/\chi,\alpha}\) is unscaled. Testedness comes
from the \hyperlink{proof.weighted-marked-period-reduction}
{rank-one marked period}.

Now \(F_\psi=v^{2z_p}\eta^{\sum_\tau m_\tau}\sum_{c\in C}K_\psi(g_cz)\),
since class translation preserves the unweighted sum. Substitution
in the pairing identity cancels \(v^{2z_p-12}\), leaves \(h_1^{-1}\),
and gives exactly
\[
 \frac{h_1^2\langle\Xi,F_\psi\rangle}{\rho_{q_2}(\psi)}
                  =A_\phi B'_\nu A_2(1) B_\chi .
\]
Choose \(V_{jp}\) before any embedding by
\[
 V_{jp}^{|\mathcal X|}
   =\prod_{\nu\in\mathcal X}
         t_p A_2(1) A_\nu B'_\nu .
\]
Every factor is nonzero. For fixed \(\chi\), the products
\(t_p A_2(1)A_{\chi/\nu}B'_\nu\) have that same total product.
Their largest absolute value at \(a\) is at least \(|a(V_{jp})|_\ell\).
Use the corresponding two auxiliary characters in the point
test. Each complex conjugate of \(V_{jp}\) is bounded by the product
of the uniform archimedean bounds for the three flat parts and
\(t_p\). The flat parts contribute \(3x^{34}\) and \(t_p\) contributes
\(fmb/\kappa\), whose sum is at most \(B_0\) for large \(x\).
This proves \eqref{eq:weighted-rank-one-test} on either reading
of the same \(\psi\). The characters and monomial may vary
with the reading, while \(V_{jp},\psi\) and the ramified pair
remain fixed.\footnote{\weightedpin
{grh-compact/src/New/Rows/MonomialCoefficientVector.lean}{18-L36}{Flat archimedean bound},
and \weightedpin{grh-compact/src/New/Rows/HeckeValueHeight.lean}{154-L216}{Assembly on the fitted twist}.}

Multiply the two inequalities \eqref{eq:weighted-rank-one-test}.
The product of the pairings, with the second conjugated, equals
\(\|F_\psi\|^2\langle\Xi_{j_\iota},P_{\mathbf C F_\psi}
 \Xi_{j_\iota\circ\operatorname{conj}}\rangle\). Write
\(J_{jp}(\psi)=h_1\langle\Xi_{j_\iota},P_{\mathbf C F_\psi}
 \Xi_{j_\iota\circ\operatorname{conj}}\rangle\).
It is nonzero by the two nonzero tests. The fitted rank-one
estimate gives \(|j_\iota(J_{jp}(\psi))|_\ell\le e^{\beta_{j,\ell}}\).
The ratios \(\rho_{q_2}(\psi)\) are real. Hence
\[
 \left|j_\iota\left(V_{jp}\overline{V_{jp}}
          \frac{\rho_{q_2}(\psi)^2}{h_1^3\|F_\psi\|^2}\right)\right|_\ell
     \le \ell^{2P_\ell}|j_\iota(J_{jp}(\psi))|_\ell
     \le \exp(\beta_{j,\ell}+2P_\ell\log\ell).
\]
Use the bound with \(J_{jp}(\psi)\) and multiply by \(h_1^3 I_{jp}\),
by the unit power of the conductor norm and by
\(\rho_{q_2}(\psi)^{-2}\). The first factor is integral
and the last has local logarithm at most \(2v_\ell\) by the fit.
By the Rallis normalization of the rank-one algebraic part,
\(\rho_p(\psi)^{-1}\) is \(\|F_1\|^2\|F_\psi\|^{-2}\) times that
conductor power.\footnote{\weightedpin{grh-compact/src/New/Rows/Separation.lean}{80-L95}{Untwisted norm multiplier},
\weightedpin{grh-compact/src/New/Exchange/AlgebraicPartDefs.lean}{249-L255}{Normalization of the algebraic part}, and
\weightedpin{grh-compact/src/New/Exchange/TriangleLevelExponent.lean}{122-L134}{Its ratio identity}.}
We obtain
\[
 \log|j_\iota(A_{jp}/\rho_p(\theta^j))|_\ell
       \le\log|j_\iota(J_{jp}(\theta^j))|_\ell
                           +2P_\ell\log\ell+2v_\ell .
\]

The factorization of \(\Gamma_j(\psi)/\Gamma_j(1)\) has
\(\kappa/2\) such rank-one factors and \(5m\) ramified factors, all squared,
and a conductor norm power of absolute value one at \(j_\iota\).
The rank-one factors occupy
\(7m/2<z\le7m/2+\kappa/2\) in the unsigned list. The other factors
come from the first \(7m/2\) places in that list and the two signed
lists of lengths \(m\) and \(m/2\). Put
\[
 K_\Gamma=\prod_{j,p} A_{jp}^{\,2},\qquad
 Q_\Gamma(\theta)=\prod_{j=1}^m
            \frac{\Gamma_j(1)}{\Gamma_j(\theta^j)} ,
\]
with \(p\) ranging over the sharp pairs in \(K_\Gamma\).
The sharp bounds and the \(5m\) ramified factors per stage give%
\footnote{\weightedpin{grh-compact/src/New/Rows/Separation.lean}{43-L65}{Factor list}, and
\weightedpin{grh-compact/src/New/Rows/HeckeValueHeight.lean}{218-L236}{Factorization into the rank-one and ramified ratios}.}
\begin{equation}
\label{eq:weighted-aggregate-hecke}
 \log|j_\iota(K_\Gamma Q_\Gamma(\theta))|_\ell
 \le \kappa\sum_j\beta_{j,\ell}+2m\kappa P_\ell\log\ell
             +(2m\kappa+10m^2)v_\ell .
\end{equation}

\paragraph{The joint fitted product.}
\teamanchor{proof.weighted-fit-joint-det}%
In the reciprocal volume ratio the projected pairings and the
compressed determinants occur through the sum
\[
 \frac1R\log\left|j_\iota\left(\prod_j d_j(\theta^j)\right)\right|_\ell
        +2\sum_{j,p}\log|j_\iota(J_{jp}(\theta^j))|_\ell .
\]
The first term uses the dimension-\(R\) test at each stage. The
second uses the \(\kappa/2\) sharp pairs at each stage. Their sum
is at most \(E_\ell=(r+\kappa)\sum_j\beta_{j,\ell}\),
whose finite subsums are bounded by \((r+\kappa)m B_0\).
The coefficient-place comparison uses this joint upper bound.

\paragraph{The determinant along the triangle.}
\teamanchor{proof.weighted-chain-determinant}%
For a fixed \(m\)-fold twist consider the source triangle. Write
\(a_{i,n}:F_{i,n}\to F_{i,n+1}\) for the lifts,
\(d_{i,n}:F_{i,n+1}\to F_{i+1,n}\) for the descents, and \(e_i\)
for the isometric inclusions of principal-level forms in \(F_{i,0}\).
Denote the positive triangle weights by \(q_n\) and the successive
weighted lifts in row \(i\) by \(M_{i,n}\), with \(M_{i,0}=I\) and
\(M_{i,n+1}=q_n a_{i,n}M_{i,n}\).

\paragraph{The arrows on the two triangles.}
\teamanchor{proof.weighted-realized-triangle}%
Use the two actual grids \(V\) and \(Z\), with indices in the order of
\(V_{a,b}\) and \(Z_{a,b}\) below. Put \(z=0,1\), respectively. An
interval of length \(n\) in the triangle of grid \(z\) starts at
offset \(i\). Its space is the node \((0,n)\) of that grid with the
blocks numbered \(i+1,\ldots,i+n\). Its rank is
\(r+z+n(t-z)\). At length zero the space is the \(i\)-fold opposite
of the original space. The isometry \(e_i\) on base-level forms is
pullback by this identification of the unitary groups, with their
quotient measures. The base level is the chosen principal level for
\(V\), and the principal level of \(\mathfrak b=(\bar\pi_G)^{b_1}\)
for \(Z\). Write \(\mathfrak M=(\bar\pi_G)\). The finite ring of the
triangle is \(\mathcal O_E/(\mathfrak M Q_\vartheta)\) for \(V\) and
\(\mathcal O_E/\mathfrak b\) for \(Z\). In the latter all the arrows in
the triangle are untwisted. The order \(q_n\) of the general linear
group of the node of length \(n\) is the weight of its outgoing lift.
It turns the quotient measure on its source into the measure with
principal level of volume one.

Here are the real types and the characters that give the edges. If
\(o\) is the offset, use the doubled parameter with stem entries
\(2\lambda_s+z\), blocks centered at \(2w_1,\ldots,2w_a\) of length
\(t+z\), and blocks centered at \(-2w_{o+1},\ldots,-2w_{o+b}\) of
length \(t-z\). A block of length \(d\) has displacements
\(d-1,d-3,\ldots,1-d\). Lower every entry by one in \(V\) if \(a+b\)
is odd. On the stem the sign at position \(s\) is
\((-1)^{o+A+\min((A-o)_+,b)}\), where \(A=(s+1-\kappa)_+\). All
block signs are positive. In \(V\) a step increasing \(b\) across
the \(l\)-th block uses \(\mu_l^0\vartheta^l\upsilon^{\lfloor N/2\rfloor}\)
at a space of rank \(N\). A step increasing \(a\) uses
\(\mu_l^0\upsilon^{\lceil N/2\rceil}\). In \(Z\) the twist is on the
step increasing \(a\), and both roundings equal \(N/2\).

For an edge increasing \(b\), subtract the weight of the character
at the opposite end in the two parameters. In doubled coordinates
the half-rank terms cancel the lowering, and the common coordinates
are translated by \(2w_l\). The new block is consequently centered
at zero. Among existing stem coordinates the signs flip precisely for
\(A\ge l\), namely the coordinates below \(-2w_l\). All earlier
blocks in this column are above that center. For an edge increasing
\(a\), use the negative of the parameters. The ceiling cancels the
lowering and the common coordinates become \(2w_l-\widetilde\lambda\),
where \(\widetilde\lambda\) is a doubled entry before lowering. They
are all positive, and their signs agree at the two ends. For the
triangle edge deleting its first block, compare the negative parameters
at \((i+1,n)\) and \((i,n+1)\), with character
\(b_i(N)=\alpha_{i,0}(N)^{-1}\upsilon^N\), where \(\alpha_{i,0}(N)\)
is the character of the first upward arrow at rank \(N\).
The shifted common entries are \(-\widetilde\lambda-2w_{i+1}\).
The signs at offsets \(i+1\) and \(i\) are opposite on the stem for
\(A\le i\), where this entry is negative. For \(A\ge i+1\) use
\(\min(A-i,n+1)=1+\min(A-i-1,n)\) to see that they agree. The
surviving blocks lie below \(-2w_{i+1}\), so their shifted entries
are positive. This argument is symmetric in the ends of that edge.
The gaps between opposite signs exceed \(500(N+N')^3\), since a
sign change separates stem positions with different \(A\), or occurs
across distinct block centers. The minimum separation in either case
is bounded below using \(h-t-2\) with \(h=2r^8\), and
\(N,N'\le 2n_m\). The central block and these sign checks give the
Fock matching rule, which supplies the joint harmonic on each edge
with these weights. The harmonics depend on the infinite types and
weights, which are fixed as the finite twist varies.
\footnote{\weightedpin{grh/src/Lib/DiscreteSeries/Theta/BlockRule.lean}{45-L80}{Doubled block rule},
\weightedpin{grh/src/New/UpperBound/Objects/Parameters.lean}{854-L987}{The edge calculations and separated centers}.}

The finite datum of every upward arrow \(a_{i,n}\) is the indicator
of the standard inclusion \(\binom{I}{0}\) in the triangle ring.
Choose the polynomial for a left edge once for each \((i,n)\).
At \(n=0\), conjugating the upward kernel and changing its source to
the opposite space changes its characters to \(b_i\). On the residue
datum this sends \(p(x)\) to \(\overline{p(\epsilon^{-1}x)}\), up to a
scalar, with \(\epsilon\) the integral unit used in the opposite-space
identification. The polynomial is conjugated. In a Petersson pairing
the convolution projection onto the chosen left-edge type replaces
that polynomial by its left-edge harmonic up to a scalar independent
of the forms. Absorbing the scalar into this first descent datum
\(p_{i,0}\) gives
\(\langle x,a_{i,0}e_i y\rangle=\langle d_{i,0}x,e_{i+1}y\rangle\).
This is an identity on the realized form spaces. It follows by
substituting the conjugate kernel into the Petersson integral, then
interchanging the two integrations. The convolution replacement uses
the same joint type at the two ends.
\footnote{\weightedpin{grh/src/New/Determinant/Chain/TriangleLifts.lean}{161-L213}{Conjugating the first kernel and taking its adjoint}.}

\teamanchor{proof.weighted-triangle-point-kernel}%
In either triangle let \(q\) be the base rank and \(D'\) the residue
ring of the first arrow. Thus \(q=r\) and \(D'=\mathcal O_E/
(\mathfrak M Q_\vartheta)\) in \(V\), while \(q=r+1\) and
\(D'=\mathcal O_E/\mathfrak b\) in \(Z\). The finite point datum is
the indicator of the inclusion \(J=\binom{I_q}{0}\) in this ring.
At each of its split places the projection onto the marked rows is
a common left inverse of all matrices in the residue support. The
\hyperlink{proof.split-primitive-support}{split support calculation}
with the identity subgroup at the stated congruence depth gives
the principal subgroup as the only possible support of the
coefficient. That subgroup fixes the residue indicator and its
tame phase is one, so the coefficient there equals the squared
norm of the indicator. With probability Haar on the compact
source its integral is the squared norm divided by
\(|{\rm GL}_q(D')|\). Off the ramified places the datum is the vacuum.

A compact element \(k\) of the base level can be matched on the target
by lifting the block diagonal matrix with blocks \(k\bmod D'\) and
\(I_{r+t-q}\). This follows by inserting \(J\) in the residue action on
both sides. The resulting intertwining of translations is up to a
unitary Weil phase, which cancels in the Gram pairing. Testing the Gram
image of a \(k\)-fixed vector against all translated vectors shows that
the Gram operator preserves the base level.
\footnote{\weightedpin{grh/src/New/Determinant/Chain/TriangleLifts.lean}{305-L449}{Sharp datum, stable range, and preservation of base level},
\weightedpin{grh-compact/src/Lib/SiegelWeil/RallisInner/SharpData.lean}
{451-L470}{Primitive support and exact finite mass}, and
\weightedpin{grh-compact/src/New/Exchange/TriangleArrows.lean}{77-L109}
{The covariance of the actual first arrow}.}

\paragraph{Injectivity of the sharp-point lift.}
\teamanchor{proof.weighted-sharp-base-injectivity}%
Write \(a=r+t=6r\) for the target rank in this calculation and
\(n\in\{r,r+1\}\) for the source rank. Every target here has global
Witt index \(w\le a/2\), so \(a-w\ge a/2>n\). In Yamana's
notation the parabolic isotropic rank of the doubled source is \(n\)
and the evaluation point is \((a-n)/2\). His convergent
Siegel--Weil formula therefore applies. The primitive point support gives the
\hyperlink{proof.unitary-doubling-identity}{graph-orbit unfolding}.
With the same highest polynomial in both copies of the section it
expresses the mixed Petersson pairing of two lifts as the source
pairing with the product of the local vacuum convolutions off the
modulus, times the real integrals, positive point masses at the
specified split places and the target quotient volume. At the split
places in the modulus the principal subgroups act trivially on both
source forms, so the coefficient gives exactly the mass computed
above. At the other finite places use the volume-one compact and
include the prescribed splitting characters in the coefficient. Its
compact-coset contribution is then the identity on the retained
source vectors.

\teamanchor{proof.weighted-sharp-real-pairing}%
At a real place let \(h\) be the actual nonzero joint harmonic and put
\(f(g)=\langle h,\omega(g)h\rangle\). Let \(m\) be the corresponding
unit highest coefficient of the oscillator-side discrete series.
This series is contragredient to the bilinear automorphic input,
so the compression to its lowest-\(K\) highest line pairs \(f\)
against \(\overline m\). Its restriction
\(m_K\) to the maximal compact has
\(c=\int_K|m_K|^2>0\) with probability Haar. Formula
\eqref{eq:weighted-row-coefficient-polar} and compact Schur
orthogonality give, for positive Hermitian \(z\) in the Cartan
decomposition, the \(K\times K\) average of \(f\overline m\) as
\[
 c\int_K M_{-\mu}(u^*zu)^{-1} f(u^{-1}zu)\,du .
\]
Indeed substitute \(p_1u,p_2u\) for the two compact variables in
\eqref{eq:weighted-row-coefficient-polar}. The two compact integrals
project onto the respective copies of the highest line of \(h\),
each with factor \(c\). The real integral is absolutely convergent
by the gap on these edges and boundedness of \(f\). For such \(z\),
\(H=I_a\otimes z^{-\mathsf t}\) is positive Hermitian and
\(\det z=1\). In the normally ordered Fock Gaussian, each pure
quadratic contraction couples opposite signs in the same row. It
therefore kills the chosen harmonic polynomial \(P\) on either side.
The remaining coefficient is the positive scalar
\(\det(H_{--})^{-1}\) times the Fock pairing of
\(P(C^{-1}\,\cdot)\) with \(P\), where
\(C=P_+HP_++P_-\overline H P_-\) and \(P_\pm\) project onto the two
signs. This pairing is the squared norm obtained by substituting
the positive square root of \(C^{-1}\) into \(P\). It is positive
because the substitution is invertible. The same argument applies
to \(u^{-1}zu\), and the real integral is positive.
\footnote{\href{https://msp.org/pjm/2013/264-1/pjm-v264-n1-p11-s.pdf\#page=4}
{Yamana, Theorem~2}, \cite[Theorem~6.3 and (25)--(26)]{real-branch-sun},
\weightedpin{grh/src/Lib/Weil/Lift.lean}{590-L618}{same-vector compact pairing},
and \weightedpin{grh/src/Lib/Weil/Fock/OscillatorPair.lean}
{424-L464}{the Fock Gaussian on the actual signed rows}.}

For the finite convolutions let \(\Lambda\) be the self-dual
oscillator lattice at a place off the modulus, with unit vacuum
\(v\), and put \(v_g=\omega(g)v\). The Weyl translates of \(v_g\)
by representatives of \(\Lambda/(\Lambda\cap g\Lambda)\) are
orthonormal. For two distinct representatives the difference is
outside \(g\Lambda\), so a shift in \(g\Lambda\) with nontrivial
commutator makes their pairing zero. Their pairings with \(v\)
all have the same modulus because \(v\) is an eigenvector for
translations in \(\Lambda\). Bessel's inequality yields
\[
 |\langle v,v_g\rangle|\le [\Lambda:\Lambda\cap g\Lambda]^{-1/2}.
\]
At a split place of residue size \(q\) in one half, let \(L\) be
the source lattice, \(L'=gL\), and set
\(u=\operatorname{length}(L/(L\cap L'))\),
\(w=\operatorname{length}(L'/(L\cap L'))\).
The opposite half has intersection length \(w\) by duality.
The oscillator lattice has \(a\) copies in each half, so the index
is \(q^{a(u+w)}\). For each \(u,w\) there are at most \(q^{n(u+w)}\)
possibilities for \(L'\). Choose the intersection by \(u\) successive
index-\(q\) sublattices and its extension by \(w\) index-\(q\)
overlattices. Each step costs at most \((q^n-1)/(q-1)\le q^n\).
At a nonsplit place let \(Q\) be the residue size over \(E\).
The two lengths are equal by self-duality and equal covolumes.
The oscillator index is \(Q^{au}\), with at most \(Q^{2nu}\)
choices for \(L'\). The compact stabilizer is the source lattice
stabilizer, so each such right coset has mass one.

Consequently the convolution is \(I+T_v\) on the corresponding
source fixed space. Unitarity bounds \(\|T_v\|\) by the total
absolute coefficient mass of the nonidentity cosets, at most
\[
 \bigl(1-q^{n-a/2}\bigr)^{-2}-1
 \quad\hbox{or}\quad
 \frac{Q^{2n-a/2}}{1-Q^{2n-a/2}},
\]
at split and nonsplit places, respectively. Both bounds are less
than one and summable over the places since \(a/2-2n\ge r-2\ge2\).
The coefficients are bi-invariant under the compact subgroup because
the off-modulus datum is its vacuum and the splitting characters are
unramified. The convolutions preserve the source fixed space and
commute across different places.
Use the Petersson Hilbert closure of the fixed finite level and
real type. If \(\delta_v=\|T_v\|\), then
\(\|(I+T_v)^{-1}-I\|\le\delta_v/(1-\delta_v)\), and these bounds
are summable. The operator products and their inverse products
converge in norm and remain inverse. The mixed Petersson relation
tested against all forms of the fixed level and type, which are
dense in this closure, shows that a vector whose lift is zero is
killed by the invertible product. This
proves injectivity on the full base level for each of the two
arrows, including the deep core level in \(Z\). With the realized
first adjoint this gives \(d_{i,0}a_{i,0}e_i=e_{i+1}B_i\) with
\(B_i\) injective.

\teamanchor{proof.weighted-auxiliary-injection}%
At the trivial twist the auxiliary lift \(W_j\to I_{6r}\) has this
same point datum at the specified modulus \(\mathfrak M\), as
follows by Petersson adjunction with the definite class-sum metric.
The auxiliary construction below has \(Q_\psi=\varnothing\)
and no source translates to add at that twist. Its off-modulus
datum is the vacuum, its source rank is \(r\), and the real
scalar uses the actual joint harmonic for this edge. The argument
just proved gives injectivity on the full principal-level
source, which supplies the descent of a good form in the
\hyperlink{proof.weighted-comparison-carrier-approximation}{holomorphic
peak argument}.\footnote{\weightedpin
{grh-compact/src/New/Exchange/CharacterLiftNorm.lean}
{33-L75}{The auxiliary point datum at the trivial twist}.}

\paragraph{Complementary square integrals and the realized commutations.}
\teamanchor{proof.weighted-triangle-square-laws}%
In a square let the two exterior spaces have orthogonal sum \(W_*\)
of even rank \(N\), and let the intermediate spaces have ranks
\(a,b\) with \(a+b=N\) and \(e=a-N/2\). Their negative sign counts
add to the count \(q_*\) of \(W_*\). To define partners over the residue
ring \(R_*\) in the order \(b\oplus a\), take \(k\in\mathrm{GL}_N(R_*)\)
and write \(Y_k=k\binom{I_b}{0}\), \(X_k=(0\ I_a)k^{-1}\). The two finite
functions are supported on these frames, with
\(\Phi(X_k)=\chi(\det k)\Psi(Y_k)\), where \(\chi\) is the residue
character of \(\nu\). The product of the edge characters at spaces
whose ranks add to \(N\) must be \(\nu^{-1}\upsilon^{N/2}\).
The weight \(\mu\) of \(\nu\) satisfies \(N\le2\mu+2q_*\) and
\(N\ge2q_*+2\). Let \(S\) be the nonempty split support of the
square modulus, disjoint from its conjugate. The local determinant
characters on \(S\) are nontrivial of conductor one. For \(e\ge0\)
put
\[
 \Pi_{\nu,S}(e)=\Bigl(\prod_{v\in S}q_v\Bigr)^{e^2}
                  \prod_{l=1}^{e}|L^S(l,\nu)|^2 .
\]
Complementary integration interchanges the two kernels on matched
data with scalar \(c\) satisfying \(|c|^2=C\Pi_{\nu,S}(e)\), where
\(C>0\) is fixed in the twist. The constant for the map between form
spaces is \(\kappa_*/c\), with \(\kappa_*\) independent of the twist.
The product for \(e=0\) is one. For \(e<0\) take
\(\Pi_{\nu,S}(e)=\Pi_{\nu,S}(-e)^{-1}\).

\paragraph{Primitive support and integration in the split model.}
\teamanchor{proof.weighted-grid-primitive}%
All measures in this calculation are the measures of the fixed Weil
datum. In the finite named model take a nonempty set \(S\) of split
places disjoint from \(\bar S\). Fix the finite Bruhat--Schwartz function
\(\Psi\) in that model and the vacuum at the other places. Write \(G\)
for the external hermitian space, \(H\) for the integrated one, and
\(Y\colon G\to H\) for a matrix in the named coordinates. Suppose
\(Y\) is row-surjective at one named place wherever \(\Psi(Y)\ne0\).
Put \(h_G(g),h_H(h)\ge1\) for the heights from the moved lattices and
write \(M_j(v)=\int |v(u)|(1+\|u\|)^j\,du\) for the Fock moment, with
the normalized coherent-state measure. Then the kernel is bounded by
\[
 |\Theta_{\Psi,v}(g,h)|\,h_H(h)^B\le C M_j(v)h_G(g)^A
\]
for every prescribed \(B\ge0\), for suitable \(A,j,C\).
The constants here may depend on the fixed finite function.
For the initial bound use reduction with the inverse first
minimum as height. A real move of operator norm at most two
halves heights above the fixed cutoff by expanding the short
isotropic subspace and contracting its dual. After
\(O(\log h_G+\log h_H)\) moves, Mahler compactness gives
rational translates in fixed compact sets. Operators of
the compact set of real moves have a common norm bound
in the coherent model, so the norm of the reduced
operator is bounded by a fixed power of the heights.
The rational translates leave the theta functional
unchanged. This gives the initial bound for
\(\Lambda\) on the coherent datum norm.

Here is the short-shift argument for this estimate. For large
\(h_H(h)\), take a nonzero short rational vector \(d\) in the moved
integral lattice of \(H\), of real length \(O(h_H(h)^{-1})\).
Its hermitian norm is integral and has all real absolute values less
than one, so \(d\) is isotropic. Take an \(E\)-basis \(C\) of \(G\)
in the moved integral lattice with real lengths \(O(h_G(g)^{A_0})\).
Indeed the moved lattices have fixed covolume, and the product bound
on the successive Euclidean minima bounds the last minimum by a
fixed power of the inverse first minimum. A full-rank set over
\(\mathbb Q\) contains an \(E\)-basis. In the original rational
coordinates the shifts are \(d z C^*J_G\), with \(z\) an adelic row.
All lie in one isotropic space of \(\operatorname{Hom}(G,H)\).
The pullback \(\Lambda=\Theta\circ\omega(g,h)\) is invariant under
rational parameters \(z\), and the shift operators \(T_z\) satisfy
\(T_zT_{z'}=T_{z+z'}\), since the Heisenberg pairing on this space is
zero. Define \(\mathcal P\) on these invariant functionals by
\(\mathcal P\Lambda(f)=\int_{[z]}\Lambda(T_z f)\,dz\) with probability
Haar on the additive quotient. It is idempotent and
\(\mathcal P\Lambda\) is invariant under every adelic parameter.

At a selected split pair, restrict \(z\) to the half conjugate to the
matrix coordinates. In the fixed named model this multiplies the
finite function by an additive character of the pairing with \(Y\).
In the transported coordinates \(d',C'\) its linear coefficient is
\(d'{}^*Y C'\), up to the fixed invertible Weil and frame factors.
It is nonzero because \(d'\ne0\), \(C'\) is invertible and \(Y\) is
surjective at the selected place. In particular
\(\mathcal P\Lambda(f)=0\) for the datum in question. To see this
directly, each point of the compact support has a modulation whose
character is different from one on an open neighborhood. Refine a
finite subcover to a locally constant partition. The invariant
functional vanishes on each summand by that modulation. This
argument also works with the columns \(C\) restricted to a frame
on which \(YC'\) is surjective.

Choose representatives for \([z]\) in a fixed bounded real set with
finite part in \(D\widehat{\mathcal O}_E\), where \(D\) is a fixed
integer large enough to fix the given finite vector under the
shifts. Such representatives exist by additive approximation and
the compact real quotient of the resulting fractional lattice.
The moved finite columns of \(d,C\) are integral before this
multiplication by \(D\). Their real shifts have length at most
\(C_0 h_G(g)^{A_0}/h_H(h)\). Since
\((1-\mathcal P)^m=1-\mathcal P\) for \(m\ge1\), the value
\(\Lambda(f)\) equals the integral of
\(\Lambda((1-T_{z_1})\cdots(1-T_{z_m})f)\) over these representatives.
For shifts of length at most one, expand each coherent vector
\(\rho_\infty(u)v_{\rm vac}\) using the Heisenberg law and the
fundamental theorem of calculus in each shift. The norm of the
mixed derivative of order \(m\) is bounded by \(C_m(1+\|u\|)^m\).
The bound follows by differentiating the Gaussian translate and
its phase, with translation lengths at most \(m\) and all directions
normalized. Integrating the coherent expansion gives the product
of the shift lengths times \(C_m M_m(v)\). Use the moderate-growth
bound for the theta functional in this coherent-state norm, with
a fixed height exponent \(A_1\). Taking \(m>A_1+B\) absorbs the
height of \(h\). If the bound on a shift exceeds one, then
\(h_H(h)\ll h_G(g)^{A_0}\), and the initial moderate-growth bound
gives the assertion with a larger \(A\).

For a Fourier coefficient along a fixed isotropic frame \(K\) with
retained frame \(I\), one also needs the sum of absolute values of
the rational slice terms. Use finite sums of the slice functionals
with unit phases chosen to make their values nonnegative on the fixed
datum. The initial operator bound for these sums is uniform in the
finite subset and the phases. In fact an adapted coherent
vacuum has slice norms bounded by a Gaussian in the real
variables and a box indicator in the finite variables.
Their translates have uniformly bounded sums on the rational
lattice, by periodizing the Gaussian and the compact box.
The coherent integral gives the same bound times the datum
norm for an arbitrary vector. Rational reduction in the
integrated variable permutes the rational slices, giving the
initial height bound for the sum of their absolute values.
For this functional take \(C\) to be a common
integral multiple of \((K\mid I)\), in rational coordinates. The
shifts kill \(K\), so the rational invariance of each slice
functional follows from rational invariance of the retained theta
functional. The preceding argument bounds every finite sum
uniformly in \(h\) if the restriction of \(Y\) to this transported
frame is row-surjective. This gives the absolute slice sum bound
on the finite-volume quotient. For the top row datum this restriction
is invertible by the stacked determinant. For the top column datum
its first square block is invertible at the conjugate places.
At a parabolic point the frame changes by an invertible block
triangular matrix. For the lower row test the full row rank follows
from invertibility of the stack of \(P\) and \(Y\). These are the tests
on which termwise integration and the cone estimate below are used.
\footnote{\weightedpin{grh-compact/src/Lib/Weil/DualPair/KernelDecay.lean}
{39-L148}{Short shifts, split support and decay},
\weightedpin{grh/src/Lib/Weil/DualPair/ThetaLiftConvergence.lean}
{143-L210}{Reduction with logarithmically many bounded moves},
\weightedpin{grh/src/Lib/Weil/MixedModel/ZeroTerm.lean}
{689-L717}{Summing norms of the rational slices},
\weightedpin{grh/src/Lib/Weil/DualPair/ShortIsotropicVectors.lean}
{544-L598}{Repeated differences in the coherent model},
\weightedpin{grh/src/Lib/Weil/DualPair/ThetaIntegrals.lean}
{174-L185}{Rational invariance on each slice}
and \weightedpin{grh/src/Lib/Weil/DualPair/ThetaIntegrals.lean}
{230-L326}{Unit phases and uniform absolute slice sums}.}

\paragraph{The two realized line models.}
\teamanchor{proof.weighted-grid-plane-geometry}%
In the ordered frames \(W^+=L\oplus W\oplus L^*\) and
\(V^+=B\oplus V\oplus B^*\), dual isotropic generators pair to one.
The center parameter of the outer line radical multiplies the term
indexed by \(x\in V^+(E)\) by \(\psi(z\,x^*J_{V^+}x)\).
The radical coordinates here are those with central parameter
twice the matrix parameter of the general center formula. The
central mean is therefore the indicator \(x^*J_{V^+}x=0\).
For \(x=0\) the remaining mean is the retained theta functional
at the zero slice. For \(x=\xi b_0\ne0\), where \(b_0\) generates
\(B\) in a rational Witt frame, the other radical coordinate
\(u\) translates the first slice by \(\xi b_0 u^*\) in the
retained hom space. The map \(u\mapsto \xi u^*\) is an automorphism
of the corresponding additive quotient. The normalized zero term
of the seesaw calculation therefore identifies this mean with the
second zero projection, along \(\operatorname{Hom}(W,B)\).
At fixed integrated point the slice norms are summable and the
radical phases have modulus one, so the expansion is termwise
integrable over the compact radical. The primitive-support bound
is uniform in the integrated group on this compact family of
external points.

For clarity the two zero projections can be read in the same
hom space before using the finite tests. The trace hermitian
form pairs a map from block \(L_1\) to block \(B_1\) with the map
from its dual block to the dual of \(B_1\). In the lines just
chosen, the common isotropic summand is
\[
\operatorname{Hom}(W\oplus L^*,B)\oplus\operatorname{Hom}(L^*,V).
\]
Its dual is
\(\operatorname{Hom}(W\oplus L,B^*)\oplus\operatorname{Hom}(L,V)\),
as one sees by pairing the matrix units. The orthogonal retained
piece is \(\operatorname{Hom}(W,V)\) together with the hyperbolic
plane \(\operatorname{Hom}(L,B)\oplus\operatorname{Hom}(L^*,B^*)\).
After the common projection the last two models keep opposite
coordinates of this plane. Thus they are related by the
self-dual Fourier transform on \(\mathbb A_E\), with the
dual variable read through conjugation. This is the plane on which
the Tate calculation below takes place. The split basis is the
original adapted integral frame, so its finite scaling and the
real Gaussian are the ones in the normalized zero-term formula.

For the standard row datum \(P I_A'\), the zero lower matrix
annihilates the added first line together with the fixed rows,
so its full stacked determinant is zero. Nonzero isotropic
vectors form the one orbit of rational isotropic lines by
Witt extension. The stabilizer of \(B\) has radical and Levi
\(E^\times\times U(V)\) on rational points. The line terms are
invariant under its adelic radical because the two zero terms
are invariant under the corresponding unipotent operators. The
line-orbit sum unfolds to the quotient by that stabilizer.
In Iwasawa coordinates the Levi sends the generator of \(B\) to
\(\lambda^{-1}\) times itself. Its action on the radical column
has modulus \(|\lambda|^{-a}\) and that on the central
\(E^+\)-coordinate has modulus \(|\lambda|^{-1}\).
This gives the quotient factor \(|\lambda|^{a+1}\), with one
fixed positive change of Haar measure. The two line arguments
are respectively \(\lambda u\) and
\(\overline{\lambda u}^{-1}\) at an external line scalar \(u\).
The two determinant characters and the zero-term dilation
on the ambient side give the Weil multipliers of the next
calculation. In particular the formula at arbitrary \(u\) is
used here only when the lower two ranks are equal, so that the
extra modulus \(|u|^{a+1}\) combines with the Iwasawa factor
and \(u\) can be removed by multiplicative translation.
\footnote{\weightedpin{grh/src/Lib/SiegelWeil/Tower/ConstantTerm.lean}
{62-L207}{The central filter and the nonzero-vector slice},
\weightedpin{grh/src/Lib/SiegelWeil/Tower/ConstantTerm.lean}
{643-L717}{Grouping and unfolding the line orbit},
\weightedpin{grh/src/Lib/SiegelWeil/Tower/LineTerms.lean}
{109-L180}{Levi action in the two coordinates}
and \weightedpin{grh/src/Lib/SiegelWeil/Tower/LineTerms.lean}
{190-L274}{The common projection and Poisson on the plane}.}

On the column side the upper function is read on transposes at
\(\bar S\), with the \(W\)-projection of the adapted frame. In this
reading the transpose of the dual projection at \(S\) is the
inclusion at \(\bar S\), by the hermitian frame relations. The
finite column datum after a slice at \(x\) is consequently the
lower \(\chi^{-1}(\det(PY^{\mathsf t}))\) on integral \(Y\),
multiplied by the vacuum factor if the finite component of the
scaled slice argument is integral, and zero otherwise. In the
nonzero isotropic term the remaining theta functional is invariant
under translations along that line of the integrated space,
after transport by the same integrated point \(h\). Its
local action at a split pair is a modulation with coefficient
\(x'{}^*Y\), where \(x'\ne0\) is the moved line vector and
\(Y\) has full row rank on the retained named datum. The
partition argument above makes this functional zero there.
Thus precisely the zero slice remains. The rank twist loses
the splitting character \(\upsilon\) for the removed plane,
and the constant term at \((1,g)\) is the lower integral on
the real zero slice times the fixed nonzero factor
\(\kappa_0\) of the adapted vacuum.
\footnote{\weightedpin{grh/src/New/Determinant/Exchange/RemovingPlanes.lean}
{321-L405}{Projection at conjugate places and primitive killing}
and \weightedpin{grh/src/New/Determinant/Exchange/RemovingPlanes.lean}
{413-L464}{Normalization at the zero vector}.}

\paragraph{The signed coefficient at the top of the tower.}
\teamanchor{proof.weighted-grid-signed-top}%
Use \(W_1=\mathbb H\perp W_*\), of even rank \(n_1=2j\), and
\(X=\mathbb H^b\perp W_1\), with adapted splitting \(D\) over
\(W_1\). For the small space at the top take \(U=V_{\rm small}^{-}\)
with its form negated. Thus the other space at the top is
\(U\perp W_1\). If \(V_{\rm small}\) has \(q_b\) negative signs,
\(U\) has \(b-q_b\), and this top intermediate space has
\(b-q_b+q_*+1\) negative signs. Removing \(b+1\) planes on the
large side leaves \(q_*-q_b\), the sign count of \(V_{\rm large}\).
The rank difference at the top is \(n_1\), so the signed excess
is \(j\).

In the integral frame \((L,K,I)\) of \(D\), put \(P=L'\), the
dual row block. Right multiplication of the stacked row matrix
by the frame makes its determinant the determinant on the
\((K,I)\) columns, times an integral unit from the frame and
the prescribed row order. Call its character phase \(\epsilon\).
On the small side the same \(L'\) reads the \(K\) columns at
the conjugate split places when the data are transposed.
At a point \(n(c,z)m(q,g)\), the central character of the
rational slice at a frame \(x\) reads its hermitian Gram matrix
\(x^*Jx\). The coefficient with index \(\jmath T\), for the
chosen purely imaginary Weil element \(\jmath\), is the sum
of the orbit integrals of the slices with this Gram matrix.
Here \(\psi(\tfrac12\operatorname{tr}(zT))\) is the common
radical phase in the unscaled matrix parameter. Its normalization
is fixed by the normalized mixed-model zero term. The absolute
slice estimate just proved gives the needed integrability.

Write \(J_U\) for the form of \(U\), \(x_U\) for its adelic frame
and \(\alpha\) for the fixed scale idele. A self-dual test is
\[
 T=J_U,\qquad q=\alpha x_U B,\qquad
 B_f k\in U(1_b)(\mathbb A_f)
 \quad\text{for some }k\in{\rm GL}_b(\widehat{\mathcal O}_E).
\]
Here \(U(1_b)\) is the group of the finite form \(1_b\).
Witt extension gives a single rational orbit on each side,
represented by \(\binom{1_b}{0}\) and \(1_b\). Their stabilizers
are \(U(W_1)\) and the trivial group, respectively.
On the large side unfold first through
\(U(U)\times U(W_1)\) in \(U(U\perp W_1)\).
If a finite translate does not lie in the product of this
block group with the standard integral subgroup, the term
vanishes. Indeed a nonzero value requires that the moved
first columns be integral. Right multiplication by \(k\)
then gives an integral orthonormal frame. Integral Witt
extension carries it to the first standard columns,
and an isometry fixing those columns acts in their orthogonal
complement. This is the stated block-group support.

After the allowed finite translate we can use a block point
\(y\oplus y_1\). Put \(B_y=y^{-1}B\) and \(x=(B_y)_f\).
Both slices require integral \(x\). Since \(xk\) is unitary
when \(x\) is integral, the determinants of \(x\) and \(x^{-1}\)
are both integral. In particular \(\det x\) is a finite unit.
On the large side the determinant of
\(\bigl(\begin{smallmatrix}x&*\\0&Y_1\end{smallmatrix}\bigr)\)
factors as \(\det x\det Y_1\). The retained finite vector is
the vacuum for \((W_1,U)\) times the signed determinant datum
for \((W_1,W_1)\), multiplied by
\(\epsilon\chi^{-1}(\det x)\). On the small side it is the
vacuum times \(\chi^{-1}(\overline{\det x})\).
The finite unit identity relating the two readings is
\[
 \chi^{-1}(\det x)\widetilde\nu_f(\det x)
       =\chi^{-1}(\overline{\det x}),\qquad
 \widetilde\nu(\lambda)=\nu(\lambda/\bar\lambda).
\]
The integrated characters on \(U(U)\) differ by
\(\nu(\det y)^2\). In the product Weil action this is the factor
from \(\widetilde\nu(\det y)\), since \(\det y\) has norm one.
The common radical phase and the retained Heisenberg
translate \(x_U^{-1}c^*J_{W_1}\) agree on the first block.

\emph{Real slices.}
Let \(v\) be the prescribed real vector on \((W_1,U)\)
and \(w\) a real vector on \((W_1,W_1)\). In the real chart put
\(J=\operatorname{diag}(J_U,J_{W_1})\) in sign coordinates and
\(Z_0=(1_b,0)^{\mathsf t}\). Near its nondegenerate \(b\)-plane
use the smooth projection
\[
 P(Z)=Z(Z^*JZ)^{-1}Z^*JZ_0 .
\]
It depends only on the plane of \(Z\). Near \(Z_0\), choose
a smooth unitary \(k(P)\) taking the plane of \(Z_0\) to that
of \(P\). This follows by the local section of the orbit
submersion on nondegenerate planes. Take a nonnegative smooth
compactly supported bump \(\sigma\) near \(Z_0\), with
\(\sigma(Z_0)=1\). For the small vector prescribe the real slice
at \(B\) to be
\[
 \kappa\,\sigma(Z_0 B)
       \left|\prod_\tau(\det B_\tau)^{j+\mu_\tau}\right|^2 v .
\]
The scalar \(\kappa\) is chosen using the integral of the
displayed bump and determinant weight, which is positive.
With the Fourier measure of this mixed model the slice at zero
in the opposite splitting is then \(v\).
For the large vector put \(Z'=k(P(Z))^{-1}Z\) and prescribe
the slices
\[
 \kappa\,\sigma(P(Z))\,\sigma(Z')\,
       \prod_\tau\overline{\det Z'_{1,\tau}}^{\,2\mu_\tau}
       \omega_\infty(1,k(P(Z)))(v\boxtimes w).
\]
Here \(Z'_1\) is the first square row block and the
real target weight in \(\omega_\infty\) is \(j-\mu_\tau\).
Define the slices by zero outside the chart. The two cutoffs
make the support compact in \(Z\), with rank bounded away
from degeneration, so the extension is smooth. In the
mixed model these are compactly supported smooth families
of retained Schwartz vectors. Inverse partial Fourier
transformation produces the required vectors with all
Fock moments finite. Every nonzero large slice has real
Gram matrix congruent to \(J_U\).

To compare the periods, set \(Z=h_\infty^{-1}Z_0\). If the
bump at \(P(Z)\) is nonzero, \(h_\infty k(P(Z))\) preserves
the plane of \(Z_0\) and is block diagonal. A left translate
by its inverse in \(U(U)\times U(W_1)\) cancels it. This
translate also aligns the slice at \(Z B_y\) with the slice
at \(Z_0 (B_y)_\infty\), since \(P(ZB_y)=P(Z)\), with the
first block of the left translate included in the former
slice argument. Hence the two real factors have ratio
\[
 \frac{\rho(h_\infty)}
 {\widetilde\nu_\infty(\det (B_y)_\infty)
            |\det (B_y)_\infty|_{E,\infty}^{j}},
 \qquad
 \rho(h_\infty)=\sigma(P(h_\infty^{-1}Z_0)).
\]
The determinant relation follows place by place from
\[
 \overline{d}^{\,2\mu_\tau}
 =\frac{|d|^{2(j+\mu_\tau)}}{(d/|d|)^{2\mu_\tau}|d|^{2j}} .
\]
In the remaining block the integral of \(w\) is the constant
equal-rank period \(I_\nu(w)\) on \(W_1\). Its constancy
uses the one-plane formula as explained below.

The quotient stages supply a fixed positive factor \(R\).
Indeed the integral over the block subgroup, after quotient
by the rational stabilizer, is unchanged by its left block
translates. The weight \(\rho\) depends only on the plane
of \(h_\infty^{-1}Z_0\), and the finite indicator is also
left block invariant. On the real quotient the support of
\(\rho\) is in the image of the compact set of representatives
\(k(P)^{-1}\) over the bump. The finite support is the image
of the standard compact subgroup. Consequently the quotient
weight has finite mass, positive because it is positive near
the identity coset. Its mass is \(R\) for every character,
including a zero value of the common small coefficient.
Writing \(d_0=\det(\alpha x_U)\) and \(\kappa_D^+,\kappa_D^-\)
for the two fixed vacuum factors, the multiplier in the
original coefficient identity is
\[
 C_{\rm top}=
 R\,\frac{\kappa_D^+}{\kappa_D^-}\,|d_0|_E^{j}\,
     \epsilon\,\widetilde\nu(d_0)\, I_\nu(w).
\]
In fact the Levi modulus ratio is \(|\det q|_E^j\) and its
character ratio contributes \(\widetilde\nu(\det q)\);
\(\det q=d_0\det B\). The real determinant quotient, the
finite unit identity and the \(U(U)\)-character adjustment
cancel the factors containing \(B_y\) and \(y\). This gives
the displayed scalar including the conjugate finite reading.
Its squared modulus is a fixed positive multiple of
\(\Pi_{\nu,S}(j)\mathcal R_\nu(w)\) with
\(\mathcal R_\nu(w)=|I_\nu(w)|^2/\Pi_{\nu,S}(j)\).
\footnote{\weightedpin{grh-compact/src/New/Exchange/RealPair.lean}
{59-L132}{Signed real section},
\weightedpin{grh/src/New/Determinant/Exchange/LargeTriple.lean}
{485-L609}{Integral orbit and the two finite slices},
\weightedpin{grh/src/New/Determinant/Exchange/LargeTriple.lean}
{839-L943}{Character and Levi factors}
and \weightedpin{grh/src/New/Determinant/Exchange/LargeTriple.lean}
{1031-L1078}{Coefficient measures and staging}.}

For an invertible \(T\) at a real place outside the class of
\(U\), every small orbit is empty. With the supported large
vector just constructed, every large slice at that index is
zero. One also needs the all-negative auxiliary case for
arbitrary real input during the equal-rank nonvanishing
argument. At a place where the Gram form of a frame of rank
\(b\) has a positive direction, move that direction to a
positive coordinate by a real isometry. The projection of
the frame to the negative rows has rank at most \(b-1\).
There are \(b+q(W_1)\) negative rows. For a coherent-state
label, require in addition orthogonality to its
\(q(W_1)\) columns indexed by negative signs of \(W_1\).
A unit negative row remains. Rotation of its line by
\(\bar\eta\) fixes the frame. Its vacuum exponent on the
remaining oscillator is
\[
 n_1-q(W_1)-(j-\mu_\tau)=j-q(W_1)+\mu_\tau>0
\]
since \(q(W_1)\le j\) and \(\mu_\tau>0\).
The coherent label has zero coordinates in the terms
on which the rotation would contribute negative powers
of \(\eta\). Its circle mean therefore vanishes after that
vacuum factor. The compact stabilizer average is zero
on all coherent vectors and hence on every Schwartz vector
by the coherent-state integral. Haar on the stabilizer is
right invariant under its compact subgroup, so the unfolded
orbit integral vanishes. This also covers the alternative
sign hypothesis in the top coefficient calculation when
the real bump restriction is absent.
\footnote{\weightedpin{grh-compact/src/New/Exchange/NonnegativeVanishing.lean}
{21-L79}{Compact stabilizer and circle mean}
and \weightedpin{grh/src/New/Determinant/Exchange/NegativeIndices.lean}
{50-L104}{Positive coordinate and the remaining negative row}.}

\paragraph{Propagation from the self-dual tests.}
\teamanchor{proof.weighted-grid-standard-extension}%
Let \(\mathcal D\) be the difference of the large integral and
\(C_{\rm top}\) times the small integral. Write \(\Lambda\) and
\(I_{\rm small}\) for those integrals. At a place \(w\in S\) let
\(\tau_w(x)=\nu((x,x^{-1}))=\widetilde\nu_w(x)\), in the chosen
split coordinate. Its restriction to integral units is \(\chi_w\).
Use the full character, with its value at the uniformizer, and put
\[
 \eta_w(h)=\tau_w(\det h)|\det h|_w^{N/2}.
\]
Here \(N\) is the ambient rank for the two complementary integrals
being compared. Multiplying over the selected pairs gives
\(\eta=\prod_{w\in S}\eta_w\).
The factor \(|\det h|_w^{N/2}\) is the square root of the
additive Jacobian of \(X\mapsto hX\) on \(M_{a,N}(E_w)\).
In the split vacuum convention the unitary factor of the underlying
action on the position space \(T_w\) is
\(\upsilon_{E,w}(\det T_w)^{-1}\).
Indeed, the determinant valuations of the intersection lattices
for \(T^{-1}\) in the two halves differ by
\(-v_w(\det T_w)\), and the two local values of
\(\upsilon_E\) are inverse under the split identification.
On the large integrated group, postcomposition by \(h'\) has
phase \(\upsilon_{E,w}(\det h')^{-N}\).
Multiplication by \((\nu^{-1}\upsilon^{N/2})(\det h')\) leaves
\(\tau_w(\det h')^{-1}\), since
\(\upsilon((x,x^{-1}))=\upsilon_{E,w}(x)^2\).
The action with pullback \(\Phi(hX)\) uses \(h'=h^{-1}\).
On the small group the character adjustment by \(\nu^2\) in the
opposite coordinates gives \(\nu\upsilon^{N/2}\) on the integrated
group. The element with coordinate \(l\) in the chosen half has
coordinate \(l^{-\mathsf t}\) in the other half. Transposing the
pullback in that half gives \(\Psi(Yl)\), with factor
\(\eta_w(l)\). Right invariance of the two quotient measures
therefore gives
\[
 \eta(h)\Lambda(\Phi(h\,\cdot))=\Lambda(\Phi),\qquad
 \eta(l)I_{\rm small}(\Psi(\,\cdot\, l))=I_{\rm small}(\Psi).
\]
These identities use the original quotient measures.

Write \(E_S=\prod_{w\in S}E_w\) and
\(\mathcal O_S=\prod_{w\in S}\mathcal O_w\),
with \(K_j=\mathrm{GL}_j(\mathcal O_S)\).
Use Haar measures of compact mass one in the following
pushforwards of the finite sections:
\[
 f_X(k)=\int_{\mathrm{GL}_a(E_S)}\eta(h)\Phi(hX_k)\,dh,
 \qquad
 f_Y(k)=\int_{\mathrm{GL}_b(E_S)}\eta(l)\Psi(Y_kl)\,dl.
\]
The integrals are over compact subsets for each \(k\), since the
supports are compact in the full-rank loci.
On primitive frames the pushforwards give the primitive reductions.
For example \(X\mathcal O_S^N\) is the lattice generated by
the columns of \(X\). On the support it equals \(h\mathcal O_S^a\) for
finitely many right cosets \(hK_a\). The part on each such coset
is moved to primitive \(X\) by \(X\mapsto hX\) and weighted by
\(\eta(h)\), by the measured row covariance. The compact average
with the same covariance yields \(f_X\); the column calculation
uses right multiplication on the basis of the column frame.
Thus the integrals with \(\Phi,\Psi\) equal the integrals with
their primitive reductions, with no extra measure factors.

For the block upper triangular parabolic \(P_{b,a}\) stabilizing
the first \(b\) coordinates, let \(u,v\) be the two diagonal
blocks of \(p\). Then \(X_{kp}=v^{-1}X_k\), \(Y_{kp}=Y_ku\), so
\[
 f_X(kp)=\eta(v) f_X(k),\qquad
 f_Y(kp)=\eta(u)^{-1} f_Y(k).
\]
Put \(\tau=\prod_{w\in S}\tau_w\) and
\(|\det g|_S=\prod_{w\in S}|\det g_w|_w\).
Equivalently, the sections
\[
 |\det g|_S^{a/2} f_X(g^{-1}),\qquad
 |\det g|_S^{-b/2}\tau(\det g)^{-1} f_Y(g^{-1})
\]
have the same left transformation law
\(\delta_{P_{b,a}}(p)^{1/2}\tau(\det v)^{-1}\), where
\(\delta_{P_{b,a}}(p)=|\det u|_S^a|\det v|_S^{-b}\).
Consequently the integral partner relation on \(K_N\) extends,
by \(\mathrm{GL}_N(E_S)=K_NP_{b,a}(E_S)\), to
\(f_X(k)=\eta(k) f_Y(k)\). For any full-rank column test
the relation specifies its primitive row partner from \(f_Y\)
on \(K_N\), and defines \(\mathcal D\) linearly in that test.

To see covariance under the outer group, let \(\beta_X,\beta_Y\)
be the exterior characters before the opposite-space reading,
with \(\beta_X\beta_Y=\nu^{-1}\upsilon^{N/2}\).
In the latter reading the column exterior character is
\(\beta_Y^{-1}\upsilon^b\). For \(g\) at a chosen half the
pullbacks \(X\mapsto Xg\) and \(Y\mapsto g^{-1}Y\)
have Weil multipliers of modulus \(|\det g|_w^{a/2}\) and
\(|\det g|_w^{-b/2}\), respectively. The ratio of their unitary
factors is \(\upsilon_{E,w}^{a-b}(\beta_Y^{-1}\upsilon^b)/\beta_X\),
evaluated in the split coordinate. This is \(\tau_w(\det g)\).
Their full ratio is therefore \(\eta_w(g)\), which is exactly the
factor obtained by replacing \(k\) by \(g^{-1}k\) in the
normalized partner relation. The kernel of \(\mathcal D\) is
stable under the outer split group. Its image retains the
standard integral level off the selected pairs. The
determinant-one shears below have Weil multiplier one.
For the scalar translate by \(-\varpi^{-1}\) and the
dilations \(D_d k\), divide each normalized column translate
by its nonzero Weil multiplier before summing. The multiplier
may depend on \(k\). The linear kernel therefore contains the
bare pullbacks used in the support and multiplicity calculations.

On the exterior blocks of the triangle the character product
uses the opening character \(\alpha_{i,0}^{-1}\upsilon^{N'}\)
at rank \(N'\) and the upper character of step \(i+n+2\).
The exponent from their rank terms is
\(\lfloor a/2\rfloor+\lceil b/2\rceil\), or the same with \(a,b\)
interchanged according to the exterior block. The exponents
equal \(N/2\) since \(a+b=N\) and \(a,b\) have the same parity.
This leaves
\(\nu^{-1}=(\upsilon^{-h}\chi_0\vartheta)^{n+1}\), with
\(\vartheta=1\) on the core triangle.
On a square of the full bordered grid the horizontal and
vertical rank exponents give the same floor and ceiling in
\(V\). In \(Z\) both ranks are even, so their integral halves
also sum to \(N/2\). The remaining product is
\(\mu^0_{j_1}\mu^0_{j_2}\vartheta^\ell\), for horizontal and
vertical step indices \(j_1,j_2\), with \(\ell\) the twisted
step. This calculation holds on each exterior block.
For the full ambient group take exterior characters of product
\(\nu^{-1}\upsilon^{N/2}\). The difference between these
characters and the original block characters contributes the
same multiplier on both paths after the opposite-space reading.
It therefore cancels on return to the original periods.
In the auxiliary comparisons, put
\(\delta_N=\nu^{-1}\upsilon^{N/2}\). For the row pair
\(W_N\to H_a\) write its two characters \((A_R,B_R)\);
for the column pair \(H_b\to W_N\) before taking the adjoint
write \((A_C,B_C)\). Use \(A_R=B_C=\delta_N\) on the
integrated groups and \(A_CB_R=\delta_N\) on the exterior.
On raising the rank by a plane use
\(\delta_{N+2}=\upsilon\delta_N\),
increase \(B_R\) by the factor \(\upsilon\), and keep \(A_C\)
fixed. The zero-term action on descent restores the lower
characters. Taking the opposite small space gives integrated
character \(\upsilon^N B_C^{-1}=\nu\upsilon^{N/2}\)
and exterior character \(\upsilon^b A_C^{-1}\). Thus at the
top, with \(a-b=2j\), the row exterior character divided by
the opposite small one is \(\delta_N\upsilon^{-b}
=\nu^{-1}\upsilon^j\). For the two rank-zero tests use
\((B_R,A_C)=(\delta_N,1)\) when \(b=0\) and
\((1,\delta_N)\) when \(a=0\), giving the same
characters \((\delta_N,\delta_N)\) on the common
equal-rank kernel.
At an old chosen pair the unramified factor of the trial
character still takes its full value on a uniformizer. At a
new pair the unramified factors of the original stage
characters are likewise retained. On \(K_N\) all modulus
factors are one and only the specified residue character
remains.

Retain the extra plane of \(W_1=\mathbb H\perp W_*\) outside
the frame of \(D\). Write \(P=L'\) and let
\(\Psi_0(Y)=1_{\rm int}(Y)\chi^{-1}(\det(PY))\), with its
zero extension off the unit determinant locus. Let \(F\)
be its image under \(\mathcal D\). At a point that fixes the
\(D\)-frame at the selected pairs consider the two partial
means of a coefficient \(F_T\), one for each line of the
extra plane. Partial averaging along the opposite line
gives invariance under \(Y\mapsto Y+u z P Y\) in the linear
functional of the finite datum. Here \(u\) is the extra
line column in the moved frame at \(S\), and \(z\) is any row.
Only determinant-one shears occur. This implies invariance
of the value at \(\Psi_0\) under \(Y\mapsto Y+u rY\) for
\(Pu=ru=0\). In detail, average both tests over
\(Y\mapsto Y+u r H P Y\), where \(H\) varies over integral
matrices. If the affine space \(Y+u Z\) meets the support,
take \(Y_0=Y+u Z_0\) in that support. Then
\(H_0=Y_0(PY)^{-1}\) is integral, \(rY_0=rY\), and the
two averages differ by \(H\mapsto H+H_0\). Otherwise both
are zero. The averaging is a finite sum after passage to a
sufficiently small open subgroup. The actual local unitary
shear has its operation \(1+ur\) at \(S\) and its adjoint
inverse at \(\bar S\). The column test reads \(PY\) at the
indicated half and \(Pu=0\). This proves the partial-mean
shear invariance on both lines of the extra plane.

Write \(\Xi=\jmath T\) for the skew-hermitian Fourier
index and \(\delta=\jmath\alpha\bar\alpha\) for the
inverse-different idele in these coordinates. Consider an
invertible index for which \(\delta^{-1}\Xi\) is unimodular
at \(S\). It is then unimodular at \(\bar S\), since this
normalized matrix and its inverse are hermitian. Use the
standard compact level over the complementary finite places.
Let \(\phi_g^+(a)\) be the mean of \(F_T(n_-(b)n_+(a)g)\)
over \(b\in E^b\backslash\mathbb A_E^b\), where \(n_\pm\)
are the two shifts of the added line by a row linking it
to the \(D\)-frame. Define \(\phi_g^-\) with the lines
exchanged. The commutator of the shifts is in the
\(D\)-center. Moving \(g\) by \(n_+(a)\) translates
\(\phi_g^+\) by \(a\) and multiplies \(\phi_g^-(b)\) by
\(\psi(-\operatorname{tr}(\Xi a^*b))\) in the coefficient
convention. Rational translates of \(a\) read all Fourier
coefficients in \(b\), because \(\Xi\) is invertible.
Thus \(\phi_g^+\) determines \(F_T(n_-(b)n_+(a)g)\)
and then \(\phi_g^-\), and scalar identities pass between
the two means.

To straighten a point with a nonzero mean, use Iwasawa for
the extra line and the \(D\)-frame. Right multiplication
by an integral unitary matrix away from the selected
pairs makes the action on this flag block upper triangular,
with the \(D\)-block denoted by \(\gamma\).
The outer radical factor may be removed from the partial
mean. For the center of the extra line this follows
from invariance under a split local shear and additive
approximation with its rational center. The partial mean
then equals the coefficient at \(0\oplus\Xi\) of the
enlarged frame. The rational radical acts on that frame
by \(\bigl(\begin{smallmatrix}1&\lambda\\0&1\end{smallmatrix}\bigr)\),
which preserves \(0\oplus\Xi\). The radical over the
selected split half acts through the shear invariance
just proved. Apply additive approximation in the radical
column and then in its center to remove the whole factor.
In particular straightening preserves a nonzero mean and
gives a representative which scales the two extra lines
by mutually dual scalars \(t,t'=\bar t^{-1}\) and acts
on the original frame by \(\gamma\) away from the selected
pairs. At those pairs it still fixes the original frame.

Among such representatives with a nonzero partial mean
in either orientation, choose one minimizing the denominator index of
\(\det Q_\gamma^{-1}\) away from the selected pairs, where
\(Q_\gamma=\delta^{-1}\gamma^*\Xi\gamma\).
The matrix \(Q_\gamma\) is integral there. Indeed a
nonzero mean gives a nonzero coefficient at a shifted
point with the same original frame. Invariance under
integral central shifts of the moved frame forces
integrality by the integral additive conductor.
The support of \(\phi^+\) now satisfies
\(\gamma^{-1}a^*t\) integral away from these pairs.
To check this, straighten \(n_+(a)g\).
Its triangular frame before straightening has lower-left
block \(-a^*t\) and lower-right block \(\gamma\).
If \(\gamma'\) is the new lower-right block, comparing
the two frames by the inverse integral right factor gives
\[
 -a^*t=\gamma' H_{21},\qquad \gamma=\gamma' H_{22},
\]
with \(H_{21},H_{22}\) integral at these places.
Nonintegrality of \(\gamma^{-1}a^*t\) would force
\(H_{22}\) to be nonunimodular. Both normalized Gram
matrices are integral and
\(Q_\gamma=H_{22}^*Q_{\gamma'}H_{22}\), so the
denominator index would strictly decrease. By symmetry
the support of \(\phi^-\) satisfies the corresponding
condition on \(\gamma^{-1}b^*t'\).
For each column \(x\) of \(Q_\gamma^{-1}\), the shear
with \(a=(\gamma x\,\overline{t'})^*\) supported at the
complementary finite places has character one on this
support of \(\phi^-\). This follows from the preceding
support condition and the integrality of \(Q_\gamma x\).
It is therefore a period of \(\phi^+\), by the transfer
of scalar identities between the two means.
Comparing one nonzero point of \(\phi^+\) with its
translate by that period shows \(x\) is integral. This
makes \(Q_\gamma\) unimodular.

There is then a point with a nonzero coefficient and
this unimodular Gram block, since \(\phi^+\ne0\) is a
mean of coefficients at unchanged original frames.
Glue the unimodular matrices at all finite places using
the fixed frame at the selected pairs. For the field
in the grid, which contains a primitive \(5040\)-th root
of unity, the finite unimodular hermitian form admits an
integral orthonormal frame. At split pairs this uses
the perfect integral pairing, and at the unramified fixed
places it uses the residue norm and integral hermitian
diagonalization. This is applied after removing the
inverse-different factor \(\delta\).
If the real class is wrong, the off-class vanishing
already applies. In the matching real class Landherr's
principle identifies the form rationally with \(J_U\),
since the finite orthonormal test gives the required
local isometries. Move by this rational isometry on
the original frame.\footnote{\weightedpin{grh-compact/src/New/Exchange/StandardCoefficients.lean}
{18-L78}{Primitive coefficients and global transport},
\weightedpin{grh-compact/src/New/Exchange/BlockPartners.lean}
{17-L29}{Common coset representatives},
\weightedpin{grh/src/New/Determinant/Exchange/MatchingPartners.lean}
{424-L492}{Linear factor through the common local sections},
\weightedpin{grh/src/New/Determinant/Exchange/StableSpacesStandardData.lean}
{686-L750}{Partial-mean shears on the named datum},
\weightedpin{grh/src/Lib/SiegelWeil/SingularForms/CentreCoefficients.lean}
{482-L581}{Two means and rational Fourier inversion},
\weightedpin{grh/src/Lib/SiegelWeil/SingularForms/CentreCoefficients.lean}
{705-L837}{Radical invariance from the split shears}
and \weightedpin{grh/src/Lib/SiegelWeil/SingularForms/UnimodularCoefficients.lean}
{675-L856}{The denominator descent}.} The new coefficient is a self-dual
test and was proved zero. This contradiction proves
vanishing at all points trivial at the selected pairs for
indices primitive there.

Now use the integral opposite frame at the selected places.
On the translates \(\Psi_z(Y)=\Psi_0(Y+K z L'Y)\), a
central Fourier functional at a point trivial at those
places is multiplied by the character of \(z\) with the
normalized Fourier index. On integral \(z\) the standard
datum is unchanged. Thus nonzero functionals can only
have integral normalized index, by the conductor.
Put \(\varpi=(\varpi_v)_{v\in S}\) and combine the tests
\(\Psi_{\varpi^{-1}\widetilde r}\) with weights
\(\chi^{-1}(\det r)\), over square matrices \(r\) in the
product of residue fields. Use zero extension of the
determinant character and take the lift of \(1\) to be \(1\).
The Fourier multiplier of this sum is a matrix Gauss
sum. At a singular frequency it is zero, since a
stabilizer on a complementary line realizes every
determinant and the character is nontrivial. At the
remaining frequencies the primitive coefficient test
gives zero. Fourier uniqueness on the center and weak
approximation away from the selected places put the
combination in the kernel.

The scalar translate by \(-\varpi^{-1}\) of this
combination is \(r(Y)\) times the standard datum for
the opposite row block, with \(r(Y)\) a nonnegative
integer. Indeed, if \(X\) is the sheared variable of
a nonzero summand, then \(Y=K\widetilde r (L'X)-\varpi X\).
The opposite row block of \(Y\) reduces to \(r\) times
the row block \(L'X\). Each nonzero summand is therefore
the value of that opposite standard datum. The term
\(r=1\) is nonzero at \(K-\varpi L\), whose opposite
standard row block is \(1\).

Such a test generates all the required locally constant
column functions in the stable kernel. Let \(Y_0\) be
its witnessing column and \(P_0Y_0=1\) its standard row
block. A column partner is covariant by
\(\chi^{-1}(\det h)\) for right \(h\in{\rm GL}_b(\mathcal O_S)\).
Choose a depth \(d\) at which the given column function
\(\Psi\) is translation invariant on integral matrices.
The map \(D_d=Y_0P_0+\varpi^{-d}(1-Y_0P_0)\) fixes
\(Y_0\) and dilates the kernel of \(P_0\).
Average the test \(p(D_d kY)\) over integral \(k\)
with the weight \(\Psi(k^{-1}Y_0)\).
On its support
\[
 Y\equiv k^{-1}Y_0(P_0D_d kY)\pmod{\varpi^d},
\]
and covariance cancels the determinant character.
Thus the average is \(\Psi(Y)\) times the integral of
\(|p(D_d kY)|\). The latter is constant on primitive
columns by transitivity, and is positive at \(Y_0\).
Compactness and local constancy make this average a finite
linear combination of translates. This proves the extension
to standard rows and columns. The last averaging argument,
starting from an ordinary standard column test in the
kernel after descent, proves the identity on all partners.
The depth \(d\) concerns the matrix functions and can exceed
the conductor of the determinant character.
\footnote{\weightedpin{grh-compact/src/New/Exchange/StableSpaces.lean}
{29-L88}{Generation after the Gauss sum},
\weightedpin{grh/src/New/Determinant/Exchange/StableSpacesStandardData.lean}
{302-L395}{The opposite standard datum}
and \weightedpin{grh/src/New/Determinant/Exchange/TheoremForPartners.lean}
{241-L272}{Extension after descent to the actual partners}.}

\paragraph{The normalized Tate factor for one plane.}
\teamanchor{proof.weighted-one-plane-tate}%
Write \(W\) for the lower ambient space of even rank \(N\) and
\(V\) for the lower intermediate space of rank \(a\). The unchanged
intermediate space has rank \(b=N-a\). Let \(A,X\) be the adapted
plane splittings of \(W^+=\mathbb H\perp W\) and
\(V^+=\mathbb H\perp V\). Use the prescribed row orders
\(\sigma\colon b\oplus a\longrightarrow N\) and
\(\sigma^+\colon b\oplus(a+2)\longrightarrow N+2\), where the integers
in these orders denote their coordinate sets. The fixed \(b\) rows
at the lower stage are an integral matrix \(P\) at the selected
split places, and the upper fixed rows are \(P I_A'\), where \(I_A'\)
projects onto the \(W\)-coordinates of the adapted frame. The
standard row datum at either stage is the inverse determinant
character of the stacked, ordered rows on the integral invertible
locus, extended by zero.

Let \(\nu\) be the square character of positive infinite weight
\(\mu\), put \(\widetilde\nu(\lambda)=\nu(\lambda/\bar\lambda)\),
and let \(S\) be the one-sided, nonempty split support of the tame
part, disjoint from \(\bar S\). Put \(Q_S=\prod_{w\in S}q_w\).
For \(w\in S\) let \(\chi_w\ne1\) be the residue-field character
read on the \(w\)-component, and identify \(E_{\bar w}\) with \(E_w\)
via conjugation. The finite character on the two unit factors
is \((\chi_w,\chi_w^{-1})\). Thus \(\widetilde\nu\) has conductor
exponent one on both factors and is unramified off \(S\cup\bar S\).
Set \(e=a-N/2\), so the upper difference is
\((a+2)-(N+2)/2=e+1\). Put \(x=e+1\) for \(e\ge0\) and
\(x=-e\) for \(e<0\).
For the two line models use \(\beta=+\) in the first case and
\(\beta=-\) in the second. With the cone vector formed using this
\((\beta,x)\), the normalized one-plane multiplier \(\gamma\) has
\[
 |\gamma|^2=C_{\beta,x}\,Q_S^{2x-1}|L^S(x,\nu)|^2
   = C_{\beta,x}\frac{\Pi_{\nu,S}(e+1)}{\Pi_{\nu,S}(e)},
 \qquad C_{\beta,x}>0 .
\]
The constant can depend on the fixed spaces, splittings, infinite
weight and measures. It is independent of the character of weight
\(\mu\), the ramified set, the integral \(P\), the real vector and
the adelic point. This scalar evaluation uses the mixed-model line
and Levi formulas and the primitive-data lower theta kernel bound
stated below.

\emph{The split line calculation.}
In an adapted frame, a line term fixes a scalar \(s\) on one pivot
and partially Fourier transforms the other removed coordinates. At
\(w\in S\), let \(M\) be the \(N\)-square matrix of the rows of \(P\)
and of the variable lower matrix \(Y\), stacked in order \(\sigma\).
The ambient and intermediate isotropic line coordinates put the
upper stacked matrix, up to fixed integral invertible frame changes
on the left and right, in the form
\[
 T_+=\begin{pmatrix}
       s & z_W & z_0\\
       0 & M & (0,z_V)^{\mathsf T}\\
       0 & 0 & t
      \end{pmatrix}.
\]
Here \(z_W\) is a row of length \(N\), \(z_0\) is a scalar, and
the last column of the middle block has \(b\) zero entries and
\(a\) entries from \(z_V\), in order \(\sigma\). The other model
uses \(T_-\) obtained by interchanging the two diagonal scalars
\(s,t\). These are the two mixed coordinates on the same removed
plane. Indeed the rows of \(P I_A'\) vanish on both plane
coordinates, the lower variable block maps \(W\) to \(V\), and
the remaining integrated maps range from \(W\) to one removed line,
from the other removed line to \(V\), and between the removed
lines. The two models exchange the fixed and integrated maps
at the two diagonal pivots. The full frames and the permutations
\(\sigma,\sigma^+\) change the determinant by an integral unit
independent of the matrix arguments. On either model
\(\det T_\beta=s\,t\,\det M\). Integrality is entrywise and unit
determinant requires each of \(s,t,\det M\) to be a unit.

Write \(D_w(t)=1_{\mathcal O_w^\times}(t)\chi_w(t)^{-1}\) and use
the additive Haar measure of volume one on \(\mathcal O_w\) for
the additive character \(\psi_w\) of conductor \(\mathcal O_w\).
In the line term at \(\alpha=\alpha_0\lambda\), where \(\alpha_0\)
is the fixed Weil scale idele, \(s=\lambda_w\) and the Fourier
frequency of \(t\) is \(\bar\lambda_w\). All other Fourier
frequencies are zero. Those variables give \(\mathcal O_w\)-volumes
one; the determinant phase from the fixed frames has modulus one.
The remaining transform is
\[
 \int_{\mathcal O_w}D_w(t)\psi_w(bt)\,dt
   =g_w\,1_{\mathcal O_w^\times}(\varpi_w b)\chi_w(\varpi_w b),
 \qquad
 g_w=q_w^{-1}\!\sum_{u\in k_w^\times}
           \chi_w(u)^{-1}\psi_w(\varpi_w^{-1}\tilde u).
\]
This is the conductor-one specialization of the local ramified
calculation in \href{https://sites.math.rutgers.edu/~alexk/2023S572/Tate1950.pdf\#page=26}
{Tate, \S2.5}, in the stipulated volume-one gauge.
For \(b\in\mathcal O_w\) the character sum of \(D_w\) vanishes.
If \(v_w(b)<-1\), orthogonality on each class modulo \(\varpi_w\)
gives zero. If \(v_w(b)=-1\), summing the residue classes and
changing variables by \(\varpi_w b\) gives the formula. For the
unnormalized finite Fourier transform of \(\chi_w^{-1}\) extended
by zero at zero, the value at frequency zero is zero and all
nonzero frequencies have the same modulus. Plancherel gives
\((q_w-1)q_w^2|g_w|^2=q_w(q_w-1)\), hence
\(|g_w|^2=q_w^{-1}\).

For \(p=(p_w)_{w\in S}\), define \(F_p(\lambda_f)\) on the
finite adeles by requiring integrality off \(S\cup\bar S\) and
requiring \(\lambda_w,p_w\bar\lambda_w\) to be units for every
\(w\in S\). On that support put
\[
 F_p(\lambda_f)=\prod_{w\in S}\chi_w(\lambda_w)^{-1}
                                  \chi_w(p_w\bar\lambda_w);
\]
give it value zero elsewhere. The fixed vacuum in the scaled
coordinates gives the integrality condition off \(S\cup\bar S\).
The split calculation expresses the line kernel at
\(\alpha_0\lambda\) as \(l_\beta G_\nu F_{\varpi}(\lambda_f)\)
times the lower theta kernel with the same \(P\) and the real
line slice at \(\lambda_\infty\). Here \(l_\beta\ne0\) is the
fixed vacuum normalization of the adapted splitting and \(G_\nu\)
is the product of the \(g_w\) with the unit determinant phase.
Thus \(|G_\nu|^2=Q_S^{-1}\). For the real slice the upper vector
is acted on by the real component of the compact argument \(k\)
of \(U(V^+)\) in the pair representation with character
\(\nu^{-1}\upsilon^{(N+2)/2}\) on both sides. The finite compact
action fixes the datum because its determinant effect on the
standard rows cancels the splitting-character effect on the
units. At each removed plane the rank twist is
\(\nu^{-1}\upsilon^{N/2}\) on the lower pair.
\footnote{\weightedpin{grh/src/Lib/SiegelWeil/Tower/LineTerms.lean}
{185-L315}{The two adapted mixed models},
\weightedpin{grh/src/Lib/SiegelWeil/Tower/LineGaussSums.lean}
{347-L500}{Ordered block frame and fixed vacuum},
\weightedpin{grh/src/Lib/SiegelWeil/Tower/LargeSideConstantTerm.lean}
{91-L117}{Finite compact invariance}.}

\emph{Unfolding the line and its finite integral.}
Write \([V]\) for the quotient by rational points of \(U(V)\)
and \(K^+\) for the fixed adelic maximal compact subgroup of
\(U(V^+)\), with its full compact Haar measure.
For a lower \(g\in U(W)(\mathbb A)\), let
\(B_\beta(\alpha;h,k)\) denote the line kernel just described,
with \(h\) in \([V]\) and \(k\in K^+\).
The line constant term at \((1,g)\) is the fixed multiple
\(c_{\rm unf}\ne0\) of the integral at \(z=e+1\) of
\[
 \widetilde\nu(\lambda)|\lambda|^z
 \sum_{\xi\in E^\times} B_+(\xi\lambda;h,k)
\]
on \(E^\times\backslash\mathbb A_E^\times\times[V]\times K^+\).
Here the line constant-term formula first selects rational
isotropic vectors. Its zero orbit on the large side vanishes
because the fixed rows kill the added first isotropic column,
and a variable block in the zero slice kills that same column.
The full stacked matrix then has a kernel. The other vectors
are grouped by their line. Witt's theorem gives one orbit of
isotropic lines, and unfolding along its stabilizer gives the
Levi, lower quotient and compact integrations. In the coordinate
\(\lambda\) the actual line Levi acts on its generator by
\(\lambda^{-1}\). Conjugation on its radical, whose coordinates
are a column in \(E^a\) and a one-dimensional center
over the real subfield, has modulus \(|\lambda|^{-(a+1)}\).
Quotient integration therefore contributes
\(|\lambda|^{a+1}\). The direct line Weil dilation contributes
\(\widetilde\nu(\lambda)|\lambda|^{-N/2}\), giving \(z=e+1\).
For the opposite model its dilation contributes
\(\widetilde\nu(\lambda)|\lambda|^{-(N+2)/2}\), with line
coordinate \(\bar\lambda^{-1}\). The measures here are the fixed
quotient measures of the two adapted splittings. For the equal-rank
period calculation the same relation holds at \((u,g)\), for
\(a=N\), with \(\lambda u\) in place of \(\lambda\). The two
line-slice determinants in that case give a numerator
\(|u|^{a+1}\) in the Weil dilation, as required to replace
\(|\lambda|^{a+1}\) by \(|\lambda u|^{a+1}\).
\footnote{\weightedpin{grh/src/Lib/SiegelWeil/Tower/ConstantTerm.lean}
{643-L717}{Line orbit and quotient integration},
\weightedpin{grh/src/Lib/SiegelWeil/Tower/LineTerms.lean}
{112-L180}{The two Levi dilations},
\weightedpin{grh/src/Lib/SiegelWeil/Tower/LargeSideConstantTerm.lean}
{122-L159}{Primitive support excludes the zero orbit}.}

The two line kernels have zero term at the origin by the split
support just computed. The mixed slices first take the common
zero projection away from the remaining hyperbolic plane. The
two scalar line variables are dual in that plane, and the lower
theta functional commutes with the line Fourier transform.
Adelic Poisson on this line, using self-dual measures, gives
\[
 \sum_{\xi\in E^\times}B_-(\xi\bar\lambda^{-1};h,k)
   =|\lambda|\sum_{\xi\in E^\times}B_+(\xi\lambda;h,k).
\]
The bilinear Fourier pairing on that line uses conjugation in the
dual variable. Under \(j_-(\lambda)=\bar\lambda^{-1}\) the norm
inverts and \(\widetilde\nu\) is unchanged, since
\(\bar\lambda^{-1}/\lambda^{-1}=\lambda/\bar\lambda\).
Consequently the Mellin integral at \(z=1-x\) uses the
\(\beta=-\) line kernel at exponent \(x\) in the
\(j_-\)-variable. The integral at \(z=x\) uses \(\beta=+\)
and \(j_+(\lambda)=\lambda\).

Choose the finite idele \(c\) with components \(1\) off
\(\bar S\) and with conjugate component \(\varpi_w\) at each
\(w\in S\). Then \(F_\varpi(\lambda_f)=F_1(\lambda_f c)\).
In the unfolded idele integral write the line argument as
\(\lambda K_\nu\), with \(K_\nu=\alpha_0 c^{-1}\).
One has
\[
 |K_\nu|=|\alpha_0|Q_S,\qquad
 |\widetilde\nu(K_\nu)|=1 .
\]
For \(x>1\), unfold the sum over \(E^\times\) using the
multiplicative quotient measure and split the idele measure in the
variable \(j_\beta\) as the chosen real multiplicative Haar
times a finite multiplicative Haar \(\rho_\beta\). This splitting
is independent of \(\nu\) and translation by \(K_\nu\) preserves
the measure. The finite integral is
\[
 \int_{\mathbb A_{E,f}^\times}F_1(\lambda_f)
       \widetilde\nu_f(\lambda_f)|\lambda_f|^x\,d\rho_\beta
   =\rho_\beta(U_f)L^S(x,\nu),
 \qquad U_f=\prod_{w<\infty}\mathcal O_w^\times .
\]
Indeed both components above a selected split prime are units,
and their character values cancel \(F_1\) on the unit ideles.
At each remaining prime the valuation runs over \(j\ge0\)
and its summand is \(\widetilde\nu_w(\varpi_w)^j q_w^{-jx}\).
The normalized coset volume is the same over every valuation
tuple, namely \(\rho_\beta(U_f)>0\). Absolute convergence
follows by majorization by the Dedekind Euler product for \(x>1\).
Here \(L^S\) omits \(S\cup\bar S\).

The real integral is the lower theta integral on the cone vector.
It averages the real line slice at \(\lambda_\infty\) over
\(K^+\times(E\otimes_{\mathbb Q}\mathbb R)^\times\) against
\[
 |\lambda_\infty|_E^x
 \prod_{\tau}\left(\frac{\lambda_\tau}{|\lambda_\tau|}\right)^{2\mu_\tau}.
\]
This fixes the real operator before the finite part of \(\nu\)
is chosen. The mixed-model slice of a Schwartz vector, after
acting by \(k\) in the compact group, has every lower Schwartz
seminorm rapidly decreasing in the line coordinate, uniformly
in \(k\). This follows by partial Fourier transformation and
restriction in the removed variables. Use the primitive-support
theta kernel bound with zero exponent on the height of \(h\).
It gives a bound by a fixed Fock moment, uniformly in \(h\) for
a fixed \(g,P,\nu\). On the support of the lower datum \(Y\)
has full row rank at the selected places, since the matrix
formed by stacking \(P\) and \(Y\) is invertible. Thus the bound
applies to this datum.
In each complex radial variable the measure is proportional to
\(dr/r\), and the weight at zero is \(r^{2x}\). The cone operator
and the integral after applying the lower theta functional
are therefore absolutely convergent and continuous locally
for \(x>0\). These bounds justify Fubini for \(x>1\).

At \(x=1\) use the common Mellin integral on the idele quotient.
For \(R>1\) its absolute convergence at \(R\) and \(1-R\)
follows from the two unfolded formulas. For \(z\) in between,
the weight \(|\lambda|^{\Re z}\) is bounded by
\(|\lambda|^R+|\lambda|^{1-R}\), giving continuity at the
interior points, including \(0,1\) for \(R>1\).
The partial Hecke function is continuous at \(1\) by
Tate's continuation and the factors at the omitted places.
In fact \(\widetilde\nu\) and \(\widetilde\nu^2\) both have
nonzero angular type and the corresponding partial Hecke
functions are holomorphic there. If \(L^S(1,\nu)=0\), then
for \(\sigma>1\) tending to one the expression
\[
 \bigl(\zeta_E^{S\cup\bar S}(\sigma)\bigr)^3
 |L^S(\sigma,\nu)|^4\,|L^S(\sigma,\nu^2)|
\]
tends to zero by the simple pole of \(\zeta_E\). It is at
least one by taking the logarithm of the Euler products and
using \(3+4\cos t+\cos(2t)\ge0\) term by term. It follows that
\(L^S(1,\nu)\ne0\). Thus the factorization at \(x=1\) uses the
right limit of the partial Euler product. On either branch
the finite and archimedean factors combine to give the
complex scalar
\[
 \gamma = c_{\rm unf}l_\beta G_\nu\,
         \widetilde\nu(K_\nu)|K_\nu|^x
         \rho_\beta(U_f)L^S(x,\nu).
\]
The factor \(Q_S^{2x-1}\) in its squared modulus is the product
of \(Q_S^{-1}\) from the additive Fourier transform and
\(Q_S^{2x}\) from the idele translate. For \(e\ge0\) the
ratio of the two \(\Pi\)'s is \(Q_S^{2e+1}|L^S(e+1,\nu)|^2\).
For \(e<0\), set \(x=-e\) and use
\(\Pi_{\nu,S}(-j)=\Pi_{\nu,S}(j)^{-1}\). The ratio becomes
\(\Pi_{\nu,S}(x)/\Pi_{\nu,S}(x-1)\) and has the same formula.
\footnote{\weightedpin{grh/src/Lib/SiegelWeil/Tower/LineTerms.lean}
{231-L274}{Poisson for the common hyperbolic plane},
\weightedpin{grh/src/Lib/Weil/DualPair/KernelDecay.lean}
{364-L378}{Primitive-support theta kernel bound},
\weightedpin{grh/src/Lib/SiegelWeil/Tower/LargeSideConstantTerm.lean}
{173-L209}{Primitive theta kernel bound and cone integration}.
For holomorphy of the two Hecke functions and the simple pole
of the Dedekind zeta function, see J.~T.~Tate,
\emph{Fourier analysis in number fields and Hecke's
zeta-functions}, \S\S4.4--4.5.}

\emph{Descent of the exchange scalar.}
Use the corresponding projected standard column datum with the
same \(P\) on the other side. Its line constant term is the
lower theta integral on the real zero slice times the fixed
nonzero vacuum factor \(\kappa_0\), by the primitive-support
small-side constant-term identity. If the upper
large integral equals \(c_+\) times the upper small integral
on the standard data, then
\(\gamma I_{\rm large}=c_+\kappa_0 I_{\rm small}\)
at the lower stage. Thus \(c_-=c_+\kappa_0/\gamma\). For an
upper scalar squared modulus proportional to
\(\Pi_{\nu,S}(e+1)\mathcal R\), the lower squared modulus
is proportional to \(\Pi_{\nu,S}(e)\mathcal R\) with the
same \(\mathcal R=\mathcal R_\nu(w)\) from the fixed determinant
test above. For the equal-rank period use \(a=N\)
and the equal-rank constant function identity. Its numerator
\(|I^+(1)|^2\) has the same factor, up to \(C_{\beta,x}\) and
the fixed radical volume, as the denominator
\(\Pi_{\nu,S}(N/2+1)\) when passing down to rank \(N\).

\paragraph{Nonvanishing and the complementary-period norm.}
\teamanchor{proof.weighted-complementary-period-norm}%
For an even-rank space \(Y\) with an adapted hyperbolic
line write \(I_Y(v)\) for its equal-rank theta integral on
the signed determinant datum. The split equivariance with
zero column rank makes \(I_Y(v)\) right invariant under
the full unitary group at a selected split place. The
derived group is the simply connected special unitary
group and is noncompact at this place. Strong approximation
away from the split place gives density of rational special
unitary elements times the special unitary group at that
place. In applying this density to \(I_Y(v)(sg)\), the
local factor can be commuted past \(g\), since its support
stays at that place. Thus \(I_Y(v)\) factors through the
determinant. In particular it equals its line radical
average. For equal lower ranks the one-plane computation
has an external scalar \(u\) on the line. Its contribution
combines with the quotient modulus to give
\(|\lambda u|^{a+1}\), so the substitution \(\lambda u\)
removes \(u\) from the constant term. The determinant of the
ambient line Levi ranges over the norm-one ideles by
Hilbert 90, choosing unit lifts at almost all places.
The integral \(I_Y(v)\) is therefore constant on the
adelic group. The same argument applies to the two
sign-compatible spaces reached by the local unitary
transport after changing the hermitian form: the standard
finite datum and the measures are transported, and the
remaining Weil factor has modulus one.
\footnote{\weightedpin{grh-compact/src/New/Exchange/LowRankIdentity.lean}
{78-L86}{Invariance of the zero-column datum},
\weightedpin{grh/src/New/Determinant/Exchange/LargeTriple.lean}
{387-L425}{Constancy on an adapted equal-rank pair}.
For strong approximation use M.~Kneser,
\emph{Starke Approximation in algebraischen Gruppen I},
J.~reine angew.~Math. 218 (1965), Satz~2, p.~191.}

We first record the two real preimages used with the
one-plane calculation. For the positive branch, the
cone slice is read along the null generator \(b_0\) of
the added plane on the integrated side. For a compact
argument \(k\) and real line variable \(\lambda\), the
moved column is \(x=k^{-1}b_0\lambda\).
Near \(b_0\), the orbit of nonzero isotropic vectors has
a smooth local unitary section \(h(x)b_0=x\). Indeed
the derivative of the action, \(X\mapsto Xb_0\) for
skew-adjoint \(X\), surjects onto the real tangent space
\(\operatorname{Re}\langle b_0,y\rangle=0\). In a
hyperbolic pair this follows by freely choosing the
column at \(b_0\), filling its adjoint row, and imposing
the stated real condition at their intersection.
Choose a smooth nonnegative bump \(\beta\) near
\(b_0\) positive there, using an ambient smooth extension
of the section from the null manifold on this chart.
Extend the lower vector by a
Schwartz scalar of value one in the second zero
projection, and prescribe the first slices to be
\(\beta(x)\) times the translates of this extension by
\(h(x)\), with target weight \(n/2-\mu_\tau\).
The real mixed model realizes these slices
by partial Fourier inversion. On their support \(kh(x)\)
multiplies \(b_0\) by \(\lambda\). The line Levi law
thus gives a lower translated vector times
\(\beta(x)|\lambda|^{n/2}_E\) and angular phase
\((\lambda/|\lambda|)^{-2\mu}\). The angular weight of
the cone cancels this phase, and translating the lower
integrated variable removes the lower unitary element.
The remaining real multiplier is the integral of
\(\beta(x)|\lambda|_E^{z+n/2}\). It is finite since the
bump confines every \(|\lambda_\tau|\) to a compact
subinterval of \((0,\infty)\), and positive by its
value on a neighborhood of \((1,1)\).
Dividing the upper input by this factor gives any
prescribed lower equal-rank integral on its standard
datum. This argument uses only the positive branch.

For the two branches in the nonvanishing argument it
suffices to have preimages on determinant vectors.
Write \(n=2(k_0+q)\) at a lower step and
\(\mu_\tau=k_0+m_\tau\) with \(m_\tau\ge0\).
Use the Gaussian times the conjugate product of
positive-coordinate principal minors of exponents
\(m_\tau\) on the positive block and \(m_\tau+2k_0\)
on the negative block. In the notation of the fixed
Fock model these minors are taken after multiplying the
positive-coordinate matrix by one fixed matrix \(A_0\).
Extend \(A_0\) by zero across the complements and by
the two marked line pivots
\(-\alpha_\infty^{-1}(K_WK_V^*+L_WL_V^*)\) in the
adapted real frames. For the real compact action each
block transforms by the determinant of its compact
matrix. The vacuum and splitting character cancel
both powers because \(n'/2=k_0+q(W^+)\) at the
upper step. In the mixed slice the equal-sign projection
of the added plane has one added \(\lambda_\tau\) pivot
in each compact sign block, with zero off-diagonal
entries relative to the old frame. This is obtained
by projecting the two dual isotropic generators onto
their equal-sign sums and taking the chosen scalar in
the retained line. The two mixed models exchange those
generators. Thus each sign minor acquires a factor
\(\lambda_\tau\) in either model, while the Gaussian
acquires the fixed line Gaussian. The angular factor
of the cone cancels the conjugate determinant power
\(\overline{\lambda_\tau}^{\,2(k_0+m_\tau)}\).
With the normalized real scale, the resulting
integrand is the lower determinant vector times
\[
 \exp\!\left(-\sum_\tau|\lambda_\tau|^2\right)
       |\lambda|_E^z\,\prod_\tau|\lambda_\tau|^{\,2(k_0+m_\tau)} .
\]
For \(z\ge1\) its integral is positive and finite.
Rescaling gives the desired preimage on each
successive plane. The rank-zero vacuum is in this
determinant family.
\footnote{\weightedpin{grh/src/New/Determinant/Exchange/ConeImage.lean}
{225-L284}{The local isotropic section},
\weightedpin{grh/src/New/Determinant/Exchange/ConeImage.lean}
{324-L376}{Cancellation inside the cone}
and \weightedpin{grh-compact/src/New/Exchange/DeterminantVectors.lean}
{25-L139}{Both branches on determinant vectors}.}

Take \(J\) positive definite of rank \(2k_0\) and
\(Y=\mathbb H\perp J\). Fix a tame character of weight
\(\mu_\tau\ge k_0\). If \(I_Y(w)(1)\) vanished for every
Schwartz \(w\), every top coefficient of
\((\mathbb H^{2k_0}\perp Y,J^-\perp Y)\) in the self-dual
test would vanish, even for arbitrary large real input.
Indeed the real tensor product first gives this for pure
tensors in the slice, and the coherent-state integral
with all moments gives it for arbitrary slices.
For invertible indices outside the negative class use
the negative-line rotation above. The coefficient
propagation gives zero on all standard rows. Descend
through the \(2k_0+1\) planes, using the determinant
preimages and the nonzero Tate factor at each stage.
On \((J,0)\) this would kill the vacuum on the standard
finite frame, whose theta integral is the positive
volume of the rank-zero quotient. This proves a
nonzero period for the chosen character.

Fix such a \(w\). Apply the descended identity first to
\((J,J,0)\) and then to \((J,0,J)\) with the same \(w\)
at the top. In the first identity the small vector is
the rank-zero vacuum. In the second use the large
vector produced in the first as the prescribed
\((J,J)\)-vector on the small side. Put \(\mathcal R_\nu
=|I_Y(w)(1)|^2/\Pi_{\nu,S}(k_0+1)\). With fixed
normalization constants \(B,C\), the two identities
have factors \(c_1,c_2\) satisfying
\[
 B|c_1|^2=\Pi_{\nu,S}(k_0)\mathcal R_\nu,\qquad
 C|c_2|^2=\Pi_{\nu,S}(-k_0)\mathcal R_\nu .
\]
Their common middle kernel has the signed determinant
datum and the same real vector on the two copies of
\(J\). Fubini, with the same Haar on both copies,
gives \(c_1c_2=\xi\), where \(\xi\) is the value of the
fixed vector on the rank-zero hom space. The positive
rank-zero quotient volume cancels. The two \(\Pi\)'s
are reciprocal, so \(\mathcal R_\nu^2=BC|\xi|^2\).
Thus \(\mathcal R_\nu\) is one positive constant
throughout the family. To pass from the positive
\(J\) to the actual \(W_*\), use the positive preimage
at each added plane. The one-plane factor and the
constant function identity multiply numerator and
denominator of \(\mathcal R_\nu\) by the same
\(\Pi\)-ratio up to a fixed positive factor. The
number of added planes is \(q_*\).

Now use the fixed vector \(w\) on
\((\mathbb H\perp W_*,\mathbb H\perp W_*)\) in the top
coefficient construction. An arbitrary Schwartz vector
on the small pair extends to the required opposite
zero slice, since in the real mixed model one can
multiply it by a Schwartz scalar of value one there.
The \(b+1\) successive removals leave exactly that
small vector. At each step the standard column is
projected using the original adapted dual frame,
and the scalar update is \(c_-\!=c_+\kappa_0/\gamma\).
The powers of \(\Pi_{\nu,S}\) telescope to the final
excess \(e=a-N/2\), including steps with \(e<0\)
in the two rank-zero tests. The top calculation uses
the opposite \(V_{\rm small}^-\); return to
\(V_{\rm small}\) by adjoint coordinates, with the finite
component at its indicated conjugate half. The
conversion of the unitary quotient is measure
preserving and the remaining determinant factor is
a unitary scalar. The rank twist
\(\nu^{-1}\upsilon^{N/2}\) on the exterior group
converts the two original integrals into the
complementary integrals used in the square.
It contributes no factor to the squared scalar.
The general output is the partner identity with
\(|c|^2=C_*\,\Pi_{\nu,S}(e)\), \(C_*>0\) independent
of the twist and of the partner tests. It holds for
the chosen real vector on the small pair, all
Schwartz---Bruhat determinant partners and the
fixed measures. A finite quotient can be of greater
depth than the tame determinant character.
\footnote{\weightedpin{grh-compact/src/New/Exchange/EqualRankValue.lean}
{97-L189}{Reciprocal comparison},
\weightedpin{grh/src/New/Determinant/Exchange/TheoremForPartners.lean}
{204-L232}{Returning from the opposite space},
\weightedpin{grh/src/New/Determinant/Exchange/TheoremForPartners.lean}
{501-L562}{Nonvanishing at the rank-zero test},
\weightedpin{grh/src/New/Determinant/Exchange/TheoremForPartners.lean}
{686-L784}{The same kernel in the reciprocal comparisons}
and \weightedpin{grh/src/New/Determinant/Exchange/TheoremForPartners.lean}
{805-L842}{The normalization at the actual signs}.}

In triangle \(z\) the two exterior ranks are both
\(r+z+(n+1)(t-z)\). The intermediate ranks in the forward
square are \(a=r+z+(n+2)(t-z)\) and \(b=r+z+n(t-z)\), so
the lower endpoint of the plane descent is \(e=t-z\). The square
character has \(\nu^{-1}=(\upsilon^{-h}\chi_0\vartheta)^{n+1}\)
for \(V\) and the same formula with \(\vartheta=1\) for \(Z\).
The finite matrix modulus is \(\mathfrak M Q_\vartheta\) for
\(V\) and \(\mathfrak M^{b_1}\) for \(Z\). In the latter case
the tame character factors through the residue fields. The
scalar uses conductor exponent one at the selected places in
both triangles. The power \(1\le n+1\le m\) preserves
nontriviality there. The two intermediate signs add to the
exterior count, and \((n+1)h\ge N/2-q_*\) at the actual ranks.
The equal-intermediate-rank exchange of the
reverse square in \(V\) uses \(e=0\). Its exchange scalar
product is \(\Pi_{\nu,S}(0)=1\). The reversal factors enter
through the separate Weil kernel comparison.
\footnote{\weightedpin{grh/src/New/Determinant/Exchange/RemovingPlanes.lean}
{413-L464}{Small-side constant term},
\weightedpin{grh/src/New/Determinant/Chain/ChainNorms.lean}
{157-L224}{The forward square and tame family},
\weightedpin{grh/src/New/Determinant/Chain/ChainNorms.lean}
{261-L294}{The reversed square and its equal intermediate ranks}.}

\paragraph{The real comparison on the realized form spaces.}
\teamanchor{proof.weighted-grid-real-transport}%
The complementary integral is an identity at every exterior point.
For a fixed small-side real vector \(v\), it pairs \(v\) with the
constructed large-side vector \(Tv\). The extension to the top
zero slice, the fixed real bumps and equal-rank vector, and the
real cone operators give a continuous linear choice of \(T\) on
Schwartz space. This choice depends on the signatures, frames
and real weights. It is independent of the finite data and
finite part of \(\nu\). Thus a pointwise use of the identity may
transport the real vector by \(T\), also in the presence of
other tensor factors.

\teamanchor{proof.weighted-grid-real-contraction}%
For the operator comparison on realized forms, first pair with
the two prescribed exterior types. In an oriented path through
an intermediate group \(H\), regard each edge as a lift from
its exterior group to \(H\), using the transpose with the
contragredient type where required by the bilinear pairing.
For each edge let \(\pi_i\) be the exterior real type and
\(\pi_H^{(i)}\) the resulting intermediate type. The signed
center and weight calculations give the same intermediate
representation on one edge and its contragredient on the
other, with the compact positive order reversed in the
contragredient factor so that the two chosen lines are paired.

Here are the real hypotheses for using the oscillator quotients
on these edges. After removing the determinant centers, every
matching adds or removes one positive-labeled central string
of length \(d=|N'-N|\). All remaining centers are beyond
\((d+1)/2\) in absolute value. The signs of the surviving
coordinates are precisely the match in the signed parameter
calculation. Paul's Theorem~2.7(b) gives the larger discrete
series, and Proposition~3.4(ii) gives the reverse first
occurrence since that central string is maximal. Negating the
parameters or taking the opposite convention uses the
contragredient identification of Proposition~2.9. In the
forward direction the smaller parameter has no central entry.
For odd \(d\) its \(k_\pi=-1\) and its first occurrence in the
indicated tower follows from Ichino's Theorem~6.1(i), with
\(l=0\), \(t=(d-1)/2\) and both near sets empty. For even \(d\)
the smaller centered entries are half-integers, \(k_\pi=0\),
and Theorem~6.1(ii) applies with \(l=0\), \(t=d/2\) and
the same empty near sets. In either case removing a
hyperbolic plane would make \(l<0\). When the smaller
Witt-tower member does not exist, no such check is needed.
These are the odd \(V\) and the even \(Z\) excesses, including
\(t-1\) and \(t+1\) in the bordered grid. In the transpose
the determinant centers are preserved upon removing a
hyperbolic plane by the normalized zero-term formula.
Wu's tower constant-term formula in both orientations
therefore gives cuspidal intermediate lifts for arbitrary
real sections and cuspidal exterior inputs, as in the column
calculation. In the realized comparisons the exterior forms
are in the cuspidal spaces at the stated levels and real
types. For the original bounded row inputs, the
\hyperlink{proof.weighted-seesaw-rapid-test}{constant-term argument}
uses their integrable discrete-series reproducer and its
vanishing real unipotent constant terms to give cuspidality.
The core comparison uses the cuspidal multiplicity space
at the full deep principal level; it contains the core
Poincar\'e tests in the corank-one separation. Subsequent
triangle inputs remain cuspidal by the tower formula.
The Petersson tests separating the square outputs range
over the corresponding cuspidal spaces. Their rapid decay
and the Schwartz estimates give continuity of the pairings.
\cite[Theorem~2.7(b) and Propositions~2.9 and~3.4(ii)]{p2paul-real-theta}
\cite[Theorem~6.1(i)--(ii)]{p2ichino-real-theta}
\cite[Proposition~3.1]{p2wu-unitary-tower}

For a single edge the smooth oscillator quotient
onto the exterior contragredient and the intermediate
representation is unique up to scalar. The prescribed
joint harmonic has nonzero image in it. In fact the
larger discrete series is in the case of Paul's
Proposition~3.4(ii). Its lowest compact type has minimal
oscillator degree by the calculation in the proof of
that proposition. The paired lowest compact type on the
smaller side also has minimal degree by the joint
harmonic correspondence. In the proof of Howe's
Lemma~4.1(c), the joint harmonic of that degree has
nonzero image in the maximal isotypic quotient. Its
image generates the exterior lowest-type isotypic
space under the enveloping algebra of the other
member. The same conclusion holds in the irreducible
paired quotient. The polynomial fixed by the signed
minor construction is the joint highest vector in
this harmonic. Pairing its exterior compact factor
with the selected conjugate line gives the nonzero
intermediate lowest vector. These statements are
unchanged by the fixed characters at infinity.
\cite[Theorem~1, Lemma~4.1(c) and (4.5)]{p2flip-howe-real}
\cite[Theorem~2.2(2) and the proof of Proposition~3.4]{p2paul-real-theta}

To apply the two quotients, fix a product of finite
tests and exterior forms on their selected real lines.
For a polynomial edge vector \(Q_i\), first integrate
against the exterior form. Pairing the lift with
any intermediate cuspidal form on the indicated real
type gives a continuous bilinear oscillator pairing.
After the exterior and intermediate representations
are specified, this factors through the single
paired quotient by Howe's uniqueness. Equivalently,
the lift has the form
\[
 m_i\otimes u_i(Q_i)\quad\text{in}\quad
 \mathcal M_i\otimes\pi_H^{(i)} ,
\]
where \(\mathcal M_i\) is the automorphic multiplicity
space, \(m_i\) contains the finite data and exterior
form in the multiplicity space, and \(u_i\) is the
fixed real quotient contracted against the exterior
lowest vector. Normalize \(u_i(P_i)\) using the
prescribed harmonic \(P_i\). Nonzero images of \(P_i\)
make this normalization common to all choices of
\(m_i\). The Petersson or bilinear contraction through
\(H\) pairs the two multiplicity vectors and contracts
\(u_1(Q_1),u_2(Q_2)\) by the invariant dual pairing
of the two intermediate types. The latter pairing
is nonzero on \(u_1(P_1),u_2(P_2)\). Thus, for
\(Q_1\otimes Q_2\), the integral through \(H\) is its
value on \(P_1\otimes P_2\) times the ratio of the
real pairings. This ratio depends only on the real
vectors and the fixed representations.

Finite residue functions of primitive support are
finite linear combinations of products of matrix
indicators, using boxes inside the primitive support.
On every resulting product both descriptions give
the same kernel by the tensor product of the two
Weil actions. The primitive-support bound above
and the rapid decay of the exterior forms justify
the group integrations in either order. The real
ratio is the same on each term of the finite
expansion. For each fixed finite datum and pair
of exterior forms the resulting period is continuous
in the Schwartz topology. If the prescribed tensor
has nonzero period at one choice of these inputs,
the period at that choice divided by its value on
the prescribed tensor gives the continuous real
functional extending the ratio on polynomial tensors.
The identity with that functional extends by density
to arbitrary real vectors for every other choice of
finite data and exterior forms. If the prescribed
tensor gives zero for all choices, polynomial
tensor factorization likewise forces zero for all
real vectors on that path.

Now take the actual small-path tensor for \(v\)
and compare \(Tv\) on the large side with its actual
tensor. The preceding factorization gives
\(\kappa_*\) on the large-side form pairing, independent
of the twist and of the finite and exterior tests.
In the case where the large-side form pairings all
vanish take \(\kappa_*=0\). This proves that the
actual small-path pairing is \(\kappa_*/c\) times
the actual large-path pairing. The resulting
identity of operators is on the prescribed form
spaces. The bordered endpoint path uses the
pointwise comparison with \(T\) on its transported vectors.
\footnote{\weightedpin{grh/src/New/Determinant/Exchange/LargeTriple.lean}{381-L431}{Constancy of the equal-rank period},
\weightedpin{grh/src/New/Determinant/Exchange/TheoremForPartners.lean}{501-L562}{Nonvanishing down to the rank-zero test},
\weightedpin{grh/src/New/Determinant/Exchange/TheoremForPartners.lean}{568-L840}{Reciprocal comparison and the equal-rank value},
\weightedpin{grh/src/New/Determinant/Exchange/TheoremForPartners.lean}{241-L279}{Stable-space extension},
\weightedpin{grh/src/New/Determinant/Exchange/OneSquare.lean}{308-L365}{Primitive-support tail estimate},
\weightedpin{grh-compact/src/New/Exchange/OneSquare.lean}{59-L134}{Contraction to the two-arrow comparison}.}

For a triangle square the intermediate ranks differ by \(2(t-z)\).
The negative-count equality follows entry by entry. The changes from
\((i,n+1),(i+1,n+1)\) to \((i,n+2),(i+1,n)\) simultaneously flip the
two signs exactly when \(A\ge i+n+2\), where the old signs are
opposite. The stable inequality follows because each exterior node
has rank \(r+z+(n+1)(t-z)\), only \(r+z\) entries can be negative,
and \(t-z>r+z\). The character product has \(\nu^{-1}\) equal to
\((\upsilon^{-h}\chi_0\vartheta)^{n+1}\) in \(V\), where
\(\chi_0\) has residue character \(\omega\). The same formula with
\(\vartheta=1\) holds in \(Z\). The opening character \(b_i\) inverts
the \((i+1)\)-st character and the upper lift uses the
\((i+n+2)\)-nd. At same-parity ranks, the two roundings add to half
the total rank. The weight \((n+1)h\) exceeds the rank of either
middle column. Ramification is at the indicated places since
\(n+1\le m\), \(\omega^{n+1}\) is nontrivial at the chosen places, and
\(\vartheta\) is an \(m\)-fold twist. The loss in the scalar thus uses
one product of at most \(t-z\) ramified ratios at the actual power
\(\vartheta^{n+1}\).

Set \(p_{i,n+1}=p_{i,n}^+\), where the new block has size \(b=t-z\)
and
\[
 p^+\!\begin{pmatrix}Y&S\\ C&D\end{pmatrix}
 =\begin{cases}p(-Y)/\#\{S\},& C=0,\ D=I_b,\\0,&\text{otherwise.}\end{cases}
\]
For the two square kernels \(L(X)\) and \(L'(Y)\), applying the
partner rule to the graph of \(\binom{I}{-W}\) and \((W\ I)\)
gives \(|\mathrm{GL}_{a_0}(R_*)|L'\binom{I}{-W}
= (\kappa_*/c)|\mathrm{GL}_{a_1}(R_*)|L(W\ I)\), with
\(a_0,a_1\) the smaller and larger intermediate ranks. This follows
by summing over their invertible frames. The inverse determinant
characters in the partner weights cancel the characters of the
kernels. Specialize \(W=\binom{Y}{0}\) after exchanging the common
blocks. The determinant-one shear \(\bigl(\begin{smallmatrix}I&S\\0&I\end{smallmatrix}\bigr)\)
on the larger intermediate space leaves its value unchanged.
Averaging over \(S\) gives exactly \(p^+\) and the two weights
\(q_n,q_{n+2}\). In operator form the identity is
\[
 q_n a_{i+1,n}d_{i,n}
       =(\kappa_*/c)q_{n+2}d_{i,n+1}a_{i,n+1}
       =\frac{q_{n+2}}{\bar s_{i,n}}d_{i,n+1}a_{i,n+1},
\]
where \(s_{i,n}=\overline{c/\kappa_*}\) if \(\kappa_*\ne0\), and is
zero otherwise, with the convention \(0^{-1}=0\). Wherever the chain
gives nonzero \(s\), \(|s_\vartheta/s_1|^{-2}=|c_\vartheta/c_1|^{-2}\).

\paragraph{The realized reversals in the source triangle.}
\teamanchor{proof.weighted-grid-kernel-reversal}%
We use the reversal on the same extended source tests that
occur in the first commutation. Let the two original ranks be
\(n_0,n_1\) and let \(D\) be their finite matrix datum over
\(R_*=\mathcal O_E/\mathfrak m\). Set
\(u=\overline{\alpha_f\varpi_\mathfrak m}/
       (\alpha_f\varpi_\mathfrak m)\), with \(\varpi_\mathfrak m\)
the prescribed finite generator. On the two finite half-spaces
the normalized adjoint of a matrix map exchanges its two
polarizations. Conjugation first passes to the opposite
source form, with a finite multiplication by a unit of
norm minus one and the conjugated polynomial. On returning
by the adjoint map the same unit occurs as the change of
Fourier variable. The identity \(c_0^{-1}=-\bar c_0\)
for the fixed unit \(c_0\) of norm minus one shows that
the resulting primitive character is the original
\(\psi_\mathfrak m\). At a split pair the scaled vacuum
support on the second half has now been moved by \(u\);
the two factors in \(u\) are the conjugate and original
scales of this support. Off the modulus the factors from
the fixed Weil scale are included in \(u\) in the
same way.

For a point indicator at \(A\) the adjoint coordinates
carry its finite translation to the modulation
\(\psi_\mathfrak m(\operatorname{tr}(YA))\), in the
transposed variables. This follows from the Heisenberg
pairing of those two polarizations. Use the same gauge
as in the diagonal line reversal. The point indicator
is normalized by the measure \(1/|R_*|^{n_0n_1}\)
of its coset in the standard matrix box. The character
function has absolute value one on that box. Since the
normalized adjoint map is unitary, the factor relating
the indicator kernel to the unnormalized Fourier
character kernel has squared modulus
\(|R_*|^{-n_0n_1}\). The global change of theta functional
has modulus one by the rational Heisenberg normalization.
At infinity the conjugation and adjoint carry a matrix
coordinate to its signed transpose, so a polynomial
is changed to the signed transpose \(P^\natural\) of
its coefficientwise conjugate. In particular vacuum
moduli introduce no further twist-dependent factor.
Linearity in the functions on the residue quotient gives
\[
 \overline{K_{D,P}(g,h)}
  =\tau_\vartheta K_{\widehat{\bar D},P^\natural}(h,gu),
 \qquad |\tau_\vartheta|^2=|R_*|^{-n_0n_1}.
\]
Here the two determinant characters exchange positions,
with the same signs as the two original groups. This formula
uses the full matrix dimension for each edge.

\paragraph{The real entry of a reversed edge.}
\teamanchor{proof.weighted-grid-reversed-entry}%
On a reversed edge take the matrix entry of the smooth
one-edge oscillator quotient on the two selected real
lines. For any polynomial \(Q\) write \(\ell_\infty(Q)\)
for that entry, with the determinant characters and
real measures of the edge. A theta pairing with the
fixed exterior source and target types factors through
the same quotient by Howe's multiplicity one, as in
the square comparison. The joint harmonic
\(P_{\rm edge}\) has nonzero entry by the minimal-degree
calculation there. Thus replacement by \(Q\) multiplies
the pairing by \(\ell_\infty(Q)/\ell_\infty(P_{\rm edge})\).
The finite part of the pairing is the same in this
calculation. Equivalently, when \(Q\) is represented on
the selected line by convolution by \(\alpha_\infty\)
on the other group, the ratio is
\[
 \int \alpha_\infty(x)\,\overline{e'(x)}/e'(1)\,dx,
\]
where \(e'\) is the prescribed lowest-vector coefficient.
The opposite matrix entry conjugates this scalar.

For the actual reversal use the transported harmonic
\(P^\natural\) on the reversed edge. The signed transpose
and coefficient conjugation interchange the mixed-sign
Laplacians of the two real actions. They exchange the
paired compact types, with the reversed compact order
on the contragredient argument. They send the nonzero
joint highest line to the corresponding joint highest
line in those orders. Thus \(P^\natural\) has nonzero
matrix entry. Using the chosen generator of that same
line changes the reversal by a fixed nonzero real-place
normalization, independent of the finite test, points
and twist.
\footnote{\weightedpin{grh-compact/src/Lib/SiegelWeil/RallisInner/SharpData.lean}
{286-L337}{Conjugation and finite Fourier in both matrix directions},
\weightedpin{grh/src/Lib/Weil/DualPair/Flip.lean}
{489-L568}{Transport of a point indicator and exact scale},
\weightedpin{grh/src/Lib/DiscreteSeries/Theta/LocalCorrespondence.lean}
{734-L802}{Oscillator projection and the real convolution}
and \weightedpin{grh/src/Lib/Weil/Lift.lean}
{624-L682}{Application to the prescribed joint harmonics}.}

For the adjoint companion in \(V\), the reversed square has
equal intermediate ranks and the same weight and character
product, so the complementary integral has \(e=0\).
Its finite identity uses the \emph{same} \(p^+\).
More precisely, for an additive
Fourier variable \(\bigl(\begin{smallmatrix}A&B\\C&D\end{smallmatrix}\bigr)\),
the transform of \(\overline{p^+}\) is
\(1_{C=0}\widehat{\bar p}(-A)\psi(\operatorname{tr}D)\).
This is the trace expansion of the block product and orthogonality
of the average over \(S\). The corner sum over \(D\) now equals
the matrix Gauss sum times the value at \(D=I_b\). To see its modulus,
work at a field factor of the squarefree ring \(R_*\). The determinant
covariance forces a function of \(D\) to vanish at singular \(D\), by
a stabilizer scaling a complementary line with any determinant.
The same argument shows that its additive Fourier transform is zero
at singular dual matrices. At invertible dual matrices the values
differ by unitary determinant characters. Finite Plancherel applied
to the function \(\chi(\det D)\) on invertible matrices gives
\(|G_b|^2=|R_*|^{b^2}\), after multiplying the field factors. The free
block adds a factor \(\#\{S\}=|R_*|^{bN_0}\), where \(N_0\) is the
rank at length \(n\). Thus the reversed comparison uses
\(|G_b\#\{S\}|^{-2}=|R_*|^{-b(b+2N_0)}\). The translations from
the two reversals move through the intervening upward lift up to
its unitary determinant phase. Taking adjoints of the reversed
edges gives a factor of squared modulus \(|R_*|^{-N_0N_1}\) for
the first reversal and \(|R_*|^{N_1N_2}\) for the inverse second
reversal. Thus the scalar \(t_{i,n}\) in the companion has squared
modulus a fixed multiple of
\(|R_*|^{N_1(N_2-N_0)-t(t+2N_0)}=|R_*|^{t^2}\), since
\(N_1=N_0+t\) and \(N_2=N_0+2t\) in \(V\).
The real type convolutions are fixed in the twist. In detail, one
first writes \(\langle d_{i,n}a_{i,n+1}^*x,y\rangle\) as a multiple
of the pairing of \(a_{i,n+1}^*x\) with the reversed lift of
\(R(u^{-1})y\). Move the upper adjoint to the other argument, commute
the two lifts by the reversed rule, move the central translation
through the lower lift, then reverse the remaining left edge back.
This gives \(t_{i,n}\langle d_{i,n+1}x,a_{i+1,n}y\rangle\).
The scalar introduces no ramified ratio. The first commutation is
available in both triangles, while the adjoint companion uses the
squarefree ring of \(V\).
\footnote{\weightedpin{grh/src/New/Determinant/Chain/SquareLaw.lean}{152-L214}{Shears and group orders},
\weightedpin{grh/src/New/Determinant/Chain/SquareLaw.lean}{394-L511}{Finite Fourier transform of the same extension},
\weightedpin{grh/src/New/Determinant/Chain/ChainNorms.lean}{157-L235}{First triangle square},
\weightedpin{grh/src/New/Determinant/Chain/ChainNorms.lean}{261-L371}{Reversed, equal-rank square},
\weightedpin{grh/src/New/Determinant/Chain/ChainNorms.lean}{539-L661}{Combining the same descents and their adjoints}.}

\paragraph{The Gram recurrence.}
The square law just proved is
\[
  q_n a_{i+1,n}d_{i,n}
    =\frac{q_{n+2}}{\overline{s_{i,n}}}d_{i,n+1}a_{i,n+1},
\]
with the reciprocal of zero read as zero. Its adjoint
companion is
\(\langle d_{i,n}a_{i,n+1}^*z,y\rangle
 =t_{i,n}\langle d_{i,n+1}z,a_{i+1,n}y\rangle\).
Here \(s_{i,n}\) costs one square and \(t_{i,n}\) costs zero.
The first descent is adjoint to the first lift on the
base forms under the inclusions, and the composite
\(d_{i,0}a_{i,0}e_i=e_{i+1} B_i\) is injective there.

Induction on \(n\) gives scalars \(s_i^{(n)},t_i^{(n)}\), products of
the first \(n\) successive square factors in row \(i\), such that
\begin{align*}
 M_{i+1,n}d_{i,0}a_{i,0}
   &=\frac{q_{n+1}}{q_1q_0\overline{s_i^{(n)}}}d_{i,n} M_{i,n+1},\\
 \langle z,a_{i,n}M_{i,n}e_i y\rangle
   &=t_i^{(n)}\langle d_{i,n}z,M_{i+1,n}e_{i+1}y\rangle.
\end{align*}
For the first formula multiply by the next weighted lift
\(q_n a_{i+1,n}\) and use the square law. For the second move
\(a_{i,n+1}\) to the first variable, apply the current identity and
then its adjoint companion, using \(q_n a_{i+1,n}M_{i+1,n}=M_{i+1,n+1}\).
For \(n=0\) the first formula follows from \(M_{i,1}=q_0a_{i,0}\)
and the second is the first-arrow adjoint identity. The first formula
also proves injectivity of the weighted chains on the base by
induction, since applying the last descent to a zero image gives a
zero image of \(B_i y\) under the next shorter chain. For any nonzero
base form \(y\), the left side of the first identity evaluated at
\(e_i y\) is therefore nonzero. Hence \(s_i^{(n)}\ne0\) for all the
factors used in the norm comparison. For the second identity take
\(z=a_{i,n}M_{i,n}e_i y\) for a nonzero base vector \(y\).
The left side is then positive by chain injectivity, so
\(t_i^{(n)}\ne0\) on the same range. This establishes the
nonvanishing of the realized comparison factors used here.

The first-arrow Gram form is
\(q_0\langle a_{i,0}e_i x,a_{i,0}e_i y\rangle\) on the good
base forms. Write \(S_i\) for its operator.

\paragraph{The first-arrow pairings on the two inner forms.}
\teamanchor{proof.weighted-first-arrow-rallis}%
Put \(j=i+1\), \(\psi=\vartheta^j\), and
\(D_\psi=\mathcal O_E/(\mathfrak M Q_\psi)\). The support of
\(Q_\psi\) is that of \(Q_\vartheta\) for \(1\le j\le m\).
Both first-arrow pairs in the comparison have ranks \(r\) and
\(r+t=6r\). The target dimension minus its global Witt index is
at least \(3r>r\) on the grid, and the auxiliary target is definite.
Thus both are in Yamana's convergent case. Both pairs have the
joint harmonics of the corresponding real edge, with weights
independent of \(\vartheta\), and the source vectors in these
types are cuspidal. Apply the convergent Siegel--Weil formula and
the \hyperlink{proof.unitary-doubling-identity}{graph-orbit unfolding}
used in the \hyperlink{proof.stage-petersson}{earlier stage pairing}.
In Yamana's notation the doubled skew Hermitian space has
dimension \(2r=2n_{\rm Yam}\) and the target has
\(m_{\rm Yam}=6r\). The normalized point is
\(s_0=(m_{\rm Yam}-n_{\rm Yam})/2=5r/2>n_{\rm Yam}/2\).
Use Theorem~2 with the convergence bound above on the mixed-sign target.
For the neutral doubled datum \((\upsilon^r,1)\), write \(D\)
for the action of a Siegel Levi element on the stabilized
diagonal isotropic subspace. The zero term uses the inverse of
the source splitting character \(\upsilon_E^{-6r}\), so its
factor at \(D\) is
\[
 \upsilon_E(\det D)^{6r}|\det D|_{\mathbb A_E}^{3r}.
\]
Here \(3r=s_0+r/2\) includes the half-modular character. Thus the
normalized inducing character is \(\chi_\square=\upsilon_E^{6r}\)
in this Levi coordinate, with restriction \(\epsilon_{E/F}^{6r}\)
to \(\mathbb A_F^\times\).
Under the polarization change, the neutral doubled action on
\((g_1,g_2)\) is the neutral action on the first copy tensored
with its dual on the second copy twisted by
\(\chi_\square(\det g_2)\).
\cite[\S11.3, p.~52 in arXiv v3]{gqt}
The two-source Petersson integral therefore multiplies the doubled
section by
\[
 \chi_\square(\det g_2)^{-1}\,b(\det(g_2^{-1}g_1))^{-1}.
\]
The second factor restores the original \(b(\det g_1)^{-1}\) and
its conjugate at \(g_2\). On the graph representative
\((\gamma',1)\) the \(\chi_\square\)-factor is one, giving
\(b(\det\gamma')^{-1}\) in the unfolded source coefficient.
The earlier stage pairing uses conjugated bilinear inputs
and its second-copy factor is written there as \(\chi_B(\det g_2)\)
for the inducing character \(\chi_B^{-1}\).
Take the fixed polynomial Fock tensor as the Schwartz
datum and evaluate its section at the two source points in
the doubled group. Primitive support at a split place forces
the rational Lagrangians that contribute to be graphs.
Rapid decrease of the source coefficients justifies the
unfolding with this fixed section. The
\hyperlink{proof.weighted-sharp-real-pairing}{same-vector real
calculation} gives a positive scalar on the joint harmonic. Include
the target quotient volume in that scalar. The opposite-space
identification of the source is isometric and intertwines
finite translations.\footnote{\href{https://msp.org/pjm/2013/264-1/pjm-v264-n1-p11-s.pdf\#page=4}
{Yamana, The Siegel--Weil formula for unitary groups, pp.~235--237 and Theorem~2};
\weightedpin{grh-compact/src/New/Exchange/GridParameters.lean}{16-L67}{Ranks, signs and harmonic};
\weightedpin{grh-compact/src/Lib/Weil/DualPair/PairSplitting.lean}
{302-L328}{The inverse splitting character in the stabilized-space Levi};
\weightedpin{grh-compact/src/Lib/Weil/DualPair/PairSumNegation.lean}
{208-L241}{The dual and the doubled tensor};
\weightedpin{grh-compact/src/Lib/SiegelWeil/RallisInner/SharpData.lean}
{93-L189}{Primitive graph support, double datum and measure}.}

For the first arrow in \(V\) use the
\hyperlink{proof.weighted-triangle-point-kernel}{principal support
and covariance just proved} with \(D'=D_\psi\) and \(q=r\).
Its point indicator has squared norm \(|D_\psi|^{-(r+t)r}\) in
the residue measure. The corresponding integral over the compact
source is \(q_0^{-1}|D_\psi|^{-(r+t)r}\), where
\(q_0=|{\rm GL}_r(D_\psi)|\).

For the auxiliary lift \(u_\psi\), Petersson adjunction with
the definite class-sum metric realizes its values by the
highest-vector theta lift. The projector on the \(Q_\psi\)-part
of the section is a sum over source translates by
\(\prod_{v\in Q_\psi}{\rm GL}_r(k_v)\). They act trivially on
the chosen-level source forms. The resulting section is the
indicator with marked rows equal to the identity multiplied by
\(\gamma=\prod_{v\in Q_\psi}|{\rm GL}_r(k_v)|\).
The unmarked entries at each of these places give \(q_v^{rt}\)
possibilities. The auxiliary mass is therefore
\[
 \frac{\gamma^2\,\operatorname N(Q_\psi)^{rt}}
 {q_0|D_\psi|^{(r+t)r}}.
\]
The conversion from Petersson to the definite class-sum norm
multiplies by the auxiliary level index and by the fixed
dimension of the coefficient representation. The resulting
factor is \(A_j\Delta_\vartheta |D_\psi|^{-(r+t)r}\), where
\(A_j\ne0\) is fixed and \(\Delta_\vartheta\) has exactly the
meaning of the auxiliary Gram scalar used in the terminal
level comparison. The finite count also gives
\[
 \Delta_\vartheta |D_\psi|^{-(r+t)r}=
 C_{\mathfrak M}\operatorname N(Q_\psi)^{rt}
 \prod_{v\in Q_\psi}\prod_{h=t+1}^{t+r}(1-q_v^{-h}),
\]
since the local Gram factor \(G_{r+t}(q_v)/G_t(q_v)\)
contains \(q_v^{r^2+2rt}\), with \(G_s(q)=|{\rm GL}_s(\mathbf F_q)|\).
This is the factor when the common
Hecke sum is written without its residue-matrix denominator.
The character of the \(r+t\) side in both lifts is
\(\mu_j^0\psi\,\upsilon^{(r+t)/2}\), since \(r+t\) is
even. At the places away from the modulus and its conjugate,
the vacuum matrix coefficient for a source element \(\gamma'\)
is consequently the same scalar
\[
 \rho_{j,\psi}(\gamma')=
 \left((\mu_j^0\psi\upsilon^{(r+t)/2})(\det\gamma')\,
                    \Delta_f(\gamma')^{r+t}\right)^{-1}.
\]
Here \(\Delta_f\) is the finite Weil multiplier in the original
vacuum convention. The change from the realized level-\(J_m\)
chain to a definite class-sum metric occurs after this
comparison.\footnote{\weightedpin
{grh-compact/src/New/Exchange/CharacterLiftNorm.lean}{33-L80}
{Lift, source translates and Gram factor};
\weightedpin{grh-compact/src/New/Exchange/CharacterLiftNorm.lean}{82-L86}
{The normalized off-level term}.}

Let \(K\) be the chosen level on either rank-\(r\) group. With
\(x,y\) fixed by \(K\), the matrix coefficient of
\(\rho_V(\gamma')\) can be replaced by that of
\[
 \frac{1}{[K:K\cap\gamma'K(\gamma')^{-1}]}
   \sum_{k\in K/(K\cap\gamma'K(\gamma')^{-1})}
                    \rho_V(k\gamma'k^{-1}).
\]
Indeed each summand has the same pairing with \(x,y\), and
changing the chosen representative of \(k\) changes the
action on \(y\) by an element of \(K\). This operator on
\(K\)-fixed vectors is the normalized double-coset operator.
The common finite adelic identification uses \(\gamma'\)
supported off the chosen places and the twist places. The
coefficients \(\rho_{j,\psi}(\gamma')/|D_\psi|^{(r+t)r}\)
of the normalized double cosets agree in the two expressions.
The unfolded integrals converge absolutely on each matrix
entry. Since both base spaces are finite dimensional, the
series lies in the closed paired image of the off-level
Hecke algebra. Its operator \(P_{j,\psi}\) thus acts on the
good spaces in both inner forms, commutes with the chosen
finite group and preserves the good subspaces. The two Gram
forms \(q_0\langle a_{i,0}e_i x,a_{i,0}e_i y\rangle\) and
\(\langle u_\psi x,u_\psi y\rangle\) are the actions of
\(P_{j,\psi}\) with scalars \(C_i\) and \(A_j\Delta_\vartheta\),
respectively, for \(C_i\ne0\) independent of \(\vartheta\).
The \hyperlink{proof.weighted-good-transfer}{good trace transfer}
gives equal determinants of \(P_{j,\psi}\) on the two good
spaces. This yields \(B_i=S_i/q_0\) in the realized first
descent and the scalar required by the chain recurrence.
\footnote{\weightedpin{grh-compact/src/New/Trace/GoodFormTransfer.lean}
{24-L49}{Common Hecke image};
\weightedpin{grh-compact/src/New/Rows/RowVolume.lean}{32-L58}
{Application to the original rows}.}

\paragraph{The determinant of the chain.}
Pairing the two images under
\(M_{i,n+1}\) and using both induction formulas reduces their pairing
to that for \(M_{i+1,n}\) on \(B_i x,y\), multiplied by
\(q_nq_1q_0\cdot t_i^{(n)}s_i^{(n)}/q_{n+1}\). Thus the full Gram form on
\(W\) for the weighted chain in row zero is that of
\(S_0\cdots S_{m-1}\), times
\(\bigl(q_0q_1^m/q_m\bigr)\prod_{i=0}^{m-1}t_i^{(m-i-1)}s_i^{(m-i-1)}\).
The cost of the latter product of square factors is at most
\(\sum_i (m-i-1)\le m^2\) squares. This recurrence takes determinants
of the product without a choice of a common eigenbasis.

The first-arrow calculation gives the precise determinant ratio
\[
  \det S_{j-1}
    =\left(\frac{C_{j-1}}{A_j\Delta_\vartheta}\right)^R
                     \lambda_j(\vartheta^j).
\]

The conversion from the weighted triangle norm to the class-sum norm
multiplies its Gram form by the final weight dimension and by the
last restriction level index. The following count identifies that
index with the factors used in the single-lift comparison.

\paragraph{The terminal level in the chain norm.}
\teamanchor{proof.weighted-chain-level-index}%
For the fixed chain let \(J_m\) be its last restriction level.
At a new place with residue field of size \(q\), the block sizes
of its parabolic are \(r,t,\ldots,t\) with \(m\) blocks of size
\(t\). Each strictly upper block is arbitrary, so its unipotent
order is \(q^{e_m}\), where \(e_m=mrt+\binom m2 t^2\).
Writing \(G_u(q)=|{\rm GL}_u(\mathbf F_q)|\), the local index is
\[
 \frac{G_{r+mt}(q)}{G_r(q)G_t(q)^m q^{e_m}}.
\]
This uses the same ordered-frame count as the earlier
two-block norm index. For the single auxiliary lift the
local target index is \(G_{r+t}(q)/(G_r(q)G_t(q)q^{rt})\).
Its section has \(q^{rt}\) free residues after fixing the marked
submatrix, and the source projection contributes \(G_r(q)\)
before squaring. The normalization divides by \(G_r(q)\).
Thus the local factor of the actual auxiliary Gram scalar
\(\Delta_\vartheta\) is \(G_{r+t}(q)/G_t(q)\). Meanwhile
the triangle weight \(q_n\) acquires \(G_{r+nt}(q)\) at
this place. The same level modulus occurs for the \(m\)
powers, by the definition of the fitted \(m\)-fold twist.
Multiplying gives
\[
 [K_{\rm full}:J_m]=C\operatorname N(Q_\vartheta)^{-e_m}
                  \Delta_\vartheta^{\,m}\frac{q_m}{q_0q_1^m},
\]
with \(C\) fixed by the untwisted chosen places. Hence the
weight product cancels on the realized chain, leaving a
conductor power of absolute value one at the fitted place.
The remaining passage to the terminal class-sum norm uses
the dimension of its compact coefficient type and the
actual realization of the weighted arrows. The square and
adjoint recurrences use the same descent datum at a node in
both applications. Use the common Rallis operator and Hecke transfer in the chain
calculation above.\footnote{\weightedpin
{grh-compact/src/New/Seesaw/GradingsOfY.lean}{96-L99}{The exact restriction level},
\weightedpin{grh-compact/src/New/Exchange/CharacterLiftNorm.lean}{58-L80}{The auxiliary Gram factor},
and \weightedpin{grh-compact/src/New/Exchange/ChainNorms.lean}{16-L47}{Level and norm on the realized chain}.}

Taking absolute values thus gives
\[
  Z_\vartheta=|c_\vartheta|^R
                      \prod_{j=1}^m\lambda_j(\vartheta^j),
  \qquad Z_\vartheta\ne0.
\]
The square-factor product controls \(c_\vartheta\) up to a fixed
constant and a conductor norm power. The latter is a unit at the
fitted place \(\iota\). Each charged square gives at most \(t\)
squared ramified inverse ratios, so the local logarithm of the
squared absolute ratio is at most \(2m^2t v_\ell\).%
\footnote{\weightedpin{grh-compact/src/New/Exchange/ChainNorms.lean}{15-L49}{Last level and norm normalization},
\weightedpin{grh-compact/src/New/Rows/RowVolume.lean}{32-L58}{Application with the actual chain and good subspaces},
\weightedpin{grh-compact/src/New/Exchange/CostOfSquares.lean}{26-L63}{Conductor and ramified factors}, and
\weightedpin{grh-compact/src/New/Exchange/TriangleNorms.lean}{36-L66}{Multiplication in the triangle}.}
Dividing by two gives
\begin{equation}
\label{eq:weighted-unsquared-chain}
 \log|j_\iota(|c_1/c_\theta|)|_\ell\le m^2t v_\ell .
\end{equation}

Fix the product of the untwisted compressed determinants
\[
 D_0=\prod_{j=1}^m d_j(1),\qquad
 \log|\sigma D_0|\le 2RrmB_0 \quad\text{for every }\sigma .
\]
The single-lift ratio identity gives
\(D_0K_\Gamma^R Z_1/Z_\theta
 =|c_1/c_\theta|^R\bigl(\prod_j d_j(\theta^j)\bigr)
                       (K_\Gamma Q_\Gamma(\theta))^R\).
Use the factor counts of \eqref{eq:weighted-aggregate-hecke},
the joint fitted product at this place and
\eqref{eq:weighted-unsquared-chain} to obtain
\begin{equation}
\label{eq:weighted-volume-lift}
\begin{aligned}
 \log|j_\iota(D_0K_\Gamma^R Z_1/Z_\theta)|_\ell
    \le{}& R(r+\kappa)\sum_j\beta_{j,\ell}+2Rm\kappa P_\ell\log\ell\\
         &+ R(2m\kappa+10m^2+m^2t)v_\ell .
\end{aligned}
\end{equation}

\paragraph{Endpoint period maps.}
\teamanchor{proof.weighted-completion-period}%
Write \(\mathcal H_0\) for the full principal-level space on \(Y_0\)
and \(\mathcal H_{\rm c}\) for the prescribed core forms on
\(Z_{0,0}=Y_0\perp\langle1\rangle\). Let \(\mathscr F\) be the level
space on the definite end of the core column. Write
\(S\colon\mathscr F\to\mathcal H_{\rm c}\) for its fixed
untwisted descent. It is onto by the
\hyperlink{proof.weighted-endpoint-first-path}{transpose calculation}
below.

For each real translate \(a_\infty\) let \(r_{a_\infty}\colon\mathcal H_{\rm c}\to
\mathcal H_0\) represent the bordered restriction pairing by Riesz
duality. Write \(K\) for the chosen principal level on \(Y_0\) and
\(K^{\rm deep}\) for the level of modulus \(\bar\pi_G^{b_1}\) there.
The restriction identity against a test of the twisted bottom level
uses the same \(r_{a_\infty}\). Indeed, for \(k\in K\) choose \(v\) supported
at the chosen places with reduction inverse to \(k\) at depth \(b_1\).
Then \(v\) fixes the test, since it is principal modulo the chosen
places and is trivial on the twist places. The element \(kv\) lies
in \(K^{\rm deep}\), so its block image fixes the core form. By
simultaneous finite translation in the period, replacing the test by its \(k\)-translate
does not change the period. Averaging over \(K\) gives the asserted
pairing against \(r_{a_\infty}\), independent of the twist. Put \(Q_{a_\infty}=r_{a_\infty}S\).
\footnote{\weightedpin{grh-compact/src/New/Exchange/GridFirstColumns.lean}{32-L69}{Simultaneous translation and the common bordered map}.}

Write \(T_\vartheta\) for the final-level chain on \(\mathcal H_0\),
with restriction \(\Theta_\vartheta\) to \(W\). At \(Y_m\) let
\(K_{\rm full}\) be the full integral terminal level. Write
\(K_\vartheta^{\rm pr}\) for its subgroup with identity
reduction at the chosen and twist places, in the specified
split coordinates. Write \(K_\vartheta^{\rm res}\) for the
subgroup with identity reduction at the chosen places and
upper block triangular reduction at the places of \(Q_\vartheta\),
with block sizes \(r,t,\ldots,t\) and \(m\) blocks of size \(t\).
The latter is the restriction level used in the
\hyperlink{proof.weighted-chain-level-index}{chain norm}.
Its tame character on an element with diagonal blocks
\(k_{ll,w}\) is
\(\prod_{w\mid Q_\vartheta}\prod_{l=1}^m
  \vartheta_w^l(\det k_{ll,w})\), trivial on
\(K_\vartheta^{\rm pr}\).
The initial path identity, exchanges and normalization of the
last path give
\begin{equation}
\label{eq:weighted-completion-period}
   [K_{\rm full}:K_\vartheta^{\rm pr}]\int\Delta_{h,v}\,
        T_\vartheta(Q_{a_\infty}\phi)
     =\frac{\alpha_{a_\infty}}{c'_\vartheta}
         \mathcal E_\vartheta(\Delta_{h,v},\phi).
\end{equation}
The integral takes the automorphic realization of the algebraic
form at level \(K_\vartheta^{\rm pr}\), with \(\Delta_{h,v}\)
the dual-weight class delta test at \(h\) of value \(v\).
The determinant uses the class sum at \(K_\vartheta^{\rm res}\).
The two readings evaluate the same lift at a given adelic argument.
The normalized path calculations below give the same \(c'\)
for all real translates, with a scalar \(\alpha_{a_\infty}\)
that may depend on the translate.
\footnote{\weightedpin{grh-compact/src/New/Seesaw/GradingsOfY.lean}
{30-L32}{The smaller pairing level} and
\weightedpin{grh-compact/src/New/Seesaw/GradingsOfY.lean}
{96-L104}{The determinant level}.}

\paragraph{The first path and its level weights.}
\teamanchor{proof.weighted-endpoint-first-path}%
The weighted lift with the inclusion datum and the descent on the
first column of a grid use the same kernel. Their transpose identity
follows by Fubini in the bilinear pairing of a form and its conjugate
type, with the factor \(q_j\) on each arrow. Each endpoint image
after a nonempty sequence is fixed by the kernel of reduction in
its triangle ring. The inclusion datum is fixed there on either
side, since both determinant characters are trivial on this kernel.
In \(Z\) take the actual untwisted ring of depth \(b_1\). The core
descent maps its terminal principal-level space to the bottom
principal-level space. If a bottom vector annihilates the image, test
the transpose identity with the conjugate of its upward image at the
terminal level. The resulting pairing is the squared norm of the
upward image. The first commutation and the injectivity of the first
arrow, applied inductively to the core triangle, make the upward
image injective, so the bottom vector is zero. This proves
surjectivity of the required descent \(S\) on the full core level.

For the initial path put \(R=\mathcal O_E/(\mathfrak bQ_\vartheta)\).
In the grid let \(I_{z,a,b}\) be the index of the residue subgroup at
node \((a,b)\). This subgroup is trivial on the first column of \(V\).
In \(Z\) and off the first column of \(V\) it is the kernel of reduction
to \(\mathcal O_E/\mathfrak b\). Thus \(I_{1,a,b}\) and
\(I_{0,a,b}\) for \(a>0\) are independent of the twist. On the first
column of \(V\), the ratio \(q_j/I_{0,0,j}\) is independent of the
twist by reduction of matrices from \(\mathfrak bQ_\vartheta\) to
\(\mathfrak M Q_\vartheta\). At a truncated discrete valuation ring
the kernel consists of all matrices congruent to the identity at the
smaller depth. Its size is the corresponding residue cardinality to
the power \(N^2\) at a node of rank \(N\). On the first column of \(Z\)
the two orders are equal.

Expand the two first-column descents by pulling back their respective
inclusion data to \(R\), and then summing point indicators. The
coefficient is either zero or one. The matrices on the twisted edges
have left inverses because they reduce to the inclusion at every maximal
ideal and hence split over each local factor of \(R\). The sum has
at most \(|R|^{(2n_m)^2(2m)}\) terms. Let \(\Phi\) be this sum as a
linear functional on functions of tuples. If \(g_p\) acts at node
\(p\) and a matrix on edge \(p\to q\) changes to
\(g_q^{-1}Xg_p\), the character of the edge is evaluated at
\(\det(g_q)/\det(g_p)\). Write \(\eta(g)\) for the product of these
unitary factors over the two legs. Initially \(\Phi\) has the required
covariance for the residue subgroup at every node because reductions
to the triangle rings act trivially there.

Choose one representative \(t_x\) for each tuple of cosets of these
subgroups at the \(2m\) nonterminal nodes. At the bottom node of
\(Z\), use the block extension of the chosen bottom representative
in \(V\). Keep the former coset index in the sum so its multiplicity
is retained. At the terminal nodes take the identity. For a function
\(F\) of the tuple write \(A_gF(X)=F(g\cdot X)\), and put
\(\Lambda_0(F)=\sum_x\eta(t_x)\Phi(A_{t_x^{-1}}F)\). To check
covariance, permute the cosets by an admissible family \(g\). The
error \(t_{gx}^{-1}gt_x\) lies in the level at each node. At the
modified bottom node of \(Z\) this holds because the error in \(V\)
is the identity and both lifts are the same block extension there.
The character of this error cancels against the changed coefficient
\(\eta(t_{gx})\). Thus \(\Lambda_0(A_gF)=\eta(g)\Lambda_0(F)\). The
coefficient in each term is a root of unity. This second sum adds at
most \(2m\) tuple powers, giving order \(4m\) in the compact model.

In the actual bilinear period at the bottom, simultaneous translation
in the two groups by \(k\) and \(k\perp1\) preserves the integral
for every real translate. Interior translations on each leg cancel,
leaving exactly \(\eta(t_x)^{-1}\) on the term read at \(t_x^{-1}\).
Hence the coset sum contributes the product of the \(2m\) vertex
indices. This converts the weight product from the expansion into
\(\prod_{j<m}(q_j/I_{0,0,j})\), a fixed constant. At the terminal
principal level \(K_\vartheta^{\rm pr}\)
the finite index is \(q_m\). The stage periods put \(I_{0,0,m}\) on
the test form before using the two grid paths. Converting this last
factor to \(q_m\) adds \(q_m/I_{0,0,m}\). The initial identity
therefore holds with \(C_0=\prod_{j=0}^{m}(q_j/I_{0,0,j})\), with
all the \(q_j\) here taken on \(V\). Insert the twist-free bordered
Riesz map \(r_{a_\infty}\) and the core descent \(S\), and transpose
the \(V\)-leg. This gives the left side of
\eqref{eq:weighted-completion-period} with \(Q_{a_\infty}=r_{a_\infty}S\)
and the right side \(C_0\mathscr P_0\). The datum \(\Lambda_0\) is
independent of \(a_\infty\).
\footnote{\weightedpin{grh-compact/src/New/Exchange/RowsReached.lean}{18-L59}{Weighted transpose at the same core level},
\weightedpin{grh-compact/src/New/Exchange/GridFirstColumns.lean}{22-L30}{Ratio of group orders to indices},
\weightedpin{grh-compact/src/New/Exchange/GridFirstColumns.lean}{199-L272}{Point data and compatible coset averaging}.}

\paragraph{Pointwise transport through the off-diagonal squares.}
\teamanchor{proof.weighted-endpoint-grid-transfers}%
\teamanchor{proof.weighted-endpoint-pointwise-transport}%
At a corner \((a,b)\) the old path between the exterior nodes
\((a,b+1)\) and \((a+1,b)\) goes through \((a,b)\). The new path
goes through \((a+1,b+1)\). Thus moving the descent past the lift
uses the complementary identity from the smaller to the larger
intermediate space. Apply it to the kernel integrated over that
intermediate group at every pair of exterior arguments.
For an off-diagonal square with corner \((a,b)\), the two intermediate
ranks differ by \(2t\). At offset zero the negative count in either
grid depends only on \(b\) and is \(\lfloor(m-b+z)/2\rfloor\).
Indeed, the only negative entries are the stem positions with
\(b<A\le m-1+z\) and \(A+b\) odd. This verifies the signed sum
condition. The exterior rank is at least \(2(r+z+t)\), and each
exterior node has at most \(r+z\) negative signs. Since \(t>r+z\),
the rank is at least twice the negative count plus two. The total
rank is at most \(4n_m\) and
the weight \(w_{a+1}+w_{b+1}\) exceeds \(2n_m\). The floor and ceiling
in \(V\) add to half the total rank because the two ranks have the
same parity. The two products in \(Z\) have integral halves. In
the square rule this gives
\(\nu^{-1}=\mu^0_{a+1}\mu^0_{b+1}\vartheta^{l}\), where \(l=b+1\)
for \(V\) and \(l=a+1\) for \(Z\). The residue exponent of
\(\omega\) at chosen places is \(a+b+2\le 2m\), and the twist is
nontrivial at its places by the \(m\)-fold condition.
Use the
\hyperlink{proof.weighted-grid-real-transport}{pointwise complementary-period identity}
with \(e=t\) on the finite functions of modulus
\(\mathfrak M^{b_1} Q_\vartheta\). Its arbitrary residue depth
covers both grids, including the ordinary \(V\)-square of the
diagonal move.

Let \(\mathcal V_0(a_\infty)\) be the real test on the initial
path. Move the right translate on the core form into the adjacent
real Weil vector by covariance of its theta kernel. This uses the
same real action at every twist, since \(\vartheta\) has weight zero.
At each square let \(T_\square\) be the small-to-large real map
from the pointwise identity. The zero-slice extension uses fixed
Schwartz factors of value one. The top construction uses the
fixed compact chart and bump and tensors with its fixed
determinant vector. The real cone maps have positive exponents
along this descent to \(e=t\). Their estimates give continuity
in the Schwartz topology. Thus \(T_\square\) depends only on the
real data of the square. The finite Tate factors remain in the
normalized comparison scalar \(c_\square(\vartheta)\), whose
squared modulus is the fixed multiple of \(\Pi_{\nu,S}(t)\)
in the complementary identity.

Use \(T_\square\otimes 1\) on the whole real test, with the other
edges as tensor factors. First take a finite sum of decomposable
real tests. At fixed exterior arguments and fixed finite tuple,
the spectator factors come out of the integral over the square
intermediate group. The finite tuple map below then extends
linearly over the partner graph with the same \(T_\square\) on
every term. The primitive-support estimate in the
\hyperlink{proof.weighted-grid-primitive}{split model} bounds
the absolute intermediate integral by a Schwartz seminorm
times a power of the exterior heights. It applies to the
complete column and row stacks of the partner data.
To pair at an exterior node, integrate the prefix of the leg
from its initial form. The
\hyperlink{proof.weighted-grid-real-contraction}{one-edge first-occurrence calculation}
applies in the traversed orientation to each pure Schwartz
factor. The output is a smooth cusp form of the prescribed
type. The theta kernel's derivative bounds and the
zero-constant-term rapid decay estimate bound its rapid-decay
seminorms by Schwartz and input cusp seminorms at the fixed
finite level. These estimates justify Fubini at the
intermediate group and the adjacent exterior node. Rapid decay
also bounds the paired border integral. A short vector in the
smaller lattice is still short in the larger lattice of the
block embedding, and a fixed translate changes the height
bound by a constant. For each fixed finite tuple this gives
continuous multilinear bounds on pure tests and hence continuous
linear periods on the Schwartz space of the direct sum. Density
of finite sums of pure tests gives the same identities for
arbitrary real tests. The estimates may depend on the fixed
finite tuple. The real map and the scalar in the pointwise
identity are common to all tuples.

The finite move on the actual tuple sums uses the following direct
description. For matrices \((X,Y)=(X_k,Y_k)\) in the partner graph
give the pair coefficient \(\chi(\det k)/|\mathrm{GL}_b(R)|\), with
\(b\) in this formula the smaller intermediate rank, and give all
other pairs coefficient zero. Completions of a pair differ by a
determinant-one block shear, so the coefficient is well defined. For
fixed primitive \(Y\), sum over its orbit \(Yg\) with weight
\(\chi(\det g)^{-1}\). The two partner sums then give the pointwise
square rule with this graph coefficient. Indeed each of the
\(|\mathrm{GL}_b(R)|\) replacements gives the same product of
weights and kernel values. If a block of \(Y\) has a left inverse,
the opposite block of \(X\) does as well. To see this, reduce at each
maximal ideal and use \(\ker X=\operatorname{im}Y\): a vector in
the kernel of the opposite block would give a vector of
\(\operatorname{im}Y\) with zero projection on the original
primitive block. That projection is injective on
\(\operatorname{im}Y\). Full column rank in all residue fields
gives a left inverse over \(R\).

Use one representative for each right \(\mathrm{GL}_b(R)\)-orbit of
\(Y\). Retaining this representative and multiplying its graph
coefficient by \(|\mathrm{GL}_b(R)|\) gives the same functional as
keeping the average. This follows by averaging over the admissible
change of basis at the removed node, extended by one on the extra
line in a diagonal move. The action is free on the primitive
frames. Under exterior changes \(g\), new-row change \(f\) and
old-column change \(v\), the completion changes from \(k\) to
\(g^{-1}k\operatorname{diag}(v,f)\). The character of its determinant
is exactly the quotient of the new and old edge characters.
Consequently the new functional has the required covariance and
root-of-unity coefficients. The two new matrix variables add at
most two to the tuple order. The finite split matching at depth
\(b_1\) uses this graph over \(R\).

At stage \(j\) move the last vertical descent past the \(j\) lifts
in each grid. Apply the square rule with all other edges fixed and
with the complementary pairing as the linear functional. These
\(2j\) moves raise the order by \(4j\) and use \(2j\) square
products. What remains is the pair of descents to the diagonal
\((j,j)\). There is an isometry
\(Z_{j,j}=V_{j,j}\perp L_j\), with line sign \((-1)^j\), by the
negative counts and the actual diagonal basis identifications. In
the off-diagonal position the groups \(Z_{j,j+1}\) and \(V_{j+1,j}\)
are identified by a relabeling, since they have equal negative
counts and the same even rank \(N_j=r+(2j+1)t\). Use this common
group for the seesaw. At the old diagonal the \(V\)-rank is odd:
adding the line preserves the ceiling of half the rank. Thus the
kernel on the \(Z\)-edge leaving that diagonal splits into the
corresponding untwisted \(V\)-kernel and the kernel from \(L_j\) of
line character one. At the new diagonal the floor increases by one.
The new twisted \(V\)-kernel is multiplied by the kernel to
\(L_{j+1}\) of line character \(\upsilon\). At the common even-rank
group the twisted and untwisted characters differ by
\(\vartheta^{j+1}\). This twist enters the line reversal only through
a unitary determinant factor at the line argument.

The first seesaw splits the old edge matrix into the columns for
\(V_{j,j}\) and the column for its signed line.

\paragraph{The normalized reversal on the diagonal line.}
\teamanchor{proof.weighted-endpoint-line-reversal}%
Fix \(j<m\). Use the identification of \(Z_{j,j+1}\) with
\(S=V_{j+1,j}\) on both kernels. Write \(N=N_j\) for its rank,
\(s_i\) for its signs, and \(\epsilon=(-1)^{j+1}\). The old and
new lines are \(L_{-\epsilon}\) and \(L_\epsilon\), respectively.
Use \(R=\mathcal O_E/\mathfrak m_\vartheta\) with
\(\mathfrak m_\vartheta=\mathfrak b Q_\vartheta\), and write
\(\lambda\) for its finite idele generator, equal to one off the
support \(A\) of the modulus. This support is disjoint from its
conjugate. Put \(u=\overline{\lambda\alpha_{\rm W}}/
(\lambda\alpha_{\rm W})\), \(k=(1,u^j)\) and \(b,a'\) equal to the
horizontal and vertical edge characters on \(S\). Denote the
modulus kernel by \(K^{U,V}_{a,b}[\Psi,P](g,h)\), with character
multiplier \(a(\det h)b(\det g)^{-1}\) on the underlying Weil
action on \(\operatorname{Hom}(U,V)\). There are a primitive
additive character \(\psi\) of \(R\) and an \(f_j=f_j(\vartheta)\ne0\)
such that, for each \(g\in U(S)(\mathbb A)\), column \(v\in R^N\)
and polynomial \(P\),
\begin{align}
 K^{L_{-\epsilon},S}_{1,b}[\mathbf 1_v,P](k,g)
 &=
 f_j\sum_{y\in R^{1\times N}}\psi(\operatorname{tr}(yv))\,
 K^{S,L_\epsilon}_{a',\upsilon}[\mathbf 1_y,P^\vee]
       (g,(1,u^{j+1})),
 & |f_j|^2&=|R|^{-N}.
 \label{eq:weighted-endpoint-line-reversal}
\end{align}
Here \(\upsilon\) is the restriction of the fixed splitting
character to the norm-one ideles. The map \(P\mapsto P^\vee\) is
complex linear. In the \(i\)-th column variable it substitutes
\(-s_i\epsilon\) times the corresponding row variable, at each
infinite place. Thus the substitution is independent of
\(\vartheta\). This is a complete-kernel identity. The
\hyperlink{proof.weighted-grid-kernel-reversal}{conjugate reversal of a chain edge}
uses coefficient conjugation and a highest-line comparison on
realized form spaces.

Let \(c_{\rm neg}\) be the fixed integral finite unit of norm
\(-1\) in the standard signed line, and put \(c_0=c_{\rm neg}^\epsilon\). The old
line point is \(x_-=(1,c_0)x_+\), where \(x_+\) is the standard
point of \(L_\epsilon\). Set \(\ell_0=(1,\bar\alpha_{\rm W}/
\alpha_{\rm W})\) and
\(\xi=(1,-c_0\alpha_{\rm W}/\bar\alpha_{\rm W})\). Then
\(x_+\ell_0=\bar\xi x_-\). In pointed coordinates let \(D_S\) and
\(D_-\) be the standard Hermitian matrices of \(S\) and
\(L_{-\epsilon}\). They give the prescribed signs at infinity
and are identity matrices at finite places. The adjoint
identification from new row coordinates to old column coordinates is
\[
       Y\longmapsto \xi D_S Y^* D_- .
\]
It identifies rational vectors for the new line pointed at
\(x_+\ell_0\) with the old rational vectors. Indeed, for an
old rational matrix \(M\) the adjoint \(J_-^{-1}M^*J_S\) is
rational, where \(J_-,J_S\) are the rational Hermitian matrices.
If \(X=x_S^{-1}M x_-\) is its old column coordinate, the new row
coordinate of the adjoint is
\((x_+\ell_0)^{-1}x_-D_-X^*D_S\). The displayed identification sends
this row to \(X\).

In the coordinates scaled by \(\alpha_{\rm W}\) on the finite
model, this map is \(Y\mapsto-c_0Y^*\). The self-dual additive
character there is
\(\psi_0(z)=\psi_f(\jmath\alpha_{\rm W}\bar\alpha_{\rm W}z)\),
where \(\bar\jmath=-\jmath\) is the fixed Weil datum. Thus
\(\psi_0(-\bar z)=\psi_0(z)\). The identity
\(\bar c_0c_0=-1\) gives preservation of the finite Heisenberg
bicharacter, including its factor \(1/2\), and of the twisted
integral-box vacuum. At infinity the maps are
\(z\mapsto r\bar z\) on coordinates of old sign
\(r=-s_i\epsilon\) and new sign \(-r\). They preserve the
Gaussian and transport the bicharacter with sign \(-r\) to that
with sign \(r\). With \(\zeta_r(z)=\bar z\) for \(r=1\) and
\(\zeta_r(z)=z\) for \(r=-1\), one has
\(\zeta_r(r\bar z)=r\zeta_{-r}(z)\). This gives precisely
\(P^\vee\) in the same Fock normalization.

For the action \(X\mapsto gXh^{-1}\) of the old pair, put
\(T=g\otimes h^{-{\mathsf T}}\), where \(h\) lies in the old
line group. The action of the exchanged pair in the new
coordinates has matrix \(T^{-{\mathsf T}}\) after the
coordinate permutation. The explicit Weil action in the
vacuum model is
\[
     \omega(T)F=\Delta(T)\int F(X)\rho(TX)e_0\,dX,
\]
where \(\rho\) is the Heisenberg action and \(e_0\) the vacuum.
The unweighted integrals intertwine by the change of
coordinates, which preserves the vacuum Haar measure. To
compare the multipliers, transposition preserves the finite
elementary divisors and the real principal minors. In the real
factor \(\Delta_r(T)=\overline{\det T_{--}}\), Jacobi's
complementary-minor identity gives
\(\overline{\Delta_r(T)}=\det(T)\Delta_{-r}(T)\).
At the finite places \(\Delta_f(T)=\chi_f(\det D)\), where
\(D\widehat{\mathcal O}_E^N=
\widehat{\mathcal O}_E^N\cap T^{-1}\widehat{\mathcal O}_E^N\)
and \(\chi_f=\Upsilon_f|\cdot|_f^{-1/2}\), with \(\Upsilon_f\)
the finite part of the full idele splitting character.
One may use \(TD\) for the corresponding lattice of \(T^{-1}\).
As \(|\det T|_f=1\), it follows that
\(\Delta_f(T^{-1})=\Upsilon_f(\det T)\Delta_f(T)\).
Together with \(\Delta_r(T^{-1})=\overline{\Delta_r(T)}\),
these relations give
\[
      \frac{\Delta_{\rm old}(T)}
                {\Delta_{\rm new}(T^{-{\mathsf T}})}
           =\upsilon(\det T)^{-1}
           =\upsilon(\det h)^N\upsilon(\det g)^{-1}.
\]
Consequently the old characters \(1,b\) become
\(b_*=b^{-1}\upsilon^N,\upsilon\) on the exchanged kernel.

The theta functional is invariant under rational Heisenberg
operators. In a rational Schr\"odinger polarization it is unique
up to scalar. Rational modulations force an invariant
distribution to be supported on the rational lattice, since
that lattice is its own annihilator. At a lattice point the
modulation factors minus one give local coordinates at infinity,
so the distribution has no derivative terms there. Rational
translations make all point masses have the same coefficient.
The adjoint preserves the vacuum and gives an isomorphism of
the rational quotients. Both theta functionals have mean square
one on translates of their vacua in Haar probability measure
on these quotients. Their proportionality factor has modulus
one. Moving the new line point from \(x_+\ell_0\) to \(x_+\)
puts \(\ell_0\) into the line argument and contributes the
unitary character \(b_*(\det\ell_0)\) in the identity for the
new kernel in terms of the old. This follows by multiplying the
underlying Weil operators and keeping the explicit determinant
character on the moved point.
\footnote{For the rational summation model and the ad\'elic
oscillator action see Weil, \emph{Sur certains groupes
d'op\'erateurs unitaires}, Acta Math.\ 111 (1964), \S\S18--19
and \S41. The conventions here are
\weightedpin{grh/src/Lib/Weil/Construction/AdelicRepresentation.lean}
{148-L166}{The Weil multiplier and vacuum measure},
\weightedpin{grh/src/Lib/Weil/Construction/CanonicalTheta.lean}
{54-L80}{The phase for mean-square-one theta functionals}, and
\weightedpin{grh/src/Lib/Weil/DualPair/FiniteAction.lean}
{496-L521}{Changing the point without changing data}.}

To compute the factor with the prescribed finite datum, start
with the Heisenberg translate of the old vacuum by
\(-\lambda^{-1}\tilde v\) at \(A\), in the scaled finite
coordinates, for an integral lift \(\tilde v\) of \(v\). The
split Schr\"odinger function of this translate is
\(\mathbf 1_{\mathcal O_A^N}(X-\lambda^{-1}\tilde v)\).
Acting by \(d=(1,\bar\lambda/\lambda)\) in the old line changes
its argument to \(\mathbf 1_{\mathcal O_A^N}
(\lambda^{-1}(X-\tilde v))\), the specified residue indicator.
On the split integral model the action on functions is
pullback by the inverse matrix times the square root of its
absolute determinant, with the unitary splitting and edge
characters. Here the inverse matrix on the position
coordinates is \(d I_N=\lambda^{-1}I_N\) at \(A\), so the
scalar \(C\) of this action satisfies
\[
      |C|^2=\prod_{w\in A}|\lambda_w|_w^{-N}=|R|^N .
\]
This computation uses the characteristic function of the
residue class on the original vacuum.

On the new side the transported translation is supported on
\(\bar A\), with coordinate
\(\overline{\lambda^{-1}(-\bar c_0)\tilde v^{\mathsf T}}\)
there. In the split Schr\"odinger model on \(A\), translation
at \(\bar A\) multiplies by the character obtained by
pairing with the new position variable. Thus the new named
function is \(\psi(\operatorname{tr}(yv))\) on the integral box,
with \(\psi(t)=\psi_0(\lambda^{-1}(-\bar c_0)\tilde t)\) on
components in \(A\). The self-dual lattice pairing and the
unit \(-\bar c_0\) make \(\psi\) primitive, including at the
prime powers of \(\mathfrak b\). This verifies the sign and
the orientation of the finite Fourier row. Since
\(\ell_0 d=(1,u)\), the new line point is exactly
\(k\ell_0 d=(1,u^{j+1})\). The adjusted kernel in terms of the
old kernel at \(k\) has multiplier \(C\) times the theta
and point phases. On inversion it has squared modulus
\(|R|^{-N}\). Finally changing \(b_*\) to \(a'\) on the first
character slot of the new kernel places the factor
\((b_*/a')(\det(1,u^{j+1}))\) on the right of
\eqref{eq:weighted-endpoint-line-reversal}. Its modulus is one
and it is independent of \(g,v,P\). Expand the character
function on the residue box in the row indicators to obtain
the unnormalized sum in \eqref{eq:weighted-endpoint-line-reversal}.
\footnote{\weightedpin{grh/src/Lib/Weil/MixedModel/SchwartzBruhat.lean}
{131-L143}{Split translation and modulation},
\weightedpin{grh/src/Lib/Weil/DualPair/FiniteAction.lean}
{196-L228}{The split change of variables},
\weightedpin{grh/src/Lib/Weil/DualPair/Flip.lean}
{459-L531}{Point and modulus conventions},
\weightedpin{grh/src/New/Determinant/AuxiliaryForms/DiagonalTurns.lean}
{410-L443}{Application to the same signed grid lines}.}

The signed real adjoint, including the fixed change of
polarization, acts continuously on Schwartz space. Hence the
identity extends by density at every fixed \(g\) and finite
column. It acts on the line factor of a real test with
identity on all other factors. The finite Fourier row is
unchanged. With \(C_j\) the normalized complementary scalar
for the ordinary diagonal \(V\)-square, the combined diagonal
denominator is \(C_j/f_j\). Its squared inverse-line factor is
\(|R|^{N_j}\). Relative to the trivial twist this factor
contributes \(\mathrm N(Q_\vartheta)^{N_j}\), a unit at the
coefficient place selected for that twist. The determinant
uses the squared modulus, in which the unitary phases cancel.
\footnote{\weightedpin{grh/src/New/Determinant/AuxiliaryForms/DiagonalTurns.lean}
{493-L524}{The common period test},
\weightedpin{grh/src/New/Determinant/AuxiliaryForms/DiagonalTurns.lean}
{592-L620}{The diagonal denominator}, and
\weightedpin{grh-compact/src/New/Exchange/CrudeRows.lean}
{52-L65}{The same factor in the stage product}.}

\paragraph{The diagonal square and stage product.}
Use the tensor product of the Weil kernels in the first seesaw.
It separates the \(Z\)-edge into the ordinary \(V\)-edge and the
old line kernel on the common exterior group \(S\), using the same
product measures. At the two exterior arguments, reverse the line
by \eqref{eq:weighted-endpoint-line-reversal} and apply
\(T_\square\) to the ordinary \(V\)-square. This square moves from
\(V_{j,j}\) to \(V_{j+1,j+1}\), of rank \(2t\) greater. The line
kernel is a spectator for the square, since it depends on the
exterior \(S\)-argument and the indicated pointed line argument.
The marked old \(Z\)-edge has a left inverse, so its restriction
to the non-line columns does also. The partner graph therefore
preserves the needed primitive block. In the second seesaw tensor
the new row for \(L_{j+1}\) with the new \(V\)-kernel. For pure
tests this is the direct sum identity of the kernel and for
general tests it is the corresponding continuous Schwartz
tensor identification.
\footnote{\weightedpin{grh-compact/src/New/Exchange/RowLadder.lean}{17-L62}{Off-diagonal rank and character check},
\weightedpin{grh-compact/src/New/Exchange/GridMove.lean}{238-L293}{Partner graph and diagonal block covariance},
\weightedpin{grh/src/New/Determinant/AuxiliaryForms/DiagonalTurns.lean}{231-L338}{Both factorizations and the actual line characters},
\weightedpin{grh/src/New/Determinant/AuxiliaryForms/DiagonalTurns.lean}{406-L620}{Line Fourier transform and the two seesaws}.}

In the split presentation the line factor multiplies the
exterior prefix on \(S\). The theta line kernel with its point
fixed is of moderate growth, with the same Schwartz seminorm
bounds for its derivatives. Its product with the rapid prefix
is rapid. The intermediate integral is still bounded by the
primitive-support estimate for the ordinary square. Thus the
two exterior integrations and the square can be grouped as
above, even with the line factor. All real operations in this
move are on the kernel before pairing with the exterior forms.

The denominator of the combined move is the product of the
\(2j\) off-diagonal square scalars and the normalized diagonal
denominator \(C_j/f_j\) of
\hyperlink{proof.weighted-endpoint-line-reversal}{the line comparison}.
The tidy order goes up by \(4j+2\). At every stage the real
maps act on the same tensor for every finite tuple and both
input forms. The finite graph and Fourier coefficients are
those of the normalized move and are independent of the real
translate. Multiplication yields one real tensor
\(\mathcal V_m(a_\infty)\) on the two last paths, one last finite
datum of order \(4m+2m^2\), and one nowhere-zero denominator
\(c'_\vartheta\) costing at most \(2m^2\) squares. The actual number
of Hecke products in this construction is
\(\sum_{j<m}(2j+1)=m^2\). Keep the larger budget in the endpoint
estimate. All modulus powers in these relative scalars are powers
of \(\operatorname{N}Q_\vartheta\) times fixed constants, hence they
have absolute value one after division by the untwisted value and
evaluation away from the conductor.

\paragraph{The last path and the single multiplier on the completion.}
\teamanchor{proof.weighted-endpoint-last-path}%
The final paths on both grids run from \((0,m)\) to \((m,m)\)
by increasing the first index. Every edge is a lift from the
smaller to the larger definite space. In the coordinates
\((a,b)\) the two lists of nodes are \(V_{j,m}\) and \(Z_{j,m}\).
The compact coefficient models for these oriented edges use
the negatives of the grid parameters.

For the terminal pairing, compare the doubled \(V\)-row
parameter with the negative of the doubled \(Z\)-row parameter.
Let \(\mathcal B_d(c)\) denote the length-\(d\) string
\(c+d-1,c+d-3,\ldots,c-d+1\). At each real place the entries are
\[
 \begin{array}{c|c|c}
 & V & \text{reflected }Z\\ \hline
 \text{above the center}
   &\mathcal B_t(2w_l)
   &\mathcal B_{t-1}(2w_l),\quad 2w_i-h-1\\
 \text{central stem}
   &\mathcal B_{\kappa-1}(0)&\mathcal B_{\kappa-1}(-1)\\
 \text{below the center}
   &\mathcal B_t(-2w_l),\quad h-2w_l
   &\mathcal B_{t+1}(-2w_l)
 \end{array}
\]
Here \(1\le l\le m\) and \(1\le i\le m+1\). The two blocks at
\(2w_l\) sit between \(2w_{l+1}-h-1\) and \(2w_l-h-1\)
and alternate with a \(V\)-entry at each end. The central
reflected entries lie one below the respective \(V\)-entries.
The last central reflected entry is above \(h-2w_1\).
After \(h-2w_l\) the two blocks at \(-2w_l\) alternate with a
reflected entry at each end. These separations use
\(w_{l+1}-w_l=h\), \(h>t+\kappa+2\) and \(2w_1-h>\kappa\).
In decreasing order the two full lists therefore strictly
alternate, beginning and ending with the reflected list.

\paragraph{An integral compact terminal pairing.}
\teamanchor{proof.weighted-endpoint-compact-branching}%
Use the two definite spaces \(V_{m,m}\) and \(Z_{m,m}\) with
their fixed block embedding. At a CM-type place let \(n\) be the
rank on the \(V\)-side. Write \(\lambda=(\lambda_1,\ldots,
\lambda_{n+1})\) for the highest weight whose doubled parameter
is the negative of the terminal \(Z\) row parameter, and
\(\nu=(\nu_1,\ldots,\nu_n)\) for the highest weight with the
terminal \(V\) row parameter. The two actual algebraic
weights are \(\lambda\) and \(-w_0\nu=(-\nu_n,\ldots,-\nu_1)\).
For descending rank \(d\) entries the \(i\)-th doubled parameter
is \(2\mu_i+d+1-2i\). The two successive doubled gaps are thus
\(2(\lambda_i-\nu_i)+1\) and \(2(\nu_i-\lambda_{i+1})+1\).
Both are positive, so \(\lambda_i\ge\nu_i\ge\lambda_{i+1}\).

We use the reader's Weyl lattice and Gram form. For an
arbitrary highest weight \(\mu\) of rank \(d\), keep the
exponent \(\mu_d\) of the determinant separately and put
\(\kappa_h=\mu_h-\mu_{h+1}\) for \(1\le h<d\). With \(X_i\)
the columns of \(X\), there is an integral coordinate map
\[
 I_\mu\colon {\rm ev}_X\longmapsto
   \bigotimes_{h=1}^{d-1}
       (X_1\wedge\cdots\wedge X_h)^{\otimes\kappa_h}.
\]
Here \({\rm ev}_X\) lists the standard flag-minor products of
this shape at \(X\). Each coordinate of the displayed tensor
is an integer linear combination of those products by
integral flag-minor straightening. Applying Cauchy--Binet to
each factor makes the pairing of two image evaluation vectors
the product of the same principal minors of \(X^*Y\) that
defines their Gram pairing. The evaluation vectors span
because the standard products are independent polynomials.
Thus \(I_\mu\) is an isometry for the Weyl norm, with the
ordinary exterior tensor norm on its target. Its action on
the representation includes the single external factor
\(\det(g)^{\mu_d}\).
For this integral straightening, Theorem~II.2.16 of Akin,
Buchsbaum and Weyman says that for a free module with an
ordered basis over a commutative ring the standard tableaux
of a given shape give a free basis of the Schur module.
Use the rank-\(d\) standard basis over \(\mathbf Z\) and
transpose the diagram from their row convention to the
column heights above. This is applied to the polynomial shape
after separating the determinant power.
\footnote{Akin, Buchsbaum and Weyman,
\href{https://jerzyweyman.com/wp-content/uploads/2022/11/Publ1-Schur-Functors-and-Schur-Complexes.pdf}
{\emph{Schur Functors and Schur Complexes}},
Adv.\ Math.\ 44 (1982), Theorem~II.2.16, p.~233.
Map each exterior factor of height \(h\) to the minor on
the first \(h\) columns of the variable array. The tableau
relations are the flag-minor relations, and the basis images
are the standard flag-minor products. The precise Gram
kernel is \weightedpin{grh-compact/src/Lib/GLn/SchurModule/Gram.lean}
{89-L96}{Weyl evaluation and the minor kernel}.}

In the \(Z\)-side tensor let \(e_L\) be the extra line basis
vector. For each height \(h\) shorten \(\gamma_h=\lambda_h-\nu_h\)
columns by taking the \(\bigwedge^{h-1}V\wedge e_L\) component
and dropping \(e_L\). Keep the remaining
\(\kappa'_h=\nu_h-\lambda_{h+1}\) columns by projecting onto
\(\bigwedge^h V\). These operations are integral contractions.
After the tensor factors are reordered and a full-height
factor \(\bigwedge^n V\) is written in its volume basis, the
resulting map \(D\) has the tensor target of \(I_\nu\).
The number of factors of height \(h<n\) is
\(\kappa'_h+\gamma_{h+1}=\nu_h-\nu_{h+1}\). For
\(\operatorname{diag}(g,t)\) the full-height factors give the
determinant exponent
\(\lambda_{n+1}+\kappa'_n=\nu_n\); the added line has exponent
\[
       s_\tau=\lambda_{n+1}+\sum_{h=1}^n\gamma_h
              =\sum_{i=1}^{n+1}\lambda_i-\sum_{i=1}^{n}\nu_i .
\]
On the \(V\)-side use the signed complementary wedge map
\(C_h\colon\bigwedge^h V\to\bigwedge^{n-h} V\). Its matrix in
the exterior bases is a signed permutation and
\(C_h\bigwedge^h g=\det(g)(\bigwedge^{n-h}g^{-{\mathsf T}})C_h\).
Applied on every factor of \(I_{-w_0\nu}\), with the factors
reordered, this gives \(C\) into the tensor target of
\(I_\nu\). The \(V\)-side determinant exponent after this
operation is \(-\nu_1+(\nu_1-\nu_n)=-\nu_n\).

Take the ordinary bilinear tensor contraction of
\(CI_{-w_0\nu}v_-\) and \(DI_\lambda v_+\). Its matrix is
integral, and its character under the simultaneous block
action is exactly \(t^{s_\tau}\), since the tensor actions
\(g^{-{\mathsf T}}\) and \(g\) cancel. This identity holds in
characteristic zero and hence in the coefficient readings
used in the transported integral gauges. The contraction
is nonzero. On the \(V\)-side evaluate at the reversed
basis flag, whose complementary factors give the standard
highest tensor of \(I_\nu\) up to sign. On the \(Z\)-side
evaluate at \(X_i=e_i+e_L\) for \(1\le i\le n\) and
\(X_{n+1}=e_L\). Each retained factor has standard highest
wedge coefficient one. Each shortened factor has coefficient
one at \(e_1\wedge\cdots\wedge e_{h-1}\), coming from the
line entry of its last column. The contraction of these two
tensors is therefore a sign. Tensor over the CM-type places,
using the specified reindexing of the block. This gives a
single integer matrix \(B_{\rm end}\) with
\[
 \|B_{\rm end}\|_{\rm bil}\le1,\qquad
 \sum_\tau |s_\tau|\le C_{\rm W}+f\lceil20 b\rceil .
\]
The first inequality is in the two unscaled Weyl Gram norms,
by the isometries above and the norm-one exterior projections.
The second follows from the bounds for the actual
\(Z\)-side diagram and determinant exponents. These are the
norm and character needed at \(h\perp u^m\).

Take any real coefficient covectors \(L_-,L_+\) in the
terminal compact models, with the translate already in the
real datum. At the finite points \(h\) and \(h\perp u^m\),
write their algebraic values as \(v_-,v_+\). The compact
integral is
\[
 \int L_-(\rho_-(k^{-1})v_-)\,
        L_+(\rho_+((k\perp1)^{-1})v_+)\,dk
      =\beta(L_-,L_+)\, B_{\rm end}(v_-,v_+).
\]
Indeed the invariant space for the common compact block is
one-dimensional. In the Schur character on \(n+1\) letters,
the boxes with letter \(n+1\) form a horizontal strip.
Removing it gives exactly one branching term for each
interlacing shape on \(n\) letters, with the remaining
tableau arbitrary. A common determinant shift treats
arbitrary integral highest weights and gives the indicated
added-line character \(t^{s_\tau}\). Pairing with
the weight \(-w_0\nu\) takes the term of weight \(\nu\)
with multiplicity one at every real place. Transport to the
conjugate Weyl models takes the invariant contraction to
the same \(B_{\rm end}\). Thus \(\beta\) is a fixed bilinear
functional of the real covectors alone.
The original finite measure is \(I_{0,m,m}^{-1}\) times the
measure giving the grid terminal level mass one. Apply this
rescaling to the tuple-averaged integrand, whose right
invariance and free classes give the finite class sum. This
is the factor \(I_{0,m,m}^{-1}\) in the last-path cancellation.
A fixed nonzero scalar multiple of \(B_{\rm end}\) changes
\(\beta(L_-,L_+)\) by the inverse scale at every twist.
\footnote{\weightedpin{grh-compact/src/New/Exchange/GridMove.lean}
{195-L235}{Actual terminal Weyl weights and block point},
\weightedpin{grh-compact/src/Lib/Weil/CompactLift/BranchingPeriod.lean}
{148-L172}{Compact and finite measures at the branching step},
and \weightedpin{grh-compact/src/New/Exchange/GridFirstColumns.lean}
{108-L134}{Terminal index cancellation}.}

\paragraph{The two definite CM jet comparisons.}
\teamanchor{proof.weighted-endpoint-edge-jets}%
Write \(z=0,1\) for the \(V,Z\) legs. The rank at step \(j\) is
\(N_{z,j}=r+z+j(t+z)+m(t-z)\) for
\(0\le j\le m\). All signs are positive and the adelic frame is
the standard one, including at \(V_{0,m}=Y_m\). At each embedding
in the CM type \(\Phi\), let \(l_j\) be the negative of the grid
parameter at this node. Write \(l_{j,i}\) in decreasing order
when assigning the highest weight
\(\lambda_{j,i}=(l_{j,i}-N_{z,j}-1+2i)/2\).
The finite matrices and the Weyl bases keep their original
indices. The decreasing order specifies the highest flags.

For one edge put \(n=N_{z,j}\), \(n'=N_{z,j+1}\) and
\(\Delta=n'-n=t+z\). Its splitting characters are
\[
 a=\gamma_{z,j,\vartheta}\upsilon^{\lfloor(n+1-z)/2\rfloor},
 \qquad
 b=\gamma_{z,j,\vartheta}\upsilon^{\lfloor(n'+1-z)/2\rfloor},
 \qquad
 \gamma_{z,j,\vartheta}=
 \begin{cases}
  \mu_{j+1}^0 & z=0,\\
  \mu_{j+1}^0\vartheta^{j+1} & z=1 .
 \end{cases}
\]
Here \(\mu_{j+1}^0\) has real weight \(-w_{j+1}\) and tame part
\(\omega^{j+1}\), and \(\upsilon\) is restricted to norm-one
ideles. Write \(a_\tau,b_\tau\) for the weights of \(a,b\) at
\(\tau\). They are independent of \(\vartheta\).

The new positive block lies above the old grid entries.
The preceding positive centres are at most \(2w_j\), separated
from \(2w_{j+1}\) by at least \(2h\). The small central base
entries have modulus at most \(\kappa\), the remaining base
entries \(2(h_2-w_i)\) are negative, and the negative blocks
are centred below zero. The inequalities for \(w_1,h,t,\kappa\)
put these entries below the first positive block. After negating
the parameter the new block is thus the last \(\Delta\) entries.
In the decreasing orders,
\[
 l_{j+1,i}-l_{j,i}=\Delta+2(a_\tau-b_\tau)\quad(i\le n),\qquad
 l_{j+1,n+k}=2a_\tau+\Delta-n+1-2k\quad(1\le k\le\Delta).
\]
For \(z=0\), \(m\) is even and \(t\) is odd. The old entries
shift by \(+1\) for even \(j\) and \(-1\) for odd \(j\).
The parity \(n\equiv j+1\bmod 2\) gives this shift from the
ceilings in \(a_\tau,b_\tau\). The new negated block has
centre \(-2w_{j+1}+(n\bmod 2)\). For \(z=1\), the old
entries are unchanged, \(n,n'\) are even, and the new
negated centre is \(-2w_{j+1}\). The highest weights therefore
satisfy \(\lambda_{j+1,i}=\lambda_{j,i}+a_\tau-b_\tau\) for
\(i\le n\) and \(\lambda_{j+1,i}=a_\tau\) for \(i>n\).

Set \(e_{z,j,\tau}=\lambda_{j,n}\). The diagram
\(\kappa_{z,\tau}\) has column multiplicity
\(\lambda_{j,i}-\lambda_{j,i+1}\) at \(i<n\), multiplicity
\(e_{z,j,\tau}-b_\tau\) at \(n\), and zero at \(n<i<n'\).
The multiplicity at \(n\) is nonnegative by the decrease
of the target weight. Successive arrows thus use truncations
of one diagram, with the determinant exponents \(e_{z,j,\tau}\)
read separately at their ranks. In this rectangle the polynomial
degree is \(D_{z,j}=\sum_{\tau}\sum_{h<n'}h\kappa_{z,\tau}(h)\).
It includes the columns of full source height \(n\).
The bound on the grid parameters gives row lengths \(O(b)\),
so \(D_{z,j}=O(fn_*b)\), with \(f=|\Phi|\) and \(n_*=2n_m\).
Since \(f=O(x^{11}\log x)\), \(n_*=O(x^{11/2})\),
\(r\asymp x^5\) and \(b_1\ge b\), this is at most
\(C_{\rm W}=b_1r^4\) for large \(x\).
\footnote{\weightedpin{grh-compact/src/New/Exchange/GridParameters.lean}
{30-L72}{Grid signs, parameters, character weights and ranks},
\weightedpin{grh-compact/src/New/Seesaw/PairTensorData.lean}
{68-L104}{The spaces and characters}, and
\weightedpin{grh-compact/src/New/Exchange/GridMove.lean}
{154-L204}{The diagrams and end kernels}.}

On the transported tuple \(S\) use its component
\(S_e=S((z,j,m),(z,j+1,m))\). Put
\(R_\vartheta=\mathcal O_E/\mathfrak m_\vartheta\), where
\(\mathfrak m_\vartheta=\mathfrak b Q_\vartheta\). The raw
finite datum is \(F_e=\operatorname{name}_{\mathfrak m_\vartheta}
(\mathbf 1_{S_e})v_0\) on \(\operatorname{Hom}(E^n,E^{n'})\).
Naming means multiplication by that residue indicator in the
specified split Schr\"odinger coordinates, with the vacuum
elsewhere. The analytic arrow uses \(F_e\), and the algebraic
arrow uses \(F_e^\circ=v^{nn'}F_e\). Here \(v=c_s\ne0\) is the
\hyperlink{proof.weighted-vacuum-constant}{rank-one vacuum multiplier}
of the same datum. The Hom space has rank \(nn'\) over \(E\),
so this is the product normalization \(c_s^{nn'}v_0\) on
its vacuum. Use the
\hyperlink{proof.weighted-cm-jet-transport}{CM Taylor scale}
\(c=\mathrm{per}=(s_0/u_0)\eta/\Omega_0\ne0\).
At the actual finite source and target points \(g,h\) put
\[
 J_e(g,h)=\mathcal J_c\bigl(\omega_{a,b,f}(g,h)F_e^\circ\bigr),
 \qquad
 (\mathcal J_c(F))_{i',i}=
     S_{\kappa_z}(\partial)_{i',i}\,T_c(F)(0).
\]
Here the polynomial entries of \(S_{\kappa_z}\) in the Weyl
bases are integral. The finite action is
\(\omega_{a,b,f}(g,h)=a(\det h)b(\det g)^{-1}\omega_f(G(g,h))\),
where \(G(g,h)\) acts by \(X\mapsto hXg^{-1}\).
The Taylor coefficients and factorial differential pairing
are those of
\hyperlink{proof.weighted-fock-jet-coefficient}{the definite Fock calculation}.

For the compact conjugate Weyl actions \(\rho,\rho'\) set
\(P_{\ell,\xi}(Z)=\ell(S_{\kappa_z}(Z)\xi)\).
The real action substitutes \(y_\tau^{\mathsf T}Z_\tau
x_\tau^{-{\mathsf T}}\) and contributes
\(\det(y_\tau)^{a_\tau}\det(x_\tau)^{-b_\tau}\).
Right substitution by \(x_\tau^{-{\mathsf T}}\) in the
rectangular Schur matrix contributes
\(\kappa_{z,\tau}(n)\) powers of \(\det(x_\tau)^{-1}\).
Thus the transformed coefficient is
\(P_{\ell\circ\rho'(y^{-1}),\,\rho(x)\xi}\). Homogeneity
and the factorial pairing give, for every finite \(F\),
\[
 c^{D_{z,j}}\,\Theta\bigl(\omega_{a,b}((x,g),(y,h))
                    (\mathsf v(P_{\ell,\xi})\otimes F)\bigr)
 =\ell\bigl(\rho'(y^{-1})\mathcal J_c(\omega_{a,b,f}(g,h)F)
                       \rho(x)\xi\bigr).
\]
Adelic invariance of \(\Theta\) then gives the rational
covariance \(J_e(\gamma g,\delta h)\rho(\gamma)
=\rho'(\delta)J_e(g,h)\) by testing every \(\ell,\xi\).

Take the two highest evaluation vectors \(u,u'\) on the
decreasing flags and put
\(\ell'=\langle u',\,\cdot\,\rangle_{\mathrm W}\).
The Gram minors of a permutation flag are one, so the
highest vectors have norm one. The shifted source entries
\(l_{j,i}-2b_\tau+n'\) equal \(l_{j+1,i}-2a_\tau+n\)
and exceed \(\Delta\), since \(\lambda_{j+1,i}\ge a_\tau\).
The new shifted string is \(\Delta+1-2k\). This is the
positive definite harmonic match. At \(\tau\), the product
\(H_\tau\) of leading minors in the ordered rectangle,
with the specified column multiplicities, is joint highest.
The opposite-sign Laplacians are vacuous. The lower row
and upper column unipotents preserve the minors, and the
diagonal characters are the row and column lengths with
the splitting weights. These triangular groups have an
open orbit, obtained by elimination when the successive
leading minors are nonzero and by diagonal rescaling.
Hence semi-invariants of the same character on this orbit
are proportional. The prescribed joint harmonic is the
product of \(u_\tau H_\tau\), with \(u_\tau\ne0\) fixed.

If \(P_\tau\) is the permutation matrix of the ordered
source flag, the rectangular Schur action takes the
evaluation at \(P_\tau\) to that at \(Z_\tau P_\tau\) after
multiplication by \(\det(P_\tau)^{\kappa_{z,\tau}(n)}\).
The target Gram covector extracts the leading minors in
its ordered rows. Thus the actual harmonic is
\(\eta_{z,j}P_{\ell',u}\),\footnote{\weightedpin{grh-compact/src/Lib/Weil/CompactLift/BranchingPeriod.lean}
{90-L139}{Rectangular covariance and homogeneity},
\weightedpin{grh-compact/src/Lib/Weil/CompactLift/RallisFormula.lean}
{42-L87}{The fixed harmonic and its highest coefficient}, and
\weightedpin{grh/src/Lib/Weil/CompactLift/DefiniteThetaLift.lean}
{127-L239}{The ordered minors and the permutation determinant}.} where
\(\eta_{z,j}=\prod_\tau
 (u_\tau\det(P_\tau)^{\kappa_{z,\tau}(n)})\).

Write \(I_{z,j,m}=[K^0_{z,j}:K_{z,j}]\) for the source
index. In the real source variable the algebraic form has
coefficient \(\langle u,\rho(x^{-1})\phi(g)\rangle_{\mathrm W}\).
Its product with the kernel averages the conjugates of
\(u\langle u,\,\cdot\,\rangle_{\mathrm W}\) over compact
probability Haar. The average has trace one and commutes
with the irreducible action, so it is the identity divided
by the source Weyl dimension \(d_{z,j}\).
After this average, the finite integrand is proportional to
\(J_e(g,h)\phi(g)\). The old measure gives \(K^0_{z,j}\)
mass one, while the algebraic class mean gives \(K_{z,j}\)
mass one. The ratio is \(I_{z,j,m}^{-1}\) on any common
refinement fixing the integrand at this tuple. On that
refinement both measures divide by the same rational
stabilizer orders. Thus the analytic lift at \((y,h)\) equals
\[
 \frac{\beta_{z,j}}{I_{z,j,m}\,v^{nn'}}\,
       \ell'\!\left(\rho'(y^{-1})(\mathscr L_e\phi)(h)\right),
 \qquad
 \beta_{z,j}=\frac{\eta_{z,j}}{d_{z,j}\,c^{D_{z,j}}},
\]
where \(\mathscr L_e\) integrates \(J_e(g,h)\phi(g)\) in the
source measure with \(K_{z,j}\) of mass one.
For raw finite datum on both sides the multiplier is
\(\beta_{z,j}/I_{z,j,m}\). In either case \(\beta_{z,j}\)
is fixed by the real coefficient models and the CM scale,
uniformly for all twists, tuples and class arguments.
\footnote{\weightedpin{grh-compact/src/New/Exchange/GridMove.lean}
{183-L231}{The levels, CM scale and normalized datum},
\weightedpin{grh-compact/src/Lib/Adelic/Hecke/ClassMean.lean}
{25-L37}{The refined class integral}, and
\weightedpin{grh/src/Lib/Adelic/AutomorphicForms/AutomorphicAsAlgebraic.lean}
{297-L367}{The finite index and compact matrix coefficient}.}

\paragraph{Compact averaging of the transported tensor.}
\teamanchor{proof.weighted-endpoint-definite-branching}%
On the final paths the input on each leg has its prescribed
compact type, with its fixed coefficient covector. For a
real place \(\tau\), the polynomial Fock space of an edge
is the rectangular polynomial space with the determinant
characters of the action above. Its Cauchy decomposition
has one pair of Schur modules for each partition with at
most \(n\) rows. This follows by applying Schur--Weyl
decomposition to the two tensor powers and taking the
diagonal symmetric-group invariants. With the compact-action
conventions of \(P_{\ell,\xi}\) above, the entering weight
forces the partition
\[
    (\lambda_{j,1}-b_\tau,\ldots,\lambda_{j,n}-b_\tau).
\]
Its last entry is nonnegative by the target weight formula.
The outgoing coefficient weight is the same partition padded
with \(\Delta\) zeros and shifted by \(a_\tau\). It is exactly
\(\lambda_{j+1}\). Thus the entering type forces the outgoing
type at every step, with no choice of intermediate compact
type. This component of the Schwartz oscillator is the span
of the vectors \(\mathsf v(P_{\ell,\xi})\) after taking all
real places. Indeed the compact character projector is
continuous on Schwartz space and the polynomial Fock
vectors are dense there. The projected polynomial component
is finite-dimensional.

For the projected edge vector \(\mathsf v(P_{\ell,\xi})\),
let \(L\) be the covector of the entering form, with algebraic
value \(\phi(g)\). Schur averaging in the real source
variable integrates
\[
 \rho(x)\xi\,L\circ\rho(x^{-1})
     \quad\hbox{to}\quad \frac{L(\xi)}{d_{z,j}}\,\mathrm{id}.
\]
The resulting automorphic coefficient at \((y,h)\) is
\[
  \frac{L(\xi)}
       {I_{z,j,m}\,v^{nn'}\,d_{z,j}\,c^{D_{z,j}}}
       \ \ell\bigl(\rho'(y^{-1})(\mathscr L_e\phi)(h)\bigr).
\]
The arithmetic operator \(\mathscr L_e\) is the very same
finite jet integral as for the joint harmonic. The outgoing
real covector is \(\ell\). The finite measure and scale
factors are precisely those of the normalized harmonic
comparison, with the real factor \(L(\xi)/\eta_{z,j}\) relative
to its \(\beta_{z,j}\).

For a pure projected tensor on both legs, write
\(\ell_{z,j}\) for the successive real covectors and
\(\xi_{z,j}\) for the entering vectors of its edge factors.
After retaining the normalized edge constants, the remaining
real coefficient of the period is
\[
 \beta(\ell_{0,m},\ell_{1,m})
             \prod_{z=0}^1\prod_{j=0}^{m-1}
                     \frac{\ell_{z,j}(\xi_{z,j})}{\eta_{z,j}} .
\]
Here \(\beta(\ell_{0,m},\ell_{1,m})\) is the compact
terminal coefficient of \(B_{\rm end}\) just computed.
This is a multilinear contraction on the fixed finite-dimensional
edge components. Composing it with the tensor product of the
compact projections gives one continuous linear functional
\(\beta(\mathcal V)\) on the real Schwartz tests of the last
paths. The period identity extends to arbitrary tensors by
Schwartz density and continuity of the integrals. In particular
for \(\mathcal V_m(a_\infty)\) the functional is the same for
each term of the finite tuple sum, for both input forms, and
for every \(\vartheta\). The finite operator in every term
is the composition of the prescribed jet matrices and the
integer pairing \(B_{\rm end}\), at its indicated block point.

\paragraph{Cancellation on the last path.}
\teamanchor{proof.weighted-endpoint-edge-indices}%
The reduction of \(K^0_{z,j}\) to
\(\operatorname{GL}_{N_{z,j}}(R_\vartheta)\) is surjective
at the chosen split places. At \(V_{0,m}\) the reduction
subgroup defining the level is trivial. At every other node
it is the kernel of reduction modulo \(\mathfrak b\).
Thus every target index and every \(Z\)-index is independent
of \(\vartheta\). At the start of \(V\), \(K_{0,0}\) is
contained in the principal delta level of modulus
\(\mathfrak M Q_\vartheta\), since \(\mathfrak b=\mathfrak M^{b_1}\).
At the start of \(Z\) the level is the core principal level
of \(\mathfrak b\). The two initial forms have these levels.

Write the arrow factor as
\[
  C_{z,j}\frac{I_{z,j+1,m}}{I_{z,j,m}},
  \qquad
  C_{z,j}=\frac{\beta_{z,j}}
        {v^{N_{z,j}N_{z,j+1}} I_{z,j+1,m}(1)} .
\]
Along a leg these factors multiply to
\(\prod_j C_{z,j}\) times \(I_{z,m,m}/I_{z,0,m}\).
The compact terminal pairing at \(h\perp F_\vartheta^m\)
contributes \(I_{0,m,m}^{-1}\) and the contraction
\(\beta(\mathcal V_m(a_\infty))\) accounts for the full real
tensor. Combining this with the initial stage factor
\(I_{0,0,m}\) gives the common multiplier
\[
   \alpha_{a_\infty}
     = C_0\,\beta(\mathcal V_m(a_\infty))\,
          \frac{I_{1,m,m}}{I_{1,0,m}}\prod_{z,j} C_{z,j}.
\]
This proves \eqref{eq:weighted-completion-period} with the
same \(c'_\vartheta\) from the pointwise path moves.
\footnote{\weightedpin{grh-compact/src/New/Seesaw/PairTensorData.lean}
{107-L113}{Reduction subgroups},
\weightedpin{grh-compact/src/New/Seesaw/PairTensorData.lean}
{173-L186}{The determinant character and fixed indices}, and
\weightedpin{grh-compact/src/New/Exchange/GridFirstColumns.lean}
{108-L134}{The stage-period cancellations}.}

\paragraph{Corank-one periods on the weighted datum.}
\teamanchor{proof.weighted-endpoint-separation}%
For each nonzero \(y\in\mathcal H_0\) we need a core form of the
prescribed real type and level and some real translate \(a_\infty\) whose
bordered period with \(y\) is nonzero. We use the
\hyperlink{candidate.fixed-level-period}{corank-one argument of the
earlier window route} at the present weights and depth.

For the row group \(H\) and its fixed-line inclusion \(\iota\) in
the core group, the Petersson pairing with core form \(u\) is
\[
 \int_{H(F)\backslash H(\mathbb A_F)}
       \overline{y(h)}\,u(\iota(h)a_\infty)\,dh .
\]
At one real place write \(\lambda_i\) for the row Harish--Chandra
entries, numbered from zero, and \(\kappa=r+1-m\).
The lower type in He's restriction is the type of \(y\), with entries
\(\lambda_i\). The upper type of \(u\) has entries
\(-\lambda_i-1/2\) at the shared indices. The shared signs
are positive for \(i<\kappa-1\), then alternate from positive.
The added upper entry at index \(r\) has positive sign since
\(m\) is even. In descending order, the upper tail from index \(r\)
to \(\kappa-1\) comes first, with alternating signs and two
positive endpoints. The middle interleaves the positive-sign row
entries \(\kappa/2-1-i\) with their upper reflections and half-unit
shift. It starts with a row entry and ends with an upper entry.
The lower tail then starts positive and alternates. The gaps from
the two tails to the middle follow because both rounded weights
\(a/x\) and \(a/x_1\) dominate every fixed power of \(r\).
Every adjacent pair in this merged list satisfies
He's signed interlacing rule. For discrete series on \(U(p,q)\)
and on its subgroup \(U(p-1,q)\) fixing a positive line,
sort the two labelled Harish--Chandra lists together in descending
order. The discrete multiplicity is one when every adjacent
pair from the same list has opposite signs and every adjacent
pair from different lists has equal signs. Otherwise it is zero.
The alternating tails and the interleaved positive middle give
exactly these two conditions, including the joins. The ranks
are consecutive and the Harish--Chandra lists have opposite
integer parities, so no ties occur. The extra positive entry
is the fixed positive line. This gives the required copy of the
row representation in the restricted core representation.
Let \(v_0\) be the unit highest vector of the upper lowest compact
type whose coefficient defines the real core kernel. Use this
vector as the initial vector in the earlier proof.
\cite[Definition~1 and Theorem~0.1]{depth-he}
\footnote{\weightedpin{grh-compact/src/New/Seesaw/IndexTypes.lean}{124-L126}{Row entries},
\weightedpin{grh-compact/src/New/Exchange/GridParameters.lean}{32-L54}{Signs and upper shift},
and \weightedpin{grh-compact/src/New/Exchange/GridParameters.lean}{81-L87}{Signature of the added line}.}

The upper ambient gaps in the compact-sign block can equal one.
For the transverse second moment of the earlier proof, use its
tensor filtration for \(\pi\otimes(\mathbf C^{r+1})^*\). Let
\(\mu_1>\cdots>\mu_{r+1}\) be the upper entries in their ambient
order. A shift \(\mu-e_i\) is singular only if
\(\mu_i-\mu_{i+1}=1\), and the root \(\alpha=e_i-e_{i+1}\) is compact by the
opposite-sign gap. Cohomologically induce in stages through the
compact rank-two Levi on these adjacent coordinates. In the
inner step the shifted parameter has root pairing zero. The
half-sum of the outer positive roots pairs to zero with
\(\alpha^\vee\), since their set is stable under this simple
reflection. Subtracting the Levi half-sum gives inducing degree
\(0-\langle\rho_{U(2)},\alpha^\vee\rangle=-1\) on its root line.
Thus the Borel--Weil--Bott calculation is that of
\(\mathcal O(-1)\) on \(\mathbf P^1\), whose cohomology vanishes.
The remaining
shifts are regular in the original chamber, with cohomology
concentrated in the same degree. Their central eigenvalues are
distinct. Thus the tensor sum in the earlier proof is over the
regular shifts. The same Capelli minors give their projectors;
the denominators \(\prod_{k\ne i}|\mu_k-\mu_i|\) are at least
one. For each surviving shift the ratios of the changed factors
in the formal degree are at most two, since both the original
and the shifted differences are positive integers in ambient
order. This proves the same transverse moment bound on the
actual upper type.
\cite[Proposition~3.6 and (14)]{real-branch-prudhon}
\cite[Chapters~V and~XI]{real-branch-knapp-vogan}

The opposite-sign gaps are at least \(r^8\) and the absolute
entries are \(O(b)\). These bounds give the hypotheses of the
earlier lowest-type decay estimate, the transversal integration
and the first-variation bound. The orbit-average translate of
the lowest-type vector uses the same reproducing scalar in all
three estimates. Write \(Q=O(fr^4\log b)\) for their common
logarithmic loss over the \(f\) real places. For the finite
part in the old depth calculation use the principal modulus
\(\mathfrak m_1=(\bar\pi_G)^{b_1}\) of the weighted core, with
\(\log\mathrm N\mathfrak m_1=fb_1\log g\). At the row level take
\(A=r^2f g\operatorname{rd}_E\), and put
\(s=(2e^Q A)^{-m}/2\) and \(R_{\mathrm{c}}=s^f\sqrt{\mathrm N\mathfrak m_1}\).
The chosen primes are completely split and their chosen halves
divide \(\mathfrak m_1\). Its intersection with its conjugate
ideal is contained in \((g)\). The norm-one block restriction
of the core level therefore lies in the row principal level.
Indeed either split-half congruence implies the conjugate
congruence by the inverse-adjoint identity on the self-dual
lattice. The bounds \(b_1\asymp b\), \(\log g\asymp\log b\) and
the polynomial bounds on \(f,r,m\) give
\[
  e^Q A(2s)^{1/m}=\tfrac12,\qquad
  \log 2+Q+(6(r+1)+9)f+f\log(1/s)
        <\tfrac12\log\mathrm N\mathfrak m_1
\]
for large \(x\). Thus \(R_{\mathrm{c}}>2e^{Q+(6(r+1)+9)f}\). These
are the requirements for the short-lattice reduction and the
unfolded stabilizer term to dominate the tail. The earlier
count and the earlier proof of absolute convergence apply at
this level. In particular the count uses the rank \(r+1\)
of the same norm-one frame, not the ranks later in the core
column.
\footnote{\hyperlink{proof.principal-short-lattice-reduction}{The earlier short-lattice estimate},
and \weightedpin{grh-compact/src/New/Exchange/GridModuli.lean}{120-L142}{The weighted core modulus}.}

The resulting Poincar\'e form is a value of a smooth cuspidal
multiplicity map at a translate of the fixed lowest-type ambient
vector. The right real span of the row coefficient vector is
dense in its irreducible Hilbert space. The bounded period
functional of the Poincar\'e witness is nonzero on the row
reproducing vector, hence on a right translate of the prescribed
nonzero row coefficient vector. Move that translate to the
ambient argument using right invariance of the row quotient
measure. The ambient vector is still a translate of the fixed
lowest-type vector. It is therefore a reading at some real
\(a_\infty\) of a form in \(\mathcal H_{\rm c}\), at the same finite
level. This proves separation with variable \(a_\infty\). In particular,
the argument applies to each nonzero vector of \(W\) orthogonal
to the original selected rows.

\paragraph{The endpoint scalar at a specified coefficient place.}
\teamanchor{proof.weighted-endpoint-local-scalar}%
Set \(\zeta_\vartheta=c'_\vartheta/c'_1\) for this common path.
For the local fit at \(\iota\), each square gives a product of
at most \(t\) squared inverse ramified ratios and a power of
the conductor norm. The latter is an \(\ell\)-unit, and the
inverse ratios have local logarithm at most \(v_\ell\).
Multiplication of the same square scalars as in
\eqref{eq:weighted-completion-period} gives%
\footnote{\weightedpin{grh-compact/src/New/Exchange/RowLadder.lean}
{16-L52}{The factors of the path squares},
\weightedpin{grh-compact/src/New/Exchange/CrudeRows.lean}
{52-L65}{The inverse flip factor}, and
\weightedpin{grh-compact/src/New/Exchange/CostOfSquares.lean}
{26-L63}{Conductor and ramified factors}.}
\begin{equation}
\label{eq:weighted-completion-scalar}
   \log|j_\iota(|\zeta_\theta|^{-2})|_\ell\le4m^2t v_\ell .
\end{equation}

\paragraph{The residue indicators in the arithmetic jet.}
\teamanchor{proof.weighted-endpoint-edge-transport}%
For \(A\in\operatorname{Mat}_{n',n}(R_\vartheta)\), put
\([A,B]=\sum_{i,j} A_{ij}B_{ij}\). Let \(\psi_\vartheta\)
be the primitive additive character at the chosen generator
\(\lambda_\vartheta\) of \(\mathfrak m_\vartheta\), and let
\(U_\vartheta\) be the support of the modulus. It meets none
of its conjugate places. For a lift \(\widetilde B\) of \(B\)
on \(U_\vartheta\), let
\(x_B=\overline{\mathbf 1_{U_\vartheta}
               \lambda_\vartheta^{-1}\widetilde B}\).
Character orthogonality gives the exact identity
\[
 F_e^\circ=\frac{1}{|R_\vartheta|^{nn'}}
       \sum_B\psi_\vartheta(-[S_e,B])\,T(x_B)v_{\rm int},
 \qquad v_{\rm int}=v^{nn'}v_0 .
\]
Here \(T\) is normalized finite Heisenberg translation.
The Fourier mode \(\psi_\vartheta([X,B])\) acts by \(T(x_B)\)
in the conjugate coordinates. Its quadratic normalization
has phase one since the support is on one side of each
conjugate pair. After every complex automorphism the
absolute coefficient sum in the displayed expansion is one.

For \(k,k'\) in the standard compact groups, the finite
Weil action on the indicator datum is
\[
 \omega_{a,b,f}(k,k')F_e^\circ
 =\chi_{e,\vartheta}\!\left(\frac{\det\widetilde k'}
                                  {\det\widetilde k}\right)
     v^{nn'}\operatorname{name}_{\mathfrak m_\vartheta}
             (\mathbf 1_{\widetilde k'S_e\widetilde k^{-1}})v_0 .
\]
The tildes on \(k,k'\) now denote reductions. Indeed the
transformed indicator is evaluated at
\((\widetilde k')^{-1}X\widetilde k\), and the splitting
characters give the displayed determinant quotient.
The components of \(\chi_{e,\vartheta}\) are
\(\omega_w^{j+1}\) in \(V\) and
\((\omega_w\vartheta_w)^{j+1}\) in \(Z\).
The indicator index thus changes by the inverse of the
action defining tidiness. The product of the determinant
factors is cancelled by the inverse character from tidiness.
This gives joint invariance in the class variables with the
fixed input levels, using \((k,k\perp1)\) at the terminal
block. The arithmetic readings preserve this exact identity
on the combined finite sum.
\footnote{\weightedpin{grh-compact/src/Lib/Weil/MixedModel/NameModulus.lean}
{163-L187}{The primitive character and named Fourier modes},
\weightedpin{grh-compact/src/New/Curve/IntegralVacuum.lean}
{17-L45}{Exact Fourier coefficients and denominator},
\weightedpin{grh-compact/src/Lib/Weil/DualPair/TwoLineTerms.lean}
{64-L82}{The finite change of arguments}, and
\weightedpin{grh-compact/src/New/Exchange/GridFirstColumns.lean}
{136-L193}{Invariance of the joint sum}.}

\paragraph{Transport of the integral gauge.}
\teamanchor{proof.weighted-endpoint-edge-gauge}%
Put \(J=g^{b_1}\), \(q_\vartheta=\mathrm N Q_\vartheta\).
For an embedding \(\iota:\mathbf C\to\mathbf C_\ell\) with
\(\ell\nmid q_\vartheta\), use the Weyl gauge at a finite
point \(g_0\) given by the maximum absolute value after
transporting by \(g_0^{-1}\) at the readings
\(\iota\circ\Phi_\tau\circ{\rm conj}\). The transport uses
the determinant exponents \(e_{z,j,\tau}\).
Use the \(\iota\)-integral span \(\mathrm{Std}_J\) of
\(T(x)v_{\rm int}\) with \(Jx\) integral above \(\ell\).
Since \(q_\vartheta J\in\mathfrak m_\vartheta\), \(Jx_B\)
is integral above \(\ell\). Also \(|R_\vartheta|\) divides
\(|\mathrm N_{E/\mathbf Q}(q_\vartheta J)|=(q_\vartheta J)^{2f}\).
The Fourier identity shows that
\(J^{2f nn'}F_e^\circ\in\mathrm{Std}_J\), using the fact that
the phases are roots of unity.

Make \(g,h\) integral at the places over \(\ell\) by
independent rational unitary changes. Weak approximation
follows by the Cayley chart on skew-Hermitian matrices and
approximation in its coordinates. Rational covariance
transports the two gauges to these new points, where
integral Weyl matrices and their integral inverses preserve
the maximum norm. The finite action there preserves
\(\mathrm{Std}_J\). Indeed, the move of the vacuum is the
finite Weil multiplier divided by its squared complex
modulus times the sum indexed by the intersection lattice
quotient. Its multiplier and index are \(\ell\)-units,
since the standard lattice and its inverse image agree at
the primes above \(\ell\). For the multiplier one uses
the product of the local determinant absolute values of
the intersection frame. The algebraic character valuation
formula gives units for both determinant characters at the
integral points. Conjugation of the translations moves
\(x\) by \(X\mapsto hXg^{-1}\), preserving the integral
condition on \(Jx\) over \(\ell\).

The
\hyperlink{proof.weighted-local-source-gauge}{normalized CM product calculation}
applies with \(\delta=J\) and the multiplier \(c_0=1\)
of \(\mathrm{per}\). The prime-to-\(\ell\) frame in that
calculation expresses the normalized data by elliptic
products at \(y_i\) with \(n_0J y_i\in\Lambda\), where \(n_0\)
is prime to \(\ell\). The coefficients
\(\iota((s_0/u_0)\overline{\Phi_\tau(w_i)})\) are integral,
and the frame-length denominator is a unit. The derivative
algebra of the marked elliptic model has integral constant
terms and is preserved by \(d\,\partial_z\).
The isogeny identity then makes the derivatives of the
translated factors integral after rescaling by
\(\overline{n_0J}d\). The chain rule on the products bounds the
factorial jet paired with a homogeneous integral polynomial
of degree \(D_{z,j}\) by \(|\iota(Jd)|_\ell^{-D_{z,j}}\).
The integral Schur entries have exactly this degree when
nonzero. Combining the entrywise estimate with the
ultrametric inequality and the rational changes, the matrix
\[
       J^{2f nn'}(Jd)^{D_{z,j}}J_e(g,h)
\]
preserves the integral gauges between the two class points.
The fixed integer \(\nu=J^{2f n_*^2}(Jd)^{C_{\rm W}}\) works
for all edges and all \(\iota\) above \(\ell\nmid q_\vartheta\).
\footnote{\weightedpin{grh-compact/src/New/Exchange/CrudeModularData.lean}
{118-L139}{Normalization of the residue data},
\weightedpin{grh-compact/src/New/Curve/FormalParameter.lean}
{23-L52}{The derivative algebra for CM products},
\weightedpin{grh-compact/src/New/Curve/ConjugatesOfJets.lean}
{25-L96}{The standard-data stability and tight frame}, and
\weightedpin{grh-compact/src/New/Curve/ConjugatesOfJets.lean}
{112-L153}{The transported gauge and denominator per box}.}

\paragraph{Complex readings of the same arithmetic edge.}
\teamanchor{proof.weighted-endpoint-edge-growth}%
The normalized Taylor series of the finite translates
have algebraic coefficients by the marked CM products and
finite recovery in the
\hyperlink{proof.weighted-cm-jet-transport}{CM Taylor transport}.
The finite Weil moves are combinations of translations
with algebraic multipliers, and the extra determinant
characters are algebraic as well. Thus the integral Schur
polynomials give algebraic jets on the same data.
The scalar \(\beta_{z,j}\) is used in the complex
automorphic identity before these arithmetic readings.

For \(\sigma\) fixing \(\mathbf Q(\sqrt{-d})\), use the
semilinear Taylor map \(L_\sigma\) on the standard Hom
space. It permutes the places by \(\Phi_\tau\mapsto
\sigma\circ\Phi_\tau\), sends \(T(x)\) to
\(T(b_\sigma x)\), and changes \(\mathrm{per}\) to
\(t_\sigma\), with \(|t_\sigma|^2=\Omega^{-1}\).
Here \(b_\sigma\bar b_\sigma=\chi_{\rm cyc}(\sigma)\)
for the finite CM Tate multiplier \(b_\sigma\).
Factor \(b_\sigma=z_\sigma a_\sigma\) with
\(z_\sigma\bar z_\sigma=1\) and \(a_\sigma\) an integral
unit equal to one on \(U_\vartheta\). At
\(\overline{U_\vartheta}\) take
\(a_\sigma=\chi_{\rm cyc}(\sigma)\), and complete the
factorization on the other places with integral units of
the specified norm. The vector \(a_\sigma x_B\) is then
the translating vector for \(\chi_{\rm cyc}(\sigma)B\),
modulo the integral lattice. Its phase is paired with
\(\sigma\psi_\vartheta(-[S_e,B])
=\psi_\vartheta(-[S_e,\chi_{\rm cyc}(\sigma)B])\).
The change of index is bijective in the same residue ring.
If \(L_\sigma v_0=c_\sigma\omega_f(z_\sigma)v_0\), the
result is
\[
 L_\sigma F_e^\circ
 =\sigma(v^{nn'})c_\sigma\omega_f(z_\sigma)
       \operatorname{name}_{\mathfrak m_\vartheta}
                       (\mathbf 1_{S_e})v_0 .
\]
Transport leaves the finite move \(G\) at the same argument.
Its multiplier is \(\sigma(\Delta_f(G))/\Delta_f(G)\).
Indeed its vacuum pairing is
\(\Delta_f(G)/|\Delta_f(G)|^2\), and the squared modulus
is a rational lattice index. Comparing the vacuum
pairings gives the claimed multiplier on the vacuum,
and translation covariance gives it on the span.
The factor \(\sigma(a(\det h)/b(\det g))\) is retained.
Both of these factors have modulus one by the
principal-ideal norm formulas. For \(\sigma\) not fixing
the quadratic field, compose on its output with complex
conjugation for the norm estimate, since the Weyl Gram
matrices are real.

The
\hyperlink{proof.weighted-cm-moved-vacuum}{moved vacuum bound}
now gives an entire majorant \(M e^{\|Z\|^2/2}\) at the
trivial twist, where \(M=\exp(C_d(N_1nn'+1)^2)\).
The absolute coefficient sum of the translated vacua
is one. Write \(D_0=fnn'\) and let \(r_{\tau,i}\) be
the row lengths in the rectangle. For unit input and dual
test vectors in the conjugate Weyl norms, the exterior
tensor model bounds the coefficient polynomial by the
product of the singular values to these row powers.
For one singular value \(s\) and row length \(r_0\),
the \(r_0\)-th term of \(e^{s^2/2}\) yields
\(s^{r_0}\le (r_0!)^{1/2}2^{r_0/2}e^{s^2/4}\).
The scale \(t_\sigma\) contributes \(\Omega^{-r_0/2}\)
for the row. The factorial pairing is Gaussian integration
against \(e^{-\|Z\|^2}\pi^{-D_0}\,dZ\).
The product of the two bounds leaves
\(e^{-\|Z\|^2/4}\), whose integral against
\(\pi^{-D_0}\,dZ\) is \(4^{D_0}\). Using \(\log 2\le1\)
gives the operator bound
\[
 \|\sigma J_e(g,h)\|_{\mathrm W}\le M4^{D_0}
   \exp\left(\frac12\sum_{\tau,i}
     \bigl[\log(r_{\tau,i}!)+(1-\log\Omega)r_{\tau,i}\bigr]\right).
\]
This includes degree zero. The scale in the squared
bound is \(\Omega^{-1}\) per box at the same period
normalization, which gives the one-step loss for the
endpoint class sum.
\footnote{\weightedpin{grh-compact/src/Lib/CM/Reciprocity/SemilinearWeilLaw.lean}
{22-L78}{The multiplier and transported indicator},
\weightedpin{grh-compact/src/New/Curve/ConjugatesOfJets.lean}
{156-L211}{Taylor transport and envelope}, and
\weightedpin{grh-compact/src/New/Rows/MemberKernels.lean}
{42-L87}{The Schur jet norm estimate}.}

\paragraph{Finite class sums and an integer that clears the endpoint period.}
\teamanchor{proof.weighted-endpoint-finite-evaluation}%
In the arithmetic realization of \(\mathcal E_\vartheta\), the last
path has two definite chains of \(m\) arrows. Write \(K_{z,j}\) for
the level at step \(j\) on side \(z\in\{-,+\}\). Put \(J=g^{b_1}\)
and keep \(\mathfrak m_\vartheta=\mathfrak bQ_\vartheta\) and
\(q_\vartheta=\#(\mathcal O_E/Q_\vartheta)\). The residue
ring for every arrow is \(\mathcal O_E/\mathfrak m_\vartheta\), and
\(q_\vartheta J\in\mathfrak m_\vartheta\). Write \(n_*=2n_m\) for
the rank bound on both chains and \(C_{\rm W}=b_1r^4\) for their
bound on the total number of polynomial Weyl boxes. Each terminal
level is read at its prescribed definite group. The \(+\) terminal
point above a \(-\) point \(h\) is the block point
\(\operatorname{diag}(h,F_\vartheta^m)\), in the given order of its
columns, where
\[
   F_\vartheta=\frac{\overline{\lambda_\vartheta\alpha_{\rm W}}}
                     {\lambda_\vartheta\alpha_{\rm W}}.
\]
Here \(\lambda_\vartheta\) is the finite idele generating the
modulus, and \(\alpha_{\rm W}\) is the fixed Weil idele. The integer
branching matrix pairs the two terminal vectors. Its character on
the added line has exponents \(s_\tau=e_\tau+\sum_h\gamma_\tau(h)\),
where \(e_\tau\) is the determinant exponent and \(\gamma_\tau(h)\)
counts the shortened columns of height \(h\). There are at most as
many shortened columns as boxes, so
\(\sum_\tau|s_\tau|\le C_{\rm W}+f\lceil20b\rceil\).
Its operator norm as a bilinear form for the two Weyl norms is at
most \((C_{\rm W}^{C_{\rm W}})^3\).
\footnote{\weightedpin{grh-compact/src/New/Exchange/GridMove.lean}{154-L235}{Definite weights, levels, and block point},
\weightedpin{grh-compact/src/New/Seesaw/PairTensorData.lean}{19-L55}{Integer branching form and its character},
\weightedpin{grh-compact/src/Lib/GLn/Branching/Branching.lean}{224-L237}{Shortened columns in the last-coordinate exponent},
and \weightedpin{grh-compact/src/Lib/SiegelWeil/RallisInner/SharpData.lean}{283-L288}{Orientation of the block scalar}.}

The last-path datum is a linear combination of at most
\[
   |\mathcal O_E/\mathfrak m_\vartheta|^{\,n_*^2(4m+2m^2)}
\]
evaluations at edge tuples, with every nonzero coefficient a root
of unity. Expand the two kernels on each tuple before summing over
classes. An intermediate level change on one side moves through the
two adjacent jets. The change in the datum is the product of its
edge characters, which cancels the reciprocal factor from the Weil
action. At the input vertices use level invariance of the two forms.
At the paired terminal vertex move the two arguments by \(k\) and
\(\operatorname{diag}(k,1)\), so the last-coordinate character is
one. The resulting sum over edge tuples is separately right
invariant in the class variable at each integration. This is the
\hyperlink{proof.weighted-short-class-sum}{joint class-sum argument}
applied to two branches with a terminal contraction. Every endpoint
level lies in the chosen principal level, so the
\hyperlink{proof.weighted-free-level}{free-level calculation}
gives class mass one. The endpoint period is the sum over \(2m+1\)
class variables of the combined bracket. Inside the bracket each
summand is one root of unity times the branching pairing of two
products of \(m\) jet matrices, applied to the original endpoint
values. The count uses both input levels and a single terminal level.
\footnote{\weightedpin{grh-compact/src/New/Exchange/GridMove.lean}{37-L76}{Level action and short-sum datum},
and \weightedpin{grh-compact/src/New/Exchange/GridFirstColumns.lean}{136-L193}{Invariance and the class expansion}.}

Use the common integer
\[
 \nu=J^{2fn_*^2}(Jd)^{C_{\rm W}} .
\]
The \hyperlink{proof.weighted-endpoint-edge-gauge}{transported edge estimate}
shows that \(\nu\) times each edge matrix preserves the integral Weyl
gauges when \(\ell\nmid q_\vartheta\).

The factor \(JN_1\) clears both \(F_\vartheta\) and its inverse at
every finite place over \(\ell\nmid q_\vartheta\). Indeed,
\(q_\vartheta J/\lambda_\vartheta\),
\(N_1\overline{\alpha_{\rm W}}\), \(\overline{\lambda_\vartheta}\)
and \(\alpha_{\rm W}^{-1}\) are integral finite adeles. Their
product is \(q_\vartheta JN_1 F_\vartheta\). Conjugation gives
the same result for \(F_\vartheta^{-1}\), and \(q_\vartheta\) is a
unit at the coefficient place. In transported Weyl coordinates the
integer branching form is evaluated on the two integral vectors
obtained after the \(m\) steps on each side. Undoing the transport
multiplies the pairing by the character of the extra line at
\(F_\vartheta^m\). The common integer
\[
  N_{\rm end}=(JN_1)^{m\sum_\tau|s_\tau|}\nu^{2m}
\]
therefore makes each summand of the combined bracket integral.
Apply the ultrametric inequality to conclude that
\(N_{\rm end}\mathcal E_\vartheta\) is integral when both input
forms are integral in their transported Weyl gauges. The integer
depends only on the datum. From \(b_1=O(b)\), \(\log J=O(b\log b)\),
\(n_m=O(x^{11/2})\) and \(f=O(x^{11}\log x)\), its logarithm is
\(O_d(x^{30}b^2)\), hence at most \(x^{99}b^2\) eventually.
\footnote{\weightedpin{grh-compact/src/Lib/Weil/MixedModel/NameModulus.lean}{96-L114}{The modulus generator},
\weightedpin{grh-compact/src/New/Data/WeilDataOfE.lean}{198-L220}{The fixed Weil idele},
and \weightedpin{grh-compact/src/New/Rows/LineArithmetic.lean}{20-L23}{The integer in the modulus}.}

At the trivial twist the residue cardinality is \(J^f\).
The \hyperlink{proof.weighted-definite-class-count}{definite class
count} with \(q\le n_*\) and \(M=J\) bounds the number of classes at
each endpoint level by \(\exp(x^{36}(1+\log J))\), since the level
contains the principal subgroup modulo \(J\).

The \hyperlink{proof.weighted-endpoint-edge-growth}{one-step bound on the same arithmetic edge}
gives logarithmic loss at most
\[
 C_d(N_1 n_*^2+1)^2+fn_*^2\log 4
 +\tfrac12 C_{\rm W}(1+\log(C_{\rm W}+1)+|\log\Omega|).
\]
Here \(\log(h!)\le h\log(h+1)\) for each row of length \(h\),
and \(|\log\Omega|=O_d(1+\log x)\) by the field and CM
normalization. Applying this bound to the \(2m\) jets, the norm
bound to the branching form and the triangle inequality to the class
and tuple sums gives a total log loss \(O_d(x^{60}b)\). If the
two input Weyl norms are bounded by \(B,B'\) on these conjugate
values, the endpoint bound is
\(|\sigma\mathcal E_1|\le \exp(x^{200}b^2)BB'\).

\paragraph{Short endpoint forms.}
\teamanchor{proof.weighted-endpoint-short-source}%
At the definite source \(\mathscr F\) use the
\hyperlink{proof.weighted-classwise-evaluation}{classwise evaluation
construction} with \(q\le n_*\) and \(D=C_{\rm W}\). Its frames,
including the ones for the determinant exponents of this source
weight, have entry bounds \(P_{\rm frame}=6dN_1(fn_*)^2\).
The resulting delta forms span \(\mathscr F\) and are integral at
all coefficient embeddings in the endpoint level.

The squared Weyl norm of an evaluation vector is the product of the
leading Gram minors of its matrix arguments, each raised to its
column multiplicity. This is the highest-weight Weyl kernel evaluated
on the matrix and its conjugate. Hadamard's inequality bounds the
squared norm including the determinant factors by
\((n_*^2P_{\rm frame})^{2(C_{\rm W}+fn_*)}\), uniformly on the conjugates,
because \(C_{\rm W}\) bounds the total number of boxes over all
embeddings. This is at most \(\exp(x^{99}b^2)\) for large \(x\).
\footnote{\weightedpin{grh-compact/src/Lib/GLn/SchurModule/Gram.lean}{43-L59}{Product of leading Gram minors},
\weightedpin{grh-compact/src/New/Columns/ChainValuation.lean}{101-L147}{Determinant corrections in the actual Weyl coordinates},
and \weightedpin{grh-compact/src/New/Exchange/CrudeRows.lean}{203-L218}{Application at the source endpoint}.}

\paragraph{The integral duality used by the principal tests.}
\teamanchor{proof.weighted-endpoint-dual-extraction}%
At \(Y_m=V_{0,m}\) the row lift and the delta test have
opposite doubled parameters. At each place their algebraic
highest weights are \(\mu\) and \(\mu^\vee=-w_0\mu\), with
\(w_0\) reversing the entries. The test multiplicity at
height \(h\) is the row multiplicity at height \(N-h\),
where \(N=\operatorname{rank}Y_m\). Apply the
\hyperlink{proof.weighted-fischer-weyl}{Jacobi and Fischer calculation}
to these diagrams, in the original Weyl coordinates.
Write \(A\) for the integral isometry from the row
to the test diagram, obtained by the complementary-minor
substitution. Its inverse is integral. Since the two
decreasing flags are reverse to each other, its sign can be
fixed to send the highest evaluation of the row to that of
the test. For a matrix \(g\), the substitution replaces the
row action by \(g^{-{\mathsf T}}\) times the determinant
raised to the number of columns. The determinant exponents
complete the two actions to exact contragredients, since
\(\mu_1-\mu_N\) is that column count and
\(\mu^\vee_N=-\mu_1\).

With \(G^\vee\) the test Gram matrix, the invariant bilinear
pairing therefore has integral matrix \(B=G^\vee A\).
Its dual norm is the row Weyl norm by the isometry. Let
\(F^\vee\) be the integral Fischer matrix of the test shape
and \(h^\vee\) the squared Fischer norm of its highest
flag-minor product. The common calculation gives
\(F^\vee G^\vee=h^\vee I\), so \(B'=A^{-1}F^\vee\)
satisfies \(B'B=h^\vee I\). The grid parameter bound also
applies to the reflected weight, whose number of polynomial
boxes \(D^\vee\) is \(O(fn_* b)\le C_{\rm W}\).
The determinant expansion and factorial norm bound give
\(1\le h^\vee\le(D^\vee!)^3\), including the empty shape.
Thus \(\log h^\vee\le3 C_{\rm W}\log C_{\rm W}\).

Use the delta form of the opposite weight at class \(h\) in
the principal test level \(K^{\rm pr}_\vartheta\), with
value \(v\). Write \(\Delta^\flat\) for its finite algebraic
form. Its highest coefficient at \((x,g)\) is
\(\Delta^\flat(g)^{\mathsf T}B\rho(x)u\), since \(Au=u^\vee\)
and \(B\) is invariant. It pairs with the row-lift
coefficient by compact Schur orthogonality.
As in the definite edge calculation, averaging the Gram-unit
highest projector divides by the row weight dimension \(d_w\).
The index in \eqref{eq:weighted-completion-period} gives
the class-sum measure on the principal test level, and
freeness gives mass one for that delta's class. The result is
\(v^{\mathsf T}B(T_\vartheta Q_{a_\infty}\phi)(h)/d_w\).
The determinant uses the same pointwise row lift at
\(K_\vartheta^{\rm res}\). Its extraction uses the integral inverse
identity even at primes dividing \(h^\vee\).
\footnote{\weightedpin{grh-compact/src/Lib/GLn/Dual/DualWeights.lean}
{14-L40}{Opposite-weight pairing and the highest coefficient},
\weightedpin{grh-compact/src/New/Exchange/CrudeRows.lean}
{219-L235}{Extraction at the principal terminal tests}.}

\paragraph{One completion of the selected rows.}
\teamanchor{proof.weighted-endpoint-completion-union}%
Use the
\hyperlink{proof.weighted-endpoint-dual-extraction}{integral pairing and principal delta tests}
just established.

If \(Q_{a_\infty}\ne0\), then \(\alpha_{a_\infty}\ne0\).
Indeed the untwisted chain on \(\mathcal H_0\) is injective by the
row triangle, and the terminal delta tests separate its values.
For each such translate choose \(s_{a_\infty}\) with
\(s_{a_\infty}\alpha_{a_\infty}d_w/c'_1=h^\vee N_{\rm end}\)
and write \(f_{a_\infty,\phi}=s_{a_\infty}Q_{a_\infty}\phi\)
for a short source form \(\phi\).
\footnote{\weightedpin{grh-compact/src/New/Exchange/ChainNorms.lean}
{39-L55}{Injectivity at the full row principal level}.}
The normalizing factors \(s_{a_\infty}\) are independent of the
twist. In the nonvanishing case the pairing of
\(\zeta_\vartheta T_\vartheta f_{a_\infty,\phi}\) with a dual vector is
\(h^\vee N_{\rm end}\mathcal E_\vartheta(\Delta_{h,v},\phi)\).
At each coefficient place use a representative \(h\) integral
above its rational prime, by weak approximation. For each row of
\(B'\), the corresponding delta test is integral in its Weyl gauge.
The endpoint clearing shows that the pairing is \(h^\vee\) times
an integral scalar. By \(B'B=h^\vee I\), cancellation proves
coordinate integrality of \(\zeta_\vartheta T_\vartheta f_{a_\infty,\phi}\),
including at primes dividing \(h^\vee\).

For an automorphism \(\sigma\) and a dual vector \(v\), take the
delta test with coefficient \(\sigma^{-1}v\). The conjugate delta
has Weyl norm at most \(\|v\|_{\rm W}\) on the entire class,
by unitarity of the rational action at every embedding. The pairing
matrix is integral and fixed by \(\sigma\). The trivial-twist
endpoint bound and norm duality therefore give
\[
 \bigl\|\sigma (T_1 f_{a_\infty,\phi})(h)\bigr\|_{\rm W}^2
 \le \exp\bigl(2\log h^\vee+(2x^{200}+3x^{99})b^2\bigr)
 \le \exp(x_1^2b^2).
\]
The determinant and norm costs are uniform in \(a_\infty\). Apply the
\hyperlink{proof.weighted-projection-transfer}{projection transfer}
with the identity translate to these rows. Write \(P_W\) for the
orthogonal projector onto \(W\). The transfer is by the fixed
nonzero multiple \(|\mathcal G| P_W\), and it preserves the local
integrality after multiplying by the same \(\zeta_\vartheta\).
The resulting rows have untwisted squared class-sum norms at most
\(e^{B_z}\), where \(B_z=x_1^2 b^2+x^{99}\).

Let \(L_y\subset W\) be the span of the original selected rows.
If a vector in \(W\) is perpendicular to both \(L_y\) and every
projected row just constructed, it is perpendicular to
\(r_{a_\infty}\mathcal H_{\rm c}\) for every \(a_\infty\), because the short source
forms span and \(S\) is onto. The corank-one period then forces the
vector to be zero. Consequently the projected short images span \(W/L_y\).

\paragraph{The untwisted frame and its determinant.}
\teamanchor{proof.weighted-completion-cost}%
Choose \(u=R-n\) representatives \(z\) from the short completion
rows whose images form a basis of \(W/L_y\). Fix this basis
\(w=(y,z)\), with its chosen row preimages, before specifying a
coefficient place.
\teamanchor{proof.weighted-mixed-gram}%
At each \(\sigma\), the squared covolume of the coefficientwise
frame \((\Theta_1w)^\sigma\) is at most \(e^{V+u B_z}\). This follows
by successively projecting each added vector orthogonally off the
previous span and using its squared norm bound. By rationality of the
Weyl Gram entries and the class weights, \(\sigma H_1\) for
\(H_1=\det\operatorname{Gram}(\Theta_1 w)\) is the mixed determinant
pairing this frame with \((\Theta_1w)^{\sigma^\#}\). The exterior
Cauchy--Schwarz inequality from the
\hyperlink{proof.fitted-row-mixed-gram}{base quotient comparison}
then gives \(|\sigma H_1|\le e^{V+u B_z}\).

\paragraph{Selection of a coefficient place.}
\teamanchor{proof.weighted-selected-coefficient}%
From the fixed untwisted determinants and rank-one multipliers put
\[
       Y=H_1D_0K_\Gamma^R .
\]
The untwisted chain is injective on \(W\) and every factor in \(D_0\)
and \(K_\Gamma\) is nonzero. The estimate for \(H_1\), the archimedean
bound in the single-lift comparison, and the \(\kappa/2\) sharp pairs
at each stage give
\begin{equation}
\label{eq:weighted-selected-place-arch}
 \log|\sigma Y|
 \le V+uB_z+2RrmB_0+Rm\kappa(3B_0-S_0)
       \qquad(\sigma\in\operatorname{Aut}(\mathbf C)).
\end{equation}
It follows that \(Y\in\overline{\mathbf Q}^{\,\times}\). If \(Y\)
were transcendental, replacing \(Y\) by \(hY\) in a transcendence
basis over \(\overline{\mathbf Q}\) would extend to an automorphism
of \(\mathbf C\) for every positive integer \(h\), contrary to
\eqref{eq:weighted-selected-place-arch}.

The datum and the row preimages already fix \(F_\ell\),
\(\beta_{j,\ell}\), \(v_\ell\) and \(P_\ell\). Put
\begin{equation}
\label{eq:weighted-selected-place-cost}
\begin{aligned}
 A_\ell={}&F_\ell+R(r+\kappa)\sum_j\beta_{j,\ell}
                  +2Rm\kappa P_\ell\log\ell\\
           &+\bigl[4u m^2t+R(2m\kappa+10m^2+m^2t)\bigr]v_\ell .
\end{aligned}
\end{equation}
Let \(f_Y\) be the monic minimal polynomial over \(\mathbf Q\), let
\(e=\deg f_Y\), and put \(q_\ell=e^{-1}\log|f_Y(0)|_\ell\).
Its positive values occur at the prime divisors of the reduced
denominator of \(f_Y(0)\). Write \(|f_Y(0)|=a_Y/b_Y\) in lowest terms.
The sum of the finite logs is \((\log b_Y-\log a_Y)/e\).
If \(a_Y\ge b_Y\), this sum is nonpositive and suffices for the lower
determinant bound. If \(a_Y<b_Y\), the denominator is greater than one.
Choose the least prime divisor \(\ell_*\) maximizing
\(q_\ell-A_\ell\) over that finite set. At \(\ell_*\), take a root
of \(f_Y\) of largest absolute value. The corresponding embedding
of \(\mathbf Q(Y)\) extends to a coefficient isomorphism
\(j_\iota\colon\mathbf C\simeq\overline{\mathbf Q}_{\ell_*}\).
One extends first to the algebraic closures of \(\mathbf Q\) and
then identifies transcendence bases, which both have the
cardinality of the continuum. The largest root norm is at least
the geometric mean of the \(e\) root norms, so
\(q_{\ell_*}\le\log|j_\iota Y|_{\ell_*}\). The inequality
\(\log|j_\iota Y|_{\ell_*}\le A_{\ell_*}\) will therefore give
\(q_\ell\le A_\ell\) at every denominator prime.

\paragraph{The local calculation at the selected place.}
\teamanchor{proof.weighted-selected-coefficient-tests}%
If \(|f_Y(0)|<1\), set \(\ell=\ell_*\) and apply
the \hyperlink{proof.weighted-fit}{local fit} at \((\ell,j_\iota)\)
to obtain one \(\theta\). Use the index recombination for \(y\)
at \(j_\iota\) and
\(j_\iota\circ\operatorname{conj}\) from
\eqref{eq:weighted-A-paired-index}. Extend it on the completion rows
by \(\zeta_\theta I_u\). The two
resulting lifted frames are integral at the paired places. The
final class stabilizers are trivial, so their mixed Gram entries
are integral. Taking determinants, writing
\(H_\vartheta=\det\operatorname{Gram}(\Theta_\vartheta w)\), gives
\[
 |j_\iota(H_\theta|\zeta_\theta|^{2u})|_\ell\le e^{F_\ell}.
\]
The endpoint scalar bound \eqref{eq:weighted-completion-scalar} gives
\begin{equation}
\label{eq:weighted-exact-completion-height}
 \log|j_\iota H_\theta|_\ell\le F_\ell+4um^2t v_\ell .
\end{equation}
On the same basis \(w\), \(H_\theta/H_1=Z_\theta/Z_1\).
Add \eqref{eq:weighted-volume-lift} and
\eqref{eq:weighted-exact-completion-height}. The twist cancels in
the product \(H_\theta D_0K_\Gamma^R Z_1/Z_\theta=Y\).
The sum of the local allowances is exactly \(A_{\ell_*}\)
in \eqref{eq:weighted-selected-place-cost}.

\paragraph{The product formula and the lower determinant bound.}
\teamanchor{proof.weighted-selected-product-formula}%
Apply the rational product formula to
\((-1)^e f_Y(0)=\mathrm N_{\mathbf Q(Y)/\mathbf Q}Y\).
Its archimedean logarithm divided by \(e\) is the average of
\(\log|\sigma Y|\) over the \(e\) conjugates. If \(|f_Y(0)|\ge1\),
the finite sum of the \(q_\ell\) is nonpositive. If \(|f_Y(0)|<1\),
each \(q_\ell\) at a denominator prime is at most its \(A_\ell\)
by the local test and maximality, and the other \(q_\ell\) are
nonpositive. In both cases the finite sum is at most \(\sum_{\ell\in S}A_\ell\)
for the finite set \(S\) of denominator primes and primes through
\(\max(N,\lceil x_1\rceil)\). It contains the support of \(P_\ell\).
The finite sums
of \(F_\ell\), \(\sum_j\beta_{j,\ell}\) and \(v_\ell\) are at most
\(2\delta\), \(mB_0\) and \(x^{999}\), respectively. The identity
\eqref{eq:weighted-forced-rank-one} combines the sum involving
\(P_\ell\) with \(-Rm\kappa S_0\) to give
\(-Rm\kappa\mathcal U=-R h_{\rm main}\). Thus
\begin{equation}
\label{eq:weighted-exact-lower}
 R(h_{\rm main}-B_*)\le V+2\delta+u(B_z+B_{\rm c}),
\end{equation}
where
\[
 \begin{aligned}
 B_{\rm c}&=4m^2t x^{999},\\
 B_*&=2rmB_0+(r+\kappa)mB_0
          +3m\kappa B_0+(2m\kappa+10m^2+m^2t)x^{999}.
 \end{aligned}
\]
Here \(B_*=O(\mathcal B)+O(x^{1005})\), since \(r/\kappa\to1\),
\(m\kappa B_0=\mathcal B\) and \(t=5r\). Thus \(B_*\le2L\)
for sufficiently large scales. The rank loss is at most
\(u/R\le4b^{-3/2}\), and \(x_1^2=b^{o(1)}\), so eventually
\begin{equation}
\label{eq:weighted-volume-rows}
             u(B_z+B_{\rm c})\le Rb.
\end{equation}
Together with \(b\le L\), these bounds give Theorem~C.
\begin{equation}
\label{eq:weighted-long-forms}
 R(h_{\rm main}-3L)\le V+2\delta.
\end{equation}

\paragraph{The signed estimate.}
\teamanchor{proof.weighted-height}%
On the selected datum \(h_{\rm main}=m\kappa\mathcal U\) is
nonnegative for all sufficiently large \(x\). For each \(\sigma\),
\(w(\sigma)\le a\) and the minimum of the weights is \(a/x_1\).
A positive summand of \(\operatorname{Ask}(w(\sigma),\sigma)\) must
therefore have \(\nu\le x_1\). Uniformly in \(\sigma\),
\[
 \operatorname{Ask}(w(\sigma),\sigma)
 \le 2w(\sigma)\sum_{\nu\le x_1}\frac{\log\nu}{\nu}
 = O\bigl(w(\sigma)(1+\log^2 x_1)\bigr)
 = o(w(\sigma)\log b).
\]
Also \(g\ge b^{20}\) and \(\Omega\le2\pi e^{\ell_d}\), using
\(\operatorname{rd}_E\ge1\). The constant here depends on the fixed
\(d=d_k\). Since \(w_0(\sigma)\ge w(\sigma)\) and the logarithm of a
factorial is nonnegative, we have
\[
 S(w_0(\sigma))\ge w_0(\sigma)(20\log b-O_d(1))
                   \ge\operatorname{Ask}(w(\sigma),\sigma)
\]
eventually. Here \((\log x_1)^2/\log b=x^{-2\varepsilon/3}\to0\)
at the fixed precision. The estimate holds at real scales, so its
threshold applies to any selection of sufficiently large scales.

Substitute \(V=n(h_{\rm main}-m\operatorname{rem}+L)\) into
\eqref{eq:weighted-exact-lower} and keep the actual \(n\)
on the side containing the signed remainder. This yields
\[
 n m\operatorname{rem}\le nL+R B_*+2\delta
                          +u(B_z+B_{\rm c}-h_{\rm main}).
\]
Since \(u\ge0\), dropping \(-u h_{\rm main}\) gives an upper bound.
For rank loss \(u/R\le4b^{-3/2}\) we have \(n\ge R/2>0\) eventually,
including when the number of selected rows is chosen adaptively on
the datum. Using \(\delta\le5R\log(1+b^2)\mathcal B\), division
by \(n\), \(B_*\le2L\), and the completion bound give
\[
 m\operatorname{rem}\le 7L+20\log(1+b^2)\mathcal B
     \le (47\log b+20\log 2)\mathcal B.
\]
The right side is at most \(\mathcal B(\log b)^2\) for sufficiently
large \(x\), giving the signed bound on that datum,
\begin{equation}
\label{eq:weighted-star}
 m\operatorname{rem}\le fa(\log b)^5.
\end{equation}

The \hyperlink{proof.weighted-prime-identity}{weighted prime identity}
uses this signed bound and the same auxiliary scale \(x_1\).

\paragraph{Simultaneous height bounds in the compact Lean source.}
The compact version assembles fits at all coefficient places and
formulates reciprocal volume bounds as heights of two families.
The statements and their proofs are in the fixed compact
version.\footnote{\weightedpin{grh-compact/src/New/Selmer/FitTwistsExist.lean}
{24-L65}{Assembly of the simultaneous fitted family}.}
The determinant comparison above uses the cost at the selected
coefficient place and the finite sums of the numerical allowances.

\paragraph{Side estimate in the compact Lean source.}
The compact proof uses \(|h_{\rm main}|+m|\mathrm{rem}|\le b^2x^{99}\)
to replace \(nh\) by an upper bound in terms of the full rank \(R\).
The signed argument above uses \(h_{\rm main}\ge0\) and divides by
the number \(n\) of selected rows. The stronger estimate is recorded
from the statement in the fixed source. Its proof awaits comparison
with this derivation.\footnote{\weightedpin
{grh-compact/src/New/Endgame/ShortAndEnough.lean}{72-L82}{The compact two-sided estimate}.}
